% ==============================================================================
% BANCO CONSOLIDADO DE PREGUNTAS DE ADMISIÓN (7,500 PREGUNTAS)
% Universidad de El Salvador / Prueba AVANZO
% Asignaturas:
% 1. Matemáticas (1.1 al 1.7) -- 2,100 Preguntas
% 2. Lenguaje (2.1 al 2.5) -- 1,500 Preguntas
% 3. Literatura (3.1 al 3.5) -- 1,500 Preguntas
% 4. Ciencias Sociales (4.1 al 4.4) -- 1,200 Preguntas
% 5. Ciencias Naturales (5.1 al 5.4) -- 1,200 Preguntas
% ==============================================================================
\documentclass[10pt,a4paper]{article}
\usepackage[utf8]{inputenc}
\usepackage[T1]{fontenc}
\usepackage[spanish,es-nodecimaldot]{babel}
\usepackage{amsmath,amssymb,amsfonts}
\usepackage{geometry}
\usepackage{multicol}
\usepackage{enumitem}
\usepackage{tcolorbox}
\usepackage{fancyhdr}
\usepackage{hyperref}
\usepackage{titlesec}

\geometry{top=2cm, bottom=2cm, left=1.8cm, right=1.8cm}
\pagestyle{fancy}
\fancyhf{}
\rhead{\small\scshape Banco Maestro de Preguntas}
\lhead{\small\bfseries UES / AVANZO}
\cfoot{\thepage}
\renewcommand{\headrulewidth}{0.4pt}

\begin{document}
\begin{titlepage}
\centering
{\scshape\LARGE UNIVERSIDAD DE EL SALVADOR / PRUEBA AVANZO \par}
\vspace{1.5cm}
{\huge\bfseries BANCO GENERAL DE PREGUNTAS (7,500 PREGUNTAS) \par}
\vspace{0.5cm}
{\Large\bfseries Compendio Maestro Oficial \par}
\vspace{1.5cm}
\begin{tcolorbox}[colback=blue!5!white,colframe=blue!60!black,title=\textbf{RESUMEN DE CONTENIDO}]
\begin{enumerate}[label=\textbf{\Roman*.}]
  \item \textbf{Matemáticas} (1.1 al 1.7) -- 2,100 Preguntas
  \item \textbf{Lenguaje} (2.1 al 2.5) -- 1,500 Preguntas
  \item \textbf{Literatura} (3.1 al 3.5) -- 1,500 Preguntas
  \item \textbf{Ciencias Sociales} (4.1 al 4.4) -- 1,200 Preguntas
  \item \textbf{Ciencias Naturales} (5.1 al 5.4) -- 1,200 Preguntas
\end{enumerate}
\end{tcolorbox}
\vfill
{\large San Salvador, El Salvador \par}
{\large \today \par}
\end{titlepage}
\tableofcontents
\newpage

% ==============================================================================
% MATERIA 1: MATEMÁTICAS
% ==============================================================================
\section{Matemáticas (Módulos 1.1 al 1.7)}

\subsection{1.1 Números Reales}

\subsubsection*{SECCIÓN I: CLASIFICACIÓN DE NÚMEROS REALES (1–60)}

\begin{enumerate}[label=\textbf{\arabic*.}]
\setcounter{enumi}{0}
\item ¿Cuál es el conjunto más específico al que pertenece \(\sqrt{2}\)?
\begin{enumerate}[label=\textbf{\Alph*)}]
\begin{multicols}{2}
  \item Natural
  \item Entero
  \item Racional
  \item Irracional
  \item Complejo no real
  \item Imaginario
\end{multicols}
\end{enumerate}
\vspace{0.15cm}
\item ¿Cuál es el conjunto más específico al que pertenece \(-7\)?
\begin{enumerate}[label=\textbf{\Alph*)}]
\begin{multicols}{2}
  \item Natural
  \item Entero
  \item Racional
  \item Irracional
  \item Complejo no real
  \item Imaginario
\end{multicols}
\end{enumerate}
\vspace{0.15cm}
\item ¿Cuál es el conjunto más específico al que pertenece \(\frac{3}{4}\)?
\begin{enumerate}[label=\textbf{\Alph*)}]
\begin{multicols}{2}
  \item Natural
  \item Entero
  \item Racional
  \item Irracional
  \item Complejo no real
  \item Imaginario
\end{multicols}
\end{enumerate}
\vspace{0.15cm}
\item ¿Cuál es el conjunto más específico al que pertenece \(0\)?
\begin{enumerate}[label=\textbf{\Alph*)}]
\begin{multicols}{2}
  \item Natural
  \item Entero
  \item Racional
  \item Irracional
  \item Complejo no real
  \item Imaginario
\end{multicols}
\end{enumerate}
\vspace{0.15cm}
\item ¿Cuál es el conjunto más específico al que pertenece \(\pi\)?
\begin{enumerate}[label=\textbf{\Alph*)}]
\begin{multicols}{2}
  \item Natural
  \item Entero
  \item Racional
  \item Irracional
  \item Complejo no real
  \item Imaginario
\end{multicols}
\end{enumerate}
\vspace{0.15cm}
\item ¿Cuál es el conjunto más específico al que pertenece \(15\)?
\begin{enumerate}[label=\textbf{\Alph*)}]
\begin{multicols}{2}
  \item Natural
  \item Entero
  \item Racional
  \item Irracional
  \item Complejo no real
  \item Imaginario
\end{multicols}
\end{enumerate}
\vspace{0.15cm}
\item ¿Cuál es el conjunto más específico al que pertenece \(0.25\)?
\begin{enumerate}[label=\textbf{\Alph*)}]
\begin{multicols}{2}
  \item Natural
  \item Entero
  \item Racional
  \item Irracional
  \item Complejo no real
  \item Imaginario
\end{multicols}
\end{enumerate}
\vspace{0.15cm}
\item ¿Cuál es el conjunto más específico al que pertenece \(\sqrt{9}\)?
\begin{enumerate}[label=\textbf{\Alph*)}]
\begin{multicols}{2}
  \item Natural
  \item Entero
  \item Racional
  \item Irracional
  \item Complejo no real
  \item Imaginario
\end{multicols}
\end{enumerate}
\vspace{0.15cm}
\item ¿Cuál es el conjunto más específico al que pertenece \(-\frac{5}{2}\)?
\begin{enumerate}[label=\textbf{\Alph*)}]
\begin{multicols}{2}
  \item Natural
  \item Entero
  \item Racional
  \item Irracional
  \item Complejo no real
  \item Imaginario
\end{multicols}
\end{enumerate}
\vspace{0.15cm}
\item ¿Cuál es el conjunto más específico al que pertenece \(e\)?
\begin{enumerate}[label=\textbf{\Alph*)}]
\begin{multicols}{2}
  \item Natural
  \item Entero
  \item Racional
  \item Irracional
  \item Complejo no real
  \item Imaginario
\end{multicols}
\end{enumerate}
\vspace{0.15cm}
\item ¿Cuál de los siguientes números es irracional?
\begin{enumerate}[label=\textbf{\Alph*)}]
\begin{multicols}{2}
  \item \frac{2}{3}
  \item -\sqrt{16}
  \item \(\sqrt{7}\)
  \item 0.75
  \item 5
  \item -1
\end{multicols}
\end{enumerate}
\vspace{0.15cm}
\item ¿Cuál de los siguientes números es racional?
\begin{enumerate}[label=\textbf{\Alph*)}]
\begin{multicols}{2}
  \item \sqrt{2}
  \item \pi
  \item e
  \item \sqrt{5}
  \item \(\frac{7}{9}\)
  \item \sqrt{11}
\end{multicols}
\end{enumerate}
\vspace{0.15cm}
\item ¿Cuál de los siguientes números es entero?
\begin{enumerate}[label=\textbf{\Alph*)}]
\begin{multicols}{2}
  \item Natural
  \item Entero
  \item Racional
  \item Irracional
  \item Complejo no real
  \item Imaginario
\end{multicols}
\end{enumerate}
\vspace{0.15cm}
\item ¿Cuál de los siguientes números es natural?
\begin{enumerate}[label=\textbf{\Alph*)}]
\begin{multicols}{2}
  \item Natural
  \item Entero
  \item Racional
  \item Irracional
  \item Complejo no real
  \item Imaginario
\end{multicols}
\end{enumerate}
\vspace{0.15cm}
\item ¿A qué conjunto pertenece el número \(-\sqrt{25}\)?
\begin{enumerate}[label=\textbf{\Alph*)}]
\begin{multicols}{2}
  \item Natural
  \item Entero
  \item Racional
  \item Irracional
  \item Complejo no real
  \item Imaginario
\end{multicols}
\end{enumerate}
\vspace{0.15cm}
\item ¿A qué conjunto pertenece el número \(\sqrt{16}\)?
\begin{enumerate}[label=\textbf{\Alph*)}]
\begin{multicols}{2}
  \item Natural
  \item Entero
  \item Racional
  \item Irracional
  \item Complejo no real
  \item Imaginario
\end{multicols}
\end{enumerate}
\vspace{0.15cm}
\item ¿Cuál de los siguientes números es irracional?
\begin{enumerate}[label=\textbf{\Alph*)}]
\begin{multicols}{2}
  \item \frac{2}{3}
  \item -\sqrt{16}
  \item \(\sqrt{3}\)
  \item 0.75
  \item 5
  \item -1
\end{multicols}
\end{enumerate}
\vspace{0.15cm}
\item ¿Cuál de los siguientes números es racional?
\begin{enumerate}[label=\textbf{\Alph*)}]
\begin{multicols}{2}
  \item \sqrt{2}
  \item \pi
  \item e
  \item \sqrt{5}
  \item \(\frac{5}{6}\)
  \item \sqrt{11}
\end{multicols}
\end{enumerate}
\vspace{0.15cm}
\item ¿A qué conjunto pertenece el número \(2.5\)?
\begin{enumerate}[label=\textbf{\Alph*)}]
\begin{multicols}{2}
  \item Natural
  \item Entero
  \item Racional
  \item Irracional
  \item Complejo no real
  \item Imaginario
\end{multicols}
\end{enumerate}
\vspace{0.15cm}
\item ¿A qué conjunto pertenece el número \(-0.75\)?
\begin{enumerate}[label=\textbf{\Alph*)}]
\begin{multicols}{2}
  \item Natural
  \item Entero
  \item Racional
  \item Irracional
  \item Complejo no real
  \item Imaginario
\end{multicols}
\end{enumerate}
\vspace{0.15cm}
\item ¿Cuál de los siguientes números NO es real?
\begin{enumerate}[label=\textbf{\Alph*)}]
\begin{multicols}{2}
  \item \sqrt{2}
  \item \pi
  \item \sqrt{-1}
  \item -5
  \item 0
  \item \frac{3}{7}
\end{multicols}
\end{enumerate}
\vspace{0.15cm}
\item ¿Cuál de los siguientes números es irracional?
\begin{enumerate}[label=\textbf{\Alph*)}]
\begin{multicols}{2}
  \item \frac{2}{3}
  \item -\sqrt{16}
  \item \sqrt{7}
  \item 0.75
  \item 5
  \item -1
\end{multicols}
\end{enumerate}
\vspace{0.15cm}
\item ¿Cuál de los siguientes números es racional?
\begin{enumerate}[label=\textbf{\Alph*)}]
\begin{multicols}{2}
  \item \sqrt{2}
  \item \pi
  \item e
  \item \sqrt{5}
  \item \frac{7}{9}
  \item \(\sqrt{9}\)
\end{multicols}
\end{enumerate}
\vspace{0.15cm}
\item ¿A qué conjunto pertenece el número \(-\frac{8}{4}\)?
\begin{enumerate}[label=\textbf{\Alph*)}]
\begin{multicols}{2}
  \item Natural
  \item Entero
  \item Racional
  \item Irracional
  \item Complejo no real
  \item Imaginario
\end{multicols}
\end{enumerate}
\vspace{0.15cm}
\item ¿A qué conjunto pertenece el número \(\frac{0}{5}\)?
\begin{enumerate}[label=\textbf{\Alph*)}]
\begin{multicols}{2}
  \item Natural
  \item Entero
  \item Racional
  \item Irracional
  \item Complejo no real
  \item Imaginario
\end{multicols}
\end{enumerate}
\vspace{0.15cm}
\item ¿A qué conjunto pertenece el número \(\frac{5}{0}\)?
\begin{enumerate}[label=\textbf{\Alph*)}]
\begin{multicols}{2}
  \item Natural
  \item Entero
  \item Racional
  \item Irracional
  \item Complejo no real
  \item Imaginario
\end{multicols}
\end{enumerate}
\vspace{0.15cm}
\item ¿Cuál de los siguientes números es irracional?
\begin{enumerate}[label=\textbf{\Alph*)}]
\begin{multicols}{2}
  \item \frac{2}{3}
  \item -\sqrt{16}
  \item \(\sqrt{10}\)
  \item 0.75
  \item 5
  \item -1
\end{multicols}
\end{enumerate}
\vspace{0.15cm}
\item ¿Cuál de los siguientes números es racional?
\begin{enumerate}[label=\textbf{\Alph*)}]
\begin{multicols}{2}
  \item \sqrt{2}
  \item \pi
  \item e
  \item \sqrt{5}
  \item \(\frac{11}{3}\)
  \item \sqrt{11}
\end{multicols}
\end{enumerate}
\vspace{0.15cm}
\item ¿A qué conjunto pertenece el número \(0.121212...\)?
\begin{enumerate}[label=\textbf{\Alph*)}]
\begin{multicols}{2}
  \item Natural
  \item Entero
  \item Racional
  \item Irracional
  \item Complejo no real
  \item Imaginario
\end{multicols}
\end{enumerate}
\vspace{0.15cm}
\item ¿A qué conjunto pertenece el número \(0.1010010001...\)?
\begin{enumerate}[label=\textbf{\Alph*)}]
\begin{multicols}{2}
  \item Natural
  \item Entero
  \item Racional
  \item Irracional
  \item Complejo no real
  \item Imaginario
\end{multicols}
\end{enumerate}
\vspace{0.15cm}
\item ¿Cuál de los siguientes números es natural?
\begin{enumerate}[label=\textbf{\Alph*)}]
\begin{multicols}{2}
  \item Natural
  \item Entero
  \item Racional
  \item Irracional
  \item Complejo no real
  \item Imaginario
\end{multicols}
\end{enumerate}
\vspace{0.15cm}
\item ¿Cuál de los siguientes números NO es entero?
\begin{enumerate}[label=\textbf{\Alph*)}]
\begin{multicols}{2}
  \item Natural
  \item Entero
  \item Racional
  \item Irracional
  \item Complejo no real
  \item Imaginario
\end{multicols}
\end{enumerate}
\vspace{0.15cm}
\item ¿A qué conjunto pertenece el número \(\frac{10}{2}\)?
\begin{enumerate}[label=\textbf{\Alph*)}]
\begin{multicols}{2}
  \item Natural
  \item Entero
  \item Racional
  \item Irracional
  \item Complejo no real
  \item Imaginario
\end{multicols}
\end{enumerate}
\vspace{0.15cm}
\item ¿Cuál de los siguientes números es irracional?
\begin{enumerate}[label=\textbf{\Alph*)}]
\begin{multicols}{2}
  \item \frac{2}{3}
  \item -\sqrt{16}
  \item \(\sqrt{2}\)
  \item 0.75
  \item 5
  \item -1
\end{multicols}
\end{enumerate}
\vspace{0.15cm}
\item ¿Cuál de los siguientes números es racional?
\begin{enumerate}[label=\textbf{\Alph*)}]
\begin{multicols}{2}
  \item \sqrt{2}
  \item \pi
  \item e
  \item \sqrt{5}
  \item \frac{7}{9}
  \item \sqrt{11}
\end{multicols}
\end{enumerate}
\vspace{0.15cm}
\item ¿A qué conjunto pertenece el número \(-1.5\)?
\begin{enumerate}[label=\textbf{\Alph*)}]
\begin{multicols}{2}
  \item Natural
  \item Entero
  \item Racional
  \item Irracional
  \item Complejo no real
  \item Imaginario
\end{multicols}
\end{enumerate}
\vspace{0.15cm}
\item ¿A qué conjunto pertenece el número \(\sqrt{6}\)?
\begin{enumerate}[label=\textbf{\Alph*)}]
\begin{multicols}{2}
  \item Natural
  \item Entero
  \item Racional
  \item Irracional
  \item Complejo no real
  \item Imaginario
\end{multicols}
\end{enumerate}
\vspace{0.15cm}
\item ¿A qué conjunto pertenece el número \(\sqrt{1}\)?
\begin{enumerate}[label=\textbf{\Alph*)}]
\begin{multicols}{2}
  \item Natural
  \item Entero
  \item Racional
  \item Irracional
  \item Complejo no real
  \item Imaginario
\end{multicols}
\end{enumerate}
\vspace{0.15cm}
\item ¿A qué conjunto pertenece el número \(-8\)?
\begin{enumerate}[label=\textbf{\Alph*)}]
\begin{multicols}{2}
  \item Natural
  \item Entero
  \item Racional
  \item Irracional
  \item Complejo no real
  \item Imaginario
\end{multicols}
\end{enumerate}
\vspace{0.15cm}
\item ¿A qué conjunto pertenece el número \(\frac{3}{4}\)?
\begin{enumerate}[label=\textbf{\Alph*)}]
\begin{multicols}{2}
  \item Natural
  \item Entero
  \item Racional
  \item Irracional
  \item Complejo no real
  \item Imaginario
\end{multicols}
\end{enumerate}
\vspace{0.15cm}
\item ¿A qué conjunto pertenece el número \(\sqrt{0}\)?
\begin{enumerate}[label=\textbf{\Alph*)}]
\begin{multicols}{2}
  \item Natural
  \item Entero
  \item Racional
  \item Irracional
  \item Complejo no real
  \item Imaginario
\end{multicols}
\end{enumerate}
\vspace{0.15cm}
\item ¿Cuál de los siguientes números es irracional?
\begin{enumerate}[label=\textbf{\Alph*)}]
\begin{multicols}{2}
  \item \frac{2}{3}
  \item -\sqrt{16}
  \item \(\sqrt{15}\)
  \item 0.75
  \item 5
  \item -1
\end{multicols}
\end{enumerate}
\vspace{0.15cm}
\item ¿A qué conjunto pertenece el número \(-\frac{9}{3}\)?
\begin{enumerate}[label=\textbf{\Alph*)}]
\begin{multicols}{2}
  \item Natural
  \item Entero
  \item Racional
  \item Irracional
  \item Complejo no real
  \item Imaginario
\end{multicols}
\end{enumerate}
\vspace{0.15cm}
\item ¿Cuál de los siguientes números es natural?
\begin{enumerate}[label=\textbf{\Alph*)}]
\begin{multicols}{2}
  \item Natural
  \item Entero
  \item Racional
  \item Irracional
  \item Complejo no real
  \item Imaginario
\end{multicols}
\end{enumerate}
\vspace{0.15cm}
\item ¿Cuál de los siguientes números es irracional?
\begin{enumerate}[label=\textbf{\Alph*)}]
\begin{multicols}{2}
  \item \frac{2}{3}
  \item -\sqrt{16}
  \item \(\sqrt{17}\)
  \item 0.75
  \item 5
  \item -1
\end{multicols}
\end{enumerate}
\vspace{0.15cm}
\item ¿A qué conjunto pertenece el número \(\pi\)?
\begin{enumerate}[label=\textbf{\Alph*)}]
\begin{multicols}{2}
  \item Natural
  \item Entero
  \item Racional
  \item Irracional
  \item Complejo no real
  \item Imaginario
\end{multicols}
\end{enumerate}
\vspace{0.15cm}
\item ¿Cuál de los siguientes números es racional?
\begin{enumerate}[label=\textbf{\Alph*)}]
\begin{multicols}{2}
  \item \sqrt{2}
  \item \pi
  \item e
  \item \sqrt{5}
  \item \frac{7}{9}
  \item \sqrt{11}
\end{multicols}
\end{enumerate}
\vspace{0.15cm}
\item ¿A qué conjunto pertenece el número \(-0.333...\)?
\begin{enumerate}[label=\textbf{\Alph*)}]
\begin{multicols}{2}
  \item Natural
  \item Entero
  \item Racional
  \item Irracional
  \item Complejo no real
  \item Imaginario
\end{multicols}
\end{enumerate}
\vspace{0.15cm}
\item ¿Cuál de los siguientes números es irracional?
\begin{enumerate}[label=\textbf{\Alph*)}]
\begin{multicols}{2}
  \item \frac{2}{3}
  \item -\sqrt{16}
  \item \(\sqrt{19}\)
  \item 0.75
  \item 5
  \item -1
\end{multicols}
\end{enumerate}
\vspace{0.15cm}
\item ¿A qué conjunto pertenece el número \(\frac{7}{7}\)?
\begin{enumerate}[label=\textbf{\Alph*)}]
\begin{multicols}{2}
  \item Natural
  \item Entero
  \item Racional
  \item Irracional
  \item Complejo no real
  \item Imaginario
\end{multicols}
\end{enumerate}
\vspace{0.15cm}
\item ¿Cuál de los siguientes números es entero?
\begin{enumerate}[label=\textbf{\Alph*)}]
\begin{multicols}{2}
  \item Natural
  \item Entero
  \item Racional
  \item Irracional
  \item Complejo no real
  \item Imaginario
\end{multicols}
\end{enumerate}
\vspace{0.15cm}
\item ¿Cuál de los siguientes números es racional?
\begin{enumerate}[label=\textbf{\Alph*)}]
\begin{multicols}{2}
  \item \sqrt{2}
  \item \pi
  \item e
  \item \sqrt{5}
  \item \frac{7}{9}
  \item \sqrt{11}
\end{multicols}
\end{enumerate}
\vspace{0.15cm}
\item ¿A qué conjunto pertenece el número \(3.14\)?
\begin{enumerate}[label=\textbf{\Alph*)}]
\begin{multicols}{2}
  \item Natural
  \item Entero
  \item Racional
  \item Irracional
  \item Complejo no real
  \item Imaginario
\end{multicols}
\end{enumerate}
\vspace{0.15cm}
\item ¿A qué conjunto pertenece el número \(\pi \approx 3.14159...\)?
\begin{enumerate}[label=\textbf{\Alph*)}]
\begin{multicols}{2}
  \item Natural
  \item Entero
  \item Racional
  \item Irracional
  \item Complejo no real
  \item Imaginario
\end{multicols}
\end{enumerate}
\vspace{0.15cm}
\item ¿Cuál de los siguientes números es irracional?
\begin{enumerate}[label=\textbf{\Alph*)}]
\begin{multicols}{2}
  \item \frac{2}{3}
  \item -\sqrt{16}
  \item \(\sqrt{23}\)
  \item 0.75
  \item 5
  \item -1
\end{multicols}
\end{enumerate}
\vspace{0.15cm}
\item ¿A qué conjunto pertenece el número \(-\sqrt{36}\)?
\begin{enumerate}[label=\textbf{\Alph*)}]
\begin{multicols}{2}
  \item Natural
  \item Entero
  \item Racional
  \item Irracional
  \item Complejo no real
  \item Imaginario
\end{multicols}
\end{enumerate}
\vspace{0.15cm}
\item ¿Cuál de los siguientes números es racional?
\begin{enumerate}[label=\textbf{\Alph*)}]
\begin{multicols}{2}
  \item \sqrt{2}
  \item \pi
  \item e
  \item \sqrt{5}
  \item \frac{7}{9}
  \item \(\sqrt{16}\)
\end{multicols}
\end{enumerate}
\vspace{0.15cm}
\item ¿A qué conjunto pertenece el número \(0.666...\)?
\begin{enumerate}[label=\textbf{\Alph*)}]
\begin{multicols}{2}
  \item Natural
  \item Entero
  \item Racional
  \item Irracional
  \item Complejo no real
  \item Imaginario
\end{multicols}
\end{enumerate}
\vspace{0.15cm}
\item ¿Cuál de los siguientes números es irracional?
\begin{enumerate}[label=\textbf{\Alph*)}]
\begin{multicols}{2}
  \item \frac{2}{3}
  \item -\sqrt{16}
  \item \(\sqrt{29}\)
  \item 0.75
  \item 5
  \item -1
\end{multicols}
\end{enumerate}
\vspace{0.15cm}
\item ¿A qué conjunto pertenece el número \(\frac{12}{4}\)?
\begin{enumerate}[label=\textbf{\Alph*)}]
\begin{multicols}{2}
  \item Natural
  \item Entero
  \item Racional
  \item Irracional
  \item Complejo no real
  \item Imaginario
\end{multicols}
\end{enumerate}
\vspace{0.15cm}
\end{enumerate}\n
\subsubsection*{SECCIÓN II: PROPIEDADES DE LOS NÚMEROS REALES (61–110)}

\begin{enumerate}[label=\textbf{\arabic*.}]
\setcounter{enumi}{60}
\item ¿Qué propiedad se ilustra en \(3+5=5+3\)?
\begin{enumerate}[label=\textbf{\Alph*)}]
\begin{multicols}{2}
  \item Asociativa
  \item Conmutativa
  \item Distributiva
  \item Elemento neutro
  \item Elemento inverso
  \item Clausurativa
\end{multicols}
\end{enumerate}
\vspace{0.15cm}
\item ¿Qué propiedad se ilustra en \((2+3)+4=2+(3+4)\)?
\begin{enumerate}[label=\textbf{\Alph*)}]
\begin{multicols}{2}
  \item Asociativa
  \item Conmutativa
  \item Distributiva
  \item Elemento neutro
  \item Elemento inverso
  \item Clausurativa
\end{multicols}
\end{enumerate}
\vspace{0.15cm}
\item ¿Qué propiedad se ilustra en \(2(3+4)=2\cdot3+2\cdot4\)?
\begin{enumerate}[label=\textbf{\Alph*)}]
\begin{multicols}{2}
  \item Asociativa
  \item Conmutativa
  \item Distributiva
  \item Elemento neutro
  \item Elemento inverso
  \item Clausurativa
\end{multicols}
\end{enumerate}
\vspace{0.15cm}
\item ¿Qué propiedad se ilustra en \(7+0=7\)?
\begin{enumerate}[label=\textbf{\Alph*)}]
\begin{multicols}{2}
  \item Asociativa
  \item Conmutativa
  \item Distributiva
  \item Elemento neutro
  \item Elemento inverso
  \item Clausurativa
\end{multicols}
\end{enumerate}
\vspace{0.15cm}
\item ¿Qué propiedad se ilustra en \(5+(-5)=0\)?
\begin{enumerate}[label=\textbf{\Alph*)}]
\begin{multicols}{2}
  \item Asociativa
  \item Conmutativa
  \item Distributiva
  \item Elemento neutro
  \item Elemento inverso
  \item Clausurativa
\end{multicols}
\end{enumerate}
\vspace{0.15cm}
\item ¿Qué propiedad se ilustra en \(3\cdot1=3\)?
\begin{enumerate}[label=\textbf{\Alph*)}]
\begin{multicols}{2}
  \item Asociativa
  \item Conmutativa
  \item Distributiva
  \item Elemento neutro
  \item Elemento inverso
  \item Clausurativa
\end{multicols}
\end{enumerate}
\vspace{0.15cm}
\item ¿Qué propiedad se ilustra en \(4\cdot\frac{1}{4}=1\)?
\begin{enumerate}[label=\textbf{\Alph*)}]
\begin{multicols}{2}
  \item Asociativa
  \item Conmutativa
  \item Distributiva
  \item Elemento neutro
  \item Elemento inverso
  \item Clausurativa
\end{multicols}
\end{enumerate}
\vspace{0.15cm}
\item ¿Qué propiedad se ilustra en \(3\cdot5=5\cdot3\)?
\begin{enumerate}[label=\textbf{\Alph*)}]
\begin{multicols}{2}
  \item Asociativa
  \item Conmutativa
  \item Distributiva
  \item Elemento neutro
  \item Elemento inverso
  \item Clausurativa
\end{multicols}
\end{enumerate}
\vspace{0.15cm}
\item ¿Qué propiedad se ilustra en \((2\cdot3)\cdot4=2\cdot(3\cdot4)\)?
\begin{enumerate}[label=\textbf{\Alph*)}]
\begin{multicols}{2}
  \item Asociativa
  \item Conmutativa
  \item Distributiva
  \item Elemento neutro
  \item Elemento inverso
  \item Clausurativa
\end{multicols}
\end{enumerate}
\vspace{0.15cm}
\item ¿Qué propiedad se ilustra en \(5(2+3)=5\cdot2+5\cdot3\)?
\begin{enumerate}[label=\textbf{\Alph*)}]
\begin{multicols}{2}
  \item Asociativa
  \item Conmutativa
  \item Distributiva
  \item Elemento neutro
  \item Elemento inverso
  \item Clausurativa
\end{multicols}
\end{enumerate}
\vspace{0.15cm}
\item ¿Qué propiedad se ilustra en \(9+0=9\)?
\begin{enumerate}[label=\textbf{\Alph*)}]
\begin{multicols}{2}
  \item Asociativa
  \item Conmutativa
  \item Distributiva
  \item Elemento neutro
  \item Elemento inverso
  \item Clausurativa
\end{multicols}
\end{enumerate}
\vspace{0.15cm}
\item ¿Qué propiedad se ilustra en \((-3)+3=0\)?
\begin{enumerate}[label=\textbf{\Alph*)}]
\begin{multicols}{2}
  \item Asociativa
  \item Conmutativa
  \item Distributiva
  \item Elemento neutro
  \item Elemento inverso
  \item Clausurativa
\end{multicols}
\end{enumerate}
\vspace{0.15cm}
\item ¿Qué propiedad se ilustra en \(8\cdot1=8\)?
\begin{enumerate}[label=\textbf{\Alph*)}]
\begin{multicols}{2}
  \item Asociativa
  \item Conmutativa
  \item Distributiva
  \item Elemento neutro
  \item Elemento inverso
  \item Clausurativa
\end{multicols}
\end{enumerate}
\vspace{0.15cm}
\item ¿Qué propiedad se ilustra en \(6\cdot\frac{1}{6}=1\)?
\begin{enumerate}[label=\textbf{\Alph*)}]
\begin{multicols}{2}
  \item Asociativa
  \item Conmutativa
  \item Distributiva
  \item Elemento neutro
  \item Elemento inverso
  \item Clausurativa
\end{multicols}
\end{enumerate}
\vspace{0.15cm}
\item ¿Qué propiedad se ilustra en \(a+b=b+a\)?
\begin{enumerate}[label=\textbf{\Alph*)}]
\begin{multicols}{2}
  \item Asociativa
  \item Conmutativa
  \item Distributiva
  \item Elemento neutro
  \item Elemento inverso
  \item Clausurativa
\end{multicols}
\end{enumerate}
\vspace{0.15cm}
\item ¿Qué propiedad se ilustra en \((a+b)+c=a+(b+c)\)?
\begin{enumerate}[label=\textbf{\Alph*)}]
\begin{multicols}{2}
  \item Asociativa
  \item Conmutativa
  \item Distributiva
  \item Elemento neutro
  \item Elemento inverso
  \item Clausurativa
\end{multicols}
\end{enumerate}
\vspace{0.15cm}
\item ¿Qué propiedad se ilustra en \(a(b+c)=ab+ac\)?
\begin{enumerate}[label=\textbf{\Alph*)}]
\begin{multicols}{2}
  \item Asociativa
  \item Conmutativa
  \item Distributiva
  \item Elemento neutro
  \item Elemento inverso
  \item Clausurativa
\end{multicols}
\end{enumerate}
\vspace{0.15cm}
\item ¿Qué propiedad se ilustra en \(a+0=a\)?
\begin{enumerate}[label=\textbf{\Alph*)}]
\begin{multicols}{2}
  \item Asociativa
  \item Conmutativa
  \item Distributiva
  \item Elemento neutro
  \item Elemento inverso
  \item Clausurativa
\end{multicols}
\end{enumerate}
\vspace{0.15cm}
\item ¿Qué propiedad se ilustra en \(a+(-a)=0\)?
\begin{enumerate}[label=\textbf{\Alph*)}]
\begin{multicols}{2}
  \item Asociativa
  \item Conmutativa
  \item Distributiva
  \item Elemento neutro
  \item Elemento inverso
  \item Clausurativa
\end{multicols}
\end{enumerate}
\vspace{0.15cm}
\item ¿Qué propiedad se ilustra en \(a\cdot1=a\)?
\begin{enumerate}[label=\textbf{\Alph*)}]
\begin{multicols}{2}
  \item Asociativa
  \item Conmutativa
  \item Distributiva
  \item Elemento neutro
  \item Elemento inverso
  \item Clausurativa
\end{multicols}
\end{enumerate}
\vspace{0.15cm}
\item ¿Qué propiedad se ilustra en \(a\cdot\frac{1}{a}=1\)?
\begin{enumerate}[label=\textbf{\Alph*)}]
\begin{multicols}{2}
  \item Asociativa
  \item Conmutativa
  \item Distributiva
  \item Elemento neutro
  \item Elemento inverso
  \item Clausurativa
\end{multicols}
\end{enumerate}
\vspace{0.15cm}
\item ¿Qué propiedad se ilustra en \(a\cdot b=b\cdot a\)?
\begin{enumerate}[label=\textbf{\Alph*)}]
\begin{multicols}{2}
  \item Asociativa
  \item Conmutativa
  \item Distributiva
  \item Elemento neutro
  \item Elemento inverso
  \item Clausurativa
\end{multicols}
\end{enumerate}
\vspace{0.15cm}
\item ¿Qué propiedad se ilustra en \((a\cdot b)\cdot c=a\cdot(b\cdot c)\)?
\begin{enumerate}[label=\textbf{\Alph*)}]
\begin{multicols}{2}
  \item Asociativa
  \item Conmutativa
  \item Distributiva
  \item Elemento neutro
  \item Elemento inverso
  \item Clausurativa
\end{multicols}
\end{enumerate}
\vspace{0.15cm}
\item ¿Qué propiedad se ilustra en \(4(5-2)=4\cdot5-4\cdot2\)?
\begin{enumerate}[label=\textbf{\Alph*)}]
\begin{multicols}{2}
  \item Asociativa
  \item Conmutativa
  \item Distributiva
  \item Elemento neutro
  \item Elemento inverso
  \item Clausurativa
\end{multicols}
\end{enumerate}
\vspace{0.15cm}
\item ¿Qué propiedad se ilustra en \((-7)+0=-7\)?
\begin{enumerate}[label=\textbf{\Alph*)}]
\begin{multicols}{2}
  \item Asociativa
  \item Conmutativa
  \item Distributiva
  \item Elemento neutro
  \item Elemento inverso
  \item Clausurativa
\end{multicols}
\end{enumerate}
\vspace{0.15cm}
\item ¿Qué propiedad se ilustra en \((-5)+5=0\)?
\begin{enumerate}[label=\textbf{\Alph*)}]
\begin{multicols}{2}
  \item Asociativa
  \item Conmutativa
  \item Distributiva
  \item Elemento neutro
  \item Elemento inverso
  \item Clausurativa
\end{multicols}
\end{enumerate}
\vspace{0.15cm}
\item ¿Qué propiedad se ilustra en \((-1)\cdot1=-1\)?
\begin{enumerate}[label=\textbf{\Alph*)}]
\begin{multicols}{2}
  \item Asociativa
  \item Conmutativa
  \item Distributiva
  \item Elemento neutro
  \item Elemento inverso
  \item Clausurativa
\end{multicols}
\end{enumerate}
\vspace{0.15cm}
\item ¿Qué propiedad se ilustra en \((-3)\cdot(-\frac{1}{3})=1\)?
\begin{enumerate}[label=\textbf{\Alph*)}]
\begin{multicols}{2}
  \item Asociativa
  \item Conmutativa
  \item Distributiva
  \item Elemento neutro
  \item Elemento inverso
  \item Clausurativa
\end{multicols}
\end{enumerate}
\vspace{0.15cm}
\item ¿Qué propiedad se ilustra en \(2\cdot(3\cdot4)=(2\cdot3)\cdot4\)?
\begin{enumerate}[label=\textbf{\Alph*)}]
\begin{multicols}{2}
  \item Asociativa
  \item Conmutativa
  \item Distributiva
  \item Elemento neutro
  \item Elemento inverso
  \item Clausurativa
\end{multicols}
\end{enumerate}
\vspace{0.15cm}
\item ¿Qué propiedad se ilustra en \(6+3=3+6\)?
\begin{enumerate}[label=\textbf{\Alph*)}]
\begin{multicols}{2}
  \item Asociativa
  \item Conmutativa
  \item Distributiva
  \item Elemento neutro
  \item Elemento inverso
  \item Clausurativa
\end{multicols}
\end{enumerate}
\vspace{0.15cm}
\item ¿Qué propiedad se ilustra en \(10(1+2)=10\cdot1+10\cdot2\)?
\begin{enumerate}[label=\textbf{\Alph*)}]
\begin{multicols}{2}
  \item Asociativa
  \item Conmutativa
  \item Distributiva
  \item Elemento neutro
  \item Elemento inverso
  \item Clausurativa
\end{multicols}
\end{enumerate}
\vspace{0.15cm}
\item ¿Qué propiedad se ilustra en \(0+8=8\)?
\begin{enumerate}[label=\textbf{\Alph*)}]
\begin{multicols}{2}
  \item Asociativa
  \item Conmutativa
  \item Distributiva
  \item Elemento neutro
  \item Elemento inverso
  \item Clausurativa
\end{multicols}
\end{enumerate}
\vspace{0.15cm}
\item ¿Qué propiedad se ilustra en \(4+(-4)=0\)?
\begin{enumerate}[label=\textbf{\Alph*)}]
\begin{multicols}{2}
  \item Asociativa
  \item Conmutativa
  \item Distributiva
  \item Elemento neutro
  \item Elemento inverso
  \item Clausurativa
\end{multicols}
\end{enumerate}
\vspace{0.15cm}
\item ¿Qué propiedad se ilustra en \(9\cdot1=9\)?
\begin{enumerate}[label=\textbf{\Alph*)}]
\begin{multicols}{2}
  \item Asociativa
  \item Conmutativa
  \item Distributiva
  \item Elemento neutro
  \item Elemento inverso
  \item Clausurativa
\end{multicols}
\end{enumerate}
\vspace{0.15cm}
\item ¿Qué propiedad se ilustra en \(\frac{2}{3}\cdot\frac{3}{2}=1\)?
\begin{enumerate}[label=\textbf{\Alph*)}]
\begin{multicols}{2}
  \item Asociativa
  \item Conmutativa
  \item Distributiva
  \item Elemento neutro
  \item Elemento inverso
  \item Clausurativa
\end{multicols}
\end{enumerate}
\vspace{0.15cm}
\item ¿Qué propiedad se ilustra en \(7\cdot3=3\cdot7\)?
\begin{enumerate}[label=\textbf{\Alph*)}]
\begin{multicols}{2}
  \item Asociativa
  \item Conmutativa
  \item Distributiva
  \item Elemento neutro
  \item Elemento inverso
  \item Clausurativa
\end{multicols}
\end{enumerate}
\vspace{0.15cm}
\item ¿Qué propiedad se ilustra en \((1+2)+3=1+(2+3)\)?
\begin{enumerate}[label=\textbf{\Alph*)}]
\begin{multicols}{2}
  \item Asociativa
  \item Conmutativa
  \item Distributiva
  \item Elemento neutro
  \item Elemento inverso
  \item Clausurativa
\end{multicols}
\end{enumerate}
\vspace{0.15cm}
\item ¿Qué propiedad se ilustra en \(3(4+5)=3\cdot4+3\cdot5\)?
\begin{enumerate}[label=\textbf{\Alph*)}]
\begin{multicols}{2}
  \item Asociativa
  \item Conmutativa
  \item Distributiva
  \item Elemento neutro
  \item Elemento inverso
  \item Clausurativa
\end{multicols}
\end{enumerate}
\vspace{0.15cm}
\item ¿Qué propiedad se ilustra en \(-6+0=-6\)?
\begin{enumerate}[label=\textbf{\Alph*)}]
\begin{multicols}{2}
  \item Asociativa
  \item Conmutativa
  \item Distributiva
  \item Elemento neutro
  \item Elemento inverso
  \item Clausurativa
\end{multicols}
\end{enumerate}
\vspace{0.15cm}
\item ¿Qué propiedad se ilustra en \(-8+8=0\)?
\begin{enumerate}[label=\textbf{\Alph*)}]
\begin{multicols}{2}
  \item Asociativa
  \item Conmutativa
  \item Distributiva
  \item Elemento neutro
  \item Elemento inverso
  \item Clausurativa
\end{multicols}
\end{enumerate}
\vspace{0.15cm}
\item La suma de dos reales es real. Esto se conoce como:
\begin{enumerate}[label=\textbf{\Alph*)}]
\begin{multicols}{2}
  \item Asociativa
  \item Conmutativa
  \item Distributiva
  \item Elemento neutro
  \item Elemento inverso
  \item Clausurativa
\end{multicols}
\end{enumerate}
\vspace{0.15cm}
\item El producto de dos reales es real. Esto se conoce como:
\begin{enumerate}[label=\textbf{\Alph*)}]
\begin{multicols}{2}
  \item Asociativa
  \item Conmutativa
  \item Distributiva
  \item Elemento neutro
  \item Elemento inverso
  \item Clausurativa
\end{multicols}
\end{enumerate}
\vspace{0.15cm}
\item ¿Qué propiedad permite reordenar sumandos sin alterar?
\begin{enumerate}[label=\textbf{\Alph*)}]
\begin{multicols}{2}
  \item Asociativa
  \item Conmutativa
  \item Distributiva
  \item Elemento neutro
  \item Elemento inverso
  \item Clausurativa
\end{multicols}
\end{enumerate}
\vspace{0.15cm}
\item ¿Qué propiedad permite reagrupar sumandos sin alterar?
\begin{enumerate}[label=\textbf{\Alph*)}]
\begin{multicols}{2}
  \item Asociativa
  \item Conmutativa
  \item Distributiva
  \item Elemento neutro
  \item Elemento inverso
  \item Clausurativa
\end{multicols}
\end{enumerate}
\vspace{0.15cm}
\item ¿Qué propiedad permite multiplicar factor por suma?
\begin{enumerate}[label=\textbf{\Alph*)}]
\begin{multicols}{2}
  \item Asociativa
  \item Conmutativa
  \item Distributiva
  \item Elemento neutro
  \item Elemento inverso
  \item Clausurativa
\end{multicols}
\end{enumerate}
\vspace{0.15cm}
\item El número \(0\) es el elemento neutro de la:
\begin{enumerate}[label=\textbf{\Alph*)}]
\begin{multicols}{2}
  \item Resta
  \item Suma
  \item Multiplicación
  \item División
  \item Potenciación
  \item Ninguna
\end{multicols}
\end{enumerate}
\vspace{0.15cm}
\item El número \(1\) es el elemento neutro de la:
\begin{enumerate}[label=\textbf{\Alph*)}]
\begin{multicols}{2}
  \item Suma
  \item Multiplicación
  \item División
  \item Potenciación
  \item Potenciación
  \item Ninguna
\end{multicols}
\end{enumerate}
\vspace{0.15cm}
\item El inverso aditivo de \(7\) es:
\begin{enumerate}[label=\textbf{\Alph*)}]
\begin{multicols}{2}
  \item 7
  \item -7
  \item \frac{1}{7}
  \item -\frac{1}{7}
  \item 0
  \item 1
\end{multicols}
\end{enumerate}
\vspace{0.15cm}
\item El inverso multiplicativo de \(5\) es:
\begin{enumerate}[label=\textbf{\Alph*)}]
\begin{multicols}{2}
  \item 5
  \item -5
  \item \frac{1}{5}
  \item -\frac{1}{5}
  \item 0
  \item 1
\end{multicols}
\end{enumerate}
\vspace{0.15cm}
\item El inverso aditivo de \(-3\) es:
\begin{enumerate}[label=\textbf{\Alph*)}]
\begin{multicols}{2}
  \item -3
  \item 3
  \item \frac{1}{3}
  \item -\frac{1}{3}
  \item 0
  \item 1
\end{multicols}
\end{enumerate}
\vspace{0.15cm}
\end{enumerate}\n
\subsubsection*{SECCIÓN III: EXPANSIÓN DECIMAL (111–150)}

\begin{enumerate}[label=\textbf{\arabic*.}]
\setcounter{enumi}{110}
\item ¿Qué tipo de expansión decimal tiene \(\frac{1}{2}\)?
\begin{enumerate}[label=\textbf{\Alph*)}]
\begin{multicols}{2}
  \item Terminante
  \item Periódica pura
  \item Periódica mixta
  \item Irracional
  \item No real
  \item Ninguna
\end{multicols}
\end{enumerate}
\vspace{0.15cm}
\item ¿Qué tipo de expansión decimal tiene \(\frac{1}{3}\)?
\begin{enumerate}[label=\textbf{\Alph*)}]
\begin{multicols}{2}
  \item Terminante
  \item Periódica pura
  \item Periódica mixta
  \item Irracional
  \item No real
  \item Ninguna
\end{multicols}
\end{enumerate}
\vspace{0.15cm}
\item ¿Qué tipo de expansión decimal tiene \(\frac{1}{6}\)?
\begin{enumerate}[label=\textbf{\Alph*)}]
\begin{multicols}{2}
  \item Terminante
  \item Periódica pura
  \item Periódica mixta
  \item Irracional
  \item No real
  \item Ninguna
\end{multicols}
\end{enumerate}
\vspace{0.15cm}
\item ¿Qué tipo de expansión decimal tiene \(\frac{1}{4}\)?
\begin{enumerate}[label=\textbf{\Alph*)}]
\begin{multicols}{2}
  \item Terminante
  \item Periódica pura
  \item Periódica mixta
  \item Irracional
  \item No real
  \item Ninguna
\end{multicols}
\end{enumerate}
\vspace{0.15cm}
\item ¿Qué tipo de expansión decimal tiene \(\frac{1}{7}\)?
\begin{enumerate}[label=\textbf{\Alph*)}]
\begin{multicols}{2}
  \item Terminante
  \item Periódica pura
  \item Periódica mixta
  \item Irracional
  \item No real
  \item Ninguna
\end{multicols}
\end{enumerate}
\vspace{0.15cm}
\item ¿Qué tipo de expansión decimal tiene \(\frac{1}{8}\)?
\begin{enumerate}[label=\textbf{\Alph*)}]
\begin{multicols}{2}
  \item Terminante
  \item Periódica pura
  \item Periódica mixta
  \item Irracional
  \item No real
  \item Ninguna
\end{multicols}
\end{enumerate}
\vspace{0.15cm}
\item ¿Qué tipo de expansión decimal tiene \(\frac{1}{9}\)?
\begin{enumerate}[label=\textbf{\Alph*)}]
\begin{multicols}{2}
  \item Terminante
  \item Periódica pura
  \item Periódica mixta
  \item Irracional
  \item No real
  \item Ninguna
\end{multicols}
\end{enumerate}
\vspace{0.15cm}
\item ¿Qué tipo de expansión decimal tiene \(\frac{1}{11}\)?
\begin{enumerate}[label=\textbf{\Alph*)}]
\begin{multicols}{2}
  \item Terminante
  \item Periódica pura
  \item Periódica mixta
  \item Irracional
  \item No real
  \item Ninguna
\end{multicols}
\end{enumerate}
\vspace{0.15cm}
\item ¿Qué tipo de expansión decimal tiene \(\frac{1}{12}\)?
\begin{enumerate}[label=\textbf{\Alph*)}]
\begin{multicols}{2}
  \item Terminante
  \item Periódica pura
  \item Periódica mixta
  \item Irracional
  \item No real
  \item Ninguna
\end{multicols}
\end{enumerate}
\vspace{0.15cm}
\item ¿Qué tipo de expansión decimal tiene \(\frac{1}{20}\)?
\begin{enumerate}[label=\textbf{\Alph*)}]
\begin{multicols}{2}
  \item Terminante
  \item Periódica pura
  \item Periódica mixta
  \item Irracional
  \item No real
  \item Ninguna
\end{multicols}
\end{enumerate}
\vspace{0.15cm}
\item La expansión decimal de \(\frac{1}{2}\) es:
\begin{enumerate}[label=\textbf{\Alph*)}]
\begin{multicols}{2}
  \item 0.5
  \item 0.333...
  \item 0.1666...
  \item 0.25
  \item 0.75
  \item 0.125
\end{multicols}
\end{enumerate}
\vspace{0.15cm}
\item La expansión decimal de \(\frac{1}{3}\) es:
\begin{enumerate}[label=\textbf{\Alph*)}]
\begin{multicols}{2}
  \item 0.1
  \item 0.333...
  \item 0.1666...
  \item 0.25
  \item 0.75
  \item 0.125
\end{multicols}
\end{enumerate}
\vspace{0.15cm}
\item La expansión decimal de \(\frac{1}{4}\) es:
\begin{enumerate}[label=\textbf{\Alph*)}]
\begin{multicols}{2}
  \item 0.1
  \item 0.333...
  \item 0.1666...
  \item 0.25
  \item 0.75
  \item 0.125
\end{multicols}
\end{enumerate}
\vspace{0.15cm}
\item La expansión decimal de \(\frac{1}{5}\) es:
\begin{enumerate}[label=\textbf{\Alph*)}]
\begin{multicols}{2}
  \item 0.1
  \item 0.2
  \item 0.1666...
  \item 0.25
  \item 0.75
  \item 0.125
\end{multicols}
\end{enumerate}
\vspace{0.15cm}
\item La expansión decimal de \(\frac{1}{8}\) es:
\begin{enumerate}[label=\textbf{\Alph*)}]
\begin{multicols}{2}
  \item 0.1
  \item 0.333...
  \item 0.1666...
  \item 0.25
  \item 0.75
  \item 0.125
\end{multicols}
\end{enumerate}
\vspace{0.15cm}
\item La expansión decimal de \(\frac{1}{10}\) es:
\begin{enumerate}[label=\textbf{\Alph*)}]
\begin{multicols}{2}
  \item 0.1
  \item 0.333...
  \item 0.1666...
  \item 0.25
  \item 0.75
  \item 0.125
\end{multicols}
\end{enumerate}
\vspace{0.15cm}
\item La expansión decimal de \(\frac{3}{4}\) es:
\begin{enumerate}[label=\textbf{\Alph*)}]
\begin{multicols}{2}
  \item 0.1
  \item 0.75
  \item 0.1666...
  \item 0.25
  \item 0.75
  \item 0.125
\end{multicols}
\end{enumerate}
\vspace{0.15cm}
\item La expansión decimal de \(\frac{2}{5}\) es:
\begin{enumerate}[label=\textbf{\Alph*)}]
\begin{multicols}{2}
  \item 0.1
  \item 0.4
  \item 0.1666...
  \item 0.25
  \item 0.75
  \item 0.125
\end{multicols}
\end{enumerate}
\vspace{0.15cm}
\item La expansión decimal de \(\frac{7}{10}\) es:
\begin{enumerate}[label=\textbf{\Alph*)}]
\begin{multicols}{2}
  \item 0.7
  \item 0.333...
  \item 0.1666...
  \item 0.25
  \item 0.75
  \item 0.125
\end{multicols}
\end{enumerate}
\vspace{0.15cm}
\item La expansión decimal de \(\frac{5}{8}\) es:
\begin{enumerate}[label=\textbf{\Alph*)}]
\begin{multicols}{2}
  \item 0.1
  \item 0.625
  \item 0.1666...
  \item 0.25
  \item 0.75
  \item 0.125
\end{multicols}
\end{enumerate}
\vspace{0.15cm}
\item La fracción generatriz de \(0.5\) es:
\begin{enumerate}[label=\textbf{\Alph*)}]
\begin{multicols}{2}
  \item \frac{1}{5}
  \item \frac{1}{2}
  \item \frac{5}{10}
  \item \frac{2}{5}
  \item \frac{1}{4}
  \item \frac{5}{100}
\end{multicols}
\end{enumerate}
\vspace{0.15cm}
\item La fracción generatriz de \(0.25\) es:
\begin{enumerate}[label=\textbf{\Alph*)}]
\begin{multicols}{2}
  \item \frac{1}{5}
  \item \frac{1}{4}
  \item \frac{5}{10}
  \item \frac{2}{5}
  \item \frac{1}{4}
  \item \frac{5}{100}
\end{multicols}
\end{enumerate}
\vspace{0.15cm}
\item La fracción generatriz de \(0.75\) es:
\begin{enumerate}[label=\textbf{\Alph*)}]
\begin{multicols}{2}
  \item \frac{1}{5}
  \item \frac{3}{4}
  \item \frac{5}{10}
  \item \frac{2}{5}
  \item \frac{1}{4}
  \item \frac{5}{100}
\end{multicols}
\end{enumerate}
\vspace{0.15cm}
\item La fracción generatriz de \(0.2\) es:
\begin{enumerate}[label=\textbf{\Alph*)}]
\begin{multicols}{2}
  \item \frac{1}{5}
  \item \frac{1}{5}
  \item \frac{5}{10}
  \item \frac{2}{5}
  \item \frac{1}{4}
  \item \frac{5}{100}
\end{multicols}
\end{enumerate}
\vspace{0.15cm}
\item La fracción generatriz de \(0.125\) es:
\begin{enumerate}[label=\textbf{\Alph*)}]
\begin{multicols}{2}
  \item \frac{1}{5}
  \item \frac{1}{2}
  \item \frac{1}{8}
  \item \frac{2}{5}
  \item \frac{1}{4}
  \item \frac{5}{100}
\end{multicols}
\end{enumerate}
\vspace{0.15cm}
\item La fracción generatriz de \(0.333...\) es:
\begin{enumerate}[label=\textbf{\Alph*)}]
\begin{multicols}{2}
  \item \frac{1}{3}
  \item \frac{1}{2}
  \item \frac{5}{10}
  \item \frac{2}{5}
  \item \frac{1}{4}
  \item \frac{5}{100}
\end{multicols}
\end{enumerate}
\vspace{0.15cm}
\item La fracción generatriz de \(0.666...\) es:
\begin{enumerate}[label=\textbf{\Alph*)}]
\begin{multicols}{2}
  \item \frac{1}{5}
  \item \frac{2}{3}
  \item \frac{5}{10}
  \item \frac{2}{5}
  \item \frac{1}{4}
  \item \frac{5}{100}
\end{multicols}
\end{enumerate}
\vspace{0.15cm}
\item La fracción generatriz de \(0.111...\) es:
\begin{enumerate}[label=\textbf{\Alph*)}]
\begin{multicols}{2}
  \item \frac{1}{9}
  \item \frac{1}{2}
  \item \frac{5}{10}
  \item \frac{2}{5}
  \item \frac{1}{4}
  \item \frac{5}{100}
\end{multicols}
\end{enumerate}
\vspace{0.15cm}
\item La fracción generatriz de \(0.0909...\) es:
\begin{enumerate}[label=\textbf{\Alph*)}]
\begin{multicols}{2}
  \item \frac{1}{5}
  \item \frac{1}{11}
  \item \frac{5}{10}
  \item \frac{2}{5}
  \item \frac{1}{4}
  \item \frac{5}{100}
\end{multicols}
\end{enumerate}
\vspace{0.15cm}
\item La fracción generatriz de \(0.142857...\) es:
\begin{enumerate}[label=\textbf{\Alph*)}]
\begin{multicols}{2}
  \item \frac{1}{7}
  \item \frac{1}{2}
  \item \frac{5}{10}
  \item \frac{2}{5}
  \item \frac{1}{4}
  \item \frac{5}{100}
\end{multicols}
\end{enumerate}
\vspace{0.15cm}
\item ¿Cuál de las siguientes fracciones tiene expansión decimal terminante?
\begin{enumerate}[label=\textbf{\Alph*)}]
\begin{multicols}{2}
  \item \frac{1}{3}
  \item \frac{1}{6}
  \item \frac{1}{7}
  \item \frac{1}{8}
  \item \frac{1}{20}
  \item \frac{1}{11}
\end{multicols}
\end{enumerate}
\vspace{0.15cm}
\item ¿Cuál de las siguientes fracciones tiene expansión decimal periódica pura?
\begin{enumerate}[label=\textbf{\Alph*)}]
\begin{multicols}{2}
  \item \frac{1}{2}
  \item \frac{1}{4}
  \item \frac{1}{3}
  \item \frac{1}{6}
  \item \frac{1}{8}
  \item \frac{1}{11}
\end{multicols}
\end{enumerate}
\vspace{0.15cm}
\item ¿Cuál de las siguientes fracciones tiene expansión decimal periódica mixta?
\begin{enumerate}[label=\textbf{\Alph*)}]
\begin{multicols}{2}
  \item \frac{1}{2}
  \item \frac{1}{3}
  \item \frac{1}{4}
  \item \frac{1}{6}
  \item \frac{1}{8}
  \item \frac{1}{9}
\end{multicols}
\end{enumerate}
\vspace{0.15cm}
\item ¿Cuál de las siguientes fracciones tiene expansión decimal terminante?
\begin{enumerate}[label=\textbf{\Alph*)}]
\begin{multicols}{2}
  \item \frac{1}{3}
  \item \frac{1}{6}
  \item \frac{1}{7}
  \item \frac{1}{8}
  \item \frac{1}{20}
  \item \frac{1}{11}
\end{multicols}
\end{enumerate}
\vspace{0.15cm}
\item ¿Cuál de las siguientes fracciones tiene expansión decimal periódica pura?
\begin{enumerate}[label=\textbf{\Alph*)}]
\begin{multicols}{2}
  \item \frac{1}{2}
  \item \frac{1}{4}
  \item \frac{1}{3}
  \item \frac{1}{6}
  \item \frac{1}{8}
  \item \frac{1}{11}
\end{multicols}
\end{enumerate}
\vspace{0.15cm}
\item ¿Cuál de las siguientes fracciones tiene expansión decimal periódica mixta?
\begin{enumerate}[label=\textbf{\Alph*)}]
\begin{multicols}{2}
  \item \frac{1}{2}
  \item \frac{1}{3}
  \item \frac{1}{4}
  \item \frac{1}{12}
  \item \frac{1}{8}
  \item \frac{1}{9}
\end{multicols}
\end{enumerate}
\vspace{0.15cm}
\item La expansión decimal de \(\frac{2}{3}\) es:
\begin{enumerate}[label=\textbf{\Alph*)}]
\begin{multicols}{2}
  \item 0.5
  \item 0.666...
  \item 0.333...
  \item 0.25
  \item 0.75
  \item 0.2
\end{multicols}
\end{enumerate}
\vspace{0.15cm}
\item La expansión decimal de \(\frac{3}{5}\) es:
\begin{enumerate}[label=\textbf{\Alph*)}]
\begin{multicols}{2}
  \item 0.5
  \item 0.666...
  \item 0.6
  \item 0.25
  \item 0.75
  \item 0.2
\end{multicols}
\end{enumerate}
\vspace{0.15cm}
\item La expansión decimal de \(\frac{7}{8}\) es:
\begin{enumerate}[label=\textbf{\Alph*)}]
\begin{multicols}{2}
  \item 0.5
  \item 0.666...
  \item 0.875
  \item 0.25
  \item 0.75
  \item 0.2
\end{multicols}
\end{enumerate}
\vspace{0.15cm}
\item La expansión decimal de \(\frac{5}{6}\) es:
\begin{enumerate}[label=\textbf{\Alph*)}]
\begin{multicols}{2}
  \item 0.5
  \item 0.8333...
  \item 0.333...
  \item 0.25
  \item 0.75
  \item 0.2
\end{multicols}
\end{enumerate}
\vspace{0.15cm}
\end{enumerate}\n
\subsubsection*{SECCIÓN IV: ORDEN Y COMPARACIÓN (151–190)}

\begin{enumerate}[label=\textbf{\arabic*.}]
\setcounter{enumi}{150}
\item ¿Cuál es el símbolo correcto para comparar \(-3\) y \(2\)?
\begin{enumerate}[label=\textbf{\Alph*)}]
\begin{multicols}{2}
  \item <
  \item >
  \item =
  \item \leq
  \item \geq
  \item \neq
\end{multicols}
\end{enumerate}
\vspace{0.15cm}
\item ¿Cuál es el símbolo correcto para comparar \(5\) y \(-7\)?
\begin{enumerate}[label=\textbf{\Alph*)}]
\begin{multicols}{2}
  \item <
  \item >
  \item =
  \item \leq
  \item \geq
  \item \neq
\end{multicols}
\end{enumerate}
\vspace{0.15cm}
\item ¿Cuál es el símbolo correcto para comparar \(-4\) y \(-1\)?
\begin{enumerate}[label=\textbf{\Alph*)}]
\begin{multicols}{2}
  \item <
  \item >
  \item =
  \item \leq
  \item \geq
  \item \neq
\end{multicols}
\end{enumerate}
\vspace{0.15cm}
\item ¿Cuál es el símbolo correcto para comparar \(0\) y \(-2\)?
\begin{enumerate}[label=\textbf{\Alph*)}]
\begin{multicols}{2}
  \item <
  \item >
  \item =
  \item \leq
  \item \geq
  \item \neq
\end{multicols}
\end{enumerate}
\vspace{0.15cm}
\item ¿Cuál es el símbolo correcto para comparar \(-5\) y \(-5\)?
\begin{enumerate}[label=\textbf{\Alph*)}]
\begin{multicols}{2}
  \item <
  \item >
  \item =
  \item \leq
  \item \geq
  \item \neq
\end{multicols}
\end{enumerate}
\vspace{0.15cm}
\item ¿Cuál es el símbolo correcto para comparar \(\frac{1}{2}\) y \(0.5\)?
\begin{enumerate}[label=\textbf{\Alph*)}]
\begin{multicols}{2}
  \item <
  \item >
  \item =
  \item \leq
  \item \geq
  \item \neq
\end{multicols}
\end{enumerate}
\vspace{0.15cm}
\item ¿Cuál es el símbolo correcto para comparar \(\frac{1}{3}\) y \(0.33\)?
\begin{enumerate}[label=\textbf{\Alph*)}]
\begin{multicols}{2}
  \item <
  \item >
  \item =
  \item \leq
  \item \geq
  \item \neq
\end{multicols}
\end{enumerate}
\vspace{0.15cm}
\item ¿Cuál es el símbolo correcto para comparar \(\sqrt{2}\) y \(1.5\)?
\begin{enumerate}[label=\textbf{\Alph*)}]
\begin{multicols}{2}
  \item <
  \item >
  \item =
  \item \leq
  \item \geq
  \item \neq
\end{multicols}
\end{enumerate}
\vspace{0.15cm}
\item ¿Cuál es el símbolo correcto para comparar \(\pi\) y \(3.14\)?
\begin{enumerate}[label=\textbf{\Alph*)}]
\begin{multicols}{2}
  \item <
  \item >
  \item =
  \item \leq
  \item \geq
  \item \neq
\end{multicols}
\end{enumerate}
\vspace{0.15cm}
\item ¿Cuál es el símbolo correcto para comparar \(-0.5\) y \(-\frac{1}{2}\)?
\begin{enumerate}[label=\textbf{\Alph*)}]
\begin{multicols}{2}
  \item <
  \item >
  \item =
  \item \leq
  \item \geq
  \item \neq
\end{multicols}
\end{enumerate}
\vspace{0.15cm}
\item Ordena de menor a mayor: -3, 0, 2, -5
\begin{enumerate}[label=\textbf{\Alph*)}]
\begin{multicols}{2}
  \item -5, -3, 0, 2
  \item 1
  \item 2
  \item 3
  \item -1
  \item -2
\end{multicols}
\end{enumerate}
\vspace{0.15cm}
\item Ordena de menor a mayor: \frac{1}{2}, \frac{1}{4}, \frac{3}{4}, \frac{1}{8}
\begin{enumerate}[label=\textbf{\Alph*)}]
\begin{multicols}{2}
  \item \frac{1}{8}, \frac{1}{4}, \frac{1}{2}, \frac{3}{4}
  \item 1
  \item 2
  \item 3
  \item -1
  \item -2
\end{multicols}
\end{enumerate}
\vspace{0.15cm}
\item Ordena de menor a mayor: -1, -0.5, 0, 0.5, 1
\begin{enumerate}[label=\textbf{\Alph*)}]
\begin{multicols}{2}
  \item 0
  \item -1, -0.5, 0, 0.5, 1
  \item 2
  \item 3
  \item -1
  \item -2
\end{multicols}
\end{enumerate}
\vspace{0.15cm}
\item Ordena de mayor a menor: 3, -2, 5, -7
\begin{enumerate}[label=\textbf{\Alph*)}]
\begin{multicols}{2}
  \item 5, 3, -2, -7
  \item 1
  \item 2
  \item 3
  \item -1
  \item -2
\end{multicols}
\end{enumerate}
\vspace{0.15cm}
\item ¿Cuál de los siguientes números es el mayor?
\begin{enumerate}[label=\textbf{\Alph*)}]
\begin{multicols}{2}
  \item -5
  \item -1
  \item 0
  \item 2
  \item -3
  \item -4
\end{multicols}
\end{enumerate}
\vspace{0.15cm}
\item ¿Cuál de los siguientes números es el menor?
\begin{enumerate}[label=\textbf{\Alph*)}]
\begin{multicols}{2}
  \item -5
  \item -1
  \item 0
  \item 2
  \item -3
  \item -4
\end{multicols}
\end{enumerate}
\vspace{0.15cm}
\item ¿Cuál de los siguientes números es el mayor?
\begin{enumerate}[label=\textbf{\Alph*)}]
\begin{multicols}{2}
  \item \frac{1}{2}
  \item \frac{1}{3}
  \item \frac{1}{4}
  \item \frac{1}{5}
  \item \frac{1}{6}
  \item \frac{1}{7}
\end{multicols}
\end{enumerate}
\vspace{0.15cm}
\item ¿Cuál de los siguientes números es el menor?
\begin{enumerate}[label=\textbf{\Alph*)}]
\begin{multicols}{2}
  \item \frac{1}{2}
  \item \frac{1}{3}
  \item \frac{1}{4}
  \item \frac{1}{5}
  \item \frac{1}{6}
  \item \frac{1}{7}
\end{multicols}
\end{enumerate}
\vspace{0.15cm}
\item ¿Cuál de los siguientes números es el mayor?
\begin{enumerate}[label=\textbf{\Alph*)}]
\begin{multicols}{2}
  \item 0.5
  \item 0.05
  \item 0.005
  \item 0.555
  \item 0.55
  \item 0.505
\end{multicols}
\end{enumerate}
\vspace{0.15cm}
\item ¿Cuál de los siguientes números es el menor?
\begin{enumerate}[label=\textbf{\Alph*)}]
\begin{multicols}{2}
  \item 0.5
  \item 0.05
  \item 0.005
  \item 0.555
  \item 0.55
  \item 0.505
\end{multicols}
\end{enumerate}
\vspace{0.15cm}
\item Si a < b y b < c, entonces:
\begin{enumerate}[label=\textbf{\Alph*)}]
\begin{multicols}{2}
  \item a > c
  \item a < c
  \item Igual a
  \item No se puede determinar
  \item Siempre positivo
  \item Siempre negativo
\end{multicols}
\end{enumerate}
\vspace{0.15cm}
\item Si a < b, entonces a + c es:
\begin{enumerate}[label=\textbf{\Alph*)}]
\begin{multicols}{2}
  \item Mayor que
  \item Menor que
  \item Igual a
  \item No se puede determinar
  \item Siempre positivo
  \item Siempre negativo
\end{multicols}
\end{enumerate}
\vspace{0.15cm}
\item Si a < b y c > 0, entonces ac es:
\begin{enumerate}[label=\textbf{\Alph*)}]
\begin{multicols}{2}
  \item Mayor que
  \item Menor que
  \item Igual a
  \item No se puede determinar
  \item Siempre positivo
  \item Siempre negativo
\end{multicols}
\end{enumerate}
\vspace{0.15cm}
\item Si a < b y c < 0, entonces ac es:
\begin{enumerate}[label=\textbf{\Alph*)}]
\begin{multicols}{2}
  \item Mayor que
  \item Menor que
  \item Igual a
  \item No se puede determinar
  \item Siempre positivo
  \item Siempre negativo
\end{multicols}
\end{enumerate}
\vspace{0.15cm}
\item ¿Cuál de los siguientes números está entre 1 y 2?
\begin{enumerate}[label=\textbf{\Alph*)}]
\begin{multicols}{2}
  \item 0
  \item 1.5
  \item 2
  \item 3
  \item -1
  \item -2
\end{multicols}
\end{enumerate}
\vspace{0.15cm}
\item ¿Cuál de los siguientes números está entre -1 y 0?
\begin{enumerate}[label=\textbf{\Alph*)}]
\begin{multicols}{2}
  \item -0.5
  \item 1
  \item 2
  \item 3
  \item -1
  \item -2
\end{multicols}
\end{enumerate}
\vspace{0.15cm}
\item ¿Cuál de los siguientes números está entre \frac{1}{4} y \frac{1}{2}?
\begin{enumerate}[label=\textbf{\Alph*)}]
\begin{multicols}{2}
  \item 0
  \item \frac{3}{8}
  \item 2
  \item 3
  \item -1
  \item -2
\end{multicols}
\end{enumerate}
\vspace{0.15cm}
\item ¿Cuál de los siguientes números está entre 0.1 y 0.2?
\begin{enumerate}[label=\textbf{\Alph*)}]
\begin{multicols}{2}
  \item 0
  \item 0.15
  \item 2
  \item 3
  \item -1
  \item -2
\end{multicols}
\end{enumerate}
\vspace{0.15cm}
\item ¿Cuál de los siguientes números está entre -2 y -1?
\begin{enumerate}[label=\textbf{\Alph*)}]
\begin{multicols}{2}
  \item -1.5
  \item 1
  \item 2
  \item 3
  \item -1
  \item -2
\end{multicols}
\end{enumerate}
\vspace{0.15cm}
\item ¿Cuál de los siguientes números está entre 3 y 4?
\begin{enumerate}[label=\textbf{\Alph*)}]
\begin{multicols}{2}
  \item 0
  \item 3.5
  \item 2
  \item 3
  \item -1
  \item -2
\end{multicols}
\end{enumerate}
\vspace{0.15cm}
\item ¿Cuál de los siguientes números está entre -\frac{1}{2} y 0?
\begin{enumerate}[label=\textbf{\Alph*)}]
\begin{multicols}{2}
  \item -\frac{1}{4}
  \item 1
  \item 2
  \item 3
  \item -1
  \item -2
\end{multicols}
\end{enumerate}
\vspace{0.15cm}
\item ¿Cuál de los siguientes números está entre \frac{2}{3} y 1?
\begin{enumerate}[label=\textbf{\Alph*)}]
\begin{multicols}{2}
  \item 0
  \item \frac{5}{6}
  \item 2
  \item 3
  \item -1
  \item -2
\end{multicols}
\end{enumerate}
\vspace{0.15cm}
\item ¿Cuál de los siguientes números está entre -0.5 y 0.5?
\begin{enumerate}[label=\textbf{\Alph*)}]
\begin{multicols}{2}
  \item 0
  \item 1
  \item 2
  \item 3
  \item -1
  \item -2
\end{multicols}
\end{enumerate}
\vspace{0.15cm}
\item ¿Cuál de los siguientes números está entre 1.5 y 2.5?
\begin{enumerate}[label=\textbf{\Alph*)}]
\begin{multicols}{2}
  \item 0
  \item 2
  \item 2
  \item 3
  \item -1
  \item -2
\end{multicols}
\end{enumerate}
\vspace{0.15cm}
\item Si x > 5, ¿cuál de los siguientes valores satisface la desigualdad?
\begin{enumerate}[label=\textbf{\Alph*)}]
\begin{multicols}{2}
  \item 0
  \item 1
  \item 2
  \item 6
  \item -1
  \item -2
\end{multicols}
\end{enumerate}
\vspace{0.15cm}
\item Si x < 3, ¿cuál de los siguientes valores satisface la desigualdad?
\begin{enumerate}[label=\textbf{\Alph*)}]
\begin{multicols}{2}
  \item 0
  \item 1
  \item 2
  \item 3
  \item -1
  \item -2
\end{multicols}
\end{enumerate}
\vspace{0.15cm}
\item Si x \geq -2, ¿cuál NO satisface la desigualdad?
\begin{enumerate}[label=\textbf{\Alph*)}]
\begin{multicols}{2}
  \item 0
  \item 1
  \item 2
  \item 3
  \item -3
  \item -2
\end{multicols}
\end{enumerate}
\vspace{0.15cm}
\item Si x \leq 4, ¿cuál NO satisface la desigualdad?
\begin{enumerate}[label=\textbf{\Alph*)}]
\begin{multicols}{2}
  \item 0
  \item 1
  \item 2
  \item 5
  \item -1
  \item -2
\end{multicols}
\end{enumerate}
\vspace{0.15cm}
\item Si -3 < x < 3, ¿cuál satisface la desigualdad?
\begin{enumerate}[label=\textbf{\Alph*)}]
\begin{multicols}{2}
  \item 0
  \item 1
  \item 2
  \item 3
  \item -1
  \item -2
\end{multicols}
\end{enumerate}
\vspace{0.15cm}
\item Si 0 < x < 1, ¿cuál satisface la desigualdad?
\begin{enumerate}[label=\textbf{\Alph*)}]
\begin{multicols}{2}
  \item 0
  \item 1
  \item 0.5
  \item 3
  \item -1
  \item -2
\end{multicols}
\end{enumerate}
\vspace{0.15cm}
\end{enumerate}\n
\subsubsection*{SECCIÓN V: OPERACIONES CON NÚMEROS REALES (191–240)}

\begin{enumerate}[label=\textbf{\arabic*.}]
\setcounter{enumi}{190}
\item ¿Cuál es el resultado de \(-3+5\)?
\begin{enumerate}[label=\textbf{\Alph*)}]
\begin{multicols}{2}
  \item 0
  \item 2
  \item 2
  \item -1
  \item -2
  \item Indefinido
\end{multicols}
\end{enumerate}
\vspace{0.15cm}
\item ¿Cuál es el resultado de \(7+(-4)\)?
\begin{enumerate}[label=\textbf{\Alph*)}]
\begin{multicols}{2}
  \item 3
  \item 1
  \item 2
  \item -1
  \item -2
  \item Indefinido
\end{multicols}
\end{enumerate}
\vspace{0.15cm}
\item ¿Cuál es el resultado de \(-6-2\)?
\begin{enumerate}[label=\textbf{\Alph*)}]
\begin{multicols}{2}
  \item 0
  \item 1
  \item -8
  \item -1
  \item -2
  \item Indefinido
\end{multicols}
\end{enumerate}
\vspace{0.15cm}
\item ¿Cuál es el resultado de \(5-(-3)\)?
\begin{enumerate}[label=\textbf{\Alph*)}]
\begin{multicols}{2}
  \item 0
  \item 1
  \item 8
  \item -1
  \item -2
  \item Indefinido
\end{multicols}
\end{enumerate}
\vspace{0.15cm}
\item ¿Cuál es el resultado de \((-4)(-3)\)?
\begin{enumerate}[label=\textbf{\Alph*)}]
\begin{multicols}{2}
  \item 0
  \item 12
  \item 2
  \item -1
  \item -2
  \item Indefinido
\end{multicols}
\end{enumerate}
\vspace{0.15cm}
\item ¿Cuál es el resultado de \((-5)(3)\)?
\begin{enumerate}[label=\textbf{\Alph*)}]
\begin{multicols}{2}
  \item -15
  \item 1
  \item 2
  \item -1
  \item -2
  \item Indefinido
\end{multicols}
\end{enumerate}
\vspace{0.15cm}
\item ¿Cuál es el resultado de \((-12)\div(-4)\)?
\begin{enumerate}[label=\textbf{\Alph*)}]
\begin{multicols}{2}
  \item 0
  \item 3
  \item 2
  \item -1
  \item -2
  \item Indefinido
\end{multicols}
\end{enumerate}
\vspace{0.15cm}
\item ¿Cuál es el resultado de \(15\div(-3)\)?
\begin{enumerate}[label=\textbf{\Alph*)}]
\begin{multicols}{2}
  \item -5
  \item 1
  \item 2
  \item -1
  \item -2
  \item Indefinido
\end{multicols}
\end{enumerate}
\vspace{0.15cm}
\item ¿Cuál es el resultado de \(-2+3-4+5\)?
\begin{enumerate}[label=\textbf{\Alph*)}]
\begin{multicols}{2}
  \item 2
  \item 1
  \item 2
  \item -1
  \item -2
  \item Indefinido
\end{multicols}
\end{enumerate}
\vspace{0.15cm}
\item ¿Cuál es el resultado de \(2\cdot3+4\cdot5\)?
\begin{enumerate}[label=\textbf{\Alph*)}]
\begin{multicols}{2}
  \item 26
  \item 1
  \item 2
  \item -1
  \item -2
  \item Indefinido
\end{multicols}
\end{enumerate}
\vspace{0.15cm}
\item ¿Cuál es el resultado de \((2+3)\cdot4\)?
\begin{enumerate}[label=\textbf{\Alph*)}]
\begin{multicols}{2}
  \item 20
  \item 1
  \item 2
  \item -1
  \item -2
  \item Indefinido
\end{multicols}
\end{enumerate}
\vspace{0.15cm}
\item ¿Cuál es el resultado de \(2+3\cdot4\)?
\begin{enumerate}[label=\textbf{\Alph*)}]
\begin{multicols}{2}
  \item 0
  \item 14
  \item 2
  \item -1
  \item -2
  \item Indefinido
\end{multicols}
\end{enumerate}
\vspace{0.15cm}
\item ¿Cuál es el resultado de \(-3\cdot(-2)+5\)?
\begin{enumerate}[label=\textbf{\Alph*)}]
\begin{multicols}{2}
  \item 11
  \item 1
  \item 2
  \item -1
  \item -2
  \item Indefinido
\end{multicols}
\end{enumerate}
\vspace{0.15cm}
\item ¿Cuál es el resultado de \((-2)^3\)?
\begin{enumerate}[label=\textbf{\Alph*)}]
\begin{multicols}{2}
  \item -8
  \item 1
  \item 2
  \item -1
  \item -2
  \item Indefinido
\end{multicols}
\end{enumerate}
\vspace{0.15cm}
\item ¿Cuál es el resultado de \((-2)^4\)?
\begin{enumerate}[label=\textbf{\Alph*)}]
\begin{multicols}{2}
  \item 0
  \item 16
  \item 2
  \item -1
  \item -2
  \item Indefinido
\end{multicols}
\end{enumerate}
\vspace{0.15cm}
\item ¿Cuál es el resultado de \(-2^2\)?
\begin{enumerate}[label=\textbf{\Alph*)}]
\begin{multicols}{2}
  \item 0
  \item -4
  \item 2
  \item -1
  \item -2
  \item Indefinido
\end{multicols}
\end{enumerate}
\vspace{0.15cm}
\item ¿Cuál es el resultado de \((-2)^2\)?
\begin{enumerate}[label=\textbf{\Alph*)}]
\begin{multicols}{2}
  \item 0
  \item 4
  \item 2
  \item -1
  \item -2
  \item Indefinido
\end{multicols}
\end{enumerate}
\vspace{0.15cm}
\item ¿Cuál es el resultado de \(\sqrt{16}\)?
\begin{enumerate}[label=\textbf{\Alph*)}]
\begin{multicols}{2}
  \item 4
  \item 1
  \item 2
  \item -1
  \item -2
  \item Indefinido
\end{multicols}
\end{enumerate}
\vspace{0.15cm}
\item ¿Cuál es el resultado de \(-\sqrt{9}\)?
\begin{enumerate}[label=\textbf{\Alph*)}]
\begin{multicols}{2}
  \item 0
  \item -3
  \item 2
  \item -1
  \item -2
  \item Indefinido
\end{multicols}
\end{enumerate}
\vspace{0.15cm}
\item ¿Cuál es el resultado de \(\sqrt{-4}\)?
\begin{enumerate}[label=\textbf{\Alph*)}]
\begin{multicols}{2}
  \item 0
  \item 1
  \item 2
  \item No es real
  \item -2
  \item Indefinido
\end{multicols}
\end{enumerate}
\vspace{0.15cm}
\item ¿Cuál es el resultado de \(2^3\cdot2^2\)?
\begin{enumerate}[label=\textbf{\Alph*)}]
\begin{multicols}{2}
  \item 2^5
  \item 1
  \item 2
  \item -1
  \item -2
  \item Indefinido
\end{multicols}
\end{enumerate}
\vspace{0.15cm}
\item ¿Cuál es el resultado de \(\frac{2^5}{2^2}\)?
\begin{enumerate}[label=\textbf{\Alph*)}]
\begin{multicols}{2}
  \item 2^3
  \item 1
  \item 2
  \item -1
  \item -2
  \item Indefinido
\end{multicols}
\end{enumerate}
\vspace{0.15cm}
\item ¿Cuál es el resultado de \((2^3)^2\)?
\begin{enumerate}[label=\textbf{\Alph*)}]
\begin{multicols}{2}
  \item 0
  \item 2^6
  \item 2
  \item -1
  \item -2
  \item Indefinido
\end{multicols}
\end{enumerate}
\vspace{0.15cm}
\item ¿Cuál es el resultado de \(2^0\)?
\begin{enumerate}[label=\textbf{\Alph*)}]
\begin{multicols}{2}
  \item 0
  \item 1
  \item 2
  \item -1
  \item -2
  \item Indefinido
\end{multicols}
\end{enumerate}
\vspace{0.15cm}
\item ¿Cuál es el resultado de \(2^{-2}\)?
\begin{enumerate}[label=\textbf{\Alph*)}]
\begin{multicols}{2}
  \item 0
  \item \frac{1}{4}
  \item 2
  \item -1
  \item -2
  \item Indefinido
\end{multicols}
\end{enumerate}
\vspace{0.15cm}
\item ¿Cuál es el resultado de \(\frac{1}{2}+\frac{1}{3}\)?
\begin{enumerate}[label=\textbf{\Alph*)}]
\begin{multicols}{2}
  \item 0
  \item \frac{5}{6}
  \item 2
  \item -1
  \item -2
  \item Indefinido
\end{multicols}
\end{enumerate}
\vspace{0.15cm}
\item ¿Cuál es el resultado de \(\frac{3}{4}-\frac{1}{2}\)?
\begin{enumerate}[label=\textbf{\Alph*)}]
\begin{multicols}{2}
  \item \frac{1}{4}
  \item 1
  \item 2
  \item -1
  \item -2
  \item Indefinido
\end{multicols}
\end{enumerate}
\vspace{0.15cm}
\item ¿Cuál es el resultado de \(\frac{2}{3}\cdot\frac{3}{4}\)?
\begin{enumerate}[label=\textbf{\Alph*)}]
\begin{multicols}{2}
  \item \frac{1}{2}
  \item 1
  \item 2
  \item -1
  \item -2
  \item Indefinido
\end{multicols}
\end{enumerate}
\vspace{0.15cm}
\item ¿Cuál es el resultado de \(\frac{2}{3}\div\frac{4}{5}\)?
\begin{enumerate}[label=\textbf{\Alph*)}]
\begin{multicols}{2}
  \item 0
  \item \frac{5}{6}
  \item 2
  \item -1
  \item -2
  \item Indefinido
\end{multicols}
\end{enumerate}
\vspace{0.15cm}
\item ¿Cuál es el resultado de \(\frac{1}{2}+\frac{1}{4}+\frac{1}{8}\)?
\begin{enumerate}[label=\textbf{\Alph*)}]
\begin{multicols}{2}
  \item 0
  \item \frac{7}{8}
  \item 2
  \item -1
  \item -2
  \item Indefinido
\end{multicols}
\end{enumerate}
\vspace{0.15cm}
\item ¿Cuál es el resultado de \(0.5+0.25\)?
\begin{enumerate}[label=\textbf{\Alph*)}]
\begin{multicols}{2}
  \item 0.75
  \item 1
  \item 2
  \item -1
  \item -2
  \item Indefinido
\end{multicols}
\end{enumerate}
\vspace{0.15cm}
\item ¿Cuál es el resultado de \(1.5-0.75\)?
\begin{enumerate}[label=\textbf{\Alph*)}]
\begin{multicols}{2}
  \item 0.75
  \item 1
  \item 2
  \item -1
  \item -2
  \item Indefinido
\end{multicols}
\end{enumerate}
\vspace{0.15cm}
\item ¿Cuál es el resultado de \(0.5\cdot0.4\)?
\begin{enumerate}[label=\textbf{\Alph*)}]
\begin{multicols}{2}
  \item 0.2
  \item 1
  \item 2
  \item -1
  \item -2
  \item Indefinido
\end{multicols}
\end{enumerate}
\vspace{0.15cm}
\item ¿Cuál es el resultado de \(0.8\div0.2\)?
\begin{enumerate}[label=\textbf{\Alph*)}]
\begin{multicols}{2}
  \item 0
  \item 4
  \item 2
  \item -1
  \item -2
  \item Indefinido
\end{multicols}
\end{enumerate}
\vspace{0.15cm}
\item ¿Cuál es el resultado de \(-5+3-2+4\)?
\begin{enumerate}[label=\textbf{\Alph*)}]
\begin{multicols}{2}
  \item 0
  \item 1
  \item 2
  \item -1
  \item -2
  \item Indefinido
\end{multicols}
\end{enumerate}
\vspace{0.15cm}
\item ¿Cuál es el resultado de \((-3)^2+(-2)^3\)?
\begin{enumerate}[label=\textbf{\Alph*)}]
\begin{multicols}{2}
  \item 1
  \item 1
  \item 2
  \item -1
  \item -2
  \item Indefinido
\end{multicols}
\end{enumerate}
\vspace{0.15cm}
\item ¿Cuál es el resultado de \(\sqrt{4}+\sqrt{9}\)?
\begin{enumerate}[label=\textbf{\Alph*)}]
\begin{multicols}{2}
  \item 0
  \item 5
  \item 2
  \item -1
  \item -2
  \item Indefinido
\end{multicols}
\end{enumerate}
\vspace{0.15cm}
\item ¿Cuál es el resultado de \(\sqrt{25}-\sqrt{16}\)?
\begin{enumerate}[label=\textbf{\Alph*)}]
\begin{multicols}{2}
  \item 1
  \item 1
  \item 2
  \item -1
  \item -2
  \item Indefinido
\end{multicols}
\end{enumerate}
\vspace{0.15cm}
\item ¿Cuál es el resultado de \(3^2+4^2\)?
\begin{enumerate}[label=\textbf{\Alph*)}]
\begin{multicols}{2}
  \item 0
  \item 25
  \item 2
  \item -1
  \item -2
  \item Indefinido
\end{multicols}
\end{enumerate}
\vspace{0.15cm}
\item ¿Cuál es el resultado de \(\sqrt{3^2+4^2}\)?
\begin{enumerate}[label=\textbf{\Alph*)}]
\begin{multicols}{2}
  \item 0
  \item 5
  \item 2
  \item -1
  \item -2
  \item Indefinido
\end{multicols}
\end{enumerate}
\vspace{0.15cm}
\item ¿Cuál es el resultado de \((-1)^{100}\)?
\begin{enumerate}[label=\textbf{\Alph*)}]
\begin{multicols}{2}
  \item 0
  \item 1
  \item 2
  \item -1
  \item -2
  \item Indefinido
\end{multicols}
\end{enumerate}
\vspace{0.15cm}
\item ¿Cuál es el resultado de \((-1)^{101}\)?
\begin{enumerate}[label=\textbf{\Alph*)}]
\begin{multicols}{2}
  \item -1
  \item 1
  \item 2
  \item -1
  \item -2
  \item Indefinido
\end{multicols}
\end{enumerate}
\vspace{0.15cm}
\item ¿Cuál es el resultado de \(2^3+2^2\)?
\begin{enumerate}[label=\textbf{\Alph*)}]
\begin{multicols}{2}
  \item 0
  \item 12
  \item 2
  \item -1
  \item -2
  \item Indefinido
\end{multicols}
\end{enumerate}
\vspace{0.15cm}
\item ¿Cuál es el resultado de \(5\cdot0\)?
\begin{enumerate}[label=\textbf{\Alph*)}]
\begin{multicols}{2}
  \item 0
  \item 0
  \item 2
  \item -1
  \item -2
  \item Indefinido
\end{multicols}
\end{enumerate}
\vspace{0.15cm}
\item ¿Cuál es el resultado de \(0\div5\)?
\begin{enumerate}[label=\textbf{\Alph*)}]
\begin{multicols}{2}
  \item 0
  \item 0
  \item 2
  \item -1
  \item -2
  \item Indefinido
\end{multicols}
\end{enumerate}
\vspace{0.15cm}
\item ¿Cuál es el resultado de \(5\div0\)?
\begin{enumerate}[label=\textbf{\Alph*)}]
\begin{multicols}{2}
  \item 0
  \item 1
  \item 2
  \item Indefinido
  \item -2
  \item Indefinido
\end{multicols}
\end{enumerate}
\vspace{0.15cm}
\item ¿Cuál es el resultado de \(\frac{2}{3}+\frac{1}{6}\)?
\begin{enumerate}[label=\textbf{\Alph*)}]
\begin{multicols}{2}
  \item 0
  \item \frac{5}{6}
  \item 2
  \item -1
  \item -2
  \item Indefinido
\end{multicols}
\end{enumerate}
\vspace{0.15cm}
\item ¿Cuál es el resultado de \(\frac{5}{8}-\frac{1}{4}\)?
\begin{enumerate}[label=\textbf{\Alph*)}]
\begin{multicols}{2}
  \item 0
  \item \frac{3}{8}
  \item 2
  \item -1
  \item -2
  \item Indefinido
\end{multicols}
\end{enumerate}
\vspace{0.15cm}
\item ¿Cuál es el resultado de \(\frac{3}{5}\cdot\frac{10}{9}\)?
\begin{enumerate}[label=\textbf{\Alph*)}]
\begin{multicols}{2}
  \item \frac{2}{3}
  \item 1
  \item 2
  \item -1
  \item -2
  \item Indefinido
\end{multicols}
\end{enumerate}
\vspace{0.15cm}
\item ¿Cuál es el resultado de \(\frac{4}{7}\div\frac{2}{7}\)?
\begin{enumerate}[label=\textbf{\Alph*)}]
\begin{multicols}{2}
  \item 0
  \item 2
  \item 2
  \item -1
  \item -2
  \item Indefinido
\end{multicols}
\end{enumerate}
\vspace{0.15cm}
\end{enumerate}\n
\subsubsection*{SECCIÓN VI: DENSIDAD E INTERVALOS (241–300)}

\begin{enumerate}[label=\textbf{\arabic*.}]
\setcounter{enumi}{240}
\item ¿Cuál de los siguientes números está entre \(2\) y \(3\)?
\begin{enumerate}[label=\textbf{\Alph*)}]
\begin{multicols}{2}
  \item 0
  \item 2.5
  \item 2
  \item -1
  \item -2
  \item 10
\end{multicols}
\end{enumerate}
\vspace{0.15cm}
\item ¿Cuál de los siguientes números está entre \(-2\) y \(-1\)?
\begin{enumerate}[label=\textbf{\Alph*)}]
\begin{multicols}{2}
  \item -1.5
  \item 1
  \item 2
  \item -1
  \item -2
  \item 10
\end{multicols}
\end{enumerate}
\vspace{0.15cm}
\item ¿Cuál de los siguientes números está entre \(0\) y \(1\)?
\begin{enumerate}[label=\textbf{\Alph*)}]
\begin{multicols}{2}
  \item 0.5
  \item 1
  \item 2
  \item -1
  \item -2
  \item 10
\end{multicols}
\end{enumerate}
\vspace{0.15cm}
\item ¿Cuál de los siguientes números está entre \(\frac{1}{2}\) y \(\frac{3}{4}\)?
\begin{enumerate}[label=\textbf{\Alph*)}]
\begin{multicols}{2}
  \item 0
  \item \frac{5}{8}
  \item 2
  \item -1
  \item -2
  \item 10
\end{multicols}
\end{enumerate}
\vspace{0.15cm}
\item ¿Cuál de los siguientes números está entre \(1.1\) y \(1.2\)?
\begin{enumerate}[label=\textbf{\Alph*)}]
\begin{multicols}{2}
  \item 1.15
  \item 1
  \item 2
  \item -1
  \item -2
  \item 10
\end{multicols}
\end{enumerate}
\vspace{0.15cm}
\item ¿Cuál de los siguientes números está entre \(-0.1\) y \(0.1\)?
\begin{enumerate}[label=\textbf{\Alph*)}]
\begin{multicols}{2}
  \item 0
  \item 1
  \item 2
  \item -1
  \item -2
  \item 10
\end{multicols}
\end{enumerate}
\vspace{0.15cm}
\item ¿Cuál de los siguientes números está entre \(\sqrt{2}\) y \(\sqrt{3}\)?
\begin{enumerate}[label=\textbf{\Alph*)}]
\begin{multicols}{2}
  \item 1.5
  \item 1
  \item 2
  \item -1
  \item -2
  \item 10
\end{multicols}
\end{enumerate}
\vspace{0.15cm}
\item ¿Cuál de los siguientes números está entre \(\pi\) y \(4\)?
\begin{enumerate}[label=\textbf{\Alph*)}]
\begin{multicols}{2}
  \item 3.5
  \item 1
  \item 2
  \item -1
  \item -2
  \item 10
\end{multicols}
\end{enumerate}
\vspace{0.15cm}
\item ¿Cuál de los siguientes números está entre \(-\frac{3}{4}\) y \(-\frac{1}{2}\)?
\begin{enumerate}[label=\textbf{\Alph*)}]
\begin{multicols}{2}
  \item -\frac{5}{8}
  \item 1
  \item 2
  \item -1
  \item -2
  \item 10
\end{multicols}
\end{enumerate}
\vspace{0.15cm}
\item ¿Cuál de los siguientes números está entre \(0.3\) y \(0.4\)?
\begin{enumerate}[label=\textbf{\Alph*)}]
\begin{multicols}{2}
  \item 0.35
  \item 1
  \item 2
  \item -1
  \item -2
  \item 10
\end{multicols}
\end{enumerate}
\vspace{0.15cm}
\item ¿Cuál de los siguientes números está entre \(\frac{5}{6}\) y \(1\)?
\begin{enumerate}[label=\textbf{\Alph*)}]
\begin{multicols}{2}
  \item \frac{11}{12}
  \item 1
  \item 2
  \item -1
  \item -2
  \item 10
\end{multicols}
\end{enumerate}
\vspace{0.15cm}
\item ¿Cuál de los siguientes números está entre \(-5\) y \(-4\)?
\begin{enumerate}[label=\textbf{\Alph*)}]
\begin{multicols}{2}
  \item -4.5
  \item 1
  \item 2
  \item -1
  \item -2
  \item 10
\end{multicols}
\end{enumerate}
\vspace{0.15cm}
\item ¿Cuál de los siguientes números está entre \(10\) y \(11\)?
\begin{enumerate}[label=\textbf{\Alph*)}]
\begin{multicols}{2}
  \item 10.5
  \item 1
  \item 2
  \item -1
  \item -2
  \item 10
\end{multicols}
\end{enumerate}
\vspace{0.15cm}
\item ¿Cuál de los siguientes números está entre \(-1\) y \(1\)?
\begin{enumerate}[label=\textbf{\Alph*)}]
\begin{multicols}{2}
  \item 0.5
  \item 1
  \item 2
  \item -1
  \item -2
  \item 10
\end{multicols}
\end{enumerate}
\vspace{0.15cm}
\item ¿Cuál de los siguientes números está entre \(0.01\) y \(0.02\)?
\begin{enumerate}[label=\textbf{\Alph*)}]
\begin{multicols}{2}
  \item 0.015
  \item 1
  \item 2
  \item -1
  \item -2
  \item 10
\end{multicols}
\end{enumerate}
\vspace{0.15cm}
\item ¿Cuál de los siguientes números está entre \(\frac{2}{5}\) y \(\frac{3}{5}\)?
\begin{enumerate}[label=\textbf{\Alph*)}]
\begin{multicols}{2}
  \item \frac{1}{2}
  \item 1
  \item 2
  \item -1
  \item -2
  \item 10
\end{multicols}
\end{enumerate}
\vspace{0.15cm}
\item ¿Cuál de los siguientes números está entre \(-\frac{1}{3}\) y \(\frac{1}{3}\)?
\begin{enumerate}[label=\textbf{\Alph*)}]
\begin{multicols}{2}
  \item 0
  \item 1
  \item 2
  \item -1
  \item -2
  \item 10
\end{multicols}
\end{enumerate}
\vspace{0.15cm}
\item ¿Cuál de los siguientes números está entre \(2.5\) y \(2.6\)?
\begin{enumerate}[label=\textbf{\Alph*)}]
\begin{multicols}{2}
  \item 2.55
  \item 1
  \item 2
  \item -1
  \item -2
  \item 10
\end{multicols}
\end{enumerate}
\vspace{0.15cm}
\item ¿Cuál de los siguientes números está entre \(-0.01\) y \(0\)?
\begin{enumerate}[label=\textbf{\Alph*)}]
\begin{multicols}{2}
  \item -0.005
  \item 1
  \item 2
  \item -1
  \item -2
  \item 10
\end{multicols}
\end{enumerate}
\vspace{0.15cm}
\item ¿Cuál de los siguientes números está entre \(100\) y \(101\)?
\begin{enumerate}[label=\textbf{\Alph*)}]
\begin{multicols}{2}
  \item 100.5
  \item 1
  \item 2
  \item -1
  \item -2
  \item 10
\end{multicols}
\end{enumerate}
\vspace{0.15cm}
\item El intervalo \((2,5)\) incluye:
\begin{enumerate}[label=\textbf{\Alph*)}]
\begin{multicols}{2}
  \item Solo 2 y 5
  \item Todos los números entre 2 y 5 sin incluirlos
  \item Todos entre 2 y 5 incluyendo 2 pero no 5
  \item Solo 2
  \item Solo 5
  \item Ninguno
\end{multicols}
\end{enumerate}
\vspace{0.15cm}
\item El intervalo \([2,5]\) incluye:
\begin{enumerate}[label=\textbf{\Alph*)}]
\begin{multicols}{2}
  \item Solo 2 y 5
  \item Todos entre 2 y 5 incluyendo 5 pero no 2
  \item Todos los números entre 2 y 5 incluyéndolos
  \item Solo 2
  \item Solo 5
  \item Ninguno
\end{multicols}
\end{enumerate}
\vspace{0.15cm}
\item El intervalo \([2,5)\) incluye:
\begin{enumerate}[label=\textbf{\Alph*)}]
\begin{multicols}{2}
  \item Solo 2 y 5
  \item Todos los números entre 2 y 5 sin incluirlos
  \item Todos entre 2 y 5 incluyendo 2 pero no 5
  \item Solo 2
  \item Solo 5
  \item Ninguno
\end{multicols}
\end{enumerate}
\vspace{0.15cm}
\item El intervalo \((2,5]\) incluye:
\begin{enumerate}[label=\textbf{\Alph*)}]
\begin{multicols}{2}
  \item Solo 2 y 5
  \item Todos entre 2 y 5 incluyendo 5 pero no 2
  \item Todos los números entre 2 y 5 incluyéndolos
  \item Solo 2
  \item Solo 5
  \item Ninguno
\end{multicols}
\end{enumerate}
\vspace{0.15cm}
\item La desigualdad \(2 < x < 5\) se escribe como:
\begin{enumerate}[label=\textbf{\Alph*)}]
\begin{multicols}{2}
  \item [2,5]
  \item (2,5)
  \item [2,5)
  \item (2,5]
  \item (-\infty,5)
  \item (2,\infty)
\end{multicols}
\end{enumerate}
\vspace{0.15cm}
\item La desigualdad \(2 \leq x \leq 5\) se escribe como:
\begin{enumerate}[label=\textbf{\Alph*)}]
\begin{multicols}{2}
  \item [2,5]
  \item (2,5)
  \item [2,5)
  \item (2,5]
  \item (-\infty,5)
  \item (2,\infty)
\end{multicols}
\end{enumerate}
\vspace{0.15cm}
\item La desigualdad \(2 \leq x < 5\) se escribe como:
\begin{enumerate}[label=\textbf{\Alph*)}]
\begin{multicols}{2}
  \item [2,5]
  \item (2,5)
  \item [2,5)
  \item (2,5]
  \item (-\infty,5)
  \item (2,\infty)
\end{multicols}
\end{enumerate}
\vspace{0.15cm}
\item La desigualdad \(2 < x \leq 5\) se escribe como:
\begin{enumerate}[label=\textbf{\Alph*)}]
\begin{multicols}{2}
  \item [2,5]
  \item (2,5)
  \item [2,5)
  \item (2,5]
  \item (-\infty,5)
  \item (2,\infty)
\end{multicols}
\end{enumerate}
\vspace{0.15cm}
\item La desigualdad \(x > 3\) se escribe como:
\begin{enumerate}[label=\textbf{\Alph*)}]
\begin{multicols}{2}
  \item (3,\infty)
  \item [3,\infty)
  \item (-\infty,3)
  \item (-\infty,3]
  \item [3,\infty]
  \item (3,\infty]
\end{multicols}
\end{enumerate}
\vspace{0.15cm}
\item La desigualdad \(x \geq 3\) se escribe como:
\begin{enumerate}[label=\textbf{\Alph*)}]
\begin{multicols}{2}
  \item (3,\infty)
  \item [3,\infty)
  \item (-\infty,3)
  \item (-\infty,3]
  \item [3,\infty]
  \item (3,\infty]
\end{multicols}
\end{enumerate}
\vspace{0.15cm}
\item La desigualdad \(x < 3\) se escribe como:
\begin{enumerate}[label=\textbf{\Alph*)}]
\begin{multicols}{2}
  \item (3,\infty)
  \item [3,\infty)
  \item (-\infty,3)
  \item (-\infty,3]
  \item [3,\infty]
  \item (3,\infty]
\end{multicols}
\end{enumerate}
\vspace{0.15cm}
\item La desigualdad \(x \leq 3\) se escribe como:
\begin{enumerate}[label=\textbf{\Alph*)}]
\begin{multicols}{2}
  \item (3,\infty)
  \item [3,\infty)
  \item (-\infty,3)
  \item (-\infty,3]
  \item [3,\infty]
  \item (3,\infty]
\end{multicols}
\end{enumerate}
\vspace{0.15cm}
\item El intervalo \((-\infty,5)\) representa:
\begin{enumerate}[label=\textbf{\Alph*)}]
\begin{multicols}{2}
  \item x > 5
  \item x < 5
  \item x \geq 5
  \item x \leq 5
  \item x = 5
  \item Todos los reales
\end{multicols}
\end{enumerate}
\vspace{0.15cm}
\item El intervalo \((-\infty,5]\) representa:
\begin{enumerate}[label=\textbf{\Alph*)}]
\begin{multicols}{2}
  \item x > 5
  \item x < 5
  \item x \geq 5
  \item x \leq 5
  \item x = 5
  \item Todos los reales
\end{multicols}
\end{enumerate}
\vspace{0.15cm}
\item El intervalo \((5,\infty)\) representa:
\begin{enumerate}[label=\textbf{\Alph*)}]
\begin{multicols}{2}
  \item x > 5
  \item x < 5
  \item x \geq 5
  \item x \leq 5
  \item x = 5
  \item Todos los reales
\end{multicols}
\end{enumerate}
\vspace{0.15cm}
\item El intervalo \([5,\infty)\) representa:
\begin{enumerate}[label=\textbf{\Alph*)}]
\begin{multicols}{2}
  \item x > 5
  \item x < 5
  \item x \geq 5
  \item x \leq 5
  \item x = 5
  \item Todos los reales
\end{multicols}
\end{enumerate}
\vspace{0.15cm}
\item ¿Cuál es la intersección de \((1,5)\) y \((3,7)\)?
\begin{enumerate}[label=\textbf{\Alph*)}]
\begin{multicols}{2}
  \item (1,7)
  \item (3,5)
  \item \emptyset
  \item (5,7)
  \item \emptyset
  \item (1,5)
\end{multicols}
\end{enumerate}
\vspace{0.15cm}
\item ¿Cuál es la unión de \((1,5)\) y \((3,7)\)?
\begin{enumerate}[label=\textbf{\Alph*)}]
\begin{multicols}{2}
  \item (1,7)
  \item (3,5)
  \item \emptyset
  \item (5,7)
  \item \emptyset
  \item (1,5)
\end{multicols}
\end{enumerate}
\vspace{0.15cm}
\item ¿Cuál es la intersección de \((-\infty,3)\) y \((1,\infty)\)?
\begin{enumerate}[label=\textbf{\Alph*)}]
\begin{multicols}{2}
  \item (1,3)
  \item (3,5)
  \item \emptyset
  \item (5,7)
  \item \emptyset
  \item (1,5)
\end{multicols}
\end{enumerate}
\vspace{0.15cm}
\item ¿Cuál es la unión de \((-\infty,3)\) y \((1,\infty)\)?
\begin{enumerate}[label=\textbf{\Alph*)}]
\begin{multicols}{2}
  \item (1,7)
  \item (-\infty,\infty)
  \item \emptyset
  \item (5,7)
  \item \emptyset
  \item (1,5)
\end{multicols}
\end{enumerate}
\vspace{0.15cm}
\item ¿Cuál es la intersección de \([1,5]\) y \([3,7]\)?
\begin{enumerate}[label=\textbf{\Alph*)}]
\begin{multicols}{2}
  \item (1,7)
  \item [3,5]
  \item \emptyset
  \item (5,7)
  \item \emptyset
  \item (1,5)
\end{multicols}
\end{enumerate}
\vspace{0.15cm}
\item ¿Cuál es la unión de \([1,5]\) y \([3,7]\)?
\begin{enumerate}[label=\textbf{\Alph*)}]
\begin{multicols}{2}
  \item [1,7]
  \item (3,5)
  \item \emptyset
  \item (5,7)
  \item \emptyset
  \item (1,5)
\end{multicols}
\end{enumerate}
\vspace{0.15cm}
\item ¿Cuál es la intersección de \((1,3)\) y \((5,7)\)?
\begin{enumerate}[label=\textbf{\Alph*)}]
\begin{multicols}{2}
  \item (-\infty,\infty)
  \item (3,5)
  \item \emptyset
  \item (5,7)
  \item \emptyset
  \item (1,5)
\end{multicols}
\end{enumerate}
\vspace{0.15cm}
\item ¿Cuál es la unión de \((1,3)\) y \((5,7)\)?
\begin{enumerate}[label=\textbf{\Alph*)}]
\begin{multicols}{2}
  \item (1,7)
  \item (1,3)\cup(5,7)
  \item \emptyset
  \item (5,7)
  \item \emptyset
  \item (1,5)
\end{multicols}
\end{enumerate}
\vspace{0.15cm}
\item ¿Cuál es la intersección de \([-2,2]\) y \([0,4]\)?
\begin{enumerate}[label=\textbf{\Alph*)}]
\begin{multicols}{2}
  \item [-2,4]
  \item [0,2]
  \item \emptyset
  \item (5,7)
  \item \emptyset
  \item (1,5)
\end{multicols}
\end{enumerate}
\vspace{0.15cm}
\item ¿Cuál es la unión de \([-2,2]\) y \([0,4]\)?
\begin{enumerate}[label=\textbf{\Alph*)}]
\begin{multicols}{2}
  \item [-2,4]
  \item (3,5)
  \item \emptyset
  \item (5,7)
  \item \emptyset
  \item (1,5)
\end{multicols}
\end{enumerate}
\vspace{0.15cm}
\item ¿Cuál es la intersección de \((-\infty,0)\) y \((0,\infty)\)?
\begin{enumerate}[label=\textbf{\Alph*)}]
\begin{multicols}{2}
  \item (-\infty,\infty)
  \item (3,5)
  \item \emptyset
  \item (5,7)
  \item \emptyset
  \item (1,5)
\end{multicols}
\end{enumerate}
\vspace{0.15cm}
\item ¿Cuál es la unión de \((-\infty,0)\) y \((0,\infty)\)?
\begin{enumerate}[label=\textbf{\Alph*)}]
\begin{multicols}{2}
  \item (1,7)
  \item \mathbb{R}\setminus\{0\}
  \item \emptyset
  \item (5,7)
  \item \emptyset
  \item (1,5)
\end{multicols}
\end{enumerate}
\vspace{0.15cm}
\item [a,b] con a > b
\begin{enumerate}[label=\textbf{\Alph*)}]
\begin{multicols}{2}
  \item Válido
  \item Vacío
  \item Infinito
  \item {a}
  \item Igual a R
  \item No definido
\end{multicols}
\end{enumerate}
\vspace{0.15cm}
\item (a,a)
\begin{enumerate}[label=\textbf{\Alph*)}]
\begin{multicols}{2}
  \item Válido
  \item Vacío
  \item Infinito
  \item {a}
  \item Igual a R
  \item No definido
\end{multicols}
\end{enumerate}
\vspace{0.15cm}
\item [a,a]
\begin{enumerate}[label=\textbf{\Alph*)}]
\begin{multicols}{2}
  \item Válido
  \item Vacío
  \item Infinito
  \item {a}
  \item Igual a R
  \item No definido
\end{multicols}
\end{enumerate}
\vspace{0.15cm}
\item ¿Cuál contiene a 0?
\begin{enumerate}[label=\textbf{\Alph*)}]
\begin{multicols}{2}
  \item (0,1)
  \item (-1,1)
  \item (2,3)
  \item (0,2)
  \item (-2,-1)
  \item (3,4)
\end{multicols}
\end{enumerate}
\vspace{0.15cm}
\item ¿Cuál contiene a -1?
\begin{enumerate}[label=\textbf{\Alph*)}]
\begin{multicols}{2}
  \item Válido
  \item (-2,0)
  \item Infinito
  \item {a}
  \item Igual a R
  \item No definido
\end{multicols}
\end{enumerate}
\vspace{0.15cm}
\item ¿Cuál contiene a 5?
\begin{enumerate}[label=\textbf{\Alph*)}]
\begin{multicols}{2}
  \item Válido
  \item (4,6)
  \item Infinito
  \item {a}
  \item Igual a R
  \item No definido
\end{multicols}
\end{enumerate}
\vspace{0.15cm}
\item ¿Cuál contiene a \frac{1}{2}?
\begin{enumerate}[label=\textbf{\Alph*)}]
\begin{multicols}{2}
  \item Válido
  \item Vacío
  \item Infinito
  \item {a}
  \item Igual a R
  \item No definido
\end{multicols}
\end{enumerate}
\vspace{0.15cm}
\item ¿Cuál contiene a \sqrt{2}?
\begin{enumerate}[label=\textbf{\Alph*)}]
\begin{multicols}{2}
  \item Válido
  \item (1,2)
  \item Infinito
  \item {a}
  \item Igual a R
  \item No definido
\end{multicols}
\end{enumerate}
\vspace{0.15cm}
\item ¿Cuál contiene a \pi?
\begin{enumerate}[label=\textbf{\Alph*)}]
\begin{multicols}{2}
  \item Válido
  \item Vacío
  \item Infinito
  \item (3,4)
  \item Igual a R
  \item No definido
\end{multicols}
\end{enumerate}
\vspace{0.15cm}
\item (0,\infty)
\begin{enumerate}[label=\textbf{\Alph*)}]
\begin{multicols}{2}
  \item (0,1)
  \item (-1,1)
  \item (2,3)
  \item (0,2)
  \item (-2,-1)
  \item (3,4)
\end{multicols}
\end{enumerate}
\vspace{0.15cm}
\item (-\infty,0)
\begin{enumerate}[label=\textbf{\Alph*)}]
\begin{multicols}{2}
  \item (0,1)
  \item x < 0
  \item (2,3)
  \item (0,2)
  \item (-2,-1)
  \item (3,4)
\end{multicols}
\end{enumerate}
\vspace{0.15cm}
\item (-\infty,\infty)
\begin{enumerate}[label=\textbf{\Alph*)}]
\begin{multicols}{2}
  \item Válido
  \item Vacío
  \item Infinito
  \item Todos los números reales
  \item Igual a R
  \item No definido
\end{multicols}
\end{enumerate}
\vspace{0.15cm}
\end{enumerate}\n
\vspace{0.4cm}
\begin{tcolorbox}[colback=gray!5!white,colframe=gray!60!black,title=\textbf{SOLUCIONARIO -- 1.1 Números Reales}]
\begin{multicols}{4}
\begin{enumerate}[label=\textbf{\arabic*.}]
  \item \textbf{D}
  \item \textbf{B}
  \item \textbf{C}
  \item \textbf{B}
  \item \textbf{D}
  \item \textbf{A}
  \item \textbf{C}
  \item \textbf{A}
  \item \textbf{C}
  \item \textbf{D}
  \item \textbf{C}
  \item \textbf{E}
  \item \textbf{C}
  \item \textbf{D}
  \item \textbf{B}
  \item \textbf{B}
  \item \textbf{C}
  \item \textbf{E}
  \item \textbf{C}
  \item \textbf{C}
  \item \textbf{C}
  \item \textbf{E}
  \item \textbf{F}
  \item \textbf{B}
  \item \textbf{B}
  \item \textbf{D}
  \item \textbf{C}
  \item \textbf{E}
  \item \textbf{C}
  \item \textbf{D}
  \item \textbf{C}
  \item \textbf{E}
  \item \textbf{A}
  \item \textbf{C}
  \item \textbf{D}
  \item \textbf{C}
  \item \textbf{C}
  \item \textbf{A}
  \item \textbf{C}
  \item \textbf{C}
  \item \textbf{A}
  \item \textbf{C}
  \item \textbf{B}
  \item \textbf{D}
  \item \textbf{C}
  \item \textbf{D}
  \item \textbf{C}
  \item \textbf{C}
  \item \textbf{C}
  \item \textbf{A}
  \item \textbf{C}
  \item \textbf{D}
  \item \textbf{C}
  \item \textbf{D}
  \item \textbf{C}
  \item \textbf{B}
  \item \textbf{F}
  \item \textbf{C}
  \item \textbf{C}
  \item \textbf{A}
  \item \textbf{B}
  \item \textbf{B}
  \item \textbf{C}
  \item \textbf{D}
  \item \textbf{E}
  \item \textbf{D}
  \item \textbf{E}
  \item \textbf{B}
  \item \textbf{B}
  \item \textbf{C}
  \item \textbf{D}
  \item \textbf{E}
  \item \textbf{D}
  \item \textbf{E}
  \item \textbf{B}
  \item \textbf{B}
  \item \textbf{C}
  \item \textbf{D}
  \item \textbf{E}
  \item \textbf{D}
  \item \textbf{E}
  \item \textbf{B}
  \item \textbf{B}
  \item \textbf{C}
  \item \textbf{D}
  \item \textbf{E}
  \item \textbf{D}
  \item \textbf{E}
  \item \textbf{B}
  \item \textbf{B}
  \item \textbf{C}
  \item \textbf{D}
  \item \textbf{E}
  \item \textbf{D}
  \item \textbf{E}
  \item \textbf{B}
  \item \textbf{B}
  \item \textbf{C}
  \item \textbf{D}
  \item \textbf{E}
  \item \textbf{F}
  \item \textbf{F}
  \item \textbf{B}
  \item \textbf{B}
  \item \textbf{C}
  \item \textbf{B}
  \item \textbf{B}
  \item \textbf{B}
  \item \textbf{C}
  \item \textbf{B}
  \item \textbf{A}
  \item \textbf{B}
  \item \textbf{C}
  \item \textbf{A}
  \item \textbf{B}
  \item \textbf{A}
  \item \textbf{B}
  \item \textbf{B}
  \item \textbf{C}
  \item \textbf{A}
  \item \textbf{A}
  \item \textbf{B}
  \item \textbf{D}
  \item \textbf{B}
  \item \textbf{F}
  \item \textbf{A}
  \item \textbf{B}
  \item \textbf{B}
  \item \textbf{A}
  \item \textbf{B}
  \item \textbf{B}
  \item \textbf{B}
  \item \textbf{B}
  \item \textbf{B}
  \item \textbf{C}
  \item \textbf{A}
  \item \textbf{B}
  \item \textbf{A}
  \item \textbf{B}
  \item \textbf{A}
  \item \textbf{D}
  \item \textbf{C}
  \item \textbf{D}
  \item \textbf{E}
  \item \textbf{F}
  \item \textbf{D}
  \item \textbf{B}
  \item \textbf{C}
  \item \textbf{C}
  \item \textbf{B}
  \item \textbf{A}
  \item \textbf{B}
  \item \textbf{A}
  \item \textbf{B}
  \item \textbf{C}
  \item \textbf{C}
  \item \textbf{B}
  \item \textbf{A}
  \item \textbf{B}
  \item \textbf{C}
  \item \textbf{A}
  \item \textbf{A}
  \item \textbf{B}
  \item \textbf{A}
  \item \textbf{D}
  \item \textbf{A}
  \item \textbf{A}
  \item \textbf{F}
  \item \textbf{D}
  \item \textbf{C}
  \item \textbf{B}
  \item \textbf{B}
  \item \textbf{B}
  \item \textbf{A}
  \item \textbf{B}
  \item \textbf{A}
  \item \textbf{B}
  \item \textbf{B}
  \item \textbf{A}
  \item \textbf{B}
  \item \textbf{A}
  \item \textbf{B}
  \item \textbf{A}
  \item \textbf{B}
  \item \textbf{D}
  \item \textbf{C}
  \item \textbf{E}
  \item \textbf{D}
  \item \textbf{C}
  \item \textbf{C}
  \item \textbf{B}
  \item \textbf{A}
  \item \textbf{C}
  \item \textbf{C}
  \item \textbf{B}
  \item \textbf{A}
  \item \textbf{B}
  \item \textbf{A}
  \item \textbf{A}
  \item \textbf{A}
  \item \textbf{A}
  \item \textbf{B}
  \item \textbf{A}
  \item \textbf{A}
  \item \textbf{B}
  \item \textbf{B}
  \item \textbf{B}
  \item \textbf{A}
  \item \textbf{B}
  \item \textbf{D}
  \item \textbf{A}
  \item \textbf{A}
  \item \textbf{B}
  \item \textbf{B}
  \item \textbf{B}
  \item \textbf{B}
  \item \textbf{A}
  \item \textbf{A}
  \item \textbf{B}
  \item \textbf{B}
  \item \textbf{A}
  \item \textbf{A}
  \item \textbf{A}
  \item \textbf{B}
  \item \textbf{A}
  \item \textbf{A}
  \item \textbf{B}
  \item \textbf{A}
  \item \textbf{B}
  \item \textbf{B}
  \item \textbf{B}
  \item \textbf{A}
  \item \textbf{B}
  \item \textbf{B}
  \item \textbf{B}
  \item \textbf{D}
  \item \textbf{B}
  \item \textbf{B}
  \item \textbf{A}
  \item \textbf{B}
  \item \textbf{B}
  \item \textbf{A}
  \item \textbf{A}
  \item \textbf{B}
  \item \textbf{A}
  \item \textbf{A}
  \item \textbf{A}
  \item \textbf{A}
  \item \textbf{A}
  \item \textbf{A}
  \item \textbf{A}
  \item \textbf{A}
  \item \textbf{A}
  \item \textbf{A}
  \item \textbf{A}
  \item \textbf{A}
  \item \textbf{A}
  \item \textbf{A}
  \item \textbf{A}
  \item \textbf{A}
  \item \textbf{B}
  \item \textbf{C}
  \item \textbf{B}
  \item \textbf{C}
  \item \textbf{B}
  \item \textbf{A}
  \item \textbf{C}
  \item \textbf{D}
  \item \textbf{A}
  \item \textbf{B}
  \item \textbf{C}
  \item \textbf{D}
  \item \textbf{B}
  \item \textbf{D}
  \item \textbf{A}
  \item \textbf{C}
  \item \textbf{B}
  \item \textbf{A}
  \item \textbf{A}
  \item \textbf{B}
  \item \textbf{B}
  \item \textbf{A}
  \item \textbf{E}
  \item \textbf{B}
  \item \textbf{B}
  \item \textbf{A}
  \item \textbf{C}
  \item \textbf{B}
  \item \textbf{B}
  \item \textbf{B}
  \item \textbf{D}
  \item \textbf{B}
  \item \textbf{B}
  \item \textbf{B}
  \item \textbf{A}
  \item \textbf{B}
  \item \textbf{D}
  \item \textbf{A}
  \item \textbf{B}
  \item \textbf{D}
\end{enumerate}
\end{multicols}
\end{tcolorbox}
\newpage\n
\subsection{1.2 La Recta Real}

\subsubsection*{SECCIÓN I: CONCEPTOS BÁSICOS DE LA RECTA REAL (1–30)}

\begin{enumerate}[label=\textbf{\arabic*.}]
\setcounter{enumi}{0}
\item ¿Cómo se llama el punto que corresponde al número \(0\) en la recta real?
\begin{enumerate}[label=\textbf{\Alph*)}]
\begin{multicols}{2}
  \item Origen
  \item Extremo
  \item Vértice
  \item 2
  \item -1
  \item Ninguna
\end{multicols}
\end{enumerate}
\vspace{0.15cm}
\item ¿Cuál es la coordenada del origen en la recta real?
\begin{enumerate}[label=\textbf{\Alph*)}]
\begin{multicols}{2}
  \item -5
  \item 0
  \item 0
  \item 2
  \item -1
  \item No existe
\end{multicols}
\end{enumerate}
\vspace{0.15cm}
\item En la recta real, los números positivos se ubican:
\begin{enumerate}[label=\textbf{\Alph*)}]
\begin{multicols}{2}
  \item -5
  \item A la derecha del origen
  \item 1
  \item 2
  \item -1
  \item Ninguna
\end{multicols}
\end{enumerate}
\vspace{0.15cm}
\item Los números negativos se ubican:
\begin{enumerate}[label=\textbf{\Alph*)}]
\begin{multicols}{2}
  \item -5
  \item A la izquierda del origen
  \item 1
  \item 2
  \item -1
  \item Ninguna
\end{multicols}
\end{enumerate}
\vspace{0.15cm}
\item ¿Qué establece la correspondencia biunívoca en la recta real?
\begin{enumerate}[label=\textbf{\Alph*)}]
\begin{multicols}{2}
  \item Cada punto corresponde a un número real y viceversa
  \item 0
  \item 1
  \item 2
  \item -1
  \item Ninguna
\end{multicols}
\end{enumerate}
\vspace{0.15cm}
\item La distancia entre dos puntos en la recta real se calcula como:
\begin{enumerate}[label=\textbf{\Alph*)}]
\begin{multicols}{2}
  \item -5
  \item Diferencia en valor absoluto
  \item 1
  \item 2
  \item -1
  \item Ninguna
\end{multicols}
\end{enumerate}
\vspace{0.15cm}
\item ¿Cuál es el sentido positivo de la recta real?
\begin{enumerate}[label=\textbf{\Alph*)}]
\begin{multicols}{2}
  \item -5
  \item Hacia la derecha
  \item 1
  \item 2
  \item -1
  \item Ninguna
\end{multicols}
\end{enumerate}
\vspace{0.15cm}
\item ¿Cuál es el sentido negativo de la recta real?
\begin{enumerate}[label=\textbf{\Alph*)}]
\begin{multicols}{2}
  \item Hacia la izquierda
  \item 0
  \item 1
  \item 2
  \item -1
  \item Ninguna
\end{multicols}
\end{enumerate}
\vspace{0.15cm}
\item La escala en la recta real debe ser:
\begin{enumerate}[label=\textbf{\Alph*)}]
\begin{multicols}{2}
  \item -5
  \item Uniforme
  \item 1
  \item 2
  \item -1
  \item Ninguna
\end{multicols}
\end{enumerate}
\vspace{0.15cm}
\item ¿Qué número está a la misma distancia del origen que \(5\)?
\begin{enumerate}[label=\textbf{\Alph*)}]
\begin{multicols}{2}
  \item -5
  \item -5
  \item 1
  \item 2
  \item -1
  \item Ninguna
\end{multicols}
\end{enumerate}
\vspace{0.15cm}
\item ¿Qué número está a la misma distancia del origen que \(-3\)?
\begin{enumerate}[label=\textbf{\Alph*)}]
\begin{multicols}{2}
  \item 3
  \item 0
  \item 1
  \item 2
  \item -1
  \item Ninguna
\end{multicols}
\end{enumerate}
\vspace{0.15cm}
\item El opuesto de un número \(a\) en la recta real es:
\begin{enumerate}[label=\textbf{\Alph*)}]
\begin{multicols}{2}
  \item -5
  \item -a
  \item 1
  \item 2
  \item -1
  \item Ninguna
\end{multicols}
\end{enumerate}
\vspace{0.15cm}
\item ¿Cómo se representa el conjunto de todos los números reales?
\begin{enumerate}[label=\textbf{\Alph*)}]
\begin{multicols}{2}
  \item -5
  \item 0
  \item 1
  \item \mathbb{R}
  \item \mathbb{C}
  \item Ninguna
\end{multicols}
\end{enumerate}
\vspace{0.15cm}
\item ¿Qué número corresponde al punto medio entre \(0\) y \(2\)?
\begin{enumerate}[label=\textbf{\Alph*)}]
\begin{multicols}{2}
  \item -5
  \item 1
  \item 1
  \item 2
  \item -1
  \item Ninguna
\end{multicols}
\end{enumerate}
\vspace{0.15cm}
\item ¿Qué número corresponde al punto medio entre \(-4\) y \(4\)?
\begin{enumerate}[label=\textbf{\Alph*)}]
\begin{multicols}{2}
  \item -5
  \item 0
  \item 1
  \item 2
  \item -1
  \item Ninguna
\end{multicols}
\end{enumerate}
\vspace{0.15cm}
\item Si un punto está a \(7\) unidades a la derecha del origen, su coordenada es:
\begin{enumerate}[label=\textbf{\Alph*)}]
\begin{multicols}{2}
  \item -5
  \item 7
  \item 1
  \item 2
  \item -1
  \item Ninguna
\end{multicols}
\end{enumerate}
\vspace{0.15cm}
\item Si un punto está a \(7\) unidades a la izquierda del origen, su coordenada es:
\begin{enumerate}[label=\textbf{\Alph*)}]
\begin{multicols}{2}
  \item -7
  \item 7
  \item 1
  \item 2
  \item -1
  \item Ninguna
\end{multicols}
\end{enumerate}
\vspace{0.15cm}
\item ¿Qué número está más a la derecha en la recta real?
\begin{enumerate}[label=\textbf{\Alph*)}]
\begin{multicols}{2}
  \item -5
  \item 0
  \item 1
  \item 3
  \item -1
  \item Ninguna
\end{multicols}
\end{enumerate}
\vspace{0.15cm}
\item ¿Qué número está más a la izquierda en la recta real?
\begin{enumerate}[label=\textbf{\Alph*)}]
\begin{multicols}{2}
  \item -5
  \item 0
  \item 1
  \item 2
  \item -1
  \item Ninguna
\end{multicols}
\end{enumerate}
\vspace{0.15cm}
\item La recta real es un conjunto:
\begin{enumerate}[label=\textbf{\Alph*)}]
\begin{multicols}{2}
  \item -5
  \item Infinito y continuo
  \item 1
  \item 2
  \item -1
  \item Ninguna
\end{multicols}
\end{enumerate}
\vspace{0.15cm}
\item ¿Qué tipo de números se pueden representar en la recta real?
\begin{enumerate}[label=\textbf{\Alph*)}]
\begin{multicols}{2}
  \item -5
  \item 0
  \item 1
  \item Todos los reales
  \item -1
  \item Ninguna
\end{multicols}
\end{enumerate}
\vspace{0.15cm}
\item La coordenada de un punto que está \(2\) unidades a la izquierda de \(5\) es:
\begin{enumerate}[label=\textbf{\Alph*)}]
\begin{multicols}{2}
  \item -7
  \item 3
  \item 1
  \item 2
  \item -1
  \item Ninguna
\end{multicols}
\end{enumerate}
\vspace{0.15cm}
\item La coordenada de un punto que está \(2\) unidades a la derecha de \(-1\) es:
\begin{enumerate}[label=\textbf{\Alph*)}]
\begin{multicols}{2}
  \item -5
  \item 1
  \item 1
  \item 2
  \item -1
  \item Ninguna
\end{multicols}
\end{enumerate}
\vspace{0.15cm}
\item ¿Cuál es la distancia entre \(0\) y \(-6\)?
\begin{enumerate}[label=\textbf{\Alph*)}]
\begin{multicols}{2}
  \item -5
  \item 6
  \item 1
  \item 2
  \item -1
  \item Ninguna
\end{multicols}
\end{enumerate}
\vspace{0.15cm}
\item ¿Cuál es la distancia entre \(3\) y \(7\)?
\begin{enumerate}[label=\textbf{\Alph*)}]
\begin{multicols}{2}
  \item -5
  \item 4
  \item 1
  \item 2
  \item -1
  \item Ninguna
\end{multicols}
\end{enumerate}
\vspace{0.15cm}
\item ¿Cuál es el punto medio entre \(1\) y \(5\)?
\begin{enumerate}[label=\textbf{\Alph*)}]
\begin{multicols}{2}
  \item -5
  \item 3
  \item 1
  \item 2
  \item -1
  \item Ninguna
\end{multicols}
\end{enumerate}
\vspace{0.15cm}
\item ¿Cuál es el punto medio entre \(-3\) y \(3\)?
\begin{enumerate}[label=\textbf{\Alph*)}]
\begin{multicols}{2}
  \item -5
  \item 0
  \item 1
  \item 2
  \item -1
  \item Ninguna
\end{multicols}
\end{enumerate}
\vspace{0.15cm}
\item Si \(a<b\), entonces en la recta real \(a\) está:
\begin{enumerate}[label=\textbf{\Alph*)}]
\begin{multicols}{2}
  \item -5
  \item A la izquierda de b
  \item 1
  \item 2
  \item -1
  \item Ninguna
\end{multicols}
\end{enumerate}
\vspace{0.15cm}
\item Si \(a>b\), entonces \(a\) está:
\begin{enumerate}[label=\textbf{\Alph*)}]
\begin{multicols}{2}
  \item A la derecha de b
  \item 0
  \item 1
  \item 2
  \item -1
  \item Ninguna
\end{multicols}
\end{enumerate}
\vspace{0.15cm}
\item La distancia entre \(a\) y \(b\) se denota como:
\begin{enumerate}[label=\textbf{\Alph*)}]
\begin{multicols}{2}
  \item -5
  \item |a-b|
  \item 1
  \item 2
  \item -1
  \item Ninguna
\end{multicols}
\end{enumerate}
\vspace{0.15cm}
\end{enumerate}\n
\subsubsection*{SECCIÓN II: ORDEN Y UBICACIÓN (31–60)}

\begin{enumerate}[label=\textbf{\arabic*.}]
\setcounter{enumi}{30}
\item ¿Cuál de los siguientes números está a la derecha de \(-2\)?
\begin{enumerate}[label=\textbf{\Alph*)}]
\begin{multicols}{2}
  \item -5
  \item 0
  \item -1
  \item 4
  \item -2
  \item -4
\end{multicols}
\end{enumerate}
\vspace{0.15cm}
\item ¿Cuál de los siguientes números está a la izquierda de \(3\)?
\begin{enumerate}[label=\textbf{\Alph*)}]
\begin{multicols}{2}
  \item -5
  \item 0
  \item 2
  \item 4
  \item -2
  \item -4
\end{multicols}
\end{enumerate}
\vspace{0.15cm}
\item ¿Cuál de los siguientes números está entre \(-1\) y \(1\)?
\begin{enumerate}[label=\textbf{\Alph*)}]
\begin{multicols}{2}
  \item -5
  \item 0
  \item 0
  \item 4
  \item -2
  \item -4
\end{multicols}
\end{enumerate}
\vspace{0.15cm}
\item ¿Cuál de los siguientes números está entre \(2\) y \(5\)?
\begin{enumerate}[label=\textbf{\Alph*)}]
\begin{multicols}{2}
  \item -5
  \item 0
  \item 3
  \item 4
  \item -2
  \item -4
\end{multicols}
\end{enumerate}
\vspace{0.15cm}
\item ¿Cuál de los siguientes números está a la derecha de \(0\)?
\begin{enumerate}[label=\textbf{\Alph*)}]
\begin{multicols}{2}
  \item -5
  \item 0
  \item 1
  \item 4
  \item -2
  \item -4
\end{multicols}
\end{enumerate}
\vspace{0.15cm}
\item ¿Cuál de los siguientes números está a la izquierda de \(0\)?
\begin{enumerate}[label=\textbf{\Alph*)}]
\begin{multicols}{2}
  \item -5
  \item 0
  \item -1
  \item 4
  \item -2
  \item -4
\end{multicols}
\end{enumerate}
\vspace{0.15cm}
\item ¿Cuál es el orden correcto de menor a mayor?
\begin{enumerate}[label=\textbf{\Alph*)}]
\begin{multicols}{2}
  \item -3, -1, 0, 2
  \item 0
  \item 2
  \item 4
  \item -2
  \item -4
\end{multicols}
\end{enumerate}
\vspace{0.15cm}
\item ¿Cuál es el orden correcto de mayor a menor?
\begin{enumerate}[label=\textbf{\Alph*)}]
\begin{multicols}{2}
  \item -5
  \item 3, 0, -2, -5
  \item 2
  \item 4
  \item -2
  \item -4
\end{multicols}
\end{enumerate}
\vspace{0.15cm}
\item ¿Qué número está más cerca del origen?
\begin{enumerate}[label=\textbf{\Alph*)}]
\begin{multicols}{2}
  \item -5
  \item 0
  \item -1
  \item 4
  \item -2
  \item -4
\end{multicols}
\end{enumerate}
\vspace{0.15cm}
\item ¿Qué número está más lejos del origen?
\begin{enumerate}[label=\textbf{\Alph*)}]
\begin{multicols}{2}
  \item -5
  \item 0
  \item 2
  \item 4
  \item -2
  \item -4
\end{multicols}
\end{enumerate}
\vspace{0.15cm}
\item ¿Cuál de los siguientes números es mayor?
\begin{enumerate}[label=\textbf{\Alph*)}]
\begin{multicols}{2}
  \item -5
  \item 0
  \item -1
  \item 4
  \item -2
  \item -4
\end{multicols}
\end{enumerate}
\vspace{0.15cm}
\item ¿Cuál de los siguientes números es menor?
\begin{enumerate}[label=\textbf{\Alph*)}]
\begin{multicols}{2}
  \item -5
  \item 0
  \item 2
  \item -1
  \item -2
  \item -4
\end{multicols}
\end{enumerate}
\vspace{0.15cm}
\item Si \(x\) está a la derecha de \(-4\), entonces:
\begin{enumerate}[label=\textbf{\Alph*)}]
\begin{multicols}{2}
  \item x < -4
  \item x > -4
  \item x = -4
  \item x \leq -4
  \item x \geq -4
  \item Ninguna
\end{multicols}
\end{enumerate}
\vspace{0.15cm}
\item Si \(x\) está a la izquierda de \(2\), entonces:
\begin{enumerate}[label=\textbf{\Alph*)}]
\begin{multicols}{2}
  \item x > 2
  \item x < 2
  \item x = 2
  \item x \leq 2
  \item x \geq 2
  \item Ninguna
\end{multicols}
\end{enumerate}
\vspace{0.15cm}
\item ¿Cuál de los siguientes números NO está entre \(-3\) y \(3\)?
\begin{enumerate}[label=\textbf{\Alph*)}]
\begin{multicols}{2}
  \item -5
  \item 0
  \item 2
  \item 4
  \item -2
  \item -4
\end{multicols}
\end{enumerate}
\vspace{0.15cm}
\item ¿Cuál de los siguientes números está entre \(-2\) y \(0\)?
\begin{enumerate}[label=\textbf{\Alph*)}]
\begin{multicols}{2}
  \item -1
  \item 0
  \item 2
  \item 4
  \item -2
  \item -4
\end{multicols}
\end{enumerate}
\vspace{0.15cm}
\item ¿Cuál de los siguientes números está entre \(1\) y \(4\)?
\begin{enumerate}[label=\textbf{\Alph*)}]
\begin{multicols}{2}
  \item -5
  \item 0
  \item 2
  \item 4
  \item -2
  \item -4
\end{multicols}
\end{enumerate}
\vspace{0.15cm}
\item ¿Cuál de los siguientes números está a la derecha de \(-5\) pero a la izquierda de \(-2\)?
\begin{enumerate}[label=\textbf{\Alph*)}]
\begin{multicols}{2}
  \item -5
  \item 0
  \item -3
  \item 4
  \item -2
  \item -4
\end{multicols}
\end{enumerate}
\vspace{0.15cm}
\item ¿Cuál de los siguientes números está a la derecha de \(2\) pero a la izquierda de \(5\)?
\begin{enumerate}[label=\textbf{\Alph*)}]
\begin{multicols}{2}
  \item -5
  \item 0
  \item 3
  \item 4
  \item -2
  \item -4
\end{multicols}
\end{enumerate}
\vspace{0.15cm}
\item ¿Cuál es el número que está exactamente a la mitad entre \(-2\) y \(4\)?
\begin{enumerate}[label=\textbf{\Alph*)}]
\begin{multicols}{2}
  \item -5
  \item 1
  \item 2
  \item 4
  \item -2
  \item -4
\end{multicols}
\end{enumerate}
\vspace{0.15cm}
\item ¿Cuál es el número que está exactamente a la mitad entre \(-6\) y \(2\)?
\begin{enumerate}[label=\textbf{\Alph*)}]
\begin{multicols}{2}
  \item -5
  \item -2
  \item 2
  \item 4
  \item -2
  \item -4
\end{multicols}
\end{enumerate}
\vspace{0.15cm}
\item ¿Cuál es el número que está exactamente a la mitad entre \(1\) y \(7\)?
\begin{enumerate}[label=\textbf{\Alph*)}]
\begin{multicols}{2}
  \item -5
  \item 4
  \item 2
  \item 4
  \item -2
  \item -4
\end{multicols}
\end{enumerate}
\vspace{0.15cm}
\item ¿Cuál de los siguientes números es el opuesto de \(-3\)?
\begin{enumerate}[label=\textbf{\Alph*)}]
\begin{multicols}{2}
  \item -5
  \item 3
  \item 2
  \item 4
  \item -2
  \item -4
\end{multicols}
\end{enumerate}
\vspace{0.15cm}
\item ¿Cuál de los siguientes números es el opuesto de \(4\)?
\begin{enumerate}[label=\textbf{\Alph*)}]
\begin{multicols}{2}
  \item -5
  \item -4
  \item 2
  \item 4
  \item -2
  \item -4
\end{multicols}
\end{enumerate}
\vspace{0.15cm}
\item ¿Cuál es el simétrico de \(-2\) respecto al origen?
\begin{enumerate}[label=\textbf{\Alph*)}]
\begin{multicols}{2}
  \item -5
  \item 2
  \item 2
  \item 4
  \item -2
  \item -4
\end{multicols}
\end{enumerate}
\vspace{0.15cm}
\item ¿Cuál es el simétrico de \(5\) respecto al origen?
\begin{enumerate}[label=\textbf{\Alph*)}]
\begin{multicols}{2}
  \item -5
  \item -5
  \item 2
  \item 4
  \item -2
  \item -4
\end{multicols}
\end{enumerate}
\vspace{0.15cm}
\item Si un punto está a \(3\) unidades del origen, ¿cuáles son sus posibles coordenadas?
\begin{enumerate}[label=\textbf{\Alph*)}]
\begin{multicols}{2}
  \item 3 y -3
  \item 3 y 0
  \item 0 y -3
  \item 3 y 6
  \item -3 y -6
  \item 0 y 3
\end{multicols}
\end{enumerate}
\vspace{0.15cm}
\item Si un punto está a \(4\) unidades del origen, ¿cuáles son sus posibles coordenadas?
\begin{enumerate}[label=\textbf{\Alph*)}]
\begin{multicols}{2}
  \item 4 y -4
  \item 4 y 0
  \item 0 y -4
  \item 4 y 8
  \item -4 y -8
  \item 0 y 4
\end{multicols}
\end{enumerate}
\vspace{0.15cm}
\item La distancia entre \(-3\) y \(2\) es:
\begin{enumerate}[label=\textbf{\Alph*)}]
\begin{multicols}{2}
  \item -5
  \item 5
  \item 2
  \item 4
  \item -2
  \item -4
\end{multicols}
\end{enumerate}
\vspace{0.15cm}
\item La distancia entre \(-1\) y \(4\) es:
\begin{enumerate}[label=\textbf{\Alph*)}]
\begin{multicols}{2}
  \item -5
  \item 5
  \item 2
  \item 4
  \item -2
  \item -4
\end{multicols}
\end{enumerate}
\vspace{0.15cm}
\end{enumerate}\n
\subsubsection*{SECCIÓN III: COMPARACIÓN EN LA RECTA REAL (61–90)}

\begin{enumerate}[label=\textbf{\arabic*.}]
\setcounter{enumi}{60}
\item ¿Cuál símbolo es correcto? \(-2 _ 3\)
\begin{enumerate}[label=\textbf{\Alph*)}]
\begin{multicols}{2}
  \item <
  \item >
  \item =
  \item \leq
  \item \geq
  \item \neq
\end{multicols}
\end{enumerate}
\vspace{0.15cm}
\item ¿Cuál símbolo es correcto? \(-5 _ -1\)
\begin{enumerate}[label=\textbf{\Alph*)}]
\begin{multicols}{2}
  \item <
  \item >
  \item =
  \item \leq
  \item \geq
  \item \neq
\end{multicols}
\end{enumerate}
\vspace{0.15cm}
\item ¿Cuál símbolo es correcto? \(0 _ -3\)
\begin{enumerate}[label=\textbf{\Alph*)}]
\begin{multicols}{2}
  \item <
  \item >
  \item =
  \item \leq
  \item \geq
  \item \neq
\end{multicols}
\end{enumerate}
\vspace{0.15cm}
\item ¿Cuál símbolo es correcto? \(2 _ 2\)
\begin{enumerate}[label=\textbf{\Alph*)}]
\begin{multicols}{2}
  \item <
  \item >
  \item =
  \item \leq
  \item \geq
  \item \neq
\end{multicols}
\end{enumerate}
\vspace{0.15cm}
\item ¿Cuál símbolo es correcto? \(-4 _ -4\)
\begin{enumerate}[label=\textbf{\Alph*)}]
\begin{multicols}{2}
  \item <
  \item >
  \item =
  \item \leq
  \item \geq
  \item \neq
\end{multicols}
\end{enumerate}
\vspace{0.15cm}
\item ¿Cuál símbolo es correcto? \(3 _ -7\)
\begin{enumerate}[label=\textbf{\Alph*)}]
\begin{multicols}{2}
  \item <
  \item >
  \item =
  \item \leq
  \item \geq
  \item \neq
\end{multicols}
\end{enumerate}
\vspace{0.15cm}
\item ¿Cuál símbolo es correcto? \(-1 _ 0\)
\begin{enumerate}[label=\textbf{\Alph*)}]
\begin{multicols}{2}
  \item <
  \item >
  \item =
  \item \leq
  \item \geq
  \item \neq
\end{multicols}
\end{enumerate}
\vspace{0.15cm}
\item ¿Cuál símbolo es correcto? \(5 _ 1\)
\begin{enumerate}[label=\textbf{\Alph*)}]
\begin{multicols}{2}
  \item <
  \item >
  \item =
  \item \leq
  \item \geq
  \item \neq
\end{multicols}
\end{enumerate}
\vspace{0.15cm}
\item ¿Cuál símbolo es correcto? \(-6 _ -10\)
\begin{enumerate}[label=\textbf{\Alph*)}]
\begin{multicols}{2}
  \item <
  \item >
  \item =
  \item \leq
  \item \geq
  \item \neq
\end{multicols}
\end{enumerate}
\vspace{0.15cm}
\item ¿Cuál símbolo es correcto? \(-2 _ 0\)
\begin{enumerate}[label=\textbf{\Alph*)}]
\begin{multicols}{2}
  \item <
  \item >
  \item =
  \item \leq
  \item \geq
  \item \neq
\end{multicols}
\end{enumerate}
\vspace{0.15cm}
\item ¿Cuál de los siguientes números es mayor que -1?
\begin{enumerate}[label=\textbf{\Alph*)}]
\begin{multicols}{2}
  \item 0
  \item 0
  \item 3
  \item -1
  \item -2
  \item 2
\end{multicols}
\end{enumerate}
\vspace{0.15cm}
\item ¿Cuál de los siguientes números es menor que -2?
\begin{enumerate}[label=\textbf{\Alph*)}]
\begin{multicols}{2}
  \item 0
  \item -3
  \item 3
  \item -1
  \item -2
  \item 2
\end{multicols}
\end{enumerate}
\vspace{0.15cm}
\item ¿Cuál de los siguientes números es mayor que 2?
\begin{enumerate}[label=\textbf{\Alph*)}]
\begin{multicols}{2}
  \item 0
  \item 3
  \item 3
  \item -1
  \item -2
  \item 2
\end{multicols}
\end{enumerate}
\vspace{0.15cm}
\item ¿Cuál de los siguientes números es menor que 1?
\begin{enumerate}[label=\textbf{\Alph*)}]
\begin{multicols}{2}
  \item 0
  \item 0
  \item 3
  \item -1
  \item -2
  \item 2
\end{multicols}
\end{enumerate}
\vspace{0.15cm}
\item Si a = -3 y b = 1, entonces:
\begin{enumerate}[label=\textbf{\Alph*)}]
\begin{multicols}{2}
  \item a > b
  \item a < b
  \item a = b
  \item a \geq b
  \item a \leq b
  \item No se puede
\end{multicols}
\end{enumerate}
\vspace{0.15cm}
\item Si a = 4 y b = -2, entonces:
\begin{enumerate}[label=\textbf{\Alph*)}]
\begin{multicols}{2}
  \item a > b
  \item a < b
  \item a = b
  \item a \geq b
  \item a \leq b
  \item No se puede
\end{multicols}
\end{enumerate}
\vspace{0.15cm}
\item Si a = -5 y b = -5, entonces:
\begin{enumerate}[label=\textbf{\Alph*)}]
\begin{multicols}{2}
  \item a > b
  \item a < b
  \item a = b
  \item a \geq b
  \item a \leq b
  \item No se puede
\end{multicols}
\end{enumerate}
\vspace{0.15cm}
\item ¿Cuál de los siguientes pares está ordenado de menor a mayor?
\begin{enumerate}[label=\textbf{\Alph*)}]
\begin{multicols}{2}
  \item (3,-1)
  \item (-2,0)
  \item (2,1)
  \item (-1,-3)
  \item (0,-2)
  \item (4,2)
\end{multicols}
\end{enumerate}
\vspace{0.15cm}
\item ¿Cuál de los siguientes pares está ordenado de mayor a menor?
\begin{enumerate}[label=\textbf{\Alph*)}]
\begin{multicols}{2}
  \item (-1,3)
  \item (0,-2)
  \item (2,5)
  \item (-3,-1)
  \item (1,4)
  \item (5,2)
\end{multicols}
\end{enumerate}
\vspace{0.15cm}
\item ¿Cuál de los siguientes números es el mayor?
\begin{enumerate}[label=\textbf{\Alph*)}]
\begin{multicols}{2}
  \item -5
  \item -1
  \item 3
  \item -1
  \item -2
  \item 2
\end{multicols}
\end{enumerate}
\vspace{0.15cm}
\item ¿Cuál de los siguientes números es el menor?
\begin{enumerate}[label=\textbf{\Alph*)}]
\begin{multicols}{2}
  \item 0
  \item 1
  \item 3
  \item -1
  \item -2
  \item 2
\end{multicols}
\end{enumerate}
\vspace{0.15cm}
\item ¿Cuál de los siguientes números está entre -3 y -1?
\begin{enumerate}[label=\textbf{\Alph*)}]
\begin{multicols}{2}
  \item 0
  \item -2
  \item 3
  \item -1
  \item -2
  \item 2
\end{multicols}
\end{enumerate}
\vspace{0.15cm}
\item ¿Cuál de los siguientes números está entre 0 y 2?
\begin{enumerate}[label=\textbf{\Alph*)}]
\begin{multicols}{2}
  \item 0
  \item 1
  \item 3
  \item -1
  \item -2
  \item 2
\end{multicols}
\end{enumerate}
\vspace{0.15cm}
\item ¿Cuál de los siguientes números está entre -1 y 2?
\begin{enumerate}[label=\textbf{\Alph*)}]
\begin{multicols}{2}
  \item 0
  \item 0
  \item 3
  \item -1
  \item -2
  \item 2
\end{multicols}
\end{enumerate}
\vspace{0.15cm}
\item ¿Cuál de los siguientes números está entre -5 y -2?
\begin{enumerate}[label=\textbf{\Alph*)}]
\begin{multicols}{2}
  \item -5
  \item -3
  \item 3
  \item -1
  \item -2
  \item 2
\end{multicols}
\end{enumerate}
\vspace{0.15cm}
\item ¿Cuál de los siguientes números está entre 3 y 6?
\begin{enumerate}[label=\textbf{\Alph*)}]
\begin{multicols}{2}
  \item 0
  \item 4
  \item 3
  \item -1
  \item -2
  \item 2
\end{multicols}
\end{enumerate}
\vspace{0.15cm}
\item ¿Cuál de los siguientes números está entre -4 y 0?
\begin{enumerate}[label=\textbf{\Alph*)}]
\begin{multicols}{2}
  \item 0
  \item -2
  \item 3
  \item -1
  \item -2
  \item 2
\end{multicols}
\end{enumerate}
\vspace{0.15cm}
\item ¿Cuál de los siguientes números está entre 1 y 3?
\begin{enumerate}[label=\textbf{\Alph*)}]
\begin{multicols}{2}
  \item 0
  \item 2
  \item 3
  \item -1
  \item -2
  \item 2
\end{multicols}
\end{enumerate}
\vspace{0.15cm}
\item ¿Cuál de los siguientes números está entre -2 y 2?
\begin{enumerate}[label=\textbf{\Alph*)}]
\begin{multicols}{2}
  \item 0
  \item 0
  \item 3
  \item -1
  \item -2
  \item 2
\end{multicols}
\end{enumerate}
\vspace{0.15cm}
\item ¿Cuál de los siguientes números está entre -10 y -8?
\begin{enumerate}[label=\textbf{\Alph*)}]
\begin{multicols}{2}
  \item 0
  \item -9
  \item 3
  \item -1
  \item -2
  \item 2
\end{multicols}
\end{enumerate}
\vspace{0.15cm}
\end{enumerate}\n
\subsubsection*{SECCIÓN IV: DISTANCIA ENTRE PUNTOS (91–120)}

\begin{enumerate}[label=\textbf{\arabic*.}]
\setcounter{enumi}{90}
\item ¿Cuál es la distancia entre \(2\) y \(5\)?
\begin{enumerate}[label=\textbf{\Alph*)}]
\begin{multicols}{2}
  \item 3
  \item -3
  \item 6
  \item 5
  \item 2
  \item 1
\end{multicols}
\end{enumerate}
\vspace{0.15cm}
\item ¿Cuál es la distancia entre \(-3\) y \(1\)?
\begin{enumerate}[label=\textbf{\Alph*)}]
\begin{multicols}{2}
  \item 4
  \item -4
  \item 8
  \item 6
  \item 3
  \item 1
\end{multicols}
\end{enumerate}
\vspace{0.15cm}
\item ¿Cuál es la distancia entre \(-5\) y \(-2\)?
\begin{enumerate}[label=\textbf{\Alph*)}]
\begin{multicols}{2}
  \item 3
  \item -3
  \item 6
  \item 5
  \item 2
  \item 1
\end{multicols}
\end{enumerate}
\vspace{0.15cm}
\item ¿Cuál es la distancia entre \(0\) y \(4\)?
\begin{enumerate}[label=\textbf{\Alph*)}]
\begin{multicols}{2}
  \item 4
  \item -4
  \item 8
  \item 6
  \item 3
  \item 1
\end{multicols}
\end{enumerate}
\vspace{0.15cm}
\item ¿Cuál es la distancia entre \(-6\) y \(0\)?
\begin{enumerate}[label=\textbf{\Alph*)}]
\begin{multicols}{2}
  \item 6
  \item -6
  \item 12
  \item 8
  \item 5
  \item 1
\end{multicols}
\end{enumerate}
\vspace{0.15cm}
\item ¿Cuál es la distancia entre \(3\) y \(-3\)?
\begin{enumerate}[label=\textbf{\Alph*)}]
\begin{multicols}{2}
  \item 6
  \item -6
  \item 12
  \item 8
  \item 5
  \item 1
\end{multicols}
\end{enumerate}
\vspace{0.15cm}
\item ¿Cuál es la distancia entre \(-1\) y \(4\)?
\begin{enumerate}[label=\textbf{\Alph*)}]
\begin{multicols}{2}
  \item 5
  \item -5
  \item 10
  \item 7
  \item 4
  \item 1
\end{multicols}
\end{enumerate}
\vspace{0.15cm}
\item ¿Cuál es la distancia entre \(7\) y \(2\)?
\begin{enumerate}[label=\textbf{\Alph*)}]
\begin{multicols}{2}
  \item 5
  \item -5
  \item 10
  \item 7
  \item 4
  \item 1
\end{multicols}
\end{enumerate}
\vspace{0.15cm}
\item ¿Cuál es la distancia entre \(-4\) y \(-7\)?
\begin{enumerate}[label=\textbf{\Alph*)}]
\begin{multicols}{2}
  \item 3
  \item -3
  \item 6
  \item 5
  \item 2
  \item 1
\end{multicols}
\end{enumerate}
\vspace{0.15cm}
\item ¿Cuál es la distancia entre \(10\) y \(3\)?
\begin{enumerate}[label=\textbf{\Alph*)}]
\begin{multicols}{2}
  \item 7
  \item -7
  \item 14
  \item 9
  \item 6
  \item 1
\end{multicols}
\end{enumerate}
\vspace{0.15cm}
\item ¿Cuál es la distancia entre \(-8\) y \(2\)?
\begin{enumerate}[label=\textbf{\Alph*)}]
\begin{multicols}{2}
  \item 10
  \item -10
  \item 20
  \item 12
  \item 9
  \item 1
\end{multicols}
\end{enumerate}
\vspace{0.15cm}
\item ¿Cuál es la distancia entre \(5\) y \(-5\)?
\begin{enumerate}[label=\textbf{\Alph*)}]
\begin{multicols}{2}
  \item 10
  \item -10
  \item 20
  \item 12
  \item 9
  \item 1
\end{multicols}
\end{enumerate}
\vspace{0.15cm}
\item ¿Cuál es la distancia entre \(-9\) y \(-1\)?
\begin{enumerate}[label=\textbf{\Alph*)}]
\begin{multicols}{2}
  \item 8
  \item -8
  \item 16
  \item 10
  \item 7
  \item 1
\end{multicols}
\end{enumerate}
\vspace{0.15cm}
\item ¿Cuál es la distancia entre \(0\) y \(-7\)?
\begin{enumerate}[label=\textbf{\Alph*)}]
\begin{multicols}{2}
  \item 7
  \item -7
  \item 14
  \item 9
  \item 6
  \item 1
\end{multicols}
\end{enumerate}
\vspace{0.15cm}
\item ¿Cuál es la distancia entre \(6\) y \(6\)?
\begin{enumerate}[label=\textbf{\Alph*)}]
\begin{multicols}{2}
  \item 0
  \item 6
  \item -6
  \item 2
  \item 1
  \item 1
\end{multicols}
\end{enumerate}
\vspace{0.15cm}
\item ¿Cuál es la distancia entre \(-2\) y \(-2\)?
\begin{enumerate}[label=\textbf{\Alph*)}]
\begin{multicols}{2}
  \item 0
  \item 6
  \item -6
  \item 2
  \item 1
  \item 1
\end{multicols}
\end{enumerate}
\vspace{0.15cm}
\item ¿Cuál es la distancia entre \(8\) y \(1\)?
\begin{enumerate}[label=\textbf{\Alph*)}]
\begin{multicols}{2}
  \item 7
  \item -7
  \item 14
  \item 9
  \item 6
  \item 1
\end{multicols}
\end{enumerate}
\vspace{0.15cm}
\item ¿Cuál es la distancia entre \(-3\) y \(-8\)?
\begin{enumerate}[label=\textbf{\Alph*)}]
\begin{multicols}{2}
  \item 5
  \item -5
  \item 10
  \item 7
  \item 4
  \item 1
\end{multicols}
\end{enumerate}
\vspace{0.15cm}
\item ¿Cuál es la distancia entre \(4\) y \(-1\)?
\begin{enumerate}[label=\textbf{\Alph*)}]
\begin{multicols}{2}
  \item 5
  \item -5
  \item 10
  \item 7
  \item 4
  \item 1
\end{multicols}
\end{enumerate}
\vspace{0.15cm}
\item ¿Cuál es la distancia entre \(-10\) y \(-5\)?
\begin{enumerate}[label=\textbf{\Alph*)}]
\begin{multicols}{2}
  \item 5
  \item -5
  \item 10
  \item 7
  \item 4
  \item 1
\end{multicols}
\end{enumerate}
\vspace{0.15cm}
\item Si \(A=2 y B=7\), ¿cuál es la distancia \(AB\)?
\begin{enumerate}[label=\textbf{\Alph*)}]
\begin{multicols}{2}
  \item 5
  \item -5
  \item 10
  \item 7
  \item 4
  \item 1
\end{multicols}
\end{enumerate}
\vspace{0.15cm}
\item Si \(A=-4 y B=3\), ¿cuál es la distancia \(AB\)?
\begin{enumerate}[label=\textbf{\Alph*)}]
\begin{multicols}{2}
  \item 7
  \item -7
  \item 14
  \item 9
  \item 6
  \item 1
\end{multicols}
\end{enumerate}
\vspace{0.15cm}
\item Si \(A=-1 y B=-6\), ¿cuál es la distancia \(AB\)?
\begin{enumerate}[label=\textbf{\Alph*)}]
\begin{multicols}{2}
  \item 5
  \item -5
  \item 10
  \item 7
  \item 4
  \item 1
\end{multicols}
\end{enumerate}
\vspace{0.15cm}
\item Si \(A=0 y B=-9\), ¿cuál es la distancia \(AB\)?
\begin{enumerate}[label=\textbf{\Alph*)}]
\begin{multicols}{2}
  \item 9
  \item -9
  \item 18
  \item 11
  \item 8
  \item 1
\end{multicols}
\end{enumerate}
\vspace{0.15cm}
\item Si \(A=5 y B=-2\), ¿cuál es la distancia \(AB\)?
\begin{enumerate}[label=\textbf{\Alph*)}]
\begin{multicols}{2}
  \item 7
  \item -7
  \item 14
  \item 9
  \item 6
  \item 1
\end{multicols}
\end{enumerate}
\vspace{0.15cm}
\item Si \(A=-7 y B=-3\), ¿cuál es la distancia \(AB\)?
\begin{enumerate}[label=\textbf{\Alph*)}]
\begin{multicols}{2}
  \item 4
  \item -4
  \item 8
  \item 6
  \item 3
  \item 1
\end{multicols}
\end{enumerate}
\vspace{0.15cm}
\item Si \(A=6 y B=2\), ¿cuál es la distancia \(AB\)?
\begin{enumerate}[label=\textbf{\Alph*)}]
\begin{multicols}{2}
  \item 4
  \item -4
  \item 8
  \item 6
  \item 3
  \item 1
\end{multicols}
\end{enumerate}
\vspace{0.15cm}
\item Si \(A=-2 y B=5\), ¿cuál es la distancia \(AB\)?
\begin{enumerate}[label=\textbf{\Alph*)}]
\begin{multicols}{2}
  \item 7
  \item -7
  \item 14
  \item 9
  \item 6
  \item 1
\end{multicols}
\end{enumerate}
\vspace{0.15cm}
\item Si \(A=1 y B=-4\), ¿cuál es la distancia \(AB\)?
\begin{enumerate}[label=\textbf{\Alph*)}]
\begin{multicols}{2}
  \item 5
  \item -5
  \item 10
  \item 7
  \item 4
  \item 1
\end{multicols}
\end{enumerate}
\vspace{0.15cm}
\item Si \(A=-5 y B=0\), ¿cuál es la distancia \(AB\)?
\begin{enumerate}[label=\textbf{\Alph*)}]
\begin{multicols}{2}
  \item 5
  \item -5
  \item 10
  \item 7
  \item 4
  \item 1
\end{multicols}
\end{enumerate}
\vspace{0.15cm}
\end{enumerate}\n
\subsubsection*{SECCIÓN V: PUNTO MEDIO (121–150)}

\begin{enumerate}[label=\textbf{\arabic*.}]
\setcounter{enumi}{120}
\item ¿Cuál es el punto medio entre \(2\) y \(6\)?
\begin{enumerate}[label=\textbf{\Alph*)}]
\begin{multicols}{2}
  \item 3
  \item 4
  \item 6
  \item -4
  \item 0
  \item 1
\end{multicols}
\end{enumerate}
\vspace{0.15cm}
\item ¿Cuál es el punto medio entre \(-2\) y \(4\)?
\begin{enumerate}[label=\textbf{\Alph*)}]
\begin{multicols}{2}
  \item 0
  \item 1
  \item 3
  \item -1
  \item 0
  \item -1
\end{multicols}
\end{enumerate}
\vspace{0.15cm}
\item ¿Cuál es el punto medio entre \(-4\) y \(2\)?
\begin{enumerate}[label=\textbf{\Alph*)}]
\begin{multicols}{2}
  \item -2
  \item -1
  \item 1
  \item 1
  \item 0
  \item 1
\end{multicols}
\end{enumerate}
\vspace{0.15cm}
\item ¿Cuál es el punto medio entre \(1\) y \(5\)?
\begin{enumerate}[label=\textbf{\Alph*)}]
\begin{multicols}{2}
  \item 2
  \item 3
  \item 5
  \item -3
  \item 0
  \item 1
\end{multicols}
\end{enumerate}
\vspace{0.15cm}
\item ¿Cuál es el punto medio entre \(-3\) y \(3\)?
\begin{enumerate}[label=\textbf{\Alph*)}]
\begin{multicols}{2}
  \item -1
  \item 0
  \item 2
  \item 6
  \item 1
  \item 1
\end{multicols}
\end{enumerate}
\vspace{0.15cm}
\item ¿Cuál es el punto medio entre \(0\) y \(8\)?
\begin{enumerate}[label=\textbf{\Alph*)}]
\begin{multicols}{2}
  \item 4
  \item 5
  \item 6
  \item -4
  \item 0
  \item 1
\end{multicols}
\end{enumerate}
\vspace{0.15cm}
\item ¿Cuál es el punto medio entre \(-6\) y \(0\)?
\begin{enumerate}[label=\textbf{\Alph*)}]
\begin{multicols}{2}
  \item -3
  \item -2
  \item -1
  \item 3
  \item 0
  \item 1
\end{multicols}
\end{enumerate}
\vspace{0.15cm}
\item ¿Cuál es el punto medio entre \(-5\) y \(-1\)?
\begin{enumerate}[label=\textbf{\Alph*)}]
\begin{multicols}{2}
  \item -3
  \item -2
  \item -1
  \item 3
  \item 0
  \item 1
\end{multicols}
\end{enumerate}
\vspace{0.15cm}
\item ¿Cuál es el punto medio entre \(3\) y \(9\)?
\begin{enumerate}[label=\textbf{\Alph*)}]
\begin{multicols}{2}
  \item 6
  \item 7
  \item 8
  \item -6
  \item 0
  \item 1
\end{multicols}
\end{enumerate}
\vspace{0.15cm}
\item ¿Cuál es el punto medio entre \(-7\) y \(1\)?
\begin{enumerate}[label=\textbf{\Alph*)}]
\begin{multicols}{2}
  \item -3
  \item -2
  \item -1
  \item 3
  \item 0
  \item 1
\end{multicols}
\end{enumerate}
\vspace{0.15cm}
\item ¿Cuál es el punto medio entre \(10\) y \(2\)?
\begin{enumerate}[label=\textbf{\Alph*)}]
\begin{multicols}{2}
  \item 6
  \item 7
  \item 8
  \item -6
  \item 0
  \item 1
\end{multicols}
\end{enumerate}
\vspace{0.15cm}
\item ¿Cuál es el punto medio entre \(-8\) y \(-4\)?
\begin{enumerate}[label=\textbf{\Alph*)}]
\begin{multicols}{2}
  \item -6
  \item -5
  \item -4
  \item 6
  \item 0
  \item 1
\end{multicols}
\end{enumerate}
\vspace{0.15cm}
\item ¿Cuál es el punto medio entre \(1\) y \(-5\)?
\begin{enumerate}[label=\textbf{\Alph*)}]
\begin{multicols}{2}
  \item -2
  \item -1
  \item 0
  \item 2
  \item 0
  \item 1
\end{multicols}
\end{enumerate}
\vspace{0.15cm}
\item ¿Cuál es el punto medio entre \(-1\) y \(7\)?
\begin{enumerate}[label=\textbf{\Alph*)}]
\begin{multicols}{2}
  \item 3
  \item 4
  \item 5
  \item -3
  \item 0
  \item 1
\end{multicols}
\end{enumerate}
\vspace{0.15cm}
\item ¿Cuál es el punto medio entre \(4\) y \(10\)?
\begin{enumerate}[label=\textbf{\Alph*)}]
\begin{multicols}{2}
  \item 7
  \item 8
  \item 9
  \item -7
  \item 0
  \item 1
\end{multicols}
\end{enumerate}
\vspace{0.15cm}
\item ¿Cuál es el punto medio entre \(-10\) y \(0\)?
\begin{enumerate}[label=\textbf{\Alph*)}]
\begin{multicols}{2}
  \item -5
  \item -4
  \item -3
  \item 5
  \item 0
  \item 1
\end{multicols}
\end{enumerate}
\vspace{0.15cm}
\item ¿Cuál es el punto medio entre \(-2\) y \(-8\)?
\begin{enumerate}[label=\textbf{\Alph*)}]
\begin{multicols}{2}
  \item -5
  \item -4
  \item -3
  \item 5
  \item 0
  \item 1
\end{multicols}
\end{enumerate}
\vspace{0.15cm}
\item ¿Cuál es el punto medio entre \(6\) y \(-2\)?
\begin{enumerate}[label=\textbf{\Alph*)}]
\begin{multicols}{2}
  \item 2
  \item 3
  \item 4
  \item -2
  \item 0
  \item 1
\end{multicols}
\end{enumerate}
\vspace{0.15cm}
\item ¿Cuál es el punto medio entre \(-3\) y \(5\)?
\begin{enumerate}[label=\textbf{\Alph*)}]
\begin{multicols}{2}
  \item 1
  \item 2
  \item 3
  \item -1
  \item 0
  \item -1
\end{multicols}
\end{enumerate}
\vspace{0.15cm}
\item ¿Cuál es el punto medio entre \(8\) y \(0\)?
\begin{enumerate}[label=\textbf{\Alph*)}]
\begin{multicols}{2}
  \item 4
  \item 5
  \item 6
  \item -4
  \item 0
  \item 1
\end{multicols}
\end{enumerate}
\vspace{0.15cm}
\item Si \(A=1 y B=9\), el punto medio es:
\begin{enumerate}[label=\textbf{\Alph*)}]
\begin{multicols}{2}
  \item 5
  \item 6
  \item 7
  \item -5
  \item 0
  \item 1
\end{multicols}
\end{enumerate}
\vspace{0.15cm}
\item Si \(A=-4 y B=6\), el punto medio es:
\begin{enumerate}[label=\textbf{\Alph*)}]
\begin{multicols}{2}
  \item 1
  \item 2
  \item 3
  \item -1
  \item 0
  \item -1
\end{multicols}
\end{enumerate}
\vspace{0.15cm}
\item Si \(A=-7 y B=-1\), el punto medio es:
\begin{enumerate}[label=\textbf{\Alph*)}]
\begin{multicols}{2}
  \item -4
  \item -3
  \item -2
  \item 4
  \item 0
  \item 1
\end{multicols}
\end{enumerate}
\vspace{0.15cm}
\item Si \(A=2 y B=-8\), el punto medio es:
\begin{enumerate}[label=\textbf{\Alph*)}]
\begin{multicols}{2}
  \item -3
  \item -2
  \item -1
  \item 3
  \item 0
  \item 1
\end{multicols}
\end{enumerate}
\vspace{0.15cm}
\item Si \(A=-5 y B=5\), el punto medio es:
\begin{enumerate}[label=\textbf{\Alph*)}]
\begin{multicols}{2}
  \item 0
  \item 1
  \item 2
  \item 6
  \item 1
  \item 1
\end{multicols}
\end{enumerate}
\vspace{0.15cm}
\item Si \(A=3 y B=3\), el punto medio es:
\begin{enumerate}[label=\textbf{\Alph*)}]
\begin{multicols}{2}
  \item 3
  \item 4
  \item 5
  \item -3
  \item 0
  \item 1
\end{multicols}
\end{enumerate}
\vspace{0.15cm}
\item Si \(A=-6 y B=-2\), el punto medio es:
\begin{enumerate}[label=\textbf{\Alph*)}]
\begin{multicols}{2}
  \item -4
  \item -3
  \item -2
  \item 4
  \item 0
  \item 1
\end{multicols}
\end{enumerate}
\vspace{0.15cm}
\item Si \(A=0 y B=10\), el punto medio es:
\begin{enumerate}[label=\textbf{\Alph*)}]
\begin{multicols}{2}
  \item 5
  \item 6
  \item 7
  \item -5
  \item 0
  \item 1
\end{multicols}
\end{enumerate}
\vspace{0.15cm}
\item Si \(A=-9 y B=3\), el punto medio es:
\begin{enumerate}[label=\textbf{\Alph*)}]
\begin{multicols}{2}
  \item -3
  \item -2
  \item -1
  \item 3
  \item 0
  \item 1
\end{multicols}
\end{enumerate}
\vspace{0.15cm}
\item Si \(A=7 y B=-1\), el punto medio es:
\begin{enumerate}[label=\textbf{\Alph*)}]
\begin{multicols}{2}
  \item 3
  \item 4
  \item 5
  \item -3
  \item 0
  \item 1
\end{multicols}
\end{enumerate}
\vspace{0.15cm}
\end{enumerate}\n
\subsubsection*{SECCIÓN VI: DESPLAZAMIENTOS EN LA RECTA REAL (151–180)}

\begin{enumerate}[label=\textbf{\arabic*.}]
\setcounter{enumi}{150}
\item ¿Qué número está 3 unidades a la derecha de 2?
\begin{enumerate}[label=\textbf{\Alph*)}]
\begin{multicols}{2}
  \item 5
  \item -5
  \item 0
  \item 7
  \item 1
  \item -1
\end{multicols}
\end{enumerate}
\vspace{0.15cm}
\item ¿Qué número está 4 unidades a la izquierda de 1?
\begin{enumerate}[label=\textbf{\Alph*)}]
\begin{multicols}{2}
  \item -3
  \item 3
  \item 0
  \item -1
  \item 1
  \item -1
\end{multicols}
\end{enumerate}
\vspace{0.15cm}
\item ¿Qué número está 2 unidades a la derecha de -3?
\begin{enumerate}[label=\textbf{\Alph*)}]
\begin{multicols}{2}
  \item -1
  \item 1
  \item 0
  \item 1
  \item 1
  \item -1
\end{multicols}
\end{enumerate}
\vspace{0.15cm}
\item ¿Qué número está 5 unidades a la izquierda de 2?
\begin{enumerate}[label=\textbf{\Alph*)}]
\begin{multicols}{2}
  \item -3
  \item 3
  \item 0
  \item -1
  \item 1
  \item -1
\end{multicols}
\end{enumerate}
\vspace{0.15cm}
\item ¿Qué número está 6 unidades a la derecha de -4?
\begin{enumerate}[label=\textbf{\Alph*)}]
\begin{multicols}{2}
  \item 2
  \item -2
  \item 0
  \item 4
  \item 1
  \item -1
\end{multicols}
\end{enumerate}
\vspace{0.15cm}
\item ¿Qué número está 3 unidades a la izquierda de -2?
\begin{enumerate}[label=\textbf{\Alph*)}]
\begin{multicols}{2}
  \item -5
  \item 5
  \item 0
  \item -3
  \item 1
  \item -1
\end{multicols}
\end{enumerate}
\vspace{0.15cm}
\item ¿Qué número está 7 unidades a la derecha de 0?
\begin{enumerate}[label=\textbf{\Alph*)}]
\begin{multicols}{2}
  \item 7
  \item -7
  \item 0
  \item 9
  \item 1
  \item -1
\end{multicols}
\end{enumerate}
\vspace{0.15cm}
\item ¿Qué número está 7 unidades a la izquierda de 0?
\begin{enumerate}[label=\textbf{\Alph*)}]
\begin{multicols}{2}
  \item -7
  \item 7
  \item 0
  \item -5
  \item 1
  \item -1
\end{multicols}
\end{enumerate}
\vspace{0.15cm}
\item ¿Qué número está 2 unidades a la derecha de 5?
\begin{enumerate}[label=\textbf{\Alph*)}]
\begin{multicols}{2}
  \item 7
  \item -7
  \item 0
  \item 9
  \item 1
  \item -1
\end{multicols}
\end{enumerate}
\vspace{0.15cm}
\item ¿Qué número está 4 unidades a la izquierda de 3?
\begin{enumerate}[label=\textbf{\Alph*)}]
\begin{multicols}{2}
  \item -1
  \item 1
  \item 0
  \item 1
  \item 1
  \item -1
\end{multicols}
\end{enumerate}
\vspace{0.15cm}
\item Si estás en -2 y te mueves 5 unidades a la derecha, ¿dónde quedas?
\begin{enumerate}[label=\textbf{\Alph*)}]
\begin{multicols}{2}
  \item 3
  \item -3
  \item 0
  \item 5
  \item 1
  \item -1
\end{multicols}
\end{enumerate}
\vspace{0.15cm}
\item Si estás en 4 y te mueves 6 unidades a la izquierda, ¿dónde quedas?
\begin{enumerate}[label=\textbf{\Alph*)}]
\begin{multicols}{2}
  \item -2
  \item 2
  \item 0
  \item 0
  \item 1
  \item -1
\end{multicols}
\end{enumerate}
\vspace{0.15cm}
\item Si estás en -1 y te mueves 3 unidades a la izquierda, ¿dónde quedas?
\begin{enumerate}[label=\textbf{\Alph*)}]
\begin{multicols}{2}
  \item -4
  \item 4
  \item 0
  \item -2
  \item 1
  \item -1
\end{multicols}
\end{enumerate}
\vspace{0.15cm}
\item Si estás en 0 y te mueves 8 unidades a la derecha, ¿dónde quedas?
\begin{enumerate}[label=\textbf{\Alph*)}]
\begin{multicols}{2}
  \item 8
  \item -8
  \item 0
  \item 10
  \item 1
  \item -1
\end{multicols}
\end{enumerate}
\vspace{0.15cm}
\item Si estás en 0 y te mueves 8 unidades a la izquierda, ¿dónde quedas?
\begin{enumerate}[label=\textbf{\Alph*)}]
\begin{multicols}{2}
  \item -8
  \item 8
  \item 0
  \item -6
  \item 1
  \item -1
\end{multicols}
\end{enumerate}
\vspace{0.15cm}
\item Si estás en -5 y te mueves 10 unidades a la derecha, ¿dónde quedas?
\begin{enumerate}[label=\textbf{\Alph*)}]
\begin{multicols}{2}
  \item 5
  \item -5
  \item 0
  \item 7
  \item 1
  \item -1
\end{multicols}
\end{enumerate}
\vspace{0.15cm}
\item Si estás en 7 y te mueves 10 unidades a la izquierda, ¿dónde quedas?
\begin{enumerate}[label=\textbf{\Alph*)}]
\begin{multicols}{2}
  \item -3
  \item 3
  \item 0
  \item -1
  \item 1
  \item -1
\end{multicols}
\end{enumerate}
\vspace{0.15cm}
\item Si estás en -3 y te mueves 2 unidades a la derecha, ¿dónde quedas?
\begin{enumerate}[label=\textbf{\Alph*)}]
\begin{multicols}{2}
  \item -1
  \item 1
  \item 0
  \item 1
  \item 1
  \item -1
\end{multicols}
\end{enumerate}
\vspace{0.15cm}
\item Si estás en 6 y te mueves 9 unidades a la izquierda, ¿dónde quedas?
\begin{enumerate}[label=\textbf{\Alph*)}]
\begin{multicols}{2}
  \item -3
  \item 3
  \item 0
  \item -1
  \item 1
  \item -1
\end{multicols}
\end{enumerate}
\vspace{0.15cm}
\item Si estás en -8 y te mueves 3 unidades a la derecha, ¿dónde quedas?
\begin{enumerate}[label=\textbf{\Alph*)}]
\begin{multicols}{2}
  \item -5
  \item 5
  \item 0
  \item -3
  \item 1
  \item -1
\end{multicols}
\end{enumerate}
\vspace{0.15cm}
\item ¿Cuántas unidades te mueves de -2 a 3?
\begin{enumerate}[label=\textbf{\Alph*)}]
\begin{multicols}{2}
  \item 5
  \item -5
  \item 0
  \item 7
  \item 1
  \item -1
\end{multicols}
\end{enumerate}
\vspace{0.15cm}
\item ¿Cuántas unidades te mueves de 4 a -1?
\begin{enumerate}[label=\textbf{\Alph*)}]
\begin{multicols}{2}
  \item 5
  \item -5
  \item 0
  \item 7
  \item 1
  \item -1
\end{multicols}
\end{enumerate}
\vspace{0.15cm}
\item ¿Cuántas unidades te mueves de -5 a -1?
\begin{enumerate}[label=\textbf{\Alph*)}]
\begin{multicols}{2}
  \item 4
  \item -4
  \item 0
  \item 6
  \item 1
  \item -1
\end{multicols}
\end{enumerate}
\vspace{0.15cm}
\item ¿Cuántas unidades te mueves de 0 a 6?
\begin{enumerate}[label=\textbf{\Alph*)}]
\begin{multicols}{2}
  \item 6
  \item -6
  \item 0
  \item 8
  \item 1
  \item -1
\end{multicols}
\end{enumerate}
\vspace{0.15cm}
\item ¿Cuántas unidades te mueves de -3 a 0?
\begin{enumerate}[label=\textbf{\Alph*)}]
\begin{multicols}{2}
  \item 3
  \item -3
  \item 0
  \item 5
  \item 1
  \item -1
\end{multicols}
\end{enumerate}
\vspace{0.15cm}
\item ¿Cuántas unidades te mueves de 2 a 2?
\begin{enumerate}[label=\textbf{\Alph*)}]
\begin{multicols}{2}
  \item 0
  \item 2
  \item 1
  \item 4
  \item 1
  \item -1
\end{multicols}
\end{enumerate}
\vspace{0.15cm}
\item ¿Cuántas unidades te mueves de -7 a 7?
\begin{enumerate}[label=\textbf{\Alph*)}]
\begin{multicols}{2}
  \item 14
  \item -14
  \item 0
  \item 16
  \item 1
  \item -1
\end{multicols}
\end{enumerate}
\vspace{0.15cm}
\item ¿Cuántas unidades te mueves de 5 a -5?
\begin{enumerate}[label=\textbf{\Alph*)}]
\begin{multicols}{2}
  \item 10
  \item -10
  \item 0
  \item 12
  \item 1
  \item -1
\end{multicols}
\end{enumerate}
\vspace{0.15cm}
\item ¿Cuántas unidades te mueves de -10 a -2?
\begin{enumerate}[label=\textbf{\Alph*)}]
\begin{multicols}{2}
  \item 8
  \item -8
  \item 0
  \item 10
  \item 1
  \item -1
\end{multicols}
\end{enumerate}
\vspace{0.15cm}
\item ¿Cuántas unidades te mueves de 3 a -4?
\begin{enumerate}[label=\textbf{\Alph*)}]
\begin{multicols}{2}
  \item 7
  \item -7
  \item 0
  \item 9
  \item 1
  \item -1
\end{multicols}
\end{enumerate}
\vspace{0.15cm}
\end{enumerate}\n
\subsubsection*{SECCIÓN VII: OPUESTOS Y SIMÉTRICOS (181–210)}

\begin{enumerate}[label=\textbf{\arabic*.}]
\setcounter{enumi}{180}
\item ¿Cuál es el opuesto de 8?
\begin{enumerate}[label=\textbf{\Alph*)}]
\begin{multicols}{2}
  \item -8
  \item 8
  \item 0
  \item 10
  \item -10
  \item 1
\end{multicols}
\end{enumerate}
\vspace{0.15cm}
\item ¿Cuál es el opuesto de -6?
\begin{enumerate}[label=\textbf{\Alph*)}]
\begin{multicols}{2}
  \item 6
  \item -6
  \item 0
  \item 10
  \item -10
  \item 1
\end{multicols}
\end{enumerate}
\vspace{0.15cm}
\item ¿Cuál es el opuesto de 0?
\begin{enumerate}[label=\textbf{\Alph*)}]
\begin{multicols}{2}
  \item 0
  \item 1
  \item 0
  \item 10
  \item -10
  \item 1
\end{multicols}
\end{enumerate}
\vspace{0.15cm}
\item ¿Cuál es el opuesto de -1?
\begin{enumerate}[label=\textbf{\Alph*)}]
\begin{multicols}{2}
  \item 1
  \item -1
  \item 0
  \item 10
  \item -10
  \item 1
\end{multicols}
\end{enumerate}
\vspace{0.15cm}
\item ¿Cuál es el opuesto de 3.5?
\begin{enumerate}[label=\textbf{\Alph*)}]
\begin{multicols}{2}
  \item -3.5
  \item 3.5
  \item 0
  \item 10
  \item -10
  \item 1
\end{multicols}
\end{enumerate}
\vspace{0.15cm}
\item ¿Cuál es el opuesto de -2.5?
\begin{enumerate}[label=\textbf{\Alph*)}]
\begin{multicols}{2}
  \item 2.5
  \item -2.5
  \item 0
  \item 10
  \item -10
  \item 1
\end{multicols}
\end{enumerate}
\vspace{0.15cm}
\item ¿Cuál es el simétrico de 4 respecto al origen?
\begin{enumerate}[label=\textbf{\Alph*)}]
\begin{multicols}{2}
  \item -4
  \item 4
  \item 0
  \item 10
  \item -10
  \item 1
\end{multicols}
\end{enumerate}
\vspace{0.15cm}
\item ¿Cuál es el simétrico de -9 respecto al origen?
\begin{enumerate}[label=\textbf{\Alph*)}]
\begin{multicols}{2}
  \item 9
  \item -9
  \item 0
  \item 10
  \item -10
  \item 1
\end{multicols}
\end{enumerate}
\vspace{0.15cm}
\item ¿Cuál es el simétrico de 0 respecto al origen?
\begin{enumerate}[label=\textbf{\Alph*)}]
\begin{multicols}{2}
  \item 0
  \item 1
  \item 0
  \item 10
  \item -10
  \item 1
\end{multicols}
\end{enumerate}
\vspace{0.15cm}
\item ¿Cuál es el simétrico de 7 respecto al origen?
\begin{enumerate}[label=\textbf{\Alph*)}]
\begin{multicols}{2}
  \item -7
  \item 7
  \item 0
  \item 10
  \item -10
  \item 1
\end{multicols}
\end{enumerate}
\vspace{0.15cm}
\item ¿Cuál es el simétrico de -3 respecto al origen?
\begin{enumerate}[label=\textbf{\Alph*)}]
\begin{multicols}{2}
  \item 3
  \item -3
  \item 0
  \item 10
  \item -10
  \item 1
\end{multicols}
\end{enumerate}
\vspace{0.15cm}
\item Si a=-5, ¿cuál es -a?
\begin{enumerate}[label=\textbf{\Alph*)}]
\begin{multicols}{2}
  \item 5
  \item -5
  \item 0
  \item 10
  \item -10
  \item 1
\end{multicols}
\end{enumerate}
\vspace{0.15cm}
\item Si a=2, ¿cuál es -a?
\begin{enumerate}[label=\textbf{\Alph*)}]
\begin{multicols}{2}
  \item -2
  \item 2
  \item 0
  \item 10
  \item -10
  \item 1
\end{multicols}
\end{enumerate}
\vspace{0.15cm}
\item Si a=-7, ¿cuál es -a?
\begin{enumerate}[label=\textbf{\Alph*)}]
\begin{multicols}{2}
  \item 7
  \item -7
  \item 0
  \item 10
  \item -10
  \item 1
\end{multicols}
\end{enumerate}
\vspace{0.15cm}
\item Si a=0, ¿cuál es -a?
\begin{enumerate}[label=\textbf{\Alph*)}]
\begin{multicols}{2}
  \item 0
  \item 1
  \item 0
  \item 10
  \item -10
  \item 1
\end{multicols}
\end{enumerate}
\vspace{0.15cm}
\item ¿Qué número está a la misma distancia del origen que -4?
\begin{enumerate}[label=\textbf{\Alph*)}]
\begin{multicols}{2}
  \item 4
  \item -4
  \item 0
  \item 10
  \item -10
  \item 1
\end{multicols}
\end{enumerate}
\vspace{0.15cm}
\item ¿Qué número está a la misma distancia del origen que 6?
\begin{enumerate}[label=\textbf{\Alph*)}]
\begin{multicols}{2}
  \item -6
  \item 6
  \item 0
  \item 10
  \item -10
  \item 1
\end{multicols}
\end{enumerate}
\vspace{0.15cm}
\item ¿Qué número está a la misma distancia del origen que -10?
\begin{enumerate}[label=\textbf{\Alph*)}]
\begin{multicols}{2}
  \item 10
  \item -10
  \item 0
  \item 10
  \item -10
  \item 1
\end{multicols}
\end{enumerate}
\vspace{0.15cm}
\item ¿Qué número está a la misma distancia del origen que 1.5?
\begin{enumerate}[label=\textbf{\Alph*)}]
\begin{multicols}{2}
  \item -1.5
  \item 1.5
  \item 0
  \item 10
  \item -10
  \item 1
\end{multicols}
\end{enumerate}
\vspace{0.15cm}
\item ¿Qué número está a la misma distancia del origen que -0.5?
\begin{enumerate}[label=\textbf{\Alph*)}]
\begin{multicols}{2}
  \item 0.5
  \item -0.5
  \item 0
  \item 10
  \item -10
  \item 1
\end{multicols}
\end{enumerate}
\vspace{0.15cm}
\item El opuesto de un número positivo es:
\begin{enumerate}[label=\textbf{\Alph*)}]
\begin{multicols}{2}
  \item Positivo
  \item Negativo
  \item Cero
  \item Igual
  \item Mayor
  \item Menor
\end{multicols}
\end{enumerate}
\vspace{0.15cm}
\item El opuesto de un número negativo es:
\begin{enumerate}[label=\textbf{\Alph*)}]
\begin{multicols}{2}
  \item Positivo
  \item Negativo
  \item Cero
  \item Igual
  \item Mayor
  \item Menor
\end{multicols}
\end{enumerate}
\vspace{0.15cm}
\item El opuesto de cero es:
\begin{enumerate}[label=\textbf{\Alph*)}]
\begin{multicols}{2}
  \item Cero
  \item Uno
  \item Menos uno
  \item Infinito
  \item Mayor
  \item Menor
\end{multicols}
\end{enumerate}
\vspace{0.15cm}
\item Si x=-2, ¿cuál es el opuesto de x?
\begin{enumerate}[label=\textbf{\Alph*)}]
\begin{multicols}{2}
  \item 2
  \item -2
  \item 0
  \item 10
  \item -10
  \item 1
\end{multicols}
\end{enumerate}
\vspace{0.15cm}
\item Si x=5, ¿cuál es el opuesto de x?
\begin{enumerate}[label=\textbf{\Alph*)}]
\begin{multicols}{2}
  \item -5
  \item 5
  \item 0
  \item 10
  \item -10
  \item 1
\end{multicols}
\end{enumerate}
\vspace{0.15cm}
\item Si x=-8, ¿cuál es el opuesto de x?
\begin{enumerate}[label=\textbf{\Alph*)}]
\begin{multicols}{2}
  \item 8
  \item -8
  \item 0
  \item 10
  \item -10
  \item 1
\end{multicols}
\end{enumerate}
\vspace{0.15cm}
\item Si x=0, ¿cuál es el opuesto de x?
\begin{enumerate}[label=\textbf{\Alph*)}]
\begin{multicols}{2}
  \item 0
  \item 1
  \item 0
  \item 10
  \item -10
  \item 1
\end{multicols}
\end{enumerate}
\vspace{0.15cm}
\item ¿Cuál es el simétrico de -6 respecto al origen?
\begin{enumerate}[label=\textbf{\Alph*)}]
\begin{multicols}{2}
  \item 6
  \item -6
  \item 0
  \item 10
  \item -10
  \item 1
\end{multicols}
\end{enumerate}
\vspace{0.15cm}
\item ¿Cuál es el simétrico de 10 respecto al origen?
\begin{enumerate}[label=\textbf{\Alph*)}]
\begin{multicols}{2}
  \item -10
  \item 10
  \item 0
  \item 10
  \item -10
  \item 1
\end{multicols}
\end{enumerate}
\vspace{0.15cm}
\item ¿Cuál es el simétrico de -1 respecto al origen?
\begin{enumerate}[label=\textbf{\Alph*)}]
\begin{multicols}{2}
  \item 1
  \item -1
  \item 0
  \item 10
  \item -10
  \item 1
\end{multicols}
\end{enumerate}
\vspace{0.15cm}
\end{enumerate}\n
\subsubsection*{SECCIÓN VIII: PUNTOS ENTRE DOS COORDENADAS (211–240)}

\begin{enumerate}[label=\textbf{\arabic*.}]
\setcounter{enumi}{210}
\item ¿Cuál de los siguientes números está entre \(1\) y \(4\)?
\begin{enumerate}[label=\textbf{\Alph*)}]
\begin{multicols}{2}
  \item 0
  \item 1
  \item 2
  \item -1
  \item -2
  \item 3
\end{multicols}
\end{enumerate}
\vspace{0.15cm}
\item ¿Cuál de los siguientes números está entre \(-3\) y \(0\)?
\begin{enumerate}[label=\textbf{\Alph*)}]
\begin{multicols}{2}
  \item 0
  \item -1
  \item 2
  \item -1
  \item -2
  \item 3
\end{multicols}
\end{enumerate}
\vspace{0.15cm}
\item ¿Cuál de los siguientes números está entre \(-2\) y \(2\)?
\begin{enumerate}[label=\textbf{\Alph*)}]
\begin{multicols}{2}
  \item 0
  \item 0
  \item 2
  \item -1
  \item -2
  \item 3
\end{multicols}
\end{enumerate}
\vspace{0.15cm}
\item ¿Cuál de los siguientes números está entre \(0\) y \(5\)?
\begin{enumerate}[label=\textbf{\Alph*)}]
\begin{multicols}{2}
  \item 0
  \item 3
  \item 2
  \item -1
  \item -2
  \item 3
\end{multicols}
\end{enumerate}
\vspace{0.15cm}
\item ¿Cuál de los siguientes números está entre \(-5\) y \(-2\)?
\begin{enumerate}[label=\textbf{\Alph*)}]
\begin{multicols}{2}
  \item 0
  \item -3
  \item 2
  \item -1
  \item -2
  \item 3
\end{multicols}
\end{enumerate}
\vspace{0.15cm}
\item ¿Cuál de los siguientes números está entre \(2\) y \(8\)?
\begin{enumerate}[label=\textbf{\Alph*)}]
\begin{multicols}{2}
  \item 0
  \item 5
  \item 2
  \item -1
  \item -2
  \item 3
\end{multicols}
\end{enumerate}
\vspace{0.15cm}
\item ¿Cuál de los siguientes números está entre \(-4\) y \(4\)?
\begin{enumerate}[label=\textbf{\Alph*)}]
\begin{multicols}{2}
  \item 0
  \item 0
  \item 2
  \item -1
  \item -2
  \item 3
\end{multicols}
\end{enumerate}
\vspace{0.15cm}
\item ¿Cuál de los siguientes números está entre \(-1\) y \(3\)?
\begin{enumerate}[label=\textbf{\Alph*)}]
\begin{multicols}{2}
  \item 0
  \item 1
  \item 2
  \item -1
  \item -2
  \item 3
\end{multicols}
\end{enumerate}
\vspace{0.15cm}
\item ¿Cuál de los siguientes números está entre \(3\) y \(7\)?
\begin{enumerate}[label=\textbf{\Alph*)}]
\begin{multicols}{2}
  \item 0
  \item 5
  \item 2
  \item -1
  \item -2
  \item 3
\end{multicols}
\end{enumerate}
\vspace{0.15cm}
\item ¿Cuál de los siguientes números está entre \(-7\) y \(-3\)?
\begin{enumerate}[label=\textbf{\Alph*)}]
\begin{multicols}{2}
  \item 0
  \item -5
  \item 2
  \item -1
  \item -2
  \item 3
\end{multicols}
\end{enumerate}
\vspace{0.15cm}
\item ¿Cuál de los siguientes números está entre \(0.5\) y \(1.5\)?
\begin{enumerate}[label=\textbf{\Alph*)}]
\begin{multicols}{2}
  \item 0
  \item 1
  \item 2
  \item -1
  \item -2
  \item 3
\end{multicols}
\end{enumerate}
\vspace{0.15cm}
\item ¿Cuál de los siguientes números está entre \(-1.5\) y \(-0.5\)?
\begin{enumerate}[label=\textbf{\Alph*)}]
\begin{multicols}{2}
  \item 0
  \item -1
  \item 2
  \item -1
  \item -2
  \item 3
\end{multicols}
\end{enumerate}
\vspace{0.15cm}
\item ¿Cuál de los siguientes números está entre \(2.5\) y \(3.5\)?
\begin{enumerate}[label=\textbf{\Alph*)}]
\begin{multicols}{2}
  \item 0
  \item 3
  \item 2
  \item -1
  \item -2
  \item 3
\end{multicols}
\end{enumerate}
\vspace{0.15cm}
\item ¿Cuál de los siguientes números está entre \(-3.5\) y \(-2.5\)?
\begin{enumerate}[label=\textbf{\Alph*)}]
\begin{multicols}{2}
  \item 0
  \item -3
  \item 2
  \item -1
  \item -2
  \item 3
\end{multicols}
\end{enumerate}
\vspace{0.15cm}
\item ¿Cuál de los siguientes números está entre \(10\) y \(20\)?
\begin{enumerate}[label=\textbf{\Alph*)}]
\begin{multicols}{2}
  \item 0
  \item 15
  \item 2
  \item -1
  \item -2
  \item 3
\end{multicols}
\end{enumerate}
\vspace{0.15cm}
\item ¿Cuál de los siguientes números está entre \(-20\) y \(-10\)?
\begin{enumerate}[label=\textbf{\Alph*)}]
\begin{multicols}{2}
  \item 0
  \item -15
  \item 2
  \item -1
  \item -2
  \item 3
\end{multicols}
\end{enumerate}
\vspace{0.15cm}
\item ¿Cuál de los siguientes números está entre \(100\) y \(200\)?
\begin{enumerate}[label=\textbf{\Alph*)}]
\begin{multicols}{2}
  \item 0
  \item 150
  \item 2
  \item -1
  \item -2
  \item 3
\end{multicols}
\end{enumerate}
\vspace{0.15cm}
\item ¿Cuál de los siguientes números está entre \(-100\) y \(-50\)?
\begin{enumerate}[label=\textbf{\Alph*)}]
\begin{multicols}{2}
  \item 0
  \item -75
  \item 2
  \item -1
  \item -2
  \item 3
\end{multicols}
\end{enumerate}
\vspace{0.15cm}
\item ¿Cuál de los siguientes números está entre \(\frac{1}{2}\) y \(\frac{3}{4}\)?
\begin{enumerate}[label=\textbf{\Alph*)}]
\begin{multicols}{2}
  \item 0
  \item \frac{5}{8}
  \item 2
  \item -1
  \item -2
  \item 3
\end{multicols}
\end{enumerate}
\vspace{0.15cm}
\item ¿Cuál de los siguientes números está entre \(\frac{1}{3}\) y \(\frac{2}{3}\)?
\begin{enumerate}[label=\textbf{\Alph*)}]
\begin{multicols}{2}
  \item 0
  \item \frac{1}{2}
  \item 2
  \item -1
  \item -2
  \item 3
\end{multicols}
\end{enumerate}
\vspace{0.15cm}
\item ¿Cuál de los siguientes números está entre \(-\frac{1}{2}\) y \(\frac{1}{2}\)?
\begin{enumerate}[label=\textbf{\Alph*)}]
\begin{multicols}{2}
  \item 0
  \item 0
  \item 2
  \item -1
  \item -2
  \item 3
\end{multicols}
\end{enumerate}
\vspace{0.15cm}
\item ¿Cuál de los siguientes números está entre \(-1\) y \(1\)?
\begin{enumerate}[label=\textbf{\Alph*)}]
\begin{multicols}{2}
  \item 0
  \item 0
  \item 2
  \item -1
  \item -2
  \item 3
\end{multicols}
\end{enumerate}
\vspace{0.15cm}
\item ¿Cuál de los siguientes números está entre \(-10\) y \(10\)?
\begin{enumerate}[label=\textbf{\Alph*)}]
\begin{multicols}{2}
  \item 0
  \item 0
  \item 2
  \item -1
  \item -2
  \item 3
\end{multicols}
\end{enumerate}
\vspace{0.15cm}
\item ¿Cuál de los siguientes números está entre \(5\) y \(15\)?
\begin{enumerate}[label=\textbf{\Alph*)}]
\begin{multicols}{2}
  \item 0
  \item 10
  \item 2
  \item -1
  \item -2
  \item 3
\end{multicols}
\end{enumerate}
\vspace{0.15cm}
\item ¿Cuál de los siguientes números está entre \(-15\) y \(-5\)?
\begin{enumerate}[label=\textbf{\Alph*)}]
\begin{multicols}{2}
  \item 0
  \item -10
  \item 2
  \item -1
  \item -2
  \item 3
\end{multicols}
\end{enumerate}
\vspace{0.15cm}
\item ¿Cuál de los siguientes números está entre \(0\) y \(0.5\)?
\begin{enumerate}[label=\textbf{\Alph*)}]
\begin{multicols}{2}
  \item 0.25
  \item 1
  \item 2
  \item -1
  \item -2
  \item 3
\end{multicols}
\end{enumerate}
\vspace{0.15cm}
\item ¿Cuál de los siguientes números está entre \(-0.5\) y \(0\)?
\begin{enumerate}[label=\textbf{\Alph*)}]
\begin{multicols}{2}
  \item -0.25
  \item 1
  \item 2
  \item -1
  \item -2
  \item 3
\end{multicols}
\end{enumerate}
\vspace{0.15cm}
\item ¿Cuál de los siguientes números está entre \(1.1\) y \(1.2\)?
\begin{enumerate}[label=\textbf{\Alph*)}]
\begin{multicols}{2}
  \item 1.15
  \item 1
  \item 2
  \item -1
  \item -2
  \item 3
\end{multicols}
\end{enumerate}
\vspace{0.15cm}
\item ¿Cuál de los siguientes números está entre \(-2.1\) y \(-2.0\)?
\begin{enumerate}[label=\textbf{\Alph*)}]
\begin{multicols}{2}
  \item -2.05
  \item 1
  \item 2
  \item -1
  \item -2
  \item 3
\end{multicols}
\end{enumerate}
\vspace{0.15cm}
\item ¿Cuál de los siguientes números está entre \(\sqrt{2}\) y \(\sqrt{3}\)?
\begin{enumerate}[label=\textbf{\Alph*)}]
\begin{multicols}{2}
  \item 1.5
  \item 1
  \item 2
  \item -1
  \item -2
  \item 3
\end{multicols}
\end{enumerate}
\vspace{0.15cm}
\end{enumerate}\n
\subsubsection*{SECCIÓN IX: DISTANCIA AL ORIGEN (241–270)}

\begin{enumerate}[label=\textbf{\arabic*.}]
\setcounter{enumi}{240}
\item ¿Cuál es la distancia de \(5\) al origen?
\begin{enumerate}[label=\textbf{\Alph*)}]
\begin{multicols}{2}
  \item 5
  \item -5
  \item 0
  \item 10
  \item -10
  \item 1
\end{multicols}
\end{enumerate}
\vspace{0.15cm}
\item ¿Cuál es la distancia de \(-3\) al origen?
\begin{enumerate}[label=\textbf{\Alph*)}]
\begin{multicols}{2}
  \item 3
  \item -3
  \item 0
  \item 10
  \item -10
  \item 1
\end{multicols}
\end{enumerate}
\vspace{0.15cm}
\item ¿Cuál es la distancia de \(0\) al origen?
\begin{enumerate}[label=\textbf{\Alph*)}]
\begin{multicols}{2}
  \item 0
  \item 1
  \item 0
  \item 10
  \item -10
  \item 1
\end{multicols}
\end{enumerate}
\vspace{0.15cm}
\item ¿Cuál es la distancia de \(-7\) al origen?
\begin{enumerate}[label=\textbf{\Alph*)}]
\begin{multicols}{2}
  \item 7
  \item -7
  \item 0
  \item 10
  \item -10
  \item 1
\end{multicols}
\end{enumerate}
\vspace{0.15cm}
\item ¿Cuál es la distancia de \(2.5\) al origen?
\begin{enumerate}[label=\textbf{\Alph*)}]
\begin{multicols}{2}
  \item 2.5
  \item -2.5
  \item 0
  \item 10
  \item -10
  \item 1
\end{multicols}
\end{enumerate}
\vspace{0.15cm}
\item ¿Cuál es la distancia de \(-4.5\) al origen?
\begin{enumerate}[label=\textbf{\Alph*)}]
\begin{multicols}{2}
  \item 4.5
  \item -4.5
  \item 0
  \item 10
  \item -10
  \item 1
\end{multicols}
\end{enumerate}
\vspace{0.15cm}
\item ¿Cuál es la distancia de \(10\) al origen?
\begin{enumerate}[label=\textbf{\Alph*)}]
\begin{multicols}{2}
  \item 10
  \item -10
  \item 0
  \item 10
  \item -10
  \item 1
\end{multicols}
\end{enumerate}
\vspace{0.15cm}
\item ¿Cuál es la distancia de \(-10\) al origen?
\begin{enumerate}[label=\textbf{\Alph*)}]
\begin{multicols}{2}
  \item 10
  \item -10
  \item 0
  \item 10
  \item -10
  \item 1
\end{multicols}
\end{enumerate}
\vspace{0.15cm}
\item ¿Cuál es la distancia de \(1\) al origen?
\begin{enumerate}[label=\textbf{\Alph*)}]
\begin{multicols}{2}
  \item 1
  \item -1
  \item 0
  \item 10
  \item -10
  \item 1
\end{multicols}
\end{enumerate}
\vspace{0.15cm}
\item ¿Cuál es la distancia de \(-1\) al origen?
\begin{enumerate}[label=\textbf{\Alph*)}]
\begin{multicols}{2}
  \item 1
  \item -1
  \item 0
  \item 10
  \item -10
  \item 1
\end{multicols}
\end{enumerate}
\vspace{0.15cm}
\item Si |x|=4, ¿cuáles son los posibles valores de x?
\begin{enumerate}[label=\textbf{\Alph*)}]
\begin{multicols}{2}
  \item 4 y -4
  \item 1
  \item 0
  \item 10
  \item -10
  \item 1
\end{multicols}
\end{enumerate}
\vspace{0.15cm}
\item Si |x|=7, ¿cuáles son los posibles valores de x?
\begin{enumerate}[label=\textbf{\Alph*)}]
\begin{multicols}{2}
  \item 7 y -7
  \item 1
  \item 0
  \item 10
  \item -10
  \item 1
\end{multicols}
\end{enumerate}
\vspace{0.15cm}
\item Si |x|=0, ¿cuál es el valor de x?
\begin{enumerate}[label=\textbf{\Alph*)}]
\begin{multicols}{2}
  \item 0
  \item 1
  \item 0
  \item 10
  \item -10
  \item 1
\end{multicols}
\end{enumerate}
\vspace{0.15cm}
\item Si |x|=3, ¿cuáles son los posibles valores de x?
\begin{enumerate}[label=\textbf{\Alph*)}]
\begin{multicols}{2}
  \item 3 y -3
  \item 1
  \item 0
  \item 10
  \item -10
  \item 1
\end{multicols}
\end{enumerate}
\vspace{0.15cm}
\item Si |x|=5, ¿cuáles son los posibles valores de x?
\begin{enumerate}[label=\textbf{\Alph*)}]
\begin{multicols}{2}
  \item 5 y -5
  \item 1
  \item 0
  \item 10
  \item -10
  \item 1
\end{multicols}
\end{enumerate}
\vspace{0.15cm}
\item Si |x|=2, ¿cuáles son los posibles valores de x?
\begin{enumerate}[label=\textbf{\Alph*)}]
\begin{multicols}{2}
  \item 2 y -2
  \item 1
  \item 0
  \item 10
  \item -10
  \item 1
\end{multicols}
\end{enumerate}
\vspace{0.15cm}
\item Si |x|=8, ¿cuáles son los posibles valores de x?
\begin{enumerate}[label=\textbf{\Alph*)}]
\begin{multicols}{2}
  \item 8 y -8
  \item 1
  \item 0
  \item 10
  \item -10
  \item 1
\end{multicols}
\end{enumerate}
\vspace{0.15cm}
\item Si |x|=1, ¿cuáles son los posibles valores de x?
\begin{enumerate}[label=\textbf{\Alph*)}]
\begin{multicols}{2}
  \item 1 y -1
  \item 1
  \item 0
  \item 10
  \item -10
  \item 1
\end{multicols}
\end{enumerate}
\vspace{0.15cm}
\item Si |x|=6, ¿cuáles son los posibles valores de x?
\begin{enumerate}[label=\textbf{\Alph*)}]
\begin{multicols}{2}
  \item 6 y -6
  \item 1
  \item 0
  \item 10
  \item -10
  \item 1
\end{multicols}
\end{enumerate}
\vspace{0.15cm}
\item Si |x|=9, ¿cuáles son los posibles valores de x?
\begin{enumerate}[label=\textbf{\Alph*)}]
\begin{multicols}{2}
  \item 9 y -9
  \item 1
  \item 0
  \item 10
  \item -10
  \item 1
\end{multicols}
\end{enumerate}
\vspace{0.15cm}
\item ¿Cuál es la distancia de \(-2\) al origen?
\begin{enumerate}[label=\textbf{\Alph*)}]
\begin{multicols}{2}
  \item 2
  \item -2
  \item 0
  \item 10
  \item -10
  \item 1
\end{multicols}
\end{enumerate}
\vspace{0.15cm}
\item ¿Cuál es la distancia de \(6\) al origen?
\begin{enumerate}[label=\textbf{\Alph*)}]
\begin{multicols}{2}
  \item 6
  \item -6
  \item 0
  \item 10
  \item -10
  \item 1
\end{multicols}
\end{enumerate}
\vspace{0.15cm}
\item ¿Cuál es la distancia de \(-8\) al origen?
\begin{enumerate}[label=\textbf{\Alph*)}]
\begin{multicols}{2}
  \item 8
  \item -8
  \item 0
  \item 10
  \item -10
  \item 1
\end{multicols}
\end{enumerate}
\vspace{0.15cm}
\item ¿Cuál es la distancia de \(3.5\) al origen?
\begin{enumerate}[label=\textbf{\Alph*)}]
\begin{multicols}{2}
  \item 3.5
  \item -3.5
  \item 0
  \item 10
  \item -10
  \item 1
\end{multicols}
\end{enumerate}
\vspace{0.15cm}
\item ¿Cuál es la distancia de \(-0.5\) al origen?
\begin{enumerate}[label=\textbf{\Alph*)}]
\begin{multicols}{2}
  \item 0.5
  \item -0.5
  \item 0
  \item 10
  \item -10
  \item 1
\end{multicols}
\end{enumerate}
\vspace{0.15cm}
\item ¿Cuál es la distancia de \(100\) al origen?
\begin{enumerate}[label=\textbf{\Alph*)}]
\begin{multicols}{2}
  \item 100
  \item -100
  \item 0
  \item 10
  \item -10
  \item 1
\end{multicols}
\end{enumerate}
\vspace{0.15cm}
\item ¿Cuál es la distancia de \(-100\) al origen?
\begin{enumerate}[label=\textbf{\Alph*)}]
\begin{multicols}{2}
  \item 100
  \item -100
  \item 0
  \item 10
  \item -10
  \item 1
\end{multicols}
\end{enumerate}
\vspace{0.15cm}
\item ¿Cuál es la distancia de \(\frac{1}{2}\) al origen?
\begin{enumerate}[label=\textbf{\Alph*)}]
\begin{multicols}{2}
  \item \frac{1}{2}
  \item -\frac{1}{2}
  \item 0
  \item 10
  \item -10
  \item 1
\end{multicols}
\end{enumerate}
\vspace{0.15cm}
\item ¿Cuál es la distancia de \(-\frac{3}{4}\) al origen?
\begin{enumerate}[label=\textbf{\Alph*)}]
\begin{multicols}{2}
  \item \frac{3}{4}
  \item -\frac{3}{4}
  \item 0
  \item 10
  \item -10
  \item 1
\end{multicols}
\end{enumerate}
\vspace{0.15cm}
\item ¿Cuál es la distancia de \(\sqrt{2}\) al origen?
\begin{enumerate}[label=\textbf{\Alph*)}]
\begin{multicols}{2}
  \item \sqrt{2}
  \item -\sqrt{2}
  \item 0
  \item 10
  \item -10
  \item 1
\end{multicols}
\end{enumerate}
\vspace{0.15cm}
\end{enumerate}\n
\subsubsection*{SECCIÓN X: DESIGUALDADES Y REPRESENTACIÓN (271–300)}

\begin{enumerate}[label=\textbf{\arabic*.}]
\setcounter{enumi}{270}
\item La representación de \(x > 2\) es:
\begin{enumerate}[label=\textbf{\Alph*)}]
\begin{multicols}{2}
  \item Punto lleno en 2 hacia derecha
  \item Punto vacío en 2 hacia derecha
  \item Punto lleno en 2 hacia izquierda
  \item Punto vacío en 2 hacia izquierda
  \item Solo el punto
  \item Toda la recta
\end{multicols}
\end{enumerate}
\vspace{0.15cm}
\item La representación de \(x >= 2\) es:
\begin{enumerate}[label=\textbf{\Alph*)}]
\begin{multicols}{2}
  \item Punto lleno en 2 hacia derecha
  \item Punto vacío en 2 hacia derecha
  \item Punto lleno en 2 hacia izquierda
  \item Punto vacío en 2 hacia izquierda
  \item Solo el punto
  \item Toda la recta
\end{multicols}
\end{enumerate}
\vspace{0.15cm}
\item La representación de \(x < 3\) es:
\begin{enumerate}[label=\textbf{\Alph*)}]
\begin{multicols}{2}
  \item Punto lleno en 2 hacia derecha
  \item Punto vacío en 2 hacia derecha
  \item Punto lleno en 2 hacia izquierda
  \item Punto vacío en 3 hacia izquierda
  \item Solo el punto
  \item Toda la recta
\end{multicols}
\end{enumerate}
\vspace{0.15cm}
\item La representación de \(x <= 3\) es:
\begin{enumerate}[label=\textbf{\Alph*)}]
\begin{multicols}{2}
  \item Punto lleno en 2 hacia derecha
  \item Punto vacío en 2 hacia derecha
  \item Punto lleno en 3 hacia izquierda
  \item Punto vacío en 2 hacia izquierda
  \item Solo el punto
  \item Toda la recta
\end{multicols}
\end{enumerate}
\vspace{0.15cm}
\item La representación de \(-1 < x < 2\) es:
\begin{enumerate}[label=\textbf{\Alph*)}]
\begin{multicols}{2}
  \item Puntos vacíos en -1 y 2, segmento entre ellos
  \item Punto vacío en 2 hacia derecha
  \item Punto lleno en 2 hacia izquierda
  \item Punto vacío en 2 hacia izquierda
  \item Solo el punto
  \item Toda la recta
\end{multicols}
\end{enumerate}
\vspace{0.15cm}
\item La representación de \(-1 <= x <= 2\) es:
\begin{enumerate}[label=\textbf{\Alph*)}]
\begin{multicols}{2}
  \item Punto lleno en 2 hacia derecha
  \item Puntos llenos en -1 y 2, segmento entre ellos
  \item Punto lleno en 2 hacia izquierda
  \item Punto vacío en 2 hacia izquierda
  \item Solo el punto
  \item Toda la recta
\end{multicols}
\end{enumerate}
\vspace{0.15cm}
\item La representación de \(-1 <= x < 2\) es:
\begin{enumerate}[label=\textbf{\Alph*)}]
\begin{multicols}{2}
  \item Punto lleno en -1, vacío en 2
  \item Punto vacío en 2 hacia derecha
  \item Punto lleno en 2 hacia izquierda
  \item Punto vacío en 2 hacia izquierda
  \item Solo el punto
  \item Toda la recta
\end{multicols}
\end{enumerate}
\vspace{0.15cm}
\item La representación de \(-1 < x <= 2\) es:
\begin{enumerate}[label=\textbf{\Alph*)}]
\begin{multicols}{2}
  \item Punto lleno en 2 hacia derecha
  \item Punto vacío en -1, lleno en 2
  \item Punto lleno en 2 hacia izquierda
  \item Punto vacío en 2 hacia izquierda
  \item Solo el punto
  \item Toda la recta
\end{multicols}
\end{enumerate}
\vspace{0.15cm}
\item ¿Qué desigualdad representa el conjunto a la derecha de -2?
\begin{enumerate}[label=\textbf{\Alph*)}]
\begin{multicols}{2}
  \item -2 < x < 3
  \item x > -2
  \item -2 < x \leq 3
  \item -2 \leq x < 3
  \item x = 0
  \item Todos los reales
\end{multicols}
\end{enumerate}
\vspace{0.15cm}
\item ¿Qué desigualdad representa el conjunto a la izquierda de 4?
\begin{enumerate}[label=\textbf{\Alph*)}]
\begin{multicols}{2}
  \item x < 4
  \item -2 \leq x \leq 3
  \item -2 < x \leq 3
  \item -2 \leq x < 3
  \item x = 0
  \item Todos los reales
\end{multicols}
\end{enumerate}
\vspace{0.15cm}
\item ¿Qué desigualdad representa el entre 1 y 5 sin incluirlos?
\begin{enumerate}[label=\textbf{\Alph*)}]
\begin{multicols}{2}
  \item 1 < x < 5
  \item -2 \leq x \leq 3
  \item -2 < x \leq 3
  \item -2 \leq x < 3
  \item x = 0
  \item Todos los reales
\end{multicols}
\end{enumerate}
\vspace{0.15cm}
\item ¿Qué desigualdad representa el entre -3 y 2 incluyendo ambos?
\begin{enumerate}[label=\textbf{\Alph*)}]
\begin{multicols}{2}
  \item -2 < x < 3
  \item -3 <= x <= 2
  \item -2 < x \leq 3
  \item -2 \leq x < 3
  \item x = 0
  \item Todos los reales
\end{multicols}
\end{enumerate}
\vspace{0.15cm}
\item ¿Qué desigualdad representa el entre 0 y 10 incluyendo 0 pero no 10?
\begin{enumerate}[label=\textbf{\Alph*)}]
\begin{multicols}{2}
  \item -2 < x < 3
  \item -2 \leq x \leq 3
  \item 0 <= x < 10
  \item -2 \leq x < 3
  \item x = 0
  \item Todos los reales
\end{multicols}
\end{enumerate}
\vspace{0.15cm}
\item ¿Qué desigualdad representa el entre -5 y 5 incluyendo 5 pero no -5?
\begin{enumerate}[label=\textbf{\Alph*)}]
\begin{multicols}{2}
  \item -2 < x < 3
  \item -2 \leq x \leq 3
  \item -5 < x <= 5
  \item -2 \leq x < 3
  \item x = 0
  \item Todos los reales
\end{multicols}
\end{enumerate}
\vspace{0.15cm}
\item ¿Qué desigualdad representa el mayores o iguales que 3?
\begin{enumerate}[label=\textbf{\Alph*)}]
\begin{multicols}{2}
  \item -2 < x < 3
  \item -2 \leq x \leq 3
  \item -2 < x \leq 3
  \item x >= 3
  \item x = 0
  \item Todos los reales
\end{multicols}
\end{enumerate}
\vspace{0.15cm}
\item ¿Qué desigualdad representa el menores o iguales que -1?
\begin{enumerate}[label=\textbf{\Alph*)}]
\begin{multicols}{2}
  \item -2 < x < 3
  \item -2 \leq x \leq 3
  \item x <= -1
  \item -2 \leq x < 3
  \item x = 0
  \item Todos los reales
\end{multicols}
\end{enumerate}
\vspace{0.15cm}
\item ¿Qué desigualdad representa el mayores que 0?
\begin{enumerate}[label=\textbf{\Alph*)}]
\begin{multicols}{2}
  \item -2 < x < 3
  \item x > 0
  \item -2 < x \leq 3
  \item -2 \leq x < 3
  \item x = 0
  \item Todos los reales
\end{multicols}
\end{enumerate}
\vspace{0.15cm}
\item ¿Qué desigualdad representa el menores que 0?
\begin{enumerate}[label=\textbf{\Alph*)}]
\begin{multicols}{2}
  \item x < 0
  \item -2 \leq x \leq 3
  \item -2 < x \leq 3
  \item -2 \leq x < 3
  \item x = 0
  \item Todos los reales
\end{multicols}
\end{enumerate}
\vspace{0.15cm}
\item Si x está en el intervalo (-2,3), entonces:
\begin{enumerate}[label=\textbf{\Alph*)}]
\begin{multicols}{2}
  \item -2 < x < 3
  \item -2 \leq x \leq 3
  \item -2 < x \leq 3
  \item -2 \leq x < 3
  \item x = 0
  \item Todos los reales
\end{multicols}
\end{enumerate}
\vspace{0.15cm}
\item Si x está en el intervalo [-2,3], entonces:
\begin{enumerate}[label=\textbf{\Alph*)}]
\begin{multicols}{2}
  \item -2 < x < 3
  \item -2 <= x <= 3
  \item -2 < x \leq 3
  \item -2 \leq x < 3
  \item x = 0
  \item Todos los reales
\end{multicols}
\end{enumerate}
\vspace{0.15cm}
\item Si x está en el intervalo [-2,3), entonces:
\begin{enumerate}[label=\textbf{\Alph*)}]
\begin{multicols}{2}
  \item -2 < x < 3
  \item -2 \leq x \leq 3
  \item -2 < x \leq 3
  \item -2 <= x < 3
  \item x = 0
  \item Todos los reales
\end{multicols}
\end{enumerate}
\vspace{0.15cm}
\item Si x está en el intervalo (-2,3], entonces:
\begin{enumerate}[label=\textbf{\Alph*)}]
\begin{multicols}{2}
  \item -2 < x < 3
  \item -2 \leq x \leq 3
  \item -2 < x <= 3
  \item -2 \leq x < 3
  \item x = 0
  \item Todos los reales
\end{multicols}
\end{enumerate}
\vspace{0.15cm}
\item La representación de \(x > 5\) es:
\begin{enumerate}[label=\textbf{\Alph*)}]
\begin{multicols}{2}
  \item Punto vacío en 5 hacia derecha
  \item -2 \leq x \leq 3
  \item -2 < x \leq 3
  \item -2 \leq x < 3
  \item x = 0
  \item Todos los reales
\end{multicols}
\end{enumerate}
\vspace{0.15cm}
\item La representación de \(x >= 5\) es:
\begin{enumerate}[label=\textbf{\Alph*)}]
\begin{multicols}{2}
  \item -2 < x < 3
  \item Punto lleno en 5 hacia derecha
  \item -2 < x \leq 3
  \item -2 \leq x < 3
  \item x = 0
  \item Todos los reales
\end{multicols}
\end{enumerate}
\vspace{0.15cm}
\item La representación de \(x < 5\) es:
\begin{enumerate}[label=\textbf{\Alph*)}]
\begin{multicols}{2}
  \item -2 < x < 3
  \item -2 \leq x \leq 3
  \item Punto vacío en 5 hacia izquierda
  \item -2 \leq x < 3
  \item x = 0
  \item Todos los reales
\end{multicols}
\end{enumerate}
\vspace{0.15cm}
\item La representación de \(x <= 5\) es:
\begin{enumerate}[label=\textbf{\Alph*)}]
\begin{multicols}{2}
  \item -2 < x < 3
  \item -2 \leq x \leq 3
  \item -2 < x \leq 3
  \item Punto lleno en 5 hacia izquierda
  \item x = 0
  \item Todos los reales
\end{multicols}
\end{enumerate}
\vspace{0.15cm}
\item El intervalo \((0, infty)\) corresponde a:
\begin{enumerate}[label=\textbf{\Alph*)}]
\begin{multicols}{2}
  \item x > 0
  \item x < 0
  \item x \geq 0
  \item x \leq 0
  \item x = 0
  \item Todos los reales
\end{multicols}
\end{enumerate}
\vspace{0.15cm}
\item El intervalo \((-infty, 0)\) corresponde a:
\begin{enumerate}[label=\textbf{\Alph*)}]
\begin{multicols}{2}
  \item x > 0
  \item x < 0
  \item x \geq 0
  \item x \leq 0
  \item x = 0
  \item Todos los reales
\end{multicols}
\end{enumerate}
\vspace{0.15cm}
\item El intervalo \((-infty, infty)\) corresponde a:
\begin{enumerate}[label=\textbf{\Alph*)}]
\begin{multicols}{2}
  \item x > 0
  \item x < 0
  \item x \geq 0
  \item Todos los números reales
  \item x = 0
  \item Todos los reales
\end{multicols}
\end{enumerate}
\vspace{0.15cm}
\item El intervalo \([0, infty)\) corresponde a:
\begin{enumerate}[label=\textbf{\Alph*)}]
\begin{multicols}{2}
  \item x > 0
  \item x < 0
  \item x >= 0
  \item x \leq 0
  \item x = 0
  \item Todos los reales
\end{multicols}
\end{enumerate}
\vspace{0.15cm}
\end{enumerate}\n
\vspace{0.4cm}
\begin{tcolorbox}[colback=gray!5!white,colframe=gray!60!black,title=\textbf{SOLUCIONARIO -- 1.2 La Recta Real}]
\begin{multicols}{4}
\begin{enumerate}[label=\textbf{\arabic*.}]
  \item \textbf{A}
  \item \textbf{C}
  \item \textbf{B}
  \item \textbf{B}
  \item \textbf{A}
  \item \textbf{B}
  \item \textbf{B}
  \item \textbf{A}
  \item \textbf{B}
  \item \textbf{B}
  \item \textbf{A}
  \item \textbf{B}
  \item \textbf{D}
  \item \textbf{B}
  \item \textbf{B}
  \item \textbf{B}
  \item \textbf{A}
  \item \textbf{D}
  \item \textbf{A}
  \item \textbf{B}
  \item \textbf{D}
  \item \textbf{B}
  \item \textbf{B}
  \item \textbf{B}
  \item \textbf{B}
  \item \textbf{B}
  \item \textbf{B}
  \item \textbf{B}
  \item \textbf{A}
  \item \textbf{B}
  \item \textbf{C}
  \item \textbf{C}
  \item \textbf{C}
  \item \textbf{C}
  \item \textbf{C}
  \item \textbf{C}
  \item \textbf{A}
  \item \textbf{B}
  \item \textbf{C}
  \item \textbf{A}
  \item \textbf{C}
  \item \textbf{D}
  \item \textbf{B}
  \item \textbf{B}
  \item \textbf{D}
  \item \textbf{A}
  \item \textbf{C}
  \item \textbf{C}
  \item \textbf{C}
  \item \textbf{B}
  \item \textbf{B}
  \item \textbf{B}
  \item \textbf{B}
  \item \textbf{B}
  \item \textbf{B}
  \item \textbf{B}
  \item \textbf{A}
  \item \textbf{A}
  \item \textbf{B}
  \item \textbf{B}
  \item \textbf{A}
  \item \textbf{A}
  \item \textbf{B}
  \item \textbf{C}
  \item \textbf{C}
  \item \textbf{B}
  \item \textbf{A}
  \item \textbf{B}
  \item \textbf{B}
  \item \textbf{A}
  \item \textbf{C}
  \item \textbf{C}
  \item \textbf{C}
  \item \textbf{C}
  \item \textbf{B}
  \item \textbf{A}
  \item \textbf{C}
  \item \textbf{B}
  \item \textbf{F}
  \item \textbf{B}
  \item \textbf{B}
  \item \textbf{B}
  \item \textbf{B}
  \item \textbf{B}
  \item \textbf{B}
  \item \textbf{B}
  \item \textbf{B}
  \item \textbf{B}
  \item \textbf{B}
  \item \textbf{B}
  \item \textbf{A}
  \item \textbf{A}
  \item \textbf{A}
  \item \textbf{A}
  \item \textbf{A}
  \item \textbf{A}
  \item \textbf{A}
  \item \textbf{A}
  \item \textbf{A}
  \item \textbf{A}
  \item \textbf{A}
  \item \textbf{A}
  \item \textbf{A}
  \item \textbf{A}
  \item \textbf{A}
  \item \textbf{A}
  \item \textbf{A}
  \item \textbf{A}
  \item \textbf{A}
  \item \textbf{A}
  \item \textbf{A}
  \item \textbf{A}
  \item \textbf{A}
  \item \textbf{A}
  \item \textbf{A}
  \item \textbf{A}
  \item \textbf{A}
  \item \textbf{A}
  \item \textbf{A}
  \item \textbf{A}
  \item \textbf{B}
  \item \textbf{B}
  \item \textbf{B}
  \item \textbf{B}
  \item \textbf{B}
  \item \textbf{A}
  \item \textbf{A}
  \item \textbf{A}
  \item \textbf{A}
  \item \textbf{A}
  \item \textbf{A}
  \item \textbf{A}
  \item \textbf{A}
  \item \textbf{A}
  \item \textbf{A}
  \item \textbf{A}
  \item \textbf{A}
  \item \textbf{A}
  \item \textbf{A}
  \item \textbf{A}
  \item \textbf{A}
  \item \textbf{A}
  \item \textbf{A}
  \item \textbf{A}
  \item \textbf{A}
  \item \textbf{A}
  \item \textbf{A}
  \item \textbf{A}
  \item \textbf{A}
  \item \textbf{A}
  \item \textbf{A}
  \item \textbf{A}
  \item \textbf{A}
  \item \textbf{A}
  \item \textbf{A}
  \item \textbf{A}
  \item \textbf{A}
  \item \textbf{A}
  \item \textbf{A}
  \item \textbf{A}
  \item \textbf{A}
  \item \textbf{A}
  \item \textbf{A}
  \item \textbf{A}
  \item \textbf{A}
  \item \textbf{A}
  \item \textbf{A}
  \item \textbf{A}
  \item \textbf{A}
  \item \textbf{A}
  \item \textbf{A}
  \item \textbf{A}
  \item \textbf{A}
  \item \textbf{A}
  \item \textbf{A}
  \item \textbf{A}
  \item \textbf{A}
  \item \textbf{A}
  \item \textbf{A}
  \item \textbf{A}
  \item \textbf{A}
  \item \textbf{A}
  \item \textbf{A}
  \item \textbf{A}
  \item \textbf{A}
  \item \textbf{A}
  \item \textbf{A}
  \item \textbf{A}
  \item \textbf{A}
  \item \textbf{A}
  \item \textbf{A}
  \item \textbf{A}
  \item \textbf{A}
  \item \textbf{A}
  \item \textbf{A}
  \item \textbf{A}
  \item \textbf{A}
  \item \textbf{A}
  \item \textbf{A}
  \item \textbf{A}
  \item \textbf{B}
  \item \textbf{A}
  \item \textbf{A}
  \item \textbf{A}
  \item \textbf{A}
  \item \textbf{A}
  \item \textbf{A}
  \item \textbf{A}
  \item \textbf{A}
  \item \textbf{A}
  \item \textbf{C}
  \item \textbf{B}
  \item \textbf{B}
  \item \textbf{B}
  \item \textbf{B}
  \item \textbf{B}
  \item \textbf{B}
  \item \textbf{B}
  \item \textbf{B}
  \item \textbf{B}
  \item \textbf{B}
  \item \textbf{B}
  \item \textbf{B}
  \item \textbf{B}
  \item \textbf{B}
  \item \textbf{B}
  \item \textbf{B}
  \item \textbf{B}
  \item \textbf{B}
  \item \textbf{B}
  \item \textbf{B}
  \item \textbf{B}
  \item \textbf{B}
  \item \textbf{B}
  \item \textbf{B}
  \item \textbf{A}
  \item \textbf{A}
  \item \textbf{A}
  \item \textbf{A}
  \item \textbf{A}
  \item \textbf{A}
  \item \textbf{A}
  \item \textbf{A}
  \item \textbf{A}
  \item \textbf{A}
  \item \textbf{A}
  \item \textbf{A}
  \item \textbf{A}
  \item \textbf{A}
  \item \textbf{A}
  \item \textbf{A}
  \item \textbf{A}
  \item \textbf{A}
  \item \textbf{A}
  \item \textbf{A}
  \item \textbf{A}
  \item \textbf{A}
  \item \textbf{A}
  \item \textbf{A}
  \item \textbf{A}
  \item \textbf{A}
  \item \textbf{A}
  \item \textbf{A}
  \item \textbf{A}
  \item \textbf{A}
  \item \textbf{A}
  \item \textbf{A}
  \item \textbf{A}
  \item \textbf{A}
  \item \textbf{A}
  \item \textbf{B}
  \item \textbf{A}
  \item \textbf{D}
  \item \textbf{C}
  \item \textbf{A}
  \item \textbf{B}
  \item \textbf{A}
  \item \textbf{B}
  \item \textbf{B}
  \item \textbf{A}
  \item \textbf{A}
  \item \textbf{B}
  \item \textbf{C}
  \item \textbf{C}
  \item \textbf{D}
  \item \textbf{C}
  \item \textbf{B}
  \item \textbf{A}
  \item \textbf{A}
  \item \textbf{B}
  \item \textbf{D}
  \item \textbf{C}
  \item \textbf{A}
  \item \textbf{B}
  \item \textbf{C}
  \item \textbf{D}
  \item \textbf{A}
  \item \textbf{B}
  \item \textbf{D}
  \item \textbf{C}
\end{enumerate}
\end{multicols}
\end{tcolorbox}
\newpage\n
\subsection{1.3 Conjuntos e Intervalos}

\subsubsection*{SECCIÓN I: CONCEPTOS BÁSICOS DE CONJUNTOS (1–50)}

\begin{enumerate}[label=\textbf{\arabic*.}]
\setcounter{enumi}{0}
\item ¿Qué es un conjunto?
\begin{enumerate}[label=\textbf{\Alph*)}]
\begin{multicols}{2}
  \item Una colección de objetos bien definida
  \item U
  \item 2
  \item 4
  \item 0
  \item Ninguna
\end{multicols}
\end{enumerate}
\vspace{0.15cm}
\item ¿Cómo se denota que \(a\) es elemento de \(A\)?
\begin{enumerate}[label=\textbf{\Alph*)}]
\begin{multicols}{2}
  \item 0
  \item a \in A
  \item 2
  \item \mathbb{R}
  \item \mathbb{C}
  \item \mathbb{I}
\end{multicols}
\end{enumerate}
\vspace{0.15cm}
\item ¿Cómo se denota que \(a\) NO es elemento de \(A\)?
\begin{enumerate}[label=\textbf{\Alph*)}]
\begin{multicols}{2}
  \item 0
  \item U
  \item 2
  \item \mathbb{R}
  \item a \notin A
  \item \mathbb{I}
\end{multicols}
\end{enumerate}
\vspace{0.15cm}
\item El conjunto vacío se denota como:
\begin{enumerate}[label=\textbf{\Alph*)}]
\begin{multicols}{2}
  \item 0
  \item \emptyset
  \item 2
  \item \mathbb{R}
  \item \mathbb{C}
  \item \mathbb{I}
\end{multicols}
\end{enumerate}
\vspace{0.15cm}
\item ¿Cuál es el conjunto vacío?
\begin{enumerate}[label=\textbf{\Alph*)}]
\begin{multicols}{2}
  \item 0
  \item U
  \item {}
  \item 4
  \item 0
  \item Ninguna
\end{multicols}
\end{enumerate}
\vspace{0.15cm}
\item El conjunto universal se denota usualmente como:
\begin{enumerate}[label=\textbf{\Alph*)}]
\begin{multicols}{2}
  \item 0
  \item U
  \item 2
  \item \mathbb{R}
  \item \mathbb{C}
  \item \mathbb{I}
\end{multicols}
\end{enumerate}
\vspace{0.15cm}
\item ¿Cuántos elementos tiene el conjunto \(\{1,2,3\}\)?
\begin{enumerate}[label=\textbf{\Alph*)}]
\begin{multicols}{2}
  \item 0
  \item U
  \item 3
  \item 4
  \item 0
  \item Ninguna
\end{multicols}
\end{enumerate}
\vspace{0.15cm}
\item ¿Cuál es la cardinalidad de \(\{a,b,c,d\}\)?
\begin{enumerate}[label=\textbf{\Alph*)}]
\begin{multicols}{2}
  \item 0
  \item U
  \item 4
  \item 4
  \item 0
  \item Ninguna
\end{multicols}
\end{enumerate}
\vspace{0.15cm}
\item El conjunto \(\{x \mid x \text{ es par}\}\) se lee como:
\begin{enumerate}[label=\textbf{\Alph*)}]
\begin{multicols}{2}
  \item Conjunto de x tal que x es par
  \item U
  \item 2
  \item 4
  \item 0
  \item Ninguna
\end{multicols}
\end{enumerate}
\vspace{0.15cm}
\item ¿Cuál de los siguientes es un conjunto finito?
\begin{enumerate}[label=\textbf{\Alph*)}]
\begin{multicols}{2}
  \item 0
  \item U
  \item {1,2,3,4,5}
  \item 4
  \item 0
  \item Ninguna
\end{multicols}
\end{enumerate}
\vspace{0.15cm}
\item ¿Cuál de los siguientes es un conjunto infinito?
\begin{enumerate}[label=\textbf{\Alph*)}]
\begin{multicols}{2}
  \item 0
  \item U
  \item \mathbb{N}
  \item 4
  \item 0
  \item Ninguna
\end{multicols}
\end{enumerate}
\vspace{0.15cm}
\item El conjunto de los números naturales se denota:
\begin{enumerate}[label=\textbf{\Alph*)}]
\begin{multicols}{2}
  \item 0
  \item U
  \item \mathbb{N}
  \item \mathbb{R}
  \item \mathbb{C}
  \item \mathbb{I}
\end{multicols}
\end{enumerate}
\vspace{0.15cm}
\item El conjunto de los números enteros se denota:
\begin{enumerate}[label=\textbf{\Alph*)}]
\begin{multicols}{2}
  \item 0
  \item \mathbb{Z}
  \item 2
  \item \mathbb{R}
  \item \mathbb{C}
  \item \mathbb{I}
\end{multicols}
\end{enumerate}
\vspace{0.15cm}
\item El conjunto de los números racionales se denota:
\begin{enumerate}[label=\textbf{\Alph*)}]
\begin{multicols}{2}
  \item 0
  \item U
  \item \mathbb{Q}
  \item \mathbb{R}
  \item \mathbb{C}
  \item \mathbb{I}
\end{multicols}
\end{enumerate}
\vspace{0.15cm}
\item El conjunto de los números reales se denota:
\begin{enumerate}[label=\textbf{\Alph*)}]
\begin{multicols}{2}
  \item 0
  \item U
  \item 2
  \item \mathbb{R}
  \item \mathbb{C}
  \item \mathbb{I}
\end{multicols}
\end{enumerate}
\vspace{0.15cm}
\item El conjunto de los números irracionales se denota:
\begin{enumerate}[label=\textbf{\Alph*)}]
\begin{multicols}{2}
  \item 0
  \item U
  \item 2
  \item \mathbb{R}
  \item \mathbb{C}
  \item \mathbb{I}
\end{multicols}
\end{enumerate}
\vspace{0.15cm}
\item ¿Cuál de los siguientes conjuntos es vacío?
\begin{enumerate}[label=\textbf{\Alph*)}]
\begin{multicols}{2}
  \item \{x \in \mathbb{N} \mid x < 1\}
  \item U
  \item 2
  \item 4
  \item 0
  \item Ninguna
\end{multicols}
\end{enumerate}
\vspace{0.15cm}
\item ¿Cuál es la cardinalidad del conjunto vacío?
\begin{enumerate}[label=\textbf{\Alph*)}]
\begin{multicols}{2}
  \item 0
  \item U
  \item 2
  \item 4
  \item 0
  \item Ninguna
\end{multicols}
\end{enumerate}
\vspace{0.15cm}
\item ¿Cuál de los siguientes conjuntos es unitario?
\begin{enumerate}[label=\textbf{\Alph*)}]
\begin{multicols}{2}
  \item 0
  \item {1}
  \item 2
  \item 4
  \item 0
  \item Ninguna
\end{multicols}
\end{enumerate}
\vspace{0.15cm}
\item ¿Cuál de los siguientes conjuntos es unitario?
\begin{enumerate}[label=\textbf{\Alph*)}]
\begin{multicols}{2}
  \item {0}
  \item U
  \item 2
  \item 4
  \item 0
  \item Ninguna
\end{multicols}
\end{enumerate}
\vspace{0.15cm}
\item El conjunto \(\{1,2,3,\ldots\}\) representa:
\begin{enumerate}[label=\textbf{\Alph*)}]
\begin{multicols}{2}
  \item 0
  \item \mathbb{N}
  \item 2
  \item 4
  \item 0
  \item Ninguna
\end{multicols}
\end{enumerate}
\vspace{0.15cm}
\item El conjunto \(\{\ldots,-2,-1,0,1,2,\ldots\}\) representa:
\begin{enumerate}[label=\textbf{\Alph*)}]
\begin{multicols}{2}
  \item 0
  \item \mathbb{Z}
  \item 2
  \item 4
  \item 0
  \item Ninguna
\end{multicols}
\end{enumerate}
\vspace{0.15cm}
\item ¿Cuál de los siguientes conjuntos es igual a \(\{1,2,3\}\)?
\begin{enumerate}[label=\textbf{\Alph*)}]
\begin{multicols}{2}
  \item {3,2,1}
  \item U
  \item 2
  \item 4
  \item 0
  \item Ninguna
\end{multicols}
\end{enumerate}
\vspace{0.15cm}
\item Dos conjuntos son iguales si:
\begin{enumerate}[label=\textbf{\Alph*)}]
\begin{multicols}{2}
  \item 0
  \item Tienen exactamente los mismos elementos
  \item 2
  \item 4
  \item 0
  \item Ninguna
\end{multicols}
\end{enumerate}
\vspace{0.15cm}
\item ¿Cuál de los siguientes conjuntos es equivalente a \(\{a,b,c\}\)?
\begin{enumerate}[label=\textbf{\Alph*)}]
\begin{multicols}{2}
  \item {1,2,3}
  \item U
  \item 2
  \item 4
  \item 0
  \item Ninguna
\end{multicols}
\end{enumerate}
\vspace{0.15cm}
\item El conjunto potencia de A se denota:
\begin{enumerate}[label=\textbf{\Alph*)}]
\begin{multicols}{2}
  \item P(A)
  \item U
  \item 2
  \item \mathbb{R}
  \item \mathbb{C}
  \item \mathbb{I}
\end{multicols}
\end{enumerate}
\vspace{0.15cm}
\item ¿Cuántos subconjuntos tiene un conjunto con 3 elementos?
\begin{enumerate}[label=\textbf{\Alph*)}]
\begin{multicols}{2}
  \item 0
  \item 1
  \item 8
  \item 4
  \item 0
  \item Ninguna
\end{multicols}
\end{enumerate}
\vspace{0.15cm}
\item ¿Cuántos subconjuntos tiene un conjunto con 2 elementos?
\begin{enumerate}[label=\textbf{\Alph*)}]
\begin{multicols}{2}
  \item 0
  \item 1
  \item 4
  \item 4
  \item 0
  \item Ninguna
\end{multicols}
\end{enumerate}
\vspace{0.15cm}
\item ¿Cuántos subconjuntos tiene un conjunto con 1 elemento?
\begin{enumerate}[label=\textbf{\Alph*)}]
\begin{multicols}{2}
  \item 0
  \item 2
  \item 2
  \item 4
  \item 0
  \item Ninguna
\end{multicols}
\end{enumerate}
\vspace{0.15cm}
\item ¿Cuántos subconjuntos tiene el conjunto vacío?
\begin{enumerate}[label=\textbf{\Alph*)}]
\begin{multicols}{2}
  \item 0
  \item 1
  \item 2
  \item 4
  \item 0
  \item Ninguna
\end{multicols}
\end{enumerate}
\vspace{0.15cm}
\item ¿Cuál es el conjunto potencia de \(\{1,2\}\)?
\begin{enumerate}[label=\textbf{\Alph*)}]
\begin{multicols}{2}
  \item \{\emptyset,\{1\},\{2\},\{1,2\}\}
  \item U
  \item 2
  \item 4
  \item 0
  \item Ninguna
\end{multicols}
\end{enumerate}
\vspace{0.15cm}
\item ¿Cuál es el conjunto potencia de \(\{a\}\)?
\begin{enumerate}[label=\textbf{\Alph*)}]
\begin{multicols}{2}
  \item \{\emptyset,\{a\}\}
  \item U
  \item 2
  \item 4
  \item 0
  \item Ninguna
\end{multicols}
\end{enumerate}
\vspace{0.15cm}
\item ¿Cuál de los siguientes es subconjunto de \(\{1,2,3\}\)?
\begin{enumerate}[label=\textbf{\Alph*)}]
\begin{multicols}{2}
  \item {1,2}
  \item U
  \item 2
  \item 4
  \item 0
  \item Ninguna
\end{multicols}
\end{enumerate}
\vspace{0.15cm}
\item ¿Cuál de los siguientes NO es subconjunto de \(\{1,2,3\}\)?
\begin{enumerate}[label=\textbf{\Alph*)}]
\begin{multicols}{2}
  \item 0
  \item U
  \item 2
  \item {4}
  \item 0
  \item Ninguna
\end{multicols}
\end{enumerate}
\vspace{0.15cm}
\item El símbolo \(\subset\) significa:
\begin{enumerate}[label=\textbf{\Alph*)}]
\begin{multicols}{2}
  \item 0
  \item Subconjunto
  \item 2
  \item 4
  \item 0
  \item Ninguna
\end{multicols}
\end{enumerate}
\vspace{0.15cm}
\item El símbolo \(\subseteq\) significa:
\begin{enumerate}[label=\textbf{\Alph*)}]
\begin{multicols}{2}
  \item 0
  \item Subconjunto o igual
  \item 2
  \item 4
  \item 0
  \item Ninguna
\end{multicols}
\end{enumerate}
\vspace{0.15cm}
\item Si \(A = \{1,2\}\) y \(B = \{1,2,3\}\), entonces:
\begin{enumerate}[label=\textbf{\Alph*)}]
\begin{multicols}{2}
  \item A \subset B
  \item U
  \item 2
  \item 4
  \item 0
  \item Ninguna
\end{multicols}
\end{enumerate}
\vspace{0.15cm}
\item Si \(A = \{1,2,3\}\) y \(B = \{1,2,3\}\), entonces:
\begin{enumerate}[label=\textbf{\Alph*)}]
\begin{multicols}{2}
  \item 0
  \item A = B
  \item 2
  \item 4
  \item 0
  \item Ninguna
\end{multicols}
\end{enumerate}
\vspace{0.15cm}
\item ¿Cuál de los siguientes conjuntos tiene cardinalidad 5?
\begin{enumerate}[label=\textbf{\Alph*)}]
\begin{multicols}{2}
  \item {1,2,3,4,5}
  \item U
  \item 2
  \item 4
  \item 0
  \item Ninguna
\end{multicols}
\end{enumerate}
\vspace{0.15cm}
\item El conjunto \(\mathbb{N} \cup \{0\}\) se denota como:
\begin{enumerate}[label=\textbf{\Alph*)}]
\begin{multicols}{2}
  \item \mathbb{N}_0
  \item U
  \item 2
  \item \mathbb{R}
  \item \mathbb{C}
  \item \mathbb{I}
\end{multicols}
\end{enumerate}
\vspace{0.15cm}
\item ¿Cuál de los siguientes es un conjunto por extensión?
\begin{enumerate}[label=\textbf{\Alph*)}]
\begin{multicols}{2}
  \item 0
  \item {1,2,3,4,5}
  \item 2
  \item 4
  \item 0
  \item Ninguna
\end{multicols}
\end{enumerate}
\vspace{0.15cm}
\item ¿Cuál de los siguientes es un conjunto por comprensión?
\begin{enumerate}[label=\textbf{\Alph*)}]
\begin{multicols}{2}
  \item 0
  \item U
  \item \{x \mid x \text{ es par}\}
  \item 4
  \item 0
  \item Ninguna
\end{multicols}
\end{enumerate}
\vspace{0.15cm}
\item El conjunto \(\{x \in \mathbb{N} \mid x < 5\}\) por extensión es:
\begin{enumerate}[label=\textbf{\Alph*)}]
\begin{multicols}{2}
  \item {1,2,3,4}
  \item U
  \item 2
  \item 4
  \item 0
  \item Ninguna
\end{multicols}
\end{enumerate}
\vspace{0.15cm}
\item El conjunto \(\{x \in \mathbb{Z} \mid -2 \leq x \leq 2\}\) por extensión es:
\begin{enumerate}[label=\textbf{\Alph*)}]
\begin{multicols}{2}
  \item {-2,-1,0,1,2}
  \item U
  \item 2
  \item 4
  \item 0
  \item Ninguna
\end{multicols}
\end{enumerate}
\vspace{0.15cm}
\item ¿Cuál es el conjunto de los números pares positivos?
\begin{enumerate}[label=\textbf{\Alph*)}]
\begin{multicols}{2}
  \item {2,4,6,8,...}
  \item U
  \item 2
  \item 4
  \item 0
  \item Ninguna
\end{multicols}
\end{enumerate}
\vspace{0.15cm}
\item ¿Cuál es el conjunto de los números impares positivos?
\begin{enumerate}[label=\textbf{\Alph*)}]
\begin{multicols}{2}
  \item 0
  \item {1,3,5,7,...}
  \item 2
  \item 4
  \item 0
  \item Ninguna
\end{multicols}
\end{enumerate}
\vspace{0.15cm}
\item El conjunto \(\{x \in \mathbb{R} \mid x^2 = 4\}\) por extensión es:
\begin{enumerate}[label=\textbf{\Alph*)}]
\begin{multicols}{2}
  \item 0
  \item U
  \item {-2,2}
  \item 4
  \item 0
  \item Ninguna
\end{multicols}
\end{enumerate}
\vspace{0.15cm}
\item El conjunto \(\{x \in \mathbb{R} \mid x^2 = -1\}\) es:
\begin{enumerate}[label=\textbf{\Alph*)}]
\begin{multicols}{2}
  \item 0
  \item U
  \item \emptyset
  \item 4
  \item 0
  \item Ninguna
\end{multicols}
\end{enumerate}
\vspace{0.15cm}
\item El conjunto \(\{x \in \mathbb{N} \mid x^2 = 9\}\) por extensión es:
\begin{enumerate}[label=\textbf{\Alph*)}]
\begin{multicols}{2}
  \item {3}
  \item U
  \item 2
  \item 4
  \item 0
  \item Ninguna
\end{multicols}
\end{enumerate}
\vspace{0.15cm}
\item Si \(A = \{1,2,3,4,5\}\), ¿cuántos subconjuntos propios tiene?
\begin{enumerate}[label=\textbf{\Alph*)}]
\begin{multicols}{2}
  \item 0
  \item 31
  \item 2
  \item 4
  \item 0
  \item Ninguna
\end{multicols}
\end{enumerate}
\vspace{0.15cm}
\end{enumerate}\n
\subsubsection*{SECCIÓN II: PERTENENCIA Y SUBCONJUNTOS (51–100)}

\begin{enumerate}[label=\textbf{\arabic*.}]
\setcounter{enumi}{50}
\item ¿Es \(3 \in \{1,2,3,4\}\)?
\begin{enumerate}[label=\textbf{\Alph*)}]
\begin{multicols}{2}
  \item Sí
  \item No
  \item A veces
  \item Depende
  \item No se puede
  \item Falta información
\end{multicols}
\end{enumerate}
\vspace{0.15cm}
\item ¿Es \(5 \in \{1,2,3,4\}\)?
\begin{enumerate}[label=\textbf{\Alph*)}]
\begin{multicols}{2}
  \item Sí
  \item No
  \item A veces
  \item Depende
  \item No se puede
  \item Falta información
\end{multicols}
\end{enumerate}
\vspace{0.15cm}
\item ¿Es \(\emptyset \subset \{1,2,3\}\)?
\begin{enumerate}[label=\textbf{\Alph*)}]
\begin{multicols}{2}
  \item Sí
  \item No
  \item A veces
  \item Depende
  \item No se puede
  \item Falta información
\end{multicols}
\end{enumerate}
\vspace{0.15cm}
\item ¿Es \(\{1,2\} \subset \{1,2,3\}\)?
\begin{enumerate}[label=\textbf{\Alph*)}]
\begin{multicols}{2}
  \item Sí
  \item No
  \item A veces
  \item Depende
  \item No se puede
  \item Falta información
\end{multicols}
\end{enumerate}
\vspace{0.15cm}
\item ¿Es \(\{1,4\} \subset \{1,2,3\}\)?
\begin{enumerate}[label=\textbf{\Alph*)}]
\begin{multicols}{2}
  \item Sí
  \item No
  \item A veces
  \item Depende
  \item No se puede
  \item Falta información
\end{multicols}
\end{enumerate}
\vspace{0.15cm}
\item ¿Es \(\{1,2,3\} \subseteq \{1,2,3\}\)?
\begin{enumerate}[label=\textbf{\Alph*)}]
\begin{multicols}{2}
  \item Sí
  \item No
  \item A veces
  \item Depende
  \item No se puede
  \item Falta información
\end{multicols}
\end{enumerate}
\vspace{0.15cm}
\item ¿Es \(\{1,2,3\} \subset \{1,2,3\}\)?
\begin{enumerate}[label=\textbf{\Alph*)}]
\begin{multicols}{2}
  \item Sí
  \item No
  \item A veces
  \item Depende
  \item No se puede
  \item Falta información
\end{multicols}
\end{enumerate}
\vspace{0.15cm}
\item Si \(A = \{a,b,c\}\), ¿es \(a \in A\)?
\begin{enumerate}[label=\textbf{\Alph*)}]
\begin{multicols}{2}
  \item Sí
  \item No
  \item A veces
  \item Indeterminada
  \item No se puede
  \item Falta información
\end{multicols}
\end{enumerate}
\vspace{0.15cm}
\item Si \(A = \{a,b,c\}\), ¿es \(\{a\} \in A\)?
\begin{enumerate}[label=\textbf{\Alph*)}]
\begin{multicols}{2}
  \item Sí
  \item No
  \item A veces
  \item Indeterminada
  \item No se puede
  \item Falta información
\end{multicols}
\end{enumerate}
\vspace{0.15cm}
\item Si \(A = \{a,b,c\}\), ¿es \(\{a\} \subset A\)?
\begin{enumerate}[label=\textbf{\Alph*)}]
\begin{multicols}{2}
  \item Sí
  \item No
  \item A veces
  \item Indeterminada
  \item No se puede
  \item Falta información
\end{multicols}
\end{enumerate}
\vspace{0.15cm}
\item ¿Cuál es el número de subconjuntos de un conjunto con 4 elementos?
\begin{enumerate}[label=\textbf{\Alph*)}]
\begin{multicols}{2}
  \item 4
  \item 8
  \item 16
  \item 32
  \item 64
  \item 12
\end{multicols}
\end{enumerate}
\vspace{0.15cm}
\item ¿Cuál es el número de subconjuntos propios de un conjunto con 4 elementos?
\begin{enumerate}[label=\textbf{\Alph*)}]
\begin{multicols}{2}
  \item 4
  \item 15
  \item 16
  \item 32
  \item 64
  \item 12
\end{multicols}
\end{enumerate}
\vspace{0.15cm}
\item ¿Cuál es el número de subconjuntos de un conjunto con 5 elementos?
\begin{enumerate}[label=\textbf{\Alph*)}]
\begin{multicols}{2}
  \item 4
  \item 32
  \item 16
  \item 32
  \item 64
  \item 12
\end{multicols}
\end{enumerate}
\vspace{0.15cm}
\item ¿Cuál es el número de subconjuntos propios de un conjunto con 5 elementos?
\begin{enumerate}[label=\textbf{\Alph*)}]
\begin{multicols}{2}
  \item 4
  \item 31
  \item 16
  \item 32
  \item 64
  \item 12
\end{multicols}
\end{enumerate}
\vspace{0.15cm}
\item Si \(A \subset B\) y \(B \subset C\), entonces:
\begin{enumerate}[label=\textbf{\Alph*)}]
\begin{multicols}{2}
  \item A \subset C
  \item \emptyset
  \item {1,2}
  \item {1,2,3}
  \item U
  \item {1}
\end{multicols}
\end{enumerate}
\vspace{0.15cm}
\item Si \(A \subseteq B\), significa que:
\begin{enumerate}[label=\textbf{\Alph*)}]
\begin{multicols}{2}
  \item Todo elemento de A está en B
  \item \emptyset
  \item {1,2}
  \item {1,2,3}
  \item U
  \item {1}
\end{multicols}
\end{enumerate}
\vspace{0.15cm}
\item Si \(A \subset B\), significa que:
\begin{enumerate}[label=\textbf{\Alph*)}]
\begin{multicols}{2}
  \item A \subseteq B y A \neq B
  \item \emptyset
  \item {1,2}
  \item {1,2,3}
  \item U
  \item {1}
\end{multicols}
\end{enumerate}
\vspace{0.15cm}
\item ¿Cuál de los siguientes conjuntos es subconjunto de \(\mathbb{N}\)?
\begin{enumerate}[label=\textbf{\Alph*)}]
\begin{multicols}{2}
  \item {3,4}
  \item {1,2,3}
  \item {1,2}
  \item {1,2,3}
  \item U
  \item {1}
\end{multicols}
\end{enumerate}
\vspace{0.15cm}
\item ¿Cuál de los siguientes conjuntos es subconjunto de \(\mathbb{Z}\)?
\begin{enumerate}[label=\textbf{\Alph*)}]
\begin{multicols}{2}
  \item {-1,0,1}
  \item \emptyset
  \item {1,2}
  \item {1,2,3}
  \item U
  \item {1}
\end{multicols}
\end{enumerate}
\vspace{0.15cm}
\item ¿Cuál de los siguientes conjuntos es subconjunto de \(\mathbb{Q}\)?
\begin{enumerate}[label=\textbf{\Alph*)}]
\begin{multicols}{2}
  \item {3,4}
  \item \emptyset
  \item \{\frac{1}{2},\frac{3}{4}\}
  \item {1,2,3}
  \item U
  \item {1}
\end{multicols}
\end{enumerate}
\vspace{0.15cm}
\item ¿Cuál de los siguientes conjuntos es subconjunto de \(\mathbb{R}\)?
\begin{enumerate}[label=\textbf{\Alph*)}]
\begin{multicols}{2}
  \item \{\pi,\sqrt{2},e\}
  \item \emptyset
  \item {1,2}
  \item {1,2,3}
  \item U
  \item {1}
\end{multicols}
\end{enumerate}
\vspace{0.15cm}
\item La relación \(\mathbb{N} \subset \mathbb{Z}\) es:
\begin{enumerate}[label=\textbf{\Alph*)}]
\begin{multicols}{2}
  \item Sí
  \item No
  \item A veces
  \item Indeterminada
  \item No se puede
  \item Falta información
\end{multicols}
\end{enumerate}
\vspace{0.15cm}
\item La relación \(\mathbb{Z} \subset \mathbb{N}\) es:
\begin{enumerate}[label=\textbf{\Alph*)}]
\begin{multicols}{2}
  \item Sí
  \item No
  \item A veces
  \item Indeterminada
  \item No se puede
  \item Falta información
\end{multicols}
\end{enumerate}
\vspace{0.15cm}
\item La relación \(\mathbb{Q} \subset \mathbb{R}\) es:
\begin{enumerate}[label=\textbf{\Alph*)}]
\begin{multicols}{2}
  \item Sí
  \item No
  \item A veces
  \item Indeterminada
  \item No se puede
  \item Falta información
\end{multicols}
\end{enumerate}
\vspace{0.15cm}
\item La relación \(\mathbb{I} \subset \mathbb{R}\) es:
\begin{enumerate}[label=\textbf{\Alph*)}]
\begin{multicols}{2}
  \item Sí
  \item No
  \item A veces
  \item Indeterminada
  \item No se puede
  \item Falta información
\end{multicols}
\end{enumerate}
\vspace{0.15cm}
\item La relación \(\mathbb{N} \subset \mathbb{Q}\) es:
\begin{enumerate}[label=\textbf{\Alph*)}]
\begin{multicols}{2}
  \item Sí
  \item No
  \item A veces
  \item Indeterminada
  \item No se puede
  \item Falta información
\end{multicols}
\end{enumerate}
\vspace{0.15cm}
\item La relación \(\mathbb{Z} \subset \mathbb{Q}\) es:
\begin{enumerate}[label=\textbf{\Alph*)}]
\begin{multicols}{2}
  \item Sí
  \item No
  \item A veces
  \item Indeterminada
  \item No se puede
  \item Falta información
\end{multicols}
\end{enumerate}
\vspace{0.15cm}
\item ¿Cuál es el complemento de \(A = \{1,2\}\) en \(U = \{1,2,3,4\}\)?
\begin{enumerate}[label=\textbf{\Alph*)}]
\begin{multicols}{2}
  \item {3,4}
  \item \emptyset
  \item {1,2}
  \item {1,2,3}
  \item U
  \item {1}
\end{multicols}
\end{enumerate}
\vspace{0.15cm}
\item ¿Cuál es el complemento de \(A = \{1,2,3\}\) en \(U = \{1,2,3\}\)?
\begin{enumerate}[label=\textbf{\Alph*)}]
\begin{multicols}{2}
  \item {3,4}
  \item \emptyset
  \item {1,2}
  \item {1,2,3}
  \item U
  \item {1}
\end{multicols}
\end{enumerate}
\vspace{0.15cm}
\item ¿Cuál es el complemento de \(\emptyset\) en \(U\)?
\begin{enumerate}[label=\textbf{\Alph*)}]
\begin{multicols}{2}
  \item U
  \item \emptyset
  \item {1,2}
  \item {1,2,3}
  \item U
  \item {1}
\end{multicols}
\end{enumerate}
\vspace{0.15cm}
\item ¿Cuál es el complemento de \(U\) en \(U\)?
\begin{enumerate}[label=\textbf{\Alph*)}]
\begin{multicols}{2}
  \item {3,4}
  \item \emptyset
  \item {1,2}
  \item {1,2,3}
  \item U
  \item {1}
\end{multicols}
\end{enumerate}
\vspace{0.15cm}
\item Si \(A = \{1,2,3,4,5\}\) y \(B = \{1,3,5\}\), entonces \(B \subset A\) es:
\begin{enumerate}[label=\textbf{\Alph*)}]
\begin{multicols}{2}
  \item Sí
  \item No
  \item A veces
  \item Indeterminada
  \item No se puede
  \item Falta información
\end{multicols}
\end{enumerate}
\vspace{0.15cm}
\item Si \(A = \{1,2,3\}\) y \(B = \{1,2,3,4\}\), entonces \(A \subset B\) es:
\begin{enumerate}[label=\textbf{\Alph*)}]
\begin{multicols}{2}
  \item Sí
  \item No
  \item A veces
  \item Indeterminada
  \item No se puede
  \item Falta información
\end{multicols}
\end{enumerate}
\vspace{0.15cm}
\item Si \(A = \{1,2,3\}\) y \(B = \{1,2,3\}\), entonces \(A \subset B\) es:
\begin{enumerate}[label=\textbf{\Alph*)}]
\begin{multicols}{2}
  \item Sí
  \item No
  \item A veces
  \item Indeterminada
  \item No se puede
  \item Falta información
\end{multicols}
\end{enumerate}
\vspace{0.15cm}
\item Si \(A = \{1,2,3\}\) y \(B = \{1,2,3\}\), entonces \(A \subseteq B\) es:
\begin{enumerate}[label=\textbf{\Alph*)}]
\begin{multicols}{2}
  \item Sí
  \item No
  \item A veces
  \item Indeterminada
  \item No se puede
  \item Falta información
\end{multicols}
\end{enumerate}
\vspace{0.15cm}
\item ¿Cuántos elementos tiene el conjunto potencia de un conjunto con n elementos?
\begin{enumerate}[label=\textbf{\Alph*)}]
\begin{multicols}{2}
  \item 4
  \item 8
  \item 2^n
  \item 32
  \item 64
  \item 12
\end{multicols}
\end{enumerate}
\vspace{0.15cm}
\item ¿Cuántos elementos tiene el conjunto potencia de \(\{1,2,3,4\}\)?
\begin{enumerate}[label=\textbf{\Alph*)}]
\begin{multicols}{2}
  \item 4
  \item 8
  \item 16
  \item 32
  \item 64
  \item 12
\end{multicols}
\end{enumerate}
\vspace{0.15cm}
\item ¿Cuántos elementos tiene el conjunto potencia de \(\emptyset\)?
\begin{enumerate}[label=\textbf{\Alph*)}]
\begin{multicols}{2}
  \item 4
  \item 1
  \item 16
  \item 32
  \item 64
  \item 12
\end{multicols}
\end{enumerate}
\vspace{0.15cm}
\item El conjunto potencia de A contiene:
\begin{enumerate}[label=\textbf{\Alph*)}]
\begin{multicols}{2}
  \item Todos los subconjuntos de A
  \item \emptyset
  \item {1,2}
  \item {1,2,3}
  \item U
  \item {1}
\end{multicols}
\end{enumerate}
\vspace{0.15cm}
\item ¿Cuál es el conjunto potencia de \(\{1\}\)?
\begin{enumerate}[label=\textbf{\Alph*)}]
\begin{multicols}{2}
  \item \{\emptyset,\{1\}\}
  \item \emptyset
  \item {1,2}
  \item {1,2,3}
  \item U
  \item {1}
\end{multicols}
\end{enumerate}
\vspace{0.15cm}
\item Si \(A = \{1,2\}\), ¿cuál es \(P(A)\)?
\begin{enumerate}[label=\textbf{\Alph*)}]
\begin{multicols}{2}
  \item \{\emptyset,\{1\},\{2\},\{1,2\}\}
  \item \emptyset
  \item {1,2}
  \item {1,2,3}
  \item U
  \item {1}
\end{multicols}
\end{enumerate}
\vspace{0.15cm}
\item Si \(A = \{a,b,c\}\), ¿cuántos elementos tiene \(P(A)\)?
\begin{enumerate}[label=\textbf{\Alph*)}]
\begin{multicols}{2}
  \item 4
  \item 8
  \item 8
  \item 32
  \item 64
  \item 12
\end{multicols}
\end{enumerate}
\vspace{0.15cm}
\item ¿Cuál de los siguientes es un subconjunto propio de \(\{1,2,3\}\)?
\begin{enumerate}[label=\textbf{\Alph*)}]
\begin{multicols}{2}
  \item {3,4}
  \item {1,2}
  \item {1,2}
  \item {1,2,3}
  \item U
  \item {1}
\end{multicols}
\end{enumerate}
\vspace{0.15cm}
\item ¿Cuál de los siguientes NO es un subconjunto propio de \(\{1,2,3\}\)?
\begin{enumerate}[label=\textbf{\Alph*)}]
\begin{multicols}{2}
  \item {3,4}
  \item \emptyset
  \item {1,2}
  \item {1,2,3}
  \item U
  \item {1}
\end{multicols}
\end{enumerate}
\vspace{0.15cm}
\item Si \(A = \{1,2,3,4\}\), ¿cuántos subconjuntos de A contienen al 1?
\begin{enumerate}[label=\textbf{\Alph*)}]
\begin{multicols}{2}
  \item 4
  \item 8
  \item 16
  \item 32
  \item 64
  \item 12
\end{multicols}
\end{enumerate}
\vspace{0.15cm}
\item Si \(A = \{1,2,3,4\}\), ¿cuántos subconjuntos de A NO contienen al 1?
\begin{enumerate}[label=\textbf{\Alph*)}]
\begin{multicols}{2}
  \item 4
  \item 8
  \item 16
  \item 32
  \item 64
  \item 12
\end{multicols}
\end{enumerate}
\vspace{0.15cm}
\item Si \(A = \{1,2,3,4\}\), ¿cuántos subconjuntos de A tienen exactamente 2 elementos?
\begin{enumerate}[label=\textbf{\Alph*)}]
\begin{multicols}{2}
  \item 4
  \item 6
  \item 16
  \item 32
  \item 64
  \item 12
\end{multicols}
\end{enumerate}
\vspace{0.15cm}
\item Si \(A = \{1,2,3,4,5\}\), ¿cuántos subconjuntos de A tienen exactamente 3 elementos?
\begin{enumerate}[label=\textbf{\Alph*)}]
\begin{multicols}{2}
  \item 4
  \item 10
  \item 16
  \item 32
  \item 64
  \item 12
\end{multicols}
\end{enumerate}
\vspace{0.15cm}
\item Si \(A = \{1,2,3,4\}\), ¿cuántos subconjuntos de A tienen exactamente 1 elemento?
\begin{enumerate}[label=\textbf{\Alph*)}]
\begin{multicols}{2}
  \item 4
  \item 8
  \item 16
  \item 4
  \item 64
  \item 12
\end{multicols}
\end{enumerate}
\vspace{0.15cm}
\item Si \(A = \{1,2,3,4\}\), ¿cuántos subconjuntos de A tienen exactamente 0 elementos?
\begin{enumerate}[label=\textbf{\Alph*)}]
\begin{multicols}{2}
  \item 4
  \item 1
  \item 16
  \item 32
  \item 64
  \item 12
\end{multicols}
\end{enumerate}
\vspace{0.15cm}
\end{enumerate}\n
\subsubsection*{SECCIÓN III: OPERACIONES CON CONJUNTOS (101–160)}

\begin{enumerate}[label=\textbf{\arabic*.}]
\setcounter{enumi}{100}
\item ¿Qué es \(A \cup B\)?
\begin{enumerate}[label=\textbf{\Alph*)}]
\begin{multicols}{2}
  \item Unión de A y B
  \item B
  \item \emptyset
  \item U
  \item \emptyset
  \item Ninguna
\end{multicols}
\end{enumerate}
\vspace{0.15cm}
\item ¿Qué es \(A \cap B\)?
\begin{enumerate}[label=\textbf{\Alph*)}]
\begin{multicols}{2}
  \item A
  \item Intersección de A y B
  \item \emptyset
  \item U
  \item \emptyset
  \item Ninguna
\end{multicols}
\end{enumerate}
\vspace{0.15cm}
\item ¿Qué es \(A - B\)?
\begin{enumerate}[label=\textbf{\Alph*)}]
\begin{multicols}{2}
  \item Unión
  \item Intersección
  \item Diferencia
  \item {4,5}
  \item \emptyset
  \item Ninguna
\end{multicols}
\end{enumerate}
\vspace{0.15cm}
\item Si \(A = \{1,2,3\}\) y \(B = \{3,4,5\}\), ¿cuál es \(A \cup B\)?
\begin{enumerate}[label=\textbf{\Alph*)}]
\begin{multicols}{2}
  \item {1,2,3,4,5}
  \item B
  \item \emptyset
  \item U
  \item \emptyset
  \item Ninguna
\end{multicols}
\end{enumerate}
\vspace{0.15cm}
\item Si \(A = \{1,2,3\}\) y \(B = \{3,4,5\}\), ¿cuál es \(A \cap B\)?
\begin{enumerate}[label=\textbf{\Alph*)}]
\begin{multicols}{2}
  \item A
  \item {3}
  \item \emptyset
  \item U
  \item \emptyset
  \item Ninguna
\end{multicols}
\end{enumerate}
\vspace{0.15cm}
\item Si \(A = \{1,2,3\}\) y \(B = \{3,4,5\}\), ¿cuál es \(A - B\)?
\begin{enumerate}[label=\textbf{\Alph*)}]
\begin{multicols}{2}
  \item {1,2}
  \item B
  \item {1,2}
  \item {4,5}
  \item \emptyset
  \item Ninguna
\end{multicols}
\end{enumerate}
\vspace{0.15cm}
\item Si \(A = \{1,2,3\}\) y \(B = \{3,4,5\}\), ¿cuál es \(B - A\)?
\begin{enumerate}[label=\textbf{\Alph*)}]
\begin{multicols}{2}
  \item {1,2}
  \item B
  \item {1,2}
  \item {4,5}
  \item \emptyset
  \item Ninguna
\end{multicols}
\end{enumerate}
\vspace{0.15cm}
\item Si \(A = \{1,2\}\) y \(B = \{3,4\}\), ¿cuál es \(A \cap B\)?
\begin{enumerate}[label=\textbf{\Alph*)}]
\begin{multicols}{2}
  \item A
  \item \emptyset
  \item \emptyset
  \item U
  \item \emptyset
  \item Ninguna
\end{multicols}
\end{enumerate}
\vspace{0.15cm}
\item Si \(A = \{1,2\}\) y \(B = \{3,4\}\), ¿cuál es \(A \cup B\)?
\begin{enumerate}[label=\textbf{\Alph*)}]
\begin{multicols}{2}
  \item {1,2,3,4}
  \item B
  \item \emptyset
  \item U
  \item \emptyset
  \item Ninguna
\end{multicols}
\end{enumerate}
\vspace{0.15cm}
\item Si \(A = \{1,2,3,4,5\}\) y \(B = \{2,4\}\), ¿cuál es \(A \cap B\)?
\begin{enumerate}[label=\textbf{\Alph*)}]
\begin{multicols}{2}
  \item {2,4}
  \item \emptyset
  \item \emptyset
  \item U
  \item \emptyset
  \item Ninguna
\end{multicols}
\end{enumerate}
\vspace{0.15cm}
\item Si \(A = \{1,2,3,4,5\}\) y \(B = \{2,4\}\), ¿cuál es \(A - B\)?
\begin{enumerate}[label=\textbf{\Alph*)}]
\begin{multicols}{2}
  \item {1,2}
  \item {1,3,5}
  \item {1,2}
  \item {4,5}
  \item \emptyset
  \item Ninguna
\end{multicols}
\end{enumerate}
\vspace{0.15cm}
\item Si \(A = \{1,2,3\}\) y \(B = \{1,2,3\}\), ¿cuál es \(A \cup B\)?
\begin{enumerate}[label=\textbf{\Alph*)}]
\begin{multicols}{2}
  \item {1,2,3}
  \item B
  \item \emptyset
  \item U
  \item \emptyset
  \item Ninguna
\end{multicols}
\end{enumerate}
\vspace{0.15cm}
\item Si \(A = \{1,2,3\}\) y \(B = \{1,2,3\}\), ¿cuál es \(A \cap B\)?
\begin{enumerate}[label=\textbf{\Alph*)}]
\begin{multicols}{2}
  \item {1,2,3}
  \item \emptyset
  \item \emptyset
  \item U
  \item \emptyset
  \item Ninguna
\end{multicols}
\end{enumerate}
\vspace{0.15cm}
\item Si \(A = \{1,2,3\}\) y \(B = \{1,2,3\}\), ¿cuál es \(A - B\)?
\begin{enumerate}[label=\textbf{\Alph*)}]
\begin{multicols}{2}
  \item {1,2}
  \item \emptyset
  \item {1,2}
  \item {4,5}
  \item \emptyset
  \item Ninguna
\end{multicols}
\end{enumerate}
\vspace{0.15cm}
\item Si \(A = \emptyset\) y \(B = \{1,2\}\), ¿cuál es \(A \cup B\)?
\begin{enumerate}[label=\textbf{\Alph*)}]
\begin{multicols}{2}
  \item {1,2}
  \item B
  \item \emptyset
  \item U
  \item \emptyset
  \item Ninguna
\end{multicols}
\end{enumerate}
\vspace{0.15cm}
\item Si \(A = \emptyset\) y \(B = \{1,2\}\), ¿cuál es \(A \cap B\)?
\begin{enumerate}[label=\textbf{\Alph*)}]
\begin{multicols}{2}
  \item A
  \item \emptyset
  \item \emptyset
  \item U
  \item \emptyset
  \item Ninguna
\end{multicols}
\end{enumerate}
\vspace{0.15cm}
\item Si \(A = \{1,2,3,4\}\) y \(B = \{3,4,5,6\}\), ¿cuál es \(A \cup B\)?
\begin{enumerate}[label=\textbf{\Alph*)}]
\begin{multicols}{2}
  \item {1,2,3,4,5,6}
  \item B
  \item \emptyset
  \item U
  \item \emptyset
  \item Ninguna
\end{multicols}
\end{enumerate}
\vspace{0.15cm}
\item Si \(A = \{1,2,3,4\}\) y \(B = \{3,4,5,6\}\), ¿cuál es \(A \cap B\)?
\begin{enumerate}[label=\textbf{\Alph*)}]
\begin{multicols}{2}
  \item A
  \item {3,4}
  \item \emptyset
  \item U
  \item \emptyset
  \item Ninguna
\end{multicols}
\end{enumerate}
\vspace{0.15cm}
\item Si \(A = \{1,2,3,4\}\) y \(B = \{3,4,5,6\}\), ¿cuál es \(A - B\)?
\begin{enumerate}[label=\textbf{\Alph*)}]
\begin{multicols}{2}
  \item {1,2}
  \item B
  \item {1,2}
  \item {4,5}
  \item \emptyset
  \item Ninguna
\end{multicols}
\end{enumerate}
\vspace{0.15cm}
\item Si \(A = \{1,2,3,4\}\) y \(B = \{3,4,5,6\}\), ¿cuál es \(B - A\)?
\begin{enumerate}[label=\textbf{\Alph*)}]
\begin{multicols}{2}
  \item {1,2}
  \item B
  \item {1,2}
  \item {5,6}
  \item \emptyset
  \item Ninguna
\end{multicols}
\end{enumerate}
\vspace{0.15cm}
\item \(A \cup \emptyset =\)
\begin{enumerate}[label=\textbf{\Alph*)}]
\begin{multicols}{2}
  \item A
  \item B
  \item \emptyset
  \item U
  \item \emptyset
  \item Ninguna
\end{multicols}
\end{enumerate}
\vspace{0.15cm}
\item \(A \cap \emptyset =\)
\begin{enumerate}[label=\textbf{\Alph*)}]
\begin{multicols}{2}
  \item A
  \item \emptyset
  \item \emptyset
  \item U
  \item \emptyset
  \item Ninguna
\end{multicols}
\end{enumerate}
\vspace{0.15cm}
\item \(A \cup A =\)
\begin{enumerate}[label=\textbf{\Alph*)}]
\begin{multicols}{2}
  \item A
  \item B
  \item \emptyset
  \item U
  \item \emptyset
  \item Ninguna
\end{multicols}
\end{enumerate}
\vspace{0.15cm}
\item \(A \cap A =\)
\begin{enumerate}[label=\textbf{\Alph*)}]
\begin{multicols}{2}
  \item A
  \item \emptyset
  \item \emptyset
  \item U
  \item \emptyset
  \item Ninguna
\end{multicols}
\end{enumerate}
\vspace{0.15cm}
\item \(A \cup U =\)
\begin{enumerate}[label=\textbf{\Alph*)}]
\begin{multicols}{2}
  \item {1,2,3,4,5}
  \item B
  \item U
  \item U
  \item \emptyset
  \item Ninguna
\end{multicols}
\end{enumerate}
\vspace{0.15cm}
\item \(A \cap U =\)
\begin{enumerate}[label=\textbf{\Alph*)}]
\begin{multicols}{2}
  \item A
  \item \emptyset
  \item \emptyset
  \item U
  \item \emptyset
  \item Ninguna
\end{multicols}
\end{enumerate}
\vspace{0.15cm}
\item La propiedad conmutativa de la unión establece:
\begin{enumerate}[label=\textbf{\Alph*)}]
\begin{multicols}{2}
  \item A \cup B = B \cup A
  \item B
  \item \emptyset
  \item U
  \item \emptyset
  \item Ninguna
\end{multicols}
\end{enumerate}
\vspace{0.15cm}
\item La propiedad asociativa de la unión establece:
\begin{enumerate}[label=\textbf{\Alph*)}]
\begin{multicols}{2}
  \item A
  \item A \cup (B \cup C) = (A \cup B) \cup C
  \item \emptyset
  \item U
  \item \emptyset
  \item Ninguna
\end{multicols}
\end{enumerate}
\vspace{0.15cm}
\item La propiedad distributiva establece:
\begin{enumerate}[label=\textbf{\Alph*)}]
\begin{multicols}{2}
  \item A \cap (B \cup C) = (A \cap B) \cup (A \cap C)
  \item B
  \item \emptyset
  \item U
  \item \emptyset
  \item Ninguna
\end{multicols}
\end{enumerate}
\vspace{0.15cm}
\item La ley de De Morgan para conjuntos establece:
\begin{enumerate}[label=\textbf{\Alph*)}]
\begin{multicols}{2}
  \item \overline{A \cup B} = \bar{A} \cap \bar{B}
  \item B
  \item \emptyset
  \item U
  \item \emptyset
  \item Ninguna
\end{multicols}
\end{enumerate}
\vspace{0.15cm}
\item La ley de De Morgan también establece:
\begin{enumerate}[label=\textbf{\Alph*)}]
\begin{multicols}{2}
  \item \overline{A \cap B} = \bar{A} \cup \bar{B}
  \item B
  \item \emptyset
  \item U
  \item \emptyset
  \item Ninguna
\end{multicols}
\end{enumerate}
\vspace{0.15cm}
\item Si \(A = \{1,2,3,4,5\}\) y \(B = \{4,5,6,7,8\}\), ¿cuál es \(A \cap B\)?
\begin{enumerate}[label=\textbf{\Alph*)}]
\begin{multicols}{2}
  \item {4,5}
  \item \emptyset
  \item \emptyset
  \item U
  \item \emptyset
  \item Ninguna
\end{multicols}
\end{enumerate}
\vspace{0.15cm}
\item Si \(A = \{1,2,3,4,5\}\) y \(B = \{4,5,6,7,8\}\), ¿cuál es \(A \cup B\)?
\begin{enumerate}[label=\textbf{\Alph*)}]
\begin{multicols}{2}
  \item {1,2,3,4,5}
  \item B
  \item \emptyset
  \item {1,2,3,4,5,6,7,8}
  \item \emptyset
  \item Ninguna
\end{multicols}
\end{enumerate}
\vspace{0.15cm}
\item Si \(A = \{1,2,3,4,5\}\) y \(B = \{4,5,6,7,8\}\), ¿cuál es \(A - B\)?
\begin{enumerate}[label=\textbf{\Alph*)}]
\begin{multicols}{2}
  \item {1,2}
  \item {1,2,3}
  \item {1,2}
  \item {4,5}
  \item \emptyset
  \item Ninguna
\end{multicols}
\end{enumerate}
\vspace{0.15cm}
\item Si \(A = \{1,2,3,4,5\}\) y \(B = \{4,5,6,7,8\}\), ¿cuál es \(B - A\)?
\begin{enumerate}[label=\textbf{\Alph*)}]
\begin{multicols}{2}
  \item {1,2}
  \item B
  \item {6,7,8}
  \item {4,5}
  \item \emptyset
  \item Ninguna
\end{multicols}
\end{enumerate}
\vspace{0.15cm}
\item Si \(U = \{1,2,3,4,5,6,7,8,9,10\}\) y \(A = \{2,4,6,8,10\}\), ¿cuál es \(\bar{A}\)?
\begin{enumerate}[label=\textbf{\Alph*)}]
\begin{multicols}{2}
  \item {1,3,5,7,9}
  \item B
  \item \emptyset
  \item U
  \item \emptyset
  \item Ninguna
\end{multicols}
\end{enumerate}
\vspace{0.15cm}
\item Si \(U = \{1,2,3,4,5,6,7,8,9,10\}\) y \(A = \{1,2,3,4,5\}\), ¿cuál es \(\bar{A}\)?
\begin{enumerate}[label=\textbf{\Alph*)}]
\begin{multicols}{2}
  \item {6,7,8,9,10}
  \item B
  \item \emptyset
  \item U
  \item \emptyset
  \item Ninguna
\end{multicols}
\end{enumerate}
\vspace{0.15cm}
\item Si \(A \subset B\), entonces \(A \cup B =\)
\begin{enumerate}[label=\textbf{\Alph*)}]
\begin{multicols}{2}
  \item {1,2,3,4,5}
  \item B
  \item \emptyset
  \item U
  \item \emptyset
  \item Ninguna
\end{multicols}
\end{enumerate}
\vspace{0.15cm}
\item Si \(A \subset B\), entonces \(A \cap B =\)
\begin{enumerate}[label=\textbf{\Alph*)}]
\begin{multicols}{2}
  \item A
  \item \emptyset
  \item \emptyset
  \item U
  \item \emptyset
  \item Ninguna
\end{multicols}
\end{enumerate}
\vspace{0.15cm}
\item Si \(A \subset B\), entonces \(A - B =\)
\begin{enumerate}[label=\textbf{\Alph*)}]
\begin{multicols}{2}
  \item {1,2}
  \item B
  \item \emptyset
  \item {4,5}
  \item \emptyset
  \item Ninguna
\end{multicols}
\end{enumerate}
\vspace{0.15cm}
\item Si A y B son disjuntos, entonces \(A \cap B =\)
\begin{enumerate}[label=\textbf{\Alph*)}]
\begin{multicols}{2}
  \item A
  \item \emptyset
  \item \emptyset
  \item U
  \item \emptyset
  \item Ninguna
\end{multicols}
\end{enumerate}
\vspace{0.15cm}
\item Si A y B son disjuntos, entonces A cup B se llama:
\begin{enumerate}[label=\textbf{\Alph*)}]
\begin{multicols}{2}
  \item Unión disjunta
  \item B
  \item \emptyset
  \item U
  \item \emptyset
  \item Ninguna
\end{multicols}
\end{enumerate}
\vspace{0.15cm}
\item ¿Cuál es la intersección de \(\{1,2,3\}\) y \(\{4,5,6\}\)?
\begin{enumerate}[label=\textbf{\Alph*)}]
\begin{multicols}{2}
  \item A
  \item \emptyset
  \item \emptyset
  \item U
  \item \emptyset
  \item Ninguna
\end{multicols}
\end{enumerate}
\vspace{0.15cm}
\item ¿Cuál es la unión de \(\{1,2,3\}\) y \(\{4,5,6\}\)?
\begin{enumerate}[label=\textbf{\Alph*)}]
\begin{multicols}{2}
  \item {1,2,3,4,5,6}
  \item B
  \item \emptyset
  \item U
  \item \emptyset
  \item Ninguna
\end{multicols}
\end{enumerate}
\vspace{0.15cm}
\item Si \(A = \{a,b,c\}\) y \(B = \{b,c,d\}\), ¿cuál es \(A \cap B\)?
\begin{enumerate}[label=\textbf{\Alph*)}]
\begin{multicols}{2}
  \item A
  \item {b,c}
  \item \emptyset
  \item U
  \item \emptyset
  \item Ninguna
\end{multicols}
\end{enumerate}
\vspace{0.15cm}
\item Si \(A = \{a,b,c\}\) y \(B = \{b,c,d\}\), ¿cuál es \(A \cup B\)?
\begin{enumerate}[label=\textbf{\Alph*)}]
\begin{multicols}{2}
  \item {a,b,c,d}
  \item B
  \item \emptyset
  \item U
  \item \emptyset
  \item Ninguna
\end{multicols}
\end{enumerate}
\vspace{0.15cm}
\item Si \(A = \{a,b,c\}\) y \(B = \{b,c,d\}\), ¿cuál es \(A - B\)?
\begin{enumerate}[label=\textbf{\Alph*)}]
\begin{multicols}{2}
  \item {1,2}
  \item B
  \item {a}
  \item {4,5}
  \item \emptyset
  \item Ninguna
\end{multicols}
\end{enumerate}
\vspace{0.15cm}
\item Si \(A = \{a,b,c\}\) y \(B = \{b,c,d\}\), ¿cuál es \(B - A\)?
\begin{enumerate}[label=\textbf{\Alph*)}]
\begin{multicols}{2}
  \item {1,2}
  \item B
  \item {1,2}
  \item {d}
  \item \emptyset
  \item Ninguna
\end{multicols}
\end{enumerate}
\vspace{0.15cm}
\item Si \(A = \{1,2,3,4,5,6\}\) y \(B = \{2,4,6\}\), ¿cuál es \(A - B\)?
\begin{enumerate}[label=\textbf{\Alph*)}]
\begin{multicols}{2}
  \item {1,3,5}
  \item B
  \item {1,2}
  \item {4,5}
  \item \emptyset
  \item Ninguna
\end{multicols}
\end{enumerate}
\vspace{0.15cm}
\item Si \(A = \{1,2,3,4,5,6\}\) y \(B = \{2,4,6\}\), ¿cuál es \(B - A\)?
\begin{enumerate}[label=\textbf{\Alph*)}]
\begin{multicols}{2}
  \item {1,2}
  \item B
  \item {1,2}
  \item \emptyset
  \item \emptyset
  \item Ninguna
\end{multicols}
\end{enumerate}
\vspace{0.15cm}
\item ¿Cuál es la diferencia simétrica de \(A = \{1,2,3\}\) y \(B = \{2,3,4\}\)?
\begin{enumerate}[label=\textbf{\Alph*)}]
\begin{multicols}{2}
  \item {1,4}
  \item B
  \item \emptyset
  \item U
  \item \emptyset
  \item Ninguna
\end{multicols}
\end{enumerate}
\vspace{0.15cm}
\item La diferencia simétrica \(A \triangle B\) se define como:
\begin{enumerate}[label=\textbf{\Alph*)}]
\begin{multicols}{2}
  \item (A - B) \cup (B - A)
  \item B
  \item \emptyset
  \item U
  \item \emptyset
  \item Ninguna
\end{multicols}
\end{enumerate}
\vspace{0.15cm}
\item Si \(A = \{1,2,3,4\}\) y \(B = \{3,4,5,6\}\), ¿cuál es \(A \triangle B\)?
\begin{enumerate}[label=\textbf{\Alph*)}]
\begin{multicols}{2}
  \item {1,2,5,6}
  \item B
  \item \emptyset
  \item U
  \item \emptyset
  \item Ninguna
\end{multicols}
\end{enumerate}
\vspace{0.15cm}
\item Si \(A = \{1,2\}\) y \(B = \{1,2\}\), ¿cuál es \(A \triangle B\)?
\begin{enumerate}[label=\textbf{\Alph*)}]
\begin{multicols}{2}
  \item A
  \item \emptyset
  \item \emptyset
  \item U
  \item \emptyset
  \item Ninguna
\end{multicols}
\end{enumerate}
\vspace{0.15cm}
\item Si \(A = \{1,2,3\}\) y \(B = \{4,5,6\}\), ¿cuál es \(A \triangle B\)?
\begin{enumerate}[label=\textbf{\Alph*)}]
\begin{multicols}{2}
  \item {1,2,3,4,5,6}
  \item B
  \item \emptyset
  \item U
  \item \emptyset
  \item Ninguna
\end{multicols}
\end{enumerate}
\vspace{0.15cm}
\item ¿Cuál es el complemento de \(A = \{1,3,5,7,9\}\) en \(U = \{1,2,\dots,10\}\)?
\begin{enumerate}[label=\textbf{\Alph*)}]
\begin{multicols}{2}
  \item {2,4,6,8,10}
  \item B
  \item \emptyset
  \item U
  \item \emptyset
  \item Ninguna
\end{multicols}
\end{enumerate}
\vspace{0.15cm}
\item ¿Cuál es el complemento de \(A = \{2,4,6,8,10\}\) en \(U = \{1,2,\dots,10\}\)?
\begin{enumerate}[label=\textbf{\Alph*)}]
\begin{multicols}{2}
  \item {1,3,5,7,9}
  \item B
  \item \emptyset
  \item U
  \item \emptyset
  \item Ninguna
\end{multicols}
\end{enumerate}
\vspace{0.15cm}
\item Si \(A = \{1,2,3\}\) y \(U = \{1,2,3,4,5\}\), ¿cuál es \(\bar{A}\)?
\begin{enumerate}[label=\textbf{\Alph*)}]
\begin{multicols}{2}
  \item {4,5}
  \item B
  \item \emptyset
  \item U
  \item \emptyset
  \item Ninguna
\end{multicols}
\end{enumerate}
\vspace{0.15cm}
\item Si \(A = \{1,2,3,4\}\) y \(U = \{1,2,3,4\}\), ¿cuál es \(\bar{A}\)?
\begin{enumerate}[label=\textbf{\Alph*)}]
\begin{multicols}{2}
  \item A
  \item \emptyset
  \item \emptyset
  \item U
  \item \emptyset
  \item Ninguna
\end{multicols}
\end{enumerate}
\vspace{0.15cm}
\item La diferencia \(A - B\) también se puede expresar como:
\begin{enumerate}[label=\textbf{\Alph*)}]
\begin{multicols}{2}
  \item A \cap \bar{B}
  \item B
  \item {1,2}
  \item {4,5}
  \item \emptyset
  \item Ninguna
\end{multicols}
\end{enumerate}
\vspace{0.15cm}
\end{enumerate}\n
\subsubsection*{SECCIÓN IV: INTERVALOS (161–220)}

\begin{enumerate}[label=\textbf{\arabic*.}]
\setcounter{enumi}{160}
\item El intervalo \((a,b)\) es:
\begin{enumerate}[label=\textbf{\Alph*)}]
\begin{multicols}{2}
  \item Abierto
  \item Cerrado
  \item Semiabierto
  \item Infinito
  \item Vacío
  \item Unitario
\end{multicols}
\end{enumerate}
\vspace{0.15cm}
\item El intervalo \([a,b]\) es:
\begin{enumerate}[label=\textbf{\Alph*)}]
\begin{multicols}{2}
  \item Abierto
  \item Cerrado
  \item Semiabierto
  \item Infinito
  \item Vacío
  \item Unitario
\end{multicols}
\end{enumerate}
\vspace{0.15cm}
\item El intervalo \([a,b)\) es:
\begin{enumerate}[label=\textbf{\Alph*)}]
\begin{multicols}{2}
  \item Abierto
  \item Cerrado
  \item Semiabierto
  \item Infinito
  \item Vacío
  \item Unitario
\end{multicols}
\end{enumerate}
\vspace{0.15cm}
\item El intervalo \((a,b]\) es:
\begin{enumerate}[label=\textbf{\Alph*)}]
\begin{multicols}{2}
  \item Abierto
  \item Cerrado
  \item Semiabierto
  \item Infinito
  \item Vacío
  \item Unitario
\end{multicols}
\end{enumerate}
\vspace{0.15cm}
\item El intervalo \((a,\infty)\) es:
\begin{enumerate}[label=\textbf{\Alph*)}]
\begin{multicols}{2}
  \item Abierto
  \item Cerrado
  \item Semiabierto
  \item Infinito
  \item Vacío
  \item Unitario
\end{multicols}
\end{enumerate}
\vspace{0.15cm}
\item El intervalo \([a,\infty)\) es:
\begin{enumerate}[label=\textbf{\Alph*)}]
\begin{multicols}{2}
  \item Abierto
  \item Cerrado
  \item Semiabierto
  \item Infinito
  \item Vacío
  \item Unitario
\end{multicols}
\end{enumerate}
\vspace{0.15cm}
\item El intervalo \((-\infty,b)\) es:
\begin{enumerate}[label=\textbf{\Alph*)}]
\begin{multicols}{2}
  \item Abierto
  \item Cerrado
  \item Semiabierto
  \item Infinito
  \item Vacío
  \item Unitario
\end{multicols}
\end{enumerate}
\vspace{0.15cm}
\item El intervalo \((-\infty,b]\) es:
\begin{enumerate}[label=\textbf{\Alph*)}]
\begin{multicols}{2}
  \item Abierto
  \item Cerrado
  \item Semiabierto
  \item Infinito
  \item Vacío
  \item Unitario
\end{multicols}
\end{enumerate}
\vspace{0.15cm}
\item (-\infty,\infty)
\begin{enumerate}[label=\textbf{\Alph*)}]
\begin{multicols}{2}
  \item Todos los reales
  \item (-2,2)
  \item \emptyset
  \item x <= 4
  \item \emptyset
  \item Todos los reales
\end{multicols}
\end{enumerate}
\vspace{0.15cm}
\item La desigualdad \(a < x < b\) corresponde al intervalo:
\begin{enumerate}[label=\textbf{\Alph*)}]
\begin{multicols}{2}
  \item (a,b)
  \item [a,b]
  \item [a,b)
  \item (a,b]
  \item (-\infty,b)
  \item (a,\infty)
\end{multicols}
\end{enumerate}
\vspace{0.15cm}
\item La desigualdad \(a \leq x \leq b\) corresponde al intervalo:
\begin{enumerate}[label=\textbf{\Alph*)}]
\begin{multicols}{2}
  \item (a,b)
  \item [a,b]
  \item [a,b)
  \item (a,b]
  \item (-\infty,b)
  \item (a,\infty)
\end{multicols}
\end{enumerate}
\vspace{0.15cm}
\item La desigualdad \(a \leq x < b\) corresponde al intervalo:
\begin{enumerate}[label=\textbf{\Alph*)}]
\begin{multicols}{2}
  \item (a,b)
  \item [a,b]
  \item [a,b)
  \item (a,b]
  \item (-\infty,b)
  \item (a,\infty)
\end{multicols}
\end{enumerate}
\vspace{0.15cm}
\item La desigualdad \(a < x \leq b\) corresponde al intervalo:
\begin{enumerate}[label=\textbf{\Alph*)}]
\begin{multicols}{2}
  \item (a,b)
  \item [a,b]
  \item [a,b)
  \item (a,b]
  \item (-\infty,b)
  \item (a,\infty)
\end{multicols}
\end{enumerate}
\vspace{0.15cm}
\item La desigualdad \(x > a\) corresponde al intervalo:
\begin{enumerate}[label=\textbf{\Alph*)}]
\begin{multicols}{2}
  \item (a,\infty)
  \item [a,\infty)
  \item (-\infty,a)
  \item (-\infty,a]
  \item (a,b)
  \item [a,b]
\end{multicols}
\end{enumerate}
\vspace{0.15cm}
\item La desigualdad \(x \geq a\) corresponde al intervalo:
\begin{enumerate}[label=\textbf{\Alph*)}]
\begin{multicols}{2}
  \item (a,\infty)
  \item [a,\infty)
  \item (-\infty,a)
  \item (-\infty,a]
  \item (a,b)
  \item [a,b]
\end{multicols}
\end{enumerate}
\vspace{0.15cm}
\item La desigualdad \(x < b\) corresponde al intervalo:
\begin{enumerate}[label=\textbf{\Alph*)}]
\begin{multicols}{2}
  \item (b,\infty)
  \item [b,\infty)
  \item (-\infty,b)
  \item (-\infty,b]
  \item (a,b)
  \item [a,b]
\end{multicols}
\end{enumerate}
\vspace{0.15cm}
\item La desigualdad \(x \leq b\) corresponde al intervalo:
\begin{enumerate}[label=\textbf{\Alph*)}]
\begin{multicols}{2}
  \item (b,\infty)
  \item [b,\infty)
  \item (-\infty,b)
  \item (-\infty,b]
  \item (a,b)
  \item [a,b]
\end{multicols}
\end{enumerate}
\vspace{0.15cm}
\item (2,5)
\begin{enumerate}[label=\textbf{\Alph*)}]
\begin{multicols}{2}
  \item Todos entre 2 y 5 sin incluirlos
  \item (-2,2)
  \item \emptyset
  \item x <= 4
  \item \emptyset
  \item Todos los reales
\end{multicols}
\end{enumerate}
\vspace{0.15cm}
\item [2,5]
\begin{enumerate}[label=\textbf{\Alph*)}]
\begin{multicols}{2}
  \item (-1,1)
  \item (-2,2)
  \item Todos entre 2 y 5 incluyéndolos
  \item x <= 4
  \item \emptyset
  \item Todos los reales
\end{multicols}
\end{enumerate}
\vspace{0.15cm}
\item [2,5)
\begin{enumerate}[label=\textbf{\Alph*)}]
\begin{multicols}{2}
  \item Incluye 2, no 5
  \item (-2,2)
  \item \emptyset
  \item x <= 4
  \item \emptyset
  \item Todos los reales
\end{multicols}
\end{enumerate}
\vspace{0.15cm}
\item (2,5]
\begin{enumerate}[label=\textbf{\Alph*)}]
\begin{multicols}{2}
  \item (-1,1)
  \item Incluye 5, no 2
  \item \emptyset
  \item x <= 4
  \item \emptyset
  \item Todos los reales
\end{multicols}
\end{enumerate}
\vspace{0.15cm}
\item (3,7)
\begin{enumerate}[label=\textbf{\Alph*)}]
\begin{multicols}{2}
  \item 3 < x < 7
  \item (-2,2)
  \item \emptyset
  \item x <= 4
  \item \emptyset
  \item Todos los reales
\end{multicols}
\end{enumerate}
\vspace{0.15cm}
\item [3,7]
\begin{enumerate}[label=\textbf{\Alph*)}]
\begin{multicols}{2}
  \item (-1,1)
  \item 3 <= x <= 7
  \item \emptyset
  \item x <= 4
  \item \emptyset
  \item Todos los reales
\end{multicols}
\end{enumerate}
\vspace{0.15cm}
\item [3,7)
\begin{enumerate}[label=\textbf{\Alph*)}]
\begin{multicols}{2}
  \item (-1,1)
  \item (-2,2)
  \item 3 <= x < 7
  \item x <= 4
  \item \emptyset
  \item Todos los reales
\end{multicols}
\end{enumerate}
\vspace{0.15cm}
\item (3,7]
\begin{enumerate}[label=\textbf{\Alph*)}]
\begin{multicols}{2}
  \item (-1,1)
  \item (-2,2)
  \item \emptyset
  \item 3 < x <= 7
  \item \emptyset
  \item Todos los reales
\end{multicols}
\end{enumerate}
\vspace{0.15cm}
\item (-\infty,4)
\begin{enumerate}[label=\textbf{\Alph*)}]
\begin{multicols}{2}
  \item x < 4
  \item (-2,2)
  \item \emptyset
  \item x <= 4
  \item \emptyset
  \item Todos los reales
\end{multicols}
\end{enumerate}
\vspace{0.15cm}
\item (-\infty,4]
\begin{enumerate}[label=\textbf{\Alph*)}]
\begin{multicols}{2}
  \item (-1,1)
  \item (-2,2)
  \item x <= 4
  \item x <= 4
  \item \emptyset
  \item Todos los reales
\end{multicols}
\end{enumerate}
\vspace{0.15cm}
\item (4,\infty)
\begin{enumerate}[label=\textbf{\Alph*)}]
\begin{multicols}{2}
  \item (-1,1)
  \item x > 4
  \item \emptyset
  \item x <= 4
  \item \emptyset
  \item Todos los reales
\end{multicols}
\end{enumerate}
\vspace{0.15cm}
\item [4,\infty)
\begin{enumerate}[label=\textbf{\Alph*)}]
\begin{multicols}{2}
  \item (-1,1)
  \item (-2,2)
  \item \emptyset
  \item x >= 4
  \item \emptyset
  \item Todos los reales
\end{multicols}
\end{enumerate}
\vspace{0.15cm}
\item (-\infty,\infty)
\begin{enumerate}[label=\textbf{\Alph*)}]
\begin{multicols}{2}
  \item (-1,1)
  \item (-2,2)
  \item \emptyset
  \item Todos los números reales
  \item \emptyset
  \item Todos los reales
\end{multicols}
\end{enumerate}
\vspace{0.15cm}
\item ¿Cuál es la intersección (1,5) y (3,7)?
\begin{enumerate}[label=\textbf{\Alph*)}]
\begin{multicols}{2}
  \item (1,7)
  \item (3,5)
  \item \emptyset
  \item x <= 4
  \item \emptyset
  \item Todos los reales
\end{multicols}
\end{enumerate}
\vspace{0.15cm}
\item ¿Cuál es la unión (1,5) y (3,7)?
\begin{enumerate}[label=\textbf{\Alph*)}]
\begin{multicols}{2}
  \item (1,7)
  \item (3,5)
  \item \emptyset
  \item x <= 4
  \item \emptyset
  \item Todos los reales
\end{multicols}
\end{enumerate}
\vspace{0.15cm}
\item ¿Cuál es la intersección [1,5] y [3,7]?
\begin{enumerate}[label=\textbf{\Alph*)}]
\begin{multicols}{2}
  \item (-1,1)
  \item [3,5]
  \item \emptyset
  \item x <= 4
  \item \emptyset
  \item Todos los reales
\end{multicols}
\end{enumerate}
\vspace{0.15cm}
\item ¿Cuál es la unión [1,5] y [3,7]?
\begin{enumerate}[label=\textbf{\Alph*)}]
\begin{multicols}{2}
  \item [1,7]
  \item (-2,2)
  \item \emptyset
  \item x <= 4
  \item \emptyset
  \item Todos los reales
\end{multicols}
\end{enumerate}
\vspace{0.15cm}
\item ¿Cuál es la intersección (-\infty,3) y (1,\infty)?
\begin{enumerate}[label=\textbf{\Alph*)}]
\begin{multicols}{2}
  \item (1,3)
  \item (-2,2)
  \item \emptyset
  \item x <= 4
  \item \emptyset
  \item Todos los reales
\end{multicols}
\end{enumerate}
\vspace{0.15cm}
\item ¿Cuál es la unión (-\infty,3) y (1,\infty)?
\begin{enumerate}[label=\textbf{\Alph*)}]
\begin{multicols}{2}
  \item (-1,1)
  \item (-\infty,\infty)
  \item \emptyset
  \item x <= 4
  \item \emptyset
  \item Todos los reales
\end{multicols}
\end{enumerate}
\vspace{0.15cm}
\item ¿Cuál es la intersección [1,3] y [5,7]?
\begin{enumerate}[label=\textbf{\Alph*)}]
\begin{multicols}{2}
  \item (-1,1)
  \item (-2,2)
  \item \emptyset
  \item x <= 4
  \item \emptyset
  \item Todos los reales
\end{multicols}
\end{enumerate}
\vspace{0.15cm}
\item ¿Cuál es la unión [1,3] y [5,7]?
\begin{enumerate}[label=\textbf{\Alph*)}]
\begin{multicols}{2}
  \item (-1,1)
  \item [1,3] \cup [5,7]
  \item \emptyset
  \item x <= 4
  \item \emptyset
  \item Todos los reales
\end{multicols}
\end{enumerate}
\vspace{0.15cm}
\item ¿Cuál es la intersección (-\infty,0) y (0,\infty)?
\begin{enumerate}[label=\textbf{\Alph*)}]
\begin{multicols}{2}
  \item (-1,1)
  \item (-2,2)
  \item \emptyset
  \item x <= 4
  \item \emptyset
  \item Todos los reales
\end{multicols}
\end{enumerate}
\vspace{0.15cm}
\item ¿Cuál es la unión (-\infty,0) y (0,\infty)?
\begin{enumerate}[label=\textbf{\Alph*)}]
\begin{multicols}{2}
  \item (-1,1)
  \item \mathbb{R}\setminus\{0\}
  \item \emptyset
  \item x <= 4
  \item \emptyset
  \item Todos los reales
\end{multicols}
\end{enumerate}
\vspace{0.15cm}
\item ¿Cuál es la intersección [-2,2] y [0,4]?
\begin{enumerate}[label=\textbf{\Alph*)}]
\begin{multicols}{2}
  \item [-2,4]
  \item [0,2]
  \item \emptyset
  \item x <= 4
  \item \emptyset
  \item Todos los reales
\end{multicols}
\end{enumerate}
\vspace{0.15cm}
\item ¿Cuál es la unión [-2,2] y [0,4]?
\begin{enumerate}[label=\textbf{\Alph*)}]
\begin{multicols}{2}
  \item [-2,4]
  \item [0,2]
  \item \emptyset
  \item x <= 4
  \item \emptyset
  \item Todos los reales
\end{multicols}
\end{enumerate}
\vspace{0.15cm}
\item ¿Cuál es la intersección (-3,1) y (1,5)?
\begin{enumerate}[label=\textbf{\Alph*)}]
\begin{multicols}{2}
  \item (-1,1)
  \item \emptyset
  \item \emptyset
  \item x <= 4
  \item \emptyset
  \item Todos los reales
\end{multicols}
\end{enumerate}
\vspace{0.15cm}
\item ¿Cuál es la unión (-3,1) y (1,5)?
\begin{enumerate}[label=\textbf{\Alph*)}]
\begin{multicols}{2}
  \item (-1,1)
  \item (-3,1) \cup (1,5)
  \item \emptyset
  \item x <= 4
  \item \emptyset
  \item Todos los reales
\end{multicols}
\end{enumerate}
\vspace{0.15cm}
\item ¿Cuál es la intersección (-\infty,5) y (2,\infty)?
\begin{enumerate}[label=\textbf{\Alph*)}]
\begin{multicols}{2}
  \item (2,5)
  \item (-2,2)
  \item \emptyset
  \item x <= 4
  \item \emptyset
  \item Todos los reales
\end{multicols}
\end{enumerate}
\vspace{0.15cm}
\item ¿Cuál es la unión (-\infty,5) y (2,\infty)?
\begin{enumerate}[label=\textbf{\Alph*)}]
\begin{multicols}{2}
  \item (-1,1)
  \item (-\infty,\infty)
  \item \emptyset
  \item x <= 4
  \item \emptyset
  \item Todos los reales
\end{multicols}
\end{enumerate}
\vspace{0.15cm}
\item ¿Cuál es la intersección [-1,4] y (0,6)?
\begin{enumerate}[label=\textbf{\Alph*)}]
\begin{multicols}{2}
  \item (-1,1)
  \item (0,4]
  \item \emptyset
  \item x <= 4
  \item \emptyset
  \item Todos los reales
\end{multicols}
\end{enumerate}
\vspace{0.15cm}
\item ¿Cuál es la unión [-1,4] y (0,6)?
\begin{enumerate}[label=\textbf{\Alph*)}]
\begin{multicols}{2}
  \item [-1,6)
  \item (-2,2)
  \item \emptyset
  \item x <= 4
  \item \emptyset
  \item Todos los reales
\end{multicols}
\end{enumerate}
\vspace{0.15cm}
\item ¿Cuál es la intersección (-\infty,2] y [2,\infty)?
\begin{enumerate}[label=\textbf{\Alph*)}]
\begin{multicols}{2}
  \item {2}
  \item (-2,2)
  \item \emptyset
  \item x <= 4
  \item \emptyset
  \item Todos los reales
\end{multicols}
\end{enumerate}
\vspace{0.15cm}
\item ¿Cuál es la unión (-\infty,2] y [2,\infty)?
\begin{enumerate}[label=\textbf{\Alph*)}]
\begin{multicols}{2}
  \item (-1,1)
  \item (-2,2)
  \item (-\infty,\infty)
  \item x <= 4
  \item \emptyset
  \item Todos los reales
\end{multicols}
\end{enumerate}
\vspace{0.15cm}
\item ¿Cuál es la intersección (-1,1) y (-2,2)?
\begin{enumerate}[label=\textbf{\Alph*)}]
\begin{multicols}{2}
  \item (-1,1)
  \item (-2,2)
  \item \emptyset
  \item x <= 4
  \item \emptyset
  \item Todos los reales
\end{multicols}
\end{enumerate}
\vspace{0.15cm}
\item ¿Cuál es la unión (-1,1) y (-2,2)?
\begin{enumerate}[label=\textbf{\Alph*)}]
\begin{multicols}{2}
  \item (-1,1)
  \item (-2,2)
  \item \emptyset
  \item x <= 4
  \item \emptyset
  \item Todos los reales
\end{multicols}
\end{enumerate}
\vspace{0.15cm}
\item ¿Cuál es la intersección [0,3] y [1,4]?
\begin{enumerate}[label=\textbf{\Alph*)}]
\begin{multicols}{2}
  \item (-1,1)
  \item [1,3]
  \item \emptyset
  \item x <= 4
  \item \emptyset
  \item Todos los reales
\end{multicols}
\end{enumerate}
\vspace{0.15cm}
\item ¿Cuál es la unión [0,3] y [1,4]?
\begin{enumerate}[label=\textbf{\Alph*)}]
\begin{multicols}{2}
  \item [0,4]
  \item (-2,2)
  \item \emptyset
  \item x <= 4
  \item \emptyset
  \item Todos los reales
\end{multicols}
\end{enumerate}
\vspace{0.15cm}
\item ¿Cuál es la intersección (-5,-2) y (-3,0)?
\begin{enumerate}[label=\textbf{\Alph*)}]
\begin{multicols}{2}
  \item (-1,1)
  \item (-3,-2)
  \item \emptyset
  \item x <= 4
  \item \emptyset
  \item Todos los reales
\end{multicols}
\end{enumerate}
\vspace{0.15cm}
\item ¿Cuál es la unión (-5,-2) y (-3,0)?
\begin{enumerate}[label=\textbf{\Alph*)}]
\begin{multicols}{2}
  \item (-5,0)
  \item (-2,2)
  \item \emptyset
  \item x <= 4
  \item \emptyset
  \item Todos los reales
\end{multicols}
\end{enumerate}
\vspace{0.15cm}
\item ¿Cuál es la intersección (-\infty,0] y [0,\infty)?
\begin{enumerate}[label=\textbf{\Alph*)}]
\begin{multicols}{2}
  \item {0}
  \item (-2,2)
  \item \emptyset
  \item x <= 4
  \item \emptyset
  \item Todos los reales
\end{multicols}
\end{enumerate}
\vspace{0.15cm}
\item ¿Cuál es la unión (-\infty,0] y [0,\infty)?
\begin{enumerate}[label=\textbf{\Alph*)}]
\begin{multicols}{2}
  \item (-1,1)
  \item (-2,2)
  \item (-\infty,\infty)
  \item x <= 4
  \item \emptyset
  \item Todos los reales
\end{multicols}
\end{enumerate}
\vspace{0.15cm}
\item ¿Cuál es la intersección (1,4) y (2,3)?
\begin{enumerate}[label=\textbf{\Alph*)}]
\begin{multicols}{2}
  \item (-1,1)
  \item (2,3)
  \item \emptyset
  \item x <= 4
  \item \emptyset
  \item Todos los reales
\end{multicols}
\end{enumerate}
\vspace{0.15cm}
\item ¿Cuál es la unión (1,4) y (2,3)?
\begin{enumerate}[label=\textbf{\Alph*)}]
\begin{multicols}{2}
  \item (1,4)
  \item (-2,2)
  \item \emptyset
  \item x <= 4
  \item \emptyset
  \item Todos los reales
\end{multicols}
\end{enumerate}
\vspace{0.15cm}
\end{enumerate}\n
\subsubsection*{SECCIÓN V: PROBLEMAS DE APLICACIÓN (221–300)}

\begin{enumerate}[label=\textbf{\arabic*.}]
\setcounter{enumi}{220}
\item Si \(A = \{x \in \mathbb{N} \mid x < 6\}\), ¿cuál es A por extensión?
\begin{enumerate}[label=\textbf{\Alph*)}]
\begin{multicols}{2}
  \item {1,2,3,4,5}
  \item 1
  \item {-3,3}
  \item \emptyset
  \item \emptyset
  \item Todos los reales
\end{multicols}
\end{enumerate}
\vspace{0.15cm}
\item Si \(A = \{x \in \mathbb{Z} \mid -3 < x < 3\}\), ¿cuál es A por extensión?
\begin{enumerate}[label=\textbf{\Alph*)}]
\begin{multicols}{2}
  \item {-2,-1,0,1,2}
  \item 1
  \item {-3,3}
  \item \emptyset
  \item \emptyset
  \item Todos los reales
\end{multicols}
\end{enumerate}
\vspace{0.15cm}
\item Si \(A = \{x \in \mathbb{N} \mid x \text{ es par}, x < 10\}\), ¿cuál es A?
\begin{enumerate}[label=\textbf{\Alph*)}]
\begin{multicols}{2}
  \item {2,4,6,8}
  \item 1
  \item {-3,3}
  \item \emptyset
  \item \emptyset
  \item Todos los reales
\end{multicols}
\end{enumerate}
\vspace{0.15cm}
\item Si \(A = \{x \in \mathbb{N} \mid x \text{ es impar}, x < 10\}\), ¿cuál es A?
\begin{enumerate}[label=\textbf{\Alph*)}]
\begin{multicols}{2}
  \item {1,3,5,7,9}
  \item 1
  \item {-3,3}
  \item \emptyset
  \item \emptyset
  \item Todos los reales
\end{multicols}
\end{enumerate}
\vspace{0.15cm}
\item Si \(A = \{x \in \mathbb{Z} \mid x^2 < 10\}\), ¿cuál es A?
\begin{enumerate}[label=\textbf{\Alph*)}]
\begin{multicols}{2}
  \item {-3,-2,-1,0,1,2,3}
  \item 1
  \item {-3,3}
  \item \emptyset
  \item \emptyset
  \item Todos los reales
\end{multicols}
\end{enumerate}
\vspace{0.15cm}
\item Si \(A = \{x \in \mathbb{N} \mid x^2 < 20\}\), ¿cuál es A?
\begin{enumerate}[label=\textbf{\Alph*)}]
\begin{multicols}{2}
  \item {1,2,3,4}
  \item 1
  \item {-3,3}
  \item \emptyset
  \item \emptyset
  \item Todos los reales
\end{multicols}
\end{enumerate}
\vspace{0.15cm}
\item Si \(A = \{x \in \mathbb{R} \mid x^2 = 9\}\), ¿cuál es A?
\begin{enumerate}[label=\textbf{\Alph*)}]
\begin{multicols}{2}
  \item {1,2,3,4,5}
  \item 1
  \item {-3,3}
  \item \emptyset
  \item \emptyset
  \item Todos los reales
\end{multicols}
\end{enumerate}
\vspace{0.15cm}
\item Si \(A = \{x \in \mathbb{R} \mid x^2 = 16\}\), ¿cuál es A?
\begin{enumerate}[label=\textbf{\Alph*)}]
\begin{multicols}{2}
  \item {1,2,3,4,5}
  \item 1
  \item {-4,4}
  \item \emptyset
  \item \emptyset
  \item Todos los reales
\end{multicols}
\end{enumerate}
\vspace{0.15cm}
\item Si \(A = \{x \in \mathbb{R} \mid x^2 = -4\}\), ¿cuál es A?
\begin{enumerate}[label=\textbf{\Alph*)}]
\begin{multicols}{2}
  \item 0
  \item 1
  \item 2
  \item \emptyset
  \item \emptyset
  \item Todos los reales
\end{multicols}
\end{enumerate}
\vspace{0.15cm}
\item Si \(A = \{x \in \mathbb{N} \mid 2x - 3 < 7\}\), ¿cuál es A?
\begin{enumerate}[label=\textbf{\Alph*)}]
\begin{multicols}{2}
  \item {1,2,3,4}
  \item 1
  \item {-3,3}
  \item \emptyset
  \item \emptyset
  \item Todos los reales
\end{multicols}
\end{enumerate}
\vspace{0.15cm}
\item Si \(A = \{x \in \mathbb{Z} \mid 3x + 1 > -5\}\), ¿cuál es el menor elemento de A?
\begin{enumerate}[label=\textbf{\Alph*)}]
\begin{multicols}{2}
  \item 0
  \item -1
  \item 2
  \item -1
  \item \emptyset
  \item Todos los reales
\end{multicols}
\end{enumerate}
\vspace{0.15cm}
\item Si \(A = \{x \in \mathbb{N} \mid x^2 - 4 = 0\}\), ¿cuál es A?
\begin{enumerate}[label=\textbf{\Alph*)}]
\begin{multicols}{2}
  \item {2}
  \item 1
  \item {-3,3}
  \item \emptyset
  \item \emptyset
  \item Todos los reales
\end{multicols}
\end{enumerate}
\vspace{0.15cm}
\item Si \(A = \{x \in \mathbb{Z} \mid |x| < 4\}\), ¿cuál es A?
\begin{enumerate}[label=\textbf{\Alph*)}]
\begin{multicols}{2}
  \item {-3,-2,-1,0,1,2,3}
  \item 1
  \item {-3,3}
  \item \emptyset
  \item \emptyset
  \item Todos los reales
\end{multicols}
\end{enumerate}
\vspace{0.15cm}
\item Si \(A = \{x \in \mathbb{Z} \mid |x| \leq 3\}\), ¿cuál es A?
\begin{enumerate}[label=\textbf{\Alph*)}]
\begin{multicols}{2}
  \item {-3,-2,-1,0,1,2,3}
  \item 1
  \item {-3,3}
  \item \emptyset
  \item \emptyset
  \item Todos los reales
\end{multicols}
\end{enumerate}
\vspace{0.15cm}
\item Si \(A = \{x \in \mathbb{N} \mid |x| < 5\}\), ¿cuál es A?
\begin{enumerate}[label=\textbf{\Alph*)}]
\begin{multicols}{2}
  \item {1,2,3,4}
  \item 1
  \item {-3,3}
  \item \emptyset
  \item \emptyset
  \item Todos los reales
\end{multicols}
\end{enumerate}
\vspace{0.15cm}
\item Si \(A = \{x \in \mathbb{R} \mid |x| < 3\}\), ¿cuál es A?
\begin{enumerate}[label=\textbf{\Alph*)}]
\begin{multicols}{2}
  \item (-3,3)
  \item 1
  \item 2
  \item -1
  \item \emptyset
  \item Todos los reales
\end{multicols}
\end{enumerate}
\vspace{0.15cm}
\item Si \(A = \{x \in \mathbb{R} \mid |x| \leq 3\}\), ¿cuál es A?
\begin{enumerate}[label=\textbf{\Alph*)}]
\begin{multicols}{2}
  \item 0
  \item [-3,3]
  \item 2
  \item -1
  \item \emptyset
  \item Todos los reales
\end{multicols}
\end{enumerate}
\vspace{0.15cm}
\item Si \(A = \{x \in \mathbb{R} \mid |x| > 3\}\), ¿cuál es A?
\begin{enumerate}[label=\textbf{\Alph*)}]
\begin{multicols}{2}
  \item (1,3)
  \item (-\infty,-3) \cup (3,\infty)
  \item 2
  \item -1
  \item \emptyset
  \item Todos los reales
\end{multicols}
\end{enumerate}
\vspace{0.15cm}
\item Si \(A = \{x \in \mathbb{R} \mid |x| \geq 3\}\), ¿cuál es A?
\begin{enumerate}[label=\textbf{\Alph*)}]
\begin{multicols}{2}
  \item (1,3)
  \item [-3,3]
  \item 2
  \item (-\infty,-3] \cup [3,\infty)
  \item \emptyset
  \item Todos los reales
\end{multicols}
\end{enumerate}
\vspace{0.15cm}
\item Si \(A = \{x \in \mathbb{R} \mid 1 < x < 4\}\), ¿cuál es A?
\begin{enumerate}[label=\textbf{\Alph*)}]
\begin{multicols}{2}
  \item (1,4)
  \item 1
  \item 2
  \item -1
  \item \emptyset
  \item Todos los reales
\end{multicols}
\end{enumerate}
\vspace{0.15cm}
\item Si \(A = \{x \in \mathbb{R} \mid 1 \leq x \leq 4\}\), ¿cuál es A?
\begin{enumerate}[label=\textbf{\Alph*)}]
\begin{multicols}{2}
  \item 0
  \item [1,4]
  \item 2
  \item -1
  \item \emptyset
  \item Todos los reales
\end{multicols}
\end{enumerate}
\vspace{0.15cm}
\item Si \(A = \{x \in \mathbb{R} \mid 1 \leq x < 4\}\), ¿cuál es A?
\begin{enumerate}[label=\textbf{\Alph*)}]
\begin{multicols}{2}
  \item 0
  \item [-3,3]
  \item [1,4)
  \item -1
  \item \emptyset
  \item Todos los reales
\end{multicols}
\end{enumerate}
\vspace{0.15cm}
\item Si \(A = \{x \in \mathbb{R} \mid 1 < x \leq 4\}\), ¿cuál es A?
\begin{enumerate}[label=\textbf{\Alph*)}]
\begin{multicols}{2}
  \item (1,3)
  \item 1
  \item 2
  \item (1,4]
  \item \emptyset
  \item Todos los reales
\end{multicols}
\end{enumerate}
\vspace{0.15cm}
\item Si \(A = \{x \in \mathbb{R} \mid x > 2\}\), ¿cuál es A?
\begin{enumerate}[label=\textbf{\Alph*)}]
\begin{multicols}{2}
  \item (2,\infty)
  \item 1
  \item 2
  \item -1
  \item \emptyset
  \item Todos los reales
\end{multicols}
\end{enumerate}
\vspace{0.15cm}
\item Si \(A = \{x \in \mathbb{R} \mid x \geq 2\}\), ¿cuál es A?
\begin{enumerate}[label=\textbf{\Alph*)}]
\begin{multicols}{2}
  \item 0
  \item [2,\infty)
  \item 2
  \item -1
  \item \emptyset
  \item Todos los reales
\end{multicols}
\end{enumerate}
\vspace{0.15cm}
\item Si \(A = \{x \in \mathbb{R} \mid x < 2\}\), ¿cuál es A?
\begin{enumerate}[label=\textbf{\Alph*)}]
\begin{multicols}{2}
  \item (1,3)
  \item 1
  \item (-\infty,2)
  \item -1
  \item \emptyset
  \item Todos los reales
\end{multicols}
\end{enumerate}
\vspace{0.15cm}
\item Si \(A = \{x \in \mathbb{R} \mid x \leq 2\}\), ¿cuál es A?
\begin{enumerate}[label=\textbf{\Alph*)}]
\begin{multicols}{2}
  \item (1,3)
  \item 1
  \item 2
  \item (-\infty,2]
  \item \emptyset
  \item Todos los reales
\end{multicols}
\end{enumerate}
\vspace{0.15cm}
\item Si \(A = (-1,3)\) y \(B = (1,5)\), ¿cuál es \(A \cap B\)?
\begin{enumerate}[label=\textbf{\Alph*)}]
\begin{multicols}{2}
  \item (1,3)
  \item 1
  \item 2
  \item -1
  \item \emptyset
  \item Todos los reales
\end{multicols}
\end{enumerate}
\vspace{0.15cm}
\item Si \(A = (-1,3)\) y \(B = (1,5)\), ¿cuál es \(A \cup B\)?
\begin{enumerate}[label=\textbf{\Alph*)}]
\begin{multicols}{2}
  \item (1,3)
  \item (-1,5)
  \item 2
  \item -1
  \item \emptyset
  \item Todos los reales
\end{multicols}
\end{enumerate}
\vspace{0.15cm}
\item Si \(A = (0,\infty)\) y \(B = (-\infty,5)\), ¿cuál es \(A \cap B\)?
\begin{enumerate}[label=\textbf{\Alph*)}]
\begin{multicols}{2}
  \item (0,5)
  \item 1
  \item 2
  \item -1
  \item \emptyset
  \item Todos los reales
\end{multicols}
\end{enumerate}
\vspace{0.15cm}
\item Si \(A = (0,\infty)\) y \(B = (-\infty,5)\), ¿cuál es \(A \cup B\)?
\begin{enumerate}[label=\textbf{\Alph*)}]
\begin{multicols}{2}
  \item (1,3)
  \item (-\infty,\infty)
  \item 2
  \item -1
  \item \emptyset
  \item Todos los reales
\end{multicols}
\end{enumerate}
\vspace{0.15cm}
\item Si \(A = (-2,2)\) y \(B = (-3,3)\), ¿cuál es \(A \cap B\)?
\begin{enumerate}[label=\textbf{\Alph*)}]
\begin{multicols}{2}
  \item (-2,2)
  \item 1
  \item 2
  \item -1
  \item \emptyset
  \item Todos los reales
\end{multicols}
\end{enumerate}
\vspace{0.15cm}
\item Si \(A = (-2,2)\) y \(B = (-3,3)\), ¿cuál es \(A \cup B\)?
\begin{enumerate}[label=\textbf{\Alph*)}]
\begin{multicols}{2}
  \item (1,3)
  \item (-3,3)
  \item 2
  \item -1
  \item \emptyset
  \item Todos los reales
\end{multicols}
\end{enumerate}
\vspace{0.15cm}
\item Si \(A = [-2,2]\) y \(B = (-\infty,-1] \cup [1,\infty)\), ¿cuál es \(A \cap B\)?
\begin{enumerate}[label=\textbf{\Alph*)}]
\begin{multicols}{2}
  \item [-2,-1] \cup [1,2]
  \item [-3,3]
  \item 2
  \item -1
  \item \emptyset
  \item Todos los reales
\end{multicols}
\end{enumerate}
\vspace{0.15cm}
\item Si \(A = [-2,2]\) y \(B = (-\infty,-1] \cup [1,\infty)\), ¿cuál es \(A \cup B\)?
\begin{enumerate}[label=\textbf{\Alph*)}]
\begin{multicols}{2}
  \item (1,3)
  \item 1
  \item (-\infty,\infty)
  \item -1
  \item \emptyset
  \item Todos los reales
\end{multicols}
\end{enumerate}
\vspace{0.15cm}
\item Si \(A = \{x \in \mathbb{N} \mid x < 10\}\) y \(B = \{x \in \mathbb{N} \mid x > 5\}\), ¿cuál es \(A \cap B\)?
\begin{enumerate}[label=\textbf{\Alph*)}]
\begin{multicols}{2}
  \item {6,7,8,9}
  \item 1
  \item {-3,3}
  \item \emptyset
  \item \emptyset
  \item Todos los reales
\end{multicols}
\end{enumerate}
\vspace{0.15cm}
\item Si \(A = \{x \in \mathbb{N} \mid x < 10\}\) y \(B = \{x \in \mathbb{N} \mid x > 5\}\), ¿cuál es \(A \cup B\)?
\begin{enumerate}[label=\textbf{\Alph*)}]
\begin{multicols}{2}
  \item {1,2,3,4,5}
  \item \{1,2,3,\ldots\}
  \item {-3,3}
  \item \emptyset
  \item \emptyset
  \item Todos los reales
\end{multicols}
\end{enumerate}
\vspace{0.15cm}
\item Si \(A = \{-3,\dots,3\}\) y \(B = \{0,\dots,5\}\), ¿cuál es \(A \cap B\)?
\begin{enumerate}[label=\textbf{\Alph*)}]
\begin{multicols}{2}
  \item {0,1,2,3}
  \item 1
  \item {-3,3}
  \item \emptyset
  \item \emptyset
  \item Todos los reales
\end{multicols}
\end{enumerate}
\vspace{0.15cm}
\item Si \(A = \{-3,\dots,3\}\) y \(B = \{0,\dots,5\}\), ¿cuál es \(A \cup B\)?
\begin{enumerate}[label=\textbf{\Alph*)}]
\begin{multicols}{2}
  \item {1,2,3,4,5}
  \item {-3,-2,-1,0,1,2,3,4,5}
  \item {-3,3}
  \item \emptyset
  \item \emptyset
  \item Todos los reales
\end{multicols}
\end{enumerate}
\vspace{0.15cm}
\item Si \(A = \{x \in \mathbb{R} \mid x^2 < 4\}\), ¿cuál es A?
\begin{enumerate}[label=\textbf{\Alph*)}]
\begin{multicols}{2}
  \item (-2,2)
  \item 1
  \item 2
  \item -1
  \item \emptyset
  \item Todos los reales
\end{multicols}
\end{enumerate}
\vspace{0.15cm}
\item Si \(A = \{x \in \mathbb{R} \mid x^2 \leq 4\}\), ¿cuál es A?
\begin{enumerate}[label=\textbf{\Alph*)}]
\begin{multicols}{2}
  \item 0
  \item [-2,2]
  \item 2
  \item -1
  \item \emptyset
  \item Todos los reales
\end{multicols}
\end{enumerate}
\vspace{0.15cm}
\item Si \(A = \{x \in \mathbb{R} \mid x^2 > 4\}\), ¿cuál es A?
\begin{enumerate}[label=\textbf{\Alph*)}]
\begin{multicols}{2}
  \item (1,3)
  \item (-\infty,-2) \cup (2,\infty)
  \item 2
  \item -1
  \item \emptyset
  \item Todos los reales
\end{multicols}
\end{enumerate}
\vspace{0.15cm}
\item Si \(A = \{x \in \mathbb{R} \mid x^2 \geq 4\}\), ¿cuál es A?
\begin{enumerate}[label=\textbf{\Alph*)}]
\begin{multicols}{2}
  \item (1,3)
  \item [-3,3]
  \item 2
  \item (-\infty,-2] \cup [2,\infty)
  \item \emptyset
  \item Todos los reales
\end{multicols}
\end{enumerate}
\vspace{0.15cm}
\item Si \(A = (-2,2)\) y \(B = (0,4)\), ¿cuál es A - B?
\begin{enumerate}[label=\textbf{\Alph*)}]
\begin{multicols}{2}
  \item (1,3)
  \item (-2,0]
  \item 2
  \item -1
  \item \emptyset
  \item Todos los reales
\end{multicols}
\end{enumerate}
\vspace{0.15cm}
\item Si \(A = (-2,2)\) y \(B = (0,4)\), ¿cuál es B - A?
\begin{enumerate}[label=\textbf{\Alph*)}]
\begin{multicols}{2}
  \item [2,4)
  \item [-3,3]
  \item 2
  \item -1
  \item \emptyset
  \item Todos los reales
\end{multicols}
\end{enumerate}
\vspace{0.15cm}
\item Si \(A = \{1,2,3,4,5\}\) y \(B = \{4,5,6,7\}\), ¿cuál es \(A \triangle B\)?
\begin{enumerate}[label=\textbf{\Alph*)}]
\begin{multicols}{2}
  \item {1,2,3,6,7}
  \item 1
  \item {-3,3}
  \item \emptyset
  \item \emptyset
  \item Todos los reales
\end{multicols}
\end{enumerate}
\vspace{0.15cm}
\item Si \(A = \{a,b,c,d\}\) y \(B = \{c,d,e,f\}\), ¿cuál es \(A \triangle B\)?
\begin{enumerate}[label=\textbf{\Alph*)}]
\begin{multicols}{2}
  \item {a,b,e,f}
  \item 1
  \item {-3,3}
  \item \emptyset
  \item \emptyset
  \item Todos los reales
\end{multicols}
\end{enumerate}
\vspace{0.15cm}
\item Si \(A = \{1,2,3\}\) y \(B = \{3,4,5\}\), ¿cuál es \(A \triangle B\)?
\begin{enumerate}[label=\textbf{\Alph*)}]
\begin{multicols}{2}
  \item {1,2,4,5}
  \item 1
  \item {-3,3}
  \item \emptyset
  \item \emptyset
  \item Todos los reales
\end{multicols}
\end{enumerate}
\vspace{0.15cm}
\item Si \(A = \{1,2,3,4\}\) y \(B = \{2,4,6,8\}\), ¿cuál es \(A \triangle B\)?
\begin{enumerate}[label=\textbf{\Alph*)}]
\begin{multicols}{2}
  \item {1,3,6,8}
  \item 1
  \item {-3,3}
  \item \emptyset
  \item \emptyset
  \item Todos los reales
\end{multicols}
\end{enumerate}
\vspace{0.15cm}
\item Si \(A = \{x \in \mathbb{N} \mid x \leq 5\}\) y \(B = \{x \in \mathbb{N} \mid 3 \leq x \leq 8\}\), ¿cuál es \(A \cap B\)?
\begin{enumerate}[label=\textbf{\Alph*)}]
\begin{multicols}{2}
  \item {3,4,5}
  \item 1
  \item {-3,3}
  \item \emptyset
  \item \emptyset
  \item Todos los reales
\end{multicols}
\end{enumerate}
\vspace{0.15cm}
\item Si \(A = \{x \in \mathbb{N} \mid x \leq 5\}\) y \(B = \{x \in \mathbb{N} \mid 3 \leq x \leq 8\}\), ¿cuál es \(A \cup B\)?
\begin{enumerate}[label=\textbf{\Alph*)}]
\begin{multicols}{2}
  \item {1,2,3,4,5}
  \item {1,2,3,4,5,6,7,8}
  \item {-3,3}
  \item \emptyset
  \item \emptyset
  \item Todos los reales
\end{multicols}
\end{enumerate}
\vspace{0.15cm}
\item Si \(A = \{x \in \mathbb{Z} \mid -5 < x < 0\}\) y \(B = \{x \in \mathbb{Z} \mid -2 < x < 3\}\), ¿cuál es \(A \cap B\)?
\begin{enumerate}[label=\textbf{\Alph*)}]
\begin{multicols}{2}
  \item {1,2,3,4,5}
  \item {-1}
  \item {-3,3}
  \item \emptyset
  \item \emptyset
  \item Todos los reales
\end{multicols}
\end{enumerate}
\vspace{0.15cm}
\item Si \(A = \{x \in \mathbb{Z} \mid -5 < x < 0\}\) y \(B = \{x \in \mathbb{Z} \mid -2 < x < 3\}\), ¿cuál es \(A \cup B\)?
\begin{enumerate}[label=\textbf{\Alph*)}]
\begin{multicols}{2}
  \item {1,2,3,4,5}
  \item {-4,-3,-2,-1,0,1,2}
  \item {-3,3}
  \item \emptyset
  \item \emptyset
  \item Todos los reales
\end{multicols}
\end{enumerate}
\vspace{0.15cm}
\item Si \(A = (0,1)\) y \(B = (0.5,2)\), ¿cuál es \(A \cap B\)?
\begin{enumerate}[label=\textbf{\Alph*)}]
\begin{multicols}{2}
  \item (0.5,1)
  \item 1
  \item 2
  \item -1
  \item \emptyset
  \item Todos los reales
\end{multicols}
\end{enumerate}
\vspace{0.15cm}
\item Si \(A = (0,1)\) y \(B = (0.5,2)\), ¿cuál es \(A \cup B\)?
\begin{enumerate}[label=\textbf{\Alph*)}]
\begin{multicols}{2}
  \item (1,3)
  \item (0,2)
  \item 2
  \item -1
  \item \emptyset
  \item Todos los reales
\end{multicols}
\end{enumerate}
\vspace{0.15cm}
\item Si \(A = [-1,2]\) y \(B = [1,4]\), ¿cuál es \(A \cap B\)?
\begin{enumerate}[label=\textbf{\Alph*)}]
\begin{multicols}{2}
  \item [1,2]
  \item [-3,3]
  \item 2
  \item -1
  \item \emptyset
  \item Todos los reales
\end{multicols}
\end{enumerate}
\vspace{0.15cm}
\item Si \(A = [-1,2]\) y \(B = [1,4]\), ¿cuál es \(A \cup B\)?
\begin{enumerate}[label=\textbf{\Alph*)}]
\begin{multicols}{2}
  \item 0
  \item [-1,4]
  \item 2
  \item -1
  \item \emptyset
  \item Todos los reales
\end{multicols}
\end{enumerate}
\vspace{0.15cm}
\item Si \(A = (-\infty,3)\) y \(B = (1,\infty)\), ¿cuál es \(A \cap B\)?
\begin{enumerate}[label=\textbf{\Alph*)}]
\begin{multicols}{2}
  \item (1,3)
  \item 1
  \item 2
  \item -1
  \item \emptyset
  \item Todos los reales
\end{multicols}
\end{enumerate}
\vspace{0.15cm}
\item Si \(A = (-\infty,3)\) y \(B = (1,\infty)\), ¿cuál es \(A \cup B\)?
\begin{enumerate}[label=\textbf{\Alph*)}]
\begin{multicols}{2}
  \item (1,3)
  \item (-\infty,\infty)
  \item 2
  \item -1
  \item \emptyset
  \item Todos los reales
\end{multicols}
\end{enumerate}
\vspace{0.15cm}
\item Si \(A = (-\infty,2]\) y \(B = [2,\infty)\), ¿cuál es \(A \cap B\)?
\begin{enumerate}[label=\textbf{\Alph*)}]
\begin{multicols}{2}
  \item {2}
  \item 1
  \item {-3,3}
  \item \emptyset
  \item \emptyset
  \item Todos los reales
\end{multicols}
\end{enumerate}
\vspace{0.15cm}
\item Si \(A = (-\infty,2]\) y \(B = [2,\infty)\), ¿cuál es \(A \cup B\)?
\begin{enumerate}[label=\textbf{\Alph*)}]
\begin{multicols}{2}
  \item (1,3)
  \item 1
  \item (-\infty,\infty)
  \item -1
  \item \emptyset
  \item Todos los reales
\end{multicols}
\end{enumerate}
\vspace{0.15cm}
\item Si \(A = \{x \in \mathbb{N} \mid x^2 \leq 9\}\), ¿cuál es A?
\begin{enumerate}[label=\textbf{\Alph*)}]
\begin{multicols}{2}
  \item {1,2,3}
  \item 1
  \item {-3,3}
  \item \emptyset
  \item \emptyset
  \item Todos los reales
\end{multicols}
\end{enumerate}
\vspace{0.15cm}
\item Si \(A = \{x \in \mathbb{Z} \mid x^2 \leq 9\}\), ¿cuál es A?
\begin{enumerate}[label=\textbf{\Alph*)}]
\begin{multicols}{2}
  \item {-3,-2,-1,0,1,2,3}
  \item 1
  \item {-3,3}
  \item \emptyset
  \item \emptyset
  \item Todos los reales
\end{multicols}
\end{enumerate}
\vspace{0.15cm}
\item Si \(A = \{x \in \mathbb{N} \mid 2 < x < 8\}\), ¿cuál es A?
\begin{enumerate}[label=\textbf{\Alph*)}]
\begin{multicols}{2}
  \item {3,4,5,6,7}
  \item 1
  \item {-3,3}
  \item \emptyset
  \item \emptyset
  \item Todos los reales
\end{multicols}
\end{enumerate}
\vspace{0.15cm}
\item Si \(A = \{x \in \mathbb{Z} \mid -4 < x < 2\}\), ¿cuál es A?
\begin{enumerate}[label=\textbf{\Alph*)}]
\begin{multicols}{2}
  \item {-3,-2,-1,0,1}
  \item 1
  \item {-3,3}
  \item \emptyset
  \item \emptyset
  \item Todos los reales
\end{multicols}
\end{enumerate}
\vspace{0.15cm}
\item Si \(A = (2,5)\) y \(B = (4,7)\), ¿cuál es \(A \cap B\)?
\begin{enumerate}[label=\textbf{\Alph*)}]
\begin{multicols}{2}
  \item (4,5)
  \item 1
  \item 2
  \item -1
  \item \emptyset
  \item Todos los reales
\end{multicols}
\end{enumerate}
\vspace{0.15cm}
\item Si \(A = (2,5)\) y \(B = (4,7)\), ¿cuál es \(A \cup B\)?
\begin{enumerate}[label=\textbf{\Alph*)}]
\begin{multicols}{2}
  \item (1,3)
  \item (2,7)
  \item 2
  \item -1
  \item \emptyset
  \item Todos los reales
\end{multicols}
\end{enumerate}
\vspace{0.15cm}
\item Si \(A = [-3,1)\) y \(B = (0,4]\), ¿cuál es \(A \cap B\)?
\begin{enumerate}[label=\textbf{\Alph*)}]
\begin{multicols}{2}
  \item (0,1)
  \item 1
  \item 2
  \item -1
  \item \emptyset
  \item Todos los reales
\end{multicols}
\end{enumerate}
\vspace{0.15cm}
\item Si \(A = [-3,1)\) y \(B = (0,4]\), ¿cuál es \(A \cup B\)?
\begin{enumerate}[label=\textbf{\Alph*)}]
\begin{multicols}{2}
  \item 0
  \item [-3,4]
  \item 2
  \item -1
  \item \emptyset
  \item Todos los reales
\end{multicols}
\end{enumerate}
\vspace{0.15cm}
\item Si \(A = \{x \in \mathbb{R} \mid |x-2| < 1\}\), ¿cuál es A?
\begin{enumerate}[label=\textbf{\Alph*)}]
\begin{multicols}{2}
  \item (1,3)
  \item 1
  \item 2
  \item -1
  \item \emptyset
  \item Todos los reales
\end{multicols}
\end{enumerate}
\vspace{0.15cm}
\item Si \(A = \{x \in \mathbb{R} \mid |x-2| \leq 1\}\), ¿cuál es A?
\begin{enumerate}[label=\textbf{\Alph*)}]
\begin{multicols}{2}
  \item 0
  \item [1,3]
  \item 2
  \item -1
  \item \emptyset
  \item Todos los reales
\end{multicols}
\end{enumerate}
\vspace{0.15cm}
\item Si \(A = \{x \in \mathbb{R} \mid |x+1| < 2\}\), ¿cuál es A?
\begin{enumerate}[label=\textbf{\Alph*)}]
\begin{multicols}{2}
  \item (-3,1)
  \item 1
  \item 2
  \item -1
  \item \emptyset
  \item Todos los reales
\end{multicols}
\end{enumerate}
\vspace{0.15cm}
\item Si \(A = \{x \in \mathbb{R} \mid |x+1| \leq 2\}\), ¿cuál es A?
\begin{enumerate}[label=\textbf{\Alph*)}]
\begin{multicols}{2}
  \item 0
  \item [-3,1]
  \item 2
  \item -1
  \item \emptyset
  \item Todos los reales
\end{multicols}
\end{enumerate}
\vspace{0.15cm}
\item Si \(A = \{x \in \mathbb{R} \mid |x-3| < 2\}\), ¿cuál es A?
\begin{enumerate}[label=\textbf{\Alph*)}]
\begin{multicols}{2}
  \item (1,5)
  \item 1
  \item 2
  \item -1
  \item \emptyset
  \item Todos los reales
\end{multicols}
\end{enumerate}
\vspace{0.15cm}
\item Si \(A = \{x \in \mathbb{R} \mid |x-3| \leq 2\}\), ¿cuál es A?
\begin{enumerate}[label=\textbf{\Alph*)}]
\begin{multicols}{2}
  \item 0
  \item [1,5]
  \item 2
  \item -1
  \item \emptyset
  \item Todos los reales
\end{multicols}
\end{enumerate}
\vspace{0.15cm}
\item Si \(A = \{x \in \mathbb{N} \mid x \text{ múltiplo de } 3, x < 15\}\), ¿cuál es A?
\begin{enumerate}[label=\textbf{\Alph*)}]
\begin{multicols}{2}
  \item {3,6,9,12}
  \item 1
  \item {-3,3}
  \item \emptyset
  \item \emptyset
  \item Todos los reales
\end{multicols}
\end{enumerate}
\vspace{0.15cm}
\item Si \(A = \{x \in \mathbb{N} \mid x \text{ divisor de } 12\}\), ¿cuál es A?
\begin{enumerate}[label=\textbf{\Alph*)}]
\begin{multicols}{2}
  \item {1,2,3,4,6,12}
  \item 1
  \item {-3,3}
  \item \emptyset
  \item \emptyset
  \item Todos los reales
\end{multicols}
\end{enumerate}
\vspace{0.15cm}
\item Si \(A = \{x \in \mathbb{N} \mid x \text{ divisor de } 18\}\), ¿cuál es A?
\begin{enumerate}[label=\textbf{\Alph*)}]
\begin{multicols}{2}
  \item {1,2,3,6,9,18}
  \item 1
  \item {-3,3}
  \item \emptyset
  \item \emptyset
  \item Todos los reales
\end{multicols}
\end{enumerate}
\vspace{0.15cm}
\item Si \(A = \{x \in \mathbb{N} \mid x \text{ múltiplo de } 4, x < 20\}\), ¿cuál es A?
\begin{enumerate}[label=\textbf{\Alph*)}]
\begin{multicols}{2}
  \item {4,8,12,16}
  \item 1
  \item {-3,3}
  \item \emptyset
  \item \emptyset
  \item Todos los reales
\end{multicols}
\end{enumerate}
\vspace{0.15cm}
\item Si \(A = \{x \in \mathbb{N} \mid x \text{ múltiplo de } 5, x < 30\}\), ¿cuál es A?
\begin{enumerate}[label=\textbf{\Alph*)}]
\begin{multicols}{2}
  \item {5,10,15,20,25}
  \item 1
  \item {-3,3}
  \item \emptyset
  \item \emptyset
  \item Todos los reales
\end{multicols}
\end{enumerate}
\vspace{0.15cm}
\end{enumerate}\n
\vspace{0.4cm}
\begin{tcolorbox}[colback=gray!5!white,colframe=gray!60!black,title=\textbf{SOLUCIONARIO -- 1.3 Conjuntos e Intervalos}]
\begin{multicols}{4}
\begin{enumerate}[label=\textbf{\arabic*.}]
  \item \textbf{A}
  \item \textbf{B}
  \item \textbf{E}
  \item \textbf{B}
  \item \textbf{C}
  \item \textbf{B}
  \item \textbf{C}
  \item \textbf{C}
  \item \textbf{A}
  \item \textbf{C}
  \item \textbf{C}
  \item \textbf{C}
  \item \textbf{B}
  \item \textbf{C}
  \item \textbf{D}
  \item \textbf{F}
  \item \textbf{A}
  \item \textbf{A}
  \item \textbf{B}
  \item \textbf{A}
  \item \textbf{B}
  \item \textbf{B}
  \item \textbf{A}
  \item \textbf{B}
  \item \textbf{A}
  \item \textbf{A}
  \item \textbf{C}
  \item \textbf{C}
  \item \textbf{B}
  \item \textbf{B}
  \item \textbf{A}
  \item \textbf{A}
  \item \textbf{A}
  \item \textbf{D}
  \item \textbf{B}
  \item \textbf{B}
  \item \textbf{A}
  \item \textbf{B}
  \item \textbf{A}
  \item \textbf{A}
  \item \textbf{B}
  \item \textbf{C}
  \item \textbf{A}
  \item \textbf{A}
  \item \textbf{A}
  \item \textbf{B}
  \item \textbf{C}
  \item \textbf{C}
  \item \textbf{A}
  \item \textbf{B}
  \item \textbf{A}
  \item \textbf{B}
  \item \textbf{A}
  \item \textbf{A}
  \item \textbf{B}
  \item \textbf{A}
  \item \textbf{B}
  \item \textbf{A}
  \item \textbf{B}
  \item \textbf{A}
  \item \textbf{C}
  \item \textbf{B}
  \item \textbf{B}
  \item \textbf{B}
  \item \textbf{A}
  \item \textbf{A}
  \item \textbf{A}
  \item \textbf{B}
  \item \textbf{A}
  \item \textbf{C}
  \item \textbf{A}
  \item \textbf{A}
  \item \textbf{B}
  \item \textbf{A}
  \item \textbf{A}
  \item \textbf{A}
  \item \textbf{A}
  \item \textbf{A}
  \item \textbf{B}
  \item \textbf{A}
  \item \textbf{B}
  \item \textbf{A}
  \item \textbf{A}
  \item \textbf{B}
  \item \textbf{A}
  \item \textbf{C}
  \item \textbf{C}
  \item \textbf{B}
  \item \textbf{A}
  \item \textbf{A}
  \item \textbf{A}
  \item \textbf{C}
  \item \textbf{B}
  \item \textbf{D}
  \item \textbf{B}
  \item \textbf{B}
  \item \textbf{B}
  \item \textbf{B}
  \item \textbf{D}
  \item \textbf{B}
  \item \textbf{A}
  \item \textbf{B}
  \item \textbf{C}
  \item \textbf{A}
  \item \textbf{B}
  \item \textbf{C}
  \item \textbf{D}
  \item \textbf{B}
  \item \textbf{A}
  \item \textbf{A}
  \item \textbf{B}
  \item \textbf{A}
  \item \textbf{A}
  \item \textbf{B}
  \item \textbf{A}
  \item \textbf{B}
  \item \textbf{A}
  \item \textbf{B}
  \item \textbf{C}
  \item \textbf{D}
  \item \textbf{A}
  \item \textbf{B}
  \item \textbf{A}
  \item \textbf{A}
  \item \textbf{C}
  \item \textbf{A}
  \item \textbf{A}
  \item \textbf{B}
  \item \textbf{A}
  \item \textbf{A}
  \item \textbf{A}
  \item \textbf{A}
  \item \textbf{D}
  \item \textbf{B}
  \item \textbf{C}
  \item \textbf{A}
  \item \textbf{A}
  \item \textbf{B}
  \item \textbf{A}
  \item \textbf{C}
  \item \textbf{C}
  \item \textbf{A}
  \item \textbf{B}
  \item \textbf{A}
  \item \textbf{B}
  \item \textbf{A}
  \item \textbf{C}
  \item \textbf{D}
  \item \textbf{A}
  \item \textbf{D}
  \item \textbf{A}
  \item \textbf{A}
  \item \textbf{A}
  \item \textbf{B}
  \item \textbf{A}
  \item \textbf{A}
  \item \textbf{A}
  \item \textbf{A}
  \item \textbf{B}
  \item \textbf{A}
  \item \textbf{A}
  \item \textbf{B}
  \item \textbf{C}
  \item \textbf{C}
  \item \textbf{D}
  \item \textbf{D}
  \item \textbf{D}
  \item \textbf{D}
  \item \textbf{A}
  \item \textbf{A}
  \item \textbf{B}
  \item \textbf{C}
  \item \textbf{D}
  \item \textbf{A}
  \item \textbf{B}
  \item \textbf{C}
  \item \textbf{D}
  \item \textbf{A}
  \item \textbf{C}
  \item \textbf{A}
  \item \textbf{B}
  \item \textbf{A}
  \item \textbf{B}
  \item \textbf{C}
  \item \textbf{D}
  \item \textbf{A}
  \item \textbf{C}
  \item \textbf{B}
  \item \textbf{D}
  \item \textbf{D}
  \item \textbf{B}
  \item \textbf{A}
  \item \textbf{B}
  \item \textbf{A}
  \item \textbf{A}
  \item \textbf{B}
  \item \textbf{E}
  \item \textbf{B}
  \item \textbf{C}
  \item \textbf{B}
  \item \textbf{B}
  \item \textbf{A}
  \item \textbf{B}
  \item \textbf{B}
  \item \textbf{A}
  \item \textbf{B}
  \item \textbf{B}
  \item \textbf{A}
  \item \textbf{A}
  \item \textbf{C}
  \item \textbf{A}
  \item \textbf{B}
  \item \textbf{B}
  \item \textbf{A}
  \item \textbf{B}
  \item \textbf{A}
  \item \textbf{A}
  \item \textbf{C}
  \item \textbf{B}
  \item \textbf{A}
  \item \textbf{A}
  \item \textbf{A}
  \item \textbf{A}
  \item \textbf{A}
  \item \textbf{A}
  \item \textbf{A}
  \item \textbf{C}
  \item \textbf{C}
  \item \textbf{D}
  \item \textbf{A}
  \item \textbf{B}
  \item \textbf{A}
  \item \textbf{A}
  \item \textbf{A}
  \item \textbf{A}
  \item \textbf{A}
  \item \textbf{B}
  \item \textbf{B}
  \item \textbf{D}
  \item \textbf{A}
  \item \textbf{B}
  \item \textbf{C}
  \item \textbf{D}
  \item \textbf{A}
  \item \textbf{B}
  \item \textbf{C}
  \item \textbf{D}
  \item \textbf{A}
  \item \textbf{B}
  \item \textbf{A}
  \item \textbf{B}
  \item \textbf{A}
  \item \textbf{B}
  \item \textbf{A}
  \item \textbf{C}
  \item \textbf{A}
  \item \textbf{B}
  \item \textbf{A}
  \item \textbf{B}
  \item \textbf{A}
  \item \textbf{B}
  \item \textbf{B}
  \item \textbf{D}
  \item \textbf{B}
  \item \textbf{A}
  \item \textbf{A}
  \item \textbf{A}
  \item \textbf{A}
  \item \textbf{A}
  \item \textbf{A}
  \item \textbf{B}
  \item \textbf{B}
  \item \textbf{B}
  \item \textbf{A}
  \item \textbf{B}
  \item \textbf{A}
  \item \textbf{B}
  \item \textbf{A}
  \item \textbf{B}
  \item \textbf{A}
  \item \textbf{C}
  \item \textbf{A}
  \item \textbf{A}
  \item \textbf{A}
  \item \textbf{A}
  \item \textbf{A}
  \item \textbf{B}
  \item \textbf{A}
  \item \textbf{B}
  \item \textbf{A}
  \item \textbf{B}
  \item \textbf{A}
  \item \textbf{B}
  \item \textbf{A}
  \item \textbf{B}
  \item \textbf{A}
  \item \textbf{A}
  \item \textbf{A}
  \item \textbf{A}
  \item \textbf{A}
\end{enumerate}
\end{multicols}
\end{tcolorbox}
\newpage\n
\subsection{1.4 Valor Absoluto}

\subsubsection*{SECCIÓN I: DEFINICIÓN Y CONCEPTOS BÁSICOS (1–40)}

\begin{enumerate}[label=\textbf{\arabic*.}]
\setcounter{enumi}{0}
\item ¿Qué es el valor absoluto de un número real \(x\)?
\begin{enumerate}[label=\textbf{\Alph*)}]
\begin{multicols}{2}
  \item Su distancia al origen en la recta real
  \item 1
  \item 0
  \item 10
  \item -2
  \item Ninguna
\end{multicols}
\end{enumerate}
\vspace{0.15cm}
\item ¿Cómo se denota el valor absoluto de \(x\)?
\begin{enumerate}[label=\textbf{\Alph*)}]
\begin{multicols}{2}
  \item x
  \item |x|
  \item 0
  \item 10
  \item -2
  \item Ninguna
\end{multicols}
\end{enumerate}
\vspace{0.15cm}
\item El valor absoluto de un número siempre es:
\begin{enumerate}[label=\textbf{\Alph*)}]
\begin{multicols}{2}
  \item x
  \item No negativo
  \item 0
  \item 10
  \item -2
  \item Ninguna
\end{multicols}
\end{enumerate}
\vspace{0.15cm}
\item ¿Cuál es el valor absoluto de \(5\)?
\begin{enumerate}[label=\textbf{\Alph*)}]
\begin{multicols}{2}
  \item -5
  \item 5
  \item 0
  \item 10
  \item -2
  \item Ninguna
\end{multicols}
\end{enumerate}
\vspace{0.15cm}
\item ¿Cuál es el valor absoluto de \(-7\)?
\begin{enumerate}[label=\textbf{\Alph*)}]
\begin{multicols}{2}
  \item -7
  \item 7
  \item 0
  \item 10
  \item -2
  \item Ninguna
\end{multicols}
\end{enumerate}
\vspace{0.15cm}
\item ¿Cuál es el valor absoluto de \(0\)?
\begin{enumerate}[label=\textbf{\Alph*)}]
\begin{multicols}{2}
  \item 0
  \item 1
  \item 0
  \item 10
  \item No existe
  \item Infinito
\end{multicols}
\end{enumerate}
\vspace{0.15cm}
\item ¿Cuál es el valor absoluto de \(-3\)?
\begin{enumerate}[label=\textbf{\Alph*)}]
\begin{multicols}{2}
  \item x
  \item 3
  \item 0
  \item 10
  \item -2
  \item Ninguna
\end{multicols}
\end{enumerate}
\vspace{0.15cm}
\item ¿Cuál es el valor absoluto de \(2.5\)?
\begin{enumerate}[label=\textbf{\Alph*)}]
\begin{multicols}{2}
  \item x
  \item 2.5
  \item 0
  \item 10
  \item -2
  \item Ninguna
\end{multicols}
\end{enumerate}
\vspace{0.15cm}
\item ¿Cuál es el valor absoluto de \(-0.5\)?
\begin{enumerate}[label=\textbf{\Alph*)}]
\begin{multicols}{2}
  \item x
  \item 0.5
  \item 0
  \item 10
  \item -2
  \item Ninguna
\end{multicols}
\end{enumerate}
\vspace{0.15cm}
\item ¿Cuál es el valor absoluto de \(\frac{3}{4}\)?
\begin{enumerate}[label=\textbf{\Alph*)}]
\begin{multicols}{2}
  \item x
  \item \frac{3}{4}
  \item 0
  \item 10
  \item -2
  \item Ninguna
\end{multicols}
\end{enumerate}
\vspace{0.15cm}
\item ¿Cuál es el valor absoluto de \(-\frac{5}{2}\)?
\begin{enumerate}[label=\textbf{\Alph*)}]
\begin{multicols}{2}
  \item x
  \item \frac{5}{2}
  \item 0
  \item 10
  \item -2
  \item Ninguna
\end{multicols}
\end{enumerate}
\vspace{0.15cm}
\item ¿Cuál es el valor absoluto de \(\sqrt{2}\)?
\begin{enumerate}[label=\textbf{\Alph*)}]
\begin{multicols}{2}
  \item x
  \item \sqrt{2}
  \item 0
  \item 10
  \item -2
  \item Ninguna
\end{multicols}
\end{enumerate}
\vspace{0.15cm}
\item ¿Cuál es el valor absoluto de \(-\pi\)?
\begin{enumerate}[label=\textbf{\Alph*)}]
\begin{multicols}{2}
  \item x
  \item \pi
  \item 0
  \item 10
  \item -2
  \item Ninguna
\end{multicols}
\end{enumerate}
\vspace{0.15cm}
\item ¿Cuál es el valor absoluto de \(10\)?
\begin{enumerate}[label=\textbf{\Alph*)}]
\begin{multicols}{2}
  \item x
  \item 10
  \item 0
  \item 10
  \item -2
  \item Ninguna
\end{multicols}
\end{enumerate}
\vspace{0.15cm}
\item ¿Cuál es el valor absoluto de \(-15\)?
\begin{enumerate}[label=\textbf{\Alph*)}]
\begin{multicols}{2}
  \item x
  \item 15
  \item 0
  \item 10
  \item -2
  \item Ninguna
\end{multicols}
\end{enumerate}
\vspace{0.15cm}
\item Si \(x>0\), entonces \(|x|=\)
\begin{enumerate}[label=\textbf{\Alph*)}]
\begin{multicols}{2}
  \item x
  \item x
  \item 0
  \item 10
  \item -2
  \item Ninguna
\end{multicols}
\end{enumerate}
\vspace{0.15cm}
\item Si \(x<0\), entonces \(|x|=\)
\begin{enumerate}[label=\textbf{\Alph*)}]
\begin{multicols}{2}
  \item -x
  \item 1
  \item 0
  \item 10
  \item -2
  \item Ninguna
\end{multicols}
\end{enumerate}
\vspace{0.15cm}
\item Si \(x=0\), entonces \(|x|=\)
\begin{enumerate}[label=\textbf{\Alph*)}]
\begin{multicols}{2}
  \item 0
  \item 0
  \item 0
  \item 10
  \item No existe
  \item Infinito
\end{multicols}
\end{enumerate}
\vspace{0.15cm}
\item La definición formal del valor absoluto es:
\begin{enumerate}[label=\textbf{\Alph*)}]
\begin{multicols}{2}
  \item |x|=x si x \geq 0, |x|=-x si x < 0
  \item 1
  \item 0
  \item 10
  \item -2
  \item Ninguna
\end{multicols}
\end{enumerate}
\vspace{0.15cm}
\item ¿Cuál es el valor absoluto de \(-1\)?
\begin{enumerate}[label=\textbf{\Alph*)}]
\begin{multicols}{2}
  \item x
  \item 1
  \item 0
  \item 10
  \item -2
  \item Ninguna
\end{multicols}
\end{enumerate}
\vspace{0.15cm}
\item ¿Cuál es el valor absoluto de \(1\)?
\begin{enumerate}[label=\textbf{\Alph*)}]
\begin{multicols}{2}
  \item x
  \item 1
  \item 0
  \item 10
  \item -2
  \item Ninguna
\end{multicols}
\end{enumerate}
\vspace{0.15cm}
\item ¿Cuál es el valor absoluto de \(-100\)?
\begin{enumerate}[label=\textbf{\Alph*)}]
\begin{multicols}{2}
  \item x
  \item 100
  \item 0
  \item 10
  \item -2
  \item Ninguna
\end{multicols}
\end{enumerate}
\vspace{0.15cm}
\item ¿Cuál es el valor absoluto de \(0.001\)?
\begin{enumerate}[label=\textbf{\Alph*)}]
\begin{multicols}{2}
  \item x
  \item 0.001
  \item 0
  \item 10
  \item -2
  \item Ninguna
\end{multicols}
\end{enumerate}
\vspace{0.15cm}
\item ¿Cuál es el valor absoluto de \(-0.001\)?
\begin{enumerate}[label=\textbf{\Alph*)}]
\begin{multicols}{2}
  \item x
  \item 0.001
  \item 0
  \item 10
  \item -2
  \item Ninguna
\end{multicols}
\end{enumerate}
\vspace{0.15cm}
\item El valor absoluto de un número real representa:
\begin{enumerate}[label=\textbf{\Alph*)}]
\begin{multicols}{2}
  \item Su magnitud sin signo
  \item 1
  \item 0
  \item 10
  \item -2
  \item Ninguna
\end{multicols}
\end{enumerate}
\vspace{0.15cm}
\item ¿Puede el valor absoluto de un número ser negativo?
\begin{enumerate}[label=\textbf{\Alph*)}]
\begin{multicols}{2}
  \item Sí, siempre
  \item No, nunca
  \item 0
  \item 10
  \item -2
  \item Ninguna
\end{multicols}
\end{enumerate}
\vspace{0.15cm}
\item Si \(|x|=5\), ¿cuáles son los posibles valores de \(x\)?
\begin{enumerate}[label=\textbf{\Alph*)}]
\begin{multicols}{2}
  \item x
  \item 1
  \item 5 y -5
  \item -1
  \item -2
  \item Ninguna
\end{multicols}
\end{enumerate}
\vspace{0.15cm}
\item Si \(|x|=3\), ¿cuáles son los posibles valores de \(x\)?
\begin{enumerate}[label=\textbf{\Alph*)}]
\begin{multicols}{2}
  \item x
  \item 1
  \item 3 y -3
  \item -1
  \item -2
  \item Ninguna
\end{multicols}
\end{enumerate}
\vspace{0.15cm}
\item Si \(|x|=0\), ¿cuál es el valor de \(x\)?
\begin{enumerate}[label=\textbf{\Alph*)}]
\begin{multicols}{2}
  \item 0
  \item 1
  \item 0
  \item 10
  \item No existe
  \item Infinito
\end{multicols}
\end{enumerate}
\vspace{0.15cm}
\item Si \(|x|=7\), ¿cuáles son los posibles valores de \(x\)?
\begin{enumerate}[label=\textbf{\Alph*)}]
\begin{multicols}{2}
  \item x
  \item 1
  \item 7 y -7
  \item -1
  \item -2
  \item Ninguna
\end{multicols}
\end{enumerate}
\vspace{0.15cm}
\item ¿Cuál es el valor absoluto de \(-2.25\)?
\begin{enumerate}[label=\textbf{\Alph*)}]
\begin{multicols}{2}
  \item x
  \item 2.25
  \item 0
  \item 10
  \item -2
  \item Ninguna
\end{multicols}
\end{enumerate}
\vspace{0.15cm}
\item ¿Cuál es el valor absoluto de \(\frac{7}{3}\)?
\begin{enumerate}[label=\textbf{\Alph*)}]
\begin{multicols}{2}
  \item x
  \item \frac{7}{3}
  \item 0
  \item 10
  \item -2
  \item Ninguna
\end{multicols}
\end{enumerate}
\vspace{0.15cm}
\item ¿Cuál es el valor absoluto de \(-\frac{2}{5}\)?
\begin{enumerate}[label=\textbf{\Alph*)}]
\begin{multicols}{2}
  \item x
  \item \frac{2}{5}
  \item 0
  \item 10
  \item -2
  \item Ninguna
\end{multicols}
\end{enumerate}
\vspace{0.15cm}
\item ¿Cuál es el valor absoluto de \(-\sqrt{3}\)?
\begin{enumerate}[label=\textbf{\Alph*)}]
\begin{multicols}{2}
  \item x
  \item \sqrt{3}
  \item 0
  \item 10
  \item -2
  \item Ninguna
\end{multicols}
\end{enumerate}
\vspace{0.15cm}
\item ¿Cuál es el valor absoluto de \(-e\)?
\begin{enumerate}[label=\textbf{\Alph*)}]
\begin{multicols}{2}
  \item x
  \item e
  \item 0
  \item 10
  \item -2
  \item Ninguna
\end{multicols}
\end{enumerate}
\vspace{0.15cm}
\item La expresión \(|x|\) se lee como:
\begin{enumerate}[label=\textbf{\Alph*)}]
\begin{multicols}{2}
  \item Valor absoluto de \(x\)
  \item 1
  \item 0
  \item 10
  \item -2
  \item Ninguna
\end{multicols}
\end{enumerate}
\vspace{0.15cm}
\item El valor absoluto también se conoce como:
\begin{enumerate}[label=\textbf{\Alph*)}]
\begin{multicols}{2}
  \item Módulo
  \item 1
  \item 0
  \item 10
  \item -2
  \item Ninguna
\end{multicols}
\end{enumerate}
\vspace{0.15cm}
\item ¿Cuál es el valor absoluto de \(-8\)?
\begin{enumerate}[label=\textbf{\Alph*)}]
\begin{multicols}{2}
  \item x
  \item 8
  \item 0
  \item 10
  \item -2
  \item Ninguna
\end{multicols}
\end{enumerate}
\vspace{0.15cm}
\item ¿Cuál es el valor absoluto de \(12\)?
\begin{enumerate}[label=\textbf{\Alph*)}]
\begin{multicols}{2}
  \item x
  \item 12
  \item 0
  \item 10
  \item -2
  \item Ninguna
\end{multicols}
\end{enumerate}
\vspace{0.15cm}
\item ¿Cuál es el valor absoluto de \(-20\)?
\begin{enumerate}[label=\textbf{\Alph*)}]
\begin{multicols}{2}
  \item x
  \item 20
  \item 0
  \item 10
  \item -2
  \item Ninguna
\end{multicols}
\end{enumerate}
\vspace{0.15cm}
\end{enumerate}\n
\subsubsection*{SECCIÓN II: PROPIEDADES DEL VALOR ABSOLUTO (41–90)}

\begin{enumerate}[label=\textbf{\arabic*.}]
\setcounter{enumi}{40}
\item ¿Cuál es el valor de \(|5|+|3|\)?
\begin{enumerate}[label=\textbf{\Alph*)}]
\begin{multicols}{2}
  \item 2
  \item 8
  \item -8
  \item 15
  \item 0
  \item 1
\end{multicols}
\end{enumerate}
\vspace{0.15cm}
\item ¿Cuál es el valor de \(|-5|+|-3|\)?
\begin{enumerate}[label=\textbf{\Alph*)}]
\begin{multicols}{2}
  \item 2
  \item 8
  \item -8
  \item 15
  \item 0
  \item 1
\end{multicols}
\end{enumerate}
\vspace{0.15cm}
\item ¿Cuál es el valor de \(|5|-|3|\)?
\begin{enumerate}[label=\textbf{\Alph*)}]
\begin{multicols}{2}
  \item 2
  \item 8
  \item -8
  \item 15
  \item 0
  \item 1
\end{multicols}
\end{enumerate}
\vspace{0.15cm}
\item ¿Cuál es el valor de \(|5-3|\)?
\begin{enumerate}[label=\textbf{\Alph*)}]
\begin{multicols}{2}
  \item 2
  \item 8
  \item -8
  \item 15
  \item 0
  \item 1
\end{multicols}
\end{enumerate}
\vspace{0.15cm}
\item ¿Cuál es el valor de \(|3-5|\)?
\begin{enumerate}[label=\textbf{\Alph*)}]
\begin{multicols}{2}
  \item 2
  \item 8
  \item -8
  \item 15
  \item 0
  \item 1
\end{multicols}
\end{enumerate}
\vspace{0.15cm}
\item ¿Cuál es el valor de \(|-5|\cdot|3|\)?
\begin{enumerate}[label=\textbf{\Alph*)}]
\begin{multicols}{2}
  \item 2
  \item 15
  \item -8
  \item 15
  \item 0
  \item 1
\end{multicols}
\end{enumerate}
\vspace{0.15cm}
\item ¿Cuál es el valor de \(|(-5)\cdot3|\)?
\begin{enumerate}[label=\textbf{\Alph*)}]
\begin{multicols}{2}
  \item 2
  \item 15
  \item -8
  \item 15
  \item 0
  \item 1
\end{multicols}
\end{enumerate}
\vspace{0.15cm}
\item ¿Cuál es el valor de \(\frac{|-12|}{|4|}\)?
\begin{enumerate}[label=\textbf{\Alph*)}]
\begin{multicols}{2}
  \item 3
  \item 8
  \item -8
  \item 15
  \item 0
  \item 1
\end{multicols}
\end{enumerate}
\vspace{0.15cm}
\item ¿Cuál es el valor de \(\left|\frac{-12}{4}\right|\)?
\begin{enumerate}[label=\textbf{\Alph*)}]
\begin{multicols}{2}
  \item 3
  \item 8
  \item -8
  \item 15
  \item 0
  \item 1
\end{multicols}
\end{enumerate}
\vspace{0.15cm}
\item La propiedad \(|a\cdot b|=|a|\cdot|b|\) se conoce como:
\begin{enumerate}[label=\textbf{\Alph*)}]
\begin{multicols}{2}
  \item Multiplicatividad
  \item Aditividad
  \item Desigualdad triangular
  \item Idempotencia
  \item Simetría
  \item Transitividad
\end{multicols}
\end{enumerate}
\vspace{0.15cm}
\item La propiedad \(|a+b|\leq|a|+|b|\) se conoce como:
\begin{enumerate}[label=\textbf{\Alph*)}]
\begin{multicols}{2}
  \item Desigualdad triangular
  \item Aditividad
  \item Desigualdad triangular
  \item Idempotencia
  \item Simetría
  \item Transitividad
\end{multicols}
\end{enumerate}
\vspace{0.15cm}
\item La propiedad \(|a|=|-a|\) se conoce como:
\begin{enumerate}[label=\textbf{\Alph*)}]
\begin{multicols}{2}
  \item Simetría
  \item Aditividad
  \item Desigualdad triangular
  \item Idempotencia
  \item Simetría
  \item Transitividad
\end{multicols}
\end{enumerate}
\vspace{0.15cm}
\item La propiedad \(|a|\geq0\) se conoce como:
\begin{enumerate}[label=\textbf{\Alph*)}]
\begin{multicols}{2}
  \item No negatividad
  \item Aditividad
  \item Desigualdad triangular
  \item Idempotencia
  \item Simetría
  \item Transitividad
\end{multicols}
\end{enumerate}
\vspace{0.15cm}
\item La propiedad \(|a|=0\Leftrightarrow a=0\) se conoce como:
\begin{enumerate}[label=\textbf{\Alph*)}]
\begin{multicols}{2}
  \item Anulación
  \item Aditividad
  \item Desigualdad triangular
  \item Idempotencia
  \item Simetría
  \item Transitividad
\end{multicols}
\end{enumerate}
\vspace{0.15cm}
\item ¿Cuál es el valor de \(|7|+|-2|\)?
\begin{enumerate}[label=\textbf{\Alph*)}]
\begin{multicols}{2}
  \item 2
  \item 9
  \item -8
  \item 15
  \item 0
  \item 1
\end{multicols}
\end{enumerate}
\vspace{0.15cm}
\item ¿Cuál es el valor de \(|7-2|\)?
\begin{enumerate}[label=\textbf{\Alph*)}]
\begin{multicols}{2}
  \item 5
  \item 8
  \item -8
  \item 15
  \item 0
  \item 1
\end{multicols}
\end{enumerate}
\vspace{0.15cm}
\item ¿Cuál es el valor de \(|2-7|\)?
\begin{enumerate}[label=\textbf{\Alph*)}]
\begin{multicols}{2}
  \item 5
  \item 8
  \item -8
  \item 15
  \item 0
  \item 1
\end{multicols}
\end{enumerate}
\vspace{0.15cm}
\item ¿Cuál es el valor de \(|-4|\cdot|-2|\)?
\begin{enumerate}[label=\textbf{\Alph*)}]
\begin{multicols}{2}
  \item 2
  \item 8
  \item -8
  \item 15
  \item 0
  \item 1
\end{multicols}
\end{enumerate}
\vspace{0.15cm}
\item ¿Cuál es el valor de \(|(-4)\cdot(-2)|\)?
\begin{enumerate}[label=\textbf{\Alph*)}]
\begin{multicols}{2}
  \item 2
  \item 8
  \item -8
  \item 15
  \item 0
  \item 1
\end{multicols}
\end{enumerate}
\vspace{0.15cm}
\item ¿Cuál es el valor de \(\frac{|-20|}{|-5|}\)?
\begin{enumerate}[label=\textbf{\Alph*)}]
\begin{multicols}{2}
  \item 4
  \item 8
  \item -8
  \item 15
  \item 0
  \item 1
\end{multicols}
\end{enumerate}
\vspace{0.15cm}
\item ¿Cuál es el valor de \(\left|\frac{-20}{-5}\right|\)?
\begin{enumerate}[label=\textbf{\Alph*)}]
\begin{multicols}{2}
  \item 4
  \item 8
  \item -8
  \item 15
  \item 0
  \item 1
\end{multicols}
\end{enumerate}
\vspace{0.15cm}
\item ¿Cuál es el valor de \(|3|^2\)?
\begin{enumerate}[label=\textbf{\Alph*)}]
\begin{multicols}{2}
  \item 2
  \item 9
  \item -8
  \item 15
  \item 0
  \item 1
\end{multicols}
\end{enumerate}
\vspace{0.15cm}
\item ¿Cuál es el valor de \(|-3|^2\)?
\begin{enumerate}[label=\textbf{\Alph*)}]
\begin{multicols}{2}
  \item 2
  \item 9
  \item -8
  \item 15
  \item 0
  \item 1
\end{multicols}
\end{enumerate}
\vspace{0.15cm}
\item ¿Cuál es el valor de \(\sqrt{(-3)^2}\)?
\begin{enumerate}[label=\textbf{\Alph*)}]
\begin{multicols}{2}
  \item 2
  \item 3
  \item -8
  \item 15
  \item 0
  \item 1
\end{multicols}
\end{enumerate}
\vspace{0.15cm}
\item ¿Cuál es el valor de \(\sqrt{(-5)^2}\)?
\begin{enumerate}[label=\textbf{\Alph*)}]
\begin{multicols}{2}
  \item 2
  \item 5
  \item -8
  \item 15
  \item 0
  \item 1
\end{multicols}
\end{enumerate}
\vspace{0.15cm}
\item ¿Cuál es el valor de \(|2|+|-3|+|-4|\)?
\begin{enumerate}[label=\textbf{\Alph*)}]
\begin{multicols}{2}
  \item 9
  \item 8
  \item -8
  \item 15
  \item 0
  \item 1
\end{multicols}
\end{enumerate}
\vspace{0.15cm}
\item ¿Cuál es el valor de \(|2-3-4|\)?
\begin{enumerate}[label=\textbf{\Alph*)}]
\begin{multicols}{2}
  \item 5
  \item 8
  \item -8
  \item 15
  \item 0
  \item 1
\end{multicols}
\end{enumerate}
\vspace{0.15cm}
\item ¿Cuál es el valor de \(|-2+3+4|\)?
\begin{enumerate}[label=\textbf{\Alph*)}]
\begin{multicols}{2}
  \item 5
  \item 8
  \item -8
  \item 15
  \item 0
  \item 1
\end{multicols}
\end{enumerate}
\vspace{0.15cm}
\item ¿Cuál es el valor de \(|(-2)\cdot3\cdot(-4)|\)?
\begin{enumerate}[label=\textbf{\Alph*)}]
\begin{multicols}{2}
  \item 24
  \item 8
  \item -8
  \item 15
  \item 0
  \item 1
\end{multicols}
\end{enumerate}
\vspace{0.15cm}
\item ¿Cuál es el valor de \(|-2|\cdot|3|\cdot|-4|\)?
\begin{enumerate}[label=\textbf{\Alph*)}]
\begin{multicols}{2}
  \item 24
  \item 8
  \item -8
  \item 15
  \item 0
  \item 1
\end{multicols}
\end{enumerate}
\vspace{0.15cm}
\item Si \(|a|=4\) y \(|b|=3\), ¿cuál es el valor de \(|a\cdot b|\)?
\begin{enumerate}[label=\textbf{\Alph*)}]
\begin{multicols}{2}
  \item 12
  \item 8
  \item -8
  \item 15
  \item 0
  \item 1
\end{multicols}
\end{enumerate}
\vspace{0.15cm}
\item Si \(|a|=6\) y \(|b|=2\), ¿cuál es el valor de \(\left|\frac{a}{b}\right|\)?
\begin{enumerate}[label=\textbf{\Alph*)}]
\begin{multicols}{2}
  \item 3
  \item 8
  \item -8
  \item 15
  \item 0
  \item 1
\end{multicols}
\end{enumerate}
\vspace{0.15cm}
\item Si \(|a|=5\), ¿cuál es el valor de \(|-a|\)?
\begin{enumerate}[label=\textbf{\Alph*)}]
\begin{multicols}{2}
  \item 5
  \item 8
  \item -8
  \item 15
  \item 0
  \item 1
\end{multicols}
\end{enumerate}
\vspace{0.15cm}
\item Si \(|a|=7\), ¿cuál es el valor de \(|a|^2\)?
\begin{enumerate}[label=\textbf{\Alph*)}]
\begin{multicols}{2}
  \item 2
  \item 49
  \item -8
  \item 15
  \item 0
  \item 1
\end{multicols}
\end{enumerate}
\vspace{0.15cm}
\item ¿Cuál es el valor de \(||-5|-|3||\)?
\begin{enumerate}[label=\textbf{\Alph*)}]
\begin{multicols}{2}
  \item 2
  \item 8
  \item -8
  \item 15
  \item 0
  \item 1
\end{multicols}
\end{enumerate}
\vspace{0.15cm}
\item ¿Cuál es el valor de \(|2|+|-2|-|0|\)?
\begin{enumerate}[label=\textbf{\Alph*)}]
\begin{multicols}{2}
  \item 4
  \item 8
  \item -8
  \item 15
  \item 0
  \item 1
\end{multicols}
\end{enumerate}
\vspace{0.15cm}
\item ¿Cuál es el valor de \(|-1|+|-1|+|-1|\)?
\begin{enumerate}[label=\textbf{\Alph*)}]
\begin{multicols}{2}
  \item 3
  \item 8
  \item -8
  \item 15
  \item 0
  \item 1
\end{multicols}
\end{enumerate}
\vspace{0.15cm}
\item ¿Cuál es el valor de \(|1-1|\)?
\begin{enumerate}[label=\textbf{\Alph*)}]
\begin{multicols}{2}
  \item 0
  \item 8
  \item -8
  \item 15
  \item 0
  \item 1
\end{multicols}
\end{enumerate}
\vspace{0.15cm}
\item ¿Cuál es el valor de \(|0-0|\)?
\begin{enumerate}[label=\textbf{\Alph*)}]
\begin{multicols}{2}
  \item 0
  \item 8
  \item -8
  \item 15
  \item 0
  \item 1
\end{multicols}
\end{enumerate}
\vspace{0.15cm}
\item ¿Cuál es el valor de \(|100-100|\)?
\begin{enumerate}[label=\textbf{\Alph*)}]
\begin{multicols}{2}
  \item 0
  \item 8
  \item -8
  \item 15
  \item 0
  \item 1
\end{multicols}
\end{enumerate}
\vspace{0.15cm}
\item Si \(|x|=|y|\), entonces:
\begin{enumerate}[label=\textbf{\Alph*)}]
\begin{multicols}{2}
  \item 2
  \item 8
  \item x = y o x = -y
  \item 15
  \item 0
  \item 1
\end{multicols}
\end{enumerate}
\vspace{0.15cm}
\item Si \(|x|\leq5\), entonces:
\begin{enumerate}[label=\textbf{\Alph*)}]
\begin{multicols}{2}
  \item 2
  \item 8
  \item -5 <= x <= 5
  \item 15
  \item 0
  \item 1
\end{multicols}
\end{enumerate}
\vspace{0.15cm}
\item Si \(|x|<3\), entonces:
\begin{enumerate}[label=\textbf{\Alph*)}]
\begin{multicols}{2}
  \item -3 < x < 3
  \item 8
  \item -8
  \item 15
  \item 0
  \item 1
\end{multicols}
\end{enumerate}
\vspace{0.15cm}
\item Si \(|x|>4\), entonces:
\begin{enumerate}[label=\textbf{\Alph*)}]
\begin{multicols}{2}
  \item 2
  \item 8
  \item x > 4 o x < -4
  \item 15
  \item 0
  \item 1
\end{multicols}
\end{enumerate}
\vspace{0.15cm}
\item Si \(|x|\geq2\), entonces:
\begin{enumerate}[label=\textbf{\Alph*)}]
\begin{multicols}{2}
  \item 2
  \item 8
  \item x >= 2 o x <= -2
  \item 15
  \item 0
  \item 1
\end{multicols}
\end{enumerate}
\vspace{0.15cm}
\item ¿Cuál es el valor de \(|a+b|\) si \(a=3\) y \(b=-7\)?
\begin{enumerate}[label=\textbf{\Alph*)}]
\begin{multicols}{2}
  \item 4
  \item 8
  \item -8
  \item 15
  \item 0
  \item 1
\end{multicols}
\end{enumerate}
\vspace{0.15cm}
\item ¿Cuál es el valor de \(|a|+|b|\) si \(a=3\) y \(b=-7\)?
\begin{enumerate}[label=\textbf{\Alph*)}]
\begin{multicols}{2}
  \item 2
  \item 8
  \item 10
  \item 15
  \item 0
  \item 1
\end{multicols}
\end{enumerate}
\vspace{0.15cm}
\item ¿Se cumple que \(|3+(-7)|\leq|3|+|-7|\)?
\begin{enumerate}[label=\textbf{\Alph*)}]
\begin{multicols}{2}
  \item Sí, \(4\leq10\)
  \item 8
  \item -8
  \item 15
  \item 0
  \item 1
\end{multicols}
\end{enumerate}
\vspace{0.15cm}
\item ¿Cuál es el valor de \(|a-b|\) si \(a=5\) y \(b=2\)?
\begin{enumerate}[label=\textbf{\Alph*)}]
\begin{multicols}{2}
  \item 3
  \item 8
  \item -8
  \item 15
  \item 0
  \item 1
\end{multicols}
\end{enumerate}
\vspace{0.15cm}
\item ¿Cuál es el valor de \(|a-b|\) si \(a=2\) y \(b=5\)?
\begin{enumerate}[label=\textbf{\Alph*)}]
\begin{multicols}{2}
  \item 3
  \item 8
  \item -8
  \item 15
  \item 0
  \item 1
\end{multicols}
\end{enumerate}
\vspace{0.15cm}
\end{enumerate}\n
\subsubsection*{SECCIÓN III: ECUACIONES CON VALOR ABSOLUTO (91–150)}

\begin{enumerate}[label=\textbf{\arabic*.}]
\setcounter{enumi}{90}
\item Resuelve \(|x|=4\).
\begin{enumerate}[label=\textbf{\Alph*)}]
\begin{multicols}{2}
  \item x = 4
  \item x = -4
  \item x = 4 o x = -4
  \item x = 2
  \item x = -2
  \item Sin solución
\end{multicols}
\end{enumerate}
\vspace{0.15cm}
\item Resuelve \(|x|=7\).
\begin{enumerate}[label=\textbf{\Alph*)}]
\begin{multicols}{2}
  \item x = 7
  \item x = -7
  \item x = 7 o x = -7
  \item x = 2
  \item x = -2
  \item Sin solución
\end{multicols}
\end{enumerate}
\vspace{0.15cm}
\item Resuelve \(|x|=0\).
\begin{enumerate}[label=\textbf{\Alph*)}]
\begin{multicols}{2}
  \item x = 0
  \item x = -4
  \item x = 0
  \item x = 2
  \item x = -2
  \item Sin solución
\end{multicols}
\end{enumerate}
\vspace{0.15cm}
\item Resuelve \(|x|=10\).
\begin{enumerate}[label=\textbf{\Alph*)}]
\begin{multicols}{2}
  \item x = 10
  \item x = -10
  \item x = 10 o x = -10
  \item x = 2
  \item x = -2
  \item Sin solución
\end{multicols}
\end{enumerate}
\vspace{0.15cm}
\item Resuelve \(|x|=3\).
\begin{enumerate}[label=\textbf{\Alph*)}]
\begin{multicols}{2}
  \item x = 3
  \item x = -3
  \item x = 3 o x = -3
  \item x = 2
  \item x = -2
  \item Sin solución
\end{multicols}
\end{enumerate}
\vspace{0.15cm}
\item Resuelve \(|x|=12\).
\begin{enumerate}[label=\textbf{\Alph*)}]
\begin{multicols}{2}
  \item x = 12
  \item x = -12
  \item x = 12 o x = -12
  \item x = 2
  \item x = -2
  \item Sin solución
\end{multicols}
\end{enumerate}
\vspace{0.15cm}
\item Resuelve \(|x|=1\).
\begin{enumerate}[label=\textbf{\Alph*)}]
\begin{multicols}{2}
  \item x = 1
  \item x = -1
  \item x = 1 o x = -1
  \item x = 2
  \item x = -2
  \item Sin solución
\end{multicols}
\end{enumerate}
\vspace{0.15cm}
\item Resuelve \(|x|=9\).
\begin{enumerate}[label=\textbf{\Alph*)}]
\begin{multicols}{2}
  \item x = 9
  \item x = -9
  \item x = 9 o x = -9
  \item x = 2
  \item x = -2
  \item Sin solución
\end{multicols}
\end{enumerate}
\vspace{0.15cm}
\item Resuelve \(|x|=15\).
\begin{enumerate}[label=\textbf{\Alph*)}]
\begin{multicols}{2}
  \item x = 15
  \item x = -15
  \item x = 15 o x = -15
  \item x = 2
  \item x = -2
  \item Sin solución
\end{multicols}
\end{enumerate}
\vspace{0.15cm}
\item Resuelve \(|x|=2.5\).
\begin{enumerate}[label=\textbf{\Alph*)}]
\begin{multicols}{2}
  \item x = 2.5
  \item x = -2.5
  \item x = 2.5 o x = -2.5
  \item x = 2
  \item x = -2
  \item Sin solución
\end{multicols}
\end{enumerate}
\vspace{0.15cm}
\item Resuelve \(|x|=6\).
\begin{enumerate}[label=\textbf{\Alph*)}]
\begin{multicols}{2}
  \item x = 6
  \item x = -6
  \item x = 6 o x = -6
  \item x = 2
  \item x = -2
  \item Sin solución
\end{multicols}
\end{enumerate}
\vspace{0.15cm}
\item Resuelve \(|x|=8\).
\begin{enumerate}[label=\textbf{\Alph*)}]
\begin{multicols}{2}
  \item x = 8
  \item x = -8
  \item x = 8 o x = -8
  \item x = 2
  \item x = -2
  \item Sin solución
\end{multicols}
\end{enumerate}
\vspace{0.15cm}
\item Resuelve \(|x-3|=5\).
\begin{enumerate}[label=\textbf{\Alph*)}]
\begin{multicols}{2}
  \item x = 8 o x = -2
  \item x = -2
  \item x = 0
  \item x = 2
  \item x = -2
  \item Sin solución
\end{multicols}
\end{enumerate}
\vspace{0.15cm}
\item Resuelve \(|x+2|=6\).
\begin{enumerate}[label=\textbf{\Alph*)}]
\begin{multicols}{2}
  \item x = 4 o x = -8
  \item x = -8
  \item x = 0
  \item x = 2
  \item x = -2
  \item Sin solución
\end{multicols}
\end{enumerate}
\vspace{0.15cm}
\item Resuelve \(|x-1|=4\).
\begin{enumerate}[label=\textbf{\Alph*)}]
\begin{multicols}{2}
  \item x = 5 o x = -3
  \item x = -3
  \item x = 0
  \item x = 2
  \item x = -2
  \item Sin solución
\end{multicols}
\end{enumerate}
\vspace{0.15cm}
\item Resuelve \(|x+5|=2\).
\begin{enumerate}[label=\textbf{\Alph*)}]
\begin{multicols}{2}
  \item x = -3 o x = -7
  \item x = -7
  \item x = 0
  \item x = 2
  \item x = -2
  \item Sin solución
\end{multicols}
\end{enumerate}
\vspace{0.15cm}
\item Resuelve \(|x-4|=0\).
\begin{enumerate}[label=\textbf{\Alph*)}]
\begin{multicols}{2}
  \item x = 4
  \item x = -4
  \item x = 0
  \item x = 2
  \item x = -2
  \item Sin solución
\end{multicols}
\end{enumerate}
\vspace{0.15cm}
\item Resuelve \(|x+1|=0\).
\begin{enumerate}[label=\textbf{\Alph*)}]
\begin{multicols}{2}
  \item x = -1
  \item x = -4
  \item x = 0
  \item x = 2
  \item x = -2
  \item Sin solución
\end{multicols}
\end{enumerate}
\vspace{0.15cm}
\item Resuelve \(|x-2|=3\).
\begin{enumerate}[label=\textbf{\Alph*)}]
\begin{multicols}{2}
  \item x = 5 o x = -1
  \item x = -1
  \item x = 0
  \item x = 2
  \item x = -2
  \item Sin solución
\end{multicols}
\end{enumerate}
\vspace{0.15cm}
\item Resuelve \(|x+3|=1\).
\begin{enumerate}[label=\textbf{\Alph*)}]
\begin{multicols}{2}
  \item x = -2 o x = -4
  \item x = -4
  \item x = 0
  \item x = 2
  \item x = -2
  \item Sin solución
\end{multicols}
\end{enumerate}
\vspace{0.15cm}
\item Resuelve \(|2x|=6\).
\begin{enumerate}[label=\textbf{\Alph*)}]
\begin{multicols}{2}
  \item x = 3 o x = -3
  \item x = -3
  \item x = 0
  \item x = 2
  \item x = -2
  \item Sin solución
\end{multicols}
\end{enumerate}
\vspace{0.15cm}
\item Resuelve \(|3x|=12\).
\begin{enumerate}[label=\textbf{\Alph*)}]
\begin{multicols}{2}
  \item x = 4 o x = -4
  \item x = -4
  \item x = 0
  \item x = 2
  \item x = -2
  \item Sin solución
\end{multicols}
\end{enumerate}
\vspace{0.15cm}
\item Resuelve \(|4x|=8\).
\begin{enumerate}[label=\textbf{\Alph*)}]
\begin{multicols}{2}
  \item x = 2 o x = -2
  \item x = -2
  \item x = 0
  \item x = 2
  \item x = -2
  \item Sin solución
\end{multicols}
\end{enumerate}
\vspace{0.15cm}
\item Resuelve \(|5x|=10\).
\begin{enumerate}[label=\textbf{\Alph*)}]
\begin{multicols}{2}
  \item x = 2 o x = -2
  \item x = -2
  \item x = 0
  \item x = 2
  \item x = -2
  \item Sin solución
\end{multicols}
\end{enumerate}
\vspace{0.15cm}
\item Resuelve \(|2x-1|=5\).
\begin{enumerate}[label=\textbf{\Alph*)}]
\begin{multicols}{2}
  \item x = 3 o x = -2
  \item x = -2
  \item x = 0
  \item x = 2
  \item x = -2
  \item Sin solución
\end{multicols}
\end{enumerate}
\vspace{0.15cm}
\item Resuelve \(|3x+2|=8\).
\begin{enumerate}[label=\textbf{\Alph*)}]
\begin{multicols}{2}
  \item x = 2 o x = -\frac{10}{3}
  \item x = -\frac{10}{3}
  \item x = 0
  \item x = 2
  \item x = -2
  \item Sin solución
\end{multicols}
\end{enumerate}
\vspace{0.15cm}
\item Resuelve \(|x-3|=|x+1|\).
\begin{enumerate}[label=\textbf{\Alph*)}]
\begin{multicols}{2}
  \item x = 1
  \item x = -4
  \item x = 0
  \item x = 2
  \item x = -2
  \item Sin solución
\end{multicols}
\end{enumerate}
\vspace{0.15cm}
\item Resuelve \(|x+2|=|x-4|\).
\begin{enumerate}[label=\textbf{\Alph*)}]
\begin{multicols}{2}
  \item x = 1
  \item x = -4
  \item x = 0
  \item x = 2
  \item x = -2
  \item Sin solución
\end{multicols}
\end{enumerate}
\vspace{0.15cm}
\item Resuelve \(|2x-1|=|x+3|\).
\begin{enumerate}[label=\textbf{\Alph*)}]
\begin{multicols}{2}
  \item x = 4 o x = -\frac{2}{3}
  \item x = -\frac{2}{3}
  \item x = 0
  \item x = 2
  \item x = -2
  \item Sin solución
\end{multicols}
\end{enumerate}
\vspace{0.15cm}
\item Resuelve \(|x|=x\).
\begin{enumerate}[label=\textbf{\Alph*)}]
\begin{multicols}{2}
  \item x \geq 0
  \item x = -4
  \item x = 0
  \item x = 2
  \item x = -2
  \item Sin solución
\end{multicols}
\end{enumerate}
\vspace{0.15cm}
\item Resuelve \(|x|=-x\).
\begin{enumerate}[label=\textbf{\Alph*)}]
\begin{multicols}{2}
  \item x \leq 0
  \item x = -4
  \item x = 0
  \item x = 2
  \item x = -2
  \item Sin solución
\end{multicols}
\end{enumerate}
\vspace{0.15cm}
\item Resuelve \(|x|=2x\).
\begin{enumerate}[label=\textbf{\Alph*)}]
\begin{multicols}{2}
  \item x = 0
  \item x = -4
  \item x = 0
  \item x = 2
  \item x = -2
  \item Sin solución
\end{multicols}
\end{enumerate}
\vspace{0.15cm}
\item Resuelve \(|x|=-2x\).
\begin{enumerate}[label=\textbf{\Alph*)}]
\begin{multicols}{2}
  \item x = 0
  \item x = -4
  \item x = 0
  \item x = 2
  \item x = -2
  \item Sin solución
\end{multicols}
\end{enumerate}
\vspace{0.15cm}
\item Resuelve \(|x-5|=0\).
\begin{enumerate}[label=\textbf{\Alph*)}]
\begin{multicols}{2}
  \item x = 5
  \item x = -4
  \item x = 0
  \item x = 2
  \item x = -2
  \item Sin solución
\end{multicols}
\end{enumerate}
\vspace{0.15cm}
\item Resuelve \(|x+7|=0\).
\begin{enumerate}[label=\textbf{\Alph*)}]
\begin{multicols}{2}
  \item x = -7
  \item x = -4
  \item x = 0
  \item x = 2
  \item x = -2
  \item Sin solución
\end{multicols}
\end{enumerate}
\vspace{0.15cm}
\item Resuelve \(|x-1|=1\).
\begin{enumerate}[label=\textbf{\Alph*)}]
\begin{multicols}{2}
  \item x = 2 o x = 0
  \item x = 0
  \item x = 0
  \item x = 2
  \item x = -2
  \item Sin solución
\end{multicols}
\end{enumerate}
\vspace{0.15cm}
\item Resuelve \(|x+4|=4\).
\begin{enumerate}[label=\textbf{\Alph*)}]
\begin{multicols}{2}
  \item x = 0 o x = -8
  \item x = -8
  \item x = 0
  \item x = 2
  \item x = -2
  \item Sin solución
\end{multicols}
\end{enumerate}
\vspace{0.15cm}
\item Resuelve \(|2x+1|=7\).
\begin{enumerate}[label=\textbf{\Alph*)}]
\begin{multicols}{2}
  \item x = 3 o x = -4
  \item x = -4
  \item x = 0
  \item x = 2
  \item x = -2
  \item Sin solución
\end{multicols}
\end{enumerate}
\vspace{0.15cm}
\item Resuelve \(|3x-2|=4\).
\begin{enumerate}[label=\textbf{\Alph*)}]
\begin{multicols}{2}
  \item x = 2 o x = -\frac{2}{3}
  \item x = -\frac{2}{3}
  \item x = 0
  \item x = 2
  \item x = -2
  \item Sin solución
\end{multicols}
\end{enumerate}
\vspace{0.15cm}
\item Resuelve \(|4x+3|=1\).
\begin{enumerate}[label=\textbf{\Alph*)}]
\begin{multicols}{2}
  \item x = -\frac{1}{2} o x = -1
  \item x = -1
  \item x = 0
  \item x = 2
  \item x = -2
  \item Sin solución
\end{multicols}
\end{enumerate}
\vspace{0.15cm}
\item Resuelve \(|x-6|=2\).
\begin{enumerate}[label=\textbf{\Alph*)}]
\begin{multicols}{2}
  \item x = 8 o x = 4
  \item x = 4
  \item x = 0
  \item x = 2
  \item x = -2
  \item Sin solución
\end{multicols}
\end{enumerate}
\vspace{0.15cm}
\item Resuelve \(|x+8|=3\).
\begin{enumerate}[label=\textbf{\Alph*)}]
\begin{multicols}{2}
  \item x = -5 o x = -11
  \item x = -11
  \item x = 0
  \item x = 2
  \item x = -2
  \item Sin solución
\end{multicols}
\end{enumerate}
\vspace{0.15cm}
\item Resuelve \(|x-2|=2\).
\begin{enumerate}[label=\textbf{\Alph*)}]
\begin{multicols}{2}
  \item x = 4 o x = 0
  \item x = 0
  \item x = 0
  \item x = 2
  \item x = -2
  \item Sin solución
\end{multicols}
\end{enumerate}
\vspace{0.15cm}
\item Resuelve \(|x+3|=3\).
\begin{enumerate}[label=\textbf{\Alph*)}]
\begin{multicols}{2}
  \item x = 0 o x = -6
  \item x = -6
  \item x = 0
  \item x = 2
  \item x = -2
  \item Sin solución
\end{multicols}
\end{enumerate}
\vspace{0.15cm}
\item Resuelve \(|x-7|=1\).
\begin{enumerate}[label=\textbf{\Alph*)}]
\begin{multicols}{2}
  \item x = 8 o x = 6
  \item x = 6
  \item x = 0
  \item x = 2
  \item x = -2
  \item Sin solución
\end{multicols}
\end{enumerate}
\vspace{0.15cm}
\item Resuelve \(|x+9|=2\).
\begin{enumerate}[label=\textbf{\Alph*)}]
\begin{multicols}{2}
  \item x = -7 o x = -11
  \item x = -11
  \item x = 0
  \item x = 2
  \item x = -2
  \item Sin solución
\end{multicols}
\end{enumerate}
\vspace{0.15cm}
\item Resuelve \(|x-4|=6\).
\begin{enumerate}[label=\textbf{\Alph*)}]
\begin{multicols}{2}
  \item x = 10 o x = -2
  \item x = -2
  \item x = 0
  \item x = 2
  \item x = -2
  \item Sin solución
\end{multicols}
\end{enumerate}
\vspace{0.15cm}
\item Resuelve \(|x+1|=5\).
\begin{enumerate}[label=\textbf{\Alph*)}]
\begin{multicols}{2}
  \item x = 4 o x = -6
  \item x = -6
  \item x = 0
  \item x = 2
  \item x = -2
  \item Sin solución
\end{multicols}
\end{enumerate}
\vspace{0.15cm}
\item Resuelve \(|x-3|=7\).
\begin{enumerate}[label=\textbf{\Alph*)}]
\begin{multicols}{2}
  \item x = 10 o x = -4
  \item x = -4
  \item x = 0
  \item x = 2
  \item x = -2
  \item Sin solución
\end{multicols}
\end{enumerate}
\vspace{0.15cm}
\item Resuelve \(|x+2|=9\).
\begin{enumerate}[label=\textbf{\Alph*)}]
\begin{multicols}{2}
  \item x = 7 o x = -11
  \item x = -11
  \item x = 0
  \item x = 2
  \item x = -2
  \item Sin solución
\end{multicols}
\end{enumerate}
\vspace{0.15cm}
\item Resuelve \(|2x-3|=1\).
\begin{enumerate}[label=\textbf{\Alph*)}]
\begin{multicols}{2}
  \item x = 2 o x = 1
  \item x = 1
  \item x = 0
  \item x = 2
  \item x = -2
  \item Sin solución
\end{multicols}
\end{enumerate}
\vspace{0.15cm}
\item Resuelve \(|3x+1|=4\).
\begin{enumerate}[label=\textbf{\Alph*)}]
\begin{multicols}{2}
  \item x = 1 o x = -\frac{5}{3}
  \item x = -\frac{5}{3}
  \item x = 0
  \item x = 2
  \item x = -2
  \item Sin solución
\end{multicols}
\end{enumerate}
\vspace{0.15cm}
\item Resuelve \(|4x-1|=3\).
\begin{enumerate}[label=\textbf{\Alph*)}]
\begin{multicols}{2}
  \item x = 1 o x = -\frac{1}{2}
  \item x = -\frac{1}{2}
  \item x = 0
  \item x = 2
  \item x = -2
  \item Sin solución
\end{multicols}
\end{enumerate}
\vspace{0.15cm}
\item Resuelve \(|5x+2|=7\).
\begin{enumerate}[label=\textbf{\Alph*)}]
\begin{multicols}{2}
  \item x = 1 o x = -\frac{9}{5}
  \item x = -\frac{9}{5}
  \item x = 0
  \item x = 2
  \item x = -2
  \item Sin solución
\end{multicols}
\end{enumerate}
\vspace{0.15cm}
\item Resuelve \(|x-1|=|x-3|\).
\begin{enumerate}[label=\textbf{\Alph*)}]
\begin{multicols}{2}
  \item x = 2
  \item x = -4
  \item x = 0
  \item x = 2
  \item x = -2
  \item Sin solución
\end{multicols}
\end{enumerate}
\vspace{0.15cm}
\item Resuelve \(|x+2|=|x-2|\).
\begin{enumerate}[label=\textbf{\Alph*)}]
\begin{multicols}{2}
  \item x = 0
  \item x = -4
  \item x = 0
  \item x = 2
  \item x = -2
  \item Sin solución
\end{multicols}
\end{enumerate}
\vspace{0.15cm}
\item Resuelve \(|x-4|=|x+2|\).
\begin{enumerate}[label=\textbf{\Alph*)}]
\begin{multicols}{2}
  \item x = 1
  \item x = -4
  \item x = 0
  \item x = 2
  \item x = -2
  \item Sin solución
\end{multicols}
\end{enumerate}
\vspace{0.15cm}
\item Resuelve \(|2x+1|=|x-2|\).
\begin{enumerate}[label=\textbf{\Alph*)}]
\begin{multicols}{2}
  \item x = -3 o x = \frac{1}{3}
  \item x = \frac{1}{3}
  \item x = 0
  \item x = 2
  \item x = -2
  \item Sin solución
\end{multicols}
\end{enumerate}
\vspace{0.15cm}
\item Resuelve \(|3x-1|=|x+1|\).
\begin{enumerate}[label=\textbf{\Alph*)}]
\begin{multicols}{2}
  \item x = 1 o x = 0
  \item x = 0
  \item x = 0
  \item x = 2
  \item x = -2
  \item Sin solución
\end{multicols}
\end{enumerate}
\vspace{0.15cm}
\item Resuelve \(|4x+2|=|2x-4|\).
\begin{enumerate}[label=\textbf{\Alph*)}]
\begin{multicols}{2}
  \item x = -3 o x = \frac{1}{3}
  \item x = \frac{1}{3}
  \item x = 0
  \item x = 2
  \item x = -2
  \item Sin solución
\end{multicols}
\end{enumerate}
\vspace{0.15cm}
\end{enumerate}\n
\subsubsection*{SECCIÓN IV: DESIGUALDADES CON VALOR ABSOLUTO (151–210)}

\begin{enumerate}[label=\textbf{\arabic*.}]
\setcounter{enumi}{150}
\item Resuelve \(|x|<3\).
\begin{enumerate}[label=\textbf{\Alph*)}]
\begin{multicols}{2}
  \item -3 < x < 3
  \item x < 3
  \item x > -3
  \item x = 0
  \item Todos los reales
  \item Sin solución
\end{multicols}
\end{enumerate}
\vspace{0.15cm}
\item Resuelve \(|x|\leq4\).
\begin{enumerate}[label=\textbf{\Alph*)}]
\begin{multicols}{2}
  \item -4 <= x <= 4
  \item x = 3
  \item x = -3
  \item x = 0
  \item Todos los reales
  \item Sin solución
\end{multicols}
\end{enumerate}
\vspace{0.15cm}
\item Resuelve \(|x|>2\).
\begin{enumerate}[label=\textbf{\Alph*)}]
\begin{multicols}{2}
  \item x > 2 o x < -2
  \item x > 3
  \item x < -3
  \item x = 0
  \item Todos los reales
  \item Sin solución
\end{multicols}
\end{enumerate}
\vspace{0.15cm}
\item Resuelve \(|x|\geq5\).
\begin{enumerate}[label=\textbf{\Alph*)}]
\begin{multicols}{2}
  \item x >= 5 o x <= -5
  \item x = 3
  \item x = -3
  \item x = 0
  \item Todos los reales
  \item Sin solución
\end{multicols}
\end{enumerate}
\vspace{0.15cm}
\item Resuelve \(|x|<1\).
\begin{enumerate}[label=\textbf{\Alph*)}]
\begin{multicols}{2}
  \item -1 < x < 1
  \item x < 3
  \item x > -3
  \item x = 0
  \item Todos los reales
  \item Sin solución
\end{multicols}
\end{enumerate}
\vspace{0.15cm}
\item Resuelve \(|x|\leq2\).
\begin{enumerate}[label=\textbf{\Alph*)}]
\begin{multicols}{2}
  \item -2 <= x <= 2
  \item x = 3
  \item x = -3
  \item x = 0
  \item Todos los reales
  \item Sin solución
\end{multicols}
\end{enumerate}
\vspace{0.15cm}
\item Resuelve \(|x|>6\).
\begin{enumerate}[label=\textbf{\Alph*)}]
\begin{multicols}{2}
  \item x > 6 o x < -6
  \item x > 3
  \item x < -3
  \item x = 0
  \item Todos los reales
  \item Sin solución
\end{multicols}
\end{enumerate}
\vspace{0.15cm}
\item Resuelve \(|x|\geq3\).
\begin{enumerate}[label=\textbf{\Alph*)}]
\begin{multicols}{2}
  \item x >= 3 o x <= -3
  \item x = 3
  \item x = -3
  \item x = 0
  \item Todos los reales
  \item Sin solución
\end{multicols}
\end{enumerate}
\vspace{0.15cm}
\item Resuelve \(|x|<7\).
\begin{enumerate}[label=\textbf{\Alph*)}]
\begin{multicols}{2}
  \item -7 < x < 7
  \item x < 3
  \item x > -3
  \item x = 0
  \item Todos los reales
  \item Sin solución
\end{multicols}
\end{enumerate}
\vspace{0.15cm}
\item Resuelve \(|x|\leq8\).
\begin{enumerate}[label=\textbf{\Alph*)}]
\begin{multicols}{2}
  \item -8 <= x <= 8
  \item x = 3
  \item x = -3
  \item x = 0
  \item Todos los reales
  \item Sin solución
\end{multicols}
\end{enumerate}
\vspace{0.15cm}
\item Resuelve \(|x-2|<3\).
\begin{enumerate}[label=\textbf{\Alph*)}]
\begin{multicols}{2}
  \item -1 < x < 5
  \item x < 3
  \item x > -3
  \item x = 0
  \item Todos los reales
  \item Sin solución
\end{multicols}
\end{enumerate}
\vspace{0.15cm}
\item Resuelve \(|x+1|\leq4\).
\begin{enumerate}[label=\textbf{\Alph*)}]
\begin{multicols}{2}
  \item -5 <= x <= 3
  \item x = 3
  \item x = -3
  \item x = 0
  \item Todos los reales
  \item Sin solución
\end{multicols}
\end{enumerate}
\vspace{0.15cm}
\item Resuelve \(|x-3|>2\).
\begin{enumerate}[label=\textbf{\Alph*)}]
\begin{multicols}{2}
  \item x > 5 o x < 1
  \item x > 3
  \item x < -3
  \item x = 0
  \item Todos los reales
  \item Sin solución
\end{multicols}
\end{enumerate}
\vspace{0.15cm}
\item Resuelve \(|x+2|\geq3\).
\begin{enumerate}[label=\textbf{\Alph*)}]
\begin{multicols}{2}
  \item x >= 1 o x <= -5
  \item x = 3
  \item x = -3
  \item x = 0
  \item Todos los reales
  \item Sin solución
\end{multicols}
\end{enumerate}
\vspace{0.15cm}
\item Resuelve \(|x-4|<1\).
\begin{enumerate}[label=\textbf{\Alph*)}]
\begin{multicols}{2}
  \item 3 < x < 5
  \item x < 3
  \item x > -3
  \item x = 0
  \item Todos los reales
  \item Sin solución
\end{multicols}
\end{enumerate}
\vspace{0.15cm}
\item Resuelve \(|x+3|\leq2\).
\begin{enumerate}[label=\textbf{\Alph*)}]
\begin{multicols}{2}
  \item -5 <= x <= -1
  \item x = 3
  \item x = -3
  \item x = 0
  \item Todos los reales
  \item Sin solución
\end{multicols}
\end{enumerate}
\vspace{0.15cm}
\item Resuelve \(|x-1|>4\).
\begin{enumerate}[label=\textbf{\Alph*)}]
\begin{multicols}{2}
  \item x > 5 o x < -3
  \item x > 3
  \item x < -3
  \item x = 0
  \item Todos los reales
  \item Sin solución
\end{multicols}
\end{enumerate}
\vspace{0.15cm}
\item Resuelve \(|x+4|\geq1\).
\begin{enumerate}[label=\textbf{\Alph*)}]
\begin{multicols}{2}
  \item x >= -3 o x <= -5
  \item x = 3
  \item x = -3
  \item x = 0
  \item Todos los reales
  \item Sin solución
\end{multicols}
\end{enumerate}
\vspace{0.15cm}
\item Resuelve \(|x-5|<2\).
\begin{enumerate}[label=\textbf{\Alph*)}]
\begin{multicols}{2}
  \item 3 < x < 7
  \item x < 3
  \item x > -3
  \item x = 0
  \item Todos los reales
  \item Sin solución
\end{multicols}
\end{enumerate}
\vspace{0.15cm}
\item Resuelve \(|x+2|\leq5\).
\begin{enumerate}[label=\textbf{\Alph*)}]
\begin{multicols}{2}
  \item -7 <= x <= 3
  \item x = 3
  \item x = -3
  \item x = 0
  \item Todos los reales
  \item Sin solución
\end{multicols}
\end{enumerate}
\vspace{0.15cm}
\item Resuelve \(|x-6|>1\).
\begin{enumerate}[label=\textbf{\Alph*)}]
\begin{multicols}{2}
  \item x > 7 o x < 5
  \item x > 3
  \item x < -3
  \item x = 0
  \item Todos los reales
  \item Sin solución
\end{multicols}
\end{enumerate}
\vspace{0.15cm}
\item Resuelve \(|x+5|\geq2\).
\begin{enumerate}[label=\textbf{\Alph*)}]
\begin{multicols}{2}
  \item x >= -3 o x <= -7
  \item x = 3
  \item x = -3
  \item x = 0
  \item Todos los reales
  \item Sin solución
\end{multicols}
\end{enumerate}
\vspace{0.15cm}
\item Resuelve \(|x-2|\leq0\).
\begin{enumerate}[label=\textbf{\Alph*)}]
\begin{multicols}{2}
  \item x = 2
  \item x = 3
  \item x = -3
  \item x = 0
  \item Todos los reales
  \item Sin solución
\end{multicols}
\end{enumerate}
\vspace{0.15cm}
\item Resuelve \(|x+1|<0\).
\begin{enumerate}[label=\textbf{\Alph*)}]
\begin{multicols}{2}
  \item Sin solución
  \item x < 3
  \item x > -3
  \item x = 0
  \item Todos los reales
  \item x = 1
\end{multicols}
\end{enumerate}
\vspace{0.15cm}
\item Resuelve \(|x-3|\geq0\).
\begin{enumerate}[label=\textbf{\Alph*)}]
\begin{multicols}{2}
  \item Todos los reales
  \item x = 3
  \item x = -3
  \item x = 0
  \item x = 0
  \item Sin solución
\end{multicols}
\end{enumerate}
\vspace{0.15cm}
\item Resuelve \(|x+2|>0\).
\begin{enumerate}[label=\textbf{\Alph*)}]
\begin{multicols}{2}
  \item Todos los reales excepto x = -2
  \item x > 3
  \item x < -3
  \item x = 0
  \item Todos los reales
  \item Sin solución
\end{multicols}
\end{enumerate}
\vspace{0.15cm}
\item Resuelve \(|2x|<6\).
\begin{enumerate}[label=\textbf{\Alph*)}]
\begin{multicols}{2}
  \item -3 < x < 3
  \item x < 3
  \item x > -3
  \item x = 0
  \item Todos los reales
  \item Sin solución
\end{multicols}
\end{enumerate}
\vspace{0.15cm}
\item Resuelve \(|3x|\leq9\).
\begin{enumerate}[label=\textbf{\Alph*)}]
\begin{multicols}{2}
  \item -3 <= x <= 3
  \item x = 3
  \item x = -3
  \item x = 0
  \item Todos los reales
  \item Sin solución
\end{multicols}
\end{enumerate}
\vspace{0.15cm}
\item Resuelve \(|4x|>8\).
\begin{enumerate}[label=\textbf{\Alph*)}]
\begin{multicols}{2}
  \item x > 2 o x < -2
  \item x > 3
  \item x < -3
  \item x = 0
  \item Todos los reales
  \item Sin solución
\end{multicols}
\end{enumerate}
\vspace{0.15cm}
\item Resuelve \(|5x|\geq10\).
\begin{enumerate}[label=\textbf{\Alph*)}]
\begin{multicols}{2}
  \item x >= 2 o x <= -2
  \item x = 3
  \item x = -3
  \item x = 0
  \item Todos los reales
  \item Sin solución
\end{multicols}
\end{enumerate}
\vspace{0.15cm}
\item Resuelve \(|x-1|<1\).
\begin{enumerate}[label=\textbf{\Alph*)}]
\begin{multicols}{2}
  \item 0 < x < 2
  \item x < 3
  \item x > -3
  \item x = 0
  \item Todos los reales
  \item Sin solución
\end{multicols}
\end{enumerate}
\vspace{0.15cm}
\item Resuelve \(|x+3|<2\).
\begin{enumerate}[label=\textbf{\Alph*)}]
\begin{multicols}{2}
  \item -5 < x < -1
  \item x < 3
  \item x > -3
  \item x = 0
  \item Todos los reales
  \item Sin solución
\end{multicols}
\end{enumerate}
\vspace{0.15cm}
\item Resuelve \(|x-4|>3\).
\begin{enumerate}[label=\textbf{\Alph*)}]
\begin{multicols}{2}
  \item x > 7 o x < 1
  \item x > 3
  \item x < -3
  \item x = 0
  \item Todos los reales
  \item Sin solución
\end{multicols}
\end{enumerate}
\vspace{0.15cm}
\item Resuelve \(|x+2|\geq4\).
\begin{enumerate}[label=\textbf{\Alph*)}]
\begin{multicols}{2}
  \item x >= 2 o x <= -6
  \item x = 3
  \item x = -3
  \item x = 0
  \item Todos los reales
  \item Sin solución
\end{multicols}
\end{enumerate}
\vspace{0.15cm}
\item Resuelve \(|x-5|\leq3\).
\begin{enumerate}[label=\textbf{\Alph*)}]
\begin{multicols}{2}
  \item 2 <= x <= 8
  \item x = 3
  \item x = -3
  \item x = 0
  \item Todos los reales
  \item Sin solución
\end{multicols}
\end{enumerate}
\vspace{0.15cm}
\item Resuelve \(|x+4|>2\).
\begin{enumerate}[label=\textbf{\Alph*)}]
\begin{multicols}{2}
  \item x > -2 o x < -6
  \item x > 3
  \item x < -3
  \item x = 0
  \item Todos los reales
  \item Sin solución
\end{multicols}
\end{enumerate}
\vspace{0.15cm}
\item Resuelve \(|x-6|\leq2\).
\begin{enumerate}[label=\textbf{\Alph*)}]
\begin{multicols}{2}
  \item 4 <= x <= 8
  \item x = 3
  \item x = -3
  \item x = 0
  \item Todos los reales
  \item Sin solución
\end{multicols}
\end{enumerate}
\vspace{0.15cm}
\item Resuelve \(|x+1|>3\).
\begin{enumerate}[label=\textbf{\Alph*)}]
\begin{multicols}{2}
  \item x > 2 o x < -4
  \item x > 3
  \item x < -3
  \item x = 0
  \item Todos los reales
  \item Sin solución
\end{multicols}
\end{enumerate}
\vspace{0.15cm}
\item Resuelve \(|x-3|\geq5\).
\begin{enumerate}[label=\textbf{\Alph*)}]
\begin{multicols}{2}
  \item x >= 8 o x <= -2
  \item x = 3
  \item x = -3
  \item x = 0
  \item Todos los reales
  \item Sin solución
\end{multicols}
\end{enumerate}
\vspace{0.15cm}
\item Resuelve \(|x+5|<1\).
\begin{enumerate}[label=\textbf{\Alph*)}]
\begin{multicols}{2}
  \item -6 < x < -4
  \item x < 3
  \item x > -3
  \item x = 0
  \item Todos los reales
  \item Sin solución
\end{multicols}
\end{enumerate}
\vspace{0.15cm}
\item Resuelve \(|2x-1|<5\).
\begin{enumerate}[label=\textbf{\Alph*)}]
\begin{multicols}{2}
  \item -2 < x < 3
  \item x < 3
  \item x > -3
  \item x = 0
  \item Todos los reales
  \item Sin solución
\end{multicols}
\end{enumerate}
\vspace{0.15cm}
\item Resuelve \(|3x+2|\leq7\).
\begin{enumerate}[label=\textbf{\Alph*)}]
\begin{multicols}{2}
  \item -3 \leq x \leq \frac{5}{3}
  \item x = 3
  \item x = -3
  \item x = 0
  \item Todos los reales
  \item Sin solución
\end{multicols}
\end{enumerate}
\vspace{0.15cm}
\item Resuelve \(|4x-3|>1\).
\begin{enumerate}[label=\textbf{\Alph*)}]
\begin{multicols}{2}
  \item x > 1 o x < \frac{1}{2}
  \item x > 3
  \item x < -3
  \item x = 0
  \item Todos los reales
  \item Sin solución
\end{multicols}
\end{enumerate}
\vspace{0.15cm}
\item Resuelve \(|2x+3|\geq1\).
\begin{enumerate}[label=\textbf{\Alph*)}]
\begin{multicols}{2}
  \item x >= -1 o x <= -2
  \item x = 3
  \item x = -3
  \item x = 0
  \item Todos los reales
  \item Sin solución
\end{multicols}
\end{enumerate}
\vspace{0.15cm}
\item Resuelve \(|3x-1|<2\).
\begin{enumerate}[label=\textbf{\Alph*)}]
\begin{multicols}{2}
  \item -\frac{1}{3} < x < 1
  \item x < 3
  \item x > -3
  \item x = 0
  \item Todos los reales
  \item Sin solución
\end{multicols}
\end{enumerate}
\vspace{0.15cm}
\item Resuelve \(|5x+2|\leq3\).
\begin{enumerate}[label=\textbf{\Alph*)}]
\begin{multicols}{2}
  \item -1 \leq x \leq \frac{1}{5}
  \item x = 3
  \item x = -3
  \item x = 0
  \item Todos los reales
  \item Sin solución
\end{multicols}
\end{enumerate}
\vspace{0.15cm}
\item Resuelve \(|2x-5|>3\).
\begin{enumerate}[label=\textbf{\Alph*)}]
\begin{multicols}{2}
  \item x > 4 o x < 1
  \item x > 3
  \item x < -3
  \item x = 0
  \item Todos los reales
  \item Sin solución
\end{multicols}
\end{enumerate}
\vspace{0.15cm}
\item Resuelve \(|3x+1|\geq5\).
\begin{enumerate}[label=\textbf{\Alph*)}]
\begin{multicols}{2}
  \item x \geq \frac{4}{3} o x \leq -2
  \item x = 3
  \item x = -3
  \item x = 0
  \item Todos los reales
  \item Sin solución
\end{multicols}
\end{enumerate}
\vspace{0.15cm}
\item Resuelve \(|4x-2|<6\).
\begin{enumerate}[label=\textbf{\Alph*)}]
\begin{multicols}{2}
  \item -1 < x < 2
  \item x < 3
  \item x > -3
  \item x = 0
  \item Todos los reales
  \item Sin solución
\end{multicols}
\end{enumerate}
\vspace{0.15cm}
\item Resuelve \(|2x+1|\leq3\).
\begin{enumerate}[label=\textbf{\Alph*)}]
\begin{multicols}{2}
  \item -2 <= x <= 1
  \item x = 3
  \item x = -3
  \item x = 0
  \item Todos los reales
  \item Sin solución
\end{multicols}
\end{enumerate}
\vspace{0.15cm}
\item Resuelve \(|x-2|>0\).
\begin{enumerate}[label=\textbf{\Alph*)}]
\begin{multicols}{2}
  \item x \neq 2
  \item x > 3
  \item x < -3
  \item x = 0
  \item Todos los reales
  \item Sin solución
\end{multicols}
\end{enumerate}
\vspace{0.15cm}
\item Resuelve \(|x+3|\geq0\).
\begin{enumerate}[label=\textbf{\Alph*)}]
\begin{multicols}{2}
  \item Todos los reales
  \item x = 3
  \item x = -3
  \item x = 0
  \item x = 0
  \item Sin solución
\end{multicols}
\end{enumerate}
\vspace{0.15cm}
\item Resuelve \(|x-4|<0\).
\begin{enumerate}[label=\textbf{\Alph*)}]
\begin{multicols}{2}
  \item Sin solución
  \item x < 3
  \item x > -3
  \item x = 0
  \item Todos los reales
  \item x = 1
\end{multicols}
\end{enumerate}
\vspace{0.15cm}
\item Resuelve \(|2x-6|\leq0\).
\begin{enumerate}[label=\textbf{\Alph*)}]
\begin{multicols}{2}
  \item x = 3
  \item x = 3
  \item x = -3
  \item x = 0
  \item Todos los reales
  \item Sin solución
\end{multicols}
\end{enumerate}
\vspace{0.15cm}
\item Resuelve \(|3x+9|<0\).
\begin{enumerate}[label=\textbf{\Alph*)}]
\begin{multicols}{2}
  \item Sin solución
  \item x < 3
  \item x > -3
  \item x = 0
  \item Todos los reales
  \item x = 1
\end{multicols}
\end{enumerate}
\vspace{0.15cm}
\item Resuelve \(|x-1|\geq0\).
\begin{enumerate}[label=\textbf{\Alph*)}]
\begin{multicols}{2}
  \item Todos los reales
  \item x = 3
  \item x = -3
  \item x = 0
  \item x = 0
  \item Sin solución
\end{multicols}
\end{enumerate}
\vspace{0.15cm}
\item Resuelve \(|x+7|>0\).
\begin{enumerate}[label=\textbf{\Alph*)}]
\begin{multicols}{2}
  \item x \neq -7
  \item x > 3
  \item x < -3
  \item x = 0
  \item Todos los reales
  \item Sin solución
\end{multicols}
\end{enumerate}
\vspace{0.15cm}
\item Resuelve \(|4x-8|\leq0\).
\begin{enumerate}[label=\textbf{\Alph*)}]
\begin{multicols}{2}
  \item x = 2
  \item x = 3
  \item x = -3
  \item x = 0
  \item Todos los reales
  \item Sin solución
\end{multicols}
\end{enumerate}
\vspace{0.15cm}
\item Resuelve \(|2x+4|<0\).
\begin{enumerate}[label=\textbf{\Alph*)}]
\begin{multicols}{2}
  \item Sin solución
  \item x < 3
  \item x > -3
  \item x = 0
  \item Todos los reales
  \item x = 1
\end{multicols}
\end{enumerate}
\vspace{0.15cm}
\item Resuelve \(|5x-10|\geq0\).
\begin{enumerate}[label=\textbf{\Alph*)}]
\begin{multicols}{2}
  \item Todos los reales
  \item x = 3
  \item x = -3
  \item x = 0
  \item x = 0
  \item Sin solución
\end{multicols}
\end{enumerate}
\vspace{0.15cm}
\end{enumerate}\n
\subsubsection*{SECCIÓN V: DISTANCIA Y APLICACIONES (211–260)}

\begin{enumerate}[label=\textbf{\arabic*.}]
\setcounter{enumi}{210}
\item ¿Cuál es la distancia entre \(2\) y \(7\)?
\begin{enumerate}[label=\textbf{\Alph*)}]
\begin{multicols}{2}
  \item 5
  \item 1
  \item 0
  \item -1
  \item 5
  \item Sin solución
\end{multicols}
\end{enumerate}
\vspace{0.15cm}
\item ¿Cuál es la distancia entre \(-3\) y \(4\)?
\begin{enumerate}[label=\textbf{\Alph*)}]
\begin{multicols}{2}
  \item 7
  \item 1
  \item 0
  \item -1
  \item 5
  \item Sin solución
\end{multicols}
\end{enumerate}
\vspace{0.15cm}
\item ¿Cuál es la distancia entre \(-5\) y \(-1\)?
\begin{enumerate}[label=\textbf{\Alph*)}]
\begin{multicols}{2}
  \item 4
  \item 1
  \item 0
  \item -1
  \item 5
  \item Sin solución
\end{multicols}
\end{enumerate}
\vspace{0.15cm}
\item ¿Cuál es la distancia entre \(0\) y \(-6\)?
\begin{enumerate}[label=\textbf{\Alph*)}]
\begin{multicols}{2}
  \item 6
  \item 1
  \item 0
  \item -1
  \item 5
  \item Sin solución
\end{multicols}
\end{enumerate}
\vspace{0.15cm}
\item ¿Cuál es la distancia entre \(3\) y \(3\)?
\begin{enumerate}[label=\textbf{\Alph*)}]
\begin{multicols}{2}
  \item 0
  \item 1
  \item 3
  \item -1
  \item 5
  \item Sin solución
\end{multicols}
\end{enumerate}
\vspace{0.15cm}
\item ¿Cuál es la distancia entre \(-2\) y \(5\)?
\begin{enumerate}[label=\textbf{\Alph*)}]
\begin{multicols}{2}
  \item 7
  \item 1
  \item 0
  \item -1
  \item 5
  \item Sin solución
\end{multicols}
\end{enumerate}
\vspace{0.15cm}
\item ¿Cuál es la distancia entre \(1\) y \(-4\)?
\begin{enumerate}[label=\textbf{\Alph*)}]
\begin{multicols}{2}
  \item 5
  \item 1
  \item 0
  \item -1
  \item 5
  \item Sin solución
\end{multicols}
\end{enumerate}
\vspace{0.15cm}
\item ¿Cuál es la distancia entre \(-7\) y \(-2\)?
\begin{enumerate}[label=\textbf{\Alph*)}]
\begin{multicols}{2}
  \item 5
  \item 1
  \item 0
  \item -1
  \item 5
  \item Sin solución
\end{multicols}
\end{enumerate}
\vspace{0.15cm}
\item ¿Cuál es la distancia entre \(8\) y \(2\)?
\begin{enumerate}[label=\textbf{\Alph*)}]
\begin{multicols}{2}
  \item 6
  \item 1
  \item 0
  \item -1
  \item 5
  \item Sin solución
\end{multicols}
\end{enumerate}
\vspace{0.15cm}
\item ¿Cuál es la distancia entre \(-1\) y \(6\)?
\begin{enumerate}[label=\textbf{\Alph*)}]
\begin{multicols}{2}
  \item 7
  \item 1
  \item 0
  \item -1
  \item 5
  \item Sin solución
\end{multicols}
\end{enumerate}
\vspace{0.15cm}
\item La distancia entre a y b se calcula como:
\begin{enumerate}[label=\textbf{\Alph*)}]
\begin{multicols}{2}
  \item |a-b|
  \item 1
  \item 0
  \item -1
  \item 5
  \item Sin solución
\end{multicols}
\end{enumerate}
\vspace{0.15cm}
\item La distancia entre -3 y 2 expresada con valor absoluto es:
\begin{enumerate}[label=\textbf{\Alph*)}]
\begin{multicols}{2}
  \item |-3-2|
  \item 1
  \item 0
  \item -1
  \item 5
  \item Sin solución
\end{multicols}
\end{enumerate}
\vspace{0.15cm}
\item Si |x-3|=5, ¿cuál es la distancia de x a 3?
\begin{enumerate}[label=\textbf{\Alph*)}]
\begin{multicols}{2}
  \item 5
  \item 1
  \item 0
  \item -1
  \item 5
  \item Sin solución
\end{multicols}
\end{enumerate}
\vspace{0.15cm}
\item Si |x+2|=4, ¿cuál es la distancia de x a -2?
\begin{enumerate}[label=\textbf{\Alph*)}]
\begin{multicols}{2}
  \item 4
  \item 1
  \item 0
  \item -1
  \item 5
  \item Sin solución
\end{multicols}
\end{enumerate}
\vspace{0.15cm}
\item Si |x-5|=2, ¿cuáles son los valores de x?
\begin{enumerate}[label=\textbf{\Alph*)}]
\begin{multicols}{2}
  \item 7 y 3
  \item 1
  \item 0
  \item -1
  \item 5
  \item Sin solución
\end{multicols}
\end{enumerate}
\vspace{0.15cm}
\item Si |x+1|=3, ¿cuáles son los valores de x?
\begin{enumerate}[label=\textbf{\Alph*)}]
\begin{multicols}{2}
  \item 2 y -4
  \item 1
  \item 0
  \item -1
  \item 5
  \item Sin solución
\end{multicols}
\end{enumerate}
\vspace{0.15cm}
\item Si |x-4|=1, ¿cuáles son los valores de x?
\begin{enumerate}[label=\textbf{\Alph*)}]
\begin{multicols}{2}
  \item 5 y 3
  \item 1
  \item 0
  \item -1
  \item 5
  \item Sin solución
\end{multicols}
\end{enumerate}
\vspace{0.15cm}
\item Si |x+3|=2, ¿cuáles son los valores de x?
\begin{enumerate}[label=\textbf{\Alph*)}]
\begin{multicols}{2}
  \item -1 y -5
  \item 1
  \item 0
  \item -1
  \item 5
  \item Sin solución
\end{multicols}
\end{enumerate}
\vspace{0.15cm}
\item Si |x-2|=6, ¿cuáles son los valores de x?
\begin{enumerate}[label=\textbf{\Alph*)}]
\begin{multicols}{2}
  \item 8 y -4
  \item 1
  \item 0
  \item -1
  \item 5
  \item Sin solución
\end{multicols}
\end{enumerate}
\vspace{0.15cm}
\item Si |x+5|=1, ¿cuáles son los valores de x?
\begin{enumerate}[label=\textbf{\Alph*)}]
\begin{multicols}{2}
  \item -4 y -6
  \item 1
  \item 0
  \item -1
  \item 5
  \item Sin solución
\end{multicols}
\end{enumerate}
\vspace{0.15cm}
\item ¿Cuál es la distancia entre \(a=3\) y \(b=-2\)?
\begin{enumerate}[label=\textbf{\Alph*)}]
\begin{multicols}{2}
  \item 5
  \item 1
  \item 0
  \item -1
  \item 5
  \item Sin solución
\end{multicols}
\end{enumerate}
\vspace{0.15cm}
\item ¿Cuál es la distancia entre \(a=-4\) y \(b=4\)?
\begin{enumerate}[label=\textbf{\Alph*)}]
\begin{multicols}{2}
  \item 8
  \item 1
  \item 0
  \item -1
  \item 5
  \item Sin solución
\end{multicols}
\end{enumerate}
\vspace{0.15cm}
\item ¿Cuál es la distancia entre \(a=-6\) y \(b=-3\)?
\begin{enumerate}[label=\textbf{\Alph*)}]
\begin{multicols}{2}
  \item 3
  \item 1
  \item 0
  \item -1
  \item 5
  \item Sin solución
\end{multicols}
\end{enumerate}
\vspace{0.15cm}
\item ¿Cuál es la distancia entre \(a=0\) y \(b=-9\)?
\begin{enumerate}[label=\textbf{\Alph*)}]
\begin{multicols}{2}
  \item 9
  \item 1
  \item 0
  \item -1
  \item 5
  \item Sin solución
\end{multicols}
\end{enumerate}
\vspace{0.15cm}
\item ¿Cuál es la distancia entre \(a=5\) y \(b=-5\)?
\begin{enumerate}[label=\textbf{\Alph*)}]
\begin{multicols}{2}
  \item 10
  \item 1
  \item 0
  \item -1
  \item 5
  \item Sin solución
\end{multicols}
\end{enumerate}
\vspace{0.15cm}
\item Si \(|x-7|\leq2\), ¿cuál es el intervalo de x?
\begin{enumerate}[label=\textbf{\Alph*)}]
\begin{multicols}{2}
  \item 5 \leq x \leq 9
  \item 1
  \item 0
  \item -1
  \item 5
  \item Sin solución
\end{multicols}
\end{enumerate}
\vspace{0.15cm}
\item Si \(|x+2|\leq3\), ¿cuál es el intervalo de x?
\begin{enumerate}[label=\textbf{\Alph*)}]
\begin{multicols}{2}
  \item -5 \leq x \leq 1
  \item 1
  \item 0
  \item -1
  \item 5
  \item Sin solución
\end{multicols}
\end{enumerate}
\vspace{0.15cm}
\item Si |x-1|<4, ¿cuál es el intervalo de x?
\begin{enumerate}[label=\textbf{\Alph*)}]
\begin{multicols}{2}
  \item -3 < x < 5
  \item 1
  \item 0
  \item -1
  \item 5
  \item Sin solución
\end{multicols}
\end{enumerate}
\vspace{0.15cm}
\item Si |x+3|>2, ¿cuál es el conjunto solución?
\begin{enumerate}[label=\textbf{\Alph*)}]
\begin{multicols}{2}
  \item x > -1 o x < -5
  \item 1
  \item 0
  \item -1
  \item 5
  \item Sin solución
\end{multicols}
\end{enumerate}
\vspace{0.15cm}
\item Si \(|x-4|\geq3\), ¿cuál es el conjunto solución?
\begin{enumerate}[label=\textbf{\Alph*)}]
\begin{multicols}{2}
  \item x >= 7 o x <= 1
  \item 1
  \item 0
  \item -1
  \item 5
  \item Sin solución
\end{multicols}
\end{enumerate}
\vspace{0.15cm}
\item Un número x está a distancia 3 de 5. ¿Cuáles son los posibles valores de x?
\begin{enumerate}[label=\textbf{\Alph*)}]
\begin{multicols}{2}
  \item 8 y 2
  \item 1
  \item 0
  \item -1
  \item 5
  \item Sin solución
\end{multicols}
\end{enumerate}
\vspace{0.15cm}
\item Un número x está a distancia 4 de -1. ¿Cuáles son los posibles valores de x?
\begin{enumerate}[label=\textbf{\Alph*)}]
\begin{multicols}{2}
  \item 3 y -5
  \item 1
  \item 0
  \item -1
  \item 5
  \item Sin solución
\end{multicols}
\end{enumerate}
\vspace{0.15cm}
\item Un número x está a distancia 2 de 0. ¿Cuáles son los posibles valores de x?
\begin{enumerate}[label=\textbf{\Alph*)}]
\begin{multicols}{2}
  \item 2 y -2
  \item 1
  \item 0
  \item -1
  \item 5
  \item Sin solución
\end{multicols}
\end{enumerate}
\vspace{0.15cm}
\item Un número x está a distancia 6 de 3. ¿Cuáles son los posibles valores de x?
\begin{enumerate}[label=\textbf{\Alph*)}]
\begin{multicols}{2}
  \item 9 y -3
  \item 1
  \item 0
  \item -1
  \item 5
  \item Sin solución
\end{multicols}
\end{enumerate}
\vspace{0.15cm}
\item Un número x está a distancia 1 de -4. ¿Cuáles son los posibles valores de x?
\begin{enumerate}[label=\textbf{\Alph*)}]
\begin{multicols}{2}
  \item -3 y -5
  \item 1
  \item 0
  \item -1
  \item 5
  \item Sin solución
\end{multicols}
\end{enumerate}
\vspace{0.15cm}
\item Si |x-2|<1, ¿cuáles son los valores enteros de x?
\begin{enumerate}[label=\textbf{\Alph*)}]
\begin{multicols}{2}
  \item 2
  \item 1
  \item 0
  \item -1
  \item 5
  \item Sin solución
\end{multicols}
\end{enumerate}
\vspace{0.15cm}
\item Si \(|x+1|\leq2\), ¿cuáles son los valores enteros de x?
\begin{enumerate}[label=\textbf{\Alph*)}]
\begin{multicols}{2}
  \item -3,-2,-1,0,1
  \item 1
  \item 0
  \item -1
  \item 5
  \item Sin solución
\end{multicols}
\end{enumerate}
\vspace{0.15cm}
\item Si \(|x-3|\leq2\), ¿cuáles son los valores enteros de x?
\begin{enumerate}[label=\textbf{\Alph*)}]
\begin{multicols}{2}
  \item 1,2,3,4,5
  \item 1
  \item 0
  \item -1
  \item 5
  \item Sin solución
\end{multicols}
\end{enumerate}
\vspace{0.15cm}
\item Si |x+2|<3, ¿cuáles son los valores enteros de x?
\begin{enumerate}[label=\textbf{\Alph*)}]
\begin{multicols}{2}
  \item -4,-3,-2,-1,0
  \item 1
  \item 0
  \item -1
  \item 5
  \item Sin solución
\end{multicols}
\end{enumerate}
\vspace{0.15cm}
\item Si |x-1|>2, ¿cuáles son los valores enteros de x menores que 5?
\begin{enumerate}[label=\textbf{\Alph*)}]
\begin{multicols}{2}
  \item 4
  \item 1
  \item 0
  \item -1
  \item 5
  \item Sin solución
\end{multicols}
\end{enumerate}
\vspace{0.15cm}
\item La distancia entre dos puntos a y b en la recta real es:
\begin{enumerate}[label=\textbf{\Alph*)}]
\begin{multicols}{2}
  \item |a-b|
  \item 1
  \item 0
  \item -1
  \item 5
  \item Sin solución
\end{multicols}
\end{enumerate}
\vspace{0.15cm}
\item Si |x-5|=0, ¿cuál es el valor de x?
\begin{enumerate}[label=\textbf{\Alph*)}]
\begin{multicols}{2}
  \item 5
  \item 1
  \item 0
  \item -1
  \item 5
  \item Sin solución
\end{multicols}
\end{enumerate}
\vspace{0.15cm}
\item Si |x+3|=0, ¿cuál es el valor de x?
\begin{enumerate}[label=\textbf{\Alph*)}]
\begin{multicols}{2}
  \item -3
  \item 1
  \item 0
  \item -1
  \item 5
  \item Sin solución
\end{multicols}
\end{enumerate}
\vspace{0.15cm}
\item Si |x-2|=|x+4|, ¿cuál es el valor de x?
\begin{enumerate}[label=\textbf{\Alph*)}]
\begin{multicols}{2}
  \item -1
  \item 1
  \item 0
  \item -1
  \item 5
  \item Sin solución
\end{multicols}
\end{enumerate}
\vspace{0.15cm}
\item Si |x+1|=|x-5|, ¿cuál es el valor de x?
\begin{enumerate}[label=\textbf{\Alph*)}]
\begin{multicols}{2}
  \item 2
  \item 1
  \item 0
  \item -1
  \item 5
  \item Sin solución
\end{multicols}
\end{enumerate}
\vspace{0.15cm}
\item Si |2x-1|=|x+2|, ¿cuáles son los valores de x?
\begin{enumerate}[label=\textbf{\Alph*)}]
\begin{multicols}{2}
  \item 3 y -\frac{1}{3}
  \item 1
  \item 0
  \item -1
  \item 5
  \item Sin solución
\end{multicols}
\end{enumerate}
\vspace{0.15cm}
\item Si |3x+1|=|x-3|, ¿cuáles son los valores de x?
\begin{enumerate}[label=\textbf{\Alph*)}]
\begin{multicols}{2}
  \item -2 y \frac{1}{2}
  \item 1
  \item 0
  \item -1
  \item 5
  \item Sin solución
\end{multicols}
\end{enumerate}
\vspace{0.15cm}
\item Si |x-4|=|2x+1|, ¿cuáles son los valores de x?
\begin{enumerate}[label=\textbf{\Alph*)}]
\begin{multicols}{2}
  \item 1 y -5
  \item 1
  \item 0
  \item -1
  \item 5
  \item Sin solución
\end{multicols}
\end{enumerate}
\vspace{0.15cm}
\item Si |x+2|=|3x-1|, ¿cuáles son los valores de x?
\begin{enumerate}[label=\textbf{\Alph*)}]
\begin{multicols}{2}
  \item \frac{3}{2} y -\frac{1}{4}
  \item 1
  \item 0
  \item -1
  \item 5
  \item Sin solución
\end{multicols}
\end{enumerate}
\vspace{0.15cm}
\item Si |2x+3|=|x-2|, ¿cuáles son los valores de x?
\begin{enumerate}[label=\textbf{\Alph*)}]
\begin{multicols}{2}
  \item -\frac{1}{3} y -5
  \item 1
  \item 0
  \item -1
  \item 5
  \item Sin solución
\end{multicols}
\end{enumerate}
\vspace{0.15cm}
\end{enumerate}\n
\subsubsection*{SECCIÓN VI: PROBLEMAS AVANZADOS (261–300)}

\begin{enumerate}[label=\textbf{\arabic*.}]
\setcounter{enumi}{260}
\item ¿Cuál es el valor de \(|3-7|+|2-5|\)?
\begin{enumerate}[label=\textbf{\Alph*)}]
\begin{multicols}{2}
  \item 7
  \item 1
  \item 0
  \item -1
  \item 4
  \item Sin solución
\end{multicols}
\end{enumerate}
\vspace{0.15cm}
\item ¿Cuál es el valor de \(||-3|-|-5||\)?
\begin{enumerate}[label=\textbf{\Alph*)}]
\begin{multicols}{2}
  \item 2
  \item 1
  \item 0
  \item -1
  \item 4
  \item Sin solución
\end{multicols}
\end{enumerate}
\vspace{0.15cm}
\item ¿Cuál es el valor de \(|2|^3+|-3|^2\)?
\begin{enumerate}[label=\textbf{\Alph*)}]
\begin{multicols}{2}
  \item 17
  \item 1
  \item 0
  \item -1
  \item 4
  \item Sin solución
\end{multicols}
\end{enumerate}
\vspace{0.15cm}
\item ¿Cuál es el valor de \(\sqrt{(-7)^2}\)?
\begin{enumerate}[label=\textbf{\Alph*)}]
\begin{multicols}{2}
  \item 7
  \item 1
  \item 0
  \item -1
  \item 4
  \item Sin solución
\end{multicols}
\end{enumerate}
\vspace{0.15cm}
\item ¿Cuál es el valor de \(|(-2)^3|\)?
\begin{enumerate}[label=\textbf{\Alph*)}]
\begin{multicols}{2}
  \item 8
  \item 1
  \item 0
  \item -1
  \item 4
  \item Sin solución
\end{multicols}
\end{enumerate}
\vspace{0.15cm}
\item ¿Cuál es el valor de \(|-2^3|\)?
\begin{enumerate}[label=\textbf{\Alph*)}]
\begin{multicols}{2}
  \item 8
  \item 1
  \item 0
  \item -1
  \item 4
  \item Sin solución
\end{multicols}
\end{enumerate}
\vspace{0.15cm}
\item ¿Cuál es el valor de \(|(-2)^2|\)?
\begin{enumerate}[label=\textbf{\Alph*)}]
\begin{multicols}{2}
  \item 4
  \item 1
  \item 0
  \item -1
  \item 4
  \item Sin solución
\end{multicols}
\end{enumerate}
\vspace{0.15cm}
\item ¿Cuál es el valor de \(|-2^2|\)?
\begin{enumerate}[label=\textbf{\Alph*)}]
\begin{multicols}{2}
  \item 4
  \item 1
  \item 0
  \item -1
  \item 4
  \item Sin solución
\end{multicols}
\end{enumerate}
\vspace{0.15cm}
\item Si \(|x|=3\) y \(|y|=4\), ¿cuál es el valor de \(|x+y|\)?
\begin{enumerate}[label=\textbf{\Alph*)}]
\begin{multicols}{2}
  \item 7 o 1
  \item 1
  \item 0
  \item -1
  \item 4
  \item Sin solución
\end{multicols}
\end{enumerate}
\vspace{0.15cm}
\item Si \(|x|=3\) y \(|y|=4\), ¿cuál es el valor de \(|x\cdot y|\)?
\begin{enumerate}[label=\textbf{\Alph*)}]
\begin{multicols}{2}
  \item 12
  \item 1
  \item 0
  \item -1
  \item 4
  \item Sin solución
\end{multicols}
\end{enumerate}
\vspace{0.15cm}
\item Si \(|x|=5\), ¿cuál es el valor de \(|x-3|\)?
\begin{enumerate}[label=\textbf{\Alph*)}]
\begin{multicols}{2}
  \item 2 u 8
  \item 1
  \item 0
  \item -1
  \item 4
  \item Sin solución
\end{multicols}
\end{enumerate}
\vspace{0.15cm}
\item Si \(|x|=2\), ¿cuál es el valor de \(|x+4|\)?
\begin{enumerate}[label=\textbf{\Alph*)}]
\begin{multicols}{2}
  \item 6 o 2
  \item 1
  \item 0
  \item -1
  \item 4
  \item Sin solución
\end{multicols}
\end{enumerate}
\vspace{0.15cm}
\item Si \(|x|=1\), ¿cuál es el valor de \(|2x-1|\)?
\begin{enumerate}[label=\textbf{\Alph*)}]
\begin{multicols}{2}
  \item 1 o 3
  \item 1
  \item 0
  \item -1
  \item 4
  \item Sin solución
\end{multicols}
\end{enumerate}
\vspace{0.15cm}
\item ¿Cuál es el valor de \(|a|+|b|\) si \(a=-3\) y \(b=5\)?
\begin{enumerate}[label=\textbf{\Alph*)}]
\begin{multicols}{2}
  \item 8
  \item 1
  \item 0
  \item -1
  \item 4
  \item Sin solución
\end{multicols}
\end{enumerate}
\vspace{0.15cm}
\item ¿Cuál es el valor de \(|a+b|\) si \(a=-3\) y \(b=5\)?
\begin{enumerate}[label=\textbf{\Alph*)}]
\begin{multicols}{2}
  \item 2
  \item 1
  \item 0
  \item -1
  \item 4
  \item Sin solución
\end{multicols}
\end{enumerate}
\vspace{0.15cm}
\item ¿Se cumple \(|a+b|\leq|a|+|b|\) para \(a=-3\) y \(b=5\)?
\begin{enumerate}[label=\textbf{\Alph*)}]
\begin{multicols}{2}
  \item Sí, \(2\leq8\)
  \item 1
  \item 0
  \item -1
  \item 4
  \item Sin solución
\end{multicols}
\end{enumerate}
\vspace{0.15cm}
\item ¿Cuál es el valor de \(|a|-|b|\) si \(a=-3\) y \(b=5\)?
\begin{enumerate}[label=\textbf{\Alph*)}]
\begin{multicols}{2}
  \item -2
  \item 1
  \item 0
  \item -1
  \item 4
  \item Sin solución
\end{multicols}
\end{enumerate}
\vspace{0.15cm}
\item ¿Cuál es el valor de \(|a-b|\) si \(a=-3\) y \(b=5\)?
\begin{enumerate}[label=\textbf{\Alph*)}]
\begin{multicols}{2}
  \item 8
  \item 1
  \item 0
  \item -1
  \item 4
  \item Sin solución
\end{multicols}
\end{enumerate}
\vspace{0.15cm}
\item ¿Cuál es el valor de \(|a\cdot b|\) si \(a=-3\) y \(b=5\)?
\begin{enumerate}[label=\textbf{\Alph*)}]
\begin{multicols}{2}
  \item 15
  \item 1
  \item 0
  \item -1
  \item 4
  \item Sin solución
\end{multicols}
\end{enumerate}
\vspace{0.15cm}
\item ¿Cuál es el valor de \(\left|\frac{a}{b}\right|\) si \(a=-3\) y \(b=5\)?
\begin{enumerate}[label=\textbf{\Alph*)}]
\begin{multicols}{2}
  \item \frac{3}{5}
  \item 1
  \item 0
  \item -1
  \item 4
  \item Sin solución
\end{multicols}
\end{enumerate}
\vspace{0.15cm}
\item Resuelve \(|x-2|+|x+1|=5\).
\begin{enumerate}[label=\textbf{\Alph*)}]
\begin{multicols}{2}
  \item x = 3 o x = -2
  \item 1
  \item 0
  \item -1
  \item 4
  \item Sin solución
\end{multicols}
\end{enumerate}
\vspace{0.15cm}
\item Resuelve \(|x-1|+|x-3|=2\).
\begin{enumerate}[label=\textbf{\Alph*)}]
\begin{multicols}{2}
  \item 1 \leq x \leq 3
  \item 1
  \item 0
  \item -1
  \item 4
  \item Sin solución
\end{multicols}
\end{enumerate}
\vspace{0.15cm}
\item Resuelve \(|x+2|+|x-2|=4\).
\begin{enumerate}[label=\textbf{\Alph*)}]
\begin{multicols}{2}
  \item -2 \leq x \leq 2
  \item 1
  \item 0
  \item -1
  \item 4
  \item Sin solución
\end{multicols}
\end{enumerate}
\vspace{0.15cm}
\item Resuelve \(|x-1|+|x+1|=4\).
\begin{enumerate}[label=\textbf{\Alph*)}]
\begin{multicols}{2}
  \item x = 2 o x = -2
  \item 1
  \item 0
  \item -1
  \item 4
  \item Sin solución
\end{multicols}
\end{enumerate}
\vspace{0.15cm}
\item Resuelve \(|x-3|+|x+2|=7\).
\begin{enumerate}[label=\textbf{\Alph*)}]
\begin{multicols}{2}
  \item x = 4 o x = -3
  \item 1
  \item 0
  \item -1
  \item 4
  \item Sin solución
\end{multicols}
\end{enumerate}
\vspace{0.15cm}
\item Resuelve \(|2x-1|+|x+2|=5\).
\begin{enumerate}[label=\textbf{\Alph*)}]
\begin{multicols}{2}
  \item x = \frac{4}{3} o x = -2
  \item 1
  \item 0
  \item -1
  \item 4
  \item Sin solución
\end{multicols}
\end{enumerate}
\vspace{0.15cm}
\item Resuelve \(|x-4|+|x+1|=6\).
\begin{enumerate}[label=\textbf{\Alph*)}]
\begin{multicols}{2}
  \item x = 4.5 o x = -1.5
  \item 1
  \item 0
  \item -1
  \item 4
  \item Sin solución
\end{multicols}
\end{enumerate}
\vspace{0.15cm}
\item Resuelve \(|x+3|+|x-1|=6\).
\begin{enumerate}[label=\textbf{\Alph*)}]
\begin{multicols}{2}
  \item x = 2 o x = -4
  \item 1
  \item 0
  \item -1
  \item 4
  \item Sin solución
\end{multicols}
\end{enumerate}
\vspace{0.15cm}
\item Resuelve \(|x-2|+|x+4|=8\).
\begin{enumerate}[label=\textbf{\Alph*)}]
\begin{multicols}{2}
  \item x = 3 o x = -5
  \item 1
  \item 0
  \item -1
  \item 4
  \item Sin solución
\end{multicols}
\end{enumerate}
\vspace{0.15cm}
\item Resuelve \(|x-1|+|x-5|=4\).
\begin{enumerate}[label=\textbf{\Alph*)}]
\begin{multicols}{2}
  \item 1 \leq x \leq 5
  \item 1
  \item 0
  \item -1
  \item 4
  \item Sin solución
\end{multicols}
\end{enumerate}
\vspace{0.15cm}
\item ¿Para qué valores de \(x\) se cumple \(|x-2|=x-2\)?
\begin{enumerate}[label=\textbf{\Alph*)}]
\begin{multicols}{2}
  \item x \geq 2
  \item 1
  \item 0
  \item -1
  \item 4
  \item Sin solución
\end{multicols}
\end{enumerate}
\vspace{0.15cm}
\item ¿Para qué valores de \(x\) se cumple \(|x+3|=-(x+3)\)?
\begin{enumerate}[label=\textbf{\Alph*)}]
\begin{multicols}{2}
  \item x \leq -3
  \item 1
  \item 0
  \item -1
  \item 4
  \item Sin solución
\end{multicols}
\end{enumerate}
\vspace{0.15cm}
\item ¿Para qué valores de \(x\) se cumple \(|2x-1|=2x-1\)?
\begin{enumerate}[label=\textbf{\Alph*)}]
\begin{multicols}{2}
  \item x \geq \frac{1}{2}
  \item 1
  \item 0
  \item -1
  \item 4
  \item Sin solución
\end{multicols}
\end{enumerate}
\vspace{0.15cm}
\item ¿Para qué valores de \(x\) se cumple \(|3x+6|=-(3x+6)\)?
\begin{enumerate}[label=\textbf{\Alph*)}]
\begin{multicols}{2}
  \item x \leq -2
  \item 1
  \item 0
  \item -1
  \item 4
  \item Sin solución
\end{multicols}
\end{enumerate}
\vspace{0.15cm}
\item ¿Cuál es el valor mínimo de \(|x-2|\)?
\begin{enumerate}[label=\textbf{\Alph*)}]
\begin{multicols}{2}
  \item 0
  \item 1
  \item 3
  \item -1
  \item 4
  \item Sin solución
\end{multicols}
\end{enumerate}
\vspace{0.15cm}
\item ¿Cuál es el valor mínimo de \(|x+3|+2\)?
\begin{enumerate}[label=\textbf{\Alph*)}]
\begin{multicols}{2}
  \item 2
  \item 1
  \item 0
  \item -1
  \item 4
  \item Sin solución
\end{multicols}
\end{enumerate}
\vspace{0.15cm}
\item ¿Cuál es el valor mínimo de \(|x-1|-4\)?
\begin{enumerate}[label=\textbf{\Alph*)}]
\begin{multicols}{2}
  \item -4
  \item 1
  \item 0
  \item -1
  \item 4
  \item Sin solución
\end{multicols}
\end{enumerate}
\vspace{0.15cm}
\item ¿Cuál es el valor mínimo de \(2|x-3|+1\)?
\begin{enumerate}[label=\textbf{\Alph*)}]
\begin{multicols}{2}
  \item 1
  \item 2
  \item 0
  \item -1
  \item 4
  \item Sin solución
\end{multicols}
\end{enumerate}
\vspace{0.15cm}
\item ¿Para qué valor de \(x\) se alcanza el mínimo de \(|x-5|+3\)?
\begin{enumerate}[label=\textbf{\Alph*)}]
\begin{multicols}{2}
  \item x = 5
  \item 1
  \item 0
  \item -1
  \item 4
  \item Sin solución
\end{multicols}
\end{enumerate}
\vspace{0.15cm}
\item ¿Para qué valor de \(x\) se alcanza el mínimo de \(|x+2|-1\)?
\begin{enumerate}[label=\textbf{\Alph*)}]
\begin{multicols}{2}
  \item x = -2
  \item 1
  \item 0
  \item -1
  \item 4
  \item Sin solución
\end{multicols}
\end{enumerate}
\vspace{0.15cm}
\end{enumerate}\n
\vspace{0.4cm}
\begin{tcolorbox}[colback=gray!5!white,colframe=gray!60!black,title=\textbf{SOLUCIONARIO -- 1.4 Valor Absoluto}]
\begin{multicols}{4}
\begin{enumerate}[label=\textbf{\arabic*.}]
  \item \textbf{A}
  \item \textbf{B}
  \item \textbf{B}
  \item \textbf{B}
  \item \textbf{B}
  \item \textbf{A}
  \item \textbf{B}
  \item \textbf{B}
  \item \textbf{B}
  \item \textbf{B}
  \item \textbf{B}
  \item \textbf{B}
  \item \textbf{B}
  \item \textbf{B}
  \item \textbf{B}
  \item \textbf{B}
  \item \textbf{A}
  \item \textbf{B}
  \item \textbf{A}
  \item \textbf{B}
  \item \textbf{B}
  \item \textbf{B}
  \item \textbf{B}
  \item \textbf{B}
  \item \textbf{A}
  \item \textbf{B}
  \item \textbf{C}
  \item \textbf{C}
  \item \textbf{A}
  \item \textbf{C}
  \item \textbf{B}
  \item \textbf{B}
  \item \textbf{B}
  \item \textbf{B}
  \item \textbf{B}
  \item \textbf{A}
  \item \textbf{A}
  \item \textbf{B}
  \item \textbf{B}
  \item \textbf{B}
  \item \textbf{B}
  \item \textbf{B}
  \item \textbf{A}
  \item \textbf{A}
  \item \textbf{A}
  \item \textbf{B}
  \item \textbf{B}
  \item \textbf{A}
  \item \textbf{A}
  \item \textbf{A}
  \item \textbf{A}
  \item \textbf{A}
  \item \textbf{A}
  \item \textbf{A}
  \item \textbf{B}
  \item \textbf{A}
  \item \textbf{A}
  \item \textbf{B}
  \item \textbf{B}
  \item \textbf{A}
  \item \textbf{A}
  \item \textbf{B}
  \item \textbf{B}
  \item \textbf{B}
  \item \textbf{B}
  \item \textbf{A}
  \item \textbf{A}
  \item \textbf{A}
  \item \textbf{A}
  \item \textbf{A}
  \item \textbf{A}
  \item \textbf{A}
  \item \textbf{A}
  \item \textbf{B}
  \item \textbf{A}
  \item \textbf{A}
  \item \textbf{A}
  \item \textbf{A}
  \item \textbf{A}
  \item \textbf{A}
  \item \textbf{C}
  \item \textbf{C}
  \item \textbf{A}
  \item \textbf{C}
  \item \textbf{C}
  \item \textbf{A}
  \item \textbf{C}
  \item \textbf{A}
  \item \textbf{A}
  \item \textbf{A}
  \item \textbf{C}
  \item \textbf{C}
  \item \textbf{A}
  \item \textbf{C}
  \item \textbf{C}
  \item \textbf{C}
  \item \textbf{C}
  \item \textbf{C}
  \item \textbf{C}
  \item \textbf{C}
  \item \textbf{C}
  \item \textbf{C}
  \item \textbf{A}
  \item \textbf{A}
  \item \textbf{A}
  \item \textbf{A}
  \item \textbf{A}
  \item \textbf{A}
  \item \textbf{A}
  \item \textbf{A}
  \item \textbf{A}
  \item \textbf{A}
  \item \textbf{A}
  \item \textbf{A}
  \item \textbf{A}
  \item \textbf{A}
  \item \textbf{A}
  \item \textbf{A}
  \item \textbf{A}
  \item \textbf{A}
  \item \textbf{A}
  \item \textbf{A}
  \item \textbf{A}
  \item \textbf{A}
  \item \textbf{A}
  \item \textbf{A}
  \item \textbf{A}
  \item \textbf{A}
  \item \textbf{A}
  \item \textbf{A}
  \item \textbf{A}
  \item \textbf{A}
  \item \textbf{A}
  \item \textbf{A}
  \item \textbf{A}
  \item \textbf{A}
  \item \textbf{A}
  \item \textbf{A}
  \item \textbf{A}
  \item \textbf{A}
  \item \textbf{A}
  \item \textbf{A}
  \item \textbf{A}
  \item \textbf{A}
  \item \textbf{A}
  \item \textbf{A}
  \item \textbf{A}
  \item \textbf{A}
  \item \textbf{A}
  \item \textbf{A}
  \item \textbf{A}
  \item \textbf{A}
  \item \textbf{A}
  \item \textbf{A}
  \item \textbf{A}
  \item \textbf{A}
  \item \textbf{A}
  \item \textbf{A}
  \item \textbf{A}
  \item \textbf{A}
  \item \textbf{A}
  \item \textbf{A}
  \item \textbf{A}
  \item \textbf{A}
  \item \textbf{A}
  \item \textbf{A}
  \item \textbf{A}
  \item \textbf{A}
  \item \textbf{A}
  \item \textbf{A}
  \item \textbf{A}
  \item \textbf{A}
  \item \textbf{A}
  \item \textbf{A}
  \item \textbf{A}
  \item \textbf{A}
  \item \textbf{A}
  \item \textbf{A}
  \item \textbf{A}
  \item \textbf{A}
  \item \textbf{A}
  \item \textbf{A}
  \item \textbf{A}
  \item \textbf{A}
  \item \textbf{A}
  \item \textbf{A}
  \item \textbf{A}
  \item \textbf{A}
  \item \textbf{A}
  \item \textbf{A}
  \item \textbf{A}
  \item \textbf{A}
  \item \textbf{A}
  \item \textbf{A}
  \item \textbf{A}
  \item \textbf{A}
  \item \textbf{A}
  \item \textbf{A}
  \item \textbf{A}
  \item \textbf{A}
  \item \textbf{A}
  \item \textbf{A}
  \item \textbf{A}
  \item \textbf{A}
  \item \textbf{A}
  \item \textbf{A}
  \item \textbf{A}
  \item \textbf{A}
  \item \textbf{A}
  \item \textbf{A}
  \item \textbf{A}
  \item \textbf{A}
  \item \textbf{A}
  \item \textbf{A}
  \item \textbf{A}
  \item \textbf{A}
  \item \textbf{A}
  \item \textbf{A}
  \item \textbf{A}
  \item \textbf{A}
  \item \textbf{A}
  \item \textbf{A}
  \item \textbf{A}
  \item \textbf{A}
  \item \textbf{A}
  \item \textbf{A}
  \item \textbf{A}
  \item \textbf{A}
  \item \textbf{A}
  \item \textbf{A}
  \item \textbf{A}
  \item \textbf{A}
  \item \textbf{A}
  \item \textbf{A}
  \item \textbf{A}
  \item \textbf{A}
  \item \textbf{A}
  \item \textbf{A}
  \item \textbf{A}
  \item \textbf{A}
  \item \textbf{A}
  \item \textbf{A}
  \item \textbf{A}
  \item \textbf{A}
  \item \textbf{A}
  \item \textbf{A}
  \item \textbf{A}
  \item \textbf{A}
  \item \textbf{A}
  \item \textbf{A}
  \item \textbf{A}
  \item \textbf{A}
  \item \textbf{A}
  \item \textbf{A}
  \item \textbf{A}
  \item \textbf{A}
  \item \textbf{A}
  \item \textbf{A}
  \item \textbf{A}
  \item \textbf{A}
  \item \textbf{A}
  \item \textbf{A}
  \item \textbf{A}
  \item \textbf{A}
  \item \textbf{A}
  \item \textbf{A}
  \item \textbf{A}
  \item \textbf{A}
  \item \textbf{A}
  \item \textbf{A}
  \item \textbf{A}
  \item \textbf{A}
  \item \textbf{A}
  \item \textbf{A}
  \item \textbf{A}
  \item \textbf{A}
  \item \textbf{A}
  \item \textbf{A}
  \item \textbf{A}
  \item \textbf{A}
  \item \textbf{A}
  \item \textbf{A}
  \item \textbf{A}
  \item \textbf{A}
  \item \textbf{A}
  \item \textbf{A}
  \item \textbf{A}
  \item \textbf{A}
  \item \textbf{A}
  \item \textbf{A}
  \item \textbf{A}
  \item \textbf{A}
  \item \textbf{A}
  \item \textbf{A}
  \item \textbf{A}
  \item \textbf{A}
  \item \textbf{A}
  \item \textbf{A}
  \item \textbf{A}
  \item \textbf{A}
\end{enumerate}
\end{multicols}
\end{tcolorbox}
\newpage\n
\subsection{1.5 Álgebra}

\subsubsection*{SECCIÓN I: EXPRESIONES ALGEBRAICAS BÁSICAS (1–40)}

\begin{enumerate}[label=\textbf{\arabic*.}]
\setcounter{enumi}{0}
\item ¿Cuál es el valor numérico de \(2x+3 si x=4\)?
\begin{enumerate}[label=\textbf{\Alph*)}]
\begin{multicols}{2}
  \item 8
  \item 11
  \item 14
  \item 5
  \item 7
  \item 10
\end{multicols}
\end{enumerate}
\vspace{0.15cm}
\item ¿Cuál es el valor numérico de \(3x-5 si x=2\)?
\begin{enumerate}[label=\textbf{\Alph*)}]
\begin{multicols}{2}
  \item 1
  \item 11
  \item 14
  \item 5
  \item 7
  \item 10
\end{multicols}
\end{enumerate}
\vspace{0.15cm}
\item ¿Cuál es el valor numérico de \(x^2+2x si x=3\)?
\begin{enumerate}[label=\textbf{\Alph*)}]
\begin{multicols}{2}
  \item 15
  \item 11
  \item 14
  \item 5
  \item 7
  \item 10
\end{multicols}
\end{enumerate}
\vspace{0.15cm}
\item ¿Cuál es el valor numérico de \(2x^2-1 si x=3\)?
\begin{enumerate}[label=\textbf{\Alph*)}]
\begin{multicols}{2}
  \item 17
  \item 11
  \item 14
  \item 5
  \item 7
  \item 10
\end{multicols}
\end{enumerate}
\vspace{0.15cm}
\item ¿Cuál es el valor numérico de \(\frac{x+1}{2} si x=5\)?
\begin{enumerate}[label=\textbf{\Alph*)}]
\begin{multicols}{2}
  \item 3
  \item 11
  \item 14
  \item 5
  \item 7
  \item 10
\end{multicols}
\end{enumerate}
\vspace{0.15cm}
\item ¿Cuál es el valor numérico de \(x^3 si x=2\)?
\begin{enumerate}[label=\textbf{\Alph*)}]
\begin{multicols}{2}
  \item 8
  \item 8
  \item 14
  \item 5
  \item 7
  \item 10
\end{multicols}
\end{enumerate}
\vspace{0.15cm}
\item ¿Cuál es el valor numérico de \(3x+2y si x=2 e y=3\)?
\begin{enumerate}[label=\textbf{\Alph*)}]
\begin{multicols}{2}
  \item 12
  \item 11
  \item 14
  \item 5
  \item 7
  \item 10
\end{multicols}
\end{enumerate}
\vspace{0.15cm}
\item ¿Cuál es el valor numérico de \(x^2-y^2 si x=5 e y=3\)?
\begin{enumerate}[label=\textbf{\Alph*)}]
\begin{multicols}{2}
  \item 16
  \item 11
  \item 14
  \item 5
  \item 7
  \item 10
\end{multicols}
\end{enumerate}
\vspace{0.15cm}
\item ¿Cuál es el valor numérico de \(2x^2+3x-1 si x=2\)?
\begin{enumerate}[label=\textbf{\Alph*)}]
\begin{multicols}{2}
  \item 13
  \item 11
  \item 14
  \item 5
  \item 7
  \item 10
\end{multicols}
\end{enumerate}
\vspace{0.15cm}
\item ¿Cuál es el valor numérico de \(\frac{x+y}{2} si x=3 e y=7\)?
\begin{enumerate}[label=\textbf{\Alph*)}]
\begin{multicols}{2}
  \item 5
  \item 11
  \item 14
  \item 5
  \item 7
  \item 10
\end{multicols}
\end{enumerate}
\vspace{0.15cm}
\item El término 5x tiene coeficiente:
\begin{enumerate}[label=\textbf{\Alph*)}]
\begin{multicols}{2}
  \item 8
  \item 5
  \item 14
  \item 5
  \item 7
  \item 10
\end{multicols}
\end{enumerate}
\vspace{0.15cm}
\item El término -3x^2 tiene coeficiente:
\begin{enumerate}[label=\textbf{\Alph*)}]
\begin{multicols}{2}
  \item 8
  \item -3
  \item 14
  \item 5
  \item 7
  \item 10
\end{multicols}
\end{enumerate}
\vspace{0.15cm}
\item El grado del polinomio 4x^3+2x-1 es:
\begin{enumerate}[label=\textbf{\Alph*)}]
\begin{multicols}{2}
  \item 8
  \item 11
  \item 3
  \item 5
  \item 7
  \item 10
\end{multicols}
\end{enumerate}
\vspace{0.15cm}
\item El grado del polinomio x^2+5x^4-3 es:
\begin{enumerate}[label=\textbf{\Alph*)}]
\begin{multicols}{2}
  \item 8
  \item 4
  \item 14
  \item 5
  \item 7
  \item 10
\end{multicols}
\end{enumerate}
\vspace{0.15cm}
\item El término independiente de 3x^2+2x-7 es:
\begin{enumerate}[label=\textbf{\Alph*)}]
\begin{multicols}{2}
  \item 8
  \item 11
  \item -7
  \item 5
  \item 7
  \item 10
\end{multicols}
\end{enumerate}
\vspace{0.15cm}
\item ¿Cuántos términos tiene 2x^3-x+5?
\begin{enumerate}[label=\textbf{\Alph*)}]
\begin{multicols}{2}
  \item 8
  \item 11
  \item 3
  \item 5
  \item 7
  \item 10
\end{multicols}
\end{enumerate}
\vspace{0.15cm}
\item El coeficiente principal de 5x^2-3x+1 es:
\begin{enumerate}[label=\textbf{\Alph*)}]
\begin{multicols}{2}
  \item 5
  \item 11
  \item 14
  \item 5
  \item 7
  \item 10
\end{multicols}
\end{enumerate}
\vspace{0.15cm}
\item El grado del monomio -7x^5 es:
\begin{enumerate}[label=\textbf{\Alph*)}]
\begin{multicols}{2}
  \item 8
  \item 5
  \item 14
  \item 5
  \item 7
  \item 10
\end{multicols}
\end{enumerate}
\vspace{0.15cm}
\item ¿Cuál es el resultado de \(3x+5x\)?
\begin{enumerate}[label=\textbf{\Alph*)}]
\begin{multicols}{2}
  \item 8x
  \item 11
  \item 14
  \item 5
  \item 7
  \item 10
\end{multicols}
\end{enumerate}
\vspace{0.15cm}
\item ¿Cuál es el resultado de \(7x^2-2x^2\)?
\begin{enumerate}[label=\textbf{\Alph*)}]
\begin{multicols}{2}
  \item 5x^2
  \item 11
  \item 14
  \item 5
  \item 7
  \item 10
\end{multicols}
\end{enumerate}
\vspace{0.15cm}
\item ¿Cuál es el resultado de \(4x^2\cdot2x^3\)?
\begin{enumerate}[label=\textbf{\Alph*)}]
\begin{multicols}{2}
  \item 8x^5
  \item 11
  \item 14
  \item 5
  \item 7
  \item 10
\end{multicols}
\end{enumerate}
\vspace{0.15cm}
\item ¿Cuál es el resultado de \(\frac{10x^5}{2x^2}\)?
\begin{enumerate}[label=\textbf{\Alph*)}]
\begin{multicols}{2}
  \item 5x^3
  \item 11
  \item 14
  \item 5
  \item 7
  \item 10
\end{multicols}
\end{enumerate}
\vspace{0.15cm}
\item ¿Cuál es el resultado de \((3x^2)^3\)?
\begin{enumerate}[label=\textbf{\Alph*)}]
\begin{multicols}{2}
  \item 27x^6
  \item 11
  \item 14
  \item 5
  \item 7
  \item 10
\end{multicols}
\end{enumerate}
\vspace{0.15cm}
\item ¿Cuál es el resultado de \(2x(3x+4)\)?
\begin{enumerate}[label=\textbf{\Alph*)}]
\begin{multicols}{2}
  \item 6x^2+8x
  \item 11
  \item 14
  \item 5
  \item 7
  \item 10
\end{multicols}
\end{enumerate}
\vspace{0.15cm}
\item ¿Cuál es el resultado de \(-3x(2x-5)\)?
\begin{enumerate}[label=\textbf{\Alph*)}]
\begin{multicols}{2}
  \item -6x^2+15x
  \item 11
  \item 14
  \item 5
  \item 7
  \item 10
\end{multicols}
\end{enumerate}
\vspace{0.15cm}
\item ¿Cuál es el resultado de \((x+2)+(x+3)\)?
\begin{enumerate}[label=\textbf{\Alph*)}]
\begin{multicols}{2}
  \item 2x+5
  \item 11
  \item 14
  \item 5
  \item 7
  \item 10
\end{multicols}
\end{enumerate}
\vspace{0.15cm}
\item ¿Cuál es el resultado de \((3x+1)-(x-4)\)?
\begin{enumerate}[label=\textbf{\Alph*)}]
\begin{multicols}{2}
  \item 2x+5
  \item 11
  \item 14
  \item 5
  \item 7
  \item 10
\end{multicols}
\end{enumerate}
\vspace{0.15cm}
\item ¿Cuál es el resultado de \((x^2+2x-1)+(2x^2-x+3)\)?
\begin{enumerate}[label=\textbf{\Alph*)}]
\begin{multicols}{2}
  \item 3x^2+x+2
  \item 11
  \item 14
  \item 5
  \item 7
  \item 10
\end{multicols}
\end{enumerate}
\vspace{0.15cm}
\item ¿Cuál es el resultado de \((5x^2-3x+2)-(2x^2+x-1)\)?
\begin{enumerate}[label=\textbf{\Alph*)}]
\begin{multicols}{2}
  \item 3x^2-4x+3
  \item 11
  \item 14
  \item 5
  \item 7
  \item 10
\end{multicols}
\end{enumerate}
\vspace{0.15cm}
\item ¿Cuál es el resultado de \(Si P(x)=2x+1 y Q(x)=x-3, P(x)+Q(x) es:\)?
\begin{enumerate}[label=\textbf{\Alph*)}]
\begin{multicols}{2}
  \item 3x-2
  \item 11
  \item 14
  \item 5
  \item 7
  \item 10
\end{multicols}
\end{enumerate}
\vspace{0.15cm}
\item ¿Cuál es el resultado de \(Si P(x)=x^2+1 y Q(x)=x^2-1, P(x)-Q(x) es:\)?
\begin{enumerate}[label=\textbf{\Alph*)}]
\begin{multicols}{2}
  \item 2
  \item 11
  \item 14
  \item 5
  \item 7
  \item 10
\end{multicols}
\end{enumerate}
\vspace{0.15cm}
\item ¿Cuál es el resultado de \(x^3\cdot x^4\)?
\begin{enumerate}[label=\textbf{\Alph*)}]
\begin{multicols}{2}
  \item 8
  \item x^7
  \item 14
  \item 5
  \item 7
  \item 10
\end{multicols}
\end{enumerate}
\vspace{0.15cm}
\item ¿Cuál es el resultado de \(\frac{x^8}{x^3}\)?
\begin{enumerate}[label=\textbf{\Alph*)}]
\begin{multicols}{2}
  \item 8
  \item x^5
  \item 14
  \item 5
  \item 7
  \item 10
\end{multicols}
\end{enumerate}
\vspace{0.15cm}
\item ¿Cuál es el resultado de \((x^4)^2\)?
\begin{enumerate}[label=\textbf{\Alph*)}]
\begin{multicols}{2}
  \item 8
  \item x^8
  \item 14
  \item 5
  \item 7
  \item 10
\end{multicols}
\end{enumerate}
\vspace{0.15cm}
\item ¿Cuál es el resultado de \((2x^3)(-4x^2)\)?
\begin{enumerate}[label=\textbf{\Alph*)}]
\begin{multicols}{2}
  \item -8x^5
  \item 11
  \item 14
  \item 5
  \item 7
  \item 10
\end{multicols}
\end{enumerate}
\vspace{0.15cm}
\item ¿Cuál es el resultado de \(\frac{-12x^6}{4x^2}\)?
\begin{enumerate}[label=\textbf{\Alph*)}]
\begin{multicols}{2}
  \item -3x^4
  \item 11
  \item 14
  \item 5
  \item 7
  \item 10
\end{multicols}
\end{enumerate}
\vspace{0.15cm}
\item ¿Cuál es el resultado de \(3x^2(2x^2-4x+1)\)?
\begin{enumerate}[label=\textbf{\Alph*)}]
\begin{multicols}{2}
  \item 6x^4-12x^3+3x^2
  \item 11
  \item 14
  \item 5
  \item 7
  \item 10
\end{multicols}
\end{enumerate}
\vspace{0.15cm}
\item ¿Cuál es el resultado de \(-2x(x^2-3x+2)\)?
\begin{enumerate}[label=\textbf{\Alph*)}]
\begin{multicols}{2}
  \item -2x^3+6x^2-4x
  \item 11
  \item 14
  \item 5
  \item 7
  \item 10
\end{multicols}
\end{enumerate}
\vspace{0.15cm}
\item ¿Cuál es el valor numérico de \(Si x=-2, el valor de x^2-3x+1 es:\)?
\begin{enumerate}[label=\textbf{\Alph*)}]
\begin{multicols}{2}
  \item 11
  \item 11
  \item 14
  \item 5
  \item 7
  \item 10
\end{multicols}
\end{enumerate}
\vspace{0.15cm}
\item ¿Cuál es el valor numérico de \(Si x=-1, el valor de 2x^3-x^2+4 es:\)?
\begin{enumerate}[label=\textbf{\Alph*)}]
\begin{multicols}{2}
  \item 1
  \item 11
  \item 14
  \item 5
  \item 7
  \item 10
\end{multicols}
\end{enumerate}
\vspace{0.15cm}
\end{enumerate}\n
\subsubsection*{SECCIÓN II: PRODUCTOS NOTABLES (41–90)}

\begin{enumerate}[label=\textbf{\arabic*.}]
\setcounter{enumi}{40}
\item ¿Cuál es el desarrollo de \((x+3)^2\)?
\begin{enumerate}[label=\textbf{\Alph*)}]
\begin{multicols}{2}
  \item x^2+6x+9
  \item x^2-6x-9
  \item x+6x+9
  \item x^2+1
  \item 2x
  \item 0
\end{multicols}
\end{enumerate}
\vspace{0.15cm}
\item ¿Cuál es el desarrollo de \((x-4)^2\)?
\begin{enumerate}[label=\textbf{\Alph*)}]
\begin{multicols}{2}
  \item x^2-8x+16
  \item x^2-8x-16
  \item x-8x+16
  \item x^2+1
  \item 2x
  \item 0
\end{multicols}
\end{enumerate}
\vspace{0.15cm}
\item ¿Cuál es el desarrollo de \((x+5)^2\)?
\begin{enumerate}[label=\textbf{\Alph*)}]
\begin{multicols}{2}
  \item x^2+10x+25
  \item x^2-10x-25
  \item x+10x+25
  \item x^2+1
  \item 2x
  \item 0
\end{multicols}
\end{enumerate}
\vspace{0.15cm}
\item ¿Cuál es el desarrollo de \((x-2)^2\)?
\begin{enumerate}[label=\textbf{\Alph*)}]
\begin{multicols}{2}
  \item x^2-4x+4
  \item x^2-4x-4
  \item x-4x+4
  \item x^2+1
  \item 2x
  \item 0
\end{multicols}
\end{enumerate}
\vspace{0.15cm}
\item ¿Cuál es el desarrollo de \((2x+1)^2\)?
\begin{enumerate}[label=\textbf{\Alph*)}]
\begin{multicols}{2}
  \item 4x^2+4x+1
  \item 4x^2-4x-1
  \item 4x+4x+1
  \item x^2+1
  \item 2x
  \item 0
\end{multicols}
\end{enumerate}
\vspace{0.15cm}
\item ¿Cuál es el desarrollo de \((3x-2)^2\)?
\begin{enumerate}[label=\textbf{\Alph*)}]
\begin{multicols}{2}
  \item 9x^2-12x+4
  \item 9x^2-12x-4
  \item 9x-12x+4
  \item x^2+1
  \item 2x
  \item 0
\end{multicols}
\end{enumerate}
\vspace{0.15cm}
\item ¿Cuál es el desarrollo de \((x+3)(x-3)\)?
\begin{enumerate}[label=\textbf{\Alph*)}]
\begin{multicols}{2}
  \item x^2-9
  \item x^2+9
  \item x-9
  \item x^2+1
  \item 2x
  \item 0
\end{multicols}
\end{enumerate}
\vspace{0.15cm}
\item ¿Cuál es el desarrollo de \((x+5)(x-5)\)?
\begin{enumerate}[label=\textbf{\Alph*)}]
\begin{multicols}{2}
  \item x^2-25
  \item x^2+25
  \item x-25
  \item x^2+1
  \item 2x
  \item 0
\end{multicols}
\end{enumerate}
\vspace{0.15cm}
\item ¿Cuál es el desarrollo de \((2x+3)(2x-3)\)?
\begin{enumerate}[label=\textbf{\Alph*)}]
\begin{multicols}{2}
  \item 4x^2-9
  \item 4x^2+9
  \item 4x-9
  \item x^2+1
  \item 2x
  \item 0
\end{multicols}
\end{enumerate}
\vspace{0.15cm}
\item ¿Cuál es el desarrollo de \((x+2)(x+5)\)?
\begin{enumerate}[label=\textbf{\Alph*)}]
\begin{multicols}{2}
  \item x^2+7x+10
  \item x^2-7x-10
  \item x+7x+10
  \item x^2+1
  \item 2x
  \item 0
\end{multicols}
\end{enumerate}
\vspace{0.15cm}
\item ¿Cuál es el desarrollo de \((x-3)(x-4)\)?
\begin{enumerate}[label=\textbf{\Alph*)}]
\begin{multicols}{2}
  \item x^2-7x+12
  \item x^2-7x-12
  \item x-7x+12
  \item x^2+1
  \item 2x
  \item 0
\end{multicols}
\end{enumerate}
\vspace{0.15cm}
\item ¿Cuál es el desarrollo de \((x+4)(x-2)\)?
\begin{enumerate}[label=\textbf{\Alph*)}]
\begin{multicols}{2}
  \item x^2+2x-8
  \item x^2-2x-8
  \item x+2x-8
  \item x^2+1
  \item 2x
  \item 0
\end{multicols}
\end{enumerate}
\vspace{0.15cm}
\item ¿Cuál es el desarrollo de \((x-5)(x+3)\)?
\begin{enumerate}[label=\textbf{\Alph*)}]
\begin{multicols}{2}
  \item x^2-2x-15
  \item x^2+2x+15
  \item x-2x-15
  \item x^2+1
  \item 2x
  \item 0
\end{multicols}
\end{enumerate}
\vspace{0.15cm}
\item ¿Cuál es el desarrollo de \((2x+1)(x+3)\)?
\begin{enumerate}[label=\textbf{\Alph*)}]
\begin{multicols}{2}
  \item 2x^2+7x+3
  \item 2x^2-7x-3
  \item 2x+7x+3
  \item x^2+1
  \item 2x
  \item 0
\end{multicols}
\end{enumerate}
\vspace{0.15cm}
\item ¿Cuál es el desarrollo de \((3x-2)(x+4)\)?
\begin{enumerate}[label=\textbf{\Alph*)}]
\begin{multicols}{2}
  \item 3x^2+10x-8
  \item 3x^2-10x-8
  \item 3x+10x-8
  \item x^2+1
  \item 2x
  \item 0
\end{multicols}
\end{enumerate}
\vspace{0.15cm}
\item ¿Cuál es el desarrollo de \((x+1)^3\)?
\begin{enumerate}[label=\textbf{\Alph*)}]
\begin{multicols}{2}
  \item x^3+3x^2+3x+1
  \item x^3-3x^2-3x-1
  \item x^3+3x+3x+1
  \item x^2+1
  \item 2x
  \item 0
\end{multicols}
\end{enumerate}
\vspace{0.15cm}
\item ¿Cuál es el desarrollo de \((x-1)^3\)?
\begin{enumerate}[label=\textbf{\Alph*)}]
\begin{multicols}{2}
  \item x^3-3x^2+3x-1
  \item x^3-3x^2-3x-1
  \item x^3-3x+3x-1
  \item x^2+1
  \item 2x
  \item 0
\end{multicols}
\end{enumerate}
\vspace{0.15cm}
\item ¿Cuál es el desarrollo de \((x+2)^3\)?
\begin{enumerate}[label=\textbf{\Alph*)}]
\begin{multicols}{2}
  \item x^3+6x^2+12x+8
  \item x^3-6x^2-12x-8
  \item x^3+6x+12x+8
  \item x^2+1
  \item 2x
  \item 0
\end{multicols}
\end{enumerate}
\vspace{0.15cm}
\item ¿Cuál es el desarrollo de \((x-2)^3\)?
\begin{enumerate}[label=\textbf{\Alph*)}]
\begin{multicols}{2}
  \item x^3-6x^2+12x-8
  \item x^3-6x^2-12x-8
  \item x^3-6x+12x-8
  \item x^2+1
  \item 2x
  \item 0
\end{multicols}
\end{enumerate}
\vspace{0.15cm}
\item ¿Cuál es el desarrollo de \((2x+3)^2\)?
\begin{enumerate}[label=\textbf{\Alph*)}]
\begin{multicols}{2}
  \item 4x^2+12x+9
  \item 4x^2-12x-9
  \item 4x+12x+9
  \item x^2+1
  \item 2x
  \item 0
\end{multicols}
\end{enumerate}
\vspace{0.15cm}
\item ¿Cuál es el desarrollo de \((3x-1)^2\)?
\begin{enumerate}[label=\textbf{\Alph*)}]
\begin{multicols}{2}
  \item 9x^2-6x+1
  \item 9x^2-6x-1
  \item 9x-6x+1
  \item x^2+1
  \item 2x
  \item 0
\end{multicols}
\end{enumerate}
\vspace{0.15cm}
\item ¿Cuál es el desarrollo de \((x+6)^2\)?
\begin{enumerate}[label=\textbf{\Alph*)}]
\begin{multicols}{2}
  \item x^2+12x+36
  \item x^2-12x-36
  \item x+12x+36
  \item x^2+1
  \item 2x
  \item 0
\end{multicols}
\end{enumerate}
\vspace{0.15cm}
\item ¿Cuál es el desarrollo de \((x-7)^2\)?
\begin{enumerate}[label=\textbf{\Alph*)}]
\begin{multicols}{2}
  \item x^2-14x+49
  \item x^2-14x-49
  \item x-14x+49
  \item x^2+1
  \item 2x
  \item 0
\end{multicols}
\end{enumerate}
\vspace{0.15cm}
\item ¿Cuál es el desarrollo de \((x+8)(x-8)\)?
\begin{enumerate}[label=\textbf{\Alph*)}]
\begin{multicols}{2}
  \item x^2-64
  \item x^2+64
  \item x-64
  \item x^2+1
  \item 2x
  \item 0
\end{multicols}
\end{enumerate}
\vspace{0.15cm}
\item ¿Cuál es el desarrollo de \((4x+1)(4x-1)\)?
\begin{enumerate}[label=\textbf{\Alph*)}]
\begin{multicols}{2}
  \item 16x^2-1
  \item 16x^2+1
  \item 16x-1
  \item x^2+1
  \item 2x
  \item 0
\end{multicols}
\end{enumerate}
\vspace{0.15cm}
\item ¿Cuál es el desarrollo de \((x+10)^2\)?
\begin{enumerate}[label=\textbf{\Alph*)}]
\begin{multicols}{2}
  \item x^2+20x+100
  \item x^2-20x-100
  \item x+20x+100
  \item x^2+1
  \item 2x
  \item 0
\end{multicols}
\end{enumerate}
\vspace{0.15cm}
\item ¿Cuál es el desarrollo de \((x-9)^2\)?
\begin{enumerate}[label=\textbf{\Alph*)}]
\begin{multicols}{2}
  \item x^2-18x+81
  \item x^2-18x-81
  \item x-18x+81
  \item x^2+1
  \item 2x
  \item 0
\end{multicols}
\end{enumerate}
\vspace{0.15cm}
\item ¿Cuál es el desarrollo de \((x+1)(x+2)\)?
\begin{enumerate}[label=\textbf{\Alph*)}]
\begin{multicols}{2}
  \item x^2+3x+2
  \item x^2-3x-2
  \item x+3x+2
  \item x^2+1
  \item 2x
  \item 0
\end{multicols}
\end{enumerate}
\vspace{0.15cm}
\item ¿Cuál es el desarrollo de \((x-1)(x-2)\)?
\begin{enumerate}[label=\textbf{\Alph*)}]
\begin{multicols}{2}
  \item x^2-3x+2
  \item x^2-3x-2
  \item x-3x+2
  \item x^2+1
  \item 2x
  \item 0
\end{multicols}
\end{enumerate}
\vspace{0.15cm}
\item ¿Cuál es el desarrollo de \((x+7)(x-2)\)?
\begin{enumerate}[label=\textbf{\Alph*)}]
\begin{multicols}{2}
  \item x^2+5x-14
  \item x^2-5x-14
  \item x+5x-14
  \item x^2+1
  \item 2x
  \item 0
\end{multicols}
\end{enumerate}
\vspace{0.15cm}
\item ¿Cuál es el desarrollo de \((x-6)(x+4)\)?
\begin{enumerate}[label=\textbf{\Alph*)}]
\begin{multicols}{2}
  \item x^2-2x-24
  \item x^2+2x+24
  \item x-2x-24
  \item x^2+1
  \item 2x
  \item 0
\end{multicols}
\end{enumerate}
\vspace{0.15cm}
\item ¿Cuál es el desarrollo de \((2x+5)(2x-5)\)?
\begin{enumerate}[label=\textbf{\Alph*)}]
\begin{multicols}{2}
  \item 4x^2-25
  \item 4x^2+25
  \item 4x-25
  \item x^2+1
  \item 2x
  \item 0
\end{multicols}
\end{enumerate}
\vspace{0.15cm}
\item ¿Cuál es el desarrollo de \((3x+2)(3x-2)\)?
\begin{enumerate}[label=\textbf{\Alph*)}]
\begin{multicols}{2}
  \item 9x^2-4
  \item 9x^2+4
  \item 9x-4
  \item x^2+1
  \item 2x
  \item 0
\end{multicols}
\end{enumerate}
\vspace{0.15cm}
\item ¿Cuál es el desarrollo de \((x+2)(x+2)\)?
\begin{enumerate}[label=\textbf{\Alph*)}]
\begin{multicols}{2}
  \item x^2+4x+4
  \item x^2-4x-4
  \item x+4x+4
  \item x^2+1
  \item 2x
  \item 0
\end{multicols}
\end{enumerate}
\vspace{0.15cm}
\item ¿Cuál es el desarrollo de \((x-3)(x-3)\)?
\begin{enumerate}[label=\textbf{\Alph*)}]
\begin{multicols}{2}
  \item x^2-6x+9
  \item x^2-6x-9
  \item x-6x+9
  \item x^2+1
  \item 2x
  \item 0
\end{multicols}
\end{enumerate}
\vspace{0.15cm}
\item ¿Cuál es el desarrollo de \((a+b)^2\)?
\begin{enumerate}[label=\textbf{\Alph*)}]
\begin{multicols}{2}
  \item a^2+2ab+b^2
  \item a^2-2ab-b^2
  \item x^2
  \item x^2+1
  \item 2x
  \item 0
\end{multicols}
\end{enumerate}
\vspace{0.15cm}
\item ¿Cuál es el desarrollo de \((a-b)^2\)?
\begin{enumerate}[label=\textbf{\Alph*)}]
\begin{multicols}{2}
  \item a^2-2ab+b^2
  \item a^2-2ab-b^2
  \item x^2
  \item x^2+1
  \item 2x
  \item 0
\end{multicols}
\end{enumerate}
\vspace{0.15cm}
\item ¿Cuál es el desarrollo de \((a+b)(a-b)\)?
\begin{enumerate}[label=\textbf{\Alph*)}]
\begin{multicols}{2}
  \item a^2-b^2
  \item a^2+b^2
  \item x^2
  \item x^2+1
  \item 2x
  \item 0
\end{multicols}
\end{enumerate}
\vspace{0.15cm}
\item ¿Cuál es el desarrollo de \((a+b)^3\)?
\begin{enumerate}[label=\textbf{\Alph*)}]
\begin{multicols}{2}
  \item a^3+3a^2b+3ab^2+b^3
  \item a^3-3a^2b-3ab^2-b^3
  \item x^2
  \item x^2+1
  \item 2x
  \item 0
\end{multicols}
\end{enumerate}
\vspace{0.15cm}
\item ¿Cuál es el desarrollo de \((a-b)^3\)?
\begin{enumerate}[label=\textbf{\Alph*)}]
\begin{multicols}{2}
  \item a^3-3a^2b+3ab^2-b^3
  \item a^3-3a^2b-3ab^2-b^3
  \item x^2
  \item x^2+1
  \item 2x
  \item 0
\end{multicols}
\end{enumerate}
\vspace{0.15cm}
\item ¿Cuál es el desarrollo de \((x+2)(x^2-2x+4)\)?
\begin{enumerate}[label=\textbf{\Alph*)}]
\begin{multicols}{2}
  \item x^3+8
  \item x^3-8
  \item x^2
  \item x^2+1
  \item 2x
  \item 0
\end{multicols}
\end{enumerate}
\vspace{0.15cm}
\item ¿Cuál es el desarrollo de \((x-3)(x^2+3x+9)\)?
\begin{enumerate}[label=\textbf{\Alph*)}]
\begin{multicols}{2}
  \item x^3-27
  \item x^3+27
  \item x^2
  \item x^2+1
  \item 2x
  \item 0
\end{multicols}
\end{enumerate}
\vspace{0.15cm}
\item ¿Cuál es el desarrollo de \((x+1)(x^2-x+1)\)?
\begin{enumerate}[label=\textbf{\Alph*)}]
\begin{multicols}{2}
  \item x^3+1
  \item x^3-1
  \item x^2
  \item x^2+1
  \item 2x
  \item 0
\end{multicols}
\end{enumerate}
\vspace{0.15cm}
\item ¿Cuál es el desarrollo de \((x-2)(x^2+2x+4)\)?
\begin{enumerate}[label=\textbf{\Alph*)}]
\begin{multicols}{2}
  \item x^3-8
  \item x^3+8
  \item x^2
  \item x^2+1
  \item 2x
  \item 0
\end{multicols}
\end{enumerate}
\vspace{0.15cm}
\item ¿Cuál es el desarrollo de \((2x+1)(4x^2-2x+1)\)?
\begin{enumerate}[label=\textbf{\Alph*)}]
\begin{multicols}{2}
  \item 8x^3+1
  \item 8x^3-1
  \item x^2
  \item x^2+1
  \item 2x
  \item 0
\end{multicols}
\end{enumerate}
\vspace{0.15cm}
\item ¿Cuál es el desarrollo de \((3x-2)(9x^2+6x+4)\)?
\begin{enumerate}[label=\textbf{\Alph*)}]
\begin{multicols}{2}
  \item 27x^3-8
  \item 27x^3+8
  \item x^2
  \item x^2+1
  \item 2x
  \item 0
\end{multicols}
\end{enumerate}
\vspace{0.15cm}
\item ¿Cuál es el desarrollo de \((x+4)^2-(x-4)^2\)?
\begin{enumerate}[label=\textbf{\Alph*)}]
\begin{multicols}{2}
  \item 16x
  \item 16x
  \item x^2
  \item x^2+1
  \item 2x
  \item 0
\end{multicols}
\end{enumerate}
\vspace{0.15cm}
\item ¿Cuál es el desarrollo de \((x+3)^2+(x-3)^2\)?
\begin{enumerate}[label=\textbf{\Alph*)}]
\begin{multicols}{2}
  \item 2x^2+18
  \item 2x^2-18
  \item 2x+18
  \item x^2+1
  \item 2x
  \item 0
\end{multicols}
\end{enumerate}
\vspace{0.15cm}
\item ¿Cuál es el desarrollo de \((x+5)^2-(x+5)(x-5)\)?
\begin{enumerate}[label=\textbf{\Alph*)}]
\begin{multicols}{2}
  \item 10x+50
  \item 10x-50
  \item x^2
  \item x^2+1
  \item 2x
  \item 0
\end{multicols}
\end{enumerate}
\vspace{0.15cm}
\item ¿Cuál es el desarrollo de \((x-2)^2+(x+2)(x-2)\)?
\begin{enumerate}[label=\textbf{\Alph*)}]
\begin{multicols}{2}
  \item 2x^2-4x-4
  \item 2x^2+4x+4
  \item 2x-4x-4
  \item x^2+1
  \item 2x
  \item 0
\end{multicols}
\end{enumerate}
\vspace{0.15cm}
\end{enumerate}\n
\subsubsection*{SECCIÓN III: FACTORIZACIÓN (91–150)}

\begin{enumerate}[label=\textbf{\arabic*.}]
\setcounter{enumi}{90}
\item Factoriza \(x^2+5x\).
\begin{enumerate}[label=\textbf{\Alph*)}]
\begin{multicols}{2}
  \item x(x+5)
  \item x(x-5)
  \item (x+1)^2
  \item (x+2)^2
  \item (x+3)(x-3)
  \item No factorizable
\end{multicols}
\end{enumerate}
\vspace{0.15cm}
\item Factoriza \(x^2-7x\).
\begin{enumerate}[label=\textbf{\Alph*)}]
\begin{multicols}{2}
  \item x(x-7)
  \item x(x+7)
  \item x(x-1)
  \item (x-1)^2
  \item (x+3)(x-3)
  \item No factorizable
\end{multicols}
\end{enumerate}
\vspace{0.15cm}
\item Factoriza \(3x^2+6x\).
\begin{enumerate}[label=\textbf{\Alph*)}]
\begin{multicols}{2}
  \item 3x(x+2)
  \item 3x(x-2)
  \item (x+1)^2
  \item (x+2)^2
  \item (x+3)(x-3)
  \item No factorizable
\end{multicols}
\end{enumerate}
\vspace{0.15cm}
\item Factoriza \(4x^3-8x^2\).
\begin{enumerate}[label=\textbf{\Alph*)}]
\begin{multicols}{2}
  \item 4x^2(x-2)
  \item 4x^2(x+2)
  \item x(x-1)
  \item (x-1)^2
  \item (x+3)(x-3)
  \item No factorizable
\end{multicols}
\end{enumerate}
\vspace{0.15cm}
\item Factoriza \(5x^2+10x+15\).
\begin{enumerate}[label=\textbf{\Alph*)}]
\begin{multicols}{2}
  \item 5(x^2+2x+3)
  \item 5(x^2-2x-3)
  \item x(x-1)
  \item (x+2)^2
  \item (x+3)(x-3)
  \item No factorizable
\end{multicols}
\end{enumerate}
\vspace{0.15cm}
\item Factoriza \(x^2+6x+9\).
\begin{enumerate}[label=\textbf{\Alph*)}]
\begin{multicols}{2}
  \item (x+3)^2
  \item (x-3)^2
  \item (x+1)^2
  \item (x+2)^2
  \item (x+3)(x-3)
  \item No factorizable
\end{multicols}
\end{enumerate}
\vspace{0.15cm}
\item Factoriza \(x^2-10x+25\).
\begin{enumerate}[label=\textbf{\Alph*)}]
\begin{multicols}{2}
  \item (x-5)^2
  \item (x+5)^2
  \item x(x-1)
  \item (x-1)^2
  \item (x+3)(x-3)
  \item No factorizable
\end{multicols}
\end{enumerate}
\vspace{0.15cm}
\item Factoriza \(x^2+8x+16\).
\begin{enumerate}[label=\textbf{\Alph*)}]
\begin{multicols}{2}
  \item (x+4)^2
  \item (x-4)^2
  \item (x+1)^2
  \item (x+2)^2
  \item (x+3)(x-3)
  \item No factorizable
\end{multicols}
\end{enumerate}
\vspace{0.15cm}
\item Factoriza \(x^2-4x+4\).
\begin{enumerate}[label=\textbf{\Alph*)}]
\begin{multicols}{2}
  \item (x-2)^2
  \item (x+2)^2
  \item x(x-1)
  \item (x-1)^2
  \item (x+3)(x-3)
  \item No factorizable
\end{multicols}
\end{enumerate}
\vspace{0.15cm}
\item Factoriza \(x^2+2x+1\).
\begin{enumerate}[label=\textbf{\Alph*)}]
\begin{multicols}{2}
  \item (x+1)^2
  \item (x-1)^2
  \item (x+1)^2
  \item (x+2)^2
  \item (x+3)(x-3)
  \item No factorizable
\end{multicols}
\end{enumerate}
\vspace{0.15cm}
\item Factoriza \(x^2-1\).
\begin{enumerate}[label=\textbf{\Alph*)}]
\begin{multicols}{2}
  \item (x+1)(x-1)
  \item (x-1)(x-1)
  \item (x+1)^2
  \item (x-1)^2
  \item (x+3)(x-3)
  \item No factorizable
\end{multicols}
\end{enumerate}
\vspace{0.15cm}
\item Factoriza \(x^2-9\).
\begin{enumerate}[label=\textbf{\Alph*)}]
\begin{multicols}{2}
  \item (x+3)(x-3)
  \item (x-3)(x-3)
  \item (x+1)^2
  \item (x-1)^2
  \item (x+3)(x-3)
  \item No factorizable
\end{multicols}
\end{enumerate}
\vspace{0.15cm}
\item Factoriza \(x^2-16\).
\begin{enumerate}[label=\textbf{\Alph*)}]
\begin{multicols}{2}
  \item (x+4)(x-4)
  \item (x-4)(x-4)
  \item (x+1)^2
  \item (x-1)^2
  \item (x+3)(x-3)
  \item No factorizable
\end{multicols}
\end{enumerate}
\vspace{0.15cm}
\item Factoriza \(x^2-25\).
\begin{enumerate}[label=\textbf{\Alph*)}]
\begin{multicols}{2}
  \item (x+5)(x-5)
  \item (x-5)(x-5)
  \item (x+1)^2
  \item (x-1)^2
  \item (x+3)(x-3)
  \item No factorizable
\end{multicols}
\end{enumerate}
\vspace{0.15cm}
\item Factoriza \(4x^2-9\).
\begin{enumerate}[label=\textbf{\Alph*)}]
\begin{multicols}{2}
  \item (2x+3)(2x-3)
  \item (2x-3)(2x-3)
  \item x(x-1)
  \item (x+2)^2
  \item (x+3)(x-3)
  \item No factorizable
\end{multicols}
\end{enumerate}
\vspace{0.15cm}
\item Factoriza \(9x^2-16\).
\begin{enumerate}[label=\textbf{\Alph*)}]
\begin{multicols}{2}
  \item (3x+4)(3x-4)
  \item (3x-4)(3x-4)
  \item x(x-1)
  \item (x+2)^2
  \item (x+3)(x-3)
  \item No factorizable
\end{multicols}
\end{enumerate}
\vspace{0.15cm}
\item Factoriza \(x^2+7x+12\).
\begin{enumerate}[label=\textbf{\Alph*)}]
\begin{multicols}{2}
  \item (x+3)(x+4)
  \item (x-3)(x-4)
  \item (x+1)^2
  \item (x+2)^2
  \item (x+3)(x-3)
  \item No factorizable
\end{multicols}
\end{enumerate}
\vspace{0.15cm}
\item Factoriza \(x^2+8x+15\).
\begin{enumerate}[label=\textbf{\Alph*)}]
\begin{multicols}{2}
  \item (x+3)(x+5)
  \item (x-3)(x-5)
  \item (x+1)^2
  \item (x+2)^2
  \item (x+3)(x-3)
  \item No factorizable
\end{multicols}
\end{enumerate}
\vspace{0.15cm}
\item Factoriza \(x^2+9x+20\).
\begin{enumerate}[label=\textbf{\Alph*)}]
\begin{multicols}{2}
  \item (x+4)(x+5)
  \item (x-4)(x-5)
  \item (x+1)^2
  \item (x+2)^2
  \item (x+3)(x-3)
  \item No factorizable
\end{multicols}
\end{enumerate}
\vspace{0.15cm}
\item Factoriza \(x^2-7x+10\).
\begin{enumerate}[label=\textbf{\Alph*)}]
\begin{multicols}{2}
  \item (x-2)(x-5)
  \item (x+2)(x+5)
  \item x(x-1)
  \item (x-1)^2
  \item (x+3)(x-3)
  \item No factorizable
\end{multicols}
\end{enumerate}
\vspace{0.15cm}
\item Factoriza \(x^2-5x+6\).
\begin{enumerate}[label=\textbf{\Alph*)}]
\begin{multicols}{2}
  \item (x-2)(x-3)
  \item (x+2)(x+3)
  \item x(x-1)
  \item (x-1)^2
  \item (x+3)(x-3)
  \item No factorizable
\end{multicols}
\end{enumerate}
\vspace{0.15cm}
\item Factoriza \(x^2+x-12\).
\begin{enumerate}[label=\textbf{\Alph*)}]
\begin{multicols}{2}
  \item (x+4)(x-3)
  \item (x-4)(x-3)
  \item (x+1)^2
  \item (x-1)^2
  \item (x+3)(x-3)
  \item No factorizable
\end{multicols}
\end{enumerate}
\vspace{0.15cm}
\item Factoriza \(x^2-2x-15\).
\begin{enumerate}[label=\textbf{\Alph*)}]
\begin{multicols}{2}
  \item (x-5)(x+3)
  \item (x-5)(x-3)
  \item (x+1)^2
  \item (x-1)^2
  \item (x+3)(x-3)
  \item No factorizable
\end{multicols}
\end{enumerate}
\vspace{0.15cm}
\item Factoriza \(x^2+3x-10\).
\begin{enumerate}[label=\textbf{\Alph*)}]
\begin{multicols}{2}
  \item (x+5)(x-2)
  \item (x-5)(x-2)
  \item (x+1)^2
  \item (x-1)^2
  \item (x+3)(x-3)
  \item No factorizable
\end{multicols}
\end{enumerate}
\vspace{0.15cm}
\item Factoriza \(x^2-6x-16\).
\begin{enumerate}[label=\textbf{\Alph*)}]
\begin{multicols}{2}
  \item (x-8)(x+2)
  \item (x-8)(x-2)
  \item (x+1)^2
  \item (x-1)^2
  \item (x+3)(x-3)
  \item No factorizable
\end{multicols}
\end{enumerate}
\vspace{0.15cm}
\item Factoriza \(x^2+10x+21\).
\begin{enumerate}[label=\textbf{\Alph*)}]
\begin{multicols}{2}
  \item (x+3)(x+7)
  \item (x-3)(x-7)
  \item (x+1)^2
  \item (x+2)^2
  \item (x+3)(x-3)
  \item No factorizable
\end{multicols}
\end{enumerate}
\vspace{0.15cm}
\item Factoriza \(x^2-11x+24\).
\begin{enumerate}[label=\textbf{\Alph*)}]
\begin{multicols}{2}
  \item (x-3)(x-8)
  \item (x+3)(x+8)
  \item x(x-1)
  \item (x-1)^2
  \item (x+3)(x-3)
  \item No factorizable
\end{multicols}
\end{enumerate}
\vspace{0.15cm}
\item Factoriza \(x^2+2x-8\).
\begin{enumerate}[label=\textbf{\Alph*)}]
\begin{multicols}{2}
  \item (x+4)(x-2)
  \item (x-4)(x-2)
  \item (x+1)^2
  \item (x-1)^2
  \item (x+3)(x-3)
  \item No factorizable
\end{multicols}
\end{enumerate}
\vspace{0.15cm}
\item Factoriza \(x^2-4x-12\).
\begin{enumerate}[label=\textbf{\Alph*)}]
\begin{multicols}{2}
  \item (x-6)(x+2)
  \item (x-6)(x-2)
  \item (x+1)^2
  \item (x-1)^2
  \item (x+3)(x-3)
  \item No factorizable
\end{multicols}
\end{enumerate}
\vspace{0.15cm}
\item Factoriza \(x^2+12x+35\).
\begin{enumerate}[label=\textbf{\Alph*)}]
\begin{multicols}{2}
  \item (x+5)(x+7)
  \item (x-5)(x-7)
  \item (x+1)^2
  \item (x+2)^2
  \item (x+3)(x-3)
  \item No factorizable
\end{multicols}
\end{enumerate}
\vspace{0.15cm}
\item Factoriza \(2x^2+7x+3\).
\begin{enumerate}[label=\textbf{\Alph*)}]
\begin{multicols}{2}
  \item (2x+1)(x+3)
  \item (2x-1)(x-3)
  \item (x+1)^2
  \item (x+2)^2
  \item (x+3)(x-3)
  \item No factorizable
\end{multicols}
\end{enumerate}
\vspace{0.15cm}
\item Factoriza \(3x^2+10x+8\).
\begin{enumerate}[label=\textbf{\Alph*)}]
\begin{multicols}{2}
  \item (3x+4)(x+2)
  \item (3x-4)(x-2)
  \item (x+1)^2
  \item (x+2)^2
  \item (x+3)(x-3)
  \item No factorizable
\end{multicols}
\end{enumerate}
\vspace{0.15cm}
\item Factoriza \(2x^2-5x+3\).
\begin{enumerate}[label=\textbf{\Alph*)}]
\begin{multicols}{2}
  \item (2x-3)(x-1)
  \item (2x+3)(x+1)
  \item x(x-1)
  \item (x-1)^2
  \item (x+3)(x-3)
  \item No factorizable
\end{multicols}
\end{enumerate}
\vspace{0.15cm}
\item Factoriza \(6x^2+7x+2\).
\begin{enumerate}[label=\textbf{\Alph*)}]
\begin{multicols}{2}
  \item (3x+2)(2x+1)
  \item (3x-2)(2x-1)
  \item x(x-1)
  \item (x+2)^2
  \item (x+3)(x-3)
  \item No factorizable
\end{multicols}
\end{enumerate}
\vspace{0.15cm}
\item Factoriza \(4x^2-4x+1\).
\begin{enumerate}[label=\textbf{\Alph*)}]
\begin{multicols}{2}
  \item (2x-1)^2
  \item (2x+1)^2
  \item x(x-1)
  \item (x+2)^2
  \item (x+3)(x-3)
  \item No factorizable
\end{multicols}
\end{enumerate}
\vspace{0.15cm}
\item Factoriza \(9x^2+6x+1\).
\begin{enumerate}[label=\textbf{\Alph*)}]
\begin{multicols}{2}
  \item (3x+1)^2
  \item (3x-1)^2
  \item x(x-1)
  \item (x+2)^2
  \item (x+3)(x-3)
  \item No factorizable
\end{multicols}
\end{enumerate}
\vspace{0.15cm}
\item Factoriza \(25x^2-10x+1\).
\begin{enumerate}[label=\textbf{\Alph*)}]
\begin{multicols}{2}
  \item (5x-1)^2
  \item (5x+1)^2
  \item x(x-1)
  \item (x+2)^2
  \item (x+3)(x-3)
  \item No factorizable
\end{multicols}
\end{enumerate}
\vspace{0.15cm}
\item Factoriza \(16x^2+8x+1\).
\begin{enumerate}[label=\textbf{\Alph*)}]
\begin{multicols}{2}
  \item (4x+1)^2
  \item (4x-1)^2
  \item x(x-1)
  \item (x+2)^2
  \item (x+3)(x-3)
  \item No factorizable
\end{multicols}
\end{enumerate}
\vspace{0.15cm}
\item Factoriza \(x^3+8\).
\begin{enumerate}[label=\textbf{\Alph*)}]
\begin{multicols}{2}
  \item (x+2)(x^2-2x+4)
  \item (x-2)(x^2-2x-4)
  \item (x+1)^2
  \item (x+2)^2
  \item (x+3)(x-3)
  \item No factorizable
\end{multicols}
\end{enumerate}
\vspace{0.15cm}
\item Factoriza \(x^3-27\).
\begin{enumerate}[label=\textbf{\Alph*)}]
\begin{multicols}{2}
  \item (x-3)(x^2+3x+9)
  \item (x-3)(x^2-3x-9)
  \item x(x-1)
  \item (x-1)^2
  \item (x+3)(x-3)
  \item No factorizable
\end{multicols}
\end{enumerate}
\vspace{0.15cm}
\item Factoriza \(x^3+1\).
\begin{enumerate}[label=\textbf{\Alph*)}]
\begin{multicols}{2}
  \item (x+1)(x^2-x+1)
  \item (x-1)(x^2-x-1)
  \item (x+1)^2
  \item (x+2)^2
  \item (x+3)(x-3)
  \item No factorizable
\end{multicols}
\end{enumerate}
\vspace{0.15cm}
\item Factoriza \(x^3-1\).
\begin{enumerate}[label=\textbf{\Alph*)}]
\begin{multicols}{2}
  \item (x-1)(x^2+x+1)
  \item (x-1)(x^2-x-1)
  \item x(x-1)
  \item (x-1)^2
  \item (x+3)(x-3)
  \item No factorizable
\end{multicols}
\end{enumerate}
\vspace{0.15cm}
\item Factoriza \(8x^3+27\).
\begin{enumerate}[label=\textbf{\Alph*)}]
\begin{multicols}{2}
  \item (2x+3)(4x^2-6x+9)
  \item (2x-3)(4x^2-6x-9)
  \item x(x-1)
  \item (x+2)^2
  \item (x+3)(x-3)
  \item No factorizable
\end{multicols}
\end{enumerate}
\vspace{0.15cm}
\item Factoriza \(27x^3-8\).
\begin{enumerate}[label=\textbf{\Alph*)}]
\begin{multicols}{2}
  \item (3x-2)(9x^2+6x+4)
  \item (3x-2)(9x^2-6x-4)
  \item x(x-1)
  \item (x+2)^2
  \item (x+3)(x-3)
  \item No factorizable
\end{multicols}
\end{enumerate}
\vspace{0.15cm}
\item Factoriza \(x^4-16\).
\begin{enumerate}[label=\textbf{\Alph*)}]
\begin{multicols}{2}
  \item (x^2+4)(x+2)(x-2)
  \item (x^2-4)(x-2)(x-2)
  \item (x+1)^2
  \item (x-1)^2
  \item (x+3)(x-3)
  \item No factorizable
\end{multicols}
\end{enumerate}
\vspace{0.15cm}
\item Factoriza \(x^4-1\).
\begin{enumerate}[label=\textbf{\Alph*)}]
\begin{multicols}{2}
  \item (x^2+1)(x+1)(x-1)
  \item (x^2-1)(x-1)(x-1)
  \item (x+1)^2
  \item (x-1)^2
  \item (x+3)(x-3)
  \item No factorizable
\end{multicols}
\end{enumerate}
\vspace{0.15cm}
\item Factoriza \(x^2+5x+6\).
\begin{enumerate}[label=\textbf{\Alph*)}]
\begin{multicols}{2}
  \item (x+2)(x+3)
  \item (x-2)(x-3)
  \item (x+1)^2
  \item (x+2)^2
  \item (x+3)(x-3)
  \item No factorizable
\end{multicols}
\end{enumerate}
\vspace{0.15cm}
\item Factoriza \(x^2-8x+15\).
\begin{enumerate}[label=\textbf{\Alph*)}]
\begin{multicols}{2}
  \item (x-3)(x-5)
  \item (x+3)(x+5)
  \item x(x-1)
  \item (x-1)^2
  \item (x+3)(x-3)
  \item No factorizable
\end{multicols}
\end{enumerate}
\vspace{0.15cm}
\item Factoriza \(x^2+4x-21\).
\begin{enumerate}[label=\textbf{\Alph*)}]
\begin{multicols}{2}
  \item (x+7)(x-3)
  \item (x-7)(x-3)
  \item (x+1)^2
  \item (x-1)^2
  \item (x+3)(x-3)
  \item No factorizable
\end{multicols}
\end{enumerate}
\vspace{0.15cm}
\item Factoriza \(x^2-3x-18\).
\begin{enumerate}[label=\textbf{\Alph*)}]
\begin{multicols}{2}
  \item (x-6)(x+3)
  \item (x-6)(x-3)
  \item (x+1)^2
  \item (x-1)^2
  \item (x+3)(x-3)
  \item No factorizable
\end{multicols}
\end{enumerate}
\vspace{0.15cm}
\item Factoriza \(3x^2-12\).
\begin{enumerate}[label=\textbf{\Alph*)}]
\begin{multicols}{2}
  \item 3(x+2)(x-2)
  \item 3(x-2)(x-2)
  \item (x+1)^2
  \item (x-1)^2
  \item (x+3)(x-3)
  \item No factorizable
\end{multicols}
\end{enumerate}
\vspace{0.15cm}
\item Factoriza \(2x^2-18\).
\begin{enumerate}[label=\textbf{\Alph*)}]
\begin{multicols}{2}
  \item 2(x+3)(x-3)
  \item 2(x-3)(x-3)
  \item (x+1)^2
  \item (x-1)^2
  \item (x+3)(x-3)
  \item No factorizable
\end{multicols}
\end{enumerate}
\vspace{0.15cm}
\item Factoriza \(5x^2-20\).
\begin{enumerate}[label=\textbf{\Alph*)}]
\begin{multicols}{2}
  \item 5(x+2)(x-2)
  \item 5(x-2)(x-2)
  \item (x+1)^2
  \item (x-1)^2
  \item (x+3)(x-3)
  \item No factorizable
\end{multicols}
\end{enumerate}
\vspace{0.15cm}
\item Factoriza \(x^3-4x\).
\begin{enumerate}[label=\textbf{\Alph*)}]
\begin{multicols}{2}
  \item x(x+2)(x-2)
  \item x(x-2)(x-2)
  \item (x+1)^2
  \item (x-1)^2
  \item (x+3)(x-3)
  \item No factorizable
\end{multicols}
\end{enumerate}
\vspace{0.15cm}
\item Factoriza \(x^3-9x\).
\begin{enumerate}[label=\textbf{\Alph*)}]
\begin{multicols}{2}
  \item x(x+3)(x-3)
  \item x(x-3)(x-3)
  \item (x+1)^2
  \item (x-1)^2
  \item (x+3)(x-3)
  \item No factorizable
\end{multicols}
\end{enumerate}
\vspace{0.15cm}
\item Factoriza \(x^3+2x^2-3x\).
\begin{enumerate}[label=\textbf{\Alph*)}]
\begin{multicols}{2}
  \item x(x+3)(x-1)
  \item x(x-3)(x-1)
  \item (x+1)^2
  \item (x-1)^2
  \item (x+3)(x-3)
  \item No factorizable
\end{multicols}
\end{enumerate}
\vspace{0.15cm}
\item Factoriza \(x^3-6x^2+9x\).
\begin{enumerate}[label=\textbf{\Alph*)}]
\begin{multicols}{2}
  \item x(x-3)^2
  \item x(x+3)^2
  \item x(x-1)
  \item (x-1)^2
  \item (x+3)(x-3)
  \item No factorizable
\end{multicols}
\end{enumerate}
\vspace{0.15cm}
\item Factoriza \(2x^3-8x\).
\begin{enumerate}[label=\textbf{\Alph*)}]
\begin{multicols}{2}
  \item 2x(x+2)(x-2)
  \item 2x(x-2)(x-2)
  \item (x+1)^2
  \item (x-1)^2
  \item (x+3)(x-3)
  \item No factorizable
\end{multicols}
\end{enumerate}
\vspace{0.15cm}
\item Factoriza \(x^4-5x^2+4\).
\begin{enumerate}[label=\textbf{\Alph*)}]
\begin{multicols}{2}
  \item (x+1)(x-1)(x+2)(x-2)
  \item (x-1)(x-1)(x-2)(x-2)
  \item (x+1)^2
  \item (x-1)^2
  \item (x+3)(x-3)
  \item No factorizable
\end{multicols}
\end{enumerate}
\vspace{0.15cm}
\item Factoriza \(x^4-13x^2+36\).
\begin{enumerate}[label=\textbf{\Alph*)}]
\begin{multicols}{2}
  \item (x+2)(x-2)(x+3)(x-3)
  \item (x-2)(x-2)(x-3)(x-3)
  \item (x+1)^2
  \item (x-1)^2
  \item (x+3)(x-3)
  \item No factorizable
\end{multicols}
\end{enumerate}
\vspace{0.15cm}
\end{enumerate}\n
\subsubsection*{SECCIÓN IV: ECUACIONES LINEALES (151–200)}

\begin{enumerate}[label=\textbf{\arabic*.}]
\setcounter{enumi}{150}
\item Resuelve \(2x+3=7\).
\begin{enumerate}[label=\textbf{\Alph*)}]
\begin{multicols}{2}
  \item x = 2
  \item x = 5
  \item x = 1
  \item x = -2
  \item x = 0
  \item Sin solución
\end{multicols}
\end{enumerate}
\vspace{0.15cm}
\item Resuelve \(3x-5=10\).
\begin{enumerate}[label=\textbf{\Alph*)}]
\begin{multicols}{2}
  \item x = 5
  \item x = 2
  \item x = 1
  \item x = -2
  \item x = 0
  \item Sin solución
\end{multicols}
\end{enumerate}
\vspace{0.15cm}
\item Resuelve \(x+4=9\).
\begin{enumerate}[label=\textbf{\Alph*)}]
\begin{multicols}{2}
  \item x = 5
  \item x = 2
  \item x = 1
  \item x = -2
  \item x = 0
  \item Sin solución
\end{multicols}
\end{enumerate}
\vspace{0.15cm}
\item Resuelve \(x-7=3\).
\begin{enumerate}[label=\textbf{\Alph*)}]
\begin{multicols}{2}
  \item x = 10
  \item x = 5
  \item x = 1
  \item x = -2
  \item x = 0
  \item Sin solución
\end{multicols}
\end{enumerate}
\vspace{0.15cm}
\item Resuelve \(5x=20\).
\begin{enumerate}[label=\textbf{\Alph*)}]
\begin{multicols}{2}
  \item x = 4
  \item x = 5
  \item x = 1
  \item x = -2
  \item x = 0
  \item Sin solución
\end{multicols}
\end{enumerate}
\vspace{0.15cm}
\item Resuelve \(\frac{x}{3}=4\).
\begin{enumerate}[label=\textbf{\Alph*)}]
\begin{multicols}{2}
  \item x = 12
  \item x = 5
  \item x = 1
  \item x = -2
  \item x = 0
  \item Sin solución
\end{multicols}
\end{enumerate}
\vspace{0.15cm}
\item Resuelve \(2x+1=x+5\).
\begin{enumerate}[label=\textbf{\Alph*)}]
\begin{multicols}{2}
  \item x = 4
  \item x = 5
  \item x = 1
  \item x = -2
  \item x = 0
  \item Sin solución
\end{multicols}
\end{enumerate}
\vspace{0.15cm}
\item Resuelve \(3x-2=2x+4\).
\begin{enumerate}[label=\textbf{\Alph*)}]
\begin{multicols}{2}
  \item x = 6
  \item x = 5
  \item x = 1
  \item x = -2
  \item x = 0
  \item Sin solución
\end{multicols}
\end{enumerate}
\vspace{0.15cm}
\item Resuelve \(4x+3=2x+11\).
\begin{enumerate}[label=\textbf{\Alph*)}]
\begin{multicols}{2}
  \item x = 4
  \item x = 5
  \item x = 1
  \item x = -2
  \item x = 0
  \item Sin solución
\end{multicols}
\end{enumerate}
\vspace{0.15cm}
\item Resuelve \(5x-4=3x+6\).
\begin{enumerate}[label=\textbf{\Alph*)}]
\begin{multicols}{2}
  \item x = 5
  \item x = 2
  \item x = 1
  \item x = -2
  \item x = 0
  \item Sin solución
\end{multicols}
\end{enumerate}
\vspace{0.15cm}
\item Resuelve \(2(x+3)=10\).
\begin{enumerate}[label=\textbf{\Alph*)}]
\begin{multicols}{2}
  \item x = 2
  \item x = 5
  \item x = 1
  \item x = -2
  \item x = 0
  \item Sin solución
\end{multicols}
\end{enumerate}
\vspace{0.15cm}
\item Resuelve \(3(x-2)=9\).
\begin{enumerate}[label=\textbf{\Alph*)}]
\begin{multicols}{2}
  \item x = 5
  \item x = 2
  \item x = 1
  \item x = -2
  \item x = 0
  \item Sin solución
\end{multicols}
\end{enumerate}
\vspace{0.15cm}
\item Resuelve \(4(x+1)=2(x+5)\).
\begin{enumerate}[label=\textbf{\Alph*)}]
\begin{multicols}{2}
  \item x = 3
  \item x = 5
  \item x = 1
  \item x = -2
  \item x = 0
  \item Sin solución
\end{multicols}
\end{enumerate}
\vspace{0.15cm}
\item Resuelve \(5(x-1)=3(x+1)\).
\begin{enumerate}[label=\textbf{\Alph*)}]
\begin{multicols}{2}
  \item x = 4
  \item x = 5
  \item x = 1
  \item x = -2
  \item x = 0
  \item Sin solución
\end{multicols}
\end{enumerate}
\vspace{0.15cm}
\item Resuelve \(2(x+1)+3(x-1)=11\).
\begin{enumerate}[label=\textbf{\Alph*)}]
\begin{multicols}{2}
  \item x = 2
  \item x = 5
  \item x = 1
  \item x = -2
  \item x = 0
  \item Sin solución
\end{multicols}
\end{enumerate}
\vspace{0.15cm}
\item Resuelve \(3(x+2)-2(x-1)=10\).
\begin{enumerate}[label=\textbf{\Alph*)}]
\begin{multicols}{2}
  \item x = 2
  \item x = 5
  \item x = 1
  \item x = -2
  \item x = 0
  \item Sin solución
\end{multicols}
\end{enumerate}
\vspace{0.15cm}
\item Resuelve \(\frac{x}{2}+3=5\).
\begin{enumerate}[label=\textbf{\Alph*)}]
\begin{multicols}{2}
  \item x = 4
  \item x = 5
  \item x = 1
  \item x = -2
  \item x = 0
  \item Sin solución
\end{multicols}
\end{enumerate}
\vspace{0.15cm}
\item Resuelve \(\frac{x}{3}-2=1\).
\begin{enumerate}[label=\textbf{\Alph*)}]
\begin{multicols}{2}
  \item x = 9
  \item x = 5
  \item x = 1
  \item x = -2
  \item x = 0
  \item Sin solución
\end{multicols}
\end{enumerate}
\vspace{0.15cm}
\item Resuelve \(\frac{x+1}{2}=3\).
\begin{enumerate}[label=\textbf{\Alph*)}]
\begin{multicols}{2}
  \item x = 5
  \item x = 2
  \item x = 1
  \item x = -2
  \item x = 0
  \item Sin solución
\end{multicols}
\end{enumerate}
\vspace{0.15cm}
\item Resuelve \(\frac{x-2}{3}=4\).
\begin{enumerate}[label=\textbf{\Alph*)}]
\begin{multicols}{2}
  \item x = 14
  \item x = 5
  \item x = 1
  \item x = -2
  \item x = 0
  \item Sin solución
\end{multicols}
\end{enumerate}
\vspace{0.15cm}
\item Resuelve \(\frac{2x+1}{3}=3\).
\begin{enumerate}[label=\textbf{\Alph*)}]
\begin{multicols}{2}
  \item x = 4
  \item x = 5
  \item x = 1
  \item x = -2
  \item x = 0
  \item Sin solución
\end{multicols}
\end{enumerate}
\vspace{0.15cm}
\item Resuelve \(\frac{3x-1}{2}=4\).
\begin{enumerate}[label=\textbf{\Alph*)}]
\begin{multicols}{2}
  \item x = 3
  \item x = 5
  \item x = 1
  \item x = -2
  \item x = 0
  \item Sin solución
\end{multicols}
\end{enumerate}
\vspace{0.15cm}
\item Resuelve \(2x+3=3x-2\).
\begin{enumerate}[label=\textbf{\Alph*)}]
\begin{multicols}{2}
  \item x = 5
  \item x = 2
  \item x = 1
  \item x = -2
  \item x = 0
  \item Sin solución
\end{multicols}
\end{enumerate}
\vspace{0.15cm}
\item Resuelve \(4x-1=x+8\).
\begin{enumerate}[label=\textbf{\Alph*)}]
\begin{multicols}{2}
  \item x = 3
  \item x = 5
  \item x = 1
  \item x = -2
  \item x = 0
  \item Sin solución
\end{multicols}
\end{enumerate}
\vspace{0.15cm}
\item Resuelve \(6x+2=2x+10\).
\begin{enumerate}[label=\textbf{\Alph*)}]
\begin{multicols}{2}
  \item x = 2
  \item x = 5
  \item x = 1
  \item x = -2
  \item x = 0
  \item Sin solución
\end{multicols}
\end{enumerate}
\vspace{0.15cm}
\item Resuelve \(7x-3=4x+6\).
\begin{enumerate}[label=\textbf{\Alph*)}]
\begin{multicols}{2}
  \item x = 3
  \item x = 5
  \item x = 1
  \item x = -2
  \item x = 0
  \item Sin solución
\end{multicols}
\end{enumerate}
\vspace{0.15cm}
\item Resuelve \(8x+1=5x+10\).
\begin{enumerate}[label=\textbf{\Alph*)}]
\begin{multicols}{2}
  \item x = 3
  \item x = 5
  \item x = 1
  \item x = -2
  \item x = 0
  \item Sin solución
\end{multicols}
\end{enumerate}
\vspace{0.15cm}
\item Resuelve \(9x-4=6x+5\).
\begin{enumerate}[label=\textbf{\Alph*)}]
\begin{multicols}{2}
  \item x = 3
  \item x = 5
  \item x = 1
  \item x = -2
  \item x = 0
  \item Sin solución
\end{multicols}
\end{enumerate}
\vspace{0.15cm}
\item Resuelve \(10x+3=7x+12\).
\begin{enumerate}[label=\textbf{\Alph*)}]
\begin{multicols}{2}
  \item x = 3
  \item x = 5
  \item x = 1
  \item x = -2
  \item x = 0
  \item Sin solución
\end{multicols}
\end{enumerate}
\vspace{0.15cm}
\item Resuelve \(2x-5=7\).
\begin{enumerate}[label=\textbf{\Alph*)}]
\begin{multicols}{2}
  \item x = 6
  \item x = 5
  \item x = 1
  \item x = -2
  \item x = 0
  \item Sin solución
\end{multicols}
\end{enumerate}
\vspace{0.15cm}
\item Resuelve \(-3x=12\).
\begin{enumerate}[label=\textbf{\Alph*)}]
\begin{multicols}{2}
  \item x = -4
  \item x = 5
  \item x = 1
  \item x = -2
  \item x = 0
  \item Sin solución
\end{multicols}
\end{enumerate}
\vspace{0.15cm}
\item Resuelve \(-2x+5=11\).
\begin{enumerate}[label=\textbf{\Alph*)}]
\begin{multicols}{2}
  \item x = -3
  \item x = 5
  \item x = 1
  \item x = -2
  \item x = 0
  \item Sin solución
\end{multicols}
\end{enumerate}
\vspace{0.15cm}
\item Resuelve \(-4x-1=7\).
\begin{enumerate}[label=\textbf{\Alph*)}]
\begin{multicols}{2}
  \item x = -2
  \item x = 5
  \item x = 1
  \item x = -2
  \item x = 0
  \item Sin solución
\end{multicols}
\end{enumerate}
\vspace{0.15cm}
\item Resuelve \(-5x+2=-8\).
\begin{enumerate}[label=\textbf{\Alph*)}]
\begin{multicols}{2}
  \item x = 2
  \item x = 5
  \item x = 1
  \item x = -2
  \item x = 0
  \item Sin solución
\end{multicols}
\end{enumerate}
\vspace{0.15cm}
\item Resuelve \(-x+3=7\).
\begin{enumerate}[label=\textbf{\Alph*)}]
\begin{multicols}{2}
  \item x = -4
  \item x = 5
  \item x = 1
  \item x = -2
  \item x = 0
  \item Sin solución
\end{multicols}
\end{enumerate}
\vspace{0.15cm}
\item Resuelve \(-2(x+1)=6\).
\begin{enumerate}[label=\textbf{\Alph*)}]
\begin{multicols}{2}
  \item x = -4
  \item x = 5
  \item x = 1
  \item x = -2
  \item x = 0
  \item Sin solución
\end{multicols}
\end{enumerate}
\vspace{0.15cm}
\item Resuelve \(-3(x-2)=9\).
\begin{enumerate}[label=\textbf{\Alph*)}]
\begin{multicols}{2}
  \item x = -1
  \item x = 5
  \item x = 1
  \item x = -2
  \item x = 0
  \item Sin solución
\end{multicols}
\end{enumerate}
\vspace{0.15cm}
\item Resuelve \(2(x-3)=-8\).
\begin{enumerate}[label=\textbf{\Alph*)}]
\begin{multicols}{2}
  \item x = -1
  \item x = 5
  \item x = 1
  \item x = -2
  \item x = 0
  \item Sin solución
\end{multicols}
\end{enumerate}
\vspace{0.15cm}
\item Resuelve \(5(2-x)=10\).
\begin{enumerate}[label=\textbf{\Alph*)}]
\begin{multicols}{2}
  \item x = 0
  \item x = 5
  \item x = 1
  \item x = -2
  \item x = 0
  \item Sin solución
\end{multicols}
\end{enumerate}
\vspace{0.15cm}
\item Resuelve \(3(4-x)=6\).
\begin{enumerate}[label=\textbf{\Alph*)}]
\begin{multicols}{2}
  \item x = 2
  \item x = 5
  \item x = 1
  \item x = -2
  \item x = 0
  \item Sin solución
\end{multicols}
\end{enumerate}
\vspace{0.15cm}
\item Resuelve \(\frac{2x}{3}=4\).
\begin{enumerate}[label=\textbf{\Alph*)}]
\begin{multicols}{2}
  \item x = 6
  \item x = 5
  \item x = 1
  \item x = -2
  \item x = 0
  \item Sin solución
\end{multicols}
\end{enumerate}
\vspace{0.15cm}
\item Resuelve \(\frac{3x}{4}=6\).
\begin{enumerate}[label=\textbf{\Alph*)}]
\begin{multicols}{2}
  \item x = 8
  \item x = 5
  \item x = 1
  \item x = -2
  \item x = 0
  \item Sin solución
\end{multicols}
\end{enumerate}
\vspace{0.15cm}
\item Resuelve \(\frac{5x}{2}=10\).
\begin{enumerate}[label=\textbf{\Alph*)}]
\begin{multicols}{2}
  \item x = 4
  \item x = 5
  \item x = 1
  \item x = -2
  \item x = 0
  \item Sin solución
\end{multicols}
\end{enumerate}
\vspace{0.15cm}
\item Resuelve \(\frac{x}{2}+\frac{x}{3}=5\).
\begin{enumerate}[label=\textbf{\Alph*)}]
\begin{multicols}{2}
  \item x = 6
  \item x = 5
  \item x = 1
  \item x = -2
  \item x = 0
  \item Sin solución
\end{multicols}
\end{enumerate}
\vspace{0.15cm}
\item Resuelve \(\frac{x}{4}+\frac{x}{2}=3\).
\begin{enumerate}[label=\textbf{\Alph*)}]
\begin{multicols}{2}
  \item x = 4
  \item x = 5
  \item x = 1
  \item x = -2
  \item x = 0
  \item Sin solución
\end{multicols}
\end{enumerate}
\vspace{0.15cm}
\item Resuelve \(\frac{x+1}{2}+\frac{x-1}{3}=2\).
\begin{enumerate}[label=\textbf{\Alph*)}]
\begin{multicols}{2}
  \item x = \frac{11}{5}
  \item x = 5
  \item x = 1
  \item x = -2
  \item x = 0
  \item Sin solución
\end{multicols}
\end{enumerate}
\vspace{0.15cm}
\item Resuelve \(\frac{2x-1}{3}+\frac{x+2}{4}=3\).
\begin{enumerate}[label=\textbf{\Alph*)}]
\begin{multicols}{2}
  \item x = \frac{34}{11}
  \item x = 5
  \item x = 1
  \item x = -2
  \item x = 0
  \item Sin solución
\end{multicols}
\end{enumerate}
\vspace{0.15cm}
\item Resuelve \(\frac{x}{2}-\frac{x}{3}=1\).
\begin{enumerate}[label=\textbf{\Alph*)}]
\begin{multicols}{2}
  \item x = 6
  \item x = 5
  \item x = 1
  \item x = -2
  \item x = 0
  \item Sin solución
\end{multicols}
\end{enumerate}
\vspace{0.15cm}
\item Resuelve \(\frac{x}{3}-\frac{x}{4}=2\).
\begin{enumerate}[label=\textbf{\Alph*)}]
\begin{multicols}{2}
  \item x = 24
  \item x = 5
  \item x = 1
  \item x = -2
  \item x = 0
  \item Sin solución
\end{multicols}
\end{enumerate}
\vspace{0.15cm}
\item Resuelve \(\frac{2x}{5}-\frac{x}{2}=1\).
\begin{enumerate}[label=\textbf{\Alph*)}]
\begin{multicols}{2}
  \item x = -10
  \item x = 5
  \item x = 1
  \item x = -2
  \item x = 0
  \item Sin solución
\end{multicols}
\end{enumerate}
\vspace{0.15cm}
\end{enumerate}\n
\subsubsection*{SECCIÓN V: SISTEMAS DE ECUACIONES (201–250)}

\begin{enumerate}[label=\textbf{\arabic*.}]
\setcounter{enumi}{200}
\item Resuelve el sistema: \(x+y=7, x-y=3\).
\begin{enumerate}[label=\textbf{\Alph*)}]
\begin{multicols}{2}
  \item x=5, y=2
  \item x=2, y=4
  \item x=4, y=1
  \item x=1, y=5
  \item x=5, y=0
  \item x=0, y=5
\end{multicols}
\end{enumerate}
\vspace{0.15cm}
\item Resuelve el sistema: \(x+y=10, x-y=4\).
\begin{enumerate}[label=\textbf{\Alph*)}]
\begin{multicols}{2}
  \item x=7, y=3
  \item x=2, y=4
  \item x=4, y=1
  \item x=1, y=5
  \item x=5, y=0
  \item x=0, y=5
\end{multicols}
\end{enumerate}
\vspace{0.15cm}
\item Resuelve el sistema: \(x+y=5, x-y=1\).
\begin{enumerate}[label=\textbf{\Alph*)}]
\begin{multicols}{2}
  \item x=3, y=2
  \item x=2, y=4
  \item x=4, y=1
  \item x=1, y=5
  \item x=5, y=0
  \item x=0, y=5
\end{multicols}
\end{enumerate}
\vspace{0.15cm}
\item Resuelve el sistema: \(2x+y=8, x-y=1\).
\begin{enumerate}[label=\textbf{\Alph*)}]
\begin{multicols}{2}
  \item x=3, y=2
  \item x=2, y=4
  \item x=4, y=1
  \item x=1, y=5
  \item x=5, y=0
  \item x=0, y=5
\end{multicols}
\end{enumerate}
\vspace{0.15cm}
\item Resuelve el sistema: \(3x+y=10, x+y=4\).
\begin{enumerate}[label=\textbf{\Alph*)}]
\begin{multicols}{2}
  \item x=3, y=1
  \item x=2, y=4
  \item x=4, y=1
  \item x=1, y=5
  \item x=5, y=0
  \item x=0, y=5
\end{multicols}
\end{enumerate}
\vspace{0.15cm}
\item Resuelve el sistema: \(x+2y=8, x-y=2\).
\begin{enumerate}[label=\textbf{\Alph*)}]
\begin{multicols}{2}
  \item x=4, y=2
  \item x=2, y=4
  \item x=4, y=1
  \item x=1, y=5
  \item x=5, y=0
  \item x=0, y=5
\end{multicols}
\end{enumerate}
\vspace{0.15cm}
\item Resuelve el sistema: \(2x+3y=12, x-y=1\).
\begin{enumerate}[label=\textbf{\Alph*)}]
\begin{multicols}{2}
  \item x=3, y=2
  \item x=2, y=4
  \item x=4, y=1
  \item x=1, y=5
  \item x=5, y=0
  \item x=0, y=5
\end{multicols}
\end{enumerate}
\vspace{0.15cm}
\item Resuelve el sistema: \(3x+2y=16, x+y=6\).
\begin{enumerate}[label=\textbf{\Alph*)}]
\begin{multicols}{2}
  \item x=4, y=2
  \item x=2, y=4
  \item x=4, y=1
  \item x=1, y=5
  \item x=5, y=0
  \item x=0, y=5
\end{multicols}
\end{enumerate}
\vspace{0.15cm}
\item Resuelve el sistema: \(4x+y=14, 2x-y=4\).
\begin{enumerate}[label=\textbf{\Alph*)}]
\begin{multicols}{2}
  \item x=3, y=2
  \item x=2, y=4
  \item x=4, y=1
  \item x=1, y=5
  \item x=5, y=0
  \item x=0, y=5
\end{multicols}
\end{enumerate}
\vspace{0.15cm}
\item Resuelve el sistema: \(5x+2y=16, x+y=5\).
\begin{enumerate}[label=\textbf{\Alph*)}]
\begin{multicols}{2}
  \item x=2, y=3
  \item x=2, y=4
  \item x=4, y=1
  \item x=1, y=5
  \item x=5, y=0
  \item x=0, y=5
\end{multicols}
\end{enumerate}
\vspace{0.15cm}
\item Resuelve el sistema: \(x+y=6, 2x-y=3\).
\begin{enumerate}[label=\textbf{\Alph*)}]
\begin{multicols}{2}
  \item x=3, y=3
  \item x=2, y=4
  \item x=4, y=1
  \item x=1, y=5
  \item x=5, y=0
  \item x=0, y=5
\end{multicols}
\end{enumerate}
\vspace{0.15cm}
\item Resuelve el sistema: \(3x-y=5, x+y=3\).
\begin{enumerate}[label=\textbf{\Alph*)}]
\begin{multicols}{2}
  \item x=2, y=1
  \item x=2, y=4
  \item x=4, y=1
  \item x=1, y=5
  \item x=5, y=0
  \item x=0, y=5
\end{multicols}
\end{enumerate}
\vspace{0.15cm}
\item Resuelve el sistema: \(2x+y=7, x+2y=8\).
\begin{enumerate}[label=\textbf{\Alph*)}]
\begin{multicols}{2}
  \item x=2, y=3
  \item x=2, y=4
  \item x=4, y=1
  \item x=1, y=5
  \item x=5, y=0
  \item x=0, y=5
\end{multicols}
\end{enumerate}
\vspace{0.15cm}
\item Resuelve el sistema: \(3x+2y=13, 2x-y=4\).
\begin{enumerate}[label=\textbf{\Alph*)}]
\begin{multicols}{2}
  \item x=3, y=2
  \item x=2, y=4
  \item x=4, y=1
  \item x=1, y=5
  \item x=5, y=0
  \item x=0, y=5
\end{multicols}
\end{enumerate}
\vspace{0.15cm}
\item Resuelve el sistema: \(x+3y=10, 2x-y=6\).
\begin{enumerate}[label=\textbf{\Alph*)}]
\begin{multicols}{2}
  \item x=4, y=2
  \item x=2, y=4
  \item x=4, y=1
  \item x=1, y=5
  \item x=5, y=0
  \item x=0, y=5
\end{multicols}
\end{enumerate}
\vspace{0.15cm}
\item Resuelve el sistema: \(4x+3y=18, 2x+y=8\).
\begin{enumerate}[label=\textbf{\Alph*)}]
\begin{multicols}{2}
  \item x=3, y=2
  \item x=2, y=4
  \item x=4, y=1
  \item x=1, y=5
  \item x=5, y=0
  \item x=0, y=5
\end{multicols}
\end{enumerate}
\vspace{0.15cm}
\item Resuelve el sistema: \(5x-2y=11, x+y=5\).
\begin{enumerate}[label=\textbf{\Alph*)}]
\begin{multicols}{2}
  \item x=3, y=2
  \item x=2, y=4
  \item x=4, y=1
  \item x=1, y=5
  \item x=5, y=0
  \item x=0, y=5
\end{multicols}
\end{enumerate}
\vspace{0.15cm}
\item Resuelve el sistema: \(2x+5y=16, x-y=1\).
\begin{enumerate}[label=\textbf{\Alph*)}]
\begin{multicols}{2}
  \item x=3, y=2
  \item x=2, y=4
  \item x=4, y=1
  \item x=1, y=5
  \item x=5, y=0
  \item x=0, y=5
\end{multicols}
\end{enumerate}
\vspace{0.15cm}
\item Resuelve el sistema: \(3x+4y=18, x+y=5\).
\begin{enumerate}[label=\textbf{\Alph*)}]
\begin{multicols}{2}
  \item x=2, y=3
  \item x=2, y=4
  \item x=4, y=1
  \item x=1, y=5
  \item x=5, y=0
  \item x=0, y=5
\end{multicols}
\end{enumerate}
\vspace{0.15cm}
\item Resuelve el sistema: \(x+4y=14, 2x-y=5\).
\begin{enumerate}[label=\textbf{\Alph*)}]
\begin{multicols}{2}
  \item x=3, y=2
  \item x=2, y=4
  \item x=4, y=1
  \item x=1, y=5
  \item x=5, y=0
  \item x=0, y=5
\end{multicols}
\end{enumerate}
\vspace{0.15cm}
\item Resuelve el sistema: \(y=2x, x+y=9\).
\begin{enumerate}[label=\textbf{\Alph*)}]
\begin{multicols}{2}
  \item x=3, y=6
  \item x=2, y=4
  \item x=4, y=1
  \item x=1, y=5
  \item x=5, y=0
  \item x=0, y=5
\end{multicols}
\end{enumerate}
\vspace{0.15cm}
\item Resuelve el sistema: \(y=x+1, 2x+y=7\).
\begin{enumerate}[label=\textbf{\Alph*)}]
\begin{multicols}{2}
  \item x=2, y=3
  \item x=2, y=4
  \item x=4, y=1
  \item x=1, y=5
  \item x=5, y=0
  \item x=0, y=5
\end{multicols}
\end{enumerate}
\vspace{0.15cm}
\item Resuelve el sistema: \(y=3x, 2x+y=10\).
\begin{enumerate}[label=\textbf{\Alph*)}]
\begin{multicols}{2}
  \item x=2, y=6
  \item x=2, y=4
  \item x=4, y=1
  \item x=1, y=5
  \item x=5, y=0
  \item x=0, y=5
\end{multicols}
\end{enumerate}
\vspace{0.15cm}
\item Resuelve el sistema: \(y=x-2, x+y=8\).
\begin{enumerate}[label=\textbf{\Alph*)}]
\begin{multicols}{2}
  \item x=5, y=3
  \item x=2, y=4
  \item x=4, y=1
  \item x=1, y=5
  \item x=5, y=0
  \item x=0, y=5
\end{multicols}
\end{enumerate}
\vspace{0.15cm}
\item Resuelve el sistema: \(x=2y, x+y=9\).
\begin{enumerate}[label=\textbf{\Alph*)}]
\begin{multicols}{2}
  \item x=6, y=3
  \item x=2, y=4
  \item x=4, y=1
  \item x=1, y=5
  \item x=5, y=0
  \item x=0, y=5
\end{multicols}
\end{enumerate}
\vspace{0.15cm}
\item Resuelve el sistema: \(y=2x+1, y=x+4\).
\begin{enumerate}[label=\textbf{\Alph*)}]
\begin{multicols}{2}
  \item x=3, y=7
  \item x=2, y=4
  \item x=4, y=1
  \item x=1, y=5
  \item x=5, y=0
  \item x=0, y=5
\end{multicols}
\end{enumerate}
\vspace{0.15cm}
\item Resuelve el sistema: \(y=3x-2, y=x+2\).
\begin{enumerate}[label=\textbf{\Alph*)}]
\begin{multicols}{2}
  \item x=2, y=4
  \item x=2, y=4
  \item x=4, y=1
  \item x=1, y=5
  \item x=5, y=0
  \item x=0, y=5
\end{multicols}
\end{enumerate}
\vspace{0.15cm}
\item Resuelve el sistema: \(y=-x+5, y=2x-1\).
\begin{enumerate}[label=\textbf{\Alph*)}]
\begin{multicols}{2}
  \item x=2, y=3
  \item x=2, y=4
  \item x=4, y=1
  \item x=1, y=5
  \item x=5, y=0
  \item x=0, y=5
\end{multicols}
\end{enumerate}
\vspace{0.15cm}
\item Resuelve el sistema: \(x+y=8, x-y=2\).
\begin{enumerate}[label=\textbf{\Alph*)}]
\begin{multicols}{2}
  \item x=5, y=3
  \item x=2, y=4
  \item x=4, y=1
  \item x=1, y=5
  \item x=5, y=0
  \item x=0, y=5
\end{multicols}
\end{enumerate}
\vspace{0.15cm}
\item Resuelve el sistema: \(2x+y=9, x-y=0\).
\begin{enumerate}[label=\textbf{\Alph*)}]
\begin{multicols}{2}
  \item x=3, y=3
  \item x=2, y=4
  \item x=4, y=1
  \item x=1, y=5
  \item x=5, y=0
  \item x=0, y=5
\end{multicols}
\end{enumerate}
\vspace{0.15cm}
\item Resuelve el sistema: \(3x+y=11, x+y=5\).
\begin{enumerate}[label=\textbf{\Alph*)}]
\begin{multicols}{2}
  \item x=3, y=2
  \item x=2, y=4
  \item x=4, y=1
  \item x=1, y=5
  \item x=5, y=0
  \item x=0, y=5
\end{multicols}
\end{enumerate}
\vspace{0.15cm}
\item Resuelve el sistema: \(x+2y=7, 2x+y=8\).
\begin{enumerate}[label=\textbf{\Alph*)}]
\begin{multicols}{2}
  \item x=3, y=2
  \item x=2, y=4
  \item x=4, y=1
  \item x=1, y=5
  \item x=5, y=0
  \item x=0, y=5
\end{multicols}
\end{enumerate}
\vspace{0.15cm}
\item Resuelve el sistema: \(4x-y=7, 2x+y=5\).
\begin{enumerate}[label=\textbf{\Alph*)}]
\begin{multicols}{2}
  \item x=2, y=1
  \item x=2, y=4
  \item x=4, y=1
  \item x=1, y=5
  \item x=5, y=0
  \item x=0, y=5
\end{multicols}
\end{enumerate}
\vspace{0.15cm}
\item Resuelve el sistema: \(5x+2y=16, 3x-2y=0\).
\begin{enumerate}[label=\textbf{\Alph*)}]
\begin{multicols}{2}
  \item x=2, y=3
  \item x=2, y=4
  \item x=4, y=1
  \item x=1, y=5
  \item x=5, y=0
  \item x=0, y=5
\end{multicols}
\end{enumerate}
\vspace{0.15cm}
\item Resuelve el sistema: \(2x+3y=13, 4x-y=5\).
\begin{enumerate}[label=\textbf{\Alph*)}]
\begin{multicols}{2}
  \item x=2, y=3
  \item x=2, y=4
  \item x=4, y=1
  \item x=1, y=5
  \item x=5, y=0
  \item x=0, y=5
\end{multicols}
\end{enumerate}
\vspace{0.15cm}
\item Resuelve el sistema: \(x+3y=11, 3x+2y=12\).
\begin{enumerate}[label=\textbf{\Alph*)}]
\begin{multicols}{2}
  \item x=2, y=3
  \item x=2, y=4
  \item x=4, y=1
  \item x=1, y=5
  \item x=5, y=0
  \item x=0, y=5
\end{multicols}
\end{enumerate}
\vspace{0.15cm}
\item Resuelve el sistema: \(3x+2y=16, 2x+3y=14\).
\begin{enumerate}[label=\textbf{\Alph*)}]
\begin{multicols}{2}
  \item x=4, y=2
  \item x=2, y=4
  \item x=4, y=1
  \item x=1, y=5
  \item x=5, y=0
  \item x=0, y=5
\end{multicols}
\end{enumerate}
\vspace{0.15cm}
\item Resuelve el sistema: \(5x-3y=1, 2x+y=7\).
\begin{enumerate}[label=\textbf{\Alph*)}]
\begin{multicols}{2}
  \item x=2, y=3
  \item x=2, y=4
  \item x=4, y=1
  \item x=1, y=5
  \item x=5, y=0
  \item x=0, y=5
\end{multicols}
\end{enumerate}
\vspace{0.15cm}
\item Resuelve el sistema: \(4x+5y=22, 2x-y=1\).
\begin{enumerate}[label=\textbf{\Alph*)}]
\begin{multicols}{2}
  \item x=2, y=3
  \item x=2, y=4
  \item x=4, y=1
  \item x=1, y=5
  \item x=5, y=0
  \item x=0, y=5
\end{multicols}
\end{enumerate}
\vspace{0.15cm}
\item Resuelve el sistema: \(x+2y=9, 3x-y=4\).
\begin{enumerate}[label=\textbf{\Alph*)}]
\begin{multicols}{2}
  \item x=3, y=3
  \item x=2, y=4
  \item x=4, y=1
  \item x=1, y=5
  \item x=5, y=0
  \item x=0, y=5
\end{multicols}
\end{enumerate}
\vspace{0.15cm}
\item Resuelve el sistema: \(2x+y=5, x+2y=7\).
\begin{enumerate}[label=\textbf{\Alph*)}]
\begin{multicols}{2}
  \item x=1, y=3
  \item x=2, y=4
  \item x=4, y=1
  \item x=1, y=5
  \item x=5, y=0
  \item x=0, y=5
\end{multicols}
\end{enumerate}
\vspace{0.15cm}
\item Resuelve el sistema: \(3x-y=7, x+2y=7\).
\begin{enumerate}[label=\textbf{\Alph*)}]
\begin{multicols}{2}
  \item x=3, y=2
  \item x=2, y=4
  \item x=4, y=1
  \item x=1, y=5
  \item x=5, y=0
  \item x=0, y=5
\end{multicols}
\end{enumerate}
\vspace{0.15cm}
\item Resuelve el sistema: \(x+y=4, 2x+3y=11\).
\begin{enumerate}[label=\textbf{\Alph*)}]
\begin{multicols}{2}
  \item x=1, y=3
  \item x=2, y=4
  \item x=4, y=1
  \item x=1, y=5
  \item x=5, y=0
  \item x=0, y=5
\end{multicols}
\end{enumerate}
\vspace{0.15cm}
\item Resuelve el sistema: \(3x+2y=12, x+3y=11\).
\begin{enumerate}[label=\textbf{\Alph*)}]
\begin{multicols}{2}
  \item x=2, y=3
  \item x=2, y=4
  \item x=4, y=1
  \item x=1, y=5
  \item x=5, y=0
  \item x=0, y=5
\end{multicols}
\end{enumerate}
\vspace{0.15cm}
\item Resuelve el sistema: \(4x+y=9, 2x+3y=13\).
\begin{enumerate}[label=\textbf{\Alph*)}]
\begin{multicols}{2}
  \item x=1, y=5
  \item x=2, y=4
  \item x=4, y=1
  \item x=1, y=5
  \item x=5, y=0
  \item x=0, y=5
\end{multicols}
\end{enumerate}
\vspace{0.15cm}
\item Resuelve el sistema: \(5x+3y=21, 2x-y=1\).
\begin{enumerate}[label=\textbf{\Alph*)}]
\begin{multicols}{2}
  \item x=2, y=3
  \item x=2, y=4
  \item x=4, y=1
  \item x=1, y=5
  \item x=5, y=0
  \item x=0, y=5
\end{multicols}
\end{enumerate}
\vspace{0.15cm}
\item Resuelve el sistema: \(x+4y=13, 3x+2y=11\).
\begin{enumerate}[label=\textbf{\Alph*)}]
\begin{multicols}{2}
  \item x=1, y=3
  \item x=2, y=4
  \item x=4, y=1
  \item x=1, y=5
  \item x=5, y=0
  \item x=0, y=5
\end{multicols}
\end{enumerate}
\vspace{0.15cm}
\item Resuelve el sistema: \(2x+5y=19, 3x-2y=0\).
\begin{enumerate}[label=\textbf{\Alph*)}]
\begin{multicols}{2}
  \item x=2, y=3
  \item x=2, y=4
  \item x=4, y=1
  \item x=1, y=5
  \item x=5, y=0
  \item x=0, y=5
\end{multicols}
\end{enumerate}
\vspace{0.15cm}
\item Resuelve el sistema: \(3x+4y=17, 2x+3y=12\).
\begin{enumerate}[label=\textbf{\Alph*)}]
\begin{multicols}{2}
  \item x=3, y=2
  \item x=2, y=4
  \item x=4, y=1
  \item x=1, y=5
  \item x=5, y=0
  \item x=0, y=5
\end{multicols}
\end{enumerate}
\vspace{0.15cm}
\item Resuelve el sistema: \(4x+3y=19, 2x+5y=17\).
\begin{enumerate}[label=\textbf{\Alph*)}]
\begin{multicols}{2}
  \item x=4, y=1
  \item x=2, y=4
  \item x=4, y=1
  \item x=1, y=5
  \item x=5, y=0
  \item x=0, y=5
\end{multicols}
\end{enumerate}
\vspace{0.15cm}
\end{enumerate}\n
\subsubsection*{SECCIÓN VI: ECUACIONES CUADRÁTICAS (251–300)}

\begin{enumerate}[label=\textbf{\arabic*.}]
\setcounter{enumi}{250}
\item Resuelve \(x^2-5x+6=0\).
\begin{enumerate}[label=\textbf{\Alph*)}]
\begin{multicols}{2}
  \item x=2, x=3
  \item x=-2, x=-3
  \item x=1, x=6
  \item x=0
  \item x=2
  \item Sin solución
\end{multicols}
\end{enumerate}
\vspace{0.15cm}
\item Resuelve \(x^2-7x+12=0\).
\begin{enumerate}[label=\textbf{\Alph*)}]
\begin{multicols}{2}
  \item x=3, x=4
  \item x=-2, x=-3
  \item x=1, x=6
  \item x=0
  \item x=2
  \item Sin solución
\end{multicols}
\end{enumerate}
\vspace{0.15cm}
\item Resuelve \(x^2-9x+20=0\).
\begin{enumerate}[label=\textbf{\Alph*)}]
\begin{multicols}{2}
  \item x=4, x=5
  \item x=-2, x=-3
  \item x=1, x=6
  \item x=0
  \item x=2
  \item Sin solución
\end{multicols}
\end{enumerate}
\vspace{0.15cm}
\item Resuelve \(x^2-6x+8=0\).
\begin{enumerate}[label=\textbf{\Alph*)}]
\begin{multicols}{2}
  \item x=2, x=4
  \item x=-2, x=-3
  \item x=1, x=6
  \item x=0
  \item x=2
  \item Sin solución
\end{multicols}
\end{enumerate}
\vspace{0.15cm}
\item Resuelve \(x^2-8x+15=0\).
\begin{enumerate}[label=\textbf{\Alph*)}]
\begin{multicols}{2}
  \item x=3, x=5
  \item x=-2, x=-3
  \item x=1, x=6
  \item x=0
  \item x=2
  \item Sin solución
\end{multicols}
\end{enumerate}
\vspace{0.15cm}
\item Resuelve \(x^2+5x+6=0\).
\begin{enumerate}[label=\textbf{\Alph*)}]
\begin{multicols}{2}
  \item x=-2, x=-3
  \item x=-2, x=-3
  \item x=1, x=6
  \item x=0
  \item x=2
  \item Sin solución
\end{multicols}
\end{enumerate}
\vspace{0.15cm}
\item Resuelve \(x^2+7x+12=0\).
\begin{enumerate}[label=\textbf{\Alph*)}]
\begin{multicols}{2}
  \item x=-3, x=-4
  \item x=-2, x=-3
  \item x=1, x=6
  \item x=0
  \item x=2
  \item Sin solución
\end{multicols}
\end{enumerate}
\vspace{0.15cm}
\item Resuelve \(x^2+9x+20=0\).
\begin{enumerate}[label=\textbf{\Alph*)}]
\begin{multicols}{2}
  \item x=-4, x=-5
  \item x=-2, x=-3
  \item x=1, x=6
  \item x=0
  \item x=2
  \item Sin solución
\end{multicols}
\end{enumerate}
\vspace{0.15cm}
\item Resuelve \(x^2+6x+8=0\).
\begin{enumerate}[label=\textbf{\Alph*)}]
\begin{multicols}{2}
  \item x=-2, x=-4
  \item x=-2, x=-3
  \item x=1, x=6
  \item x=0
  \item x=2
  \item Sin solución
\end{multicols}
\end{enumerate}
\vspace{0.15cm}
\item Resuelve \(x^2+8x+15=0\).
\begin{enumerate}[label=\textbf{\Alph*)}]
\begin{multicols}{2}
  \item x=-3, x=-5
  \item x=-2, x=-3
  \item x=1, x=6
  \item x=0
  \item x=2
  \item Sin solución
\end{multicols}
\end{enumerate}
\vspace{0.15cm}
\item Resuelve \(x^2-4=0\).
\begin{enumerate}[label=\textbf{\Alph*)}]
\begin{multicols}{2}
  \item x=2, x=-2
  \item x=-2, x=-3
  \item x=1, x=6
  \item x=0
  \item x=2
  \item Sin solución
\end{multicols}
\end{enumerate}
\vspace{0.15cm}
\item Resuelve \(x^2-9=0\).
\begin{enumerate}[label=\textbf{\Alph*)}]
\begin{multicols}{2}
  \item x=3, x=-3
  \item x=-2, x=-3
  \item x=1, x=6
  \item x=0
  \item x=2
  \item Sin solución
\end{multicols}
\end{enumerate}
\vspace{0.15cm}
\item Resuelve \(x^2-16=0\).
\begin{enumerate}[label=\textbf{\Alph*)}]
\begin{multicols}{2}
  \item x=4, x=-4
  \item x=-2, x=-3
  \item x=1, x=6
  \item x=0
  \item x=2
  \item Sin solución
\end{multicols}
\end{enumerate}
\vspace{0.15cm}
\item Resuelve \(x^2-25=0\).
\begin{enumerate}[label=\textbf{\Alph*)}]
\begin{multicols}{2}
  \item x=5, x=-5
  \item x=-2, x=-3
  \item x=1, x=6
  \item x=0
  \item x=2
  \item Sin solución
\end{multicols}
\end{enumerate}
\vspace{0.15cm}
\item Resuelve \(x^2-1=0\).
\begin{enumerate}[label=\textbf{\Alph*)}]
\begin{multicols}{2}
  \item x=1, x=-1
  \item x=-2, x=-3
  \item x=1, x=6
  \item x=0
  \item x=2
  \item Sin solución
\end{multicols}
\end{enumerate}
\vspace{0.15cm}
\item Resuelve \(x^2=4\).
\begin{enumerate}[label=\textbf{\Alph*)}]
\begin{multicols}{2}
  \item x=2, x=-2
  \item x=-2, x=-3
  \item x=1, x=6
  \item x=0
  \item x=2
  \item Sin solución
\end{multicols}
\end{enumerate}
\vspace{0.15cm}
\item Resuelve \(x^2=9\).
\begin{enumerate}[label=\textbf{\Alph*)}]
\begin{multicols}{2}
  \item x=3, x=-3
  \item x=-2, x=-3
  \item x=1, x=6
  \item x=0
  \item x=2
  \item Sin solución
\end{multicols}
\end{enumerate}
\vspace{0.15cm}
\item Resuelve \(x^2=16\).
\begin{enumerate}[label=\textbf{\Alph*)}]
\begin{multicols}{2}
  \item x=4, x=-4
  \item x=-2, x=-3
  \item x=1, x=6
  \item x=0
  \item x=2
  \item Sin solución
\end{multicols}
\end{enumerate}
\vspace{0.15cm}
\item Resuelve \(x^2=25\).
\begin{enumerate}[label=\textbf{\Alph*)}]
\begin{multicols}{2}
  \item x=5, x=-5
  \item x=-2, x=-3
  \item x=1, x=6
  \item x=0
  \item x=2
  \item Sin solución
\end{multicols}
\end{enumerate}
\vspace{0.15cm}
\item Resuelve \(x^2=1\).
\begin{enumerate}[label=\textbf{\Alph*)}]
\begin{multicols}{2}
  \item x=1, x=-1
  \item x=-2, x=-3
  \item x=1, x=6
  \item x=0
  \item x=2
  \item Sin solución
\end{multicols}
\end{enumerate}
\vspace{0.15cm}
\item Resuelve \(2x^2-8=0\).
\begin{enumerate}[label=\textbf{\Alph*)}]
\begin{multicols}{2}
  \item x=2, x=-2
  \item x=-2, x=-3
  \item x=1, x=6
  \item x=0
  \item x=2
  \item Sin solución
\end{multicols}
\end{enumerate}
\vspace{0.15cm}
\item Resuelve \(3x^2-12=0\).
\begin{enumerate}[label=\textbf{\Alph*)}]
\begin{multicols}{2}
  \item x=2, x=-2
  \item x=-2, x=-3
  \item x=1, x=6
  \item x=0
  \item x=2
  \item Sin solución
\end{multicols}
\end{enumerate}
\vspace{0.15cm}
\item Resuelve \(4x^2-16=0\).
\begin{enumerate}[label=\textbf{\Alph*)}]
\begin{multicols}{2}
  \item x=2, x=-2
  \item x=-2, x=-3
  \item x=1, x=6
  \item x=0
  \item x=2
  \item Sin solución
\end{multicols}
\end{enumerate}
\vspace{0.15cm}
\item Resuelve \(5x^2-20=0\).
\begin{enumerate}[label=\textbf{\Alph*)}]
\begin{multicols}{2}
  \item x=2, x=-2
  \item x=-2, x=-3
  \item x=1, x=6
  \item x=0
  \item x=2
  \item Sin solución
\end{multicols}
\end{enumerate}
\vspace{0.15cm}
\item Resuelve \(x^2+2x=0\).
\begin{enumerate}[label=\textbf{\Alph*)}]
\begin{multicols}{2}
  \item x=0, x=-2
  \item x=-2, x=-3
  \item x=1, x=6
  \item x=0
  \item x=2
  \item Sin solución
\end{multicols}
\end{enumerate}
\vspace{0.15cm}
\item Resuelve \(x^2-3x=0\).
\begin{enumerate}[label=\textbf{\Alph*)}]
\begin{multicols}{2}
  \item x=0, x=3
  \item x=-2, x=-3
  \item x=1, x=6
  \item x=0
  \item x=2
  \item Sin solución
\end{multicols}
\end{enumerate}
\vspace{0.15cm}
\item Resuelve \(x^2+4x=0\).
\begin{enumerate}[label=\textbf{\Alph*)}]
\begin{multicols}{2}
  \item x=0, x=-4
  \item x=-2, x=-3
  \item x=1, x=6
  \item x=0
  \item x=2
  \item Sin solución
\end{multicols}
\end{enumerate}
\vspace{0.15cm}
\item Resuelve \(x^2-5x=0\).
\begin{enumerate}[label=\textbf{\Alph*)}]
\begin{multicols}{2}
  \item x=0, x=5
  \item x=-2, x=-3
  \item x=1, x=6
  \item x=0
  \item x=2
  \item Sin solución
\end{multicols}
\end{enumerate}
\vspace{0.15cm}
\item Resuelve \(2x^2+4x=0\).
\begin{enumerate}[label=\textbf{\Alph*)}]
\begin{multicols}{2}
  \item x=0, x=-2
  \item x=-2, x=-3
  \item x=1, x=6
  \item x=0
  \item x=2
  \item Sin solución
\end{multicols}
\end{enumerate}
\vspace{0.15cm}
\item Resuelve \(3x^2-9x=0\).
\begin{enumerate}[label=\textbf{\Alph*)}]
\begin{multicols}{2}
  \item x=0, x=3
  \item x=-2, x=-3
  \item x=1, x=6
  \item x=0
  \item x=2
  \item Sin solución
\end{multicols}
\end{enumerate}
\vspace{0.15cm}
\item Resuelve \(x^2-4x+3=0\).
\begin{enumerate}[label=\textbf{\Alph*)}]
\begin{multicols}{2}
  \item x=1, x=3
  \item x=-2, x=-3
  \item x=1, x=6
  \item x=0
  \item x=2
  \item Sin solución
\end{multicols}
\end{enumerate}
\vspace{0.15cm}
\item Resuelve \(x^2-6x+5=0\).
\begin{enumerate}[label=\textbf{\Alph*)}]
\begin{multicols}{2}
  \item x=1, x=5
  \item x=-2, x=-3
  \item x=1, x=6
  \item x=0
  \item x=2
  \item Sin solución
\end{multicols}
\end{enumerate}
\vspace{0.15cm}
\item Resuelve \(x^2-2x-3=0\).
\begin{enumerate}[label=\textbf{\Alph*)}]
\begin{multicols}{2}
  \item x=3, x=-1
  \item x=-2, x=-3
  \item x=1, x=6
  \item x=0
  \item x=2
  \item Sin solución
\end{multicols}
\end{enumerate}
\vspace{0.15cm}
\item Resuelve \(x^2+2x-3=0\).
\begin{enumerate}[label=\textbf{\Alph*)}]
\begin{multicols}{2}
  \item x=1, x=-3
  \item x=-2, x=-3
  \item x=1, x=6
  \item x=0
  \item x=2
  \item Sin solución
\end{multicols}
\end{enumerate}
\vspace{0.15cm}
\item Resuelve \(x^2-5x+4=0\).
\begin{enumerate}[label=\textbf{\Alph*)}]
\begin{multicols}{2}
  \item x=1, x=4
  \item x=-2, x=-3
  \item x=1, x=6
  \item x=0
  \item x=2
  \item Sin solución
\end{multicols}
\end{enumerate}
\vspace{0.15cm}
\item Resuelve \(x^2+3x-10=0\).
\begin{enumerate}[label=\textbf{\Alph*)}]
\begin{multicols}{2}
  \item x=2, x=-5
  \item x=-2, x=-3
  \item x=1, x=6
  \item x=0
  \item x=2
  \item Sin solución
\end{multicols}
\end{enumerate}
\vspace{0.15cm}
\item Resuelve \(x^2-3x-10=0\).
\begin{enumerate}[label=\textbf{\Alph*)}]
\begin{multicols}{2}
  \item x=5, x=-2
  \item x=-2, x=-3
  \item x=1, x=6
  \item x=0
  \item x=2
  \item Sin solución
\end{multicols}
\end{enumerate}
\vspace{0.15cm}
\item Resuelve \(x^2+5x+4=0\).
\begin{enumerate}[label=\textbf{\Alph*)}]
\begin{multicols}{2}
  \item x=-1, x=-4
  \item x=-2, x=-3
  \item x=1, x=6
  \item x=0
  \item x=2
  \item Sin solución
\end{multicols}
\end{enumerate}
\vspace{0.15cm}
\item Resuelve \(x^2-7x+10=0\).
\begin{enumerate}[label=\textbf{\Alph*)}]
\begin{multicols}{2}
  \item x=2, x=5
  \item x=-2, x=-3
  \item x=1, x=6
  \item x=0
  \item x=2
  \item Sin solución
\end{multicols}
\end{enumerate}
\vspace{0.15cm}
\item Resuelve \(x^2+6x+5=0\).
\begin{enumerate}[label=\textbf{\Alph*)}]
\begin{multicols}{2}
  \item x=-1, x=-5
  \item x=-2, x=-3
  \item x=1, x=6
  \item x=0
  \item x=2
  \item Sin solución
\end{multicols}
\end{enumerate}
\vspace{0.15cm}
\item Resuelve \(2x^2-7x+3=0\).
\begin{enumerate}[label=\textbf{\Alph*)}]
\begin{multicols}{2}
  \item x=3, x=\frac{1}{2}
  \item x=-2, x=-3
  \item x=1, x=6
  \item x=0
  \item x=2
  \item Sin solución
\end{multicols}
\end{enumerate}
\vspace{0.15cm}
\item Resuelve \(3x^2-10x+8=0\).
\begin{enumerate}[label=\textbf{\Alph*)}]
\begin{multicols}{2}
  \item x=2, x=\frac{4}{3}
  \item x=-2, x=-3
  \item x=1, x=6
  \item x=0
  \item x=2
  \item Sin solución
\end{multicols}
\end{enumerate}
\vspace{0.15cm}
\item Resuelve \(2x^2+5x-3=0\).
\begin{enumerate}[label=\textbf{\Alph*)}]
\begin{multicols}{2}
  \item x=\frac{1}{2}, x=-3
  \item x=-2, x=-3
  \item x=1, x=6
  \item x=0
  \item x=2
  \item Sin solución
\end{multicols}
\end{enumerate}
\vspace{0.15cm}
\item Resuelve \(6x^2+7x+2=0\).
\begin{enumerate}[label=\textbf{\Alph*)}]
\begin{multicols}{2}
  \item x=-\frac{1}{2}, x=-\frac{2}{3}
  \item x=-2, x=-3
  \item x=1, x=6
  \item x=0
  \item x=2
  \item Sin solución
\end{multicols}
\end{enumerate}
\vspace{0.15cm}
\item Resuelve \(4x^2-4x+1=0\).
\begin{enumerate}[label=\textbf{\Alph*)}]
\begin{multicols}{2}
  \item x=\frac{1}{2} (doble)
  \item x=-2, x=-3
  \item x=1, x=6
  \item x=0
  \item x=2
  \item Sin solución
\end{multicols}
\end{enumerate}
\vspace{0.15cm}
\item Resuelve \(9x^2+6x+1=0\).
\begin{enumerate}[label=\textbf{\Alph*)}]
\begin{multicols}{2}
  \item x=-\frac{1}{3} (doble)
  \item x=-2, x=-3
  \item x=1, x=6
  \item x=0
  \item x=2
  \item Sin solución
\end{multicols}
\end{enumerate}
\vspace{0.15cm}
\item Resuelve \(x^2-6x+9=0\).
\begin{enumerate}[label=\textbf{\Alph*)}]
\begin{multicols}{2}
  \item x=3 (doble)
  \item x=-2, x=-3
  \item x=1, x=6
  \item x=0
  \item x=2
  \item Sin solución
\end{multicols}
\end{enumerate}
\vspace{0.15cm}
\item Resuelve \(x^2+8x+16=0\).
\begin{enumerate}[label=\textbf{\Alph*)}]
\begin{multicols}{2}
  \item x=-4 (doble)
  \item x=-2, x=-3
  \item x=1, x=6
  \item x=0
  \item x=2
  \item Sin solución
\end{multicols}
\end{enumerate}
\vspace{0.15cm}
\item ¿Cuál es el discriminante de x^2-4x+3=0?
\begin{enumerate}[label=\textbf{\Alph*)}]
\begin{multicols}{2}
  \item 4
  \item 1
  \item 2
  \item 3
  \item x=2
  \item Sin solución
\end{multicols}
\end{enumerate}
\vspace{0.15cm}
\item ¿Cuántas soluciones reales tiene una ecuación con discriminante negativo?
\begin{enumerate}[label=\textbf{\Alph*)}]
\begin{multicols}{2}
  \item 0
  \item 1
  \item 2
  \item 3
  \item x=2
  \item Sin solución
\end{multicols}
\end{enumerate}
\vspace{0.15cm}
\end{enumerate}\n
\vspace{0.4cm}
\begin{tcolorbox}[colback=gray!5!white,colframe=gray!60!black,title=\textbf{SOLUCIONARIO -- 1.5 Álgebra}]
\begin{multicols}{4}
\begin{enumerate}[label=\textbf{\arabic*.}]
  \item \textbf{B}
  \item \textbf{A}
  \item \textbf{A}
  \item \textbf{A}
  \item \textbf{A}
  \item \textbf{B}
  \item \textbf{A}
  \item \textbf{A}
  \item \textbf{A}
  \item \textbf{A}
  \item \textbf{B}
  \item \textbf{B}
  \item \textbf{C}
  \item \textbf{B}
  \item \textbf{C}
  \item \textbf{C}
  \item \textbf{A}
  \item \textbf{B}
  \item \textbf{A}
  \item \textbf{A}
  \item \textbf{A}
  \item \textbf{A}
  \item \textbf{A}
  \item \textbf{A}
  \item \textbf{A}
  \item \textbf{A}
  \item \textbf{A}
  \item \textbf{A}
  \item \textbf{A}
  \item \textbf{A}
  \item \textbf{A}
  \item \textbf{B}
  \item \textbf{B}
  \item \textbf{B}
  \item \textbf{A}
  \item \textbf{A}
  \item \textbf{A}
  \item \textbf{A}
  \item \textbf{A}
  \item \textbf{A}
  \item \textbf{A}
  \item \textbf{A}
  \item \textbf{A}
  \item \textbf{A}
  \item \textbf{A}
  \item \textbf{A}
  \item \textbf{A}
  \item \textbf{A}
  \item \textbf{A}
  \item \textbf{A}
  \item \textbf{A}
  \item \textbf{A}
  \item \textbf{A}
  \item \textbf{A}
  \item \textbf{A}
  \item \textbf{A}
  \item \textbf{A}
  \item \textbf{A}
  \item \textbf{A}
  \item \textbf{A}
  \item \textbf{A}
  \item \textbf{A}
  \item \textbf{A}
  \item \textbf{A}
  \item \textbf{A}
  \item \textbf{A}
  \item \textbf{A}
  \item \textbf{A}
  \item \textbf{A}
  \item \textbf{A}
  \item \textbf{A}
  \item \textbf{A}
  \item \textbf{A}
  \item \textbf{A}
  \item \textbf{A}
  \item \textbf{A}
  \item \textbf{A}
  \item \textbf{A}
  \item \textbf{A}
  \item \textbf{A}
  \item \textbf{A}
  \item \textbf{A}
  \item \textbf{A}
  \item \textbf{A}
  \item \textbf{A}
  \item \textbf{A}
  \item \textbf{A}
  \item \textbf{A}
  \item \textbf{A}
  \item \textbf{A}
  \item \textbf{A}
  \item \textbf{A}
  \item \textbf{A}
  \item \textbf{A}
  \item \textbf{A}
  \item \textbf{A}
  \item \textbf{A}
  \item \textbf{A}
  \item \textbf{A}
  \item \textbf{A}
  \item \textbf{A}
  \item \textbf{A}
  \item \textbf{A}
  \item \textbf{A}
  \item \textbf{A}
  \item \textbf{A}
  \item \textbf{A}
  \item \textbf{A}
  \item \textbf{A}
  \item \textbf{A}
  \item \textbf{A}
  \item \textbf{A}
  \item \textbf{A}
  \item \textbf{A}
  \item \textbf{A}
  \item \textbf{A}
  \item \textbf{A}
  \item \textbf{A}
  \item \textbf{A}
  \item \textbf{A}
  \item \textbf{A}
  \item \textbf{A}
  \item \textbf{A}
  \item \textbf{A}
  \item \textbf{A}
  \item \textbf{A}
  \item \textbf{A}
  \item \textbf{A}
  \item \textbf{A}
  \item \textbf{A}
  \item \textbf{A}
  \item \textbf{A}
  \item \textbf{A}
  \item \textbf{A}
  \item \textbf{A}
  \item \textbf{A}
  \item \textbf{A}
  \item \textbf{A}
  \item \textbf{A}
  \item \textbf{A}
  \item \textbf{A}
  \item \textbf{A}
  \item \textbf{A}
  \item \textbf{A}
  \item \textbf{A}
  \item \textbf{A}
  \item \textbf{A}
  \item \textbf{A}
  \item \textbf{A}
  \item \textbf{A}
  \item \textbf{A}
  \item \textbf{A}
  \item \textbf{A}
  \item \textbf{A}
  \item \textbf{A}
  \item \textbf{A}
  \item \textbf{A}
  \item \textbf{A}
  \item \textbf{A}
  \item \textbf{A}
  \item \textbf{A}
  \item \textbf{A}
  \item \textbf{A}
  \item \textbf{A}
  \item \textbf{A}
  \item \textbf{A}
  \item \textbf{A}
  \item \textbf{A}
  \item \textbf{A}
  \item \textbf{A}
  \item \textbf{A}
  \item \textbf{A}
  \item \textbf{A}
  \item \textbf{A}
  \item \textbf{A}
  \item \textbf{A}
  \item \textbf{A}
  \item \textbf{A}
  \item \textbf{A}
  \item \textbf{A}
  \item \textbf{A}
  \item \textbf{A}
  \item \textbf{A}
  \item \textbf{A}
  \item \textbf{A}
  \item \textbf{A}
  \item \textbf{A}
  \item \textbf{A}
  \item \textbf{A}
  \item \textbf{A}
  \item \textbf{A}
  \item \textbf{A}
  \item \textbf{A}
  \item \textbf{A}
  \item \textbf{A}
  \item \textbf{A}
  \item \textbf{A}
  \item \textbf{A}
  \item \textbf{A}
  \item \textbf{A}
  \item \textbf{A}
  \item \textbf{A}
  \item \textbf{A}
  \item \textbf{A}
  \item \textbf{A}
  \item \textbf{A}
  \item \textbf{A}
  \item \textbf{A}
  \item \textbf{A}
  \item \textbf{A}
  \item \textbf{A}
  \item \textbf{A}
  \item \textbf{A}
  \item \textbf{A}
  \item \textbf{A}
  \item \textbf{A}
  \item \textbf{A}
  \item \textbf{A}
  \item \textbf{A}
  \item \textbf{A}
  \item \textbf{A}
  \item \textbf{A}
  \item \textbf{A}
  \item \textbf{A}
  \item \textbf{A}
  \item \textbf{A}
  \item \textbf{A}
  \item \textbf{A}
  \item \textbf{A}
  \item \textbf{A}
  \item \textbf{A}
  \item \textbf{A}
  \item \textbf{A}
  \item \textbf{A}
  \item \textbf{A}
  \item \textbf{A}
  \item \textbf{A}
  \item \textbf{A}
  \item \textbf{A}
  \item \textbf{A}
  \item \textbf{A}
  \item \textbf{A}
  \item \textbf{A}
  \item \textbf{A}
  \item \textbf{A}
  \item \textbf{A}
  \item \textbf{A}
  \item \textbf{A}
  \item \textbf{A}
  \item \textbf{A}
  \item \textbf{A}
  \item \textbf{A}
  \item \textbf{A}
  \item \textbf{A}
  \item \textbf{A}
  \item \textbf{A}
  \item \textbf{A}
  \item \textbf{A}
  \item \textbf{A}
  \item \textbf{A}
  \item \textbf{A}
  \item \textbf{A}
  \item \textbf{A}
  \item \textbf{A}
  \item \textbf{A}
  \item \textbf{A}
  \item \textbf{A}
  \item \textbf{A}
  \item \textbf{A}
  \item \textbf{A}
  \item \textbf{A}
  \item \textbf{A}
  \item \textbf{A}
  \item \textbf{A}
  \item \textbf{A}
  \item \textbf{A}
  \item \textbf{A}
  \item \textbf{A}
  \item \textbf{A}
  \item \textbf{A}
  \item \textbf{A}
  \item \textbf{A}
  \item \textbf{A}
  \item \textbf{A}
  \item \textbf{A}
  \item \textbf{A}
  \item \textbf{A}
  \item \textbf{A}
  \item \textbf{A}
  \item \textbf{A}
  \item \textbf{A}
  \item \textbf{A}
  \item \textbf{A}
  \item \textbf{A}
  \item \textbf{A}
  \item \textbf{A}
  \item \textbf{A}
  \item \textbf{A}
  \item \textbf{A}
  \item \textbf{A}
\end{enumerate}
\end{multicols}
\end{tcolorbox}
\newpage\n
\subsection{1.6 Geometría y Trigonometría}

\subsubsection*{SECCIÓN I: CONCEPTOS BÁSICOS DE GEOMETRÍA (1–40)}

\begin{enumerate}[label=\textbf{\arabic*.}]
\setcounter{enumi}{0}
\item ¿Cuántos grados mide un ángulo recto?
\begin{enumerate}[label=\textbf{\Alph*)}]
\begin{multicols}{2}
  \item 45^\circ
  \item 90^\circ
  \item 180^\circ
  \item 360^\circ
  \item 60^\circ
  \item 270^\circ
\end{multicols}
\end{enumerate}
\vspace{0.15cm}
\item ¿Cuántos grados mide un ángulo llano?
\begin{enumerate}[label=\textbf{\Alph*)}]
\begin{multicols}{2}
  \item 45^\circ
  \item 90^\circ
  \item 180^\circ
  \item 360^\circ
  \item 60^\circ
  \item 270^\circ
\end{multicols}
\end{enumerate}
\vspace{0.15cm}
\item ¿Cuántos grados mide un ángulo completo?
\begin{enumerate}[label=\textbf{\Alph*)}]
\begin{multicols}{2}
  \item 45^\circ
  \item 90^\circ
  \item 180^\circ
  \item 360^\circ
  \item 60^\circ
  \item 270^\circ
\end{multicols}
\end{enumerate}
\vspace{0.15cm}
\item Un ángulo agudo mide:
\begin{enumerate}[label=\textbf{\Alph*)}]
\begin{multicols}{2}
  \item Menos de 90^\circ
  \item 90^\circ
  \item 180^\circ
  \item 360^\circ
  \item 60^\circ
  \item 270^\circ
\end{multicols}
\end{enumerate}
\vspace{0.15cm}
\item Un ángulo obtuso mide:
\begin{enumerate}[label=\textbf{\Alph*)}]
\begin{multicols}{2}
  \item 45^\circ
  \item 90^\circ
  \item Más de 90^\circ y menos de 180^\circ
  \item 360^\circ
  \item 60^\circ
  \item 270^\circ
\end{multicols}
\end{enumerate}
\vspace{0.15cm}
\item Dos ángulos son complementarios si suman:
\begin{enumerate}[label=\textbf{\Alph*)}]
\begin{multicols}{2}
  \item 3
  \item 90^\circ
  \item 5
  \item 6
  \item 7
  \item 8
\end{multicols}
\end{enumerate}
\vspace{0.15cm}
\item Dos ángulos son suplementarios si suman:
\begin{enumerate}[label=\textbf{\Alph*)}]
\begin{multicols}{2}
  \item 3
  \item 4
  \item 180^\circ
  \item 6
  \item 7
  \item 8
\end{multicols}
\end{enumerate}
\vspace{0.15cm}
\item El complemento de \(30^\circ\) es:
\begin{enumerate}[label=\textbf{\Alph*)}]
\begin{multicols}{2}
  \item 3
  \item 60^\circ
  \item 5
  \item 6
  \item 7
  \item 8
\end{multicols}
\end{enumerate}
\vspace{0.15cm}
\item El suplemento de \(120^\circ\) es:
\begin{enumerate}[label=\textbf{\Alph*)}]
\begin{multicols}{2}
  \item 3
  \item 60^\circ
  \item 5
  \item 6
  \item 7
  \item 8
\end{multicols}
\end{enumerate}
\vspace{0.15cm}
\item El complemento de \(45^\circ\) es:
\begin{enumerate}[label=\textbf{\Alph*)}]
\begin{multicols}{2}
  \item 3
  \item 45^\circ
  \item 5
  \item 6
  \item 7
  \item 8
\end{multicols}
\end{enumerate}
\vspace{0.15cm}
\item El suplemento de \(60^\circ\) es:
\begin{enumerate}[label=\textbf{\Alph*)}]
\begin{multicols}{2}
  \item 3
  \item 4
  \item 120^\circ
  \item 6
  \item 7
  \item 8
\end{multicols}
\end{enumerate}
\vspace{0.15cm}
\item ¿Cómo se llama un triángulo con tres lados iguales?
\begin{enumerate}[label=\textbf{\Alph*)}]
\begin{multicols}{2}
  \item Escaleno
  \item Isósceles
  \item Equilátero
  \item 6
  \item 7
  \item 8
\end{multicols}
\end{enumerate}
\vspace{0.15cm}
\item ¿Cómo se llama un triángulo con dos lados iguales?
\begin{enumerate}[label=\textbf{\Alph*)}]
\begin{multicols}{2}
  \item Escaleno
  \item Isósceles
  \item Equilátero
  \item 6
  \item 7
  \item 8
\end{multicols}
\end{enumerate}
\vspace{0.15cm}
\item ¿Cómo se llama un triángulo con tres lados diferentes?
\begin{enumerate}[label=\textbf{\Alph*)}]
\begin{multicols}{2}
  \item Escaleno
  \item Isósceles
  \item Equilátero
  \item 6
  \item 7
  \item 8
\end{multicols}
\end{enumerate}
\vspace{0.15cm}
\item En un triángulo equilátero, cada ángulo mide:
\begin{enumerate}[label=\textbf{\Alph*)}]
\begin{multicols}{2}
  \item 45^\circ
  \item 90^\circ
  \item 60^\circ
  \item 360^\circ
  \item 60^\circ
  \item 270^\circ
\end{multicols}
\end{enumerate}
\vspace{0.15cm}
\item La suma de los ángulos internos de un triángulo es:
\begin{enumerate}[label=\textbf{\Alph*)}]
\begin{multicols}{2}
  \item Escaleno
  \item 180^\circ
  \item Equilátero
  \item 6
  \item 7
  \item 8
\end{multicols}
\end{enumerate}
\vspace{0.15cm}
\item La suma de los ángulos internos de un cuadrilátero es:
\begin{enumerate}[label=\textbf{\Alph*)}]
\begin{multicols}{2}
  \item 3
  \item 4
  \item 5
  \item 360^\circ
  \item 7
  \item 8
\end{multicols}
\end{enumerate}
\vspace{0.15cm}
\item Un triángulo rectángulo tiene:
\begin{enumerate}[label=\textbf{\Alph*)}]
\begin{multicols}{2}
  \item Un ángulo recto
  \item Isósceles
  \item Equilátero
  \item 6
  \item 7
  \item 8
\end{multicols}
\end{enumerate}
\vspace{0.15cm}
\item En un triángulo rectángulo, el lado opuesto al ángulo recto se llama:
\begin{enumerate}[label=\textbf{\Alph*)}]
\begin{multicols}{2}
  \item Escaleno
  \item Hipotenusa
  \item Equilátero
  \item 6
  \item 7
  \item 8
\end{multicols}
\end{enumerate}
\vspace{0.15cm}
\item En un triángulo rectángulo, los lados que forman el ángulo recto se llaman:
\begin{enumerate}[label=\textbf{\Alph*)}]
\begin{multicols}{2}
  \item Escaleno
  \item Catetos
  \item Equilátero
  \item 6
  \item 7
  \item 8
\end{multicols}
\end{enumerate}
\vspace{0.15cm}
\item Si un triángulo tiene ángulos de \(30^\circ\) y \(60^\circ\), el tercer ángulo mide:
\begin{enumerate}[label=\textbf{\Alph*)}]
\begin{multicols}{2}
  \item 45^\circ
  \item 90^\circ
  \item 90^\circ
  \item 360^\circ
  \item 60^\circ
  \item 270^\circ
\end{multicols}
\end{enumerate}
\vspace{0.15cm}
\item Si un triángulo tiene ángulos de \(45^\circ\) y \(45^\circ\), el tercer ángulo mide:
\begin{enumerate}[label=\textbf{\Alph*)}]
\begin{multicols}{2}
  \item 45^\circ
  \item 90^\circ
  \item 90^\circ
  \item 360^\circ
  \item 60^\circ
  \item 270^\circ
\end{multicols}
\end{enumerate}
\vspace{0.15cm}
\item Si un triángulo tiene ángulos de \(50^\circ\) y \(70^\circ\), el tercer ángulo mide:
\begin{enumerate}[label=\textbf{\Alph*)}]
\begin{multicols}{2}
  \item 45^\circ
  \item 60^\circ
  \item 180^\circ
  \item 360^\circ
  \item 60^\circ
  \item 270^\circ
\end{multicols}
\end{enumerate}
\vspace{0.15cm}
\item Si un triángulo tiene ángulos de \(100^\circ\) y \(40^\circ\), el tercer ángulo mide:
\begin{enumerate}[label=\textbf{\Alph*)}]
\begin{multicols}{2}
  \item 45^\circ
  \item 40^\circ
  \item 180^\circ
  \item 360^\circ
  \item 60^\circ
  \item 270^\circ
\end{multicols}
\end{enumerate}
\vspace{0.15cm}
\item ¿Cuántos lados tiene un pentágono?
\begin{enumerate}[label=\textbf{\Alph*)}]
\begin{multicols}{2}
  \item 3
  \item 4
  \item 5
  \item 6
  \item 7
  \item 8
\end{multicols}
\end{enumerate}
\vspace{0.15cm}
\item ¿Cuántos lados tiene un hexágono?
\begin{enumerate}[label=\textbf{\Alph*)}]
\begin{multicols}{2}
  \item 3
  \item 4
  \item 5
  \item 6
  \item 7
  \item 8
\end{multicols}
\end{enumerate}
\vspace{0.15cm}
\item ¿Cuántos lados tiene un octágono?
\begin{enumerate}[label=\textbf{\Alph*)}]
\begin{multicols}{2}
  \item 3
  \item 4
  \item 5
  \item 8
  \item 7
  \item 8
\end{multicols}
\end{enumerate}
\vspace{0.15cm}
\item Un ángulo de \(270^\circ\) se clasifica como:
\begin{enumerate}[label=\textbf{\Alph*)}]
\begin{multicols}{2}
  \item 3
  \item 4
  \item 5
  \item 6
  \item Cóncavo
  \item 8
\end{multicols}
\end{enumerate}
\vspace{0.15cm}
\item Dos rectas que nunca se intersectan se llaman:
\begin{enumerate}[label=\textbf{\Alph*)}]
\begin{multicols}{2}
  \item 3
  \item Paralelas
  \item 5
  \item 6
  \item 7
  \item 8
\end{multicols}
\end{enumerate}
\vspace{0.15cm}
\item Dos rectas que se intersectan formando \(90^\circ\) se llaman:
\begin{enumerate}[label=\textbf{\Alph*)}]
\begin{multicols}{2}
  \item 3
  \item Perpendiculares
  \item 5
  \item 6
  \item 7
  \item 8
\end{multicols}
\end{enumerate}
\vspace{0.15cm}
\item ¿Cuánto mide cada ángulo de un cuadrado?
\begin{enumerate}[label=\textbf{\Alph*)}]
\begin{multicols}{2}
  \item 45^\circ
  \item 90^\circ
  \item 90^\circ
  \item 360^\circ
  \item 60^\circ
  \item 270^\circ
\end{multicols}
\end{enumerate}
\vspace{0.15cm}
\item ¿Cuánto mide cada ángulo de un triángulo equilátero?
\begin{enumerate}[label=\textbf{\Alph*)}]
\begin{multicols}{2}
  \item 45^\circ
  \item 90^\circ
  \item 60^\circ
  \item 360^\circ
  \item 60^\circ
  \item 270^\circ
\end{multicols}
\end{enumerate}
\vspace{0.15cm}
\item Un ángulo que mide \(0^\circ\) se llama:
\begin{enumerate}[label=\textbf{\Alph*)}]
\begin{multicols}{2}
  \item Nulo
  \item 90^\circ
  \item 180^\circ
  \item 360^\circ
  \item 60^\circ
  \item 270^\circ
\end{multicols}
\end{enumerate}
\vspace{0.15cm}
\item El perímetro de un polígono es:
\begin{enumerate}[label=\textbf{\Alph*)}]
\begin{multicols}{2}
  \item 3
  \item La suma de sus lados
  \item 5
  \item 6
  \item 7
  \item 8
\end{multicols}
\end{enumerate}
\vspace{0.15cm}
\item La mediana de un triángulo es el segmento que une:
\begin{enumerate}[label=\textbf{\Alph*)}]
\begin{multicols}{2}
  \item Un vértice con el punto medio del lado opuesto
  \item Isósceles
  \item Equilátero
  \item 6
  \item 7
  \item 8
\end{multicols}
\end{enumerate}
\vspace{0.15cm}
\item La altura de un triángulo es el segmento que va:
\begin{enumerate}[label=\textbf{\Alph*)}]
\begin{multicols}{2}
  \item De un vértice al lado opuesto perpendicularmente
  \item Isósceles
  \item Equilátero
  \item 6
  \item 7
  \item 8
\end{multicols}
\end{enumerate}
\vspace{0.15cm}
\item La bisectriz de un ángulo es la recta que:
\begin{enumerate}[label=\textbf{\Alph*)}]
\begin{multicols}{2}
  \item Divide el ángulo en dos partes iguales
  \item 4
  \item 5
  \item 6
  \item 7
  \item 8
\end{multicols}
\end{enumerate}
\vspace{0.15cm}
\item ¿Cuántas diagonales tiene un cuadrado?
\begin{enumerate}[label=\textbf{\Alph*)}]
\begin{multicols}{2}
  \item 3
  \item 4
  \item 2
  \item 6
  \item 7
  \item 8
\end{multicols}
\end{enumerate}
\vspace{0.15cm}
\item ¿Cuántas diagonales tiene un pentágono?
\begin{enumerate}[label=\textbf{\Alph*)}]
\begin{multicols}{2}
  \item 3
  \item 4
  \item 5
  \item 6
  \item 7
  \item 8
\end{multicols}
\end{enumerate}
\vspace{0.15cm}
\item La suma de los ángulos externos de cualquier polígono convexo es:
\begin{enumerate}[label=\textbf{\Alph*)}]
\begin{multicols}{2}
  \item 3
  \item 4
  \item 5
  \item 360^\circ
  \item 7
  \item 8
\end{multicols}
\end{enumerate}
\vspace{0.15cm}
\end{enumerate}\n
\subsubsection*{SECCIÓN II: TEOREMA DE PITÁGORAS (41–90)}

\begin{enumerate}[label=\textbf{\arabic*.}]
\setcounter{enumi}{40}
\item El teorema de Pitágoras establece que:
\begin{enumerate}[label=\textbf{\Alph*)}]
\begin{multicols}{2}
  \item 5
  \item a^2+b^2=c^2
  \item Obtusángulo
  \item 15
  \item 25
  \item 3
\end{multicols}
\end{enumerate}
\vspace{0.15cm}
\item En el teorema de Pitágoras, \(c\) representa:
\begin{enumerate}[label=\textbf{\Alph*)}]
\begin{multicols}{2}
  \item 5
  \item La hipotenusa
  \item Obtusángulo
  \item 15
  \item 25
  \item 3
\end{multicols}
\end{enumerate}
\vspace{0.15cm}
\item Si los catetos miden 3 y 4, la hipotenusa mide:
\begin{enumerate}[label=\textbf{\Alph*)}]
\begin{multicols}{2}
  \item 5
  \item 10
  \item 12
  \item 15
  \item 25
  \item 3
\end{multicols}
\end{enumerate}
\vspace{0.15cm}
\item Si los catetos miden 6 y 8, la hipotenusa mide:
\begin{enumerate}[label=\textbf{\Alph*)}]
\begin{multicols}{2}
  \item 10
  \item 10
  \item 12
  \item 15
  \item 25
  \item 3
\end{multicols}
\end{enumerate}
\vspace{0.15cm}
\item Si los catetos miden 5 y 12, la hipotenusa mide:
\begin{enumerate}[label=\textbf{\Alph*)}]
\begin{multicols}{2}
  \item 13
  \item 10
  \item 12
  \item 15
  \item 25
  \item 3
\end{multicols}
\end{enumerate}
\vspace{0.15cm}
\item Si los catetos miden 8 y 15, la hipotenusa mide:
\begin{enumerate}[label=\textbf{\Alph*)}]
\begin{multicols}{2}
  \item 17
  \item 10
  \item 12
  \item 15
  \item 25
  \item 3
\end{multicols}
\end{enumerate}
\vspace{0.15cm}
\item Si los catetos miden 9 y 12, la hipotenusa mide:
\begin{enumerate}[label=\textbf{\Alph*)}]
\begin{multicols}{2}
  \item 15
  \item 10
  \item 12
  \item 15
  \item 25
  \item 3
\end{multicols}
\end{enumerate}
\vspace{0.15cm}
\item Si los catetos miden 7 y 24, la hipotenusa mide:
\begin{enumerate}[label=\textbf{\Alph*)}]
\begin{multicols}{2}
  \item 25
  \item 10
  \item 12
  \item 15
  \item 25
  \item 3
\end{multicols}
\end{enumerate}
\vspace{0.15cm}
\item Si un cateto mide 3 y la hipotenusa 5, el otro cateto mide:
\begin{enumerate}[label=\textbf{\Alph*)}]
\begin{multicols}{2}
  \item 5
  \item 4
  \item 12
  \item 15
  \item 25
  \item 3
\end{multicols}
\end{enumerate}
\vspace{0.15cm}
\item Si un cateto mide 6 y la hipotenusa 10, el otro cateto mide:
\begin{enumerate}[label=\textbf{\Alph*)}]
\begin{multicols}{2}
  \item 5
  \item 8
  \item 12
  \item 15
  \item 25
  \item 3
\end{multicols}
\end{enumerate}
\vspace{0.15cm}
\item Si un cateto mide 5 y la hipotenusa 13, el otro cateto mide:
\begin{enumerate}[label=\textbf{\Alph*)}]
\begin{multicols}{2}
  \item 5
  \item 12
  \item 12
  \item 15
  \item 25
  \item 3
\end{multicols}
\end{enumerate}
\vspace{0.15cm}
\item Si un cateto mide 8 y la hipotenusa 17, el otro cateto mide:
\begin{enumerate}[label=\textbf{\Alph*)}]
\begin{multicols}{2}
  \item 5
  \item 15
  \item 12
  \item 15
  \item 25
  \item 3
\end{multicols}
\end{enumerate}
\vspace{0.15cm}
\item Si un cateto mide 9 y la hipotenusa 15, el otro cateto mide:
\begin{enumerate}[label=\textbf{\Alph*)}]
\begin{multicols}{2}
  \item 5
  \item 12
  \item 12
  \item 15
  \item 25
  \item 3
\end{multicols}
\end{enumerate}
\vspace{0.15cm}
\item Si un cateto mide 7 y la hipotenusa 25, el otro cateto mide:
\begin{enumerate}[label=\textbf{\Alph*)}]
\begin{multicols}{2}
  \item 5
  \item 24
  \item 12
  \item 15
  \item 25
  \item 3
\end{multicols}
\end{enumerate}
\vspace{0.15cm}
\item Si los catetos miden 1 y 1, la hipotenusa mide:
\begin{enumerate}[label=\textbf{\Alph*)}]
\begin{multicols}{2}
  \item 5
  \item \sqrt{2}
  \item 12
  \item 15
  \item 25
  \item 3
\end{multicols}
\end{enumerate}
\vspace{0.15cm}
\item Si los catetos miden 1 y \sqrt{3}, la hipotenusa mide:
\begin{enumerate}[label=\textbf{\Alph*)}]
\begin{multicols}{2}
  \item 5
  \item 2
  \item 12
  \item 15
  \item 25
  \item 3
\end{multicols}
\end{enumerate}
\vspace{0.15cm}
\item Si los catetos miden 2 y 2, la hipotenusa mide:
\begin{enumerate}[label=\textbf{\Alph*)}]
\begin{multicols}{2}
  \item 5
  \item 2\sqrt{2}
  \item 12
  \item 15
  \item 25
  \item 3
\end{multicols}
\end{enumerate}
\vspace{0.15cm}
\item Si los catetos miden \sqrt{2} y \sqrt{2}, la hipotenusa mide:
\begin{enumerate}[label=\textbf{\Alph*)}]
\begin{multicols}{2}
  \item 5
  \item 2
  \item 12
  \item 15
  \item 25
  \item 3
\end{multicols}
\end{enumerate}
\vspace{0.15cm}
\item Un triángulo con lados 3, 4, 5 es:
\begin{enumerate}[label=\textbf{\Alph*)}]
\begin{multicols}{2}
  \item 5
  \item Rectángulo
  \item Obtusángulo
  \item 15
  \item 25
  \item 3
\end{multicols}
\end{enumerate}
\vspace{0.15cm}
\item Un triángulo con lados 5, 12, 13 es:
\begin{enumerate}[label=\textbf{\Alph*)}]
\begin{multicols}{2}
  \item 5
  \item Rectángulo
  \item Obtusángulo
  \item 15
  \item 25
  \item 3
\end{multicols}
\end{enumerate}
\vspace{0.15cm}
\item Un triángulo con lados 8, 15, 17 es:
\begin{enumerate}[label=\textbf{\Alph*)}]
\begin{multicols}{2}
  \item 5
  \item Rectángulo
  \item Obtusángulo
  \item 15
  \item 25
  \item 3
\end{multicols}
\end{enumerate}
\vspace{0.15cm}
\item Un triángulo con lados 7, 24, 25 es:
\begin{enumerate}[label=\textbf{\Alph*)}]
\begin{multicols}{2}
  \item 5
  \item Rectángulo
  \item Obtusángulo
  \item 15
  \item 25
  \item 3
\end{multicols}
\end{enumerate}
\vspace{0.15cm}
\item Un triángulo con lados 9, 12, 15 es:
\begin{enumerate}[label=\textbf{\Alph*)}]
\begin{multicols}{2}
  \item 5
  \item Rectángulo
  \item Obtusángulo
  \item 15
  \item 25
  \item 3
\end{multicols}
\end{enumerate}
\vspace{0.15cm}
\item Un triángulo con lados 6, 8, 10 es:
\begin{enumerate}[label=\textbf{\Alph*)}]
\begin{multicols}{2}
  \item 5
  \item Rectángulo
  \item Obtusángulo
  \item 15
  \item 25
  \item 3
\end{multicols}
\end{enumerate}
\vspace{0.15cm}
\item Un triángulo con lados 5, 5, 6 es:
\begin{enumerate}[label=\textbf{\Alph*)}]
\begin{multicols}{2}
  \item Acutángulo
  \item Rectángulo
  \item Obtusángulo
  \item 15
  \item 25
  \item 3
\end{multicols}
\end{enumerate}
\vspace{0.15cm}
\item Si \(a^2+b^2<c^2\), el triángulo es:
\begin{enumerate}[label=\textbf{\Alph*)}]
\begin{multicols}{2}
  \item 5
  \item Rectángulo
  \item Obtusángulo
  \item 15
  \item 25
  \item 3
\end{multicols}
\end{enumerate}
\vspace{0.15cm}
\item Si \(a^2+b^2>c^2\), el triángulo es:
\begin{enumerate}[label=\textbf{\Alph*)}]
\begin{multicols}{2}
  \item Acutángulo
  \item Rectángulo
  \item Obtusángulo
  \item 15
  \item 25
  \item 3
\end{multicols}
\end{enumerate}
\vspace{0.15cm}
\item La diagonal de un cuadrado de lado 1 mide:
\begin{enumerate}[label=\textbf{\Alph*)}]
\begin{multicols}{2}
  \item 5
  \item \sqrt{2}
  \item Obtusángulo
  \item 15
  \item 25
  \item 3
\end{multicols}
\end{enumerate}
\vspace{0.15cm}
\item La diagonal de un cuadrado de lado 5 mide:
\begin{enumerate}[label=\textbf{\Alph*)}]
\begin{multicols}{2}
  \item 5
  \item 5\sqrt{2}
  \item Obtusángulo
  \item 15
  \item 25
  \item 3
\end{multicols}
\end{enumerate}
\vspace{0.15cm}
\item La diagonal de un rectángulo de lados 3 y 4 mide:
\begin{enumerate}[label=\textbf{\Alph*)}]
\begin{multicols}{2}
  \item 5
  \item Rectángulo
  \item Obtusángulo
  \item 15
  \item 25
  \item 3
\end{multicols}
\end{enumerate}
\vspace{0.15cm}
\item La diagonal de un rectángulo de lados 6 y 8 mide:
\begin{enumerate}[label=\textbf{\Alph*)}]
\begin{multicols}{2}
  \item 10
  \item Rectángulo
  \item Obtusángulo
  \item 15
  \item 25
  \item 3
\end{multicols}
\end{enumerate}
\vspace{0.15cm}
\item La diagonal de un rectángulo de lados 5 y 12 mide:
\begin{enumerate}[label=\textbf{\Alph*)}]
\begin{multicols}{2}
  \item 13
  \item Rectángulo
  \item Obtusángulo
  \item 15
  \item 25
  \item 3
\end{multicols}
\end{enumerate}
\vspace{0.15cm}
\item La altura de un triángulo equilátero de lado 2 mide:
\begin{enumerate}[label=\textbf{\Alph*)}]
\begin{multicols}{2}
  \item 5
  \item \sqrt{3}
  \item Obtusángulo
  \item 15
  \item 25
  \item 3
\end{multicols}
\end{enumerate}
\vspace{0.15cm}
\item La altura de un triángulo equilátero de lado 6 mide:
\begin{enumerate}[label=\textbf{\Alph*)}]
\begin{multicols}{2}
  \item 5
  \item 3\sqrt{3}
  \item Obtusángulo
  \item 15
  \item 25
  \item 3
\end{multicols}
\end{enumerate}
\vspace{0.15cm}
\item La altura de un triángulo equilátero de lado 10 mide:
\begin{enumerate}[label=\textbf{\Alph*)}]
\begin{multicols}{2}
  \item 5
  \item 5\sqrt{3}
  \item Obtusángulo
  \item 15
  \item 25
  \item 3
\end{multicols}
\end{enumerate}
\vspace{0.15cm}
\item Si la hipotenusa mide 10 y un cateto mide 6, el otro cateto mide:
\begin{enumerate}[label=\textbf{\Alph*)}]
\begin{multicols}{2}
  \item 5
  \item 8
  \item 12
  \item 15
  \item 25
  \item 3
\end{multicols}
\end{enumerate}
\vspace{0.15cm}
\item Si la hipotenusa mide 13 y un cateto mide 5, el otro cateto mide:
\begin{enumerate}[label=\textbf{\Alph*)}]
\begin{multicols}{2}
  \item 5
  \item 12
  \item 12
  \item 15
  \item 25
  \item 3
\end{multicols}
\end{enumerate}
\vspace{0.15cm}
\item Si la hipotenusa mide 17 y un cateto mide 8, el otro cateto mide:
\begin{enumerate}[label=\textbf{\Alph*)}]
\begin{multicols}{2}
  \item 5
  \item 15
  \item 12
  \item 15
  \item 25
  \item 3
\end{multicols}
\end{enumerate}
\vspace{0.15cm}
\item Si la hipotenusa mide 25 y un cateto mide 7, el otro cateto mide:
\begin{enumerate}[label=\textbf{\Alph*)}]
\begin{multicols}{2}
  \item 5
  \item 24
  \item 12
  \item 15
  \item 25
  \item 3
\end{multicols}
\end{enumerate}
\vspace{0.15cm}
\item Si la hipotenusa mide 15 y un cateto mide 9, el otro cateto mide:
\begin{enumerate}[label=\textbf{\Alph*)}]
\begin{multicols}{2}
  \item 5
  \item 12
  \item 12
  \item 15
  \item 25
  \item 3
\end{multicols}
\end{enumerate}
\vspace{0.15cm}
\item La distancia entre los puntos (0,0) y (3,4) es:
\begin{enumerate}[label=\textbf{\Alph*)}]
\begin{multicols}{2}
  \item 5
  \item Rectángulo
  \item Obtusángulo
  \item 15
  \item 25
  \item 3
\end{multicols}
\end{enumerate}
\vspace{0.15cm}
\item La distancia entre los puntos (0,0) y (6,8) es:
\begin{enumerate}[label=\textbf{\Alph*)}]
\begin{multicols}{2}
  \item 10
  \item Rectángulo
  \item Obtusángulo
  \item 15
  \item 25
  \item 3
\end{multicols}
\end{enumerate}
\vspace{0.15cm}
\item La distancia entre los puntos (0,0) y (5,12) es:
\begin{enumerate}[label=\textbf{\Alph*)}]
\begin{multicols}{2}
  \item 13
  \item Rectángulo
  \item Obtusángulo
  \item 15
  \item 25
  \item 3
\end{multicols}
\end{enumerate}
\vspace{0.15cm}
\item La distancia entre los puntos (1,1) y (4,5) es:
\begin{enumerate}[label=\textbf{\Alph*)}]
\begin{multicols}{2}
  \item 5
  \item Rectángulo
  \item Obtusángulo
  \item 15
  \item 25
  \item 3
\end{multicols}
\end{enumerate}
\vspace{0.15cm}
\item La distancia entre los puntos (-1,-1) y (2,3) es:
\begin{enumerate}[label=\textbf{\Alph*)}]
\begin{multicols}{2}
  \item 5
  \item Rectángulo
  \item Obtusángulo
  \item 15
  \item 25
  \item 3
\end{multicols}
\end{enumerate}
\vspace{0.15cm}
\item La distancia entre los puntos (2,3) y (5,7) es:
\begin{enumerate}[label=\textbf{\Alph*)}]
\begin{multicols}{2}
  \item 5
  \item Rectángulo
  \item Obtusángulo
  \item 15
  \item 25
  \item 3
\end{multicols}
\end{enumerate}
\vspace{0.15cm}
\item Si un poste de 3 m proyecta una sombra de 4 m, la distancia de la punta a la sombra es:
\begin{enumerate}[label=\textbf{\Alph*)}]
\begin{multicols}{2}
  \item 5 m
  \item Rectángulo
  \item Obtusángulo
  \item 15
  \item 25
  \item 3
\end{multicols}
\end{enumerate}
\vspace{0.15cm}
\item Una escalera de 5 m se apoya a 3 m de la pared. ¿A qué altura llega?
\begin{enumerate}[label=\textbf{\Alph*)}]
\begin{multicols}{2}
  \item 4 m
  \item Rectángulo
  \item Obtusángulo
  \item 15
  \item 25
  \item 3
\end{multicols}
\end{enumerate}
\vspace{0.15cm}
\item Una escalera de 13 m se apoya a 5 m de la pared. ¿A qué altura llega?
\begin{enumerate}[label=\textbf{\Alph*)}]
\begin{multicols}{2}
  \item 12 m
  \item Rectángulo
  \item Obtusángulo
  \item 15
  \item 25
  \item 3
\end{multicols}
\end{enumerate}
\vspace{0.15cm}
\item Un triángulo rectángulo tiene catetos 9 y 12. La hipotenusa mide:
\begin{enumerate}[label=\textbf{\Alph*)}]
\begin{multicols}{2}
  \item 15
  \item 10
  \item 12
  \item 15
  \item 25
  \item 3
\end{multicols}
\end{enumerate}
\vspace{0.15cm}
\end{enumerate}\n
\subsubsection*{SECCIÓN III: RAZONES TRIGONOMÉTRICAS (91–160)}

\begin{enumerate}[label=\textbf{\arabic*.}]
\setcounter{enumi}{90}
\item En un triángulo rectángulo, \(\sen\theta\) se define como:
\begin{enumerate}[label=\textbf{\Alph*)}]
\begin{multicols}{2}
  \item Cateto opuesto / Hipotenusa
  \item 1
  \item -1
  \item -\frac{1}{2}
  \item -\frac{\sqrt{3}}{2}
  \item 0
\end{multicols}
\end{enumerate}
\vspace{0.15cm}
\item ¿A qué es igual \(\cos\theta se define como:\)?
\begin{enumerate}[label=\textbf{\Alph*)}]
\begin{multicols}{2}
  \item 0
  \item Cateto adyacente / Hipotenusa
  \item -1
  \item -\frac{1}{2}
  \item -\frac{\sqrt{3}}{2}
  \item 0
\end{multicols}
\end{enumerate}
\vspace{0.15cm}
\item ¿A qué es igual \(\tan\theta se define como:\)?
\begin{enumerate}[label=\textbf{\Alph*)}]
\begin{multicols}{2}
  \item 0
  \item 1
  \item Cateto opuesto / Cateto adyacente
  \item -\frac{1}{2}
  \item -\frac{\sqrt{3}}{2}
  \item 0
\end{multicols}
\end{enumerate}
\vspace{0.15cm}
\item ¿A qué es igual \(\csc\theta es el recíproco de:\)?
\begin{enumerate}[label=\textbf{\Alph*)}]
\begin{multicols}{2}
  \item \sen\theta
  \item 1
  \item -1
  \item -\frac{1}{2}
  \item -\frac{\sqrt{3}}{2}
  \item 0
\end{multicols}
\end{enumerate}
\vspace{0.15cm}
\item ¿A qué es igual \(\sec\theta es el recíproco de:\)?
\begin{enumerate}[label=\textbf{\Alph*)}]
\begin{multicols}{2}
  \item 0
  \item \cos\theta
  \item -1
  \item -\frac{1}{2}
  \item -\frac{\sqrt{3}}{2}
  \item 0
\end{multicols}
\end{enumerate}
\vspace{0.15cm}
\item ¿A qué es igual \(\cot\theta es el recíproco de:\)?
\begin{enumerate}[label=\textbf{\Alph*)}]
\begin{multicols}{2}
  \item 0
  \item 1
  \item \tan\theta
  \item -\frac{1}{2}
  \item -\frac{\sqrt{3}}{2}
  \item 0
\end{multicols}
\end{enumerate}
\vspace{0.15cm}
\item ¿A qué es igual \(\sen 30^\circ =\)?
\begin{enumerate}[label=\textbf{\Alph*)}]
\begin{multicols}{2}
  \item \frac{1}{2}
  \item \frac{\sqrt{2}}{2}
  \item \frac{\sqrt{3}}{2}
  \item -\frac{1}{2}
  \item -\frac{\sqrt{3}}{2}
  \item 0
\end{multicols}
\end{enumerate}
\vspace{0.15cm}
\item ¿A qué es igual \(\cos 30^\circ =\)?
\begin{enumerate}[label=\textbf{\Alph*)}]
\begin{multicols}{2}
  \item \frac{1}{2}
  \item \frac{\sqrt{2}}{2}
  \item \frac{\sqrt{3}}{2}
  \item -\frac{1}{2}
  \item -\frac{\sqrt{3}}{2}
  \item 0
\end{multicols}
\end{enumerate}
\vspace{0.15cm}
\item ¿A qué es igual \(\tan 30^\circ =\)?
\begin{enumerate}[label=\textbf{\Alph*)}]
\begin{multicols}{2}
  \item \frac{1}{2}
  \item \frac{\sqrt{2}}{2}
  \item \frac{\sqrt{3}}{2}
  \item \frac{\sqrt{3}}{3}
  \item -\frac{\sqrt{3}}{2}
  \item 0
\end{multicols}
\end{enumerate}
\vspace{0.15cm}
\item ¿A qué es igual \(\sen 45^\circ =\)?
\begin{enumerate}[label=\textbf{\Alph*)}]
\begin{multicols}{2}
  \item \frac{1}{2}
  \item \frac{\sqrt{2}}{2}
  \item \frac{\sqrt{3}}{2}
  \item -\frac{1}{2}
  \item -\frac{\sqrt{3}}{2}
  \item 0
\end{multicols}
\end{enumerate}
\vspace{0.15cm}
\item ¿A qué es igual \(\cos 45^\circ =\)?
\begin{enumerate}[label=\textbf{\Alph*)}]
\begin{multicols}{2}
  \item \frac{1}{2}
  \item \frac{\sqrt{2}}{2}
  \item \frac{\sqrt{3}}{2}
  \item -\frac{1}{2}
  \item -\frac{\sqrt{3}}{2}
  \item 0
\end{multicols}
\end{enumerate}
\vspace{0.15cm}
\item ¿A qué es igual \(\tan 45^\circ =\)?
\begin{enumerate}[label=\textbf{\Alph*)}]
\begin{multicols}{2}
  \item 0
  \item 1
  \item -1
  \item 1
  \item -\frac{\sqrt{3}}{2}
  \item 0
\end{multicols}
\end{enumerate}
\vspace{0.15cm}
\item ¿A qué es igual \(\sen 60^\circ =\)?
\begin{enumerate}[label=\textbf{\Alph*)}]
\begin{multicols}{2}
  \item \frac{1}{2}
  \item \frac{\sqrt{2}}{2}
  \item \frac{\sqrt{3}}{2}
  \item -\frac{1}{2}
  \item -\frac{\sqrt{3}}{2}
  \item 0
\end{multicols}
\end{enumerate}
\vspace{0.15cm}
\item ¿A qué es igual \(\cos 60^\circ =\)?
\begin{enumerate}[label=\textbf{\Alph*)}]
\begin{multicols}{2}
  \item \frac{1}{2}
  \item \frac{\sqrt{2}}{2}
  \item \frac{\sqrt{3}}{2}
  \item -\frac{1}{2}
  \item -\frac{\sqrt{3}}{2}
  \item 0
\end{multicols}
\end{enumerate}
\vspace{0.15cm}
\item ¿A qué es igual \(\tan 60^\circ =\)?
\begin{enumerate}[label=\textbf{\Alph*)}]
\begin{multicols}{2}
  \item 0
  \item 1
  \item -1
  \item \sqrt{3}
  \item -\frac{\sqrt{3}}{2}
  \item 0
\end{multicols}
\end{enumerate}
\vspace{0.15cm}
\item ¿A qué es igual \(\sen 0^\circ =\)?
\begin{enumerate}[label=\textbf{\Alph*)}]
\begin{multicols}{2}
  \item 0
  \item 1
  \item -1
  \item -\frac{1}{2}
  \item -\frac{\sqrt{3}}{2}
  \item 0
\end{multicols}
\end{enumerate}
\vspace{0.15cm}
\item ¿A qué es igual \(\cos 0^\circ =\)?
\begin{enumerate}[label=\textbf{\Alph*)}]
\begin{multicols}{2}
  \item 0
  \item 1
  \item -1
  \item -\frac{1}{2}
  \item -\frac{\sqrt{3}}{2}
  \item 0
\end{multicols}
\end{enumerate}
\vspace{0.15cm}
\item ¿A qué es igual \(\tan 0^\circ =\)?
\begin{enumerate}[label=\textbf{\Alph*)}]
\begin{multicols}{2}
  \item 0
  \item 1
  \item -1
  \item -\frac{1}{2}
  \item -\frac{\sqrt{3}}{2}
  \item 0
\end{multicols}
\end{enumerate}
\vspace{0.15cm}
\item ¿A qué es igual \(\sen 90^\circ =\)?
\begin{enumerate}[label=\textbf{\Alph*)}]
\begin{multicols}{2}
  \item 0
  \item 1
  \item -1
  \item -\frac{1}{2}
  \item -\frac{\sqrt{3}}{2}
  \item 0
\end{multicols}
\end{enumerate}
\vspace{0.15cm}
\item ¿A qué es igual \(\cos 90^\circ =\)?
\begin{enumerate}[label=\textbf{\Alph*)}]
\begin{multicols}{2}
  \item 0
  \item 1
  \item -1
  \item -\frac{1}{2}
  \item -\frac{\sqrt{3}}{2}
  \item 0
\end{multicols}
\end{enumerate}
\vspace{0.15cm}
\item ¿A qué es igual \(\tan 90^\circ =\)?
\begin{enumerate}[label=\textbf{\Alph*)}]
\begin{multicols}{2}
  \item 0
  \item 1
  \item -1
  \item -\frac{1}{2}
  \item -\frac{\sqrt{3}}{2}
  \item Indefinido
\end{multicols}
\end{enumerate}
\vspace{0.15cm}
\item ¿A qué es igual \(\sen 180^\circ =\)?
\begin{enumerate}[label=\textbf{\Alph*)}]
\begin{multicols}{2}
  \item 0
  \item 1
  \item -1
  \item -\frac{1}{2}
  \item -\frac{\sqrt{3}}{2}
  \item 0
\end{multicols}
\end{enumerate}
\vspace{0.15cm}
\item ¿A qué es igual \(\cos 180^\circ =\)?
\begin{enumerate}[label=\textbf{\Alph*)}]
\begin{multicols}{2}
  \item 0
  \item 1
  \item -1
  \item -\frac{1}{2}
  \item -\frac{\sqrt{3}}{2}
  \item 0
\end{multicols}
\end{enumerate}
\vspace{0.15cm}
\item ¿A qué es igual \(\tan 180^\circ =\)?
\begin{enumerate}[label=\textbf{\Alph*)}]
\begin{multicols}{2}
  \item 0
  \item 1
  \item -1
  \item -\frac{1}{2}
  \item -\frac{\sqrt{3}}{2}
  \item 0
\end{multicols}
\end{enumerate}
\vspace{0.15cm}
\item Si \sen\theta = \frac{3}{5}, entonces \cos\theta =
\begin{enumerate}[label=\textbf{\Alph*)}]
\begin{multicols}{2}
  \item \frac{4}{5}
  \item \frac{\sqrt{2}}{2}
  \item \frac{\sqrt{3}}{2}
  \item -\frac{1}{2}
  \item -\frac{\sqrt{3}}{2}
  \item 0
\end{multicols}
\end{enumerate}
\vspace{0.15cm}
\item Si \cos\theta = \frac{4}{5}, entonces \sen\theta =
\begin{enumerate}[label=\textbf{\Alph*)}]
\begin{multicols}{2}
  \item \frac{3}{5}
  \item \frac{\sqrt{2}}{2}
  \item \frac{\sqrt{3}}{2}
  \item -\frac{1}{2}
  \item -\frac{\sqrt{3}}{2}
  \item 0
\end{multicols}
\end{enumerate}
\vspace{0.15cm}
\item Si \sen\theta = \frac{5}{13}, entonces \cos\theta =
\begin{enumerate}[label=\textbf{\Alph*)}]
\begin{multicols}{2}
  \item \frac{12}{13}
  \item \frac{\sqrt{2}}{2}
  \item \frac{\sqrt{3}}{2}
  \item -\frac{1}{2}
  \item -\frac{\sqrt{3}}{2}
  \item 0
\end{multicols}
\end{enumerate}
\vspace{0.15cm}
\item Si \cos\theta = \frac{12}{13}, entonces \sen\theta =
\begin{enumerate}[label=\textbf{\Alph*)}]
\begin{multicols}{2}
  \item \frac{5}{13}
  \item \frac{\sqrt{2}}{2}
  \item \frac{\sqrt{3}}{2}
  \item -\frac{1}{2}
  \item -\frac{\sqrt{3}}{2}
  \item 0
\end{multicols}
\end{enumerate}
\vspace{0.15cm}
\item Si \sen\theta = \frac{8}{17}, entonces \cos\theta =
\begin{enumerate}[label=\textbf{\Alph*)}]
\begin{multicols}{2}
  \item \frac{15}{17}
  \item \frac{\sqrt{2}}{2}
  \item \frac{\sqrt{3}}{2}
  \item -\frac{1}{2}
  \item -\frac{\sqrt{3}}{2}
  \item 0
\end{multicols}
\end{enumerate}
\vspace{0.15cm}
\item Si \cos\theta = \frac{15}{17}, entonces \sen\theta =
\begin{enumerate}[label=\textbf{\Alph*)}]
\begin{multicols}{2}
  \item \frac{8}{17}
  \item \frac{\sqrt{2}}{2}
  \item \frac{\sqrt{3}}{2}
  \item -\frac{1}{2}
  \item -\frac{\sqrt{3}}{2}
  \item 0
\end{multicols}
\end{enumerate}
\vspace{0.15cm}
\item Si \tan\theta = \frac{3}{4}, entonces \sen\theta =
\begin{enumerate}[label=\textbf{\Alph*)}]
\begin{multicols}{2}
  \item \frac{3}{5}
  \item \frac{\sqrt{2}}{2}
  \item \frac{\sqrt{3}}{2}
  \item -\frac{1}{2}
  \item -\frac{\sqrt{3}}{2}
  \item 0
\end{multicols}
\end{enumerate}
\vspace{0.15cm}
\item Si \tan\theta = \frac{3}{4}, entonces \cos\theta =
\begin{enumerate}[label=\textbf{\Alph*)}]
\begin{multicols}{2}
  \item \frac{1}{2}
  \item \frac{4}{5}
  \item \frac{\sqrt{3}}{2}
  \item -\frac{1}{2}
  \item -\frac{\sqrt{3}}{2}
  \item 0
\end{multicols}
\end{enumerate}
\vspace{0.15cm}
\item Si \tan\theta = \frac{5}{12}, entonces \sen\theta =
\begin{enumerate}[label=\textbf{\Alph*)}]
\begin{multicols}{2}
  \item \frac{5}{13}
  \item \frac{\sqrt{2}}{2}
  \item \frac{\sqrt{3}}{2}
  \item -\frac{1}{2}
  \item -\frac{\sqrt{3}}{2}
  \item 0
\end{multicols}
\end{enumerate}
\vspace{0.15cm}
\item Si \tan\theta = \frac{5}{12}, entonces \cos\theta =
\begin{enumerate}[label=\textbf{\Alph*)}]
\begin{multicols}{2}
  \item \frac{1}{2}
  \item \frac{12}{13}
  \item \frac{\sqrt{3}}{2}
  \item -\frac{1}{2}
  \item -\frac{\sqrt{3}}{2}
  \item 0
\end{multicols}
\end{enumerate}
\vspace{0.15cm}
\item Si \tan\theta = \frac{8}{15}, entonces \sen\theta =
\begin{enumerate}[label=\textbf{\Alph*)}]
\begin{multicols}{2}
  \item \frac{8}{17}
  \item \frac{\sqrt{2}}{2}
  \item \frac{\sqrt{3}}{2}
  \item -\frac{1}{2}
  \item -\frac{\sqrt{3}}{2}
  \item 0
\end{multicols}
\end{enumerate}
\vspace{0.15cm}
\item Si \tan\theta = \frac{8}{15}, entonces \cos\theta =
\begin{enumerate}[label=\textbf{\Alph*)}]
\begin{multicols}{2}
  \item \frac{1}{2}
  \item \frac{15}{17}
  \item \frac{\sqrt{3}}{2}
  \item -\frac{1}{2}
  \item -\frac{\sqrt{3}}{2}
  \item 0
\end{multicols}
\end{enumerate}
\vspace{0.15cm}
\item ¿A qué es igual \(\sen^2\theta + \cos^2\theta =\)?
\begin{enumerate}[label=\textbf{\Alph*)}]
\begin{multicols}{2}
  \item 0
  \item 1
  \item -1
  \item -\frac{1}{2}
  \item -\frac{\sqrt{3}}{2}
  \item 0
\end{multicols}
\end{enumerate}
\vspace{0.15cm}
\item ¿A qué es igual \(\tan\theta =\)?
\begin{enumerate}[label=\textbf{\Alph*)}]
\begin{multicols}{2}
  \item \frac{\sen\theta}{\cos\theta}
  \item \frac{\sqrt{2}}{2}
  \item \frac{\sqrt{3}}{2}
  \item -\frac{1}{2}
  \item -\frac{\sqrt{3}}{2}
  \item 0
\end{multicols}
\end{enumerate}
\vspace{0.15cm}
\item ¿A qué es igual \(\cot\theta =\)?
\begin{enumerate}[label=\textbf{\Alph*)}]
\begin{multicols}{2}
  \item \frac{1}{2}
  \item \frac{\cos\theta}{\sen\theta}
  \item \frac{\sqrt{3}}{2}
  \item -\frac{1}{2}
  \item -\frac{\sqrt{3}}{2}
  \item 0
\end{multicols}
\end{enumerate}
\vspace{0.15cm}
\item ¿A qué es igual \(\sec\theta =\)?
\begin{enumerate}[label=\textbf{\Alph*)}]
\begin{multicols}{2}
  \item \frac{1}{2}
  \item \frac{1}{\cos\theta}
  \item \frac{\sqrt{3}}{2}
  \item -\frac{1}{2}
  \item -\frac{\sqrt{3}}{2}
  \item 0
\end{multicols}
\end{enumerate}
\vspace{0.15cm}
\item ¿A qué es igual \(\csc\theta =\)?
\begin{enumerate}[label=\textbf{\Alph*)}]
\begin{multicols}{2}
  \item \frac{1}{\sen\theta}
  \item \frac{\sqrt{2}}{2}
  \item \frac{\sqrt{3}}{2}
  \item -\frac{1}{2}
  \item -\frac{\sqrt{3}}{2}
  \item 0
\end{multicols}
\end{enumerate}
\vspace{0.15cm}
\item ¿A qué es igual \(1+\tan^2\theta =\)?
\begin{enumerate}[label=\textbf{\Alph*)}]
\begin{multicols}{2}
  \item \sec^2\theta
  \item 1
  \item -1
  \item -\frac{1}{2}
  \item -\frac{\sqrt{3}}{2}
  \item 0
\end{multicols}
\end{enumerate}
\vspace{0.15cm}
\item ¿A qué es igual \(1+\cot^2\theta =\)?
\begin{enumerate}[label=\textbf{\Alph*)}]
\begin{multicols}{2}
  \item 0
  \item \csc^2\theta
  \item -1
  \item -\frac{1}{2}
  \item -\frac{\sqrt{3}}{2}
  \item 0
\end{multicols}
\end{enumerate}
\vspace{0.15cm}
\item ¿A qué es igual \(\sen(90^\circ-\theta) =\)?
\begin{enumerate}[label=\textbf{\Alph*)}]
\begin{multicols}{2}
  \item 0
  \item \cos\theta
  \item -1
  \item -\frac{1}{2}
  \item -\frac{\sqrt{3}}{2}
  \item 0
\end{multicols}
\end{enumerate}
\vspace{0.15cm}
\item ¿A qué es igual \(\cos(90^\circ-\theta) =\)?
\begin{enumerate}[label=\textbf{\Alph*)}]
\begin{multicols}{2}
  \item \sen\theta
  \item 1
  \item -1
  \item -\frac{1}{2}
  \item -\frac{\sqrt{3}}{2}
  \item 0
\end{multicols}
\end{enumerate}
\vspace{0.15cm}
\item ¿A qué es igual \(\tan(90^\circ-\theta) =\)?
\begin{enumerate}[label=\textbf{\Alph*)}]
\begin{multicols}{2}
  \item 0
  \item 1
  \item -1
  \item \cot\theta
  \item -\frac{\sqrt{3}}{2}
  \item 0
\end{multicols}
\end{enumerate}
\vspace{0.15cm}
\item ¿A qué es igual \(\sen(-\theta) =\)?
\begin{enumerate}[label=\textbf{\Alph*)}]
\begin{multicols}{2}
  \item 0
  \item -\sen\theta
  \item -1
  \item -\frac{1}{2}
  \item -\frac{\sqrt{3}}{2}
  \item 0
\end{multicols}
\end{enumerate}
\vspace{0.15cm}
\item ¿A qué es igual \(\cos(-\theta) =\)?
\begin{enumerate}[label=\textbf{\Alph*)}]
\begin{multicols}{2}
  \item 0
  \item 1
  \item \cos\theta
  \item -\frac{1}{2}
  \item -\frac{\sqrt{3}}{2}
  \item 0
\end{multicols}
\end{enumerate}
\vspace{0.15cm}
\item ¿A qué es igual \(\tan(-\theta) =\)?
\begin{enumerate}[label=\textbf{\Alph*)}]
\begin{multicols}{2}
  \item 0
  \item 1
  \item -1
  \item -\frac{1}{2}
  \item -\frac{\sqrt{3}}{2}
  \item 0
\end{multicols}
\end{enumerate}
\vspace{0.15cm}
\item ¿A qué es igual \(\sen(180^\circ-\theta) =\)?
\begin{enumerate}[label=\textbf{\Alph*)}]
\begin{multicols}{2}
  \item \sen\theta
  \item 1
  \item -1
  \item -\frac{1}{2}
  \item -\frac{\sqrt{3}}{2}
  \item 0
\end{multicols}
\end{enumerate}
\vspace{0.15cm}
\item ¿A qué es igual \(\cos(180^\circ-\theta) =\)?
\begin{enumerate}[label=\textbf{\Alph*)}]
\begin{multicols}{2}
  \item 0
  \item 1
  \item -1
  \item -\cos\theta
  \item -\frac{\sqrt{3}}{2}
  \item 0
\end{multicols}
\end{enumerate}
\vspace{0.15cm}
\item ¿A qué es igual \(\tan(180^\circ-\theta) =\)?
\begin{enumerate}[label=\textbf{\Alph*)}]
\begin{multicols}{2}
  \item 0
  \item -\tan\theta
  \item -1
  \item -\frac{1}{2}
  \item -\frac{\sqrt{3}}{2}
  \item 0
\end{multicols}
\end{enumerate}
\vspace{0.15cm}
\item ¿A qué es igual \(\sen(180^\circ+\theta) =\)?
\begin{enumerate}[label=\textbf{\Alph*)}]
\begin{multicols}{2}
  \item 0
  \item -\sen\theta
  \item -1
  \item -\frac{1}{2}
  \item -\frac{\sqrt{3}}{2}
  \item 0
\end{multicols}
\end{enumerate}
\vspace{0.15cm}
\item ¿A qué es igual \(\cos(180^\circ+\theta) =\)?
\begin{enumerate}[label=\textbf{\Alph*)}]
\begin{multicols}{2}
  \item 0
  \item 1
  \item -1
  \item -\cos\theta
  \item -\frac{\sqrt{3}}{2}
  \item 0
\end{multicols}
\end{enumerate}
\vspace{0.15cm}
\item ¿A qué es igual \(\sen(360^\circ-\theta) =\)?
\begin{enumerate}[label=\textbf{\Alph*)}]
\begin{multicols}{2}
  \item 0
  \item -\sen\theta
  \item -1
  \item -\frac{1}{2}
  \item -\frac{\sqrt{3}}{2}
  \item 0
\end{multicols}
\end{enumerate}
\vspace{0.15cm}
\item ¿A qué es igual \(\cos(360^\circ-\theta) =\)?
\begin{enumerate}[label=\textbf{\Alph*)}]
\begin{multicols}{2}
  \item 0
  \item 1
  \item \cos\theta
  \item -\frac{1}{2}
  \item -\frac{\sqrt{3}}{2}
  \item 0
\end{multicols}
\end{enumerate}
\vspace{0.15cm}
\item ¿A qué es igual \(\sen 120^\circ =\)?
\begin{enumerate}[label=\textbf{\Alph*)}]
\begin{multicols}{2}
  \item \frac{1}{2}
  \item \frac{\sqrt{2}}{2}
  \item \frac{\sqrt{3}}{2}
  \item -\frac{1}{2}
  \item -\frac{\sqrt{3}}{2}
  \item 0
\end{multicols}
\end{enumerate}
\vspace{0.15cm}
\item ¿A qué es igual \(\cos 120^\circ =\)?
\begin{enumerate}[label=\textbf{\Alph*)}]
\begin{multicols}{2}
  \item \frac{1}{2}
  \item -\frac{1}{2}
  \item \frac{\sqrt{3}}{2}
  \item -\frac{1}{2}
  \item -\frac{\sqrt{3}}{2}
  \item 0
\end{multicols}
\end{enumerate}
\vspace{0.15cm}
\item ¿A qué es igual \(\sen 150^\circ =\)?
\begin{enumerate}[label=\textbf{\Alph*)}]
\begin{multicols}{2}
  \item \frac{1}{2}
  \item \frac{\sqrt{2}}{2}
  \item \frac{\sqrt{3}}{2}
  \item -\frac{1}{2}
  \item -\frac{\sqrt{3}}{2}
  \item 0
\end{multicols}
\end{enumerate}
\vspace{0.15cm}
\item ¿A qué es igual \(\cos 150^\circ =\)?
\begin{enumerate}[label=\textbf{\Alph*)}]
\begin{multicols}{2}
  \item \frac{1}{2}
  \item \frac{\sqrt{2}}{2}
  \item \frac{\sqrt{3}}{2}
  \item -\frac{\sqrt{3}}{2}
  \item -\frac{\sqrt{3}}{2}
  \item 0
\end{multicols}
\end{enumerate}
\vspace{0.15cm}
\item ¿A qué es igual \(\sen 210^\circ =\)?
\begin{enumerate}[label=\textbf{\Alph*)}]
\begin{multicols}{2}
  \item \frac{1}{2}
  \item -\frac{1}{2}
  \item \frac{\sqrt{3}}{2}
  \item -\frac{1}{2}
  \item -\frac{\sqrt{3}}{2}
  \item 0
\end{multicols}
\end{enumerate}
\vspace{0.15cm}
\item ¿A qué es igual \(\cos 210^\circ =\)?
\begin{enumerate}[label=\textbf{\Alph*)}]
\begin{multicols}{2}
  \item \frac{1}{2}
  \item \frac{\sqrt{2}}{2}
  \item \frac{\sqrt{3}}{2}
  \item -\frac{\sqrt{3}}{2}
  \item -\frac{\sqrt{3}}{2}
  \item 0
\end{multicols}
\end{enumerate}
\vspace{0.15cm}
\item ¿A qué es igual \(\sen 240^\circ =\)?
\begin{enumerate}[label=\textbf{\Alph*)}]
\begin{multicols}{2}
  \item \frac{1}{2}
  \item \frac{\sqrt{2}}{2}
  \item \frac{\sqrt{3}}{2}
  \item -\frac{\sqrt{3}}{2}
  \item -\frac{\sqrt{3}}{2}
  \item 0
\end{multicols}
\end{enumerate}
\vspace{0.15cm}
\item ¿A qué es igual \(\cos 240^\circ =\)?
\begin{enumerate}[label=\textbf{\Alph*)}]
\begin{multicols}{2}
  \item \frac{1}{2}
  \item -\frac{1}{2}
  \item \frac{\sqrt{3}}{2}
  \item -\frac{1}{2}
  \item -\frac{\sqrt{3}}{2}
  \item 0
\end{multicols}
\end{enumerate}
\vspace{0.15cm}
\item ¿A qué es igual \(\sen 300^\circ =\)?
\begin{enumerate}[label=\textbf{\Alph*)}]
\begin{multicols}{2}
  \item \frac{1}{2}
  \item \frac{\sqrt{2}}{2}
  \item \frac{\sqrt{3}}{2}
  \item -\frac{\sqrt{3}}{2}
  \item -\frac{\sqrt{3}}{2}
  \item 0
\end{multicols}
\end{enumerate}
\vspace{0.15cm}
\item ¿A qué es igual \(\cos 300^\circ =\)?
\begin{enumerate}[label=\textbf{\Alph*)}]
\begin{multicols}{2}
  \item \frac{1}{2}
  \item \frac{\sqrt{2}}{2}
  \item \frac{\sqrt{3}}{2}
  \item -\frac{1}{2}
  \item -\frac{\sqrt{3}}{2}
  \item 0
\end{multicols}
\end{enumerate}
\vspace{0.15cm}
\item ¿A qué es igual \(\sen 330^\circ =\)?
\begin{enumerate}[label=\textbf{\Alph*)}]
\begin{multicols}{2}
  \item \frac{1}{2}
  \item -\frac{1}{2}
  \item \frac{\sqrt{3}}{2}
  \item -\frac{1}{2}
  \item -\frac{\sqrt{3}}{2}
  \item 0
\end{multicols}
\end{enumerate}
\vspace{0.15cm}
\item ¿A qué es igual \(\cos 330^\circ =\)?
\begin{enumerate}[label=\textbf{\Alph*)}]
\begin{multicols}{2}
  \item \frac{1}{2}
  \item \frac{\sqrt{2}}{2}
  \item \frac{\sqrt{3}}{2}
  \item -\frac{1}{2}
  \item -\frac{\sqrt{3}}{2}
  \item 0
\end{multicols}
\end{enumerate}
\vspace{0.15cm}
\item ¿A qué es igual \(\sen 135^\circ =\)?
\begin{enumerate}[label=\textbf{\Alph*)}]
\begin{multicols}{2}
  \item \frac{1}{2}
  \item \frac{\sqrt{2}}{2}
  \item \frac{\sqrt{3}}{2}
  \item -\frac{1}{2}
  \item -\frac{\sqrt{3}}{2}
  \item 0
\end{multicols}
\end{enumerate}
\vspace{0.15cm}
\item ¿A qué es igual \(\cos 135^\circ =\)?
\begin{enumerate}[label=\textbf{\Alph*)}]
\begin{multicols}{2}
  \item \frac{1}{2}
  \item \frac{\sqrt{2}}{2}
  \item -\frac{\sqrt{2}}{2}
  \item -\frac{1}{2}
  \item -\frac{\sqrt{3}}{2}
  \item 0
\end{multicols}
\end{enumerate}
\vspace{0.15cm}
\end{enumerate}\n
\subsubsection*{SECCIÓN IV: RESOLUCIÓN DE TRIÁNGULOS RECTÁNGULOS (161–210)}

\begin{enumerate}[label=\textbf{\arabic*.}]
\setcounter{enumi}{160}
\item Si \sen\theta=3/5, cateto opuesto=3, hipotenusa=5. El cateto adyacente mide:
\begin{enumerate}[label=\textbf{\Alph*)}]
\begin{multicols}{2}
  \item 5
  \item 4
  \item 3
  \item 10
  \item 2
  \item 1
\end{multicols}
\end{enumerate}
\vspace{0.15cm}
\item Si \cos\theta=4/5, cateto adyacente=4, hipotenusa=5. El cateto opuesto mide:
\begin{enumerate}[label=\textbf{\Alph*)}]
\begin{multicols}{2}
  \item 3
  \item 4
  \item 3
  \item 10
  \item 2
  \item 1
\end{multicols}
\end{enumerate}
\vspace{0.15cm}
\item Si \tan\theta=3/4, catetos 3 y 4. La hipotenusa mide:
\begin{enumerate}[label=\textbf{\Alph*)}]
\begin{multicols}{2}
  \item 5
  \item 4
  \item 3
  \item 10
  \item 2
  \item 1
\end{multicols}
\end{enumerate}
\vspace{0.15cm}
\item Hipotenusa 10 y ángulo 30^\circ, el cateto opuesto mide:
\begin{enumerate}[label=\textbf{\Alph*)}]
\begin{multicols}{2}
  \item 5
  \item 4
  \item 3
  \item 10
  \item 2
  \item 1
\end{multicols}
\end{enumerate}
\vspace{0.15cm}
\item Hipotenusa 10 y ángulo 30^\circ, el cateto adyacente mide:
\begin{enumerate}[label=\textbf{\Alph*)}]
\begin{multicols}{2}
  \item 5
  \item 5\sqrt{3}
  \item 3
  \item 10
  \item 2
  \item 1
\end{multicols}
\end{enumerate}
\vspace{0.15cm}
\item Hipotenusa 10 y ángulo 60^\circ, el cateto opuesto mide:
\begin{enumerate}[label=\textbf{\Alph*)}]
\begin{multicols}{2}
  \item 5
  \item 5\sqrt{3}
  \item 3
  \item 10
  \item 2
  \item 1
\end{multicols}
\end{enumerate}
\vspace{0.15cm}
\item Hipotenusa 10 y ángulo 60^\circ, el cateto adyacente mide:
\begin{enumerate}[label=\textbf{\Alph*)}]
\begin{multicols}{2}
  \item 5
  \item 4
  \item 3
  \item 10
  \item 2
  \item 1
\end{multicols}
\end{enumerate}
\vspace{0.15cm}
\item Hipotenusa 8 y ángulo 45^\circ, cada cateto mide:
\begin{enumerate}[label=\textbf{\Alph*)}]
\begin{multicols}{2}
  \item 5
  \item 4\sqrt{2}
  \item 3
  \item 10
  \item 2
  \item 1
\end{multicols}
\end{enumerate}
\vspace{0.15cm}
\item Cateto opuesto 6 a ángulo 30^\circ, la hipotenusa mide:
\begin{enumerate}[label=\textbf{\Alph*)}]
\begin{multicols}{2}
  \item 12
  \item 4
  \item 3
  \item 10
  \item 2
  \item 1
\end{multicols}
\end{enumerate}
\vspace{0.15cm}
\item Cateto opuesto 6 a ángulo 30^\circ, el otro cateto mide:
\begin{enumerate}[label=\textbf{\Alph*)}]
\begin{multicols}{2}
  \item 5
  \item 6\sqrt{3}
  \item 3
  \item 10
  \item 2
  \item 1
\end{multicols}
\end{enumerate}
\vspace{0.15cm}
\item Cateto opuesto 5 a ángulo 45^\circ, la hipotenusa mide:
\begin{enumerate}[label=\textbf{\Alph*)}]
\begin{multicols}{2}
  \item 5
  \item 5\sqrt{2}
  \item 3
  \item 10
  \item 2
  \item 1
\end{multicols}
\end{enumerate}
\vspace{0.15cm}
\item Cateto opuesto 5 a ángulo 45^\circ, el otro cateto mide:
\begin{enumerate}[label=\textbf{\Alph*)}]
\begin{multicols}{2}
  \item 5
  \item 4
  \item 3
  \item 10
  \item 2
  \item 1
\end{multicols}
\end{enumerate}
\vspace{0.15cm}
\item Cateto adyacente 4 a ángulo 60^\circ, la hipotenusa mide:
\begin{enumerate}[label=\textbf{\Alph*)}]
\begin{multicols}{2}
  \item 8
  \item 4
  \item 3
  \item 10
  \item 2
  \item 1
\end{multicols}
\end{enumerate}
\vspace{0.15cm}
\item Cateto adyacente 4 a ángulo 60^\circ, el otro cateto mide:
\begin{enumerate}[label=\textbf{\Alph*)}]
\begin{multicols}{2}
  \item 5
  \item 4\sqrt{3}
  \item 3
  \item 10
  \item 2
  \item 1
\end{multicols}
\end{enumerate}
\vspace{0.15cm}
\item Si \sen\theta=1/2, entonces \theta =
\begin{enumerate}[label=\textbf{\Alph*)}]
\begin{multicols}{2}
  \item 30^\circ
  \item 30^\circ
  \item 60^\circ
  \item 90^\circ
  \item 180^\circ
  \item 1
\end{multicols}
\end{enumerate}
\vspace{0.15cm}
\item Si \cos\theta=1/2, entonces \theta =
\begin{enumerate}[label=\textbf{\Alph*)}]
\begin{multicols}{2}
  \item 30^\circ
  \item 45^\circ
  \item 60^\circ
  \item 60^\circ
  \item 180^\circ
  \item 1
\end{multicols}
\end{enumerate}
\vspace{0.15cm}
\item Si \tan\theta=1, entonces \theta =
\begin{enumerate}[label=\textbf{\Alph*)}]
\begin{multicols}{2}
  \item 30^\circ
  \item 45^\circ
  \item 45^\circ
  \item 90^\circ
  \item 180^\circ
  \item 1
\end{multicols}
\end{enumerate}
\vspace{0.15cm}
\item Si \sen\theta=\sqrt{3}/2, entonces \theta =
\begin{enumerate}[label=\textbf{\Alph*)}]
\begin{multicols}{2}
  \item 30^\circ
  \item 45^\circ
  \item 60^\circ
  \item 60^\circ
  \item 180^\circ
  \item 1
\end{multicols}
\end{enumerate}
\vspace{0.15cm}
\item Si \cos\theta=\sqrt{3}/2, entonces \theta =
\begin{enumerate}[label=\textbf{\Alph*)}]
\begin{multicols}{2}
  \item 30^\circ
  \item 30^\circ
  \item 60^\circ
  \item 90^\circ
  \item 180^\circ
  \item 1
\end{multicols}
\end{enumerate}
\vspace{0.15cm}
\item Si \sen\theta=\sqrt{2}/2, entonces \theta =
\begin{enumerate}[label=\textbf{\Alph*)}]
\begin{multicols}{2}
  \item 30^\circ
  \item 45^\circ
  \item 45^\circ
  \item 90^\circ
  \item 180^\circ
  \item 1
\end{multicols}
\end{enumerate}
\vspace{0.15cm}
\item Si \cos\theta=\sqrt{2}/2, entonces \theta =
\begin{enumerate}[label=\textbf{\Alph*)}]
\begin{multicols}{2}
  \item 30^\circ
  \item 45^\circ
  \item 45^\circ
  \item 90^\circ
  \item 180^\circ
  \item 1
\end{multicols}
\end{enumerate}
\vspace{0.15cm}
\item Si \tan\theta=\sqrt{3}, entonces \theta =
\begin{enumerate}[label=\textbf{\Alph*)}]
\begin{multicols}{2}
  \item 30^\circ
  \item 45^\circ
  \item 60^\circ
  \item 60^\circ
  \item 180^\circ
  \item 1
\end{multicols}
\end{enumerate}
\vspace{0.15cm}
\item Si \tan\theta=\sqrt{3}/3, entonces \theta =
\begin{enumerate}[label=\textbf{\Alph*)}]
\begin{multicols}{2}
  \item 30^\circ
  \item 30^\circ
  \item 60^\circ
  \item 90^\circ
  \item 180^\circ
  \item 1
\end{multicols}
\end{enumerate}
\vspace{0.15cm}
\item Si \sen\theta=0, entonces \theta =
\begin{enumerate}[label=\textbf{\Alph*)}]
\begin{multicols}{2}
  \item 0^\circ o 180^\circ
  \item 45^\circ
  \item 60^\circ
  \item 90^\circ
  \item 180^\circ
  \item 1
\end{multicols}
\end{enumerate}
\vspace{0.15cm}
\item Si \cos\theta=0, entonces \theta =
\begin{enumerate}[label=\textbf{\Alph*)}]
\begin{multicols}{2}
  \item 30^\circ
  \item 45^\circ
  \item 60^\circ
  \item 90^\circ
  \item 90^\circ
  \item 1
\end{multicols}
\end{enumerate}
\vspace{0.15cm}
\item Si \sen\theta=1, entonces \theta =
\begin{enumerate}[label=\textbf{\Alph*)}]
\begin{multicols}{2}
  \item 30^\circ
  \item 45^\circ
  \item 60^\circ
  \item 90^\circ
  \item 90^\circ
  \item 1
\end{multicols}
\end{enumerate}
\vspace{0.15cm}
\item Si \cos\theta=1, entonces \theta =
\begin{enumerate}[label=\textbf{\Alph*)}]
\begin{multicols}{2}
  \item 0^\circ
  \item 45^\circ
  \item 60^\circ
  \item 90^\circ
  \item 180^\circ
  \item 1
\end{multicols}
\end{enumerate}
\vspace{0.15cm}
\item Si \cos\theta=-1, entonces \theta =
\begin{enumerate}[label=\textbf{\Alph*)}]
\begin{multicols}{2}
  \item 30^\circ
  \item 45^\circ
  \item 60^\circ
  \item 90^\circ
  \item 180^\circ
  \item 180^\circ
\end{multicols}
\end{enumerate}
\vspace{0.15cm}
\item Ángulo de elevación 30° a 100 m de base, altura edificio:
\begin{enumerate}[label=\textbf{\Alph*)}]
\begin{multicols}{2}
  \item 5
  \item \frac{100\sqrt{3}}{3} m
  \item 3
  \item 10
  \item 2
  \item 1
\end{multicols}
\end{enumerate}
\vspace{0.15cm}
\item Ángulo de elevación 45° a 50 m de base, altura edificio:
\begin{enumerate}[label=\textbf{\Alph*)}]
\begin{multicols}{2}
  \item 5
  \item 50 m
  \item 3
  \item 10
  \item 2
  \item 1
\end{multicols}
\end{enumerate}
\vspace{0.15cm}
\item Ángulo de elevación 60° a 30 m de base, altura edificio:
\begin{enumerate}[label=\textbf{\Alph*)}]
\begin{multicols}{2}
  \item 30\sqrt{3} m
  \item 4
  \item 3
  \item 10
  \item 2
  \item 1
\end{multicols}
\end{enumerate}
\vspace{0.15cm}
\item Edificio 40 m, ángulo de depresión 45°, distancia horizontal:
\begin{enumerate}[label=\textbf{\Alph*)}]
\begin{multicols}{2}
  \item 5
  \item 40 m
  \item 3
  \item 10
  \item 2
  \item 1
\end{multicols}
\end{enumerate}
\vspace{0.15cm}
\item Edificio 30 m, ángulo de depresión 30°, distancia horizontal:
\begin{enumerate}[label=\textbf{\Alph*)}]
\begin{multicols}{2}
  \item 30\sqrt{3} m
  \item 4
  \item 3
  \item 10
  \item 2
  \item 1
\end{multicols}
\end{enumerate}
\vspace{0.15cm}
\item Edificio 60 m, ángulo de depresión 60°, distancia horizontal:
\begin{enumerate}[label=\textbf{\Alph*)}]
\begin{multicols}{2}
  \item 5
  \item 20\sqrt{3} m
  \item 3
  \item 10
  \item 2
  \item 1
\end{multicols}
\end{enumerate}
\vspace{0.15cm}
\item Rampa inclinación 30°, longitud 10 m, altura que alcanza:
\begin{enumerate}[label=\textbf{\Alph*)}]
\begin{multicols}{2}
  \item 5 m
  \item 4
  \item 3
  \item 10
  \item 2
  \item 1
\end{multicols}
\end{enumerate}
\vspace{0.15cm}
\item Rampa inclinación 45°, longitud 10 m, altura que alcanza:
\begin{enumerate}[label=\textbf{\Alph*)}]
\begin{multicols}{2}
  \item 5
  \item 5\sqrt{2} m
  \item 3
  \item 10
  \item 2
  \item 1
\end{multicols}
\end{enumerate}
\vspace{0.15cm}
\item Rampa inclinación 60°, longitud 10 m, altura que alcanza:
\begin{enumerate}[label=\textbf{\Alph*)}]
\begin{multicols}{2}
  \item 5
  \item 5\sqrt{3} m
  \item 3
  \item 10
  \item 2
  \item 1
\end{multicols}
\end{enumerate}
\vspace{0.15cm}
\item Cable 20 m sostiene poste con ángulo 30°, altura poste:
\begin{enumerate}[label=\textbf{\Alph*)}]
\begin{multicols}{2}
  \item 10 m
  \item 4
  \item 3
  \item 10
  \item 2
  \item 1
\end{multicols}
\end{enumerate}
\vspace{0.15cm}
\item Cable 20 m sostiene poste con ángulo 60°, altura poste:
\begin{enumerate}[label=\textbf{\Alph*)}]
\begin{multicols}{2}
  \item 5
  \item 10\sqrt{3} m
  \item 3
  \item 10
  \item 2
  \item 1
\end{multicols}
\end{enumerate}
\vspace{0.15cm}
\item Cable 20 m sostiene poste con ángulo 45°, altura poste:
\begin{enumerate}[label=\textbf{\Alph*)}]
\begin{multicols}{2}
  \item 5
  \item 10\sqrt{2} m
  \item 3
  \item 10
  \item 2
  \item 1
\end{multicols}
\end{enumerate}
\vspace{0.15cm}
\item Si \sen\theta=7/25, entonces \cos\theta =
\begin{enumerate}[label=\textbf{\Alph*)}]
\begin{multicols}{2}
  \item \frac{24}{25}
  \item 4
  \item 3
  \item 10
  \item 2
  \item 1
\end{multicols}
\end{enumerate}
\vspace{0.15cm}
\item Si \cos\theta=24/25, entonces \sen\theta =
\begin{enumerate}[label=\textbf{\Alph*)}]
\begin{multicols}{2}
  \item \frac{7}{25}
  \item 4
  \item 3
  \item 10
  \item 2
  \item 1
\end{multicols}
\end{enumerate}
\vspace{0.15cm}
\item Si \tan\theta=7/24, entonces \sen\theta =
\begin{enumerate}[label=\textbf{\Alph*)}]
\begin{multicols}{2}
  \item \frac{7}{25}
  \item 4
  \item 3
  \item 10
  \item 2
  \item 1
\end{multicols}
\end{enumerate}
\vspace{0.15cm}
\item Si \tan\theta=7/24, entonces \cos\theta =
\begin{enumerate}[label=\textbf{\Alph*)}]
\begin{multicols}{2}
  \item 5
  \item \frac{24}{25}
  \item 3
  \item 10
  \item 2
  \item 1
\end{multicols}
\end{enumerate}
\vspace{0.15cm}
\item Si \sen\theta=9/41, entonces \cos\theta =
\begin{enumerate}[label=\textbf{\Alph*)}]
\begin{multicols}{2}
  \item \frac{40}{41}
  \item 4
  \item 3
  \item 10
  \item 2
  \item 1
\end{multicols}
\end{enumerate}
\vspace{0.15cm}
\item Si \cos\theta=40/41, entonces \sen\theta =
\begin{enumerate}[label=\textbf{\Alph*)}]
\begin{multicols}{2}
  \item \frac{9}{41}
  \item 4
  \item 3
  \item 10
  \item 2
  \item 1
\end{multicols}
\end{enumerate}
\vspace{0.15cm}
\item Si \tan\theta=9/40, entonces \sen\theta =
\begin{enumerate}[label=\textbf{\Alph*)}]
\begin{multicols}{2}
  \item \frac{9}{41}
  \item 4
  \item 3
  \item 10
  \item 2
  \item 1
\end{multicols}
\end{enumerate}
\vspace{0.15cm}
\item Si \tan\theta=9/40, entonces \cos\theta =
\begin{enumerate}[label=\textbf{\Alph*)}]
\begin{multicols}{2}
  \item 5
  \item \frac{40}{41}
  \item 3
  \item 10
  \item 2
  \item 1
\end{multicols}
\end{enumerate}
\vspace{0.15cm}
\item Si \sen\theta=20/29, entonces \cos\theta =
\begin{enumerate}[label=\textbf{\Alph*)}]
\begin{multicols}{2}
  \item \frac{21}{29}
  \item 4
  \item 3
  \item 10
  \item 2
  \item 1
\end{multicols}
\end{enumerate}
\vspace{0.15cm}
\item Si \cos\theta=21/29, entonces \sen\theta =
\begin{enumerate}[label=\textbf{\Alph*)}]
\begin{multicols}{2}
  \item \frac{20}{29}
  \item 4
  \item 3
  \item 10
  \item 2
  \item 1
\end{multicols}
\end{enumerate}
\vspace{0.15cm}
\end{enumerate}\n
\subsubsection*{SECCIÓN V: LEY DE SENOS Y COSENOS (211–250)}

\begin{enumerate}[label=\textbf{\arabic*.}]
\setcounter{enumi}{210}
\item La ley de senos establece que:
\begin{enumerate}[label=\textbf{\Alph*)}]
\begin{multicols}{2}
  \item \frac{a}{\sen A}=\frac{b}{\sen B}=\frac{c}{\sen C}
  \item 5
  \item 12
  \item 7
  \item 4
  \item 1
\end{multicols}
\end{enumerate}
\vspace{0.15cm}
\item La ley de cosenos establece que:
\begin{enumerate}[label=\textbf{\Alph*)}]
\begin{multicols}{2}
  \item a^2=b^2+c^2-2bc\cos A
  \item 5
  \item 12
  \item 7
  \item 4
  \item 1
\end{multicols}
\end{enumerate}
\vspace{0.15cm}
\item Triángulo a=5, A=30°, B=45°, b es aprox:
\begin{enumerate}[label=\textbf{\Alph*)}]
\begin{multicols}{2}
  \item 7.07
  \item 5
  \item 12
  \item 7
  \item 4
  \item 1
\end{multicols}
\end{enumerate}
\vspace{0.15cm}
\item Triángulo a=8, A=45°, B=60°, b es aprox:
\begin{enumerate}[label=\textbf{\Alph*)}]
\begin{multicols}{2}
  \item 9.80
  \item 5
  \item 12
  \item 7
  \item 4
  \item 1
\end{multicols}
\end{enumerate}
\vspace{0.15cm}
\item Triángulo a=10, A=30°, B=60°, b es aprox:
\begin{enumerate}[label=\textbf{\Alph*)}]
\begin{multicols}{2}
  \item 17.32
  \item 5
  \item 12
  \item 7
  \item 4
  \item 1
\end{multicols}
\end{enumerate}
\vspace{0.15cm}
\item Triángulo a=6, A=60°, B=45°, b es aprox:
\begin{enumerate}[label=\textbf{\Alph*)}]
\begin{multicols}{2}
  \item 4.90
  \item 5
  \item 12
  \item 7
  \item 4
  \item 1
\end{multicols}
\end{enumerate}
\vspace{0.15cm}
\item Triángulo b=7, B=45°, C=60°, c es aprox:
\begin{enumerate}[label=\textbf{\Alph*)}]
\begin{multicols}{2}
  \item 8.57
  \item 5
  \item 12
  \item 7
  \item 4
  \item 1
\end{multicols}
\end{enumerate}
\vspace{0.15cm}
\item Triángulo a=5, b=7, C=60°, c es aprox:
\begin{enumerate}[label=\textbf{\Alph*)}]
\begin{multicols}{2}
  \item 6.24
  \item 5
  \item 12
  \item 7
  \item 4
  \item 1
\end{multicols}
\end{enumerate}
\vspace{0.15cm}
\item Triángulo a=8, b=6, C=90°, c es:
\begin{enumerate}[label=\textbf{\Alph*)}]
\begin{multicols}{2}
  \item 10
  \item 5
  \item 12
  \item 7
  \item 4
  \item 1
\end{multicols}
\end{enumerate}
\vspace{0.15cm}
\item Triángulo a=5, b=5, C=60°, c es:
\begin{enumerate}[label=\textbf{\Alph*)}]
\begin{multicols}{2}
  \item 5
  \item 10
  \item 12
  \item 7
  \item 4
  \item 1
\end{multicols}
\end{enumerate}
\vspace{0.15cm}
\item Triángulo a=4, b=4, C=90°, c es:
\begin{enumerate}[label=\textbf{\Alph*)}]
\begin{multicols}{2}
  \item 4\sqrt{2}
  \item 5
  \item 12
  \item 7
  \item 4
  \item 1
\end{multicols}
\end{enumerate}
\vspace{0.15cm}
\item Triángulo a=3, b=4, C=90°, c es:
\begin{enumerate}[label=\textbf{\Alph*)}]
\begin{multicols}{2}
  \item 5
  \item 10
  \item 12
  \item 7
  \item 4
  \item 1
\end{multicols}
\end{enumerate}
\vspace{0.15cm}
\item Triángulo a=6, b=8, C=90°, c es:
\begin{enumerate}[label=\textbf{\Alph*)}]
\begin{multicols}{2}
  \item 10
  \item 5
  \item 12
  \item 7
  \item 4
  \item 1
\end{multicols}
\end{enumerate}
\vspace{0.15cm}
\item Triángulo a=5, b=12, C=90°, c es:
\begin{enumerate}[label=\textbf{\Alph*)}]
\begin{multicols}{2}
  \item 13
  \item 5
  \item 12
  \item 7
  \item 4
  \item 1
\end{multicols}
\end{enumerate}
\vspace{0.15cm}
\item Triángulo a=7, b=24, C=90°, c es:
\begin{enumerate}[label=\textbf{\Alph*)}]
\begin{multicols}{2}
  \item 25
  \item 5
  \item 12
  \item 7
  \item 4
  \item 1
\end{multicols}
\end{enumerate}
\vspace{0.15cm}
\item Triángulo a=9, b=12, C=90°, c es:
\begin{enumerate}[label=\textbf{\Alph*)}]
\begin{multicols}{2}
  \item 15
  \item 5
  \item 12
  \item 7
  \item 4
  \item 1
\end{multicols}
\end{enumerate}
\vspace{0.15cm}
\item Triángulo a=8, b=15, C=90°, c es:
\begin{enumerate}[label=\textbf{\Alph*)}]
\begin{multicols}{2}
  \item 17
  \item 5
  \item 12
  \item 7
  \item 4
  \item 1
\end{multicols}
\end{enumerate}
\vspace{0.15cm}
\item Triángulo a=10, b=24, C=90°, c es:
\begin{enumerate}[label=\textbf{\Alph*)}]
\begin{multicols}{2}
  \item 26
  \item 5
  \item 12
  \item 7
  \item 4
  \item 1
\end{multicols}
\end{enumerate}
\vspace{0.15cm}
\item Triángulo a=6, b=6, C=60°, c es:
\begin{enumerate}[label=\textbf{\Alph*)}]
\begin{multicols}{2}
  \item 6
  \item 5
  \item 12
  \item 7
  \item 4
  \item 1
\end{multicols}
\end{enumerate}
\vspace{0.15cm}
\item Triángulo a=2, b=3, C=60°, c es:
\begin{enumerate}[label=\textbf{\Alph*)}]
\begin{multicols}{2}
  \item \sqrt{7}
  \item 5
  \item 12
  \item 7
  \item 4
  \item 1
\end{multicols}
\end{enumerate}
\vspace{0.15cm}
\item Triángulo a=1, b=1, C=60°, c es:
\begin{enumerate}[label=\textbf{\Alph*)}]
\begin{multicols}{2}
  \item 1
  \item 5
  \item 12
  \item 7
  \item 4
  \item 1
\end{multicols}
\end{enumerate}
\vspace{0.15cm}
\item Triángulo a=1, b=1, C=90°, c es:
\begin{enumerate}[label=\textbf{\Alph*)}]
\begin{multicols}{2}
  \item \sqrt{2}
  \item 5
  \item 12
  \item 7
  \item 4
  \item 1
\end{multicols}
\end{enumerate}
\vspace{0.15cm}
\item Triángulo a=1, b=1, C=120°, c es:
\begin{enumerate}[label=\textbf{\Alph*)}]
\begin{multicols}{2}
  \item \sqrt{3}
  \item 5
  \item 12
  \item 7
  \item 4
  \item 1
\end{multicols}
\end{enumerate}
\vspace{0.15cm}
\item Triángulo a=2, b=2, C=120°, c es:
\begin{enumerate}[label=\textbf{\Alph*)}]
\begin{multicols}{2}
  \item 2\sqrt{3}
  \item 5
  \item 12
  \item 7
  \item 4
  \item 1
\end{multicols}
\end{enumerate}
\vspace{0.15cm}
\item Triángulo a=3, b=5, C=120°, c es:
\begin{enumerate}[label=\textbf{\Alph*)}]
\begin{multicols}{2}
  \item 7
  \item 5
  \item 12
  \item 7
  \item 4
  \item 1
\end{multicols}
\end{enumerate}
\vspace{0.15cm}
\item Triángulo a=4, b=5, C=60°, c es:
\begin{enumerate}[label=\textbf{\Alph*)}]
\begin{multicols}{2}
  \item \sqrt{21}
  \item 5
  \item 12
  \item 7
  \item 4
  \item 1
\end{multicols}
\end{enumerate}
\vspace{0.15cm}
\item Triángulo a=5, b=8, C=60°, c es:
\begin{enumerate}[label=\textbf{\Alph*)}]
\begin{multicols}{2}
  \item 7
  \item 5
  \item 12
  \item 7
  \item 4
  \item 1
\end{multicols}
\end{enumerate}
\vspace{0.15cm}
\item Triángulo a=6, b=10, C=60°, c es:
\begin{enumerate}[label=\textbf{\Alph*)}]
\begin{multicols}{2}
  \item 2\sqrt{19}
  \item 5
  \item 12
  \item 7
  \item 4
  \item 1
\end{multicols}
\end{enumerate}
\vspace{0.15cm}
\item Triángulo a=7, b=8, C=60°, c es:
\begin{enumerate}[label=\textbf{\Alph*)}]
\begin{multicols}{2}
  \item \sqrt{57}
  \item 5
  \item 12
  \item 7
  \item 4
  \item 1
\end{multicols}
\end{enumerate}
\vspace{0.15cm}
\item Triángulo a=3, b=3, C=90°, c es:
\begin{enumerate}[label=\textbf{\Alph*)}]
\begin{multicols}{2}
  \item 3\sqrt{2}
  \item 5
  \item 12
  \item 7
  \item 4
  \item 1
\end{multicols}
\end{enumerate}
\vspace{0.15cm}
\item Triángulo a=8, b=8, C=90°, c es:
\begin{enumerate}[label=\textbf{\Alph*)}]
\begin{multicols}{2}
  \item 8\sqrt{2}
  \item 5
  \item 12
  \item 7
  \item 4
  \item 1
\end{multicols}
\end{enumerate}
\vspace{0.15cm}
\item Triángulo a=10, b=10, C=90°, c es:
\begin{enumerate}[label=\textbf{\Alph*)}]
\begin{multicols}{2}
  \item 10\sqrt{2}
  \item 5
  \item 12
  \item 7
  \item 4
  \item 1
\end{multicols}
\end{enumerate}
\vspace{0.15cm}
\item Triángulo a=5, b=6, C=30°, c es aprox:
\begin{enumerate}[label=\textbf{\Alph*)}]
\begin{multicols}{2}
  \item 3.03
  \item 5
  \item 12
  \item 7
  \item 4
  \item 1
\end{multicols}
\end{enumerate}
\vspace{0.15cm}
\item Triángulo a=4, b=7, C=45°, c es aprox:
\begin{enumerate}[label=\textbf{\Alph*)}]
\begin{multicols}{2}
  \item 4.95
  \item 5
  \item 12
  \item 7
  \item 4
  \item 1
\end{multicols}
\end{enumerate}
\vspace{0.15cm}
\item Triángulo a=6, b=9, C=120°, c es aprox:
\begin{enumerate}[label=\textbf{\Alph*)}]
\begin{multicols}{2}
  \item 13.08
  \item 5
  \item 12
  \item 7
  \item 4
  \item 1
\end{multicols}
\end{enumerate}
\vspace{0.15cm}
\item Triángulo a=10, b=10, C=120°, c es:
\begin{enumerate}[label=\textbf{\Alph*)}]
\begin{multicols}{2}
  \item 10\sqrt{3}
  \item 5
  \item 12
  \item 7
  \item 4
  \item 1
\end{multicols}
\end{enumerate}
\vspace{0.15cm}
\item Triángulo a=2, b=4, C=60°, c es:
\begin{enumerate}[label=\textbf{\Alph*)}]
\begin{multicols}{2}
  \item 2\sqrt{3}
  \item 5
  \item 12
  \item 7
  \item 4
  \item 1
\end{multicols}
\end{enumerate}
\vspace{0.15cm}
\item Triángulo a=3, b=6, C=60°, c es:
\begin{enumerate}[label=\textbf{\Alph*)}]
\begin{multicols}{2}
  \item 3\sqrt{3}
  \item 5
  \item 12
  \item 7
  \item 4
  \item 1
\end{multicols}
\end{enumerate}
\vspace{0.15cm}
\item Triángulo a=4, b=8, C=60°, c es:
\begin{enumerate}[label=\textbf{\Alph*)}]
\begin{multicols}{2}
  \item 4\sqrt{3}
  \item 5
  \item 12
  \item 7
  \item 4
  \item 1
\end{multicols}
\end{enumerate}
\vspace{0.15cm}
\item Triángulo a=5, b=10, C=60°, c es:
\begin{enumerate}[label=\textbf{\Alph*)}]
\begin{multicols}{2}
  \item 5\sqrt{3}
  \item 5
  \item 12
  \item 7
  \item 4
  \item 1
\end{multicols}
\end{enumerate}
\vspace{0.15cm}
\end{enumerate}\n
\subsubsection*{SECCIÓN VI: ÁREAS Y PERÍMETROS (251–300)}

\begin{enumerate}[label=\textbf{\arabic*.}]
\setcounter{enumi}{250}
\item Área de triángulo base 6 y altura 4:
\begin{enumerate}[label=\textbf{\Alph*)}]
\begin{multicols}{2}
  \item 12
  \item 24
  \item 36
  \item 12
  \item 10
  \item 1
\end{multicols}
\end{enumerate}
\vspace{0.15cm}
\item Área de triángulo base 10 y altura 5:
\begin{enumerate}[label=\textbf{\Alph*)}]
\begin{multicols}{2}
  \item 25
  \item 24
  \item 36
  \item 12
  \item 10
  \item 1
\end{multicols}
\end{enumerate}
\vspace{0.15cm}
\item Área de triángulo base 8 y altura 3:
\begin{enumerate}[label=\textbf{\Alph*)}]
\begin{multicols}{2}
  \item 12
  \item 24
  \item 36
  \item 12
  \item 10
  \item 1
\end{multicols}
\end{enumerate}
\vspace{0.15cm}
\item Área de triángulo base 12 y altura 7:
\begin{enumerate}[label=\textbf{\Alph*)}]
\begin{multicols}{2}
  \item 42
  \item 24
  \item 36
  \item 12
  \item 10
  \item 1
\end{multicols}
\end{enumerate}
\vspace{0.15cm}
\item Área de triángulo base 5 y altura 8:
\begin{enumerate}[label=\textbf{\Alph*)}]
\begin{multicols}{2}
  \item 20
  \item 24
  \item 36
  \item 12
  \item 10
  \item 1
\end{multicols}
\end{enumerate}
\vspace{0.15cm}
\item Área de cuadrado de lado 5:
\begin{enumerate}[label=\textbf{\Alph*)}]
\begin{multicols}{2}
  \item 25
  \item 24
  \item 36
  \item 12
  \item 10
  \item 1
\end{multicols}
\end{enumerate}
\vspace{0.15cm}
\item Área de cuadrado de lado 8:
\begin{enumerate}[label=\textbf{\Alph*)}]
\begin{multicols}{2}
  \item 64
  \item 24
  \item 36
  \item 12
  \item 10
  \item 1
\end{multicols}
\end{enumerate}
\vspace{0.15cm}
\item Perímetro de cuadrado de lado 6:
\begin{enumerate}[label=\textbf{\Alph*)}]
\begin{multicols}{2}
  \item 24
  \item 48
  \item 36
  \item 12
  \item 10
  \item 1
\end{multicols}
\end{enumerate}
\vspace{0.15cm}
\item Perímetro de cuadrado de lado 9:
\begin{enumerate}[label=\textbf{\Alph*)}]
\begin{multicols}{2}
  \item 36
  \item 24
  \item 36
  \item 12
  \item 10
  \item 1
\end{multicols}
\end{enumerate}
\vspace{0.15cm}
\item Área de rectángulo base 8 y altura 5:
\begin{enumerate}[label=\textbf{\Alph*)}]
\begin{multicols}{2}
  \item 40
  \item 24
  \item 36
  \item 12
  \item 10
  \item 1
\end{multicols}
\end{enumerate}
\vspace{0.15cm}
\item Área de rectángulo base 10 y altura 6:
\begin{enumerate}[label=\textbf{\Alph*)}]
\begin{multicols}{2}
  \item 60
  \item 24
  \item 36
  \item 12
  \item 10
  \item 1
\end{multicols}
\end{enumerate}
\vspace{0.15cm}
\item Perímetro de rectángulo base 7 y altura 3:
\begin{enumerate}[label=\textbf{\Alph*)}]
\begin{multicols}{2}
  \item 20
  \item 24
  \item 36
  \item 12
  \item 10
  \item 1
\end{multicols}
\end{enumerate}
\vspace{0.15cm}
\item Perímetro de rectángulo base 12 y altura 5:
\begin{enumerate}[label=\textbf{\Alph*)}]
\begin{multicols}{2}
  \item 34
  \item 24
  \item 36
  \item 12
  \item 10
  \item 1
\end{multicols}
\end{enumerate}
\vspace{0.15cm}
\item Área de círculo de radio 3:
\begin{enumerate}[label=\textbf{\Alph*)}]
\begin{multicols}{2}
  \item 9\pi
  \item 24
  \item 36
  \item 12
  \item 10
  \item 1
\end{multicols}
\end{enumerate}
\vspace{0.15cm}
\item Área de círculo de radio 5:
\begin{enumerate}[label=\textbf{\Alph*)}]
\begin{multicols}{2}
  \item 25\pi
  \item 24
  \item 36
  \item 12
  \item 10
  \item 1
\end{multicols}
\end{enumerate}
\vspace{0.15cm}
\item Área de círculo de radio 10:
\begin{enumerate}[label=\textbf{\Alph*)}]
\begin{multicols}{2}
  \item 100\pi
  \item 24
  \item 36
  \item 12
  \item 10
  \item 1
\end{multicols}
\end{enumerate}
\vspace{0.15cm}
\item Circunferencia de círculo de radio 4:
\begin{enumerate}[label=\textbf{\Alph*)}]
\begin{multicols}{2}
  \item 8\pi
  \item 24
  \item 36
  \item 12
  \item 10
  \item 1
\end{multicols}
\end{enumerate}
\vspace{0.15cm}
\item Circunferencia de círculo de radio 6:
\begin{enumerate}[label=\textbf{\Alph*)}]
\begin{multicols}{2}
  \item 12\pi
  \item 24
  \item 36
  \item 12
  \item 10
  \item 1
\end{multicols}
\end{enumerate}
\vspace{0.15cm}
\item Circunferencia de círculo de radio 10:
\begin{enumerate}[label=\textbf{\Alph*)}]
\begin{multicols}{2}
  \item 20\pi
  \item 24
  \item 36
  \item 12
  \item 10
  \item 1
\end{multicols}
\end{enumerate}
\vspace{0.15cm}
\item Área triángulo rectángulo catetos 3 y 4:
\begin{enumerate}[label=\textbf{\Alph*)}]
\begin{multicols}{2}
  \item 6
  \item 24
  \item 36
  \item 12
  \item 10
  \item 1
\end{multicols}
\end{enumerate}
\vspace{0.15cm}
\item Área triángulo rectángulo catetos 6 y 8:
\begin{enumerate}[label=\textbf{\Alph*)}]
\begin{multicols}{2}
  \item 24
  \item 48
  \item 36
  \item 12
  \item 10
  \item 1
\end{multicols}
\end{enumerate}
\vspace{0.15cm}
\item Área triángulo rectángulo catetos 5 y 12:
\begin{enumerate}[label=\textbf{\Alph*)}]
\begin{multicols}{2}
  \item 30
  \item 24
  \item 36
  \item 12
  \item 10
  \item 1
\end{multicols}
\end{enumerate}
\vspace{0.15cm}
\item Área triángulo equilátero lado 4:
\begin{enumerate}[label=\textbf{\Alph*)}]
\begin{multicols}{2}
  \item 4\sqrt{3}
  \item 24
  \item 36
  \item 12
  \item 10
  \item 1
\end{multicols}
\end{enumerate}
\vspace{0.15cm}
\item Área triángulo equilátero lado 6:
\begin{enumerate}[label=\textbf{\Alph*)}]
\begin{multicols}{2}
  \item 9\sqrt{3}
  \item 24
  \item 36
  \item 12
  \item 10
  \item 1
\end{multicols}
\end{enumerate}
\vspace{0.15cm}
\item Área triángulo equilátero lado 2:
\begin{enumerate}[label=\textbf{\Alph*)}]
\begin{multicols}{2}
  \item \sqrt{3}
  \item 24
  \item 36
  \item 12
  \item 10
  \item 1
\end{multicols}
\end{enumerate}
\vspace{0.15cm}
\item Área trapecio bases 6 y 4, altura 3:
\begin{enumerate}[label=\textbf{\Alph*)}]
\begin{multicols}{2}
  \item 15
  \item 24
  \item 36
  \item 12
  \item 10
  \item 1
\end{multicols}
\end{enumerate}
\vspace{0.15cm}
\item Área trapecio bases 10 y 6, altura 4:
\begin{enumerate}[label=\textbf{\Alph*)}]
\begin{multicols}{2}
  \item 32
  \item 24
  \item 36
  \item 12
  \item 10
  \item 1
\end{multicols}
\end{enumerate}
\vspace{0.15cm}
\item Área rombo diagonales 6 y 8:
\begin{enumerate}[label=\textbf{\Alph*)}]
\begin{multicols}{2}
  \item 24
  \item 48
  \item 36
  \item 12
  \item 10
  \item 1
\end{multicols}
\end{enumerate}
\vspace{0.15cm}
\item Área rombo diagonales 10 y 12:
\begin{enumerate}[label=\textbf{\Alph*)}]
\begin{multicols}{2}
  \item 60
  \item 24
  \item 36
  \item 12
  \item 10
  \item 1
\end{multicols}
\end{enumerate}
\vspace{0.15cm}
\item Área paralelogramo base 8 y altura 5:
\begin{enumerate}[label=\textbf{\Alph*)}]
\begin{multicols}{2}
  \item 40
  \item 24
  \item 36
  \item 12
  \item 10
  \item 1
\end{multicols}
\end{enumerate}
\vspace{0.15cm}
\item Área paralelogramo base 12 y altura 7:
\begin{enumerate}[label=\textbf{\Alph*)}]
\begin{multicols}{2}
  \item 84
  \item 24
  \item 36
  \item 12
  \item 10
  \item 1
\end{multicols}
\end{enumerate}
\vspace{0.15cm}
\item Perímetro triángulo lados 3, 4, 5:
\begin{enumerate}[label=\textbf{\Alph*)}]
\begin{multicols}{2}
  \item 12
  \item 24
  \item 36
  \item 12
  \item 10
  \item 1
\end{multicols}
\end{enumerate}
\vspace{0.15cm}
\item Perímetro triángulo lados 6, 8, 10:
\begin{enumerate}[label=\textbf{\Alph*)}]
\begin{multicols}{2}
  \item 24
  \item 48
  \item 36
  \item 12
  \item 10
  \item 1
\end{multicols}
\end{enumerate}
\vspace{0.15cm}
\item Perímetro triángulo equilátero lado 7:
\begin{enumerate}[label=\textbf{\Alph*)}]
\begin{multicols}{2}
  \item 21
  \item 24
  \item 36
  \item 12
  \item 10
  \item 1
\end{multicols}
\end{enumerate}
\vspace{0.15cm}
\item Perímetro triángulo equilátero lado 12:
\begin{enumerate}[label=\textbf{\Alph*)}]
\begin{multicols}{2}
  \item 36
  \item 24
  \item 36
  \item 12
  \item 10
  \item 1
\end{multicols}
\end{enumerate}
\vspace{0.15cm}
\item Área hexágono regular lado 2:
\begin{enumerate}[label=\textbf{\Alph*)}]
\begin{multicols}{2}
  \item 6\sqrt{3}
  \item 24
  \item 36
  \item 12
  \item 10
  \item 1
\end{multicols}
\end{enumerate}
\vspace{0.15cm}
\item Área hexágono regular lado 3:
\begin{enumerate}[label=\textbf{\Alph*)}]
\begin{multicols}{2}
  \item \frac{27\sqrt{3}}{2}
  \item 24
  \item 36
  \item 12
  \item 10
  \item 1
\end{multicols}
\end{enumerate}
\vspace{0.15cm}
\item Área sector circular radio 6, ángulo 60°:
\begin{enumerate}[label=\textbf{\Alph*)}]
\begin{multicols}{2}
  \item 6\pi
  \item 24
  \item 36
  \item 12
  \item 10
  \item 1
\end{multicols}
\end{enumerate}
\vspace{0.15cm}
\item Área sector circular radio 4, ángulo 90°:
\begin{enumerate}[label=\textbf{\Alph*)}]
\begin{multicols}{2}
  \item 4\pi
  \item 24
  \item 36
  \item 12
  \item 10
  \item 1
\end{multicols}
\end{enumerate}
\vspace{0.15cm}
\item Área sector circular radio 10, ángulo 45°:
\begin{enumerate}[label=\textbf{\Alph*)}]
\begin{multicols}{2}
  \item \frac{25\pi}{2}
  \item 24
  \item 36
  \item 12
  \item 10
  \item 1
\end{multicols}
\end{enumerate}
\vspace{0.15cm}
\item Área triángulo lados 5, 6, 7 (Herón):
\begin{enumerate}[label=\textbf{\Alph*)}]
\begin{multicols}{2}
  \item 14.70
  \item 24
  \item 36
  \item 12
  \item 10
  \item 1
\end{multicols}
\end{enumerate}
\vspace{0.15cm}
\item Área triángulo lados 3, 4, 5 (Herón):
\begin{enumerate}[label=\textbf{\Alph*)}]
\begin{multicols}{2}
  \item 6
  \item 24
  \item 36
  \item 12
  \item 10
  \item 1
\end{multicols}
\end{enumerate}
\vspace{0.15cm}
\item Área triángulo lados 6, 8, 10 (Herón):
\begin{enumerate}[label=\textbf{\Alph*)}]
\begin{multicols}{2}
  \item 24
  \item 48
  \item 36
  \item 12
  \item 10
  \item 1
\end{multicols}
\end{enumerate}
\vspace{0.15cm}
\item Área triángulo lados 9, 12, 15 (Herón):
\begin{enumerate}[label=\textbf{\Alph*)}]
\begin{multicols}{2}
  \item 54
  \item 24
  \item 36
  \item 12
  \item 10
  \item 1
\end{multicols}
\end{enumerate}
\vspace{0.15cm}
\item Área triángulo lados 7, 24, 25 (Herón):
\begin{enumerate}[label=\textbf{\Alph*)}]
\begin{multicols}{2}
  \item 84
  \item 24
  \item 36
  \item 12
  \item 10
  \item 1
\end{multicols}
\end{enumerate}
\vspace{0.15cm}
\item Volumen cubo arista 3:
\begin{enumerate}[label=\textbf{\Alph*)}]
\begin{multicols}{2}
  \item 27
  \item 24
  \item 36
  \item 12
  \item 10
  \item 1
\end{multicols}
\end{enumerate}
\vspace{0.15cm}
\item Volumen cubo arista 5:
\begin{enumerate}[label=\textbf{\Alph*)}]
\begin{multicols}{2}
  \item 125
  \item 24
  \item 36
  \item 12
  \item 10
  \item 1
\end{multicols}
\end{enumerate}
\vspace{0.15cm}
\item Volumen cilindro radio 3, altura 5:
\begin{enumerate}[label=\textbf{\Alph*)}]
\begin{multicols}{2}
  \item 45\pi
  \item 24
  \item 36
  \item 12
  \item 10
  \item 1
\end{multicols}
\end{enumerate}
\vspace{0.15cm}
\item Volumen esfera radio 3:
\begin{enumerate}[label=\textbf{\Alph*)}]
\begin{multicols}{2}
  \item 36\pi
  \item 24
  \item 36
  \item 12
  \item 10
  \item 1
\end{multicols}
\end{enumerate}
\vspace{0.15cm}
\item Volumen cono radio 3, altura 4:
\begin{enumerate}[label=\textbf{\Alph*)}]
\begin{multicols}{2}
  \item 12\pi
  \item 24
  \item 36
  \item 12
  \item 10
  \item 1
\end{multicols}
\end{enumerate}
\vspace{0.15cm}
\end{enumerate}\n
\vspace{0.4cm}
\begin{tcolorbox}[colback=gray!5!white,colframe=gray!60!black,title=\textbf{SOLUCIONARIO -- 1.6 Geometría y Trigonometría}]
\begin{multicols}{4}
\begin{enumerate}[label=\textbf{\arabic*.}]
  \item \textbf{B}
  \item \textbf{C}
  \item \textbf{D}
  \item \textbf{A}
  \item \textbf{C}
  \item \textbf{B}
  \item \textbf{C}
  \item \textbf{B}
  \item \textbf{B}
  \item \textbf{B}
  \item \textbf{C}
  \item \textbf{C}
  \item \textbf{B}
  \item \textbf{A}
  \item \textbf{C}
  \item \textbf{B}
  \item \textbf{D}
  \item \textbf{A}
  \item \textbf{B}
  \item \textbf{B}
  \item \textbf{C}
  \item \textbf{C}
  \item \textbf{B}
  \item \textbf{B}
  \item \textbf{C}
  \item \textbf{D}
  \item \textbf{D}
  \item \textbf{E}
  \item \textbf{B}
  \item \textbf{B}
  \item \textbf{C}
  \item \textbf{C}
  \item \textbf{A}
  \item \textbf{B}
  \item \textbf{A}
  \item \textbf{A}
  \item \textbf{A}
  \item \textbf{C}
  \item \textbf{C}
  \item \textbf{D}
  \item \textbf{B}
  \item \textbf{B}
  \item \textbf{A}
  \item \textbf{A}
  \item \textbf{A}
  \item \textbf{A}
  \item \textbf{A}
  \item \textbf{A}
  \item \textbf{B}
  \item \textbf{B}
  \item \textbf{B}
  \item \textbf{B}
  \item \textbf{B}
  \item \textbf{B}
  \item \textbf{B}
  \item \textbf{B}
  \item \textbf{B}
  \item \textbf{B}
  \item \textbf{B}
  \item \textbf{B}
  \item \textbf{B}
  \item \textbf{B}
  \item \textbf{B}
  \item \textbf{B}
  \item \textbf{A}
  \item \textbf{C}
  \item \textbf{A}
  \item \textbf{B}
  \item \textbf{B}
  \item \textbf{A}
  \item \textbf{A}
  \item \textbf{A}
  \item \textbf{B}
  \item \textbf{B}
  \item \textbf{B}
  \item \textbf{B}
  \item \textbf{B}
  \item \textbf{B}
  \item \textbf{B}
  \item \textbf{B}
  \item \textbf{A}
  \item \textbf{A}
  \item \textbf{A}
  \item \textbf{A}
  \item \textbf{A}
  \item \textbf{A}
  \item \textbf{A}
  \item \textbf{A}
  \item \textbf{A}
  \item \textbf{A}
  \item \textbf{A}
  \item \textbf{B}
  \item \textbf{C}
  \item \textbf{A}
  \item \textbf{B}
  \item \textbf{C}
  \item \textbf{A}
  \item \textbf{C}
  \item \textbf{D}
  \item \textbf{B}
  \item \textbf{B}
  \item \textbf{D}
  \item \textbf{C}
  \item \textbf{A}
  \item \textbf{D}
  \item \textbf{A}
  \item \textbf{B}
  \item \textbf{A}
  \item \textbf{B}
  \item \textbf{A}
  \item \textbf{F}
  \item \textbf{A}
  \item \textbf{C}
  \item \textbf{A}
  \item \textbf{A}
  \item \textbf{A}
  \item \textbf{A}
  \item \textbf{A}
  \item \textbf{A}
  \item \textbf{A}
  \item \textbf{A}
  \item \textbf{B}
  \item \textbf{A}
  \item \textbf{B}
  \item \textbf{A}
  \item \textbf{B}
  \item \textbf{B}
  \item \textbf{A}
  \item \textbf{B}
  \item \textbf{B}
  \item \textbf{A}
  \item \textbf{A}
  \item \textbf{B}
  \item \textbf{B}
  \item \textbf{A}
  \item \textbf{D}
  \item \textbf{B}
  \item \textbf{C}
  \item \textbf{F}
  \item \textbf{A}
  \item \textbf{D}
  \item \textbf{B}
  \item \textbf{B}
  \item \textbf{D}
  \item \textbf{B}
  \item \textbf{C}
  \item \textbf{C}
  \item \textbf{B}
  \item \textbf{A}
  \item \textbf{D}
  \item \textbf{B}
  \item \textbf{D}
  \item \textbf{D}
  \item \textbf{B}
  \item \textbf{D}
  \item \textbf{A}
  \item \textbf{B}
  \item \textbf{C}
  \item \textbf{B}
  \item \textbf{C}
  \item \textbf{B}
  \item \textbf{A}
  \item \textbf{A}
  \item \textbf{A}
  \item \textbf{B}
  \item \textbf{B}
  \item \textbf{A}
  \item \textbf{B}
  \item \textbf{A}
  \item \textbf{B}
  \item \textbf{B}
  \item \textbf{A}
  \item \textbf{A}
  \item \textbf{B}
  \item \textbf{B}
  \item \textbf{D}
  \item \textbf{C}
  \item \textbf{D}
  \item \textbf{B}
  \item \textbf{C}
  \item \textbf{C}
  \item \textbf{D}
  \item \textbf{B}
  \item \textbf{A}
  \item \textbf{E}
  \item \textbf{E}
  \item \textbf{A}
  \item \textbf{F}
  \item \textbf{B}
  \item \textbf{B}
  \item \textbf{A}
  \item \textbf{B}
  \item \textbf{A}
  \item \textbf{B}
  \item \textbf{A}
  \item \textbf{B}
  \item \textbf{B}
  \item \textbf{A}
  \item \textbf{B}
  \item \textbf{B}
  \item \textbf{A}
  \item \textbf{A}
  \item \textbf{A}
  \item \textbf{B}
  \item \textbf{A}
  \item \textbf{A}
  \item \textbf{A}
  \item \textbf{B}
  \item \textbf{A}
  \item \textbf{A}
  \item \textbf{A}
  \item \textbf{B}
  \item \textbf{A}
  \item \textbf{A}
  \item \textbf{A}
  \item \textbf{A}
  \item \textbf{A}
  \item \textbf{A}
  \item \textbf{A}
  \item \textbf{A}
  \item \textbf{A}
  \item \textbf{A}
  \item \textbf{A}
  \item \textbf{A}
  \item \textbf{A}
  \item \textbf{A}
  \item \textbf{A}
  \item \textbf{A}
  \item \textbf{A}
  \item \textbf{A}
  \item \textbf{A}
  \item \textbf{A}
  \item \textbf{A}
  \item \textbf{A}
  \item \textbf{A}
  \item \textbf{A}
  \item \textbf{A}
  \item \textbf{A}
  \item \textbf{A}
  \item \textbf{A}
  \item \textbf{A}
  \item \textbf{A}
  \item \textbf{A}
  \item \textbf{A}
  \item \textbf{A}
  \item \textbf{A}
  \item \textbf{A}
  \item \textbf{A}
  \item \textbf{A}
  \item \textbf{A}
  \item \textbf{A}
  \item \textbf{A}
  \item \textbf{A}
  \item \textbf{A}
  \item \textbf{A}
  \item \textbf{A}
  \item \textbf{A}
  \item \textbf{A}
  \item \textbf{A}
  \item \textbf{A}
  \item \textbf{A}
  \item \textbf{A}
  \item \textbf{A}
  \item \textbf{A}
  \item \textbf{A}
  \item \textbf{A}
  \item \textbf{A}
  \item \textbf{A}
  \item \textbf{A}
  \item \textbf{A}
  \item \textbf{A}
  \item \textbf{A}
  \item \textbf{A}
  \item \textbf{A}
  \item \textbf{A}
  \item \textbf{A}
  \item \textbf{A}
  \item \textbf{A}
  \item \textbf{A}
  \item \textbf{A}
  \item \textbf{A}
  \item \textbf{A}
  \item \textbf{A}
  \item \textbf{A}
  \item \textbf{A}
  \item \textbf{A}
  \item \textbf{A}
  \item \textbf{A}
  \item \textbf{A}
  \item \textbf{A}
  \item \textbf{A}
  \item \textbf{A}
  \item \textbf{A}
  \item \textbf{A}
  \item \textbf{A}
  \item \textbf{A}
  \item \textbf{A}
  \item \textbf{A}
  \item \textbf{A}
  \item \textbf{A}
\end{enumerate}
\end{multicols}
\end{tcolorbox}
\newpage\n
\subsection{1.7 Funciones}

\subsubsection*{SECCIÓN I: CONCEPTOS BÁSICOS DE FUNCIONES (1–40)}

\begin{enumerate}[label=\textbf{\arabic*.}]
\setcounter{enumi}{0}
\item ¿Qué es una función?
\begin{enumerate}[label=\textbf{\Alph*)}]
\begin{multicols}{2}
  \item Una regla que asigna a cada elemento del dominio un único elemento del rango
  \item 8
  \item 14
  \item 7
  \item 5
  \item 10
\end{multicols}
\end{enumerate}
\vspace{0.15cm}
\item El dominio de una función es:
\begin{enumerate}[label=\textbf{\Alph*)}]
\begin{multicols}{2}
  \item Par
  \item El conjunto de valores de entrada
  \item 14
  \item 7
  \item 5
  \item 10
\end{multicols}
\end{enumerate}
\vspace{0.15cm}
\item El rango de una función es:
\begin{enumerate}[label=\textbf{\Alph*)}]
\begin{multicols}{2}
  \item Par
  \item El conjunto de valores de salida
  \item 14
  \item 7
  \item 5
  \item 10
\end{multicols}
\end{enumerate}
\vspace{0.15cm}
\item ¿Cuál de las siguientes relaciones NO es una función?
\begin{enumerate}[label=\textbf{\Alph*)}]
\begin{multicols}{2}
  \item 11
  \item 8
  \item x^2+y^2=25
  \item 7
  \item 5
  \item 10
\end{multicols}
\end{enumerate}
\vspace{0.15cm}
\item La notación \(f(x)\) se lee como:
\begin{enumerate}[label=\textbf{\Alph*)}]
\begin{multicols}{2}
  \item 11
  \item f evaluada en x
  \item 14
  \item 7
  \item 5
  \item 10
\end{multicols}
\end{enumerate}
\vspace{0.15cm}
\item Si \(f(x)=2x+3\), ¿cuál es \(f(4)\)?
\begin{enumerate}[label=\textbf{\Alph*)}]
\begin{multicols}{2}
  \item 11
  \item 8
  \item 14
  \item 7
  \item 5
  \item 10
\end{multicols}
\end{enumerate}
\vspace{0.15cm}
\item Si \(f(x)=x^2-1\), ¿cuál es \(f(3)\)?
\begin{enumerate}[label=\textbf{\Alph*)}]
\begin{multicols}{2}
  \item 8
  \item 8
  \item 14
  \item 7
  \item 5
  \item 10
\end{multicols}
\end{enumerate}
\vspace{0.15cm}
\item Si \(f(x)=3x-2\), ¿cuál es \(f(0)\)?
\begin{enumerate}[label=\textbf{\Alph*)}]
\begin{multicols}{2}
  \item -2
  \item 8
  \item 14
  \item 7
  \item 5
  \item 10
\end{multicols}
\end{enumerate}
\vspace{0.15cm}
\item Si \(f(x)=x^2+2x\), ¿cuál es \(f(-1)\)?
\begin{enumerate}[label=\textbf{\Alph*)}]
\begin{multicols}{2}
  \item -1
  \item 8
  \item 14
  \item 7
  \item 5
  \item 10
\end{multicols}
\end{enumerate}
\vspace{0.15cm}
\item Si \(f(x)=\frac{1}{x}\), ¿cuál es \(f(2)\)?
\begin{enumerate}[label=\textbf{\Alph*)}]
\begin{multicols}{2}
  \item \frac{1}{2}
  \item 8
  \item 14
  \item 7
  \item 5
  \item 10
\end{multicols}
\end{enumerate}
\vspace{0.15cm}
\item Si \(f(x)=\sqrt{x}\), ¿cuál es \(f(9)\)?
\begin{enumerate}[label=\textbf{\Alph*)}]
\begin{multicols}{2}
  \item 3
  \item 8
  \item 14
  \item 7
  \item 5
  \item 10
\end{multicols}
\end{enumerate}
\vspace{0.15cm}
\item Si \(f(x)=|x|\), ¿cuál es \(f(-5)\)?
\begin{enumerate}[label=\textbf{\Alph*)}]
\begin{multicols}{2}
  \item 11
  \item 5
  \item 14
  \item 7
  \item 5
  \item 10
\end{multicols}
\end{enumerate}
\vspace{0.15cm}
\item ¿Cuál es el dominio de \(f(x)=\frac{1}{x}\)?
\begin{enumerate}[label=\textbf{\Alph*)}]
\begin{multicols}{2}
  \item 11
  \item Todos los reales excepto \(x=0\)
  \item 14
  \item 7
  \item 5
  \item 10
\end{multicols}
\end{enumerate}
\vspace{0.15cm}
\item ¿Cuál es el dominio de \(f(x)=\sqrt{x}\)?
\begin{enumerate}[label=\textbf{\Alph*)}]
\begin{multicols}{2}
  \item 11
  \item x \geq 0
  \item 14
  \item 7
  \item 5
  \item 10
\end{multicols}
\end{enumerate}
\vspace{0.15cm}
\item ¿Cuál es el dominio de \(f(x)=\frac{1}{x-2}\)?
\begin{enumerate}[label=\textbf{\Alph*)}]
\begin{multicols}{2}
  \item 11
  \item Todos los reales excepto \(x=2\)
  \item 14
  \item 7
  \item 5
  \item 10
\end{multicols}
\end{enumerate}
\vspace{0.15cm}
\item ¿Cuál es el dominio de \(f(x)=\sqrt{x-3}\)?
\begin{enumerate}[label=\textbf{\Alph*)}]
\begin{multicols}{2}
  \item x \geq 3
  \item 8
  \item 14
  \item 7
  \item 5
  \item 10
\end{multicols}
\end{enumerate}
\vspace{0.15cm}
\item ¿Cuál es el dominio de \(f(x)=x^2+1\)?
\begin{enumerate}[label=\textbf{\Alph*)}]
\begin{multicols}{2}
  \item Todos los reales
  \item 8
  \item 14
  \item 7
  \item 5
  \item 10
\end{multicols}
\end{enumerate}
\vspace{0.15cm}
\item ¿Cuál es el rango de \(f(x)=x^2\)?
\begin{enumerate}[label=\textbf{\Alph*)}]
\begin{multicols}{2}
  \item 11
  \item y \geq 0
  \item 14
  \item 7
  \item 5
  \item 10
\end{multicols}
\end{enumerate}
\vspace{0.15cm}
\item ¿Cuál es el rango de \(f(x)=|x|\)?
\begin{enumerate}[label=\textbf{\Alph*)}]
\begin{multicols}{2}
  \item 11
  \item y \geq 0
  \item 14
  \item 7
  \item 5
  \item 10
\end{multicols}
\end{enumerate}
\vspace{0.15cm}
\item ¿Cuál es el rango de \(f(x)=x^2+2\)?
\begin{enumerate}[label=\textbf{\Alph*)}]
\begin{multicols}{2}
  \item y \geq 2
  \item 8
  \item 14
  \item 7
  \item 5
  \item 10
\end{multicols}
\end{enumerate}
\vspace{0.15cm}
\item Si \(f(x)=2x+1\), ¿cuál es \(f(a+1)\)?
\begin{enumerate}[label=\textbf{\Alph*)}]
\begin{multicols}{2}
  \item 2a+3
  \item 8
  \item 14
  \item 7
  \item 5
  \item 10
\end{multicols}
\end{enumerate}
\vspace{0.15cm}
\item Si \(f(x)=x^2\), ¿cuál es \(f(x+1)\)?
\begin{enumerate}[label=\textbf{\Alph*)}]
\begin{multicols}{2}
  \item 11
  \item x^2+2x+1
  \item 14
  \item 7
  \item 5
  \item 10
\end{multicols}
\end{enumerate}
\vspace{0.15cm}
\item Si \(f(x)=3x-2\), ¿cuál es \(f(2x)\)?
\begin{enumerate}[label=\textbf{\Alph*)}]
\begin{multicols}{2}
  \item 11
  \item 6x-2
  \item 14
  \item 7
  \item 5
  \item 10
\end{multicols}
\end{enumerate}
\vspace{0.15cm}
\item Si \(f(x)=x+1\) y \(g(x)=2x\), ¿cuál es \(f(g(3))\)?
\begin{enumerate}[label=\textbf{\Alph*)}]
\begin{multicols}{2}
  \item 7
  \item 8
  \item 14
  \item 7
  \item 5
  \item 10
\end{multicols}
\end{enumerate}
\vspace{0.15cm}
\item Si \(f(x)=2x\) y \(g(x)=x+3\), ¿cuál es \(f(g(2))\)?
\begin{enumerate}[label=\textbf{\Alph*)}]
\begin{multicols}{2}
  \item 10
  \item 8
  \item 14
  \item 7
  \item 5
  \item 10
\end{multicols}
\end{enumerate}
\vspace{0.15cm}
\item Si \(f(x)=x^2\) y \(g(x)=x+1\), ¿cuál es \(f(g(1))\)?
\begin{enumerate}[label=\textbf{\Alph*)}]
\begin{multicols}{2}
  \item 4
  \item 8
  \item 14
  \item 7
  \item 5
  \item 10
\end{multicols}
\end{enumerate}
\vspace{0.15cm}
\item Si \(f(x)=x-2\) y \(g(x)=3x\), ¿cuál es \(g(f(5))\)?
\begin{enumerate}[label=\textbf{\Alph*)}]
\begin{multicols}{2}
  \item 9
  \item 8
  \item 14
  \item 7
  \item 5
  \item 10
\end{multicols}
\end{enumerate}
\vspace{0.15cm}
\item La función identidad es:
\begin{enumerate}[label=\textbf{\Alph*)}]
\begin{multicols}{2}
  \item Par
  \item Impar
  \item f(x)=x
  \item 7
  \item 5
  \item 10
\end{multicols}
\end{enumerate}
\vspace{0.15cm}
\item La función constante es de la forma:
\begin{enumerate}[label=\textbf{\Alph*)}]
\begin{multicols}{2}
  \item 11
  \item f(x)=c
  \item 14
  \item 7
  \item 5
  \item 10
\end{multicols}
\end{enumerate}
\vspace{0.15cm}
\item La prueba de la línea vertical determina si una gráfica representa:
\begin{enumerate}[label=\textbf{\Alph*)}]
\begin{multicols}{2}
  \item Una función
  \item 8
  \item 14
  \item 7
  \item 5
  \item 10
\end{multicols}
\end{enumerate}
\vspace{0.15cm}
\item Si una gráfica interseca una línea vertical en más de un punto, entonces:
\begin{enumerate}[label=\textbf{\Alph*)}]
\begin{multicols}{2}
  \item Par
  \item No es una función
  \item 14
  \item 7
  \item 5
  \item 10
\end{multicols}
\end{enumerate}
\vspace{0.15cm}
\item La función \(f(x)=x^3\) es:
\begin{enumerate}[label=\textbf{\Alph*)}]
\begin{multicols}{2}
  \item Par
  \item Impar
  \item Ni par ni impar
  \item 7
  \item 5
  \item 10
\end{multicols}
\end{enumerate}
\vspace{0.15cm}
\item La función \(f(x)=x^2\) es:
\begin{enumerate}[label=\textbf{\Alph*)}]
\begin{multicols}{2}
  \item Par
  \item Impar
  \item 14
  \item 7
  \item 5
  \item 10
\end{multicols}
\end{enumerate}
\vspace{0.15cm}
\item Una función es par si:
\begin{enumerate}[label=\textbf{\Alph*)}]
\begin{multicols}{2}
  \item f(-x)=f(x)
  \item 8
  \item 14
  \item 7
  \item 5
  \item 10
\end{multicols}
\end{enumerate}
\vspace{0.15cm}
\item Una función es impar si:
\begin{enumerate}[label=\textbf{\Alph*)}]
\begin{multicols}{2}
  \item 11
  \item f(-x)=-f(x)
  \item 14
  \item 7
  \item 5
  \item 10
\end{multicols}
\end{enumerate}
\vspace{0.15cm}
\item La función \(f(x)=x^2+1\) es:
\begin{enumerate}[label=\textbf{\Alph*)}]
\begin{multicols}{2}
  \item Par
  \item Impar
  \item 14
  \item 7
  \item 5
  \item 10
\end{multicols}
\end{enumerate}
\vspace{0.15cm}
\item La función \(f(x)=x^3+x\) es:
\begin{enumerate}[label=\textbf{\Alph*)}]
\begin{multicols}{2}
  \item Par
  \item Impar
  \item Ni par ni impar
  \item 7
  \item 5
  \item 10
\end{multicols}
\end{enumerate}
\vspace{0.15cm}
\item La función \(f(x)=x^2+x\) es:
\begin{enumerate}[label=\textbf{\Alph*)}]
\begin{multicols}{2}
  \item Par
  \item Impar
  \item Ni par ni impar
  \item 7
  \item 5
  \item 10
\end{multicols}
\end{enumerate}
\vspace{0.15cm}
\item Si \(f(x)=2x+1\), ¿cuál es \(f^{-1}(x)\)?
\begin{enumerate}[label=\textbf{\Alph*)}]
\begin{multicols}{2}
  \item \frac{x-1}{2}
  \item 8
  \item 14
  \item 7
  \item 5
  \item 10
\end{multicols}
\end{enumerate}
\vspace{0.15cm}
\item Si \(f(x)=x-3\), ¿cuál es \(f^{-1}(x)\)?
\begin{enumerate}[label=\textbf{\Alph*)}]
\begin{multicols}{2}
  \item x+3
  \item 8
  \item 14
  \item 7
  \item 5
  \item 10
\end{multicols}
\end{enumerate}
\vspace{0.15cm}
\end{enumerate}\n
\subsubsection*{SECCIÓN II: FUNCIÓN LINEAL (41–100)}

\begin{enumerate}[label=\textbf{\arabic*.}]
\setcounter{enumi}{40}
\item La forma general de una función lineal es:
\begin{enumerate}[label=\textbf{\Alph*)}]
\begin{multicols}{2}
  \item f(x)=mx+b
  \item Decreciente
  \item Constante
  \item -2
  \item 0
  \item 1
\end{multicols}
\end{enumerate}
\vspace{0.15cm}
\item En f(x)=mx+b, m representa:
\begin{enumerate}[label=\textbf{\Alph*)}]
\begin{multicols}{2}
  \item La pendiente
  \item 3
  \item 4
  \item -2
  \item 0
  \item 1
\end{multicols}
\end{enumerate}
\vspace{0.15cm}
\item En f(x)=mx+b, b representa:
\begin{enumerate}[label=\textbf{\Alph*)}]
\begin{multicols}{2}
  \item 2
  \item El intercepto con y
  \item 4
  \item -2
  \item 0
  \item 1
\end{multicols}
\end{enumerate}
\vspace{0.15cm}
\item La pendiente de f(x)=3x+2 es:
\begin{enumerate}[label=\textbf{\Alph*)}]
\begin{multicols}{2}
  \item Creciente
  \item 3
  \item Constante
  \item -2
  \item 0
  \item 1
\end{multicols}
\end{enumerate}
\vspace{0.15cm}
\item La pendiente de f(x)=-2x+5 es:
\begin{enumerate}[label=\textbf{\Alph*)}]
\begin{multicols}{2}
  \item Creciente
  \item -2
  \item Constante
  \item -2
  \item 0
  \item 1
\end{multicols}
\end{enumerate}
\vspace{0.15cm}
\item La pendiente de f(x)=x-1 es:
\begin{enumerate}[label=\textbf{\Alph*)}]
\begin{multicols}{2}
  \item 1
  \item Decreciente
  \item Constante
  \item -2
  \item 0
  \item 1
\end{multicols}
\end{enumerate}
\vspace{0.15cm}
\item El intercepto con y de f(x)=2x+4 es:
\begin{enumerate}[label=\textbf{\Alph*)}]
\begin{multicols}{2}
  \item 4
  \item Decreciente
  \item Constante
  \item -2
  \item 0
  \item 1
\end{multicols}
\end{enumerate}
\vspace{0.15cm}
\item El intercepto con y de f(x)=-3x+1 es:
\begin{enumerate}[label=\textbf{\Alph*)}]
\begin{multicols}{2}
  \item 1
  \item Decreciente
  \item Constante
  \item -2
  \item 0
  \item 1
\end{multicols}
\end{enumerate}
\vspace{0.15cm}
\item El intercepto con y de f(x)=5x es:
\begin{enumerate}[label=\textbf{\Alph*)}]
\begin{multicols}{2}
  \item Creciente
  \item 0
  \item Constante
  \item -2
  \item 0
  \item 1
\end{multicols}
\end{enumerate}
\vspace{0.15cm}
\item La pendiente de una recta horizontal es:
\begin{enumerate}[label=\textbf{\Alph*)}]
\begin{multicols}{2}
  \item 0
  \item Decreciente
  \item Constante
  \item -2
  \item 0
  \item 1
\end{multicols}
\end{enumerate}
\vspace{0.15cm}
\item La pendiente de una recta vertical es:
\begin{enumerate}[label=\textbf{\Alph*)}]
\begin{multicols}{2}
  \item Creciente
  \item Decreciente
  \item Constante
  \item -2
  \item 0
  \item 1
\end{multicols}
\end{enumerate}
\vspace{0.15cm}
\item Pendiente de la recta que pasa por (1,2) y (3,6):
\begin{enumerate}[label=\textbf{\Alph*)}]
\begin{multicols}{2}
  \item 2
  \item 3
  \item 4
  \item -2
  \item 0
  \item 1
\end{multicols}
\end{enumerate}
\vspace{0.15cm}
\item Pendiente de la recta que pasa por (0,0) y (2,4):
\begin{enumerate}[label=\textbf{\Alph*)}]
\begin{multicols}{2}
  \item 2
  \item 3
  \item 4
  \item -2
  \item 0
  \item 1
\end{multicols}
\end{enumerate}
\vspace{0.15cm}
\item Pendiente de la recta que pasa por (1,5) y (3,1):
\begin{enumerate}[label=\textbf{\Alph*)}]
\begin{multicols}{2}
  \item -2
  \item 3
  \item 4
  \item -2
  \item 0
  \item 1
\end{multicols}
\end{enumerate}
\vspace{0.15cm}
\item Pendiente de la recta que pasa por (-1,2) y (1,6):
\begin{enumerate}[label=\textbf{\Alph*)}]
\begin{multicols}{2}
  \item 2
  \item 3
  \item 4
  \item -2
  \item 0
  \item 1
\end{multicols}
\end{enumerate}
\vspace{0.15cm}
\item Ecuación de la recta con pendiente 2 e intercepto 3:
\begin{enumerate}[label=\textbf{\Alph*)}]
\begin{multicols}{2}
  \item y=2x+3
  \item 3
  \item 4
  \item -2
  \item 0
  \item 1
\end{multicols}
\end{enumerate}
\vspace{0.15cm}
\item Ecuación de la recta con pendiente -1 e intercepto 4:
\begin{enumerate}[label=\textbf{\Alph*)}]
\begin{multicols}{2}
  \item y=-x+4
  \item 3
  \item 4
  \item -2
  \item 0
  \item 1
\end{multicols}
\end{enumerate}
\vspace{0.15cm}
\item Ecuación de la recta con pendiente 0 e intercepto 5:
\begin{enumerate}[label=\textbf{\Alph*)}]
\begin{multicols}{2}
  \item y=5
  \item 3
  \item 4
  \item -2
  \item 0
  \item 1
\end{multicols}
\end{enumerate}
\vspace{0.15cm}
\item Ecuación de la recta por (0,1) con pendiente 2:
\begin{enumerate}[label=\textbf{\Alph*)}]
\begin{multicols}{2}
  \item y=2x+1
  \item 3
  \item 4
  \item -2
  \item 0
  \item 1
\end{multicols}
\end{enumerate}
\vspace{0.15cm}
\item Ecuación de la recta por (1,3) con pendiente 2:
\begin{enumerate}[label=\textbf{\Alph*)}]
\begin{multicols}{2}
  \item y=2x+1
  \item 3
  \item 4
  \item -2
  \item 0
  \item 1
\end{multicols}
\end{enumerate}
\vspace{0.15cm}
\item Ecuación de la recta por (2,5) con pendiente 3:
\begin{enumerate}[label=\textbf{\Alph*)}]
\begin{multicols}{2}
  \item y=3x-1
  \item 3
  \item 4
  \item -2
  \item 0
  \item 1
\end{multicols}
\end{enumerate}
\vspace{0.15cm}
\item Ecuación de la recta por (1,2) y (3,8):
\begin{enumerate}[label=\textbf{\Alph*)}]
\begin{multicols}{2}
  \item y=3x-1
  \item 3
  \item 4
  \item -2
  \item 0
  \item 1
\end{multicols}
\end{enumerate}
\vspace{0.15cm}
\item Ecuación de la recta por (0,0) y (2,6):
\begin{enumerate}[label=\textbf{\Alph*)}]
\begin{multicols}{2}
  \item y=3x
  \item 3
  \item 4
  \item -2
  \item 0
  \item 1
\end{multicols}
\end{enumerate}
\vspace{0.15cm}
\item La función lineal f(x)=2x+1 es:
\begin{enumerate}[label=\textbf{\Alph*)}]
\begin{multicols}{2}
  \item Creciente
  \item Decreciente
  \item Constante
  \item -2
  \item 0
  \item 1
\end{multicols}
\end{enumerate}
\vspace{0.15cm}
\item La función lineal f(x)=-3x+2 es:
\begin{enumerate}[label=\textbf{\Alph*)}]
\begin{multicols}{2}
  \item Creciente
  \item Decreciente
  \item Constante
  \item -2
  \item 0
  \item 1
\end{multicols}
\end{enumerate}
\vspace{0.15cm}
\item La función lineal f(x)=5 es:
\begin{enumerate}[label=\textbf{\Alph*)}]
\begin{multicols}{2}
  \item Creciente
  \item Decreciente
  \item Constante
  \item -2
  \item 0
  \item 1
\end{multicols}
\end{enumerate}
\vspace{0.15cm}
\item Si m > 0, la función lineal es:
\begin{enumerate}[label=\textbf{\Alph*)}]
\begin{multicols}{2}
  \item Creciente
  \item Decreciente
  \item Constante
  \item -2
  \item 0
  \item 1
\end{multicols}
\end{enumerate}
\vspace{0.15cm}
\item Si m < 0, la función lineal es:
\begin{enumerate}[label=\textbf{\Alph*)}]
\begin{multicols}{2}
  \item Creciente
  \item Decreciente
  \item Constante
  \item -2
  \item 0
  \item 1
\end{multicols}
\end{enumerate}
\vspace{0.15cm}
\item Si m = 0, la función lineal es:
\begin{enumerate}[label=\textbf{\Alph*)}]
\begin{multicols}{2}
  \item Creciente
  \item Decreciente
  \item Constante
  \item -2
  \item 0
  \item 1
\end{multicols}
\end{enumerate}
\vspace{0.15cm}
\item El intercepto con x de f(x)=2x-4 es:
\begin{enumerate}[label=\textbf{\Alph*)}]
\begin{multicols}{2}
  \item 2
  \item Decreciente
  \item Constante
  \item -2
  \item 0
  \item 1
\end{multicols}
\end{enumerate}
\vspace{0.15cm}
\item El intercepto con x de f(x)=3x+6 es:
\begin{enumerate}[label=\textbf{\Alph*)}]
\begin{multicols}{2}
  \item -2
  \item Decreciente
  \item Constante
  \item -2
  \item 0
  \item 1
\end{multicols}
\end{enumerate}
\vspace{0.15cm}
\item El intercepto con x de f(x)=x-5 es:
\begin{enumerate}[label=\textbf{\Alph*)}]
\begin{multicols}{2}
  \item 5
  \item Decreciente
  \item Constante
  \item -2
  \item 0
  \item 1
\end{multicols}
\end{enumerate}
\vspace{0.15cm}
\item La recta y=2x+3 y la recta y=2x-1 son:
\begin{enumerate}[label=\textbf{\Alph*)}]
\begin{multicols}{2}
  \item Paralelas
  \item 3
  \item 4
  \item -2
  \item 0
  \item 1
\end{multicols}
\end{enumerate}
\vspace{0.15cm}
\item La recta y=2x+3 y la recta y=-\frac{1}{2}x+1 son:
\begin{enumerate}[label=\textbf{\Alph*)}]
\begin{multicols}{2}
  \item 2
  \item Perpendiculares
  \item 4
  \item -2
  \item 0
  \item 1
\end{multicols}
\end{enumerate}
\vspace{0.15cm}
\item Dos rectas son paralelas si:
\begin{enumerate}[label=\textbf{\Alph*)}]
\begin{multicols}{2}
  \item Tienen la misma pendiente
  \item 3
  \item 4
  \item -2
  \item 0
  \item 1
\end{multicols}
\end{enumerate}
\vspace{0.15cm}
\item Dos rectas son perpendiculares si:
\begin{enumerate}[label=\textbf{\Alph*)}]
\begin{multicols}{2}
  \item 2
  \item Sus pendientes son recíprocas negativas
  \item 4
  \item -2
  \item 0
  \item 1
\end{multicols}
\end{enumerate}
\vspace{0.15cm}
\item Si f(x)=2x+1, ¿cuál es f(3)?
\begin{enumerate}[label=\textbf{\Alph*)}]
\begin{multicols}{2}
  \item 7
  \item 3
  \item 4
  \item -2
  \item 0
  \item 1
\end{multicols}
\end{enumerate}
\vspace{0.15cm}
\item Si f(x)=-x+4, ¿cuál es f(2)?
\begin{enumerate}[label=\textbf{\Alph*)}]
\begin{multicols}{2}
  \item 2
  \item 3
  \item 4
  \item -2
  \item 0
  \item 1
\end{multicols}
\end{enumerate}
\vspace{0.15cm}
\item Si f(x)=3x-1, ¿cuál es f(-2)?
\begin{enumerate}[label=\textbf{\Alph*)}]
\begin{multicols}{2}
  \item -7
  \item 3
  \item 4
  \item -2
  \item 0
  \item 1
\end{multicols}
\end{enumerate}
\vspace{0.15cm}
\item Si f(x)=\frac{1}{2}x+3, ¿cuál es f(4)?
\begin{enumerate}[label=\textbf{\Alph*)}]
\begin{multicols}{2}
  \item 5
  \item 3
  \item 4
  \item -2
  \item 0
  \item 1
\end{multicols}
\end{enumerate}
\vspace{0.15cm}
\item Si f(x)=2x+1 y g(x)=x-3, f(x)+g(x) es:
\begin{enumerate}[label=\textbf{\Alph*)}]
\begin{multicols}{2}
  \item 3x-2
  \item Decreciente
  \item Constante
  \item -2
  \item 0
  \item 1
\end{multicols}
\end{enumerate}
\vspace{0.15cm}
\item Si f(x)=x+2 y g(x)=2x-1, f(x)*g(x) es:
\begin{enumerate}[label=\textbf{\Alph*)}]
\begin{multicols}{2}
  \item 2x^2+3x-2
  \item Decreciente
  \item Constante
  \item -2
  \item 0
  \item 1
\end{multicols}
\end{enumerate}
\vspace{0.15cm}
\item La función lineal f(x)=mx+b tiene como gráfica:
\begin{enumerate}[label=\textbf{\Alph*)}]
\begin{multicols}{2}
  \item 2
  \item Una recta
  \item 4
  \item -2
  \item 0
  \item 1
\end{multicols}
\end{enumerate}
\vspace{0.15cm}
\item Punto de intersección de y=x+1 con y=-x+3:
\begin{enumerate}[label=\textbf{\Alph*)}]
\begin{multicols}{2}
  \item (1,2)
  \item 3
  \item 4
  \item -2
  \item 0
  \item 1
\end{multicols}
\end{enumerate}
\vspace{0.15cm}
\item Punto de intersección de y=2x con y=x+3:
\begin{enumerate}[label=\textbf{\Alph*)}]
\begin{multicols}{2}
  \item (3,6)
  \item 3
  \item 4
  \item -2
  \item 0
  \item 1
\end{multicols}
\end{enumerate}
\vspace{0.15cm}
\item Punto de intersección de y=3x-1 con y=x+3:
\begin{enumerate}[label=\textbf{\Alph*)}]
\begin{multicols}{2}
  \item (2,5)
  \item 3
  \item 4
  \item -2
  \item 0
  \item 1
\end{multicols}
\end{enumerate}
\vspace{0.15cm}
\item Si f(x)=2x+1, ¿para qué x se cumple f(x)=9?
\begin{enumerate}[label=\textbf{\Alph*)}]
\begin{multicols}{2}
  \item 4
  \item 3
  \item 4
  \item -2
  \item 0
  \item 1
\end{multicols}
\end{enumerate}
\vspace{0.15cm}
\item Si f(x)=3x-2, ¿para qué x se cumple f(x)=7?
\begin{enumerate}[label=\textbf{\Alph*)}]
\begin{multicols}{2}
  \item 3
  \item 3
  \item 4
  \item -2
  \item 0
  \item 1
\end{multicols}
\end{enumerate}
\vspace{0.15cm}
\item Si f(x)=-x+5, ¿para qué x se cumple f(x)=0?
\begin{enumerate}[label=\textbf{\Alph*)}]
\begin{multicols}{2}
  \item 5
  \item 3
  \item 4
  \item -2
  \item 0
  \item 1
\end{multicols}
\end{enumerate}
\vspace{0.15cm}
\item La pendiente de 2x+3y=6 es:
\begin{enumerate}[label=\textbf{\Alph*)}]
\begin{multicols}{2}
  \item -\frac{2}{3}
  \item Decreciente
  \item Constante
  \item -2
  \item 0
  \item 1
\end{multicols}
\end{enumerate}
\vspace{0.15cm}
\item La pendiente de 4x-2y=8 es:
\begin{enumerate}[label=\textbf{\Alph*)}]
\begin{multicols}{2}
  \item 2
  \item Decreciente
  \item Constante
  \item -2
  \item 0
  \item 1
\end{multicols}
\end{enumerate}
\vspace{0.15cm}
\item El intercepto con y de 2x+3y=6 es:
\begin{enumerate}[label=\textbf{\Alph*)}]
\begin{multicols}{2}
  \item 2
  \item Decreciente
  \item Constante
  \item -2
  \item 0
  \item 1
\end{multicols}
\end{enumerate}
\vspace{0.15cm}
\item El intercepto con x de 2x+3y=6 es:
\begin{enumerate}[label=\textbf{\Alph*)}]
\begin{multicols}{2}
  \item 3
  \item Decreciente
  \item Constante
  \item -2
  \item 0
  \item 1
\end{multicols}
\end{enumerate}
\vspace{0.15cm}
\item Recta que pasa por (1,1) y es paralela a y=2x+3:
\begin{enumerate}[label=\textbf{\Alph*)}]
\begin{multicols}{2}
  \item y=2x-1
  \item 3
  \item 4
  \item -2
  \item 0
  \item 1
\end{multicols}
\end{enumerate}
\vspace{0.15cm}
\item Recta que pasa por (2,1) y es perpendicular a y=2x+1:
\begin{enumerate}[label=\textbf{\Alph*)}]
\begin{multicols}{2}
  \item y=-\frac{1}{2}x+2
  \item 3
  \item 4
  \item -2
  \item 0
  \item 1
\end{multicols}
\end{enumerate}
\vspace{0.15cm}
\item ¿Cuál de las siguientes funciones lineales es creciente?
\begin{enumerate}[label=\textbf{\Alph*)}]
\begin{multicols}{2}
  \item 2
  \item f(x)=3x-4
  \item 4
  \item -2
  \item 0
  \item 1
\end{multicols}
\end{enumerate}
\vspace{0.15cm}
\item ¿Cuál de las siguientes funciones lineales es decreciente?
\begin{enumerate}[label=\textbf{\Alph*)}]
\begin{multicols}{2}
  \item 2
  \item f(x)=-3x+2
  \item 4
  \item -2
  \item 0
  \item 1
\end{multicols}
\end{enumerate}
\vspace{0.15cm}
\item Si f(x)=2x+1, ¿cuál es f(f(1))?
\begin{enumerate}[label=\textbf{\Alph*)}]
\begin{multicols}{2}
  \item 7
  \item 3
  \item 4
  \item -2
  \item 0
  \item 1
\end{multicols}
\end{enumerate}
\vspace{0.15cm}
\item Si f(x)=x+2, ¿cuál es f(f(0))?
\begin{enumerate}[label=\textbf{\Alph*)}]
\begin{multicols}{2}
  \item 4
  \item 3
  \item 4
  \item -2
  \item 0
  \item 1
\end{multicols}
\end{enumerate}
\vspace{0.15cm}
\item Si f(x)=3x-1, ¿cuál es f(f(1))?
\begin{enumerate}[label=\textbf{\Alph*)}]
\begin{multicols}{2}
  \item 5
  \item 3
  \item 4
  \item -2
  \item 0
  \item 1
\end{multicols}
\end{enumerate}
\vspace{0.15cm}
\end{enumerate}\n
\subsubsection*{SECCIÓN III: FUNCIÓN CUADRÁTICA (101–170)}

\begin{enumerate}[label=\textbf{\arabic*.}]
\setcounter{enumi}{100}
\item La forma general de una función cuadrática es:
\begin{enumerate}[label=\textbf{\Alph*)}]
\begin{multicols}{2}
  \item f(x)=ax^2+bx+c
  \item (1,1)
  \item 0
  \item 2
  \item -1
  \item Todos los reales
\end{multicols}
\end{enumerate}
\vspace{0.15cm}
\item La gráfica de una función cuadrática es:
\begin{enumerate}[label=\textbf{\Alph*)}]
\begin{multicols}{2}
  \item Una parábola
  \item 1
  \item 0
  \item 2
  \item -1
  \item Todos los reales
\end{multicols}
\end{enumerate}
\vspace{0.15cm}
\item En f(x)=ax^2+bx+c, si a > 0, la parábola abre:
\begin{enumerate}[label=\textbf{\Alph*)}]
\begin{multicols}{2}
  \item Hacia arriba
  \item 1
  \item 0
  \item 2
  \item -1
  \item Todos los reales
\end{multicols}
\end{enumerate}
\vspace{0.15cm}
\item En f(x)=ax^2+bx+c, si a < 0, la parábola abre:
\begin{enumerate}[label=\textbf{\Alph*)}]
\begin{multicols}{2}
  \item Hacia abajo
  \item 1
  \item 0
  \item 2
  \item -1
  \item Todos los reales
\end{multicols}
\end{enumerate}
\vspace{0.15cm}
\item El vértice de una parábola es:
\begin{enumerate}[label=\textbf{\Alph*)}]
\begin{multicols}{2}
  \item El punto más alto o más bajo
  \item 1
  \item 0
  \item 2
  \item -1
  \item Todos los reales
\end{multicols}
\end{enumerate}
\vspace{0.15cm}
\item El eje de simetría de una parábola es:
\begin{enumerate}[label=\textbf{\Alph*)}]
\begin{multicols}{2}
  \item Una recta vertical que pasa por el vértice
  \item 1
  \item 0
  \item 2
  \item -1
  \item Todos los reales
\end{multicols}
\end{enumerate}
\vspace{0.15cm}
\item La coordenada x del vértice de ax^2+bx+c es:
\begin{enumerate}[label=\textbf{\Alph*)}]
\begin{multicols}{2}
  \item x=-\frac{b}{2a}
  \item 1
  \item 0
  \item 2
  \item -1
  \item Todos los reales
\end{multicols}
\end{enumerate}
\vspace{0.15cm}
\item El vértice de f(x)=x^2 es:
\begin{enumerate}[label=\textbf{\Alph*)}]
\begin{multicols}{2}
  \item (0,0)
  \item (1,1)
  \item 0
  \item 2
  \item -1
  \item Todos los reales
\end{multicols}
\end{enumerate}
\vspace{0.15cm}
\item El vértice de f(x)=x^2+1 es:
\begin{enumerate}[label=\textbf{\Alph*)}]
\begin{multicols}{2}
  \item (0,1)
  \item (1,1)
  \item 0
  \item 2
  \item -1
  \item Todos los reales
\end{multicols}
\end{enumerate}
\vspace{0.15cm}
\item El vértice de f(x)=x^2-1 es:
\begin{enumerate}[label=\textbf{\Alph*)}]
\begin{multicols}{2}
  \item (0,-1)
  \item (1,1)
  \item 0
  \item 2
  \item -1
  \item Todos los reales
\end{multicols}
\end{enumerate}
\vspace{0.15cm}
\item El vértice de f(x)=(x-1)^2 es:
\begin{enumerate}[label=\textbf{\Alph*)}]
\begin{multicols}{2}
  \item (1,0)
  \item (1,1)
  \item 0
  \item 2
  \item -1
  \item Todos los reales
\end{multicols}
\end{enumerate}
\vspace{0.15cm}
\item El vértice de f(x)=(x+2)^2 es:
\begin{enumerate}[label=\textbf{\Alph*)}]
\begin{multicols}{2}
  \item (-2,0)
  \item (1,1)
  \item 0
  \item 2
  \item -1
  \item Todos los reales
\end{multicols}
\end{enumerate}
\vspace{0.15cm}
\item El vértice de f(x)=(x-3)^2+2 es:
\begin{enumerate}[label=\textbf{\Alph*)}]
\begin{multicols}{2}
  \item (3,2)
  \item (1,1)
  \item 0
  \item 2
  \item -1
  \item Todos los reales
\end{multicols}
\end{enumerate}
\vspace{0.15cm}
\item El vértice de f(x)=x^2-4x+3 es:
\begin{enumerate}[label=\textbf{\Alph*)}]
\begin{multicols}{2}
  \item (2,-1)
  \item (1,1)
  \item 0
  \item 2
  \item -1
  \item Todos los reales
\end{multicols}
\end{enumerate}
\vspace{0.15cm}
\item El vértice de f(x)=x^2-6x+5 es:
\begin{enumerate}[label=\textbf{\Alph*)}]
\begin{multicols}{2}
  \item (3,-4)
  \item (1,1)
  \item 0
  \item 2
  \item -1
  \item Todos los reales
\end{multicols}
\end{enumerate}
\vspace{0.15cm}
\item El vértice de f(x)=x^2+2x+1 es:
\begin{enumerate}[label=\textbf{\Alph*)}]
\begin{multicols}{2}
  \item (-1,0)
  \item (1,1)
  \item 0
  \item 2
  \item -1
  \item Todos los reales
\end{multicols}
\end{enumerate}
\vspace{0.15cm}
\item El vértice de f(x)=2x^2-4x+1 es:
\begin{enumerate}[label=\textbf{\Alph*)}]
\begin{multicols}{2}
  \item (1,-1)
  \item (1,1)
  \item 0
  \item 2
  \item -1
  \item Todos los reales
\end{multicols}
\end{enumerate}
\vspace{0.15cm}
\item El vértice de f(x)=-x^2+2x+3 es:
\begin{enumerate}[label=\textbf{\Alph*)}]
\begin{multicols}{2}
  \item (1,4)
  \item (1,1)
  \item 0
  \item 2
  \item -1
  \item Todos los reales
\end{multicols}
\end{enumerate}
\vspace{0.15cm}
\item El eje de simetría de f(x)=x^2-4x+3 es:
\begin{enumerate}[label=\textbf{\Alph*)}]
\begin{multicols}{2}
  \item x = 2
  \item 1
  \item 0
  \item 2
  \item -1
  \item Todos los reales
\end{multicols}
\end{enumerate}
\vspace{0.15cm}
\item El eje de simetría de f(x)=x^2+6x+5 es:
\begin{enumerate}[label=\textbf{\Alph*)}]
\begin{multicols}{2}
  \item x = -3
  \item 1
  \item 0
  \item 2
  \item -1
  \item Todos los reales
\end{multicols}
\end{enumerate}
\vspace{0.15cm}
\item El eje de simetría de f(x)=2x^2-8x+1 es:
\begin{enumerate}[label=\textbf{\Alph*)}]
\begin{multicols}{2}
  \item x = 2
  \item 1
  \item 0
  \item 2
  \item -1
  \item Todos los reales
\end{multicols}
\end{enumerate}
\vspace{0.15cm}
\item Los interceptos con x de f(x)=x^2-5x+6 son:
\begin{enumerate}[label=\textbf{\Alph*)}]
\begin{multicols}{2}
  \item x=2, x=3
  \item 1
  \item 0
  \item 2
  \item -1
  \item Todos los reales
\end{multicols}
\end{enumerate}
\vspace{0.15cm}
\item Los interceptos con x de f(x)=x^2-4 son:
\begin{enumerate}[label=\textbf{\Alph*)}]
\begin{multicols}{2}
  \item x=2, x=-2
  \item 1
  \item 0
  \item 2
  \item -1
  \item Todos los reales
\end{multicols}
\end{enumerate}
\vspace{0.15cm}
\item Los interceptos con x de f(x)=x^2-9 son:
\begin{enumerate}[label=\textbf{\Alph*)}]
\begin{multicols}{2}
  \item x=3, x=-3
  \item 1
  \item 0
  \item 2
  \item -1
  \item Todos los reales
\end{multicols}
\end{enumerate}
\vspace{0.15cm}
\item Los interceptos con x de f(x)=x^2+5x+6 son:
\begin{enumerate}[label=\textbf{\Alph*)}]
\begin{multicols}{2}
  \item x=-2, x=-3
  \item 1
  \item 0
  \item 2
  \item -1
  \item Todos los reales
\end{multicols}
\end{enumerate}
\vspace{0.15cm}
\item El intercepto con y de f(x)=x^2-4x+3 es:
\begin{enumerate}[label=\textbf{\Alph*)}]
\begin{multicols}{2}
  \item 3
  \item 1
  \item 0
  \item 2
  \item -1
  \item Todos los reales
\end{multicols}
\end{enumerate}
\vspace{0.15cm}
\item El intercepto con y de f(x)=2x^2-3x+1 es:
\begin{enumerate}[label=\textbf{\Alph*)}]
\begin{multicols}{2}
  \item 1
  \item 1
  \item 0
  \item 2
  \item -1
  \item Todos los reales
\end{multicols}
\end{enumerate}
\vspace{0.15cm}
\item El intercepto con y de f(x)=x^2+2x es:
\begin{enumerate}[label=\textbf{\Alph*)}]
\begin{multicols}{2}
  \item 0
  \item 1
  \item 0
  \item 2
  \item -1
  \item Todos los reales
\end{multicols}
\end{enumerate}
\vspace{0.15cm}
\item Si f(x)=x^2-2x-3, ¿cuál es f(0)?
\begin{enumerate}[label=\textbf{\Alph*)}]
\begin{multicols}{2}
  \item -3
  \item 1
  \item 0
  \item 2
  \item -1
  \item Todos los reales
\end{multicols}
\end{enumerate}
\vspace{0.15cm}
\item Si f(x)=x^2-2x-3, ¿cuál es f(1)?
\begin{enumerate}[label=\textbf{\Alph*)}]
\begin{multicols}{2}
  \item -4
  \item 1
  \item 0
  \item 2
  \item -1
  \item Todos los reales
\end{multicols}
\end{enumerate}
\vspace{0.15cm}
\item Si f(x)=x^2-2x-3, ¿cuál es f(3)?
\begin{enumerate}[label=\textbf{\Alph*)}]
\begin{multicols}{2}
  \item 0
  \item 1
  \item 0
  \item 2
  \item -1
  \item Todos los reales
\end{multicols}
\end{enumerate}
\vspace{0.15cm}
\item Si f(x)=x^2-2x-3, ¿cuál es f(-1)?
\begin{enumerate}[label=\textbf{\Alph*)}]
\begin{multicols}{2}
  \item 0
  \item 1
  \item 0
  \item 2
  \item -1
  \item Todos los reales
\end{multicols}
\end{enumerate}
\vspace{0.15cm}
\item El discriminante de f(x)=x^2-4x+3 es:
\begin{enumerate}[label=\textbf{\Alph*)}]
\begin{multicols}{2}
  \item 4
  \item 1
  \item 0
  \item 2
  \item -1
  \item Todos los reales
\end{multicols}
\end{enumerate}
\vspace{0.15cm}
\item El discriminante de f(x)=x^2-6x+9 es:
\begin{enumerate}[label=\textbf{\Alph*)}]
\begin{multicols}{2}
  \item 0
  \item 1
  \item 0
  \item 2
  \item -1
  \item Todos los reales
\end{multicols}
\end{enumerate}
\vspace{0.15cm}
\item El discriminante de f(x)=x^2+x+1 es:
\begin{enumerate}[label=\textbf{\Alph*)}]
\begin{multicols}{2}
  \item -3
  \item 1
  \item 0
  \item 2
  \item -1
  \item Todos los reales
\end{multicols}
\end{enumerate}
\vspace{0.15cm}
\item Si el discriminante es positivo, la parábola tiene:
\begin{enumerate}[label=\textbf{\Alph*)}]
\begin{multicols}{2}
  \item Dos interceptos con x
  \item 1
  \item 0
  \item 2
  \item -1
  \item Todos los reales
\end{multicols}
\end{enumerate}
\vspace{0.15cm}
\item Si el discriminante es cero, la parábola tiene:
\begin{enumerate}[label=\textbf{\Alph*)}]
\begin{multicols}{2}
  \item Un intercepto con x (doble)
  \item (1,1)
  \item 0
  \item 2
  \item -1
  \item Todos los reales
\end{multicols}
\end{enumerate}
\vspace{0.15cm}
\item Si el discriminante es negativo, la parábola tiene:
\begin{enumerate}[label=\textbf{\Alph*)}]
\begin{multicols}{2}
  \item Ningún intercepto con x
  \item 1
  \item 0
  \item 2
  \item -1
  \item Todos los reales
\end{multicols}
\end{enumerate}
\vspace{0.15cm}
\item La forma vértice de una parábola es:
\begin{enumerate}[label=\textbf{\Alph*)}]
\begin{multicols}{2}
  \item f(x)=a(x-h)^2+k
  \item (1,1)
  \item 0
  \item 2
  \item -1
  \item Todos los reales
\end{multicols}
\end{enumerate}
\vspace{0.15cm}
\item En f(x)=a(x-h)^2+k, el vértice es:
\begin{enumerate}[label=\textbf{\Alph*)}]
\begin{multicols}{2}
  \item (h,k)
  \item (1,1)
  \item 0
  \item 2
  \item -1
  \item Todos los reales
\end{multicols}
\end{enumerate}
\vspace{0.15cm}
\item La función f(x)=x^2 tiene un:
\begin{enumerate}[label=\textbf{\Alph*)}]
\begin{multicols}{2}
  \item Mínimo en (0,0)
  \item (1,1)
  \item 0
  \item 2
  \item -1
  \item Todos los reales
\end{multicols}
\end{enumerate}
\vspace{0.15cm}
\item La función f(x)=-x^2 tiene un:
\begin{enumerate}[label=\textbf{\Alph*)}]
\begin{multicols}{2}
  \item Máximo en (0,0)
  \item (1,1)
  \item 0
  \item 2
  \item -1
  \item Todos los reales
\end{multicols}
\end{enumerate}
\vspace{0.15cm}
\item La función f(x)=x^2-4x+3 tiene un:
\begin{enumerate}[label=\textbf{\Alph*)}]
\begin{multicols}{2}
  \item Mínimo en (2,-1)
  \item (1,1)
  \item 0
  \item 2
  \item -1
  \item Todos los reales
\end{multicols}
\end{enumerate}
\vspace{0.15cm}
\item La función f(x)=-x^2+2x+3 tiene un:
\begin{enumerate}[label=\textbf{\Alph*)}]
\begin{multicols}{2}
  \item Máximo en (1,4)
  \item (1,1)
  \item 0
  \item 2
  \item -1
  \item Todos los reales
\end{multicols}
\end{enumerate}
\vspace{0.15cm}
\item ¿Cuál es el valor mínimo de f(x)=x^2-6x+5?
\begin{enumerate}[label=\textbf{\Alph*)}]
\begin{multicols}{2}
  \item -4
  \item 1
  \item 0
  \item 2
  \item -1
  \item Todos los reales
\end{multicols}
\end{enumerate}
\vspace{0.15cm}
\item ¿Cuál es el valor máximo de f(x)=-x^2+4x-1?
\begin{enumerate}[label=\textbf{\Alph*)}]
\begin{multicols}{2}
  \item 3
  \item 1
  \item 0
  \item 2
  \item -1
  \item Todos los reales
\end{multicols}
\end{enumerate}
\vspace{0.15cm}
\item ¿Cuál es el valor mínimo de f(x)=x^2+2x+1?
\begin{enumerate}[label=\textbf{\Alph*)}]
\begin{multicols}{2}
  \item 0
  \item 1
  \item 0
  \item 2
  \item -1
  \item Todos los reales
\end{multicols}
\end{enumerate}
\vspace{0.15cm}
\item ¿Cuál es el valor mínimo de f(x)=2x^2-4x+1?
\begin{enumerate}[label=\textbf{\Alph*)}]
\begin{multicols}{2}
  \item -1
  \item 1
  \item 0
  \item 2
  \item -1
  \item Todos los reales
\end{multicols}
\end{enumerate}
\vspace{0.15cm}
\item La parábola f(x)=x^2-4x+3 es:
\begin{enumerate}[label=\textbf{\Alph*)}]
\begin{multicols}{2}
  \item Cóncava hacia arriba
  \item 1
  \item 0
  \item 2
  \item -1
  \item Todos los reales
\end{multicols}
\end{enumerate}
\vspace{0.15cm}
\item La parábola f(x)=-x^2+2x+3 es:
\begin{enumerate}[label=\textbf{\Alph*)}]
\begin{multicols}{2}
  \item Cóncava hacia abajo
  \item 1
  \item 0
  \item 2
  \item -1
  \item Todos los reales
\end{multicols}
\end{enumerate}
\vspace{0.15cm}
\item Si f(x)=x^2-2x-3, ¿cuáles son los interceptos con x?
\begin{enumerate}[label=\textbf{\Alph*)}]
\begin{multicols}{2}
  \item x=3, x=-1
  \item 1
  \item 0
  \item 2
  \item -1
  \item Todos los reales
\end{multicols}
\end{enumerate}
\vspace{0.15cm}
\item Si f(x)=x^2+3x-10, ¿cuáles son los interceptos con x?
\begin{enumerate}[label=\textbf{\Alph*)}]
\begin{multicols}{2}
  \item x=2, x=-5
  \item 1
  \item 0
  \item 2
  \item -1
  \item Todos los reales
\end{multicols}
\end{enumerate}
\vspace{0.15cm}
\item Si f(x)=x^2-3x-10, ¿cuáles son los interceptos con x?
\begin{enumerate}[label=\textbf{\Alph*)}]
\begin{multicols}{2}
  \item x=5, x=-2
  \item 1
  \item 0
  \item 2
  \item -1
  \item Todos los reales
\end{multicols}
\end{enumerate}
\vspace{0.15cm}
\item Si f(x)=2x^2-7x+3, ¿cuáles son los interceptos con x?
\begin{enumerate}[label=\textbf{\Alph*)}]
\begin{multicols}{2}
  \item x=3, x=\frac{1}{2}
  \item 1
  \item 0
  \item 2
  \item -1
  \item Todos los reales
\end{multicols}
\end{enumerate}
\vspace{0.15cm}
\item Si f(x)=3x^2-10x+8, ¿cuáles son los interceptos con x?
\begin{enumerate}[label=\textbf{\Alph*)}]
\begin{multicols}{2}
  \item x=2, x=\frac{4}{3}
  \item 1
  \item 0
  \item 2
  \item -1
  \item Todos los reales
\end{multicols}
\end{enumerate}
\vspace{0.15cm}
\item La función f(x)=x^2+1 tiene:
\begin{enumerate}[label=\textbf{\Alph*)}]
\begin{multicols}{2}
  \item Ningún intercepto con x
  \item 1
  \item 0
  \item 2
  \item -1
  \item Todos los reales
\end{multicols}
\end{enumerate}
\vspace{0.15cm}
\item La función f(x)=x^2-6x+9 tiene:
\begin{enumerate}[label=\textbf{\Alph*)}]
\begin{multicols}{2}
  \item Un intercepto con x (doble)
  \item (1,1)
  \item 0
  \item 2
  \item -1
  \item Todos los reales
\end{multicols}
\end{enumerate}
\vspace{0.15cm}
\item La función f(x)=x^2+x+1 tiene:
\begin{enumerate}[label=\textbf{\Alph*)}]
\begin{multicols}{2}
  \item Ningún intercepto con x
  \item 1
  \item 0
  \item 2
  \item -1
  \item Todos los reales
\end{multicols}
\end{enumerate}
\vspace{0.15cm}
\item Si f(x)=x^2-4x+3, ¿cuál es el valor de f(2)?
\begin{enumerate}[label=\textbf{\Alph*)}]
\begin{multicols}{2}
  \item -1
  \item 1
  \item 0
  \item 2
  \item -1
  \item Todos los reales
\end{multicols}
\end{enumerate}
\vspace{0.15cm}
\item Si f(x)=x^2-6x+5, ¿cuál es el valor de f(3)?
\begin{enumerate}[label=\textbf{\Alph*)}]
\begin{multicols}{2}
  \item -4
  \item 1
  \item 0
  \item 2
  \item -1
  \item Todos los reales
\end{multicols}
\end{enumerate}
\vspace{0.15cm}
\item Si f(x)=x^2+2x+1, ¿cuál es el valor de f(-1)?
\begin{enumerate}[label=\textbf{\Alph*)}]
\begin{multicols}{2}
  \item 0
  \item 1
  \item 0
  \item 2
  \item -1
  \item Todos los reales
\end{multicols}
\end{enumerate}
\vspace{0.15cm}
\item Si f(x)=2x^2-4x+1, ¿cuál es el valor de f(1)?
\begin{enumerate}[label=\textbf{\Alph*)}]
\begin{multicols}{2}
  \item -1
  \item 1
  \item 0
  \item 2
  \item -1
  \item Todos los reales
\end{multicols}
\end{enumerate}
\vspace{0.15cm}
\item Si f(x)=-x^2+2x+3, ¿cuál es el valor de f(1)?
\begin{enumerate}[label=\textbf{\Alph*)}]
\begin{multicols}{2}
  \item 4
  \item 1
  \item 0
  \item 2
  \item -1
  \item Todos los reales
\end{multicols}
\end{enumerate}
\vspace{0.15cm}
\item La función f(x)=x^2-4x+3 es simétrica respecto a:
\begin{enumerate}[label=\textbf{\Alph*)}]
\begin{multicols}{2}
  \item x = 2
  \item 1
  \item 0
  \item 2
  \item -1
  \item Todos los reales
\end{multicols}
\end{enumerate}
\vspace{0.15cm}
\item La función f(x)=x^2+6x+5 es simétrica respecto a:
\begin{enumerate}[label=\textbf{\Alph*)}]
\begin{multicols}{2}
  \item x = -3
  \item 1
  \item 0
  \item 2
  \item -1
  \item Todos los reales
\end{multicols}
\end{enumerate}
\vspace{0.15cm}
\item La función f(x)=2x^2-8x+1 es simétrica respecto a:
\begin{enumerate}[label=\textbf{\Alph*)}]
\begin{multicols}{2}
  \item x = 2
  \item 1
  \item 0
  \item 2
  \item -1
  \item Todos los reales
\end{multicols}
\end{enumerate}
\vspace{0.15cm}
\item Si f(x)=x^2-4x+3, ¿cuál es el rango?
\begin{enumerate}[label=\textbf{\Alph*)}]
\begin{multicols}{2}
  \item y \geq -1
  \item 1
  \item 0
  \item 2
  \item -1
  \item Todos los reales
\end{multicols}
\end{enumerate}
\vspace{0.15cm}
\item Si f(x)=-x^2+2x+3, ¿cuál es el rango?
\begin{enumerate}[label=\textbf{\Alph*)}]
\begin{multicols}{2}
  \item y \leq 4
  \item 1
  \item 0
  \item 2
  \item -1
  \item Todos los reales
\end{multicols}
\end{enumerate}
\vspace{0.15cm}
\item Si f(x)=x^2+2x+1, ¿cuál es el rango?
\begin{enumerate}[label=\textbf{\Alph*)}]
\begin{multicols}{2}
  \item y \geq 0
  \item 1
  \item 0
  \item 2
  \item -1
  \item Todos los reales
\end{multicols}
\end{enumerate}
\vspace{0.15cm}
\item Si f(x)=2x^2-4x+1, ¿cuál es el rango?
\begin{enumerate}[label=\textbf{\Alph*)}]
\begin{multicols}{2}
  \item y \geq -1
  \item 1
  \item 0
  \item 2
  \item -1
  \item Todos los reales
\end{multicols}
\end{enumerate}
\vspace{0.15cm}
\end{enumerate}\n
\subsubsection*{SECCIÓN IV: OPERACIONES CON FUNCIONES (171–220)}

\begin{enumerate}[label=\textbf{\arabic*.}]
\setcounter{enumi}{170}
\item f(x)=2x+1 y g(x)=x-3, (f+g)(x):
\begin{enumerate}[label=\textbf{\Alph*)}]
\begin{multicols}{2}
  \item 3x-2
  \item 2x+1
  \item x
  \item x^2
  \item 1
  \item 0
\end{multicols}
\end{enumerate}
\vspace{0.15cm}
\item f(x)=x+2 y g(x)=3x-1, (f+g)(x):
\begin{enumerate}[label=\textbf{\Alph*)}]
\begin{multicols}{2}
  \item 4x+1
  \item 2x+1
  \item x
  \item x^2
  \item 1
  \item 0
\end{multicols}
\end{enumerate}
\vspace{0.15cm}
\item f(x)=x^2 y g(x)=x+1, (f+g)(x):
\begin{enumerate}[label=\textbf{\Alph*)}]
\begin{multicols}{2}
  \item x^2+x+1
  \item 2x+1
  \item x
  \item x^2
  \item 1
  \item 0
\end{multicols}
\end{enumerate}
\vspace{0.15cm}
\item f(x)=2x+1 y g(x)=x-3, (f-g)(x):
\begin{enumerate}[label=\textbf{\Alph*)}]
\begin{multicols}{2}
  \item x+4
  \item 2x+1
  \item x
  \item x^2
  \item 1
  \item 0
\end{multicols}
\end{enumerate}
\vspace{0.15cm}
\item f(x)=x^2 y g(x)=x+1, (f-g)(x):
\begin{enumerate}[label=\textbf{\Alph*)}]
\begin{multicols}{2}
  \item x^2-x-1
  \item 2x+1
  \item x
  \item x^2
  \item 1
  \item 0
\end{multicols}
\end{enumerate}
\vspace{0.15cm}
\item f(x)=2x+1 y g(x)=x-3, (f*g)(x):
\begin{enumerate}[label=\textbf{\Alph*)}]
\begin{multicols}{2}
  \item 2x^2-5x-3
  \item 2x+1
  \item x
  \item x^2
  \item 1
  \item 0
\end{multicols}
\end{enumerate}
\vspace{0.15cm}
\item f(x)=x+1 y g(x)=x-1, (f*g)(x):
\begin{enumerate}[label=\textbf{\Alph*)}]
\begin{multicols}{2}
  \item x^2-1
  \item 2x+1
  \item x
  \item x^2
  \item 1
  \item 0
\end{multicols}
\end{enumerate}
\vspace{0.15cm}
\item f(x)=x^2 y g(x)=x+1, (f*g)(x):
\begin{enumerate}[label=\textbf{\Alph*)}]
\begin{multicols}{2}
  \item x^3+x^2
  \item 2x+1
  \item x
  \item x^2
  \item 1
  \item 0
\end{multicols}
\end{enumerate}
\vspace{0.15cm}
\item f(x)=2x y g(x)=x+3, (f o g)(x):
\begin{enumerate}[label=\textbf{\Alph*)}]
\begin{multicols}{2}
  \item 2x+6
  \item 2x+1
  \item x
  \item x^2
  \item 1
  \item 0
\end{multicols}
\end{enumerate}
\vspace{0.15cm}
\item f(x)=x+1 y g(x)=2x, (f o g)(x):
\begin{enumerate}[label=\textbf{\Alph*)}]
\begin{multicols}{2}
  \item 2x+1
  \item 2x+1
  \item x
  \item x^2
  \item 1
  \item 0
\end{multicols}
\end{enumerate}
\vspace{0.15cm}
\item f(x)=x^2 y g(x)=x+1, (f o g)(x):
\begin{enumerate}[label=\textbf{\Alph*)}]
\begin{multicols}{2}
  \item (x+1)^2
  \item 2x+1
  \item x
  \item x^2
  \item 1
  \item 0
\end{multicols}
\end{enumerate}
\vspace{0.15cm}
\item f(x)=x-2 y g(x)=3x, (f o g)(x):
\begin{enumerate}[label=\textbf{\Alph*)}]
\begin{multicols}{2}
  \item 3x-2
  \item 2x+1
  \item x
  \item x^2
  \item 1
  \item 0
\end{multicols}
\end{enumerate}
\vspace{0.15cm}
\item f(x)=2x+1 y g(x)=x-3, (f o g)(x):
\begin{enumerate}[label=\textbf{\Alph*)}]
\begin{multicols}{2}
  \item 2x-5
  \item 2x+1
  \item x
  \item x^2
  \item 1
  \item 0
\end{multicols}
\end{enumerate}
\vspace{0.15cm}
\item f(x)=x^2 y g(x)=2x, (f o g)(x):
\begin{enumerate}[label=\textbf{\Alph*)}]
\begin{multicols}{2}
  \item 4x^2
  \item 2x+1
  \item x
  \item x^2
  \item 1
  \item 0
\end{multicols}
\end{enumerate}
\vspace{0.15cm}
\item f(x)=3x y g(x)=x^2, (f o g)(x):
\begin{enumerate}[label=\textbf{\Alph*)}]
\begin{multicols}{2}
  \item 3x^2
  \item 2x+1
  \item x
  \item x^2
  \item 1
  \item 0
\end{multicols}
\end{enumerate}
\vspace{0.15cm}
\item f(x)=x+2 y g(x)=x-2, (f o g)(x):
\begin{enumerate}[label=\textbf{\Alph*)}]
\begin{multicols}{2}
  \item x
  \item 2x+1
  \item 2x
  \item x^2
  \item 1
  \item 0
\end{multicols}
\end{enumerate}
\vspace{0.15cm}
\item f(x)=2x y g(x)=x/2, (f o g)(x):
\begin{enumerate}[label=\textbf{\Alph*)}]
\begin{multicols}{2}
  \item x
  \item 2x+1
  \item 2x
  \item x^2
  \item 1
  \item 0
\end{multicols}
\end{enumerate}
\vspace{0.15cm}
\item f(x)=x^2 y g(x)=\sqrt{x}, (f o g)(x):
\begin{enumerate}[label=\textbf{\Alph*)}]
\begin{multicols}{2}
  \item x
  \item 2x+1
  \item 2x
  \item x^2
  \item 1
  \item 0
\end{multicols}
\end{enumerate}
\vspace{0.15cm}
\item f(x)=\sqrt{x} y g(x)=x^2, (f o g)(x):
\begin{enumerate}[label=\textbf{\Alph*)}]
\begin{multicols}{2}
  \item |x|
  \item 2x+1
  \item x
  \item x^2
  \item 1
  \item 0
\end{multicols}
\end{enumerate}
\vspace{0.15cm}
\item Si f(x)=x+1, f^{-1}(x) es:
\begin{enumerate}[label=\textbf{\Alph*)}]
\begin{multicols}{2}
  \item x-1
  \item 2x+1
  \item x
  \item x^2
  \item 1
  \item 0
\end{multicols}
\end{enumerate}
\vspace{0.15cm}
\item Si f(x)=2x, f^{-1}(x) es:
\begin{enumerate}[label=\textbf{\Alph*)}]
\begin{multicols}{2}
  \item x/2
  \item 2x+1
  \item x
  \item x^2
  \item 1
  \item 0
\end{multicols}
\end{enumerate}
\vspace{0.15cm}
\item Si f(x)=3x-2, f^{-1}(x) es:
\begin{enumerate}[label=\textbf{\Alph*)}]
\begin{multicols}{2}
  \item (x+2)/3
  \item 2x+1
  \item x
  \item x^2
  \item 1
  \item 0
\end{multicols}
\end{enumerate}
\vspace{0.15cm}
\item Si f(x)=x^3, f^{-1}(x) es:
\begin{enumerate}[label=\textbf{\Alph*)}]
\begin{multicols}{2}
  \item \sqrt[3]{x}
  \item 2x+1
  \item x
  \item x^2
  \item 1
  \item 0
\end{multicols}
\end{enumerate}
\vspace{0.15cm}
\item Si f(x)=1/x, f^{-1}(x) es:
\begin{enumerate}[label=\textbf{\Alph*)}]
\begin{multicols}{2}
  \item 1/x
  \item 2x+1
  \item x
  \item x^2
  \item 1
  \item 0
\end{multicols}
\end{enumerate}
\vspace{0.15cm}
\item Si f(x)=x+3, f^{-1}(2) es:
\begin{enumerate}[label=\textbf{\Alph*)}]
\begin{multicols}{2}
  \item -1
  \item 2
  \item x
  \item x^2
  \item 1
  \item 0
\end{multicols}
\end{enumerate}
\vspace{0.15cm}
\item Si f(x)=2x-1, f^{-1}(3) es:
\begin{enumerate}[label=\textbf{\Alph*)}]
\begin{multicols}{2}
  \item 2
  \item 2
  \item x
  \item x^2
  \item 1
  \item 0
\end{multicols}
\end{enumerate}
\vspace{0.15cm}
\item Si f(x)=3x+2, f^{-1}(8) es:
\begin{enumerate}[label=\textbf{\Alph*)}]
\begin{multicols}{2}
  \item 2
  \item 2
  \item x
  \item x^2
  \item 1
  \item 0
\end{multicols}
\end{enumerate}
\vspace{0.15cm}
\item Si f(x)=x^2 con x>=0, f^{-1}(x) es:
\begin{enumerate}[label=\textbf{\Alph*)}]
\begin{multicols}{2}
  \item \sqrt{x}
  \item 2x+1
  \item x
  \item x^2
  \item 1
  \item 0
\end{multicols}
\end{enumerate}
\vspace{0.15cm}
\item Si f(x)=\sqrt{x}, f^{-1}(x) es:
\begin{enumerate}[label=\textbf{\Alph*)}]
\begin{multicols}{2}
  \item x^2
  \item 2x+1
  \item x
  \item x^2
  \item 1
  \item 0
\end{multicols}
\end{enumerate}
\vspace{0.15cm}
\item Si f(x)=x^2 con x>=0, f^{-1}(9) es:
\begin{enumerate}[label=\textbf{\Alph*)}]
\begin{multicols}{2}
  \item 3
  \item 2
  \item x
  \item x^2
  \item 1
  \item 0
\end{multicols}
\end{enumerate}
\vspace{0.15cm}
\item Si f(x)=x^2 con x>=0, f^{-1}(16) es:
\begin{enumerate}[label=\textbf{\Alph*)}]
\begin{multicols}{2}
  \item 4
  \item 2
  \item x
  \item x^2
  \item 1
  \item 0
\end{multicols}
\end{enumerate}
\vspace{0.15cm}
\item Si f(x)=x^2 con x>=0, f^{-1}(25) es:
\begin{enumerate}[label=\textbf{\Alph*)}]
\begin{multicols}{2}
  \item 5
  \item 2
  \item x
  \item x^2
  \item 1
  \item 0
\end{multicols}
\end{enumerate}
\vspace{0.15cm}
\item La función f(x)=2x+1 es:
\begin{enumerate}[label=\textbf{\Alph*)}]
\begin{multicols}{2}
  \item Inyectiva
  \item 2
  \item x
  \item x^2
  \item 1
  \item 0
\end{multicols}
\end{enumerate}
\vspace{0.15cm}
\item La función f(x)=x^2 con dominio R es:
\begin{enumerate}[label=\textbf{\Alph*)}]
\begin{multicols}{2}
  \item No inyectiva
  \item 2
  \item x
  \item x^2
  \item 1
  \item 0
\end{multicols}
\end{enumerate}
\vspace{0.15cm}
\item La función f(x)=x^2 con dominio x>=0 es:
\begin{enumerate}[label=\textbf{\Alph*)}]
\begin{multicols}{2}
  \item Inyectiva
  \item 2
  \item x
  \item x^2
  \item 1
  \item 0
\end{multicols}
\end{enumerate}
\vspace{0.15cm}
\item Una función es inyectiva si:
\begin{enumerate}[label=\textbf{\Alph*)}]
\begin{multicols}{2}
  \item Cada valor del rango corresponde a un solo valor del dominio
  \item 2
  \item x
  \item x^2
  \item 1
  \item 0
\end{multicols}
\end{enumerate}
\vspace{0.15cm}
\item Una función es sobreyectiva si:
\begin{enumerate}[label=\textbf{\Alph*)}]
\begin{multicols}{2}
  \item El rango es igual al codominio
  \item 2
  \item x
  \item x^2
  \item 1
  \item 0
\end{multicols}
\end{enumerate}
\vspace{0.15cm}
\item Una función es biyectiva si:
\begin{enumerate}[label=\textbf{\Alph*)}]
\begin{multicols}{2}
  \item Es inyectiva y sobreyectiva
  \item 2
  \item x
  \item x^2
  \item 1
  \item 0
\end{multicols}
\end{enumerate}
\vspace{0.15cm}
\item Si f(x)=2x+1, f(f^{-1}(5)) es:
\begin{enumerate}[label=\textbf{\Alph*)}]
\begin{multicols}{2}
  \item 5
  \item 2
  \item x
  \item x^2
  \item 1
  \item 0
\end{multicols}
\end{enumerate}
\vspace{0.15cm}
\item Si f(x)=3x-2, f^{-1}(f(4)) es:
\begin{enumerate}[label=\textbf{\Alph*)}]
\begin{multicols}{2}
  \item 4
  \item 2
  \item x
  \item x^2
  \item 1
  \item 0
\end{multicols}
\end{enumerate}
\vspace{0.15cm}
\item Si f(x)=x+1 y g(x)=2x, (g o f)(x) es:
\begin{enumerate}[label=\textbf{\Alph*)}]
\begin{multicols}{2}
  \item 2x+2
  \item 2x+1
  \item x
  \item x^2
  \item 1
  \item 0
\end{multicols}
\end{enumerate}
\vspace{0.15cm}
\item Si f(x)=x^2 y g(x)=x-1, (g o f)(x) es:
\begin{enumerate}[label=\textbf{\Alph*)}]
\begin{multicols}{2}
  \item x^2-1
  \item 2x+1
  \item x
  \item x^2
  \item 1
  \item 0
\end{multicols}
\end{enumerate}
\vspace{0.15cm}
\item Si f(x)=2x+1 y g(x)=x^2, (g o f)(x) es:
\begin{enumerate}[label=\textbf{\Alph*)}]
\begin{multicols}{2}
  \item (2x+1)^2
  \item 2x+1
  \item x
  \item x^2
  \item 1
  \item 0
\end{multicols}
\end{enumerate}
\vspace{0.15cm}
\item Si f(x)=x+2 y g(x)=x^2, (f o g)(x) es:
\begin{enumerate}[label=\textbf{\Alph*)}]
\begin{multicols}{2}
  \item x^2+2
  \item 2x+1
  \item x
  \item x^2
  \item 1
  \item 0
\end{multicols}
\end{enumerate}
\vspace{0.15cm}
\item Si f(x)=3x y g(x)=x+1, (f o g)(x) es:
\begin{enumerate}[label=\textbf{\Alph*)}]
\begin{multicols}{2}
  \item 3x+3
  \item 2x+1
  \item x
  \item x^2
  \item 1
  \item 0
\end{multicols}
\end{enumerate}
\vspace{0.15cm}
\item Si f(x)=x-1 y g(x)=x+1, (f o g)(x) es:
\begin{enumerate}[label=\textbf{\Alph*)}]
\begin{multicols}{2}
  \item x
  \item 2x+1
  \item 2x
  \item x^2
  \item 1
  \item 0
\end{multicols}
\end{enumerate}
\vspace{0.15cm}
\item Si f(x)=1/x y g(x)=x+1, (f o g)(x) es:
\begin{enumerate}[label=\textbf{\Alph*)}]
\begin{multicols}{2}
  \item 1/(x+1)
  \item 2x+1
  \item x
  \item x^2
  \item 1
  \item 0
\end{multicols}
\end{enumerate}
\vspace{0.15cm}
\item Si f(x)=x^2+1, f(f(1)) es:
\begin{enumerate}[label=\textbf{\Alph*)}]
\begin{multicols}{2}
  \item 5
  \item 2
  \item x
  \item x^2
  \item 1
  \item 0
\end{multicols}
\end{enumerate}
\vspace{0.15cm}
\item Si f(x)=2x-3, f(f(2)) es:
\begin{enumerate}[label=\textbf{\Alph*)}]
\begin{multicols}{2}
  \item -1
  \item 2
  \item x
  \item x^2
  \item 1
  \item 0
\end{multicols}
\end{enumerate}
\vspace{0.15cm}
\item Si f(x)=x+1, f(f(f(1))) es:
\begin{enumerate}[label=\textbf{\Alph*)}]
\begin{multicols}{2}
  \item 4
  \item 2
  \item x
  \item x^2
  \item 1
  \item 0
\end{multicols}
\end{enumerate}
\vspace{0.15cm}
\end{enumerate}\n
\subsubsection*{SECCIÓN V: PROBLEMAS DE APLICACIÓN (221–270)}

\begin{enumerate}[label=\textbf{\arabic*.}]
\setcounter{enumi}{220}
\item Un taxi cobra $2 base más $1.5 por km. Costo:
\begin{enumerate}[label=\textbf{\Alph*)}]
\begin{multicols}{2}
  \item C(x)=1.5x+2
  \item 2
  \item 10
  \item 15
  \item 25
  \item 0
\end{multicols}
\end{enumerate}
\vspace{0.15cm}
\item Costo por 10 km taxi anterior:
\begin{enumerate}[label=\textbf{\Alph*)}]
\begin{multicols}{2}
  \item $17
  \item $20
  \item 10
  \item 15
  \item 25
  \item 0
\end{multicols}
\end{enumerate}
\vspace{0.15cm}
\item Plan celular $10 fijos más $0.5 por minuto:
\begin{enumerate}[label=\textbf{\Alph*)}]
\begin{multicols}{2}
  \item C(x)=0.5x+10
  \item 2
  \item 10
  \item 15
  \item 25
  \item 0
\end{multicols}
\end{enumerate}
\vspace{0.15cm}
\item Costo plan celular por 60 minutos:
\begin{enumerate}[label=\textbf{\Alph*)}]
\begin{multicols}{2}
  \item $40
  \item $20
  \item 10
  \item 15
  \item 25
  \item 0
\end{multicols}
\end{enumerate}
\vspace{0.15cm}
\item 100°C a Fahrenheit F(C)=9/5 C + 32:
\begin{enumerate}[label=\textbf{\Alph*)}]
\begin{multicols}{2}
  \item 212°F
  \item 2
  \item 10
  \item 15
  \item 25
  \item 0
\end{multicols}
\end{enumerate}
\vspace{0.15cm}
\item 0°C a Fahrenheit:
\begin{enumerate}[label=\textbf{\Alph*)}]
\begin{multicols}{2}
  \item 32°F
  \item 2
  \item 10
  \item 15
  \item 25
  \item 0
\end{multicols}
\end{enumerate}
\vspace{0.15cm}
\item 37°C a Fahrenheit:
\begin{enumerate}[label=\textbf{\Alph*)}]
\begin{multicols}{2}
  \item 98.6°F
  \item 2
  \item 10
  \item 15
  \item 25
  \item 0
\end{multicols}
\end{enumerate}
\vspace{0.15cm}
\item 212°F a Celsius C(F)=5/9(F-32):
\begin{enumerate}[label=\textbf{\Alph*)}]
\begin{multicols}{2}
  \item 100°C
  \item 2
  \item 10
  \item 15
  \item 25
  \item 0
\end{multicols}
\end{enumerate}
\vspace{0.15cm}
\item 32°F a Celsius:
\begin{enumerate}[label=\textbf{\Alph*)}]
\begin{multicols}{2}
  \item 0°C
  \item 2
  \item 10
  \item 15
  \item 25
  \item 0
\end{multicols}
\end{enumerate}
\vspace{0.15cm}
\item 98.6°F a Celsius:
\begin{enumerate}[label=\textbf{\Alph*)}]
\begin{multicols}{2}
  \item 37°C
  \item 2
  \item 10
  \item 15
  \item 25
  \item 0
\end{multicols}
\end{enumerate}
\vspace{0.15cm}
\item Área de un círculo en función del radio:
\begin{enumerate}[label=\textbf{\Alph*)}]
\begin{multicols}{2}
  \item A(r)=\pi r^2
  \item 2
  \item 10
  \item 15
  \item 25
  \item 0
\end{multicols}
\end{enumerate}
\vspace{0.15cm}
\item Área círculo radio 5:
\begin{enumerate}[label=\textbf{\Alph*)}]
\begin{multicols}{2}
  \item 25\pi
  \item 2
  \item 10
  \item 15
  \item 25
  \item 0
\end{multicols}
\end{enumerate}
\vspace{0.15cm}
\item Perímetro de cuadrado en función de su lado:
\begin{enumerate}[label=\textbf{\Alph*)}]
\begin{multicols}{2}
  \item P(l)=4l
  \item 2
  \item 10
  \item 15
  \item 25
  \item 0
\end{multicols}
\end{enumerate}
\vspace{0.15cm}
\item Perímetro cuadrado lado 6:
\begin{enumerate}[label=\textbf{\Alph*)}]
\begin{multicols}{2}
  \item 24
  \item 2
  \item 10
  \item 15
  \item 25
  \item 0
\end{multicols}
\end{enumerate}
\vspace{0.15cm}
\item Área cuadrado función de su lado:
\begin{enumerate}[label=\textbf{\Alph*)}]
\begin{multicols}{2}
  \item A(l)=l^2
  \item 2
  \item 10
  \item 15
  \item 25
  \item 0
\end{multicols}
\end{enumerate}
\vspace{0.15cm}
\item Área cuadrado lado 5:
\begin{enumerate}[label=\textbf{\Alph*)}]
\begin{multicols}{2}
  \item 25
  \item 2
  \item 10
  \item 15
  \item 25
  \item 0
\end{multicols}
\end{enumerate}
\vspace{0.15cm}
\item Volumen cubo en función de su arista:
\begin{enumerate}[label=\textbf{\Alph*)}]
\begin{multicols}{2}
  \item V(a)=a^3
  \item 2
  \item 10
  \item 15
  \item 25
  \item 0
\end{multicols}
\end{enumerate}
\vspace{0.15cm}
\item Volumen cubo arista 3:
\begin{enumerate}[label=\textbf{\Alph*)}]
\begin{multicols}{2}
  \item 27
  \item 2
  \item 10
  \item 15
  \item 25
  \item 0
\end{multicols}
\end{enumerate}
\vspace{0.15cm}
\item Distancia función del tiempo velocidad v:
\begin{enumerate}[label=\textbf{\Alph*)}]
\begin{multicols}{2}
  \item d(t)=vt
  \item 2
  \item 10
  \item 15
  \item 25
  \item 0
\end{multicols}
\end{enumerate}
\vspace{0.15cm}
\item Auto a 60 km/h en 3 horas recorre:
\begin{enumerate}[label=\textbf{\Alph*)}]
\begin{multicols}{2}
  \item 180 km
  \item 2
  \item 10
  \item 15
  \item 25
  \item 0
\end{multicols}
\end{enumerate}
\vspace{0.15cm}
\item Ganancia G(x)=5x-200. Unidades para no perder:
\begin{enumerate}[label=\textbf{\Alph*)}]
\begin{multicols}{2}
  \item 40
  \item 2
  \item 10
  \item 15
  \item 25
  \item 0
\end{multicols}
\end{enumerate}
\vspace{0.15cm}
\item Ganancia G(x)=5x-200 por 100 unidades:
\begin{enumerate}[label=\textbf{\Alph*)}]
\begin{multicols}{2}
  \item $300
  \item $20
  \item 10
  \item 15
  \item 25
  \item 0
\end{multicols}
\end{enumerate}
\vspace{0.15cm}
\item Altura h(t)=-5t^2+20t. Altura máxima:
\begin{enumerate}[label=\textbf{\Alph*)}]
\begin{multicols}{2}
  \item 20 m
  \item 2
  \item 10
  \item 15
  \item 25
  \item 0
\end{multicols}
\end{enumerate}
\vspace{0.15cm}
\item Tiempo altura máxima proyectil anterior:
\begin{enumerate}[label=\textbf{\Alph*)}]
\begin{multicols}{2}
  \item 2 s
  \item 2
  \item 10
  \item 15
  \item 25
  \item 0
\end{multicols}
\end{enumerate}
\vspace{0.15cm}
\item Altura h(t)=-5t^2+30t. Altura máxima:
\begin{enumerate}[label=\textbf{\Alph*)}]
\begin{multicols}{2}
  \item 45 m
  \item 2
  \item 10
  \item 15
  \item 25
  \item 0
\end{multicols}
\end{enumerate}
\vspace{0.15cm}
\item Tiempo altura máxima anterior:
\begin{enumerate}[label=\textbf{\Alph*)}]
\begin{multicols}{2}
  \item 3 s
  \item 2
  \item 10
  \item 15
  \item 25
  \item 0
\end{multicols}
\end{enumerate}
\vspace{0.15cm}
\item Utilidad U(x)=-x^2+100x-1000. Utilidad máxima:
\begin{enumerate}[label=\textbf{\Alph*)}]
\begin{multicols}{2}
  \item $1500
  \item $20
  \item 10
  \item 15
  \item 25
  \item 0
\end{multicols}
\end{enumerate}
\vspace{0.15cm}
\item Valor de x para utilidad máxima anterior:
\begin{enumerate}[label=\textbf{\Alph*)}]
\begin{multicols}{2}
  \item 50
  \item 2
  \item 10
  \item 15
  \item 25
  \item 0
\end{multicols}
\end{enumerate}
\vspace{0.15cm}
\item Costo C(x)=2x+500, venta a $5. Ingreso I(x):
\begin{enumerate}[label=\textbf{\Alph*)}]
\begin{multicols}{2}
  \item I(x)=5x
  \item 2
  \item 10
  \item 15
  \item 25
  \item 0
\end{multicols}
\end{enumerate}
\vspace{0.15cm}
\item Ganancia del problema anterior:
\begin{enumerate}[label=\textbf{\Alph*)}]
\begin{multicols}{2}
  \item G(x)=3x-500
  \item 2
  \item 10
  \item 15
  \item 25
  \item 0
\end{multicols}
\end{enumerate}
\vspace{0.15cm}
\item Artículos para no perder problema anterior:
\begin{enumerate}[label=\textbf{\Alph*)}]
\begin{multicols}{2}
  \item 167
  \item 2
  \item 10
  \item 15
  \item 25
  \item 0
\end{multicols}
\end{enumerate}
\vspace{0.15cm}
\item Población P(t)=10000(1+0.02t) en 5 años:
\begin{enumerate}[label=\textbf{\Alph*)}]
\begin{multicols}{2}
  \item 11000
  \item 2
  \item 10
  \item 15
  \item 25
  \item 0
\end{multicols}
\end{enumerate}
\vspace{0.15cm}
\item Población P(t)=5000(1+0.03t) en 10 años:
\begin{enumerate}[label=\textbf{\Alph*)}]
\begin{multicols}{2}
  \item 6500
  \item 2
  \item 10
  \item 15
  \item 25
  \item 0
\end{multicols}
\end{enumerate}
\vspace{0.15cm}
\item Interés simple C=1000, r=0.05, t=3:
\begin{enumerate}[label=\textbf{\Alph*)}]
\begin{multicols}{2}
  \item $150
  \item $20
  \item 10
  \item 15
  \item 25
  \item 0
\end{multicols}
\end{enumerate}
\vspace{0.15cm}
\item Monto final C=2000, r=0.04, t=5:
\begin{enumerate}[label=\textbf{\Alph*)}]
\begin{multicols}{2}
  \item $2400
  \item $20
  \item 10
  \item 15
  \item 25
  \item 0
\end{multicols}
\end{enumerate}
\vspace{0.15cm}
\item Depreciación V(t)=20000-2000t a los 5 años:
\begin{enumerate}[label=\textbf{\Alph*)}]
\begin{multicols}{2}
  \item $10000
  \item $20
  \item 10
  \item 15
  \item 25
  \item 0
\end{multicols}
\end{enumerate}
\vspace{0.15cm}
\item Años para que auto anterior valga 0:
\begin{enumerate}[label=\textbf{\Alph*)}]
\begin{multicols}{2}
  \item 10 años
  \item 2
  \item 10
  \item 15
  \item 25
  \item 0
\end{multicols}
\end{enumerate}
\vspace{0.15cm}
\item Temperatura T(t)=100-5t a los 10 minutos:
\begin{enumerate}[label=\textbf{\Alph*)}]
\begin{multicols}{2}
  \item 50°C
  \item 2
  \item 10
  \item 15
  \item 25
  \item 0
\end{multicols}
\end{enumerate}
\vspace{0.15cm}
\item Presión P(h)=1+0.1h a 50 m de profundidad:
\begin{enumerate}[label=\textbf{\Alph*)}]
\begin{multicols}{2}
  \item 6 atm
  \item 2
  \item 10
  \item 15
  \item 25
  \item 0
\end{multicols}
\end{enumerate}
\vspace{0.15cm}
\item Longitud resorte L(w)=10+0.5w con 20 kg:
\begin{enumerate}[label=\textbf{\Alph*)}]
\begin{multicols}{2}
  \item 20 cm
  \item 2
  \item 10
  \item 15
  \item 25
  \item 0
\end{multicols}
\end{enumerate}
\vspace{0.15cm}
\item Salario S(x)=300+0.1x con ventas $5000:
\begin{enumerate}[label=\textbf{\Alph*)}]
\begin{multicols}{2}
  \item $800
  \item $20
  \item 10
  \item 15
  \item 25
  \item 0
\end{multicols}
\end{enumerate}
\vspace{0.15cm}
\item Frenado d(v)=0.05v^2 a 100 km/h:
\begin{enumerate}[label=\textbf{\Alph*)}]
\begin{multicols}{2}
  \item 500 m
  \item 2
  \item 10
  \item 15
  \item 25
  \item 0
\end{multicols}
\end{enumerate}
\vspace{0.15cm}
\item Frenado d(v)=0.05v^2 a 60 km/h:
\begin{enumerate}[label=\textbf{\Alph*)}]
\begin{multicols}{2}
  \item 180 m
  \item 2
  \item 10
  \item 15
  \item 25
  \item 0
\end{multicols}
\end{enumerate}
\vspace{0.15cm}
\item Área rectángulo perímetro 20 función base x:
\begin{enumerate}[label=\textbf{\Alph*)}]
\begin{multicols}{2}
  \item A(x)=x(10-x)
  \item 2
  \item 10
  \item 15
  \item 25
  \item 0
\end{multicols}
\end{enumerate}
\vspace{0.15cm}
\item Área máxima rectángulo anterior:
\begin{enumerate}[label=\textbf{\Alph*)}]
\begin{multicols}{2}
  \item 25
  \item 2
  \item 10
  \item 15
  \item 25
  \item 0
\end{multicols}
\end{enumerate}
\vspace{0.15cm}
\item Ingreso I(x)=50x-0.5x^2. Ingreso máximo:
\begin{enumerate}[label=\textbf{\Alph*)}]
\begin{multicols}{2}
  \item $1250
  \item $20
  \item 10
  \item 15
  \item 25
  \item 0
\end{multicols}
\end{enumerate}
\vspace{0.15cm}
\item x para maximizar ingreso anterior:
\begin{enumerate}[label=\textbf{\Alph*)}]
\begin{multicols}{2}
  \item 50
  \item 2
  \item 10
  \item 15
  \item 25
  \item 0
\end{multicols}
\end{enumerate}
\vspace{0.15cm}
\item Pelota h(t)=-4t^2+16t+1. Altura máxima:
\begin{enumerate}[label=\textbf{\Alph*)}]
\begin{multicols}{2}
  \item 17 m
  \item 2
  \item 10
  \item 15
  \item 25
  \item 0
\end{multicols}
\end{enumerate}
\vspace{0.15cm}
\item Tiempo para altura máxima pelota anterior:
\begin{enumerate}[label=\textbf{\Alph*)}]
\begin{multicols}{2}
  \item 2 s
  \item 2
  \item 10
  \item 15
  \item 25
  \item 0
\end{multicols}
\end{enumerate}
\vspace{0.15cm}
\item Demanda p(x)=100-2x, Ingreso máximo:
\begin{enumerate}[label=\textbf{\Alph*)}]
\begin{multicols}{2}
  \item $1250
  \item $20
  \item 10
  \item 15
  \item 25
  \item 0
\end{multicols}
\end{enumerate}
\vspace{0.15cm}
\end{enumerate}\n
\subsubsection*{SECCIÓN VI: FUNCIONES ESPECIALES (271–300)}

\begin{enumerate}[label=\textbf{\arabic*.}]
\setcounter{enumi}{270}
\item La función valor absoluto es:
\begin{enumerate}[label=\textbf{\Alph*)}]
\begin{multicols}{2}
  \item f(x)=|x|
  \item x > 0
  \item Línea recta
  \item Constante
  \item Todos los reales
  \item 0
\end{multicols}
\end{enumerate}
\vspace{0.15cm}
\item La gráfica de f(x)=|x| tiene forma de:
\begin{enumerate}[label=\textbf{\Alph*)}]
\begin{multicols}{2}
  \item V
  \item U
  \item Línea recta
  \item Constante
  \item Todos los reales
  \item 0
\end{multicols}
\end{enumerate}
\vspace{0.15cm}
\item El dominio de f(x)=|x| es:
\begin{enumerate}[label=\textbf{\Alph*)}]
\begin{multicols}{2}
  \item Todos los reales
  \item x > 0
  \item Línea recta
  \item Constante
  \item Todos los reales
  \item 0
\end{multicols}
\end{enumerate}
\vspace{0.15cm}
\item El rango de f(x)=|x| es:
\begin{enumerate}[label=\textbf{\Alph*)}]
\begin{multicols}{2}
  \item y \geq 0
  \item x > 0
  \item Línea recta
  \item Constante
  \item Todos los reales
  \item 0
\end{multicols}
\end{enumerate}
\vspace{0.15cm}
\item La función f(x)=|x| es:
\begin{enumerate}[label=\textbf{\Alph*)}]
\begin{multicols}{2}
  \item Par
  \item x > 0
  \item Impar
  \item Constante
  \item Todos los reales
  \item 0
\end{multicols}
\end{enumerate}
\vspace{0.15cm}
\item La función f(x)=|x-2| tiene su vértice en:
\begin{enumerate}[label=\textbf{\Alph*)}]
\begin{multicols}{2}
  \item (2,0)
  \item x > 0
  \item Línea recta
  \item Constante
  \item Todos los reales
  \item 0
\end{multicols}
\end{enumerate}
\vspace{0.15cm}
\item El dominio de f(x)=\sqrt{x-1} es:
\begin{enumerate}[label=\textbf{\Alph*)}]
\begin{multicols}{2}
  \item x \geq 1
  \item x > 0
  \item Línea recta
  \item Constante
  \item Todos los reales
  \item 0
\end{multicols}
\end{enumerate}
\vspace{0.15cm}
\item El rango de f(x)=\sqrt{x} es:
\begin{enumerate}[label=\textbf{\Alph*)}]
\begin{multicols}{2}
  \item y \geq 0
  \item x > 0
  \item Línea recta
  \item Constante
  \item Todos los reales
  \item 0
\end{multicols}
\end{enumerate}
\vspace{0.15cm}
\item La función f(x)=\sqrt{x} es:
\begin{enumerate}[label=\textbf{\Alph*)}]
\begin{multicols}{2}
  \item Creciente
  \item x > 0
  \item Línea recta
  \item Constante
  \item Todos los reales
  \item 0
\end{multicols}
\end{enumerate}
\vspace{0.15cm}
\item La función f(x)=1/x tiene asíntota vertical en:
\begin{enumerate}[label=\textbf{\Alph*)}]
\begin{multicols}{2}
  \item x = 0
  \item x > 0
  \item Línea recta
  \item Constante
  \item Todos los reales
  \item 0
\end{multicols}
\end{enumerate}
\vspace{0.15cm}
\item La función f(x)=1/x tiene asíntota horizontal en:
\begin{enumerate}[label=\textbf{\Alph*)}]
\begin{multicols}{2}
  \item y = 0
  \item x > 0
  \item Línea recta
  \item Constante
  \item Todos los reales
  \item 0
\end{multicols}
\end{enumerate}
\vspace{0.15cm}
\item El dominio de f(x)=1/(x-3) es:
\begin{enumerate}[label=\textbf{\Alph*)}]
\begin{multicols}{2}
  \item Todos los reales excepto x = 3
  \item x > 0
  \item Línea recta
  \item Constante
  \item Todos los reales
  \item 0
\end{multicols}
\end{enumerate}
\vspace{0.15cm}
\item El rango de f(x)=1/x es:
\begin{enumerate}[label=\textbf{\Alph*)}]
\begin{multicols}{2}
  \item Todos los reales excepto y = 0
  \item x > 0
  \item Línea recta
  \item Constante
  \item Todos los reales
  \item 0
\end{multicols}
\end{enumerate}
\vspace{0.15cm}
\item La función f(x)=e^x es:
\begin{enumerate}[label=\textbf{\Alph*)}]
\begin{multicols}{2}
  \item Creciente
  \item x > 0
  \item Línea recta
  \item Constante
  \item Todos los reales
  \item 0
\end{multicols}
\end{enumerate}
\vspace{0.15cm}
\item El dominio de f(x)=e^x es:
\begin{enumerate}[label=\textbf{\Alph*)}]
\begin{multicols}{2}
  \item Todos los reales
  \item x > 0
  \item Línea recta
  \item Constante
  \item Todos los reales
  \item 0
\end{multicols}
\end{enumerate}
\vspace{0.15cm}
\item El rango de f(x)=e^x es:
\begin{enumerate}[label=\textbf{\Alph*)}]
\begin{multicols}{2}
  \item y > 0
  \item x > 0
  \item Línea recta
  \item Constante
  \item Todos los reales
  \item 0
\end{multicols}
\end{enumerate}
\vspace{0.15cm}
\item La función f(x)=ln(x) es la inversa de:
\begin{enumerate}[label=\textbf{\Alph*)}]
\begin{multicols}{2}
  \item e^x
  \item x > 0
  \item Línea recta
  \item Constante
  \item Todos los reales
  \item 0
\end{multicols}
\end{enumerate}
\vspace{0.15cm}
\item El dominio de f(x)=ln(x) es:
\begin{enumerate}[label=\textbf{\Alph*)}]
\begin{multicols}{2}
  \item x > 0
  \item x > 0
  \item Línea recta
  \item Constante
  \item Todos los reales
  \item 0
\end{multicols}
\end{enumerate}
\vspace{0.15cm}
\item El rango de f(x)=ln(x) es:
\begin{enumerate}[label=\textbf{\Alph*)}]
\begin{multicols}{2}
  \item Todos los reales
  \item x > 0
  \item Línea recta
  \item Constante
  \item Todos los reales
  \item 0
\end{multicols}
\end{enumerate}
\vspace{0.15cm}
\item La función f(x)=x^3 es:
\begin{enumerate}[label=\textbf{\Alph*)}]
\begin{multicols}{2}
  \item Impar
  \item Par
  \item Línea recta
  \item Constante
  \item Todos los reales
  \item 0
\end{multicols}
\end{enumerate}
\vspace{0.15cm}
\item La función f(x)=x^4 es:
\begin{enumerate}[label=\textbf{\Alph*)}]
\begin{multicols}{2}
  \item Par
  \item x > 0
  \item Impar
  \item Constante
  \item Todos los reales
  \item 0
\end{multicols}
\end{enumerate}
\vspace{0.15cm}
\item La función f(x)=1/x es:
\begin{enumerate}[label=\textbf{\Alph*)}]
\begin{multicols}{2}
  \item Impar
  \item Par
  \item Línea recta
  \item Constante
  \item Todos los reales
  \item 0
\end{multicols}
\end{enumerate}
\vspace{0.15cm}
\item La función f(x)=1/x^2 es:
\begin{enumerate}[label=\textbf{\Alph*)}]
\begin{multicols}{2}
  \item Par
  \item x > 0
  \item Impar
  \item Constante
  \item Todos los reales
  \item 0
\end{multicols}
\end{enumerate}
\vspace{0.15cm}
\item La función f(x)=x^3+x es:
\begin{enumerate}[label=\textbf{\Alph*)}]
\begin{multicols}{2}
  \item Impar
  \item Par
  \item Línea recta
  \item Constante
  \item Todos los reales
  \item 0
\end{multicols}
\end{enumerate}
\vspace{0.15cm}
\item La función f(x)=x^2+x es:
\begin{enumerate}[label=\textbf{\Alph*)}]
\begin{multicols}{2}
  \item Ni par ni impar
  \item x > 0
  \item Línea recta
  \item Constante
  \item Todos los reales
  \item 0
\end{multicols}
\end{enumerate}
\vspace{0.15cm}
\item La función f(x)=2 es:
\begin{enumerate}[label=\textbf{\Alph*)}]
\begin{multicols}{2}
  \item Constante
  \item x > 0
  \item Línea recta
  \item Constante
  \item Todos los reales
  \item 0
\end{multicols}
\end{enumerate}
\vspace{0.15cm}
\item La función f(x)=x es:
\begin{enumerate}[label=\textbf{\Alph*)}]
\begin{multicols}{2}
  \item Impar
  \item Par
  \item Línea recta
  \item Constante
  \item Todos los reales
  \item 0
\end{multicols}
\end{enumerate}
\vspace{0.15cm}
\item La función f(x)=|x|+1 tiene su mínimo en:
\begin{enumerate}[label=\textbf{\Alph*)}]
\begin{multicols}{2}
  \item (0,1)
  \item x > 0
  \item Línea recta
  \item Constante
  \item Todos los reales
  \item 0
\end{multicols}
\end{enumerate}
\vspace{0.15cm}
\item La función f(x)=|x-3|+2 tiene su mínimo en:
\begin{enumerate}[label=\textbf{\Alph*)}]
\begin{multicols}{2}
  \item (3,2)
  \item x > 0
  \item Línea recta
  \item Constante
  \item Todos los reales
  \item 0
\end{multicols}
\end{enumerate}
\vspace{0.15cm}
\item La función f(x)=-|x| tiene su máximo en:
\begin{enumerate}[label=\textbf{\Alph*)}]
\begin{multicols}{2}
  \item (0,0)
  \item x > 0
  \item Línea recta
  \item Constante
  \item Todos los reales
  \item 0
\end{multicols}
\end{enumerate}
\vspace{0.15cm}
\end{enumerate}\n
\vspace{0.4cm}
\begin{tcolorbox}[colback=gray!5!white,colframe=gray!60!black,title=\textbf{SOLUCIONARIO -- 1.7 Funciones}]
\begin{multicols}{4}
\begin{enumerate}[label=\textbf{\arabic*.}]
  \item \textbf{A}
  \item \textbf{B}
  \item \textbf{B}
  \item \textbf{C}
  \item \textbf{B}
  \item \textbf{A}
  \item \textbf{A}
  \item \textbf{A}
  \item \textbf{A}
  \item \textbf{A}
  \item \textbf{A}
  \item \textbf{B}
  \item \textbf{B}
  \item \textbf{B}
  \item \textbf{B}
  \item \textbf{A}
  \item \textbf{A}
  \item \textbf{B}
  \item \textbf{B}
  \item \textbf{A}
  \item \textbf{A}
  \item \textbf{B}
  \item \textbf{B}
  \item \textbf{A}
  \item \textbf{A}
  \item \textbf{A}
  \item \textbf{A}
  \item \textbf{C}
  \item \textbf{B}
  \item \textbf{A}
  \item \textbf{B}
  \item \textbf{B}
  \item \textbf{A}
  \item \textbf{A}
  \item \textbf{B}
  \item \textbf{A}
  \item \textbf{B}
  \item \textbf{C}
  \item \textbf{A}
  \item \textbf{A}
  \item \textbf{A}
  \item \textbf{A}
  \item \textbf{B}
  \item \textbf{B}
  \item \textbf{B}
  \item \textbf{A}
  \item \textbf{A}
  \item \textbf{A}
  \item \textbf{B}
  \item \textbf{A}
  \item \textbf{D}
  \item \textbf{A}
  \item \textbf{A}
  \item \textbf{A}
  \item \textbf{A}
  \item \textbf{A}
  \item \textbf{A}
  \item \textbf{A}
  \item \textbf{A}
  \item \textbf{A}
  \item \textbf{A}
  \item \textbf{A}
  \item \textbf{A}
  \item \textbf{A}
  \item \textbf{B}
  \item \textbf{C}
  \item \textbf{A}
  \item \textbf{B}
  \item \textbf{C}
  \item \textbf{A}
  \item \textbf{A}
  \item \textbf{A}
  \item \textbf{A}
  \item \textbf{B}
  \item \textbf{A}
  \item \textbf{B}
  \item \textbf{A}
  \item \textbf{A}
  \item \textbf{A}
  \item \textbf{A}
  \item \textbf{A}
  \item \textbf{A}
  \item \textbf{B}
  \item \textbf{A}
  \item \textbf{A}
  \item \textbf{A}
  \item \textbf{A}
  \item \textbf{A}
  \item \textbf{A}
  \item \textbf{A}
  \item \textbf{A}
  \item \textbf{A}
  \item \textbf{A}
  \item \textbf{A}
  \item \textbf{A}
  \item \textbf{B}
  \item \textbf{B}
  \item \textbf{A}
  \item \textbf{A}
  \item \textbf{A}
  \item \textbf{B}
  \item \textbf{B}
  \item \textbf{A}
  \item \textbf{B}
  \item \textbf{A}
  \item \textbf{A}
  \item \textbf{A}
  \item \textbf{A}
  \item \textbf{A}
  \item \textbf{A}
  \item \textbf{A}
  \item \textbf{A}
  \item \textbf{A}
  \item \textbf{A}
  \item \textbf{A}
  \item \textbf{A}
  \item \textbf{A}
  \item \textbf{A}
  \item \textbf{A}
  \item \textbf{A}
  \item \textbf{A}
  \item \textbf{A}
  \item \textbf{A}
  \item \textbf{A}
  \item \textbf{A}
  \item \textbf{A}
  \item \textbf{A}
  \item \textbf{A}
  \item \textbf{A}
  \item \textbf{A}
  \item \textbf{A}
  \item \textbf{A}
  \item \textbf{A}
  \item \textbf{A}
  \item \textbf{A}
  \item \textbf{A}
  \item \textbf{B}
  \item \textbf{C}
  \item \textbf{A}
  \item \textbf{A}
  \item \textbf{A}
  \item \textbf{A}
  \item \textbf{A}
  \item \textbf{A}
  \item \textbf{A}
  \item \textbf{A}
  \item \textbf{A}
  \item \textbf{A}
  \item \textbf{A}
  \item \textbf{A}
  \item \textbf{A}
  \item \textbf{A}
  \item \textbf{A}
  \item \textbf{A}
  \item \textbf{A}
  \item \textbf{A}
  \item \textbf{A}
  \item \textbf{A}
  \item \textbf{A}
  \item \textbf{A}
  \item \textbf{A}
  \item \textbf{A}
  \item \textbf{A}
  \item \textbf{A}
  \item \textbf{A}
  \item \textbf{A}
  \item \textbf{A}
  \item \textbf{A}
  \item \textbf{A}
  \item \textbf{A}
  \item \textbf{A}
  \item \textbf{A}
  \item \textbf{A}
  \item \textbf{A}
  \item \textbf{A}
  \item \textbf{A}
  \item \textbf{A}
  \item \textbf{A}
  \item \textbf{A}
  \item \textbf{A}
  \item \textbf{A}
  \item \textbf{A}
  \item \textbf{A}
  \item \textbf{A}
  \item \textbf{A}
  \item \textbf{A}
  \item \textbf{A}
  \item \textbf{A}
  \item \textbf{A}
  \item \textbf{A}
  \item \textbf{A}
  \item \textbf{A}
  \item \textbf{A}
  \item \textbf{A}
  \item \textbf{A}
  \item \textbf{A}
  \item \textbf{A}
  \item \textbf{A}
  \item \textbf{A}
  \item \textbf{A}
  \item \textbf{A}
  \item \textbf{A}
  \item \textbf{A}
  \item \textbf{B}
  \item \textbf{A}
  \item \textbf{A}
  \item \textbf{A}
  \item \textbf{A}
  \item \textbf{A}
  \item \textbf{A}
  \item \textbf{A}
  \item \textbf{A}
  \item \textbf{A}
  \item \textbf{A}
  \item \textbf{A}
  \item \textbf{A}
  \item \textbf{A}
  \item \textbf{A}
  \item \textbf{A}
  \item \textbf{A}
  \item \textbf{B}
  \item \textbf{A}
  \item \textbf{A}
  \item \textbf{A}
  \item \textbf{A}
  \item \textbf{A}
  \item \textbf{A}
  \item \textbf{A}
  \item \textbf{A}
  \item \textbf{A}
  \item \textbf{A}
  \item \textbf{A}
  \item \textbf{A}
  \item \textbf{A}
  \item \textbf{A}
  \item \textbf{A}
  \item \textbf{A}
  \item \textbf{A}
  \item \textbf{A}
  \item \textbf{A}
  \item \textbf{A}
  \item \textbf{A}
  \item \textbf{A}
  \item \textbf{A}
  \item \textbf{A}
  \item \textbf{A}
  \item \textbf{A}
  \item \textbf{A}
  \item \textbf{A}
  \item \textbf{A}
  \item \textbf{A}
  \item \textbf{A}
  \item \textbf{A}
  \item \textbf{A}
  \item \textbf{A}
  \item \textbf{A}
  \item \textbf{A}
  \item \textbf{A}
  \item \textbf{A}
  \item \textbf{A}
  \item \textbf{A}
  \item \textbf{A}
  \item \textbf{A}
  \item \textbf{A}
  \item \textbf{A}
  \item \textbf{A}
  \item \textbf{A}
  \item \textbf{A}
  \item \textbf{A}
  \item \textbf{A}
  \item \textbf{A}
  \item \textbf{A}
  \item \textbf{A}
  \item \textbf{A}
  \item \textbf{A}
  \item \textbf{A}
  \item \textbf{A}
  \item \textbf{A}
  \item \textbf{A}
  \item \textbf{A}
  \item \textbf{A}
  \item \textbf{A}
  \item \textbf{A}
  \item \textbf{A}
  \item \textbf{A}
  \item \textbf{A}
  \item \textbf{A}
  \item \textbf{A}
  \item \textbf{A}
  \item \textbf{A}
  \item \textbf{A}
  \item \textbf{A}
  \item \textbf{A}
  \item \textbf{A}
  \item \textbf{A}
  \item \textbf{A}
  \item \textbf{A}
  \item \textbf{A}
  \item \textbf{A}
  \item \textbf{A}
\end{enumerate}
\end{multicols}
\end{tcolorbox}
\newpage\n
% ==============================================================================
% MATERIA 2: LENGUAJE
% ==============================================================================
\section{Lenguaje (Módulos 2.1 al 2.5)}

\subsection{2.1 ORTOGRAFÍA Y GRAMÁTICA}

\subsubsection*{SECCIÓN I: ACENTUACIÓN – PALABRAS AGUDAS (1–40)}

\begin{enumerate}[label=\textbf{\arabic*.}]
\setcounter{enumi}{0}
\item ¿Cuál de las siguientes palabras es aguda?
\begin{enumerate}[label=\textbf{\Alph*)}]
\begin{multicols}{2}
  \item árbol
  \item cantar
  \item lápiz
  \item música
  \item fácil
  \item césped
\end{multicols}
\end{enumerate}
\vspace{0.15cm}
\item ¿Cuál de las siguientes palabras es aguda?
\begin{enumerate}[label=\textbf{\Alph*)}]
\begin{multicols}{2}
  \item mesa
  \item reloj
  \item silla
  \item libro
  \item casa
  \item pared
\end{multicols}
\end{enumerate}
\vspace{0.15cm}
\item ¿Cuál de las siguientes palabras es aguda?
\begin{enumerate}[label=\textbf{\Alph*)}]
\begin{multicols}{2}
  \item corazón
  \item ventana
  \item teléfono
  \item pájaro
  \item cárcel
  \item lunes
\end{multicols}
\end{enumerate}
\vspace{0.15cm}
\item ¿Cuál de las siguientes palabras es aguda?
\begin{enumerate}[label=\textbf{\Alph*)}]
\begin{multicols}{2}
  \item papel
  \item cuaderno
  \item lápiz
  \item árbol
  \item azúcar
  \item mármol
\end{multicols}
\end{enumerate}
\vspace{0.15cm}
\item ¿Cuál de las siguientes palabras es aguda?
\begin{enumerate}[label=\textbf{\Alph*)}]
\begin{multicols}{2}
  \item café
  \item mesa
  \item silla
  \item libro
  \item casa
  \item ventana
\end{multicols}
\end{enumerate}
\vspace{0.15cm}
\item ¿Cuál de las siguientes palabras es aguda?
\begin{enumerate}[label=\textbf{\Alph*)}]
\begin{multicols}{2}
  \item sofá
  \item cama
  \item puerta
  \item ventana
  \item techo
  \item piso
\end{multicols}
\end{enumerate}
\vspace{0.15cm}
\item ¿Cuál de las siguientes palabras es aguda?
\begin{enumerate}[label=\textbf{\Alph*)}]
\begin{multicols}{2}
  \item mamá
  \item papá
  \item casa
  \item mesa
  \item silla
  \item libro
\end{multicols}
\end{enumerate}
\vspace{0.15cm}
\item ¿Cuál de las siguientes palabras es aguda?
\begin{enumerate}[label=\textbf{\Alph*)}]
\begin{multicols}{2}
  \item compás
  \item lápiz
  \item árbol
  \item cárcel
  \item fácil
  \item césped
\end{multicols}
\end{enumerate}
\vspace{0.15cm}
\item ¿Cuál de las siguientes palabras es aguda?
\begin{enumerate}[label=\textbf{\Alph*)}]
\begin{multicols}{2}
  \item jamás
  \item mesa
  \item silla
  \item libro
  \item casa
  \item ventana
\end{multicols}
\end{enumerate}
\vspace{0.15cm}
\item ¿Cuál de las siguientes palabras es aguda?
\begin{enumerate}[label=\textbf{\Alph*)}]
\begin{multicols}{2}
  \item atrás
  \item delante
  \item cerca
  \item lejos
  \item arriba
  \item abajo
\end{multicols}
\end{enumerate}
\vspace{0.15cm}
\item ¿Cuál de las siguientes palabras es aguda?
\begin{enumerate}[label=\textbf{\Alph*)}]
\begin{multicols}{2}
  \item después
  \item antes
  \item ahora
  \item luego
  \item siempre
  \item nunca
\end{multicols}
\end{enumerate}
\vspace{0.15cm}
\item ¿Cuál de las siguientes palabras es aguda?
\begin{enumerate}[label=\textbf{\Alph*)}]
\begin{multicols}{2}
  \item también
  \item tampoco
  \item quizás
  \item acaso
  \item tal vez
  \item quizá
\end{multicols}
\end{enumerate}
\vspace{0.15cm}
\item ¿Cuál de las siguientes palabras es aguda?
\begin{enumerate}[label=\textbf{\Alph*)}]
\begin{multicols}{2}
  \item canción
  \item música
  \item lápiz
  \item árbol
  \item fácil
  \item césped
\end{multicols}
\end{enumerate}
\vspace{0.15cm}
\item ¿Cuál de las siguientes palabras es aguda?
\begin{enumerate}[label=\textbf{\Alph*)}]
\begin{multicols}{2}
  \item razón
  \item corazón
  \item mesa
  \item silla
  \item libro
  \item casa
\end{multicols}
\end{enumerate}
\vspace{0.15cm}
\item ¿Cuál de las siguientes palabras es aguda?
\begin{enumerate}[label=\textbf{\Alph*)}]
\begin{multicols}{2}
  \item balón
  \item pelota
  \item juego
  \item cancha
  \item equipo
  \item gol
\end{multicols}
\end{enumerate}
\vspace{0.15cm}
\item ¿Cuál de las siguientes palabras es aguda?
\begin{enumerate}[label=\textbf{\Alph*)}]
\begin{multicols}{2}
  \item camión
  \item carro
  \item auto
  \item moto
  \item bus
  \item tren
\end{multicols}
\end{enumerate}
\vspace{0.15cm}
\item ¿Cuál de las siguientes palabras es aguda?
\begin{enumerate}[label=\textbf{\Alph*)}]
\begin{multicols}{2}
  \item avión
  \item barco
  \item lancha
  \item velero
  \item yate
  \item bote
\end{multicols}
\end{enumerate}
\vspace{0.15cm}
\item ¿Cuál de las siguientes palabras es aguda?
\begin{enumerate}[label=\textbf{\Alph*)}]
\begin{multicols}{2}
  \item ratón
  \item gato
  \item perro
  \item ratón
  \item conejo
  \item hámster
\end{multicols}
\end{enumerate}
\vspace{0.15cm}
\item ¿Cuál de las siguientes palabras es aguda?
\begin{enumerate}[label=\textbf{\Alph*)}]
\begin{multicols}{2}
  \item ratón
  \item ratón
  \item ratón
  \item ratón
  \item ratón
  \item ratón
\end{multicols}
\end{enumerate}
\vspace{0.15cm}
\item ¿Cuál de las siguientes palabras es aguda?
\begin{enumerate}[label=\textbf{\Alph*)}]
\begin{multicols}{2}
  \item jabón
  \item agua
  \item espuma
  \item baño
  \item ducha
  \item tina
\end{multicols}
\end{enumerate}
\vspace{0.15cm}
\item ¿Cuál de las siguientes palabras es aguda?
\begin{enumerate}[label=\textbf{\Alph*)}]
\begin{multicols}{2}
  \item botón
  \item camisa
  \item pantalón
  \item zapato
  \item media
  \item blusa
\end{multicols}
\end{enumerate}
\vspace{0.15cm}
\item ¿Cuál de las siguientes palabras es aguda?
\begin{enumerate}[label=\textbf{\Alph*)}]
\begin{multicols}{2}
  \item pantalón
  \item camisa
  \item blusa
  \item falda
  \item vestido
  \item traje
\end{multicols}
\end{enumerate}
\vspace{0.15cm}
\item ¿Cuál de las siguientes palabras es aguda?
\begin{enumerate}[label=\textbf{\Alph*)}]
\begin{multicols}{2}
  \item sillón
  \item silla
  \item mesa
  \item cama
  \item sofá
  \item banco
\end{multicols}
\end{enumerate}
\vspace{0.15cm}
\item ¿Cuál de las siguientes palabras es aguda?
\begin{enumerate}[label=\textbf{\Alph*)}]
\begin{multicols}{2}
  \item calor
  \item frío
  \item templado
  \item fresco
  \item helado
  \item tibio
\end{multicols}
\end{enumerate}
\vspace{0.15cm}
\item ¿Cuál de las siguientes palabras es aguda?
\begin{enumerate}[label=\textbf{\Alph*)}]
\begin{multicols}{2}
  \item dolor
  \item alegría
  \item tristeza
  \item pena
  \item gozo
  \item risa
\end{multicols}
\end{enumerate}
\vspace{0.15cm}
\item ¿Cuál de las siguientes palabras es aguda?
\begin{enumerate}[label=\textbf{\Alph*)}]
\begin{multicols}{2}
  \item amor
  \item odio
  \item cariño
  \item afecto
  \item pasión
  \item ternura
\end{multicols}
\end{enumerate}
\vspace{0.15cm}
\item ¿Cuál de las siguientes palabras es aguda?
\begin{enumerate}[label=\textbf{\Alph*)}]
\begin{multicols}{2}
  \item sabor
  \item dulce
  \item salado
  \item ácido
  \item amargo
  \item picante
\end{multicols}
\end{enumerate}
\vspace{0.15cm}
\item ¿Cuál de las siguientes palabras es aguda?
\begin{enumerate}[label=\textbf{\Alph*)}]
\begin{multicols}{2}
  \item olor
  \item aroma
  \item fragancia
  \item perfume
  \item esencia
  \item hedor
\end{multicols}
\end{enumerate}
\vspace{0.15cm}
\item ¿Cuál de las siguientes palabras es aguda?
\begin{enumerate}[label=\textbf{\Alph*)}]
\begin{multicols}{2}
  \item color
  \item tono
  \item matiz
  \item sombra
  \item luz
  \item brillo
\end{multicols}
\end{enumerate}
\vspace{0.15cm}
\item ¿Cuál de las siguientes palabras es aguda?
\begin{enumerate}[label=\textbf{\Alph*)}]
\begin{multicols}{2}
  \item valor
  \item precio
  \item costo
  \item tasa
  \item tarifa
  \item cuota
\end{multicols}
\end{enumerate}
\vspace{0.15cm}
\item ¿Cuál de las siguientes palabras es aguda?
\begin{enumerate}[label=\textbf{\Alph*)}]
\begin{multicols}{2}
  \item temor
  \item miedo
  \item pánico
  \item susto
  \item terror
  \item fobia
\end{multicols}
\end{enumerate}
\vspace{0.15cm}
\item ¿Cuál de las siguientes palabras es aguda?
\begin{enumerate}[label=\textbf{\Alph*)}]
\begin{multicols}{2}
  \item rumor
  \item noticia
  \item chisme
  \item comentario
  \item opinión
  \item crítica
\end{multicols}
\end{enumerate}
\vspace{0.15cm}
\item ¿Cuál de las siguientes palabras es aguda?
\begin{enumerate}[label=\textbf{\Alph*)}]
\begin{multicols}{2}
  \item humor
  \item gracia
  \item chiste
  \item broma
  \item risa
  \item sonrisa
\end{multicols}
\end{enumerate}
\vspace{0.15cm}
\item ¿Cuál de las siguientes palabras es aguda?
\begin{enumerate}[label=\textbf{\Alph*)}]
\begin{multicols}{2}
  \item rigor
  \item disciplina
  \item orden
  \item control
  \item norma
  \item regla
\end{multicols}
\end{enumerate}
\vspace{0.15cm}
\item ¿Cuál de las siguientes palabras es aguda?
\begin{enumerate}[label=\textbf{\Alph*)}]
\begin{multicols}{2}
  \item favor
  \item ayuda
  \item apoyo
  \item servicio
  \item beneficio
  \item ventaja
\end{multicols}
\end{enumerate}
\vspace{0.15cm}
\item ¿Cuál de las siguientes palabras es aguda?
\begin{enumerate}[label=\textbf{\Alph*)}]
\begin{multicols}{2}
  \item error
  \item acierto
  \item fallo
  \item equivocación
  \item descuido
  \item falta
\end{multicols}
\end{enumerate}
\vspace{0.15cm}
\item ¿Cuál de las siguientes palabras es aguda?
\begin{enumerate}[label=\textbf{\Alph*)}]
\begin{multicols}{2}
  \item temblor
  \item temblor
  \item temblor
  \item temblor
  \item temblor
  \item temblor
\end{multicols}
\end{enumerate}
\vspace{0.15cm}
\item ¿Cuál de las siguientes palabras es aguda?
\begin{enumerate}[label=\textbf{\Alph*)}]
\begin{multicols}{2}
  \item sudor
  \item transpiración
  \item humedad
  \item calor
  \item frío
  \item frescura
\end{multicols}
\end{enumerate}
\vspace{0.15cm}
\item ¿Cuál de las siguientes palabras es aguda?
\begin{enumerate}[label=\textbf{\Alph*)}]
\begin{multicols}{2}
  \item rencor
  \item perdón
  \item olvido
  \item paz
  \item calma
  \item tranquilidad
\end{multicols}
\end{enumerate}
\vspace{0.15cm}
\item ¿Cuál de las siguientes palabras es aguda?
\begin{enumerate}[label=\textbf{\Alph*)}]
\begin{multicols}{2}
  \item candor
  \item pureza
  \item inocencia
  \item bondad
  \item virtud
  \item honradez ---
\end{multicols}
\end{enumerate}
\vspace{0.15cm}
\end{enumerate}\n
\subsubsection*{SECCIÓN II: ACENTUACIÓN – PALABRAS GRAVES O LLANAS (41–80)}

\begin{enumerate}[label=\textbf{\arabic*.}]
\setcounter{enumi}{40}
\item ¿Cuál de las siguientes palabras es grave?
\begin{enumerate}[label=\textbf{\Alph*)}]
\begin{multicols}{2}
  \item árbol
  \item cantar
  \item corazón
  \item papel
  \item café
  \item sofá
\end{multicols}
\end{enumerate}
\vspace{0.15cm}
\item ¿Cuál de las siguientes palabras es grave?
\begin{enumerate}[label=\textbf{\Alph*)}]
\begin{multicols}{2}
  \item lápiz
  \item reloj
  \item jamás
  \item atrás
  \item después
  \item también
\end{multicols}
\end{enumerate}
\vspace{0.15cm}
\item ¿Cuál de las siguientes palabras es grave?
\begin{enumerate}[label=\textbf{\Alph*)}]
\begin{multicols}{2}
  \item mesa
  \item mamá
  \item papá
  \item canción
  \item razón
  \item balón
\end{multicols}
\end{enumerate}
\vspace{0.15cm}
\item ¿Cuál de las siguientes palabras es grave?
\begin{enumerate}[label=\textbf{\Alph*)}]
\begin{multicols}{2}
  \item silla
  \item café
  \item sofá
  \item compás
  \item ratón
  \item jabón
\end{multicols}
\end{enumerate}
\vspace{0.15cm}
\item ¿Cuál de las siguientes palabras es grave?
\begin{enumerate}[label=\textbf{\Alph*)}]
\begin{multicols}{2}
  \item libro
  \item corazón
  \item papel
  \item reloj
  \item jamás
  \item atrás
\end{multicols}
\end{enumerate}
\vspace{0.15cm}
\item ¿Cuál de las siguientes palabras es grave?
\begin{enumerate}[label=\textbf{\Alph*)}]
\begin{multicols}{2}
  \item casa
  \item mamá
  \item papá
  \item café
  \item sofá
  \item compás
\end{multicols}
\end{enumerate}
\vspace{0.15cm}
\item ¿Cuál de las siguientes palabras es grave?
\begin{enumerate}[label=\textbf{\Alph*)}]
\begin{multicols}{2}
  \item ventana
  \item canción
  \item razón
  \item balón
  \item camión
  \item avión
\end{multicols}
\end{enumerate}
\vspace{0.15cm}
\item ¿Cuál de las siguientes palabras es grave?
\begin{enumerate}[label=\textbf{\Alph*)}]
\begin{multicols}{2}
  \item teléfono
  \item cantar
  \item corazón
  \item papel
  \item reloj
  \item jamás
\end{multicols}
\end{enumerate}
\vspace{0.15cm}
\item ¿Cuál de las siguientes palabras es grave?
\begin{enumerate}[label=\textbf{\Alph*)}]
\begin{multicols}{2}
  \item pájaro
  \item mamá
  \item papá
  \item café
  \item sofá
  \item compás
\end{multicols}
\end{enumerate}
\vspace{0.15cm}
\item ¿Cuál de las siguientes palabras es grave?
\begin{enumerate}[label=\textbf{\Alph*)}]
\begin{multicols}{2}
  \item cárcel
  \item canción
  \item razón
  \item balón
  \item camión
  \item avión
\end{multicols}
\end{enumerate}
\vspace{0.15cm}
\item ¿Cuál de las siguientes palabras es grave?
\begin{enumerate}[label=\textbf{\Alph*)}]
\begin{multicols}{2}
  \item lunes
  \item cantar
  \item corazón
  \item papel
  \item reloj
  \item jamás
\end{multicols}
\end{enumerate}
\vspace{0.15cm}
\item ¿Cuál de las siguientes palabras es grave?
\begin{enumerate}[label=\textbf{\Alph*)}]
\begin{multicols}{2}
  \item azúcar
  \item mamá
  \item papá
  \item café
  \item sofá
  \item compás
\end{multicols}
\end{enumerate}
\vspace{0.15cm}
\item ¿Cuál de las siguientes palabras es grave?
\begin{enumerate}[label=\textbf{\Alph*)}]
\begin{multicols}{2}
  \item mármol
  \item canción
  \item razón
  \item balón
  \item camión
  \item avión
\end{multicols}
\end{enumerate}
\vspace{0.15cm}
\item ¿Cuál de las siguientes palabras es grave?
\begin{enumerate}[label=\textbf{\Alph*)}]
\begin{multicols}{2}
  \item césped
  \item cantar
  \item corazón
  \item papel
  \item reloj
  \item jamás
\end{multicols}
\end{enumerate}
\vspace{0.15cm}
\item ¿Cuál de las siguientes palabras es grave?
\begin{enumerate}[label=\textbf{\Alph*)}]
\begin{multicols}{2}
  \item fácil
  \item mamá
  \item papá
  \item café
  \item sofá
  \item compás
\end{multicols}
\end{enumerate}
\vspace{0.15cm}
\item ¿Cuál de las siguientes palabras es grave?
\begin{enumerate}[label=\textbf{\Alph*)}]
\begin{multicols}{2}
  \item útil
  \item canción
  \item razón
  \item balón
  \item camión
  \item avión
\end{multicols}
\end{enumerate}
\vspace{0.15cm}
\item ¿Cuál de las siguientes palabras es grave?
\begin{enumerate}[label=\textbf{\Alph*)}]
\begin{multicols}{2}
  \item débil
  \item cantar
  \item corazón
  \item papel
  \item reloj
  \item jamás
\end{multicols}
\end{enumerate}
\vspace{0.15cm}
\item ¿Cuál de las siguientes palabras es grave?
\begin{enumerate}[label=\textbf{\Alph*)}]
\begin{multicols}{2}
  \item móvil
  \item mamá
  \item papá
  \item café
  \item sofá
  \item compás
\end{multicols}
\end{enumerate}
\vspace{0.15cm}
\item ¿Cuál de las siguientes palabras es grave?
\begin{enumerate}[label=\textbf{\Alph*)}]
\begin{multicols}{2}
  \item túnel
  \item canción
  \item razón
  \item balón
  \item camión
  \item avión
\end{multicols}
\end{enumerate}
\vspace{0.15cm}
\item ¿Cuál de las siguientes palabras es grave?
\begin{enumerate}[label=\textbf{\Alph*)}]
\begin{multicols}{2}
  \item nivel
  \item cantar
  \item corazón
  \item papel
  \item reloj
  \item jamás
\end{multicols}
\end{enumerate}
\vspace{0.15cm}
\item ¿Cuál de las siguientes palabras es grave?
\begin{enumerate}[label=\textbf{\Alph*)}]
\begin{multicols}{2}
  \item revista
  \item mamá
  \item papá
  \item café
  \item sofá
  \item compás
\end{multicols}
\end{enumerate}
\vspace{0.15cm}
\item ¿Cuál de las siguientes palabras es grave?
\begin{enumerate}[label=\textbf{\Alph*)}]
\begin{multicols}{2}
  \item cuaderno
  \item canción
  \item razón
  \item balón
  \item camión
  \item avión
\end{multicols}
\end{enumerate}
\vspace{0.15cm}
\item ¿Cuál de las siguientes palabras es grave?
\begin{enumerate}[label=\textbf{\Alph*)}]
\begin{multicols}{2}
  \item ventana
  \item cantar
  \item corazón
  \item papel
  \item reloj
  \item jamás
\end{multicols}
\end{enumerate}
\vspace{0.15cm}
\item ¿Cuál de las siguientes palabras es grave?
\begin{enumerate}[label=\textbf{\Alph*)}]
\begin{multicols}{2}
  \item puerta
  \item mamá
  \item papá
  \item café
  \item sofá
  \item compás
\end{multicols}
\end{enumerate}
\vspace{0.15cm}
\item ¿Cuál de las siguientes palabras es grave?
\begin{enumerate}[label=\textbf{\Alph*)}]
\begin{multicols}{2}
  \item techo
  \item canción
  \item razón
  \item balón
  \item camión
  \item avión
\end{multicols}
\end{enumerate}
\vspace{0.15cm}
\item ¿Cuál de las siguientes palabras es grave?
\begin{enumerate}[label=\textbf{\Alph*)}]
\begin{multicols}{2}
  \item piso
  \item cantar
  \item corazón
  \item papel
  \item reloj
  \item jamás
\end{multicols}
\end{enumerate}
\vspace{0.15cm}
\item ¿Cuál de las siguientes palabras es grave?
\begin{enumerate}[label=\textbf{\Alph*)}]
\begin{multicols}{2}
  \item pared
  \item mamá
  \item papá
  \item café
  \item sofá
  \item compás
\end{multicols}
\end{enumerate}
\vspace{0.15cm}
\item ¿Cuál de las siguientes palabras es grave?
\begin{enumerate}[label=\textbf{\Alph*)}]
\begin{multicols}{2}
  \item mantel
  \item canción
  \item razón
  \item balón
  \item camión
  \item avión
\end{multicols}
\end{enumerate}
\vspace{0.15cm}
\item ¿Cuál de las siguientes palabras es grave?
\begin{enumerate}[label=\textbf{\Alph*)}]
\begin{multicols}{2}
  \item pastel
  \item cantar
  \item corazón
  \item papel
  \item reloj
  \item jamás
\end{multicols}
\end{enumerate}
\vspace{0.15cm}
\item ¿Cuál de las siguientes palabras es grave?
\begin{enumerate}[label=\textbf{\Alph*)}]
\begin{multicols}{2}
  \item clavel
  \item mamá
  \item papá
  \item café
  \item sofá
  \item compás
\end{multicols}
\end{enumerate}
\vspace{0.15cm}
\item ¿Cuál de las siguientes palabras es grave?
\begin{enumerate}[label=\textbf{\Alph*)}]
\begin{multicols}{2}
  \item examen
  \item canción
  \item razón
  \item balón
  \item camión
  \item avión
\end{multicols}
\end{enumerate}
\vspace{0.15cm}
\item ¿Cuál de las siguientes palabras es grave?
\begin{enumerate}[label=\textbf{\Alph*)}]
\begin{multicols}{2}
  \item joven
  \item cantar
  \item corazón
  \item papel
  \item reloj
  \item jamás
\end{multicols}
\end{enumerate}
\vspace{0.15cm}
\item ¿Cuál de las siguientes palabras es grave?
\begin{enumerate}[label=\textbf{\Alph*)}]
\begin{multicols}{2}
  \item orden
  \item mamá
  \item papá
  \item café
  \item sofá
  \item compás
\end{multicols}
\end{enumerate}
\vspace{0.15cm}
\item ¿Cuál de las siguientes palabras es grave?
\begin{enumerate}[label=\textbf{\Alph*)}]
\begin{multicols}{2}
  \item imagen
  \item canción
  \item razón
  \item balón
  \item camión
  \item avión
\end{multicols}
\end{enumerate}
\vspace{0.15cm}
\item ¿Cuál de las siguientes palabras es grave?
\begin{enumerate}[label=\textbf{\Alph*)}]
\begin{multicols}{2}
  \item virgen
  \item cantar
  \item corazón
  \item papel
  \item reloj
  \item jamás
\end{multicols}
\end{enumerate}
\vspace{0.15cm}
\item ¿Cuál de las siguientes palabras es grave?
\begin{enumerate}[label=\textbf{\Alph*)}]
\begin{multicols}{2}
  \item margen
  \item mamá
  \item papá
  \item café
  \item sofá
  \item compás
\end{multicols}
\end{enumerate}
\vspace{0.15cm}
\item ¿Cuál de las siguientes palabras es grave?
\begin{enumerate}[label=\textbf{\Alph*)}]
\begin{multicols}{2}
  \item origen
  \item canción
  \item razón
  \item balón
  \item camión
  \item avión
\end{multicols}
\end{enumerate}
\vspace{0.15cm}
\item ¿Cuál de las siguientes palabras es grave?
\begin{enumerate}[label=\textbf{\Alph*)}]
\begin{multicols}{2}
  \item certamen
  \item cantar
  \item corazón
  \item papel
  \item reloj
  \item jamás
\end{multicols}
\end{enumerate}
\vspace{0.15cm}
\item ¿Cuál de las siguientes palabras es grave?
\begin{enumerate}[label=\textbf{\Alph*)}]
\begin{multicols}{2}
  \item resumen
  \item mamá
  \item papá
  \item café
  \item sofá
  \item compás
\end{multicols}
\end{enumerate}
\vspace{0.15cm}
\item ¿Cuál de las siguientes palabras es grave?
\begin{enumerate}[label=\textbf{\Alph*)}]
\begin{multicols}{2}
  \item volumen
  \item canción
  \item razón
  \item balón
  \item camión
  \item avión ---
\end{multicols}
\end{enumerate}
\vspace{0.15cm}
\end{enumerate}\n
\subsubsection*{SECCIÓN III: ACENTUACIÓN – PALABRAS ESDRÚJULAS Y SOBRESDRÚJULAS (81–120)}

\begin{enumerate}[label=\textbf{\arabic*.}]
\setcounter{enumi}{80}
\item ¿Cuál de las siguientes palabras es esdrújula?
\begin{enumerate}[label=\textbf{\Alph*)}]
\begin{multicols}{2}
  \item música
  \item cantar
  \item corazón
  \item papel
  \item reloj
  \item jamás
\end{multicols}
\end{enumerate}
\vspace{0.15cm}
\item ¿Cuál de las siguientes palabras es esdrújula?
\begin{enumerate}[label=\textbf{\Alph*)}]
\begin{multicols}{2}
  \item lápiz
  \item mamá
  \item papá
  \item café
  \item sofá
  \item compás
\end{multicols}
\end{enumerate}
\vspace{0.15cm}
\item ¿Cuál de las siguientes palabras es esdrújula?
\begin{enumerate}[label=\textbf{\Alph*)}]
\begin{multicols}{2}
  \item árbol
  \item canción
  \item razón
  \item balón
  \item camión
  \item avión
\end{multicols}
\end{enumerate}
\vspace{0.15cm}
\item ¿Cuál de las siguientes palabras es esdrújula?
\begin{enumerate}[label=\textbf{\Alph*)}]
\begin{multicols}{2}
  \item teléfono
  \item cantar
  \item corazón
  \item papel
  \item reloj
  \item jamás
\end{multicols}
\end{enumerate}
\vspace{0.15cm}
\item ¿Cuál de las siguientes palabras es esdrújula?
\begin{enumerate}[label=\textbf{\Alph*)}]
\begin{multicols}{2}
  \item pájaro
  \item mamá
  \item papá
  \item café
  \item sofá
  \item compás
\end{multicols}
\end{enumerate}
\vspace{0.15cm}
\item ¿Cuál de las siguientes palabras es esdrújula?
\begin{enumerate}[label=\textbf{\Alph*)}]
\begin{multicols}{2}
  \item cárcel
  \item canción
  \item razón
  \item balón
  \item camión
  \item avión
\end{multicols}
\end{enumerate}
\vspace{0.15cm}
\item ¿Cuál de las siguientes palabras es esdrújula?
\begin{enumerate}[label=\textbf{\Alph*)}]
\begin{multicols}{2}
  \item azúcar
  \item cantar
  \item corazón
  \item papel
  \item reloj
  \item jamás
\end{multicols}
\end{enumerate}
\vspace{0.15cm}
\item ¿Cuál de las siguientes palabras es esdrújula?
\begin{enumerate}[label=\textbf{\Alph*)}]
\begin{multicols}{2}
  \item mármol
  \item mamá
  \item papá
  \item café
  \item sofá
  \item compás
\end{multicols}
\end{enumerate}
\vspace{0.15cm}
\item ¿Cuál de las siguientes palabras es esdrújula?
\begin{enumerate}[label=\textbf{\Alph*)}]
\begin{multicols}{2}
  \item césped
  \item canción
  \item razón
  \item balón
  \item camión
  \item avión
\end{multicols}
\end{enumerate}
\vspace{0.15cm}
\item ¿Cuál de las siguientes palabras es esdrújula?
\begin{enumerate}[label=\textbf{\Alph*)}]
\begin{multicols}{2}
  \item fácil
  \item cantar
  \item corazón
  \item papel
  \item reloj
  \item jamás
\end{multicols}
\end{enumerate}
\vspace{0.15cm}
\item ¿Cuál de las siguientes palabras es esdrújula?
\begin{enumerate}[label=\textbf{\Alph*)}]
\begin{multicols}{2}
  \item útil
  \item mamá
  \item papá
  \item café
  \item sofá
  \item compás
\end{multicols}
\end{enumerate}
\vspace{0.15cm}
\item ¿Cuál de las siguientes palabras es esdrújula?
\begin{enumerate}[label=\textbf{\Alph*)}]
\begin{multicols}{2}
  \item débil
  \item canción
  \item razón
  \item balón
  \item camión
  \item avión
\end{multicols}
\end{enumerate}
\vspace{0.15cm}
\item ¿Cuál de las siguientes palabras es esdrújula?
\begin{enumerate}[label=\textbf{\Alph*)}]
\begin{multicols}{2}
  \item móvil
  \item cantar
  \item corazón
  \item papel
  \item reloj
  \item jamás
\end{multicols}
\end{enumerate}
\vspace{0.15cm}
\item ¿Cuál de las siguientes palabras es esdrújula?
\begin{enumerate}[label=\textbf{\Alph*)}]
\begin{multicols}{2}
  \item túnel
  \item mamá
  \item papá
  \item café
  \item sofá
  \item compás
\end{multicols}
\end{enumerate}
\vspace{0.15cm}
\item ¿Cuál de las siguientes palabras es esdrújula?
\begin{enumerate}[label=\textbf{\Alph*)}]
\begin{multicols}{2}
  \item médico
  \item canción
  \item razón
  \item balón
  \item camión
  \item avión
\end{multicols}
\end{enumerate}
\vspace{0.15cm}
\item ¿Cuál de las siguientes palabras es esdrújula?
\begin{enumerate}[label=\textbf{\Alph*)}]
\begin{multicols}{2}
  \item médico
  \item cantar
  \item corazón
  \item papel
  \item reloj
  \item jamás
\end{multicols}
\end{enumerate}
\vspace{0.15cm}
\item ¿Cuál de las siguientes palabras es esdrújula?
\begin{enumerate}[label=\textbf{\Alph*)}]
\begin{multicols}{2}
  \item brújula
  \item mamá
  \item papá
  \item café
  \item sofá
  \item compás
\end{multicols}
\end{enumerate}
\vspace{0.15cm}
\item ¿Cuál de las siguientes palabras es esdrújula?
\begin{enumerate}[label=\textbf{\Alph*)}]
\begin{multicols}{2}
  \item cámara
  \item canción
  \item razón
  \item balón
  \item camión
  \item avión
\end{multicols}
\end{enumerate}
\vspace{0.15cm}
\item ¿Cuál de las siguientes palabras es esdrújula?
\begin{enumerate}[label=\textbf{\Alph*)}]
\begin{multicols}{2}
  \item sábado
  \item cantar
  \item corazón
  \item papel
  \item reloj
  \item jamás
\end{multicols}
\end{enumerate}
\vspace{0.15cm}
\item ¿Cuál de las siguientes palabras es esdrújula?
\begin{enumerate}[label=\textbf{\Alph*)}]
\begin{multicols}{2}
  \item miércoles
  \item mamá
  \item papá
  \item café
  \item sofá
  \item compás
\end{multicols}
\end{enumerate}
\vspace{0.15cm}
\item ¿Cuál de las siguientes palabras es esdrújula?
\begin{enumerate}[label=\textbf{\Alph*)}]
\begin{multicols}{2}
  \item murciélago
  \item canción
  \item razón
  \item balón
  \item camión
  \item avión
\end{multicols}
\end{enumerate}
\vspace{0.15cm}
\item ¿Cuál de las siguientes palabras es esdrújula?
\begin{enumerate}[label=\textbf{\Alph*)}]
\begin{multicols}{2}
  \item parásito
  \item cantar
  \item corazón
  \item papel
  \item reloj
  \item jamás
\end{multicols}
\end{enumerate}
\vspace{0.15cm}
\item ¿Cuál de las siguientes palabras es esdrújula?
\begin{enumerate}[label=\textbf{\Alph*)}]
\begin{multicols}{2}
  \item pelícano
  \item mamá
  \item papá
  \item café
  \item sofá
  \item compás
\end{multicols}
\end{enumerate}
\vspace{0.15cm}
\item ¿Cuál de las siguientes palabras es esdrújula?
\begin{enumerate}[label=\textbf{\Alph*)}]
\begin{multicols}{2}
  \item plátano
  \item canción
  \item razón
  \item balón
  \item camión
  \item avión
\end{multicols}
\end{enumerate}
\vspace{0.15cm}
\item ¿Cuál de las siguientes palabras es esdrújula?
\begin{enumerate}[label=\textbf{\Alph*)}]
\begin{multicols}{2}
  \item prójimo
  \item cantar
  \item corazón
  \item papel
  \item reloj
  \item jamás
\end{multicols}
\end{enumerate}
\vspace{0.15cm}
\item ¿Cuál de las siguientes palabras es esdrújula?
\begin{enumerate}[label=\textbf{\Alph*)}]
\begin{multicols}{2}
  \item ráfaga
  \item mamá
  \item papá
  \item café
  \item sofá
  \item compás
\end{multicols}
\end{enumerate}
\vspace{0.15cm}
\item ¿Cuál de las siguientes palabras es esdrújula?
\begin{enumerate}[label=\textbf{\Alph*)}]
\begin{multicols}{2}
  \item relámpago
  \item canción
  \item razón
  \item balón
  \item camión
  \item avión
\end{multicols}
\end{enumerate}
\vspace{0.15cm}
\item ¿Cuál de las siguientes palabras es esdrújula?
\begin{enumerate}[label=\textbf{\Alph*)}]
\begin{multicols}{2}
  \item síntoma
  \item cantar
  \item corazón
  \item papel
  \item reloj
  \item jamás
\end{multicols}
\end{enumerate}
\vspace{0.15cm}
\item ¿Cuál de las siguientes palabras es esdrújula?
\begin{enumerate}[label=\textbf{\Alph*)}]
\begin{multicols}{2}
  \item teléfono
  \item mamá
  \item papá
  \item café
  \item sofá
  \item compás
\end{multicols}
\end{enumerate}
\vspace{0.15cm}
\item ¿Cuál de las siguientes palabras es esdrújula?
\begin{enumerate}[label=\textbf{\Alph*)}]
\begin{multicols}{2}
  \item término
  \item canción
  \item razón
  \item balón
  \item camión
  \item avión
\end{multicols}
\end{enumerate}
\vspace{0.15cm}
\item ¿Cuál de las siguientes palabras es esdrújula?
\begin{enumerate}[label=\textbf{\Alph*)}]
\begin{multicols}{2}
  \item último
  \item cantar
  \item corazón
  \item papel
  \item reloj
  \item jamás
\end{multicols}
\end{enumerate}
\vspace{0.15cm}
\item ¿Cuál de las siguientes palabras es esdrújula?
\begin{enumerate}[label=\textbf{\Alph*)}]
\begin{multicols}{2}
  \item útilmente
  \item mamá
  \item papá
  \item café
  \item sofá
  \item compás
\end{multicols}
\end{enumerate}
\vspace{0.15cm}
\item ¿Cuál de las siguientes palabras es esdrújula?
\begin{enumerate}[label=\textbf{\Alph*)}]
\begin{multicols}{2}
  \item ábaco
  \item canción
  \item razón
  \item balón
  \item camión
  \item avión
\end{multicols}
\end{enumerate}
\vspace{0.15cm}
\item ¿Cuál de las siguientes palabras es esdrújula?
\begin{enumerate}[label=\textbf{\Alph*)}]
\begin{multicols}{2}
  \item álgebra
  \item cantar
  \item corazón
  \item papel
  \item reloj
  \item jamás
\end{multicols}
\end{enumerate}
\vspace{0.15cm}
\item ¿Cuál de las siguientes palabras es esdrújula?
\begin{enumerate}[label=\textbf{\Alph*)}]
\begin{multicols}{2}
  \item ánimo
  \item mamá
  \item papá
  \item café
  \item sofá
  \item compás
\end{multicols}
\end{enumerate}
\vspace{0.15cm}
\item ¿Cuál de las siguientes palabras es esdrújula?
\begin{enumerate}[label=\textbf{\Alph*)}]
\begin{multicols}{2}
  \item área
  \item canción
  \item razón
  \item balón
  \item camión
  \item avión
\end{multicols}
\end{enumerate}
\vspace{0.15cm}
\item ¿Cuál de las siguientes palabras es esdrújula?
\begin{enumerate}[label=\textbf{\Alph*)}]
\begin{multicols}{2}
  \item bárbaro
  \item cantar
  \item corazón
  \item papel
  \item reloj
  \item jamás
\end{multicols}
\end{enumerate}
\vspace{0.15cm}
\item ¿Cuál de las siguientes palabras es esdrújula?
\begin{enumerate}[label=\textbf{\Alph*)}]
\begin{multicols}{2}
  \item cálculo
  \item mamá
  \item papá
  \item café
  \item sofá
  \item compás
\end{multicols}
\end{enumerate}
\vspace{0.15cm}
\item ¿Cuál de las siguientes palabras es esdrújula?
\begin{enumerate}[label=\textbf{\Alph*)}]
\begin{multicols}{2}
  \item capítulo
  \item canción
  \item razón
  \item balón
  \item camión
  \item avión
\end{multicols}
\end{enumerate}
\vspace{0.15cm}
\item ¿Cuál de las siguientes palabras es sobresdrújula?
\begin{enumerate}[label=\textbf{\Alph*)}]
\begin{multicols}{2}
  \item cómpramelo
  \item cantar
  \item corazón
  \item papel
  \item reloj
  \item jamás ---
\end{multicols}
\end{enumerate}
\vspace{0.15cm}
\end{enumerate}\n
\subsubsection*{SECCIÓN IV: TILDE DIACRÍTICA (121–170)}

\begin{enumerate}[label=\textbf{\arabic*.}]
\setcounter{enumi}{120}
\item ¿Cuál es la forma correcta?
\begin{enumerate}[label=\textbf{\Alph*)}]
\begin{multicols}{2}
  \item Él niño
  \item Él niño
  \item Él niño
\end{multicols}
\end{enumerate}
\vspace{0.15cm}
\item ¿Cuál es la forma correcta?
\begin{enumerate}[label=\textbf{\Alph*)}]
\begin{multicols}{2}
  \item Mi casa
  \item Mí casa
  \item Mi casa
  \item Mí casa
  \item Mi casa
  \item Mí casa
\end{multicols}
\end{enumerate}
\vspace{0.15cm}
\item ¿Cuál es la forma correcta?
\begin{enumerate}[label=\textbf{\Alph*)}]
\begin{multicols}{2}
  \item Tu libro
  \item Tú libro
  \item Tu libro
  \item Tú libro
  \item Tu libro
  \item Tú libro
\end{multicols}
\end{enumerate}
\vspace{0.15cm}
\item ¿Cuál es la forma correcta?
\begin{enumerate}[label=\textbf{\Alph*)}]
\begin{multicols}{2}
\end{multicols}
\end{enumerate}
\vspace{0.15cm}
\item ¿Cuál es la forma correcta?
\begin{enumerate}[label=\textbf{\Alph*)}]
\begin{multicols}{2}
  \item Mi
  \item Mí
  \item Mi
  \item Mí
  \item Mi
  \item Mí
\end{multicols}
\end{enumerate}
\vspace{0.15cm}
\item ¿Cuál es la forma correcta?
\begin{enumerate}[label=\textbf{\Alph*)}]
\begin{multicols}{2}
  \item Sé
  \item Se
  \item Sé
  \item Se
  \item Sé
  \item Se
\end{multicols}
\end{enumerate}
\vspace{0.15cm}
\item ¿Cuál es la forma correcta?
\begin{enumerate}[label=\textbf{\Alph*)}]
\begin{multicols}{2}
\end{multicols}
\end{enumerate}
\vspace{0.15cm}
\item ¿Cuál es la forma correcta?
\begin{enumerate}[label=\textbf{\Alph*)}]
\begin{multicols}{2}
  \item Más
  \item Mas
  \item Más
  \item Mas
  \item Más
  \item Mas
\end{multicols}
\end{enumerate}
\vspace{0.15cm}
\item ¿Cuál es la forma correcta?
\begin{enumerate}[label=\textbf{\Alph*)}]
\begin{multicols}{2}
  \item Sí
  \item Si
  \item Sí
  \item Si
  \item Sí
  \item Si
\end{multicols}
\end{enumerate}
\vspace{0.15cm}
\item ¿Cuál es la forma correcta?
\begin{enumerate}[label=\textbf{\Alph*)}]
\begin{multicols}{2}
\end{multicols}
\end{enumerate}
\vspace{0.15cm}
\item ¿Cuál es la forma correcta?
\begin{enumerate}[label=\textbf{\Alph*)}]
\begin{multicols}{2}
  \item Él
  \item Él
  \item Él
\end{multicols}
\end{enumerate}
\vspace{0.15cm}
\item ¿Cuál es la forma correcta?
\begin{enumerate}[label=\textbf{\Alph*)}]
\begin{multicols}{2}
  \item Tú
  \item Tu
  \item Tú
  \item Tu
  \item Tú
  \item Tu
\end{multicols}
\end{enumerate}
\vspace{0.15cm}
\item ¿Cuál es la forma correcta?
\begin{enumerate}[label=\textbf{\Alph*)}]
\begin{multicols}{2}
  \item Mí
  \item Mi
  \item Mí
  \item Mi
  \item Mí
  \item Mi
\end{multicols}
\end{enumerate}
\vspace{0.15cm}
\item ¿Cuál es la forma correcta?
\begin{enumerate}[label=\textbf{\Alph*)}]
\begin{multicols}{2}
  \item Sé
  \item Se
  \item Sé
  \item Se
  \item Sé
  \item Se
\end{multicols}
\end{enumerate}
\vspace{0.15cm}
\item ¿Cuál es la forma correcta?
\begin{enumerate}[label=\textbf{\Alph*)}]
\begin{multicols}{2}
\end{multicols}
\end{enumerate}
\vspace{0.15cm}
\item ¿Cuál es la forma correcta?
\begin{enumerate}[label=\textbf{\Alph*)}]
\begin{multicols}{2}
  \item Más
  \item Mas
  \item Más
  \item Mas
  \item Más
  \item Mas
\end{multicols}
\end{enumerate}
\vspace{0.15cm}
\item ¿Cuál es la forma correcta?
\begin{enumerate}[label=\textbf{\Alph*)}]
\begin{multicols}{2}
  \item Sí
  \item Si
  \item Sí
  \item Si
  \item Sí
  \item Si
\end{multicols}
\end{enumerate}
\vspace{0.15cm}
\item ¿Cuál es la forma correcta?
\begin{enumerate}[label=\textbf{\Alph*)}]
\begin{multicols}{2}
\end{multicols}
\end{enumerate}
\vspace{0.15cm}
\item ¿Cuál es la forma correcta?
\begin{enumerate}[label=\textbf{\Alph*)}]
\begin{multicols}{2}
  \item Él
  \item Él
  \item Él
\end{multicols}
\end{enumerate}
\vspace{0.15cm}
\item ¿Cuál es la forma correcta?
\begin{enumerate}[label=\textbf{\Alph*)}]
\begin{multicols}{2}
  \item Tu
  \item Tú
  \item Tu
  \item Tú
  \item Tu
  \item Tú
\end{multicols}
\end{enumerate}
\vspace{0.15cm}
\item ¿Cuál es la forma correcta?
\begin{enumerate}[label=\textbf{\Alph*)}]
\begin{multicols}{2}
  \item Mi
  \item Mí
  \item Mi
  \item Mí
  \item Mi
  \item Mí
\end{multicols}
\end{enumerate}
\vspace{0.15cm}
\item ¿Cuál es la forma correcta?
\begin{enumerate}[label=\textbf{\Alph*)}]
\begin{multicols}{2}
  \item Se
  \item Sé
  \item Se
  \item Sé
  \item Se
  \item Sé
\end{multicols}
\end{enumerate}
\vspace{0.15cm}
\item ¿Cuál es la forma correcta?
\begin{enumerate}[label=\textbf{\Alph*)}]
\begin{multicols}{2}
\end{multicols}
\end{enumerate}
\vspace{0.15cm}
\item ¿Cuál es la forma correcta?
\begin{enumerate}[label=\textbf{\Alph*)}]
\begin{multicols}{2}
  \item Mas
  \item Más
  \item Mas
  \item Más
  \item Mas
  \item Más
\end{multicols}
\end{enumerate}
\vspace{0.15cm}
\item ¿Cuál es la forma correcta?
\begin{enumerate}[label=\textbf{\Alph*)}]
\begin{multicols}{2}
  \item Si
  \item Sí
  \item Si
  \item Sí
  \item Si
  \item Sí
\end{multicols}
\end{enumerate}
\vspace{0.15cm}
\item ¿Cuál es la forma correcta?
\begin{enumerate}[label=\textbf{\Alph*)}]
\begin{multicols}{2}
\end{multicols}
\end{enumerate}
\vspace{0.15cm}
\item ¿Cuál es la forma correcta?
\begin{enumerate}[label=\textbf{\Alph*)}]
\begin{multicols}{2}
  \item Él
  \item Él
  \item Él
\end{multicols}
\end{enumerate}
\vspace{0.15cm}
\item ¿Cuál es la forma correcta?
\begin{enumerate}[label=\textbf{\Alph*)}]
\begin{multicols}{2}
  \item Tú
  \item Tu
  \item Tú
  \item Tu
  \item Tú
  \item Tu
\end{multicols}
\end{enumerate}
\vspace{0.15cm}
\item ¿Cuál es la forma correcta?
\begin{enumerate}[label=\textbf{\Alph*)}]
\begin{multicols}{2}
  \item Mí
  \item Mi
  \item Mí
  \item Mi
  \item Mí
  \item Mi
\end{multicols}
\end{enumerate}
\vspace{0.15cm}
\item ¿Cuál es la forma correcta?
\begin{enumerate}[label=\textbf{\Alph*)}]
\begin{multicols}{2}
  \item Sé
  \item Se
  \item Sé
  \item Se
  \item Sé
  \item Se
\end{multicols}
\end{enumerate}
\vspace{0.15cm}
\item ¿Cuál es la forma correcta?
\begin{enumerate}[label=\textbf{\Alph*)}]
\begin{multicols}{2}
\end{multicols}
\end{enumerate}
\vspace{0.15cm}
\item ¿Cuál es la forma correcta?
\begin{enumerate}[label=\textbf{\Alph*)}]
\begin{multicols}{2}
  \item Más
  \item Mas
  \item Más
  \item Mas
  \item Más
  \item Mas
\end{multicols}
\end{enumerate}
\vspace{0.15cm}
\item ¿Cuál es la forma correcta?
\begin{enumerate}[label=\textbf{\Alph*)}]
\begin{multicols}{2}
  \item Sí
  \item Si
  \item Sí
  \item Si
  \item Sí
  \item Si
\end{multicols}
\end{enumerate}
\vspace{0.15cm}
\item ¿Cuál es la forma correcta?
\begin{enumerate}[label=\textbf{\Alph*)}]
\begin{multicols}{2}
\end{multicols}
\end{enumerate}
\vspace{0.15cm}
\item ¿Cuál es la forma correcta?
\begin{enumerate}[label=\textbf{\Alph*)}]
\begin{multicols}{2}
  \item Él
  \item Él
  \item Él
\end{multicols}
\end{enumerate}
\vspace{0.15cm}
\item ¿Cuál es la forma correcta?
\begin{enumerate}[label=\textbf{\Alph*)}]
\begin{multicols}{2}
  \item Tú
  \item Tu
  \item Tú
  \item Tu
  \item Tú
  \item Tu
\end{multicols}
\end{enumerate}
\vspace{0.15cm}
\item ¿Cuál es la forma correcta?
\begin{enumerate}[label=\textbf{\Alph*)}]
\begin{multicols}{2}
  \item Mí
  \item Mi
  \item Mí
  \item Mi
  \item Mí
  \item Mi
\end{multicols}
\end{enumerate}
\vspace{0.15cm}
\item ¿Cuál es la forma correcta?
\begin{enumerate}[label=\textbf{\Alph*)}]
\begin{multicols}{2}
  \item Sé
  \item Se
  \item Sé
  \item Se
  \item Sé
  \item Se
\end{multicols}
\end{enumerate}
\vspace{0.15cm}
\item ¿Cuál es la forma correcta?
\begin{enumerate}[label=\textbf{\Alph*)}]
\begin{multicols}{2}
\end{multicols}
\end{enumerate}
\vspace{0.15cm}
\item ¿Cuál es la forma correcta?
\begin{enumerate}[label=\textbf{\Alph*)}]
\begin{multicols}{2}
  \item Más
  \item Mas
  \item Más
  \item Mas
  \item Más
  \item Mas
\end{multicols}
\end{enumerate}
\vspace{0.15cm}
\item ¿Cuál es la forma correcta?
\begin{enumerate}[label=\textbf{\Alph*)}]
\begin{multicols}{2}
  \item Sí
  \item Si
  \item Sí
  \item Si
  \item Sí
  \item Si
\end{multicols}
\end{enumerate}
\vspace{0.15cm}
\item ¿Cuál es la forma correcta?
\begin{enumerate}[label=\textbf{\Alph*)}]
\begin{multicols}{2}
\end{multicols}
\end{enumerate}
\vspace{0.15cm}
\item ¿Cuál es la forma correcta?
\begin{enumerate}[label=\textbf{\Alph*)}]
\begin{multicols}{2}
  \item Él
  \item Él
  \item Él
\end{multicols}
\end{enumerate}
\vspace{0.15cm}
\item ¿Cuál es la forma correcta?
\begin{enumerate}[label=\textbf{\Alph*)}]
\begin{multicols}{2}
  \item Tú
  \item Tu
  \item Tú
  \item Tu
  \item Tú
  \item Tu
\end{multicols}
\end{enumerate}
\vspace{0.15cm}
\item ¿Cuál es la forma correcta?
\begin{enumerate}[label=\textbf{\Alph*)}]
\begin{multicols}{2}
  \item Mí
  \item Mi
  \item Mí
  \item Mi
  \item Mí
  \item Mi
\end{multicols}
\end{enumerate}
\vspace{0.15cm}
\item ¿Cuál es la forma correcta?
\begin{enumerate}[label=\textbf{\Alph*)}]
\begin{multicols}{2}
  \item Sé
  \item Se
  \item Sé
  \item Se
  \item Sé
  \item Se
\end{multicols}
\end{enumerate}
\vspace{0.15cm}
\item ¿Cuál es la forma correcta?
\begin{enumerate}[label=\textbf{\Alph*)}]
\begin{multicols}{2}
\end{multicols}
\end{enumerate}
\vspace{0.15cm}
\item ¿Cuál es la forma correcta?
\begin{enumerate}[label=\textbf{\Alph*)}]
\begin{multicols}{2}
  \item Más
  \item Mas
  \item Más
  \item Mas
  \item Más
  \item Mas
\end{multicols}
\end{enumerate}
\vspace{0.15cm}
\item ¿Cuál es la forma correcta?
\begin{enumerate}[label=\textbf{\Alph*)}]
\begin{multicols}{2}
  \item Sí
  \item Si
  \item Sí
  \item Si
  \item Sí
  \item Si
\end{multicols}
\end{enumerate}
\vspace{0.15cm}
\item ¿Cuál es la forma correcta?
\begin{enumerate}[label=\textbf{\Alph*)}]
\begin{multicols}{2}
\end{multicols}
\end{enumerate}
\vspace{0.15cm}
\end{enumerate}\n
\subsubsection*{SECCIÓN V: TILDE DISOLVENTE (171–200)}

\begin{enumerate}[label=\textbf{\arabic*.}]
\setcounter{enumi}{170}
\item ¿Cuál es la forma correcta?
\begin{enumerate}[label=\textbf{\Alph*)}]
\begin{multicols}{2}
\end{multicols}
\end{enumerate}
\vspace{0.15cm}
\item ¿Cuál es la forma correcta?
\begin{enumerate}[label=\textbf{\Alph*)}]
\begin{multicols}{2}
  \item Sólo
  \item Solo
  \item Sólo
  \item Solo
  \item Sólo
  \item Solo
\end{multicols}
\end{enumerate}
\vspace{0.15cm}
\item ¿Cuál es la forma correcta?
\begin{enumerate}[label=\textbf{\Alph*)}]
\begin{multicols}{2}
  \item Éste
  \item Éste
  \item Éste
\end{multicols}
\end{enumerate}
\vspace{0.15cm}
\item ¿Cuál es la forma correcta?
\begin{enumerate}[label=\textbf{\Alph*)}]
\begin{multicols}{2}
\end{multicols}
\end{enumerate}
\vspace{0.15cm}
\item ¿Cuál es la forma correcta?
\begin{enumerate}[label=\textbf{\Alph*)}]
\begin{multicols}{2}
\end{multicols}
\end{enumerate}
\vspace{0.15cm}
\item ¿Cuál es la forma correcta?
\begin{enumerate}[label=\textbf{\Alph*)}]
\begin{multicols}{2}
  \item Sé
  \item Se
  \item Sé
  \item Se
  \item Sé
  \item Se
\end{multicols}
\end{enumerate}
\vspace{0.15cm}
\item ¿Cuál es la forma correcta?
\begin{enumerate}[label=\textbf{\Alph*)}]
\begin{multicols}{2}
  \item Té
  \item Te
  \item Té
  \item Te
  \item Té
  \item Te
\end{multicols}
\end{enumerate}
\vspace{0.15cm}
\item ¿Cuál es la forma correcta?
\begin{enumerate}[label=\textbf{\Alph*)}]
\begin{multicols}{2}
  \item Más
  \item Mas
  \item Más
  \item Mas
  \item Más
  \item Mas
\end{multicols}
\end{enumerate}
\vspace{0.15cm}
\item ¿Cuál es la forma correcta?
\begin{enumerate}[label=\textbf{\Alph*)}]
\begin{multicols}{2}
  \item Sí
  \item Si
  \item Sí
  \item Si
  \item Sí
  \item Si
\end{multicols}
\end{enumerate}
\vspace{0.15cm}
\item ¿Cuál es la forma correcta?
\begin{enumerate}[label=\textbf{\Alph*)}]
\begin{multicols}{2}
  \item Mí
  \item Mi
  \item Mí
  \item Mi
  \item Mí
  \item Mi
\end{multicols}
\end{enumerate}
\vspace{0.15cm}
\item ¿Cuál es la forma correcta?
\begin{enumerate}[label=\textbf{\Alph*)}]
\begin{multicols}{2}
  \item Tú
  \item Tu
  \item Tú
  \item Tu
  \item Tú
  \item Tu
\end{multicols}
\end{enumerate}
\vspace{0.15cm}
\item ¿Cuál es la forma correcta?
\begin{enumerate}[label=\textbf{\Alph*)}]
\begin{multicols}{2}
  \item Él
  \item Él
  \item Él
\end{multicols}
\end{enumerate}
\vspace{0.15cm}
\item ¿Cuál es la forma correcta?
\begin{enumerate}[label=\textbf{\Alph*)}]
\begin{multicols}{2}
\end{multicols}
\end{enumerate}
\vspace{0.15cm}
\item ¿Cuál es la forma correcta?
\begin{enumerate}[label=\textbf{\Alph*)}]
\begin{multicols}{2}
  \item Sólo
  \item Solo
  \item Sólo
  \item Solo
  \item Sólo
  \item Solo
\end{multicols}
\end{enumerate}
\vspace{0.15cm}
\item ¿Cuál es la forma correcta?
\begin{enumerate}[label=\textbf{\Alph*)}]
\begin{multicols}{2}
  \item Éste
  \item Éste
  \item Éste
\end{multicols}
\end{enumerate}
\vspace{0.15cm}
\item ¿Cuál es la forma correcta?
\begin{enumerate}[label=\textbf{\Alph*)}]
\begin{multicols}{2}
\end{multicols}
\end{enumerate}
\vspace{0.15cm}
\item ¿Cuál es la forma correcta?
\begin{enumerate}[label=\textbf{\Alph*)}]
\begin{multicols}{2}
\end{multicols}
\end{enumerate}
\vspace{0.15cm}
\item ¿Cuál es la forma correcta?
\begin{enumerate}[label=\textbf{\Alph*)}]
\begin{multicols}{2}
  \item Sé
  \item Se
  \item Sé
  \item Se
  \item Sé
  \item Se
\end{multicols}
\end{enumerate}
\vspace{0.15cm}
\item ¿Cuál es la forma correcta?
\begin{enumerate}[label=\textbf{\Alph*)}]
\begin{multicols}{2}
  \item Té
  \item Te
  \item Té
  \item Te
  \item Té
  \item Te
\end{multicols}
\end{enumerate}
\vspace{0.15cm}
\item ¿Cuál es la forma correcta?
\begin{enumerate}[label=\textbf{\Alph*)}]
\begin{multicols}{2}
  \item Más
  \item Mas
  \item Más
  \item Mas
  \item Más
  \item Mas
\end{multicols}
\end{enumerate}
\vspace{0.15cm}
\item ¿Cuál es la forma correcta?
\begin{enumerate}[label=\textbf{\Alph*)}]
\begin{multicols}{2}
  \item Sí
  \item Si
  \item Sí
  \item Si
  \item Sí
  \item Si
\end{multicols}
\end{enumerate}
\vspace{0.15cm}
\item ¿Cuál es la forma correcta?
\begin{enumerate}[label=\textbf{\Alph*)}]
\begin{multicols}{2}
  \item Mí
  \item Mi
  \item Mí
  \item Mi
  \item Mí
  \item Mi
\end{multicols}
\end{enumerate}
\vspace{0.15cm}
\item ¿Cuál es la forma correcta?
\begin{enumerate}[label=\textbf{\Alph*)}]
\begin{multicols}{2}
  \item Tú
  \item Tu
  \item Tú
  \item Tu
  \item Tú
  \item Tu
\end{multicols}
\end{enumerate}
\vspace{0.15cm}
\item ¿Cuál es la forma correcta?
\begin{enumerate}[label=\textbf{\Alph*)}]
\begin{multicols}{2}
  \item Él
  \item Él
  \item Él
\end{multicols}
\end{enumerate}
\vspace{0.15cm}
\item ¿Cuál es la forma correcta?
\begin{enumerate}[label=\textbf{\Alph*)}]
\begin{multicols}{2}
\end{multicols}
\end{enumerate}
\vspace{0.15cm}
\item ¿Cuál es la forma correcta?
\begin{enumerate}[label=\textbf{\Alph*)}]
\begin{multicols}{2}
  \item Sólo
  \item Solo
  \item Sólo
  \item Solo
  \item Sólo
  \item Solo
\end{multicols}
\end{enumerate}
\vspace{0.15cm}
\item ¿Cuál es la forma correcta?
\begin{enumerate}[label=\textbf{\Alph*)}]
\begin{multicols}{2}
  \item Éste
  \item Éste
  \item Éste
\end{multicols}
\end{enumerate}
\vspace{0.15cm}
\item ¿Cuál es la forma correcta?
\begin{enumerate}[label=\textbf{\Alph*)}]
\begin{multicols}{2}
\end{multicols}
\end{enumerate}
\vspace{0.15cm}
\item ¿Cuál es la forma correcta?
\begin{enumerate}[label=\textbf{\Alph*)}]
\begin{multicols}{2}
\end{multicols}
\end{enumerate}
\vspace{0.15cm}
\item ¿Cuál es la forma correcta?
\begin{enumerate}[label=\textbf{\Alph*)}]
\begin{multicols}{2}
  \item Sé
  \item Se
  \item Sé
  \item Se
  \item Sé
  \item Se ---
\end{multicols}
\end{enumerate}
\vspace{0.15cm}
\end{enumerate}\n
\subsubsection*{SECCIÓN VI: USO DE GRAFÍAS (201–250)}

\begin{enumerate}[label=\textbf{\arabic*.}]
\setcounter{enumi}{200}
\item ¿Cuál es la forma correcta?
\begin{enumerate}[label=\textbf{\Alph*)}]
\begin{multicols}{2}
  \item Vaca
  \item Vaca
  \item Vaca
\end{multicols}
\end{enumerate}
\vspace{0.15cm}
\item ¿Cuál es la forma correcta?
\begin{enumerate}[label=\textbf{\Alph*)}]
\begin{multicols}{2}
  \item Vello
  \item Vello
  \item Vello
\end{multicols}
\end{enumerate}
\vspace{0.15cm}
\item ¿Cuál es la forma correcta?
\begin{enumerate}[label=\textbf{\Alph*)}]
\begin{multicols}{2}
\end{multicols}
\end{enumerate}
\vspace{0.15cm}
\item ¿Cuál es la forma correcta?
\begin{enumerate}[label=\textbf{\Alph*)}]
\begin{multicols}{2}
  \item Grabar
  \item Gravar
  \item Grabar
  \item Gravar
  \item Grabar
  \item Gravar
\end{multicols}
\end{enumerate}
\vspace{0.15cm}
\item ¿Cuál es la forma correcta?
\begin{enumerate}[label=\textbf{\Alph*)}]
\begin{multicols}{2}
  \item Haber
  \item Haber
  \item Haber
\end{multicols}
\end{enumerate}
\vspace{0.15cm}
\item ¿Cuál es la forma correcta?
\begin{enumerate}[label=\textbf{\Alph*)}]
\begin{multicols}{2}
  \item Iba
  \item Iva
  \item Iba
  \item Iva
  \item Iba
  \item Iva
\end{multicols}
\end{enumerate}
\vspace{0.15cm}
\item ¿Cuál es la forma correcta?
\begin{enumerate}[label=\textbf{\Alph*)}]
\begin{multicols}{2}
  \item Rebelar
  \item Revelar
  \item Rebelar
  \item Revelar
  \item Rebelar
  \item Revelar
\end{multicols}
\end{enumerate}
\vspace{0.15cm}
\item ¿Cuál es la forma correcta?
\begin{enumerate}[label=\textbf{\Alph*)}]
\begin{multicols}{2}
  \item Sabia
  \item Savia
  \item Sabia
  \item Savia
  \item Sabia
  \item Savia
\end{multicols}
\end{enumerate}
\vspace{0.15cm}
\item ¿Cuál es la forma correcta?
\begin{enumerate}[label=\textbf{\Alph*)}]
\begin{multicols}{2}
  \item Tubo
  \item Tuvo
  \item Tubo
  \item Tuvo
  \item Tubo
  \item Tuvo
\end{multicols}
\end{enumerate}
\vspace{0.15cm}
\item ¿Cuál es la forma correcta?
\begin{enumerate}[label=\textbf{\Alph*)}]
\begin{multicols}{2}
  \item Valla
  \item Vaya
  \item Valla
  \item Vaya
  \item Valla
  \item Vaya
\end{multicols}
\end{enumerate}
\vspace{0.15cm}
\item ¿Cuál es la forma correcta?
\begin{enumerate}[label=\textbf{\Alph*)}]
\begin{multicols}{2}
\end{multicols}
\end{enumerate}
\vspace{0.15cm}
\item ¿Cuál es la forma correcta?
\begin{enumerate}[label=\textbf{\Alph*)}]
\begin{multicols}{2}
  \item Seda
  \item Seda
  \item Seda
\end{multicols}
\end{enumerate}
\vspace{0.15cm}
\item ¿Cuál es la forma correcta?
\begin{enumerate}[label=\textbf{\Alph*)}]
\begin{multicols}{2}
  \item Sima
  \item Sima
  \item Sima
\end{multicols}
\end{enumerate}
\vspace{0.15cm}
\item ¿Cuál es la forma correcta?
\begin{enumerate}[label=\textbf{\Alph*)}]
\begin{multicols}{2}
\end{multicols}
\end{enumerate}
\vspace{0.15cm}
\item ¿Cuál es la forma correcta?
\begin{enumerate}[label=\textbf{\Alph*)}]
\begin{multicols}{2}
\end{multicols}
\end{enumerate}
\vspace{0.15cm}
\item ¿Cuál es la forma correcta?
\begin{enumerate}[label=\textbf{\Alph*)}]
\begin{multicols}{2}
  \item Sien
  \item Sien
  \item Sien
\end{multicols}
\end{enumerate}
\vspace{0.15cm}
\item ¿Cuál es la forma correcta?
\begin{enumerate}[label=\textbf{\Alph*)}]
\begin{multicols}{2}
  \item Siervo
  \item Siervo
  \item Siervo
\end{multicols}
\end{enumerate}
\vspace{0.15cm}
\item ¿Cuál es la forma correcta?
\begin{enumerate}[label=\textbf{\Alph*)}]
\begin{multicols}{2}
  \item Sima
  \item Sima
  \item Sima
\end{multicols}
\end{enumerate}
\vspace{0.15cm}
\item ¿Cuál es la forma correcta?
\begin{enumerate}[label=\textbf{\Alph*)}]
\begin{multicols}{2}
\end{multicols}
\end{enumerate}
\vspace{0.15cm}
\item ¿Cuál es la forma correcta?
\begin{enumerate}[label=\textbf{\Alph*)}]
\begin{multicols}{2}
\end{multicols}
\end{enumerate}
\vspace{0.15cm}
\item ¿Cuál es la forma correcta?
\begin{enumerate}[label=\textbf{\Alph*)}]
\begin{multicols}{2}
  \item Sien
  \item Sien
  \item Sien
\end{multicols}
\end{enumerate}
\vspace{0.15cm}
\item ¿Cuál es la forma correcta?
\begin{enumerate}[label=\textbf{\Alph*)}]
\begin{multicols}{2}
  \item Siervo
  \item Siervo
  \item Siervo
\end{multicols}
\end{enumerate}
\vspace{0.15cm}
\item ¿Cuál es la forma correcta?
\begin{enumerate}[label=\textbf{\Alph*)}]
\begin{multicols}{2}
  \item Sima
  \item Sima
  \item Sima
\end{multicols}
\end{enumerate}
\vspace{0.15cm}
\item ¿Cuál es la forma correcta?
\begin{enumerate}[label=\textbf{\Alph*)}]
\begin{multicols}{2}
\end{multicols}
\end{enumerate}
\vspace{0.15cm}
\item ¿Cuál es la forma correcta?
\begin{enumerate}[label=\textbf{\Alph*)}]
\begin{multicols}{2}
\end{multicols}
\end{enumerate}
\vspace{0.15cm}
\item ¿Cuál es la forma correcta?
\begin{enumerate}[label=\textbf{\Alph*)}]
\begin{multicols}{2}
  \item Sien
  \item Sien
  \item Sien
\end{multicols}
\end{enumerate}
\vspace{0.15cm}
\item ¿Cuál es la forma correcta?
\begin{enumerate}[label=\textbf{\Alph*)}]
\begin{multicols}{2}
  \item Siervo
  \item Siervo
  \item Siervo
\end{multicols}
\end{enumerate}
\vspace{0.15cm}
\item ¿Cuál es la forma correcta?
\begin{enumerate}[label=\textbf{\Alph*)}]
\begin{multicols}{2}
  \item Sima
  \item Sima
  \item Sima
\end{multicols}
\end{enumerate}
\vspace{0.15cm}
\item ¿Cuál es la forma correcta?
\begin{enumerate}[label=\textbf{\Alph*)}]
\begin{multicols}{2}
\end{multicols}
\end{enumerate}
\vspace{0.15cm}
\item ¿Cuál es la forma correcta?
\begin{enumerate}[label=\textbf{\Alph*)}]
\begin{multicols}{2}
\end{multicols}
\end{enumerate}
\vspace{0.15cm}
\item ¿Cuál es la forma correcta?
\begin{enumerate}[label=\textbf{\Alph*)}]
\begin{multicols}{2}
  \item Sien
  \item Sien
  \item Sien
\end{multicols}
\end{enumerate}
\vspace{0.15cm}
\item ¿Cuál es la forma correcta?
\begin{enumerate}[label=\textbf{\Alph*)}]
\begin{multicols}{2}
  \item Siervo
  \item Siervo
  \item Siervo
\end{multicols}
\end{enumerate}
\vspace{0.15cm}
\item ¿Cuál es la forma correcta?
\begin{enumerate}[label=\textbf{\Alph*)}]
\begin{multicols}{2}
  \item Sima
  \item Sima
  \item Sima
\end{multicols}
\end{enumerate}
\vspace{0.15cm}
\item ¿Cuál es la forma correcta?
\begin{enumerate}[label=\textbf{\Alph*)}]
\begin{multicols}{2}
\end{multicols}
\end{enumerate}
\vspace{0.15cm}
\item ¿Cuál es la forma correcta?
\begin{enumerate}[label=\textbf{\Alph*)}]
\begin{multicols}{2}
\end{multicols}
\end{enumerate}
\vspace{0.15cm}
\item ¿Cuál es la forma correcta?
\begin{enumerate}[label=\textbf{\Alph*)}]
\begin{multicols}{2}
  \item Sien
  \item Sien
  \item Sien
\end{multicols}
\end{enumerate}
\vspace{0.15cm}
\item ¿Cuál es la forma correcta?
\begin{enumerate}[label=\textbf{\Alph*)}]
\begin{multicols}{2}
  \item Siervo
  \item Siervo
  \item Siervo
\end{multicols}
\end{enumerate}
\vspace{0.15cm}
\item ¿Cuál es la forma correcta?
\begin{enumerate}[label=\textbf{\Alph*)}]
\begin{multicols}{2}
  \item Sima
  \item Sima
  \item Sima
\end{multicols}
\end{enumerate}
\vspace{0.15cm}
\item ¿Cuál es la forma correcta?
\begin{enumerate}[label=\textbf{\Alph*)}]
\begin{multicols}{2}
\end{multicols}
\end{enumerate}
\vspace{0.15cm}
\item ¿Cuál es la forma correcta?
\begin{enumerate}[label=\textbf{\Alph*)}]
\begin{multicols}{2}
\end{multicols}
\end{enumerate}
\vspace{0.15cm}
\item ¿Cuál es la forma correcta?
\begin{enumerate}[label=\textbf{\Alph*)}]
\begin{multicols}{2}
  \item Sien
  \item Sien
  \item Sien
\end{multicols}
\end{enumerate}
\vspace{0.15cm}
\item ¿Cuál es la forma correcta?
\begin{enumerate}[label=\textbf{\Alph*)}]
\begin{multicols}{2}
  \item Siervo
  \item Siervo
  \item Siervo
\end{multicols}
\end{enumerate}
\vspace{0.15cm}
\item ¿Cuál es la forma correcta?
\begin{enumerate}[label=\textbf{\Alph*)}]
\begin{multicols}{2}
  \item Sima
  \item Sima
  \item Sima
\end{multicols}
\end{enumerate}
\vspace{0.15cm}
\item ¿Cuál es la forma correcta?
\begin{enumerate}[label=\textbf{\Alph*)}]
\begin{multicols}{2}
\end{multicols}
\end{enumerate}
\vspace{0.15cm}
\item ¿Cuál es la forma correcta?
\begin{enumerate}[label=\textbf{\Alph*)}]
\begin{multicols}{2}
\end{multicols}
\end{enumerate}
\vspace{0.15cm}
\item ¿Cuál es la forma correcta?
\begin{enumerate}[label=\textbf{\Alph*)}]
\begin{multicols}{2}
  \item Sien
  \item Sien
  \item Sien
\end{multicols}
\end{enumerate}
\vspace{0.15cm}
\item ¿Cuál es la forma correcta?
\begin{enumerate}[label=\textbf{\Alph*)}]
\begin{multicols}{2}
  \item Siervo
  \item Siervo
  \item Siervo
\end{multicols}
\end{enumerate}
\vspace{0.15cm}
\item ¿Cuál es la forma correcta?
\begin{enumerate}[label=\textbf{\Alph*)}]
\begin{multicols}{2}
  \item Sima
  \item Sima
  \item Sima
\end{multicols}
\end{enumerate}
\vspace{0.15cm}
\item ¿Cuál es la forma correcta?
\begin{enumerate}[label=\textbf{\Alph*)}]
\begin{multicols}{2}
\end{multicols}
\end{enumerate}
\vspace{0.15cm}
\item ¿Cuál es la forma correcta?
\begin{enumerate}[label=\textbf{\Alph*)}]
\begin{multicols}{2}
\end{multicols}
\end{enumerate}
\vspace{0.15cm}
\end{enumerate}\n
\subsubsection*{SECCIÓN VII: SIGNOS DE PUNTUACIÓN (251–300)}

\begin{enumerate}[label=\textbf{\arabic*.}]
\setcounter{enumi}{250}
\item ¿Qué signo se usa para separar elementos de una enumeración?
\begin{enumerate}[label=\textbf{\Alph*)}]
\begin{multicols}{2}
  \item Punto
  \item Punto y coma
  \item Puntos suspensivos
  \item Interrogación
\end{multicols}
\end{enumerate}
\vspace{0.15cm}
\item ¿Qué signo se usa al final de una oración?
\begin{enumerate}[label=\textbf{\Alph*)}]
\begin{multicols}{2}
  \item Punto
  \item Punto y coma
  \item Puntos suspensivos
  \item Interrogación
\end{multicols}
\end{enumerate}
\vspace{0.15cm}
\item ¿Qué signo se usa para indicar una pausa mayor que la coma?
\begin{enumerate}[label=\textbf{\Alph*)}]
\begin{multicols}{2}
  \item Punto
  \item Punto y coma
  \item Puntos suspensivos
  \item Interrogación
\end{multicols}
\end{enumerate}
\vspace{0.15cm}
\item ¿Qué signo se usa antes de una enumeración?
\begin{enumerate}[label=\textbf{\Alph*)}]
\begin{multicols}{2}
  \item Punto
  \item Punto y coma
  \item Puntos suspensivos
  \item Interrogación
\end{multicols}
\end{enumerate}
\vspace{0.15cm}
\item ¿Qué signo se usa para indicar duda o suspenso?
\begin{enumerate}[label=\textbf{\Alph*)}]
\begin{multicols}{2}
  \item Punto
  \item Punto y coma
  \item Puntos suspensivos
  \item Interrogación
\end{multicols}
\end{enumerate}
\vspace{0.15cm}
\item ¿Qué signo se usa al final de una pregunta?
\begin{enumerate}[label=\textbf{\Alph*)}]
\begin{multicols}{2}
  \item Punto
  \item Punto y coma
  \item Puntos suspensivos
  \item Interrogación
\end{multicols}
\end{enumerate}
\vspace{0.15cm}
\item ¿Qué signo se usa al final de una exclamación?
\begin{enumerate}[label=\textbf{\Alph*)}]
\begin{multicols}{2}
  \item Punto
  \item Punto y coma
  \item Puntos suspensivos
\end{multicols}
\end{enumerate}
\vspace{0.15cm}
\item ¿Qué signo se usa para introducir una cita textual?
\begin{enumerate}[label=\textbf{\Alph*)}]
\begin{multicols}{2}
  \item Punto
  \item Punto y coma
  \item Puntos suspensivos
\end{multicols}
\end{enumerate}
\vspace{0.15cm}
\item ¿Qué signo se usa para separar los elementos de una fecha?
\begin{enumerate}[label=\textbf{\Alph*)}]
\begin{multicols}{2}
  \item Punto
  \item Guion
  \item Paréntesis
\end{multicols}
\end{enumerate}
\vspace{0.15cm}
\item ¿Qué signo se usa para encerrar aclaraciones?
\begin{enumerate}[label=\textbf{\Alph*)}]
\begin{multicols}{2}
  \item Punto
  \item Punto y coma
  \item Paréntesis
  \item Interrogación
\end{multicols}
\end{enumerate}
\vspace{0.15cm}
\item ¿Cuál es el uso correcto de la coma en: "Juan, ven aquí"?
\begin{enumerate}[label=\textbf{\Alph*)}]
\begin{multicols}{2}
  \item Vocativo
  \item Inciso
  \item Nexo
\end{multicols}
\end{enumerate}
\vspace{0.15cm}
\item ¿Cuál es el uso correcto de la coma en: "Compré manzanas, peras y uvas"?
\begin{enumerate}[label=\textbf{\Alph*)}]
\begin{multicols}{2}
  \item Vocativo
  \item Inciso
  \item Nexo
\end{multicols}
\end{enumerate}
\vspace{0.15cm}
\item ¿Cuál es el uso correcto de la coma en: "Madrid, capital de España, es una ciudad hermosa"?
\begin{enumerate}[label=\textbf{\Alph*)}]
\begin{multicols}{2}
  \item Vocativo
  \item Inciso
  \item Nexo
\end{multicols}
\end{enumerate}
\vspace{0.15cm}
\item ¿Cuál es el uso correcto de la coma en: "Sin embargo, no fui"?
\begin{enumerate}[label=\textbf{\Alph*)}]
\begin{multicols}{2}
  \item Vocativo
  \item Inciso
  \item Nexo
\end{multicols}
\end{enumerate}
\vspace{0.15cm}
\item ¿Cuál es el uso correcto de la coma en: "Yo estudio; tú, trabajas"?
\begin{enumerate}[label=\textbf{\Alph*)}]
\begin{multicols}{2}
  \item Vocativo
  \item Inciso
  \item Nexo
\end{multicols}
\end{enumerate}
\vspace{0.15cm}
\item ¿Cuál es el uso correcto del punto y coma en: "Llegó tarde; sin embargo, entró"?
\begin{enumerate}[label=\textbf{\Alph*)}]
\begin{multicols}{2}
  \item Separar oraciones
  \item Inciso
  \item Nexo
\end{multicols}
\end{enumerate}
\vspace{0.15cm}
\item ¿Cuál es el uso correcto de los dos puntos en: "Tengo tres hermanos: Ana, Luis y Pedro"?
\begin{enumerate}[label=\textbf{\Alph*)}]
\begin{multicols}{2}
  \item Inciso
  \item Nexo
\end{multicols}
\end{enumerate}
\vspace{0.15cm}
\item ¿Cuál es el uso correcto de los dos puntos en: "Dijo: 'No iré'"?
\begin{enumerate}[label=\textbf{\Alph*)}]
\begin{multicols}{2}
  \item Inciso
  \item Nexo
\end{multicols}
\end{enumerate}
\vspace{0.15cm}
\item ¿Cuál es el uso correcto de los puntos suspensivos en: "No sé si..."?
\begin{enumerate}[label=\textbf{\Alph*)}]
\begin{multicols}{2}
  \item Inciso
  \item Nexo
\end{multicols}
\end{enumerate}
\vspace{0.15cm}
\item ¿Cuál es el uso correcto de las comillas en: "Dijo 'hola'"?
\begin{enumerate}[label=\textbf{\Alph*)}]
\begin{multicols}{2}
  \item Inciso
  \item Nexo
\end{multicols}
\end{enumerate}
\vspace{0.15cm}
\item ¿Cuál es el uso correcto del guion en: "veintidós"?
\begin{enumerate}[label=\textbf{\Alph*)}]
\begin{multicols}{2}
  \item Inciso
  \item Nexo
\end{multicols}
\end{enumerate}
\vspace{0.15cm}
\item ¿Cuál es el uso correcto del guion en: "teórico-práctico"?
\begin{enumerate}[label=\textbf{\Alph*)}]
\begin{multicols}{2}
  \item Inciso
  \item Nexo
\end{multicols}
\end{enumerate}
\vspace{0.15cm}
\item ¿Cuál es el uso correcto de la raya en: "—Ven —dijo—"?
\begin{enumerate}[label=\textbf{\Alph*)}]
\begin{multicols}{2}
  \item Inciso
  \item Nexo
\end{multicols}
\end{enumerate}
\vspace{0.15cm}
\item ¿Cuál es el uso correcto del paréntesis en: "Llegó (por fin) a casa"?
\begin{enumerate}[label=\textbf{\Alph*)}]
\begin{multicols}{2}
  \item Inciso
  \item Nexo
\end{multicols}
\end{enumerate}
\vspace{0.15cm}
\item ¿Cuál es el uso correcto de los signos de interrogación en: "¿Vienes?"?
\begin{enumerate}[label=\textbf{\Alph*)}]
\begin{multicols}{2}
  \item Pregunta
  \item Inciso
  \item Nexo
\end{multicols}
\end{enumerate}
\vspace{0.15cm}
\item ¿Cuál es el uso correcto de los signos de exclamación en: "¡Ven!"?
\begin{enumerate}[label=\textbf{\Alph*)}]
\begin{multicols}{2}
  \item Inciso
  \item Nexo
\end{multicols}
\end{enumerate}
\vspace{0.15cm}
\item ¿Cuál es el uso correcto de la diéresis en: "pingüino"?
\begin{enumerate}[label=\textbf{\Alph*)}]
\begin{multicols}{2}
  \item Pronunciar la u
  \item Inciso
  \item Nexo
\end{multicols}
\end{enumerate}
\vspace{0.15cm}
\item ¿Cuál es el uso correcto de la diéresis en: "vergüenza"?
\begin{enumerate}[label=\textbf{\Alph*)}]
\begin{multicols}{2}
  \item Pronunciar la u
  \item Inciso
  \item Nexo
\end{multicols}
\end{enumerate}
\vspace{0.15cm}
\item ¿Cuál es el uso correcto de la diéresis en: "güero"?
\begin{enumerate}[label=\textbf{\Alph*)}]
\begin{multicols}{2}
  \item Pronunciar la u
  \item Inciso
  \item Nexo
\end{multicols}
\end{enumerate}
\vspace{0.15cm}
\item ¿Cuál es el uso correcto de la diéresis en: "cigüeña"?
\begin{enumerate}[label=\textbf{\Alph*)}]
\begin{multicols}{2}
  \item Pronunciar la u
  \item Inciso
  \item Nexo
\end{multicols}
\end{enumerate}
\vspace{0.15cm}
\item ¿Cuál es el uso correcto de la diéresis en: "agüita"?
\begin{enumerate}[label=\textbf{\Alph*)}]
\begin{multicols}{2}
  \item Pronunciar la u
  \item Inciso
  \item Nexo
\end{multicols}
\end{enumerate}
\vspace{0.15cm}
\item ¿Cuál es el uso correcto de la diéresis en: "lingüística"?
\begin{enumerate}[label=\textbf{\Alph*)}]
\begin{multicols}{2}
  \item Pronunciar la u
  \item Inciso
  \item Nexo
\end{multicols}
\end{enumerate}
\vspace{0.15cm}
\item ¿Cuál es el uso correcto de la diéresis en: "piragüismo"?
\begin{enumerate}[label=\textbf{\Alph*)}]
\begin{multicols}{2}
  \item Pronunciar la u
  \item Inciso
  \item Nexo
\end{multicols}
\end{enumerate}
\vspace{0.15cm}
\item ¿Cuál es el uso correcto de la diéresis en: "bilingüe"?
\begin{enumerate}[label=\textbf{\Alph*)}]
\begin{multicols}{2}
  \item Pronunciar la u
  \item Inciso
  \item Nexo
\end{multicols}
\end{enumerate}
\vspace{0.15cm}
\item ¿Cuál es el uso correcto de la diéresis en: "antigüedad"?
\begin{enumerate}[label=\textbf{\Alph*)}]
\begin{multicols}{2}
  \item Pronunciar la u
  \item Inciso
  \item Nexo
\end{multicols}
\end{enumerate}
\vspace{0.15cm}
\item ¿Cuál es el uso correcto de la diéresis en: "desagüe"?
\begin{enumerate}[label=\textbf{\Alph*)}]
\begin{multicols}{2}
  \item Pronunciar la u
  \item Inciso
  \item Nexo
\end{multicols}
\end{enumerate}
\vspace{0.15cm}
\item ¿Cuál es el uso correcto de la coma en: "Pedro, el médico, llegó"?
\begin{enumerate}[label=\textbf{\Alph*)}]
\begin{multicols}{2}
  \item Vocativo
  \item Inciso
  \item Nexo
\end{multicols}
\end{enumerate}
\vspace{0.15cm}
\item ¿Cuál es el uso correcto de la coma en: "Compré pan, leche y huevos"?
\begin{enumerate}[label=\textbf{\Alph*)}]
\begin{multicols}{2}
  \item Vocativo
  \item Inciso
  \item Nexo
\end{multicols}
\end{enumerate}
\vspace{0.15cm}
\item ¿Cuál es el uso correcto de la coma en: "María, te llaman"?
\begin{enumerate}[label=\textbf{\Alph*)}]
\begin{multicols}{2}
  \item Vocativo
  \item Inciso
  \item Nexo
\end{multicols}
\end{enumerate}
\vspace{0.15cm}
\item ¿Cuál es el uso correcto de la coma en: "Bogotá, capital de Colombia, es grande"?
\begin{enumerate}[label=\textbf{\Alph*)}]
\begin{multicols}{2}
  \item Vocativo
  \item Inciso
  \item Nexo
\end{multicols}
\end{enumerate}
\vspace{0.15cm}
\item ¿Cuál es el uso correcto de la coma en: "Por lo tanto, no iré"?
\begin{enumerate}[label=\textbf{\Alph*)}]
\begin{multicols}{2}
  \item Vocativo
  \item Inciso
  \item Nexo
\end{multicols}
\end{enumerate}
\vspace{0.15cm}
\item ¿Cuál es el uso correcto del punto y coma en: "Unos cantan; otros, bailan"?
\begin{enumerate}[label=\textbf{\Alph*)}]
\begin{multicols}{2}
  \item Separar oraciones
  \item Inciso
  \item Nexo
\end{multicols}
\end{enumerate}
\vspace{0.15cm}
\item ¿Cuál es el uso correcto de los dos puntos en: "Tengo dos opciones: quedarme o irme"?
\begin{enumerate}[label=\textbf{\Alph*)}]
\begin{multicols}{2}
  \item Inciso
  \item Nexo
\end{multicols}
\end{enumerate}
\vspace{0.15cm}
\item ¿Cuál es el uso correcto de los dos puntos en: "Explicó: 'No puedo'"?
\begin{enumerate}[label=\textbf{\Alph*)}]
\begin{multicols}{2}
  \item Inciso
  \item Nexo
\end{multicols}
\end{enumerate}
\vspace{0.15cm}
\item ¿Cuál es el uso correcto de los puntos suspensivos en: "Tal vez..."?
\begin{enumerate}[label=\textbf{\Alph*)}]
\begin{multicols}{2}
  \item Inciso
  \item Nexo
\end{multicols}
\end{enumerate}
\vspace{0.15cm}
\item ¿Cuál es el uso correcto de las comillas en: "Le dijo 'adiós'"?
\begin{enumerate}[label=\textbf{\Alph*)}]
\begin{multicols}{2}
  \item Inciso
  \item Nexo
\end{multicols}
\end{enumerate}
\vspace{0.15cm}
\item ¿Cuál es el uso correcto del guion en: "físico-químico"?
\begin{enumerate}[label=\textbf{\Alph*)}]
\begin{multicols}{2}
  \item Inciso
  \item Nexo
\end{multicols}
\end{enumerate}
\vspace{0.15cm}
\item ¿Cuál es el uso correcto de la raya en: "—Hola —saludó—"?
\begin{enumerate}[label=\textbf{\Alph*)}]
\begin{multicols}{2}
  \item Inciso
  \item Nexo
\end{multicols}
\end{enumerate}
\vspace{0.15cm}
\item ¿Cuál es el uso correcto del paréntesis en: "Vino (creo) tarde"?
\begin{enumerate}[label=\textbf{\Alph*)}]
\begin{multicols}{2}
  \item Inciso
  \item Nexo
\end{multicols}
\end{enumerate}
\vspace{0.15cm}
\item ¿Cuál es el uso correcto de los signos de interrogación en: "¿Qué hora es?"?
\begin{enumerate}[label=\textbf{\Alph*)}]
\begin{multicols}{2}
  \item Pregunta
  \item Inciso
  \item Nexo
\end{multicols}
\end{enumerate}
\vspace{0.15cm}
\end{enumerate}\n
\vspace{0.4cm}
\begin{tcolorbox}[colback=gray!5!white,colframe=gray!60!black,title=\textbf{SOLUCIONARIO -- 2.1 ORTOGRAFÍA Y GRAMÁTICA}]
\begin{multicols}{4}
\begin{enumerate}[label=\textbf{\arabic*.}]
  \item \textbf{B}
  \item \textbf{B}
  \item \textbf{A}
  \item \textbf{A}
  \item \textbf{A}
  \item \textbf{A}
  \item \textbf{A}
  \item \textbf{A}
  \item \textbf{A}
  \item \textbf{A}
  \item \textbf{A}
  \item \textbf{A}
  \item \textbf{A}
  \item \textbf{A}
  \item \textbf{A}
  \item \textbf{A}
  \item \textbf{A}
  \item \textbf{A}
  \item \textbf{A}
  \item \textbf{A}
  \item \textbf{A}
  \item \textbf{A}
  \item \textbf{A}
  \item \textbf{A}
  \item \textbf{A}
  \item \textbf{A}
  \item \textbf{A}
  \item \textbf{A}
  \item \textbf{A}
  \item \textbf{A}
  \item \textbf{A}
  \item \textbf{A}
  \item \textbf{A}
  \item \textbf{A}
  \item \textbf{A}
  \item \textbf{A}
  \item \textbf{A}
  \item \textbf{A}
  \item \textbf{A}
  \item \textbf{A}
  \item \textbf{A}
  \item \textbf{A}
  \item \textbf{A}
  \item \textbf{A}
  \item \textbf{A}
  \item \textbf{A}
  \item \textbf{A}
  \item \textbf{A}
  \item \textbf{A}
  \item \textbf{A}
  \item \textbf{A}
  \item \textbf{A}
  \item \textbf{A}
  \item \textbf{A}
  \item \textbf{A}
  \item \textbf{A}
  \item \textbf{A}
  \item \textbf{A}
  \item \textbf{A}
  \item \textbf{A}
  \item \textbf{A}
  \item \textbf{A}
  \item \textbf{A}
  \item \textbf{A}
  \item \textbf{A}
  \item \textbf{A}
  \item \textbf{A}
  \item \textbf{A}
  \item \textbf{A}
  \item \textbf{A}
  \item \textbf{A}
  \item \textbf{A}
  \item \textbf{A}
  \item \textbf{A}
  \item \textbf{A}
  \item \textbf{A}
  \item \textbf{A}
  \item \textbf{A}
  \item \textbf{A}
  \item \textbf{A}
  \item \textbf{A}
  \item \textbf{A}
  \item \textbf{A}
  \item \textbf{A}
  \item \textbf{A}
  \item \textbf{A}
  \item \textbf{A}
  \item \textbf{A}
  \item \textbf{A}
  \item \textbf{A}
  \item \textbf{A}
  \item \textbf{A}
  \item \textbf{A}
  \item \textbf{A}
  \item \textbf{A}
  \item \textbf{A}
  \item \textbf{A}
  \item \textbf{A}
  \item \textbf{A}
  \item \textbf{A}
  \item \textbf{A}
  \item \textbf{A}
  \item \textbf{A}
  \item \textbf{A}
  \item \textbf{A}
  \item \textbf{A}
  \item \textbf{A}
  \item \textbf{A}
  \item \textbf{A}
  \item \textbf{A}
  \item \textbf{A}
  \item \textbf{A}
  \item \textbf{A}
  \item \textbf{A}
  \item \textbf{A}
  \item \textbf{A}
  \item \textbf{A}
  \item \textbf{A}
  \item \textbf{A}
  \item \textbf{A}
  \item \textbf{A}
  \item \textbf{A}
  \item \textbf{A}
  \item \textbf{B}
  \item \textbf{A}
  \item \textbf{A}
  \item \textbf{A}
  \item \textbf{A}
  \item \textbf{A}
  \item \textbf{A}
  \item \textbf{A}
  \item \textbf{A}
  \item \textbf{A}
  \item \textbf{A}
  \item \textbf{A}
  \item \textbf{A}
  \item \textbf{A}
  \item \textbf{A}
  \item \textbf{B}
  \item \textbf{B}
  \item \textbf{B}
  \item \textbf{B}
  \item \textbf{B}
  \item \textbf{B}
  \item \textbf{B}
  \item \textbf{B}
  \item \textbf{A}
  \item \textbf{A}
  \item \textbf{A}
  \item \textbf{A}
  \item \textbf{A}
  \item \textbf{A}
  \item \textbf{A}
  \item \textbf{A}
  \item \textbf{A}
  \item \textbf{A}
  \item \textbf{A}
  \item \textbf{A}
  \item \textbf{A}
  \item \textbf{A}
  \item \textbf{A}
  \item \textbf{A}
  \item \textbf{A}
  \item \textbf{A}
  \item \textbf{A}
  \item \textbf{A}
  \item \textbf{A}
  \item \textbf{A}
  \item \textbf{A}
  \item \textbf{A}
  \item \textbf{A}
  \item \textbf{A}
  \item \textbf{A}
  \item \textbf{A}
  \item \textbf{A}
  \item \textbf{A}
  \item \textbf{A}
  \item \textbf{A}
  \item \textbf{A}
  \item \textbf{A}
  \item \textbf{A}
  \item \textbf{A}
  \item \textbf{A}
  \item \textbf{A}
  \item \textbf{A}
  \item \textbf{A}
  \item \textbf{A}
  \item \textbf{A}
  \item \textbf{A}
  \item \textbf{A}
  \item \textbf{A}
  \item \textbf{A}
  \item \textbf{A}
  \item \textbf{A}
  \item \textbf{A}
  \item \textbf{A}
  \item \textbf{A}
  \item \textbf{A}
  \item \textbf{A}
  \item \textbf{A}
  \item \textbf{B}
  \item \textbf{A}
  \item \textbf{A}
  \item \textbf{A}
  \item \textbf{A}
  \item \textbf{A}
  \item \textbf{B}
  \item \textbf{C}
  \item \textbf{A}
  \item \textbf{A}
  \item \textbf{A}
  \item \textbf{A}
  \item \textbf{A}
  \item \textbf{A}
  \item \textbf{A}
  \item \textbf{A}
  \item \textbf{A}
  \item \textbf{A}
  \item \textbf{A}
  \item \textbf{A}
  \item \textbf{A}
  \item \textbf{A}
  \item \textbf{A}
  \item \textbf{A}
  \item \textbf{A}
  \item \textbf{A}
  \item \textbf{A}
  \item \textbf{A}
  \item \textbf{A}
  \item \textbf{A}
  \item \textbf{A}
  \item \textbf{A}
  \item \textbf{A}
  \item \textbf{A}
  \item \textbf{A}
  \item \textbf{A}
  \item \textbf{A}
  \item \textbf{A}
  \item \textbf{A}
  \item \textbf{A}
  \item \textbf{A}
  \item \textbf{A}
  \item \textbf{A}
  \item \textbf{A}
  \item \textbf{A}
  \item \textbf{A}
  \item \textbf{A}
  \item \textbf{A}
  \item \textbf{A}
  \item \textbf{A}
  \item \textbf{B}
  \item \textbf{B}
  \item \textbf{C}
  \item \textbf{D}
  \item \textbf{E}
  \item \textbf{F}
  \item \textbf{F}
  \item \textbf{F}
  \item \textbf{C}
  \item \textbf{E}
  \item \textbf{A}
  \item \textbf{B}
  \item \textbf{C}
  \item \textbf{E}
  \item \textbf{F}
  \item \textbf{A}
  \item \textbf{B}
  \item \textbf{A}
  \item \textbf{A}
  \item \textbf{A}
  \item \textbf{A}
  \item \textbf{A}
  \item \textbf{A}
  \item \textbf{A}
  \item \textbf{A}
  \item \textbf{A}
  \item \textbf{A}
  \item \textbf{A}
  \item \textbf{A}
  \item \textbf{A}
  \item \textbf{A}
  \item \textbf{A}
  \item \textbf{A}
  \item \textbf{A}
  \item \textbf{A}
  \item \textbf{A}
  \item \textbf{C}
  \item \textbf{B}
  \item \textbf{A}
  \item \textbf{C}
  \item \textbf{E}
  \item \textbf{A}
  \item \textbf{B}
  \item \textbf{A}
  \item \textbf{A}
  \item \textbf{A}
  \item \textbf{A}
  \item \textbf{A}
  \item \textbf{A}
  \item \textbf{A}
\end{enumerate}
\end{multicols}
\end{tcolorbox}
\newpage\n
\subsection{2.2 COHERENCIA Y ESTRUCTURA TEXTUAL}

\subsubsection*{SECCIÓN I: CONCEPTOS BÁSICOS DE COHERENCIA Y COHESIÓN (1–40)}

\begin{enumerate}[label=\textbf{\arabic*.}]
\setcounter{enumi}{0}
\item ¿Qué es la coherencia textual?
\begin{enumerate}[label=\textbf{\Alph*)}]
\begin{multicols}{2}
  \item La propiedad que da unidad de sentido al texto
  \item La suma de oraciones sin relación
  \item La repetición de palabras
  \item La extensión del texto
  \item La cantidad de párrafos
\end{multicols}
\end{enumerate}
\vspace{0.15cm}
\item ¿Qué es la cohesión textual?
\begin{enumerate}[label=\textbf{\Alph*)}]
\begin{multicols}{2}
  \item La unión superficial entre elementos del texto mediante mecanismos lingüísticos
  \item La idea principal del texto
  \item La estructura de los párrafos
  \item La intención del autor
  \item La extensión del texto
\end{multicols}
\end{enumerate}
\vspace{0.15cm}
\item La coherencia se refiere a:
\begin{enumerate}[label=\textbf{\Alph*)}]
\begin{multicols}{2}
  \item La unidad temática y lógica del texto
  \item La ortografía del texto
  \item La puntuación del texto
  \item La caligrafía del texto
  \item La cantidad de palabras
  \item La sintaxis del texto
\end{multicols}
\end{enumerate}
\vspace{0.15cm}
\item La cohesión se manifiesta a través de:
\begin{enumerate}[label=\textbf{\Alph*)}]
\begin{multicols}{2}
  \item Solo signos de puntuación
  \item Solo la ortografía
  \item Solo la caligrafía
  \item Solo la extensión
  \item Solo el tema
\end{multicols}
\end{enumerate}
\vspace{0.15cm}
\item Un texto coherente debe tener:
\begin{enumerate}[label=\textbf{\Alph*)}]
\begin{multicols}{2}
  \item Unidad temática
  \item Muchas palabras
  \item Muchos párrafos
  \item Muchos conectores
  \item Mucha puntuación
  \item Mucha extensión
\end{multicols}
\end{enumerate}
\vspace{0.15cm}
\item La progresión temática se refiere a:
\begin{enumerate}[label=\textbf{\Alph*)}]
\begin{multicols}{2}
  \item La repetición de la misma idea
  \item La cantidad de párrafos
  \item La extensión del texto
  \item La cantidad de conectores
  \item La ortografía
\end{multicols}
\end{enumerate}
\vspace{0.15cm}
\item ¿Cuál de los siguientes elementos NO es un mecanismo de cohesión?
\begin{enumerate}[label=\textbf{\Alph*)}]
\begin{multicols}{2}
  \item Pronombres
  \item Sinónimos
  \item Repetición
\end{multicols}
\end{enumerate}
\vspace{0.15cm}
\item La unidad temática significa que:
\begin{enumerate}[label=\textbf{\Alph*)}]
\begin{multicols}{2}
  \item Todo el texto gira en torno a un mismo tema
\end{multicols}
\end{enumerate}
\vspace{0.15cm}
\item La coherencia semántica se refiere a:
\begin{enumerate}[label=\textbf{\Alph*)}]
\begin{multicols}{2}
  \item La relación de significado entre las partes del texto
  \item La ortografía del texto
  \item La puntuación del texto
  \item La caligrafía del texto
  \item La sintaxis del texto
  \item La extensión del texto
\end{multicols}
\end{enumerate}
\vspace{0.15cm}
\item La coherencia pragmática se refiere a:
\begin{enumerate}[label=\textbf{\Alph*)}]
\begin{multicols}{2}
  \item La adecuación del texto al contexto comunicativo
  \item La ortografía del texto
  \item La puntuación del texto
  \item La caligrafía del texto
  \item La sintaxis del texto
  \item La extensión del texto
\end{multicols}
\end{enumerate}
\vspace{0.15cm}
\item Un texto incoherente se caracteriza por:
\begin{enumerate}[label=\textbf{\Alph*)}]
\begin{multicols}{2}
\end{multicols}
\end{enumerate}
\vspace{0.15cm}
\item La cohesión léxica se logra mediante:
\begin{enumerate}[label=\textbf{\Alph*)}]
\begin{multicols}{2}
  \item Repetición, sinonimia, antonimia, hiperonimia
  \item Solo conectores
  \item Solo puntuación
  \item Solo ortografía
  \item Solo caligrafía
  \item Solo extensión
\end{multicols}
\end{enumerate}
\vspace{0.15cm}
\item La cohesión gramatical se logra mediante:
\begin{enumerate}[label=\textbf{\Alph*)}]
\begin{multicols}{2}
  \item Pronombres, elipsis, determinantes, concordancia
  \item Solo conectores
  \item Solo puntuación
  \item Solo ortografía
  \item Solo caligrafía
  \item Solo extensión
\end{multicols}
\end{enumerate}
\vspace{0.15cm}
\item ¿Qué es un conector?
\begin{enumerate}[label=\textbf{\Alph*)}]
\begin{multicols}{2}
  \item Un elemento que une ideas o partes del texto
  \item Un signo de puntuación
  \item Una palabra sin significado
  \item Un tipo de oración
  \item Un tipo de párrafo
  \item Un tipo de texto
\end{multicols}
\end{enumerate}
\vspace{0.15cm}
\item ¿Qué es la elipsis?
\begin{enumerate}[label=\textbf{\Alph*)}]
\begin{multicols}{2}
  \item Omisión de elementos que se sobreentienden
  \item Repetición de elementos
  \item Uso excesivo de conectores
  \item Uso excesivo de puntuación
  \item Uso excesivo de ortografía
  \item Uso excesivo de caligrafía
\end{multicols}
\end{enumerate}
\vspace{0.15cm}
\item La referencia anafórica se produce cuando:
\begin{enumerate}[label=\textbf{\Alph*)}]
\begin{multicols}{2}
  \item Un elemento remite a otro mencionado antes
  \item Un elemento remite a otro mencionado después
  \item No hay referencia
  \item Se repite la misma palabra
  \item Se usan sinónimos
  \item Se usan antónimos
\end{multicols}
\end{enumerate}
\vspace{0.15cm}
\item La referencia catafórica se produce cuando:
\begin{enumerate}[label=\textbf{\Alph*)}]
\begin{multicols}{2}
  \item Un elemento remite a otro mencionado después
  \item Un elemento remite a otro mencionado antes
  \item No hay referencia
  \item Se repite la misma palabra
  \item Se usan sinónimos
  \item Se usan antónimos
\end{multicols}
\end{enumerate}
\vspace{0.15cm}
\item La deixis se refiere a:
\begin{enumerate}[label=\textbf{\Alph*)}]
\begin{multicols}{2}
  \item La ortografía del texto
  \item La puntuación del texto
  \item La caligrafía del texto
  \item La sintaxis del texto
  \item La extensión del texto
\end{multicols}
\end{enumerate}
\vspace{0.15cm}
\item Un texto bien cohesionado:
\begin{enumerate}[label=\textbf{\Alph*)}]
\begin{multicols}{2}
  \item Usa adecuadamente los mecanismos de cohesión
  \item Usa muchos conectores sin criterio
  \item Repite palabras sin necesidad
  \item No usa pronombres
  \item No usa sinónimos
  \item No usa elipsis
\end{multicols}
\end{enumerate}
\vspace{0.15cm}
\item La repetición como mecanismo de cohesión:
\begin{enumerate}[label=\textbf{\Alph*)}]
\begin{multicols}{2}
  \item Puede ser útil si no es excesiva
  \item Siempre es incorrecta
  \item Siempre es correcta
  \item Nunca se usa
  \item Solo se usa en poesía
  \item Solo se usa en ciencia
\end{multicols}
\end{enumerate}
\vspace{0.15cm}
\item La sinonimia como mecanismo de cohesión:
\begin{enumerate}[label=\textbf{\Alph*)}]
\begin{multicols}{2}
  \item Nunca se usa
  \item Solo se usa en poesía
  \item Solo se usa en ciencia
  \item Solo se usa en periodismo
\end{multicols}
\end{enumerate}
\vspace{0.15cm}
\item La antonimia como mecanismo de cohesión:
\begin{enumerate}[label=\textbf{\Alph*)}]
\begin{multicols}{2}
  \item Relaciona ideas mediante palabras de significado opuesto
  \item Nunca se usa
  \item Solo se usa en poesía
  \item Solo se usa en ciencia
  \item Solo se usa en periodismo
\end{multicols}
\end{enumerate}
\vspace{0.15cm}
\item La hiperonimia como mecanismo de cohesión:
\begin{enumerate}[label=\textbf{\Alph*)}]
\begin{multicols}{2}
  \item Relaciona palabras de significado más general con otras más específicas
  \item Nunca se usa
  \item Solo se usa en poesía
  \item Solo se usa en ciencia
  \item Solo se usa en periodismo
\end{multicols}
\end{enumerate}
\vspace{0.15cm}
\item La hiponimia como mecanismo de cohesión:
\begin{enumerate}[label=\textbf{\Alph*)}]
\begin{multicols}{2}
  \item Relaciona palabras de significado más específico con otras más generales
  \item Nunca se usa
  \item Solo se usa en poesía
  \item Solo se usa en ciencia
  \item Solo se usa en periodismo
\end{multicols}
\end{enumerate}
\vspace{0.15cm}
\item Los marcadores textuales son:
\begin{enumerate}[label=\textbf{\Alph*)}]
\begin{multicols}{2}
  \item Signos de puntuación
  \item Tipos de oraciones
  \item Tipos de párrafos
  \item Tipos de textos
  \item Tipos de letras
\end{multicols}
\end{enumerate}
\vspace{0.15cm}
\item La adecuación textual se refiere a:
\begin{enumerate}[label=\textbf{\Alph*)}]
\begin{multicols}{2}
  \item La adaptación del texto al contexto y al destinatario
  \item La ortografía del texto
  \item La puntuación del texto
  \item La caligrafía del texto
  \item La sintaxis del texto
  \item La extensión del texto
\end{multicols}
\end{enumerate}
\vspace{0.15cm}
\item Un texto adecuado debe:
\begin{enumerate}[label=\textbf{\Alph*)}]
\begin{multicols}{2}
  \item Usar el registro lingüístico apropiado
  \item Usar siempre lenguaje formal
  \item Usar siempre lenguaje informal
  \item Usar siempre tecnicismos
  \item Usar siempre coloquialismos
  \item Usar siempre cultismos
\end{multicols}
\end{enumerate}
\vspace{0.15cm}
\item La estructura textual se refiere a:
\begin{enumerate}[label=\textbf{\Alph*)}]
\begin{multicols}{2}
  \item La organización de las ideas en el texto
  \item La ortografía del texto
  \item La puntuación del texto
  \item La caligrafía del texto
  \item La sintaxis del texto
  \item La extensión del texto
\end{multicols}
\end{enumerate}
\vspace{0.15cm}
\item La introducción de un texto tiene como función:
\begin{enumerate}[label=\textbf{\Alph*)}]
\begin{multicols}{2}
  \item Presentar el tema y captar la atención
  \item Repetir el tema
\end{multicols}
\end{enumerate}
\vspace{0.15cm}
\item El desarrollo de un texto tiene como función:
\begin{enumerate}[label=\textbf{\Alph*)}]
\begin{multicols}{2}
  \item Presentar el tema
  \item Repetir el tema
\end{multicols}
\end{enumerate}
\vspace{0.15cm}
\item La conclusión de un texto tiene como función:
\begin{enumerate}[label=\textbf{\Alph*)}]
\begin{multicols}{2}
  \item Sintetizar o cerrar las ideas presentadas
  \item Presentar el tema
  \item Repetir el tema
\end{multicols}
\end{enumerate}
\vspace{0.15cm}
\item Un párrafo es:
\begin{enumerate}[label=\textbf{\Alph*)}]
\begin{multicols}{2}
  \item Un conjunto de oraciones que desarrollan una idea
  \item Un conjunto de palabras sin relación
  \item Un conjunto de textos
  \item Un conjunto de párrafos
  \item Un conjunto de conectores
  \item Un conjunto de signos
\end{multicols}
\end{enumerate}
\vspace{0.15cm}
\item La idea principal de un párrafo:
\begin{enumerate}[label=\textbf{\Alph*)}]
\begin{multicols}{2}
\end{multicols}
\end{enumerate}
\vspace{0.15cm}
\item Las ideas secundarias de un párrafo:
\begin{enumerate}[label=\textbf{\Alph*)}]
\begin{multicols}{2}
  \item Son más importantes que la principal
  \item No tienen relación con la principal
  \item Son la conclusión
  \item Son la introducción
  \item Son el título
\end{multicols}
\end{enumerate}
\vspace{0.15cm}
\item La oración temática de un párrafo:
\begin{enumerate}[label=\textbf{\Alph*)}]
\begin{multicols}{2}
\end{multicols}
\end{enumerate}
\vspace{0.15cm}
\item Un texto expositivo tiene como propósito:
\begin{enumerate}[label=\textbf{\Alph*)}]
\begin{multicols}{2}
  \item Informar o explicar un tema
\end{multicols}
\end{enumerate}
\vspace{0.15cm}
\item Un texto argumentativo tiene como propósito:
\begin{enumerate}[label=\textbf{\Alph*)}]
\begin{multicols}{2}
  \item Informar o explicar un tema
\end{multicols}
\end{enumerate}
\vspace{0.15cm}
\item Un texto narrativo tiene como propósito:
\begin{enumerate}[label=\textbf{\Alph*)}]
\begin{multicols}{2}
  \item Informar o explicar un tema
\end{multicols}
\end{enumerate}
\vspace{0.15cm}
\item Un texto descriptivo tiene como propósito:
\begin{enumerate}[label=\textbf{\Alph*)}]
\begin{multicols}{2}
  \item Representar características de personas, objetos o lugares
  \item Informar o explicar un tema
\end{multicols}
\end{enumerate}
\vspace{0.15cm}
\item La tesis de un texto argumentativo es:
\begin{enumerate}[label=\textbf{\Alph*)}]
\begin{multicols}{2}
  \item La idea principal que se defiende
  \item La conclusión
  \item La introducción
  \item La idea secundaria ---
\end{multicols}
\end{enumerate}
\vspace{0.15cm}
\end{enumerate}\n
\subsubsection*{SECCIÓN II: COHESIÓN – MECANISMOS DE REFERENCIA (41–90)}

\begin{enumerate}[label=\textbf{\arabic*.}]
\setcounter{enumi}{40}
\item En "Juan llegó tarde. Él siempre llega tarde", el pronombre "él" es un mecanismo de:
\begin{enumerate}[label=\textbf{\Alph*)}]
\begin{multicols}{2}
  \item Referencia anafórica
  \item Referencia catafórica
  \item Sinonimia
  \item Hiperonimia
\end{multicols}
\end{enumerate}
\vspace{0.15cm}
\item En "Cuando lo vi, Pedro corría", el pronombre "lo" es un mecanismo de:
\begin{enumerate}[label=\textbf{\Alph*)}]
\begin{multicols}{2}
  \item Referencia catafórica
  \item Referencia anafórica
  \item Sinonimia
  \item Hiperonimia
\end{multicols}
\end{enumerate}
\vspace{0.15cm}
\item En "Compré flores. Las puse en un jarrón", el artículo "las" es un mecanismo de:
\begin{enumerate}[label=\textbf{\Alph*)}]
\begin{multicols}{2}
  \item Referencia anafórica
  \item Referencia catafórica
  \item Sinonimia
  \item Hiperonimia
\end{multicols}
\end{enumerate}
\vspace{0.15cm}
\item En "María estudia medicina. Ella quiere ser doctora", el pronombre "ella" es un mecanismo de:
\begin{enumerate}[label=\textbf{\Alph*)}]
\begin{multicols}{2}
  \item Referencia anafórica
  \item Referencia catafórica
  \item Sinonimia
  \item Hiperonimia
\end{multicols}
\end{enumerate}
\vspace{0.15cm}
\item En "Pedro y Ana llegaron. Ambos estaban cansados", el pronombre "ambos" es un mecanismo de:
\begin{enumerate}[label=\textbf{\Alph*)}]
\begin{multicols}{2}
  \item Referencia anafórica
  \item Referencia catafórica
  \item Sinonimia
  \item Hiperonimia
\end{multicols}
\end{enumerate}
\vspace{0.15cm}
\item En "Tengo un perro. El animal es muy fiel", "el animal" es un mecanismo de:
\begin{enumerate}[label=\textbf{\Alph*)}]
\begin{multicols}{2}
  \item Sinonimia
  \item Referencia catafórica
  \item Hiperonimia
  \item Hiponimia
\end{multicols}
\end{enumerate}
\vspace{0.15cm}
\item En "Compré un automóvil. El vehículo es rojo", "el vehículo" es un mecanismo de:
\begin{enumerate}[label=\textbf{\Alph*)}]
\begin{multicols}{2}
  \item Hiperonimia
  \item Referencia catafórica
  \item Sinonimia
  \item Hiponimia
\end{multicols}
\end{enumerate}
\vspace{0.15cm}
\item En "Llegó el niño. El pequeño sonreía", "el pequeño" es un mecanismo de:
\begin{enumerate}[label=\textbf{\Alph*)}]
\begin{multicols}{2}
  \item Sinonimia
  \item Referencia catafórica
  \item Hiperonimia
  \item Hiponimia
\end{multicols}
\end{enumerate}
\vspace{0.15cm}
\item En "Ana compró rosas. Las flores eran hermosas", "las flores" es un mecanismo de:
\begin{enumerate}[label=\textbf{\Alph*)}]
\begin{multicols}{2}
  \item Hiperonimia
  \item Referencia catafórica
  \item Sinonimia
  \item Hiponimia
\end{multicols}
\end{enumerate}
\vspace{0.15cm}
\item En "El árbol creció. El roble es fuerte", "el roble" es un mecanismo de:
\begin{enumerate}[label=\textbf{\Alph*)}]
\begin{multicols}{2}
  \item Hiponimia
  \item Referencia catafórica
  \item Sinonimia
  \item Hiperonimia
\end{multicols}
\end{enumerate}
\vspace{0.15cm}
\item En "Pedro estudia y trabaja. También ayuda en casa", "también" es un mecanismo de:
\begin{enumerate}[label=\textbf{\Alph*)}]
\begin{multicols}{2}
  \item Referencia catafórica
  \item Sinonimia
  \item Hiperonimia
\end{multicols}
\end{enumerate}
\vspace{0.15cm}
\item En "No estudié. Por eso reprobé", "por eso" es un mecanismo de:
\begin{enumerate}[label=\textbf{\Alph*)}]
\begin{multicols}{2}
  \item Referencia catafórica
  \item Sinonimia
  \item Hiperonimia
\end{multicols}
\end{enumerate}
\vspace{0.15cm}
\item En "Llegué tarde. Sin embargo, entré", "sin embargo" es un mecanismo de:
\begin{enumerate}[label=\textbf{\Alph*)}]
\begin{multicols}{2}
  \item Referencia catafórica
  \item Sinonimia
  \item Hiperonimia
\end{multicols}
\end{enumerate}
\vspace{0.15cm}
\item En "Primero, estudia. Luego, descansa", "primero" y "luego" son:
\begin{enumerate}[label=\textbf{\Alph*)}]
\begin{multicols}{2}
  \item Referencia catafórica
  \item Sinonimia
  \item Hiperonimia
\end{multicols}
\end{enumerate}
\vspace{0.15cm}
\item En "En conclusión, todo salió bien", "en conclusión" es un:
\begin{enumerate}[label=\textbf{\Alph*)}]
\begin{multicols}{2}
  \item Referencia catafórica
  \item Sinonimia
  \item Hiperonimia
\end{multicols}
\end{enumerate}
\vspace{0.15cm}
\item En "Por ejemplo, los mamíferos son vertebrados", "por ejemplo" es un:
\begin{enumerate}[label=\textbf{\Alph*)}]
\begin{multicols}{2}
  \item Referencia catafórica
  \item Sinonimia
  \item Hiperonimia
\end{multicols}
\end{enumerate}
\vspace{0.15cm}
\item En "Es decir, no vendrás", "es decir" es un:
\begin{enumerate}[label=\textbf{\Alph*)}]
\begin{multicols}{2}
  \item Referencia catafórica
  \item Sinonimia
  \item Hiperonimia
\end{multicols}
\end{enumerate}
\vspace{0.15cm}
\item En "Además, trajo postres", "además" es un:
\begin{enumerate}[label=\textbf{\Alph*)}]
\begin{multicols}{2}
  \item Referencia catafórica
  \item Sinonimia
  \item Hiperonimia
\end{multicols}
\end{enumerate}
\vspace{0.15cm}
\item En "No obstante, insistió", "no obstante" es un:
\begin{enumerate}[label=\textbf{\Alph*)}]
\begin{multicols}{2}
  \item Referencia catafórica
  \item Sinonimia
  \item Hiperonimia
\end{multicols}
\end{enumerate}
\vspace{0.15cm}
\item En "Por lo tanto, aprobó", "por lo tanto" es un:
\begin{enumerate}[label=\textbf{\Alph*)}]
\begin{multicols}{2}
  \item Referencia catafórica
  \item Sinonimia
  \item Hiperonimia
\end{multicols}
\end{enumerate}
\vspace{0.15cm}
\item En "Mientras tanto, esperaba", "mientras tanto" es un:
\begin{enumerate}[label=\textbf{\Alph*)}]
\begin{multicols}{2}
  \item Referencia catafórica
  \item Sinonimia
  \item Hiperonimia
\end{multicols}
\end{enumerate}
\vspace{0.15cm}
\item En "En cambio, él no fue", "en cambio" es un:
\begin{enumerate}[label=\textbf{\Alph*)}]
\begin{multicols}{2}
  \item Referencia catafórica
  \item Sinonimia
  \item Hiperonimia
\end{multicols}
\end{enumerate}
\vspace{0.15cm}
\item En "Así mismo, lo haré", "así mismo" es un:
\begin{enumerate}[label=\textbf{\Alph*)}]
\begin{multicols}{2}
  \item Referencia catafórica
  \item Sinonimia
  \item Hiperonimia
\end{multicols}
\end{enumerate}
\vspace{0.15cm}
\item En "Por consiguiente, renunció", "por consiguiente" es un:
\begin{enumerate}[label=\textbf{\Alph*)}]
\begin{multicols}{2}
  \item Referencia catafórica
  \item Sinonimia
  \item Hiperonimia
\end{multicols}
\end{enumerate}
\vspace{0.15cm}
\item En "Finalmente, llegó", "finalmente" es un:
\begin{enumerate}[label=\textbf{\Alph*)}]
\begin{multicols}{2}
  \item Referencia catafórica
  \item Sinonimia
  \item Hiperonimia
\end{multicols}
\end{enumerate}
\vspace{0.15cm}
\item En "En primer lugar, hablemos del tema", "en primer lugar" es un:
\begin{enumerate}[label=\textbf{\Alph*)}]
\begin{multicols}{2}
  \item Referencia catafórica
  \item Sinonimia
  \item Hiperonimia
\end{multicols}
\end{enumerate}
\vspace{0.15cm}
\item En "Por otro lado, existe otra teoría", "por otro lado" es un:
\begin{enumerate}[label=\textbf{\Alph*)}]
\begin{multicols}{2}
  \item Referencia catafórica
  \item Sinonimia
  \item Hiperonimia
\end{multicols}
\end{enumerate}
\vspace{0.15cm}
\item En "En resumen, todo fue positivo", "en resumen" es un:
\begin{enumerate}[label=\textbf{\Alph*)}]
\begin{multicols}{2}
  \item Referencia catafórica
  \item Sinonimia
  \item Hiperonimia
\end{multicols}
\end{enumerate}
\vspace{0.15cm}
\item En "Dado que llovía, no salimos", "dado que" es un:
\begin{enumerate}[label=\textbf{\Alph*)}]
\begin{multicols}{2}
  \item Referencia catafórica
  \item Sinonimia
  \item Hiperonimia
\end{multicols}
\end{enumerate}
\vspace{0.15cm}
\item En "Aunque insistió, no cedí", "aunque" es un:
\begin{enumerate}[label=\textbf{\Alph*)}]
\begin{multicols}{2}
  \item Referencia catafórica
  \item Sinonimia
  \item Hiperonimia
\end{multicols}
\end{enumerate}
\vspace{0.15cm}
\item En "Si estudias, aprobarás", "si" es un:
\begin{enumerate}[label=\textbf{\Alph*)}]
\begin{multicols}{2}
  \item Referencia catafórica
  \item Sinonimia
  \item Hiperonimia
\end{multicols}
\end{enumerate}
\vspace{0.15cm}
\item En "Para que entiendas, te explico", "para que" es un:
\begin{enumerate}[label=\textbf{\Alph*)}]
\begin{multicols}{2}
  \item Referencia catafórica
  \item Sinonimia
  \item Hiperonimia
\end{multicols}
\end{enumerate}
\vspace{0.15cm}
\item En "Tan pronto como llegó, se fue", "tan pronto como" es un:
\begin{enumerate}[label=\textbf{\Alph*)}]
\begin{multicols}{2}
  \item Referencia catafórica
  \item Sinonimia
  \item Hiperonimia
\end{multicols}
\end{enumerate}
\vspace{0.15cm}
\item En "Ora estudia, ora trabaja", "ora... ora" es un:
\begin{enumerate}[label=\textbf{\Alph*)}]
\begin{multicols}{2}
  \item Referencia catafórica
  \item Sinonimia
  \item Hiperonimia
\end{multicols}
\end{enumerate}
\vspace{0.15cm}
\item En "Es decir, no es cierto", "es decir" es un:
\begin{enumerate}[label=\textbf{\Alph*)}]
\begin{multicols}{2}
  \item Referencia catafórica
  \item Sinonimia
  \item Hiperonimia
\end{multicols}
\end{enumerate}
\vspace{0.15cm}
\item En "Mejor dicho, no iré", "mejor dicho" es un:
\begin{enumerate}[label=\textbf{\Alph*)}]
\begin{multicols}{2}
  \item Referencia catafórica
  \item Sinonimia
  \item Hiperonimia
\end{multicols}
\end{enumerate}
\vspace{0.15cm}
\item En "O sea, no estás de acuerdo", "o sea" es un:
\begin{enumerate}[label=\textbf{\Alph*)}]
\begin{multicols}{2}
  \item Referencia catafórica
  \item Sinonimia
  \item Hiperonimia
\end{multicols}
\end{enumerate}
\vspace{0.15cm}
\item En "En efecto, así fue", "en efecto" es un:
\begin{enumerate}[label=\textbf{\Alph*)}]
\begin{multicols}{2}
  \item Referencia catafórica
  \item Sinonimia
  \item Hiperonimia
\end{multicols}
\end{enumerate}
\vspace{0.15cm}
\item En "Por supuesto, iré", "por supuesto" es un:
\begin{enumerate}[label=\textbf{\Alph*)}]
\begin{multicols}{2}
  \item Referencia catafórica
  \item Sinonimia
  \item Hiperonimia
\end{multicols}
\end{enumerate}
\vspace{0.15cm}
\item En "Sin duda, es correcto", "sin duda" es un:
\begin{enumerate}[label=\textbf{\Alph*)}]
\begin{multicols}{2}
  \item Referencia catafórica
  \item Sinonimia
  \item Hiperonimia
\end{multicols}
\end{enumerate}
\vspace{0.15cm}
\item En "Quizás venga", "quizás" es un:
\begin{enumerate}[label=\textbf{\Alph*)}]
\begin{multicols}{2}
  \item Referencia catafórica
  \item Sinonimia
  \item Hiperonimia
\end{multicols}
\end{enumerate}
\vspace{0.15cm}
\item En "Tal vez llueva", "tal vez" es un:
\begin{enumerate}[label=\textbf{\Alph*)}]
\begin{multicols}{2}
  \item Referencia catafórica
  \item Sinonimia
  \item Hiperonimia
\end{multicols}
\end{enumerate}
\vspace{0.15cm}
\item En "Posiblemente llegue tarde", "posiblemente" es un:
\begin{enumerate}[label=\textbf{\Alph*)}]
\begin{multicols}{2}
  \item Referencia catafórica
  \item Sinonimia
  \item Hiperonimia
\end{multicols}
\end{enumerate}
\vspace{0.15cm}
\item En "Obviamente, ganó", "obviamente" es un:
\begin{enumerate}[label=\textbf{\Alph*)}]
\begin{multicols}{2}
  \item Referencia catafórica
  \item Sinonimia
  \item Hiperonimia
\end{multicols}
\end{enumerate}
\vspace{0.15cm}
\item En "Evidentemente, es así", "evidentemente" es un:
\begin{enumerate}[label=\textbf{\Alph*)}]
\begin{multicols}{2}
  \item Referencia catafórica
  \item Sinonimia
  \item Hiperonimia
\end{multicols}
\end{enumerate}
\vspace{0.15cm}
\item En "Claramente, no es correcto", "claramente" es un:
\begin{enumerate}[label=\textbf{\Alph*)}]
\begin{multicols}{2}
  \item Referencia catafórica
  \item Sinonimia
  \item Hiperonimia
\end{multicols}
\end{enumerate}
\vspace{0.15cm}
\item En "Resumiendo, todo está bien", "resumiendo" es un:
\begin{enumerate}[label=\textbf{\Alph*)}]
\begin{multicols}{2}
  \item Referencia catafórica
  \item Sinonimia
  \item Hiperonimia
\end{multicols}
\end{enumerate}
\vspace{0.15cm}
\item En "En síntesis, es viable", "en síntesis" es un:
\begin{enumerate}[label=\textbf{\Alph*)}]
\begin{multicols}{2}
  \item Referencia catafórica
  \item Sinonimia
  \item Hiperonimia
\end{multicols}
\end{enumerate}
\vspace{0.15cm}
\item En "Para finalizar, gracias", "para finalizar" es un:
\begin{enumerate}[label=\textbf{\Alph*)}]
\begin{multicols}{2}
  \item Referencia catafórica
  \item Sinonimia
  \item Hiperonimia
\end{multicols}
\end{enumerate}
\vspace{0.15cm}
\item En "Por último, recordemos", "por último" es un:
\begin{enumerate}[label=\textbf{\Alph*)}]
\begin{multicols}{2}
  \item Referencia catafórica
  \item Sinonimia
  \item Hiperonimia ---
\end{multicols}
\end{enumerate}
\vspace{0.15cm}
\end{enumerate}\n
\subsubsection*{SECCIÓN III: CONECTORES Y MARCADORES TEXTUALES (91–150)}

\begin{enumerate}[label=\textbf{\arabic*.}]
\setcounter{enumi}{90}
\item ¿Qué tipo de conector es "y"?
\begin{enumerate}[label=\textbf{\Alph*)}]
\begin{multicols}{2}
  \item Temporal
\end{multicols}
\end{enumerate}
\vspace{0.15cm}
\item ¿Qué tipo de conector es "pero"?
\begin{enumerate}[label=\textbf{\Alph*)}]
\begin{multicols}{2}
  \item Temporal
\end{multicols}
\end{enumerate}
\vspace{0.15cm}
\item ¿Qué tipo de conector es "porque"?
\begin{enumerate}[label=\textbf{\Alph*)}]
\begin{multicols}{2}
  \item Temporal
\end{multicols}
\end{enumerate}
\vspace{0.15cm}
\item ¿Qué tipo de conector es "por lo tanto"?
\begin{enumerate}[label=\textbf{\Alph*)}]
\begin{multicols}{2}
  \item Temporal
\end{multicols}
\end{enumerate}
\vspace{0.15cm}
\item ¿Qué tipo de conector es "mientras"?
\begin{enumerate}[label=\textbf{\Alph*)}]
\begin{multicols}{2}
  \item Temporal
\end{multicols}
\end{enumerate}
\vspace{0.15cm}
\item ¿Qué tipo de conector es "si"?
\begin{enumerate}[label=\textbf{\Alph*)}]
\begin{multicols}{2}
  \item Temporal
\end{multicols}
\end{enumerate}
\vspace{0.15cm}
\item ¿Qué tipo de conector es "además"?
\begin{enumerate}[label=\textbf{\Alph*)}]
\begin{multicols}{2}
  \item Temporal
\end{multicols}
\end{enumerate}
\vspace{0.15cm}
\item ¿Qué tipo de conector es "sin embargo"?
\begin{enumerate}[label=\textbf{\Alph*)}]
\begin{multicols}{2}
  \item Temporal
\end{multicols}
\end{enumerate}
\vspace{0.15cm}
\item ¿Qué tipo de conector es "puesto que"?
\begin{enumerate}[label=\textbf{\Alph*)}]
\begin{multicols}{2}
  \item Temporal
\end{multicols}
\end{enumerate}
\vspace{0.15cm}
\item ¿Qué tipo de conector es "en consecuencia"?
\begin{enumerate}[label=\textbf{\Alph*)}]
\begin{multicols}{2}
  \item Temporal
\end{multicols}
\end{enumerate}
\vspace{0.15cm}
\item ¿Qué tipo de conector es "antes de que"?
\begin{enumerate}[label=\textbf{\Alph*)}]
\begin{multicols}{2}
  \item Temporal
\end{multicols}
\end{enumerate}
\vspace{0.15cm}
\item ¿Qué tipo de conector es "a menos que"?
\begin{enumerate}[label=\textbf{\Alph*)}]
\begin{multicols}{2}
  \item Temporal
\end{multicols}
\end{enumerate}
\vspace{0.15cm}
\item ¿Qué tipo de conector es "también"?
\begin{enumerate}[label=\textbf{\Alph*)}]
\begin{multicols}{2}
  \item Temporal
\end{multicols}
\end{enumerate}
\vspace{0.15cm}
\item ¿Qué tipo de conector es "no obstante"?
\begin{enumerate}[label=\textbf{\Alph*)}]
\begin{multicols}{2}
  \item Temporal
\end{multicols}
\end{enumerate}
\vspace{0.15cm}
\item ¿Qué tipo de conector es "ya que"?
\begin{enumerate}[label=\textbf{\Alph*)}]
\begin{multicols}{2}
  \item Temporal
\end{multicols}
\end{enumerate}
\vspace{0.15cm}
\item ¿Qué tipo de conector es "así que"?
\begin{enumerate}[label=\textbf{\Alph*)}]
\begin{multicols}{2}
  \item Temporal
\end{multicols}
\end{enumerate}
\vspace{0.15cm}
\item ¿Qué tipo de conector es "después de que"?
\begin{enumerate}[label=\textbf{\Alph*)}]
\begin{multicols}{2}
  \item Temporal
\end{multicols}
\end{enumerate}
\vspace{0.15cm}
\item ¿Qué tipo de conector es "con tal de que"?
\begin{enumerate}[label=\textbf{\Alph*)}]
\begin{multicols}{2}
  \item Temporal
\end{multicols}
\end{enumerate}
\vspace{0.15cm}
\item ¿Qué tipo de conector es "incluso"?
\begin{enumerate}[label=\textbf{\Alph*)}]
\begin{multicols}{2}
  \item Temporal
\end{multicols}
\end{enumerate}
\vspace{0.15cm}
\item ¿Qué tipo de conector es "aunque"?
\begin{enumerate}[label=\textbf{\Alph*)}]
\begin{multicols}{2}
  \item Temporal
\end{multicols}
\end{enumerate}
\vspace{0.15cm}
\item ¿Qué tipo de conector es "dado que"?
\begin{enumerate}[label=\textbf{\Alph*)}]
\begin{multicols}{2}
  \item Temporal
\end{multicols}
\end{enumerate}
\vspace{0.15cm}
\item ¿Qué tipo de conector es "por consiguiente"?
\begin{enumerate}[label=\textbf{\Alph*)}]
\begin{multicols}{2}
  \item Temporal
\end{multicols}
\end{enumerate}
\vspace{0.15cm}
\item ¿Qué tipo de conector es "tan pronto como"?
\begin{enumerate}[label=\textbf{\Alph*)}]
\begin{multicols}{2}
  \item Temporal
\end{multicols}
\end{enumerate}
\vspace{0.15cm}
\item ¿Qué tipo de conector es "siempre que"?
\begin{enumerate}[label=\textbf{\Alph*)}]
\begin{multicols}{2}
  \item Temporal
\end{multicols}
\end{enumerate}
\vspace{0.15cm}
\item ¿Qué tipo de conector es "encima"?
\begin{enumerate}[label=\textbf{\Alph*)}]
\begin{multicols}{2}
  \item Temporal
\end{multicols}
\end{enumerate}
\vspace{0.15cm}
\item ¿Qué tipo de conector es "empero"?
\begin{enumerate}[label=\textbf{\Alph*)}]
\begin{multicols}{2}
  \item Temporal
\end{multicols}
\end{enumerate}
\vspace{0.15cm}
\item ¿Qué tipo de conector es "debido a que"?
\begin{enumerate}[label=\textbf{\Alph*)}]
\begin{multicols}{2}
  \item Temporal
\end{multicols}
\end{enumerate}
\vspace{0.15cm}
\item ¿Qué tipo de conector es "de modo que"?
\begin{enumerate}[label=\textbf{\Alph*)}]
\begin{multicols}{2}
  \item Temporal
\end{multicols}
\end{enumerate}
\vspace{0.15cm}
\item ¿Qué tipo de conector es "en cuanto"?
\begin{enumerate}[label=\textbf{\Alph*)}]
\begin{multicols}{2}
  \item Temporal
\end{multicols}
\end{enumerate}
\vspace{0.15cm}
\item ¿Qué tipo de conector es "a condición de que"?
\begin{enumerate}[label=\textbf{\Alph*)}]
\begin{multicols}{2}
  \item Temporal
\end{multicols}
\end{enumerate}
\vspace{0.15cm}
\item ¿Qué tipo de conector es "asimismo"?
\begin{enumerate}[label=\textbf{\Alph*)}]
\begin{multicols}{2}
  \item Temporal
\end{multicols}
\end{enumerate}
\vspace{0.15cm}
\item ¿Qué tipo de conector es "mas"?
\begin{enumerate}[label=\textbf{\Alph*)}]
\begin{multicols}{2}
  \item Temporal
\end{multicols}
\end{enumerate}
\vspace{0.15cm}
\item ¿Qué tipo de conector es "pues"?
\begin{enumerate}[label=\textbf{\Alph*)}]
\begin{multicols}{2}
  \item Temporal
\end{multicols}
\end{enumerate}
\vspace{0.15cm}
\item ¿Qué tipo de conector es "luego"?
\begin{enumerate}[label=\textbf{\Alph*)}]
\begin{multicols}{2}
  \item Temporal
\end{multicols}
\end{enumerate}
\vspace{0.15cm}
\item ¿Qué tipo de conector es "ahora bien"?
\begin{enumerate}[label=\textbf{\Alph*)}]
\begin{multicols}{2}
  \item Temporal
\end{multicols}
\end{enumerate}
\vspace{0.15cm}
\item ¿Qué tipo de conector es "por el contrario"?
\begin{enumerate}[label=\textbf{\Alph*)}]
\begin{multicols}{2}
  \item Temporal
\end{multicols}
\end{enumerate}
\vspace{0.15cm}
\item ¿Qué tipo de conector es "en efecto"?
\begin{enumerate}[label=\textbf{\Alph*)}]
\begin{multicols}{2}
\end{multicols}
\end{enumerate}
\vspace{0.15cm}
\item ¿Qué tipo de conector es "en resumen"?
\begin{enumerate}[label=\textbf{\Alph*)}]
\begin{multicols}{2}
  \item Temporal
\end{multicols}
\end{enumerate}
\vspace{0.15cm}
\item ¿Qué tipo de conector es "por ejemplo"?
\begin{enumerate}[label=\textbf{\Alph*)}]
\begin{multicols}{2}
\end{multicols}
\end{enumerate}
\vspace{0.15cm}
\item ¿Qué tipo de conector es "es decir"?
\begin{enumerate}[label=\textbf{\Alph*)}]
\begin{multicols}{2}
\end{multicols}
\end{enumerate}
\vspace{0.15cm}
\item ¿Qué tipo de conector es "en primer lugar"?
\begin{enumerate}[label=\textbf{\Alph*)}]
\begin{multicols}{2}
  \item Ordenador
\end{multicols}
\end{enumerate}
\vspace{0.15cm}
\item ¿Qué tipo de conector es "finalmente"?
\begin{enumerate}[label=\textbf{\Alph*)}]
\begin{multicols}{2}
\end{multicols}
\end{enumerate}
\vspace{0.15cm}
\item ¿Qué tipo de conector es "por supuesto"?
\begin{enumerate}[label=\textbf{\Alph*)}]
\begin{multicols}{2}
\end{multicols}
\end{enumerate}
\vspace{0.15cm}
\item ¿Qué tipo de conector es "sin duda"?
\begin{enumerate}[label=\textbf{\Alph*)}]
\begin{multicols}{2}
\end{multicols}
\end{enumerate}
\vspace{0.15cm}
\item ¿Qué tipo de conector es "quizás"?
\begin{enumerate}[label=\textbf{\Alph*)}]
\begin{multicols}{2}
\end{multicols}
\end{enumerate}
\vspace{0.15cm}
\item ¿Qué tipo de conector es "tal vez"?
\begin{enumerate}[label=\textbf{\Alph*)}]
\begin{multicols}{2}
\end{multicols}
\end{enumerate}
\vspace{0.15cm}
\item ¿Qué tipo de conector es "obviamente"?
\begin{enumerate}[label=\textbf{\Alph*)}]
\begin{multicols}{2}
\end{multicols}
\end{enumerate}
\vspace{0.15cm}
\item ¿Qué tipo de conector es "evidentemente"?
\begin{enumerate}[label=\textbf{\Alph*)}]
\begin{multicols}{2}
\end{multicols}
\end{enumerate}
\vspace{0.15cm}
\item ¿Qué tipo de conector es "en síntesis"?
\begin{enumerate}[label=\textbf{\Alph*)}]
\begin{multicols}{2}
  \item Temporal
\end{multicols}
\end{enumerate}
\vspace{0.15cm}
\item ¿Qué tipo de conector es "para finalizar"?
\begin{enumerate}[label=\textbf{\Alph*)}]
\begin{multicols}{2}
\end{multicols}
\end{enumerate}
\vspace{0.15cm}
\item ¿Qué tipo de conector es "por último"?
\begin{enumerate}[label=\textbf{\Alph*)}]
\begin{multicols}{2}
\end{multicols}
\end{enumerate}
\vspace{0.15cm}
\item ¿Qué tipo de conector es "en cambio"?
\begin{enumerate}[label=\textbf{\Alph*)}]
\begin{multicols}{2}
  \item Temporal
\end{multicols}
\end{enumerate}
\vspace{0.15cm}
\item ¿Qué tipo de conector es "mientras tanto"?
\begin{enumerate}[label=\textbf{\Alph*)}]
\begin{multicols}{2}
  \item Temporal
\end{multicols}
\end{enumerate}
\vspace{0.15cm}
\item ¿Qué tipo de conector es "a fin de que"?
\begin{enumerate}[label=\textbf{\Alph*)}]
\begin{multicols}{2}
\end{multicols}
\end{enumerate}
\vspace{0.15cm}
\item ¿Qué tipo de conector es "con el propósito de"?
\begin{enumerate}[label=\textbf{\Alph*)}]
\begin{multicols}{2}
\end{multicols}
\end{enumerate}
\vspace{0.15cm}
\item ¿Qué tipo de conector es "a pesar de que"?
\begin{enumerate}[label=\textbf{\Alph*)}]
\begin{multicols}{2}
  \item Temporal
\end{multicols}
\end{enumerate}
\vspace{0.15cm}
\item ¿Qué tipo de conector es "por más que"?
\begin{enumerate}[label=\textbf{\Alph*)}]
\begin{multicols}{2}
  \item Temporal
\end{multicols}
\end{enumerate}
\vspace{0.15cm}
\item ¿Qué tipo de conector es "si bien"?
\begin{enumerate}[label=\textbf{\Alph*)}]
\begin{multicols}{2}
  \item Temporal
\end{multicols}
\end{enumerate}
\vspace{0.15cm}
\item ¿Qué tipo de conector es "así pues"?
\begin{enumerate}[label=\textbf{\Alph*)}]
\begin{multicols}{2}
  \item Temporal
\end{multicols}
\end{enumerate}
\vspace{0.15cm}
\item ¿Qué tipo de conector es "de ahí que"?
\begin{enumerate}[label=\textbf{\Alph*)}]
\begin{multicols}{2}
  \item Temporal
\end{multicols}
\end{enumerate}
\vspace{0.15cm}
\end{enumerate}\n
\subsubsection*{SECCIÓN IV: ESTRUCTURA DE TEXTOS EXPOSITIVOS (151–200)}

\begin{enumerate}[label=\textbf{\arabic*.}]
\setcounter{enumi}{150}
\item Un texto expositivo tiene como función principal:
\begin{enumerate}[label=\textbf{\Alph*)}]
\begin{multicols}{2}
  \item Informar y explicar un tema
\end{multicols}
\end{enumerate}
\vspace{0.15cm}
\item La estructura básica de un texto expositivo es:
\begin{enumerate}[label=\textbf{\Alph*)}]
\begin{multicols}{2}
  \item Introducción, desarrollo y conclusión
  \item Planteamiento, nudo y desenlace
  \item Tesis, argumentos y conclusión
  \item Inicio, conflicto y resolución
  \item Título, cuerpo y firma
\end{multicols}
\end{enumerate}
\vspace{0.15cm}
\item La introducción de un texto expositivo debe:
\begin{enumerate}[label=\textbf{\Alph*)}]
\begin{multicols}{2}
  \item Presentar el tema y despertar el interés
  \item Repetir el título
\end{multicols}
\end{enumerate}
\vspace{0.15cm}
\item El desarrollo de un texto expositivo debe:
\begin{enumerate}[label=\textbf{\Alph*)}]
\begin{multicols}{2}
  \item Presentar solo el tema
  \item Repetir la introducción
\end{multicols}
\end{enumerate}
\vspace{0.15cm}
\item La conclusión de un texto expositivo debe:
\begin{enumerate}[label=\textbf{\Alph*)}]
\begin{multicols}{2}
  \item Sintetizar la información presentada
  \item Presentar el tema
  \item Repetir la introducción
\end{multicols}
\end{enumerate}
\vspace{0.15cm}
\item ¿Cuál de los siguientes es un texto expositivo?
\begin{enumerate}[label=\textbf{\Alph*)}]
\begin{multicols}{2}
  \item Un artículo de enciclopedia
  \item Un cuento
  \item Un poema
  \item Una novela
  \item Una obra de teatro
  \item Una fábula
\end{multicols}
\end{enumerate}
\vspace{0.15cm}
\item ¿Cuál de los siguientes es un texto expositivo?
\begin{enumerate}[label=\textbf{\Alph*)}]
\begin{multicols}{2}
  \item Un manual de instrucciones
  \item Un cuento
  \item Un poema
  \item Una novela
  \item Una obra de teatro
  \item Una fábula
\end{multicols}
\end{enumerate}
\vspace{0.15cm}
\item ¿Cuál de los siguientes es un texto expositivo?
\begin{enumerate}[label=\textbf{\Alph*)}]
\begin{multicols}{2}
  \item Un texto científico
  \item Un cuento
  \item Un poema
  \item Una novela
  \item Una obra de teatro
  \item Una fábula
\end{multicols}
\end{enumerate}
\vspace{0.15cm}
\item La exposición puede ser:
\begin{enumerate}[label=\textbf{\Alph*)}]
\begin{multicols}{2}
  \item Solo divulgativa
  \item Solo especializada
  \item Solo oral
  \item Solo escrita
  \item Solo visual
\end{multicols}
\end{enumerate}
\vspace{0.15cm}
\item La exposición divulgativa se dirige a:
\begin{enumerate}[label=\textbf{\Alph*)}]
\begin{multicols}{2}
  \item Público general
  \item Niños
\end{multicols}
\end{enumerate}
\vspace{0.15cm}
\item La exposición especializada se dirige a:
\begin{enumerate}[label=\textbf{\Alph*)}]
\begin{multicols}{2}
  \item Público general
  \item Niños
\end{multicols}
\end{enumerate}
\vspace{0.15cm}
\item El lenguaje de un texto expositivo divulgativo debe ser:
\begin{enumerate}[label=\textbf{\Alph*)}]
\begin{multicols}{2}
  \item Muy técnico
  \item Informal
\end{multicols}
\end{enumerate}
\vspace{0.15cm}
\item El lenguaje de un texto expositivo especializado debe ser:
\begin{enumerate}[label=\textbf{\Alph*)}]
\begin{multicols}{2}
  \item Técnico y preciso
  \item Informal
\end{multicols}
\end{enumerate}
\vspace{0.15cm}
\item La objetividad en un texto expositivo implica:
\begin{enumerate}[label=\textbf{\Alph*)}]
\begin{multicols}{2}
  \item Presentar la información sin opiniones personales
  \item Presentar opiniones personales
  \item Presentar solo sentimientos
  \item Presentar solo emociones
  \item Presentar solo creencias
  \item Presentar solo suposiciones
\end{multicols}
\end{enumerate}
\vspace{0.15cm}
\item La claridad en un texto expositivo implica:
\begin{enumerate}[label=\textbf{\Alph*)}]
\begin{multicols}{2}
  \item Usar un lenguaje comprensible
  \item Usar un lenguaje confuso
  \item Usar un lenguaje ambiguo
  \item Usar un lenguaje informal
  \item Usar un lenguaje coloquial
  \item Usar un lenguaje poético
\end{multicols}
\end{enumerate}
\vspace{0.15cm}
\item La precisión en un texto expositivo implica:
\begin{enumerate}[label=\textbf{\Alph*)}]
\begin{multicols}{2}
  \item Usar términos exactos y adecuados
  \item Usar términos vagos
  \item Usar términos ambiguos
  \item Usar términos informales
  \item Usar términos coloquiales
  \item Usar términos poéticos
\end{multicols}
\end{enumerate}
\vspace{0.15cm}
\item La ordenación en un texto expositivo implica:
\begin{enumerate}[label=\textbf{\Alph*)}]
\begin{multicols}{2}
  \item Organizar las ideas de forma lógica
  \item Organizar las ideas de forma caótica
  \item Organizar las ideas de forma aleatoria
  \item Organizar las ideas de forma confusa
  \item Organizar las ideas de forma ambigua
  \item Organizar las ideas de forma informal
\end{multicols}
\end{enumerate}
\vspace{0.15cm}
\item ¿Qué tipo de estructura expositiva presenta el tema, lo analiza y luego lo sintetiza?
\begin{enumerate}[label=\textbf{\Alph*)}]
\begin{multicols}{2}
  \item Sintetizante
  \item Paralela
\end{multicols}
\end{enumerate}
\vspace{0.15cm}
\item ¿Qué tipo de estructura expositiva presenta primero la conclusión y luego la explica?
\begin{enumerate}[label=\textbf{\Alph*)}]
\begin{multicols}{2}
  \item Sintetizante
  \item Paralela
\end{multicols}
\end{enumerate}
\vspace{0.15cm}
\item ¿Qué tipo de estructura expositiva presenta el tema al inicio y al final?
\begin{enumerate}[label=\textbf{\Alph*)}]
\begin{multicols}{2}
  \item Sintetizante
  \item Paralela
\end{multicols}
\end{enumerate}
\vspace{0.15cm}
\item ¿Qué tipo de estructura expositiva presenta ideas en orden temporal?
\begin{enumerate}[label=\textbf{\Alph*)}]
\begin{multicols}{2}
  \item Sintetizante
  \item Paralela
\end{multicols}
\end{enumerate}
\vspace{0.15cm}
\item ¿Qué tipo de estructura expositiva presenta causas y consecuencias?
\begin{enumerate}[label=\textbf{\Alph*)}]
\begin{multicols}{2}
  \item Sintetizante
  \item Paralela
\end{multicols}
\end{enumerate}
\vspace{0.15cm}
\item ¿Qué tipo de estructura expositiva presenta ideas sin jerarquía?
\begin{enumerate}[label=\textbf{\Alph*)}]
\begin{multicols}{2}
  \item Paralela
  \item Sintetizante
\end{multicols}
\end{enumerate}
\vspace{0.15cm}
\item La definición en un texto expositivo sirve para:
\begin{enumerate}[label=\textbf{\Alph*)}]
\begin{multicols}{2}
  \item Repetir ideas
\end{multicols}
\end{enumerate}
\vspace{0.15cm}
\item La clasificación en un texto expositivo sirve para:
\begin{enumerate}[label=\textbf{\Alph*)}]
\begin{multicols}{2}
  \item Organizar elementos en categorías
  \item Repetir ideas
\end{multicols}
\end{enumerate}
\vspace{0.15cm}
\item La comparación en un texto expositivo sirve para:
\begin{enumerate}[label=\textbf{\Alph*)}]
\begin{multicols}{2}
  \item Repetir ideas
\end{multicols}
\end{enumerate}
\vspace{0.15cm}
\item La ejemplificación en un texto expositivo sirve para:
\begin{enumerate}[label=\textbf{\Alph*)}]
\begin{multicols}{2}
  \item Ilustrar una idea con casos concretos
  \item Repetir ideas
\end{multicols}
\end{enumerate}
\vspace{0.15cm}
\item La reformulación en un texto expositivo sirve para:
\begin{enumerate}[label=\textbf{\Alph*)}]
\begin{multicols}{2}
  \item Repetir ideas
\end{multicols}
\end{enumerate}
\vspace{0.15cm}
\item La enumeración en un texto expositivo sirve para:
\begin{enumerate}[label=\textbf{\Alph*)}]
\begin{multicols}{2}
  \item Presentar elementos de una serie
  \item Repetir ideas
\end{multicols}
\end{enumerate}
\vspace{0.15cm}
\item La descripción en un texto expositivo sirve para:
\begin{enumerate}[label=\textbf{\Alph*)}]
\begin{multicols}{2}
  \item Repetir ideas
\end{multicols}
\end{enumerate}
\vspace{0.15cm}
\item La narración en un texto expositivo sirve para:
\begin{enumerate}[label=\textbf{\Alph*)}]
\begin{multicols}{2}
  \item Relatar hechos o procesos
  \item Repetir ideas
\end{multicols}
\end{enumerate}
\vspace{0.15cm}
\item La argumentación en un texto expositivo sirve para:
\begin{enumerate}[label=\textbf{\Alph*)}]
\begin{multicols}{2}
  \item Repetir ideas
\end{multicols}
\end{enumerate}
\vspace{0.15cm}
\item Un texto expositivo debe evitar:
\begin{enumerate}[label=\textbf{\Alph*)}]
\begin{multicols}{2}
  \item Las opiniones personales
  \item Los datos objetivos
  \item Las definiciones
  \item Las clasificaciones
  \item Las comparaciones
  \item Los ejemplos
\end{multicols}
\end{enumerate}
\vspace{0.15cm}
\item Un texto expositivo debe incluir:
\begin{enumerate}[label=\textbf{\Alph*)}]
\begin{multicols}{2}
  \item Información verificable
  \item Opiniones personales
  \item Sentimientos
  \item Suposiciones
\end{multicols}
\end{enumerate}
\vspace{0.15cm}
\item La función referencial del lenguaje predomina en:
\begin{enumerate}[label=\textbf{\Alph*)}]
\begin{multicols}{2}
  \item Textos expositivos
  \item Textos narrativos
  \item Textos poéticos
  \item Textos dramáticos
  \item Textos argumentativos
  \item Textos descriptivos
\end{multicols}
\end{enumerate}
\vspace{0.15cm}
\item La función apelativa del lenguaje predomina en:
\begin{enumerate}[label=\textbf{\Alph*)}]
\begin{multicols}{2}
  \item Textos argumentativos
  \item Textos expositivos
  \item Textos narrativos
  \item Textos poéticos
  \item Textos dramáticos
  \item Textos descriptivos
\end{multicols}
\end{enumerate}
\vspace{0.15cm}
\item La función poética del lenguaje predomina en:
\begin{enumerate}[label=\textbf{\Alph*)}]
\begin{multicols}{2}
  \item Textos literarios
  \item Textos expositivos
  \item Textos argumentativos
  \item Textos científicos
  \item Textos periodísticos
  \item Textos académicos
\end{multicols}
\end{enumerate}
\vspace{0.15cm}
\item La función emotiva del lenguaje predomina en:
\begin{enumerate}[label=\textbf{\Alph*)}]
\begin{multicols}{2}
  \item Textos líricos
  \item Textos expositivos
  \item Textos argumentativos
  \item Textos científicos
  \item Textos periodísticos
  \item Textos académicos
\end{multicols}
\end{enumerate}
\vspace{0.15cm}
\item La función metalingüística del lenguaje predomina en:
\begin{enumerate}[label=\textbf{\Alph*)}]
\begin{multicols}{2}
  \item Textos que hablan sobre el lenguaje
  \item Textos narrativos
  \item Textos poéticos
  \item Textos dramáticos
  \item Textos descriptivos
  \item Textos argumentativos
\end{multicols}
\end{enumerate}
\vspace{0.15cm}
\item La función fática del lenguaje predomina en:
\begin{enumerate}[label=\textbf{\Alph*)}]
\begin{multicols}{2}
  \item Textos que verifican la comunicación
  \item Textos narrativos
  \item Textos poéticos
  \item Textos dramáticos
  \item Textos descriptivos
  \item Textos argumentativos
\end{multicols}
\end{enumerate}
\vspace{0.15cm}
\item ¿Cuál es el propósito de un texto expositivo?
\begin{enumerate}[label=\textbf{\Alph*)}]
\begin{multicols}{2}
  \item Transmitir información objetiva
\end{multicols}
\end{enumerate}
\vspace{0.15cm}
\item ¿Cuál es el propósito de un texto argumentativo?
\begin{enumerate}[label=\textbf{\Alph*)}]
\begin{multicols}{2}
  \item Transmitir información objetiva
\end{multicols}
\end{enumerate}
\vspace{0.15cm}
\item ¿Cuál es el propósito de un texto narrativo?
\begin{enumerate}[label=\textbf{\Alph*)}]
\begin{multicols}{2}
  \item Transmitir información objetiva
\end{multicols}
\end{enumerate}
\vspace{0.15cm}
\item ¿Cuál es el propósito de un texto descriptivo?
\begin{enumerate}[label=\textbf{\Alph*)}]
\begin{multicols}{2}
  \item Representar características
  \item Transmitir información objetiva
\end{multicols}
\end{enumerate}
\vspace{0.15cm}
\item ¿Cuál es el propósito de un texto literario?
\begin{enumerate}[label=\textbf{\Alph*)}]
\begin{multicols}{2}
  \item Transmitir información objetiva
\end{multicols}
\end{enumerate}
\vspace{0.15cm}
\item La estructura "introducción, desarrollo y conclusión" es típica de:
\begin{enumerate}[label=\textbf{\Alph*)}]
\begin{multicols}{2}
  \item Textos expositivos y argumentativos
  \item Textos narrativos
  \item Textos poéticos
  \item Textos dramáticos
  \item Textos descriptivos
  \item Textos literarios
\end{multicols}
\end{enumerate}
\vspace{0.15cm}
\item La estructura "planteamiento, nudo y desenlace" es típica de:
\begin{enumerate}[label=\textbf{\Alph*)}]
\begin{multicols}{2}
  \item Textos narrativos
  \item Textos expositivos
  \item Textos argumentativos
  \item Textos científicos
  \item Textos académicos
  \item Textos periodísticos
\end{multicols}
\end{enumerate}
\vspace{0.15cm}
\item La estructura "tesis, argumentos y conclusión" es típica de:
\begin{enumerate}[label=\textbf{\Alph*)}]
\begin{multicols}{2}
  \item Textos argumentativos
  \item Textos narrativos
  \item Textos expositivos
  \item Textos poéticos
  \item Textos dramáticos
  \item Textos descriptivos
\end{multicols}
\end{enumerate}
\vspace{0.15cm}
\item Los textos expositivos suelen usar:
\begin{enumerate}[label=\textbf{\Alph*)}]
\begin{multicols}{2}
  \item Tercera persona y lenguaje objetivo
  \item Primera persona y lenguaje subjetivo
  \item Segunda persona y lenguaje coloquial
  \item Primera persona y lenguaje poético
  \item Segunda persona y lenguaje informal
  \item Tercera persona y lenguaje poético
\end{multicols}
\end{enumerate}
\vspace{0.15cm}
\item Los textos argumentativos suelen usar:
\begin{enumerate}[label=\textbf{\Alph*)}]
\begin{multicols}{2}
  \item Primera persona y lenguaje subjetivo
  \item Tercera persona y lenguaje objetivo
  \item Segunda persona y lenguaje coloquial
  \item Tercera persona y lenguaje poético
  \item Primera persona y lenguaje científico
  \item Segunda persona y lenguaje formal ---
\end{multicols}
\end{enumerate}
\vspace{0.15cm}
\end{enumerate}\n
\subsubsection*{SECCIÓN V: ESTRUCTURA DE TEXTOS ARGUMENTATIVOS (201–250)}

\begin{enumerate}[label=\textbf{\arabic*.}]
\setcounter{enumi}{200}
\item Un texto argumentativo tiene como propósito:
\begin{enumerate}[label=\textbf{\Alph*)}]
\begin{multicols}{2}
  \item Informar objetivamente
\end{multicols}
\end{enumerate}
\vspace{0.15cm}
\item La tesis de un texto argumentativo es:
\begin{enumerate}[label=\textbf{\Alph*)}]
\begin{multicols}{2}
  \item La idea principal que se defiende
  \item La conclusión
  \item La introducción
  \item La idea secundaria
\end{multicols}
\end{enumerate}
\vspace{0.15cm}
\item Los argumentos de un texto argumentativo son:
\begin{enumerate}[label=\textbf{\Alph*)}]
\begin{multicols}{2}
  \item Las razones que apoyan la tesis
  \item Las ideas secundarias
  \item La conclusión
  \item La introducción
\end{multicols}
\end{enumerate}
\vspace{0.15cm}
\item La conclusión de un texto argumentativo:
\begin{enumerate}[label=\textbf{\Alph*)}]
\begin{multicols}{2}
  \item Reafirma la tesis y resume los argumentos
  \item Presenta el tema
  \item Repite la introducción
\end{multicols}
\end{enumerate}
\vspace{0.15cm}
\item ¿Cuál de los siguientes es un texto argumentativo?
\begin{enumerate}[label=\textbf{\Alph*)}]
\begin{multicols}{2}
  \item Un artículo de opinión
  \item Un cuento
  \item Un poema
  \item Una noticia
  \item Una receta
  \item Un manual
\end{multicols}
\end{enumerate}
\vspace{0.15cm}
\item ¿Cuál de los siguientes es un texto argumentativo?
\begin{enumerate}[label=\textbf{\Alph*)}]
\begin{multicols}{2}
  \item Un ensayo
  \item Un cuento
  \item Un poema
  \item Una noticia
  \item Una receta
  \item Un manual
\end{multicols}
\end{enumerate}
\vspace{0.15cm}
\item ¿Cuál de los siguientes es un texto argumentativo?
\begin{enumerate}[label=\textbf{\Alph*)}]
\begin{multicols}{2}
  \item Un discurso político
  \item Un cuento
  \item Un poema
  \item Una noticia
  \item Una receta
  \item Un manual
\end{multicols}
\end{enumerate}
\vspace{0.15cm}
\item ¿Cuál de los siguientes es un texto argumentativo?
\begin{enumerate}[label=\textbf{\Alph*)}]
\begin{multicols}{2}
  \item Una reseña crítica
  \item Un cuento
  \item Un poema
  \item Una noticia
  \item Una receta
  \item Un manual
\end{multicols}
\end{enumerate}
\vspace{0.15cm}
\item La argumentación se basa en:
\begin{enumerate}[label=\textbf{\Alph*)}]
\begin{multicols}{2}
  \item Razonamientos lógicos
  \item Sentimientos
  \item Suposiciones
  \item Opiniones sin fundamento
\end{multicols}
\end{enumerate}
\vspace{0.15cm}
\item La argumentación puede ser:
\begin{enumerate}[label=\textbf{\Alph*)}]
\begin{multicols}{2}
  \item Inductiva o deductiva
  \item Solo inductiva
  \item Solo deductiva
  \item Solo oral
  \item Solo escrita
  \item Solo visual
\end{multicols}
\end{enumerate}
\vspace{0.15cm}
\item La argumentación inductiva parte de:
\begin{enumerate}[label=\textbf{\Alph*)}]
\begin{multicols}{2}
  \item Una conclusión general para llegar a casos particulares
  \item Una premisa falsa
  \item Una premisa verdadera
  \item Una conclusión falsa
  \item Una conclusión verdadera
\end{multicols}
\end{enumerate}
\vspace{0.15cm}
\item La argumentación deductiva parte de:
\begin{enumerate}[label=\textbf{\Alph*)}]
\begin{multicols}{2}
  \item Una conclusión general para llegar a casos particulares
  \item Una premisa falsa
  \item Una premisa verdadera
  \item Una conclusión falsa
  \item Una conclusión verdadera
\end{multicols}
\end{enumerate}
\vspace{0.15cm}
\item La falacia argumentativa es:
\begin{enumerate}[label=\textbf{\Alph*)}]
\begin{multicols}{2}
  \item Un razonamiento incorrecto que parece válido
  \item Un razonamiento correcto
  \item Una premisa verdadera
  \item Una conclusión verdadera
  \item Un argumento sólido
  \item Una tesis correcta
\end{multicols}
\end{enumerate}
\vspace{0.15cm}
\item La falacia "ad hominem" consiste en:
\begin{enumerate}[label=\textbf{\Alph*)}]
\begin{multicols}{2}
\end{multicols}
\end{enumerate}
\vspace{0.15cm}
\item La falacia "ad populum" consiste en:
\begin{enumerate}[label=\textbf{\Alph*)}]
\begin{multicols}{2}
\end{multicols}
\end{enumerate}
\vspace{0.15cm}
\item La falacia "ad verecundiam" consiste en:
\begin{enumerate}[label=\textbf{\Alph*)}]
\begin{multicols}{2}
\end{multicols}
\end{enumerate}
\vspace{0.15cm}
\item La falacia "ad ignorantiam" consiste en:
\begin{enumerate}[label=\textbf{\Alph*)}]
\begin{multicols}{2}
\end{multicols}
\end{enumerate}
\vspace{0.15cm}
\item La falacia "ad baculum" consiste en:
\begin{enumerate}[label=\textbf{\Alph*)}]
\begin{multicols}{2}
\end{multicols}
\end{enumerate}
\vspace{0.15cm}
\item La falacia "ad misericordiam" consiste en:
\begin{enumerate}[label=\textbf{\Alph*)}]
\begin{multicols}{2}
\end{multicols}
\end{enumerate}
\vspace{0.15cm}
\item La falacia "pendiente resbaladiza" consiste en:
\begin{enumerate}[label=\textbf{\Alph*)}]
\begin{multicols}{2}
\end{multicols}
\end{enumerate}
\vspace{0.15cm}
\item La falacia "falsa dicotomía" consiste en:
\begin{enumerate}[label=\textbf{\Alph*)}]
\begin{multicols}{2}
  \item Presentar solo dos opciones cuando hay más
\end{multicols}
\end{enumerate}
\vspace{0.15cm}
\item La falacia "generalización apresurada" consiste en:
\begin{enumerate}[label=\textbf{\Alph*)}]
\begin{multicols}{2}
  \item Sacar una conclusión general de pocos casos
\end{multicols}
\end{enumerate}
\vspace{0.15cm}
\item La falacia "causa falsa" consiste en:
\begin{enumerate}[label=\textbf{\Alph*)}]
\begin{multicols}{2}
\end{multicols}
\end{enumerate}
\vspace{0.15cm}
\item La falacia "hombre de paja" consiste en:
\begin{enumerate}[label=\textbf{\Alph*)}]
\begin{multicols}{2}
\end{multicols}
\end{enumerate}
\vspace{0.15cm}
\item La falacia "circular" consiste en:
\begin{enumerate}[label=\textbf{\Alph*)}]
\begin{multicols}{2}
  \item Usar la conclusión como premisa
\end{multicols}
\end{enumerate}
\vspace{0.15cm}
\item La falacia "equívoco" consiste en:
\begin{enumerate}[label=\textbf{\Alph*)}]
\begin{multicols}{2}
  \item Usar una palabra con distintos significados en el mismo argumento
\end{multicols}
\end{enumerate}
\vspace{0.15cm}
\item La falacia "anfibología" consiste en:
\begin{enumerate}[label=\textbf{\Alph*)}]
\begin{multicols}{2}
  \item Usar construcciones gramaticales ambiguas
\end{multicols}
\end{enumerate}
\vspace{0.15cm}
\item La falacia "composición" consiste en:
\begin{enumerate}[label=\textbf{\Alph*)}]
\begin{multicols}{2}
\end{multicols}
\end{enumerate}
\vspace{0.15cm}
\item La falacia "división" consiste en:
\begin{enumerate}[label=\textbf{\Alph*)}]
\begin{multicols}{2}
\end{multicols}
\end{enumerate}
\vspace{0.15cm}
\item La falacia "non sequitur" consiste en:
\begin{enumerate}[label=\textbf{\Alph*)}]
\begin{multicols}{2}
  \item Llegar a una conclusión que no se deriva de las premisas
\end{multicols}
\end{enumerate}
\vspace{0.15cm}
\item El ensayo argumentativo tiene una estructura de:
\begin{enumerate}[label=\textbf{\Alph*)}]
\begin{multicols}{2}
  \item Introducción, desarrollo y conclusión
  \item Planteamiento, nudo y desenlace
  \item Inicio, conflicto y resolución
  \item Título, cuerpo y firma
  \item Tesis, antítesis y síntesis
\end{multicols}
\end{enumerate}
\vspace{0.15cm}
\item La introducción de un ensayo argumentativo debe:
\begin{enumerate}[label=\textbf{\Alph*)}]
\begin{multicols}{2}
  \item Presentar la tesis y captar la atención
  \item Repetir el título
\end{multicols}
\end{enumerate}
\vspace{0.15cm}
\item El desarrollo de un ensayo argumentativo debe:
\begin{enumerate}[label=\textbf{\Alph*)}]
\begin{multicols}{2}
  \item Presentar los argumentos que apoyan la tesis
  \item Presentar solo la tesis
  \item Repetir la introducción
\end{multicols}
\end{enumerate}
\vspace{0.15cm}
\item La conclusión de un ensayo argumentativo debe:
\begin{enumerate}[label=\textbf{\Alph*)}]
\begin{multicols}{2}
  \item Reafirmar la tesis y resumir los argumentos
  \item Presentar la tesis
  \item Repetir la introducción
\end{multicols}
\end{enumerate}
\vspace{0.15cm}
\item Los argumentos de autoridad se basan en:
\begin{enumerate}[label=\textbf{\Alph*)}]
\begin{multicols}{2}
  \item La opinión de expertos en la materia
  \item La opinión de la mayoría
  \item La emoción
  \item La fuerza
  \item La ignorancia
  \item La tradición
\end{multicols}
\end{enumerate}
\vspace{0.15cm}
\item Los argumentos de hecho se basan en:
\begin{enumerate}[label=\textbf{\Alph*)}]
\begin{multicols}{2}
  \item Opiniones personales
  \item Suposiciones
  \item Tradiciones
\end{multicols}
\end{enumerate}
\vspace{0.15cm}
\item Los argumentos de ejemplificación se basan en:
\begin{enumerate}[label=\textbf{\Alph*)}]
\begin{multicols}{2}
  \item Opiniones personales
  \item Suposiciones
  \item Tradiciones
\end{multicols}
\end{enumerate}
\vspace{0.15cm}
\item Los argumentos de analogía se basan en:
\begin{enumerate}[label=\textbf{\Alph*)}]
\begin{multicols}{2}
  \item Opiniones personales
  \item Suposiciones
  \item Tradiciones
\end{multicols}
\end{enumerate}
\vspace{0.15cm}
\item Los argumentos de causa-efecto se basan en:
\begin{enumerate}[label=\textbf{\Alph*)}]
\begin{multicols}{2}
  \item Relaciones de causalidad
  \item Opiniones personales
  \item Suposiciones
  \item Tradiciones
\end{multicols}
\end{enumerate}
\vspace{0.15cm}
\item La refutación en un texto argumentativo consiste en:
\begin{enumerate}[label=\textbf{\Alph*)}]
\begin{multicols}{2}
  \item Rebartir los argumentos del oponente
  \item Ignorar los argumentos del oponente
  \item Repetir los argumentos del oponente
\end{multicols}
\end{enumerate}
\vspace{0.15cm}
\item La contraargumentación consiste en:
\begin{enumerate}[label=\textbf{\Alph*)}]
\begin{multicols}{2}
  \item Presentar argumentos contrarios a la tesis para refutarlos
  \item Ignorar la tesis
  \item Repetir la tesis
\end{multicols}
\end{enumerate}
\vspace{0.15cm}
\item La concesión en un texto argumentativo consiste en:
\begin{enumerate}[label=\textbf{\Alph*)}]
\begin{multicols}{2}
  \item Rechazar totalmente un argumento contrario
  \item Ignorar un argumento contrario
  \item Repetir un argumento contrario
\end{multicols}
\end{enumerate}
\vspace{0.15cm}
\item La persuasión en un texto argumentativo busca:
\begin{enumerate}[label=\textbf{\Alph*)}]
\begin{multicols}{2}
  \item Informar al lector
\end{multicols}
\end{enumerate}
\vspace{0.15cm}
\item La disuasión en un texto argumentativo busca:
\begin{enumerate}[label=\textbf{\Alph*)}]
\begin{multicols}{2}
  \item Informar al lector
\end{multicols}
\end{enumerate}
\vspace{0.15cm}
\item El lenguaje de un texto argumentativo debe ser:
\begin{enumerate}[label=\textbf{\Alph*)}]
\begin{multicols}{2}
  \item Informal y coloquial
  \item Poético y metafórico
  \item Técnico y especializado
  \item Neutro y objetivo
\end{multicols}
\end{enumerate}
\vspace{0.15cm}
\item Los conectores más usados en textos argumentativos son:
\begin{enumerate}[label=\textbf{\Alph*)}]
\begin{multicols}{2}
  \item Ordenadores y de cierre
\end{multicols}
\end{enumerate}
\vspace{0.15cm}
\item La tesis de un texto argumentativo debe ser:
\begin{enumerate}[label=\textbf{\Alph*)}]
\begin{multicols}{2}
  \item Informal y coloquial
  \item Poética y metafórica
  \item Técnica y especializada
  \item Neutra y objetiva
\end{multicols}
\end{enumerate}
\vspace{0.15cm}
\item Un texto argumentativo sólido debe:
\begin{enumerate}[label=\textbf{\Alph*)}]
\begin{multicols}{2}
  \item Tener argumentos bien fundamentados
  \item Tener argumentos sin fundamento
  \item Tener argumentos contradictorios
  \item Tener argumentos falsos
  \item Tener argumentos ambiguos
  \item Tener argumentos confusos
\end{multicols}
\end{enumerate}
\vspace{0.15cm}
\item La falacia "tu quoque" consiste en:
\begin{enumerate}[label=\textbf{\Alph*)}]
\begin{multicols}{2}
  \item Responder a un ataque con otro ataque similar
\end{multicols}
\end{enumerate}
\vspace{0.15cm}
\item La falacia "post hoc" consiste en:
\begin{enumerate}[label=\textbf{\Alph*)}]
\begin{multicols}{2}
\end{multicols}
\end{enumerate}
\vspace{0.15cm}
\end{enumerate}\n
\subsubsection*{SECCIÓN VI: TIPOS DE PÁRRAFOS Y PROGRESIÓN TEMÁTICA (251–300)}

\begin{enumerate}[label=\textbf{\arabic*.}]
\setcounter{enumi}{250}
\item Un párrafo es:
\begin{enumerate}[label=\textbf{\Alph*)}]
\begin{multicols}{2}
  \item Un conjunto de oraciones que desarrollan una idea
  \item Un conjunto de palabras sin relación
  \item Un conjunto de textos
  \item Un conjunto de párrafos
  \item Un conjunto de conectores
  \item Un conjunto de signos
\end{multicols}
\end{enumerate}
\vspace{0.15cm}
\item La idea principal de un párrafo:
\begin{enumerate}[label=\textbf{\Alph*)}]
\begin{multicols}{2}
\end{multicols}
\end{enumerate}
\vspace{0.15cm}
\item Las ideas secundarias de un párrafo:
\begin{enumerate}[label=\textbf{\Alph*)}]
\begin{multicols}{2}
  \item Son más importantes que la principal
  \item No tienen relación con la principal
  \item Son la conclusión
  \item Son la introducción
  \item Son el título
\end{multicols}
\end{enumerate}
\vspace{0.15cm}
\item La oración temática de un párrafo:
\begin{enumerate}[label=\textbf{\Alph*)}]
\begin{multicols}{2}
\end{multicols}
\end{enumerate}
\vspace{0.15cm}
\item Un párrafo de introducción tiene como función:
\begin{enumerate}[label=\textbf{\Alph*)}]
\begin{multicols}{2}
  \item Presentar el tema y captar la atención
  \item Repetir el tema
\end{multicols}
\end{enumerate}
\vspace{0.15cm}
\item Un párrafo de desarrollo tiene como función:
\begin{enumerate}[label=\textbf{\Alph*)}]
\begin{multicols}{2}
  \item Presentar el tema
  \item Repetir el tema
\end{multicols}
\end{enumerate}
\vspace{0.15cm}
\item Un párrafo de conclusión tiene como función:
\begin{enumerate}[label=\textbf{\Alph*)}]
\begin{multicols}{2}
  \item Sintetizar o cerrar las ideas presentadas
  \item Presentar el tema
  \item Repetir el tema
\end{multicols}
\end{enumerate}
\vspace{0.15cm}
\item Un párrafo expositivo tiene como función:
\begin{enumerate}[label=\textbf{\Alph*)}]
\begin{multicols}{2}
\end{multicols}
\end{enumerate}
\vspace{0.15cm}
\item Un párrafo argumentativo tiene como función:
\begin{enumerate}[label=\textbf{\Alph*)}]
\begin{multicols}{2}
  \item Informar objetivamente
\end{multicols}
\end{enumerate}
\vspace{0.15cm}
\item Un párrafo narrativo tiene como función:
\begin{enumerate}[label=\textbf{\Alph*)}]
\begin{multicols}{2}
  \item Informar objetivamente
\end{multicols}
\end{enumerate}
\vspace{0.15cm}
\item Un párrafo descriptivo tiene como función:
\begin{enumerate}[label=\textbf{\Alph*)}]
\begin{multicols}{2}
  \item Representar características de personas, objetos o lugares
  \item Informar objetivamente
\end{multicols}
\end{enumerate}
\vspace{0.15cm}
\item Un párrafo de comparación tiene como función:
\begin{enumerate}[label=\textbf{\Alph*)}]
\begin{multicols}{2}
  \item Informar objetivamente
\end{multicols}
\end{enumerate}
\vspace{0.15cm}
\item Un párrafo de causa-efecto tiene como función:
\begin{enumerate}[label=\textbf{\Alph*)}]
\begin{multicols}{2}
  \item Informar objetivamente
\end{multicols}
\end{enumerate}
\vspace{0.15cm}
\item Un párrafo de definición tiene como función:
\begin{enumerate}[label=\textbf{\Alph*)}]
\begin{multicols}{2}
  \item Informar objetivamente
\end{multicols}
\end{enumerate}
\vspace{0.15cm}
\item Un párrafo de clasificación tiene como función:
\begin{enumerate}[label=\textbf{\Alph*)}]
\begin{multicols}{2}
  \item Organizar elementos en categorías
  \item Informar objetivamente
\end{multicols}
\end{enumerate}
\vspace{0.15cm}
\item Un párrafo de ejemplificación tiene como función:
\begin{enumerate}[label=\textbf{\Alph*)}]
\begin{multicols}{2}
  \item Ilustrar una idea con casos concretos
  \item Informar objetivamente
\end{multicols}
\end{enumerate}
\vspace{0.15cm}
\item La progresión temática se refiere a:
\begin{enumerate}[label=\textbf{\Alph*)}]
\begin{multicols}{2}
  \item La repetición de la misma idea
  \item La cantidad de párrafos
  \item La extensión del texto
  \item La cantidad de conectores
  \item La ortografía
\end{multicols}
\end{enumerate}
\vspace{0.15cm}
\item La progresión temática puede ser:
\begin{enumerate}[label=\textbf{\Alph*)}]
\begin{multicols}{2}
  \item Lineal, con tema constante o derivada
  \item Solo lineal
  \item Solo con tema constante
  \item Solo derivada
  \item Solo circular
  \item Solo repetitiva
\end{multicols}
\end{enumerate}
\vspace{0.15cm}
\item La progresión lineal se caracteriza porque:
\begin{enumerate}[label=\textbf{\Alph*)}]
\begin{multicols}{2}
  \item No hay progresión
  \item No hay tema
  \item No hay rema
\end{multicols}
\end{enumerate}
\vspace{0.15cm}
\item La progresión con tema constante se caracteriza porque:
\begin{enumerate}[label=\textbf{\Alph*)}]
\begin{multicols}{2}
  \item No hay progresión
  \item No hay tema
  \item No hay rema
\end{multicols}
\end{enumerate}
\vspace{0.15cm}
\item La progresión derivada se caracteriza porque:
\begin{enumerate}[label=\textbf{\Alph*)}]
\begin{multicols}{2}
  \item No hay progresión
  \item No hay tema
  \item No hay rema
\end{multicols}
\end{enumerate}
\vspace{0.15cm}
\item El tema de un texto es:
\begin{enumerate}[label=\textbf{\Alph*)}]
\begin{multicols}{2}
  \item La idea principal
  \item La conclusión
  \item La introducción
\end{multicols}
\end{enumerate}
\vspace{0.15cm}
\item El rema de un texto es:
\begin{enumerate}[label=\textbf{\Alph*)}]
\begin{multicols}{2}
  \item La información nueva que se aporta sobre el tema
  \item La conclusión
  \item La introducción
\end{multicols}
\end{enumerate}
\vspace{0.15cm}
\item La estructura "tema-rema" se refiere a:
\begin{enumerate}[label=\textbf{\Alph*)}]
\begin{multicols}{2}
  \item La organización de la información en el texto
  \item La ortografía del texto
  \item La puntuación del texto
  \item La caligrafía del texto
  \item La sintaxis del texto
  \item La extensión del texto
\end{multicols}
\end{enumerate}
\vspace{0.15cm}
\item Un texto bien estructurado debe tener:
\begin{enumerate}[label=\textbf{\Alph*)}]
\begin{multicols}{2}
  \item Introducción, desarrollo y conclusión
  \item Solo introducción
  \item Solo desarrollo
  \item Solo conclusión
  \item Sin estructura
\end{multicols}
\end{enumerate}
\vspace{0.15cm}
\item La introducción de un texto debe:
\begin{enumerate}[label=\textbf{\Alph*)}]
\begin{multicols}{2}
  \item Presentar el tema y captar la atención
  \item Repetir el título
\end{multicols}
\end{enumerate}
\vspace{0.15cm}
\item El desarrollo de un texto debe:
\begin{enumerate}[label=\textbf{\Alph*)}]
\begin{multicols}{2}
  \item Presentar solo el tema
  \item Repetir la introducción
\end{multicols}
\end{enumerate}
\vspace{0.15cm}
\item La conclusión de un texto debe:
\begin{enumerate}[label=\textbf{\Alph*)}]
\begin{multicols}{2}
  \item Sintetizar o cerrar las ideas presentadas
  \item Presentar el tema
  \item Repetir la introducción
\end{multicols}
\end{enumerate}
\vspace{0.15cm}
\item La cohesión textual se logra mediante:
\begin{enumerate}[label=\textbf{\Alph*)}]
\begin{multicols}{2}
  \item Solo signos de puntuación
  \item Solo la ortografía
  \item Solo la caligrafía
  \item Solo la extensión
  \item Solo el tema
\end{multicols}
\end{enumerate}
\vspace{0.15cm}
\item La coherencia textual se logra mediante:
\begin{enumerate}[label=\textbf{\Alph*)}]
\begin{multicols}{2}
  \item Unidad temática y lógica
  \item Muchas palabras
  \item Muchos párrafos
  \item Muchos conectores
  \item Mucha puntuación
  \item Mucha extensión
\end{multicols}
\end{enumerate}
\vspace{0.15cm}
\item Un texto coherente y cohesionado debe:
\begin{enumerate}[label=\textbf{\Alph*)}]
\begin{multicols}{2}
  \item Tener unidad temática y usar adecuadamente los mecanismos de cohesión
  \item Tener muchos temas
  \item No tener tema
  \item No usar conectores
  \item No usar pronombres
\end{multicols}
\end{enumerate}
\vspace{0.15cm}
\item La adecuación textual se refiere a:
\begin{enumerate}[label=\textbf{\Alph*)}]
\begin{multicols}{2}
  \item La adaptación del texto al contexto y al destinatario
  \item La ortografía del texto
  \item La puntuación del texto
  \item La caligrafía del texto
  \item La sintaxis del texto
  \item La extensión del texto
\end{multicols}
\end{enumerate}
\vspace{0.15cm}
\item Un texto adecuado debe:
\begin{enumerate}[label=\textbf{\Alph*)}]
\begin{multicols}{2}
  \item Usar el registro lingüístico apropiado
  \item Usar siempre lenguaje formal
  \item Usar siempre lenguaje informal
  \item Usar siempre tecnicismos
  \item Usar siempre coloquialismos
  \item Usar siempre cultismos
\end{multicols}
\end{enumerate}
\vspace{0.15cm}
\item La estructura textual se refiere a:
\begin{enumerate}[label=\textbf{\Alph*)}]
\begin{multicols}{2}
  \item La organización de las ideas en el texto
  \item La ortografía del texto
  \item La puntuación del texto
  \item La caligrafía del texto
  \item La sintaxis del texto
  \item La extensión del texto
\end{multicols}
\end{enumerate}
\vspace{0.15cm}
\item La macroestructura de un texto se refiere a:
\begin{enumerate}[label=\textbf{\Alph*)}]
\begin{multicols}{2}
  \item La organización global de las ideas
  \item La organización de las oraciones
  \item La organización de las palabras
  \item La organización de las sílabas
  \item La organización de las letras
  \item La organización de los signos
\end{multicols}
\end{enumerate}
\vspace{0.15cm}
\item La microestructura de un texto se refiere a:
\begin{enumerate}[label=\textbf{\Alph*)}]
\begin{multicols}{2}
  \item La organización de las oraciones y párrafos
  \item La organización global de las ideas
  \item La organización de las palabras
  \item La organización de las sílabas
  \item La organización de las letras
  \item La organización de los signos
\end{multicols}
\end{enumerate}
\vspace{0.15cm}
\item La superestructura de un texto se refiere a:
\begin{enumerate}[label=\textbf{\Alph*)}]
\begin{multicols}{2}
  \item La organización de las oraciones
  \item La organización de las palabras
  \item La organización de las sílabas
  \item La organización de las letras
  \item La organización de los signos
\end{multicols}
\end{enumerate}
\vspace{0.15cm}
\item La superestructura de un texto narrativo es:
\begin{enumerate}[label=\textbf{\Alph*)}]
\begin{multicols}{2}
  \item Planteamiento, nudo y desenlace
  \item Introducción, desarrollo y conclusión
  \item Tesis, argumentos y conclusión
  \item Título, cuerpo y firma
  \item Inicio, conflicto y resolución
\end{multicols}
\end{enumerate}
\vspace{0.15cm}
\item La superestructura de un texto expositivo es:
\begin{enumerate}[label=\textbf{\Alph*)}]
\begin{multicols}{2}
  \item Introducción, desarrollo y conclusión
  \item Planteamiento, nudo y desenlace
  \item Tesis, argumentos y conclusión
  \item Título, cuerpo y firma
  \item Inicio, conflicto y resolución
\end{multicols}
\end{enumerate}
\vspace{0.15cm}
\item La superestructura de un texto argumentativo es:
\begin{enumerate}[label=\textbf{\Alph*)}]
\begin{multicols}{2}
  \item Tesis, argumentos y conclusión
  \item Planteamiento, nudo y desenlace
  \item Introducción, desarrollo y conclusión
  \item Título, cuerpo y firma
  \item Inicio, conflicto y resolución
\end{multicols}
\end{enumerate}
\vspace{0.15cm}
\item La macroestructura de un texto se construye a partir de:
\begin{enumerate}[label=\textbf{\Alph*)}]
\begin{multicols}{2}
  \item Las ideas principales de cada párrafo
  \item Las palabras de cada oración
  \item Las sílabas de cada palabra
  \item Las letras de cada sílaba
  \item Los signos de puntuación
  \item Los conectores
\end{multicols}
\end{enumerate}
\vspace{0.15cm}
\item La microestructura de un texto se construye a partir de:
\begin{enumerate}[label=\textbf{\Alph*)}]
\begin{multicols}{2}
  \item Las oraciones y su organización en párrafos
  \item Las ideas principales
  \item Las palabras clave
  \item Las sílabas
  \item Las letras
  \item Los signos
\end{multicols}
\end{enumerate}
\vspace{0.15cm}
\item La progresión temática de un texto debe ser:
\begin{enumerate}[label=\textbf{\Alph*)}]
\begin{multicols}{2}
  \item Repetitiva
  \item Incoherente
\end{multicols}
\end{enumerate}
\vspace{0.15cm}
\item La unidad temática de un texto implica:
\begin{enumerate}[label=\textbf{\Alph*)}]
\begin{multicols}{2}
  \item Que todo el texto gire en torno a un mismo tema
  \item Que el texto tenga muchos temas
  \item Que el texto no tenga tema
  \item Que el texto cambie de tema constantemente
  \item Que el texto solo tenga un párrafo
  \item Que el texto solo tenga una oración
\end{multicols}
\end{enumerate}
\vspace{0.15cm}
\item La estructura de un texto debe ser:
\begin{enumerate}[label=\textbf{\Alph*)}]
\begin{multicols}{2}
  \item Lógica y ordenada
  \item Incoherente
  \item Repetitiva
\end{multicols}
\end{enumerate}
\vspace{0.15cm}
\item La introducción de un texto debe captar:
\begin{enumerate}[label=\textbf{\Alph*)}]
\begin{multicols}{2}
  \item La atención del lector
  \item La paciencia del lector
  \item La curiosidad del lector
\end{multicols}
\end{enumerate}
\vspace{0.15cm}
\item El desarrollo de un texto debe mantener:
\begin{enumerate}[label=\textbf{\Alph*)}]
\begin{multicols}{2}
  \item La paciencia del lector
  \item La curiosidad del lector
\end{multicols}
\end{enumerate}
\vspace{0.15cm}
\item La conclusión de un texto debe dejar:
\begin{enumerate}[label=\textbf{\Alph*)}]
\begin{multicols}{2}
  \item Una idea clara y completa
  \item Una idea confusa
  \item Una idea ambigua
  \item Una idea incoherente
  \item Una idea repetitiva
  \item Una idea caótica
\end{multicols}
\end{enumerate}
\vspace{0.15cm}
\item Un texto bien escrito debe ser:
\begin{enumerate}[label=\textbf{\Alph*)}]
\begin{multicols}{2}
  \item Incoherente, sin cohesión e inadecuado
  \item Repetitivo, desordenado y confuso
  \item Incoherente, desordenado y ambiguo
\end{multicols}
\end{enumerate}
\vspace{0.15cm}
\item La revisión de un texto debe verificar:
\begin{enumerate}[label=\textbf{\Alph*)}]
\begin{multicols}{2}
  \item Solo la ortografía
  \item Solo la puntuación
  \item Solo la caligrafía
  \item Solo la extensión
  \item Solo el tema ---
\end{multicols}
\end{enumerate}
\vspace{0.15cm}
\end{enumerate}\n
\vspace{0.4cm}
\begin{tcolorbox}[colback=gray!5!white,colframe=gray!60!black,title=\textbf{SOLUCIONARIO -- 2.2 COHERENCIA Y ESTRUCTURA TEXTUAL}]
\begin{multicols}{4}
\begin{enumerate}[label=\textbf{\arabic*.}]
  \item \textbf{A}
  \item \textbf{A}
  \item \textbf{A}
  \item \textbf{A}
  \item \textbf{A}
  \item \textbf{A}
  \item \textbf{D}
  \item \textbf{A}
  \item \textbf{A}
  \item \textbf{A}
  \item \textbf{A}
  \item \textbf{A}
  \item \textbf{A}
  \item \textbf{A}
  \item \textbf{A}
  \item \textbf{A}
  \item \textbf{A}
  \item \textbf{A}
  \item \textbf{A}
  \item \textbf{A}
  \item \textbf{A}
  \item \textbf{A}
  \item \textbf{A}
  \item \textbf{A}
  \item \textbf{A}
  \item \textbf{A}
  \item \textbf{A}
  \item \textbf{A}
  \item \textbf{A}
  \item \textbf{A}
  \item \textbf{A}
  \item \textbf{A}
  \item \textbf{A}
  \item \textbf{A}
  \item \textbf{A}
  \item \textbf{A}
  \item \textbf{A}
  \item \textbf{A}
  \item \textbf{A}
  \item \textbf{A}
  \item \textbf{A}
  \item \textbf{A}
  \item \textbf{A}
  \item \textbf{A}
  \item \textbf{A}
  \item \textbf{A}
  \item \textbf{A}
  \item \textbf{A}
  \item \textbf{A}
  \item \textbf{A}
  \item \textbf{A}
  \item \textbf{A}
  \item \textbf{A}
  \item \textbf{A}
  \item \textbf{A}
  \item \textbf{A}
  \item \textbf{A}
  \item \textbf{A}
  \item \textbf{A}
  \item \textbf{A}
  \item \textbf{A}
  \item \textbf{A}
  \item \textbf{A}
  \item \textbf{A}
  \item \textbf{A}
  \item \textbf{A}
  \item \textbf{A}
  \item \textbf{A}
  \item \textbf{A}
  \item \textbf{A}
  \item \textbf{A}
  \item \textbf{A}
  \item \textbf{A}
  \item \textbf{A}
  \item \textbf{A}
  \item \textbf{A}
  \item \textbf{A}
  \item \textbf{A}
  \item \textbf{A}
  \item \textbf{A}
  \item \textbf{A}
  \item \textbf{A}
  \item \textbf{A}
  \item \textbf{A}
  \item \textbf{A}
  \item \textbf{A}
  \item \textbf{A}
  \item \textbf{A}
  \item \textbf{A}
  \item \textbf{A}
  \item \textbf{A}
  \item \textbf{A}
  \item \textbf{A}
  \item \textbf{A}
  \item \textbf{A}
  \item \textbf{A}
  \item \textbf{A}
  \item \textbf{A}
  \item \textbf{A}
  \item \textbf{A}
  \item \textbf{A}
  \item \textbf{A}
  \item \textbf{A}
  \item \textbf{A}
  \item \textbf{A}
  \item \textbf{A}
  \item \textbf{A}
  \item \textbf{A}
  \item \textbf{A}
  \item \textbf{A}
  \item \textbf{A}
  \item \textbf{A}
  \item \textbf{A}
  \item \textbf{A}
  \item \textbf{A}
  \item \textbf{A}
  \item \textbf{A}
  \item \textbf{A}
  \item \textbf{A}
  \item \textbf{A}
  \item \textbf{A}
  \item \textbf{A}
  \item \textbf{A}
  \item \textbf{A}
  \item \textbf{A}
  \item \textbf{A}
  \item \textbf{A}
  \item \textbf{A}
  \item \textbf{A}
  \item \textbf{A}
  \item \textbf{A}
  \item \textbf{A}
  \item \textbf{A}
  \item \textbf{A}
  \item \textbf{A}
  \item \textbf{A}
  \item \textbf{A}
  \item \textbf{A}
  \item \textbf{A}
  \item \textbf{A}
  \item \textbf{A}
  \item \textbf{A}
  \item \textbf{A}
  \item \textbf{A}
  \item \textbf{A}
  \item \textbf{A}
  \item \textbf{A}
  \item \textbf{A}
  \item \textbf{A}
  \item \textbf{A}
  \item \textbf{A}
  \item \textbf{A}
  \item \textbf{A}
  \item \textbf{A}
  \item \textbf{A}
  \item \textbf{A}
  \item \textbf{A}
  \item \textbf{A}
  \item \textbf{A}
  \item \textbf{A}
  \item \textbf{A}
  \item \textbf{A}
  \item \textbf{A}
  \item \textbf{A}
  \item \textbf{A}
  \item \textbf{A}
  \item \textbf{A}
  \item \textbf{A}
  \item \textbf{A}
  \item \textbf{A}
  \item \textbf{A}
  \item \textbf{A}
  \item \textbf{A}
  \item \textbf{A}
  \item \textbf{A}
  \item \textbf{A}
  \item \textbf{A}
  \item \textbf{A}
  \item \textbf{A}
  \item \textbf{A}
  \item \textbf{A}
  \item \textbf{A}
  \item \textbf{A}
  \item \textbf{A}
  \item \textbf{A}
  \item \textbf{A}
  \item \textbf{A}
  \item \textbf{A}
  \item \textbf{A}
  \item \textbf{A}
  \item \textbf{A}
  \item \textbf{A}
  \item \textbf{A}
  \item \textbf{A}
  \item \textbf{A}
  \item \textbf{A}
  \item \textbf{A}
  \item \textbf{A}
  \item \textbf{A}
  \item \textbf{A}
  \item \textbf{A}
  \item \textbf{A}
  \item \textbf{A}
  \item \textbf{A}
  \item \textbf{A}
  \item \textbf{A}
  \item \textbf{A}
  \item \textbf{A}
  \item \textbf{A}
  \item \textbf{A}
  \item \textbf{A}
  \item \textbf{A}
  \item \textbf{A}
  \item \textbf{A}
  \item \textbf{A}
  \item \textbf{A}
  \item \textbf{A}
  \item \textbf{A}
  \item \textbf{A}
  \item \textbf{A}
  \item \textbf{A}
  \item \textbf{A}
  \item \textbf{A}
  \item \textbf{A}
  \item \textbf{A}
  \item \textbf{A}
  \item \textbf{A}
  \item \textbf{A}
  \item \textbf{A}
  \item \textbf{A}
  \item \textbf{A}
  \item \textbf{A}
  \item \textbf{A}
  \item \textbf{A}
  \item \textbf{A}
  \item \textbf{A}
  \item \textbf{A}
  \item \textbf{A}
  \item \textbf{A}
  \item \textbf{A}
  \item \textbf{A}
  \item \textbf{A}
  \item \textbf{A}
  \item \textbf{A}
  \item \textbf{A}
  \item \textbf{A}
  \item \textbf{A}
  \item \textbf{A}
  \item \textbf{A}
  \item \textbf{A}
  \item \textbf{A}
  \item \textbf{A}
  \item \textbf{A}
  \item \textbf{A}
  \item \textbf{A}
  \item \textbf{A}
  \item \textbf{A}
  \item \textbf{A}
  \item \textbf{A}
  \item \textbf{A}
  \item \textbf{A}
  \item \textbf{A}
  \item \textbf{A}
  \item \textbf{A}
  \item \textbf{A}
  \item \textbf{A}
  \item \textbf{A}
  \item \textbf{A}
  \item \textbf{A}
  \item \textbf{A}
  \item \textbf{A}
  \item \textbf{A}
  \item \textbf{A}
  \item \textbf{A}
  \item \textbf{A}
  \item \textbf{A}
  \item \textbf{A}
  \item \textbf{A}
  \item \textbf{A}
  \item \textbf{A}
  \item \textbf{A}
  \item \textbf{A}
  \item \textbf{A}
  \item \textbf{A}
  \item \textbf{A}
  \item \textbf{A}
  \item \textbf{A}
  \item \textbf{A}
  \item \textbf{A}
  \item \textbf{A}
  \item \textbf{A}
  \item \textbf{A}
  \item \textbf{A}
  \item \textbf{A}
  \item \textbf{A}
  \item \textbf{A}
  \item \textbf{A}
  \item \textbf{A}
  \item \textbf{A}
  \item \textbf{A}
\end{enumerate}
\end{multicols}
\end{tcolorbox}
\newpage\n
\subsection{2.3 SINTAXIS}

\subsubsection*{SECCIÓN I: CONCEPTOS BÁSICOS DE SINTAXIS (1–40)}

\begin{enumerate}[label=\textbf{\arabic*.}]
\setcounter{enumi}{0}
\item ¿Qué estudia la sintaxis?
\begin{enumerate}[label=\textbf{\Alph*)}]
\begin{multicols}{2}
  \item La estructura de las oraciones y las relaciones entre sus constituyentes
  \item La ortografía de las palabras
  \item La pronunciación de las palabras
  \item La formación de las palabras
\end{multicols}
\end{enumerate}
\vspace{0.15cm}
\item ¿Qué es una oración?
\begin{enumerate}[label=\textbf{\Alph*)}]
\begin{multicols}{2}
  \item Una unidad lingüística con sentido completo que contiene verbo conjugado
  \item Un conjunto de palabras sin sentido
  \item Una palabra aislada
  \item Un conjunto de letras
  \item Un conjunto de sílabas
  \item Un conjunto de fonemas
\end{multicols}
\end{enumerate}
\vspace{0.15cm}
\item La oración simple se caracteriza por:
\begin{enumerate}[label=\textbf{\Alph*)}]
\begin{multicols}{2}
  \item Tener un solo verbo conjugado
  \item Tener dos o más verbos conjugados
  \item No tener verbo
  \item Tener solo sustantivos
  \item Tener solo adjetivos
  \item Tener solo adverbios
\end{multicols}
\end{enumerate}
\vspace{0.15cm}
\item La oración compuesta se caracteriza por:
\begin{enumerate}[label=\textbf{\Alph*)}]
\begin{multicols}{2}
  \item Tener dos o más verbos conjugados
  \item Tener un solo verbo conjugado
  \item No tener verbo
  \item Tener solo sustantivos
  \item Tener solo adjetivos
  \item Tener solo adverbios
\end{multicols}
\end{enumerate}
\vspace{0.15cm}
\item ¿Qué es un sintagma?
\begin{enumerate}[label=\textbf{\Alph*)}]
\begin{multicols}{2}
  \item Un grupo de palabras que desempeña una función sintáctica
  \item Una palabra aislada
  \item Un conjunto de oraciones
  \item Un conjunto de párrafos
  \item Un conjunto de textos
  \item Un conjunto de letras
\end{multicols}
\end{enumerate}
\vspace{0.15cm}
\item ¿Cuántos tipos principales de sintagmas existen?
\begin{enumerate}[label=\textbf{\Alph*)}]
\begin{multicols}{2}
  \item Tres: nominal, verbal y adjetival
  \item Uno: nominal
  \item Seis: nominal, verbal, adjetival, adverbial, preposicional y conjuntivo
\end{multicols}
\end{enumerate}
\vspace{0.15cm}
\item El núcleo del sintagma nominal es:
\begin{enumerate}[label=\textbf{\Alph*)}]
\begin{multicols}{2}
  \item Un sustantivo o pronombre
  \item Un verbo
  \item Un adjetivo
  \item Un adverbio
  \item Una preposición
  \item Una conjunción
\end{multicols}
\end{enumerate}
\vspace{0.15cm}
\item El núcleo del sintagma verbal es:
\begin{enumerate}[label=\textbf{\Alph*)}]
\begin{multicols}{2}
  \item Un verbo
  \item Un sustantivo
  \item Un adjetivo
  \item Un adverbio
  \item Una preposición
  \item Una conjunción
\end{multicols}
\end{enumerate}
\vspace{0.15cm}
\item El núcleo del sintagma adjetival es:
\begin{enumerate}[label=\textbf{\Alph*)}]
\begin{multicols}{2}
  \item Un adjetivo
  \item Un sustantivo
  \item Un verbo
  \item Un adverbio
  \item Una preposición
  \item Una conjunción
\end{multicols}
\end{enumerate}
\vspace{0.15cm}
\item El núcleo del sintagma adverbial es:
\begin{enumerate}[label=\textbf{\Alph*)}]
\begin{multicols}{2}
  \item Un adverbio
  \item Un sustantivo
  \item Un verbo
  \item Un adjetivo
  \item Una preposición
  \item Una conjunción
\end{multicols}
\end{enumerate}
\vspace{0.15cm}
\item El sintagma preposicional se caracteriza por:
\begin{enumerate}[label=\textbf{\Alph*)}]
\begin{multicols}{2}
\end{multicols}
\end{enumerate}
\vspace{0.15cm}
\item La función sintáctica se refiere a:
\begin{enumerate}[label=\textbf{\Alph*)}]
\begin{multicols}{2}
  \item La ortografía de la palabra
  \item La pronunciación de la palabra
  \item La formación de la palabra
\end{multicols}
\end{enumerate}
\vspace{0.15cm}
\item Los constituyentes inmediatos de la oración son:
\begin{enumerate}[label=\textbf{\Alph*)}]
\begin{multicols}{2}
  \item Sujeto y predicado
  \item Sustantivo y adjetivo
  \item Verbo y adverbio
  \item Preposición y conjunción
  \item Pronombre y verbo
\end{multicols}
\end{enumerate}
\vspace{0.15cm}
\item El sujeto de una oración es:
\begin{enumerate}[label=\textbf{\Alph*)}]
\begin{multicols}{2}
  \item Lo que se dice del sujeto
\end{multicols}
\end{enumerate}
\vspace{0.15cm}
\item El predicado de una oración es:
\begin{enumerate}[label=\textbf{\Alph*)}]
\begin{multicols}{2}
  \item Lo que se dice del sujeto
\end{multicols}
\end{enumerate}
\vspace{0.15cm}
\item La concordancia entre sujeto y predicado se da en:
\begin{enumerate}[label=\textbf{\Alph*)}]
\begin{multicols}{2}
  \item Número y persona
  \item Género y número
  \item Persona y género
  \item Tiempo y modo
  \item Modo y número
\end{multicols}
\end{enumerate}
\vspace{0.15cm}
\item El sujeto léxico es:
\begin{enumerate}[label=\textbf{\Alph*)}]
\begin{multicols}{2}
\end{multicols}
\end{enumerate}
\vspace{0.15cm}
\item El sujeto gramatical o elidido es:
\begin{enumerate}[label=\textbf{\Alph*)}]
\begin{multicols}{2}
\end{multicols}
\end{enumerate}
\vspace{0.15cm}
\item El sujeto puede estar formado por:
\begin{enumerate}[label=\textbf{\Alph*)}]
\begin{multicols}{2}
  \item Un sustantivo, pronombre, infinitivo o una oración
  \item Solo un sustantivo
  \item Solo un pronombre
  \item Solo un infinitivo
  \item Solo una oración
  \item Solo un adjetivo
\end{multicols}
\end{enumerate}
\vspace{0.15cm}
\item El predicado puede ser:
\begin{enumerate}[label=\textbf{\Alph*)}]
\begin{multicols}{2}
  \item Nominal o verbal
  \item Solo nominal
  \item Solo verbal
  \item Solo adjetival
  \item Solo adverbial
  \item Solo preposicional
\end{multicols}
\end{enumerate}
\vspace{0.15cm}
\item El predicado nominal se caracteriza por:
\begin{enumerate}[label=\textbf{\Alph*)}]
\begin{multicols}{2}
  \item Llevar verbos predicativos
  \item No llevar verbo
  \item Llevar solo adverbios
  \item Llevar solo adjetivos
  \item Llevar solo sustantivos
\end{multicols}
\end{enumerate}
\vspace{0.15cm}
\item El predicado verbal se caracteriza por:
\begin{enumerate}[label=\textbf{\Alph*)}]
\begin{multicols}{2}
  \item Llevar verbos predicativos
  \item Llevar verbos copulativos
  \item No llevar verbo
  \item Llevar solo adverbios
  \item Llevar solo adjetivos
  \item Llevar solo sustantivos
\end{multicols}
\end{enumerate}
\vspace{0.15cm}
\item Los verbos copulativos son:
\begin{enumerate}[label=\textbf{\Alph*)}]
\begin{multicols}{2}
  \item Ser, estar y parecer
  \item Ser, estar y tener
  \item Ser, haber y tener
  \item Ser, estar y haber
  \item Parecer, tener y hacer
\end{multicols}
\end{enumerate}
\vspace{0.15cm}
\item El atributo es:
\begin{enumerate}[label=\textbf{\Alph*)}]
\begin{multicols}{2}
\end{multicols}
\end{enumerate}
\vspace{0.15cm}
\item El complemento directo es:
\begin{enumerate}[label=\textbf{\Alph*)}]
\begin{multicols}{2}
\end{multicols}
\end{enumerate}
\vspace{0.15cm}
\item El complemento indirecto es:
\begin{enumerate}[label=\textbf{\Alph*)}]
\begin{multicols}{2}
\end{multicols}
\end{enumerate}
\vspace{0.15cm}
\item El complemento circunstancial es:
\begin{enumerate}[label=\textbf{\Alph*)}]
\begin{multicols}{2}
\end{multicols}
\end{enumerate}
\vspace{0.15cm}
\item El complemento agente es:
\begin{enumerate}[label=\textbf{\Alph*)}]
\begin{multicols}{2}
\end{multicols}
\end{enumerate}
\vspace{0.15cm}
\item El complemento predicativo es:
\begin{enumerate}[label=\textbf{\Alph*)}]
\begin{multicols}{2}
\end{multicols}
\end{enumerate}
\vspace{0.15cm}
\item El complemento del nombre es:
\begin{enumerate}[label=\textbf{\Alph*)}]
\begin{multicols}{2}
\end{multicols}
\end{enumerate}
\vspace{0.15cm}
\item El complemento del adjetivo es:
\begin{enumerate}[label=\textbf{\Alph*)}]
\begin{multicols}{2}
\end{multicols}
\end{enumerate}
\vspace{0.15cm}
\item El complemento del adverbio es:
\begin{enumerate}[label=\textbf{\Alph*)}]
\begin{multicols}{2}
\end{multicols}
\end{enumerate}
\vspace{0.15cm}
\item La aposición es:
\begin{enumerate}[label=\textbf{\Alph*)}]
\begin{multicols}{2}
  \item Un complemento del nombre que explica o identifica al sustantivo
  \item Un complemento del verbo
  \item Un complemento del adjetivo
  \item Un complemento del adverbio
  \item Un complemento de la oración
  \item Un complemento del pronombre
\end{multicols}
\end{enumerate}
\vspace{0.15cm}
\item Los elementos extraoracionales son:
\begin{enumerate}[label=\textbf{\Alph*)}]
\begin{multicols}{2}
  \item Vocativos, interjecciones y marcadores del discurso
  \item Sujeto y predicado
\end{multicols}
\end{enumerate}
\vspace{0.15cm}
\item El vocativo es:
\begin{enumerate}[label=\textbf{\Alph*)}]
\begin{multicols}{2}
\end{multicols}
\end{enumerate}
\vspace{0.15cm}
\item Las interjecciones son:
\begin{enumerate}[label=\textbf{\Alph*)}]
\begin{multicols}{2}
  \item Palabras que expresan emociones o reacciones
  \item Palabras que modifican al sustantivo
  \item Palabras que modifican al verbo
  \item Palabras que modifican al adjetivo
  \item Palabras que unen oraciones
  \item Palabras que indican circunstancias
\end{multicols}
\end{enumerate}
\vspace{0.15cm}
\item Los marcadores del discurso son:
\begin{enumerate}[label=\textbf{\Alph*)}]
\begin{multicols}{2}
\end{multicols}
\end{enumerate}
\vspace{0.15cm}
\item La oración impersonal se caracteriza por:
\begin{enumerate}[label=\textbf{\Alph*)}]
\begin{multicols}{2}
  \item No tener sujeto léxico ni gramatical
  \item Tener sujeto léxico
  \item Tener sujeto gramatical
  \item Tener dos sujetos
  \item No tener predicado
  \item No tener verbo
\end{multicols}
\end{enumerate}
\vspace{0.15cm}
\item La oración pasiva se caracteriza por:
\begin{enumerate}[label=\textbf{\Alph*)}]
\begin{multicols}{2}
  \item No tener sujeto
  \item No tener verbo
  \item No tener predicado
  \item Tener dos verbos
\end{multicols}
\end{enumerate}
\vspace{0.15cm}
\item La oración activa se caracteriza por:
\begin{enumerate}[label=\textbf{\Alph*)}]
\begin{multicols}{2}
  \item No tener sujeto
  \item No tener verbo
  \item No tener predicado
  \item Tener dos verbos ---
\end{multicols}
\end{enumerate}
\vspace{0.15cm}
\end{enumerate}\n
\subsubsection*{SECCIÓN II: EL SINTAGMA NOMINAL Y SUS COMPONENTES (41–90)}

\begin{enumerate}[label=\textbf{\arabic*.}]
\setcounter{enumi}{40}
\item El sintagma nominal tiene como núcleo:
\begin{enumerate}[label=\textbf{\Alph*)}]
\begin{multicols}{2}
  \item Un sustantivo o pronombre
  \item Un verbo
  \item Un adjetivo
  \item Un adverbio
  \item Una preposición
  \item Una conjunción
\end{multicols}
\end{enumerate}
\vspace{0.15cm}
\item Los determinantes del sintagma nominal son:
\begin{enumerate}[label=\textbf{\Alph*)}]
\begin{multicols}{2}
  \item Solo artículos
  \item Solo demostrativos
  \item Solo posesivos
  \item Solo numerales
  \item Solo indefinidos
\end{multicols}
\end{enumerate}
\vspace{0.15cm}
\item El artículo determinado es:
\begin{enumerate}[label=\textbf{\Alph*)}]
\begin{multicols}{2}
  \item Un, una, unos, unas
  \item Mi, tu, su
  \item Uno, dos, tres
\end{multicols}
\end{enumerate}
\vspace{0.15cm}
\item El artículo indeterminado es:
\begin{enumerate}[label=\textbf{\Alph*)}]
\begin{multicols}{2}
  \item Un, una, unos, unas
  \item Mi, tu, su
  \item Uno, dos, tres
\end{multicols}
\end{enumerate}
\vspace{0.15cm}
\item Los demostrativos son:
\begin{enumerate}[label=\textbf{\Alph*)}]
\begin{multicols}{2}
  \item Un, una, unos, unas
  \item Mi, tu, su
  \item Uno, dos, tres
\end{multicols}
\end{enumerate}
\vspace{0.15cm}
\item Los posesivos son:
\begin{enumerate}[label=\textbf{\Alph*)}]
\begin{multicols}{2}
  \item Mi, tu, su, nuestro, vuestro y sus variantes
  \item Un, una, unos, unas
  \item Uno, dos, tres
\end{multicols}
\end{enumerate}
\vspace{0.15cm}
\item Los numerales pueden ser:
\begin{enumerate}[label=\textbf{\Alph*)}]
\begin{multicols}{2}
  \item Solo cardinales
  \item Solo ordinales
  \item Solo partitivos
  \item Solo multiplicativos
  \item Solo romanos
\end{multicols}
\end{enumerate}
\vspace{0.15cm}
\item Los indefinidos son:
\begin{enumerate}[label=\textbf{\Alph*)}]
\begin{multicols}{2}
  \item Un, una, unos, unas
  \item Mi, tu, su
  \item Uno, dos, tres
\end{multicols}
\end{enumerate}
\vspace{0.15cm}
\item El adjetivo dentro del sintagma nominal funciona como:
\begin{enumerate}[label=\textbf{\Alph*)}]
\begin{multicols}{2}
  \item Modificador o complemento del nombre
  \item Núcleo del sintagma
  \item Verbo
  \item Preposición
\end{multicols}
\end{enumerate}
\vspace{0.15cm}
\item La concordancia entre sustantivo y adjetivo se da en:
\begin{enumerate}[label=\textbf{\Alph*)}]
\begin{multicols}{2}
  \item Género y número
  \item Persona y número
  \item Persona y género
  \item Tiempo y modo
  \item Modo y número
\end{multicols}
\end{enumerate}
\vspace{0.15cm}
\item La aposición puede ser:
\begin{enumerate}[label=\textbf{\Alph*)}]
\begin{multicols}{2}
  \item Solo explicativa
  \item Solo especificativa
  \item Solo restrictiva
  \item Solo no restrictiva
  \item Solo determinativa
\end{multicols}
\end{enumerate}
\vspace{0.15cm}
\item La aposición explicativa se separa mediante:
\begin{enumerate}[label=\textbf{\Alph*)}]
\begin{multicols}{2}
  \item Punto y coma
  \item Puntos suspensivos
  \item Paréntesis
\end{multicols}
\end{enumerate}
\vspace{0.15cm}
\item El complemento del nombre puede ser:
\begin{enumerate}[label=\textbf{\Alph*)}]
\begin{multicols}{2}
  \item Un sintagma preposicional, un adjetivo o una aposición
  \item Solo un sintagma preposicional
  \item Solo un adjetivo
  \item Solo una aposición
  \item Solo un adverbio
  \item Solo un verbo
\end{multicols}
\end{enumerate}
\vspace{0.15cm}
\item En "el libro de Juan", "de Juan" es:
\begin{enumerate}[label=\textbf{\Alph*)}]
\begin{multicols}{2}
  \item Predicativo
\end{multicols}
\end{enumerate}
\vspace{0.15cm}
\item En "la casa blanca", "blanca" es:
\begin{enumerate}[label=\textbf{\Alph*)}]
\begin{multicols}{2}
  \item Núcleo del sintagma
\end{multicols}
\end{enumerate}
\vspace{0.15cm}
\item En "mi amigo Pedro", "Pedro" es:
\begin{enumerate}[label=\textbf{\Alph*)}]
\begin{multicols}{2}
  \item Predicativo
\end{multicols}
\end{enumerate}
\vspace{0.15cm}
\item El núcleo del sintagma nominal puede ser:
\begin{enumerate}[label=\textbf{\Alph*)}]
\begin{multicols}{2}
  \item Un sustantivo, pronombre, infinitivo o una oración sustantivada
  \item Solo un sustantivo
  \item Solo un pronombre
  \item Solo un infinitivo
  \item Solo una oración
  \item Solo un adjetivo
\end{multicols}
\end{enumerate}
\vspace{0.15cm}
\item En "los tres libros", "los" y "tres" son:
\begin{enumerate}[label=\textbf{\Alph*)}]
\begin{multicols}{2}
  \item Núcleos
  \item Verbos
\end{multicols}
\end{enumerate}
\vspace{0.15cm}
\item En "este coche rojo", "este" es:
\begin{enumerate}[label=\textbf{\Alph*)}]
\begin{multicols}{2}
  \item Sustantivo
  \item Pronombre
  \item Verbo
\end{multicols}
\end{enumerate}
\vspace{0.15cm}
\item En "este coche rojo", "rojo" es:
\begin{enumerate}[label=\textbf{\Alph*)}]
\begin{multicols}{2}
  \item Sustantivo
  \item Pronombre
  \item Verbo
\end{multicols}
\end{enumerate}
\vspace{0.15cm}
\item En "muchos alumnos", "muchos" es:
\begin{enumerate}[label=\textbf{\Alph*)}]
\begin{multicols}{2}
  \item Sustantivo
  \item Pronombre
  \item Verbo
\end{multicols}
\end{enumerate}
\vspace{0.15cm}
\item En "el primer día", "primer" es:
\begin{enumerate}[label=\textbf{\Alph*)}]
\begin{multicols}{2}
  \item Sustantivo
  \item Pronombre
  \item Verbo
\end{multicols}
\end{enumerate}
\vspace{0.15cm}
\item En "dos casas", "dos" es:
\begin{enumerate}[label=\textbf{\Alph*)}]
\begin{multicols}{2}
  \item Sustantivo
  \item Pronombre
  \item Verbo
\end{multicols}
\end{enumerate}
\vspace{0.15cm}
\item El sintagma nominal puede desempeñar la función de:
\begin{enumerate}[label=\textbf{\Alph*)}]
\begin{multicols}{2}
  \item Sujeto, complemento directo, complemento indirecto, etc.
  \item Solo sujeto
  \item Solo complemento directo
  \item Solo complemento indirecto
  \item Solo complemento circunstancial
  \item Solo atributo
\end{multicols}
\end{enumerate}
\vspace{0.15cm}
\item En "Juan compró pan", "Juan" es:
\begin{enumerate}[label=\textbf{\Alph*)}]
\begin{multicols}{2}
  \item Sujeto
  \item Predicado
\end{multicols}
\end{enumerate}
\vspace{0.15cm}
\item En "Juan compró pan", "pan" es:
\begin{enumerate}[label=\textbf{\Alph*)}]
\begin{multicols}{2}
  \item Sujeto
  \item Predicativo
\end{multicols}
\end{enumerate}
\vspace{0.15cm}
\item En "Juan dio el libro a María", "a María" es:
\begin{enumerate}[label=\textbf{\Alph*)}]
\begin{multicols}{2}
  \item Sujeto
  \item Predicativo
\end{multicols}
\end{enumerate}
\vspace{0.15cm}
\item En "Juan dio el libro a María", "el libro" es:
\begin{enumerate}[label=\textbf{\Alph*)}]
\begin{multicols}{2}
  \item Sujeto
  \item Predicativo
\end{multicols}
\end{enumerate}
\vspace{0.15cm}
\item El complemento directo puede sustituirse por:
\begin{enumerate}[label=\textbf{\Alph*)}]
\begin{multicols}{2}
  \item Lo, la, los, las
  \item Le, les
  \item Me, te, se
  \item Nos, os
  \item Él, ella, ellos
\end{multicols}
\end{enumerate}
\vspace{0.15cm}
\item El complemento indirecto puede sustituirse por:
\begin{enumerate}[label=\textbf{\Alph*)}]
\begin{multicols}{2}
  \item Le, les
  \item Lo, la, los, las
  \item Me, te, se
  \item Nos, os
  \item Él, ella, ellos
\end{multicols}
\end{enumerate}
\vspace{0.15cm}
\item El sujeto puede sustituirse por:
\begin{enumerate}[label=\textbf{\Alph*)}]
\begin{multicols}{2}
  \item Él, ella, ellos, ellas
  \item Lo, la, los, las
  \item Le, les
  \item Me, te, se
  \item Nos, os
\end{multicols}
\end{enumerate}
\vspace{0.15cm}
\item En "Se lo dije", "se" sustituye a:
\begin{enumerate}[label=\textbf{\Alph*)}]
\begin{multicols}{2}
  \item Sujeto
  \item Predicativo
\end{multicols}
\end{enumerate}
\vspace{0.15cm}
\item En "Se lo dije", "lo" sustituye a:
\begin{enumerate}[label=\textbf{\Alph*)}]
\begin{multicols}{2}
  \item Sujeto
  \item Predicativo
\end{multicols}
\end{enumerate}
\vspace{0.15cm}
\item El atributo puede sustituirse por:
\begin{enumerate}[label=\textbf{\Alph*)}]
\begin{multicols}{2}
  \item Lo
  \item Le
  \item La
  \item Los
  \item Las
  \item Les
\end{multicols}
\end{enumerate}
\vspace{0.15cm}
\item En "Juan es médico", "médico" es:
\begin{enumerate}[label=\textbf{\Alph*)}]
\begin{multicols}{2}
  \item Predicativo
  \item Sujeto
\end{multicols}
\end{enumerate}
\vspace{0.15cm}
\item En "Juan está cansado", "cansado" es:
\begin{enumerate}[label=\textbf{\Alph*)}]
\begin{multicols}{2}
  \item Predicativo
  \item Sujeto
\end{multicols}
\end{enumerate}
\vspace{0.15cm}
\item En "Juan parece inteligente", "inteligente" es:
\begin{enumerate}[label=\textbf{\Alph*)}]
\begin{multicols}{2}
  \item Predicativo
  \item Sujeto
\end{multicols}
\end{enumerate}
\vspace{0.15cm}
\item El atributo concuerda con el sujeto en:
\begin{enumerate}[label=\textbf{\Alph*)}]
\begin{multicols}{2}
  \item Género y número
  \item Persona y número
  \item Persona y género
  \item Tiempo y modo
  \item Modo y número
\end{multicols}
\end{enumerate}
\vspace{0.15cm}
\item El predicativo concuerda con:
\begin{enumerate}[label=\textbf{\Alph*)}]
\begin{multicols}{2}
  \item Solo el sujeto
  \item Solo el complemento directo
  \item Solo el complemento indirecto
  \item Solo el verbo
  \item Solo el adverbio
\end{multicols}
\end{enumerate}
\vspace{0.15cm}
\item En "Juan llegó cansado", "cansado" es:
\begin{enumerate}[label=\textbf{\Alph*)}]
\begin{multicols}{2}
  \item Predicativo
  \item Sujeto
\end{multicols}
\end{enumerate}
\vspace{0.15cm}
\item En "Considero a Juan inteligente", "inteligente" es:
\begin{enumerate}[label=\textbf{\Alph*)}]
\begin{multicols}{2}
  \item Predicativo del complemento directo
  \item Sujeto
\end{multicols}
\end{enumerate}
\vspace{0.15cm}
\item El complemento circunstancial puede ser de:
\begin{enumerate}[label=\textbf{\Alph*)}]
\begin{multicols}{2}
  \item Lugar, tiempo, modo, causa, finalidad, etc.
  \item Solo lugar
  \item Solo tiempo
  \item Solo modo
  \item Solo causa
  \item Solo finalidad
\end{multicols}
\end{enumerate}
\vspace{0.15cm}
\item En "Juan estudia en la biblioteca", "en la biblioteca" es:
\begin{enumerate}[label=\textbf{\Alph*)}]
\begin{multicols}{2}
  \item Predicativo
  \item Sujeto
\end{multicols}
\end{enumerate}
\vspace{0.15cm}
\item En "Juan llegó ayer", "ayer" es:
\begin{enumerate}[label=\textbf{\Alph*)}]
\begin{multicols}{2}
  \item Predicativo
  \item Sujeto
\end{multicols}
\end{enumerate}
\vspace{0.15cm}
\item En "Juan estudia con atención", "con atención" es:
\begin{enumerate}[label=\textbf{\Alph*)}]
\begin{multicols}{2}
  \item Predicativo
  \item Sujeto
\end{multicols}
\end{enumerate}
\vspace{0.15cm}
\item En "Juan no vino por la lluvia", "por la lluvia" es:
\begin{enumerate}[label=\textbf{\Alph*)}]
\begin{multicols}{2}
  \item Predicativo
  \item Sujeto
\end{multicols}
\end{enumerate}
\vspace{0.15cm}
\item En "Juan estudia para aprobar", "para aprobar" es:
\begin{enumerate}[label=\textbf{\Alph*)}]
\begin{multicols}{2}
  \item Predicativo
  \item Sujeto
\end{multicols}
\end{enumerate}
\vspace{0.15cm}
\item El complemento agente aparece en:
\begin{enumerate}[label=\textbf{\Alph*)}]
\begin{multicols}{2}
  \item Oraciones pasivas
  \item Oraciones activas
  \item Oraciones impersonales
  \item Oraciones copulativas
  \item Oraciones transitivas
  \item Oraciones intransitivas
\end{multicols}
\end{enumerate}
\vspace{0.15cm}
\item En "El libro fue escrito por Juan", "por Juan" es:
\begin{enumerate}[label=\textbf{\Alph*)}]
\begin{multicols}{2}
  \item Predicativo
\end{multicols}
\end{enumerate}
\vspace{0.15cm}
\item En "El libro fue escrito por Juan", "El libro" es:
\begin{enumerate}[label=\textbf{\Alph*)}]
\begin{multicols}{2}
  \item Sujeto paciente
  \item Sujeto agente
\end{multicols}
\end{enumerate}
\vspace{0.15cm}
\end{enumerate}\n
\subsubsection*{SECCIÓN III: EL SINTAGMA VERBAL Y SUS COMPONENTES (91–150)}

\begin{enumerate}[label=\textbf{\arabic*.}]
\setcounter{enumi}{90}
\item El sintagma verbal tiene como núcleo:
\begin{enumerate}[label=\textbf{\Alph*)}]
\begin{multicols}{2}
  \item Un verbo
  \item Un sustantivo
  \item Un adjetivo
  \item Un adverbio
  \item Una preposición
  \item Una conjunción
\end{multicols}
\end{enumerate}
\vspace{0.15cm}
\item El verbo expresa:
\begin{enumerate}[label=\textbf{\Alph*)}]
\begin{multicols}{2}
  \item Solo acción
  \item Solo proceso
  \item Solo estado
  \item Solo tiempo
  \item Solo modo
\end{multicols}
\end{enumerate}
\vspace{0.15cm}
\item Las formas no personales del verbo son:
\begin{enumerate}[label=\textbf{\Alph*)}]
\begin{multicols}{2}
  \item Infinitivo, gerundio y participio
  \item Solo infinitivo
  \item Solo gerundio
  \item Solo participio
  \item Solo imperativo
  \item Solo subjuntivo
\end{multicols}
\end{enumerate}
\vspace{0.15cm}
\item El infinitivo termina en:
\begin{enumerate}[label=\textbf{\Alph*)}]
\begin{multicols}{2}
  \item -ar, -er, -ir
  \item -ando, -iendo
  \item -ado, -ido
  \item -o, -as, -a
  \item -e, -es, -e
  \item -ría, -rías
\end{multicols}
\end{enumerate}
\vspace{0.15cm}
\item El gerundio termina en:
\begin{enumerate}[label=\textbf{\Alph*)}]
\begin{multicols}{2}
  \item -ando, -iendo
  \item -ar, -er, -ir
  \item -ado, -ido
  \item -o, -as, -a
  \item -e, -es, -e
  \item -ría, -rías
\end{multicols}
\end{enumerate}
\vspace{0.15cm}
\item El participio termina en:
\begin{enumerate}[label=\textbf{\Alph*)}]
\begin{multicols}{2}
  \item -ado, -ido
  \item -ar, -er, -ir
  \item -ando, -iendo
  \item -o, -as, -a
  \item -e, -es, -e
  \item -ría, -rías
\end{multicols}
\end{enumerate}
\vspace{0.15cm}
\item Los verbos copulativos son:
\begin{enumerate}[label=\textbf{\Alph*)}]
\begin{multicols}{2}
  \item Ser, estar, parecer
  \item Ser, estar, tener
  \item Ser, haber, tener
  \item Ser, estar, haber
  \item Parecer, tener, hacer
\end{multicols}
\end{enumerate}
\vspace{0.15cm}
\item Los verbos predicativos son:
\begin{enumerate}[label=\textbf{\Alph*)}]
\begin{multicols}{2}
  \item Todos los que no son copulativos
  \item Solo ser, estar, parecer
  \item Solo los transitivos
  \item Solo los intransitivos
  \item Solo los reflexivos
  \item Solo los impersonales
\end{multicols}
\end{enumerate}
\vspace{0.15cm}
\item Los verbos transitivos son:
\begin{enumerate}[label=\textbf{\Alph*)}]
\begin{multicols}{2}
  \item Los que necesitan complemento directo
  \item Los que no necesitan complemento directo
  \item Los que necesitan complemento indirecto
  \item Los que necesitan complemento circunstancial
  \item Los que necesitan atributo
  \item Los que necesitan predicativo
\end{multicols}
\end{enumerate}
\vspace{0.15cm}
\item Los verbos intransitivos son:
\begin{enumerate}[label=\textbf{\Alph*)}]
\begin{multicols}{2}
  \item Los que no necesitan complemento directo
  \item Los que necesitan complemento directo
  \item Los que necesitan complemento indirecto
  \item Los que necesitan complemento circunstancial
  \item Los que necesitan atributo
  \item Los que necesitan predicativo
\end{multicols}
\end{enumerate}
\vspace{0.15cm}
\item Los verbos reflexivos son:
\begin{enumerate}[label=\textbf{\Alph*)}]
\begin{multicols}{2}
  \item Los que llevan pronombre reflexivo
  \item Los que no llevan pronombre
  \item Los que llevan complemento directo
  \item Los que llevan complemento indirecto
  \item Los que llevan atributo
  \item Los que llevan predicativo
\end{multicols}
\end{enumerate}
\vspace{0.15cm}
\item Los verbos recíprocos son:
\begin{enumerate}[label=\textbf{\Alph*)}]
\begin{multicols}{2}
  \item Los que expresan acción mutua
  \item Los que expresan acción reflexiva
  \item Los que expresan acción transitiva
  \item Los que expresan acción intransitiva
  \item Los que expresan acción copulativa
  \item Los que expresan acción impersonal
\end{multicols}
\end{enumerate}
\vspace{0.15cm}
\item Los verbos impersonales son:
\begin{enumerate}[label=\textbf{\Alph*)}]
\begin{multicols}{2}
  \item Los que no tienen sujeto
  \item Los que tienen sujeto
  \item Los que tienen dos sujetos
  \item Los que tienen sujeto elidido
  \item Los que tienen sujeto paciente
  \item Los que tienen sujeto agente
\end{multicols}
\end{enumerate}
\vspace{0.15cm}
\item Los verbos auxiliares son:
\begin{enumerate}[label=\textbf{\Alph*)}]
\begin{multicols}{2}
  \item Haber, ser, estar y otros que ayudan a formar tiempos compuestos
  \item Solo haber
  \item Solo ser
  \item Solo estar
  \item Solo tener
  \item Solo hacer
\end{multicols}
\end{enumerate}
\vspace{0.15cm}
\item Los tiempos compuestos se forman con:
\begin{enumerate}[label=\textbf{\Alph*)}]
\begin{multicols}{2}
  \item Haber + participio
  \item Ser + participio
  \item Tener + participio
  \item Ir + infinitivo
  \item Poder + infinitivo
\end{multicols}
\end{enumerate}
\vspace{0.15cm}
\item Las perífrasis verbales se forman con:
\begin{enumerate}[label=\textbf{\Alph*)}]
\begin{multicols}{2}
  \item Verbo auxiliar + verbo en forma no personal
  \item Un verbo y un adverbio
  \item Un verbo y un sustantivo
  \item Un verbo y un adjetivo
  \item Un verbo y una preposición
\end{multicols}
\end{enumerate}
\vspace{0.15cm}
\item "Voy a estudiar" es una perífrasis de:
\begin{enumerate}[label=\textbf{\Alph*)}]
\begin{multicols}{2}
  \item Infinitivo
  \item Gerundio
  \item Participio
  \item Subjuntivo
  \item Imperativo
  \item Indicativo
\end{multicols}
\end{enumerate}
\vspace{0.15cm}
\item "Estoy estudiando" es una perífrasis de:
\begin{enumerate}[label=\textbf{\Alph*)}]
\begin{multicols}{2}
  \item Gerundio
  \item Infinitivo
  \item Participio
  \item Subjuntivo
  \item Imperativo
  \item Indicativo
\end{multicols}
\end{enumerate}
\vspace{0.15cm}
\item "Tengo estudiado" es una perífrasis de:
\begin{enumerate}[label=\textbf{\Alph*)}]
\begin{multicols}{2}
  \item Participio
  \item Infinitivo
  \item Gerundio
  \item Subjuntivo
  \item Imperativo
  \item Indicativo
\end{multicols}
\end{enumerate}
\vspace{0.15cm}
\item El modo indicativo expresa:
\begin{enumerate}[label=\textbf{\Alph*)}]
\begin{multicols}{2}
  \item Realidad o certeza
  \item Mandato
  \item Irrealidad
  \item Hipótesis
\end{multicols}
\end{enumerate}
\vspace{0.15cm}
\item El modo subjuntivo expresa:
\begin{enumerate}[label=\textbf{\Alph*)}]
\begin{multicols}{2}
  \item Realidad o certeza
  \item Mandato
  \item Negación
  \item Pregunta
\end{multicols}
\end{enumerate}
\vspace{0.15cm}
\item El modo imperativo expresa:
\begin{enumerate}[label=\textbf{\Alph*)}]
\begin{multicols}{2}
  \item Mandato o ruego
  \item Realidad o certeza
  \item Irrealidad
  \item Hipótesis
\end{multicols}
\end{enumerate}
\vspace{0.15cm}
\item El tiempo verbal indica:
\begin{enumerate}[label=\textbf{\Alph*)}]
\begin{multicols}{2}
  \item La voz de la acción
  \item La persona de la acción
\end{multicols}
\end{enumerate}
\vspace{0.15cm}
\item El aspecto verbal indica:
\begin{enumerate}[label=\textbf{\Alph*)}]
\begin{multicols}{2}
  \item Si la acción está terminada o en desarrollo
  \item La voz de la acción
  \item La persona de la acción
\end{multicols}
\end{enumerate}
\vspace{0.15cm}
\item La voz activa indica que:
\begin{enumerate}[label=\textbf{\Alph*)}]
\begin{multicols}{2}
  \item No hay sujeto
  \item Hay dos sujetos
\end{multicols}
\end{enumerate}
\vspace{0.15cm}
\item La voz pasiva indica que:
\begin{enumerate}[label=\textbf{\Alph*)}]
\begin{multicols}{2}
  \item No hay sujeto
  \item Hay dos sujetos
\end{multicols}
\end{enumerate}
\vspace{0.15cm}
\item La pasiva se forma con:
\begin{enumerate}[label=\textbf{\Alph*)}]
\begin{multicols}{2}
  \item Ser + participio + por + agente
  \item Haber + participio
  \item Tener + participio
  \item Ir + infinitivo
  \item Poder + infinitivo
\end{multicols}
\end{enumerate}
\vspace{0.15cm}
\item En "La casa fue construida por Juan", la oración es:
\begin{enumerate}[label=\textbf{\Alph*)}]
\begin{multicols}{2}
  \item Pasiva
  \item Impersonal
  \item Reflexiva
  \item Recíproca
\end{multicols}
\end{enumerate}
\vspace{0.15cm}
\item En "Juan construyó la casa", la oración es:
\begin{enumerate}[label=\textbf{\Alph*)}]
\begin{multicols}{2}
  \item Pasiva
  \item Impersonal
  \item Reflexiva
  \item Recíproca
\end{multicols}
\end{enumerate}
\vspace{0.15cm}
\item La pasiva refleja se forma con:
\begin{enumerate}[label=\textbf{\Alph*)}]
\begin{multicols}{2}
  \item Se + verbo en tercera persona
  \item Ser + participio
  \item Haber + participio
  \item Ir + infinitivo
  \item Tener + participio
\end{multicols}
\end{enumerate}
\vspace{0.15cm}
\item En "Se venden casas", la oración es:
\begin{enumerate}[label=\textbf{\Alph*)}]
\begin{multicols}{2}
  \item Pasiva refleja
  \item Pasiva perifrástica
  \item Impersonal
  \item Reflexiva
\end{multicols}
\end{enumerate}
\vspace{0.15cm}
\item En "Se vive bien aquí", la oración es:
\begin{enumerate}[label=\textbf{\Alph*)}]
\begin{multicols}{2}
  \item Impersonal refleja
  \item Pasiva refleja
  \item Reflexiva
  \item Recíproca
\end{multicols}
\end{enumerate}
\vspace{0.15cm}
\item Los verbos de la tercera conjugación terminan en:
\begin{enumerate}[label=\textbf{\Alph*)}]
\begin{multicols}{2}
  \item -ir
  \item -ar
  \item -er
  \item -or
  \item -ur
  \item -ando
\end{multicols}
\end{enumerate}
\vspace{0.15cm}
\item Los verbos de la segunda conjugación terminan en:
\begin{enumerate}[label=\textbf{\Alph*)}]
\begin{multicols}{2}
  \item -er
  \item -ar
  \item -ir
  \item -or
  \item -ur
  \item -iendo
\end{multicols}
\end{enumerate}
\vspace{0.15cm}
\item Los verbos de la primera conjugación terminan en:
\begin{enumerate}[label=\textbf{\Alph*)}]
\begin{multicols}{2}
  \item -ar
  \item -er
  \item -ir
  \item -or
  \item -ur
  \item -ado
\end{multicols}
\end{enumerate}
\vspace{0.15cm}
\item El verbo "ser" es:
\begin{enumerate}[label=\textbf{\Alph*)}]
\begin{multicols}{2}
  \item Predicativo
  \item Transitivo
  \item Intransitivo
  \item Reflexivo
  \item Impersonal
\end{multicols}
\end{enumerate}
\vspace{0.15cm}
\item El verbo "estar" es:
\begin{enumerate}[label=\textbf{\Alph*)}]
\begin{multicols}{2}
  \item Predicativo
  \item Transitivo
  \item Intransitivo
  \item Reflexivo
  \item Impersonal
\end{multicols}
\end{enumerate}
\vspace{0.15cm}
\item El verbo "parecer" es:
\begin{enumerate}[label=\textbf{\Alph*)}]
\begin{multicols}{2}
  \item Predicativo
  \item Transitivo
  \item Intransitivo
  \item Reflexivo
  \item Impersonal
\end{multicols}
\end{enumerate}
\vspace{0.15cm}
\item El verbo "tener" es:
\begin{enumerate}[label=\textbf{\Alph*)}]
\begin{multicols}{2}
  \item Predicativo
  \item Impersonal
  \item Reflexivo
  \item Recíproco
\end{multicols}
\end{enumerate}
\vspace{0.15cm}
\item El verbo "haber" es:
\begin{enumerate}[label=\textbf{\Alph*)}]
\begin{multicols}{2}
  \item Predicativo
  \item Impersonal
  \item Reflexivo
  \item Recíproco
\end{multicols}
\end{enumerate}
\vspace{0.15cm}
\item En "Hay muchos libros", "hay" es:
\begin{enumerate}[label=\textbf{\Alph*)}]
\begin{multicols}{2}
  \item Impersonal
  \item Predicativo
  \item Reflexivo
  \item Recíproco
  \item Transitivo
\end{multicols}
\end{enumerate}
\vspace{0.15cm}
\item En "Llueve mucho", "llueve" es:
\begin{enumerate}[label=\textbf{\Alph*)}]
\begin{multicols}{2}
  \item Impersonal
  \item Predicativo
  \item Reflexivo
  \item Recíproco
  \item Transitivo
\end{multicols}
\end{enumerate}
\vspace{0.15cm}
\item En "Juan se lava", "se" es:
\begin{enumerate}[label=\textbf{\Alph*)}]
\begin{multicols}{2}
  \item Reflexivo
  \item Recíproco
  \item Impersonal
  \item Pasivo
\end{multicols}
\end{enumerate}
\vspace{0.15cm}
\item En "Juan y María se escriben", "se" es:
\begin{enumerate}[label=\textbf{\Alph*)}]
\begin{multicols}{2}
  \item Recíproco
  \item Reflexivo
  \item Impersonal
  \item Pasivo
\end{multicols}
\end{enumerate}
\vspace{0.15cm}
\item En "Se come bien en este restaurante", "se" es:
\begin{enumerate}[label=\textbf{\Alph*)}]
\begin{multicols}{2}
  \item Impersonal
  \item Reflexivo
  \item Recíproco
  \item Pasivo
\end{multicols}
\end{enumerate}
\vspace{0.15cm}
\item En "Se venden libros", "se" es:
\begin{enumerate}[label=\textbf{\Alph*)}]
\begin{multicols}{2}
  \item Pasivo reflejo
  \item Reflexivo
  \item Recíproco
  \item Impersonal
\end{multicols}
\end{enumerate}
\vspace{0.15cm}
\item El complemento directo puede ser:
\begin{enumerate}[label=\textbf{\Alph*)}]
\begin{multicols}{2}
  \item Un sintagma nominal, un pronombre o una oración
  \item Solo un sintagma nominal
  \item Solo un pronombre
  \item Solo una oración
  \item Solo un infinitivo
  \item Solo un gerundio
\end{multicols}
\end{enumerate}
\vspace{0.15cm}
\item El complemento indirecto puede ser:
\begin{enumerate}[label=\textbf{\Alph*)}]
\begin{multicols}{2}
  \item Un sintagma preposicional con "a" o "para"
  \item Solo un sintagma nominal
  \item Solo un pronombre
  \item Solo una oración
  \item Solo un infinitivo
  \item Solo un gerundio
\end{multicols}
\end{enumerate}
\vspace{0.15cm}
\item En "Dije la verdad a Juan", "la verdad" es:
\begin{enumerate}[label=\textbf{\Alph*)}]
\begin{multicols}{2}
  \item Sujeto
  \item Predicativo
\end{multicols}
\end{enumerate}
\vspace{0.15cm}
\item En "Dije la verdad a Juan", "a Juan" es:
\begin{enumerate}[label=\textbf{\Alph*)}]
\begin{multicols}{2}
  \item Sujeto
  \item Predicativo
\end{multicols}
\end{enumerate}
\vspace{0.15cm}
\item En "Compré flores para María", "para María" es:
\begin{enumerate}[label=\textbf{\Alph*)}]
\begin{multicols}{2}
  \item Sujeto
  \item Predicativo
\end{multicols}
\end{enumerate}
\vspace{0.15cm}
\item En "Compré flores en el mercado", "en el mercado" es:
\begin{enumerate}[label=\textbf{\Alph*)}]
\begin{multicols}{2}
  \item Predicativo
  \item Sujeto
\end{multicols}
\end{enumerate}
\vspace{0.15cm}
\item En "Estudio porque quiero aprender", "porque quiero aprender" es:
\begin{enumerate}[label=\textbf{\Alph*)}]
\begin{multicols}{2}
  \item Predicativo
  \item Sujeto
\end{multicols}
\end{enumerate}
\vspace{0.15cm}
\item En "Vine para ayudarte", "para ayudarte" es:
\begin{enumerate}[label=\textbf{\Alph*)}]
\begin{multicols}{2}
  \item Predicativo
  \item Sujeto
\end{multicols}
\end{enumerate}
\vspace{0.15cm}
\item En "Lo hice con mucho cuidado", "con mucho cuidado" es:
\begin{enumerate}[label=\textbf{\Alph*)}]
\begin{multicols}{2}
  \item Predicativo
  \item Sujeto
\end{multicols}
\end{enumerate}
\vspace{0.15cm}
\item En "Llegaré mañana", "mañana" es:
\begin{enumerate}[label=\textbf{\Alph*)}]
\begin{multicols}{2}
  \item Predicativo
  \item Sujeto
\end{multicols}
\end{enumerate}
\vspace{0.15cm}
\item En "Juan es más alto que Pedro", "que Pedro" es:
\begin{enumerate}[label=\textbf{\Alph*)}]
\begin{multicols}{2}
  \item Predicativo
\end{multicols}
\end{enumerate}
\vspace{0.15cm}
\item En "Muy lejos de aquí", "de aquí" es:
\begin{enumerate}[label=\textbf{\Alph*)}]
\begin{multicols}{2}
  \item Predicativo
\end{multicols}
\end{enumerate}
\vspace{0.15cm}
\item En "Cansado de esperar", "de esperar" es:
\begin{enumerate}[label=\textbf{\Alph*)}]
\begin{multicols}{2}
  \item Predicativo
\end{multicols}
\end{enumerate}
\vspace{0.15cm}
\item En "El hombre de la barba", "de la barba" es:
\begin{enumerate}[label=\textbf{\Alph*)}]
\begin{multicols}{2}
  \item Predicativo ---
\end{multicols}
\end{enumerate}
\vspace{0.15cm}
\end{enumerate}\n
\subsubsection*{SECCIÓN IV: SUJETO Y PREDICADO (151–200)}

\begin{enumerate}[label=\textbf{\arabic*.}]
\setcounter{enumi}{150}
\item El sujeto de una oración es:
\begin{enumerate}[label=\textbf{\Alph*)}]
\begin{multicols}{2}
  \item Lo que se dice del sujeto
\end{multicols}
\end{enumerate}
\vspace{0.15cm}
\item El predicado de una oración es:
\begin{enumerate}[label=\textbf{\Alph*)}]
\begin{multicols}{2}
  \item Lo que se dice del sujeto
\end{multicols}
\end{enumerate}
\vspace{0.15cm}
\item En "El niño juega", el sujeto es:
\begin{enumerate}[label=\textbf{\Alph*)}]
\begin{multicols}{2}
  \item juega
  \item niño juega
  \item No hay sujeto
\end{multicols}
\end{enumerate}
\vspace{0.15cm}
\item En "El niño juega", el predicado es:
\begin{enumerate}[label=\textbf{\Alph*)}]
\begin{multicols}{2}
  \item juega
  \item niño
  \item No hay predicado
\end{multicols}
\end{enumerate}
\vspace{0.15cm}
\item En "María y Juan estudian", el sujeto es:
\begin{enumerate}[label=\textbf{\Alph*)}]
\begin{multicols}{2}
  \item María y Juan
  \item estudian
  \item María
  \item Juan
  \item María y Juan estudian
  \item No hay sujeto
\end{multicols}
\end{enumerate}
\vspace{0.15cm}
\item En "María y Juan estudian", el predicado es:
\begin{enumerate}[label=\textbf{\Alph*)}]
\begin{multicols}{2}
  \item estudian
  \item María y Juan
  \item María
  \item Juan
  \item María y Juan estudian
  \item No hay predicado
\end{multicols}
\end{enumerate}
\vspace{0.15cm}
\item El sujeto puede ser:
\begin{enumerate}[label=\textbf{\Alph*)}]
\begin{multicols}{2}
  \item Un sintagma nominal o una oración
  \item Solo un sintagma nominal
  \item Solo una oración
  \item Solo un pronombre
  \item Solo un infinitivo
  \item Solo un gerundio
\end{multicols}
\end{enumerate}
\vspace{0.15cm}
\item El sujeto puede sustituirse por:
\begin{enumerate}[label=\textbf{\Alph*)}]
\begin{multicols}{2}
  \item Él, ella, ellos, ellas
  \item Lo, la, los, las
  \item Le, les
  \item Me, te, se
  \item Nos, os
\end{multicols}
\end{enumerate}
\vspace{0.15cm}
\item La concordancia entre sujeto y verbo se da en:
\begin{enumerate}[label=\textbf{\Alph*)}]
\begin{multicols}{2}
  \item Número y persona
  \item Género y número
  \item Persona y género
  \item Tiempo y modo
  \item Modo y número
\end{multicols}
\end{enumerate}
\vspace{0.15cm}
\item En "Los niños juegan", el sujeto es:
\begin{enumerate}[label=\textbf{\Alph*)}]
\begin{multicols}{2}
  \item Los niños
  \item juegan
  \item Los
  \item niños
  \item Los niños juegan
  \item No hay sujeto
\end{multicols}
\end{enumerate}
\vspace{0.15cm}
\item En "Los niños juegan", el verbo concuerda en:
\begin{enumerate}[label=\textbf{\Alph*)}]
\begin{multicols}{2}
  \item Tercera persona del plural
  \item Tercera persona del singular
  \item Primera persona del plural
  \item Primera persona del singular
  \item Segunda persona del plural
  \item Segunda persona del singular
\end{multicols}
\end{enumerate}
\vspace{0.15cm}
\item En "Yo estudio", el sujeto es:
\begin{enumerate}[label=\textbf{\Alph*)}]
\begin{multicols}{2}
  \item Yo
  \item estudio
  \item Yo estudio
  \item No hay sujeto
  \item estudio yo
  \item Yo y estudio
\end{multicols}
\end{enumerate}
\vspace{0.15cm}
\item En "Yo estudio", el sujeto es:
\begin{enumerate}[label=\textbf{\Alph*)}]
\begin{multicols}{2}
  \item Léxico
  \item Gramatical
  \item Impersonal
  \item Paciente
\end{multicols}
\end{enumerate}
\vspace{0.15cm}
\item En "Estudio", el sujeto es:
\begin{enumerate}[label=\textbf{\Alph*)}]
\begin{multicols}{2}
  \item Gramatical o elidido
  \item Léxico
  \item Impersonal
  \item Paciente
  \item No hay sujeto
\end{multicols}
\end{enumerate}
\vspace{0.15cm}
\item En "Llueve", el sujeto es:
\begin{enumerate}[label=\textbf{\Alph*)}]
\begin{multicols}{2}
  \item Léxico
  \item Gramatical
  \item Paciente
\end{multicols}
\end{enumerate}
\vspace{0.15cm}
\item En "Hay libros", el sujeto es:
\begin{enumerate}[label=\textbf{\Alph*)}]
\begin{multicols}{2}
  \item Léxico
  \item Gramatical
  \item Paciente
\end{multicols}
\end{enumerate}
\vspace{0.15cm}
\item El predicado nominal lleva:
\begin{enumerate}[label=\textbf{\Alph*)}]
\begin{multicols}{2}
  \item Verbo copulativo + atributo
  \item Verbo predicativo + complementos
  \item Solo verbo
  \item Solo atributo
  \item Solo complemento directo
  \item Solo complemento indirecto
\end{multicols}
\end{enumerate}
\vspace{0.15cm}
\item El predicado verbal lleva:
\begin{enumerate}[label=\textbf{\Alph*)}]
\begin{multicols}{2}
  \item Verbo predicativo + complementos
  \item Verbo copulativo + atributo
  \item Solo verbo
  \item Solo atributo
  \item Solo complemento directo
  \item Solo complemento indirecto
\end{multicols}
\end{enumerate}
\vspace{0.15cm}
\item En "Juan es médico", el predicado es:
\begin{enumerate}[label=\textbf{\Alph*)}]
\begin{multicols}{2}
  \item Nominal
  \item Verbal
  \item Impersonal
  \item Reflexivo
  \item Recíproco
  \item Pasivo
\end{multicols}
\end{enumerate}
\vspace{0.15cm}
\item En "Juan corre", el predicado es:
\begin{enumerate}[label=\textbf{\Alph*)}]
\begin{multicols}{2}
  \item Verbal
  \item Nominal
  \item Impersonal
  \item Reflexivo
  \item Recíproco
  \item Pasivo
\end{multicols}
\end{enumerate}
\vspace{0.15cm}
\item En "Juan está cansado", el predicado es:
\begin{enumerate}[label=\textbf{\Alph*)}]
\begin{multicols}{2}
  \item Nominal
  \item Verbal
  \item Impersonal
  \item Reflexivo
  \item Recíproco
  \item Pasivo
\end{multicols}
\end{enumerate}
\vspace{0.15cm}
\item En "Juan parece inteligente", el predicado es:
\begin{enumerate}[label=\textbf{\Alph*)}]
\begin{multicols}{2}
  \item Nominal
  \item Verbal
  \item Impersonal
  \item Reflexivo
  \item Recíproco
  \item Pasivo
\end{multicols}
\end{enumerate}
\vspace{0.15cm}
\item En "Juan compró un libro", el predicado es:
\begin{enumerate}[label=\textbf{\Alph*)}]
\begin{multicols}{2}
  \item Verbal
  \item Nominal
  \item Impersonal
  \item Reflexivo
  \item Recíproco
  \item Pasivo
\end{multicols}
\end{enumerate}
\vspace{0.15cm}
\item El atributo en el predicado nominal puede ser:
\begin{enumerate}[label=\textbf{\Alph*)}]
\begin{multicols}{2}
  \item Un sintagma nominal, adjetival o preposicional
  \item Solo un sintagma nominal
  \item Solo un sintagma adjetival
  \item Solo un sintagma preposicional
  \item Solo un adverbio
  \item Solo un verbo
\end{multicols}
\end{enumerate}
\vspace{0.15cm}
\item En "Juan es profesor", "profesor" es:
\begin{enumerate}[label=\textbf{\Alph*)}]
\begin{multicols}{2}
  \item Predicativo
  \item Sujeto
\end{multicols}
\end{enumerate}
\vspace{0.15cm}
\item En "Juan es alto", "alto" es:
\begin{enumerate}[label=\textbf{\Alph*)}]
\begin{multicols}{2}
  \item Predicativo
  \item Sujeto
\end{multicols}
\end{enumerate}
\vspace{0.15cm}
\item En "Juan es de Madrid", "de Madrid" es:
\begin{enumerate}[label=\textbf{\Alph*)}]
\begin{multicols}{2}
  \item Predicativo
  \item Sujeto
\end{multicols}
\end{enumerate}
\vspace{0.15cm}
\item El atributo puede sustituirse por:
\begin{enumerate}[label=\textbf{\Alph*)}]
\begin{multicols}{2}
  \item Lo
  \item Le
  \item La
  \item Los
  \item Las
  \item Les
\end{multicols}
\end{enumerate}
\vspace{0.15cm}
\item En "Juan es médico y Pedro también lo es", "lo" sustituye a:
\begin{enumerate}[label=\textbf{\Alph*)}]
\begin{multicols}{2}
  \item Predicativo
  \item Sujeto
\end{multicols}
\end{enumerate}
\vspace{0.15cm}
\item La oración copulativa es la que:
\begin{enumerate}[label=\textbf{\Alph*)}]
\begin{multicols}{2}
  \item Lleva verbo copulativo
  \item Lleva verbo predicativo
  \item No lleva verbo
  \item Lleva dos verbos
\end{multicols}
\end{enumerate}
\vspace{0.15cm}
\item La oración predicativa es la que:
\begin{enumerate}[label=\textbf{\Alph*)}]
\begin{multicols}{2}
  \item Lleva verbo predicativo
  \item Lleva verbo copulativo
  \item No lleva verbo
  \item Lleva dos verbos
\end{multicols}
\end{enumerate}
\vspace{0.15cm}
\item La oración transitiva es la que:
\begin{enumerate}[label=\textbf{\Alph*)}]
\begin{multicols}{2}
  \item Lleva complemento directo
  \item No lleva complemento directo
  \item Lleva complemento indirecto
  \item Lleva complemento circunstancial
  \item Lleva atributo
  \item Lleva predicativo
\end{multicols}
\end{enumerate}
\vspace{0.15cm}
\item La oración intransitiva es la que:
\begin{enumerate}[label=\textbf{\Alph*)}]
\begin{multicols}{2}
  \item No lleva complemento directo
  \item Lleva complemento directo
  \item Lleva complemento indirecto
  \item Lleva complemento circunstancial
  \item Lleva atributo
  \item Lleva predicativo
\end{multicols}
\end{enumerate}
\vspace{0.15cm}
\item La oración reflexiva es la que:
\begin{enumerate}[label=\textbf{\Alph*)}]
\begin{multicols}{2}
  \item Lleva pronombre reflexivo
  \item No lleva pronombre
  \item Lleva complemento directo
  \item Lleva complemento indirecto
  \item Lleva atributo
  \item Lleva predicativo
\end{multicols}
\end{enumerate}
\vspace{0.15cm}
\item La oración recíproca es la que:
\begin{enumerate}[label=\textbf{\Alph*)}]
\begin{multicols}{2}
\end{multicols}
\end{enumerate}
\vspace{0.15cm}
\item La oración impersonal es la que:
\begin{enumerate}[label=\textbf{\Alph*)}]
\begin{multicols}{2}
  \item No tiene sujeto
  \item Tiene sujeto
  \item Tiene dos sujetos
  \item Tiene sujeto elidido
  \item Tiene sujeto paciente
  \item Tiene sujeto agente
\end{multicols}
\end{enumerate}
\vspace{0.15cm}
\item La oración pasiva es la que:
\begin{enumerate}[label=\textbf{\Alph*)}]
\begin{multicols}{2}
  \item No tiene sujeto
  \item Tiene dos sujetos
  \item Tiene sujeto agente
  \item Tiene sujeto y agente
\end{multicols}
\end{enumerate}
\vspace{0.15cm}
\item La oración activa es la que:
\begin{enumerate}[label=\textbf{\Alph*)}]
\begin{multicols}{2}
  \item No tiene sujeto
  \item Tiene dos sujetos
  \item Tiene sujeto paciente
  \item Tiene sujeto y paciente
\end{multicols}
\end{enumerate}
\vspace{0.15cm}
\item En "La casa fue construida por Juan", el sujeto es:
\begin{enumerate}[label=\textbf{\Alph*)}]
\begin{multicols}{2}
  \item La casa
  \item Juan
  \item fue construida
  \item por Juan
  \item La casa fue construida
  \item No hay sujeto
\end{multicols}
\end{enumerate}
\vspace{0.15cm}
\item En "La casa fue construida por Juan", el complemento agente es:
\begin{enumerate}[label=\textbf{\Alph*)}]
\begin{multicols}{2}
  \item por Juan
  \item La casa
  \item fue construida
  \item La casa fue construida
  \item No hay complemento agente
  \item Juan
\end{multicols}
\end{enumerate}
\vspace{0.15cm}
\item En "Juan construyó la casa", el sujeto es:
\begin{enumerate}[label=\textbf{\Alph*)}]
\begin{multicols}{2}
  \item Juan
  \item la casa
  \item construyó
  \item Juan construyó
  \item No hay sujeto
  \item la casa construida
\end{multicols}
\end{enumerate}
\vspace{0.15cm}
\item En "Juan construyó la casa", el complemento directo es:
\begin{enumerate}[label=\textbf{\Alph*)}]
\begin{multicols}{2}
  \item la casa
  \item Juan
  \item construyó
  \item Juan construyó
  \item No hay complemento directo
  \item la casa construida
\end{multicols}
\end{enumerate}
\vspace{0.15cm}
\item En "Se venden casas", el sujeto es:
\begin{enumerate}[label=\textbf{\Alph*)}]
\begin{multicols}{2}
  \item casas
  \item Se
  \item venden
  \item No hay sujeto
  \item Se venden
  \item casas venden
\end{multicols}
\end{enumerate}
\vspace{0.15cm}
\item En "Se vive bien aquí", el sujeto es:
\begin{enumerate}[label=\textbf{\Alph*)}]
\begin{multicols}{2}
  \item Se
  \item vive
  \item aquí
  \item bien
  \item Se vive
\end{multicols}
\end{enumerate}
\vspace{0.15cm}
\item En "Hay muchos libros", el sujeto es:
\begin{enumerate}[label=\textbf{\Alph*)}]
\begin{multicols}{2}
  \item muchos libros
  \item Hay
  \item muchos
  \item libros
  \item Hay muchos
\end{multicols}
\end{enumerate}
\vspace{0.15cm}
\item En "Llueve mucho", el sujeto es:
\begin{enumerate}[label=\textbf{\Alph*)}]
\begin{multicols}{2}
  \item Llueve
  \item mucho
  \item Llueve mucho
  \item No hay predicado
  \item No hay oración
\end{multicols}
\end{enumerate}
\vspace{0.15cm}
\item La oración "Juan estudia" es:
\begin{enumerate}[label=\textbf{\Alph*)}]
\begin{multicols}{2}
  \item Simple, activa, transitiva
  \item Simple, activa, intransitiva
  \item Simple, pasiva, transitiva
  \item Simple, copulativa, transitiva
\end{multicols}
\end{enumerate}
\vspace{0.15cm}
\item La oración "Juan es médico" es:
\begin{enumerate}[label=\textbf{\Alph*)}]
\begin{multicols}{2}
  \item Simple, copulativa
  \item Simple, predicativa
  \item Simple, pasiva
\end{multicols}
\end{enumerate}
\vspace{0.15cm}
\item La oración "Juan compró un libro" es:
\begin{enumerate}[label=\textbf{\Alph*)}]
\begin{multicols}{2}
  \item Simple, predicativa, transitiva
  \item Simple, copulativa, transitiva
  \item Simple, predicativa, intransitiva
  \item Simple, pasiva, transitiva
\end{multicols}
\end{enumerate}
\vspace{0.15cm}
\item La oración "Juan llegó tarde" es:
\begin{enumerate}[label=\textbf{\Alph*)}]
\begin{multicols}{2}
  \item Simple, predicativa, intransitiva
  \item Simple, copulativa, intransitiva
  \item Simple, predicativa, transitiva
  \item Simple, pasiva, intransitiva ---
\end{multicols}
\end{enumerate}
\vspace{0.15cm}
\end{enumerate}\n
\subsubsection*{SECCIÓN V: COMPLEMENTOS DEL VERBO (201–250)}

\begin{enumerate}[label=\textbf{\arabic*.}]
\setcounter{enumi}{200}
\item El complemento directo es:
\begin{enumerate}[label=\textbf{\Alph*)}]
\begin{multicols}{2}
\end{multicols}
\end{enumerate}
\vspace{0.15cm}
\item El complemento directo puede sustituirse por:
\begin{enumerate}[label=\textbf{\Alph*)}]
\begin{multicols}{2}
  \item Lo, la, los, las
  \item Le, les
  \item Me, te, se
  \item Nos, os
  \item Él, ella, ellos
\end{multicols}
\end{enumerate}
\vspace{0.15cm}
\item El complemento indirecto puede sustituirse por:
\begin{enumerate}[label=\textbf{\Alph*)}]
\begin{multicols}{2}
  \item Le, les
  \item Lo, la, los, las
  \item Me, te, se
  \item Nos, os
  \item Él, ella, ellos
\end{multicols}
\end{enumerate}
\vspace{0.15cm}
\item El complemento directo puede ser:
\begin{enumerate}[label=\textbf{\Alph*)}]
\begin{multicols}{2}
  \item Un sintagma nominal, un pronombre o una oración
  \item Solo un sintagma nominal
  \item Solo un pronombre
  \item Solo una oración
  \item Solo un infinitivo
  \item Solo un gerundio
\end{multicols}
\end{enumerate}
\vspace{0.15cm}
\item El complemento indirecto puede ser:
\begin{enumerate}[label=\textbf{\Alph*)}]
\begin{multicols}{2}
  \item Un sintagma preposicional con "a" o "para"
  \item Solo un sintagma nominal
  \item Solo un pronombre
  \item Solo una oración
  \item Solo un infinitivo
  \item Solo un gerundio
\end{multicols}
\end{enumerate}
\vspace{0.15cm}
\item En "Compré el libro", "el libro" es:
\begin{enumerate}[label=\textbf{\Alph*)}]
\begin{multicols}{2}
  \item Sujeto
  \item Predicativo
\end{multicols}
\end{enumerate}
\vspace{0.15cm}
\item En "Di el libro a Juan", "el libro" es:
\begin{enumerate}[label=\textbf{\Alph*)}]
\begin{multicols}{2}
  \item Sujeto
  \item Predicativo
\end{multicols}
\end{enumerate}
\vspace{0.15cm}
\item En "Di el libro a Juan", "a Juan" es:
\begin{enumerate}[label=\textbf{\Alph*)}]
\begin{multicols}{2}
  \item Sujeto
  \item Predicativo
\end{multicols}
\end{enumerate}
\vspace{0.15cm}
\item En "Compré flores para María", "flores" es:
\begin{enumerate}[label=\textbf{\Alph*)}]
\begin{multicols}{2}
  \item Sujeto
  \item Predicativo
\end{multicols}
\end{enumerate}
\vspace{0.15cm}
\item En "Compré flores para María", "para María" es:
\begin{enumerate}[label=\textbf{\Alph*)}]
\begin{multicols}{2}
  \item Sujeto
  \item Predicativo
\end{multicols}
\end{enumerate}
\vspace{0.15cm}
\item El complemento circunstancial puede ser de:
\begin{enumerate}[label=\textbf{\Alph*)}]
\begin{multicols}{2}
  \item Lugar, tiempo, modo, causa, finalidad, etc.
  \item Solo lugar
  \item Solo tiempo
  \item Solo modo
  \item Solo causa
  \item Solo finalidad
\end{multicols}
\end{enumerate}
\vspace{0.15cm}
\item En "Estudio en la biblioteca", "en la biblioteca" es:
\begin{enumerate}[label=\textbf{\Alph*)}]
\begin{multicols}{2}
  \item Predicativo
  \item Sujeto
\end{multicols}
\end{enumerate}
\vspace{0.15cm}
\item En "Llegué ayer", "ayer" es:
\begin{enumerate}[label=\textbf{\Alph*)}]
\begin{multicols}{2}
  \item Predicativo
  \item Sujeto
\end{multicols}
\end{enumerate}
\vspace{0.15cm}
\item En "Lo hice con cuidado", "con cuidado" es:
\begin{enumerate}[label=\textbf{\Alph*)}]
\begin{multicols}{2}
  \item Predicativo
  \item Sujeto
\end{multicols}
\end{enumerate}
\vspace{0.15cm}
\item En "No vino por la lluvia", "por la lluvia" es:
\begin{enumerate}[label=\textbf{\Alph*)}]
\begin{multicols}{2}
  \item Predicativo
  \item Sujeto
\end{multicols}
\end{enumerate}
\vspace{0.15cm}
\item En "Estudio para aprobar", "para aprobar" es:
\begin{enumerate}[label=\textbf{\Alph*)}]
\begin{multicols}{2}
  \item Predicativo
  \item Sujeto
\end{multicols}
\end{enumerate}
\vspace{0.15cm}
\item En "Vino con Juan", "con Juan" es:
\begin{enumerate}[label=\textbf{\Alph*)}]
\begin{multicols}{2}
  \item Predicativo
  \item Sujeto
\end{multicols}
\end{enumerate}
\vspace{0.15cm}
\item En "Lo hizo con un martillo", "con un martillo" es:
\begin{enumerate}[label=\textbf{\Alph*)}]
\begin{multicols}{2}
  \item Predicativo
  \item Sujeto
\end{multicols}
\end{enumerate}
\vspace{0.15cm}
\item En "Cuesta cien pesos", "cien pesos" es:
\begin{enumerate}[label=\textbf{\Alph*)}]
\begin{multicols}{2}
  \item Predicativo
  \item Sujeto
\end{multicols}
\end{enumerate}
\vspace{0.15cm}
\item El complemento agente aparece en:
\begin{enumerate}[label=\textbf{\Alph*)}]
\begin{multicols}{2}
  \item Oraciones pasivas
  \item Oraciones activas
  \item Oraciones impersonales
  \item Oraciones copulativas
  \item Oraciones transitivas
  \item Oraciones intransitivas
\end{multicols}
\end{enumerate}
\vspace{0.15cm}
\item En "El libro fue escrito por Juan", "por Juan" es:
\begin{enumerate}[label=\textbf{\Alph*)}]
\begin{multicols}{2}
  \item Predicativo
\end{multicols}
\end{enumerate}
\vspace{0.15cm}
\item El complemento agente puede sustituirse por:
\begin{enumerate}[label=\textbf{\Alph*)}]
\begin{multicols}{2}
  \item No se sustituye por pronombre
  \item Lo
  \item Le
  \item La
  \item Los
  \item Las
\end{multicols}
\end{enumerate}
\vspace{0.15cm}
\item El complemento predicativo concuerda con:
\begin{enumerate}[label=\textbf{\Alph*)}]
\begin{multicols}{2}
  \item Solo el sujeto
  \item Solo el complemento directo
  \item Solo el complemento indirecto
  \item Solo el verbo
  \item Solo el adverbio
\end{multicols}
\end{enumerate}
\vspace{0.15cm}
\item En "Juan llegó cansado", "cansado" es:
\begin{enumerate}[label=\textbf{\Alph*)}]
\begin{multicols}{2}
  \item Predicativo
  \item Sujeto
\end{multicols}
\end{enumerate}
\vspace{0.15cm}
\item En "Considero a Juan inteligente", "inteligente" es:
\begin{enumerate}[label=\textbf{\Alph*)}]
\begin{multicols}{2}
  \item Predicativo del complemento directo
  \item Sujeto
\end{multicols}
\end{enumerate}
\vspace{0.15cm}
\item En "Nombraron a Juan director", "director" es:
\begin{enumerate}[label=\textbf{\Alph*)}]
\begin{multicols}{2}
  \item Predicativo del complemento directo
  \item Sujeto
\end{multicols}
\end{enumerate}
\vspace{0.15cm}
\item En "Juan es médico", "médico" es:
\begin{enumerate}[label=\textbf{\Alph*)}]
\begin{multicols}{2}
  \item Predicativo
  \item Sujeto
\end{multicols}
\end{enumerate}
\vspace{0.15cm}
\item En "Juan está cansado", "cansado" es:
\begin{enumerate}[label=\textbf{\Alph*)}]
\begin{multicols}{2}
  \item Predicativo
  \item Sujeto
\end{multicols}
\end{enumerate}
\vspace{0.15cm}
\item El atributo puede sustituirse por:
\begin{enumerate}[label=\textbf{\Alph*)}]
\begin{multicols}{2}
  \item Lo
  \item Le
  \item La
  \item Los
  \item Las
  \item Les
\end{multicols}
\end{enumerate}
\vspace{0.15cm}
\item En "Juan es médico y Pedro también lo es", "lo" sustituye a:
\begin{enumerate}[label=\textbf{\Alph*)}]
\begin{multicols}{2}
  \item Predicativo
  \item Sujeto
\end{multicols}
\end{enumerate}
\vspace{0.15cm}
\item El complemento del nombre puede ser:
\begin{enumerate}[label=\textbf{\Alph*)}]
\begin{multicols}{2}
  \item Un sintagma preposicional, un adjetivo o una aposición
  \item Solo un sintagma preposicional
  \item Solo un adjetivo
  \item Solo una aposición
  \item Solo un adverbio
  \item Solo un verbo
\end{multicols}
\end{enumerate}
\vspace{0.15cm}
\item En "el libro de Juan", "de Juan" es:
\begin{enumerate}[label=\textbf{\Alph*)}]
\begin{multicols}{2}
  \item Predicativo
\end{multicols}
\end{enumerate}
\vspace{0.15cm}
\item En "la casa blanca", "blanca" es:
\begin{enumerate}[label=\textbf{\Alph*)}]
\begin{multicols}{2}
  \item Núcleo del sintagma
\end{multicols}
\end{enumerate}
\vspace{0.15cm}
\item En "mi amigo Pedro", "Pedro" es:
\begin{enumerate}[label=\textbf{\Alph*)}]
\begin{multicols}{2}
  \item Predicativo
\end{multicols}
\end{enumerate}
\vspace{0.15cm}
\item El complemento del adjetivo puede ser:
\begin{enumerate}[label=\textbf{\Alph*)}]
\begin{multicols}{2}
  \item Un sintagma preposicional o una oración
  \item Solo un sintagma preposicional
  \item Solo una oración
  \item Solo un adverbio
  \item Solo un sustantivo
  \item Solo un verbo
\end{multicols}
\end{enumerate}
\vspace{0.15cm}
\item En "Cansado de esperar", "de esperar" es:
\begin{enumerate}[label=\textbf{\Alph*)}]
\begin{multicols}{2}
  \item Predicativo
\end{multicols}
\end{enumerate}
\vspace{0.15cm}
\item En "Muy lejos de aquí", "de aquí" es:
\begin{enumerate}[label=\textbf{\Alph*)}]
\begin{multicols}{2}
  \item Predicativo
\end{multicols}
\end{enumerate}
\vspace{0.15cm}
\item En "Juan es más alto que Pedro", "que Pedro" es:
\begin{enumerate}[label=\textbf{\Alph*)}]
\begin{multicols}{2}
  \item Predicativo
\end{multicols}
\end{enumerate}
\vspace{0.15cm}
\item El complemento del adverbio puede ser:
\begin{enumerate}[label=\textbf{\Alph*)}]
\begin{multicols}{2}
  \item Un sintagma preposicional
  \item Solo un sustantivo
  \item Solo un adjetivo
  \item Solo un verbo
  \item Solo una oración
  \item Solo un pronombre
\end{multicols}
\end{enumerate}
\vspace{0.15cm}
\item En "El hombre de la barba", "de la barba" es:
\begin{enumerate}[label=\textbf{\Alph*)}]
\begin{multicols}{2}
  \item Predicativo
\end{multicols}
\end{enumerate}
\vspace{0.15cm}
\item El vocativo es:
\begin{enumerate}[label=\textbf{\Alph*)}]
\begin{multicols}{2}
\end{multicols}
\end{enumerate}
\vspace{0.15cm}
\item En "Juan, ven aquí", "Juan" es:
\begin{enumerate}[label=\textbf{\Alph*)}]
\begin{multicols}{2}
  \item Vocativo
  \item Sujeto
\end{multicols}
\end{enumerate}
\vspace{0.15cm}
\item En "María, te llaman", "María" es:
\begin{enumerate}[label=\textbf{\Alph*)}]
\begin{multicols}{2}
  \item Vocativo
  \item Sujeto
\end{multicols}
\end{enumerate}
\vspace{0.15cm}
\item Las interjecciones son:
\begin{enumerate}[label=\textbf{\Alph*)}]
\begin{multicols}{2}
  \item Palabras que expresan emociones o reacciones
  \item Palabras que modifican al sustantivo
  \item Palabras que modifican al verbo
  \item Palabras que modifican al adjetivo
  \item Palabras que unen oraciones
  \item Palabras que indican circunstancias
\end{multicols}
\end{enumerate}
\vspace{0.15cm}
\item En "¡Ay!, me duele", "¡Ay!" es:
\begin{enumerate}[label=\textbf{\Alph*)}]
\begin{multicols}{2}
  \item Interjección
  \item Vocativo
  \item Sujeto
\end{multicols}
\end{enumerate}
\vspace{0.15cm}
\item Los marcadores del discurso son:
\begin{enumerate}[label=\textbf{\Alph*)}]
\begin{multicols}{2}
\end{multicols}
\end{enumerate}
\vspace{0.15cm}
\item En "Sin embargo, no fui", "sin embargo" es:
\begin{enumerate}[label=\textbf{\Alph*)}]
\begin{multicols}{2}
  \item Marcador del discurso
  \item Vocativo
  \item Sujeto
\end{multicols}
\end{enumerate}
\vspace{0.15cm}
\item En "En conclusión, todo salió bien", "en conclusión" es:
\begin{enumerate}[label=\textbf{\Alph*)}]
\begin{multicols}{2}
  \item Marcador del discurso
  \item Vocativo
  \item Sujeto
\end{multicols}
\end{enumerate}
\vspace{0.15cm}
\item El complemento directo puede llevar la preposición "a" cuando:
\begin{enumerate}[label=\textbf{\Alph*)}]
\begin{multicols}{2}
\end{multicols}
\end{enumerate}
\vspace{0.15cm}
\item En "Vi a Juan", "a Juan" es:
\begin{enumerate}[label=\textbf{\Alph*)}]
\begin{multicols}{2}
  \item Sujeto
  \item Predicativo ---
\end{multicols}
\end{enumerate}
\vspace{0.15cm}
\end{enumerate}\n
\subsubsection*{SECCIÓN VI: TIPOS DE ORACIONES Y COORDINACIÓN/SUBORDINACIÓN (251–300)}

\begin{enumerate}[label=\textbf{\arabic*.}]
\setcounter{enumi}{250}
\item La oración simple se caracteriza por:
\begin{enumerate}[label=\textbf{\Alph*)}]
\begin{multicols}{2}
  \item Tener un solo verbo conjugado
  \item Tener dos o más verbos conjugados
  \item No tener verbo
  \item Tener solo sustantivos
  \item Tener solo adjetivos
  \item Tener solo adverbios
\end{multicols}
\end{enumerate}
\vspace{0.15cm}
\item La oración compuesta se caracteriza por:
\begin{enumerate}[label=\textbf{\Alph*)}]
\begin{multicols}{2}
  \item Tener dos o más verbos conjugados
  \item Tener un solo verbo conjugado
  \item No tener verbo
  \item Tener solo sustantivos
  \item Tener solo adjetivos
  \item Tener solo adverbios
\end{multicols}
\end{enumerate}
\vspace{0.15cm}
\item Las oraciones compuestas pueden ser:
\begin{enumerate}[label=\textbf{\Alph*)}]
\begin{multicols}{2}
  \item Solo coordinadas
  \item Solo subordinadas
  \item Solo yuxtapuestas
  \item Solo copulativas
  \item Solo adversativas
\end{multicols}
\end{enumerate}
\vspace{0.15cm}
\item Las oraciones coordinadas se caracterizan por:
\begin{enumerate}[label=\textbf{\Alph*)}]
\begin{multicols}{2}
  \item Ser independientes sintácticamente
  \item No tener verbo
  \item Tener un solo verbo
  \item Ser impersonales
  \item Ser pasivas
\end{multicols}
\end{enumerate}
\vspace{0.15cm}
\item Las oraciones subordinadas se caracterizan por:
\begin{enumerate}[label=\textbf{\Alph*)}]
\begin{multicols}{2}
  \item Ser independientes
  \item No tener verbo
  \item Tener un solo verbo
  \item Ser impersonales
  \item Ser pasivas
\end{multicols}
\end{enumerate}
\vspace{0.15cm}
\item Las oraciones yuxtapuestas se caracterizan por:
\begin{enumerate}[label=\textbf{\Alph*)}]
\begin{multicols}{2}
  \item Unirse sin conjunciones
  \item Unirse con conjunciones
  \item Ser independientes
  \item Ser dependientes
  \item No tener verbo
  \item Tener un solo verbo
\end{multicols}
\end{enumerate}
\vspace{0.15cm}
\item Las oraciones coordinadas copulativas se unen con:
\begin{enumerate}[label=\textbf{\Alph*)}]
\begin{multicols}{2}
  \item Y, e, ni
  \item O, u
  \item Pero, mas, sino
  \item Pues, porque
  \item Luego, así que
\end{multicols}
\end{enumerate}
\vspace{0.15cm}
\item Las oraciones coordinadas disyuntivas se unen con:
\begin{enumerate}[label=\textbf{\Alph*)}]
\begin{multicols}{2}
  \item O, u
  \item Y, e, ni
  \item Pero, mas, sino
  \item Pues, porque
  \item Luego, así que
\end{multicols}
\end{enumerate}
\vspace{0.15cm}
\item Las oraciones coordinadas adversativas se unen con:
\begin{enumerate}[label=\textbf{\Alph*)}]
\begin{multicols}{2}
  \item Pero, mas, sino
  \item Y, e, ni
  \item O, u
  \item Pues, porque
  \item Luego, así que
\end{multicols}
\end{enumerate}
\vspace{0.15cm}
\item Las oraciones coordinadas explicativas se unen con:
\begin{enumerate}[label=\textbf{\Alph*)}]
\begin{multicols}{2}
  \item Y, e, ni
  \item O, u
  \item Pero, mas, sino
  \item Pues, porque
  \item Luego, así que
\end{multicols}
\end{enumerate}
\vspace{0.15cm}
\item En "Juan estudia y María trabaja", la oración es:
\begin{enumerate}[label=\textbf{\Alph*)}]
\begin{multicols}{2}
  \item Subordinada
  \item Yuxtapuesta
  \item Simple
\end{multicols}
\end{enumerate}
\vspace{0.15cm}
\item En "¿Vienes o te quedas?", la oración es:
\begin{enumerate}[label=\textbf{\Alph*)}]
\begin{multicols}{2}
  \item Subordinada
  \item Yuxtapuesta
  \item Simple
\end{multicols}
\end{enumerate}
\vspace{0.15cm}
\item En "Estudié, pero no aprobé", la oración es:
\begin{enumerate}[label=\textbf{\Alph*)}]
\begin{multicols}{2}
  \item Subordinada
  \item Yuxtapuesta
  \item Simple
\end{multicols}
\end{enumerate}
\vspace{0.15cm}
\item En "Llegó tarde, es decir, no pudo entrar", la oración es:
\begin{enumerate}[label=\textbf{\Alph*)}]
\begin{multicols}{2}
  \item Subordinada
  \item Yuxtapuesta
  \item Simple
\end{multicols}
\end{enumerate}
\vspace{0.15cm}
\item En "Juan estudia, María trabaja", la oración es:
\begin{enumerate}[label=\textbf{\Alph*)}]
\begin{multicols}{2}
  \item Yuxtapuesta
  \item Subordinada
  \item Simple
\end{multicols}
\end{enumerate}
\vspace{0.15cm}
\item Las oraciones subordinadas sustantivas funcionan como:
\begin{enumerate}[label=\textbf{\Alph*)}]
\begin{multicols}{2}
  \item Sujeto, complemento directo, etc.
  \item Verbo
  \item Preposición
\end{multicols}
\end{enumerate}
\vspace{0.15cm}
\item Las oraciones subordinadas adjetivas funcionan como:
\begin{enumerate}[label=\textbf{\Alph*)}]
\begin{multicols}{2}
  \item Sustantivo
  \item Verbo
  \item Preposición
\end{multicols}
\end{enumerate}
\vspace{0.15cm}
\item Las oraciones subordinadas adverbiales funcionan como:
\begin{enumerate}[label=\textbf{\Alph*)}]
\begin{multicols}{2}
  \item Sustantivo
  \item Verbo
  \item Preposición
\end{multicols}
\end{enumerate}
\vspace{0.15cm}
\item En "Quiero que vengas", la oración subordinada es:
\begin{enumerate}[label=\textbf{\Alph*)}]
\begin{multicols}{2}
  \item Sustantiva
  \item Yuxtapuesta
  \item Simple
\end{multicols}
\end{enumerate}
\vspace{0.15cm}
\item En "El libro que compré es interesante", la oración subordinada es:
\begin{enumerate}[label=\textbf{\Alph*)}]
\begin{multicols}{2}
  \item Sustantiva
  \item Yuxtapuesta
  \item Simple
\end{multicols}
\end{enumerate}
\vspace{0.15cm}
\item En "Estudio cuando puedo", la oración subordinada es:
\begin{enumerate}[label=\textbf{\Alph*)}]
\begin{multicols}{2}
  \item Sustantiva
  \item Yuxtapuesta
  \item Simple
\end{multicols}
\end{enumerate}
\vspace{0.15cm}
\item Las oraciones subordinadas sustantivas pueden ir introducidas por:
\begin{enumerate}[label=\textbf{\Alph*)}]
\begin{multicols}{2}
  \item Que, si, pronombres interrogativos
  \item Que, cual, quien
  \item Y, o, pero
  \item Pues, porque
\end{multicols}
\end{enumerate}
\vspace{0.15cm}
\item Las oraciones subordinadas adjetivas pueden ir introducidas por:
\begin{enumerate}[label=\textbf{\Alph*)}]
\begin{multicols}{2}
  \item Que, cual, quien, cuyo
  \item Que, si, pronombres interrogativos
  \item Y, o, pero
  \item Pues, porque
\end{multicols}
\end{enumerate}
\vspace{0.15cm}
\item Las oraciones subordinadas adverbiales pueden ir introducidas por:
\begin{enumerate}[label=\textbf{\Alph*)}]
\begin{multicols}{2}
  \item Que, si, pronombres interrogativos
  \item Que, cual, quien, cuyo
  \item Y, o, pero
  \item Pues, porque
\end{multicols}
\end{enumerate}
\vspace{0.15cm}
\item En "Dijo que vendría", la oración subordinada es:
\begin{enumerate}[label=\textbf{\Alph*)}]
\begin{multicols}{2}
  \item Sustantiva de complemento directo
  \item Yuxtapuesta
  \item Simple
\end{multicols}
\end{enumerate}
\vspace{0.15cm}
\item En "El hombre que vino es mi tío", la oración subordinada es:
\begin{enumerate}[label=\textbf{\Alph*)}]
\begin{multicols}{2}
  \item Sustantiva
  \item Yuxtapuesta
  \item Simple
\end{multicols}
\end{enumerate}
\vspace{0.15cm}
\item En "Cuando llegues, avísame", la oración subordinada es:
\begin{enumerate}[label=\textbf{\Alph*)}]
\begin{multicols}{2}
  \item Sustantiva
  \item Yuxtapuesta
  \item Simple
\end{multicols}
\end{enumerate}
\vspace{0.15cm}
\item En "Estudio porque quiero aprobar", la oración subordinada es:
\begin{enumerate}[label=\textbf{\Alph*)}]
\begin{multicols}{2}
  \item Sustantiva
  \item Yuxtapuesta
  \item Simple
\end{multicols}
\end{enumerate}
\vspace{0.15cm}
\item En "Estudio para que me aprueben", la oración subordinada es:
\begin{enumerate}[label=\textbf{\Alph*)}]
\begin{multicols}{2}
  \item Sustantiva
  \item Yuxtapuesta
  \item Simple
\end{multicols}
\end{enumerate}
\vspace{0.15cm}
\item En "Si estudias, aprobarás", la oración subordinada es:
\begin{enumerate}[label=\textbf{\Alph*)}]
\begin{multicols}{2}
  \item Sustantiva
  \item Yuxtapuesta
  \item Simple
\end{multicols}
\end{enumerate}
\vspace{0.15cm}
\item En "Aunque estudies, no aprobarás", la oración subordinada es:
\begin{enumerate}[label=\textbf{\Alph*)}]
\begin{multicols}{2}
  \item Sustantiva
  \item Yuxtapuesta
  \item Simple
\end{multicols}
\end{enumerate}
\vspace{0.15cm}
\item En "Estudia tanto que se enfermó", la oración subordinada es:
\begin{enumerate}[label=\textbf{\Alph*)}]
\begin{multicols}{2}
  \item Sustantiva
  \item Yuxtapuesta
  \item Simple
\end{multicols}
\end{enumerate}
\vspace{0.15cm}
\item En "Hice lo que me pediste", la oración subordinada es:
\begin{enumerate}[label=\textbf{\Alph*)}]
\begin{multicols}{2}
  \item Sustantiva
  \item Yuxtapuesta
  \item Simple
\end{multicols}
\end{enumerate}
\vspace{0.15cm}
\item En "Quien estudia, aprueba", la oración subordinada es:
\begin{enumerate}[label=\textbf{\Alph*)}]
\begin{multicols}{2}
  \item Sustantiva
  \item Yuxtapuesta
  \item Simple
\end{multicols}
\end{enumerate}
\vspace{0.15cm}
\item En "No sé si vendrá", la oración subordinada es:
\begin{enumerate}[label=\textbf{\Alph*)}]
\begin{multicols}{2}
  \item Sustantiva
  \item Yuxtapuesta
  \item Simple
\end{multicols}
\end{enumerate}
\vspace{0.15cm}
\item En "Me preguntó dónde vivía", la oración subordinada es:
\begin{enumerate}[label=\textbf{\Alph*)}]
\begin{multicols}{2}
  \item Sustantiva
  \item Yuxtapuesta
  \item Simple
\end{multicols}
\end{enumerate}
\vspace{0.15cm}
\item En "El hecho de que no viniera me molestó", la oración subordinada es:
\begin{enumerate}[label=\textbf{\Alph*)}]
\begin{multicols}{2}
  \item Sustantiva
  \item Yuxtapuesta
  \item Simple
\end{multicols}
\end{enumerate}
\vspace{0.15cm}
\item En "Vive donde termina la calle", la oración subordinada es:
\begin{enumerate}[label=\textbf{\Alph*)}]
\begin{multicols}{2}
  \item Sustantiva
  \item Yuxtapuesta
  \item Simple
\end{multicols}
\end{enumerate}
\vspace{0.15cm}
\item En "Lo hizo como le dijeron", la oración subordinada es:
\begin{enumerate}[label=\textbf{\Alph*)}]
\begin{multicols}{2}
  \item Sustantiva
  \item Yuxtapuesta
  \item Simple
\end{multicols}
\end{enumerate}
\vspace{0.15cm}
\item En "Es más alto que su padre", la oración subordinada es:
\begin{enumerate}[label=\textbf{\Alph*)}]
\begin{multicols}{2}
  \item Sustantiva
  \item Yuxtapuesta
  \item Simple
\end{multicols}
\end{enumerate}
\vspace{0.15cm}
\item La oración compuesta por subordinación tiene:
\begin{enumerate}[label=\textbf{\Alph*)}]
\begin{multicols}{2}
  \item Una oración principal y una o más subordinadas
  \item Solo oraciones coordinadas
  \item Solo oraciones yuxtapuestas
  \item Solo una oración
  \item Solo dos oraciones independientes
  \item Solo tres oraciones independientes
\end{multicols}
\end{enumerate}
\vspace{0.15cm}
\item La oración compuesta por coordinación tiene:
\begin{enumerate}[label=\textbf{\Alph*)}]
\begin{multicols}{2}
  \item Una oración principal y una subordinada
  \item Solo una oración
  \item Solo dos oraciones dependientes
  \item Solo tres oraciones dependientes
  \item Solo oraciones yuxtapuestas
\end{multicols}
\end{enumerate}
\vspace{0.15cm}
\item La oración compuesta por yuxtaposición tiene:
\begin{enumerate}[label=\textbf{\Alph*)}]
\begin{multicols}{2}
  \item Una oración principal y una subordinada
  \item Solo una oración
  \item Solo dos oraciones dependientes
  \item Solo tres oraciones dependientes
  \item Solo oraciones coordinadas
\end{multicols}
\end{enumerate}
\vspace{0.15cm}
\item En "Llegué, vi, vencí", las oraciones son:
\begin{enumerate}[label=\textbf{\Alph*)}]
\begin{multicols}{2}
  \item Yuxtapuestas
  \item Subordinadas
  \item Simple
\end{multicols}
\end{enumerate}
\vspace{0.15cm}
\item En "Estudia y trabaja", las oraciones son:
\begin{enumerate}[label=\textbf{\Alph*)}]
\begin{multicols}{2}
  \item Yuxtapuestas
  \item Subordinadas
  \item Simple
\end{multicols}
\end{enumerate}
\vspace{0.15cm}
\item En "¿Estudias o trabajas?", las oraciones son:
\begin{enumerate}[label=\textbf{\Alph*)}]
\begin{multicols}{2}
  \item Yuxtapuestas
  \item Subordinadas
  \item Simple
\end{multicols}
\end{enumerate}
\vspace{0.15cm}
\item En "Estudia, pero no aprende", las oraciones son:
\begin{enumerate}[label=\textbf{\Alph*)}]
\begin{multicols}{2}
  \item Yuxtapuestas
  \item Subordinadas
  \item Simple
\end{multicols}
\end{enumerate}
\vspace{0.15cm}
\item En "Quiero que estudies", la oración es:
\begin{enumerate}[label=\textbf{\Alph*)}]
\begin{multicols}{2}
  \item Simple
  \item Impersonal
  \item Pasiva
\end{multicols}
\end{enumerate}
\vspace{0.15cm}
\item En "Estudia y trabaja", la oración es:
\begin{enumerate}[label=\textbf{\Alph*)}]
\begin{multicols}{2}
  \item Simple
  \item Impersonal
  \item Pasiva
\end{multicols}
\end{enumerate}
\vspace{0.15cm}
\item En "Llegué, vi, vencí", la oración es:
\begin{enumerate}[label=\textbf{\Alph*)}]
\begin{multicols}{2}
  \item Simple
  \item Impersonal
  \item Pasiva ---
\end{multicols}
\end{enumerate}
\vspace{0.15cm}
\end{enumerate}\n
\vspace{0.4cm}
\begin{tcolorbox}[colback=gray!5!white,colframe=gray!60!black,title=\textbf{SOLUCIONARIO -- 2.3 SINTAXIS}]
\begin{multicols}{4}
\begin{enumerate}[label=\textbf{\arabic*.}]
  \item \textbf{A}
  \item \textbf{A}
  \item \textbf{A}
  \item \textbf{A}
  \item \textbf{A}
  \item \textbf{A}
  \item \textbf{A}
  \item \textbf{A}
  \item \textbf{A}
  \item \textbf{A}
  \item \textbf{A}
  \item \textbf{A}
  \item \textbf{A}
  \item \textbf{A}
  \item \textbf{A}
  \item \textbf{A}
  \item \textbf{A}
  \item \textbf{A}
  \item \textbf{A}
  \item \textbf{A}
  \item \textbf{A}
  \item \textbf{A}
  \item \textbf{A}
  \item \textbf{A}
  \item \textbf{A}
  \item \textbf{A}
  \item \textbf{A}
  \item \textbf{A}
  \item \textbf{A}
  \item \textbf{A}
  \item \textbf{A}
  \item \textbf{A}
  \item \textbf{A}
  \item \textbf{A}
  \item \textbf{A}
  \item \textbf{A}
  \item \textbf{A}
  \item \textbf{A}
  \item \textbf{A}
  \item \textbf{A}
  \item \textbf{A}
  \item \textbf{A}
  \item \textbf{A}
  \item \textbf{A}
  \item \textbf{A}
  \item \textbf{A}
  \item \textbf{A}
  \item \textbf{A}
  \item \textbf{A}
  \item \textbf{A}
  \item \textbf{A}
  \item \textbf{A}
  \item \textbf{A}
  \item \textbf{A}
  \item \textbf{A}
  \item \textbf{A}
  \item \textbf{A}
  \item \textbf{A}
  \item \textbf{A}
  \item \textbf{A}
  \item \textbf{A}
  \item \textbf{A}
  \item \textbf{A}
  \item \textbf{A}
  \item \textbf{A}
  \item \textbf{A}
  \item \textbf{A}
  \item \textbf{A}
  \item \textbf{A}
  \item \textbf{A}
  \item \textbf{A}
  \item \textbf{A}
  \item \textbf{A}
  \item \textbf{A}
  \item \textbf{A}
  \item \textbf{A}
  \item \textbf{A}
  \item \textbf{A}
  \item \textbf{A}
  \item \textbf{A}
  \item \textbf{A}
  \item \textbf{A}
  \item \textbf{A}
  \item \textbf{A}
  \item \textbf{A}
  \item \textbf{A}
  \item \textbf{A}
  \item \textbf{A}
  \item \textbf{A}
  \item \textbf{A}
  \item \textbf{A}
  \item \textbf{A}
  \item \textbf{A}
  \item \textbf{A}
  \item \textbf{A}
  \item \textbf{A}
  \item \textbf{A}
  \item \textbf{A}
  \item \textbf{A}
  \item \textbf{A}
  \item \textbf{A}
  \item \textbf{A}
  \item \textbf{A}
  \item \textbf{A}
  \item \textbf{A}
  \item \textbf{A}
  \item \textbf{A}
  \item \textbf{A}
  \item \textbf{A}
  \item \textbf{A}
  \item \textbf{A}
  \item \textbf{A}
  \item \textbf{A}
  \item \textbf{A}
  \item \textbf{A}
  \item \textbf{A}
  \item \textbf{A}
  \item \textbf{A}
  \item \textbf{A}
  \item \textbf{A}
  \item \textbf{A}
  \item \textbf{A}
  \item \textbf{A}
  \item \textbf{A}
  \item \textbf{A}
  \item \textbf{A}
  \item \textbf{A}
  \item \textbf{A}
  \item \textbf{A}
  \item \textbf{A}
  \item \textbf{A}
  \item \textbf{A}
  \item \textbf{A}
  \item \textbf{A}
  \item \textbf{A}
  \item \textbf{A}
  \item \textbf{A}
  \item \textbf{A}
  \item \textbf{A}
  \item \textbf{A}
  \item \textbf{A}
  \item \textbf{A}
  \item \textbf{A}
  \item \textbf{A}
  \item \textbf{A}
  \item \textbf{A}
  \item \textbf{A}
  \item \textbf{A}
  \item \textbf{A}
  \item \textbf{A}
  \item \textbf{A}
  \item \textbf{A}
  \item \textbf{A}
  \item \textbf{A}
  \item \textbf{A}
  \item \textbf{A}
  \item \textbf{A}
  \item \textbf{A}
  \item \textbf{A}
  \item \textbf{A}
  \item \textbf{A}
  \item \textbf{A}
  \item \textbf{A}
  \item \textbf{A}
  \item \textbf{A}
  \item \textbf{A}
  \item \textbf{A}
  \item \textbf{A}
  \item \textbf{A}
  \item \textbf{A}
  \item \textbf{A}
  \item \textbf{A}
  \item \textbf{A}
  \item \textbf{A}
  \item \textbf{A}
  \item \textbf{A}
  \item \textbf{A}
  \item \textbf{A}
  \item \textbf{A}
  \item \textbf{A}
  \item \textbf{A}
  \item \textbf{A}
  \item \textbf{A}
  \item \textbf{A}
  \item \textbf{A}
  \item \textbf{A}
  \item \textbf{A}
  \item \textbf{A}
  \item \textbf{A}
  \item \textbf{A}
  \item \textbf{A}
  \item \textbf{A}
  \item \textbf{A}
  \item \textbf{A}
  \item \textbf{A}
  \item \textbf{A}
  \item \textbf{A}
  \item \textbf{A}
  \item \textbf{A}
  \item \textbf{A}
  \item \textbf{A}
  \item \textbf{A}
  \item \textbf{A}
  \item \textbf{A}
  \item \textbf{A}
  \item \textbf{A}
  \item \textbf{A}
  \item \textbf{A}
  \item \textbf{A}
  \item \textbf{A}
  \item \textbf{A}
  \item \textbf{A}
  \item \textbf{A}
  \item \textbf{A}
  \item \textbf{A}
  \item \textbf{A}
  \item \textbf{A}
  \item \textbf{A}
  \item \textbf{A}
  \item \textbf{A}
  \item \textbf{A}
  \item \textbf{A}
  \item \textbf{A}
  \item \textbf{A}
  \item \textbf{A}
  \item \textbf{A}
  \item \textbf{A}
  \item \textbf{A}
  \item \textbf{A}
  \item \textbf{A}
  \item \textbf{A}
  \item \textbf{A}
  \item \textbf{A}
  \item \textbf{A}
  \item \textbf{A}
  \item \textbf{A}
  \item \textbf{A}
  \item \textbf{A}
  \item \textbf{A}
  \item \textbf{A}
  \item \textbf{A}
  \item \textbf{A}
  \item \textbf{A}
  \item \textbf{A}
  \item \textbf{A}
  \item \textbf{A}
  \item \textbf{A}
  \item \textbf{A}
  \item \textbf{A}
  \item \textbf{A}
  \item \textbf{A}
  \item \textbf{A}
  \item \textbf{A}
  \item \textbf{A}
  \item \textbf{A}
  \item \textbf{A}
  \item \textbf{A}
  \item \textbf{A}
  \item \textbf{A}
  \item \textbf{A}
  \item \textbf{A}
  \item \textbf{A}
  \item \textbf{A}
  \item \textbf{A}
  \item \textbf{A}
  \item \textbf{A}
  \item \textbf{A}
  \item \textbf{A}
  \item \textbf{A}
  \item \textbf{A}
  \item \textbf{A}
  \item \textbf{A}
  \item \textbf{A}
  \item \textbf{A}
  \item \textbf{A}
  \item \textbf{A}
  \item \textbf{A}
  \item \textbf{A}
  \item \textbf{A}
  \item \textbf{A}
  \item \textbf{A}
  \item \textbf{A}
  \item \textbf{A}
  \item \textbf{A}
  \item \textbf{A}
  \item \textbf{A}
  \item \textbf{A}
  \item \textbf{A}
  \item \textbf{A}
  \item \textbf{A}
  \item \textbf{A}
  \item \textbf{A}
  \item \textbf{A}
  \item \textbf{A}
  \item \textbf{A}
  \item \textbf{A}
  \item \textbf{A}
  \item \textbf{A}
  \item \textbf{A}
  \item \textbf{A}
\end{enumerate}
\end{multicols}
\end{tcolorbox}
\newpage\n
\subsection{2.4 METODOLOGÍA Y ESCRITURA ACADÉMICA}

\subsubsection*{SECCIÓN I: FUNDAMENTOS DE INVESTIGACIÓN CIENTÍFICA (1–40)}

\begin{enumerate}[label=\textbf{\arabic*.}]
\setcounter{enumi}{0}
\item ¿Qué es la investigación científica?
\begin{enumerate}[label=\textbf{\Alph*)}]
\begin{multicols}{2}
  \item Un proceso sistemático, organizado y objetivo para obtener conocimiento
  \item Una opinión personal sin fundamento
  \item Una creencia popular
  \item Una actividad aleatoria
  \item Una suposición sin verificar
  \item Una tradición oral
\end{multicols}
\end{enumerate}
\vspace{0.15cm}
\item ¿Cuál es el objetivo principal de la investigación científica?
\begin{enumerate}[label=\textbf{\Alph*)}]
\begin{multicols}{2}
  \item Generar conocimiento verificable y confiable
  \item Repetir información existente
  \item Ganar dinero
\end{multicols}
\end{enumerate}
\vspace{0.15cm}
\item ¿Qué caracteriza al método científico?
\begin{enumerate}[label=\textbf{\Alph*)}]
\begin{multicols}{2}
  \item Ser sistemático, objetivo, verificable y reproducible
  \item Ser subjetivo y arbitrario
  \item Ser aleatorio
  \item Ser improvisado
\end{multicols}
\end{enumerate}
\vspace{0.15cm}
\item ¿Cuáles son los pasos del método científico?
\begin{enumerate}[label=\textbf{\Alph*)}]
\begin{multicols}{2}
  \item Observación, hipótesis, experimentación, análisis y conclusión
  \item Opinión, creencia, tradición
  \item Improvisación, suposición, conclusión
  \item Lectura, escritura, publicación
  \item Pregunta, respuesta, fin
  \item Introducción, nudo, desenlace
\end{multicols}
\end{enumerate}
\vspace{0.15cm}
\item ¿Qué es una hipótesis?
\begin{enumerate}[label=\textbf{\Alph*)}]
\begin{multicols}{2}
  \item Una explicación tentativa que se puede verificar o refutar
  \item Una verdad absoluta
  \item Una opinión personal
  \item Una creencia popular
  \item Una conclusión final
  \item Un dato estadístico
\end{multicols}
\end{enumerate}
\vspace{0.15cm}
\item ¿Qué es una variable en investigación?
\begin{enumerate}[label=\textbf{\Alph*)}]
\begin{multicols}{2}
  \item Una característica o propiedad que puede variar y ser medida
  \item Un dato fijo
  \item Una conclusión
  \item Una hipótesis
  \item Un objetivo
  \item Una pregunta
\end{multicols}
\end{enumerate}
\vspace{0.15cm}
\item ¿Qué es una variable independiente?
\begin{enumerate}[label=\textbf{\Alph*)}]
\begin{multicols}{2}
  \item La que el investigador manipula
  \item La que se mide como resultado
  \item La que no cambia
  \item La que confunde
  \item La que se ignora
  \item La que se repite
\end{multicols}
\end{enumerate}
\vspace{0.15cm}
\item ¿Qué es una variable dependiente?
\begin{enumerate}[label=\textbf{\Alph*)}]
\begin{multicols}{2}
  \item La que se mide como resultado de la manipulación
  \item La que el investigador manipula
  \item La que no cambia
  \item La que confunde
  \item La que se ignora
  \item La que se repite
\end{multicols}
\end{enumerate}
\vspace{0.15cm}
\item ¿Qué es una variable control?
\begin{enumerate}[label=\textbf{\Alph*)}]
\begin{multicols}{2}
  \item La que se mantiene constante para evitar interferencias
  \item La que se manipula
  \item La que se mide
  \item La que cambia
  \item La que se ignora
  \item La que se repite
\end{multicols}
\end{enumerate}
\vspace{0.15cm}
\item ¿Qué es la población en investigación?
\begin{enumerate}[label=\textbf{\Alph*)}]
\begin{multicols}{2}
  \item Una parte del conjunto
  \item Un individuo
  \item Un dato
  \item Una variable
  \item Una hipótesis
\end{multicols}
\end{enumerate}
\vspace{0.15cm}
\item ¿Qué es la muestra en investigación?
\begin{enumerate}[label=\textbf{\Alph*)}]
\begin{multicols}{2}
  \item Un subconjunto representativo de la población
  \item Un individuo
  \item Un dato
  \item Una variable
  \item Una hipótesis
\end{multicols}
\end{enumerate}
\vspace{0.15cm}
\item ¿Qué es el muestreo?
\begin{enumerate}[label=\textbf{\Alph*)}]
\begin{multicols}{2}
  \item La técnica para seleccionar una muestra representativa
  \item La técnica para analizar datos
  \item La técnica para formular hipótesis
  \item La técnica para escribir
  \item La técnica para leer
  \item La técnica para exponer
\end{multicols}
\end{enumerate}
\vspace{0.15cm}
\item ¿Qué es un dato?
\begin{enumerate}[label=\textbf{\Alph*)}]
\begin{multicols}{2}
  \item La información recogida en la investigación
  \item Una opinión
  \item Una creencia
  \item Una suposición
  \item Una hipótesis
  \item Una conclusión
\end{multicols}
\end{enumerate}
\vspace{0.15cm}
\item ¿Qué es una fuente primaria?
\begin{enumerate}[label=\textbf{\Alph*)}]
\begin{multicols}{2}
  \item La que proporciona información original y directa
  \item La que interpreta información de otros
  \item La que resume información
  \item La que copia información
  \item La que traduce información
  \item La que critica información
\end{multicols}
\end{enumerate}
\vspace{0.15cm}
\item ¿Qué es una fuente secundaria?
\begin{enumerate}[label=\textbf{\Alph*)}]
\begin{multicols}{2}
  \item La que interpreta, analiza o resume información de fuentes primarias
  \item La que proporciona información original
  \item La que copia información
  \item La que traduce información
  \item La que critica información
  \item La que ignora información
\end{multicols}
\end{enumerate}
\vspace{0.15cm}
\item ¿Qué es una fuente terciaria?
\begin{enumerate}[label=\textbf{\Alph*)}]
\begin{multicols}{2}
  \item La que compila y organiza información de fuentes primarias y secundarias
  \item La que proporciona información original
  \item La que interpreta información
  \item La que copia información
  \item La que traduce información
  \item La que ignora información
\end{multicols}
\end{enumerate}
\vspace{0.15cm}
\item ¿Qué es la confiabilidad de una fuente?
\begin{enumerate}[label=\textbf{\Alph*)}]
\begin{multicols}{2}
  \item La cantidad de información
  \item La extensión del texto
\end{multicols}
\end{enumerate}
\vspace{0.15cm}
\item ¿Qué es la validez en investigación?
\begin{enumerate}[label=\textbf{\Alph*)}]
\begin{multicols}{2}
  \item La cantidad de datos
  \item La extensión del informe
\end{multicols}
\end{enumerate}
\vspace{0.15cm}
\item ¿Qué es la confiabilidad en investigación?
\begin{enumerate}[label=\textbf{\Alph*)}]
\begin{multicols}{2}
  \item La cantidad de datos
  \item La extensión del informe
\end{multicols}
\end{enumerate}
\vspace{0.15cm}
\item ¿Qué es el marco teórico?
\begin{enumerate}[label=\textbf{\Alph*)}]
\begin{multicols}{2}
  \item La conclusión de la investigación
  \item La introducción de la investigación
  \item La metodología de la investigación
  \item Los resultados de la investigación
  \item Las recomendaciones de la investigación
\end{multicols}
\end{enumerate}
\vspace{0.15cm}
\item ¿Qué es el marco metodológico?
\begin{enumerate}[label=\textbf{\Alph*)}]
\begin{multicols}{2}
  \item La descripción de los métodos y técnicas utilizados
  \item La conclusión de la investigación
  \item La introducción de la investigación
  \item Los resultados de la investigación
  \item Las recomendaciones de la investigación
\end{multicols}
\end{enumerate}
\vspace{0.15cm}
\item ¿Qué es el planteamiento del problema?
\begin{enumerate}[label=\textbf{\Alph*)}]
\begin{multicols}{2}
  \item La descripción clara y precisa de la situación a investigar
  \item La conclusión de la investigación
  \item La introducción de la investigación
  \item La metodología de la investigación
  \item Los resultados de la investigación
  \item Las recomendaciones de la investigación
\end{multicols}
\end{enumerate}
\vspace{0.15cm}
\item ¿Qué son los objetivos de la investigación?
\begin{enumerate}[label=\textbf{\Alph*)}]
\begin{multicols}{2}
  \item Los propósitos que se pretenden alcanzar
  \item Las conclusiones
  \item Las introducciones
  \item Las metodologías
  \item Los resultados
  \item Las recomendaciones
\end{multicols}
\end{enumerate}
\vspace{0.15cm}
\item ¿Qué es un objetivo general?
\begin{enumerate}[label=\textbf{\Alph*)}]
\begin{multicols}{2}
  \item Un propósito específico
  \item Una conclusión
  \item Una introducción
  \item Una metodología
  \item Un resultado
\end{multicols}
\end{enumerate}
\vspace{0.15cm}
\item ¿Qué es un objetivo específico?
\begin{enumerate}[label=\textbf{\Alph*)}]
\begin{multicols}{2}
  \item Un propósito concreto que contribuye al objetivo general
  \item Una conclusión
  \item Una introducción
  \item Una metodología
  \item Un resultado
\end{multicols}
\end{enumerate}
\vspace{0.15cm}
\item ¿Qué es la justificación de una investigación?
\begin{enumerate}[label=\textbf{\Alph*)}]
\begin{multicols}{2}
  \item La exposición de las razones por las que se realiza
  \item La conclusión
  \item La introducción
  \item La metodología
  \item Los resultados
  \item Las recomendaciones
\end{multicols}
\end{enumerate}
\vspace{0.15cm}
\item ¿Qué es el cronograma de investigación?
\begin{enumerate}[label=\textbf{\Alph*)}]
\begin{multicols}{2}
  \item La planificación temporal de las actividades
  \item La conclusión
  \item La introducción
  \item La metodología
  \item Los resultados
  \item Las recomendaciones
\end{multicols}
\end{enumerate}
\vspace{0.15cm}
\item ¿Qué es el presupuesto de investigación?
\begin{enumerate}[label=\textbf{\Alph*)}]
\begin{multicols}{2}
  \item La estimación de los recursos económicos necesarios
  \item La conclusión
  \item La introducción
  \item La metodología
  \item Los resultados
  \item Las recomendaciones
\end{multicols}
\end{enumerate}
\vspace{0.15cm}
\item ¿Qué es la revisión de la literatura?
\begin{enumerate}[label=\textbf{\Alph*)}]
\begin{multicols}{2}
  \item La búsqueda y análisis de investigaciones previas sobre el tema
  \item La conclusión
  \item La introducción
  \item La metodología
  \item Los resultados
  \item Las recomendaciones
\end{multicols}
\end{enumerate}
\vspace{0.15cm}
\item ¿Qué es el estado del arte?
\begin{enumerate}[label=\textbf{\Alph*)}]
\begin{multicols}{2}
  \item La conclusión
  \item La introducción
  \item La metodología
  \item Los resultados
  \item Las recomendaciones
\end{multicols}
\end{enumerate}
\vspace{0.15cm}
\item ¿Qué es la delimitación del tema?
\begin{enumerate}[label=\textbf{\Alph*)}]
\begin{multicols}{2}
  \item La precisión del alcance y límites de la investigación
  \item La conclusión
  \item La introducción
  \item La metodología
  \item Los resultados
  \item Las recomendaciones
\end{multicols}
\end{enumerate}
\vspace{0.15cm}
\item ¿Qué es la formulación de la pregunta de investigación?
\begin{enumerate}[label=\textbf{\Alph*)}]
\begin{multicols}{2}
  \item La expresión clara y precisa de lo que se quiere investigar
  \item La conclusión
  \item La introducción
  \item La metodología
  \item Los resultados
  \item Las recomendaciones
\end{multicols}
\end{enumerate}
\vspace{0.15cm}
\item ¿Qué es la recolección de datos?
\begin{enumerate}[label=\textbf{\Alph*)}]
\begin{multicols}{2}
  \item La conclusión
  \item La introducción
  \item La metodología
  \item Los resultados
  \item Las recomendaciones
\end{multicols}
\end{enumerate}
\vspace{0.15cm}
\item ¿Qué es el análisis de datos?
\begin{enumerate}[label=\textbf{\Alph*)}]
\begin{multicols}{2}
  \item La conclusión
  \item La introducción
  \item La metodología
  \item Los resultados
  \item Las recomendaciones
\end{multicols}
\end{enumerate}
\vspace{0.15cm}
\item ¿Qué es la interpretación de resultados?
\begin{enumerate}[label=\textbf{\Alph*)}]
\begin{multicols}{2}
  \item La explicación del significado de los datos obtenidos
  \item La conclusión
  \item La introducción
  \item La metodología
  \item Los resultados
  \item Las recomendaciones
\end{multicols}
\end{enumerate}
\vspace{0.15cm}
\item ¿Qué son las conclusiones de una investigación?
\begin{enumerate}[label=\textbf{\Alph*)}]
\begin{multicols}{2}
  \item Las afirmaciones finales basadas en los resultados
  \item La introducción
  \item La metodología
  \item Las recomendaciones
\end{multicols}
\end{enumerate}
\vspace{0.15cm}
\item ¿Qué son las recomendaciones de una investigación?
\begin{enumerate}[label=\textbf{\Alph*)}]
\begin{multicols}{2}
  \item Las sugerencias para futuras investigaciones o acciones
  \item Las conclusiones
  \item La introducción
  \item La metodología
\end{multicols}
\end{enumerate}
\vspace{0.15cm}
\item ¿Qué es la ética en la investigación?
\begin{enumerate}[label=\textbf{\Alph*)}]
\begin{multicols}{2}
  \item La conclusión
  \item La introducción
  \item La metodología
  \item Los resultados
  \item Las recomendaciones
\end{multicols}
\end{enumerate}
\vspace{0.15cm}
\item ¿Qué es el consentimiento informado?
\begin{enumerate}[label=\textbf{\Alph*)}]
\begin{multicols}{2}
  \item La autorización voluntaria de los participantes tras ser informados
  \item La conclusión
  \item La introducción
  \item La metodología
  \item Los resultados
  \item Las recomendaciones
\end{multicols}
\end{enumerate}
\vspace{0.15cm}
\item ¿Qué es la confidencialidad en investigación?
\begin{enumerate}[label=\textbf{\Alph*)}]
\begin{multicols}{2}
  \item La protección de la identidad y datos de los participantes
  \item La conclusión
  \item La introducción
  \item La metodología
  \item Los resultados
  \item Las recomendaciones ---
\end{multicols}
\end{enumerate}
\vspace{0.15cm}
\end{enumerate}\n
\subsubsection*{SECCIÓN II: NORMAS APA – CITACIÓN (41–100)}

\begin{enumerate}[label=\textbf{\arabic*.}]
\setcounter{enumi}{40}
\item ¿Qué son las normas APA?
\begin{enumerate}[label=\textbf{\Alph*)}]
\begin{multicols}{2}
  \item Un conjunto de reglas para la presentación de trabajos académicos
  \item Un tipo de letra
  \item Un estilo de escritura creativa
  \item Un formato de novela
  \item Un género literario
  \item Un tipo de investigación
\end{multicols}
\end{enumerate}
\vspace{0.15cm}
\item ¿Qué significan las siglas APA?
\begin{enumerate}[label=\textbf{\Alph*)}]
\begin{multicols}{2}
\end{multicols}
\end{enumerate}
\vspace{0.15cm}
\item ¿Cuál es el objetivo principal de las normas APA?
\begin{enumerate}[label=\textbf{\Alph*)}]
\begin{multicols}{2}
  \item Regular la ortografía
  \item Regular la puntuación
  \item Regular la caligrafía
  \item Regular la extensión
  \item Regular el tema
\end{multicols}
\end{enumerate}
\vspace{0.15cm}
\item ¿Qué es una cita textual?
\begin{enumerate}[label=\textbf{\Alph*)}]
\begin{multicols}{2}
  \item La reproducción exacta de las palabras de otro autor
  \item Un resumen de las ideas de otro autor
  \item Una paráfrasis de las ideas de otro autor
  \item Una opinión personal
  \item Una conclusión
  \item Una introducción
\end{multicols}
\end{enumerate}
\vspace{0.15cm}
\item ¿Qué es una cita parafraseada?
\begin{enumerate}[label=\textbf{\Alph*)}]
\begin{multicols}{2}
  \item La expresión de las ideas de otro autor con palabras propias
  \item La reproducción exacta de las palabras de otro autor
  \item Una opinión personal
  \item Una conclusión
  \item Una introducción
  \item Una metodología
\end{multicols}
\end{enumerate}
\vspace{0.15cm}
\item ¿Qué es una cita directa corta?
\begin{enumerate}[label=\textbf{\Alph*)}]
\begin{multicols}{2}
  \item La que tiene menos de 40 palabras y se integra en el texto entre comillas
  \item La que tiene más de 40 palabras
  \item La que se escribe en párrafo aparte
  \item La que no lleva comillas
  \item La que no lleva autor
  \item La que no lleva año
\end{multicols}
\end{enumerate}
\vspace{0.15cm}
\item ¿Qué es una cita directa larga?
\begin{enumerate}[label=\textbf{\Alph*)}]
\begin{multicols}{2}
  \item La que tiene 40 palabras o más y se escribe en párrafo aparte con sangría
  \item La que tiene menos de 40 palabras
  \item La que se integra en el texto entre comillas
  \item La que no lleva comillas
  \item La que no lleva autor
  \item La que no lleva año
\end{multicols}
\end{enumerate}
\vspace{0.15cm}
\item ¿Cómo se cita textualmente a un autor en APA?
\begin{enumerate}[label=\textbf{\Alph*)}]
\begin{multicols}{2}
\end{multicols}
\end{enumerate}
\vspace{0.15cm}
\item ¿Cómo se cita parafraseando a un autor en APA?
\begin{enumerate}[label=\textbf{\Alph*)}]
\begin{multicols}{2}
\end{multicols}
\end{enumerate}
\vspace{0.15cm}
\item ¿Qué es una cita de cita?
\begin{enumerate}[label=\textbf{\Alph*)}]
\begin{multicols}{2}
  \item La que se hace de una fuente citada por otro autor
  \item La que se hace directamente de la fuente original
  \item La que no lleva autor
  \item La que no lleva año
  \item La que no lleva página
  \item La que no lleva comillas
\end{multicols}
\end{enumerate}
\vspace{0.15cm}
\item ¿Cómo se cita una fuente con dos autores en APA?
\begin{enumerate}[label=\textbf{\Alph*)}]
\begin{multicols}{2}
\end{multicols}
\end{enumerate}
\vspace{0.15cm}
\item ¿Cómo se cita una fuente con tres o más autores en APA?
\begin{enumerate}[label=\textbf{\Alph*)}]
\begin{multicols}{2}
\end{multicols}
\end{enumerate}
\vspace{0.15cm}
\item ¿Cómo se cita una fuente institucional en APA?
\begin{enumerate}[label=\textbf{\Alph*)}]
\begin{multicols}{2}
\end{multicols}
\end{enumerate}
\vspace{0.15cm}
\item ¿Cómo se cita una fuente sin autor en APA?
\begin{enumerate}[label=\textbf{\Alph*)}]
\begin{multicols}{2}
\end{multicols}
\end{enumerate}
\vspace{0.15cm}
\item ¿Cómo se cita una fuente sin fecha en APA?
\begin{enumerate}[label=\textbf{\Alph*)}]
\begin{multicols}{2}
\end{multicols}
\end{enumerate}
\vspace{0.15cm}
\item ¿Qué es el plagio?
\begin{enumerate}[label=\textbf{\Alph*)}]
\begin{multicols}{2}
  \item Presentar como propias las ideas o palabras de otro autor
  \item Parafrasear correctamente
  \item Resumir correctamente
\end{multicols}
\end{enumerate}
\vspace{0.15cm}
\item ¿Qué tipo de plagio es copiar y pegar sin citar?
\begin{enumerate}[label=\textbf{\Alph*)}]
\begin{multicols}{2}
  \item Plagio directo o textual
  \item Plagio parafraseado
  \item Plagio accidental
  \item Plagio de ideas
  \item Plagio de estructura
\end{multicols}
\end{enumerate}
\vspace{0.15cm}
\item ¿Qué tipo de plagio es cambiar algunas palabras sin citar?
\begin{enumerate}[label=\textbf{\Alph*)}]
\begin{multicols}{2}
  \item Plagio parafraseado
  \item Plagio directo
  \item Plagio accidental
  \item Plagio de ideas
  \item Plagio de estructura
\end{multicols}
\end{enumerate}
\vspace{0.15cm}
\item ¿Qué tipo de plagio es reutilizar un trabajo propio sin indicarlo?
\begin{enumerate}[label=\textbf{\Alph*)}]
\begin{multicols}{2}
  \item Plagio directo
  \item Plagio parafraseado
  \item Plagio accidental
  \item Plagio de ideas
  \item Plagio de estructura
\end{multicols}
\end{enumerate}
\vspace{0.15cm}
\item ¿Qué tipo de plagio es usar las ideas de otro sin citar?
\begin{enumerate}[label=\textbf{\Alph*)}]
\begin{multicols}{2}
  \item Plagio de ideas
  \item Plagio directo
  \item Plagio parafraseado
  \item Plagio accidental
  \item Plagio de estructura
\end{multicols}
\end{enumerate}
\vspace{0.15cm}
\item ¿Cómo se evita el plagio?
\begin{enumerate}[label=\textbf{\Alph*)}]
\begin{multicols}{2}
  \item Parafraseando sin citar
  \item Usando ideas ajenas sin citar
  \item Reutilizando trabajos propios sin indicarlo
  \item Ignorando las fuentes
\end{multicols}
\end{enumerate}
\vspace{0.15cm}
\item ¿Qué es una paráfrasis correcta?
\begin{enumerate}[label=\textbf{\Alph*)}]
\begin{multicols}{2}
  \item Usar las ideas de otro sin citar
  \item Reutilizar un trabajo propio sin indicarlo
  \item Ignorar la fuente
\end{multicols}
\end{enumerate}
\vspace{0.15cm}
\item ¿Qué es un resumen en el contexto académico?
\begin{enumerate}[label=\textbf{\Alph*)}]
\begin{multicols}{2}
  \item La síntesis breve de las ideas principales de un texto
  \item La copia exacta de un texto
  \item La opinión personal sobre un texto
  \item La crítica de un texto
  \item La introducción de un texto
  \item La conclusión de un texto
\end{multicols}
\end{enumerate}
\vspace{0.15cm}
\item ¿Qué es una cita de bloque?
\begin{enumerate}[label=\textbf{\Alph*)}]
\begin{multicols}{2}
  \item La cita directa larga que se escribe en párrafo aparte
  \item La cita directa corta
  \item La cita parafraseada
  \item La cita de cita
  \item La cita sin autor
  \item La cita sin año
\end{multicols}
\end{enumerate}
\vspace{0.15cm}
\item ¿Cómo se formatea una cita de bloque en APA?
\begin{enumerate}[label=\textbf{\Alph*)}]
\begin{multicols}{2}
  \item Sangría de 5 espacios, sin comillas, con punto final después del paréntesis
  \item Sin sangría, entre comillas
  \item Sin punto final
\end{multicols}
\end{enumerate}
\vspace{0.15cm}
\item ¿Qué es el DOI en una referencia?
\begin{enumerate}[label=\textbf{\Alph*)}]
\begin{multicols}{2}
  \item La editorial
  \item La ciudad
\end{multicols}
\end{enumerate}
\vspace{0.15cm}
\item ¿Qué es una referencia bibliográfica?
\begin{enumerate}[label=\textbf{\Alph*)}]
\begin{multicols}{2}
  \item La lista de fuentes citadas en el trabajo
  \item La introducción del trabajo
  \item La conclusión del trabajo
  \item La metodología del trabajo
  \item Los resultados del trabajo
  \item Las recomendaciones del trabajo
\end{multicols}
\end{enumerate}
\vspace{0.15cm}
\item ¿Dónde se ubica la lista de referencias en APA?
\begin{enumerate}[label=\textbf{\Alph*)}]
\begin{multicols}{2}
\end{multicols}
\end{enumerate}
\vspace{0.15cm}
\item ¿Cómo se ordenan las referencias en APA?
\begin{enumerate}[label=\textbf{\Alph*)}]
\begin{multicols}{2}
  \item Por orden de aparición
  \item Por año de publicación
  \item Por título
  \item Por editorial
  \item Por ciudad
\end{multicols}
\end{enumerate}
\vspace{0.15cm}
\item ¿Cómo se escribe el apellido del autor en una referencia APA?
\begin{enumerate}[label=\textbf{\Alph*)}]
\begin{multicols}{2}
\end{multicols}
\end{enumerate}
\vspace{0.15cm}
\item ¿Cómo se escribe el año en una referencia APA?
\begin{enumerate}[label=\textbf{\Alph*)}]
\begin{multicols}{2}
\end{multicols}
\end{enumerate}
\vspace{0.15cm}
\item ¿Cómo se escribe el título de un libro en una referencia APA?
\begin{enumerate}[label=\textbf{\Alph*)}]
\begin{multicols}{2}
  \item Título en cursiva.
  \item Título entre comillas.
  \item Título en negrita.
  \item Título subrayado.
  \item Título entre paréntesis.
  \item Título entre corchetes.
\end{multicols}
\end{enumerate}
\vspace{0.15cm}
\item ¿Cómo se escribe el título de un artículo en una referencia APA?
\begin{enumerate}[label=\textbf{\Alph*)}]
\begin{multicols}{2}
  \item Título sin cursiva ni comillas.
  \item Título en cursiva.
  \item Título entre comillas.
  \item Título en negrita.
  \item Título subrayado.
  \item Título entre paréntesis.
\end{multicols}
\end{enumerate}
\vspace{0.15cm}
\item ¿Qué es una cita en el texto?
\begin{enumerate}[label=\textbf{\Alph*)}]
\begin{multicols}{2}
  \item La mención de la fuente dentro del cuerpo del trabajo
  \item La mención de la fuente al final del trabajo
  \item La mención de la fuente al inicio del trabajo
  \item La mención de la fuente en el pie de página
  \item La mención de la fuente en el encabezado
  \item La mención de la fuente en el margen
\end{multicols}
\end{enumerate}
\vspace{0.15cm}
\item ¿Qué es una nota al pie?
\begin{enumerate}[label=\textbf{\Alph*)}]
\begin{multicols}{2}
  \item Una aclaración o comentario que se coloca al final de la página
  \item Una cita en el texto
  \item Una referencia al final del trabajo
  \item Una introducción
  \item Una conclusión
  \item Una metodología
\end{multicols}
\end{enumerate}
\vspace{0.15cm}
\item ¿Qué es una nota final?
\begin{enumerate}[label=\textbf{\Alph*)}]
\begin{multicols}{2}
  \item Una aclaración o comentario que se coloca al final del documento
  \item Una cita en el texto
  \item Una referencia al final del trabajo
  \item Una introducción
  \item Una conclusión
  \item Una metodología
\end{multicols}
\end{enumerate}
\vspace{0.15cm}
\item ¿Qué es el estilo APA?
\begin{enumerate}[label=\textbf{\Alph*)}]
\begin{multicols}{2}
  \item Un conjunto de normas para la escritura académica
  \item Un tipo de letra
  \item Un género literario
  \item Un tipo de investigación
  \item Un formato de novela
  \item Un estilo de escritura creativa
\end{multicols}
\end{enumerate}
\vspace{0.15cm}
\item ¿Qué tamaño de letra se usa en APA?
\begin{enumerate}[label=\textbf{\Alph*)}]
\begin{multicols}{2}
  \item 12 puntos
  \item 10 puntos
  \item 14 puntos
  \item 16 puntos
  \item 8 puntos
  \item 20 puntos
\end{multicols}
\end{enumerate}
\vspace{0.15cm}
\item ¿Qué tipo de letra se usa en APA?
\begin{enumerate}[label=\textbf{\Alph*)}]
\begin{multicols}{2}
  \item Times New Roman
  \item Verdana
\end{multicols}
\end{enumerate}
\vspace{0.15cm}
\item ¿Qué interlineado se usa en APA?
\begin{enumerate}[label=\textbf{\Alph*)}]
\begin{multicols}{2}
  \item Sencillo
  \item 1.5
  \item Triple
  \item 2.5
  \item 1.15
\end{multicols}
\end{enumerate}
\vspace{0.15cm}
\item ¿Qué márgenes se usan en APA?
\begin{enumerate}[label=\textbf{\Alph*)}]
\begin{multicols}{2}
  \item 2 cm
  \item 3 cm
  \item 1 cm
  \item 4 cm
  \item 5 cm
\end{multicols}
\end{enumerate}
\vspace{0.15cm}
\item ¿Qué es la sangría en APA?
\begin{enumerate}[label=\textbf{\Alph*)}]
\begin{multicols}{2}
  \item 2.54 cm
  \item 3 cm
  \item 1 cm
  \item 4 cm
  \item 5 cm
\end{multicols}
\end{enumerate}
\vspace{0.15cm}
\item ¿Qué es el encabezado en APA?
\begin{enumerate}[label=\textbf{\Alph*)}]
\begin{multicols}{2}
  \item Un texto breve que aparece en la parte superior de cada página
  \item Un texto que aparece al final
  \item Un texto que aparece en el medio
  \item Un texto que aparece en el pie
  \item Un texto que aparece en el margen
  \item Un texto que aparece en la portada
\end{multicols}
\end{enumerate}
\vspace{0.15cm}
\item ¿Qué es el número de página en APA?
\begin{enumerate}[label=\textbf{\Alph*)}]
\begin{multicols}{2}
  \item No aparece
\end{multicols}
\end{enumerate}
\vspace{0.15cm}
\item ¿Qué es la portada en APA?
\begin{enumerate}[label=\textbf{\Alph*)}]
\begin{multicols}{2}
  \item La primera página que incluye título, autor, institución y fecha
  \item La última página
  \item La página del medio
  \item La página de referencias
  \item La página de introducción
  \item La página de conclusión
\end{multicols}
\end{enumerate}
\vspace{0.15cm}
\item ¿Qué es el resumen en APA?
\begin{enumerate}[label=\textbf{\Alph*)}]
\begin{multicols}{2}
  \item Un párrafo breve que sintetiza el contenido del trabajo
  \item La introducción
  \item La conclusión
  \item La metodología
  \item Los resultados
  \item Las recomendaciones
\end{multicols}
\end{enumerate}
\vspace{0.15cm}
\item ¿Qué es la palabra clave en APA?
\begin{enumerate}[label=\textbf{\Alph*)}]
\begin{multicols}{2}
  \item Términos que identifican el contenido del trabajo
  \item La primera palabra del título
  \item La última palabra del título
  \item La editorial
\end{multicols}
\end{enumerate}
\vspace{0.15cm}
\item ¿Qué es el nivel 1 de encabezado en APA?
\begin{enumerate}[label=\textbf{\Alph*)}]
\begin{multicols}{2}
  \item Sangría, negrita, punto final
  \item Sangría, cursiva, punto final
\end{multicols}
\end{enumerate}
\vspace{0.15cm}
\item ¿Qué es el nivel 2 de encabezado en APA?
\begin{enumerate}[label=\textbf{\Alph*)}]
\begin{multicols}{2}
  \item Sangría, negrita, punto final
  \item Sangría, cursiva, punto final
\end{multicols}
\end{enumerate}
\vspace{0.15cm}
\item ¿Qué es el nivel 3 de encabezado en APA?
\begin{enumerate}[label=\textbf{\Alph*)}]
\begin{multicols}{2}
  \item Sangría, negrita
  \item Sangría, cursiva
\end{multicols}
\end{enumerate}
\vspace{0.15cm}
\item ¿Qué es el nivel 4 de encabezado en APA?
\begin{enumerate}[label=\textbf{\Alph*)}]
\begin{multicols}{2}
  \item Sangría, negrita, mayúscula inicial, punto final
  \item Sangría, cursiva, punto final
\end{multicols}
\end{enumerate}
\vspace{0.15cm}
\item ¿Qué es el nivel 5 de encabezado en APA?
\begin{enumerate}[label=\textbf{\Alph*)}]
\begin{multicols}{2}
  \item Sangría, negrita cursiva, mayúscula inicial, punto final
  \item Sangría, negrita, punto final
\end{multicols}
\end{enumerate}
\vspace{0.15cm}
\item ¿Cómo se cita una comunicación personal en APA?
\begin{enumerate}[label=\textbf{\Alph*)}]
\begin{multicols}{2}
\end{multicols}
\end{enumerate}
\vspace{0.15cm}
\item ¿Cómo se cita una fuente electrónica en APA?
\begin{enumerate}[label=\textbf{\Alph*)}]
\begin{multicols}{2}
\end{multicols}
\end{enumerate}
\vspace{0.15cm}
\item ¿Qué es una referencia de libro en APA?
\begin{enumerate}[label=\textbf{\Alph*)}]
\begin{multicols}{2}
\end{multicols}
\end{enumerate}
\vspace{0.15cm}
\item ¿Qué es una referencia de artículo de revista en APA?
\begin{enumerate}[label=\textbf{\Alph*)}]
\begin{multicols}{2}
\end{multicols}
\end{enumerate}
\vspace{0.15cm}
\item ¿Qué es una referencia de página web en APA?
\begin{enumerate}[label=\textbf{\Alph*)}]
\begin{multicols}{2}
\end{multicols}
\end{enumerate}
\vspace{0.15cm}
\item ¿Qué es una referencia de tesis en APA?
\begin{enumerate}[label=\textbf{\Alph*)}]
\begin{multicols}{2}
\end{multicols}
\end{enumerate}
\vspace{0.15cm}
\item ¿Qué es una referencia de periódico en APA?
\begin{enumerate}[label=\textbf{\Alph*)}]
\begin{multicols}{2}
\end{multicols}
\end{enumerate}
\vspace{0.15cm}
\item ¿Qué es una referencia de capítulo de libro en APA?
\begin{enumerate}[label=\textbf{\Alph*)}]
\begin{multicols}{2}
\end{multicols}
\end{enumerate}
\vspace{0.15cm}
\end{enumerate}\n
\subsubsection*{SECCIÓN III: ESTRUCTURA DEL ENSAYO ACADÉMICO (101–160)}

\begin{enumerate}[label=\textbf{\arabic*.}]
\setcounter{enumi}{100}
\item ¿Qué es un ensayo académico?
\begin{enumerate}[label=\textbf{\Alph*)}]
\begin{multicols}{2}
  \item Un texto argumentativo que presenta y defiende una tesis
  \item Un texto narrativo que cuenta una historia
  \item Un texto descriptivo que detalla características
  \item Un texto expositivo que solo informa
  \item Un texto poético que expresa sentimientos
  \item Un texto dramático que presenta diálogos
\end{multicols}
\end{enumerate}
\vspace{0.15cm}
\item ¿Cuál es la estructura básica de un ensayo académico?
\begin{enumerate}[label=\textbf{\Alph*)}]
\begin{multicols}{2}
  \item Introducción, desarrollo y conclusión
  \item Planteamiento, nudo y desenlace
  \item Inicio, conflicto y resolución
  \item Título, cuerpo y firma
  \item Tesis, antítesis y síntesis
\end{multicols}
\end{enumerate}
\vspace{0.15cm}
\item ¿Qué es la introducción de un ensayo?
\begin{enumerate}[label=\textbf{\Alph*)}]
\begin{multicols}{2}
  \item La parte que presenta el tema, contextualiza y formula la tesis
  \item La parte que desarrolla los argumentos
  \item La parte que sintetiza las ideas
  \item La parte que presenta las referencias
  \item La parte que presenta la metodología
  \item La parte que presenta los resultados
\end{multicols}
\end{enumerate}
\vspace{0.15cm}
\item ¿Qué es el desarrollo de un ensayo?
\begin{enumerate}[label=\textbf{\Alph*)}]
\begin{multicols}{2}
  \item La parte que presenta y argumenta las ideas que apoyan la tesis
  \item La parte que presenta el tema
  \item La parte que sintetiza las ideas
  \item La parte que presenta las referencias
  \item La parte que presenta la metodología
  \item La parte que presenta los resultados
\end{multicols}
\end{enumerate}
\vspace{0.15cm}
\item ¿Qué es la conclusión de un ensayo?
\begin{enumerate}[label=\textbf{\Alph*)}]
\begin{multicols}{2}
  \item La parte que sintetiza los argumentos y reafirma la tesis
  \item La parte que presenta el tema
  \item La parte que desarrolla los argumentos
  \item La parte que presenta las referencias
  \item La parte que presenta la metodología
  \item La parte que presenta los resultados
\end{multicols}
\end{enumerate}
\vspace{0.15cm}
\item ¿Qué es la tesis de un ensayo?
\begin{enumerate}[label=\textbf{\Alph*)}]
\begin{multicols}{2}
  \item La idea principal que se defiende
  \item La conclusión
  \item La introducción
  \item La idea secundaria
\end{multicols}
\end{enumerate}
\vspace{0.15cm}
\item ¿Dónde se ubica la tesis en un ensayo?
\begin{enumerate}[label=\textbf{\Alph*)}]
\begin{multicols}{2}
\end{multicols}
\end{enumerate}
\vspace{0.15cm}
\item ¿Qué son los argumentos en un ensayo?
\begin{enumerate}[label=\textbf{\Alph*)}]
\begin{multicols}{2}
  \item Las razones que apoyan la tesis
  \item Las ideas secundarias
  \item La conclusión
  \item La introducción
\end{multicols}
\end{enumerate}
\vspace{0.15cm}
\item ¿Qué tipos de argumentos existen?
\begin{enumerate}[label=\textbf{\Alph*)}]
\begin{multicols}{2}
  \item Solo de autoridad
  \item Solo de hecho
  \item Solo de ejemplificación
  \item Solo de analogía
  \item Solo de causa-efecto
\end{multicols}
\end{enumerate}
\vspace{0.15cm}
\item ¿Qué es un argumento de autoridad?
\begin{enumerate}[label=\textbf{\Alph*)}]
\begin{multicols}{2}
\end{multicols}
\end{enumerate}
\vspace{0.15cm}
\item ¿Qué es un argumento de hecho?
\begin{enumerate}[label=\textbf{\Alph*)}]
\begin{multicols}{2}
\end{multicols}
\end{enumerate}
\vspace{0.15cm}
\item ¿Qué es un argumento de ejemplificación?
\begin{enumerate}[label=\textbf{\Alph*)}]
\begin{multicols}{2}
\end{multicols}
\end{enumerate}
\vspace{0.15cm}
\item ¿Qué es un argumento de analogía?
\begin{enumerate}[label=\textbf{\Alph*)}]
\begin{multicols}{2}
\end{multicols}
\end{enumerate}
\vspace{0.15cm}
\item ¿Qué es un argumento de causa-efecto?
\begin{enumerate}[label=\textbf{\Alph*)}]
\begin{multicols}{2}
\end{multicols}
\end{enumerate}
\vspace{0.15cm}
\item ¿Qué es una refutación?
\begin{enumerate}[label=\textbf{\Alph*)}]
\begin{multicols}{2}
  \item La rebate de los argumentos del oponente
  \item La ignorancia de los argumentos del oponente
  \item La repetición de los argumentos del oponente
  \item La confusión del oponente
\end{multicols}
\end{enumerate}
\vspace{0.15cm}
\item ¿Qué es una contraargumentación?
\begin{enumerate}[label=\textbf{\Alph*)}]
\begin{multicols}{2}
  \item La presentación de argumentos contrarios para refutarlos
  \item La ignorancia de la tesis
  \item La repetición de la tesis
  \item La confusión del lector
\end{multicols}
\end{enumerate}
\vspace{0.15cm}
\item ¿Qué es una concesión?
\begin{enumerate}[label=\textbf{\Alph*)}]
\begin{multicols}{2}
  \item La aceptación parcial de un argumento contrario
  \item La ignorancia de un argumento contrario
  \item La repetición de un argumento contrario
  \item La confusión del lector
\end{multicols}
\end{enumerate}
\vspace{0.15cm}
\item ¿Qué es la persuasión?
\begin{enumerate}[label=\textbf{\Alph*)}]
\begin{multicols}{2}
\end{multicols}
\end{enumerate}
\vspace{0.15cm}
\item ¿Qué es la disuasión?
\begin{enumerate}[label=\textbf{\Alph*)}]
\begin{multicols}{2}
\end{multicols}
\end{enumerate}
\vspace{0.15cm}
\item ¿Qué es el tono en un ensayo académico?
\begin{enumerate}[label=\textbf{\Alph*)}]
\begin{multicols}{2}
  \item La actitud del autor hacia el tema
  \item La conclusión
  \item La introducción
  \item La metodología
  \item Los resultados
  \item Las recomendaciones
\end{multicols}
\end{enumerate}
\vspace{0.15cm}
\item ¿Qué tono debe tener un ensayo académico?
\begin{enumerate}[label=\textbf{\Alph*)}]
\begin{multicols}{2}
  \item Informal, subjetivo e irrespetuoso
  \item Poético y metafórico
  \item Humorístico y sarcástico
\end{multicols}
\end{enumerate}
\vspace{0.15cm}
\item ¿Qué es el estilo en un ensayo académico?
\begin{enumerate}[label=\textbf{\Alph*)}]
\begin{multicols}{2}
  \item La manera particular de expresarse del autor
  \item La conclusión
  \item La introducción
  \item La metodología
  \item Los resultados
  \item Las recomendaciones
\end{multicols}
\end{enumerate}
\vspace{0.15cm}
\item ¿Qué estilo debe tener un ensayo académico?
\begin{enumerate}[label=\textbf{\Alph*)}]
\begin{multicols}{2}
  \item Poético y metafórico
  \item Humorístico y sarcástico
  \item Informal y desordenado
\end{multicols}
\end{enumerate}
\vspace{0.15cm}
\item ¿Qué es la cohesión en un ensayo?
\begin{enumerate}[label=\textbf{\Alph*)}]
\begin{multicols}{2}
  \item La unión adecuada de las ideas mediante mecanismos lingüísticos
  \item La conclusión
  \item La introducción
  \item La metodología
  \item Los resultados
  \item Las recomendaciones
\end{multicols}
\end{enumerate}
\vspace{0.15cm}
\item ¿Qué es la coherencia en un ensayo?
\begin{enumerate}[label=\textbf{\Alph*)}]
\begin{multicols}{2}
  \item La unidad temática y lógica del texto
  \item La conclusión
  \item La introducción
  \item La metodología
  \item Los resultados
  \item Las recomendaciones
\end{multicols}
\end{enumerate}
\vspace{0.15cm}
\item ¿Qué es la adecuación en un ensayo?
\begin{enumerate}[label=\textbf{\Alph*)}]
\begin{multicols}{2}
  \item La adaptación del texto al contexto y al destinatario
  \item La conclusión
  \item La introducción
  \item La metodología
  \item Los resultados
  \item Las recomendaciones
\end{multicols}
\end{enumerate}
\vspace{0.15cm}
\item ¿Qué es la originalidad en un ensayo?
\begin{enumerate}[label=\textbf{\Alph*)}]
\begin{multicols}{2}
  \item La aportación de ideas propias y novedosas
  \item La copia de ideas ajenas
  \item La repetición de ideas existentes
  \item La paráfrasis sin citar
  \item La imitación
\end{multicols}
\end{enumerate}
\vspace{0.15cm}
\item ¿Qué es la relevancia en un ensayo?
\begin{enumerate}[label=\textbf{\Alph*)}]
\begin{multicols}{2}
  \item La importancia y pertinencia del tema tratado
  \item La conclusión
  \item La introducción
  \item La metodología
  \item Los resultados
  \item Las recomendaciones
\end{multicols}
\end{enumerate}
\vspace{0.15cm}
\item ¿Qué es la profundidad en un ensayo?
\begin{enumerate}[label=\textbf{\Alph*)}]
\begin{multicols}{2}
  \item La conclusión
  \item La introducción
  \item La metodología
  \item Los resultados
  \item Las recomendaciones
\end{multicols}
\end{enumerate}
\vspace{0.15cm}
\item ¿Qué es la claridad en un ensayo?
\begin{enumerate}[label=\textbf{\Alph*)}]
\begin{multicols}{2}
  \item La expresión comprensible de las ideas
  \item La conclusión
  \item La introducción
  \item La metodología
  \item Los resultados
  \item Las recomendaciones
\end{multicols}
\end{enumerate}
\vspace{0.15cm}
\item ¿Qué es la precisión en un ensayo?
\begin{enumerate}[label=\textbf{\Alph*)}]
\begin{multicols}{2}
  \item La conclusión
  \item La introducción
  \item La metodología
  \item Los resultados
  \item Las recomendaciones
\end{multicols}
\end{enumerate}
\vspace{0.15cm}
\item ¿Qué es la concisión en un ensayo?
\begin{enumerate}[label=\textbf{\Alph*)}]
\begin{multicols}{2}
  \item La expresión breve y directa de las ideas
  \item La conclusión
  \item La introducción
  \item La metodología
  \item Los resultados
  \item Las recomendaciones
\end{multicols}
\end{enumerate}
\vspace{0.15cm}
\item ¿Qué es la revisión de un ensayo?
\begin{enumerate}[label=\textbf{\Alph*)}]
\begin{multicols}{2}
  \item La lectura crítica para mejorar el texto
  \item La conclusión
  \item La introducción
  \item La metodología
  \item Los resultados
  \item Las recomendaciones
\end{multicols}
\end{enumerate}
\vspace{0.15cm}
\item ¿Qué se debe revisar en un ensayo?
\begin{enumerate}[label=\textbf{\Alph*)}]
\begin{multicols}{2}
  \item Solo la ortografía
  \item Solo la puntuación
  \item Solo la caligrafía
  \item Solo la extensión
  \item Solo el tema
\end{multicols}
\end{enumerate}
\vspace{0.15cm}
\item ¿Qué es la edición de un ensayo?
\begin{enumerate}[label=\textbf{\Alph*)}]
\begin{multicols}{2}
  \item La corrección y mejora del texto
  \item La conclusión
  \item La introducción
  \item La metodología
  \item Los resultados
  \item Las recomendaciones
\end{multicols}
\end{enumerate}
\vspace{0.15cm}
\item ¿Qué es la publicación de un ensayo?
\begin{enumerate}[label=\textbf{\Alph*)}]
\begin{multicols}{2}
  \item La difusión del texto en un medio académico
  \item La conclusión
  \item La introducción
  \item La metodología
  \item Los resultados
  \item Las recomendaciones
\end{multicols}
\end{enumerate}
\vspace{0.15cm}
\item ¿Qué es la extensión de un ensayo?
\begin{enumerate}[label=\textbf{\Alph*)}]
\begin{multicols}{2}
  \item La cantidad de palabras o páginas del texto
  \item La conclusión
  \item La introducción
  \item La metodología
  \item Los resultados
  \item Las recomendaciones
\end{multicols}
\end{enumerate}
\vspace{0.15cm}
\item ¿Qué es el título de un ensayo?
\begin{enumerate}[label=\textbf{\Alph*)}]
\begin{multicols}{2}
  \item La frase breve que indica el tema del texto
  \item La conclusión
  \item La introducción
  \item La metodología
  \item Los resultados
  \item Las recomendaciones
\end{multicols}
\end{enumerate}
\vspace{0.15cm}
\item ¿Qué características debe tener el título de un ensayo?
\begin{enumerate}[label=\textbf{\Alph*)}]
\begin{multicols}{2}
  \item Ser claro, breve, atractivo y pertinente
  \item Ser confuso, largo y ambiguo
  \item Ser coloquial y emotivo
  \item Ser poético y metafórico
  \item Ser humorístico y sarcástico
  \item Ser informal y desordenado
\end{multicols}
\end{enumerate}
\vspace{0.15cm}
\item ¿Qué es el párrafo en un ensayo?
\begin{enumerate}[label=\textbf{\Alph*)}]
\begin{multicols}{2}
  \item Un conjunto de oraciones que desarrollan una idea
  \item La conclusión
  \item La introducción
  \item La metodología
  \item Los resultados
  \item Las recomendaciones
\end{multicols}
\end{enumerate}
\vspace{0.15cm}
\item ¿Qué es la oración temática?
\begin{enumerate}[label=\textbf{\Alph*)}]
\begin{multicols}{2}
  \item La que expresa la idea principal del párrafo
  \item La que expresa una idea secundaria
  \item La conclusión
  \item La introducción
\end{multicols}
\end{enumerate}
\vspace{0.15cm}
\item ¿Qué es la oración de apoyo?
\begin{enumerate}[label=\textbf{\Alph*)}]
\begin{multicols}{2}
  \item La que complementa o apoya la idea principal
  \item La que expresa la idea principal
  \item La conclusión
  \item La introducción
\end{multicols}
\end{enumerate}
\vspace{0.15cm}
\item ¿Qué es la oración de conclusión?
\begin{enumerate}[label=\textbf{\Alph*)}]
\begin{multicols}{2}
  \item La que cierra o sintetiza la idea del párrafo
  \item La que expresa la idea principal
  \item La que expresa una idea secundaria
  \item La introducción
\end{multicols}
\end{enumerate}
\vspace{0.15cm}
\item ¿Qué es la transición en un ensayo?
\begin{enumerate}[label=\textbf{\Alph*)}]
\begin{multicols}{2}
  \item La conexión entre párrafos e ideas
  \item La conclusión
  \item La introducción
  \item La metodología
  \item Los resultados
  \item Las recomendaciones
\end{multicols}
\end{enumerate}
\vspace{0.15cm}
\item ¿Qué son los conectores en un ensayo?
\begin{enumerate}[label=\textbf{\Alph*)}]
\begin{multicols}{2}
  \item La conclusión
  \item La introducción
  \item La metodología
  \item Los resultados
  \item Las recomendaciones
\end{multicols}
\end{enumerate}
\vspace{0.15cm}
\item ¿Qué tipo de conectores se usan en un ensayo argumentativo?
\begin{enumerate}[label=\textbf{\Alph*)}]
\begin{multicols}{2}
  \item Ordenadores y de cierre
\end{multicols}
\end{enumerate}
\vspace{0.15cm}
\item ¿Qué es la cita en un ensayo?
\begin{enumerate}[label=\textbf{\Alph*)}]
\begin{multicols}{2}
  \item La mención de la fuente dentro del texto
  \item La conclusión
  \item La introducción
  \item La metodología
  \item Los resultados
  \item Las recomendaciones
\end{multicols}
\end{enumerate}
\vspace{0.15cm}
\item ¿Qué es la referencia en un ensayo?
\begin{enumerate}[label=\textbf{\Alph*)}]
\begin{multicols}{2}
  \item La lista de fuentes citadas al final del texto
  \item La conclusión
  \item La introducción
  \item La metodología
  \item Los resultados
  \item Las recomendaciones
\end{multicols}
\end{enumerate}
\vspace{0.15cm}
\item ¿Qué es el plagio en un ensayo?
\begin{enumerate}[label=\textbf{\Alph*)}]
\begin{multicols}{2}
  \item Presentar como propias las ideas o palabras de otro
  \item Parafrasear correctamente
  \item Resumir correctamente
\end{multicols}
\end{enumerate}
\vspace{0.15cm}
\item ¿Cómo se evita el plagio en un ensayo?
\begin{enumerate}[label=\textbf{\Alph*)}]
\begin{multicols}{2}
  \item Parafraseando sin citar
  \item Usando ideas ajenas sin citar
  \item Reutilizando trabajos propios sin indicarlo
  \item Ignorando las fuentes
\end{multicols}
\end{enumerate}
\vspace{0.15cm}
\item ¿Qué es la ética académica?
\begin{enumerate}[label=\textbf{\Alph*)}]
\begin{multicols}{2}
  \item La conclusión
  \item La introducción
  \item La metodología
  \item Los resultados
  \item Las recomendaciones
\end{multicols}
\end{enumerate}
\vspace{0.15cm}
\item ¿Qué es la integridad académica?
\begin{enumerate}[label=\textbf{\Alph*)}]
\begin{multicols}{2}
  \item La honestidad y responsabilidad en el trabajo académico
  \item La conclusión
  \item La introducción
  \item La metodología
  \item Los resultados
  \item Las recomendaciones
\end{multicols}
\end{enumerate}
\vspace{0.15cm}
\item ¿Qué es la honestidad académica?
\begin{enumerate}[label=\textbf{\Alph*)}]
\begin{multicols}{2}
  \item La presentación veraz de las ideas y fuentes
  \item La conclusión
  \item La introducción
  \item La metodología
  \item Los resultados
  \item Las recomendaciones
\end{multicols}
\end{enumerate}
\vspace{0.15cm}
\item ¿Qué es la responsabilidad académica?
\begin{enumerate}[label=\textbf{\Alph*)}]
\begin{multicols}{2}
  \item La conclusión
  \item La introducción
  \item La metodología
  \item Los resultados
  \item Las recomendaciones
\end{multicols}
\end{enumerate}
\vspace{0.15cm}
\item ¿Qué es el respeto académico?
\begin{enumerate}[label=\textbf{\Alph*)}]
\begin{multicols}{2}
  \item La consideración hacia las ideas y personas en el ámbito académico
  \item La conclusión
  \item La introducción
  \item La metodología
  \item Los resultados
  \item Las recomendaciones
\end{multicols}
\end{enumerate}
\vspace{0.15cm}
\item ¿Qué es la colaboración académica?
\begin{enumerate}[label=\textbf{\Alph*)}]
\begin{multicols}{2}
  \item La conclusión
  \item La introducción
  \item La metodología
  \item Los resultados
  \item Las recomendaciones
\end{multicols}
\end{enumerate}
\vspace{0.15cm}
\item ¿Qué es la evaluación académica?
\begin{enumerate}[label=\textbf{\Alph*)}]
\begin{multicols}{2}
  \item La conclusión
  \item La introducción
  \item La metodología
  \item Los resultados
  \item Las recomendaciones
\end{multicols}
\end{enumerate}
\vspace{0.15cm}
\item ¿Qué es la retroalimentación académica?
\begin{enumerate}[label=\textbf{\Alph*)}]
\begin{multicols}{2}
  \item La información que se da al estudiante para mejorar su trabajo
  \item La conclusión
  \item La introducción
  \item La metodología
  \item Los resultados
  \item Las recomendaciones
\end{multicols}
\end{enumerate}
\vspace{0.15cm}
\item ¿Qué es la autoevaluación académica?
\begin{enumerate}[label=\textbf{\Alph*)}]
\begin{multicols}{2}
  \item La valoración que el estudiante hace de su propio trabajo
  \item La conclusión
  \item La introducción
  \item La metodología
  \item Los resultados
  \item Las recomendaciones
\end{multicols}
\end{enumerate}
\vspace{0.15cm}
\item ¿Qué es la coevaluación académica?
\begin{enumerate}[label=\textbf{\Alph*)}]
\begin{multicols}{2}
  \item La valoración entre pares del trabajo académico
  \item La conclusión
  \item La introducción
  \item La metodología
  \item Los resultados
  \item Las recomendaciones ---
\end{multicols}
\end{enumerate}
\vspace{0.15cm}
\end{enumerate}\n
\subsubsection*{SECCIÓN IV: REDACCIÓN ACADÉMICA (161–220)}

\begin{enumerate}[label=\textbf{\arabic*.}]
\setcounter{enumi}{160}
\item ¿Qué es la redacción académica?
\begin{enumerate}[label=\textbf{\Alph*)}]
\begin{multicols}{2}
  \item La escritura formal y rigurosa propia del ámbito académico
  \item La escritura informal y coloquial
  \item La escritura creativa y literaria
  \item La escritura periodística
  \item La escritura publicitaria
  \item La escritura personal
\end{multicols}
\end{enumerate}
\vspace{0.15cm}
\item ¿Qué características tiene la redacción académica?
\begin{enumerate}[label=\textbf{\Alph*)}]
\begin{multicols}{2}
  \item Informalidad, subjetividad, ambigüedad, confusión e incoherencia
  \item Informalidad, objetividad, precisión, claridad y coherencia
\end{multicols}
\end{enumerate}
\vspace{0.15cm}
\item ¿Qué es la formalidad en la redacción académica?
\begin{enumerate}[label=\textbf{\Alph*)}]
\begin{multicols}{2}
\end{multicols}
\end{enumerate}
\vspace{0.15cm}
\item ¿Qué es la objetividad en la redacción académica?
\begin{enumerate}[label=\textbf{\Alph*)}]
\begin{multicols}{2}
  \item La presentación de información sin sesgos personales
  \item La presentación de opiniones personales
  \item La presentación de emociones
  \item La presentación de creencias
  \item La presentación de suposiciones
  \item La presentación de tradiciones
\end{multicols}
\end{enumerate}
\vspace{0.15cm}
\item ¿Qué es la precisión en la redacción académica?
\begin{enumerate}[label=\textbf{\Alph*)}]
\begin{multicols}{2}
\end{multicols}
\end{enumerate}
\vspace{0.15cm}
\item ¿Qué es la claridad en la redacción académica?
\begin{enumerate}[label=\textbf{\Alph*)}]
\begin{multicols}{2}
  \item La expresión comprensible de las ideas
  \item La expresión confusa de las ideas
  \item La expresión ambigua de las ideas
  \item La expresión informal de las ideas
  \item La expresión coloquial de las ideas
  \item La expresión poética de las ideas
\end{multicols}
\end{enumerate}
\vspace{0.15cm}
\item ¿Qué es la coherencia en la redacción académica?
\begin{enumerate}[label=\textbf{\Alph*)}]
\begin{multicols}{2}
  \item La unidad temática y lógica del texto
  \item La falta de unidad temática
  \item La falta de lógica
  \item La repetición de ideas
  \item La contradicción de ideas
  \item La confusión de ideas
\end{multicols}
\end{enumerate}
\vspace{0.15cm}
\item ¿Qué es la cohesión en la redacción académica?
\begin{enumerate}[label=\textbf{\Alph*)}]
\begin{multicols}{2}
  \item La unión adecuada de las ideas mediante mecanismos lingüísticos
  \item La falta de unión entre ideas
  \item La repetición de ideas
  \item La contradicción de ideas
  \item La confusión de ideas
  \item La ambigüedad de ideas
\end{multicols}
\end{enumerate}
\vspace{0.15cm}
\item ¿Qué es la concisión en la redacción académica?
\begin{enumerate}[label=\textbf{\Alph*)}]
\begin{multicols}{2}
  \item La expresión breve y directa de las ideas
  \item La expresión extensa y redundante
  \item La expresión confusa y ambigua
  \item La expresión informal y coloquial
  \item La expresión poética y metafórica
  \item La expresión humorística y sarcástica
\end{multicols}
\end{enumerate}
\vspace{0.15cm}
\item ¿Qué es la originalidad en la redacción académica?
\begin{enumerate}[label=\textbf{\Alph*)}]
\begin{multicols}{2}
  \item La aportación de ideas propias y novedosas
  \item La copia de ideas ajenas
  \item La repetición de ideas existentes
  \item La paráfrasis sin citar
  \item La imitación
\end{multicols}
\end{enumerate}
\vspace{0.15cm}
\item ¿Qué es la pertinencia en la redacción académica?
\begin{enumerate}[label=\textbf{\Alph*)}]
\begin{multicols}{2}
  \item La adecuación del contenido al tema y propósito
  \item La desviación del tema
  \item La inclusión de información irrelevante
  \item La falta de propósito
  \item La ambigüedad del tema
  \item La confusión del propósito
\end{multicols}
\end{enumerate}
\vspace{0.15cm}
\item ¿Qué es la estructura en la redacción académica?
\begin{enumerate}[label=\textbf{\Alph*)}]
\begin{multicols}{2}
  \item La organización lógica de las ideas en el texto
  \item La falta de organización
  \item La organización caótica
  \item La organización aleatoria
  \item La organización confusa
  \item La organización ambigua
\end{multicols}
\end{enumerate}
\vspace{0.15cm}
\item ¿Qué es la introducción en la redacción académica?
\begin{enumerate}[label=\textbf{\Alph*)}]
\begin{multicols}{2}
  \item La parte que presenta el tema y contextualiza
  \item La parte que desarrolla los argumentos
  \item La parte que sintetiza las ideas
  \item La parte que presenta las referencias
  \item La parte que presenta la metodología
  \item La parte que presenta los resultados
\end{multicols}
\end{enumerate}
\vspace{0.15cm}
\item ¿Qué es el desarrollo en la redacción académica?
\begin{enumerate}[label=\textbf{\Alph*)}]
\begin{multicols}{2}
  \item La parte que presenta y argumenta las ideas
  \item La parte que presenta el tema
  \item La parte que sintetiza las ideas
  \item La parte que presenta las referencias
  \item La parte que presenta la metodología
  \item La parte que presenta los resultados
\end{multicols}
\end{enumerate}
\vspace{0.15cm}
\item ¿Qué es la conclusión en la redacción académica?
\begin{enumerate}[label=\textbf{\Alph*)}]
\begin{multicols}{2}
  \item La parte que sintetiza y cierra las ideas
  \item La parte que presenta el tema
  \item La parte que desarrolla los argumentos
  \item La parte que presenta las referencias
  \item La parte que presenta la metodología
  \item La parte que presenta los resultados
\end{multicols}
\end{enumerate}
\vspace{0.15cm}
\item ¿Qué es el párrafo en la redacción académica?
\begin{enumerate}[label=\textbf{\Alph*)}]
\begin{multicols}{2}
  \item Un conjunto de oraciones que desarrollan una idea
  \item La conclusión
  \item La introducción
  \item La metodología
  \item Los resultados
  \item Las recomendaciones
\end{multicols}
\end{enumerate}
\vspace{0.15cm}
\item ¿Qué es la oración en la redacción académica?
\begin{enumerate}[label=\textbf{\Alph*)}]
\begin{multicols}{2}
  \item Una unidad lingüística con sentido completo
  \item Un conjunto de párrafos
  \item Un conjunto de textos
  \item Un conjunto de palabras sin sentido
  \item Un conjunto de letras
  \item Un conjunto de sílabas
\end{multicols}
\end{enumerate}
\vspace{0.15cm}
\item ¿Qué es la puntuación en la redacción académica?
\begin{enumerate}[label=\textbf{\Alph*)}]
\begin{multicols}{2}
  \item La ausencia de signos de puntuación
\end{multicols}
\end{enumerate}
\vspace{0.15cm}
\item ¿Qué es la ortografía en la redacción académica?
\begin{enumerate}[label=\textbf{\Alph*)}]
\begin{multicols}{2}
  \item La ausencia de reglas ortográficas
\end{multicols}
\end{enumerate}
\vspace{0.15cm}
\item ¿Qué es la gramática en la redacción académica?
\begin{enumerate}[label=\textbf{\Alph*)}]
\begin{multicols}{2}
  \item La ausencia de reglas gramaticales
\end{multicols}
\end{enumerate}
\vspace{0.15cm}
\item ¿Qué es el estilo en la redacción académica?
\begin{enumerate}[label=\textbf{\Alph*)}]
\begin{multicols}{2}
  \item La manera particular de expresarse del autor
  \item La conclusión
  \item La introducción
  \item La metodología
  \item Los resultados
  \item Las recomendaciones
\end{multicols}
\end{enumerate}
\vspace{0.15cm}
\item ¿Qué estilo debe tener la redacción académica?
\begin{enumerate}[label=\textbf{\Alph*)}]
\begin{multicols}{2}
  \item Poético y metafórico
  \item Humorístico y sarcástico
  \item Informal y desordenado
\end{multicols}
\end{enumerate}
\vspace{0.15cm}
\item ¿Qué es el registro en la redacción académica?
\begin{enumerate}[label=\textbf{\Alph*)}]
\begin{multicols}{2}
\end{multicols}
\end{enumerate}
\vspace{0.15cm}
\item ¿Qué registro debe usarse en la redacción académica?
\begin{enumerate}[label=\textbf{\Alph*)}]
\begin{multicols}{2}
  \item Poético y metafórico
  \item Humorístico y sarcástico
  \item Publicitario y persuasivo
  \item Personal y emotivo
\end{multicols}
\end{enumerate}
\vspace{0.15cm}
\item ¿Qué es la voz en la redacción académica?
\begin{enumerate}[label=\textbf{\Alph*)}]
\begin{multicols}{2}
  \item La relación entre el sujeto y el verbo
  \item La conclusión
  \item La introducción
  \item La metodología
  \item Los resultados
  \item Las recomendaciones
\end{multicols}
\end{enumerate}
\vspace{0.15cm}
\item ¿Qué voz se prefiere en la redacción académica?
\begin{enumerate}[label=\textbf{\Alph*)}]
\begin{multicols}{2}
  \item Pasiva
  \item Media
  \item Reflexiva
  \item Impersonal
  \item Recíproca
\end{multicols}
\end{enumerate}
\vspace{0.15cm}
\item ¿Qué es la persona gramatical en la redacción académica?
\begin{enumerate}[label=\textbf{\Alph*)}]
\begin{multicols}{2}
  \item La conclusión
  \item La introducción
  \item La metodología
  \item Los resultados
  \item Las recomendaciones
\end{multicols}
\end{enumerate}
\vspace{0.15cm}
\item ¿Qué persona gramatical se prefiere en la redacción académica?
\begin{enumerate}[label=\textbf{\Alph*)}]
\begin{multicols}{2}
  \item Tercera persona
  \item Primera persona
  \item Segunda persona
  \item Primera persona del plural
  \item Segunda persona del plural
  \item Impersonal
\end{multicols}
\end{enumerate}
\vspace{0.15cm}
\item ¿Qué es el tiempo verbal en la redacción académica?
\begin{enumerate}[label=\textbf{\Alph*)}]
\begin{multicols}{2}
  \item La conclusión
  \item La introducción
  \item La metodología
  \item Los resultados
  \item Las recomendaciones
\end{multicols}
\end{enumerate}
\vspace{0.15cm}
\item ¿Qué tiempo verbal se prefiere en la redacción académica?
\begin{enumerate}[label=\textbf{\Alph*)}]
\begin{multicols}{2}
  \item Presente
  \item Pasado
  \item Imperativo
  \item Subjuntivo
\end{multicols}
\end{enumerate}
\vspace{0.15cm}
\item ¿Qué es la cita en la redacción académica?
\begin{enumerate}[label=\textbf{\Alph*)}]
\begin{multicols}{2}
  \item La mención de la fuente dentro del texto
  \item La conclusión
  \item La introducción
  \item La metodología
  \item Los resultados
  \item Las recomendaciones
\end{multicols}
\end{enumerate}
\vspace{0.15cm}
\item ¿Qué es la paráfrasis en la redacción académica?
\begin{enumerate}[label=\textbf{\Alph*)}]
\begin{multicols}{2}
  \item La expresión de las ideas de otro con palabras propias
  \item La copia exacta de las palabras de otro
  \item La opinión personal
  \item La conclusión
  \item La introducción
  \item La metodología
\end{multicols}
\end{enumerate}
\vspace{0.15cm}
\item ¿Qué es el resumen en la redacción académica?
\begin{enumerate}[label=\textbf{\Alph*)}]
\begin{multicols}{2}
  \item La síntesis breve de las ideas principales
  \item La copia exacta de un texto
  \item La opinión personal
  \item La crítica
  \item La introducción
  \item La conclusión
\end{multicols}
\end{enumerate}
\vspace{0.15cm}
\item ¿Qué es la síntesis en la redacción académica?
\begin{enumerate}[label=\textbf{\Alph*)}]
\begin{multicols}{2}
  \item La combinación de ideas de varias fuentes
  \item La copia exacta de un texto
  \item La opinión personal
  \item La crítica
  \item La introducción
  \item La conclusión
\end{multicols}
\end{enumerate}
\vspace{0.15cm}
\item ¿Qué es el análisis en la redacción académica?
\begin{enumerate}[label=\textbf{\Alph*)}]
\begin{multicols}{2}
  \item La descomposición de un todo en sus partes
  \item La copia exacta de un texto
  \item La opinión personal
  \item La crítica
  \item La introducción
  \item La conclusión
\end{multicols}
\end{enumerate}
\vspace{0.15cm}
\item ¿Qué es la interpretación en la redacción académica?
\begin{enumerate}[label=\textbf{\Alph*)}]
\begin{multicols}{2}
  \item La explicación del significado de algo
  \item La copia exacta de un texto
  \item La opinión personal
  \item La crítica
  \item La introducción
  \item La conclusión
\end{multicols}
\end{enumerate}
\vspace{0.15cm}
\item ¿Qué es la argumentación en la redacción académica?
\begin{enumerate}[label=\textbf{\Alph*)}]
\begin{multicols}{2}
  \item La defensa de una idea con razones
  \item La copia exacta de un texto
  \item La opinión personal
  \item La crítica
  \item La introducción
  \item La conclusión
\end{multicols}
\end{enumerate}
\vspace{0.15cm}
\item ¿Qué es la exposición en la redacción académica?
\begin{enumerate}[label=\textbf{\Alph*)}]
\begin{multicols}{2}
  \item La presentación de información de manera objetiva
  \item La copia exacta de un texto
  \item La opinión personal
  \item La crítica
  \item La introducción
  \item La conclusión
\end{multicols}
\end{enumerate}
\vspace{0.15cm}
\item ¿Qué es la descripción en la redacción académica?
\begin{enumerate}[label=\textbf{\Alph*)}]
\begin{multicols}{2}
  \item La representación de características de algo
  \item La copia exacta de un texto
  \item La opinión personal
  \item La crítica
  \item La introducción
  \item La conclusión
\end{multicols}
\end{enumerate}
\vspace{0.15cm}
\item ¿Qué es la narración en la redacción académica?
\begin{enumerate}[label=\textbf{\Alph*)}]
\begin{multicols}{2}
  \item La copia exacta de un texto
  \item La opinión personal
  \item La crítica
  \item La introducción
  \item La conclusión
\end{multicols}
\end{enumerate}
\vspace{0.15cm}
\item ¿Qué es la revisión en la redacción académica?
\begin{enumerate}[label=\textbf{\Alph*)}]
\begin{multicols}{2}
  \item La lectura crítica para mejorar el texto
  \item La copia exacta de un texto
  \item La opinión personal
  \item La crítica
  \item La introducción
  \item La conclusión
\end{multicols}
\end{enumerate}
\vspace{0.15cm}
\item ¿Qué es la edición en la redacción académica?
\begin{enumerate}[label=\textbf{\Alph*)}]
\begin{multicols}{2}
  \item La corrección y mejora del texto
  \item La copia exacta de un texto
  \item La opinión personal
  \item La crítica
  \item La introducción
  \item La conclusión
\end{multicols}
\end{enumerate}
\vspace{0.15cm}
\item ¿Qué es la publicación en la redacción académica?
\begin{enumerate}[label=\textbf{\Alph*)}]
\begin{multicols}{2}
  \item La difusión del texto en un medio académico
  \item La copia exacta de un texto
  \item La opinión personal
  \item La crítica
  \item La introducción
  \item La conclusión
\end{multicols}
\end{enumerate}
\vspace{0.15cm}
\item ¿Qué es el título en la redacción académica?
\begin{enumerate}[label=\textbf{\Alph*)}]
\begin{multicols}{2}
  \item La frase breve que indica el tema
  \item La copia exacta de un texto
  \item La opinión personal
  \item La crítica
  \item La introducción
  \item La conclusión
\end{multicols}
\end{enumerate}
\vspace{0.15cm}
\item ¿Qué características debe tener el título en la redacción académica?
\begin{enumerate}[label=\textbf{\Alph*)}]
\begin{multicols}{2}
  \item Ser claro, breve, atractivo y pertinente
  \item Ser confuso, largo y ambiguo
  \item Ser coloquial y emotivo
  \item Ser poético y metafórico
  \item Ser humorístico y sarcástico
  \item Ser informal y desordenado
\end{multicols}
\end{enumerate}
\vspace{0.15cm}
\item ¿Qué es el resumen en la redacción académica?
\begin{enumerate}[label=\textbf{\Alph*)}]
\begin{multicols}{2}
  \item Un párrafo breve que sintetiza el contenido
  \item La copia exacta de un texto
  \item La opinión personal
  \item La crítica
  \item La introducción
  \item La conclusión
\end{multicols}
\end{enumerate}
\vspace{0.15cm}
\item ¿Qué es la palabra clave en la redacción académica?
\begin{enumerate}[label=\textbf{\Alph*)}]
\begin{multicols}{2}
  \item Términos que identifican el contenido
  \item La primera palabra del título
  \item La última palabra del título
  \item La editorial
\end{multicols}
\end{enumerate}
\vspace{0.15cm}
\item ¿Qué es la portada en la redacción académica?
\begin{enumerate}[label=\textbf{\Alph*)}]
\begin{multicols}{2}
  \item La primera página que incluye título, autor, institución y fecha
  \item La última página
  \item La página del medio
  \item La página de referencias
  \item La página de introducción
  \item La página de conclusión
\end{multicols}
\end{enumerate}
\vspace{0.15cm}
\item ¿Qué es el índice en la redacción académica?
\begin{enumerate}[label=\textbf{\Alph*)}]
\begin{multicols}{2}
  \item La lista de contenidos y páginas
  \item La lista de referencias
  \item La introducción
  \item La conclusión
  \item La metodología
  \item Los resultados
\end{multicols}
\end{enumerate}
\vspace{0.15cm}
\item ¿Qué es el apéndice en la redacción académica?
\begin{enumerate}[label=\textbf{\Alph*)}]
\begin{multicols}{2}
  \item La sección que incluye material complementario
  \item La lista de referencias
  \item La introducción
  \item La conclusión
  \item La metodología
  \item Los resultados
\end{multicols}
\end{enumerate}
\vspace{0.15cm}
\item ¿Qué es el anexo en la redacción académica?
\begin{enumerate}[label=\textbf{\Alph*)}]
\begin{multicols}{2}
  \item La sección que incluye material adicional
  \item La lista de referencias
  \item La introducción
  \item La conclusión
  \item La metodología
  \item Los resultados
\end{multicols}
\end{enumerate}
\vspace{0.15cm}
\item ¿Qué es la nota al pie en la redacción académica?
\begin{enumerate}[label=\textbf{\Alph*)}]
\begin{multicols}{2}
  \item Una aclaración o comentario al final de la página
  \item Una cita en el texto
  \item Una referencia al final del trabajo
  \item Una introducción
  \item Una conclusión
  \item Una metodología
\end{multicols}
\end{enumerate}
\vspace{0.15cm}
\item ¿Qué es la nota final en la redacción académica?
\begin{enumerate}[label=\textbf{\Alph*)}]
\begin{multicols}{2}
  \item Una aclaración o comentario al final del documento
  \item Una cita en el texto
  \item Una referencia al final del trabajo
  \item Una introducción
  \item Una conclusión
  \item Una metodología
\end{multicols}
\end{enumerate}
\vspace{0.15cm}
\item ¿Qué es la tabla en la redacción académica?
\begin{enumerate}[label=\textbf{\Alph*)}]
\begin{multicols}{2}
  \item La presentación de datos en filas y columnas
  \item La conclusión
  \item La introducción
  \item La metodología
  \item Los resultados
  \item Las recomendaciones
\end{multicols}
\end{enumerate}
\vspace{0.15cm}
\item ¿Qué es la figura en la redacción académica?
\begin{enumerate}[label=\textbf{\Alph*)}]
\begin{multicols}{2}
  \item La representación visual de información
  \item La conclusión
  \item La introducción
  \item La metodología
  \item Los resultados
  \item Las recomendaciones
\end{multicols}
\end{enumerate}
\vspace{0.15cm}
\item ¿Qué es el gráfico en la redacción académica?
\begin{enumerate}[label=\textbf{\Alph*)}]
\begin{multicols}{2}
  \item La representación visual de datos
  \item La conclusión
  \item La introducción
  \item La metodología
  \item Los resultados
  \item Las recomendaciones
\end{multicols}
\end{enumerate}
\vspace{0.15cm}
\item ¿Qué es la ecuación en la redacción académica?
\begin{enumerate}[label=\textbf{\Alph*)}]
\begin{multicols}{2}
  \item La expresión matemática de una relación
  \item La conclusión
  \item La introducción
  \item La metodología
  \item Los resultados
  \item Las recomendaciones
\end{multicols}
\end{enumerate}
\vspace{0.15cm}
\item ¿Qué es la fórmula en la redacción académica?
\begin{enumerate}[label=\textbf{\Alph*)}]
\begin{multicols}{2}
  \item La expresión matemática de una regla o principio
  \item La conclusión
  \item La introducción
  \item La metodología
  \item Los resultados
  \item Las recomendaciones
\end{multicols}
\end{enumerate}
\vspace{0.15cm}
\item ¿Qué es el símbolo en la redacción académica?
\begin{enumerate}[label=\textbf{\Alph*)}]
\begin{multicols}{2}
  \item La representación convencional de un concepto
  \item La conclusión
  \item La introducción
  \item La metodología
  \item Los resultados
  \item Las recomendaciones
\end{multicols}
\end{enumerate}
\vspace{0.15cm}
\item ¿Qué es la unidad en la redacción académica?
\begin{enumerate}[label=\textbf{\Alph*)}]
\begin{multicols}{2}
  \item La medida estándar de una magnitud
  \item La conclusión
  \item La introducción
  \item La metodología
  \item Los resultados
  \item Las recomendaciones ---
\end{multicols}
\end{enumerate}
\vspace{0.15cm}
\end{enumerate}\n
\subsubsection*{SECCIÓN V: METODOLOGÍA DE LA INVESTIGACIÓN (221–270)}

\begin{enumerate}[label=\textbf{\arabic*.}]
\setcounter{enumi}{220}
\item ¿Qué es la metodología de la investigación?
\begin{enumerate}[label=\textbf{\Alph*)}]
\begin{multicols}{2}
  \item La conclusión
  \item La introducción
  \item Los resultados
  \item Las recomendaciones
\end{multicols}
\end{enumerate}
\vspace{0.15cm}
\item ¿Qué es el método científico?
\begin{enumerate}[label=\textbf{\Alph*)}]
\begin{multicols}{2}
  \item La conclusión
  \item La introducción
  \item Los resultados
  \item Las recomendaciones
\end{multicols}
\end{enumerate}
\vspace{0.15cm}
\item ¿Qué es el método deductivo?
\begin{enumerate}[label=\textbf{\Alph*)}]
\begin{multicols}{2}
\end{multicols}
\end{enumerate}
\vspace{0.15cm}
\item ¿Qué es el método inductivo?
\begin{enumerate}[label=\textbf{\Alph*)}]
\begin{multicols}{2}
\end{multicols}
\end{enumerate}
\vspace{0.15cm}
\item ¿Qué es el método analítico?
\begin{enumerate}[label=\textbf{\Alph*)}]
\begin{multicols}{2}
\end{multicols}
\end{enumerate}
\vspace{0.15cm}
\item ¿Qué es el método sintético?
\begin{enumerate}[label=\textbf{\Alph*)}]
\begin{multicols}{2}
\end{multicols}
\end{enumerate}
\vspace{0.15cm}
\item ¿Qué es el método comparativo?
\begin{enumerate}[label=\textbf{\Alph*)}]
\begin{multicols}{2}
\end{multicols}
\end{enumerate}
\vspace{0.15cm}
\item ¿Qué es el método experimental?
\begin{enumerate}[label=\textbf{\Alph*)}]
\begin{multicols}{2}
\end{multicols}
\end{enumerate}
\vspace{0.15cm}
\item ¿Qué es el método observacional?
\begin{enumerate}[label=\textbf{\Alph*)}]
\begin{multicols}{2}
\end{multicols}
\end{enumerate}
\vspace{0.15cm}
\item ¿Qué es la investigación cuantitativa?
\begin{enumerate}[label=\textbf{\Alph*)}]
\begin{multicols}{2}
  \item La que se basa en datos numéricos y estadísticos
  \item La que se basa en datos no numéricos
  \item La que no tiene dirección
  \item La que es aleatoria
  \item La que es confusa
  \item La que es ambigua
\end{multicols}
\end{enumerate}
\vspace{0.15cm}
\item ¿Qué es la investigación cualitativa?
\begin{enumerate}[label=\textbf{\Alph*)}]
\begin{multicols}{2}
  \item La que se basa en datos no numéricos y descriptivos
  \item La que se basa en datos numéricos
  \item La que no tiene dirección
  \item La que es aleatoria
  \item La que es confusa
  \item La que es ambigua
\end{multicols}
\end{enumerate}
\vspace{0.15cm}
\item ¿Qué es la investigación mixta?
\begin{enumerate}[label=\textbf{\Alph*)}]
\begin{multicols}{2}
  \item La que combina métodos cuantitativos y cualitativos
  \item La que solo usa métodos cuantitativos
  \item La que solo usa métodos cualitativos
  \item La que no tiene dirección
  \item La que es aleatoria
  \item La que es confusa
\end{multicols}
\end{enumerate}
\vspace{0.15cm}
\item ¿Qué es la investigación exploratoria?
\begin{enumerate}[label=\textbf{\Alph*)}]
\begin{multicols}{2}
  \item La que explora un tema poco estudiado
  \item La que describe un tema
  \item La que explica un tema
  \item La que correlaciona variables
  \item La que experimenta
  \item La que compara
\end{multicols}
\end{enumerate}
\vspace{0.15cm}
\item ¿Qué es la investigación descriptiva?
\begin{enumerate}[label=\textbf{\Alph*)}]
\begin{multicols}{2}
  \item La que describe características de un fenómeno
  \item La que explora un tema
  \item La que explica un tema
  \item La que correlaciona variables
  \item La que experimenta
  \item La que compara
\end{multicols}
\end{enumerate}
\vspace{0.15cm}
\item ¿Qué es la investigación explicativa?
\begin{enumerate}[label=\textbf{\Alph*)}]
\begin{multicols}{2}
  \item La que explica causas y efectos de un fenómeno
  \item La que explora un tema
  \item La que describe un tema
  \item La que correlaciona variables
  \item La que experimenta
  \item La que compara
\end{multicols}
\end{enumerate}
\vspace{0.15cm}
\item ¿Qué es la investigación correlacional?
\begin{enumerate}[label=\textbf{\Alph*)}]
\begin{multicols}{2}
  \item La que establece relaciones entre variables
  \item La que explora un tema
  \item La que describe un tema
  \item La que explica un tema
  \item La que experimenta
  \item La que compara
\end{multicols}
\end{enumerate}
\vspace{0.15cm}
\item ¿Qué es la investigación experimental?
\begin{enumerate}[label=\textbf{\Alph*)}]
\begin{multicols}{2}
  \item La que manipula variables para establecer causas
  \item La que explora un tema
  \item La que describe un tema
  \item La que explica un tema
  \item La que correlaciona variables
  \item La que compara
\end{multicols}
\end{enumerate}
\vspace{0.15cm}
\item ¿Qué es la investigación no experimental?
\begin{enumerate}[label=\textbf{\Alph*)}]
\begin{multicols}{2}
  \item La que observa sin manipular variables
  \item La que manipula variables
  \item La que explora un tema
  \item La que describe un tema
  \item La que explica un tema
  \item La que correlaciona variables
\end{multicols}
\end{enumerate}
\vspace{0.15cm}
\item ¿Qué es la investigación transversal?
\begin{enumerate}[label=\textbf{\Alph*)}]
\begin{multicols}{2}
  \item La que recoge datos en un solo momento
  \item La que recoge datos en varios momentos
  \item La que explora un tema
  \item La que describe un tema
  \item La que explica un tema
  \item La que correlaciona variables
\end{multicols}
\end{enumerate}
\vspace{0.15cm}
\item ¿Qué es la investigación longitudinal?
\begin{enumerate}[label=\textbf{\Alph*)}]
\begin{multicols}{2}
  \item La que recoge datos en varios momentos
  \item La que recoge datos en un solo momento
  \item La que explora un tema
  \item La que describe un tema
  \item La que explica un tema
  \item La que correlaciona variables
\end{multicols}
\end{enumerate}
\vspace{0.15cm}
\item ¿Qué es la técnica de investigación?
\begin{enumerate}[label=\textbf{\Alph*)}]
\begin{multicols}{2}
  \item La conclusión
  \item La introducción
  \item Los resultados
\end{multicols}
\end{enumerate}
\vspace{0.15cm}
\item ¿Qué es la encuesta?
\begin{enumerate}[label=\textbf{\Alph*)}]
\begin{multicols}{2}
  \item Una técnica de recolección de datos mediante preguntas
  \item Una técnica de análisis
  \item Una técnica de experimentación
  \item Una técnica de observación
  \item Una técnica de entrevista
  \item Una técnica de documentación
\end{multicols}
\end{enumerate}
\vspace{0.15cm}
\item ¿Qué es la entrevista?
\begin{enumerate}[label=\textbf{\Alph*)}]
\begin{multicols}{2}
  \item Una técnica de recolección de datos mediante diálogo
  \item Una técnica de análisis
  \item Una técnica de experimentación
  \item Una técnica de observación
  \item Una técnica de encuesta
  \item Una técnica de documentación
\end{multicols}
\end{enumerate}
\vspace{0.15cm}
\item ¿Qué es la observación?
\begin{enumerate}[label=\textbf{\Alph*)}]
\begin{multicols}{2}
  \item Una técnica de recolección de datos mediante la percepción directa
  \item Una técnica de análisis
  \item Una técnica de experimentación
  \item Una técnica de encuesta
  \item Una técnica de entrevista
  \item Una técnica de documentación
\end{multicols}
\end{enumerate}
\vspace{0.15cm}
\item ¿Qué es el experimento?
\begin{enumerate}[label=\textbf{\Alph*)}]
\begin{multicols}{2}
  \item Una técnica de recolección de datos mediante la manipulación de variables
  \item Una técnica de análisis
  \item Una técnica de observación
  \item Una técnica de encuesta
  \item Una técnica de entrevista
  \item Una técnica de documentación
\end{multicols}
\end{enumerate}
\vspace{0.15cm}
\item ¿Qué es el cuestionario?
\begin{enumerate}[label=\textbf{\Alph*)}]
\begin{multicols}{2}
  \item Un instrumento de recolección de datos con preguntas escritas
  \item Una técnica de análisis
  \item Una técnica de experimentación
  \item Una técnica de observación
  \item Una técnica de entrevista
  \item Una técnica de documentación
\end{multicols}
\end{enumerate}
\vspace{0.15cm}
\item ¿Qué es la escala de Likert?
\begin{enumerate}[label=\textbf{\Alph*)}]
\begin{multicols}{2}
  \item Un instrumento para medir actitudes y opiniones
  \item Una técnica de análisis
  \item Una técnica de experimentación
  \item Una técnica de observación
  \item Una técnica de entrevista
  \item Una técnica de documentación
\end{multicols}
\end{enumerate}
\vspace{0.15cm}
\item ¿Qué es la validez de un instrumento?
\begin{enumerate}[label=\textbf{\Alph*)}]
\begin{multicols}{2}
  \item La cantidad de datos
  \item La extensión del informe
\end{multicols}
\end{enumerate}
\vspace{0.15cm}
\item ¿Qué es la confiabilidad de un instrumento?
\begin{enumerate}[label=\textbf{\Alph*)}]
\begin{multicols}{2}
  \item La cantidad de datos
  \item La extensión del informe
\end{multicols}
\end{enumerate}
\vspace{0.15cm}
\item ¿Qué es la triangulación?
\begin{enumerate}[label=\textbf{\Alph*)}]
\begin{multicols}{2}
\end{multicols}
\end{enumerate}
\vspace{0.15cm}
\item ¿Qué es la población en investigación?
\begin{enumerate}[label=\textbf{\Alph*)}]
\begin{multicols}{2}
  \item Una parte del conjunto
  \item Un individuo
  \item Un dato
  \item Una variable
  \item Una hipótesis
\end{multicols}
\end{enumerate}
\vspace{0.15cm}
\item ¿Qué es la muestra en investigación?
\begin{enumerate}[label=\textbf{\Alph*)}]
\begin{multicols}{2}
  \item Un subconjunto representativo de la población
  \item Un individuo
  \item Un dato
  \item Una variable
  \item Una hipótesis
\end{multicols}
\end{enumerate}
\vspace{0.15cm}
\item ¿Qué es el muestreo probabilístico?
\begin{enumerate}[label=\textbf{\Alph*)}]
\begin{multicols}{2}
\end{multicols}
\end{enumerate}
\vspace{0.15cm}
\item ¿Qué es el muestreo no probabilístico?
\begin{enumerate}[label=\textbf{\Alph*)}]
\begin{multicols}{2}
\end{multicols}
\end{enumerate}
\vspace{0.15cm}
\item ¿Qué es el análisis estadístico?
\begin{enumerate}[label=\textbf{\Alph*)}]
\begin{multicols}{2}
  \item La conclusión
  \item La introducción
  \item Los resultados
\end{multicols}
\end{enumerate}
\vspace{0.15cm}
\item ¿Qué es la media aritmética?
\begin{enumerate}[label=\textbf{\Alph*)}]
\begin{multicols}{2}
  \item La diferencia entre el máximo y el mínimo
  \item La desviación estándar
  \item La varianza
\end{multicols}
\end{enumerate}
\vspace{0.15cm}
\item ¿Qué es la mediana?
\begin{enumerate}[label=\textbf{\Alph*)}]
\begin{multicols}{2}
  \item La diferencia entre el máximo y el mínimo
  \item La desviación estándar
  \item La varianza
\end{multicols}
\end{enumerate}
\vspace{0.15cm}
\item ¿Qué es la moda?
\begin{enumerate}[label=\textbf{\Alph*)}]
\begin{multicols}{2}
  \item La diferencia entre el máximo y el mínimo
  \item La desviación estándar
  \item La varianza
\end{multicols}
\end{enumerate}
\vspace{0.15cm}
\item ¿Qué es el rango?
\begin{enumerate}[label=\textbf{\Alph*)}]
\begin{multicols}{2}
  \item La diferencia entre el máximo y el mínimo
  \item La desviación estándar
  \item La varianza
\end{multicols}
\end{enumerate}
\vspace{0.15cm}
\item ¿Qué es la desviación estándar?
\begin{enumerate}[label=\textbf{\Alph*)}]
\begin{multicols}{2}
  \item La medida de dispersión de los datos respecto a la media
  \item La diferencia entre el máximo y el mínimo
  \item La varianza
\end{multicols}
\end{enumerate}
\vspace{0.15cm}
\item ¿Qué es la varianza?
\begin{enumerate}[label=\textbf{\Alph*)}]
\begin{multicols}{2}
  \item La diferencia entre el máximo y el mínimo
  \item La desviación estándar
\end{multicols}
\end{enumerate}
\vspace{0.15cm}
\item ¿Qué es una tabla de frecuencia?
\begin{enumerate}[label=\textbf{\Alph*)}]
\begin{multicols}{2}
  \item La presentación de datos con sus frecuencias
  \item La conclusión
  \item La introducción
  \item Los resultados
  \item Las recomendaciones
\end{multicols}
\end{enumerate}
\vspace{0.15cm}
\item ¿Qué es un gráfico de barras?
\begin{enumerate}[label=\textbf{\Alph*)}]
\begin{multicols}{2}
  \item La representación visual de datos con barras
  \item La representación visual de datos con líneas
  \item La representación visual de datos con sectores
  \item La representación visual de datos con puntos
  \item La representación visual de datos con áreas
  \item La representación visual de datos con histogramas
\end{multicols}
\end{enumerate}
\vspace{0.15cm}
\item ¿Qué es un gráfico de líneas?
\begin{enumerate}[label=\textbf{\Alph*)}]
\begin{multicols}{2}
  \item La representación visual de datos con líneas
  \item La representación visual de datos con barras
  \item La representación visual de datos con sectores
  \item La representación visual de datos con puntos
  \item La representación visual de datos con áreas
  \item La representación visual de datos con histogramas
\end{multicols}
\end{enumerate}
\vspace{0.15cm}
\item ¿Qué es un gráfico circular?
\begin{enumerate}[label=\textbf{\Alph*)}]
\begin{multicols}{2}
  \item La representación visual de datos con sectores
  \item La representación visual de datos con barras
  \item La representación visual de datos con líneas
  \item La representación visual de datos con puntos
  \item La representación visual de datos con áreas
  \item La representación visual de datos con histogramas
\end{multicols}
\end{enumerate}
\vspace{0.15cm}
\item ¿Qué es un histograma?
\begin{enumerate}[label=\textbf{\Alph*)}]
\begin{multicols}{2}
  \item La representación visual de datos con barras contiguas
  \item La representación visual de datos con barras separadas
  \item La representación visual de datos con líneas
  \item La representación visual de datos con sectores
  \item La representación visual de datos con puntos
  \item La representación visual de datos con áreas
\end{multicols}
\end{enumerate}
\vspace{0.15cm}
\item ¿Qué es un polígono de frecuencia?
\begin{enumerate}[label=\textbf{\Alph*)}]
\begin{multicols}{2}
  \item La representación visual de datos con líneas que unen puntos
  \item La representación visual de datos con barras
  \item La representación visual de datos con sectores
  \item La representación visual de datos con áreas
  \item La representación visual de datos con histogramas
  \item La representación visual de datos con puntos
\end{multicols}
\end{enumerate}
\vspace{0.15cm}
\item ¿Qué es una ojiva?
\begin{enumerate}[label=\textbf{\Alph*)}]
\begin{multicols}{2}
  \item La representación visual de frecuencias acumuladas
  \item La representación visual de frecuencias simples
  \item La representación visual de datos con barras
  \item La representación visual de datos con sectores
  \item La representación visual de datos con áreas
  \item La representación visual de datos con puntos
\end{multicols}
\end{enumerate}
\vspace{0.15cm}
\item ¿Qué es la correlación?
\begin{enumerate}[label=\textbf{\Alph*)}]
\begin{multicols}{2}
  \item La relación entre dos o más variables
  \item La conclusión
  \item La introducción
  \item Los resultados
  \item Las recomendaciones
\end{multicols}
\end{enumerate}
\vspace{0.15cm}
\item ¿Qué es la regresión?
\begin{enumerate}[label=\textbf{\Alph*)}]
\begin{multicols}{2}
  \item La conclusión
  \item La introducción
  \item Los resultados
  \item Las recomendaciones ---
\end{multicols}
\end{enumerate}
\vspace{0.15cm}
\end{enumerate}\n
\subsubsection*{SECCIÓN VI: ÉTICA Y RIGOR ACADÉMICO (271–300)}

\begin{enumerate}[label=\textbf{\arabic*.}]
\setcounter{enumi}{270}
\item ¿Qué es la ética académica?
\begin{enumerate}[label=\textbf{\Alph*)}]
\begin{multicols}{2}
  \item La conclusión
  \item La introducción
  \item La metodología
  \item Los resultados
  \item Las recomendaciones
\end{multicols}
\end{enumerate}
\vspace{0.15cm}
\item ¿Qué es la integridad académica?
\begin{enumerate}[label=\textbf{\Alph*)}]
\begin{multicols}{2}
  \item La honestidad y responsabilidad en el trabajo académico
  \item La conclusión
  \item La introducción
  \item La metodología
  \item Los resultados
  \item Las recomendaciones
\end{multicols}
\end{enumerate}
\vspace{0.15cm}
\item ¿Qué es la honestidad académica?
\begin{enumerate}[label=\textbf{\Alph*)}]
\begin{multicols}{2}
  \item La presentación veraz de las ideas y fuentes
  \item La conclusión
  \item La introducción
  \item La metodología
  \item Los resultados
  \item Las recomendaciones
\end{multicols}
\end{enumerate}
\vspace{0.15cm}
\item ¿Qué es la responsabilidad académica?
\begin{enumerate}[label=\textbf{\Alph*)}]
\begin{multicols}{2}
  \item La conclusión
  \item La introducción
  \item La metodología
  \item Los resultados
  \item Las recomendaciones
\end{multicols}
\end{enumerate}
\vspace{0.15cm}
\item ¿Qué es el respeto académico?
\begin{enumerate}[label=\textbf{\Alph*)}]
\begin{multicols}{2}
  \item La consideración hacia las ideas y personas en el ámbito académico
  \item La conclusión
  \item La introducción
  \item La metodología
  \item Los resultados
  \item Las recomendaciones
\end{multicols}
\end{enumerate}
\vspace{0.15cm}
\item ¿Qué es la colaboración académica?
\begin{enumerate}[label=\textbf{\Alph*)}]
\begin{multicols}{2}
  \item La conclusión
  \item La introducción
  \item La metodología
  \item Los resultados
  \item Las recomendaciones
\end{multicols}
\end{enumerate}
\vspace{0.15cm}
\item ¿Qué es la evaluación académica?
\begin{enumerate}[label=\textbf{\Alph*)}]
\begin{multicols}{2}
  \item La conclusión
  \item La introducción
  \item La metodología
  \item Los resultados
  \item Las recomendaciones
\end{multicols}
\end{enumerate}
\vspace{0.15cm}
\item ¿Qué es la retroalimentación académica?
\begin{enumerate}[label=\textbf{\Alph*)}]
\begin{multicols}{2}
  \item La información que se da al estudiante para mejorar su trabajo
  \item La conclusión
  \item La introducción
  \item La metodología
  \item Los resultados
  \item Las recomendaciones
\end{multicols}
\end{enumerate}
\vspace{0.15cm}
\item ¿Qué es la autoevaluación académica?
\begin{enumerate}[label=\textbf{\Alph*)}]
\begin{multicols}{2}
  \item La valoración que el estudiante hace de su propio trabajo
  \item La conclusión
  \item La introducción
  \item La metodología
  \item Los resultados
  \item Las recomendaciones
\end{multicols}
\end{enumerate}
\vspace{0.15cm}
\item ¿Qué es la coevaluación académica?
\begin{enumerate}[label=\textbf{\Alph*)}]
\begin{multicols}{2}
  \item La valoración entre pares del trabajo académico
  \item La conclusión
  \item La introducción
  \item La metodología
  \item Los resultados
  \item Las recomendaciones
\end{multicols}
\end{enumerate}
\vspace{0.15cm}
\item ¿Qué es el plagio?
\begin{enumerate}[label=\textbf{\Alph*)}]
\begin{multicols}{2}
  \item Presentar como propias las ideas o palabras de otro
  \item Parafrasear correctamente
  \item Resumir correctamente
\end{multicols}
\end{enumerate}
\vspace{0.15cm}
\item ¿Qué tipo de plagio es copiar y pegar sin citar?
\begin{enumerate}[label=\textbf{\Alph*)}]
\begin{multicols}{2}
  \item Plagio directo o textual
  \item Plagio parafraseado
  \item Plagio accidental
  \item Plagio de ideas
  \item Plagio de estructura
\end{multicols}
\end{enumerate}
\vspace{0.15cm}
\item ¿Qué tipo de plagio es cambiar algunas palabras sin citar?
\begin{enumerate}[label=\textbf{\Alph*)}]
\begin{multicols}{2}
  \item Plagio parafraseado
  \item Plagio directo
  \item Plagio accidental
  \item Plagio de ideas
  \item Plagio de estructura
\end{multicols}
\end{enumerate}
\vspace{0.15cm}
\item ¿Qué tipo de plagio es reutilizar un trabajo propio sin indicarlo?
\begin{enumerate}[label=\textbf{\Alph*)}]
\begin{multicols}{2}
  \item Plagio directo
  \item Plagio parafraseado
  \item Plagio accidental
  \item Plagio de ideas
  \item Plagio de estructura
\end{multicols}
\end{enumerate}
\vspace{0.15cm}
\item ¿Qué tipo de plagio es usar las ideas de otro sin citar?
\begin{enumerate}[label=\textbf{\Alph*)}]
\begin{multicols}{2}
  \item Plagio de ideas
  \item Plagio directo
  \item Plagio parafraseado
  \item Plagio accidental
  \item Plagio de estructura
\end{multicols}
\end{enumerate}
\vspace{0.15cm}
\item ¿Cómo se evita el plagio?
\begin{enumerate}[label=\textbf{\Alph*)}]
\begin{multicols}{2}
  \item Parafraseando sin citar
  \item Usando ideas ajenas sin citar
  \item Reutilizando trabajos propios sin indicarlo
  \item Ignorando las fuentes
\end{multicols}
\end{enumerate}
\vspace{0.15cm}
\item ¿Qué es una paráfrasis correcta?
\begin{enumerate}[label=\textbf{\Alph*)}]
\begin{multicols}{2}
  \item Usar las ideas de otro sin citar
  \item Reutilizar un trabajo propio sin indicarlo
  \item Ignorar la fuente
\end{multicols}
\end{enumerate}
\vspace{0.15cm}
\item ¿Qué es un resumen en el contexto académico?
\begin{enumerate}[label=\textbf{\Alph*)}]
\begin{multicols}{2}
  \item La síntesis breve de las ideas principales de un texto
  \item La copia exacta de un texto
  \item La opinión personal sobre un texto
  \item La crítica de un texto
  \item La introducción de un texto
  \item La conclusión de un texto
\end{multicols}
\end{enumerate}
\vspace{0.15cm}
\item ¿Qué es una cita textual?
\begin{enumerate}[label=\textbf{\Alph*)}]
\begin{multicols}{2}
  \item La reproducción exacta de las palabras de otro autor
  \item Un resumen de las ideas de otro autor
  \item Una paráfrasis de las ideas de otro autor
  \item Una opinión personal
  \item Una conclusión
  \item Una introducción
\end{multicols}
\end{enumerate}
\vspace{0.15cm}
\item ¿Qué es una cita parafraseada?
\begin{enumerate}[label=\textbf{\Alph*)}]
\begin{multicols}{2}
  \item La expresión de las ideas de otro autor con palabras propias
  \item La reproducción exacta de las palabras de otro autor
  \item Una opinión personal
  \item Una conclusión
  \item Una introducción
  \item Una metodología
\end{multicols}
\end{enumerate}
\vspace{0.15cm}
\item ¿Qué es una cita directa corta?
\begin{enumerate}[label=\textbf{\Alph*)}]
\begin{multicols}{2}
  \item La que tiene menos de 40 palabras y se integra en el texto entre comillas
  \item La que tiene más de 40 palabras
  \item La que se escribe en párrafo aparte
  \item La que no lleva comillas
  \item La que no lleva autor
  \item La que no lleva año
\end{multicols}
\end{enumerate}
\vspace{0.15cm}
\item ¿Qué es una cita directa larga?
\begin{enumerate}[label=\textbf{\Alph*)}]
\begin{multicols}{2}
  \item La que tiene 40 palabras o más y se escribe en párrafo aparte con sangría
  \item La que tiene menos de 40 palabras
  \item La que se integra en el texto entre comillas
  \item La que no lleva comillas
  \item La que no lleva autor
  \item La que no lleva año
\end{multicols}
\end{enumerate}
\vspace{0.15cm}
\item ¿Cómo se cita textualmente a un autor en APA?
\begin{enumerate}[label=\textbf{\Alph*)}]
\begin{multicols}{2}
\end{multicols}
\end{enumerate}
\vspace{0.15cm}
\item ¿Cómo se cita parafraseando a un autor en APA?
\begin{enumerate}[label=\textbf{\Alph*)}]
\begin{multicols}{2}
\end{multicols}
\end{enumerate}
\vspace{0.15cm}
\item ¿Qué es una cita de cita?
\begin{enumerate}[label=\textbf{\Alph*)}]
\begin{multicols}{2}
  \item La que se hace de una fuente citada por otro autor
  \item La que se hace directamente de la fuente original
  \item La que no lleva autor
  \item La que no lleva año
  \item La que no lleva página
  \item La que no lleva comillas
\end{multicols}
\end{enumerate}
\vspace{0.15cm}
\item ¿Cómo se cita una fuente con dos autores en APA?
\begin{enumerate}[label=\textbf{\Alph*)}]
\begin{multicols}{2}
\end{multicols}
\end{enumerate}
\vspace{0.15cm}
\item ¿Cómo se cita una fuente con tres o más autores en APA?
\begin{enumerate}[label=\textbf{\Alph*)}]
\begin{multicols}{2}
\end{multicols}
\end{enumerate}
\vspace{0.15cm}
\item ¿Cómo se cita una fuente institucional en APA?
\begin{enumerate}[label=\textbf{\Alph*)}]
\begin{multicols}{2}
\end{multicols}
\end{enumerate}
\vspace{0.15cm}
\item ¿Cómo se cita una fuente sin autor en APA?
\begin{enumerate}[label=\textbf{\Alph*)}]
\begin{multicols}{2}
\end{multicols}
\end{enumerate}
\vspace{0.15cm}
\item ¿Cómo se cita una fuente sin fecha en APA?
\begin{enumerate}[label=\textbf{\Alph*)}]
\begin{multicols}{2}
\end{multicols}
\end{enumerate}
\vspace{0.15cm}
\end{enumerate}\n
\vspace{0.4cm}
\begin{tcolorbox}[colback=gray!5!white,colframe=gray!60!black,title=\textbf{SOLUCIONARIO -- 2.4 METODOLOGÍA Y ESCRITURA ACADÉMICA}]
\begin{multicols}{4}
\begin{enumerate}[label=\textbf{\arabic*.}]
  \item \textbf{A}
  \item \textbf{A}
  \item \textbf{A}
  \item \textbf{A}
  \item \textbf{A}
  \item \textbf{A}
  \item \textbf{A}
  \item \textbf{A}
  \item \textbf{A}
  \item \textbf{A}
  \item \textbf{A}
  \item \textbf{A}
  \item \textbf{A}
  \item \textbf{A}
  \item \textbf{A}
  \item \textbf{A}
  \item \textbf{A}
  \item \textbf{A}
  \item \textbf{A}
  \item \textbf{A}
  \item \textbf{A}
  \item \textbf{A}
  \item \textbf{A}
  \item \textbf{A}
  \item \textbf{A}
  \item \textbf{A}
  \item \textbf{A}
  \item \textbf{A}
  \item \textbf{A}
  \item \textbf{A}
  \item \textbf{A}
  \item \textbf{A}
  \item \textbf{A}
  \item \textbf{A}
  \item \textbf{A}
  \item \textbf{A}
  \item \textbf{A}
  \item \textbf{A}
  \item \textbf{A}
  \item \textbf{A}
  \item \textbf{A}
  \item \textbf{A}
  \item \textbf{A}
  \item \textbf{A}
  \item \textbf{A}
  \item \textbf{A}
  \item \textbf{A}
  \item \textbf{A}
  \item \textbf{A}
  \item \textbf{A}
  \item \textbf{A}
  \item \textbf{A}
  \item \textbf{A}
  \item \textbf{A}
  \item \textbf{A}
  \item \textbf{A}
  \item \textbf{A}
  \item \textbf{A}
  \item \textbf{A}
  \item \textbf{A}
  \item \textbf{A}
  \item \textbf{A}
  \item \textbf{A}
  \item \textbf{A}
  \item \textbf{A}
  \item \textbf{A}
  \item \textbf{A}
  \item \textbf{A}
  \item \textbf{A}
  \item \textbf{A}
  \item \textbf{A}
  \item \textbf{A}
  \item \textbf{A}
  \item \textbf{A}
  \item \textbf{A}
  \item \textbf{A}
  \item \textbf{A}
  \item \textbf{A}
  \item \textbf{A}
  \item \textbf{A}
  \item \textbf{A}
  \item \textbf{A}
  \item \textbf{A}
  \item \textbf{A}
  \item \textbf{A}
  \item \textbf{A}
  \item \textbf{A}
  \item \textbf{A}
  \item \textbf{A}
  \item \textbf{A}
  \item \textbf{A}
  \item \textbf{A}
  \item \textbf{A}
  \item \textbf{A}
  \item \textbf{A}
  \item \textbf{A}
  \item \textbf{A}
  \item \textbf{A}
  \item \textbf{A}
  \item \textbf{A}
  \item \textbf{A}
  \item \textbf{A}
  \item \textbf{A}
  \item \textbf{A}
  \item \textbf{A}
  \item \textbf{A}
  \item \textbf{A}
  \item \textbf{A}
  \item \textbf{A}
  \item \textbf{A}
  \item \textbf{A}
  \item \textbf{A}
  \item \textbf{A}
  \item \textbf{A}
  \item \textbf{A}
  \item \textbf{A}
  \item \textbf{A}
  \item \textbf{A}
  \item \textbf{A}
  \item \textbf{A}
  \item \textbf{A}
  \item \textbf{A}
  \item \textbf{A}
  \item \textbf{A}
  \item \textbf{A}
  \item \textbf{A}
  \item \textbf{A}
  \item \textbf{A}
  \item \textbf{A}
  \item \textbf{A}
  \item \textbf{A}
  \item \textbf{A}
  \item \textbf{A}
  \item \textbf{A}
  \item \textbf{A}
  \item \textbf{A}
  \item \textbf{A}
  \item \textbf{A}
  \item \textbf{A}
  \item \textbf{A}
  \item \textbf{A}
  \item \textbf{A}
  \item \textbf{A}
  \item \textbf{A}
  \item \textbf{A}
  \item \textbf{A}
  \item \textbf{A}
  \item \textbf{A}
  \item \textbf{A}
  \item \textbf{A}
  \item \textbf{A}
  \item \textbf{A}
  \item \textbf{A}
  \item \textbf{A}
  \item \textbf{A}
  \item \textbf{A}
  \item \textbf{A}
  \item \textbf{A}
  \item \textbf{A}
  \item \textbf{A}
  \item \textbf{A}
  \item \textbf{A}
  \item \textbf{A}
  \item \textbf{A}
  \item \textbf{A}
  \item \textbf{A}
  \item \textbf{A}
  \item \textbf{A}
  \item \textbf{A}
  \item \textbf{A}
  \item \textbf{A}
  \item \textbf{A}
  \item \textbf{A}
  \item \textbf{A}
  \item \textbf{A}
  \item \textbf{A}
  \item \textbf{A}
  \item \textbf{A}
  \item \textbf{A}
  \item \textbf{A}
  \item \textbf{A}
  \item \textbf{A}
  \item \textbf{A}
  \item \textbf{A}
  \item \textbf{A}
  \item \textbf{A}
  \item \textbf{A}
  \item \textbf{A}
  \item \textbf{A}
  \item \textbf{A}
  \item \textbf{A}
  \item \textbf{A}
  \item \textbf{A}
  \item \textbf{A}
  \item \textbf{A}
  \item \textbf{A}
  \item \textbf{A}
  \item \textbf{A}
  \item \textbf{A}
  \item \textbf{A}
  \item \textbf{A}
  \item \textbf{A}
  \item \textbf{A}
  \item \textbf{A}
  \item \textbf{A}
  \item \textbf{A}
  \item \textbf{A}
  \item \textbf{A}
  \item \textbf{A}
  \item \textbf{A}
  \item \textbf{A}
  \item \textbf{A}
  \item \textbf{A}
  \item \textbf{A}
  \item \textbf{A}
  \item \textbf{A}
  \item \textbf{A}
  \item \textbf{A}
  \item \textbf{A}
  \item \textbf{A}
  \item \textbf{A}
  \item \textbf{A}
  \item \textbf{A}
  \item \textbf{A}
  \item \textbf{A}
  \item \textbf{A}
  \item \textbf{A}
  \item \textbf{A}
  \item \textbf{A}
  \item \textbf{A}
  \item \textbf{A}
  \item \textbf{A}
  \item \textbf{A}
  \item \textbf{A}
  \item \textbf{A}
  \item \textbf{A}
  \item \textbf{A}
  \item \textbf{A}
  \item \textbf{A}
  \item \textbf{A}
  \item \textbf{A}
  \item \textbf{A}
  \item \textbf{A}
  \item \textbf{A}
  \item \textbf{A}
  \item \textbf{A}
  \item \textbf{A}
  \item \textbf{A}
  \item \textbf{A}
  \item \textbf{A}
  \item \textbf{A}
  \item \textbf{A}
  \item \textbf{A}
  \item \textbf{A}
  \item \textbf{A}
  \item \textbf{A}
  \item \textbf{A}
  \item \textbf{A}
  \item \textbf{A}
  \item \textbf{A}
  \item \textbf{A}
  \item \textbf{A}
  \item \textbf{A}
  \item \textbf{A}
  \item \textbf{A}
  \item \textbf{A}
  \item \textbf{A}
  \item \textbf{A}
  \item \textbf{A}
  \item \textbf{A}
  \item \textbf{A}
  \item \textbf{A}
  \item \textbf{A}
  \item \textbf{A}
  \item \textbf{A}
  \item \textbf{A}
  \item \textbf{A}
  \item \textbf{A}
  \item \textbf{A}
  \item \textbf{A}
  \item \textbf{A}
  \item \textbf{A}
  \item \textbf{A}
  \item \textbf{A}
  \item \textbf{A}
  \item \textbf{A}
  \item \textbf{A}
  \item \textbf{A}
  \item \textbf{A}
  \item \textbf{A}
  \item \textbf{A}
  \item \textbf{A}
  \item \textbf{A}
  \item \textbf{A}
  \item \textbf{A}
  \item \textbf{A}
  \item \textbf{A}
  \item \textbf{A}
  \item \textbf{A}
  \item \textbf{A}
\end{enumerate}
\end{multicols}
\end{tcolorbox}
\newpage\n
\subsection{2.5 PRAGMÁTICA}

\subsubsection*{SECCIÓN I: CONCEPTOS BÁSICOS DE PRAGMÁTICA (1–40)}

\begin{enumerate}[label=\textbf{\arabic*.}]
\setcounter{enumi}{0}
\item ¿Qué estudia la pragmática?
\begin{enumerate}[label=\textbf{\Alph*)}]
\begin{multicols}{2}
  \item La estructura de las oraciones
  \item La ortografía de las palabras
  \item La pronunciación de las palabras
\end{multicols}
\end{enumerate}
\vspace{0.15cm}
\item La pragmática se diferencia de la semántica porque:
\begin{enumerate}[label=\textbf{\Alph*)}]
\begin{multicols}{2}
\end{multicols}
\end{enumerate}
\vspace{0.15cm}
\item ¿Qué es el contexto en pragmática?
\begin{enumerate}[label=\textbf{\Alph*)}]
\begin{multicols}{2}
  \item Solo el lugar donde se habla
  \item Solo el tiempo en que se habla
  \item Solo el tema del que se habla
  \item Solo el canal de comunicación
  \item Solo el código utilizado
\end{multicols}
\end{enumerate}
\vspace{0.15cm}
\item El contexto situacional se refiere a:
\begin{enumerate}[label=\textbf{\Alph*)}]
\begin{multicols}{2}
  \item Las circunstancias físicas y sociales de la comunicación
  \item La estructura gramatical
  \item La ortografía del texto
  \item La caligrafía del texto
  \item La puntuación del texto
\end{multicols}
\end{enumerate}
\vspace{0.15cm}
\item El contexto lingüístico se refiere a:
\begin{enumerate}[label=\textbf{\Alph*)}]
\begin{multicols}{2}
  \item La intención del emisor
\end{multicols}
\end{enumerate}
\vspace{0.15cm}
\item El contexto sociocultural se refiere a:
\begin{enumerate}[label=\textbf{\Alph*)}]
\begin{multicols}{2}
  \item Las normas y valores de la comunidad donde se comunica
\end{multicols}
\end{enumerate}
\vspace{0.15cm}
\item ¿Qué es un enunciado?
\begin{enumerate}[label=\textbf{\Alph*)}]
\begin{multicols}{2}
  \item La unidad mínima de comunicación en contexto
  \item Una palabra aislada
  \item Una oración gramatical
  \item Un texto completo
  \item Un párrafo
  \item Un sintagma
\end{multicols}
\end{enumerate}
\vspace{0.15cm}
\item ¿Qué es la intención comunicativa?
\begin{enumerate}[label=\textbf{\Alph*)}]
\begin{multicols}{2}
  \item La estructura gramatical
  \item La ortografía del texto
  \item La puntuación del texto
  \item La caligrafía del texto
\end{multicols}
\end{enumerate}
\vspace{0.15cm}
\item ¿Qué es el significado literal?
\begin{enumerate}[label=\textbf{\Alph*)}]
\begin{multicols}{2}
  \item La intención del emisor
  \item La presuposición
  \item La implicatura
\end{multicols}
\end{enumerate}
\vspace{0.15cm}
\item ¿Qué es el significado inferencial?
\begin{enumerate}[label=\textbf{\Alph*)}]
\begin{multicols}{2}
  \item La estructura gramatical
  \item La ortografía del texto
  \item La puntuación del texto
  \item La caligrafía del texto
\end{multicols}
\end{enumerate}
\vspace{0.15cm}
\item ¿Qué es la comunicación no verbal?
\begin{enumerate}[label=\textbf{\Alph*)}]
\begin{multicols}{2}
  \item La que se realiza sin palabras, mediante gestos, tono, etc.
  \item La que se realiza solo con palabras
  \item La que se realiza por escrito
  \item La que se realiza por señas
  \item La que se realiza por teléfono
  \item La que se realiza por correo
\end{multicols}
\end{enumerate}
\vspace{0.15cm}
\item ¿Qué es la deixis?
\begin{enumerate}[label=\textbf{\Alph*)}]
\begin{multicols}{2}
  \item La referencia a elementos del contexto mediante el lenguaje
  \item La estructura de las oraciones
  \item La ortografía del texto
  \item La puntuación del texto
  \item La caligrafía del texto
\end{multicols}
\end{enumerate}
\vspace{0.15cm}
\item La deixis personal se refiere a:
\begin{enumerate}[label=\textbf{\Alph*)}]
\begin{multicols}{2}
  \item Pronombres personales y posesivos
  \item Preposiciones
\end{multicols}
\end{enumerate}
\vspace{0.15cm}
\item La deixis espacial se refiere a:
\begin{enumerate}[label=\textbf{\Alph*)}]
\begin{multicols}{2}
  \item Pronombres personales
  \item Preposiciones
\end{multicols}
\end{enumerate}
\vspace{0.15cm}
\item La deixis temporal se refiere a:
\begin{enumerate}[label=\textbf{\Alph*)}]
\begin{multicols}{2}
  \item Pronombres personales
  \item Preposiciones
\end{multicols}
\end{enumerate}
\vspace{0.15cm}
\item La deixis social se refiere a:
\begin{enumerate}[label=\textbf{\Alph*)}]
\begin{multicols}{2}
  \item Las formas de tratamiento y cortesía
  \item Los pronombres personales
  \item Los adverbios de lugar
  \item Los adverbios de tiempo
  \item Los conectores
  \item Las preposiciones
\end{multicols}
\end{enumerate}
\vspace{0.15cm}
\item ¿Qué es la referencia?
\begin{enumerate}[label=\textbf{\Alph*)}]
\begin{multicols}{2}
  \item La relación entre un elemento lingüístico y aquello a lo que alude
  \item La estructura de las oraciones
  \item La ortografía del texto
  \item La puntuación del texto
  \item La caligrafía del texto
\end{multicols}
\end{enumerate}
\vspace{0.15cm}
\item ¿Qué es la anáfora?
\begin{enumerate}[label=\textbf{\Alph*)}]
\begin{multicols}{2}
  \item La referencia a un elemento mencionado antes
  \item La referencia a un elemento mencionado después
  \item La ausencia de referencia
  \item La repetición de palabras
\end{multicols}
\end{enumerate}
\vspace{0.15cm}
\item ¿Qué es la catáfora?
\begin{enumerate}[label=\textbf{\Alph*)}]
\begin{multicols}{2}
  \item La referencia a un elemento mencionado después
  \item La referencia a un elemento mencionado antes
  \item La ausencia de referencia
  \item La repetición de palabras
\end{multicols}
\end{enumerate}
\vspace{0.15cm}
\item ¿Qué es la elipsis?
\begin{enumerate}[label=\textbf{\Alph*)}]
\begin{multicols}{2}
  \item La omisión de elementos que se sobreentienden
  \item La repetición de elementos
  \item La ausencia de referencia
\end{multicols}
\end{enumerate}
\vspace{0.15cm}
\item ¿Qué es la presuposición?
\begin{enumerate}[label=\textbf{\Alph*)}]
\begin{multicols}{2}
  \item La información que se da por supuesta en un enunciado
  \item La información nueva que se aporta
  \item La intención del emisor
\end{multicols}
\end{enumerate}
\vspace{0.15cm}
\item ¿Qué es la implicatura?
\begin{enumerate}[label=\textbf{\Alph*)}]
\begin{multicols}{2}
  \item La estructura gramatical
  \item La ortografía del texto
  \item La puntuación del texto
  \item La caligrafía del texto
\end{multicols}
\end{enumerate}
\vspace{0.15cm}
\item ¿Qué es el principio de cooperación?
\begin{enumerate}[label=\textbf{\Alph*)}]
\begin{multicols}{2}
\end{multicols}
\end{enumerate}
\vspace{0.15cm}
\item ¿Quién propuso el principio de cooperación?
\begin{enumerate}[label=\textbf{\Alph*)}]
\begin{multicols}{2}
  \item Paul Grice
  \item Roman Jakobson
  \item Michael Halliday
\end{multicols}
\end{enumerate}
\vspace{0.15cm}
\item Las máximas de Grice son:
\begin{enumerate}[label=\textbf{\Alph*)}]
\begin{multicols}{2}
\end{multicols}
\end{enumerate}
\vspace{0.15cm}
\item La máxima de cantidad se refiere a:
\begin{enumerate}[label=\textbf{\Alph*)}]
\begin{multicols}{2}
  \item Ser relevante
  \item Ser claro
  \item Ser cortés
  \item Ser breve
\end{multicols}
\end{enumerate}
\vspace{0.15cm}
\item La máxima de calidad se refiere a:
\begin{enumerate}[label=\textbf{\Alph*)}]
\begin{multicols}{2}
  \item Ser relevante
  \item Ser claro
  \item Ser cortés
  \item Ser breve
\end{multicols}
\end{enumerate}
\vspace{0.15cm}
\item La máxima de relación se refiere a:
\begin{enumerate}[label=\textbf{\Alph*)}]
\begin{multicols}{2}
  \item Ser relevante y pertinente
  \item Ser claro
  \item Ser cortés
  \item Ser breve
\end{multicols}
\end{enumerate}
\vspace{0.15cm}
\item La máxima de manera se refiere a:
\begin{enumerate}[label=\textbf{\Alph*)}]
\begin{multicols}{2}
  \item Ser claro, breve y ordenado
  \item Ser relevante
  \item Ser cortés
  \item Ser preciso
\end{multicols}
\end{enumerate}
\vspace{0.15cm}
\item ¿Qué es la cortesía verbal?
\begin{enumerate}[label=\textbf{\Alph*)}]
\begin{multicols}{2}
\end{multicols}
\end{enumerate}
\vspace{0.15cm}
\item ¿Qué es la imagen pública o "face"?
\begin{enumerate}[label=\textbf{\Alph*)}]
\begin{multicols}{2}
  \item La imagen social que cada persona proyecta
  \item La imagen física de una persona
  \item La imagen digital de una persona
  \item La imagen profesional de una persona
  \item La imagen académica de una persona
  \item La imagen personal de una persona
\end{multicols}
\end{enumerate}
\vspace{0.15cm}
\item La cortesía positiva se refiere a:
\begin{enumerate}[label=\textbf{\Alph*)}]
\begin{multicols}{2}
  \item Mostrar aprecio y solidaridad hacia el interlocutor
  \item Mostrar distancia y respeto
  \item Ignorar al interlocutor
  \item Ofender al interlocutor
\end{multicols}
\end{enumerate}
\vspace{0.15cm}
\item La cortesía negativa se refiere a:
\begin{enumerate}[label=\textbf{\Alph*)}]
\begin{multicols}{2}
  \item Mostrar respeto y evitar imposiciones
  \item Mostrar aprecio y solidaridad
  \item Ignorar al interlocutor
  \item Ofender al interlocutor
\end{multicols}
\end{enumerate}
\vspace{0.15cm}
\item ¿Qué es un acto de habla?
\begin{enumerate}[label=\textbf{\Alph*)}]
\begin{multicols}{2}
  \item La acción que se realiza al emitir un enunciado
  \item La estructura gramatical de una oración
  \item La ortografía del texto
  \item La puntuación del texto
  \item La caligrafía del texto
\end{multicols}
\end{enumerate}
\vspace{0.15cm}
\item ¿Quién propuso la teoría de los actos de habla?
\begin{enumerate}[label=\textbf{\Alph*)}]
\begin{multicols}{2}
  \item Paul Grice
  \item Roman Jakobson
  \item Michael Halliday
\end{multicols}
\end{enumerate}
\vspace{0.15cm}
\item Los actos de habla según Austin son:
\begin{enumerate}[label=\textbf{\Alph*)}]
\begin{multicols}{2}
  \item Locutivo, ilocutivo y perlocutivo
  \item Locutivo, ilocutivo y proposicional
  \item Locutivo, proposicional y perlocutivo
  \item Ilocutivo, proposicional y perlocutivo
  \item Locutivo, ilocutivo y expresivo
  \item Locutivo, ilocutivo y directivo
\end{multicols}
\end{enumerate}
\vspace{0.15cm}
\item El acto locutivo se refiere a:
\begin{enumerate}[label=\textbf{\Alph*)}]
\begin{multicols}{2}
  \item La intención del hablante
  \item La estructura gramatical
  \item La ortografía del texto
  \item La puntuación del texto
\end{multicols}
\end{enumerate}
\vspace{0.15cm}
\item El acto ilocutivo se refiere a:
\begin{enumerate}[label=\textbf{\Alph*)}]
\begin{multicols}{2}
  \item La intención o fuerza del acto al decir algo
  \item La estructura gramatical
  \item La ortografía del texto
  \item La puntuación del texto
\end{multicols}
\end{enumerate}
\vspace{0.15cm}
\item El acto perlocutivo se refiere a:
\begin{enumerate}[label=\textbf{\Alph*)}]
\begin{multicols}{2}
  \item La intención del hablante
  \item La estructura gramatical
  \item La ortografía del texto
  \item La puntuación del texto
\end{multicols}
\end{enumerate}
\vspace{0.15cm}
\item Según Searle, los actos de habla se clasifican en:
\begin{enumerate}[label=\textbf{\Alph*)}]
\begin{multicols}{2}
  \item Locutivos, ilocutivos y perlocutivos
  \item Literales y figurados
  \item Verbales y no verbales
  \item Orales y escritos ---
\end{multicols}
\end{enumerate}
\vspace{0.15cm}
\end{enumerate}\n
\subsubsection*{SECCIÓN II: ACTOS DE HABLA (41–100)}

\begin{enumerate}[label=\textbf{\arabic*.}]
\setcounter{enumi}{40}
\item Los actos de habla asertivos o representativos:
\begin{enumerate}[label=\textbf{\Alph*)}]
\begin{multicols}{2}
  \item Intentan que el oyente haga algo
  \item Preguntan algo
\end{multicols}
\end{enumerate}
\vspace{0.15cm}
\item Los actos de habla directivos:
\begin{enumerate}[label=\textbf{\Alph*)}]
\begin{multicols}{2}
  \item Intentan que el oyente haga algo
\end{multicols}
\end{enumerate}
\vspace{0.15cm}
\item Los actos de habla compromisivos:
\begin{enumerate}[label=\textbf{\Alph*)}]
\begin{multicols}{2}
  \item Intentan que el oyente haga algo
\end{multicols}
\end{enumerate}
\vspace{0.15cm}
\item Los actos de habla expresivos:
\begin{enumerate}[label=\textbf{\Alph*)}]
\begin{multicols}{2}
  \item Intentan que el oyente haga algo
\end{multicols}
\end{enumerate}
\vspace{0.15cm}
\item Los actos de habla declarativos:
\begin{enumerate}[label=\textbf{\Alph*)}]
\begin{multicols}{2}
  \item Intentan que el oyente haga algo
  \item Preguntan algo
\end{multicols}
\end{enumerate}
\vspace{0.15cm}
\item "Te prometo que iré" es un acto de habla:
\begin{enumerate}[label=\textbf{\Alph*)}]
\begin{multicols}{2}
  \item Interrogativo
\end{multicols}
\end{enumerate}
\vspace{0.15cm}
\item "¿Puedes cerrar la puerta?" es un acto de habla:
\begin{enumerate}[label=\textbf{\Alph*)}]
\begin{multicols}{2}
\end{multicols}
\end{enumerate}
\vspace{0.15cm}
\item "Lo siento mucho" es un acto de habla:
\begin{enumerate}[label=\textbf{\Alph*)}]
\begin{multicols}{2}
  \item Interrogativo
\end{multicols}
\end{enumerate}
\vspace{0.15cm}
\item "Declaro inaugurada la sesión" es un acto de habla:
\begin{enumerate}[label=\textbf{\Alph*)}]
\begin{multicols}{2}
  \item Interrogativo
\end{multicols}
\end{enumerate}
\vspace{0.15cm}
\item "El cielo es azul" es un acto de habla:
\begin{enumerate}[label=\textbf{\Alph*)}]
\begin{multicols}{2}
  \item Interrogativo
\end{multicols}
\end{enumerate}
\vspace{0.15cm}
\item "Te aconsejo que estudies" es un acto de habla:
\begin{enumerate}[label=\textbf{\Alph*)}]
\begin{multicols}{2}
  \item Interrogativo
\end{multicols}
\end{enumerate}
\vspace{0.15cm}
\item "¡Qué alegría verte!" es un acto de habla:
\begin{enumerate}[label=\textbf{\Alph*)}]
\begin{multicols}{2}
  \item Interrogativo
\end{multicols}
\end{enumerate}
\vspace{0.15cm}
\item "Juro decir la verdad" es un acto de habla:
\begin{enumerate}[label=\textbf{\Alph*)}]
\begin{multicols}{2}
  \item Interrogativo
\end{multicols}
\end{enumerate}
\vspace{0.15cm}
\item "Bautizo este barco como Aurora" es un acto de habla:
\begin{enumerate}[label=\textbf{\Alph*)}]
\begin{multicols}{2}
  \item Interrogativo
\end{multicols}
\end{enumerate}
\vspace{0.15cm}
\item Un acto de habla directo es:
\begin{enumerate}[label=\textbf{\Alph*)}]
\begin{multicols}{2}
\end{multicols}
\end{enumerate}
\vspace{0.15cm}
\item Un acto de habla indirecto es:
\begin{enumerate}[label=\textbf{\Alph*)}]
\begin{multicols}{2}
\end{multicols}
\end{enumerate}
\vspace{0.15cm}
\item "¿Tienes hora?" es un acto de habla:
\begin{enumerate}[label=\textbf{\Alph*)}]
\begin{multicols}{2}
\end{multicols}
\end{enumerate}
\vspace{0.15cm}
\item "Hace frío aquí" puede ser un acto de habla:
\begin{enumerate}[label=\textbf{\Alph*)}]
\begin{multicols}{2}
\end{multicols}
\end{enumerate}
\vspace{0.15cm}
\item "¿Podrías pasarme la sal?" es un acto de habla:
\begin{enumerate}[label=\textbf{\Alph*)}]
\begin{multicols}{2}
\end{multicols}
\end{enumerate}
\vspace{0.15cm}
\item "Ya es tarde" puede ser un acto de habla:
\begin{enumerate}[label=\textbf{\Alph*)}]
\begin{multicols}{2}
\end{multicols}
\end{enumerate}
\vspace{0.15cm}
\item La fuerza ilocutiva se refiere a:
\begin{enumerate}[label=\textbf{\Alph*)}]
\begin{multicols}{2}
  \item La intención comunicativa del hablante
  \item La estructura gramatical
  \item La ortografía del texto
  \item La puntuación del texto
\end{multicols}
\end{enumerate}
\vspace{0.15cm}
\item El efecto perlocutivo se refiere a:
\begin{enumerate}[label=\textbf{\Alph*)}]
\begin{multicols}{2}
  \item La consecuencia o reacción que produce el enunciado
  \item La intención del hablante
  \item La estructura gramatical
  \item La ortografía del texto
  \item La puntuación del texto
\end{multicols}
\end{enumerate}
\vspace{0.15cm}
\item Los verbos realizativos son:
\begin{enumerate}[label=\textbf{\Alph*)}]
\begin{multicols}{2}
  \item Verbos que al ser pronunciados realizan la acción
  \item Verbos que describen acciones
  \item Verbos que no tienen significado
  \item Verbos que no tienen sujeto
  \item Verbos que no tienen objeto
  \item Verbos que no tienen tiempo
\end{multicols}
\end{enumerate}
\vspace{0.15cm}
\item "Prometo" es un verbo:
\begin{enumerate}[label=\textbf{\Alph*)}]
\begin{multicols}{2}
  \item Realizativo
  \item Impersonal
  \item Reflexivo
  \item Recíproco
\end{multicols}
\end{enumerate}
\vspace{0.15cm}
\item "Juro" es un verbo:
\begin{enumerate}[label=\textbf{\Alph*)}]
\begin{multicols}{2}
  \item Realizativo
  \item Impersonal
  \item Reflexivo
  \item Recíproco
\end{multicols}
\end{enumerate}
\vspace{0.15cm}
\item "Declaro" es un verbo:
\begin{enumerate}[label=\textbf{\Alph*)}]
\begin{multicols}{2}
  \item Realizativo
  \item Impersonal
  \item Reflexivo
  \item Recíproco
\end{multicols}
\end{enumerate}
\vspace{0.15cm}
\item Las condiciones de felicidad de un acto de habla son:
\begin{enumerate}[label=\textbf{\Alph*)}]
\begin{multicols}{2}
  \item Las condiciones que deben cumplirse para que el acto sea válido
  \item Las condiciones gramaticales
  \item Las condiciones ortográficas
  \item Las condiciones de puntuación
  \item Las condiciones de caligrafía
  \item Las condiciones sintácticas
\end{multicols}
\end{enumerate}
\vspace{0.15cm}
\item Para que una promesa sea válida, se requiere:
\begin{enumerate}[label=\textbf{\Alph*)}]
\begin{multicols}{2}
  \item Que el hablante tenga intención de cumplirla
  \item Que el hablante no tenga intención de cumplirla
  \item Que el oyente no la escuche
  \item Que no haya contexto
  \item Que no haya intención
  \item Que no haya efecto
\end{multicols}
\end{enumerate}
\vspace{0.15cm}
\item La cortesía en los actos de habla busca:
\begin{enumerate}[label=\textbf{\Alph*)}]
\begin{multicols}{2}
  \item Mantener la relación social y evitar conflictos
  \item Ofender al interlocutor
  \item Ignorar al interlocutor
\end{multicols}
\end{enumerate}
\vspace{0.15cm}
\item La atenuación en la cortesía consiste en:
\begin{enumerate}[label=\textbf{\Alph*)}]
\begin{multicols}{2}
  \item Suavizar la fuerza ilocutiva de un acto
  \item Ignorar la fuerza ilocutiva
  \item Repetir la fuerza ilocutiva
\end{multicols}
\end{enumerate}
\vspace{0.15cm}
\item "¿Serías tan amable de...?" es una estrategia de:
\begin{enumerate}[label=\textbf{\Alph*)}]
\begin{multicols}{2}
  \item Ironía
  \item Sarcasmo
\end{multicols}
\end{enumerate}
\vspace{0.15cm}
\item "¡Qué guapo estás hoy!" es una estrategia de:
\begin{enumerate}[label=\textbf{\Alph*)}]
\begin{multicols}{2}
  \item Ironía
  \item Sarcasmo
\end{multicols}
\end{enumerate}
\vspace{0.15cm}
\item La ironía es:
\begin{enumerate}[label=\textbf{\Alph*)}]
\begin{multicols}{2}
  \item No decir nada
\end{multicols}
\end{enumerate}
\vspace{0.15cm}
\item "¡Qué puntual eres!" dicho a alguien que llega tarde es un ejemplo de:
\begin{enumerate}[label=\textbf{\Alph*)}]
\begin{multicols}{2}
  \item Ironía
  \item Sarcasmo
\end{multicols}
\end{enumerate}
\vspace{0.15cm}
\item La metáfora es un fenómeno:
\begin{enumerate}[label=\textbf{\Alph*)}]
\begin{multicols}{2}
  \item Sintáctico
  \item Ortográfico
  \item Morfológico
\end{multicols}
\end{enumerate}
\vspace{0.15cm}
\item "Eres un sol" es un ejemplo de:
\begin{enumerate}[label=\textbf{\Alph*)}]
\begin{multicols}{2}
  \item Metáfora
  \item Metonimia
  \item Ironía
  \item Hipérbole
  \item Sarcasmo
\end{multicols}
\end{enumerate}
\vspace{0.15cm}
\item "Me muero de hambre" es un ejemplo de:
\begin{enumerate}[label=\textbf{\Alph*)}]
\begin{multicols}{2}
  \item Hipérbole
  \item Metáfora
  \item Metonimia
  \item Ironía
  \item Sarcasmo
\end{multicols}
\end{enumerate}
\vspace{0.15cm}
\item "Tiene el corazón de piedra" es un ejemplo de:
\begin{enumerate}[label=\textbf{\Alph*)}]
\begin{multicols}{2}
  \item Metáfora
  \item Metonimia
  \item Hipérbole
  \item Ironía
  \item Sarcasmo
\end{multicols}
\end{enumerate}
\vspace{0.15cm}
\item La pragmática intercultural estudia:
\begin{enumerate}[label=\textbf{\Alph*)}]
\begin{multicols}{2}
  \item Las diferencias en el uso del lenguaje entre culturas
  \item Las diferencias gramaticales entre lenguas
  \item Las diferencias ortográficas entre lenguas
  \item Las diferencias fonéticas entre lenguas
  \item Las diferencias morfológicas entre lenguas
  \item Las diferencias sintácticas entre lenguas
\end{multicols}
\end{enumerate}
\vspace{0.15cm}
\item La pragmática del español estudia:
\begin{enumerate}[label=\textbf{\Alph*)}]
\begin{multicols}{2}
  \item La gramática del español
  \item La ortografía del español
  \item La fonética del español
  \item La morfología del español
  \item La sintaxis del español
\end{multicols}
\end{enumerate}
\vspace{0.15cm}
\item El "tú" y el "usted" en español son ejemplos de:
\begin{enumerate}[label=\textbf{\Alph*)}]
\begin{multicols}{2}
\end{multicols}
\end{enumerate}
\vspace{0.15cm}
\item El uso de "usted" indica:
\begin{enumerate}[label=\textbf{\Alph*)}]
\begin{multicols}{2}
  \item Ironía
  \item Sarcasmo
\end{multicols}
\end{enumerate}
\vspace{0.15cm}
\item El uso de "tú" indica:
\begin{enumerate}[label=\textbf{\Alph*)}]
\begin{multicols}{2}
  \item Ironía
  \item Sarcasmo
\end{multicols}
\end{enumerate}
\vspace{0.15cm}
\item El voseo es:
\begin{enumerate}[label=\textbf{\Alph*)}]
\begin{multicols}{2}
\end{multicols}
\end{enumerate}
\vspace{0.15cm}
\item La pragmática del discurso estudia:
\begin{enumerate}[label=\textbf{\Alph*)}]
\begin{multicols}{2}
  \item La organización del lenguaje en unidades mayores que la oración
  \item La estructura de las oraciones
  \item La ortografía del texto
  \item La puntuación del texto
  \item La caligrafía del texto
\end{multicols}
\end{enumerate}
\vspace{0.15cm}
\item El análisis de la conversación estudia:
\begin{enumerate}[label=\textbf{\Alph*)}]
\begin{multicols}{2}
  \item La organización secuencial de la interacción verbal
  \item La estructura gramatical
  \item La ortografía del texto
  \item La puntuación del texto
  \item La caligrafía del texto
\end{multicols}
\end{enumerate}
\vspace{0.15cm}
\item Los turnos de habla son:
\begin{enumerate}[label=\textbf{\Alph*)}]
\begin{multicols}{2}
  \item Las intervenciones de cada participante en una conversación
  \item Las oraciones gramaticales
  \item Las palabras aisladas
  \item Los párrafos
  \item Los textos
  \item Los sintagmas
\end{multicols}
\end{enumerate}
\vspace{0.15cm}
\item La toma de turno se refiere a:
\begin{enumerate}[label=\textbf{\Alph*)}]
\begin{multicols}{2}
  \item La estructura gramatical
  \item La ortografía del texto
  \item La puntuación del texto
  \item La caligrafía del texto
\end{multicols}
\end{enumerate}
\vspace{0.15cm}
\item Los pares adyacentes son:
\begin{enumerate}[label=\textbf{\Alph*)}]
\begin{multicols}{2}
  \item Secuencias de tres turnos
  \item Secuencias de cuatro turnos
  \item Secuencias de un solo turno
  \item Secuencias sin relación
  \item Secuencias gramaticales
\end{multicols}
\end{enumerate}
\vspace{0.15cm}
\item "¿Cómo estás?" - "Bien, gracias" es un ejemplo de:
\begin{enumerate}[label=\textbf{\Alph*)}]
\begin{multicols}{2}
  \item Par adyacente
  \item Turno único
  \item Secuencia de tres
  \item Secuencia de cuatro
  \item Secuencia sin relación
  \item Secuencia gramatical
\end{multicols}
\end{enumerate}
\vspace{0.15cm}
\item La reparación en la conversación consiste en:
\begin{enumerate}[label=\textbf{\Alph*)}]
\begin{multicols}{2}
  \item Ignorar un error
  \item Repetir un error
\end{multicols}
\end{enumerate}
\vspace{0.15cm}
\item La superposición en la conversación se produce cuando:
\begin{enumerate}[label=\textbf{\Alph*)}]
\begin{multicols}{2}
  \item Un participante habla solo
  \item No habla nadie
  \item Se guarda silencio
  \item Se cambia de tema
  \item Se alarga la conversación
\end{multicols}
\end{enumerate}
\vspace{0.15cm}
\item La pragmática textual estudia:
\begin{enumerate}[label=\textbf{\Alph*)}]
\begin{multicols}{2}
  \item La organización y uso del lenguaje en textos
  \item La estructura de las oraciones
  \item La ortografía del texto
  \item La puntuación del texto
  \item La caligrafía del texto
\end{multicols}
\end{enumerate}
\vspace{0.15cm}
\item La coherencia pragmática se refiere a:
\begin{enumerate}[label=\textbf{\Alph*)}]
\begin{multicols}{2}
  \item La adecuación del texto al contexto comunicativo
  \item La unidad temática
  \item La cohesión gramatical
  \item La ortografía del texto
  \item La puntuación del texto
  \item La caligrafía del texto
\end{multicols}
\end{enumerate}
\vspace{0.15cm}
\item La relevancia en pragmática se refiere a:
\begin{enumerate}[label=\textbf{\Alph*)}]
\begin{multicols}{2}
  \item La pertinencia de la información en el contexto
  \item La estructura gramatical
  \item La ortografía del texto
  \item La puntuación del texto
  \item La caligrafía del texto
\end{multicols}
\end{enumerate}
\vspace{0.15cm}
\item La teoría de la relevancia fue propuesta por:
\begin{enumerate}[label=\textbf{\Alph*)}]
\begin{multicols}{2}
  \item Sperber y Wilson
  \item Grice
  \item Searle
  \item Saussure
\end{multicols}
\end{enumerate}
\vspace{0.15cm}
\item Según la teoría de la relevancia, la comunicación busca:
\begin{enumerate}[label=\textbf{\Alph*)}]
\begin{multicols}{2}
  \item Maximizar la relevancia con el mínimo esfuerzo
  \item Maximizar el esfuerzo
  \item Minimizar la relevancia
  \item Ignorar el contexto
  \item Ignorar la intención
  \item Ignorar el efecto
\end{multicols}
\end{enumerate}
\vspace{0.15cm}
\item El contexto cognitivo en la teoría de la relevancia se refiere a:
\begin{enumerate}[label=\textbf{\Alph*)}]
\begin{multicols}{2}
\end{multicols}
\end{enumerate}
\vspace{0.15cm}
\item El efecto cognitivo en la teoría de la relevancia se refiere a:
\begin{enumerate}[label=\textbf{\Alph*)}]
\begin{multicols}{2}
  \item La modificación de los supuestos del oyente
  \item La estructura gramatical
  \item La ortografía del texto
  \item La puntuación del texto
  \item La caligrafía del texto
\end{multicols}
\end{enumerate}
\vspace{0.15cm}
\item El esfuerzo de procesamiento en la teoría de la relevancia se refiere a:
\begin{enumerate}[label=\textbf{\Alph*)}]
\begin{multicols}{2}
  \item La estructura gramatical
  \item La ortografía del texto
  \item La puntuación del texto
  \item La caligrafía del texto ---
\end{multicols}
\end{enumerate}
\vspace{0.15cm}
\end{enumerate}\n
\subsubsection*{SECCIÓN III: DEIXIS Y REFERENCIA (101–150)}

\begin{enumerate}[label=\textbf{\arabic*.}]
\setcounter{enumi}{100}
\item La deixis personal se expresa mediante:
\begin{enumerate}[label=\textbf{\Alph*)}]
\begin{multicols}{2}
  \item Pronombres personales y posesivos
  \item Preposiciones
\end{multicols}
\end{enumerate}
\vspace{0.15cm}
\item "Yo" es un ejemplo de:
\begin{enumerate}[label=\textbf{\Alph*)}]
\begin{multicols}{2}
\end{multicols}
\end{enumerate}
\vspace{0.15cm}
\item "Tú" es un ejemplo de:
\begin{enumerate}[label=\textbf{\Alph*)}]
\begin{multicols}{2}
\end{multicols}
\end{enumerate}
\vspace{0.15cm}
\item "Él" es un ejemplo de:
\begin{enumerate}[label=\textbf{\Alph*)}]
\begin{multicols}{2}
\end{multicols}
\end{enumerate}
\vspace{0.15cm}
\item "Mi" es un ejemplo de:
\begin{enumerate}[label=\textbf{\Alph*)}]
\begin{multicols}{2}
\end{multicols}
\end{enumerate}
\vspace{0.15cm}
\item "Tu" es un ejemplo de:
\begin{enumerate}[label=\textbf{\Alph*)}]
\begin{multicols}{2}
\end{multicols}
\end{enumerate}
\vspace{0.15cm}
\item La deixis espacial se expresa mediante:
\begin{enumerate}[label=\textbf{\Alph*)}]
\begin{multicols}{2}
  \item Pronombres personales
  \item Preposiciones
\end{multicols}
\end{enumerate}
\vspace{0.15cm}
\item "Aquí" es un ejemplo de:
\begin{enumerate}[label=\textbf{\Alph*)}]
\begin{multicols}{2}
\end{multicols}
\end{enumerate}
\vspace{0.15cm}
\item "Allí" es un ejemplo de:
\begin{enumerate}[label=\textbf{\Alph*)}]
\begin{multicols}{2}
\end{multicols}
\end{enumerate}
\vspace{0.15cm}
\item "Este" es un ejemplo de:
\begin{enumerate}[label=\textbf{\Alph*)}]
\begin{multicols}{2}
\end{multicols}
\end{enumerate}
\vspace{0.15cm}
\item "Ese" es un ejemplo de:
\begin{enumerate}[label=\textbf{\Alph*)}]
\begin{multicols}{2}
\end{multicols}
\end{enumerate}
\vspace{0.15cm}
\item "Aquel" es un ejemplo de:
\begin{enumerate}[label=\textbf{\Alph*)}]
\begin{multicols}{2}
\end{multicols}
\end{enumerate}
\vspace{0.15cm}
\item La deixis temporal se expresa mediante:
\begin{enumerate}[label=\textbf{\Alph*)}]
\begin{multicols}{2}
  \item Pronombres personales
  \item Preposiciones
\end{multicols}
\end{enumerate}
\vspace{0.15cm}
\item "Ahora" es un ejemplo de:
\begin{enumerate}[label=\textbf{\Alph*)}]
\begin{multicols}{2}
\end{multicols}
\end{enumerate}
\vspace{0.15cm}
\item "Ayer" es un ejemplo de:
\begin{enumerate}[label=\textbf{\Alph*)}]
\begin{multicols}{2}
\end{multicols}
\end{enumerate}
\vspace{0.15cm}
\item "Mañana" es un ejemplo de:
\begin{enumerate}[label=\textbf{\Alph*)}]
\begin{multicols}{2}
\end{multicols}
\end{enumerate}
\vspace{0.15cm}
\item "Hoy" es un ejemplo de:
\begin{enumerate}[label=\textbf{\Alph*)}]
\begin{multicols}{2}
\end{multicols}
\end{enumerate}
\vspace{0.15cm}
\item "Pronto" es un ejemplo de:
\begin{enumerate}[label=\textbf{\Alph*)}]
\begin{multicols}{2}
\end{multicols}
\end{enumerate}
\vspace{0.15cm}
\item La deixis social se expresa mediante:
\begin{enumerate}[label=\textbf{\Alph*)}]
\begin{multicols}{2}
  \item Pronombres personales
  \item Preposiciones
\end{multicols}
\end{enumerate}
\vspace{0.15cm}
\item "Usted" es un ejemplo de:
\begin{enumerate}[label=\textbf{\Alph*)}]
\begin{multicols}{2}
\end{multicols}
\end{enumerate}
\vspace{0.15cm}
\item "Señor" es un ejemplo de:
\begin{enumerate}[label=\textbf{\Alph*)}]
\begin{multicols}{2}
\end{multicols}
\end{enumerate}
\vspace{0.15cm}
\item "Don" es un ejemplo de:
\begin{enumerate}[label=\textbf{\Alph*)}]
\begin{multicols}{2}
\end{multicols}
\end{enumerate}
\vspace{0.15cm}
\item La referencia es:
\begin{enumerate}[label=\textbf{\Alph*)}]
\begin{multicols}{2}
  \item La relación entre un elemento lingüístico y aquello a lo que alude
  \item La estructura de las oraciones
  \item La ortografía del texto
  \item La puntuación del texto
  \item La caligrafía del texto
\end{multicols}
\end{enumerate}
\vspace{0.15cm}
\item La referencia exofórica se produce cuando:
\begin{enumerate}[label=\textbf{\Alph*)}]
\begin{multicols}{2}
  \item No hay referente
  \item Se repite la misma palabra
  \item Se usan sinónimos
  \item Se usan antónimos
\end{multicols}
\end{enumerate}
\vspace{0.15cm}
\item La referencia endofórica se produce cuando:
\begin{enumerate}[label=\textbf{\Alph*)}]
\begin{multicols}{2}
  \item No hay referente
  \item Se repite la misma palabra
  \item Se usan sinónimos
  \item Se usan antónimos
\end{multicols}
\end{enumerate}
\vspace{0.15cm}
\item La anáfora es un tipo de:
\begin{enumerate}[label=\textbf{\Alph*)}]
\begin{multicols}{2}
  \item Referencia endofórica
  \item Referencia exofórica
  \item Sinonimia
\end{multicols}
\end{enumerate}
\vspace{0.15cm}
\item La catáfora es un tipo de:
\begin{enumerate}[label=\textbf{\Alph*)}]
\begin{multicols}{2}
  \item Referencia endofórica
  \item Referencia exofórica
  \item Sinonimia
\end{multicols}
\end{enumerate}
\vspace{0.15cm}
\item En "Juan llegó. Él estaba cansado", "él" es:
\begin{enumerate}[label=\textbf{\Alph*)}]
\begin{multicols}{2}
  \item Sinonimia
\end{multicols}
\end{enumerate}
\vspace{0.15cm}
\item En "Cuando lo vi, Pedro corría", "lo" es:
\begin{enumerate}[label=\textbf{\Alph*)}]
\begin{multicols}{2}
  \item Sinonimia
\end{multicols}
\end{enumerate}
\vspace{0.15cm}
\item En "Compré flores. Las puse en un jarrón", "las" es:
\begin{enumerate}[label=\textbf{\Alph*)}]
\begin{multicols}{2}
  \item Sinonimia
\end{multicols}
\end{enumerate}
\vspace{0.15cm}
\item La elipsis es:
\begin{enumerate}[label=\textbf{\Alph*)}]
\begin{multicols}{2}
  \item La omisión de elementos que se sobreentienden
  \item La repetición de elementos
  \item La ausencia de referencia
\end{multicols}
\end{enumerate}
\vspace{0.15cm}
\item En "Juan estudia y María también", hay:
\begin{enumerate}[label=\textbf{\Alph*)}]
\begin{multicols}{2}
  \item Sinonimia
\end{multicols}
\end{enumerate}
\vspace{0.15cm}
\item En "Yo compré el libro y ella también", hay:
\begin{enumerate}[label=\textbf{\Alph*)}]
\begin{multicols}{2}
  \item Sinonimia
\end{multicols}
\end{enumerate}
\vspace{0.15cm}
\item La sustitución en la cohesión textual consiste en:
\begin{enumerate}[label=\textbf{\Alph*)}]
\begin{multicols}{2}
  \item Reemplazar un elemento por otro
  \item Repetir el mismo elemento
  \item Omitir el elemento
  \item Usar sinónimos
  \item Usar antónimos
  \item Usar conectores
\end{multicols}
\end{enumerate}
\vspace{0.15cm}
\item En "Juan llegó. El hombre venía cansado", "el hombre" es:
\begin{enumerate}[label=\textbf{\Alph*)}]
\begin{multicols}{2}
  \item Sustitución por sinonimia
  \item Repetición
\end{multicols}
\end{enumerate}
\vspace{0.15cm}
\item En "Compré un auto. El vehículo es rojo", "el vehículo" es:
\begin{enumerate}[label=\textbf{\Alph*)}]
\begin{multicols}{2}
  \item Sustitución por hiperonimia
  \item Repetición
\end{multicols}
\end{enumerate}
\vspace{0.15cm}
\item En "Me gustan las rosas. Las flores son hermosas", "las flores" es:
\begin{enumerate}[label=\textbf{\Alph*)}]
\begin{multicols}{2}
  \item Sustitución por hiperonimia
  \item Repetición
\end{multicols}
\end{enumerate}
\vspace{0.15cm}
\item La cohesión léxica se logra mediante:
\begin{enumerate}[label=\textbf{\Alph*)}]
\begin{multicols}{2}
  \item Repetición, sinonimia, antonimia, hiperonimia
  \item Solo conectores
  \item Solo puntuación
  \item Solo ortografía
  \item Solo caligrafía
  \item Solo sintaxis
\end{multicols}
\end{enumerate}
\vspace{0.15cm}
\item La cohesión gramatical se logra mediante:
\begin{enumerate}[label=\textbf{\Alph*)}]
\begin{multicols}{2}
  \item Pronombres, elipsis, determinantes, concordancia
  \item Solo conectores
  \item Solo puntuación
  \item Solo ortografía
  \item Solo caligrafía
  \item Solo sintaxis
\end{multicols}
\end{enumerate}
\vspace{0.15cm}
\item Los conectores son:
\begin{enumerate}[label=\textbf{\Alph*)}]
\begin{multicols}{2}
  \item Signos de puntuación
  \item Palabras sin significado
  \item Tipos de oraciones
  \item Tipos de párrafos
  \item Tipos de textos
\end{multicols}
\end{enumerate}
\vspace{0.15cm}
\item Los conectores aditivos indican:
\begin{enumerate}[label=\textbf{\Alph*)}]
\begin{multicols}{2}
  \item Suma o adición de ideas
  \item Tiempo
\end{multicols}
\end{enumerate}
\vspace{0.15cm}
\item Los conectores adversativos indican:
\begin{enumerate}[label=\textbf{\Alph*)}]
\begin{multicols}{2}
  \item Suma de ideas
  \item Tiempo
\end{multicols}
\end{enumerate}
\vspace{0.15cm}
\item Los conectores causales indican:
\begin{enumerate}[label=\textbf{\Alph*)}]
\begin{multicols}{2}
  \item Suma de ideas
  \item Tiempo
\end{multicols}
\end{enumerate}
\vspace{0.15cm}
\item Los conectores conclusivos indican:
\begin{enumerate}[label=\textbf{\Alph*)}]
\begin{multicols}{2}
  \item Suma de ideas
  \item Tiempo
\end{multicols}
\end{enumerate}
\vspace{0.15cm}
\item Los conectores temporales indican:
\begin{enumerate}[label=\textbf{\Alph*)}]
\begin{multicols}{2}
  \item Orden temporal
  \item Suma de ideas
\end{multicols}
\end{enumerate}
\vspace{0.15cm}
\item Los conectores condicionales indican:
\begin{enumerate}[label=\textbf{\Alph*)}]
\begin{multicols}{2}
  \item Suma de ideas
  \item Tiempo
\end{multicols}
\end{enumerate}
\vspace{0.15cm}
\item "Y" es un conector:
\begin{enumerate}[label=\textbf{\Alph*)}]
\begin{multicols}{2}
  \item Temporal
\end{multicols}
\end{enumerate}
\vspace{0.15cm}
\item "Pero" es un conector:
\begin{enumerate}[label=\textbf{\Alph*)}]
\begin{multicols}{2}
  \item Temporal
\end{multicols}
\end{enumerate}
\vspace{0.15cm}
\item "Porque" es un conector:
\begin{enumerate}[label=\textbf{\Alph*)}]
\begin{multicols}{2}
  \item Temporal
\end{multicols}
\end{enumerate}
\vspace{0.15cm}
\item "Por lo tanto" es un conector:
\begin{enumerate}[label=\textbf{\Alph*)}]
\begin{multicols}{2}
  \item Temporal
\end{multicols}
\end{enumerate}
\vspace{0.15cm}
\end{enumerate}\n
\subsubsection*{SECCIÓN IV: IMPLICATURA Y PRESUPOSICIÓN (151–200)}

\begin{enumerate}[label=\textbf{\arabic*.}]
\setcounter{enumi}{150}
\item La implicatura es:
\begin{enumerate}[label=\textbf{\Alph*)}]
\begin{multicols}{2}
  \item La estructura gramatical
  \item La ortografía del texto
  \item La puntuación del texto
  \item La caligrafía del texto
\end{multicols}
\end{enumerate}
\vspace{0.15cm}
\item La implicatura conversacional se basa en:
\begin{enumerate}[label=\textbf{\Alph*)}]
\begin{multicols}{2}
  \item La estructura gramatical
  \item La ortografía del texto
  \item La puntuación del texto
  \item La caligrafía del texto
\end{multicols}
\end{enumerate}
\vspace{0.15cm}
\item La implicatura convencional se basa en:
\begin{enumerate}[label=\textbf{\Alph*)}]
\begin{multicols}{2}
  \item Las máximas de Grice
  \item La intención del hablante
\end{multicols}
\end{enumerate}
\vspace{0.15cm}
\item En "Juan es alto pero no sabe jugar baloncesto", la implicatura es:
\begin{enumerate}[label=\textbf{\Alph*)}]
\begin{multicols}{2}
  \item Que los altos suelen saber jugar baloncesto
  \item Que Juan es bajo
  \item Que Juan sabe jugar baloncesto
  \item Que Juan no es alto
  \item Que Juan no juega baloncesto
  \item Que Juan es alto
\end{multicols}
\end{enumerate}
\vspace{0.15cm}
\item En "Algunos estudiantes aprobaron", la implicatura es:
\begin{enumerate}[label=\textbf{\Alph*)}]
\begin{multicols}{2}
  \item Que no todos aprobaron
  \item Que todos aprobaron
  \item Que ninguno aprobó
  \item Que la mayoría aprobó
  \item Que pocos aprobaron
  \item Que todos reprobaron
\end{multicols}
\end{enumerate}
\vspace{0.15cm}
\item La presuposición es:
\begin{enumerate}[label=\textbf{\Alph*)}]
\begin{multicols}{2}
  \item La información que se da por supuesta en un enunciado
  \item La información nueva que se aporta
  \item La intención del emisor
\end{multicols}
\end{enumerate}
\vspace{0.15cm}
\item En "El rey de Francia es calvo", la presuposición es:
\begin{enumerate}[label=\textbf{\Alph*)}]
\begin{multicols}{2}
\end{multicols}
\end{enumerate}
\vspace{0.15cm}
\item En "Juan dejó de fumar", la presuposición es:
\begin{enumerate}[label=\textbf{\Alph*)}]
\begin{multicols}{2}
  \item Que Juan fumaba antes
  \item Que Juan nunca fumó
  \item Que Juan fuma ahora
  \item Que Juan no fuma
  \item Que Juan fumará
  \item Que Juan odia fumar
\end{multicols}
\end{enumerate}
\vspace{0.15cm}
\item En "María se arrepiente de haber mentido", la presuposición es:
\begin{enumerate}[label=\textbf{\Alph*)}]
\begin{multicols}{2}
  \item Que María mintió
  \item Que María no mintió
  \item Que María dice la verdad
  \item Que María no se arrepiente
  \item Que María mentirá
  \item Que María odia mentir
\end{multicols}
\end{enumerate}
\vspace{0.15cm}
\item En "El actual presidente de El Salvador nació en San Salvador", la presuposición es:
\begin{enumerate}[label=\textbf{\Alph*)}]
\begin{multicols}{2}
  \item Que no existe presidente
  \item Que el presidente nació en otro lugar
  \item Que el presidente es falso
  \item Que el presidente no nació
  \item Que el presidente es extranjero
\end{multicols}
\end{enumerate}
\vspace{0.15cm}
\item La presuposición existencial se refiere a:
\begin{enumerate}[label=\textbf{\Alph*)}]
\begin{multicols}{2}
  \item La existencia de un referente
  \item La estructura gramatical
  \item La ortografía del texto
  \item La puntuación del texto
  \item La caligrafía del texto
\end{multicols}
\end{enumerate}
\vspace{0.15cm}
\item La presuposición factiva se refiere a:
\begin{enumerate}[label=\textbf{\Alph*)}]
\begin{multicols}{2}
  \item La verdad de una información presentada como hecho
  \item La existencia de un referente
  \item La estructura gramatical
  \item La ortografía del texto
  \item La puntuación del texto
\end{multicols}
\end{enumerate}
\vspace{0.15cm}
\item La presuposición léxica se refiere a:
\begin{enumerate}[label=\textbf{\Alph*)}]
\begin{multicols}{2}
  \item La información presupuesta por el significado de una palabra
  \item La existencia de un referente
  \item La estructura gramatical
  \item La ortografía del texto
  \item La puntuación del texto
\end{multicols}
\end{enumerate}
\vspace{0.15cm}
\item La presuposición estructural se refiere a:
\begin{enumerate}[label=\textbf{\Alph*)}]
\begin{multicols}{2}
  \item La información presupuesta por la estructura gramatical
  \item La existencia de un referente
  \item La ortografía del texto
  \item La puntuación del texto
  \item La caligrafía del texto
\end{multicols}
\end{enumerate}
\vspace{0.15cm}
\item La presuposición no se cancela fácilmente porque:
\begin{enumerate}[label=\textbf{\Alph*)}]
\begin{multicols}{2}
\end{multicols}
\end{enumerate}
\vspace{0.15cm}
\item La implicatura sí se puede cancelar porque:
\begin{enumerate}[label=\textbf{\Alph*)}]
\begin{multicols}{2}
\end{multicols}
\end{enumerate}
\vspace{0.15cm}
\item En "Pedro es soltero", la presuposición es:
\begin{enumerate}[label=\textbf{\Alph*)}]
\begin{multicols}{2}
  \item Que Pedro no está casado
  \item Que Pedro está casado
  \item Que Pedro es viudo
  \item Que Pedro es divorciado
  \item Que Pedro tiene novia
  \item Que Pedro no tiene novia
\end{multicols}
\end{enumerate}
\vspace{0.15cm}
\item En "El libro que compré es interesante", la presuposición es:
\begin{enumerate}[label=\textbf{\Alph*)}]
\begin{multicols}{2}
  \item Que compré un libro
  \item Que no compré un libro
  \item Que el libro es aburrido
  \item Que el libro es nuevo
  \item Que el libro es viejo
  \item Que el libro es caro
\end{multicols}
\end{enumerate}
\vspace{0.15cm}
\item En "Lamento que hayas llegado tarde", la presuposición es:
\begin{enumerate}[label=\textbf{\Alph*)}]
\begin{multicols}{2}
  \item Que llegaste tarde
  \item Que llegaste temprano
  \item Que no llegaste
  \item Que llegarás tarde
  \item Que llegas tarde
  \item Que nunca llegas
\end{multicols}
\end{enumerate}
\vspace{0.15cm}
\item En "Es triste que no haya venido", la presuposición es:
\begin{enumerate}[label=\textbf{\Alph*)}]
\begin{multicols}{2}
  \item Que no vino
  \item Que vino
  \item Que vendrá
  \item Que viene
  \item Que nunca viene
  \item Que siempre viene
\end{multicols}
\end{enumerate}
\vspace{0.15cm}
\item La máxima de cantidad se viola cuando:
\begin{enumerate}[label=\textbf{\Alph*)}]
\begin{multicols}{2}
  \item Se da más o menos información de la necesaria
  \item Se dice la verdad
  \item Se es relevante
  \item Se es claro
  \item Se es cortés
  \item Se es breve
\end{multicols}
\end{enumerate}
\vspace{0.15cm}
\item La máxima de calidad se viola cuando:
\begin{enumerate}[label=\textbf{\Alph*)}]
\begin{multicols}{2}
  \item Se dice algo falso o sin evidencia
  \item Se da la información necesaria
  \item Se es relevante
  \item Se es claro
  \item Se es cortés
  \item Se es breve
\end{multicols}
\end{enumerate}
\vspace{0.15cm}
\item La máxima de relación se viola cuando:
\begin{enumerate}[label=\textbf{\Alph*)}]
\begin{multicols}{2}
  \item Se dice algo irrelevante
  \item Se da la información necesaria
  \item Se dice la verdad
  \item Se es claro
  \item Se es cortés
  \item Se es breve
\end{multicols}
\end{enumerate}
\vspace{0.15cm}
\item La máxima de manera se viola cuando:
\begin{enumerate}[label=\textbf{\Alph*)}]
\begin{multicols}{2}
  \item Se es ambiguo, oscuro o desordenado
  \item Se da la información necesaria
  \item Se dice la verdad
  \item Se es relevante
  \item Se es cortés
  \item Se es breve
\end{multicols}
\end{enumerate}
\vspace{0.15cm}
\item El "flouting" de una máxima consiste en:
\begin{enumerate}[label=\textbf{\Alph*)}]
\begin{multicols}{2}
  \item Violarla deliberadamente para generar una implicatura
  \item Respetarla siempre
  \item Ignorarla siempre
  \item Repetirla
\end{multicols}
\end{enumerate}
\vspace{0.15cm}
\item En "¿Me prestas tu coche? Es que mi suegra está de visita", la implicatura es:
\begin{enumerate}[label=\textbf{\Alph*)}]
\begin{multicols}{2}
  \item Que el hablante no quiere estar con su suegra
  \item Que el hablante quiere estar con su suegra
  \item Que el hablante no tiene coche
  \item Que el hablante tiene coche
  \item Que el hablante odia a su suegra
  \item Que el hablante ama a su suegra
\end{multicols}
\end{enumerate}
\vspace{0.15cm}
\item En "A: ¿Vamos al cine? B: Tengo mucho trabajo", la implicatura es:
\begin{enumerate}[label=\textbf{\Alph*)}]
\begin{multicols}{2}
\end{multicols}
\end{enumerate}
\vspace{0.15cm}
\item En "A: ¿Te gustó mi comida? B: Nunca había probado algo así", la implicatura es:
\begin{enumerate}[label=\textbf{\Alph*)}]
\begin{multicols}{2}
  \item Que la comida no le gustó
  \item Que la comida le encantó
  \item Que la comida estaba deliciosa
  \item Que la comida era normal
  \item Que la comida era exquisita
  \item Que la comida era la mejor
\end{multicols}
\end{enumerate}
\vspace{0.15cm}
\item La ironía es un tipo de:
\begin{enumerate}[label=\textbf{\Alph*)}]
\begin{multicols}{2}
  \item Implicatura
  \item Presuposición
\end{multicols}
\end{enumerate}
\vspace{0.15cm}
\item La metáfora es un tipo de:
\begin{enumerate}[label=\textbf{\Alph*)}]
\begin{multicols}{2}
  \item Implicatura
  \item Presuposición
\end{multicols}
\end{enumerate}
\vspace{0.15cm}
\item La metonimia es un tipo de:
\begin{enumerate}[label=\textbf{\Alph*)}]
\begin{multicols}{2}
  \item Implicatura
  \item Presuposición
\end{multicols}
\end{enumerate}
\vspace{0.15cm}
\item En "Leí a Cervantes", hay:
\begin{enumerate}[label=\textbf{\Alph*)}]
\begin{multicols}{2}
  \item Metáfora
  \item Ironía
  \item Hipérbole
  \item Presuposición
\end{multicols}
\end{enumerate}
\vspace{0.15cm}
\item En "Tiene una pluma excelente", hay:
\begin{enumerate}[label=\textbf{\Alph*)}]
\begin{multicols}{2}
  \item Metáfora
  \item Ironía
  \item Hipérbole
  \item Presuposición
\end{multicols}
\end{enumerate}
\vspace{0.15cm}
\item En "¡Eres un genio!" dicho a alguien que cometió un error, hay:
\begin{enumerate}[label=\textbf{\Alph*)}]
\begin{multicols}{2}
  \item Ironía
  \item Metáfora
  \item Metonimia
  \item Hipérbole
  \item Presuposición
\end{multicols}
\end{enumerate}
\vspace{0.15cm}
\item En "Te lo he dicho un millón de veces", hay:
\begin{enumerate}[label=\textbf{\Alph*)}]
\begin{multicols}{2}
  \item Hipérbole
  \item Metáfora
  \item Metonimia
  \item Ironía
  \item Presuposición
\end{multicols}
\end{enumerate}
\vspace{0.15cm}
\item La hipérbole es:
\begin{enumerate}[label=\textbf{\Alph*)}]
\begin{multicols}{2}
  \item Una exageración
  \item Una comparación
  \item Una sustitución
  \item Una omisión
  \item Una repetición
  \item Una negación
\end{multicols}
\end{enumerate}
\vspace{0.15cm}
\item La metáfora es:
\begin{enumerate}[label=\textbf{\Alph*)}]
\begin{multicols}{2}
  \item Una comparación implícita
  \item Una exageración
  \item Una sustitución
  \item Una omisión
  \item Una repetición
  \item Una negación
\end{multicols}
\end{enumerate}
\vspace{0.15cm}
\item La metonimia es:
\begin{enumerate}[label=\textbf{\Alph*)}]
\begin{multicols}{2}
  \item Una sustitución basada en una relación de contigüidad
  \item Una comparación implícita
  \item Una exageración
  \item Una omisión
  \item Una repetición
  \item Una negación
\end{multicols}
\end{enumerate}
\vspace{0.15cm}
\item El sarcasmo es:
\begin{enumerate}[label=\textbf{\Alph*)}]
\begin{multicols}{2}
  \item Una ironía mordaz o cruel
  \item Una metáfora
  \item Una metonimia
  \item Una hipérbole
  \item Una presuposición
  \item Una deixis
\end{multicols}
\end{enumerate}
\vspace{0.15cm}
\item La atenuación es:
\begin{enumerate}[label=\textbf{\Alph*)}]
\begin{multicols}{2}
  \item Una estrategia para suavizar la fuerza ilocutiva
  \item Una estrategia para aumentar la fuerza ilocutiva
  \item Una estrategia para ignorar la fuerza ilocutiva
  \item Una estrategia para confundir
  \item Una estrategia para alargar
  \item Una estrategia para cambiar de tema
\end{multicols}
\end{enumerate}
\vspace{0.15cm}
\item La intensificación es:
\begin{enumerate}[label=\textbf{\Alph*)}]
\begin{multicols}{2}
  \item Una estrategia para aumentar la fuerza ilocutiva
  \item Una estrategia para suavizar la fuerza ilocutiva
  \item Una estrategia para ignorar la fuerza ilocutiva
  \item Una estrategia para confundir
  \item Una estrategia para alargar
  \item Una estrategia para cambiar de tema
\end{multicols}
\end{enumerate}
\vspace{0.15cm}
\item "¿Podrías hacerme un favor?" es una estrategia de:
\begin{enumerate}[label=\textbf{\Alph*)}]
\begin{multicols}{2}
  \item Intensificación
  \item Ironía
  \item Sarcasmo
\end{multicols}
\end{enumerate}
\vspace{0.15cm}
\item "¡Hazlo ahora mismo!" es una estrategia de:
\begin{enumerate}[label=\textbf{\Alph*)}]
\begin{multicols}{2}
  \item Intensificación
  \item Ironía
  \item Sarcasmo
\end{multicols}
\end{enumerate}
\vspace{0.15cm}
\item La cortesía positiva busca:
\begin{enumerate}[label=\textbf{\Alph*)}]
\begin{multicols}{2}
  \item Mostrar aprecio y solidaridad
  \item Mostrar distancia y respeto
  \item Ofender
  \item Ignorar
\end{multicols}
\end{enumerate}
\vspace{0.15cm}
\item La cortesía negativa busca:
\begin{enumerate}[label=\textbf{\Alph*)}]
\begin{multicols}{2}
  \item Mostrar respeto y evitar imposiciones
  \item Mostrar aprecio y solidaridad
  \item Ofender
  \item Ignorar
\end{multicols}
\end{enumerate}
\vspace{0.15cm}
\item "Eres muy inteligente, seguro puedes ayudarme" es una estrategia de:
\begin{enumerate}[label=\textbf{\Alph*)}]
\begin{multicols}{2}
  \item Ironía
  \item Sarcasmo
\end{multicols}
\end{enumerate}
\vspace{0.15cm}
\item "No quisiera molestarte, pero ¿podrías...?" es una estrategia de:
\begin{enumerate}[label=\textbf{\Alph*)}]
\begin{multicols}{2}
  \item Ironía
  \item Sarcasmo
\end{multicols}
\end{enumerate}
\vspace{0.15cm}
\item La imagen negativa se refiere a:
\begin{enumerate}[label=\textbf{\Alph*)}]
\begin{multicols}{2}
\end{multicols}
\end{enumerate}
\vspace{0.15cm}
\item La imagen positiva se refiere a:
\begin{enumerate}[label=\textbf{\Alph*)}]
\begin{multicols}{2}
\end{multicols}
\end{enumerate}
\vspace{0.15cm}
\item Un acto de habla que amenaza la imagen es:
\begin{enumerate}[label=\textbf{\Alph*)}]
\begin{multicols}{2}
  \item Una orden, una crítica o una disculpa
  \item Un cumplido
  \item Un agradecimiento
  \item Una felicitación
  \item Un saludo
  \item Una despedida ---
\end{multicols}
\end{enumerate}
\vspace{0.15cm}
\end{enumerate}\n
\subsubsection*{SECCIÓN V: PRAGMÁTICA DEL DISCURSO Y CONVERSACIÓN (201–250)}

\begin{enumerate}[label=\textbf{\arabic*.}]
\setcounter{enumi}{200}
\item El análisis del discurso estudia:
\begin{enumerate}[label=\textbf{\Alph*)}]
\begin{multicols}{2}
  \item La organización del lenguaje en unidades mayores que la oración
  \item La estructura de las oraciones
  \item La ortografía del texto
  \item La puntuación del texto
  \item La caligrafía del texto
\end{multicols}
\end{enumerate}
\vspace{0.15cm}
\item El análisis de la conversación estudia:
\begin{enumerate}[label=\textbf{\Alph*)}]
\begin{multicols}{2}
  \item La organización secuencial de la interacción verbal
  \item La estructura gramatical
  \item La ortografía del texto
  \item La puntuación del texto
  \item La caligrafía del texto
\end{multicols}
\end{enumerate}
\vspace{0.15cm}
\item Los turnos de habla son:
\begin{enumerate}[label=\textbf{\Alph*)}]
\begin{multicols}{2}
  \item Las intervenciones de cada participante en una conversación
  \item Las oraciones gramaticales
  \item Las palabras aisladas
  \item Los párrafos
  \item Los textos
  \item Los sintagmas
\end{multicols}
\end{enumerate}
\vspace{0.15cm}
\item La toma de turno se refiere a:
\begin{enumerate}[label=\textbf{\Alph*)}]
\begin{multicols}{2}
  \item La estructura gramatical
  \item La ortografía del texto
  \item La puntuación del texto
  \item La caligrafía del texto
\end{multicols}
\end{enumerate}
\vspace{0.15cm}
\item Los pares adyacentes son:
\begin{enumerate}[label=\textbf{\Alph*)}]
\begin{multicols}{2}
  \item Secuencias de tres turnos
  \item Secuencias de cuatro turnos
  \item Secuencias de un solo turno
  \item Secuencias sin relación
  \item Secuencias gramaticales
\end{multicols}
\end{enumerate}
\vspace{0.15cm}
\item "¿Cómo estás?" - "Bien, gracias" es un ejemplo de:
\begin{enumerate}[label=\textbf{\Alph*)}]
\begin{multicols}{2}
  \item Par adyacente
  \item Turno único
  \item Secuencia de tres
  \item Secuencia de cuatro
  \item Secuencia sin relación
  \item Secuencia gramatical
\end{multicols}
\end{enumerate}
\vspace{0.15cm}
\item La reparación en la conversación consiste en:
\begin{enumerate}[label=\textbf{\Alph*)}]
\begin{multicols}{2}
  \item Ignorar un error
  \item Repetir un error
\end{multicols}
\end{enumerate}
\vspace{0.15cm}
\item La superposición en la conversación se produce cuando:
\begin{enumerate}[label=\textbf{\Alph*)}]
\begin{multicols}{2}
  \item Un participante habla solo
  \item No habla nadie
  \item Se guarda silencio
  \item Se cambia de tema
  \item Se alarga la conversación
\end{multicols}
\end{enumerate}
\vspace{0.15cm}
\item La pragmática textual estudia:
\begin{enumerate}[label=\textbf{\Alph*)}]
\begin{multicols}{2}
  \item La organización y uso del lenguaje en textos
  \item La estructura de las oraciones
  \item La ortografía del texto
  \item La puntuación del texto
  \item La caligrafía del texto
\end{multicols}
\end{enumerate}
\vspace{0.15cm}
\item La coherencia pragmática se refiere a:
\begin{enumerate}[label=\textbf{\Alph*)}]
\begin{multicols}{2}
  \item La adecuación del texto al contexto comunicativo
  \item La unidad temática
  \item La cohesión gramatical
  \item La ortografía del texto
  \item La puntuación del texto
  \item La caligrafía del texto
\end{multicols}
\end{enumerate}
\vspace{0.15cm}
\item La relevancia en pragmática se refiere a:
\begin{enumerate}[label=\textbf{\Alph*)}]
\begin{multicols}{2}
  \item La pertinencia de la información en el contexto
  \item La estructura gramatical
  \item La ortografía del texto
  \item La puntuación del texto
  \item La caligrafía del texto
\end{multicols}
\end{enumerate}
\vspace{0.15cm}
\item La teoría de la relevancia fue propuesta por:
\begin{enumerate}[label=\textbf{\Alph*)}]
\begin{multicols}{2}
  \item Sperber y Wilson
  \item Grice
  \item Searle
  \item Saussure
\end{multicols}
\end{enumerate}
\vspace{0.15cm}
\item Según la teoría de la relevancia, la comunicación busca:
\begin{enumerate}[label=\textbf{\Alph*)}]
\begin{multicols}{2}
  \item Maximizar la relevancia con el mínimo esfuerzo
  \item Maximizar el esfuerzo
  \item Minimizar la relevancia
  \item Ignorar el contexto
  \item Ignorar la intención
  \item Ignorar el efecto
\end{multicols}
\end{enumerate}
\vspace{0.15cm}
\item El contexto cognitivo en la teoría de la relevancia se refiere a:
\begin{enumerate}[label=\textbf{\Alph*)}]
\begin{multicols}{2}
\end{multicols}
\end{enumerate}
\vspace{0.15cm}
\item El efecto cognitivo en la teoría de la relevancia se refiere a:
\begin{enumerate}[label=\textbf{\Alph*)}]
\begin{multicols}{2}
  \item La modificación de los supuestos del oyente
  \item La estructura gramatical
  \item La ortografía del texto
  \item La puntuación del texto
  \item La caligrafía del texto
\end{multicols}
\end{enumerate}
\vspace{0.15cm}
\item El esfuerzo de procesamiento en la teoría de la relevancia se refiere a:
\begin{enumerate}[label=\textbf{\Alph*)}]
\begin{multicols}{2}
  \item La estructura gramatical
  \item La ortografía del texto
  \item La puntuación del texto
  \item La caligrafía del texto
\end{multicols}
\end{enumerate}
\vspace{0.15cm}
\item La comunicación ostensiva-inferencial es:
\begin{enumerate}[label=\textbf{\Alph*)}]
\begin{multicols}{2}
  \item La que combina la intención del hablante con la inferencia del oyente
  \item La que solo tiene intención
  \item La que solo tiene inferencia
  \item La que no tiene intención
  \item La que no tiene inferencia
  \item La que no tiene contexto
\end{multicols}
\end{enumerate}
\vspace{0.15cm}
\item La ostensión es:
\begin{enumerate}[label=\textbf{\Alph*)}]
\begin{multicols}{2}
  \item La manifestación de la intención comunicativa
  \item La inferencia del oyente
  \item La estructura gramatical
  \item La ortografía del texto
  \item La puntuación del texto
\end{multicols}
\end{enumerate}
\vspace{0.15cm}
\item La inferencia es:
\begin{enumerate}[label=\textbf{\Alph*)}]
\begin{multicols}{2}
  \item La manifestación de la intención
  \item La estructura gramatical
  \item La ortografía del texto
  \item La puntuación del texto
\end{multicols}
\end{enumerate}
\vspace{0.15cm}
\item El entorno cognitivo es:
\begin{enumerate}[label=\textbf{\Alph*)}]
\begin{multicols}{2}
\end{multicols}
\end{enumerate}
\vspace{0.15cm}
\item La teoría de la cortesía fue propuesta por:
\begin{enumerate}[label=\textbf{\Alph*)}]
\begin{multicols}{2}
  \item Grice
  \item Searle
  \item Sperber y Wilson
  \item Saussure
\end{multicols}
\end{enumerate}
\vspace{0.15cm}
\item La imagen pública o "face" es:
\begin{enumerate}[label=\textbf{\Alph*)}]
\begin{multicols}{2}
  \item La imagen social que cada persona proyecta
  \item La imagen física
  \item La imagen digital
  \item La imagen profesional
  \item La imagen académica
  \item La imagen personal
\end{multicols}
\end{enumerate}
\vspace{0.15cm}
\item La cortesía positiva se orienta hacia:
\begin{enumerate}[label=\textbf{\Alph*)}]
\begin{multicols}{2}
  \item La imagen positiva del interlocutor
  \item La imagen negativa del interlocutor
  \item La imagen física
  \item La imagen digital
  \item La imagen profesional
  \item La imagen académica
\end{multicols}
\end{enumerate}
\vspace{0.15cm}
\item La cortesía negativa se orienta hacia:
\begin{enumerate}[label=\textbf{\Alph*)}]
\begin{multicols}{2}
  \item La imagen negativa del interlocutor
  \item La imagen positiva del interlocutor
  \item La imagen física
  \item La imagen digital
  \item La imagen profesional
  \item La imagen académica
\end{multicols}
\end{enumerate}
\vspace{0.15cm}
\item La descortesía es:
\begin{enumerate}[label=\textbf{\Alph*)}]
\begin{multicols}{2}
  \item La violación deliberada de las normas de cortesía
  \item La indiferencia a las normas
  \item La ignorancia de las normas
  \item La confusión de las normas
  \item La repetición de las normas
\end{multicols}
\end{enumerate}
\vspace{0.15cm}
\item La pragmática intercultural estudia:
\begin{enumerate}[label=\textbf{\Alph*)}]
\begin{multicols}{2}
  \item Las diferencias en el uso del lenguaje entre culturas
  \item Las diferencias gramaticales entre lenguas
  \item Las diferencias ortográficas entre lenguas
  \item Las diferencias fonéticas entre lenguas
  \item Las diferencias morfológicas entre lenguas
  \item Las diferencias sintácticas entre lenguas
\end{multicols}
\end{enumerate}
\vspace{0.15cm}
\item La pragmática del español estudia:
\begin{enumerate}[label=\textbf{\Alph*)}]
\begin{multicols}{2}
  \item La gramática del español
  \item La ortografía del español
  \item La fonética del español
  \item La morfología del español
  \item La sintaxis del español
\end{multicols}
\end{enumerate}
\vspace{0.15cm}
\item El "tú" y el "usted" en español son ejemplos de:
\begin{enumerate}[label=\textbf{\Alph*)}]
\begin{multicols}{2}
\end{multicols}
\end{enumerate}
\vspace{0.15cm}
\item El uso de "usted" indica:
\begin{enumerate}[label=\textbf{\Alph*)}]
\begin{multicols}{2}
  \item Ironía
  \item Sarcasmo
\end{multicols}
\end{enumerate}
\vspace{0.15cm}
\item El uso de "tú" indica:
\begin{enumerate}[label=\textbf{\Alph*)}]
\begin{multicols}{2}
  \item Ironía
  \item Sarcasmo
\end{multicols}
\end{enumerate}
\vspace{0.15cm}
\item El voseo es:
\begin{enumerate}[label=\textbf{\Alph*)}]
\begin{multicols}{2}
\end{multicols}
\end{enumerate}
\vspace{0.15cm}
\item La pragmática del discurso estudia:
\begin{enumerate}[label=\textbf{\Alph*)}]
\begin{multicols}{2}
  \item La organización del lenguaje en unidades mayores que la oración
  \item La estructura de las oraciones
  \item La ortografía del texto
  \item La puntuación del texto
  \item La caligrafía del texto
\end{multicols}
\end{enumerate}
\vspace{0.15cm}
\item El análisis de la conversación estudia:
\begin{enumerate}[label=\textbf{\Alph*)}]
\begin{multicols}{2}
  \item La organización secuencial de la interacción verbal
  \item La estructura gramatical
  \item La ortografía del texto
  \item La puntuación del texto
  \item La caligrafía del texto
\end{multicols}
\end{enumerate}
\vspace{0.15cm}
\item Los turnos de habla son:
\begin{enumerate}[label=\textbf{\Alph*)}]
\begin{multicols}{2}
  \item Las intervenciones de cada participante en una conversación
  \item Las oraciones gramaticales
  \item Las palabras aisladas
  \item Los párrafos
  \item Los textos
  \item Los sintagmas
\end{multicols}
\end{enumerate}
\vspace{0.15cm}
\item La toma de turno se refiere a:
\begin{enumerate}[label=\textbf{\Alph*)}]
\begin{multicols}{2}
  \item La estructura gramatical
  \item La ortografía del texto
  \item La puntuación del texto
  \item La caligrafía del texto
\end{multicols}
\end{enumerate}
\vspace{0.15cm}
\item Los pares adyacentes son:
\begin{enumerate}[label=\textbf{\Alph*)}]
\begin{multicols}{2}
  \item Secuencias de tres turnos
  \item Secuencias de cuatro turnos
  \item Secuencias de un solo turno
  \item Secuencias sin relación
  \item Secuencias gramaticales
\end{multicols}
\end{enumerate}
\vspace{0.15cm}
\item "¿Cómo estás?" - "Bien, gracias" es un ejemplo de:
\begin{enumerate}[label=\textbf{\Alph*)}]
\begin{multicols}{2}
  \item Par adyacente
  \item Turno único
  \item Secuencia de tres
  \item Secuencia de cuatro
  \item Secuencia sin relación
  \item Secuencia gramatical
\end{multicols}
\end{enumerate}
\vspace{0.15cm}
\item La reparación en la conversación consiste en:
\begin{enumerate}[label=\textbf{\Alph*)}]
\begin{multicols}{2}
  \item Ignorar un error
  \item Repetir un error
\end{multicols}
\end{enumerate}
\vspace{0.15cm}
\item La superposición en la conversación se produce cuando:
\begin{enumerate}[label=\textbf{\Alph*)}]
\begin{multicols}{2}
  \item Un participante habla solo
  \item No habla nadie
  \item Se guarda silencio
  \item Se cambia de tema
  \item Se alarga la conversación
\end{multicols}
\end{enumerate}
\vspace{0.15cm}
\item La pragmática textual estudia:
\begin{enumerate}[label=\textbf{\Alph*)}]
\begin{multicols}{2}
  \item La organización y uso del lenguaje en textos
  \item La estructura de las oraciones
  \item La ortografía del texto
  \item La puntuación del texto
  \item La caligrafía del texto
\end{multicols}
\end{enumerate}
\vspace{0.15cm}
\item La coherencia pragmática se refiere a:
\begin{enumerate}[label=\textbf{\Alph*)}]
\begin{multicols}{2}
  \item La adecuación del texto al contexto comunicativo
  \item La unidad temática
  \item La cohesión gramatical
  \item La ortografía del texto
  \item La puntuación del texto
  \item La caligrafía del texto
\end{multicols}
\end{enumerate}
\vspace{0.15cm}
\item La relevancia en pragmática se refiere a:
\begin{enumerate}[label=\textbf{\Alph*)}]
\begin{multicols}{2}
  \item La pertinencia de la información en el contexto
  \item La estructura gramatical
  \item La ortografía del texto
  \item La puntuación del texto
  \item La caligrafía del texto
\end{multicols}
\end{enumerate}
\vspace{0.15cm}
\item La teoría de la relevancia fue propuesta por:
\begin{enumerate}[label=\textbf{\Alph*)}]
\begin{multicols}{2}
  \item Sperber y Wilson
  \item Grice
  \item Searle
  \item Saussure
\end{multicols}
\end{enumerate}
\vspace{0.15cm}
\item Según la teoría de la relevancia, la comunicación busca:
\begin{enumerate}[label=\textbf{\Alph*)}]
\begin{multicols}{2}
  \item Maximizar la relevancia con el mínimo esfuerzo
  \item Maximizar el esfuerzo
  \item Minimizar la relevancia
  \item Ignorar el contexto
  \item Ignorar la intención
  \item Ignorar el efecto
\end{multicols}
\end{enumerate}
\vspace{0.15cm}
\item El contexto cognitivo en la teoría de la relevancia se refiere a:
\begin{enumerate}[label=\textbf{\Alph*)}]
\begin{multicols}{2}
\end{multicols}
\end{enumerate}
\vspace{0.15cm}
\item El efecto cognitivo en la teoría de la relevancia se refiere a:
\begin{enumerate}[label=\textbf{\Alph*)}]
\begin{multicols}{2}
  \item La modificación de los supuestos del oyente
  \item La estructura gramatical
  \item La ortografía del texto
  \item La puntuación del texto
  \item La caligrafía del texto
\end{multicols}
\end{enumerate}
\vspace{0.15cm}
\item El esfuerzo de procesamiento en la teoría de la relevancia se refiere a:
\begin{enumerate}[label=\textbf{\Alph*)}]
\begin{multicols}{2}
  \item La estructura gramatical
  \item La ortografía del texto
  \item La puntuación del texto
  \item La caligrafía del texto
\end{multicols}
\end{enumerate}
\vspace{0.15cm}
\item La comunicación ostensiva-inferencial es:
\begin{enumerate}[label=\textbf{\Alph*)}]
\begin{multicols}{2}
  \item La que combina la intención del hablante con la inferencia del oyente
  \item La que solo tiene intención
  \item La que solo tiene inferencia
  \item La que no tiene intención
  \item La que no tiene inferencia
  \item La que no tiene contexto
\end{multicols}
\end{enumerate}
\vspace{0.15cm}
\item La ostensión es:
\begin{enumerate}[label=\textbf{\Alph*)}]
\begin{multicols}{2}
  \item La manifestación de la intención comunicativa
  \item La inferencia del oyente
  \item La estructura gramatical
  \item La ortografía del texto
  \item La puntuación del texto
\end{multicols}
\end{enumerate}
\vspace{0.15cm}
\item La inferencia es:
\begin{enumerate}[label=\textbf{\Alph*)}]
\begin{multicols}{2}
  \item La manifestación de la intención
  \item La estructura gramatical
  \item La ortografía del texto
  \item La puntuación del texto ---
\end{multicols}
\end{enumerate}
\vspace{0.15cm}
\end{enumerate}\n
\subsubsection*{SECCIÓN VI: PRAGMÁTICA Y SOCIEDAD (251–300)}

\begin{enumerate}[label=\textbf{\arabic*.}]
\setcounter{enumi}{250}
\item La pragmática estudia el lenguaje:
\begin{enumerate}[label=\textbf{\Alph*)}]
\begin{multicols}{2}
  \item Sin intención
  \item Sin efecto
  \item Sin significado
\end{multicols}
\end{enumerate}
\vspace{0.15cm}
\item El significado pragmático depende de:
\begin{enumerate}[label=\textbf{\Alph*)}]
\begin{multicols}{2}
  \item Solo la gramática
  \item Solo la ortografía
  \item Solo la puntuación
  \item Solo la caligrafía
  \item Solo la sintaxis
\end{multicols}
\end{enumerate}
\vspace{0.15cm}
\item La comunicación eficaz requiere:
\begin{enumerate}[label=\textbf{\Alph*)}]
\begin{multicols}{2}
  \item Ignorar el contexto
  \item Ignorar las intenciones
  \item Ignorar al interlocutor
  \item Ignorar el mensaje
  \item Ignorar el canal
\end{multicols}
\end{enumerate}
\vspace{0.15cm}
\item Los malentendidos se producen cuando:
\begin{enumerate}[label=\textbf{\Alph*)}]
\begin{multicols}{2}
  \item No se comparte el contexto o la intención
  \item Se comparte el contexto
  \item Se comparte la intención
  \item Se comparte el mensaje
  \item Se comparte el canal
  \item Se comparte el código
\end{multicols}
\end{enumerate}
\vspace{0.15cm}
\item La pragmática es importante en la enseñanza de lenguas porque:
\begin{enumerate}[label=\textbf{\Alph*)}]
\begin{multicols}{2}
\end{multicols}
\end{enumerate}
\vspace{0.15cm}
\item La competencia pragmática es:
\begin{enumerate}[label=\textbf{\Alph*)}]
\begin{multicols}{2}
  \item La capacidad de usar el lenguaje adecuadamente en contexto
  \item La capacidad de memorizar gramática
  \item La capacidad de memorizar ortografía
  \item La capacidad de memorizar puntuación
  \item La capacidad de memorizar caligrafía
  \item La capacidad de memorizar sintaxis
\end{multicols}
\end{enumerate}
\vspace{0.15cm}
\item La competencia pragmática incluye:
\begin{enumerate}[label=\textbf{\Alph*)}]
\begin{multicols}{2}
  \item La capacidad de interpretar y producir actos de habla
  \item Solo la gramática
  \item Solo la ortografía
  \item Solo la puntuación
  \item Solo la caligrafía
  \item Solo la sintaxis
\end{multicols}
\end{enumerate}
\vspace{0.15cm}
\item La pragmática intercultural es importante porque:
\begin{enumerate}[label=\textbf{\Alph*)}]
\begin{multicols}{2}
\end{multicols}
\end{enumerate}
\vspace{0.15cm}
\item La pragmática del español es importante porque:
\begin{enumerate}[label=\textbf{\Alph*)}]
\begin{multicols}{2}
\end{multicols}
\end{enumerate}
\vspace{0.15cm}
\item La pragmática en la traducción es importante porque:
\begin{enumerate}[label=\textbf{\Alph*)}]
\begin{multicols}{2}
\end{multicols}
\end{enumerate}
\vspace{0.15cm}
\item La pragmática en la publicidad es importante porque:
\begin{enumerate}[label=\textbf{\Alph*)}]
\begin{multicols}{2}
\end{multicols}
\end{enumerate}
\vspace{0.15cm}
\item La pragmática en la política es importante porque:
\begin{enumerate}[label=\textbf{\Alph*)}]
\begin{multicols}{2}
\end{multicols}
\end{enumerate}
\vspace{0.15cm}
\item La pragmática en la educación es importante porque:
\begin{enumerate}[label=\textbf{\Alph*)}]
\begin{multicols}{2}
\end{multicols}
\end{enumerate}
\vspace{0.15cm}
\item La pragmática en la vida cotidiana es importante porque:
\begin{enumerate}[label=\textbf{\Alph*)}]
\begin{multicols}{2}
\end{multicols}
\end{enumerate}
\vspace{0.15cm}
\item La pragmática estudia:
\begin{enumerate}[label=\textbf{\Alph*)}]
\begin{multicols}{2}
\end{multicols}
\end{enumerate}
\vspace{0.15cm}
\item La pragmática se relaciona con:
\begin{enumerate}[label=\textbf{\Alph*)}]
\begin{multicols}{2}
  \item La semántica, la sintaxis y la fonología
  \item Solo la semántica
  \item Solo la sintaxis
  \item Solo la fonología
  \item Solo la morfología
  \item Solo la ortografía
\end{multicols}
\end{enumerate}
\vspace{0.15cm}
\item La pragmática se diferencia de la semántica porque:
\begin{enumerate}[label=\textbf{\Alph*)}]
\begin{multicols}{2}
\end{multicols}
\end{enumerate}
\vspace{0.15cm}
\item La pragmática se diferencia de la sintaxis porque:
\begin{enumerate}[label=\textbf{\Alph*)}]
\begin{multicols}{2}
\end{multicols}
\end{enumerate}
\vspace{0.15cm}
\item La pragmática se diferencia de la fonología porque:
\begin{enumerate}[label=\textbf{\Alph*)}]
\begin{multicols}{2}
\end{multicols}
\end{enumerate}
\vspace{0.15cm}
\item La pragmática es una disciplina:
\begin{enumerate}[label=\textbf{\Alph*)}]
\begin{multicols}{2}
  \item Obsoleta
  \item Inútil
\end{multicols}
\end{enumerate}
\vspace{0.15cm}
\item La pragmática estudia:
\begin{enumerate}[label=\textbf{\Alph*)}]
\begin{multicols}{2}
  \item La relación entre el lenguaje y sus usuarios
  \item La relación entre las palabras
  \item La relación entre los sonidos
  \item La relación entre las letras
  \item La relación entre las sílabas
  \item La relación entre las oraciones
\end{multicols}
\end{enumerate}
\vspace{0.15cm}
\item La pragmática se interesa por:
\begin{enumerate}[label=\textbf{\Alph*)}]
\begin{multicols}{2}
  \item La intención del hablante y el efecto en el oyente
  \item Solo la estructura gramatical
  \item Solo el significado literal
  \item Solo la ortografía
  \item Solo la puntuación
  \item Solo la caligrafía
\end{multicols}
\end{enumerate}
\vspace{0.15cm}
\item La pragmática analiza:
\begin{enumerate}[label=\textbf{\Alph*)}]
\begin{multicols}{2}
  \item Los actos de habla, la deixis, la implicatura y la presuposición
  \item Solo los actos de habla
  \item Solo la deixis
  \item Solo la implicatura
  \item Solo la presuposición
  \item Solo la gramática
\end{multicols}
\end{enumerate}
\vspace{0.15cm}
\item La pragmática es esencial para:
\begin{enumerate}[label=\textbf{\Alph*)}]
\begin{multicols}{2}
  \item La comunicación eficaz
  \item La memorización de datos
  \item La repetición de información
  \item La copia de información
  \item La ignorancia de información
  \item La confusión de información
\end{multicols}
\end{enumerate}
\vspace{0.15cm}
\item La pragmática en la literatura es importante porque:
\begin{enumerate}[label=\textbf{\Alph*)}]
\begin{multicols}{2}
\end{multicols}
\end{enumerate}
\vspace{0.15cm}
\item La pragmática en el cine es importante porque:
\begin{enumerate}[label=\textbf{\Alph*)}]
\begin{multicols}{2}
\end{multicols}
\end{enumerate}
\vspace{0.15cm}
\item La pragmática en las redes sociales es importante porque:
\begin{enumerate}[label=\textbf{\Alph*)}]
\begin{multicols}{2}
\end{multicols}
\end{enumerate}
\vspace{0.15cm}
\item La pragmática en la inteligencia artificial es importante porque:
\begin{enumerate}[label=\textbf{\Alph*)}]
\begin{multicols}{2}
\end{multicols}
\end{enumerate}
\vspace{0.15cm}
\item La pragmática en la traducción automática es importante porque:
\begin{enumerate}[label=\textbf{\Alph*)}]
\begin{multicols}{2}
\end{multicols}
\end{enumerate}
\vspace{0.15cm}
\item La pragmática en la enseñanza de lenguas es importante porque:
\begin{enumerate}[label=\textbf{\Alph*)}]
\begin{multicols}{2}
\end{multicols}
\end{enumerate}
\vspace{0.15cm}
\item La pragmática en la evaluación de lenguas es importante porque:
\begin{enumerate}[label=\textbf{\Alph*)}]
\begin{multicols}{2}
\end{multicols}
\end{enumerate}
\vspace{0.15cm}
\item La pragmática en la investigación es importante porque:
\begin{enumerate}[label=\textbf{\Alph*)}]
\begin{multicols}{2}
\end{multicols}
\end{enumerate}
\vspace{0.15cm}
\item La pragmática en la psicología es importante porque:
\begin{enumerate}[label=\textbf{\Alph*)}]
\begin{multicols}{2}
\end{multicols}
\end{enumerate}
\vspace{0.15cm}
\item La pragmática en la sociología es importante porque:
\begin{enumerate}[label=\textbf{\Alph*)}]
\begin{multicols}{2}
\end{multicols}
\end{enumerate}
\vspace{0.15cm}
\item La pragmática en la antropología es importante porque:
\begin{enumerate}[label=\textbf{\Alph*)}]
\begin{multicols}{2}
\end{multicols}
\end{enumerate}
\vspace{0.15cm}
\item La pragmática en la filosofía es importante porque:
\begin{enumerate}[label=\textbf{\Alph*)}]
\begin{multicols}{2}
\end{multicols}
\end{enumerate}
\vspace{0.15cm}
\item La pragmática en la lingüística es importante porque:
\begin{enumerate}[label=\textbf{\Alph*)}]
\begin{multicols}{2}
\end{multicols}
\end{enumerate}
\vspace{0.15cm}
\item La pragmática en la comunicación es importante porque:
\begin{enumerate}[label=\textbf{\Alph*)}]
\begin{multicols}{2}
\end{multicols}
\end{enumerate}
\vspace{0.15cm}
\item La pragmática en la educación es importante porque:
\begin{enumerate}[label=\textbf{\Alph*)}]
\begin{multicols}{2}
\end{multicols}
\end{enumerate}
\vspace{0.15cm}
\item La pragmática en el trabajo es importante porque:
\begin{enumerate}[label=\textbf{\Alph*)}]
\begin{multicols}{2}
\end{multicols}
\end{enumerate}
\vspace{0.15cm}
\item La pragmática en la salud es importante porque:
\begin{enumerate}[label=\textbf{\Alph*)}]
\begin{multicols}{2}
\end{multicols}
\end{enumerate}
\vspace{0.15cm}
\item La pragmática en el derecho es importante porque:
\begin{enumerate}[label=\textbf{\Alph*)}]
\begin{multicols}{2}
\end{multicols}
\end{enumerate}
\vspace{0.15cm}
\item La pragmática en la política es importante porque:
\begin{enumerate}[label=\textbf{\Alph*)}]
\begin{multicols}{2}
\end{multicols}
\end{enumerate}
\vspace{0.15cm}
\item La pragmática en los medios es importante porque:
\begin{enumerate}[label=\textbf{\Alph*)}]
\begin{multicols}{2}
\end{multicols}
\end{enumerate}
\vspace{0.15cm}
\item La pragmática en la publicidad es importante porque:
\begin{enumerate}[label=\textbf{\Alph*)}]
\begin{multicols}{2}
\end{multicols}
\end{enumerate}
\vspace{0.15cm}
\item La pragmática en las relaciones interpersonales es importante porque:
\begin{enumerate}[label=\textbf{\Alph*)}]
\begin{multicols}{2}
\end{multicols}
\end{enumerate}
\vspace{0.15cm}
\item La pragmática en la familia es importante porque:
\begin{enumerate}[label=\textbf{\Alph*)}]
\begin{multicols}{2}
\end{multicols}
\end{enumerate}
\vspace{0.15cm}
\item La pragmática en la amistad es importante porque:
\begin{enumerate}[label=\textbf{\Alph*)}]
\begin{multicols}{2}
\end{multicols}
\end{enumerate}
\vspace{0.15cm}
\item La pragmática en el amor es importante porque:
\begin{enumerate}[label=\textbf{\Alph*)}]
\begin{multicols}{2}
\end{multicols}
\end{enumerate}
\vspace{0.15cm}
\item La pragmática es fundamental para:
\begin{enumerate}[label=\textbf{\Alph*)}]
\begin{multicols}{2}
  \item Memorizar gramática
  \item Memorizar ortografía
  \item Memorizar puntuación
  \item Memorizar caligrafía
  \item Memorizar sintaxis ---
\end{multicols}
\end{enumerate}
\vspace{0.15cm}
\end{enumerate}\n
\vspace{0.4cm}
\begin{tcolorbox}[colback=gray!5!white,colframe=gray!60!black,title=\textbf{SOLUCIONARIO -- 2.5 PRAGMÁTICA}]
\begin{multicols}{4}
\begin{enumerate}[label=\textbf{\arabic*.}]
  \item \textbf{A}
  \item \textbf{A}
  \item \textbf{A}
  \item \textbf{A}
  \item \textbf{A}
  \item \textbf{A}
  \item \textbf{A}
  \item \textbf{A}
  \item \textbf{A}
  \item \textbf{A}
  \item \textbf{A}
  \item \textbf{A}
  \item \textbf{A}
  \item \textbf{A}
  \item \textbf{A}
  \item \textbf{A}
  \item \textbf{A}
  \item \textbf{A}
  \item \textbf{A}
  \item \textbf{A}
  \item \textbf{A}
  \item \textbf{A}
  \item \textbf{A}
  \item \textbf{A}
  \item \textbf{A}
  \item \textbf{A}
  \item \textbf{A}
  \item \textbf{A}
  \item \textbf{A}
  \item \textbf{A}
  \item \textbf{A}
  \item \textbf{A}
  \item \textbf{A}
  \item \textbf{A}
  \item \textbf{A}
  \item \textbf{A}
  \item \textbf{A}
  \item \textbf{A}
  \item \textbf{A}
  \item \textbf{A}
  \item \textbf{A}
  \item \textbf{A}
  \item \textbf{A}
  \item \textbf{A}
  \item \textbf{A}
  \item \textbf{A}
  \item \textbf{A}
  \item \textbf{A}
  \item \textbf{A}
  \item \textbf{A}
  \item \textbf{A}
  \item \textbf{A}
  \item \textbf{A}
  \item \textbf{A}
  \item \textbf{A}
  \item \textbf{A}
  \item \textbf{A}
  \item \textbf{A}
  \item \textbf{A}
  \item \textbf{A}
  \item \textbf{A}
  \item \textbf{A}
  \item \textbf{A}
  \item \textbf{A}
  \item \textbf{A}
  \item \textbf{A}
  \item \textbf{A}
  \item \textbf{A}
  \item \textbf{A}
  \item \textbf{A}
  \item \textbf{A}
  \item \textbf{A}
  \item \textbf{A}
  \item \textbf{A}
  \item \textbf{A}
  \item \textbf{A}
  \item \textbf{A}
  \item \textbf{A}
  \item \textbf{A}
  \item \textbf{A}
  \item \textbf{A}
  \item \textbf{A}
  \item \textbf{A}
  \item \textbf{A}
  \item \textbf{A}
  \item \textbf{A}
  \item \textbf{A}
  \item \textbf{A}
  \item \textbf{A}
  \item \textbf{A}
  \item \textbf{A}
  \item \textbf{A}
  \item \textbf{A}
  \item \textbf{A}
  \item \textbf{A}
  \item \textbf{A}
  \item \textbf{A}
  \item \textbf{A}
  \item \textbf{A}
  \item \textbf{A}
  \item \textbf{A}
  \item \textbf{A}
  \item \textbf{A}
  \item \textbf{A}
  \item \textbf{A}
  \item \textbf{A}
  \item \textbf{A}
  \item \textbf{A}
  \item \textbf{A}
  \item \textbf{A}
  \item \textbf{A}
  \item \textbf{A}
  \item \textbf{A}
  \item \textbf{A}
  \item \textbf{A}
  \item \textbf{A}
  \item \textbf{A}
  \item \textbf{A}
  \item \textbf{A}
  \item \textbf{A}
  \item \textbf{A}
  \item \textbf{A}
  \item \textbf{A}
  \item \textbf{A}
  \item \textbf{A}
  \item \textbf{A}
  \item \textbf{A}
  \item \textbf{A}
  \item \textbf{A}
  \item \textbf{A}
  \item \textbf{A}
  \item \textbf{A}
  \item \textbf{A}
  \item \textbf{A}
  \item \textbf{A}
  \item \textbf{A}
  \item \textbf{A}
  \item \textbf{A}
  \item \textbf{A}
  \item \textbf{A}
  \item \textbf{A}
  \item \textbf{A}
  \item \textbf{A}
  \item \textbf{A}
  \item \textbf{A}
  \item \textbf{A}
  \item \textbf{A}
  \item \textbf{A}
  \item \textbf{A}
  \item \textbf{A}
  \item \textbf{A}
  \item \textbf{A}
  \item \textbf{A}
  \item \textbf{A}
  \item \textbf{A}
  \item \textbf{A}
  \item \textbf{A}
  \item \textbf{A}
  \item \textbf{A}
  \item \textbf{A}
  \item \textbf{A}
  \item \textbf{A}
  \item \textbf{A}
  \item \textbf{A}
  \item \textbf{A}
  \item \textbf{A}
  \item \textbf{A}
  \item \textbf{A}
  \item \textbf{A}
  \item \textbf{A}
  \item \textbf{A}
  \item \textbf{A}
  \item \textbf{A}
  \item \textbf{A}
  \item \textbf{A}
  \item \textbf{A}
  \item \textbf{A}
  \item \textbf{A}
  \item \textbf{A}
  \item \textbf{A}
  \item \textbf{A}
  \item \textbf{A}
  \item \textbf{A}
  \item \textbf{A}
  \item \textbf{A}
  \item \textbf{A}
  \item \textbf{A}
  \item \textbf{A}
  \item \textbf{A}
  \item \textbf{A}
  \item \textbf{A}
  \item \textbf{A}
  \item \textbf{A}
  \item \textbf{A}
  \item \textbf{A}
  \item \textbf{A}
  \item \textbf{A}
  \item \textbf{A}
  \item \textbf{A}
  \item \textbf{A}
  \item \textbf{A}
  \item \textbf{A}
  \item \textbf{A}
  \item \textbf{A}
  \item \textbf{A}
  \item \textbf{A}
  \item \textbf{A}
  \item \textbf{A}
  \item \textbf{A}
  \item \textbf{A}
  \item \textbf{A}
  \item \textbf{A}
  \item \textbf{A}
  \item \textbf{A}
  \item \textbf{A}
  \item \textbf{A}
  \item \textbf{A}
  \item \textbf{A}
  \item \textbf{A}
  \item \textbf{A}
  \item \textbf{A}
  \item \textbf{A}
  \item \textbf{A}
  \item \textbf{A}
  \item \textbf{A}
  \item \textbf{A}
  \item \textbf{A}
  \item \textbf{A}
  \item \textbf{A}
  \item \textbf{A}
  \item \textbf{A}
  \item \textbf{A}
  \item \textbf{A}
  \item \textbf{A}
  \item \textbf{A}
  \item \textbf{A}
  \item \textbf{A}
  \item \textbf{A}
  \item \textbf{A}
  \item \textbf{A}
  \item \textbf{A}
  \item \textbf{A}
  \item \textbf{A}
  \item \textbf{A}
  \item \textbf{A}
  \item \textbf{A}
  \item \textbf{A}
  \item \textbf{A}
  \item \textbf{A}
  \item \textbf{A}
  \item \textbf{A}
  \item \textbf{A}
  \item \textbf{A}
  \item \textbf{A}
  \item \textbf{A}
  \item \textbf{A}
  \item \textbf{A}
  \item \textbf{A}
  \item \textbf{A}
  \item \textbf{A}
  \item \textbf{A}
  \item \textbf{A}
  \item \textbf{A}
  \item \textbf{A}
  \item \textbf{A}
  \item \textbf{A}
  \item \textbf{A}
  \item \textbf{A}
  \item \textbf{A}
  \item \textbf{A}
  \item \textbf{A}
  \item \textbf{A}
  \item \textbf{A}
  \item \textbf{A}
  \item \textbf{A}
  \item \textbf{A}
  \item \textbf{A}
  \item \textbf{A}
  \item \textbf{A}
  \item \textbf{A}
  \item \textbf{A}
  \item \textbf{A}
  \item \textbf{A}
  \item \textbf{A}
  \item \textbf{A}
  \item \textbf{A}
  \item \textbf{A}
  \item \textbf{A}
  \item \textbf{A}
  \item \textbf{A}
  \item \textbf{A}
  \item \textbf{A}
  \item \textbf{A}
  \item \textbf{A}
  \item \textbf{A}
  \item \textbf{A}
  \item \textbf{A}
  \item \textbf{A}
  \item \textbf{A}
  \item \textbf{A}
\end{enumerate}
\end{multicols}
\end{tcolorbox}
\newpage\n
% ==============================================================================
% MATERIA 3: LITERATURA
% ==============================================================================
\section{Literatura (Módulos 3.1 al 3.5)}

\subsection{3.1 Literatura Universal y Europea}

\subsubsection*{SECCIÓN I: LITERATURA CLÁSICA GRIEGA (1–60)}

\begin{enumerate}[label=\textbf{\arabic*.}]
\setcounter{enumi}{0}
\item ¿Quién es el autor de la Ilíada y la Odisea?
\begin{enumerate}[label=\textbf{\Alph*)}]
\begin{multicols}{2}
  \item Sófocles
  \item Eurípides
  \item Homero
  \item Esquilo
  \item Aristófanes
  \item Hesíodo
\end{multicols}
\end{enumerate}
\vspace{0.15cm}
\item ¿Cuál es el tema central de la Ilíada?
\begin{enumerate}[label=\textbf{\Alph*)}]
\begin{multicols}{2}
  \item El regreso de Odiseo a Ítaca
  \item La cólera de Aquiles durante la guerra de Troya
  \item Los viajes de Eneas
  \item La fundación de Roma
  \item Las aventuras de Jasón
  \item Los trabajos de Hércules
\end{multicols}
\end{enumerate}
\vspace{0.15cm}
\item ¿Cuál es el tema central de la Odisea?
\begin{enumerate}[label=\textbf{\Alph*)}]
\begin{multicols}{2}
  \item La guerra de Troya
  \item El regreso de Odiseo a Ítaca
  \item La cólera de Aquiles
  \item La fundación de Tebas
  \item Las aventuras de Eneas
  \item Los trabajos de Hércules
\end{multicols}
\end{enumerate}
\vspace{0.15cm}
\item ¿Quién es el protagonista de la Odisea?
\begin{enumerate}[label=\textbf{\Alph*)}]
\begin{multicols}{2}
  \item Aquiles
  \item Héctor
  \item Odiseo (Ulises)
  \item Agamenón
  \item Menelao
  \item Paris
\end{multicols}
\end{enumerate}
\vspace{0.15cm}
\item ¿Quién es el protagonista de la Ilíada?
\begin{enumerate}[label=\textbf{\Alph*)}]
\begin{multicols}{2}
  \item Odiseo
  \item Aquiles
  \item Héctor
  \item Agamenón
  \item Menelao
  \item Paris
\end{multicols}
\end{enumerate}
\vspace{0.15cm}
\item ¿Qué epopeya narra la fundación de Roma?
\begin{enumerate}[label=\textbf{\Alph*)}]
\begin{multicols}{2}
  \item La Ilíada
  \item La Odisea
  \item La Eneida
  \item La Teogonía
  \item Los Trabajos y los Días
  \item Las Argonáuticas
\end{multicols}
\end{enumerate}
\vspace{0.15cm}
\item ¿Quién es el autor de la Eneida?
\begin{enumerate}[label=\textbf{\Alph*)}]
\begin{multicols}{2}
  \item Homero
  \item Virgilio
  \item Ovidio
  \item Horacio
  \item Catulo
  \item Propercio
\end{multicols}
\end{enumerate}
\vspace{0.15cm}
\item ¿Quién es el protagonista de la Eneida?
\begin{enumerate}[label=\textbf{\Alph*)}]
\begin{multicols}{2}
  \item Odiseo
  \item Aquiles
  \item Eneas
  \item Héctor
  \item Paris
  \item Agamenón
\end{multicols}
\end{enumerate}
\vspace{0.15cm}
\item ¿Qué poeta griego escribió la Teogonía?
\begin{enumerate}[label=\textbf{\Alph*)}]
\begin{multicols}{2}
  \item Homero
  \item Hesíodo
  \item Safo
  \item Píndaro
  \item Anacreonte
  \item Arquíloco
\end{multicols}
\end{enumerate}
\vspace{0.15cm}
\item ¿Qué poeta griego escribió Los Trabajos y los Días?
\begin{enumerate}[label=\textbf{\Alph*)}]
\begin{multicols}{2}
  \item Homero
  \item Hesíodo
  \item Safo
  \item Píndaro
  \item Anacreonte
  \item Arquíloco
\end{multicols}
\end{enumerate}
\vspace{0.15cm}
\item ¿Quién es el autor de la tragedia Edipo Rey?
\begin{enumerate}[label=\textbf{\Alph*)}]
\begin{multicols}{2}
  \item Esquilo
  \item Sófocles
  \item Eurípides
  \item Aristófanes
  \item Menandro
  \item Sófocles
\end{multicols}
\end{enumerate}
\vspace{0.15cm}
\item ¿Quién es el autor de la tragedia Medea?
\begin{enumerate}[label=\textbf{\Alph*)}]
\begin{multicols}{2}
  \item Esquilo
  \item Sófocles
  \item Eurípides
  \item Aristófanes
  \item Menandro
  \item Sófocles
\end{multicols}
\end{enumerate}
\vspace{0.15cm}
\item ¿Quién es el autor de la tragedia Prometeo Encadenado?
\begin{enumerate}[label=\textbf{\Alph*)}]
\begin{multicols}{2}
  \item Esquilo
  \item Sófocles
  \item Eurípides
  \item Aristófanes
  \item Menandro
  \item Sófocles
\end{multicols}
\end{enumerate}
\vspace{0.15cm}
\item ¿Quién es el autor de la comedia Lisístrata?
\begin{enumerate}[label=\textbf{\Alph*)}]
\begin{multicols}{2}
  \item Esquilo
  \item Sófocles
  \item Eurípides
  \item Aristófanes
  \item Menandro
  \item Sófocles
\end{multicols}
\end{enumerate}
\vspace{0.15cm}
\item ¿Cuál es el tema central de Edipo Rey?
\begin{enumerate}[label=\textbf{\Alph*)}]
\begin{multicols}{2}
  \item El regreso de Odiseo
  \item El destino inevitable de Edipo
  \item La guerra de Troya
  \item Los viajes de Eneas
  \item La fundación de Roma
  \item Las aventuras de Jasón
\end{multicols}
\end{enumerate}
\vspace{0.15cm}
\item ¿Cuál es el tema central de Medea?
\begin{enumerate}[label=\textbf{\Alph*)}]
\begin{multicols}{2}
  \item La venganza de Medea contra Jasón
  \item El regreso de Odiseo
  \item La guerra de Troya
  \item Los viajes de Eneas
  \item La fundación de Roma
  \item Las aventuras de Aquiles
\end{multicols}
\end{enumerate}
\vspace{0.15cm}
\item ¿Cuál es el tema central de Prometeo Encadenado?
\begin{enumerate}[label=\textbf{\Alph*)}]
\begin{multicols}{2}
  \item El castigo de Prometeo por robar el fuego
  \item El regreso de Odiseo
  \item La guerra de Troya
  \item Los viajes de Eneas
  \item La fundación de Roma
  \item Las aventuras de Aquiles
\end{multicols}
\end{enumerate}
\vspace{0.15cm}
\item ¿Qué es una tragedia griega?
\begin{enumerate}[label=\textbf{\Alph*)}]
\begin{multicols}{2}
  \item Una obra teatral con final desdichado
  \item Una obra teatral con final feliz
  \item Una epopeya
  \item Una lírica
  \item Una comedia
  \item Un drama satírico
\end{multicols}
\end{enumerate}
\vspace{0.15cm}
\item ¿Qué es una comedia griega?
\begin{enumerate}[label=\textbf{\Alph*)}]
\begin{multicols}{2}
  \item Una obra teatral con final desdichado
  \item Una obra teatral con final feliz y humor
  \item Una epopeya
  \item Una lírica
  \item Una tragedia
  \item Un drama satírico
\end{multicols}
\end{enumerate}
\vspace{0.15cm}
\item ¿Qué es la catarsis en la tragedia griega?
\begin{enumerate}[label=\textbf{\Alph*)}]
\begin{multicols}{2}
  \item La purificación emocional del espectador
  \item La risa del espectador
  \item El llanto del espectador
  \item El miedo del espectador
  \item La alegría del espectador
  \item La indiferencia del espectador
\end{multicols}
\end{enumerate}
\vspace{0.15cm}
\item ¿Qué es la hybris en la tragedia griega?
\begin{enumerate}[label=\textbf{\Alph*)}]
\begin{multicols}{2}
  \item La desmesura o soberbia del héroe
  \item La humildad del héroe
  \item La valentía del héroe
  \item La cobardía del héroe
  \item La sabiduría del héroe
  \item La ignorancia del héroe
\end{multicols}
\end{enumerate}
\vspace{0.15cm}
\item ¿Qué es el destino o moira en la tragedia griega?
\begin{enumerate}[label=\textbf{\Alph*)}]
\begin{multicols}{2}
  \item La fuerza que determina la vida de los personajes
  \item La libertad de los personajes
  \item La casualidad
  \item La voluntad de los dioses
  \item La fortuna
  \item La suerte
\end{multicols}
\end{enumerate}
\vspace{0.15cm}
\item ¿Quién es el autor de la poesía lírica griega?
\begin{enumerate}[label=\textbf{\Alph*)}]
\begin{multicols}{2}
  \item Safo
  \item Homero
  \item Hesíodo
  \item Esquilo
  \item Sófocles
  \item Eurípides
\end{multicols}
\end{enumerate}
\vspace{0.15cm}
\item ¿Quién es el autor de la poesía lírica griega?
\begin{enumerate}[label=\textbf{\Alph*)}]
\begin{multicols}{2}
  \item Píndaro
  \item Homero
  \item Hesíodo
  \item Esquilo
  \item Sófocles
  \item Eurípides
\end{multicols}
\end{enumerate}
\vspace{0.15cm}
\item ¿Quién es el autor de la poesía lírica griega?
\begin{enumerate}[label=\textbf{\Alph*)}]
\begin{multicols}{2}
  \item Anacreonte
  \item Homero
  \item Hesíodo
  \item Esquilo
  \item Sófocles
  \item Eurípides
\end{multicols}
\end{enumerate}
\vspace{0.15cm}
\item ¿Qué es la épica?
\begin{enumerate}[label=\textbf{\Alph*)}]
\begin{multicols}{2}
  \item Un género narrativo en verso que narra las hazañas de héroes
  \item Un género teatral
  \item Un género lírico
  \item Un género didáctico
  \item Un género periodístico
  \item Un género ensayístico
\end{multicols}
\end{enumerate}
\vspace{0.15cm}
\item ¿Qué es la lírica?
\begin{enumerate}[label=\textbf{\Alph*)}]
\begin{multicols}{2}
  \item Un género poético que expresa sentimientos
  \item Un género narrativo
  \item Un género teatral
  \item Un género didáctico
  \item Un género periodístico
  \item Un género ensayístico
\end{multicols}
\end{enumerate}
\vspace{0.15cm}
\item ¿Qué es la tragedia?
\begin{enumerate}[label=\textbf{\Alph*)}]
\begin{multicols}{2}
  \item Un género teatral con final desdichado
  \item Un género narrativo
  \item Un género lírico
  \item Un género didáctico
  \item Un género periodístico
  \item Un género ensayístico
\end{multicols}
\end{enumerate}
\vspace{0.15cm}
\item ¿Qué es la comedia?
\begin{enumerate}[label=\textbf{\Alph*)}]
\begin{multicols}{2}
  \item Un género teatral con final feliz y humor
  \item Un género narrativo
  \item Un género lírico
  \item Un género didáctico
  \item Un género periodístico
  \item Un género ensayístico
\end{multicols}
\end{enumerate}
\vspace{0.15cm}
\item ¿Qué es la epopeya?
\begin{enumerate}[label=\textbf{\Alph*)}]
\begin{multicols}{2}
  \item Un poema épico extenso que narra hazañas heroicas
  \item Un poema lírico
  \item Una obra teatral
  \item Un ensayo
  \item Una novela
  \item Un cuento
\end{multicols}
\end{enumerate}
\vspace{0.15cm}
\item ¿Cuál es el tema central de la Ilíada?
\begin{enumerate}[label=\textbf{\Alph*)}]
\begin{multicols}{2}
  \item La cólera de Aquiles
  \item El regreso de Odiseo
  \item La fundación de Roma
  \item Los viajes de Eneas
  \item Las aventuras de Jasón
  \item Los trabajos de Hércules
\end{multicols}
\end{enumerate}
\vspace{0.15cm}
\item ¿Cuál es el tema central de la Odisea?
\begin{enumerate}[label=\textbf{\Alph*)}]
\begin{multicols}{2}
  \item El regreso de Odiseo a Ítaca
  \item La cólera de Aquiles
  \item La fundación de Roma
  \item Los viajes de Eneas
  \item Las aventuras de Jasón
  \item Los trabajos de Hércules
\end{multicols}
\end{enumerate}
\vspace{0.15cm}
\item ¿Cuál es el tema central de la Eneida?
\begin{enumerate}[label=\textbf{\Alph*)}]
\begin{multicols}{2}
  \item La fundación de Roma
  \item El regreso de Odiseo
  \item La cólera de Aquiles
  \item Las aventuras de Jasón
  \item Los trabajos de Hércules
  \item La guerra de Tebas
\end{multicols}
\end{enumerate}
\vspace{0.15cm}
\item ¿Quién es el autor de la Ilíada?
\begin{enumerate}[label=\textbf{\Alph*)}]
\begin{multicols}{2}
  \item Homero
  \item Virgilio
  \item Hesíodo
  \item Safo
  \item Sófocles
  \item Eurípides
\end{multicols}
\end{enumerate}
\vspace{0.15cm}
\item ¿Quién es el autor de la Odisea?
\begin{enumerate}[label=\textbf{\Alph*)}]
\begin{multicols}{2}
  \item Homero
  \item Virgilio
  \item Hesíodo
  \item Safo
  \item Sófocles
  \item Eurípides
\end{multicols}
\end{enumerate}
\vspace{0.15cm}
\item ¿Quién es el autor de la Eneida?
\begin{enumerate}[label=\textbf{\Alph*)}]
\begin{multicols}{2}
  \item Homero
  \item Virgilio
  \item Hesíodo
  \item Safo
  \item Sófocles
  \item Eurípides
\end{multicols}
\end{enumerate}
\vspace{0.15cm}
\item ¿Quién es el autor de Edipo Rey?
\begin{enumerate}[label=\textbf{\Alph*)}]
\begin{multicols}{2}
  \item Sófocles
  \item Homero
  \item Virgilio
  \item Hesíodo
  \item Safo
  \item Eurípides
\end{multicols}
\end{enumerate}
\vspace{0.15cm}
\item ¿Quién es el autor de Medea?
\begin{enumerate}[label=\textbf{\Alph*)}]
\begin{multicols}{2}
  \item Eurípides
  \item Homero
  \item Virgilio
  \item Hesíodo
  \item Safo
  \item Sófocles
\end{multicols}
\end{enumerate}
\vspace{0.15cm}
\item ¿Quién es el autor de Prometeo Encadenado?
\begin{enumerate}[label=\textbf{\Alph*)}]
\begin{multicols}{2}
  \item Esquilo
  \item Homero
  \item Virgilio
  \item Hesíodo
  \item Safo
  \item Sófocles
\end{multicols}
\end{enumerate}
\vspace{0.15cm}
\item ¿Quién es el autor de Lisístrata?
\begin{enumerate}[label=\textbf{\Alph*)}]
\begin{multicols}{2}
  \item Aristófanes
  \item Homero
  \item Virgilio
  \item Hesíodo
  \item Safo
  \item Sófocles
\end{multicols}
\end{enumerate}
\vspace{0.15cm}
\item ¿Qué guerra se narra en la Ilíada?
\begin{enumerate}[label=\textbf{\Alph*)}]
\begin{multicols}{2}
  \item La guerra de Troya
  \item La guerra de Tebas
  \item Las guerras médicas
  \item La guerra del Peloponeso
  \item Las guerras púnicas
  \item La guerra de las Galias
\end{multicols}
\end{enumerate}
\vspace{0.15cm}
\item ¿Qué ciudad es el objetivo de Odiseo en la Odisea?
\begin{enumerate}[label=\textbf{\Alph*)}]
\begin{multicols}{2}
  \item Troya
  \item Ítaca
  \item Esparta
  \item Atenas
  \item Tebas
  \item Micenas
\end{multicols}
\end{enumerate}
\vspace{0.15cm}
\item ¿Qué ciudad funda Eneas en la Eneida?
\begin{enumerate}[label=\textbf{\Alph*)}]
\begin{multicols}{2}
  \item Roma
  \item Cartago
  \item Troya
  \item Ítaca
  \item Atenas
  \item Esparta
\end{multicols}
\end{enumerate}
\vspace{0.15cm}
\item ¿Qué personaje mata a su padre y se casa con su madre en la tragedia griega?
\begin{enumerate}[label=\textbf{\Alph*)}]
\begin{multicols}{2}
  \item Edipo
  \item Orestes
  \item Medea
  \item Prometeo
  \item Aquiles
  \item Odiseo
\end{multicols}
\end{enumerate}
\vspace{0.15cm}
\item ¿Qué personaje mata a sus hijos para vengarse de su esposo?
\begin{enumerate}[label=\textbf{\Alph*)}]
\begin{multicols}{2}
  \item Medea
  \item Clitemnestra
  \item Antígona
  \item Electra
  \item Helena
  \item Penélope
\end{multicols}
\end{enumerate}
\vspace{0.15cm}
\item ¿Qué personaje roba el fuego a los dioses?
\begin{enumerate}[label=\textbf{\Alph*)}]
\begin{multicols}{2}
  \item Prometeo
  \item Epimeteo
  \item Pandora
  \item Zeus
  \item Apolo
  \item Hermes
\end{multicols}
\end{enumerate}
\vspace{0.15cm}
\item ¿Qué personaje es castigado por robar el fuego?
\begin{enumerate}[label=\textbf{\Alph*)}]
\begin{multicols}{2}
  \item Prometeo
  \item Epimeteo
  \item Pandora
  \item Zeus
  \item Apolo
  \item Hermes
\end{multicols}
\end{enumerate}
\vspace{0.15cm}
\item ¿Qué es la Ilíada?
\begin{enumerate}[label=\textbf{\Alph*)}]
\begin{multicols}{2}
  \item Una epopeya griega
  \item Una tragedia griega
  \item Una comedia griega
  \item Una lírica griega
  \item Un drama satírico
  \item Un ensayo griego
\end{multicols}
\end{enumerate}
\vspace{0.15cm}
\item ¿Qué es la Odisea?
\begin{enumerate}[label=\textbf{\Alph*)}]
\begin{multicols}{2}
  \item Una epopeya griega
  \item Una tragedia griega
  \item Una comedia griega
  \item Una lírica griega
  \item Un drama satírico
  \item Un ensayo griego
\end{multicols}
\end{enumerate}
\vspace{0.15cm}
\item ¿Qué es la Eneida?
\begin{enumerate}[label=\textbf{\Alph*)}]
\begin{multicols}{2}
  \item Una epopeya latina
  \item Una tragedia latina
  \item Una comedia latina
  \item Una lírica latina
  \item Un drama satírico
  \item Un ensayo latino
\end{multicols}
\end{enumerate}
\vspace{0.15cm}
\item ¿Cuál es el tema central de la Teogonía?
\begin{enumerate}[label=\textbf{\Alph*)}]
\begin{multicols}{2}
  \item El origen de los dioses griegos
  \item La guerra de Troya
  \item El regreso de Odiseo
  \item La fundación de Roma
  \item Las aventuras de Jasón
  \item Los trabajos de Hércules
\end{multicols}
\end{enumerate}
\vspace{0.15cm}
\item ¿Cuál es el tema central de Los Trabajos y los Días?
\begin{enumerate}[label=\textbf{\Alph*)}]
\begin{multicols}{2}
  \item La vida campesina y la justicia
  \item La guerra de Troya
  \item El regreso de Odiseo
  \item La fundación de Roma
  \item Las aventuras de Jasón
  \item Los trabajos de Hércules
\end{multicols}
\end{enumerate}
\vspace{0.15cm}
\item ¿Qué poeta griego es conocido por su poesía amorosa?
\begin{enumerate}[label=\textbf{\Alph*)}]
\begin{multicols}{2}
  \item Safo
  \item Homero
  \item Hesíodo
  \item Esquilo
  \item Sófocles
  \item Eurípides
\end{multicols}
\end{enumerate}
\vspace{0.15cm}
\item ¿Qué poeta griego es conocido por sus odas triunfales?
\begin{enumerate}[label=\textbf{\Alph*)}]
\begin{multicols}{2}
  \item Píndaro
  \item Homero
  \item Hesíodo
  \item Esquilo
  \item Sófocles
  \item Eurípides
\end{multicols}
\end{enumerate}
\vspace{0.15cm}
\item ¿Qué poeta griego es conocido por sus poemas sobre el vino y el amor?
\begin{enumerate}[label=\textbf{\Alph*)}]
\begin{multicols}{2}
  \item Anacreonte
  \item Homero
  \item Hesíodo
  \item Esquilo
  \item Sófocles
  \item Eurípides
\end{multicols}
\end{enumerate}
\vspace{0.15cm}
\item ¿Qué es un drama satírico?
\begin{enumerate}[label=\textbf{\Alph*)}]
\begin{multicols}{2}
  \item Una obra teatral cómica que se representaba después de las tragedias
  \item Una tragedia
  \item Una comedia
  \item Una epopeya
  \item Una lírica
  \item Un ensayo
\end{multicols}
\end{enumerate}
\vspace{0.15cm}
\item ¿Qué es la épica?
\begin{enumerate}[label=\textbf{\Alph*)}]
\begin{multicols}{2}
  \item Un género narrativo en verso
  \item Un género teatral
  \item Un género lírico
  \item Un género didáctico
  \item Un género periodístico
  \item Un género ensayístico
\end{multicols}
\end{enumerate}
\vspace{0.15cm}
\item ¿Qué es la lírica?
\begin{enumerate}[label=\textbf{\Alph*)}]
\begin{multicols}{2}
  \item Un género poético que expresa sentimientos
  \item Un género narrativo
  \item Un género teatral
  \item Un género didáctico
  \item Un género periodístico
  \item Un género ensayístico
\end{multicols}
\end{enumerate}
\vspace{0.15cm}
\item ¿Qué es la tragedia?
\begin{enumerate}[label=\textbf{\Alph*)}]
\begin{multicols}{2}
  \item Un género teatral con final desdichado
  \item Un género narrativo
  \item Un género lírico
  \item Un género didáctico
  \item Un género periodístico
  \item Un género ensayístico
\end{multicols}
\end{enumerate}
\vspace{0.15cm}
\item ¿Qué es la comedia?
\begin{enumerate}[label=\textbf{\Alph*)}]
\begin{multicols}{2}
  \item Un género teatral con final feliz y humor
  \item Un género narrativo
  \item Un género lírico
  \item Un género didáctico
  \item Un género periodístico
  \item Un género ensayístico
\end{multicols}
\end{enumerate}
\vspace{0.15cm}
\end{enumerate}\n
\subsubsection*{SECCIÓN II: EDAD MEDIA (61–120)}

\begin{enumerate}[label=\textbf{\arabic*.}]
\setcounter{enumi}{60}
\item ¿Qué es el Cantar de Mio Cid?
\begin{enumerate}[label=\textbf{\Alph*)}]
\begin{multicols}{2}
  \item Un cantar de gesta español
  \item Una novela de caballerías
  \item Una obra teatral
  \item Una lírica
  \item Un ensayo
  \item Una epopeya griega
\end{multicols}
\end{enumerate}
\vspace{0.15cm}
\item ¿Quién es el protagonista del Cantar de Mio Cid?
\begin{enumerate}[label=\textbf{\Alph*)}]
\begin{multicols}{2}
  \item El Cid Campeador
  \item Carlomagno
  \item Roldán
  \item Arturo
  \item Lancelot
  \item Tristán
\end{multicols}
\end{enumerate}
\vspace{0.15cm}
\item ¿Qué es el Cantar de Roldán?
\begin{enumerate}[label=\textbf{\Alph*)}]
\begin{multicols}{2}
  \item Un cantar de gesta francés
  \item Una novela de caballerías
  \item Una obra teatral
  \item Una lírica
  \item Un ensayo
  \item Una epopeya griega
\end{multicols}
\end{enumerate}
\vspace{0.15cm}
\item ¿Quién es el protagonista del Cantar de Roldán?
\begin{enumerate}[label=\textbf{\Alph*)}]
\begin{multicols}{2}
  \item Roldán
  \item Carlomagno
  \item El Cid
  \item Arturo
  \item Lancelot
  \item Tristán
\end{multicols}
\end{enumerate}
\vspace{0.15cm}
\item ¿Qué es el Cantar de los Nibelungos?
\begin{enumerate}[label=\textbf{\Alph*)}]
\begin{multicols}{2}
  \item Un cantar de gesta alemán
  \item Una novela de caballerías
  \item Una obra teatral
  \item Una lírica
  \item Un ensayo
  \item Una epopeya griega
\end{multicols}
\end{enumerate}
\vspace{0.15cm}
\item ¿Quién es el protagonista del Cantar de los Nibelungos?
\begin{enumerate}[label=\textbf{\Alph*)}]
\begin{multicols}{2}
  \item Sigfrido
  \item Carlomagno
  \item Roldán
  \item Arturo
  \item Lancelot
  \item Tristán
\end{multicols}
\end{enumerate}
\vspace{0.15cm}
\item ¿Qué es la Divina Comedia?
\begin{enumerate}[label=\textbf{\Alph*)}]
\begin{multicols}{2}
  \item Un poema épico-alegórico de Dante
  \item Una novela de caballerías
  \item Una obra teatral
  \item Una lírica
  \item Un ensayo
  \item Una epopeya griega
\end{multicols}
\end{enumerate}
\vspace{0.15cm}
\item ¿Quién es el autor de la Divina Comedia?
\begin{enumerate}[label=\textbf{\Alph*)}]
\begin{multicols}{2}
  \item Dante Alighieri
  \item Petrarca
  \item Boccaccio
  \item Virgilio
  \item Homero
  \item Chaucer
\end{multicols}
\end{enumerate}
\vspace{0.15cm}
\item ¿Cuántas partes tiene la Divina Comedia?
\begin{enumerate}[label=\textbf{\Alph*)}]
\begin{multicols}{2}
  \item Tres: Infierno, Purgatorio y Paraíso
  \item Dos: Infierno y Paraíso
  \item Cuatro: Infierno, Purgatorio, Paraíso y Limbo
  \item Una: Infierno
  \item Cinco: Infierno, Purgatorio, Paraíso, Limbo y Cielo
  \item Seis: Infierno, Purgatorio, Paraíso, Limbo, Cielo y Tierra
\end{multicols}
\end{enumerate}
\vspace{0.15cm}
\item ¿Quién guía a Dante en el Infierno y el Purgatorio?
\begin{enumerate}[label=\textbf{\Alph*)}]
\begin{multicols}{2}
  \item Virgilio
  \item Beatriz
  \item Homero
  \item Petrarca
  \item Boccaccio
  \item Santo Tomás
\end{multicols}
\end{enumerate}
\vspace{0.15cm}
\item ¿Quién guía a Dante en el Paraíso?
\begin{enumerate}[label=\textbf{\Alph*)}]
\begin{multicols}{2}
  \item Beatriz
  \item Virgilio
  \item Homero
  \item Petrarca
  \item Boccaccio
  \item Santo Tomás
\end{multicols}
\end{enumerate}
\vspace{0.15cm}
\item ¿Qué es el Decamerón?
\begin{enumerate}[label=\textbf{\Alph*)}]
\begin{multicols}{2}
  \item Una colección de cuentos de Boccaccio
  \item Una novela de caballerías
  \item Una obra teatral
  \item Una lírica
  \item Un ensayo
  \item Una epopeya
\end{multicols}
\end{enumerate}
\vspace{0.15cm}
\item ¿Quién es el autor del Decamerón?
\begin{enumerate}[label=\textbf{\Alph*)}]
\begin{multicols}{2}
  \item Boccaccio
  \item Dante
  \item Petrarca
  \item Chaucer
  \item Villon
  \item Rojas
\end{multicols}
\end{enumerate}
\vspace{0.15cm}
\item ¿Qué es el Cancionero?
\begin{enumerate}[label=\textbf{\Alph*)}]
\begin{multicols}{2}
  \item Una colección de poemas de Petrarca
  \item Una novela de caballerías
  \item Una obra teatral
  \item Un ensayo
  \item Una epopeya
  \item Un cantar de gesta
\end{multicols}
\end{enumerate}
\vspace{0.15cm}
\item ¿Quién es el autor del Cancionero?
\begin{enumerate}[label=\textbf{\Alph*)}]
\begin{multicols}{2}
  \item Petrarca
  \item Dante
  \item Boccaccio
  \item Chaucer
  \item Villon
  \item Rojas
\end{multicols}
\end{enumerate}
\vspace{0.15cm}
\item ¿Qué es La Celestina?
\begin{enumerate}[label=\textbf{\Alph*)}]
\begin{multicols}{2}
  \item Una obra teatral española
  \item Una novela de caballerías
  \item Una lírica
  \item Un ensayo
  \item Una epopeya
  \item Un cantar de gesta
\end{multicols}
\end{enumerate}
\vspace{0.15cm}
\item ¿Quién es el autor de La Celestina?
\begin{enumerate}[label=\textbf{\Alph*)}]
\begin{multicols}{2}
  \item Fernando de Rojas
  \item Cervantes
  \item Lope de Vega
  \item Calderón
  \item Tirso de Molina
  \item Quevedo
\end{multicols}
\end{enumerate}
\vspace{0.15cm}
\item ¿Qué son los Cuentos de Canterbury?
\begin{enumerate}[label=\textbf{\Alph*)}]
\begin{multicols}{2}
  \item Una colección de cuentos de Chaucer
  \item Una novela de caballerías
  \item Una obra teatral
  \item Una lírica
  \item Un ensayo
  \item Una epopeya
\end{multicols}
\end{enumerate}
\vspace{0.15cm}
\item ¿Quién es el autor de los Cuentos de Canterbury?
\begin{enumerate}[label=\textbf{\Alph*)}]
\begin{multicols}{2}
  \item Chaucer
  \item Dante
  \item Petrarca
  \item Boccaccio
  \item Villon
  \item Rojas
\end{multicols}
\end{enumerate}
\vspace{0.15cm}
\item ¿Qué es el Roman de la Rose?
\begin{enumerate}[label=\textbf{\Alph*)}]
\begin{multicols}{2}
  \item Una novela alegórica francesa
  \item Una novela de caballerías
  \item Una obra teatral
  \item Una lírica
  \item Un ensayo
  \item Una epopeya
\end{multicols}
\end{enumerate}
\vspace{0.15cm}
\item ¿Qué es el Roman de Renart?
\begin{enumerate}[label=\textbf{\Alph*)}]
\begin{multicols}{2}
  \item Una colección de fábulas francesas
  \item Una novela de caballerías
  \item Una obra teatral
  \item Una lírica
  \item Un ensayo
  \item Una epopeya
\end{multicols}
\end{enumerate}
\vspace{0.15cm}
\item ¿Qué es la lírica provenzal?
\begin{enumerate}[label=\textbf{\Alph*)}]
\begin{multicols}{2}
  \item La poesía de los trovadores
  \item La poesía de los juglares
  \item La poesía de los clérigos
  \item La poesía de los caballeros
  \item La poesía de los campesinos
  \item La poesía de los reyes
\end{multicols}
\end{enumerate}
\vspace{0.15cm}
\item ¿Qué es el amor cortés?
\begin{enumerate}[label=\textbf{\Alph*)}]
\begin{multicols}{2}
  \item Un código amoroso de la Edad Media
  \item Un código legal
  \item Un código religioso
  \item Un código militar
  \item Un código comercial
  \item Un código agrícola
\end{multicols}
\end{enumerate}
\vspace{0.15cm}
\item ¿Qué es la novela de caballerías?
\begin{enumerate}[label=\textbf{\Alph*)}]
\begin{multicols}{2}
  \item Un género narrativo sobre las aventuras de caballeros
  \item Un género teatral
  \item Un género lírico
  \item Un género didáctico
  \item Un género periodístico
  \item Un género ensayístico
\end{multicols}
\end{enumerate}
\vspace{0.15cm}
\item ¿Cuál es la novela de caballerías más famosa?
\begin{enumerate}[label=\textbf{\Alph*)}]
\begin{multicols}{2}
  \item Amadís de Gaula
  \item Tirant lo Blanc
  \item El Quijote
  \item La Celestina
  \item El Cid
  \item Roldán
\end{multicols}
\end{enumerate}
\vspace{0.15cm}
\item ¿Qué es el Tirant lo Blanc?
\begin{enumerate}[label=\textbf{\Alph*)}]
\begin{multicols}{2}
  \item Una novela de caballerías catalana
  \item Una novela de caballerías española
  \item Una novela de caballerías francesa
  \item Una novela de caballerías inglesa
  \item Una novela de caballerías alemana
  \item Una novela de caballerías italiana
\end{multicols}
\end{enumerate}
\vspace{0.15cm}
\item ¿Qué es el Amadís de Gaula?
\begin{enumerate}[label=\textbf{\Alph*)}]
\begin{multicols}{2}
  \item Una novela de caballerías española
  \item Una novela de caballerías catalana
  \item Una novela de caballerías francesa
  \item Una novela de caballerías inglesa
  \item Una novela de caballerías alemana
  \item Una novela de caballerías italiana
\end{multicols}
\end{enumerate}
\vspace{0.15cm}
\item ¿Qué es el teatro medieval?
\begin{enumerate}[label=\textbf{\Alph*)}]
\begin{multicols}{2}
  \item Un teatro de carácter religioso y didáctico
  \item Un teatro de carácter profano
  \item Un teatro de carácter cómico
  \item Un teatro de carácter trágico
  \item Un teatro de carácter político
  \item Un teatro de carácter amoroso
\end{multicols}
\end{enumerate}
\vspace{0.15cm}
\item ¿Qué son los autos sacramentales?
\begin{enumerate}[label=\textbf{\Alph*)}]
\begin{multicols}{2}
  \item Obras teatrales de tema religioso
  \item Obras teatrales de tema profano
  \item Obras teatrales de tema cómico
  \item Obras teatrales de tema trágico
  \item Obras teatrales de tema político
  \item Obras teatrales de tema amoroso
\end{multicols}
\end{enumerate}
\vspace{0.15cm}
\item ¿Qué son los misterios?
\begin{enumerate}[label=\textbf{\Alph*)}]
\begin{multicols}{2}
  \item Obras teatrales de tema religioso
  \item Obras teatrales de tema profano
  \item Obras teatrales de tema cómico
  \item Obras teatrales de tema trágico
  \item Obras teatrales de tema político
  \item Obras teatrales de tema amoroso
\end{multicols}
\end{enumerate}
\vspace{0.15cm}
\item ¿Qué son los milagros?
\begin{enumerate}[label=\textbf{\Alph*)}]
\begin{multicols}{2}
  \item Obras teatrales de tema religioso
  \item Obras teatrales de tema profano
  \item Obras teatrales de tema cómico
  \item Obras teatrales de tema trágico
  \item Obras teatrales de tema político
  \item Obras teatrales de tema amoroso
\end{multicols}
\end{enumerate}
\vspace{0.15cm}
\item ¿Qué son las moralidades?
\begin{enumerate}[label=\textbf{\Alph*)}]
\begin{multicols}{2}
  \item Obras teatrales de tema moralizante
  \item Obras teatrales de tema profano
  \item Obras teatrales de tema cómico
  \item Obras teatrales de tema trágico
  \item Obras teatrales de tema político
  \item Obras teatrales de tema amoroso
\end{multicols}
\end{enumerate}
\vspace{0.15cm}
\item ¿Qué son los entremeses?
\begin{enumerate}[label=\textbf{\Alph*)}]
\begin{multicols}{2}
  \item Obras teatrales breves de tema cómico
  \item Obras teatrales de tema religioso
  \item Obras teatrales de tema profano
  \item Obras teatrales de tema trágico
  \item Obras teatrales de tema político
  \item Obras teatrales de tema amoroso
\end{multicols}
\end{enumerate}
\vspace{0.15cm}
\item ¿Qué es el mester de juglaría?
\begin{enumerate}[label=\textbf{\Alph*)}]
\begin{multicols}{2}
  \item El oficio de los juglares que recitaban cantares de gesta
  \item El oficio de los clérigos que escribían en latín
  \item El oficio de los trovadores que cantaban poesía amorosa
  \item El oficio de los caballeros que escribían novelas
  \item El oficio de los campesinos que escribían canciones
  \item El oficio de los reyes que escribían leyes
\end{multicols}
\end{enumerate}
\vspace{0.15cm}
\item ¿Qué es el mester de clerecía?
\begin{enumerate}[label=\textbf{\Alph*)}]
\begin{multicols}{2}
  \item El oficio de los clérigos que escribían poesía culta
  \item El oficio de los juglares que recitaban cantares de gesta
  \item El oficio de los trovadores que cantaban poesía amorosa
  \item El oficio de los caballeros que escribían novelas
  \item El oficio de los campesinos que escribían canciones
  \item El oficio de los reyes que escribían leyes
\end{multicols}
\end{enumerate}
\vspace{0.15cm}
\item ¿Qué es el Libro de Buen Amor?
\begin{enumerate}[label=\textbf{\Alph*)}]
\begin{multicols}{2}
  \item Una obra del mester de clerecía
  \item Un cantar de gesta
  \item Una novela de caballerías
  \item Una obra teatral
  \item Una lírica
  \item Un ensayo
\end{multicols}
\end{enumerate}
\vspace{0.15cm}
\item ¿Quién es el autor del Libro de Buen Amor?
\begin{enumerate}[label=\textbf{\Alph*)}]
\begin{multicols}{2}
  \item Juan Ruiz, Arcipreste de Hita
  \item Gonzalo de Berceo
  \item Alfonso X el Sabio
  \item Don Juan Manuel
  \item Fernando de Rojas
  \item Cervantes
\end{multicols}
\end{enumerate}
\vspace{0.15cm}
\item ¿Qué es el Conde Lucanor?
\begin{enumerate}[label=\textbf{\Alph*)}]
\begin{multicols}{2}
  \item Una colección de cuentos de Don Juan Manuel
  \item Un cantar de gesta
  \item Una novela de caballerías
  \item Una obra teatral
  \item Una lírica
  \item Un ensayo
\end{multicols}
\end{enumerate}
\vspace{0.15cm}
\item ¿Quién es el autor del Conde Lucanor?
\begin{enumerate}[label=\textbf{\Alph*)}]
\begin{multicols}{2}
  \item Don Juan Manuel
  \item Juan Ruiz
  \item Gonzalo de Berceo
  \item Alfonso X el Sabio
  \item Fernando de Rojas
  \item Cervantes
\end{multicols}
\end{enumerate}
\vspace{0.15cm}
\item ¿Qué son los Milagros de Nuestra Señora?
\begin{enumerate}[label=\textbf{\Alph*)}]
\begin{multicols}{2}
  \item Una obra de Gonzalo de Berceo
  \item Un cantar de gesta
  \item Una novela de caballerías
  \item Una obra teatral
  \item Una lírica
  \item Un ensayo
\end{multicols}
\end{enumerate}
\vspace{0.15cm}
\item ¿Quién es el autor de los Milagros de Nuestra Señora?
\begin{enumerate}[label=\textbf{\Alph*)}]
\begin{multicols}{2}
  \item Gonzalo de Berceo
  \item Juan Ruiz
  \item Don Juan Manuel
  \item Alfonso X el Sabio
  \item Fernando de Rojas
  \item Cervantes
\end{multicols}
\end{enumerate}
\vspace{0.15cm}
\item ¿Qué es la Escuela de Traductores de Toledo?
\begin{enumerate}[label=\textbf{\Alph*)}]
\begin{multicols}{2}
  \item Un centro de traducción de obras clásicas
  \item Un centro de traducción de obras religiosas
  \item Un centro de traducción de obras científicas
  \item Un centro de traducción de obras literarias
  \item Un centro de traducción de obras filosóficas
  \item Un centro de traducción de obras médicas
\end{multicols}
\end{enumerate}
\vspace{0.15cm}
\item ¿Qué rey impulsó la Escuela de Traductores de Toledo?
\begin{enumerate}[label=\textbf{\Alph*)}]
\begin{multicols}{2}
  \item Alfonso X el Sabio
  \item Fernando III el Santo
  \item Carlos I
  \item Felipe II
  \item Isabel la Católica
  \item Juan II
\end{multicols}
\end{enumerate}
\vspace{0.15cm}
\item ¿Qué es la poesía trovadoresca?
\begin{enumerate}[label=\textbf{\Alph*)}]
\begin{multicols}{2}
  \item La poesía de los trovadores provenzales
  \item La poesía de los juglares
  \item La poesía de los clérigos
  \item La poesía de los caballeros
  \item La poesía de los campesinos
  \item La poesía de los reyes
\end{multicols}
\end{enumerate}
\vspace{0.15cm}
\item ¿Qué es la poesía goliardesca?
\begin{enumerate}[label=\textbf{\Alph*)}]
\begin{multicols}{2}
  \item La poesía de los clérigos errantes
  \item La poesía de los trovadores
  \item La poesía de los juglares
  \item La poesía de los caballeros
  \item La poesía de los campesinos
  \item La poesía de los reyes
\end{multicols}
\end{enumerate}
\vspace{0.15cm}
\item ¿Qué son los Carmina Burana?
\begin{enumerate}[label=\textbf{\Alph*)}]
\begin{multicols}{2}
  \item Una colección de poemas goliardescos
  \item Una colección de cantares de gesta
  \item Una colección de novelas de caballerías
  \item Una colección de obras teatrales
  \item Una colección de líricas
  \item Una colección de ensayos
\end{multicols}
\end{enumerate}
\vspace{0.15cm}
\item ¿Qué es el Dolce Stil Novo?
\begin{enumerate}[label=\textbf{\Alph*)}]
\begin{multicols}{2}
  \item Un movimiento poético italiano
  \item Un movimiento poético francés
  \item Un movimiento poético español
  \item Un movimiento poético inglés
  \item Un movimiento poético alemán
  \item Un movimiento poético portugués
\end{multicols}
\end{enumerate}
\vspace{0.15cm}
\item ¿Quiénes son los principales representantes del Dolce Stil Novo?
\begin{enumerate}[label=\textbf{\Alph*)}]
\begin{multicols}{2}
  \item Dante y Petrarca
  \item Dante y Boccaccio
  \item Petrarca y Boccaccio
  \item Dante y Chaucer
  \item Petrarca y Chaucer
  \item Boccaccio y Chaucer
\end{multicols}
\end{enumerate}
\vspace{0.15cm}
\item ¿Qué es la cuaderna vía?
\begin{enumerate}[label=\textbf{\Alph*)}]
\begin{multicols}{2}
  \item Una estrofa de cuatro versos alejandrinos con rima consonante
  \item Una estrofa de tres versos
  \item Una estrofa de cinco versos
  \item Una estrofa de seis versos
  \item Una estrofa de ocho versos
  \item Una estrofa de diez versos
\end{multicols}
\end{enumerate}
\vspace{0.15cm}
\item ¿Qué es el alejandrino?
\begin{enumerate}[label=\textbf{\Alph*)}]
\begin{multicols}{2}
  \item Un verso de catorce sílabas
  \item Un verso de doce sílabas
  \item Un verso de diez sílabas
  \item Un verso de ocho sílabas
  \item Un verso de seis sílabas
  \item Un verso de dieciséis sílabas
\end{multicols}
\end{enumerate}
\vspace{0.15cm}
\item ¿Qué es el cantar de gesta?
\begin{enumerate}[label=\textbf{\Alph*)}]
\begin{multicols}{2}
  \item Un poema épico medieval que narra hazañas heroicas
  \item Una novela de caballerías
  \item Una obra teatral
  \item Una lírica
  \item Un ensayo
  \item Una epopeya griega
\end{multicols}
\end{enumerate}
\vspace{0.15cm}
\item ¿Qué es el Romancero?
\begin{enumerate}[label=\textbf{\Alph*)}]
\begin{multicols}{2}
  \item Una colección de romances
  \item Una colección de cantares de gesta
  \item Una colección de novelas de caballerías
  \item Una colección de obras teatrales
  \item Una colección de líricas
  \item Una colección de ensayos
\end{multicols}
\end{enumerate}
\vspace{0.15cm}
\item ¿Qué es un romance?
\begin{enumerate}[label=\textbf{\Alph*)}]
\begin{multicols}{2}
  \item Un poema épico-lírico de origen medieval
  \item Una novela de caballerías
  \item Una obra teatral
  \item Una lírica
  \item Un ensayo
  \item Una epopeya griega
\end{multicols}
\end{enumerate}
\vspace{0.15cm}
\item ¿Qué es el romancero viejo?
\begin{enumerate}[label=\textbf{\Alph*)}]
\begin{multicols}{2}
  \item El conjunto de romances anónimos de la Edad Media
  \item El conjunto de romances de autores conocidos
  \item El conjunto de romances del siglo XVI
  \item El conjunto de romances del siglo XVII
  \item El conjunto de romances del siglo XVIII
  \item El conjunto de romances del siglo XIX
\end{multicols}
\end{enumerate}
\vspace{0.15cm}
\item ¿Qué es el romancero nuevo?
\begin{enumerate}[label=\textbf{\Alph*)}]
\begin{multicols}{2}
  \item El conjunto de romances de autores conocidos del siglo XVI y XVII
  \item El conjunto de romances anónimos de la Edad Media
  \item El conjunto de romances del siglo XVIII
  \item El conjunto de romances del siglo XIX
  \item El conjunto de romances del siglo XX
  \item El conjunto de romances del siglo XXI
\end{multicols}
\end{enumerate}
\vspace{0.15cm}
\item ¿Qué es la novela sentimental?
\begin{enumerate}[label=\textbf{\Alph*)}]
\begin{multicols}{2}
  \item Un género narrativo que combina amor y aventuras
  \item Un género teatral
  \item Un género lírico
  \item Un género didáctico
  \item Un género periodístico
  \item Un género ensayístico
\end{multicols}
\end{enumerate}
\vspace{0.15cm}
\item ¿Qué es la novela pastoril?
\begin{enumerate}[label=\textbf{\Alph*)}]
\begin{multicols}{2}
  \item Un género narrativo que idealiza la vida pastoril
  \item Un género teatral
  \item Un género lírico
  \item Un género didáctico
  \item Un género periodístico
  \item Un género ensayístico
\end{multicols}
\end{enumerate}
\vspace{0.15cm}
\item ¿Qué es la novela morisca?
\begin{enumerate}[label=\textbf{\Alph*)}]
\begin{multicols}{2}
  \item Un género narrativo que idealiza la vida de los moros
  \item Un género teatral
  \item Un género lírico
  \item Un género didáctico
  \item Un género periodístico
  \item Un género ensayístico
\end{multicols}
\end{enumerate}
\vspace{0.15cm}
\item ¿Qué es la novela bizantina?
\begin{enumerate}[label=\textbf{\Alph*)}]
\begin{multicols}{2}
  \item Un género narrativo de aventuras y viajes
  \item Un género teatral
  \item Un género lírico
  \item Un género didáctico
  \item Un género periodístico
  \item Un género ensayístico
\end{multicols}
\end{enumerate}
\vspace{0.15cm}
\item ¿Qué es la novela picaresca?
\begin{enumerate}[label=\textbf{\Alph*)}]
\begin{multicols}{2}
  \item Un género narrativo que narra la vida de un pícaro
  \item Un género teatral
  \item Un género lírico
  \item Un género didáctico
  \item Un género periodístico
  \item Un género ensayístico
\end{multicols}
\end{enumerate}
\vspace{0.15cm}
\end{enumerate}\n
\subsubsection*{SECCIÓN III: RENACIMIENTO (121–180)}

\begin{enumerate}[label=\textbf{\arabic*.}]
\setcounter{enumi}{120}
\item ¿Qué es el Renacimiento?
\begin{enumerate}[label=\textbf{\Alph*)}]
\begin{multicols}{2}
  \item Un movimiento cultural que recupera la cultura clásica
  \item Un movimiento que rechaza la cultura clásica
  \item Un movimiento que ignora la cultura clásica
  \item Un movimiento que imita la cultura medieval
  \item Un movimiento que rechaza la cultura medieval
  \item Un movimiento que ignora la cultura medieval
\end{multicols}
\end{enumerate}
\vspace{0.15cm}
\item ¿Dónde surge el Renacimiento?
\begin{enumerate}[label=\textbf{\Alph*)}]
\begin{multicols}{2}
  \item En Italia
  \item En Francia
  \item En España
  \item En Inglaterra
  \item En Alemania
  \item En Portugal
\end{multicols}
\end{enumerate}
\vspace{0.15cm}
\item ¿Qué es el humanismo?
\begin{enumerate}[label=\textbf{\Alph*)}]
\begin{multicols}{2}
  \item Un movimiento intelectual que pone al ser humano en el centro
  \item Un movimiento que pone a Dios en el centro
  \item Un movimiento que pone a la naturaleza en el centro
  \item Un movimiento que pone a la ciencia en el centro
  \item Un movimiento que pone al arte en el centro
  \item Un movimiento que pone a la política en el centro
\end{multicols}
\end{enumerate}
\vspace{0.15cm}
\item ¿Quién es el autor de Elogio de la locura?
\begin{enumerate}[label=\textbf{\Alph*)}]
\begin{multicols}{2}
  \item Erasmo de Rotterdam
  \item Tomás Moro
  \item Maquiavelo
  \item Montaigne
  \item Rabelais
  \item Shakespeare
\end{multicols}
\end{enumerate}
\vspace{0.15cm}
\item ¿Quién es el autor de Utopía?
\begin{enumerate}[label=\textbf{\Alph*)}]
\begin{multicols}{2}
  \item Tomás Moro
  \item Erasmo de Rotterdam
  \item Maquiavelo
  \item Montaigne
  \item Rabelais
  \item Shakespeare
\end{multicols}
\end{enumerate}
\vspace{0.15cm}
\item ¿Quién es el autor de El Príncipe?
\begin{enumerate}[label=\textbf{\Alph*)}]
\begin{multicols}{2}
  \item Maquiavelo
  \item Erasmo de Rotterdam
  \item Tomás Moro
  \item Montaigne
  \item Rabelais
  \item Shakespeare
\end{multicols}
\end{enumerate}
\vspace{0.15cm}
\item ¿Quién es el autor de Gargantúa y Pantagruel?
\begin{enumerate}[label=\textbf{\Alph*)}]
\begin{multicols}{2}
  \item Rabelais
  \item Erasmo de Rotterdam
  \item Tomás Moro
  \item Maquiavelo
  \item Montaigne
  \item Shakespeare
\end{multicols}
\end{enumerate}
\vspace{0.15cm}
\item ¿Quién es el autor de Ensayos?
\begin{enumerate}[label=\textbf{\Alph*)}]
\begin{multicols}{2}
  \item Montaigne
  \item Erasmo de Rotterdam
  \item Tomás Moro
  \item Maquiavelo
  \item Rabelais
  \item Shakespeare
\end{multicols}
\end{enumerate}
\vspace{0.15cm}
\item ¿Quién es el autor de Romeo y Julieta?
\begin{enumerate}[label=\textbf{\Alph*)}]
\begin{multicols}{2}
  \item Shakespeare
  \item Erasmo de Rotterdam
  \item Tomás Moro
  \item Maquiavelo
  \item Rabelais
  \item Montaigne
\end{multicols}
\end{enumerate}
\vspace{0.15cm}
\item ¿Quién es el autor de Hamlet?
\begin{enumerate}[label=\textbf{\Alph*)}]
\begin{multicols}{2}
  \item Shakespeare
  \item Erasmo de Rotterdam
  \item Tomás Moro
  \item Maquiavelo
  \item Rabelais
  \item Montaigne
\end{multicols}
\end{enumerate}
\vspace{0.15cm}
\item ¿Quién es el autor de Don Quijote de la Mancha?
\begin{enumerate}[label=\textbf{\Alph*)}]
\begin{multicols}{2}
  \item Cervantes
  \item Shakespeare
  \item Erasmo de Rotterdam
  \item Tomás Moro
  \item Maquiavelo
  \item Rabelais
\end{multicols}
\end{enumerate}
\vspace{0.15cm}
\item ¿Quién es el autor de La Galatea?
\begin{enumerate}[label=\textbf{\Alph*)}]
\begin{multicols}{2}
  \item Cervantes
  \item Shakespeare
  \item Erasmo de Rotterdam
  \item Tomás Moro
  \item Maquiavelo
  \item Rabelais
\end{multicols}
\end{enumerate}
\vspace{0.15cm}
\item ¿Quién es el autor de Novelas ejemplares?
\begin{enumerate}[label=\textbf{\Alph*)}]
\begin{multicols}{2}
  \item Cervantes
  \item Shakespeare
  \item Erasmo de Rotterdam
  \item Tomás Moro
  \item Maquiavelo
  \item Rabelais
\end{multicols}
\end{enumerate}
\vspace{0.15cm}
\item ¿Quién es el autor de La vida es sueño?
\begin{enumerate}[label=\textbf{\Alph*)}]
\begin{multicols}{2}
  \item Calderón de la Barca
  \item Lope de Vega
  \item Tirso de Molina
  \item Cervantes
  \item Shakespeare
  \item Quevedo
\end{multicols}
\end{enumerate}
\vspace{0.15cm}
\item ¿Quién es el autor de Fuenteovejuna?
\begin{enumerate}[label=\textbf{\Alph*)}]
\begin{multicols}{2}
  \item Lope de Vega
  \item Calderón de la Barca
  \item Tirso de Molina
  \item Cervantes
  \item Shakespeare
  \item Quevedo
\end{multicols}
\end{enumerate}
\vspace{0.15cm}
\item ¿Quién es el autor de El burlador de Sevilla?
\begin{enumerate}[label=\textbf{\Alph*)}]
\begin{multicols}{2}
  \item Tirso de Molina
  \item Lope de Vega
  \item Calderón de la Barca
  \item Cervantes
  \item Shakespeare
  \item Quevedo
\end{multicols}
\end{enumerate}
\vspace{0.15cm}
\item ¿Quién es el autor de La vida del Buscón?
\begin{enumerate}[label=\textbf{\Alph*)}]
\begin{multicols}{2}
  \item Quevedo
  \item Lope de Vega
  \item Calderón de la Barca
  \item Tirso de Molina
  \item Cervantes
  \item Shakespeare
\end{multicols}
\end{enumerate}
\vspace{0.15cm}
\item ¿Qué es el Renacimiento?
\begin{enumerate}[label=\textbf{\Alph*)}]
\begin{multicols}{2}
  \item Un movimiento cultural que recupera la cultura clásica
  \item Un movimiento que rechaza la cultura clásica
  \item Un movimiento que ignora la cultura clásica
  \item Un movimiento que imita la cultura medieval
  \item Un movimiento que rechaza la cultura medieval
  \item Un movimiento que ignora la cultura medieval
\end{multicols}
\end{enumerate}
\vspace{0.15cm}
\item ¿Dónde surge el Renacimiento?
\begin{enumerate}[label=\textbf{\Alph*)}]
\begin{multicols}{2}
  \item En Italia
  \item En Francia
  \item En España
  \item En Inglaterra
  \item En Alemania
  \item En Portugal
\end{multicols}
\end{enumerate}
\vspace{0.15cm}
\item ¿Qué es el humanismo?
\begin{enumerate}[label=\textbf{\Alph*)}]
\begin{multicols}{2}
  \item Un movimiento intelectual que pone al ser humano en el centro
  \item Un movimiento que pone a Dios en el centro
  \item Un movimiento que pone a la naturaleza en el centro
  \item Un movimiento que pone a la ciencia en el centro
  \item Un movimiento que pone al arte en el centro
  \item Un movimiento que pone a la política en el centro
\end{multicols}
\end{enumerate}
\vspace{0.15cm}
\item ¿Quién es el autor de Elogio de la locura?
\begin{enumerate}[label=\textbf{\Alph*)}]
\begin{multicols}{2}
  \item Erasmo de Rotterdam
  \item Tomás Moro
  \item Maquiavelo
  \item Montaigne
  \item Rabelais
  \item Shakespeare
\end{multicols}
\end{enumerate}
\vspace{0.15cm}
\item ¿Quién es el autor de Utopía?
\begin{enumerate}[label=\textbf{\Alph*)}]
\begin{multicols}{2}
  \item Tomás Moro
  \item Erasmo de Rotterdam
  \item Maquiavelo
  \item Montaigne
  \item Rabelais
  \item Shakespeare
\end{multicols}
\end{enumerate}
\vspace{0.15cm}
\item ¿Quién es el autor de El Príncipe?
\begin{enumerate}[label=\textbf{\Alph*)}]
\begin{multicols}{2}
  \item Maquiavelo
  \item Erasmo de Rotterdam
  \item Tomás Moro
  \item Montaigne
  \item Rabelais
  \item Shakespeare
\end{multicols}
\end{enumerate}
\vspace{0.15cm}
\item ¿Quién es el autor de Gargantúa y Pantagruel?
\begin{enumerate}[label=\textbf{\Alph*)}]
\begin{multicols}{2}
  \item Rabelais
  \item Erasmo de Rotterdam
  \item Tomás Moro
  \item Maquiavelo
  \item Montaigne
  \item Shakespeare
\end{multicols}
\end{enumerate}
\vspace{0.15cm}
\item ¿Quién es el autor de Ensayos?
\begin{enumerate}[label=\textbf{\Alph*)}]
\begin{multicols}{2}
  \item Montaigne
  \item Erasmo de Rotterdam
  \item Tomás Moro
  \item Maquiavelo
  \item Rabelais
  \item Shakespeare
\end{multicols}
\end{enumerate}
\vspace{0.15cm}
\item ¿Quién es el autor de Romeo y Julieta?
\begin{enumerate}[label=\textbf{\Alph*)}]
\begin{multicols}{2}
  \item Shakespeare
  \item Erasmo de Rotterdam
  \item Tomás Moro
  \item Maquiavelo
  \item Rabelais
  \item Montaigne
\end{multicols}
\end{enumerate}
\vspace{0.15cm}
\item ¿Quién es el autor de Hamlet?
\begin{enumerate}[label=\textbf{\Alph*)}]
\begin{multicols}{2}
  \item Shakespeare
  \item Erasmo de Rotterdam
  \item Tomás Moro
  \item Maquiavelo
  \item Rabelais
  \item Montaigne
\end{multicols}
\end{enumerate}
\vspace{0.15cm}
\item ¿Quién es el autor de Don Quijote de la Mancha?
\begin{enumerate}[label=\textbf{\Alph*)}]
\begin{multicols}{2}
  \item Cervantes
  \item Shakespeare
  \item Erasmo de Rotterdam
  \item Tomás Moro
  \item Maquiavelo
  \item Rabelais
\end{multicols}
\end{enumerate}
\vspace{0.15cm}
\item ¿Quién es el autor de La Galatea?
\begin{enumerate}[label=\textbf{\Alph*)}]
\begin{multicols}{2}
  \item Cervantes
  \item Shakespeare
  \item Erasmo de Rotterdam
  \item Tomás Moro
  \item Maquiavelo
  \item Rabelais
\end{multicols}
\end{enumerate}
\vspace{0.15cm}
\item ¿Quién es el autor de Novelas ejemplares?
\begin{enumerate}[label=\textbf{\Alph*)}]
\begin{multicols}{2}
  \item Cervantes
  \item Shakespeare
  \item Erasmo de Rotterdam
  \item Tomás Moro
  \item Maquiavelo
  \item Rabelais
\end{multicols}
\end{enumerate}
\vspace{0.15cm}
\item ¿Quién es el autor de La vida es sueño?
\begin{enumerate}[label=\textbf{\Alph*)}]
\begin{multicols}{2}
  \item Calderón de la Barca
  \item Lope de Vega
  \item Tirso de Molina
  \item Cervantes
  \item Shakespeare
  \item Quevedo
\end{multicols}
\end{enumerate}
\vspace{0.15cm}
\item ¿Quién es el autor de Fuenteovejuna?
\begin{enumerate}[label=\textbf{\Alph*)}]
\begin{multicols}{2}
  \item Lope de Vega
  \item Calderón de la Barca
  \item Tirso de Molina
  \item Cervantes
  \item Shakespeare
  \item Quevedo
\end{multicols}
\end{enumerate}
\vspace{0.15cm}
\item ¿Quién es el autor de El burlador de Sevilla?
\begin{enumerate}[label=\textbf{\Alph*)}]
\begin{multicols}{2}
  \item Tirso de Molina
  \item Lope de Vega
  \item Calderón de la Barca
  \item Cervantes
  \item Shakespeare
  \item Quevedo
\end{multicols}
\end{enumerate}
\vspace{0.15cm}
\item ¿Quién es el autor de La vida del Buscón?
\begin{enumerate}[label=\textbf{\Alph*)}]
\begin{multicols}{2}
  \item Quevedo
  \item Lope de Vega
  \item Calderón de la Barca
  \item Tirso de Molina
  \item Cervantes
  \item Shakespeare
\end{multicols}
\end{enumerate}
\vspace{0.15cm}
\item ¿Qué es el Renacimiento?
\begin{enumerate}[label=\textbf{\Alph*)}]
\begin{multicols}{2}
  \item Un movimiento cultural que recupera la cultura clásica
  \item Un movimiento que rechaza la cultura clásica
  \item Un movimiento que ignora la cultura clásica
  \item Un movimiento que imita la cultura medieval
  \item Un movimiento que rechaza la cultura medieval
  \item Un movimiento que ignora la cultura medieval
\end{multicols}
\end{enumerate}
\vspace{0.15cm}
\item ¿Dónde surge el Renacimiento?
\begin{enumerate}[label=\textbf{\Alph*)}]
\begin{multicols}{2}
  \item En Italia
  \item En Francia
  \item En España
  \item En Inglaterra
  \item En Alemania
  \item En Portugal
\end{multicols}
\end{enumerate}
\vspace{0.15cm}
\item ¿Qué es el humanismo?
\begin{enumerate}[label=\textbf{\Alph*)}]
\begin{multicols}{2}
  \item Un movimiento intelectual que pone al ser humano en el centro
  \item Un movimiento que pone a Dios en el centro
  \item Un movimiento que pone a la naturaleza en el centro
  \item Un movimiento que pone a la ciencia en el centro
  \item Un movimiento que pone al arte en el centro
  \item Un movimiento que pone a la política en el centro
\end{multicols}
\end{enumerate}
\vspace{0.15cm}
\item ¿Quién es el autor de Elogio de la locura?
\begin{enumerate}[label=\textbf{\Alph*)}]
\begin{multicols}{2}
  \item Erasmo de Rotterdam
  \item Tomás Moro
  \item Maquiavelo
  \item Montaigne
  \item Rabelais
  \item Shakespeare
\end{multicols}
\end{enumerate}
\vspace{0.15cm}
\item ¿Quién es el autor de Utopía?
\begin{enumerate}[label=\textbf{\Alph*)}]
\begin{multicols}{2}
  \item Tomás Moro
  \item Erasmo de Rotterdam
  \item Maquiavelo
  \item Montaigne
  \item Rabelais
  \item Shakespeare
\end{multicols}
\end{enumerate}
\vspace{0.15cm}
\item ¿Quién es el autor de El Príncipe?
\begin{enumerate}[label=\textbf{\Alph*)}]
\begin{multicols}{2}
  \item Maquiavelo
  \item Erasmo de Rotterdam
  \item Tomás Moro
  \item Montaigne
  \item Rabelais
  \item Shakespeare
\end{multicols}
\end{enumerate}
\vspace{0.15cm}
\item ¿Quién es el autor de Gargantúa y Pantagruel?
\begin{enumerate}[label=\textbf{\Alph*)}]
\begin{multicols}{2}
  \item Rabelais
  \item Erasmo de Rotterdam
  \item Tomás Moro
  \item Maquiavelo
  \item Montaigne
  \item Shakespeare
\end{multicols}
\end{enumerate}
\vspace{0.15cm}
\item ¿Quién es el autor de Ensayos?
\begin{enumerate}[label=\textbf{\Alph*)}]
\begin{multicols}{2}
  \item Montaigne
  \item Erasmo de Rotterdam
  \item Tomás Moro
  \item Maquiavelo
  \item Rabelais
  \item Shakespeare
\end{multicols}
\end{enumerate}
\vspace{0.15cm}
\item ¿Quién es el autor de Romeo y Julieta?
\begin{enumerate}[label=\textbf{\Alph*)}]
\begin{multicols}{2}
  \item Shakespeare
  \item Erasmo de Rotterdam
  \item Tomás Moro
  \item Maquiavelo
  \item Rabelais
  \item Montaigne
\end{multicols}
\end{enumerate}
\vspace{0.15cm}
\item ¿Quién es el autor de Hamlet?
\begin{enumerate}[label=\textbf{\Alph*)}]
\begin{multicols}{2}
  \item Shakespeare
  \item Erasmo de Rotterdam
  \item Tomás Moro
  \item Maquiavelo
  \item Rabelais
  \item Montaigne
\end{multicols}
\end{enumerate}
\vspace{0.15cm}
\item ¿Quién es el autor de Don Quijote de la Mancha?
\begin{enumerate}[label=\textbf{\Alph*)}]
\begin{multicols}{2}
  \item Cervantes
  \item Shakespeare
  \item Erasmo de Rotterdam
  \item Tomás Moro
  \item Maquiavelo
  \item Rabelais
\end{multicols}
\end{enumerate}
\vspace{0.15cm}
\item ¿Quién es el autor de La Galatea?
\begin{enumerate}[label=\textbf{\Alph*)}]
\begin{multicols}{2}
  \item Cervantes
  \item Shakespeare
  \item Erasmo de Rotterdam
  \item Tomás Moro
  \item Maquiavelo
  \item Rabelais
\end{multicols}
\end{enumerate}
\vspace{0.15cm}
\item ¿Quién es el autor de Novelas ejemplares?
\begin{enumerate}[label=\textbf{\Alph*)}]
\begin{multicols}{2}
  \item Cervantes
  \item Shakespeare
  \item Erasmo de Rotterdam
  \item Tomás Moro
  \item Maquiavelo
  \item Rabelais
\end{multicols}
\end{enumerate}
\vspace{0.15cm}
\item ¿Quién es el autor de La vida es sueño?
\begin{enumerate}[label=\textbf{\Alph*)}]
\begin{multicols}{2}
  \item Calderón de la Barca
  \item Lope de Vega
  \item Tirso de Molina
  \item Cervantes
  \item Shakespeare
  \item Quevedo
\end{multicols}
\end{enumerate}
\vspace{0.15cm}
\item ¿Quién es el autor de Fuenteovejuna?
\begin{enumerate}[label=\textbf{\Alph*)}]
\begin{multicols}{2}
  \item Lope de Vega
  \item Calderón de la Barca
  \item Tirso de Molina
  \item Cervantes
  \item Shakespeare
  \item Quevedo
\end{multicols}
\end{enumerate}
\vspace{0.15cm}
\item ¿Quién es el autor de El burlador de Sevilla?
\begin{enumerate}[label=\textbf{\Alph*)}]
\begin{multicols}{2}
  \item Tirso de Molina
  \item Lope de Vega
  \item Calderón de la Barca
  \item Cervantes
  \item Shakespeare
  \item Quevedo
\end{multicols}
\end{enumerate}
\vspace{0.15cm}
\item ¿Quién es el autor de La vida del Buscón?
\begin{enumerate}[label=\textbf{\Alph*)}]
\begin{multicols}{2}
  \item Quevedo
  \item Lope de Vega
  \item Calderón de la Barca
  \item Tirso de Molina
  \item Cervantes
  \item Shakespeare
\end{multicols}
\end{enumerate}
\vspace{0.15cm}
\item ¿Qué es el Renacimiento?
\begin{enumerate}[label=\textbf{\Alph*)}]
\begin{multicols}{2}
  \item Un movimiento cultural que recupera la cultura clásica
  \item Un movimiento que rechaza la cultura clásica
  \item Un movimiento que ignora la cultura clásica
  \item Un movimiento que imita la cultura medieval
  \item Un movimiento que rechaza la cultura medieval
  \item Un movimiento que ignora la cultura medieval
\end{multicols}
\end{enumerate}
\vspace{0.15cm}
\item ¿Dónde surge el Renacimiento?
\begin{enumerate}[label=\textbf{\Alph*)}]
\begin{multicols}{2}
  \item En Italia
  \item En Francia
  \item En España
  \item En Inglaterra
  \item En Alemania
  \item En Portugal
\end{multicols}
\end{enumerate}
\vspace{0.15cm}
\item ¿Qué es el humanismo?
\begin{enumerate}[label=\textbf{\Alph*)}]
\begin{multicols}{2}
  \item Un movimiento intelectual que pone al ser humano en el centro
  \item Un movimiento que pone a Dios en el centro
  \item Un movimiento que pone a la naturaleza en el centro
  \item Un movimiento que pone a la ciencia en el centro
  \item Un movimiento que pone al arte en el centro
  \item Un movimiento que pone a la política en el centro
\end{multicols}
\end{enumerate}
\vspace{0.15cm}
\item ¿Quién es el autor de Elogio de la locura?
\begin{enumerate}[label=\textbf{\Alph*)}]
\begin{multicols}{2}
  \item Erasmo de Rotterdam
  \item Tomás Moro
  \item Maquiavelo
  \item Montaigne
  \item Rabelais
  \item Shakespeare
\end{multicols}
\end{enumerate}
\vspace{0.15cm}
\item ¿Quién es el autor de Utopía?
\begin{enumerate}[label=\textbf{\Alph*)}]
\begin{multicols}{2}
  \item Tomás Moro
  \item Erasmo de Rotterdam
  \item Maquiavelo
  \item Montaigne
  \item Rabelais
  \item Shakespeare
\end{multicols}
\end{enumerate}
\vspace{0.15cm}
\item ¿Quién es el autor de El Príncipe?
\begin{enumerate}[label=\textbf{\Alph*)}]
\begin{multicols}{2}
  \item Maquiavelo
  \item Erasmo de Rotterdam
  \item Tomás Moro
  \item Montaigne
  \item Rabelais
  \item Shakespeare
\end{multicols}
\end{enumerate}
\vspace{0.15cm}
\item ¿Quién es el autor de Gargantúa y Pantagruel?
\begin{enumerate}[label=\textbf{\Alph*)}]
\begin{multicols}{2}
  \item Rabelais
  \item Erasmo de Rotterdam
  \item Tomás Moro
  \item Maquiavelo
  \item Montaigne
  \item Shakespeare
\end{multicols}
\end{enumerate}
\vspace{0.15cm}
\item ¿Quién es el autor de Ensayos?
\begin{enumerate}[label=\textbf{\Alph*)}]
\begin{multicols}{2}
  \item Montaigne
  \item Erasmo de Rotterdam
  \item Tomás Moro
  \item Maquiavelo
  \item Rabelais
  \item Shakespeare
\end{multicols}
\end{enumerate}
\vspace{0.15cm}
\item ¿Quién es el autor de Romeo y Julieta?
\begin{enumerate}[label=\textbf{\Alph*)}]
\begin{multicols}{2}
  \item Shakespeare
  \item Erasmo de Rotterdam
  \item Tomás Moro
  \item Maquiavelo
  \item Rabelais
  \item Montaigne
\end{multicols}
\end{enumerate}
\vspace{0.15cm}
\end{enumerate}\n
\subsubsection*{SECCIÓN IV: BARROCO (181–240)}

\begin{enumerate}[label=\textbf{\arabic*.}]
\setcounter{enumi}{180}
\item ¿Qué es el Barroco?
\begin{enumerate}[label=\textbf{\Alph*)}]
\begin{multicols}{2}
  \item Un movimiento cultural del siglo XVII caracterizado por el contraste y la complejidad
  \item Un movimiento del siglo XVI caracterizado por la armonía
  \item Un movimiento del siglo XVIII caracterizado por la razón
  \item Un movimiento del siglo XIX
  \item Un movimiento del siglo XX
  \item Un movimiento del siglo XXI
\end{multicols}
\end{enumerate}
\vspace{0.15cm}
\item ¿Cuál es el principal representante del Barroco español?
\begin{enumerate}[label=\textbf{\Alph*)}]
\begin{multicols}{2}
  \item Quevedo y Góngora
  \item Cervantes y Lope
  \item Calderón y Tirso
  \item Shakespeare y Milton
  \item Racine y Molière
  \item Corneille y Pascal
\end{multicols}
\end{enumerate}
\vspace{0.15cm}
\item ¿Qué es el conceptismo?
\begin{enumerate}[label=\textbf{\Alph*)}]
\begin{multicols}{2}
  \item Un estilo literario que busca la concisión y el ingenio
  \item Un estilo que busca la ornamentación
  \item Un estilo que busca la claridad
  \item Un estilo que busca la sencillez
  \item Un estilo que busca la naturalidad
  \item Un estilo que busca la objetividad
\end{multicols}
\end{enumerate}
\vspace{0.15cm}
\item ¿Qué es el culteranismo?
\begin{enumerate}[label=\textbf{\Alph*)}]
\begin{multicols}{2}
  \item Un estilo literario que busca la ornamentación y la belleza formal
  \item Un estilo que busca la concisión
  \item Un estilo que busca la claridad
  \item Un estilo que busca la sencillez
  \item Un estilo que busca la naturalidad
  \item Un estilo que busca la objetividad
\end{multicols}
\end{enumerate}
\vspace{0.15cm}
\item ¿Quién es el principal representante del conceptismo?
\begin{enumerate}[label=\textbf{\Alph*)}]
\begin{multicols}{2}
  \item Quevedo
  \item Góngora
  \item Cervantes
  \item Lope de Vega
  \item Calderón
  \item Tirso de Molina
\end{multicols}
\end{enumerate}
\vspace{0.15cm}
\item ¿Quién es el principal representante del culteranismo?
\begin{enumerate}[label=\textbf{\Alph*)}]
\begin{multicols}{2}
  \item Góngora
  \item Quevedo
  \item Cervantes
  \item Lope de Vega
  \item Calderón
  \item Tirso de Molina
\end{multicols}
\end{enumerate}
\vspace{0.15cm}
\item ¿Quién es el autor de Soledades?
\begin{enumerate}[label=\textbf{\Alph*)}]
\begin{multicols}{2}
  \item Góngora
  \item Quevedo
  \item Cervantes
  \item Lope de Vega
  \item Calderón
  \item Tirso de Molina
\end{multicols}
\end{enumerate}
\vspace{0.15cm}
\item ¿Quién es el autor de Sueños?
\begin{enumerate}[label=\textbf{\Alph*)}]
\begin{multicols}{2}
  \item Quevedo
  \item Góngora
  \item Cervantes
  \item Lope de Vega
  \item Calderón
  \item Tirso de Molina
\end{multicols}
\end{enumerate}
\vspace{0.15cm}
\item ¿Quién es el autor de El Criticón?
\begin{enumerate}[label=\textbf{\Alph*)}]
\begin{multicols}{2}
  \item Baltasar Gracián
  \item Quevedo
  \item Góngora
  \item Cervantes
  \item Lope de Vega
  \item Calderón
\end{multicols}
\end{enumerate}
\vspace{0.15cm}
\item ¿Quién es el autor de Paradise Lost?
\begin{enumerate}[label=\textbf{\Alph*)}]
\begin{multicols}{2}
  \item John Milton
  \item Shakespeare
  \item Quevedo
  \item Góngora
  \item Cervantes
  \item Lope de Vega
\end{multicols}
\end{enumerate}
\vspace{0.15cm}
\item ¿Quién es el autor de El paraíso perdido?
\begin{enumerate}[label=\textbf{\Alph*)}]
\begin{multicols}{2}
  \item John Milton
  \item Shakespeare
  \item Quevedo
  \item Góngora
  \item Cervantes
  \item Lope de Vega
\end{multicols}
\end{enumerate}
\vspace{0.15cm}
\item ¿Quién es el autor de Las meninas?
\begin{enumerate}[label=\textbf{\Alph*)}]
\begin{multicols}{2}
  \item Velázquez
  \item Goya
  \item El Greco
  \item Murillo
  \item Zurbarán
  \item Ribera
\end{multicols}
\end{enumerate}
\vspace{0.15cm}
\item ¿Quién es el autor de El sueño de la razón produce monstruos?
\begin{enumerate}[label=\textbf{\Alph*)}]
\begin{multicols}{2}
  \item Goya
  \item Velázquez
  \item El Greco
  \item Murillo
  \item Zurbarán
  \item Ribera
\end{multicols}
\end{enumerate}
\vspace{0.15cm}
\item ¿Qué es el Siglo de Oro?
\begin{enumerate}[label=\textbf{\Alph*)}]
\begin{multicols}{2}
  \item El período de máximo esplendor de la literatura española
  \item El período francés
  \item El período inglés
  \item El período italiano
  \item El período alemán
  \item El período portugués
\end{multicols}
\end{enumerate}
\vspace{0.15cm}
\item ¿Qué es el teatro barroco?
\begin{enumerate}[label=\textbf{\Alph*)}]
\begin{multicols}{2}
  \item Un teatro caracterizado por la complejidad y el contraste
  \item Un teatro de sencillez
  \item Un teatro de claridad
  \item Un teatro naturalista
  \item Un teatro objetivo
  \item Un teatro de armonía
\end{multicols}
\end{enumerate}
\vspace{0.15cm}
\item ¿Qué es el Barroco?
\begin{enumerate}[label=\textbf{\Alph*)}]
\begin{multicols}{2}
  \item Un movimiento cultural del siglo XVII caracterizado por el contraste y la complejidad
  \item Un movimiento del siglo XVI caracterizado por la armonía
  \item Un movimiento del siglo XVIII caracterizado por la razón
  \item Un movimiento del siglo XIX
  \item Un movimiento del siglo XX
  \item Un movimiento del siglo XXI
\end{multicols}
\end{enumerate}
\vspace{0.15cm}
\item ¿Cuál es el principal representante del Barroco español?
\begin{enumerate}[label=\textbf{\Alph*)}]
\begin{multicols}{2}
  \item Quevedo y Góngora
  \item Cervantes y Lope
  \item Calderón y Tirso
  \item Shakespeare y Milton
  \item Racine y Molière
  \item Corneille y Pascal
\end{multicols}
\end{enumerate}
\vspace{0.15cm}
\item ¿Qué es el conceptismo?
\begin{enumerate}[label=\textbf{\Alph*)}]
\begin{multicols}{2}
  \item Un estilo literario que busca la concisión y el ingenio
  \item Un estilo que busca la ornamentación
  \item Un estilo que busca la claridad
  \item Un estilo que busca la sencillez
  \item Un estilo que busca la naturalidad
  \item Un estilo que busca la objetividad
\end{multicols}
\end{enumerate}
\vspace{0.15cm}
\item ¿Qué es el culteranismo?
\begin{enumerate}[label=\textbf{\Alph*)}]
\begin{multicols}{2}
  \item Un estilo literario que busca la ornamentación y la belleza formal
  \item Un estilo que busca la concisión
  \item Un estilo que busca la claridad
  \item Un estilo que busca la sencillez
  \item Un estilo que busca la naturalidad
  \item Un estilo que busca la objetividad
\end{multicols}
\end{enumerate}
\vspace{0.15cm}
\item ¿Quién es el principal representante del conceptismo?
\begin{enumerate}[label=\textbf{\Alph*)}]
\begin{multicols}{2}
  \item Quevedo
  \item Góngora
  \item Cervantes
  \item Lope de Vega
  \item Calderón
  \item Tirso de Molina
\end{multicols}
\end{enumerate}
\vspace{0.15cm}
\item ¿Quién es el principal representante del culteranismo?
\begin{enumerate}[label=\textbf{\Alph*)}]
\begin{multicols}{2}
  \item Góngora
  \item Quevedo
  \item Cervantes
  \item Lope de Vega
  \item Calderón
  \item Tirso de Molina
\end{multicols}
\end{enumerate}
\vspace{0.15cm}
\item ¿Quién es el autor de Soledades?
\begin{enumerate}[label=\textbf{\Alph*)}]
\begin{multicols}{2}
  \item Góngora
  \item Quevedo
  \item Cervantes
  \item Lope de Vega
  \item Calderón
  \item Tirso de Molina
\end{multicols}
\end{enumerate}
\vspace{0.15cm}
\item ¿Quién es el autor de Sueños?
\begin{enumerate}[label=\textbf{\Alph*)}]
\begin{multicols}{2}
  \item Quevedo
  \item Góngora
  \item Cervantes
  \item Lope de Vega
  \item Calderón
  \item Tirso de Molina
\end{multicols}
\end{enumerate}
\vspace{0.15cm}
\item ¿Quién es el autor de El Criticón?
\begin{enumerate}[label=\textbf{\Alph*)}]
\begin{multicols}{2}
  \item Baltasar Gracián
  \item Quevedo
  \item Góngora
  \item Cervantes
  \item Lope de Vega
  \item Calderón
\end{multicols}
\end{enumerate}
\vspace{0.15cm}
\item ¿Quién es el autor de Paradise Lost?
\begin{enumerate}[label=\textbf{\Alph*)}]
\begin{multicols}{2}
  \item John Milton
  \item Shakespeare
  \item Quevedo
  \item Góngora
  \item Cervantes
  \item Lope de Vega
\end{multicols}
\end{enumerate}
\vspace{0.15cm}
\item ¿Quién es el autor de El paraíso perdido?
\begin{enumerate}[label=\textbf{\Alph*)}]
\begin{multicols}{2}
  \item John Milton
  \item Shakespeare
  \item Quevedo
  \item Góngora
  \item Cervantes
  \item Lope de Vega
\end{multicols}
\end{enumerate}
\vspace{0.15cm}
\item ¿Quién es el autor de Las meninas?
\begin{enumerate}[label=\textbf{\Alph*)}]
\begin{multicols}{2}
  \item Velázquez
  \item Goya
  \item El Greco
  \item Murillo
  \item Zurbarán
  \item Ribera
\end{multicols}
\end{enumerate}
\vspace{0.15cm}
\item ¿Quién es el autor de El sueño de la razón produce monstruos?
\begin{enumerate}[label=\textbf{\Alph*)}]
\begin{multicols}{2}
  \item Goya
  \item Velázquez
  \item El Greco
  \item Murillo
  \item Zurbarán
  \item Ribera
\end{multicols}
\end{enumerate}
\vspace{0.15cm}
\item ¿Qué es el Siglo de Oro?
\begin{enumerate}[label=\textbf{\Alph*)}]
\begin{multicols}{2}
  \item El período de máximo esplendor de la literatura española
  \item El período francés
  \item El período inglés
  \item El período italiano
  \item El período alemán
  \item El período portugués
\end{multicols}
\end{enumerate}
\vspace{0.15cm}
\item ¿Qué es el teatro barroco?
\begin{enumerate}[label=\textbf{\Alph*)}]
\begin{multicols}{2}
  \item Un teatro caracterizado por la complejidad y el contraste
  \item Un teatro de sencillez
  \item Un teatro de claridad
  \item Un teatro naturalista
  \item Un teatro objetivo
  \item Un teatro de armonía
\end{multicols}
\end{enumerate}
\vspace{0.15cm}
\item ¿Qué es el Barroco?
\begin{enumerate}[label=\textbf{\Alph*)}]
\begin{multicols}{2}
  \item Un movimiento cultural del siglo XVII caracterizado por el contraste y la complejidad
  \item Un movimiento del siglo XVI caracterizado por la armonía
  \item Un movimiento del siglo XVIII caracterizado por la razón
  \item Un movimiento del siglo XIX
  \item Un movimiento del siglo XX
  \item Un movimiento del siglo XXI
\end{multicols}
\end{enumerate}
\vspace{0.15cm}
\item ¿Cuál es el principal representante del Barroco español?
\begin{enumerate}[label=\textbf{\Alph*)}]
\begin{multicols}{2}
  \item Quevedo y Góngora
  \item Cervantes y Lope
  \item Calderón y Tirso
  \item Shakespeare y Milton
  \item Racine y Molière
  \item Corneille y Pascal
\end{multicols}
\end{enumerate}
\vspace{0.15cm}
\item ¿Qué es el conceptismo?
\begin{enumerate}[label=\textbf{\Alph*)}]
\begin{multicols}{2}
  \item Un estilo literario que busca la concisión y el ingenio
  \item Un estilo que busca la ornamentación
  \item Un estilo que busca la claridad
  \item Un estilo que busca la sencillez
  \item Un estilo que busca la naturalidad
  \item Un estilo que busca la objetividad
\end{multicols}
\end{enumerate}
\vspace{0.15cm}
\item ¿Qué es el culteranismo?
\begin{enumerate}[label=\textbf{\Alph*)}]
\begin{multicols}{2}
  \item Un estilo literario que busca la ornamentación y la belleza formal
  \item Un estilo que busca la concisión
  \item Un estilo que busca la claridad
  \item Un estilo que busca la sencillez
  \item Un estilo que busca la naturalidad
  \item Un estilo que busca la objetividad
\end{multicols}
\end{enumerate}
\vspace{0.15cm}
\item ¿Quién es el principal representante del conceptismo?
\begin{enumerate}[label=\textbf{\Alph*)}]
\begin{multicols}{2}
  \item Quevedo
  \item Góngora
  \item Cervantes
  \item Lope de Vega
  \item Calderón
  \item Tirso de Molina
\end{multicols}
\end{enumerate}
\vspace{0.15cm}
\item ¿Quién es el principal representante del culteranismo?
\begin{enumerate}[label=\textbf{\Alph*)}]
\begin{multicols}{2}
  \item Góngora
  \item Quevedo
  \item Cervantes
  \item Lope de Vega
  \item Calderón
  \item Tirso de Molina
\end{multicols}
\end{enumerate}
\vspace{0.15cm}
\item ¿Quién es el autor de Soledades?
\begin{enumerate}[label=\textbf{\Alph*)}]
\begin{multicols}{2}
  \item Góngora
  \item Quevedo
  \item Cervantes
  \item Lope de Vega
  \item Calderón
  \item Tirso de Molina
\end{multicols}
\end{enumerate}
\vspace{0.15cm}
\item ¿Quién es el autor de Sueños?
\begin{enumerate}[label=\textbf{\Alph*)}]
\begin{multicols}{2}
  \item Quevedo
  \item Góngora
  \item Cervantes
  \item Lope de Vega
  \item Calderón
  \item Tirso de Molina
\end{multicols}
\end{enumerate}
\vspace{0.15cm}
\item ¿Quién es el autor de El Criticón?
\begin{enumerate}[label=\textbf{\Alph*)}]
\begin{multicols}{2}
  \item Baltasar Gracián
  \item Quevedo
  \item Góngora
  \item Cervantes
  \item Lope de Vega
  \item Calderón
\end{multicols}
\end{enumerate}
\vspace{0.15cm}
\item ¿Quién es el autor de Paradise Lost?
\begin{enumerate}[label=\textbf{\Alph*)}]
\begin{multicols}{2}
  \item John Milton
  \item Shakespeare
  \item Quevedo
  \item Góngora
  \item Cervantes
  \item Lope de Vega
\end{multicols}
\end{enumerate}
\vspace{0.15cm}
\item ¿Quién es el autor de El paraíso perdido?
\begin{enumerate}[label=\textbf{\Alph*)}]
\begin{multicols}{2}
  \item John Milton
  \item Shakespeare
  \item Quevedo
  \item Góngora
  \item Cervantes
  \item Lope de Vega
\end{multicols}
\end{enumerate}
\vspace{0.15cm}
\item ¿Quién es el autor de Las meninas?
\begin{enumerate}[label=\textbf{\Alph*)}]
\begin{multicols}{2}
  \item Velázquez
  \item Goya
  \item El Greco
  \item Murillo
  \item Zurbarán
  \item Ribera
\end{multicols}
\end{enumerate}
\vspace{0.15cm}
\item ¿Quién es el autor de El sueño de la razón produce monstruos?
\begin{enumerate}[label=\textbf{\Alph*)}]
\begin{multicols}{2}
  \item Goya
  \item Velázquez
  \item El Greco
  \item Murillo
  \item Zurbarán
  \item Ribera
\end{multicols}
\end{enumerate}
\vspace{0.15cm}
\item ¿Qué es el Siglo de Oro?
\begin{enumerate}[label=\textbf{\Alph*)}]
\begin{multicols}{2}
  \item El período de máximo esplendor de la literatura española
  \item El período francés
  \item El período inglés
  \item El período italiano
  \item El período alemán
  \item El período portugués
\end{multicols}
\end{enumerate}
\vspace{0.15cm}
\item ¿Qué es el teatro barroco?
\begin{enumerate}[label=\textbf{\Alph*)}]
\begin{multicols}{2}
  \item Un teatro caracterizado por la complejidad y el contraste
  \item Un teatro de sencillez
  \item Un teatro de claridad
  \item Un teatro naturalista
  \item Un teatro objetivo
  \item Un teatro de armonía
\end{multicols}
\end{enumerate}
\vspace{0.15cm}
\item ¿Qué es el Barroco?
\begin{enumerate}[label=\textbf{\Alph*)}]
\begin{multicols}{2}
  \item Un movimiento cultural del siglo XVII caracterizado por el contraste y la complejidad
  \item Un movimiento del siglo XVI caracterizado por la armonía
  \item Un movimiento del siglo XVIII caracterizado por la razón
  \item Un movimiento del siglo XIX
  \item Un movimiento del siglo XX
  \item Un movimiento del siglo XXI
\end{multicols}
\end{enumerate}
\vspace{0.15cm}
\item ¿Cuál es el principal representante del Barroco español?
\begin{enumerate}[label=\textbf{\Alph*)}]
\begin{multicols}{2}
  \item Quevedo y Góngora
  \item Cervantes y Lope
  \item Calderón y Tirso
  \item Shakespeare y Milton
  \item Racine y Molière
  \item Corneille y Pascal
\end{multicols}
\end{enumerate}
\vspace{0.15cm}
\item ¿Qué es el conceptismo?
\begin{enumerate}[label=\textbf{\Alph*)}]
\begin{multicols}{2}
  \item Un estilo literario que busca la concisión y el ingenio
  \item Un estilo que busca la ornamentación
  \item Un estilo que busca la claridad
  \item Un estilo que busca la sencillez
  \item Un estilo que busca la naturalidad
  \item Un estilo que busca la objetividad
\end{multicols}
\end{enumerate}
\vspace{0.15cm}
\item ¿Qué es el culteranismo?
\begin{enumerate}[label=\textbf{\Alph*)}]
\begin{multicols}{2}
  \item Un estilo literario que busca la ornamentación y la belleza formal
  \item Un estilo que busca la concisión
  \item Un estilo que busca la claridad
  \item Un estilo que busca la sencillez
  \item Un estilo que busca la naturalidad
  \item Un estilo que busca la objetividad
\end{multicols}
\end{enumerate}
\vspace{0.15cm}
\item ¿Quién es el principal representante del conceptismo?
\begin{enumerate}[label=\textbf{\Alph*)}]
\begin{multicols}{2}
  \item Quevedo
  \item Góngora
  \item Cervantes
  \item Lope de Vega
  \item Calderón
  \item Tirso de Molina
\end{multicols}
\end{enumerate}
\vspace{0.15cm}
\item ¿Quién es el principal representante del culteranismo?
\begin{enumerate}[label=\textbf{\Alph*)}]
\begin{multicols}{2}
  \item Góngora
  \item Quevedo
  \item Cervantes
  \item Lope de Vega
  \item Calderón
  \item Tirso de Molina
\end{multicols}
\end{enumerate}
\vspace{0.15cm}
\item ¿Quién es el autor de Soledades?
\begin{enumerate}[label=\textbf{\Alph*)}]
\begin{multicols}{2}
  \item Góngora
  \item Quevedo
  \item Cervantes
  \item Lope de Vega
  \item Calderón
  \item Tirso de Molina
\end{multicols}
\end{enumerate}
\vspace{0.15cm}
\item ¿Quién es el autor de Sueños?
\begin{enumerate}[label=\textbf{\Alph*)}]
\begin{multicols}{2}
  \item Quevedo
  \item Góngora
  \item Cervantes
  \item Lope de Vega
  \item Calderón
  \item Tirso de Molina
\end{multicols}
\end{enumerate}
\vspace{0.15cm}
\item ¿Quién es el autor de El Criticón?
\begin{enumerate}[label=\textbf{\Alph*)}]
\begin{multicols}{2}
  \item Baltasar Gracián
  \item Quevedo
  \item Góngora
  \item Cervantes
  \item Lope de Vega
  \item Calderón
\end{multicols}
\end{enumerate}
\vspace{0.15cm}
\item ¿Quién es el autor de Paradise Lost?
\begin{enumerate}[label=\textbf{\Alph*)}]
\begin{multicols}{2}
  \item John Milton
  \item Shakespeare
  \item Quevedo
  \item Góngora
  \item Cervantes
  \item Lope de Vega
\end{multicols}
\end{enumerate}
\vspace{0.15cm}
\item ¿Quién es el autor de El paraíso perdido?
\begin{enumerate}[label=\textbf{\Alph*)}]
\begin{multicols}{2}
  \item John Milton
  \item Shakespeare
  \item Quevedo
  \item Góngora
  \item Cervantes
  \item Lope de Vega
\end{multicols}
\end{enumerate}
\vspace{0.15cm}
\item ¿Quién es el autor de Las meninas?
\begin{enumerate}[label=\textbf{\Alph*)}]
\begin{multicols}{2}
  \item Velázquez
  \item Goya
  \item El Greco
  \item Murillo
  \item Zurbarán
  \item Ribera
\end{multicols}
\end{enumerate}
\vspace{0.15cm}
\item ¿Quién es el autor de El sueño de la razón produce monstruos?
\begin{enumerate}[label=\textbf{\Alph*)}]
\begin{multicols}{2}
  \item Goya
  \item Velázquez
  \item El Greco
  \item Murillo
  \item Zurbarán
  \item Ribera
\end{multicols}
\end{enumerate}
\vspace{0.15cm}
\item ¿Qué es el Siglo de Oro?
\begin{enumerate}[label=\textbf{\Alph*)}]
\begin{multicols}{2}
  \item El período de máximo esplendor de la literatura española
  \item El período francés
  \item El período inglés
  \item El período italiano
  \item El período alemán
  \item El período portugués
\end{multicols}
\end{enumerate}
\vspace{0.15cm}
\item ¿Qué es el teatro barroco?
\begin{enumerate}[label=\textbf{\Alph*)}]
\begin{multicols}{2}
  \item Un teatro caracterizado por la complejidad y el contraste
  \item Un teatro de sencillez
  \item Un teatro de claridad
  \item Un teatro naturalista
  \item Un teatro objetivo
  \item Un teatro de armonía
\end{multicols}
\end{enumerate}
\vspace{0.15cm}
\end{enumerate}\n
\subsubsection*{SECCIÓN V: ILUSTRACIÓN Y NEOCLASICISMO (241–300)}

\begin{enumerate}[label=\textbf{\arabic*.}]
\setcounter{enumi}{240}
\item ¿Qué es la Ilustración?
\begin{enumerate}[label=\textbf{\Alph*)}]
\begin{multicols}{2}
  \item Un movimiento cultural del siglo XVIII que exalta la razón
  \item Un movimiento del siglo XVII que exalta el sentimiento
  \item Un movimiento que exalta la fe
  \item Un movimiento que exalta la imaginación
  \item Un movimiento de experimentación
  \item Un movimiento de tecnología
\end{multicols}
\end{enumerate}
\vspace{0.15cm}
\item ¿Cuál es el lema de la Ilustración?
\begin{enumerate}[label=\textbf{\Alph*)}]
\begin{multicols}{2}
  \item Sapere aude (Atrévete a saber)
  \item Carpe diem (Aprovecha el día)
  \item Memento mori (Recuerda que morirás)
  \item Cogito ergo sum (Pienso, luego existo)
  \item Veni, vidi, vici (Vine, vi, vencí)
  \item Amor fati (Amor al destino)
\end{multicols}
\end{enumerate}
\vspace{0.15cm}
\item ¿Qué es el Neoclasicismo?
\begin{enumerate}[label=\textbf{\Alph*)}]
\begin{multicols}{2}
  \item Un movimiento literario que imita los modelos clásicos
  \item Un movimiento que rechaza los modelos clásicos
  \item Un movimiento que ignora los modelos clásicos
  \item Un movimiento que imita modelos medievales
  \item Un movimiento que rechaza lo medieval
  \item Un movimiento que ignora lo medieval
\end{multicols}
\end{enumerate}
\vspace{0.15cm}
\item ¿Cuál es el principio fundamental del Neoclasicismo?
\begin{enumerate}[label=\textbf{\Alph*)}]
\begin{multicols}{2}
  \item La razón y el buen gusto
  \item El sentimiento y la imaginación
  \item La fe y la tradición
  \item La experimentación y originalidad
  \item La libertad y la rebeldía
  \item La intuición y el instinto
\end{multicols}
\end{enumerate}
\vspace{0.15cm}
\item ¿Quién es el autor de El sí de las niñas?
\begin{enumerate}[label=\textbf{\Alph*)}]
\begin{multicols}{2}
  \item Leandro Fernández de Moratín
  \item José Cadalso
  \item Gaspar Melchor de Jovellanos
  \item Tomás de Iriarte
  \item Félix María Samaniego
  \item Juan Meléndez Valdés
\end{multicols}
\end{enumerate}
\vspace{0.15cm}
\item ¿Quién es el autor de Cartas marruecas?
\begin{enumerate}[label=\textbf{\Alph*)}]
\begin{multicols}{2}
  \item José Cadalso
  \item Leandro Fernández de Moratín
  \item Gaspar Melchor de Jovellanos
  \item Tomás de Iriarte
  \item Félix María Samaniego
  \item Juan Meléndez Valdés
\end{multicols}
\end{enumerate}
\vspace{0.15cm}
\item ¿Quién es el autor de Fábulas?
\begin{enumerate}[label=\textbf{\Alph*)}]
\begin{multicols}{2}
  \item Tomás de Iriarte
  \item Leandro Fernández de Moratín
  \item José Cadalso
  \item Gaspar Melchor de Jovellanos
  \item Félix María Samaniego
  \item Juan Meléndez Valdés
\end{multicols}
\end{enumerate}
\vspace{0.15cm}
\item ¿Quién es el autor de Fábulas?
\begin{enumerate}[label=\textbf{\Alph*)}]
\begin{multicols}{2}
  \item Félix María Samaniego
  \item Leandro Fernández de Moratín
  \item José Cadalso
  \item Gaspar Melchor de Jovellanos
  \item Tomás de Iriarte
  \item Juan Meléndez Valdés
\end{multicols}
\end{enumerate}
\vspace{0.15cm}
\item ¿Quién es el autor de Ensayo sobre el hombre?
\begin{enumerate}[label=\textbf{\Alph*)}]
\begin{multicols}{2}
  \item Alexander Pope
  \item Jonathan Swift
  \item Daniel Defoe
  \item Samuel Johnson
  \item Voltaire
  \item Rousseau
\end{multicols}
\end{enumerate}
\vspace{0.15cm}
\item ¿Quién es el autor de Los viajes de Gulliver?
\begin{enumerate}[label=\textbf{\Alph*)}]
\begin{multicols}{2}
  \item Jonathan Swift
  \item Alexander Pope
  \item Daniel Defoe
  \item Samuel Johnson
  \item Voltaire
  \item Rousseau
\end{multicols}
\end{enumerate}
\vspace{0.15cm}
\item ¿Quién es el autor de Robinson Crusoe?
\begin{enumerate}[label=\textbf{\Alph*)}]
\begin{multicols}{2}
  \item Daniel Defoe
  \item Jonathan Swift
  \item Alexander Pope
  \item Samuel Johnson
  \item Voltaire
  \item Rousseau
\end{multicols}
\end{enumerate}
\vspace{0.15cm}
\item ¿Quién es el autor de Cándido?
\begin{enumerate}[label=\textbf{\Alph*)}]
\begin{multicols}{2}
  \item Voltaire
  \item Rousseau
  \item Montesquieu
  \item Diderot
  \item D'Alembert
  \item Buffon
\end{multicols}
\end{enumerate}
\vspace{0.15cm}
\item ¿Quién es el autor de El contrato social?
\begin{enumerate}[label=\textbf{\Alph*)}]
\begin{multicols}{2}
  \item Rousseau
  \item Voltaire
  \item Montesquieu
  \item Diderot
  \item D'Alembert
  \item Buffon
\end{multicols}
\end{enumerate}
\vspace{0.15cm}
\item ¿Quién es el autor de El espíritu de las leyes?
\begin{enumerate}[label=\textbf{\Alph*)}]
\begin{multicols}{2}
  \item Montesquieu
  \item Voltaire
  \item Rousseau
  \item Diderot
  \item D'Alembert
  \item Buffon
\end{multicols}
\end{enumerate}
\vspace{0.15cm}
\item ¿Quién es el autor de La enciclopedia?
\begin{enumerate}[label=\textbf{\Alph*)}]
\begin{multicols}{2}
  \item Diderot y D'Alembert
  \item Voltaire y Rousseau
  \item Montesquieu y Diderot
  \item D'Alembert y Buffon
  \item Voltaire y Montesquieu
  \item Rousseau y Diderot
\end{multicols}
\end{enumerate}
\vspace{0.15cm}
\item ¿Qué es la Enciclopedia?
\begin{enumerate}[label=\textbf{\Alph*)}]
\begin{multicols}{2}
  \item Una obra que recopila todo el saber de la época
  \item Una novela
  \item Una obra teatral
  \item Una lírica
  \item Un ensayo
  \item Una epopeya
\end{multicols}
\end{enumerate}
\vspace{0.15cm}
\item ¿Qué es el despotismo ilustrado?
\begin{enumerate}[label=\textbf{\Alph*)}]
\begin{multicols}{2}
  \item Una forma de gobierno que combina el poder absoluto con reformas ilustradas
  \item Una forma democrática
  \item Una forma republicana
  \item Una forma anarquista
  \item Una forma socialista
  \item Una forma comunista
\end{multicols}
\end{enumerate}
\vspace{0.15cm}
\item ¿Qué es el teatro neoclásico?
\begin{enumerate}[label=\textbf{\Alph*)}]
\begin{multicols}{2}
  \item Un teatro que sigue las reglas clásicas de unidad de acción, lugar y tiempo
  \item Un teatro que rechaza las reglas clásicas
  \item Un teatro que ignora las reglas clásicas
  \item Un teatro medieval
  \item Un teatro barroco
  \item Un teatro romántico
\end{multicols}
\end{enumerate}
\vspace{0.15cm}
\item ¿Cuáles son las tres unidades del teatro neoclásico?
\begin{enumerate}[label=\textbf{\Alph*)}]
\begin{multicols}{2}
  \item Acción, lugar y tiempo
  \item Acción, lugar y espacio
  \item Acción, tiempo y espacio
  \item Lugar, tiempo y espacio
  \item Acción, lugar y personaje
  \item Acción, tiempo y personaje
\end{multicols}
\end{enumerate}
\vspace{0.15cm}
\item ¿Qué es la Ilustración?
\begin{enumerate}[label=\textbf{\Alph*)}]
\begin{multicols}{2}
  \item Un movimiento cultural del siglo XVIII que exalta la razón
  \item Un movimiento del siglo XVII que exalta el sentimiento
  \item Un movimiento que exalta la fe
  \item Un movimiento que exalta la imaginación
  \item Un movimiento de experimentación
  \item Un movimiento de tecnología
\end{multicols}
\end{enumerate}
\vspace{0.15cm}
\item ¿Cuál es el lema de la Ilustración?
\begin{enumerate}[label=\textbf{\Alph*)}]
\begin{multicols}{2}
  \item Sapere aude (Atrévete a saber)
  \item Carpe diem (Aprovecha el día)
  \item Memento mori (Recuerda que morirás)
  \item Cogito ergo sum (Pienso, luego existo)
  \item Veni, vidi, vici (Vine, vi, vencí)
  \item Amor fati (Amor al destino)
\end{multicols}
\end{enumerate}
\vspace{0.15cm}
\item ¿Qué es el Neoclasicismo?
\begin{enumerate}[label=\textbf{\Alph*)}]
\begin{multicols}{2}
  \item Un movimiento literario que imita los modelos clásicos
  \item Un movimiento que rechaza los modelos clásicos
  \item Un movimiento que ignora los modelos clásicos
  \item Un movimiento que imita modelos medievales
  \item Un movimiento que rechaza lo medieval
  \item Un movimiento que ignora lo medieval
\end{multicols}
\end{enumerate}
\vspace{0.15cm}
\item ¿Cuál es el principio fundamental del Neoclasicismo?
\begin{enumerate}[label=\textbf{\Alph*)}]
\begin{multicols}{2}
  \item La razón y el buen gusto
  \item El sentimiento y la imaginación
  \item La fe y la tradición
  \item La experimentación y originalidad
  \item La libertad y la rebeldía
  \item La intuición y el instinto
\end{multicols}
\end{enumerate}
\vspace{0.15cm}
\item ¿Quién es el autor de El sí de las niñas?
\begin{enumerate}[label=\textbf{\Alph*)}]
\begin{multicols}{2}
  \item Leandro Fernández de Moratín
  \item José Cadalso
  \item Gaspar Melchor de Jovellanos
  \item Tomás de Iriarte
  \item Félix María Samaniego
  \item Juan Meléndez Valdés
\end{multicols}
\end{enumerate}
\vspace{0.15cm}
\item ¿Quién es el autor de Cartas marruecas?
\begin{enumerate}[label=\textbf{\Alph*)}]
\begin{multicols}{2}
  \item José Cadalso
  \item Leandro Fernández de Moratín
  \item Gaspar Melchor de Jovellanos
  \item Tomás de Iriarte
  \item Félix María Samaniego
  \item Juan Meléndez Valdés
\end{multicols}
\end{enumerate}
\vspace{0.15cm}
\item ¿Quién es el autor de Fábulas?
\begin{enumerate}[label=\textbf{\Alph*)}]
\begin{multicols}{2}
  \item Tomás de Iriarte
  \item Leandro Fernández de Moratín
  \item José Cadalso
  \item Gaspar Melchor de Jovellanos
  \item Félix María Samaniego
  \item Juan Meléndez Valdés
\end{multicols}
\end{enumerate}
\vspace{0.15cm}
\item ¿Quién es el autor de Fábulas?
\begin{enumerate}[label=\textbf{\Alph*)}]
\begin{multicols}{2}
  \item Félix María Samaniego
  \item Leandro Fernández de Moratín
  \item José Cadalso
  \item Gaspar Melchor de Jovellanos
  \item Tomás de Iriarte
  \item Juan Meléndez Valdés
\end{multicols}
\end{enumerate}
\vspace{0.15cm}
\item ¿Quién es el autor de Ensayo sobre el hombre?
\begin{enumerate}[label=\textbf{\Alph*)}]
\begin{multicols}{2}
  \item Alexander Pope
  \item Jonathan Swift
  \item Daniel Defoe
  \item Samuel Johnson
  \item Voltaire
  \item Rousseau
\end{multicols}
\end{enumerate}
\vspace{0.15cm}
\item ¿Quién es el autor de Los viajes de Gulliver?
\begin{enumerate}[label=\textbf{\Alph*)}]
\begin{multicols}{2}
  \item Jonathan Swift
  \item Alexander Pope
  \item Daniel Defoe
  \item Samuel Johnson
  \item Voltaire
  \item Rousseau
\end{multicols}
\end{enumerate}
\vspace{0.15cm}
\item ¿Quién es el autor de Robinson Crusoe?
\begin{enumerate}[label=\textbf{\Alph*)}]
\begin{multicols}{2}
  \item Daniel Defoe
  \item Jonathan Swift
  \item Alexander Pope
  \item Samuel Johnson
  \item Voltaire
  \item Rousseau
\end{multicols}
\end{enumerate}
\vspace{0.15cm}
\item ¿Quién es el autor de Cándido?
\begin{enumerate}[label=\textbf{\Alph*)}]
\begin{multicols}{2}
  \item Voltaire
  \item Rousseau
  \item Montesquieu
  \item Diderot
  \item D'Alembert
  \item Buffon
\end{multicols}
\end{enumerate}
\vspace{0.15cm}
\item ¿Quién es el autor de El contrato social?
\begin{enumerate}[label=\textbf{\Alph*)}]
\begin{multicols}{2}
  \item Rousseau
  \item Voltaire
  \item Montesquieu
  \item Diderot
  \item D'Alembert
  \item Buffon
\end{multicols}
\end{enumerate}
\vspace{0.15cm}
\item ¿Quién es el autor de El espíritu de las leyes?
\begin{enumerate}[label=\textbf{\Alph*)}]
\begin{multicols}{2}
  \item Montesquieu
  \item Voltaire
  \item Rousseau
  \item Diderot
  \item D'Alembert
  \item Buffon
\end{multicols}
\end{enumerate}
\vspace{0.15cm}
\item ¿Quién es el autor de La enciclopedia?
\begin{enumerate}[label=\textbf{\Alph*)}]
\begin{multicols}{2}
  \item Diderot y D'Alembert
  \item Voltaire y Rousseau
  \item Montesquieu y Diderot
  \item D'Alembert y Buffon
  \item Voltaire y Montesquieu
  \item Rousseau y Diderot
\end{multicols}
\end{enumerate}
\vspace{0.15cm}
\item ¿Qué es la Enciclopedia?
\begin{enumerate}[label=\textbf{\Alph*)}]
\begin{multicols}{2}
  \item Una obra que recopila todo el saber de la época
  \item Una novela
  \item Una obra teatral
  \item Una lírica
  \item Un ensayo
  \item Una epopeya
\end{multicols}
\end{enumerate}
\vspace{0.15cm}
\item ¿Qué es el despotismo ilustrado?
\begin{enumerate}[label=\textbf{\Alph*)}]
\begin{multicols}{2}
  \item Una forma de gobierno que combina el poder absoluto con reformas ilustradas
  \item Una forma democrática
  \item Una forma republicana
  \item Una forma anarquista
  \item Una forma socialista
  \item Una forma comunista
\end{multicols}
\end{enumerate}
\vspace{0.15cm}
\item ¿Qué es el teatro neoclásico?
\begin{enumerate}[label=\textbf{\Alph*)}]
\begin{multicols}{2}
  \item Un teatro que sigue las reglas clásicas de unidad de acción, lugar y tiempo
  \item Un teatro que rechaza las reglas clásicas
  \item Un teatro que ignora las reglas clásicas
  \item Un teatro medieval
  \item Un teatro barroco
  \item Un teatro romántico
\end{multicols}
\end{enumerate}
\vspace{0.15cm}
\item ¿Cuáles son las tres unidades del teatro neoclásico?
\begin{enumerate}[label=\textbf{\Alph*)}]
\begin{multicols}{2}
  \item Acción, lugar y tiempo
  \item Acción, lugar y espacio
  \item Acción, tiempo y espacio
  \item Lugar, tiempo y espacio
  \item Acción, lugar y personaje
  \item Acción, tiempo y personaje
\end{multicols}
\end{enumerate}
\vspace{0.15cm}
\item ¿Qué es la Ilustración?
\begin{enumerate}[label=\textbf{\Alph*)}]
\begin{multicols}{2}
  \item Un movimiento cultural del siglo XVIII que exalta la razón
  \item Un movimiento del siglo XVII que exalta el sentimiento
  \item Un movimiento que exalta la fe
  \item Un movimiento que exalta la imaginación
  \item Un movimiento de experimentación
  \item Un movimiento de tecnología
\end{multicols}
\end{enumerate}
\vspace{0.15cm}
\item ¿Cuál es el lema de la Ilustración?
\begin{enumerate}[label=\textbf{\Alph*)}]
\begin{multicols}{2}
  \item Sapere aude (Atrévete a saber)
  \item Carpe diem (Aprovecha el día)
  \item Memento mori (Recuerda que morirás)
  \item Cogito ergo sum (Pienso, luego existo)
  \item Veni, vidi, vici (Vine, vi, vencí)
  \item Amor fati (Amor al destino)
\end{multicols}
\end{enumerate}
\vspace{0.15cm}
\item ¿Qué es el Neoclasicismo?
\begin{enumerate}[label=\textbf{\Alph*)}]
\begin{multicols}{2}
  \item Un movimiento literario que imita los modelos clásicos
  \item Un movimiento que rechaza los modelos clásicos
  \item Un movimiento que ignora los modelos clásicos
  \item Un movimiento que imita modelos medievales
  \item Un movimiento que rechaza lo medieval
  \item Un movimiento que ignora lo medieval
\end{multicols}
\end{enumerate}
\vspace{0.15cm}
\item ¿Cuál es el principio fundamental del Neoclasicismo?
\begin{enumerate}[label=\textbf{\Alph*)}]
\begin{multicols}{2}
  \item La razón y el buen gusto
  \item El sentimiento y la imaginación
  \item La fe y la tradición
  \item La experimentación y originalidad
  \item La libertad y la rebeldía
  \item La intuición y el instinto
\end{multicols}
\end{enumerate}
\vspace{0.15cm}
\item ¿Quién es el autor de El sí de las niñas?
\begin{enumerate}[label=\textbf{\Alph*)}]
\begin{multicols}{2}
  \item Leandro Fernández de Moratín
  \item José Cadalso
  \item Gaspar Melchor de Jovellanos
  \item Tomás de Iriarte
  \item Félix María Samaniego
  \item Juan Meléndez Valdés
\end{multicols}
\end{enumerate}
\vspace{0.15cm}
\item ¿Quién es el autor de Cartas marruecas?
\begin{enumerate}[label=\textbf{\Alph*)}]
\begin{multicols}{2}
  \item José Cadalso
  \item Leandro Fernández de Moratín
  \item Gaspar Melchor de Jovellanos
  \item Tomás de Iriarte
  \item Félix María Samaniego
  \item Juan Meléndez Valdés
\end{multicols}
\end{enumerate}
\vspace{0.15cm}
\item ¿Quién es el autor de Fábulas?
\begin{enumerate}[label=\textbf{\Alph*)}]
\begin{multicols}{2}
  \item Tomás de Iriarte
  \item Leandro Fernández de Moratín
  \item José Cadalso
  \item Gaspar Melchor de Jovellanos
  \item Félix María Samaniego
  \item Juan Meléndez Valdés
\end{multicols}
\end{enumerate}
\vspace{0.15cm}
\item ¿Quién es el autor de Fábulas?
\begin{enumerate}[label=\textbf{\Alph*)}]
\begin{multicols}{2}
  \item Félix María Samaniego
  \item Leandro Fernández de Moratín
  \item José Cadalso
  \item Gaspar Melchor de Jovellanos
  \item Tomás de Iriarte
  \item Juan Meléndez Valdés
\end{multicols}
\end{enumerate}
\vspace{0.15cm}
\item ¿Quién es el autor de Ensayo sobre el hombre?
\begin{enumerate}[label=\textbf{\Alph*)}]
\begin{multicols}{2}
  \item Alexander Pope
  \item Jonathan Swift
  \item Daniel Defoe
  \item Samuel Johnson
  \item Voltaire
  \item Rousseau
\end{multicols}
\end{enumerate}
\vspace{0.15cm}
\item ¿Quién es el autor de Los viajes de Gulliver?
\begin{enumerate}[label=\textbf{\Alph*)}]
\begin{multicols}{2}
  \item Jonathan Swift
  \item Alexander Pope
  \item Daniel Defoe
  \item Samuel Johnson
  \item Voltaire
  \item Rousseau
\end{multicols}
\end{enumerate}
\vspace{0.15cm}
\item ¿Quién es el autor de Robinson Crusoe?
\begin{enumerate}[label=\textbf{\Alph*)}]
\begin{multicols}{2}
  \item Daniel Defoe
  \item Jonathan Swift
  \item Alexander Pope
  \item Samuel Johnson
  \item Voltaire
  \item Rousseau
\end{multicols}
\end{enumerate}
\vspace{0.15cm}
\item ¿Quién es el autor de Cándido?
\begin{enumerate}[label=\textbf{\Alph*)}]
\begin{multicols}{2}
  \item Voltaire
  \item Rousseau
  \item Montesquieu
  \item Diderot
  \item D'Alembert
  \item Buffon
\end{multicols}
\end{enumerate}
\vspace{0.15cm}
\item ¿Quién es el autor de El contrato social?
\begin{enumerate}[label=\textbf{\Alph*)}]
\begin{multicols}{2}
  \item Rousseau
  \item Voltaire
  \item Montesquieu
  \item Diderot
  \item D'Alembert
  \item Buffon
\end{multicols}
\end{enumerate}
\vspace{0.15cm}
\item ¿Quién es el autor de El espíritu de las leyes?
\begin{enumerate}[label=\textbf{\Alph*)}]
\begin{multicols}{2}
  \item Montesquieu
  \item Voltaire
  \item Rousseau
  \item Diderot
  \item D'Alembert
  \item Buffon
\end{multicols}
\end{enumerate}
\vspace{0.15cm}
\item ¿Quién es el autor de La enciclopedia?
\begin{enumerate}[label=\textbf{\Alph*)}]
\begin{multicols}{2}
  \item Diderot y D'Alembert
  \item Voltaire y Rousseau
  \item Montesquieu y Diderot
  \item D'Alembert y Buffon
  \item Voltaire y Montesquieu
  \item Rousseau y Diderot
\end{multicols}
\end{enumerate}
\vspace{0.15cm}
\item ¿Qué es la Enciclopedia?
\begin{enumerate}[label=\textbf{\Alph*)}]
\begin{multicols}{2}
  \item Una obra que recopila todo el saber de la época
  \item Una novela
  \item Una obra teatral
  \item Una lírica
  \item Un ensayo
  \item Una epopeya
\end{multicols}
\end{enumerate}
\vspace{0.15cm}
\item ¿Qué es el despotismo ilustrado?
\begin{enumerate}[label=\textbf{\Alph*)}]
\begin{multicols}{2}
  \item Una forma de gobierno que combina el poder absoluto con reformas ilustradas
  \item Una forma democrática
  \item Una forma republicana
  \item Una forma anarquista
  \item Una forma socialista
  \item Una forma comunista
\end{multicols}
\end{enumerate}
\vspace{0.15cm}
\item ¿Qué es el teatro neoclásico?
\begin{enumerate}[label=\textbf{\Alph*)}]
\begin{multicols}{2}
  \item Un teatro que sigue las reglas clásicas de unidad de acción, lugar y tiempo
  \item Un teatro que rechaza las reglas clásicas
  \item Un teatro que ignora las reglas clásicas
  \item Un teatro medieval
  \item Un teatro barroco
  \item Un teatro romántico
\end{multicols}
\end{enumerate}
\vspace{0.15cm}
\item ¿Cuáles son las tres unidades del teatro neoclásico?
\begin{enumerate}[label=\textbf{\Alph*)}]
\begin{multicols}{2}
  \item Acción, lugar y tiempo
  \item Acción, lugar y espacio
  \item Acción, tiempo y espacio
  \item Lugar, tiempo y espacio
  \item Acción, lugar y personaje
  \item Acción, tiempo y personaje
\end{multicols}
\end{enumerate}
\vspace{0.15cm}
\item ¿Qué es la Ilustración?
\begin{enumerate}[label=\textbf{\Alph*)}]
\begin{multicols}{2}
  \item Un movimiento cultural del siglo XVIII que exalta la razón
  \item Un movimiento del siglo XVII que exalta el sentimiento
  \item Un movimiento que exalta la fe
  \item Un movimiento que exalta la imaginación
  \item Un movimiento de experimentación
  \item Un movimiento de tecnología
\end{multicols}
\end{enumerate}
\vspace{0.15cm}
\item ¿Cuál es el lema de la Ilustración?
\begin{enumerate}[label=\textbf{\Alph*)}]
\begin{multicols}{2}
  \item Sapere aude (Atrévete a saber)
  \item Carpe diem (Aprovecha el día)
  \item Memento mori (Recuerda que morirás)
  \item Cogito ergo sum (Pienso, luego existo)
  \item Veni, vidi, vici (Vine, vi, vencí)
  \item Amor fati (Amor al destino)
\end{multicols}
\end{enumerate}
\vspace{0.15cm}
\item ¿Qué es el Neoclasicismo?
\begin{enumerate}[label=\textbf{\Alph*)}]
\begin{multicols}{2}
  \item Un movimiento literario que imita los modelos clásicos
  \item Un movimiento que rechaza los modelos clásicos
  \item Un movimiento que ignora los modelos clásicos
  \item Un movimiento que imita modelos medievales
  \item Un movimiento que rechaza lo medieval
  \item Un movimiento que ignora lo medieval
\end{multicols}
\end{enumerate}
\vspace{0.15cm}
\end{enumerate}\n
\vspace{0.4cm}
\begin{tcolorbox}[colback=gray!5!white,colframe=gray!60!black,title=\textbf{SOLUCIONARIO -- 3.1 Literatura Universal y Europea}]
\begin{multicols}{4}
\begin{enumerate}[label=\textbf{\arabic*.}]
  \item \textbf{C}
  \item \textbf{B}
  \item \textbf{B}
  \item \textbf{C}
  \item \textbf{B}
  \item \textbf{C}
  \item \textbf{B}
  \item \textbf{C}
  \item \textbf{B}
  \item \textbf{B}
  \item \textbf{B}
  \item \textbf{C}
  \item \textbf{A}
  \item \textbf{D}
  \item \textbf{B}
  \item \textbf{A}
  \item \textbf{A}
  \item \textbf{A}
  \item \textbf{B}
  \item \textbf{A}
  \item \textbf{A}
  \item \textbf{A}
  \item \textbf{A}
  \item \textbf{A}
  \item \textbf{A}
  \item \textbf{A}
  \item \textbf{A}
  \item \textbf{A}
  \item \textbf{A}
  \item \textbf{A}
  \item \textbf{A}
  \item \textbf{A}
  \item \textbf{A}
  \item \textbf{A}
  \item \textbf{A}
  \item \textbf{B}
  \item \textbf{A}
  \item \textbf{A}
  \item \textbf{A}
  \item \textbf{A}
  \item \textbf{A}
  \item \textbf{B}
  \item \textbf{A}
  \item \textbf{A}
  \item \textbf{A}
  \item \textbf{A}
  \item \textbf{A}
  \item \textbf{A}
  \item \textbf{A}
  \item \textbf{A}
  \item \textbf{A}
  \item \textbf{A}
  \item \textbf{A}
  \item \textbf{A}
  \item \textbf{A}
  \item \textbf{A}
  \item \textbf{A}
  \item \textbf{A}
  \item \textbf{A}
  \item \textbf{A}
  \item \textbf{A}
  \item \textbf{A}
  \item \textbf{A}
  \item \textbf{A}
  \item \textbf{A}
  \item \textbf{A}
  \item \textbf{A}
  \item \textbf{A}
  \item \textbf{A}
  \item \textbf{A}
  \item \textbf{A}
  \item \textbf{A}
  \item \textbf{A}
  \item \textbf{A}
  \item \textbf{A}
  \item \textbf{A}
  \item \textbf{A}
  \item \textbf{A}
  \item \textbf{A}
  \item \textbf{A}
  \item \textbf{A}
  \item \textbf{A}
  \item \textbf{A}
  \item \textbf{A}
  \item \textbf{A}
  \item \textbf{A}
  \item \textbf{A}
  \item \textbf{A}
  \item \textbf{A}
  \item \textbf{A}
  \item \textbf{A}
  \item \textbf{A}
  \item \textbf{A}
  \item \textbf{A}
  \item \textbf{A}
  \item \textbf{A}
  \item \textbf{A}
  \item \textbf{A}
  \item \textbf{A}
  \item \textbf{A}
  \item \textbf{A}
  \item \textbf{A}
  \item \textbf{A}
  \item \textbf{A}
  \item \textbf{A}
  \item \textbf{A}
  \item \textbf{A}
  \item \textbf{A}
  \item \textbf{A}
  \item \textbf{A}
  \item \textbf{A}
  \item \textbf{A}
  \item \textbf{A}
  \item \textbf{A}
  \item \textbf{A}
  \item \textbf{A}
  \item \textbf{A}
  \item \textbf{A}
  \item \textbf{A}
  \item \textbf{A}
  \item \textbf{A}
  \item \textbf{A}
  \item \textbf{A}
  \item \textbf{A}
  \item \textbf{A}
  \item \textbf{A}
  \item \textbf{A}
  \item \textbf{A}
  \item \textbf{A}
  \item \textbf{A}
  \item \textbf{A}
  \item \textbf{A}
  \item \textbf{A}
  \item \textbf{A}
  \item \textbf{A}
  \item \textbf{A}
  \item \textbf{A}
  \item \textbf{A}
  \item \textbf{A}
  \item \textbf{A}
  \item \textbf{A}
  \item \textbf{A}
  \item \textbf{A}
  \item \textbf{A}
  \item \textbf{A}
  \item \textbf{A}
  \item \textbf{A}
  \item \textbf{A}
  \item \textbf{A}
  \item \textbf{A}
  \item \textbf{A}
  \item \textbf{A}
  \item \textbf{A}
  \item \textbf{A}
  \item \textbf{A}
  \item \textbf{A}
  \item \textbf{A}
  \item \textbf{A}
  \item \textbf{A}
  \item \textbf{A}
  \item \textbf{A}
  \item \textbf{A}
  \item \textbf{A}
  \item \textbf{A}
  \item \textbf{A}
  \item \textbf{A}
  \item \textbf{A}
  \item \textbf{A}
  \item \textbf{A}
  \item \textbf{A}
  \item \textbf{A}
  \item \textbf{A}
  \item \textbf{A}
  \item \textbf{A}
  \item \textbf{A}
  \item \textbf{A}
  \item \textbf{A}
  \item \textbf{A}
  \item \textbf{A}
  \item \textbf{A}
  \item \textbf{A}
  \item \textbf{A}
  \item \textbf{A}
  \item \textbf{A}
  \item \textbf{A}
  \item \textbf{A}
  \item \textbf{A}
  \item \textbf{A}
  \item \textbf{A}
  \item \textbf{A}
  \item \textbf{A}
  \item \textbf{A}
  \item \textbf{A}
  \item \textbf{A}
  \item \textbf{A}
  \item \textbf{A}
  \item \textbf{A}
  \item \textbf{A}
  \item \textbf{A}
  \item \textbf{A}
  \item \textbf{A}
  \item \textbf{A}
  \item \textbf{A}
  \item \textbf{A}
  \item \textbf{A}
  \item \textbf{A}
  \item \textbf{A}
  \item \textbf{A}
  \item \textbf{A}
  \item \textbf{A}
  \item \textbf{A}
  \item \textbf{A}
  \item \textbf{A}
  \item \textbf{A}
  \item \textbf{A}
  \item \textbf{A}
  \item \textbf{A}
  \item \textbf{A}
  \item \textbf{A}
  \item \textbf{A}
  \item \textbf{A}
  \item \textbf{A}
  \item \textbf{A}
  \item \textbf{A}
  \item \textbf{A}
  \item \textbf{A}
  \item \textbf{A}
  \item \textbf{A}
  \item \textbf{A}
  \item \textbf{A}
  \item \textbf{A}
  \item \textbf{A}
  \item \textbf{A}
  \item \textbf{A}
  \item \textbf{A}
  \item \textbf{A}
  \item \textbf{A}
  \item \textbf{A}
  \item \textbf{A}
  \item \textbf{A}
  \item \textbf{A}
  \item \textbf{A}
  \item \textbf{A}
  \item \textbf{A}
  \item \textbf{A}
  \item \textbf{A}
  \item \textbf{A}
  \item \textbf{A}
  \item \textbf{A}
  \item \textbf{A}
  \item \textbf{A}
  \item \textbf{A}
  \item \textbf{A}
  \item \textbf{A}
  \item \textbf{A}
  \item \textbf{A}
  \item \textbf{A}
  \item \textbf{A}
  \item \textbf{A}
  \item \textbf{A}
  \item \textbf{A}
  \item \textbf{A}
  \item \textbf{A}
  \item \textbf{A}
  \item \textbf{A}
  \item \textbf{A}
  \item \textbf{A}
  \item \textbf{A}
  \item \textbf{A}
  \item \textbf{A}
  \item \textbf{A}
  \item \textbf{A}
  \item \textbf{A}
  \item \textbf{A}
  \item \textbf{A}
  \item \textbf{A}
  \item \textbf{A}
  \item \textbf{A}
  \item \textbf{A}
  \item \textbf{A}
  \item \textbf{A}
  \item \textbf{A}
  \item \textbf{A}
  \item \textbf{A}
  \item \textbf{A}
  \item \textbf{A}
  \item \textbf{A}
  \item \textbf{A}
  \item \textbf{A}
  \item \textbf{A}
  \item \textbf{A}
  \item \textbf{A}
  \item \textbf{A}
  \item \textbf{A}
  \item \textbf{A}
  \item \textbf{A}
  \item \textbf{A}
  \item \textbf{A}
  \item \textbf{A}
  \item \textbf{A}
\end{enumerate}
\end{multicols}
\end{tcolorbox}
\newpage\n
\subsection{3.2 Movimientos del Siglo XIX}

\subsubsection*{SECCIÓN I: ROMANTICISMO – CONCEPTOS BÁSICOS (1–40)}

\begin{enumerate}[label=\textbf{\arabic*.}]
\setcounter{enumi}{0}
\item ¿Qué es el Romanticismo?
\begin{enumerate}[label=\textbf{\Alph*)}]
\begin{multicols}{2}
  \item Un movimiento cultural y literario del siglo XIX que exalta el sentimiento y la libertad
  \item Siglo XVIII
  \item Siglo XVII
  \item Siglo XX
  \item Siglo XVI
  \item Siglo XXI
\end{multicols}
\end{enumerate}
\vspace{0.15cm}
\item ¿En qué siglo surge el Romanticismo?
\begin{enumerate}[label=\textbf{\Alph*)}]
\begin{multicols}{2}
  \item Siglo XIX
  \item Siglo XVIII
  \item Siglo XVII
  \item Siglo XX
  \item Siglo XVI
  \item Siglo XXI
\end{multicols}
\end{enumerate}
\vspace{0.15cm}
\item ¿Cuál es el principal valor del Romanticismo?
\begin{enumerate}[label=\textbf{\Alph*)}]
\begin{multicols}{2}
  \item El sentimiento
  \item Siglo XVIII
  \item Siglo XVII
  \item Siglo XX
  \item Siglo XVI
  \item Siglo XXI
\end{multicols}
\end{enumerate}
\vspace{0.15cm}
\item ¿Qué es el individualismo romántico?
\begin{enumerate}[label=\textbf{\Alph*)}]
\begin{multicols}{2}
  \item La exaltación del yo y la subjetividad
  \item Siglo XVIII
  \item Siglo XVII
  \item Siglo XX
  \item Siglo XVI
  \item Siglo XXI
\end{multicols}
\end{enumerate}
\vspace{0.15cm}
\item ¿Qué es la libertad creadora en el Romanticismo?
\begin{enumerate}[label=\textbf{\Alph*)}]
\begin{multicols}{2}
  \item La ausencia de reglas y la búsqueda de la originalidad
  \item Siglo XVIII
  \item Siglo XVII
  \item Siglo XX
  \item Siglo XVI
  \item Siglo XXI
\end{multicols}
\end{enumerate}
\vspace{0.15cm}
\item ¿Qué es el subjetivismo romántico?
\begin{enumerate}[label=\textbf{\Alph*)}]
\begin{multicols}{2}
  \item La expresión de los sentimientos y emociones personales
  \item Siglo XVIII
  \item Siglo XVII
  \item Siglo XX
  \item Siglo XVI
  \item Siglo XXI
\end{multicols}
\end{enumerate}
\vspace{0.15cm}
\item ¿Qué es la evasión romántica?
\begin{enumerate}[label=\textbf{\Alph*)}]
\begin{multicols}{2}
  \item La huida de la realidad hacia mundos imaginarios o lejanos
  \item Siglo XVIII
  \item Siglo XVII
  \item Siglo XX
  \item Siglo XVI
  \item Siglo XXI
\end{multicols}
\end{enumerate}
\vspace{0.15cm}
\item ¿Qué es el amor romántico?
\begin{enumerate}[label=\textbf{\Alph*)}]
\begin{multicols}{2}
  \item Un amor idealizado, apasionado y a menudo trágico
  \item Siglo XVIII
  \item Siglo XVII
  \item Siglo XX
  \item Siglo XVI
  \item Siglo XXI
\end{multicols}
\end{enumerate}
\vspace{0.15cm}
\item ¿Qué es la naturaleza en el Romanticismo?
\begin{enumerate}[label=\textbf{\Alph*)}]
\begin{multicols}{2}
  \item Un reflejo de los sentimientos del poeta
  \item Siglo XVIII
  \item Siglo XVII
  \item Siglo XX
  \item Siglo XVI
  \item Siglo XXI
\end{multicols}
\end{enumerate}
\vspace{0.15cm}
\item ¿Qué es el paisaje romántico?
\begin{enumerate}[label=\textbf{\Alph*)}]
\begin{multicols}{2}
  \item Un paisaje melancólico, salvaje y misterioso
  \item Siglo XVIII
  \item Siglo XVII
  \item Siglo XX
  \item Siglo XVI
  \item Siglo XXI
\end{multicols}
\end{enumerate}
\vspace{0.15cm}
\item ¿Qué es el gusto por lo nocturnal en el Romanticismo?
\begin{enumerate}[label=\textbf{\Alph*)}]
\begin{multicols}{2}
  \item La preferencia por la noche, la luna y lo misterioso
  \item Siglo XVIII
  \item Siglo XVII
  \item Siglo XX
  \item Siglo XVI
  \item Siglo XXI
\end{multicols}
\end{enumerate}
\vspace{0.15cm}
\item ¿Qué es el gusto por lo macabro en el Romanticismo?
\begin{enumerate}[label=\textbf{\Alph*)}]
\begin{multicols}{2}
  \item La atracción por la muerte, las ruinas y lo sobrenatural
  \item Siglo XVIII
  \item Siglo XVII
  \item Siglo XX
  \item Siglo XVI
  \item Siglo XXI
\end{multicols}
\end{enumerate}
\vspace{0.15cm}
\item ¿Qué es el historicismo romántico?
\begin{enumerate}[label=\textbf{\Alph*)}]
\begin{multicols}{2}
  \item El interés por el pasado, especialmente la Edad Media
  \item Siglo XVIII
  \item Siglo XVII
  \item Siglo XX
  \item Siglo XVI
  \item Siglo XXI
\end{multicols}
\end{enumerate}
\vspace{0.15cm}
\item ¿Qué es el nacionalismo romántico?
\begin{enumerate}[label=\textbf{\Alph*)}]
\begin{multicols}{2}
  \item La exaltación de la identidad y las tradiciones nacionales
  \item Siglo XVIII
  \item Siglo XVII
  \item Siglo XX
  \item Siglo XVI
  \item Siglo XXI
\end{multicols}
\end{enumerate}
\vspace{0.15cm}
\item ¿Qué es el liberalismo romántico?
\begin{enumerate}[label=\textbf{\Alph*)}]
\begin{multicols}{2}
  \item La defensa de las libertades individuales y políticas
  \item Siglo XVIII
  \item Siglo XVII
  \item Siglo XX
  \item Siglo XVI
  \item Siglo XXI
\end{multicols}
\end{enumerate}
\vspace{0.15cm}
\item ¿Qué es la rebeldía romántica?
\begin{enumerate}[label=\textbf{\Alph*)}]
\begin{multicols}{2}
  \item La oposición a las normas sociales y artísticas establecidas
  \item Siglo XVIII
  \item Siglo XVII
  \item Siglo XX
  \item Siglo XVI
  \item Siglo XXI
\end{multicols}
\end{enumerate}
\vspace{0.15cm}
\item ¿Qué es el héroe romántico?
\begin{enumerate}[label=\textbf{\Alph*)}]
\begin{multicols}{2}
  \item Un personaje rebelde, solitario y atormentado
  \item Siglo XVIII
  \item Siglo XVII
  \item Siglo XX
  \item Siglo XVI
  \item Siglo XXI
\end{multicols}
\end{enumerate}
\vspace{0.15cm}
\item ¿Qué es el "mal del siglo" (mal du siècle)?
\begin{enumerate}[label=\textbf{\Alph*)}]
\begin{multicols}{2}
  \item Un sentimiento de insatisfacción y melancolía vital
  \item Siglo XVIII
  \item Siglo XVII
  \item Siglo XX
  \item Siglo XVI
  \item Siglo XXI
\end{multicols}
\end{enumerate}
\vspace{0.15cm}
\item ¿Qué es el "yo" romántico?
\begin{enumerate}[label=\textbf{\Alph*)}]
\begin{multicols}{2}
  \item El centro de la creación literaria y artística
  \item Siglo XVIII
  \item Siglo XVII
  \item Siglo XX
  \item Siglo XVI
  \item Siglo XXI
\end{multicols}
\end{enumerate}
\vspace{0.15cm}
\item ¿Qué es la imaginación en el Romanticismo?
\begin{enumerate}[label=\textbf{\Alph*)}]
\begin{multicols}{2}
  \item La facultad creadora más valorada
  \item Siglo XVIII
  \item Siglo XVII
  \item Siglo XX
  \item Siglo XVI
  \item Siglo XXI
\end{multicols}
\end{enumerate}
\vspace{0.15cm}
\item ¿Qué es la originalidad en el Romanticismo?
\begin{enumerate}[label=\textbf{\Alph*)}]
\begin{multicols}{2}
  \item La búsqueda de lo nuevo y lo propio
  \item Siglo XVIII
  \item Siglo XVII
  \item Siglo XX
  \item Siglo XVI
  \item Siglo XXI
\end{multicols}
\end{enumerate}
\vspace{0.15cm}
\item ¿Qué es la genialidad romántica?
\begin{enumerate}[label=\textbf{\Alph*)}]
\begin{multicols}{2}
  \item La capacidad creadora excepcional del artista
  \item Siglo XVIII
  \item Siglo XVII
  \item Siglo XX
  \item Siglo XVI
  \item Siglo XXI
\end{multicols}
\end{enumerate}
\vspace{0.15cm}
\item ¿Qué es el idealismo romántico?
\begin{enumerate}[label=\textbf{\Alph*)}]
\begin{multicols}{2}
  \item La búsqueda de ideales absolutos e inalcanzables
  \item Siglo XVIII
  \item Siglo XVII
  \item Siglo XX
  \item Siglo XVI
  \item Siglo XXI
\end{multicols}
\end{enumerate}
\vspace{0.15cm}
\item ¿Qué es el pesimismo romántico?
\begin{enumerate}[label=\textbf{\Alph*)}]
\begin{multicols}{2}
  \item La visión negativa y melancólica de la existencia
  \item Siglo XVIII
  \item Siglo XVII
  \item Siglo XX
  \item Siglo XVI
  \item Siglo XXI
\end{multicols}
\end{enumerate}
\vspace{0.15cm}
\item ¿Qué es el desengaño romántico?
\begin{enumerate}[label=\textbf{\Alph*)}]
\begin{multicols}{2}
  \item La decepción ante la realidad que no cumple los ideales
  \item Siglo XVIII
  \item Siglo XVII
  \item Siglo XX
  \item Siglo XVI
  \item Siglo XXI
\end{multicols}
\end{enumerate}
\vspace{0.15cm}
\item ¿Qué es la melancolía romántica?
\begin{enumerate}[label=\textbf{\Alph*)}]
\begin{multicols}{2}
  \item Un estado de tristeza y reflexión profunda
  \item Siglo XVIII
  \item Siglo XVII
  \item Siglo XX
  \item Siglo XVI
  \item Siglo XXI
\end{multicols}
\end{enumerate}
\vspace{0.15cm}
\item ¿Qué es la soledad romántica?
\begin{enumerate}[label=\textbf{\Alph*)}]
\begin{multicols}{2}
  \item La búsqueda del aislamiento y la introspección
  \item Siglo XVIII
  \item Siglo XVII
  \item Siglo XX
  \item Siglo XVI
  \item Siglo XXI
\end{multicols}
\end{enumerate}
\vspace{0.15cm}
\item ¿Qué es la muerte en el Romanticismo?
\begin{enumerate}[label=\textbf{\Alph*)}]
\begin{multicols}{2}
  \item Un tema recurrente y fascinante
  \item Siglo XVIII
  \item Siglo XVII
  \item Siglo XX
  \item Siglo XVI
  \item Siglo XXI
\end{multicols}
\end{enumerate}
\vspace{0.15cm}
\item ¿Qué es el suicidio romántico?
\begin{enumerate}[label=\textbf{\Alph*)}]
\begin{multicols}{2}
  \item Un acto trágico y literario frecuente en la época
  \item Siglo XVIII
  \item Siglo XVII
  \item Siglo XX
  \item Siglo XVI
  \item Siglo XXI
\end{multicols}
\end{enumerate}
\vspace{0.15cm}
\item ¿Qué es el sueño en el Romanticismo?
\begin{enumerate}[label=\textbf{\Alph*)}]
\begin{multicols}{2}
  \item Una vía de acceso a lo sobrenatural y lo irracional
  \item Siglo XVIII
  \item Siglo XVII
  \item Siglo XX
  \item Siglo XVI
  \item Siglo XXI
\end{multicols}
\end{enumerate}
\vspace{0.15cm}
\item ¿Qué es lo sobrenatural en el Romanticismo?
\begin{enumerate}[label=\textbf{\Alph*)}]
\begin{multicols}{2}
  \item Un tema frecuente que expresa lo misterioso e inexplicable
  \item Siglo XVIII
  \item Siglo XVII
  \item Siglo XX
  \item Siglo XVI
  \item Siglo XXI
\end{multicols}
\end{enumerate}
\vspace{0.15cm}
\item ¿Qué es el Gothic revival?
\begin{enumerate}[label=\textbf{\Alph*)}]
\begin{multicols}{2}
  \item El resurgimiento del estilo gótico en la arquitectura y la literatura
  \item Siglo XVIII
  \item Siglo XVII
  \item Siglo XX
  \item Siglo XVI
  \item Siglo XXI
\end{multicols}
\end{enumerate}
\vspace{0.15cm}
\item ¿Qué es la novela gótica?
\begin{enumerate}[label=\textbf{\Alph*)}]
\begin{multicols}{2}
  \item Un género narrativo de terror y misterio
  \item Siglo XVIII
  \item Siglo XVII
  \item Siglo XX
  \item Siglo XVI
  \item Siglo XXI
\end{multicols}
\end{enumerate}
\vspace{0.15cm}
\item ¿Quién es el autor de Frankenstein?
\begin{enumerate}[label=\textbf{\Alph*)}]
\begin{multicols}{2}
  \item Mary Shelley
  \item Siglo XVIII
  \item Siglo XVII
  \item Siglo XX
  \item Siglo XVI
  \item Siglo XXI
\end{multicols}
\end{enumerate}
\vspace{0.15cm}
\item ¿Quién es el autor de Drácula?
\begin{enumerate}[label=\textbf{\Alph*)}]
\begin{multicols}{2}
  \item Bram Stoker
  \item Siglo XVIII
  \item Siglo XVII
  \item Siglo XX
  \item Siglo XVI
  \item Siglo XXI
\end{multicols}
\end{enumerate}
\vspace{0.15cm}
\item ¿Qué es el drama romántico?
\begin{enumerate}[label=\textbf{\Alph*)}]
\begin{multicols}{2}
  \item Un género teatral que mezcla lo trágico y lo cómico
  \item Siglo XVIII
  \item Siglo XVII
  \item Siglo XX
  \item Siglo XVI
  \item Siglo XXI
\end{multicols}
\end{enumerate}
\vspace{0.15cm}
\item ¿Qué es la poesía romántica?
\begin{enumerate}[label=\textbf{\Alph*)}]
\begin{multicols}{2}
  \item La expresión lírica de los sentimientos del yo poético
  \item Siglo XVIII
  \item Siglo XVII
  \item Siglo XX
  \item Siglo XVI
  \item Siglo XXI
\end{multicols}
\end{enumerate}
\vspace{0.15cm}
\item ¿Qué es el teatro romántico?
\begin{enumerate}[label=\textbf{\Alph*)}]
\begin{multicols}{2}
  \item Un teatro que rompe las unidades clásicas y mezcla géneros
  \item Siglo XVIII
  \item Siglo XVII
  \item Siglo XX
  \item Siglo XVI
  \item Siglo XXI
\end{multicols}
\end{enumerate}
\vspace{0.15cm}
\item ¿Qué es el artículo de costumbres?
\begin{enumerate}[label=\textbf{\Alph*)}]
\begin{multicols}{2}
  \item Un género periodístico que describe las costumbres de la época
  \item Siglo XVIII
  \item Siglo XVII
  \item Siglo XX
  \item Siglo XVI
  \item Siglo XXI
\end{multicols}
\end{enumerate}
\vspace{0.15cm}
\item ¿Qué es el folletín?
\begin{enumerate}[label=\textbf{\Alph*)}]
\begin{multicols}{2}
  \item Una novela publicada por entregas en los periódicos
  \item Siglo XVIII
  \item Siglo XVII
  \item Siglo XX
  \item Siglo XVI
  \item Siglo XXI
\end{multicols}
\end{enumerate}
\vspace{0.15cm}
\end{enumerate}\n
\subsubsection*{SECCIÓN II: ROMANTICISMO – AUTORES Y OBRAS (41–90)}

\begin{enumerate}[label=\textbf{\arabic*.}]
\setcounter{enumi}{40}
\item ¿Quién es el autor de las Rimas y Leyendas?
\begin{enumerate}[label=\textbf{\Alph*)}]
\begin{multicols}{2}
  \item Gustavo Adolfo Bécquer
  \item Siglo XVIII
  \item Siglo XVII
  \item Siglo XX
  \item Siglo XVI
  \item Siglo XXI
\end{multicols}
\end{enumerate}
\vspace{0.15cm}
\item ¿Quién es el autor de Don Juan Tenorio?
\begin{enumerate}[label=\textbf{\Alph*)}]
\begin{multicols}{2}
  \item José Zorrilla
  \item Siglo XVIII
  \item Siglo XVII
  \item Siglo XX
  \item Siglo XVI
  \item Siglo XXI
\end{multicols}
\end{enumerate}
\vspace{0.15cm}
\item ¿Quién es el autor de El estudiante de Salamanca?
\begin{enumerate}[label=\textbf{\Alph*)}]
\begin{multicols}{2}
  \item José de Espronceda
  \item Siglo XVIII
  \item Siglo XVII
  \item Siglo XX
  \item Siglo XVI
  \item Siglo XXI
\end{multicols}
\end{enumerate}
\vspace{0.15cm}
\item ¿Quién es el autor de La canción del pirata?
\begin{enumerate}[label=\textbf{\Alph*)}]
\begin{multicols}{2}
  \item José de Espronceda
  \item Siglo XVIII
  \item Siglo XVII
  \item Siglo XX
  \item Siglo XVI
  \item Siglo XXI
\end{multicols}
\end{enumerate}
\vspace{0.15cm}
\item ¿Quién es el autor de Artículos de costumbres?
\begin{enumerate}[label=\textbf{\Alph*)}]
\begin{multicols}{2}
  \item Mariano José de Larra
  \item Siglo XVIII
  \item Siglo XVII
  \item Siglo XX
  \item Siglo XVI
  \item Siglo XXI
\end{multicols}
\end{enumerate}
\vspace{0.15cm}
\item ¿Quién es el autor de Don Álvaro o la fuerza del sino?
\begin{enumerate}[label=\textbf{\Alph*)}]
\begin{multicols}{2}
  \item Duque de Rivas
  \item Siglo XVIII
  \item Siglo XVII
  \item Siglo XX
  \item Siglo XVI
  \item Siglo XXI
\end{multicols}
\end{enumerate}
\vspace{0.15cm}
\item ¿Quién es la autora de Cantares gallegos?
\begin{enumerate}[label=\textbf{\Alph*)}]
\begin{multicols}{2}
  \item Rosalía de Castro
  \item Siglo XVIII
  \item Siglo XVII
  \item Siglo XX
  \item Siglo XVI
  \item Siglo XXI
\end{multicols}
\end{enumerate}
\vspace{0.15cm}
\item ¿Quién es el autor de Los amantes de Teruel?
\begin{enumerate}[label=\textbf{\Alph*)}]
\begin{multicols}{2}
  \item Juan Eugenio Hartzenbusch
  \item Siglo XVIII
  \item Siglo XVII
  \item Siglo XX
  \item Siglo XVI
  \item Siglo XXI
\end{multicols}
\end{enumerate}
\vspace{0.15cm}
\item ¿Quién es el autor de El señor de Bembibre?
\begin{enumerate}[label=\textbf{\Alph*)}]
\begin{multicols}{2}
  \item Enrique Gil y Carrasco
  \item Siglo XVIII
  \item Siglo XVII
  \item Siglo XX
  \item Siglo XVI
  \item Siglo XXI
\end{multicols}
\end{enumerate}
\vspace{0.15cm}
\item ¿Quién es el autor de Los miserables?
\begin{enumerate}[label=\textbf{\Alph*)}]
\begin{multicols}{2}
  \item Victor Hugo
  \item Siglo XVIII
  \item Siglo XVII
  \item Siglo XX
  \item Siglo XVI
  \item Siglo XXI
\end{multicols}
\end{enumerate}
\vspace{0.15cm}
\item ¿Quién es el autor de Nuestra Señora de París?
\begin{enumerate}[label=\textbf{\Alph*)}]
\begin{multicols}{2}
  \item Victor Hugo
  \item Siglo XVIII
  \item Siglo XVII
  \item Siglo XX
  \item Siglo XVI
  \item Siglo XXI
\end{multicols}
\end{enumerate}
\vspace{0.15cm}
\item ¿Quién es el autor de Los tres mosqueteros?
\begin{enumerate}[label=\textbf{\Alph*)}]
\begin{multicols}{2}
  \item Alejandro Dumas
  \item Siglo XVIII
  \item Siglo XVII
  \item Siglo XX
  \item Siglo XVI
  \item Siglo XXI
\end{multicols}
\end{enumerate}
\vspace{0.15cm}
\item ¿Quién es el autor de El conde de Montecristo?
\begin{enumerate}[label=\textbf{\Alph*)}]
\begin{multicols}{2}
  \item Alejandro Dumas
  \item Siglo XVIII
  \item Siglo XVII
  \item Siglo XX
  \item Siglo XVI
  \item Siglo XXI
\end{multicols}
\end{enumerate}
\vspace{0.15cm}
\item ¿Quién es el autor de La comedia humana?
\begin{enumerate}[label=\textbf{\Alph*)}]
\begin{multicols}{2}
  \item Honoré de Balzac
  \item Siglo XVIII
  \item Siglo XVII
  \item Siglo XX
  \item Siglo XVI
  \item Siglo XXI
\end{multicols}
\end{enumerate}
\vspace{0.15cm}
\item ¿Quién es el autor de Rojo y negro?
\begin{enumerate}[label=\textbf{\Alph*)}]
\begin{multicols}{2}
  \item Stendhal
  \item Siglo XVIII
  \item Siglo XVII
  \item Siglo XX
  \item Siglo XVI
  \item Siglo XXI
\end{multicols}
\end{enumerate}
\vspace{0.15cm}
\item ¿Quién es el autor de La cartuja de Parma?
\begin{enumerate}[label=\textbf{\Alph*)}]
\begin{multicols}{2}
  \item Stendhal
  \item Siglo XVIII
  \item Siglo XVII
  \item Siglo XX
  \item Siglo XVI
  \item Siglo XXI
\end{multicols}
\end{enumerate}
\vspace{0.15cm}
\item ¿Quién es el autor de Madame Bovary?
\begin{enumerate}[label=\textbf{\Alph*)}]
\begin{multicols}{2}
  \item Gustave Flaubert
  \item Siglo XVIII
  \item Siglo XVII
  \item Siglo XX
  \item Siglo XVI
  \item Siglo XXI
\end{multicols}
\end{enumerate}
\vspace{0.15cm}
\item ¿Quién es el autor de La educación sentimental?
\begin{enumerate}[label=\textbf{\Alph*)}]
\begin{multicols}{2}
  \item Gustave Flaubert
  \item Siglo XVIII
  \item Siglo XVII
  \item Siglo XX
  \item Siglo XVI
  \item Siglo XXI
\end{multicols}
\end{enumerate}
\vspace{0.15cm}
\item ¿Quién es el autor de Germinal?
\begin{enumerate}[label=\textbf{\Alph*)}]
\begin{multicols}{2}
  \item Émile Zola
  \item Siglo XVIII
  \item Siglo XVII
  \item Siglo XX
  \item Siglo XVI
  \item Siglo XXI
\end{multicols}
\end{enumerate}
\vspace{0.15cm}
\item ¿Quién es el autor de Naná?
\begin{enumerate}[label=\textbf{\Alph*)}]
\begin{multicols}{2}
  \item Émile Zola
  \item Siglo XVIII
  \item Siglo XVII
  \item Siglo XX
  \item Siglo XVI
  \item Siglo XXI
\end{multicols}
\end{enumerate}
\vspace{0.15cm}
\item ¿Quién es el autor de Orgullo y prejuicio?
\begin{enumerate}[label=\textbf{\Alph*)}]
\begin{multicols}{2}
  \item Jane Austen
  \item Siglo XVIII
  \item Siglo XVII
  \item Siglo XX
  \item Siglo XVI
  \item Siglo XXI
\end{multicols}
\end{enumerate}
\vspace{0.15cm}
\item ¿Quién es el autor de Cumbres borrascosas?
\begin{enumerate}[label=\textbf{\Alph*)}]
\begin{multicols}{2}
  \item Emily Brontë
  \item Siglo XVIII
  \item Siglo XVII
  \item Siglo XX
  \item Siglo XVI
  \item Siglo XXI
\end{multicols}
\end{enumerate}
\vspace{0.15cm}
\item ¿Quién es el autor de Jane Eyre?
\begin{enumerate}[label=\textbf{\Alph*)}]
\begin{multicols}{2}
  \item Charlotte Brontë
  \item Siglo XVIII
  \item Siglo XVII
  \item Siglo XX
  \item Siglo XVI
  \item Siglo XXI
\end{multicols}
\end{enumerate}
\vspace{0.15cm}
\item ¿Quién es el autor de Middlemarch?
\begin{enumerate}[label=\textbf{\Alph*)}]
\begin{multicols}{2}
  \item George Eliot
  \item Siglo XVIII
  \item Siglo XVII
  \item Siglo XX
  \item Siglo XVI
  \item Siglo XXI
\end{multicols}
\end{enumerate}
\vspace{0.15cm}
\item ¿Quién es el autor de Ivanhoe?
\begin{enumerate}[label=\textbf{\Alph*)}]
\begin{multicols}{2}
  \item Walter Scott
  \item Siglo XVIII
  \item Siglo XVII
  \item Siglo XX
  \item Siglo XVI
  \item Siglo XXI
\end{multicols}
\end{enumerate}
\vspace{0.15cm}
\item ¿Quién es el autor de Las aventuras de Robin Hood?
\begin{enumerate}[label=\textbf{\Alph*)}]
\begin{multicols}{2}
  \item Walter Scott
  \item Siglo XVIII
  \item Siglo XVII
  \item Siglo XX
  \item Siglo XVI
  \item Siglo XXI
\end{multicols}
\end{enumerate}
\vspace{0.15cm}
\item ¿Quién es el autor de Las peregrinaciones de Childe Harold?
\begin{enumerate}[label=\textbf{\Alph*)}]
\begin{multicols}{2}
  \item Lord Byron
  \item Siglo XVIII
  \item Siglo XVII
  \item Siglo XX
  \item Siglo XVI
  \item Siglo XXI
\end{multicols}
\end{enumerate}
\vspace{0.15cm}
\item ¿Quién es el autor de Prometeo liberado?
\begin{enumerate}[label=\textbf{\Alph*)}]
\begin{multicols}{2}
  \item Percy Bysshe Shelley
  \item Siglo XVIII
  \item Siglo XVII
  \item Siglo XX
  \item Siglo XVI
  \item Siglo XXI
\end{multicols}
\end{enumerate}
\vspace{0.15cm}
\item ¿Quién es el autor de Oda a una urna griega?
\begin{enumerate}[label=\textbf{\Alph*)}]
\begin{multicols}{2}
  \item John Keats
  \item Siglo XVIII
  \item Siglo XVII
  \item Siglo XX
  \item Siglo XVI
  \item Siglo XXI
\end{multicols}
\end{enumerate}
\vspace{0.15cm}
\item ¿Quién es el autor de El cuervo?
\begin{enumerate}[label=\textbf{\Alph*)}]
\begin{multicols}{2}
  \item Edgar Allan Poe
  \item Siglo XVIII
  \item Siglo XVII
  \item Siglo XX
  \item Siglo XVI
  \item Siglo XXI
\end{multicols}
\end{enumerate}
\vspace{0.15cm}
\item ¿Quién es el autor de La caída de la casa Usher?
\begin{enumerate}[label=\textbf{\Alph*)}]
\begin{multicols}{2}
  \item Edgar Allan Poe
  \item Siglo XVIII
  \item Siglo XVII
  \item Siglo XX
  \item Siglo XVI
  \item Siglo XXI
\end{multicols}
\end{enumerate}
\vspace{0.15cm}
\item ¿Quién es el autor de Fausto?
\begin{enumerate}[label=\textbf{\Alph*)}]
\begin{multicols}{2}
  \item Johann Wolfgang von Goethe
  \item Siglo XVIII
  \item Siglo XVII
  \item Siglo XX
  \item Siglo XVI
  \item Siglo XXI
\end{multicols}
\end{enumerate}
\vspace{0.15cm}
\item ¿Quién es el autor de Los bandidos?
\begin{enumerate}[label=\textbf{\Alph*)}]
\begin{multicols}{2}
  \item Friedrich Schiller
  \item Siglo XVIII
  \item Siglo XVII
  \item Siglo XX
  \item Siglo XVI
  \item Siglo XXI
\end{multicols}
\end{enumerate}
\vspace{0.15cm}
\item ¿Quién es el autor de Himnos a la noche?
\begin{enumerate}[label=\textbf{\Alph*)}]
\begin{multicols}{2}
  \item Novalis
  \item Siglo XVIII
  \item Siglo XVII
  \item Siglo XX
  \item Siglo XVI
  \item Siglo XXI
\end{multicols}
\end{enumerate}
\vspace{0.15cm}
\item ¿Quién es el autor de El hombre de la arena?
\begin{enumerate}[label=\textbf{\Alph*)}]
\begin{multicols}{2}
  \item E.T.A. Hoffmann
  \item Siglo XVIII
  \item Siglo XVII
  \item Siglo XX
  \item Siglo XVI
  \item Siglo XXI
\end{multicols}
\end{enumerate}
\vspace{0.15cm}
\item ¿Quién es el autor de Alemania, un cuento de invierno?
\begin{enumerate}[label=\textbf{\Alph*)}]
\begin{multicols}{2}
  \item Heinrich Heine
  \item Siglo XVIII
  \item Siglo XVII
  \item Siglo XX
  \item Siglo XVI
  \item Siglo XXI
\end{multicols}
\end{enumerate}
\vspace{0.15cm}
\item ¿Quién es el autor de Eugenia Grandet?
\begin{enumerate}[label=\textbf{\Alph*)}]
\begin{multicols}{2}
  \item Honoré de Balzac
  \item Siglo XVIII
  \item Siglo XVII
  \item Siglo XX
  \item Siglo XVI
  \item Siglo XXI
\end{multicols}
\end{enumerate}
\vspace{0.15cm}
\item ¿Quién es el autor de Papá Goriot?
\begin{enumerate}[label=\textbf{\Alph*)}]
\begin{multicols}{2}
  \item Honoré de Balzac
  \item Siglo XVIII
  \item Siglo XVII
  \item Siglo XX
  \item Siglo XVI
  \item Siglo XXI
\end{multicols}
\end{enumerate}
\vspace{0.15cm}
\item ¿Quién es el autor de Los trabajadores del mar?
\begin{enumerate}[label=\textbf{\Alph*)}]
\begin{multicols}{2}
  \item Victor Hugo
  \item Siglo XVIII
  \item Siglo XVII
  \item Siglo XX
  \item Siglo XVI
  \item Siglo XXI
\end{multicols}
\end{enumerate}
\vspace{0.15cm}
\item ¿Quién es el autor de La dama de las camelias?
\begin{enumerate}[label=\textbf{\Alph*)}]
\begin{multicols}{2}
  \item Alejandro Dumas (hijo)
  \item Siglo XVIII
  \item Siglo XVII
  \item Siglo XX
  \item Siglo XVI
  \item Siglo XXI
\end{multicols}
\end{enumerate}
\vspace{0.15cm}
\item ¿Quién es el autor de Salambó?
\begin{enumerate}[label=\textbf{\Alph*)}]
\begin{multicols}{2}
  \item Gustave Flaubert
  \item Siglo XVIII
  \item Siglo XVII
  \item Siglo XX
  \item Siglo XVI
  \item Siglo XXI
\end{multicols}
\end{enumerate}
\vspace{0.15cm}
\item ¿Quién es el autor de La taberna?
\begin{enumerate}[label=\textbf{\Alph*)}]
\begin{multicols}{2}
  \item Émile Zola
  \item Siglo XVIII
  \item Siglo XVII
  \item Siglo XX
  \item Siglo XVI
  \item Siglo XXI
\end{multicols}
\end{enumerate}
\vspace{0.15cm}
\item ¿Quién es el autor de La bestia humana?
\begin{enumerate}[label=\textbf{\Alph*)}]
\begin{multicols}{2}
  \item Émile Zola
  \item Siglo XVIII
  \item Siglo XVII
  \item Siglo XX
  \item Siglo XVI
  \item Siglo XXI
\end{multicols}
\end{enumerate}
\vspace{0.15cm}
\item ¿Quién es el autor de El vientre de París?
\begin{enumerate}[label=\textbf{\Alph*)}]
\begin{multicols}{2}
  \item Émile Zola
  \item Siglo XVIII
  \item Siglo XVII
  \item Siglo XX
  \item Siglo XVI
  \item Siglo XXI
\end{multicols}
\end{enumerate}
\vspace{0.15cm}
\item ¿Quién es el autor de las Rimas y Leyendas?
\begin{enumerate}[label=\textbf{\Alph*)}]
\begin{multicols}{2}
  \item Gustavo Adolfo Bécquer
  \item Siglo XVIII
  \item Siglo XVII
  \item Siglo XX
  \item Siglo XVI
  \item Siglo XXI
\end{multicols}
\end{enumerate}
\vspace{0.15cm}
\item ¿Quién es el autor de Don Juan Tenorio?
\begin{enumerate}[label=\textbf{\Alph*)}]
\begin{multicols}{2}
  \item José Zorrilla
  \item Siglo XVIII
  \item Siglo XVII
  \item Siglo XX
  \item Siglo XVI
  \item Siglo XXI
\end{multicols}
\end{enumerate}
\vspace{0.15cm}
\item ¿Quién es el autor de El estudiante de Salamanca?
\begin{enumerate}[label=\textbf{\Alph*)}]
\begin{multicols}{2}
  \item José de Espronceda
  \item Siglo XVIII
  \item Siglo XVII
  \item Siglo XX
  \item Siglo XVI
  \item Siglo XXI
\end{multicols}
\end{enumerate}
\vspace{0.15cm}
\item ¿Quién es el autor de La canción del pirata?
\begin{enumerate}[label=\textbf{\Alph*)}]
\begin{multicols}{2}
  \item José de Espronceda
  \item Siglo XVIII
  \item Siglo XVII
  \item Siglo XX
  \item Siglo XVI
  \item Siglo XXI
\end{multicols}
\end{enumerate}
\vspace{0.15cm}
\item ¿Quién es el autor de Artículos de costumbres?
\begin{enumerate}[label=\textbf{\Alph*)}]
\begin{multicols}{2}
  \item Mariano José de Larra
  \item Siglo XVIII
  \item Siglo XVII
  \item Siglo XX
  \item Siglo XVI
  \item Siglo XXI
\end{multicols}
\end{enumerate}
\vspace{0.15cm}
\item ¿Quién es el autor de Don Álvaro o la fuerza del sino?
\begin{enumerate}[label=\textbf{\Alph*)}]
\begin{multicols}{2}
  \item Duque de Rivas
  \item Siglo XVIII
  \item Siglo XVII
  \item Siglo XX
  \item Siglo XVI
  \item Siglo XXI
\end{multicols}
\end{enumerate}
\vspace{0.15cm}
\end{enumerate}\n
\subsubsection*{SECCIÓN III: REALISMO – CONCEPTOS BÁSICOS (91–140)}

\begin{enumerate}[label=\textbf{\arabic*.}]
\setcounter{enumi}{90}
\item ¿Qué es el Realismo?
\begin{enumerate}[label=\textbf{\Alph*)}]
\begin{multicols}{2}
  \item Un movimiento literario del siglo XIX que observa y describe la realidad objetivamente
  \item Siglo XVIII
  \item Siglo XVII
  \item Siglo XX
  \item Siglo XVI
  \item Siglo XXI
\end{multicols}
\end{enumerate}
\vspace{0.15cm}
\item ¿En qué siglo surge el Realismo?
\begin{enumerate}[label=\textbf{\Alph*)}]
\begin{multicols}{2}
  \item Siglo XIX
  \item Siglo XVIII
  \item Siglo XVII
  \item Siglo XX
  \item Siglo XVI
  \item Siglo XXI
\end{multicols}
\end{enumerate}
\vspace{0.15cm}
\item ¿Cuál es el principal objetivo del Realismo?
\begin{enumerate}[label=\textbf{\Alph*)}]
\begin{multicols}{2}
  \item Representar la realidad tal como es
  \item Siglo XVIII
  \item Siglo XVII
  \item Siglo XX
  \item Siglo XVI
  \item Siglo XXI
\end{multicols}
\end{enumerate}
\vspace{0.15cm}
\item ¿Qué es la objetividad realista?
\begin{enumerate}[label=\textbf{\Alph*)}]
\begin{multicols}{2}
  \item La presentación imparcial de los hechos
  \item Siglo XVIII
  \item Siglo XVII
  \item Siglo XX
  \item Siglo XVI
  \item Siglo XXI
\end{multicols}
\end{enumerate}
\vspace{0.15cm}
\item ¿Qué es la observación en el Realismo?
\begin{enumerate}[label=\textbf{\Alph*)}]
\begin{multicols}{2}
  \item El método principal para captar la realidad
  \item Siglo XVIII
  \item Siglo XVII
  \item Siglo XX
  \item Siglo XVI
  \item Siglo XXI
\end{multicols}
\end{enumerate}
\vspace{0.15cm}
\item ¿Qué es la verosimilitud realista?
\begin{enumerate}[label=\textbf{\Alph*)}]
\begin{multicols}{2}
  \item La apariencia de verdad en la narración
  \item Siglo XVIII
  \item Siglo XVII
  \item Siglo XX
  \item Siglo XVI
  \item Siglo XXI
\end{multicols}
\end{enumerate}
\vspace{0.15cm}
\item ¿Qué es la novela realista?
\begin{enumerate}[label=\textbf{\Alph*)}]
\begin{multicols}{2}
  \item Un género narrativo que describe la realidad cotidiana
  \item Siglo XVIII
  \item Siglo XVII
  \item Siglo XX
  \item Siglo XVI
  \item Siglo XXI
\end{multicols}
\end{enumerate}
\vspace{0.15cm}
\item ¿Qué es el costumbrismo?
\begin{enumerate}[label=\textbf{\Alph*)}]
\begin{multicols}{2}
  \item La descripción de las costumbres y tipos de una época
  \item Siglo XVIII
  \item Siglo XVII
  \item Siglo XX
  \item Siglo XVI
  \item Siglo XXI
\end{multicols}
\end{enumerate}
\vspace{0.15cm}
\item ¿Qué es el determinismo en el Realismo?
\begin{enumerate}[label=\textbf{\Alph*)}]
\begin{multicols}{2}
  \item La idea de que el ambiente y la herencia determinan al individuo
  \item Siglo XVIII
  \item Siglo XVII
  \item Siglo XX
  \item Siglo XVI
  \item Siglo XXI
\end{multicols}
\end{enumerate}
\vspace{0.15cm}
\item ¿Qué es el materialismo en el Realismo?
\begin{enumerate}[label=\textbf{\Alph*)}]
\begin{multicols}{2}
  \item La atención a los aspectos materiales y económicos de la vida
  \item Siglo XVIII
  \item Siglo XVII
  \item Siglo XX
  \item Siglo XVI
  \item Siglo XXI
\end{multicols}
\end{enumerate}
\vspace{0.15cm}
\item ¿Qué es la novela de tesis?
\begin{enumerate}[label=\textbf{\Alph*)}]
\begin{multicols}{2}
  \item Una novela que defiende una idea o postura
  \item Siglo XVIII
  \item Siglo XVII
  \item Siglo XX
  \item Siglo XVI
  \item Siglo XXI
\end{multicols}
\end{enumerate}
\vspace{0.15cm}
\item ¿Qué es el narrador omnisciente?
\begin{enumerate}[label=\textbf{\Alph*)}]
\begin{multicols}{2}
  \item Un narrador que conoce todo sobre los personajes y los hechos
  \item Siglo XVIII
  \item Siglo XVII
  \item Siglo XX
  \item Siglo XVI
  \item Siglo XXI
\end{multicols}
\end{enumerate}
\vspace{0.15cm}
\item ¿Qué es el estilo indirecto libre?
\begin{enumerate}[label=\textbf{\Alph*)}]
\begin{multicols}{2}
  \item Una técnica narrativa que mezcla la voz del narrador y la del personaje
  \item Siglo XVIII
  \item Siglo XVII
  \item Siglo XX
  \item Siglo XVI
  \item Siglo XXI
\end{multicols}
\end{enumerate}
\vspace{0.15cm}
\item ¿Qué es el retrato en el Realismo?
\begin{enumerate}[label=\textbf{\Alph*)}]
\begin{multicols}{2}
  \item La descripción detallada de los personajes
  \item Siglo XVIII
  \item Siglo XVII
  \item Siglo XX
  \item Siglo XVI
  \item Siglo XXI
\end{multicols}
\end{enumerate}
\vspace{0.15cm}
\item ¿Qué es el ambiente en la novela realista?
\begin{enumerate}[label=\textbf{\Alph*)}]
\begin{multicols}{2}
  \item El conjunto de circunstancias que rodean a los personajes
  \item Siglo XVIII
  \item Siglo XVII
  \item Siglo XX
  \item Siglo XVI
  \item Siglo XXI
\end{multicols}
\end{enumerate}
\vspace{0.15cm}
\item ¿Qué es el tiempo en la novela realista?
\begin{enumerate}[label=\textbf{\Alph*)}]
\begin{multicols}{2}
  \item El momento histórico en que se desarrolla la acción
  \item Siglo XVIII
  \item Siglo XVII
  \item Siglo XX
  \item Siglo XVI
  \item Siglo XXI
\end{multicols}
\end{enumerate}
\vspace{0.15cm}
\item ¿Qué es el espacio en la novela realista?
\begin{enumerate}[label=\textbf{\Alph*)}]
\begin{multicols}{2}
  \item El lugar donde se desarrolla la acción
  \item Siglo XVIII
  \item Siglo XVII
  \item Siglo XX
  \item Siglo XVI
  \item Siglo XXI
\end{multicols}
\end{enumerate}
\vspace{0.15cm}
\item ¿Qué es el personaje realista?
\begin{enumerate}[label=\textbf{\Alph*)}]
\begin{multicols}{2}
  \item Un personaje complejo y verosímil
  \item Siglo XVIII
  \item Siglo XVII
  \item Siglo XX
  \item Siglo XVI
  \item Siglo XXI
\end{multicols}
\end{enumerate}
\vspace{0.15cm}
\item ¿Qué es el protagonista en la novela realista?
\begin{enumerate}[label=\textbf{\Alph*)}]
\begin{multicols}{2}
  \item El personaje principal cuya vida se narra
  \item Siglo XVIII
  \item Siglo XVII
  \item Siglo XX
  \item Siglo XVI
  \item Siglo XXI
\end{multicols}
\end{enumerate}
\vspace{0.15cm}
\item ¿Qué es el antagonista en la novela realista?
\begin{enumerate}[label=\textbf{\Alph*)}]
\begin{multicols}{2}
  \item El personaje que se opone al protagonista
  \item Siglo XVIII
  \item Siglo XVII
  \item Siglo XX
  \item Siglo XVI
  \item Siglo XXI
\end{multicols}
\end{enumerate}
\vspace{0.15cm}
\item ¿Qué es el conflicto en la novela realista?
\begin{enumerate}[label=\textbf{\Alph*)}]
\begin{multicols}{2}
  \item La oposición de fuerzas que genera la acción
  \item Siglo XVIII
  \item Siglo XVII
  \item Siglo XX
  \item Siglo XVI
  \item Siglo XXI
\end{multicols}
\end{enumerate}
\vspace{0.15cm}
\item ¿Qué es el clímax en la novela realista?
\begin{enumerate}[label=\textbf{\Alph*)}]
\begin{multicols}{2}
  \item El momento de máxima tensión
  \item Siglo XVIII
  \item Siglo XVII
  \item Siglo XX
  \item Siglo XVI
  \item Siglo XXI
\end{multicols}
\end{enumerate}
\vspace{0.15cm}
\item ¿Qué es el desenlace en la novela realista?
\begin{enumerate}[label=\textbf{\Alph*)}]
\begin{multicols}{2}
  \item La resolución del conflicto
  \item Siglo XVIII
  \item Siglo XVII
  \item Siglo XX
  \item Siglo XVI
  \item Siglo XXI
\end{multicols}
\end{enumerate}
\vspace{0.15cm}
\item ¿Qué es el Realismo crítico?
\begin{enumerate}[label=\textbf{\Alph*)}]
\begin{multicols}{2}
  \item El que denuncia los problemas sociales
  \item Siglo XVIII
  \item Siglo XVII
  \item Siglo XX
  \item Siglo XVI
  \item Siglo XXI
\end{multicols}
\end{enumerate}
\vspace{0.15cm}
\item ¿Qué es el Realismo psicológico?
\begin{enumerate}[label=\textbf{\Alph*)}]
\begin{multicols}{2}
  \item El que profundiza en la interioridad de los personajes
  \item Siglo XVIII
  \item Siglo XVII
  \item Siglo XX
  \item Siglo XVI
  \item Siglo XXI
\end{multicols}
\end{enumerate}
\vspace{0.15cm}
\item ¿Qué es el Realismo social?
\begin{enumerate}[label=\textbf{\Alph*)}]
\begin{multicols}{2}
  \item El que retrata las condiciones de vida de las clases populares
  \item Siglo XVIII
  \item Siglo XVII
  \item Siglo XX
  \item Siglo XVI
  \item Siglo XXI
\end{multicols}
\end{enumerate}
\vspace{0.15cm}
\item ¿Qué es el Realismo mágico?
\begin{enumerate}[label=\textbf{\Alph*)}]
\begin{multicols}{2}
  \item Un movimiento literario del siglo XX que mezcla lo real y lo fantástico
  \item Siglo XVIII
  \item Siglo XVII
  \item Siglo XX
  \item Siglo XVI
  \item Siglo XXI
\end{multicols}
\end{enumerate}
\vspace{0.15cm}
\item ¿Qué es el Naturalismo?
\begin{enumerate}[label=\textbf{\Alph*)}]
\begin{multicols}{2}
  \item Una corriente del Realismo que aplica el determinismo científico
  \item Siglo XVIII
  \item Siglo XVII
  \item Siglo XX
  \item Siglo XVI
  \item Siglo XXI
\end{multicols}
\end{enumerate}
\vspace{0.15cm}
\item ¿Quién es el principal representante del Naturalismo?
\begin{enumerate}[label=\textbf{\Alph*)}]
\begin{multicols}{2}
  \item Émile Zola
  \item Siglo XVIII
  \item Siglo XVII
  \item Siglo XX
  \item Siglo XVI
  \item Siglo XXI
\end{multicols}
\end{enumerate}
\vspace{0.15cm}
\item ¿Qué es el método experimental en el Naturalismo?
\begin{enumerate}[label=\textbf{\Alph*)}]
\begin{multicols}{2}
  \item La aplicación del método científico a la literatura
  \item Siglo XVIII
  \item Siglo XVII
  \item Siglo XX
  \item Siglo XVI
  \item Siglo XXI
\end{multicols}
\end{enumerate}
\vspace{0.15cm}
\item ¿Qué es la herencia en el Naturalismo?
\begin{enumerate}[label=\textbf{\Alph*)}]
\begin{multicols}{2}
  \item La transmisión de características biológicas y psicológicas
  \item Siglo XVIII
  \item Siglo XVII
  \item Siglo XX
  \item Siglo XVI
  \item Siglo XXI
\end{multicols}
\end{enumerate}
\vspace{0.15cm}
\item ¿Qué es el medio en el Naturalismo?
\begin{enumerate}[label=\textbf{\Alph*)}]
\begin{multicols}{2}
  \item El entorno social y natural que influye en el individuo
  \item Siglo XVIII
  \item Siglo XVII
  \item Siglo XX
  \item Siglo XVI
  \item Siglo XXI
\end{multicols}
\end{enumerate}
\vspace{0.15cm}
\item ¿Qué es el pesimismo naturalista?
\begin{enumerate}[label=\textbf{\Alph*)}]
\begin{multicols}{2}
  \item La visión negativa y determinista de la existencia
  \item Siglo XVIII
  \item Siglo XVII
  \item Siglo XX
  \item Siglo XVI
  \item Siglo XXI
\end{multicols}
\end{enumerate}
\vspace{0.15cm}
\item ¿Qué es la novela naturalista?
\begin{enumerate}[label=\textbf{\Alph*)}]
\begin{multicols}{2}
  \item Una novela que aplica el determinismo y el método experimental
  \item Siglo XVIII
  \item Siglo XVII
  \item Siglo XX
  \item Siglo XVI
  \item Siglo XXI
\end{multicols}
\end{enumerate}
\vspace{0.15cm}
\item ¿Qué es el Realismo en España?
\begin{enumerate}[label=\textbf{\Alph*)}]
\begin{multicols}{2}
  \item Un movimiento literario del siglo XIX que refleja la sociedad española
  \item Siglo XVIII
  \item Siglo XVII
  \item Siglo XX
  \item Siglo XVI
  \item Siglo XXI
\end{multicols}
\end{enumerate}
\vspace{0.15cm}
\item ¿Quién es el principal representante del Realismo español?
\begin{enumerate}[label=\textbf{\Alph*)}]
\begin{multicols}{2}
  \item Benito Pérez Galdós
  \item Siglo XVIII
  \item Siglo XVII
  \item Siglo XX
  \item Siglo XVI
  \item Siglo XXI
\end{multicols}
\end{enumerate}
\vspace{0.15cm}
\item ¿Quién es el autor de Fortunata y Jacinta?
\begin{enumerate}[label=\textbf{\Alph*)}]
\begin{multicols}{2}
  \item Benito Pérez Galdós
  \item Siglo XVIII
  \item Siglo XVII
  \item Siglo XX
  \item Siglo XVI
  \item Siglo XXI
\end{multicols}
\end{enumerate}
\vspace{0.15cm}
\item ¿Quién es el autor de La Regenta?
\begin{enumerate}[label=\textbf{\Alph*)}]
\begin{multicols}{2}
  \item Leopoldo Alas "Clarín"
  \item Siglo XVIII
  \item Siglo XVII
  \item Siglo XX
  \item Siglo XVI
  \item Siglo XXI
\end{multicols}
\end{enumerate}
\vspace{0.15cm}
\item ¿Quién es el autor de Los pazos de Ulloa?
\begin{enumerate}[label=\textbf{\Alph*)}]
\begin{multicols}{2}
  \item Emilia Pardo Bazán
  \item Siglo XVIII
  \item Siglo XVII
  \item Siglo XX
  \item Siglo XVI
  \item Siglo XXI
\end{multicols}
\end{enumerate}
\vspace{0.15cm}
\item ¿Quién es el autor de Pepita Jiménez?
\begin{enumerate}[label=\textbf{\Alph*)}]
\begin{multicols}{2}
  \item Juan Valera
  \item Siglo XVIII
  \item Siglo XVII
  \item Siglo XX
  \item Siglo XVI
  \item Siglo XXI
\end{multicols}
\end{enumerate}
\vspace{0.15cm}
\item ¿Qué es el Realismo?
\begin{enumerate}[label=\textbf{\Alph*)}]
\begin{multicols}{2}
  \item Un movimiento literario del siglo XIX que observa y describe la realidad objetivamente
  \item Siglo XVIII
  \item Siglo XVII
  \item Siglo XX
  \item Siglo XVI
  \item Siglo XXI
\end{multicols}
\end{enumerate}
\vspace{0.15cm}
\item ¿En qué siglo surge el Realismo?
\begin{enumerate}[label=\textbf{\Alph*)}]
\begin{multicols}{2}
  \item Siglo XIX
  \item Siglo XVIII
  \item Siglo XVII
  \item Siglo XX
  \item Siglo XVI
  \item Siglo XXI
\end{multicols}
\end{enumerate}
\vspace{0.15cm}
\item ¿Cuál es el principal objetivo del Realismo?
\begin{enumerate}[label=\textbf{\Alph*)}]
\begin{multicols}{2}
  \item Representar la realidad tal como es
  \item Siglo XVIII
  \item Siglo XVII
  \item Siglo XX
  \item Siglo XVI
  \item Siglo XXI
\end{multicols}
\end{enumerate}
\vspace{0.15cm}
\item ¿Qué es la objetividad realista?
\begin{enumerate}[label=\textbf{\Alph*)}]
\begin{multicols}{2}
  \item La presentación imparcial de los hechos
  \item Siglo XVIII
  \item Siglo XVII
  \item Siglo XX
  \item Siglo XVI
  \item Siglo XXI
\end{multicols}
\end{enumerate}
\vspace{0.15cm}
\item ¿Qué es la observación en el Realismo?
\begin{enumerate}[label=\textbf{\Alph*)}]
\begin{multicols}{2}
  \item El método principal para captar la realidad
  \item Siglo XVIII
  \item Siglo XVII
  \item Siglo XX
  \item Siglo XVI
  \item Siglo XXI
\end{multicols}
\end{enumerate}
\vspace{0.15cm}
\item ¿Qué es la verosimilitud realista?
\begin{enumerate}[label=\textbf{\Alph*)}]
\begin{multicols}{2}
  \item La apariencia de verdad en la narración
  \item Siglo XVIII
  \item Siglo XVII
  \item Siglo XX
  \item Siglo XVI
  \item Siglo XXI
\end{multicols}
\end{enumerate}
\vspace{0.15cm}
\item ¿Qué es la novela realista?
\begin{enumerate}[label=\textbf{\Alph*)}]
\begin{multicols}{2}
  \item Un género narrativo que describe la realidad cotidiana
  \item Siglo XVIII
  \item Siglo XVII
  \item Siglo XX
  \item Siglo XVI
  \item Siglo XXI
\end{multicols}
\end{enumerate}
\vspace{0.15cm}
\item ¿Qué es el costumbrismo?
\begin{enumerate}[label=\textbf{\Alph*)}]
\begin{multicols}{2}
  \item La descripción de las costumbres y tipos de una época
  \item Siglo XVIII
  \item Siglo XVII
  \item Siglo XX
  \item Siglo XVI
  \item Siglo XXI
\end{multicols}
\end{enumerate}
\vspace{0.15cm}
\item ¿Qué es el determinismo en el Realismo?
\begin{enumerate}[label=\textbf{\Alph*)}]
\begin{multicols}{2}
  \item La idea de que el ambiente y la herencia determinan al individuo
  \item Siglo XVIII
  \item Siglo XVII
  \item Siglo XX
  \item Siglo XVI
  \item Siglo XXI
\end{multicols}
\end{enumerate}
\vspace{0.15cm}
\item ¿Qué es el materialismo en el Realismo?
\begin{enumerate}[label=\textbf{\Alph*)}]
\begin{multicols}{2}
  \item La atención a los aspectos materiales y económicos de la vida
  \item Siglo XVIII
  \item Siglo XVII
  \item Siglo XX
  \item Siglo XVI
  \item Siglo XXI
\end{multicols}
\end{enumerate}
\vspace{0.15cm}
\end{enumerate}\n
\subsubsection*{SECCIÓN IV: REALISMO – AUTORES Y OBRAS (141–190)}

\begin{enumerate}[label=\textbf{\arabic*.}]
\setcounter{enumi}{140}
\item ¿Quién es el autor de Ana Karenina?
\begin{enumerate}[label=\textbf{\Alph*)}]
\begin{multicols}{2}
  \item León Tolstói
  \item Siglo XVIII
  \item Siglo XVII
  \item Siglo XX
  \item Siglo XVI
  \item Siglo XXI
\end{multicols}
\end{enumerate}
\vspace{0.15cm}
\item ¿Quién es el autor de Guerra y paz?
\begin{enumerate}[label=\textbf{\Alph*)}]
\begin{multicols}{2}
  \item León Tolstói
  \item Siglo XVIII
  \item Siglo XVII
  \item Siglo XX
  \item Siglo XVI
  \item Siglo XXI
\end{multicols}
\end{enumerate}
\vspace{0.15cm}
\item ¿Quién es el autor de Crimen y castigo?
\begin{enumerate}[label=\textbf{\Alph*)}]
\begin{multicols}{2}
  \item Fiódor Dostoievski
  \item Siglo XVIII
  \item Siglo XVII
  \item Siglo XX
  \item Siglo XVI
  \item Siglo XXI
\end{multicols}
\end{enumerate}
\vspace{0.15cm}
\item ¿Quién es el autor de Los hermanos Karamázov?
\begin{enumerate}[label=\textbf{\Alph*)}]
\begin{multicols}{2}
  \item Fiódor Dostoievski
  \item Siglo XVIII
  \item Siglo XVII
  \item Siglo XX
  \item Siglo XVI
  \item Siglo XXI
\end{multicols}
\end{enumerate}
\vspace{0.15cm}
\item ¿Quién es el autor de El idiota?
\begin{enumerate}[label=\textbf{\Alph*)}]
\begin{multicols}{2}
  \item Fiódor Dostoievski
  \item Siglo XVIII
  \item Siglo XVII
  \item Siglo XX
  \item Siglo XVI
  \item Siglo XXI
\end{multicols}
\end{enumerate}
\vspace{0.15cm}
\item ¿Quién es el autor de Padres e hijos?
\begin{enumerate}[label=\textbf{\Alph*)}]
\begin{multicols}{2}
  \item Iván Turguénev
  \item Siglo XVIII
  \item Siglo XVII
  \item Siglo XX
  \item Siglo XVI
  \item Siglo XXI
\end{multicols}
\end{enumerate}
\vspace{0.15cm}
\item ¿Quién es el autor de Almas muertas?
\begin{enumerate}[label=\textbf{\Alph*)}]
\begin{multicols}{2}
  \item Nikolái Gógol
  \item Siglo XVIII
  \item Siglo XVII
  \item Siglo XX
  \item Siglo XVI
  \item Siglo XXI
\end{multicols}
\end{enumerate}
\vspace{0.15cm}
\item ¿Quién es el autor de La gaviota?
\begin{enumerate}[label=\textbf{\Alph*)}]
\begin{multicols}{2}
  \item Antón Chéjov
  \item Siglo XVIII
  \item Siglo XVII
  \item Siglo XX
  \item Siglo XVI
  \item Siglo XXI
\end{multicols}
\end{enumerate}
\vspace{0.15cm}
\item ¿Quién es el autor de Tío Vania?
\begin{enumerate}[label=\textbf{\Alph*)}]
\begin{multicols}{2}
  \item Antón Chéjov
  \item Siglo XVIII
  \item Siglo XVII
  \item Siglo XX
  \item Siglo XVI
  \item Siglo XXI
\end{multicols}
\end{enumerate}
\vspace{0.15cm}
\item ¿Quién es el autor de Las tres hermanas?
\begin{enumerate}[label=\textbf{\Alph*)}]
\begin{multicols}{2}
  \item Antón Chéjov
  \item Siglo XVIII
  \item Siglo XVII
  \item Siglo XX
  \item Siglo XVI
  \item Siglo XXI
\end{multicols}
\end{enumerate}
\vspace{0.15cm}
\item ¿Quién es el autor de El jardín de los cerezos?
\begin{enumerate}[label=\textbf{\Alph*)}]
\begin{multicols}{2}
  \item Antón Chéjov
  \item Siglo XVIII
  \item Siglo XVII
  \item Siglo XX
  \item Siglo XVI
  \item Siglo XXI
\end{multicols}
\end{enumerate}
\vspace{0.15cm}
\item ¿Quién es el autor de La dama del perrito?
\begin{enumerate}[label=\textbf{\Alph*)}]
\begin{multicols}{2}
  \item Antón Chéjov
  \item Siglo XVIII
  \item Siglo XVII
  \item Siglo XX
  \item Siglo XVI
  \item Siglo XXI
\end{multicols}
\end{enumerate}
\vspace{0.15cm}
\item ¿Quién es el autor de Tom Jones?
\begin{enumerate}[label=\textbf{\Alph*)}]
\begin{multicols}{2}
  \item Henry Fielding
  \item Siglo XVIII
  \item Siglo XVII
  \item Siglo XX
  \item Siglo XVI
  \item Siglo XXI
\end{multicols}
\end{enumerate}
\vspace{0.15cm}
\item ¿Quién es el autor de Oliver Twist?
\begin{enumerate}[label=\textbf{\Alph*)}]
\begin{multicols}{2}
  \item Charles Dickens
  \item Siglo XVIII
  \item Siglo XVII
  \item Siglo XX
  \item Siglo XVI
  \item Siglo XXI
\end{multicols}
\end{enumerate}
\vspace{0.15cm}
\item ¿Quién es el autor de David Copperfield?
\begin{enumerate}[label=\textbf{\Alph*)}]
\begin{multicols}{2}
  \item Charles Dickens
  \item Siglo XVIII
  \item Siglo XVII
  \item Siglo XX
  \item Siglo XVI
  \item Siglo XXI
\end{multicols}
\end{enumerate}
\vspace{0.15cm}
\item ¿Quién es el autor de Tiempos difíciles?
\begin{enumerate}[label=\textbf{\Alph*)}]
\begin{multicols}{2}
  \item Charles Dickens
  \item Siglo XVIII
  \item Siglo XVII
  \item Siglo XX
  \item Siglo XVI
  \item Siglo XXI
\end{multicols}
\end{enumerate}
\vspace{0.15cm}
\item ¿Quién es el autor de Grandes esperanzas?
\begin{enumerate}[label=\textbf{\Alph*)}]
\begin{multicols}{2}
  \item Charles Dickens
  \item Siglo XVIII
  \item Siglo XVII
  \item Siglo XX
  \item Siglo XVI
  \item Siglo XXI
\end{multicols}
\end{enumerate}
\vspace{0.15cm}
\item ¿Quién es el autor de La feria de las vanidades?
\begin{enumerate}[label=\textbf{\Alph*)}]
\begin{multicols}{2}
  \item William Thackeray
  \item Siglo XVIII
  \item Siglo XVII
  \item Siglo XX
  \item Siglo XVI
  \item Siglo XXI
\end{multicols}
\end{enumerate}
\vspace{0.15cm}
\item ¿Quién es el autor de El molino del Floss?
\begin{enumerate}[label=\textbf{\Alph*)}]
\begin{multicols}{2}
  \item George Eliot
  \item Siglo XVIII
  \item Siglo XVII
  \item Siglo XX
  \item Siglo XVI
  \item Siglo XXI
\end{multicols}
\end{enumerate}
\vspace{0.15cm}
\item ¿Quién es el autor de Tess de los d'Urberville?
\begin{enumerate}[label=\textbf{\Alph*)}]
\begin{multicols}{2}
  \item Thomas Hardy
  \item Siglo XVIII
  \item Siglo XVII
  \item Siglo XX
  \item Siglo XVI
  \item Siglo XXI
\end{multicols}
\end{enumerate}
\vspace{0.15cm}
\item ¿Quién es el autor de Jude el oscuro?
\begin{enumerate}[label=\textbf{\Alph*)}]
\begin{multicols}{2}
  \item Thomas Hardy
  \item Siglo XVIII
  \item Siglo XVII
  \item Siglo XX
  \item Siglo XVI
  \item Siglo XXI
\end{multicols}
\end{enumerate}
\vspace{0.15cm}
\item ¿Quién es el autor de Lejos del mundanal ruido?
\begin{enumerate}[label=\textbf{\Alph*)}]
\begin{multicols}{2}
  \item Thomas Hardy
  \item Siglo XVIII
  \item Siglo XVII
  \item Siglo XX
  \item Siglo XVI
  \item Siglo XXI
\end{multicols}
\end{enumerate}
\vspace{0.15cm}
\item ¿Quién es el autor de La letra escarlata?
\begin{enumerate}[label=\textbf{\Alph*)}]
\begin{multicols}{2}
  \item Nathaniel Hawthorne
  \item Siglo XVIII
  \item Siglo XVII
  \item Siglo XX
  \item Siglo XVI
  \item Siglo XXI
\end{multicols}
\end{enumerate}
\vspace{0.15cm}
\item ¿Quién es el autor de Moby Dick?
\begin{enumerate}[label=\textbf{\Alph*)}]
\begin{multicols}{2}
  \item Herman Melville
  \item Siglo XVIII
  \item Siglo XVII
  \item Siglo XX
  \item Siglo XVI
  \item Siglo XXI
\end{multicols}
\end{enumerate}
\vspace{0.15cm}
\item ¿Quién es el autor de Las aventuras de Tom Sawyer?
\begin{enumerate}[label=\textbf{\Alph*)}]
\begin{multicols}{2}
  \item Mark Twain
  \item Siglo XVIII
  \item Siglo XVII
  \item Siglo XX
  \item Siglo XVI
  \item Siglo XXI
\end{multicols}
\end{enumerate}
\vspace{0.15cm}
\item ¿Quién es el autor de Las aventuras de Huckleberry Finn?
\begin{enumerate}[label=\textbf{\Alph*)}]
\begin{multicols}{2}
  \item Mark Twain
  \item Siglo XVIII
  \item Siglo XVII
  \item Siglo XX
  \item Siglo XVI
  \item Siglo XXI
\end{multicols}
\end{enumerate}
\vspace{0.15cm}
\item ¿Quién es el autor de Retrato de una dama?
\begin{enumerate}[label=\textbf{\Alph*)}]
\begin{multicols}{2}
  \item Henry James
  \item Siglo XVIII
  \item Siglo XVII
  \item Siglo XX
  \item Siglo XVI
  \item Siglo XXI
\end{multicols}
\end{enumerate}
\vspace{0.15cm}
\item ¿Quién es el autor de La edad de la inocencia?
\begin{enumerate}[label=\textbf{\Alph*)}]
\begin{multicols}{2}
  \item Edith Wharton
  \item Siglo XVIII
  \item Siglo XVII
  \item Siglo XX
  \item Siglo XVI
  \item Siglo XXI
\end{multicols}
\end{enumerate}
\vspace{0.15cm}
\item ¿Quién es el autor de Sister Carrie?
\begin{enumerate}[label=\textbf{\Alph*)}]
\begin{multicols}{2}
  \item Theodore Dreiser
  \item Siglo XVIII
  \item Siglo XVII
  \item Siglo XX
  \item Siglo XVI
  \item Siglo XXI
\end{multicols}
\end{enumerate}
\vspace{0.15cm}
\item ¿Quién es el autor de Ana Karenina?
\begin{enumerate}[label=\textbf{\Alph*)}]
\begin{multicols}{2}
  \item León Tolstói
  \item Siglo XVIII
  \item Siglo XVII
  \item Siglo XX
  \item Siglo XVI
  \item Siglo XXI
\end{multicols}
\end{enumerate}
\vspace{0.15cm}
\item ¿Quién es el autor de Guerra y paz?
\begin{enumerate}[label=\textbf{\Alph*)}]
\begin{multicols}{2}
  \item León Tolstói
  \item Siglo XVIII
  \item Siglo XVII
  \item Siglo XX
  \item Siglo XVI
  \item Siglo XXI
\end{multicols}
\end{enumerate}
\vspace{0.15cm}
\item ¿Quién es el autor de Crimen y castigo?
\begin{enumerate}[label=\textbf{\Alph*)}]
\begin{multicols}{2}
  \item Fiódor Dostoievski
  \item Siglo XVIII
  \item Siglo XVII
  \item Siglo XX
  \item Siglo XVI
  \item Siglo XXI
\end{multicols}
\end{enumerate}
\vspace{0.15cm}
\item ¿Quién es el autor de Los hermanos Karamázov?
\begin{enumerate}[label=\textbf{\Alph*)}]
\begin{multicols}{2}
  \item Fiódor Dostoievski
  \item Siglo XVIII
  \item Siglo XVII
  \item Siglo XX
  \item Siglo XVI
  \item Siglo XXI
\end{multicols}
\end{enumerate}
\vspace{0.15cm}
\item ¿Quién es el autor de El idiota?
\begin{enumerate}[label=\textbf{\Alph*)}]
\begin{multicols}{2}
  \item Fiódor Dostoievski
  \item Siglo XVIII
  \item Siglo XVII
  \item Siglo XX
  \item Siglo XVI
  \item Siglo XXI
\end{multicols}
\end{enumerate}
\vspace{0.15cm}
\item ¿Quién es el autor de Padres e hijos?
\begin{enumerate}[label=\textbf{\Alph*)}]
\begin{multicols}{2}
  \item Iván Turguénev
  \item Siglo XVIII
  \item Siglo XVII
  \item Siglo XX
  \item Siglo XVI
  \item Siglo XXI
\end{multicols}
\end{enumerate}
\vspace{0.15cm}
\item ¿Quién es el autor de Almas muertas?
\begin{enumerate}[label=\textbf{\Alph*)}]
\begin{multicols}{2}
  \item Nikolái Gógol
  \item Siglo XVIII
  \item Siglo XVII
  \item Siglo XX
  \item Siglo XVI
  \item Siglo XXI
\end{multicols}
\end{enumerate}
\vspace{0.15cm}
\item ¿Quién es el autor de La gaviota?
\begin{enumerate}[label=\textbf{\Alph*)}]
\begin{multicols}{2}
  \item Antón Chéjov
  \item Siglo XVIII
  \item Siglo XVII
  \item Siglo XX
  \item Siglo XVI
  \item Siglo XXI
\end{multicols}
\end{enumerate}
\vspace{0.15cm}
\item ¿Quién es el autor de Tío Vania?
\begin{enumerate}[label=\textbf{\Alph*)}]
\begin{multicols}{2}
  \item Antón Chéjov
  \item Siglo XVIII
  \item Siglo XVII
  \item Siglo XX
  \item Siglo XVI
  \item Siglo XXI
\end{multicols}
\end{enumerate}
\vspace{0.15cm}
\item ¿Quién es el autor de Las tres hermanas?
\begin{enumerate}[label=\textbf{\Alph*)}]
\begin{multicols}{2}
  \item Antón Chéjov
  \item Siglo XVIII
  \item Siglo XVII
  \item Siglo XX
  \item Siglo XVI
  \item Siglo XXI
\end{multicols}
\end{enumerate}
\vspace{0.15cm}
\item ¿Quién es el autor de El jardín de los cerezos?
\begin{enumerate}[label=\textbf{\Alph*)}]
\begin{multicols}{2}
  \item Antón Chéjov
  \item Siglo XVIII
  \item Siglo XVII
  \item Siglo XX
  \item Siglo XVI
  \item Siglo XXI
\end{multicols}
\end{enumerate}
\vspace{0.15cm}
\item ¿Quién es el autor de La dama del perrito?
\begin{enumerate}[label=\textbf{\Alph*)}]
\begin{multicols}{2}
  \item Antón Chéjov
  \item Siglo XVIII
  \item Siglo XVII
  \item Siglo XX
  \item Siglo XVI
  \item Siglo XXI
\end{multicols}
\end{enumerate}
\vspace{0.15cm}
\item ¿Quién es el autor de Tom Jones?
\begin{enumerate}[label=\textbf{\Alph*)}]
\begin{multicols}{2}
  \item Henry Fielding
  \item Siglo XVIII
  \item Siglo XVII
  \item Siglo XX
  \item Siglo XVI
  \item Siglo XXI
\end{multicols}
\end{enumerate}
\vspace{0.15cm}
\item ¿Quién es el autor de Oliver Twist?
\begin{enumerate}[label=\textbf{\Alph*)}]
\begin{multicols}{2}
  \item Charles Dickens
  \item Siglo XVIII
  \item Siglo XVII
  \item Siglo XX
  \item Siglo XVI
  \item Siglo XXI
\end{multicols}
\end{enumerate}
\vspace{0.15cm}
\item ¿Quién es el autor de David Copperfield?
\begin{enumerate}[label=\textbf{\Alph*)}]
\begin{multicols}{2}
  \item Charles Dickens
  \item Siglo XVIII
  \item Siglo XVII
  \item Siglo XX
  \item Siglo XVI
  \item Siglo XXI
\end{multicols}
\end{enumerate}
\vspace{0.15cm}
\item ¿Quién es el autor de Tiempos difíciles?
\begin{enumerate}[label=\textbf{\Alph*)}]
\begin{multicols}{2}
  \item Charles Dickens
  \item Siglo XVIII
  \item Siglo XVII
  \item Siglo XX
  \item Siglo XVI
  \item Siglo XXI
\end{multicols}
\end{enumerate}
\vspace{0.15cm}
\item ¿Quién es el autor de Grandes esperanzas?
\begin{enumerate}[label=\textbf{\Alph*)}]
\begin{multicols}{2}
  \item Charles Dickens
  \item Siglo XVIII
  \item Siglo XVII
  \item Siglo XX
  \item Siglo XVI
  \item Siglo XXI
\end{multicols}
\end{enumerate}
\vspace{0.15cm}
\item ¿Quién es el autor de La feria de las vanidades?
\begin{enumerate}[label=\textbf{\Alph*)}]
\begin{multicols}{2}
  \item William Thackeray
  \item Siglo XVIII
  \item Siglo XVII
  \item Siglo XX
  \item Siglo XVI
  \item Siglo XXI
\end{multicols}
\end{enumerate}
\vspace{0.15cm}
\item ¿Quién es el autor de El molino del Floss?
\begin{enumerate}[label=\textbf{\Alph*)}]
\begin{multicols}{2}
  \item George Eliot
  \item Siglo XVIII
  \item Siglo XVII
  \item Siglo XX
  \item Siglo XVI
  \item Siglo XXI
\end{multicols}
\end{enumerate}
\vspace{0.15cm}
\item ¿Quién es el autor de Tess de los d'Urberville?
\begin{enumerate}[label=\textbf{\Alph*)}]
\begin{multicols}{2}
  \item Thomas Hardy
  \item Siglo XVIII
  \item Siglo XVII
  \item Siglo XX
  \item Siglo XVI
  \item Siglo XXI
\end{multicols}
\end{enumerate}
\vspace{0.15cm}
\item ¿Quién es el autor de Jude el oscuro?
\begin{enumerate}[label=\textbf{\Alph*)}]
\begin{multicols}{2}
  \item Thomas Hardy
  \item Siglo XVIII
  \item Siglo XVII
  \item Siglo XX
  \item Siglo XVI
  \item Siglo XXI
\end{multicols}
\end{enumerate}
\vspace{0.15cm}
\end{enumerate}\n
\subsubsection*{SECCIÓN V: NATURALISMO (191–230)}

\begin{enumerate}[label=\textbf{\arabic*.}]
\setcounter{enumi}{190}
\item ¿Qué es el Naturalismo?
\begin{enumerate}[label=\textbf{\Alph*)}]
\begin{multicols}{2}
  \item Una corriente del Realismo que aplica el determinismo científico
  \item Siglo XVIII
  \item Siglo XVII
  \item Siglo XX
  \item Siglo XVI
  \item Siglo XXI
\end{multicols}
\end{enumerate}
\vspace{0.15cm}
\item ¿Quién es el principal representante del Naturalismo?
\begin{enumerate}[label=\textbf{\Alph*)}]
\begin{multicols}{2}
  \item Émile Zola
  \item Siglo XVIII
  \item Siglo XVII
  \item Siglo XX
  \item Siglo XVI
  \item Siglo XXI
\end{multicols}
\end{enumerate}
\vspace{0.15cm}
\item ¿Qué es el método experimental en el Naturalismo?
\begin{enumerate}[label=\textbf{\Alph*)}]
\begin{multicols}{2}
  \item La aplicación del método científico a la literatura
  \item Siglo XVIII
  \item Siglo XVII
  \item Siglo XX
  \item Siglo XVI
  \item Siglo XXI
\end{multicols}
\end{enumerate}
\vspace{0.15cm}
\item ¿Qué es la herencia en el Naturalismo?
\begin{enumerate}[label=\textbf{\Alph*)}]
\begin{multicols}{2}
  \item La transmisión de características biológicas y psicológicas
  \item Siglo XVIII
  \item Siglo XVII
  \item Siglo XX
  \item Siglo XVI
  \item Siglo XXI
\end{multicols}
\end{enumerate}
\vspace{0.15cm}
\item ¿Qué es el medio en el Naturalismo?
\begin{enumerate}[label=\textbf{\Alph*)}]
\begin{multicols}{2}
  \item El entorno social y natural que influye en el individuo
  \item Siglo XVIII
  \item Siglo XVII
  \item Siglo XX
  \item Siglo XVI
  \item Siglo XXI
\end{multicols}
\end{enumerate}
\vspace{0.15cm}
\item ¿Qué es el pesimismo naturalista?
\begin{enumerate}[label=\textbf{\Alph*)}]
\begin{multicols}{2}
  \item La visión negativa y determinista de la existencia
  \item Siglo XVIII
  \item Siglo XVII
  \item Siglo XX
  \item Siglo XVI
  \item Siglo XXI
\end{multicols}
\end{enumerate}
\vspace{0.15cm}
\item ¿Qué es la novela naturalista?
\begin{enumerate}[label=\textbf{\Alph*)}]
\begin{multicols}{2}
  \item Una novela que aplica el determinismo y el método experimental
  \item Siglo XVIII
  \item Siglo XVII
  \item Siglo XX
  \item Siglo XVI
  \item Siglo XXI
\end{multicols}
\end{enumerate}
\vspace{0.15cm}
\item ¿Quién es el autor de Germinal?
\begin{enumerate}[label=\textbf{\Alph*)}]
\begin{multicols}{2}
  \item Émile Zola
  \item Siglo XVIII
  \item Siglo XVII
  \item Siglo XX
  \item Siglo XVI
  \item Siglo XXI
\end{multicols}
\end{enumerate}
\vspace{0.15cm}
\item ¿Quién es el autor de Naná?
\begin{enumerate}[label=\textbf{\Alph*)}]
\begin{multicols}{2}
  \item Émile Zola
  \item Siglo XVIII
  \item Siglo XVII
  \item Siglo XX
  \item Siglo XVI
  \item Siglo XXI
\end{multicols}
\end{enumerate}
\vspace{0.15cm}
\item ¿Quién es el autor de La taberna?
\begin{enumerate}[label=\textbf{\Alph*)}]
\begin{multicols}{2}
  \item Émile Zola
  \item Siglo XVIII
  \item Siglo XVII
  \item Siglo XX
  \item Siglo XVI
  \item Siglo XXI
\end{multicols}
\end{enumerate}
\vspace{0.15cm}
\item ¿Quién es el autor de La bestia humana?
\begin{enumerate}[label=\textbf{\Alph*)}]
\begin{multicols}{2}
  \item Émile Zola
  \item Siglo XVIII
  \item Siglo XVII
  \item Siglo XX
  \item Siglo XVI
  \item Siglo XXI
\end{multicols}
\end{enumerate}
\vspace{0.15cm}
\item ¿Quién es el autor de El vientre de París?
\begin{enumerate}[label=\textbf{\Alph*)}]
\begin{multicols}{2}
  \item Émile Zola
  \item Siglo XVIII
  \item Siglo XVII
  \item Siglo XX
  \item Siglo XVI
  \item Siglo XXI
\end{multicols}
\end{enumerate}
\vspace{0.15cm}
\item ¿Quién es el autor de La conquista de Plassans?
\begin{enumerate}[label=\textbf{\Alph*)}]
\begin{multicols}{2}
  \item Émile Zola
  \item Siglo XVIII
  \item Siglo XVII
  \item Siglo XX
  \item Siglo XVI
  \item Siglo XXI
\end{multicols}
\end{enumerate}
\vspace{0.15cm}
\item ¿Quién es el autor de La falta del abate Mouret?
\begin{enumerate}[label=\textbf{\Alph*)}]
\begin{multicols}{2}
  \item Émile Zola
  \item Siglo XVIII
  \item Siglo XVII
  \item Siglo XX
  \item Siglo XVI
  \item Siglo XXI
\end{multicols}
\end{enumerate}
\vspace{0.15cm}
\item ¿Quién es el autor de Una página de amor?
\begin{enumerate}[label=\textbf{\Alph*)}]
\begin{multicols}{2}
  \item Émile Zola
  \item Siglo XVIII
  \item Siglo XVII
  \item Siglo XX
  \item Siglo XVI
  \item Siglo XXI
\end{multicols}
\end{enumerate}
\vspace{0.15cm}
\item ¿Quién es el autor de Los pazos de Ulloa?
\begin{enumerate}[label=\textbf{\Alph*)}]
\begin{multicols}{2}
  \item Emilia Pardo Bazán
  \item Siglo XVIII
  \item Siglo XVII
  \item Siglo XX
  \item Siglo XVI
  \item Siglo XXI
\end{multicols}
\end{enumerate}
\vspace{0.15cm}
\item ¿Quién es el autor de La madre naturaleza?
\begin{enumerate}[label=\textbf{\Alph*)}]
\begin{multicols}{2}
  \item Emilia Pardo Bazán
  \item Siglo XVIII
  \item Siglo XVII
  \item Siglo XX
  \item Siglo XVI
  \item Siglo XXI
\end{multicols}
\end{enumerate}
\vspace{0.15cm}
\item ¿Quién es el autor de La barraca?
\begin{enumerate}[label=\textbf{\Alph*)}]
\begin{multicols}{2}
  \item Vicente Blasco Ibáñez
  \item Siglo XVIII
  \item Siglo XVII
  \item Siglo XX
  \item Siglo XVI
  \item Siglo XXI
\end{multicols}
\end{enumerate}
\vspace{0.15cm}
\item ¿Quién es el autor de Cañas y barro?
\begin{enumerate}[label=\textbf{\Alph*)}]
\begin{multicols}{2}
  \item Vicente Blasco Ibáñez
  \item Siglo XVIII
  \item Siglo XVII
  \item Siglo XX
  \item Siglo XVI
  \item Siglo XXI
\end{multicols}
\end{enumerate}
\vspace{0.15cm}
\item ¿Quién es el autor de Entre naranjos?
\begin{enumerate}[label=\textbf{\Alph*)}]
\begin{multicols}{2}
  \item Vicente Blasco Ibáñez
  \item Siglo XVIII
  \item Siglo XVII
  \item Siglo XX
  \item Siglo XVI
  \item Siglo XXI
\end{multicols}
\end{enumerate}
\vspace{0.15cm}
\item ¿Quién es el autor de La catedral?
\begin{enumerate}[label=\textbf{\Alph*)}]
\begin{multicols}{2}
  \item Vicente Blasco Ibáñez
  \item Siglo XVIII
  \item Siglo XVII
  \item Siglo XX
  \item Siglo XVI
  \item Siglo XXI
\end{multicols}
\end{enumerate}
\vspace{0.15cm}
\item ¿Quién es el autor de El intruso?
\begin{enumerate}[label=\textbf{\Alph*)}]
\begin{multicols}{2}
  \item Vicente Blasco Ibáñez
  \item Siglo XVIII
  \item Siglo XVII
  \item Siglo XX
  \item Siglo XVI
  \item Siglo XXI
\end{multicols}
\end{enumerate}
\vspace{0.15cm}
\item ¿Quién es el autor de La bodega?
\begin{enumerate}[label=\textbf{\Alph*)}]
\begin{multicols}{2}
  \item Vicente Blasco Ibáñez
  \item Siglo XVIII
  \item Siglo XVII
  \item Siglo XX
  \item Siglo XVI
  \item Siglo XXI
\end{multicols}
\end{enumerate}
\vspace{0.15cm}
\item ¿Quién es el autor de Sangre y arena?
\begin{enumerate}[label=\textbf{\Alph*)}]
\begin{multicols}{2}
  \item Vicente Blasco Ibáñez
  \item Siglo XVIII
  \item Siglo XVII
  \item Siglo XX
  \item Siglo XVI
  \item Siglo XXI
\end{multicols}
\end{enumerate}
\vspace{0.15cm}
\item ¿Quién es el autor de Los cuatro jinetes del Apocalipsis?
\begin{enumerate}[label=\textbf{\Alph*)}]
\begin{multicols}{2}
  \item Vicente Blasco Ibáñez
  \item Siglo XVIII
  \item Siglo XVII
  \item Siglo XX
  \item Siglo XVI
  \item Siglo XXI
\end{multicols}
\end{enumerate}
\vspace{0.15cm}
\item ¿Quién es el autor de La Regenta?
\begin{enumerate}[label=\textbf{\Alph*)}]
\begin{multicols}{2}
  \item Leopoldo Alas "Clarín"
  \item Siglo XVIII
  \item Siglo XVII
  \item Siglo XX
  \item Siglo XVI
  \item Siglo XXI
\end{multicols}
\end{enumerate}
\vspace{0.15cm}
\item ¿Quién es el autor de Su único hijo?
\begin{enumerate}[label=\textbf{\Alph*)}]
\begin{multicols}{2}
  \item Leopoldo Alas "Clarín"
  \item Siglo XVIII
  \item Siglo XVII
  \item Siglo XX
  \item Siglo XVI
  \item Siglo XXI
\end{multicols}
\end{enumerate}
\vspace{0.15cm}
\item ¿Quién es el autor de Pipá?
\begin{enumerate}[label=\textbf{\Alph*)}]
\begin{multicols}{2}
  \item Leopoldo Alas "Clarín"
  \item Siglo XVIII
  \item Siglo XVII
  \item Siglo XX
  \item Siglo XVI
  \item Siglo XXI
\end{multicols}
\end{enumerate}
\vspace{0.15cm}
\item ¿Quién es el autor de Adiós, Cordera?
\begin{enumerate}[label=\textbf{\Alph*)}]
\begin{multicols}{2}
  \item Leopoldo Alas "Clarín"
  \item Siglo XVIII
  \item Siglo XVII
  \item Siglo XX
  \item Siglo XVI
  \item Siglo XXI
\end{multicols}
\end{enumerate}
\vspace{0.15cm}
\item ¿Quién es el autor de Doña Berta?
\begin{enumerate}[label=\textbf{\Alph*)}]
\begin{multicols}{2}
  \item Leopoldo Alas "Clarín"
  \item Siglo XVIII
  \item Siglo XVII
  \item Siglo XX
  \item Siglo XVI
  \item Siglo XXI
\end{multicols}
\end{enumerate}
\vspace{0.15cm}
\item ¿Quién es el autor de Fortunata y Jacinta?
\begin{enumerate}[label=\textbf{\Alph*)}]
\begin{multicols}{2}
  \item Benito Pérez Galdós
  \item Siglo XVIII
  \item Siglo XVII
  \item Siglo XX
  \item Siglo XVI
  \item Siglo XXI
\end{multicols}
\end{enumerate}
\vspace{0.15cm}
\item ¿Quién es el autor de Misericordia?
\begin{enumerate}[label=\textbf{\Alph*)}]
\begin{multicols}{2}
  \item Benito Pérez Galdós
  \item Siglo XVIII
  \item Siglo XVII
  \item Siglo XX
  \item Siglo XVI
  \item Siglo XXI
\end{multicols}
\end{enumerate}
\vspace{0.15cm}
\item ¿Quién es el autor de Episodios nacionales?
\begin{enumerate}[label=\textbf{\Alph*)}]
\begin{multicols}{2}
  \item Benito Pérez Galdós
  \item Siglo XVIII
  \item Siglo XVII
  \item Siglo XX
  \item Siglo XVI
  \item Siglo XXI
\end{multicols}
\end{enumerate}
\vspace{0.15cm}
\item ¿Quién es el autor de Doña Perfecta?
\begin{enumerate}[label=\textbf{\Alph*)}]
\begin{multicols}{2}
  \item Benito Pérez Galdós
  \item Siglo XVIII
  \item Siglo XVII
  \item Siglo XX
  \item Siglo XVI
  \item Siglo XXI
\end{multicols}
\end{enumerate}
\vspace{0.15cm}
\item ¿Quién es el autor de La desheredada?
\begin{enumerate}[label=\textbf{\Alph*)}]
\begin{multicols}{2}
  \item Benito Pérez Galdós
  \item Siglo XVIII
  \item Siglo XVII
  \item Siglo XX
  \item Siglo XVI
  \item Siglo XXI
\end{multicols}
\end{enumerate}
\vspace{0.15cm}
\item ¿Quién es el autor de Tormento?
\begin{enumerate}[label=\textbf{\Alph*)}]
\begin{multicols}{2}
  \item Benito Pérez Galdós
  \item Siglo XVIII
  \item Siglo XVII
  \item Siglo XX
  \item Siglo XVI
  \item Siglo XXI
\end{multicols}
\end{enumerate}
\vspace{0.15cm}
\item ¿Quién es el autor de La de Bringas?
\begin{enumerate}[label=\textbf{\Alph*)}]
\begin{multicols}{2}
  \item Benito Pérez Galdós
  \item Siglo XVIII
  \item Siglo XVII
  \item Siglo XX
  \item Siglo XVI
  \item Siglo XXI
\end{multicols}
\end{enumerate}
\vspace{0.15cm}
\item ¿Quién es el autor de Tristana?
\begin{enumerate}[label=\textbf{\Alph*)}]
\begin{multicols}{2}
  \item Benito Pérez Galdós
  \item Siglo XVIII
  \item Siglo XVII
  \item Siglo XX
  \item Siglo XVI
  \item Siglo XXI
\end{multicols}
\end{enumerate}
\vspace{0.15cm}
\item ¿Quién es el autor de Nazarín?
\begin{enumerate}[label=\textbf{\Alph*)}]
\begin{multicols}{2}
  \item Benito Pérez Galdós
  \item Siglo XVIII
  \item Siglo XVII
  \item Siglo XX
  \item Siglo XVI
  \item Siglo XXI
\end{multicols}
\end{enumerate}
\vspace{0.15cm}
\item ¿Quién es el autor de Halma?
\begin{enumerate}[label=\textbf{\Alph*)}]
\begin{multicols}{2}
  \item Benito Pérez Galdós
  \item Siglo XVIII
  \item Siglo XVII
  \item Siglo XX
  \item Siglo XVI
  \item Siglo XXI
\end{multicols}
\end{enumerate}
\vspace{0.15cm}
\end{enumerate}\n
\subsubsection*{SECCIÓN VI: COMPARACIÓN ROMANTICISMO VS REALISMO (231–270)}

\begin{enumerate}[label=\textbf{\arabic*.}]
\setcounter{enumi}{230}
\item ¿Cuál es la principal diferencia entre Romanticismo y Realismo?
\begin{enumerate}[label=\textbf{\Alph*)}]
\begin{multicols}{2}
  \item El Romanticismo exalta el sentimiento; el Realismo observa la realidad
  \item Siglo XVIII
  \item Siglo XVII
  \item Siglo XX
  \item Siglo XVI
  \item Siglo XXI
\end{multicols}
\end{enumerate}
\vspace{0.15cm}
\item ¿Qué movimiento valora más la imaginación?
\begin{enumerate}[label=\textbf{\Alph*)}]
\begin{multicols}{2}
  \item Romanticismo
  \item Realismo
  \item Naturalismo
  \item Barroco
  \item Neoclasicismo
  \item Ilustración
\end{multicols}
\end{enumerate}
\vspace{0.15cm}
\item ¿Qué movimiento valora más la observación?
\begin{enumerate}[label=\textbf{\Alph*)}]
\begin{multicols}{2}
  \item Realismo
  \item Romanticismo
  \item Naturalismo
  \item Barroco
  \item Neoclasicismo
  \item Ilustración
\end{multicols}
\end{enumerate}
\vspace{0.15cm}
\item ¿Qué movimiento prefiere los temas fantásticos y sobrenaturales?
\begin{enumerate}[label=\textbf{\Alph*)}]
\begin{multicols}{2}
  \item Romanticismo
  \item Realismo
  \item Naturalismo
  \item Barroco
  \item Neoclasicismo
  \item Ilustración
\end{multicols}
\end{enumerate}
\vspace{0.15cm}
\item ¿Qué movimiento prefiere los temas cotidianos y sociales?
\begin{enumerate}[label=\textbf{\Alph*)}]
\begin{multicols}{2}
  \item Realismo
  \item Romanticismo
  \item Naturalismo
  \item Barroco
  \item Neoclasicismo
  \item Ilustración
\end{multicols}
\end{enumerate}
\vspace{0.15cm}
\item ¿Qué movimiento exalta el yo individual?
\begin{enumerate}[label=\textbf{\Alph*)}]
\begin{multicols}{2}
  \item Romanticismo
  \item Realismo
  \item Naturalismo
  \item Barroco
  \item Neoclasicismo
  \item Ilustración
\end{multicols}
\end{enumerate}
\vspace{0.15cm}
\item ¿Qué movimiento retrata la sociedad de forma objetiva?
\begin{enumerate}[label=\textbf{\Alph*)}]
\begin{multicols}{2}
  \item Realismo
  \item Romanticismo
  \item Naturalismo
  \item Barroco
  \item Neoclasicismo
  \item Ilustración
\end{multicols}
\end{enumerate}
\vspace{0.15cm}
\item ¿Qué movimiento prefiere la evasión?
\begin{enumerate}[label=\textbf{\Alph*)}]
\begin{multicols}{2}
  \item Romanticismo
  \item Realismo
  \item Naturalismo
  \item Barroco
  \item Neoclasicismo
  \item Ilustración
\end{multicols}
\end{enumerate}
\vspace{0.15cm}
\item ¿Qué movimiento prefiere el compromiso social?
\begin{enumerate}[label=\textbf{\Alph*)}]
\begin{multicols}{2}
  \item Realismo
  \item Romanticismo
  \item Naturalismo
  \item Barroco
  \item Neoclasicismo
  \item Ilustración
\end{multicols}
\end{enumerate}
\vspace{0.15cm}
\item ¿Qué movimiento exalta la libertad creadora?
\begin{enumerate}[label=\textbf{\Alph*)}]
\begin{multicols}{2}
  \item Romanticismo
  \item Realismo
  \item Naturalismo
  \item Barroco
  \item Neoclasicismo
  \item Ilustración
\end{multicols}
\end{enumerate}
\vspace{0.15cm}
\item ¿Qué movimiento respeta las reglas clásicas?
\begin{enumerate}[label=\textbf{\Alph*)}]
\begin{multicols}{2}
  \item Neoclasicismo
  \item Romanticismo
  \item Realismo
  \item Naturalismo
  \item Barroco
  \item Ilustración
\end{multicols}
\end{enumerate}
\vspace{0.15cm}
\item ¿Qué movimiento rompe las reglas clásicas?
\begin{enumerate}[label=\textbf{\Alph*)}]
\begin{multicols}{2}
  \item Romanticismo
  \item Realismo
  \item Naturalismo
  \item Barroco
  \item Neoclasicismo
  \item Ilustración
\end{multicols}
\end{enumerate}
\vspace{0.15cm}
\item ¿Qué movimiento prefiere la Edad Media?
\begin{enumerate}[label=\textbf{\Alph*)}]
\begin{multicols}{2}
  \item Romanticismo
  \item Realismo
  \item Naturalismo
  \item Barroco
  \item Neoclasicismo
  \item Ilustración
\end{multicols}
\end{enumerate}
\vspace{0.15cm}
\item ¿Qué movimiento prefiere la actualidad?
\begin{enumerate}[label=\textbf{\Alph*)}]
\begin{multicols}{2}
  \item Realismo
  \item Romanticismo
  \item Naturalismo
  \item Barroco
  \item Neoclasicismo
  \item Ilustración
\end{multicols}
\end{enumerate}
\vspace{0.15cm}
\item ¿Qué movimiento exalta la naturaleza salvaje?
\begin{enumerate}[label=\textbf{\Alph*)}]
\begin{multicols}{2}
  \item Romanticismo
  \item Realismo
  \item Naturalismo
  \item Barroco
  \item Neoclasicismo
  \item Ilustración
\end{multicols}
\end{enumerate}
\vspace{0.15cm}
\item ¿Qué movimiento describe la naturaleza urbana?
\begin{enumerate}[label=\textbf{\Alph*)}]
\begin{multicols}{2}
  \item Realismo
  \item Romanticismo
  \item Naturalismo
  \item Barroco
  \item Neoclasicismo
  \item Ilustración
\end{multicols}
\end{enumerate}
\vspace{0.15cm}
\item ¿Qué movimiento prefiere el héroe rebelde?
\begin{enumerate}[label=\textbf{\Alph*)}]
\begin{multicols}{2}
  \item Romanticismo
  \item Realismo
  \item Naturalismo
  \item Barroco
  \item Neoclasicismo
  \item Ilustración
\end{multicols}
\end{enumerate}
\vspace{0.15cm}
\item ¿Qué movimiento prefiere el personaje común?
\begin{enumerate}[label=\textbf{\Alph*)}]
\begin{multicols}{2}
  \item Realismo
  \item Romanticismo
  \item Naturalismo
  \item Barroco
  \item Neoclasicismo
  \item Ilustración
\end{multicols}
\end{enumerate}
\vspace{0.15cm}
\item ¿Qué movimiento prefiere el final trágico?
\begin{enumerate}[label=\textbf{\Alph*)}]
\begin{multicols}{2}
  \item Romanticismo
  \item Realismo
  \item Naturalismo
  \item Barroco
  \item Neoclasicismo
  \item Ilustración
\end{multicols}
\end{enumerate}
\vspace{0.15cm}
\item ¿Qué movimiento prefiere el final abierto?
\begin{enumerate}[label=\textbf{\Alph*)}]
\begin{multicols}{2}
  \item Realismo
  \item Romanticismo
  \item Naturalismo
  \item Barroco
  \item Neoclasicismo
  \item Ilustración
\end{multicols}
\end{enumerate}
\vspace{0.15cm}
\item ¿Qué movimiento prefiere el narrador omnisciente?
\begin{enumerate}[label=\textbf{\Alph*)}]
\begin{multicols}{2}
  \item Realismo
  \item Romanticismo
  \item Naturalismo
  \item Barroco
  \item Neoclasicismo
  \item Ilustración
\end{multicols}
\end{enumerate}
\vspace{0.15cm}
\item ¿Qué movimiento prefiere el narrador en primera persona?
\begin{enumerate}[label=\textbf{\Alph*)}]
\begin{multicols}{2}
  \item Romanticismo
  \item Realismo
  \item Naturalismo
  \item Barroco
  \item Neoclasicismo
  \item Ilustración
\end{multicols}
\end{enumerate}
\vspace{0.15cm}
\item ¿Qué movimiento prefiere la prosa?
\begin{enumerate}[label=\textbf{\Alph*)}]
\begin{multicols}{2}
  \item Realismo
  \item Romanticismo
  \item Naturalismo
  \item Barroco
  \item Neoclasicismo
  \item Ilustración
\end{multicols}
\end{enumerate}
\vspace{0.15cm}
\item ¿Qué movimiento prefiere el verso?
\begin{enumerate}[label=\textbf{\Alph*)}]
\begin{multicols}{2}
  \item Romanticismo
  \item Realismo
  \item Naturalismo
  \item Barroco
  \item Neoclasicismo
  \item Ilustración
\end{multicols}
\end{enumerate}
\vspace{0.15cm}
\item ¿Qué movimiento prefiere la novela?
\begin{enumerate}[label=\textbf{\Alph*)}]
\begin{multicols}{2}
  \item Realismo
  \item Romanticismo
  \item Naturalismo
  \item Barroco
  \item Neoclasicismo
  \item Ilustración
\end{multicols}
\end{enumerate}
\vspace{0.15cm}
\item ¿Qué movimiento prefiere la poesía?
\begin{enumerate}[label=\textbf{\Alph*)}]
\begin{multicols}{2}
  \item Romanticismo
  \item Realismo
  \item Naturalismo
  \item Barroco
  \item Neoclasicismo
  \item Ilustración
\end{multicols}
\end{enumerate}
\vspace{0.15cm}
\item ¿Qué movimiento prefiere el teatro?
\begin{enumerate}[label=\textbf{\Alph*)}]
\begin{multicols}{2}
  \item Romanticismo
  \item Realismo
  \item Naturalismo
  \item Barroco
  \item Neoclasicismo
  \item Ilustración
\end{multicols}
\end{enumerate}
\vspace{0.15cm}
\item ¿Qué movimiento prefiere el ensayo?
\begin{enumerate}[label=\textbf{\Alph*)}]
\begin{multicols}{2}
  \item Ilustración
  \item Romanticismo
  \item Realismo
  \item Naturalismo
  \item Barroco
  \item Neoclasicismo
\end{multicols}
\end{enumerate}
\vspace{0.15cm}
\item ¿Qué movimiento prefiere el artículo de costumbres?
\begin{enumerate}[label=\textbf{\Alph*)}]
\begin{multicols}{2}
  \item Romanticismo
  \item Realismo
  \item Naturalismo
  \item Barroco
  \item Neoclasicismo
  \item Ilustración
\end{multicols}
\end{enumerate}
\vspace{0.15cm}
\item ¿Qué movimiento prefiere la novela de tesis?
\begin{enumerate}[label=\textbf{\Alph*)}]
\begin{multicols}{2}
  \item Realismo
  \item Romanticismo
  \item Naturalismo
  \item Barroco
  \item Neoclasicismo
  \item Ilustración
\end{multicols}
\end{enumerate}
\vspace{0.15cm}
\item ¿Qué movimiento prefiere el drama histórico?
\begin{enumerate}[label=\textbf{\Alph*)}]
\begin{multicols}{2}
  \item Romanticismo
  \item Realismo
  \item Naturalismo
  \item Barroco
  \item Neoclasicismo
  \item Ilustración
\end{multicols}
\end{enumerate}
\vspace{0.15cm}
\item ¿Qué movimiento prefiere la comedia de costumbres?
\begin{enumerate}[label=\textbf{\Alph*)}]
\begin{multicols}{2}
  \item Realismo
  \item Romanticismo
  \item Naturalismo
  \item Barroco
  \item Neoclasicismo
  \item Ilustración
\end{multicols}
\end{enumerate}
\vspace{0.15cm}
\item ¿Qué movimiento prefiere la tragedia?
\begin{enumerate}[label=\textbf{\Alph*)}]
\begin{multicols}{2}
  \item Neoclasicismo
  \item Romanticismo
  \item Realismo
  \item Naturalismo
  \item Barroco
  \item Ilustración
\end{multicols}
\end{enumerate}
\vspace{0.15cm}
\item ¿Qué movimiento prefiere el melodrama?
\begin{enumerate}[label=\textbf{\Alph*)}]
\begin{multicols}{2}
  \item Romanticismo
  \item Realismo
  \item Naturalismo
  \item Barroco
  \item Neoclasicismo
  \item Ilustración
\end{multicols}
\end{enumerate}
\vspace{0.15cm}
\item ¿Qué movimiento prefiere el folletín?
\begin{enumerate}[label=\textbf{\Alph*)}]
\begin{multicols}{2}
  \item Romanticismo
  \item Realismo
  \item Naturalismo
  \item Barroco
  \item Neoclasicismo
  \item Ilustración
\end{multicols}
\end{enumerate}
\vspace{0.15cm}
\item ¿Qué movimiento prefiere la novela por entregas?
\begin{enumerate}[label=\textbf{\Alph*)}]
\begin{multicols}{2}
  \item Realismo
  \item Romanticismo
  \item Naturalismo
  \item Barroco
  \item Neoclasicismo
  \item Ilustración
\end{multicols}
\end{enumerate}
\vspace{0.15cm}
\item ¿Qué movimiento prefiere el cuento?
\begin{enumerate}[label=\textbf{\Alph*)}]
\begin{multicols}{2}
  \item Realismo
  \item Romanticismo
  \item Naturalismo
  \item Barroco
  \item Neoclasicismo
  \item Ilustración
\end{multicols}
\end{enumerate}
\vspace{0.15cm}
\item ¿Qué movimiento prefiere la leyenda?
\begin{enumerate}[label=\textbf{\Alph*)}]
\begin{multicols}{2}
  \item Romanticismo
  \item Realismo
  \item Naturalismo
  \item Barroco
  \item Neoclasicismo
  \item Ilustración
\end{multicols}
\end{enumerate}
\vspace{0.15cm}
\item ¿Qué movimiento prefiere la crónica?
\begin{enumerate}[label=\textbf{\Alph*)}]
\begin{multicols}{2}
  \item Realismo
  \item Romanticismo
  \item Naturalismo
  \item Barroco
  \item Neoclasicismo
  \item Ilustración
\end{multicols}
\end{enumerate}
\vspace{0.15cm}
\item ¿Qué movimiento prefiere la memoria?
\begin{enumerate}[label=\textbf{\Alph*)}]
\begin{multicols}{2}
  \item Romanticismo
  \item Realismo
  \item Naturalismo
  \item Barroco
  \item Neoclasicismo
  \item Ilustración
\end{multicols}
\end{enumerate}
\vspace{0.15cm}
\end{enumerate}\n
\subsubsection*{SECCIÓN VII: TEMAS Y CARACTERÍSTICAS (271–300)}

\begin{enumerate}[label=\textbf{\arabic*.}]
\setcounter{enumi}{270}
\item ¿Cuál es el tema central del Romanticismo?
\begin{enumerate}[label=\textbf{\Alph*)}]
\begin{multicols}{2}
  \item El yo y la libertad
  \item Siglo XVIII
  \item Siglo XVII
  \item Siglo XX
  \item Siglo XVI
  \item Siglo XXI
\end{multicols}
\end{enumerate}
\vspace{0.15cm}
\item ¿Cuál es el tema central del Realismo?
\begin{enumerate}[label=\textbf{\Alph*)}]
\begin{multicols}{2}
  \item La sociedad y la crítica
  \item Siglo XVIII
  \item Siglo XVII
  \item Siglo XX
  \item Siglo XVI
  \item Siglo XXI
\end{multicols}
\end{enumerate}
\vspace{0.15cm}
\item ¿Cuál es el tema central del Naturalismo?
\begin{enumerate}[label=\textbf{\Alph*)}]
\begin{multicols}{2}
  \item El determinismo y la herencia
  \item Siglo XVIII
  \item Siglo XVII
  \item Siglo XX
  \item Siglo XVI
  \item Siglo XXI
\end{multicols}
\end{enumerate}
\vspace{0.15cm}
\item ¿Qué movimiento exalta el amor trágico?
\begin{enumerate}[label=\textbf{\Alph*)}]
\begin{multicols}{2}
  \item Romanticismo
  \item Realismo
  \item Naturalismo
  \item Barroco
  \item Neoclasicismo
  \item Ilustración
\end{multicols}
\end{enumerate}
\vspace{0.15cm}
\item ¿Qué movimiento critica el matrimonio por conveniencia?
\begin{enumerate}[label=\textbf{\Alph*)}]
\begin{multicols}{2}
  \item Realismo
  \item Romanticismo
  \item Naturalismo
  \item Barroco
  \item Neoclasicismo
  \item Ilustración
\end{multicols}
\end{enumerate}
\vspace{0.15cm}
\item ¿Qué movimiento exalta la muerte?
\begin{enumerate}[label=\textbf{\Alph*)}]
\begin{multicols}{2}
  \item Romanticismo
  \item Realismo
  \item Naturalismo
  \item Barroco
  \item Neoclasicismo
  \item Ilustración
\end{multicols}
\end{enumerate}
\vspace{0.15cm}
\item ¿Qué movimiento critica la hipocresía social?
\begin{enumerate}[label=\textbf{\Alph*)}]
\begin{multicols}{2}
  \item Realismo
  \item Romanticismo
  \item Naturalismo
  \item Barroco
  \item Neoclasicismo
  \item Ilustración
\end{multicols}
\end{enumerate}
\vspace{0.15cm}
\item ¿Qué movimiento exalta la soledad?
\begin{enumerate}[label=\textbf{\Alph*)}]
\begin{multicols}{2}
  \item Romanticismo
  \item Realismo
  \item Naturalismo
  \item Barroco
  \item Neoclasicismo
  \item Ilustración
\end{multicols}
\end{enumerate}
\vspace{0.15cm}
\item ¿Qué movimiento retrata la vida de la burguesía?
\begin{enumerate}[label=\textbf{\Alph*)}]
\begin{multicols}{2}
  \item Realismo
  \item Romanticismo
  \item Naturalismo
  \item Barroco
  \item Neoclasicismo
  \item Ilustración
\end{multicols}
\end{enumerate}
\vspace{0.15cm}
\item ¿Qué movimiento retrata la vida del pueblo?
\begin{enumerate}[label=\textbf{\Alph*)}]
\begin{multicols}{2}
  \item Naturalismo
  \item Romanticismo
  \item Realismo
  \item Barroco
  \item Neoclasicismo
  \item Ilustración
\end{multicols}
\end{enumerate}
\vspace{0.15cm}
\item ¿Qué movimiento exalta la libertad política?
\begin{enumerate}[label=\textbf{\Alph*)}]
\begin{multicols}{2}
  \item Romanticismo
  \item Realismo
  \item Naturalismo
  \item Barroco
  \item Neoclasicismo
  \item Ilustración
\end{multicols}
\end{enumerate}
\vspace{0.15cm}
\item ¿Qué movimiento critica la desigualdad social?
\begin{enumerate}[label=\textbf{\Alph*)}]
\begin{multicols}{2}
  \item Realismo
  \item Romanticismo
  \item Naturalismo
  \item Barroco
  \item Neoclasicismo
  \item Ilustración
\end{multicols}
\end{enumerate}
\vspace{0.15cm}
\item ¿Qué movimiento exalta la imaginación?
\begin{enumerate}[label=\textbf{\Alph*)}]
\begin{multicols}{2}
  \item Romanticismo
  \item Realismo
  \item Naturalismo
  \item Barroco
  \item Neoclasicismo
  \item Ilustración
\end{multicols}
\end{enumerate}
\vspace{0.15cm}
\item ¿Qué movimiento exalta la observación?
\begin{enumerate}[label=\textbf{\Alph*)}]
\begin{multicols}{2}
  \item Realismo
  \item Romanticismo
  \item Naturalismo
  \item Barroco
  \item Neoclasicismo
  \item Ilustración
\end{multicols}
\end{enumerate}
\vspace{0.15cm}
\item ¿Qué movimiento exalta la razón?
\begin{enumerate}[label=\textbf{\Alph*)}]
\begin{multicols}{2}
  \item Ilustración
  \item Romanticismo
  \item Realismo
  \item Naturalismo
  \item Barroco
  \item Neoclasicismo
\end{multicols}
\end{enumerate}
\vspace{0.15cm}
\item ¿Qué movimiento exalta el sentimiento?
\begin{enumerate}[label=\textbf{\Alph*)}]
\begin{multicols}{2}
  \item Romanticismo
  \item Realismo
  \item Naturalismo
  \item Barroco
  \item Neoclasicismo
  \item Ilustración
\end{multicols}
\end{enumerate}
\vspace{0.15cm}
\item ¿Qué movimiento exalta la ciencia?
\begin{enumerate}[label=\textbf{\Alph*)}]
\begin{multicols}{2}
  \item Naturalismo
  \item Romanticismo
  \item Realismo
  \item Barroco
  \item Neoclasicismo
  \item Ilustración
\end{multicols}
\end{enumerate}
\vspace{0.15cm}
\item ¿Qué movimiento exalta la tradición?
\begin{enumerate}[label=\textbf{\Alph*)}]
\begin{multicols}{2}
  \item Neoclasicismo
  \item Romanticismo
  \item Realismo
  \item Naturalismo
  \item Barroco
  \item Ilustración
\end{multicols}
\end{enumerate}
\vspace{0.15cm}
\item ¿Qué movimiento exalta la originalidad?
\begin{enumerate}[label=\textbf{\Alph*)}]
\begin{multicols}{2}
  \item Romanticismo
  \item Realismo
  \item Naturalismo
  \item Barroco
  \item Neoclasicismo
  \item Ilustración
\end{multicols}
\end{enumerate}
\vspace{0.15cm}
\item ¿Qué movimiento exalta la imitación?
\begin{enumerate}[label=\textbf{\Alph*)}]
\begin{multicols}{2}
  \item Neoclasicismo
  \item Romanticismo
  \item Realismo
  \item Naturalismo
  \item Barroco
  \item Ilustración
\end{multicols}
\end{enumerate}
\vspace{0.15cm}
\item ¿Qué movimiento prefiere la unidad de acción, lugar y tiempo?
\begin{enumerate}[label=\textbf{\Alph*)}]
\begin{multicols}{2}
  \item Neoclasicismo
  \item Romanticismo
  \item Realismo
  \item Naturalismo
  \item Barroco
  \item Ilustración
\end{multicols}
\end{enumerate}
\vspace{0.15cm}
\item ¿Qué movimiento rompe las unidades?
\begin{enumerate}[label=\textbf{\Alph*)}]
\begin{multicols}{2}
  \item Romanticismo
  \item Realismo
  \item Naturalismo
  \item Barroco
  \item Neoclasicismo
  \item Ilustración
\end{multicols}
\end{enumerate}
\vspace{0.15cm}
\item ¿Qué movimiento prefiere la mezcla de géneros?
\begin{enumerate}[label=\textbf{\Alph*)}]
\begin{multicols}{2}
  \item Romanticismo
  \item Realismo
  \item Naturalismo
  \item Barroco
  \item Neoclasicismo
  \item Ilustración
\end{multicols}
\end{enumerate}
\vspace{0.15cm}
\item ¿Qué movimiento prefiere la pureza de géneros?
\begin{enumerate}[label=\textbf{\Alph*)}]
\begin{multicols}{2}
  \item Neoclasicismo
  \item Romanticismo
  \item Realismo
  \item Naturalismo
  \item Barroco
  \item Ilustración
\end{multicols}
\end{enumerate}
\vspace{0.15cm}
\item ¿Qué movimiento prefiere el lenguaje culto?
\begin{enumerate}[label=\textbf{\Alph*)}]
\begin{multicols}{2}
  \item Neoclasicismo
  \item Romanticismo
  \item Realismo
  \item Naturalismo
  \item Barroco
  \item Ilustración
\end{multicols}
\end{enumerate}
\vspace{0.15cm}
\item ¿Qué movimiento prefiere el lenguaje popular?
\begin{enumerate}[label=\textbf{\Alph*)}]
\begin{multicols}{2}
  \item Romanticismo
  \item Realismo
  \item Naturalismo
  \item Barroco
  \item Neoclasicismo
  \item Ilustración
\end{multicols}
\end{enumerate}
\vspace{0.15cm}
\item ¿Qué movimiento prefiere el lenguaje coloquial?
\begin{enumerate}[label=\textbf{\Alph*)}]
\begin{multicols}{2}
  \item Realismo
  \item Romanticismo
  \item Naturalismo
  \item Barroco
  \item Neoclasicismo
  \item Ilustración
\end{multicols}
\end{enumerate}
\vspace{0.15cm}
\item ¿Qué movimiento prefiere el lenguaje retórico?
\begin{enumerate}[label=\textbf{\Alph*)}]
\begin{multicols}{2}
  \item Barroco
  \item Romanticismo
  \item Realismo
  \item Naturalismo
  \item Neoclasicismo
  \item Ilustración
\end{multicols}
\end{enumerate}
\vspace{0.15cm}
\item ¿Qué movimiento prefiere el lenguaje sencillo?
\begin{enumerate}[label=\textbf{\Alph*)}]
\begin{multicols}{2}
  \item Neoclasicismo
  \item Romanticismo
  \item Realismo
  \item Naturalismo
  \item Barroco
  \item Ilustración
\end{multicols}
\end{enumerate}
\vspace{0.15cm}
\item ¿Qué movimiento prefiere el lenguaje directo?
\begin{enumerate}[label=\textbf{\Alph*)}]
\begin{multicols}{2}
  \item Realismo
  \item Romanticismo
  \item Naturalismo
  \item Barroco
  \item Neoclasicismo
  \item Ilustración
\end{multicols}
\end{enumerate}
\vspace{0.15cm}
\end{enumerate}\n
\vspace{0.4cm}
\begin{tcolorbox}[colback=gray!5!white,colframe=gray!60!black,title=\textbf{SOLUCIONARIO -- 3.2 Movimientos del Siglo XIX}]
\begin{multicols}{4}
\begin{enumerate}[label=\textbf{\arabic*.}]
  \item \textbf{A}
  \item \textbf{A}
  \item \textbf{A}
  \item \textbf{A}
  \item \textbf{A}
  \item \textbf{A}
  \item \textbf{A}
  \item \textbf{A}
  \item \textbf{A}
  \item \textbf{A}
  \item \textbf{A}
  \item \textbf{A}
  \item \textbf{A}
  \item \textbf{A}
  \item \textbf{A}
  \item \textbf{A}
  \item \textbf{A}
  \item \textbf{A}
  \item \textbf{A}
  \item \textbf{A}
  \item \textbf{A}
  \item \textbf{A}
  \item \textbf{A}
  \item \textbf{A}
  \item \textbf{A}
  \item \textbf{A}
  \item \textbf{A}
  \item \textbf{A}
  \item \textbf{A}
  \item \textbf{A}
  \item \textbf{A}
  \item \textbf{A}
  \item \textbf{A}
  \item \textbf{A}
  \item \textbf{A}
  \item \textbf{A}
  \item \textbf{A}
  \item \textbf{A}
  \item \textbf{A}
  \item \textbf{A}
  \item \textbf{A}
  \item \textbf{A}
  \item \textbf{A}
  \item \textbf{A}
  \item \textbf{A}
  \item \textbf{A}
  \item \textbf{A}
  \item \textbf{A}
  \item \textbf{A}
  \item \textbf{A}
  \item \textbf{A}
  \item \textbf{A}
  \item \textbf{A}
  \item \textbf{A}
  \item \textbf{A}
  \item \textbf{A}
  \item \textbf{A}
  \item \textbf{A}
  \item \textbf{A}
  \item \textbf{A}
  \item \textbf{A}
  \item \textbf{A}
  \item \textbf{A}
  \item \textbf{A}
  \item \textbf{A}
  \item \textbf{A}
  \item \textbf{A}
  \item \textbf{A}
  \item \textbf{A}
  \item \textbf{A}
  \item \textbf{A}
  \item \textbf{A}
  \item \textbf{A}
  \item \textbf{A}
  \item \textbf{A}
  \item \textbf{A}
  \item \textbf{A}
  \item \textbf{A}
  \item \textbf{A}
  \item \textbf{A}
  \item \textbf{A}
  \item \textbf{A}
  \item \textbf{A}
  \item \textbf{A}
  \item \textbf{A}
  \item \textbf{A}
  \item \textbf{A}
  \item \textbf{A}
  \item \textbf{A}
  \item \textbf{A}
  \item \textbf{A}
  \item \textbf{A}
  \item \textbf{A}
  \item \textbf{A}
  \item \textbf{A}
  \item \textbf{A}
  \item \textbf{A}
  \item \textbf{A}
  \item \textbf{A}
  \item \textbf{A}
  \item \textbf{A}
  \item \textbf{A}
  \item \textbf{A}
  \item \textbf{A}
  \item \textbf{A}
  \item \textbf{A}
  \item \textbf{A}
  \item \textbf{A}
  \item \textbf{A}
  \item \textbf{A}
  \item \textbf{A}
  \item \textbf{A}
  \item \textbf{A}
  \item \textbf{A}
  \item \textbf{A}
  \item \textbf{A}
  \item \textbf{A}
  \item \textbf{A}
  \item \textbf{A}
  \item \textbf{A}
  \item \textbf{A}
  \item \textbf{A}
  \item \textbf{A}
  \item \textbf{A}
  \item \textbf{A}
  \item \textbf{A}
  \item \textbf{A}
  \item \textbf{A}
  \item \textbf{A}
  \item \textbf{A}
  \item \textbf{A}
  \item \textbf{A}
  \item \textbf{A}
  \item \textbf{A}
  \item \textbf{A}
  \item \textbf{A}
  \item \textbf{A}
  \item \textbf{A}
  \item \textbf{A}
  \item \textbf{A}
  \item \textbf{A}
  \item \textbf{A}
  \item \textbf{A}
  \item \textbf{A}
  \item \textbf{A}
  \item \textbf{A}
  \item \textbf{A}
  \item \textbf{A}
  \item \textbf{A}
  \item \textbf{A}
  \item \textbf{A}
  \item \textbf{A}
  \item \textbf{A}
  \item \textbf{A}
  \item \textbf{A}
  \item \textbf{A}
  \item \textbf{A}
  \item \textbf{A}
  \item \textbf{A}
  \item \textbf{A}
  \item \textbf{A}
  \item \textbf{A}
  \item \textbf{A}
  \item \textbf{A}
  \item \textbf{A}
  \item \textbf{A}
  \item \textbf{A}
  \item \textbf{A}
  \item \textbf{A}
  \item \textbf{A}
  \item \textbf{A}
  \item \textbf{A}
  \item \textbf{A}
  \item \textbf{A}
  \item \textbf{A}
  \item \textbf{A}
  \item \textbf{A}
  \item \textbf{A}
  \item \textbf{A}
  \item \textbf{A}
  \item \textbf{A}
  \item \textbf{A}
  \item \textbf{A}
  \item \textbf{A}
  \item \textbf{A}
  \item \textbf{A}
  \item \textbf{A}
  \item \textbf{A}
  \item \textbf{A}
  \item \textbf{A}
  \item \textbf{A}
  \item \textbf{A}
  \item \textbf{A}
  \item \textbf{A}
  \item \textbf{A}
  \item \textbf{A}
  \item \textbf{A}
  \item \textbf{A}
  \item \textbf{A}
  \item \textbf{A}
  \item \textbf{A}
  \item \textbf{A}
  \item \textbf{A}
  \item \textbf{A}
  \item \textbf{A}
  \item \textbf{A}
  \item \textbf{A}
  \item \textbf{A}
  \item \textbf{A}
  \item \textbf{A}
  \item \textbf{A}
  \item \textbf{A}
  \item \textbf{A}
  \item \textbf{A}
  \item \textbf{A}
  \item \textbf{A}
  \item \textbf{A}
  \item \textbf{A}
  \item \textbf{A}
  \item \textbf{A}
  \item \textbf{A}
  \item \textbf{A}
  \item \textbf{A}
  \item \textbf{A}
  \item \textbf{A}
  \item \textbf{A}
  \item \textbf{A}
  \item \textbf{A}
  \item \textbf{A}
  \item \textbf{A}
  \item \textbf{A}
  \item \textbf{A}
  \item \textbf{A}
  \item \textbf{A}
  \item \textbf{A}
  \item \textbf{A}
  \item \textbf{A}
  \item \textbf{A}
  \item \textbf{A}
  \item \textbf{A}
  \item \textbf{A}
  \item \textbf{A}
  \item \textbf{A}
  \item \textbf{A}
  \item \textbf{A}
  \item \textbf{A}
  \item \textbf{A}
  \item \textbf{A}
  \item \textbf{A}
  \item \textbf{A}
  \item \textbf{A}
  \item \textbf{A}
  \item \textbf{A}
  \item \textbf{A}
  \item \textbf{A}
  \item \textbf{A}
  \item \textbf{A}
  \item \textbf{A}
  \item \textbf{A}
  \item \textbf{A}
  \item \textbf{A}
  \item \textbf{A}
  \item \textbf{A}
  \item \textbf{A}
  \item \textbf{A}
  \item \textbf{A}
  \item \textbf{A}
  \item \textbf{A}
  \item \textbf{A}
  \item \textbf{A}
  \item \textbf{A}
  \item \textbf{A}
  \item \textbf{A}
  \item \textbf{A}
  \item \textbf{A}
  \item \textbf{A}
  \item \textbf{A}
  \item \textbf{A}
  \item \textbf{A}
  \item \textbf{A}
  \item \textbf{A}
  \item \textbf{A}
  \item \textbf{A}
  \item \textbf{A}
  \item \textbf{A}
  \item \textbf{A}
  \item \textbf{A}
  \item \textbf{A}
  \item \textbf{A}
  \item \textbf{A}
  \item \textbf{A}
  \item \textbf{A}
  \item \textbf{A}
  \item \textbf{A}
  \item \textbf{A}
  \item \textbf{A}
  \item \textbf{A}
  \item \textbf{A}
  \item \textbf{A}
  \item \textbf{A}
\end{enumerate}
\end{multicols}
\end{tcolorbox}
\newpage\n
\subsection{3.3 Literatura del Siglo Contemporáneo}

\subsubsection*{SECCIÓN I: VANGUARDIAS ARTÍSTICAS – CONCEPTOS BÁSICOS (1–40)}

\begin{enumerate}[label=\textbf{\arabic*.}]
\setcounter{enumi}{0}
\item ¿Qué son las vanguardias artísticas?
\begin{enumerate}[label=\textbf{\Alph*)}]
\begin{multicols}{2}
  \item Movimientos artísticos de ruptura que surgen a principios del siglo XX
  \item Jorge Luis Borges
  \item Federico García Lorca
  \item Pablo Neruda
  \item Octavio Paz
  \item César Vallejo
\end{multicols}
\end{enumerate}
\vspace{0.15cm}
\item ¿En qué época surgen las vanguardias?
\begin{enumerate}[label=\textbf{\Alph*)}]
\begin{multicols}{2}
  \item A principios del siglo XX
  \item Jorge Luis Borges
  \item Federico García Lorca
  \item Pablo Neruda
  \item Octavio Paz
  \item César Vallejo
\end{multicols}
\end{enumerate}
\vspace{0.15cm}
\item ¿Cuál es la característica principal de las vanguardias?
\begin{enumerate}[label=\textbf{\Alph*)}]
\begin{multicols}{2}
  \item La ruptura con las normas y tradiciones artísticas
  \item Jorge Luis Borges
  \item Federico García Lorca
  \item Pablo Neruda
  \item Octavio Paz
  \item César Vallejo
\end{multicols}
\end{enumerate}
\vspace{0.15cm}
\item ¿Qué es el futurismo?
\begin{enumerate}[label=\textbf{\Alph*)}]
\begin{multicols}{2}
  \item Un movimiento vanguardista que exalta la velocidad y la tecnología
  \item Jorge Luis Borges
  \item Federico García Lorca
  \item Pablo Neruda
  \item Octavio Paz
  \item César Vallejo
\end{multicols}
\end{enumerate}
\vspace{0.15cm}
\item ¿Quién fundó el futurismo?
\begin{enumerate}[label=\textbf{\Alph*)}]
\begin{multicols}{2}
  \item Filippo Tommaso Marinetti
  \item Jorge Luis Borges
  \item Federico García Lorca
  \item Pablo Neruda
  \item Octavio Paz
  \item César Vallejo
\end{multicols}
\end{enumerate}
\vspace{0.15cm}
\item ¿Qué es el surrealismo?
\begin{enumerate}[label=\textbf{\Alph*)}]
\begin{multicols}{2}
  \item Un movimiento vanguardista que explora el inconsciente y los sueños
  \item Jorge Luis Borges
  \item Federico García Lorca
  \item Pablo Neruda
  \item Octavio Paz
  \item César Vallejo
\end{multicols}
\end{enumerate}
\vspace{0.15cm}
\item ¿Quién fundó el surrealismo?
\begin{enumerate}[label=\textbf{\Alph*)}]
\begin{multicols}{2}
  \item André Breton
  \item Jorge Luis Borges
  \item Federico García Lorca
  \item Pablo Neruda
  \item Octavio Paz
  \item César Vallejo
\end{multicols}
\end{enumerate}
\vspace{0.15cm}
\item ¿Qué es el dadaísmo?
\begin{enumerate}[label=\textbf{\Alph*)}]
\begin{multicols}{2}
  \item Un movimiento vanguardista que rechaza la lógica y las normas
  \item Jorge Luis Borges
  \item Federico García Lorca
  \item Pablo Neruda
  \item Octavio Paz
  \item César Vallejo
\end{multicols}
\end{enumerate}
\vspace{0.15cm}
\item ¿Quién fundó el dadaísmo?
\begin{enumerate}[label=\textbf{\Alph*)}]
\begin{multicols}{2}
  \item Tristan Tzara
  \item Jorge Luis Borges
  \item Federico García Lorca
  \item Pablo Neruda
  \item Octavio Paz
  \item César Vallejo
\end{multicols}
\end{enumerate}
\vspace{0.15cm}
\item ¿Qué es el cubismo?
\begin{enumerate}[label=\textbf{\Alph*)}]
\begin{multicols}{2}
  \item Un movimiento vanguardista que descompone la realidad en formas geométricas
  \item Jorge Luis Borges
  \item Federico García Lorca
  \item Pablo Neruda
  \item Octavio Paz
  \item César Vallejo
\end{multicols}
\end{enumerate}
\vspace{0.15cm}
\item ¿Quiénes son los principales representantes del cubismo?
\begin{enumerate}[label=\textbf{\Alph*)}]
\begin{multicols}{2}
  \item Pablo Picasso y Georges Braque
  \item Jorge Luis Borges
  \item Federico García Lorca
  \item Pablo Neruda
  \item Octavio Paz
  \item César Vallejo
\end{multicols}
\end{enumerate}
\vspace{0.15cm}
\item ¿Qué es el expresionismo?
\begin{enumerate}[label=\textbf{\Alph*)}]
\begin{multicols}{2}
  \item Un movimiento vanguardista que exalta la expresión subjetiva y emocional
  \item Jorge Luis Borges
  \item Federico García Lorca
  \item Pablo Neruda
  \item Octavio Paz
  \item César Vallejo
\end{multicols}
\end{enumerate}
\vspace{0.15cm}
\item ¿Qué es el ultraísmo?
\begin{enumerate}[label=\textbf{\Alph*)}]
\begin{multicols}{2}
  \item Un movimiento vanguardista español que busca la renovación poética
  \item Jorge Luis Borges
  \item Federico García Lorca
  \item Pablo Neruda
  \item Octavio Paz
  \item César Vallejo
\end{multicols}
\end{enumerate}
\vspace{0.15cm}
\item ¿Quién es el principal representante del ultraísmo?
\begin{enumerate}[label=\textbf{\Alph*)}]
\begin{multicols}{2}
  \item Jorge Luis Borges
  \item Federico García Lorca
  \item Pablo Neruda
  \item Octavio Paz
  \item César Vallejo
  \item Vicente Huidobro
\end{multicols}
\end{enumerate}
\vspace{0.15cm}
\item ¿Qué es el creacionismo?
\begin{enumerate}[label=\textbf{\Alph*)}]
\begin{multicols}{2}
  \item Un movimiento vanguardista que propone crear realidades nuevas en la poesía
  \item Jorge Luis Borges
  \item Federico García Lorca
  \item Pablo Neruda
  \item Octavio Paz
  \item César Vallejo
\end{multicols}
\end{enumerate}
\vspace{0.15cm}
\item ¿Quién es el principal representante del creacionismo?
\begin{enumerate}[label=\textbf{\Alph*)}]
\begin{multicols}{2}
  \item Vicente Huidobro
  \item Jorge Luis Borges
  \item Federico García Lorca
  \item Pablo Neruda
  \item Octavio Paz
  \item César Vallejo
\end{multicols}
\end{enumerate}
\vspace{0.15cm}
\item ¿Qué es el estridentismo?
\begin{enumerate}[label=\textbf{\Alph*)}]
\begin{multicols}{2}
  \item Un movimiento vanguardista mexicano
  \item Jorge Luis Borges
  \item Federico García Lorca
  \item Pablo Neruda
  \item Octavio Paz
  \item César Vallejo
\end{multicols}
\end{enumerate}
\vspace{0.15cm}
\item ¿Qué es el indigenismo?
\begin{enumerate}[label=\textbf{\Alph*)}]
\begin{multicols}{2}
  \item Un movimiento literario que reivindica la cultura indígena
  \item Jorge Luis Borges
  \item Federico García Lorca
  \item Pablo Neruda
  \item Octavio Paz
  \item César Vallejo
\end{multicols}
\end{enumerate}
\vspace{0.15cm}
\item ¿Qué es el negrismo?
\begin{enumerate}[label=\textbf{\Alph*)}]
\begin{multicols}{2}
  \item Un movimiento literario que reivindica la cultura afrocaribeña
  \item Jorge Luis Borges
  \item Federico García Lorca
  \item Pablo Neruda
  \item Octavio Paz
  \item César Vallejo
\end{multicols}
\end{enumerate}
\vspace{0.15cm}
\item ¿Qué es el realismo mágico?
\begin{enumerate}[label=\textbf{\Alph*)}]
\begin{multicols}{2}
  \item Un movimiento literario que mezcla lo real y lo fantástico
  \item Jorge Luis Borges
  \item Federico García Lorca
  \item Pablo Neruda
  \item Octavio Paz
  \item César Vallejo
\end{multicols}
\end{enumerate}
\vspace{0.15cm}
\item ¿Quién es el principal representante del realismo mágico?
\begin{enumerate}[label=\textbf{\Alph*)}]
\begin{multicols}{2}
  \item Gabriel García Márquez
  \item Jorge Luis Borges
  \item Federico García Lorca
  \item Pablo Neruda
  \item Octavio Paz
  \item César Vallejo
\end{multicols}
\end{enumerate}
\vspace{0.15cm}
\item ¿Qué es el existencialismo?
\begin{enumerate}[label=\textbf{\Alph*)}]
\begin{multicols}{2}
  \item Una corriente filosófica y literaria que reflexiona sobre la existencia humana
  \item Jorge Luis Borges
  \item Federico García Lorca
  \item Pablo Neruda
  \item Octavio Paz
  \item César Vallejo
\end{multicols}
\end{enumerate}
\vspace{0.15cm}
\item ¿Quiénes son los principales representantes del existencialismo?
\begin{enumerate}[label=\textbf{\Alph*)}]
\begin{multicols}{2}
  \item Jean-Paul Sartre y Albert Camus
  \item Jorge Luis Borges
  \item Federico García Lorca
  \item Pablo Neruda
  \item Octavio Paz
  \item César Vallejo
\end{multicols}
\end{enumerate}
\vspace{0.15cm}
\item ¿Qué es el absurdo en la literatura?
\begin{enumerate}[label=\textbf{\Alph*)}]
\begin{multicols}{2}
  \item La representación de la falta de sentido de la existencia
  \item Jorge Luis Borges
  \item Federico García Lorca
  \item Pablo Neruda
  \item Octavio Paz
  \item César Vallejo
\end{multicols}
\end{enumerate}
\vspace{0.15cm}
\item ¿Quién es el autor de El extranjero?
\begin{enumerate}[label=\textbf{\Alph*)}]
\begin{multicols}{2}
  \item Albert Camus
  \item Jorge Luis Borges
  \item Federico García Lorca
  \item Pablo Neruda
  \item Octavio Paz
  \item César Vallejo
\end{multicols}
\end{enumerate}
\vspace{0.15cm}
\item ¿Quién es el autor de La metamorfosis?
\begin{enumerate}[label=\textbf{\Alph*)}]
\begin{multicols}{2}
  \item Franz Kafka
  \item Jorge Luis Borges
  \item Federico García Lorca
  \item Pablo Neruda
  \item Octavio Paz
  \item César Vallejo
\end{multicols}
\end{enumerate}
\vspace{0.15cm}
\item ¿Quién es el autor de En busca del tiempo perdido?
\begin{enumerate}[label=\textbf{\Alph*)}]
\begin{multicols}{2}
  \item Marcel Proust
  \item Jorge Luis Borges
  \item Federico García Lorca
  \item Pablo Neruda
  \item Octavio Paz
  \item César Vallejo
\end{multicols}
\end{enumerate}
\vspace{0.15cm}
\item ¿Quién es el autor de Ulises?
\begin{enumerate}[label=\textbf{\Alph*)}]
\begin{multicols}{2}
  \item James Joyce
  \item Jorge Luis Borges
  \item Federico García Lorca
  \item Pablo Neruda
  \item Octavio Paz
  \item César Vallejo
\end{multicols}
\end{enumerate}
\vspace{0.15cm}
\item ¿Quién es la autora de Mrs. Dalloway?
\begin{enumerate}[label=\textbf{\Alph*)}]
\begin{multicols}{2}
  \item Virginia Woolf
  \item Jorge Luis Borges
  \item Federico García Lorca
  \item Pablo Neruda
  \item Octavio Paz
  \item César Vallejo
\end{multicols}
\end{enumerate}
\vspace{0.15cm}
\item ¿Qué es el monólogo interior?
\begin{enumerate}[label=\textbf{\Alph*)}]
\begin{multicols}{2}
  \item Una técnica narrativa que reproduce los pensamientos del personaje
  \item Jorge Luis Borges
  \item Federico García Lorca
  \item Pablo Neruda
  \item Octavio Paz
  \item César Vallejo
\end{multicols}
\end{enumerate}
\vspace{0.15cm}
\item ¿Qué es el stream of consciousness?
\begin{enumerate}[label=\textbf{\Alph*)}]
\begin{multicols}{2}
  \item Una técnica narrativa que reproduce el flujo de conciencia del personaje
  \item Jorge Luis Borges
  \item Federico García Lorca
  \item Pablo Neruda
  \item Octavio Paz
  \item César Vallejo
\end{multicols}
\end{enumerate}
\vspace{0.15cm}
\item ¿Qué es la novela experimental?
\begin{enumerate}[label=\textbf{\Alph*)}]
\begin{multicols}{2}
  \item Una novela que rompe con las formas narrativas tradicionales
  \item Jorge Luis Borges
  \item Federico García Lorca
  \item Pablo Neruda
  \item Octavio Paz
  \item César Vallejo
\end{multicols}
\end{enumerate}
\vspace{0.15cm}
\item ¿Qué es el teatro del absurdo?
\begin{enumerate}[label=\textbf{\Alph*)}]
\begin{multicols}{2}
  \item Un género teatral que representa la falta de sentido de la existencia
  \item Jorge Luis Borges
  \item Federico García Lorca
  \item Pablo Neruda
  \item Octavio Paz
  \item César Vallejo
\end{multicols}
\end{enumerate}
\vspace{0.15cm}
\item ¿Quién es el autor de Esperando a Godot?
\begin{enumerate}[label=\textbf{\Alph*)}]
\begin{multicols}{2}
  \item Samuel Beckett
  \item Jorge Luis Borges
  \item Federico García Lorca
  \item Pablo Neruda
  \item Octavio Paz
  \item César Vallejo
\end{multicols}
\end{enumerate}
\vspace{0.15cm}
\item ¿Quién es el autor de La cantante calva?
\begin{enumerate}[label=\textbf{\Alph*)}]
\begin{multicols}{2}
  \item Eugène Ionesco
  \item Jorge Luis Borges
  \item Federico García Lorca
  \item Pablo Neruda
  \item Octavio Paz
  \item César Vallejo
\end{multicols}
\end{enumerate}
\vspace{0.15cm}
\item ¿Quién es el autor de Seis personajes en busca de autor?
\begin{enumerate}[label=\textbf{\Alph*)}]
\begin{multicols}{2}
  \item Luigi Pirandello
  \item Jorge Luis Borges
  \item Federico García Lorca
  \item Pablo Neruda
  \item Octavio Paz
  \item César Vallejo
\end{multicols}
\end{enumerate}
\vspace{0.15cm}
\item ¿Quién es el autor de Madre Coraje?
\begin{enumerate}[label=\textbf{\Alph*)}]
\begin{multicols}{2}
  \item Bertolt Brecht
  \item Jorge Luis Borges
  \item Federico García Lorca
  \item Pablo Neruda
  \item Octavio Paz
  \item César Vallejo
\end{multicols}
\end{enumerate}
\vspace{0.15cm}
\item ¿Qué es el teatro épico?
\begin{enumerate}[label=\textbf{\Alph*)}]
\begin{multicols}{2}
  \item Un teatro que busca la reflexión crítica del espectador
  \item Jorge Luis Borges
  \item Federico García Lorca
  \item Pablo Neruda
  \item Octavio Paz
  \item César Vallejo
\end{multicols}
\end{enumerate}
\vspace{0.15cm}
\item ¿Qué es la novela existencialista?
\begin{enumerate}[label=\textbf{\Alph*)}]
\begin{multicols}{2}
  \item Una novela que reflexiona sobre la existencia humana
  \item Jorge Luis Borges
  \item Federico García Lorca
  \item Pablo Neruda
  \item Octavio Paz
  \item César Vallejo
\end{multicols}
\end{enumerate}
\vspace{0.15cm}
\item ¿Qué es la literatura del absurdo?
\begin{enumerate}[label=\textbf{\Alph*)}]
\begin{multicols}{2}
  \item Una literatura que representa la falta de sentido de la existencia
  \item Jorge Luis Borges
  \item Federico García Lorca
  \item Pablo Neruda
  \item Octavio Paz
  \item César Vallejo
\end{multicols}
\end{enumerate}
\vspace{0.15cm}
\end{enumerate}\n
\subsubsection*{SECCIÓN II: VANGUARDIAS – AUTORES Y OBRAS (41–90)}

\begin{enumerate}[label=\textbf{\arabic*.}]
\setcounter{enumi}{40}
\item ¿Quién es el autor de Manifiesto futurista?
\begin{enumerate}[label=\textbf{\Alph*)}]
\begin{multicols}{2}
  \item Filippo Tommaso Marinetti
  \item Jorge Luis Borges
  \item Federico García Lorca
  \item Pablo Neruda
  \item Octavio Paz
  \item César Vallejo
\end{multicols}
\end{enumerate}
\vspace{0.15cm}
\item ¿Quién es el autor de Manifiesto surrealista?
\begin{enumerate}[label=\textbf{\Alph*)}]
\begin{multicols}{2}
  \item André Breton
  \item Jorge Luis Borges
  \item Federico García Lorca
  \item Pablo Neruda
  \item Octavio Paz
  \item César Vallejo
\end{multicols}
\end{enumerate}
\vspace{0.15cm}
\item ¿Quién es el autor de Caligramas?
\begin{enumerate}[label=\textbf{\Alph*)}]
\begin{multicols}{2}
  \item Guillaume Apollinaire
  \item Jorge Luis Borges
  \item Federico García Lorca
  \item Pablo Neruda
  \item Octavio Paz
  \item César Vallejo
\end{multicols}
\end{enumerate}
\vspace{0.15cm}
\item ¿Quién es el autor de La persistencia de la memoria?
\begin{enumerate}[label=\textbf{\Alph*)}]
\begin{multicols}{2}
  \item Salvador Dalí
  \item Jorge Luis Borges
  \item Federico García Lorca
  \item Pablo Neruda
  \item Octavio Paz
  \item César Vallejo
\end{multicols}
\end{enumerate}
\vspace{0.15cm}
\item ¿Quién es el autor de Guernica?
\begin{enumerate}[label=\textbf{\Alph*)}]
\begin{multicols}{2}
  \item Pablo Picasso
  \item Jorge Luis Borges
  \item Federico García Lorca
  \item Pablo Neruda
  \item Octavio Paz
  \item César Vallejo
\end{multicols}
\end{enumerate}
\vspace{0.15cm}
\item ¿Quién es el autor de Altazor?
\begin{enumerate}[label=\textbf{\Alph*)}]
\begin{multicols}{2}
  \item Vicente Huidobro
  \item Jorge Luis Borges
  \item Federico García Lorca
  \item Pablo Neruda
  \item Octavio Paz
  \item César Vallejo
\end{multicols}
\end{enumerate}
\vspace{0.15cm}
\item ¿Quién es el autor de Fervor de Buenos Aires?
\begin{enumerate}[label=\textbf{\Alph*)}]
\begin{multicols}{2}
  \item Jorge Luis Borges
  \item Federico García Lorca
  \item Pablo Neruda
  \item Octavio Paz
  \item César Vallejo
  \item Vicente Huidobro
\end{multicols}
\end{enumerate}
\vspace{0.15cm}
\item ¿Quién es el autor de Poeta en Nueva York?
\begin{enumerate}[label=\textbf{\Alph*)}]
\begin{multicols}{2}
  \item Federico García Lorca
  \item Jorge Luis Borges
  \item Pablo Neruda
  \item Octavio Paz
  \item César Vallejo
  \item Vicente Huidobro
\end{multicols}
\end{enumerate}
\vspace{0.15cm}
\item ¿Quién es el autor de Veinte poemas de amor y una canción desesperada?
\begin{enumerate}[label=\textbf{\Alph*)}]
\begin{multicols}{2}
  \item Pablo Neruda
  \item Jorge Luis Borges
  \item Federico García Lorca
  \item Octavio Paz
  \item César Vallejo
  \item Vicente Huidobro
\end{multicols}
\end{enumerate}
\vspace{0.15cm}
\item ¿Quién es el autor de Trilce?
\begin{enumerate}[label=\textbf{\Alph*)}]
\begin{multicols}{2}
  \item César Vallejo
  \item Jorge Luis Borges
  \item Federico García Lorca
  \item Pablo Neruda
  \item Octavio Paz
  \item Vicente Huidobro
\end{multicols}
\end{enumerate}
\vspace{0.15cm}
\item ¿Quién es el autor de Libertad bajo palabra?
\begin{enumerate}[label=\textbf{\Alph*)}]
\begin{multicols}{2}
  \item Octavio Paz
  \item Jorge Luis Borges
  \item Federico García Lorca
  \item Pablo Neruda
  \item César Vallejo
  \item Vicente Huidobro
\end{multicols}
\end{enumerate}
\vspace{0.15cm}
\item ¿Quién es el autor de Romancero gitano?
\begin{enumerate}[label=\textbf{\Alph*)}]
\begin{multicols}{2}
  \item Federico García Lorca
  \item Jorge Luis Borges
  \item Pablo Neruda
  \item Octavio Paz
  \item César Vallejo
  \item Vicente Huidobro
\end{multicols}
\end{enumerate}
\vspace{0.15cm}
\item ¿Quién es el autor de La casa de Bernarda Alba?
\begin{enumerate}[label=\textbf{\Alph*)}]
\begin{multicols}{2}
  \item Federico García Lorca
  \item Jorge Luis Borges
  \item Pablo Neruda
  \item Octavio Paz
  \item César Vallejo
  \item Vicente Huidobro
\end{multicols}
\end{enumerate}
\vspace{0.15cm}
\item ¿Quién es el autor de Bodas de sangre?
\begin{enumerate}[label=\textbf{\Alph*)}]
\begin{multicols}{2}
  \item Federico García Lorca
  \item Jorge Luis Borges
  \item Pablo Neruda
  \item Octavio Paz
  \item César Vallejo
  \item Vicente Huidobro
\end{multicols}
\end{enumerate}
\vspace{0.15cm}
\item ¿Quién es el autor de Yerma?
\begin{enumerate}[label=\textbf{\Alph*)}]
\begin{multicols}{2}
  \item Federico García Lorca
  \item Jorge Luis Borges
  \item Pablo Neruda
  \item Octavio Paz
  \item César Vallejo
  \item Vicente Huidobro
\end{multicols}
\end{enumerate}
\vspace{0.15cm}
\item ¿Quién es el autor de Ficciones?
\begin{enumerate}[label=\textbf{\Alph*)}]
\begin{multicols}{2}
  \item Jorge Luis Borges
  \item Federico García Lorca
  \item Pablo Neruda
  \item Octavio Paz
  \item César Vallejo
  \item Vicente Huidobro
\end{multicols}
\end{enumerate}
\vspace{0.15cm}
\item ¿Quién es el autor de El Aleph?
\begin{enumerate}[label=\textbf{\Alph*)}]
\begin{multicols}{2}
  \item Jorge Luis Borges
  \item Federico García Lorca
  \item Pablo Neruda
  \item Octavio Paz
  \item César Vallejo
  \item Vicente Huidobro
\end{multicols}
\end{enumerate}
\vspace{0.15cm}
\item ¿Quién es el autor de Historia universal de la infamia?
\begin{enumerate}[label=\textbf{\Alph*)}]
\begin{multicols}{2}
  \item Jorge Luis Borges
  \item Federico García Lorca
  \item Pablo Neruda
  \item Octavio Paz
  \item César Vallejo
  \item Vicente Huidobro
\end{multicols}
\end{enumerate}
\vspace{0.15cm}
\item ¿Quién es el autor de El túnel?
\begin{enumerate}[label=\textbf{\Alph*)}]
\begin{multicols}{2}
  \item Ernesto Sábato
  \item Jorge Luis Borges
  \item Federico García Lorca
  \item Pablo Neruda
  \item Octavio Paz
  \item César Vallejo
\end{multicols}
\end{enumerate}
\vspace{0.15cm}
\item ¿Quién es el autor de Sobre héroes y tumbas?
\begin{enumerate}[label=\textbf{\Alph*)}]
\begin{multicols}{2}
  \item Ernesto Sábato
  \item Jorge Luis Borges
  \item Federico García Lorca
  \item Pablo Neruda
  \item Octavio Paz
  \item César Vallejo
\end{multicols}
\end{enumerate}
\vspace{0.15cm}
\item ¿Quién es el autor de Rayuela?
\begin{enumerate}[label=\textbf{\Alph*)}]
\begin{multicols}{2}
  \item Julio Cortázar
  \item Jorge Luis Borges
  \item Federico García Lorca
  \item Pablo Neruda
  \item Octavio Paz
  \item César Vallejo
\end{multicols}
\end{enumerate}
\vspace{0.15cm}
\item ¿Quién es el autor de Bestiario?
\begin{enumerate}[label=\textbf{\Alph*)}]
\begin{multicols}{2}
  \item Julio Cortázar
  \item Jorge Luis Borges
  \item Federico García Lorca
  \item Pablo Neruda
  \item Octavio Paz
  \item César Vallejo
\end{multicols}
\end{enumerate}
\vspace{0.15cm}
\item ¿Quién es el autor de Cien años de soledad?
\begin{enumerate}[label=\textbf{\Alph*)}]
\begin{multicols}{2}
  \item Gabriel García Márquez
  \item Jorge Luis Borges
  \item Federico García Lorca
  \item Pablo Neruda
  \item Octavio Paz
  \item César Vallejo
\end{multicols}
\end{enumerate}
\vspace{0.15cm}
\item ¿Quién es el autor de El amor en los tiempos del cólera?
\begin{enumerate}[label=\textbf{\Alph*)}]
\begin{multicols}{2}
  \item Gabriel García Márquez
  \item Jorge Luis Borges
  \item Federico García Lorca
  \item Pablo Neruda
  \item Octavio Paz
  \item César Vallejo
\end{multicols}
\end{enumerate}
\vspace{0.15cm}
\item ¿Quién es el autor de Crónica de una muerte anunciada?
\begin{enumerate}[label=\textbf{\Alph*)}]
\begin{multicols}{2}
  \item Gabriel García Márquez
  \item Jorge Luis Borges
  \item Federico García Lorca
  \item Pablo Neruda
  \item Octavio Paz
  \item César Vallejo
\end{multicols}
\end{enumerate}
\vspace{0.15cm}
\item ¿Quién es el autor de La ciudad y los perros?
\begin{enumerate}[label=\textbf{\Alph*)}]
\begin{multicols}{2}
  \item Mario Vargas Llosa
  \item Jorge Luis Borges
  \item Federico García Lorca
  \item Pablo Neruda
  \item Octavio Paz
  \item César Vallejo
\end{multicols}
\end{enumerate}
\vspace{0.15cm}
\item ¿Quién es el autor de La casa verde?
\begin{enumerate}[label=\textbf{\Alph*)}]
\begin{multicols}{2}
  \item Mario Vargas Llosa
  \item Jorge Luis Borges
  \item Federico García Lorca
  \item Pablo Neruda
  \item Octavio Paz
  \item César Vallejo
\end{multicols}
\end{enumerate}
\vspace{0.15cm}
\item ¿Quién es el autor de Conversación en La Catedral?
\begin{enumerate}[label=\textbf{\Alph*)}]
\begin{multicols}{2}
  \item Mario Vargas Llosa
  \item Jorge Luis Borges
  \item Federico García Lorca
  \item Pablo Neruda
  \item Octavio Paz
  \item César Vallejo
\end{multicols}
\end{enumerate}
\vspace{0.15cm}
\item ¿Quién es el autor de La muerte de Artemio Cruz?
\begin{enumerate}[label=\textbf{\Alph*)}]
\begin{multicols}{2}
  \item Carlos Fuentes
  \item Jorge Luis Borges
  \item Federico García Lorca
  \item Pablo Neruda
  \item Octavio Paz
  \item César Vallejo
\end{multicols}
\end{enumerate}
\vspace{0.15cm}
\item ¿Quién es el autor de Aura?
\begin{enumerate}[label=\textbf{\Alph*)}]
\begin{multicols}{2}
  \item Carlos Fuentes
  \item Jorge Luis Borges
  \item Federico García Lorca
  \item Pablo Neruda
  \item Octavio Paz
  \item César Vallejo
\end{multicols}
\end{enumerate}
\vspace{0.15cm}
\item ¿Quién es el autor de Pedro Páramo?
\begin{enumerate}[label=\textbf{\Alph*)}]
\begin{multicols}{2}
  \item Juan Rulfo
  \item Jorge Luis Borges
  \item Federico García Lorca
  \item Pablo Neruda
  \item Octavio Paz
  \item César Vallejo
\end{multicols}
\end{enumerate}
\vspace{0.15cm}
\item ¿Quién es el autor de El llano en llamas?
\begin{enumerate}[label=\textbf{\Alph*)}]
\begin{multicols}{2}
  \item Juan Rulfo
  \item Jorge Luis Borges
  \item Federico García Lorca
  \item Pablo Neruda
  \item Octavio Paz
  \item César Vallejo
\end{multicols}
\end{enumerate}
\vspace{0.15cm}
\item ¿Quién es el autor de Paradiso?
\begin{enumerate}[label=\textbf{\Alph*)}]
\begin{multicols}{2}
  \item José Lezama Lima
  \item Jorge Luis Borges
  \item Federico García Lorca
  \item Pablo Neruda
  \item Octavio Paz
  \item César Vallejo
\end{multicols}
\end{enumerate}
\vspace{0.15cm}
\item ¿Quién es el autor de Tres tristes tigres?
\begin{enumerate}[label=\textbf{\Alph*)}]
\begin{multicols}{2}
  \item Guillermo Cabrera Infante
  \item Jorge Luis Borges
  \item Federico García Lorca
  \item Pablo Neruda
  \item Octavio Paz
  \item César Vallejo
\end{multicols}
\end{enumerate}
\vspace{0.15cm}
\item ¿Quién es el autor de La vida breve?
\begin{enumerate}[label=\textbf{\Alph*)}]
\begin{multicols}{2}
  \item Juan Carlos Onetti
  \item Jorge Luis Borges
  \item Federico García Lorca
  \item Pablo Neruda
  \item Octavio Paz
  \item César Vallejo
\end{multicols}
\end{enumerate}
\vspace{0.15cm}
\item ¿Quién es el autor de El astillero?
\begin{enumerate}[label=\textbf{\Alph*)}]
\begin{multicols}{2}
  \item Juan Carlos Onetti
  \item Jorge Luis Borges
  \item Federico García Lorca
  \item Pablo Neruda
  \item Octavio Paz
  \item César Vallejo
\end{multicols}
\end{enumerate}
\vspace{0.15cm}
\item ¿Quién es el autor de Juntacadáveres?
\begin{enumerate}[label=\textbf{\Alph*)}]
\begin{multicols}{2}
  \item Juan Carlos Onetti
  \item Jorge Luis Borges
  \item Federico García Lorca
  \item Pablo Neruda
  \item Octavio Paz
  \item César Vallejo
\end{multicols}
\end{enumerate}
\vspace{0.15cm}
\item ¿Quién es el autor de La invención de Morel?
\begin{enumerate}[label=\textbf{\Alph*)}]
\begin{multicols}{2}
  \item Adolfo Bioy Casares
  \item Jorge Luis Borges
  \item Federico García Lorca
  \item Pablo Neruda
  \item Octavio Paz
  \item César Vallejo
\end{multicols}
\end{enumerate}
\vspace{0.15cm}
\item ¿Quién es el autor de El sueño de los héroes?
\begin{enumerate}[label=\textbf{\Alph*)}]
\begin{multicols}{2}
  \item Adolfo Bioy Casares
  \item Jorge Luis Borges
  \item Federico García Lorca
  \item Pablo Neruda
  \item Octavio Paz
  \item César Vallejo
\end{multicols}
\end{enumerate}
\vspace{0.15cm}
\item ¿Quién es el autor de La tregua?
\begin{enumerate}[label=\textbf{\Alph*)}]
\begin{multicols}{2}
  \item Mario Benedetti
  \item Jorge Luis Borges
  \item Federico García Lorca
  \item Pablo Neruda
  \item Octavio Paz
  \item César Vallejo
\end{multicols}
\end{enumerate}
\vspace{0.15cm}
\item ¿Quién es el autor de Gracias por el fuego?
\begin{enumerate}[label=\textbf{\Alph*)}]
\begin{multicols}{2}
  \item Mario Benedetti
  \item Jorge Luis Borges
  \item Federico García Lorca
  \item Pablo Neruda
  \item Octavio Paz
  \item César Vallejo
\end{multicols}
\end{enumerate}
\vspace{0.15cm}
\item ¿Quién es el autor de El otoño del patriarca?
\begin{enumerate}[label=\textbf{\Alph*)}]
\begin{multicols}{2}
  \item Gabriel García Márquez
  \item Jorge Luis Borges
  \item Federico García Lorca
  \item Pablo Neruda
  \item Octavio Paz
  \item César Vallejo
\end{multicols}
\end{enumerate}
\vspace{0.15cm}
\item ¿Quién es el autor de El coronel no tiene quien le escriba?
\begin{enumerate}[label=\textbf{\Alph*)}]
\begin{multicols}{2}
  \item Gabriel García Márquez
  \item Jorge Luis Borges
  \item Federico García Lorca
  \item Pablo Neruda
  \item Octavio Paz
  \item César Vallejo
\end{multicols}
\end{enumerate}
\vspace{0.15cm}
\item ¿Quién es el autor de Los funerales de la Mamá Grande?
\begin{enumerate}[label=\textbf{\Alph*)}]
\begin{multicols}{2}
  \item Gabriel García Márquez
  \item Jorge Luis Borges
  \item Federico García Lorca
  \item Pablo Neruda
  \item Octavio Paz
  \item César Vallejo
\end{multicols}
\end{enumerate}
\vspace{0.15cm}
\item ¿Quién es el autor de La increíble y triste historia de la cándida Eréndira?
\begin{enumerate}[label=\textbf{\Alph*)}]
\begin{multicols}{2}
  \item Gabriel García Márquez
  \item Jorge Luis Borges
  \item Federico García Lorca
  \item Pablo Neruda
  \item Octavio Paz
  \item César Vallejo
\end{multicols}
\end{enumerate}
\vspace{0.15cm}
\item ¿Quién es el autor de Del amor y otros demonios?
\begin{enumerate}[label=\textbf{\Alph*)}]
\begin{multicols}{2}
  \item Gabriel García Márquez
  \item Jorge Luis Borges
  \item Federico García Lorca
  \item Pablo Neruda
  \item Octavio Paz
  \item César Vallejo
\end{multicols}
\end{enumerate}
\vspace{0.15cm}
\item ¿Quién es el autor de Memoria de mis putas tristes?
\begin{enumerate}[label=\textbf{\Alph*)}]
\begin{multicols}{2}
  \item Gabriel García Márquez
  \item Jorge Luis Borges
  \item Federico García Lorca
  \item Pablo Neruda
  \item Octavio Paz
  \item César Vallejo
\end{multicols}
\end{enumerate}
\vspace{0.15cm}
\item ¿Quién es el autor de Vivir para contarla?
\begin{enumerate}[label=\textbf{\Alph*)}]
\begin{multicols}{2}
  \item Gabriel García Márquez
  \item Jorge Luis Borges
  \item Federico García Lorca
  \item Pablo Neruda
  \item Octavio Paz
  \item César Vallejo
\end{multicols}
\end{enumerate}
\vspace{0.15cm}
\item ¿Quién es el autor de Noticia de un secuestro?
\begin{enumerate}[label=\textbf{\Alph*)}]
\begin{multicols}{2}
  \item Gabriel García Márquez
  \item Jorge Luis Borges
  \item Federico García Lorca
  \item Pablo Neruda
  \item Octavio Paz
  \item César Vallejo
\end{multicols}
\end{enumerate}
\vspace{0.15cm}
\item ¿Quién es el autor de Manifiesto futurista?
\begin{enumerate}[label=\textbf{\Alph*)}]
\begin{multicols}{2}
  \item Filippo Tommaso Marinetti
  \item Jorge Luis Borges
  \item Federico García Lorca
  \item Pablo Neruda
  \item Octavio Paz
  \item César Vallejo
\end{multicols}
\end{enumerate}
\vspace{0.15cm}
\end{enumerate}\n
\subsubsection*{SECCIÓN III: EXISTENCIALISMO Y ABSURDO (91–140)}

\begin{enumerate}[label=\textbf{\arabic*.}]
\setcounter{enumi}{90}
\item ¿Qué es el existencialismo?
\begin{enumerate}[label=\textbf{\Alph*)}]
\begin{multicols}{2}
  \item Una corriente filosófica y literaria que reflexiona sobre la existencia humana
  \item Jorge Luis Borges
  \item Federico García Lorca
  \item Pablo Neruda
  \item Octavio Paz
  \item César Vallejo
\end{multicols}
\end{enumerate}
\vspace{0.15cm}
\item ¿Quiénes son los principales representantes del existencialismo?
\begin{enumerate}[label=\textbf{\Alph*)}]
\begin{multicols}{2}
  \item Jean-Paul Sartre y Albert Camus
  \item Jorge Luis Borges
  \item Federico García Lorca
  \item Pablo Neruda
  \item Octavio Paz
  \item César Vallejo
\end{multicols}
\end{enumerate}
\vspace{0.15cm}
\item ¿Qué es el absurdo en la literatura?
\begin{enumerate}[label=\textbf{\Alph*)}]
\begin{multicols}{2}
  \item La representación de la falta de sentido de la existencia
  \item Jorge Luis Borges
  \item Federico García Lorca
  \item Pablo Neruda
  \item Octavio Paz
  \item César Vallejo
\end{multicols}
\end{enumerate}
\vspace{0.15cm}
\item ¿Quién es el autor de El extranjero?
\begin{enumerate}[label=\textbf{\Alph*)}]
\begin{multicols}{2}
  \item Albert Camus
  \item Jorge Luis Borges
  \item Federico García Lorca
  \item Pablo Neruda
  \item Octavio Paz
  \item César Vallejo
\end{multicols}
\end{enumerate}
\vspace{0.15cm}
\item ¿Quién es el autor de El mito de Sísifo?
\begin{enumerate}[label=\textbf{\Alph*)}]
\begin{multicols}{2}
  \item Albert Camus
  \item Jorge Luis Borges
  \item Federico García Lorca
  \item Pablo Neruda
  \item Octavio Paz
  \item César Vallejo
\end{multicols}
\end{enumerate}
\vspace{0.15cm}
\item ¿Quién es el autor de La peste?
\begin{enumerate}[label=\textbf{\Alph*)}]
\begin{multicols}{2}
  \item Albert Camus
  \item Jorge Luis Borges
  \item Federico García Lorca
  \item Pablo Neruda
  \item Octavio Paz
  \item César Vallejo
\end{multicols}
\end{enumerate}
\vspace{0.15cm}
\item ¿Quién es el autor de La caída?
\begin{enumerate}[label=\textbf{\Alph*)}]
\begin{multicols}{2}
  \item Albert Camus
  \item Jorge Luis Borges
  \item Federico García Lorca
  \item Pablo Neruda
  \item Octavio Paz
  \item César Vallejo
\end{multicols}
\end{enumerate}
\vspace{0.15cm}
\item ¿Quién es el autor de El ser y la nada?
\begin{enumerate}[label=\textbf{\Alph*)}]
\begin{multicols}{2}
  \item Jean-Paul Sartre
  \item Jorge Luis Borges
  \item Federico García Lorca
  \item Pablo Neruda
  \item Octavio Paz
  \item César Vallejo
\end{multicols}
\end{enumerate}
\vspace{0.15cm}
\item ¿Quién es el autor de La náusea?
\begin{enumerate}[label=\textbf{\Alph*)}]
\begin{multicols}{2}
  \item Jean-Paul Sartre
  \item Jorge Luis Borges
  \item Federico García Lorca
  \item Pablo Neruda
  \item Octavio Paz
  \item César Vallejo
\end{multicols}
\end{enumerate}
\vspace{0.15cm}
\item ¿Quién es el autor de A puerta cerrada?
\begin{enumerate}[label=\textbf{\Alph*)}]
\begin{multicols}{2}
  \item Jean-Paul Sartre
  \item Jorge Luis Borges
  \item Federico García Lorca
  \item Pablo Neruda
  \item Octavio Paz
  \item César Vallejo
\end{multicols}
\end{enumerate}
\vspace{0.15cm}
\item ¿Quién es el autor de Las moscas?
\begin{enumerate}[label=\textbf{\Alph*)}]
\begin{multicols}{2}
  \item Jean-Paul Sartre
  \item Jorge Luis Borges
  \item Federico García Lorca
  \item Pablo Neruda
  \item Octavio Paz
  \item César Vallejo
\end{multicols}
\end{enumerate}
\vspace{0.15cm}
\item ¿Quién es el autor de La metamorfosis?
\begin{enumerate}[label=\textbf{\Alph*)}]
\begin{multicols}{2}
  \item Franz Kafka
  \item Jorge Luis Borges
  \item Federico García Lorca
  \item Pablo Neruda
  \item Octavio Paz
  \item César Vallejo
\end{multicols}
\end{enumerate}
\vspace{0.15cm}
\item ¿Quién es el autor de El proceso?
\begin{enumerate}[label=\textbf{\Alph*)}]
\begin{multicols}{2}
  \item Franz Kafka
  \item Jorge Luis Borges
  \item Federico García Lorca
  \item Pablo Neruda
  \item Octavio Paz
  \item César Vallejo
\end{multicols}
\end{enumerate}
\vspace{0.15cm}
\item ¿Quién es el autor de El castillo?
\begin{enumerate}[label=\textbf{\Alph*)}]
\begin{multicols}{2}
  \item Franz Kafka
  \item Jorge Luis Borges
  \item Federico García Lorca
  \item Pablo Neruda
  \item Octavio Paz
  \item César Vallejo
\end{multicols}
\end{enumerate}
\vspace{0.15cm}
\item ¿Quién es el autor de Carta al padre?
\begin{enumerate}[label=\textbf{\Alph*)}]
\begin{multicols}{2}
  \item Franz Kafka
  \item Jorge Luis Borges
  \item Federico García Lorca
  \item Pablo Neruda
  \item Octavio Paz
  \item César Vallejo
\end{multicols}
\end{enumerate}
\vspace{0.15cm}
\item ¿Quién es el autor de En busca del tiempo perdido?
\begin{enumerate}[label=\textbf{\Alph*)}]
\begin{multicols}{2}
  \item Marcel Proust
  \item Jorge Luis Borges
  \item Federico García Lorca
  \item Pablo Neruda
  \item Octavio Paz
  \item César Vallejo
\end{multicols}
\end{enumerate}
\vspace{0.15cm}
\item ¿Quién es el autor de Ulises?
\begin{enumerate}[label=\textbf{\Alph*)}]
\begin{multicols}{2}
  \item James Joyce
  \item Jorge Luis Borges
  \item Federico García Lorca
  \item Pablo Neruda
  \item Octavio Paz
  \item César Vallejo
\end{multicols}
\end{enumerate}
\vspace{0.15cm}
\item ¿Quién es el autor de Dublineses?
\begin{enumerate}[label=\textbf{\Alph*)}]
\begin{multicols}{2}
  \item James Joyce
  \item Jorge Luis Borges
  \item Federico García Lorca
  \item Pablo Neruda
  \item Octavio Paz
  \item César Vallejo
\end{multicols}
\end{enumerate}
\vspace{0.15cm}
\item ¿Quién es el autor de Retrato del artista adolescente?
\begin{enumerate}[label=\textbf{\Alph*)}]
\begin{multicols}{2}
  \item James Joyce
  \item Jorge Luis Borges
  \item Federico García Lorca
  \item Pablo Neruda
  \item Octavio Paz
  \item César Vallejo
\end{multicols}
\end{enumerate}
\vspace{0.15cm}
\item ¿Quién es la autora de Mrs. Dalloway?
\begin{enumerate}[label=\textbf{\Alph*)}]
\begin{multicols}{2}
  \item Virginia Woolf
  \item Jorge Luis Borges
  \item Federico García Lorca
  \item Pablo Neruda
  \item Octavio Paz
  \item César Vallejo
\end{multicols}
\end{enumerate}
\vspace{0.15cm}
\item ¿Quién es la autora de Al faro?
\begin{enumerate}[label=\textbf{\Alph*)}]
\begin{multicols}{2}
  \item Virginia Woolf
  \item Jorge Luis Borges
  \item Federico García Lorca
  \item Pablo Neruda
  \item Octavio Paz
  \item César Vallejo
\end{multicols}
\end{enumerate}
\vspace{0.15cm}
\item ¿Quién es la autora de Orlando?
\begin{enumerate}[label=\textbf{\Alph*)}]
\begin{multicols}{2}
  \item Virginia Woolf
  \item Jorge Luis Borges
  \item Federico García Lorca
  \item Pablo Neruda
  \item Octavio Paz
  \item César Vallejo
\end{multicols}
\end{enumerate}
\vspace{0.15cm}
\item ¿Quién es la autora de Una habitación propia?
\begin{enumerate}[label=\textbf{\Alph*)}]
\begin{multicols}{2}
  \item Virginia Woolf
  \item Jorge Luis Borges
  \item Federico García Lorca
  \item Pablo Neruda
  \item Octavio Paz
  \item César Vallejo
\end{multicols}
\end{enumerate}
\vspace{0.15cm}
\item ¿Quién es el autor de Esperando a Godot?
\begin{enumerate}[label=\textbf{\Alph*)}]
\begin{multicols}{2}
  \item Samuel Beckett
  \item Jorge Luis Borges
  \item Federico García Lorca
  \item Pablo Neruda
  \item Octavio Paz
  \item César Vallejo
\end{multicols}
\end{enumerate}
\vspace{0.15cm}
\item ¿Quién es el autor de La cantante calva?
\begin{enumerate}[label=\textbf{\Alph*)}]
\begin{multicols}{2}
  \item Eugène Ionesco
  \item Jorge Luis Borges
  \item Federico García Lorca
  \item Pablo Neruda
  \item Octavio Paz
  \item César Vallejo
\end{multicols}
\end{enumerate}
\vspace{0.15cm}
\item ¿Quién es el autor de Seis personajes en busca de autor?
\begin{enumerate}[label=\textbf{\Alph*)}]
\begin{multicols}{2}
  \item Luigi Pirandello
  \item Jorge Luis Borges
  \item Federico García Lorca
  \item Pablo Neruda
  \item Octavio Paz
  \item César Vallejo
\end{multicols}
\end{enumerate}
\vspace{0.15cm}
\item ¿Quién es el autor de Madre Coraje?
\begin{enumerate}[label=\textbf{\Alph*)}]
\begin{multicols}{2}
  \item Bertolt Brecht
  \item Jorge Luis Borges
  \item Federico García Lorca
  \item Pablo Neruda
  \item Octavio Paz
  \item César Vallejo
\end{multicols}
\end{enumerate}
\vspace{0.15cm}
\item ¿Quién es el autor de Galileo Galilei?
\begin{enumerate}[label=\textbf{\Alph*)}]
\begin{multicols}{2}
  \item Bertolt Brecht
  \item Jorge Luis Borges
  \item Federico García Lorca
  \item Pablo Neruda
  \item Octavio Paz
  \item César Vallejo
\end{multicols}
\end{enumerate}
\vspace{0.15cm}
\item ¿Quién es el autor de El círculo de tiza caucasiano?
\begin{enumerate}[label=\textbf{\Alph*)}]
\begin{multicols}{2}
  \item Bertolt Brecht
  \item Jorge Luis Borges
  \item Federico García Lorca
  \item Pablo Neruda
  \item Octavio Paz
  \item César Vallejo
\end{multicols}
\end{enumerate}
\vspace{0.15cm}
\item ¿Qué es el teatro del absurdo?
\begin{enumerate}[label=\textbf{\Alph*)}]
\begin{multicols}{2}
  \item Un género teatral que representa la falta de sentido de la existencia
  \item Jorge Luis Borges
  \item Federico García Lorca
  \item Pablo Neruda
  \item Octavio Paz
  \item César Vallejo
\end{multicols}
\end{enumerate}
\vspace{0.15cm}
\item ¿Qué es el teatro épico?
\begin{enumerate}[label=\textbf{\Alph*)}]
\begin{multicols}{2}
  \item Un teatro que busca la reflexión crítica del espectador
  \item Jorge Luis Borges
  \item Federico García Lorca
  \item Pablo Neruda
  \item Octavio Paz
  \item César Vallejo
\end{multicols}
\end{enumerate}
\vspace{0.15cm}
\item ¿Qué es la novela existencialista?
\begin{enumerate}[label=\textbf{\Alph*)}]
\begin{multicols}{2}
  \item Una novela que reflexiona sobre la existencia humana
  \item Jorge Luis Borges
  \item Federico García Lorca
  \item Pablo Neruda
  \item Octavio Paz
  \item César Vallejo
\end{multicols}
\end{enumerate}
\vspace{0.15cm}
\item ¿Qué es la literatura del absurdo?
\begin{enumerate}[label=\textbf{\Alph*)}]
\begin{multicols}{2}
  \item Una literatura que representa la falta de sentido de la existencia
  \item Jorge Luis Borges
  \item Federico García Lorca
  \item Pablo Neruda
  \item Octavio Paz
  \item César Vallejo
\end{multicols}
\end{enumerate}
\vspace{0.15cm}
\item ¿Qué es el monólogo interior?
\begin{enumerate}[label=\textbf{\Alph*)}]
\begin{multicols}{2}
  \item Una técnica narrativa que reproduce los pensamientos del personaje
  \item Jorge Luis Borges
  \item Federico García Lorca
  \item Pablo Neruda
  \item Octavio Paz
  \item César Vallejo
\end{multicols}
\end{enumerate}
\vspace{0.15cm}
\item ¿Qué es el stream of consciousness?
\begin{enumerate}[label=\textbf{\Alph*)}]
\begin{multicols}{2}
  \item Una técnica narrativa que reproduce el flujo de conciencia del personaje
  \item Jorge Luis Borges
  \item Federico García Lorca
  \item Pablo Neruda
  \item Octavio Paz
  \item César Vallejo
\end{multicols}
\end{enumerate}
\vspace{0.15cm}
\item ¿Qué es la novela experimental?
\begin{enumerate}[label=\textbf{\Alph*)}]
\begin{multicols}{2}
  \item Una novela que rompe con las formas narrativas tradicionales
  \item Jorge Luis Borges
  \item Federico García Lorca
  \item Pablo Neruda
  \item Octavio Paz
  \item César Vallejo
\end{multicols}
\end{enumerate}
\vspace{0.15cm}
\item ¿Qué es el nouveau roman?
\begin{enumerate}[label=\textbf{\Alph*)}]
\begin{multicols}{2}
  \item Un movimiento literario francés que rechaza la novela tradicional
  \item Jorge Luis Borges
  \item Federico García Lorca
  \item Pablo Neruda
  \item Octavio Paz
  \item César Vallejo
\end{multicols}
\end{enumerate}
\vspace{0.15cm}
\item ¿Quién es el principal representante del nouveau roman?
\begin{enumerate}[label=\textbf{\Alph*)}]
\begin{multicols}{2}
  \item Alain Robbe-Grillet
  \item Jorge Luis Borges
  \item Federico García Lorca
  \item Pablo Neruda
  \item Octavio Paz
  \item César Vallejo
\end{multicols}
\end{enumerate}
\vspace{0.15cm}
\item ¿Qué es la metaficción?
\begin{enumerate}[label=\textbf{\Alph*)}]
\begin{multicols}{2}
  \item Una técnica narrativa que reflexiona sobre la propia narración
  \item Jorge Luis Borges
  \item Federico García Lorca
  \item Pablo Neruda
  \item Octavio Paz
  \item César Vallejo
\end{multicols}
\end{enumerate}
\vspace{0.15cm}
\item ¿Quién es el autor de Si una noche de invierno un viajero?
\begin{enumerate}[label=\textbf{\Alph*)}]
\begin{multicols}{2}
  \item Italo Calvino
  \item Jorge Luis Borges
  \item Federico García Lorca
  \item Pablo Neruda
  \item Octavio Paz
  \item César Vallejo
\end{multicols}
\end{enumerate}
\vspace{0.15cm}
\item ¿Quién es el autor de Las ciudades invisibles?
\begin{enumerate}[label=\textbf{\Alph*)}]
\begin{multicols}{2}
  \item Italo Calvino
  \item Jorge Luis Borges
  \item Federico García Lorca
  \item Pablo Neruda
  \item Octavio Paz
  \item César Vallejo
\end{multicols}
\end{enumerate}
\vspace{0.15cm}
\item ¿Quién es el autor de El nombre de la rosa?
\begin{enumerate}[label=\textbf{\Alph*)}]
\begin{multicols}{2}
  \item Umberto Eco
  \item Jorge Luis Borges
  \item Federico García Lorca
  \item Pablo Neruda
  \item Octavio Paz
  \item César Vallejo
\end{multicols}
\end{enumerate}
\vspace{0.15cm}
\item ¿Quién es el autor de El péndulo de Foucault?
\begin{enumerate}[label=\textbf{\Alph*)}]
\begin{multicols}{2}
  \item Umberto Eco
  \item Jorge Luis Borges
  \item Federico García Lorca
  \item Pablo Neruda
  \item Octavio Paz
  \item César Vallejo
\end{multicols}
\end{enumerate}
\vspace{0.15cm}
\item ¿Quién es el autor de La insoportable levedad del ser?
\begin{enumerate}[label=\textbf{\Alph*)}]
\begin{multicols}{2}
  \item Milan Kundera
  \item Jorge Luis Borges
  \item Federico García Lorca
  \item Pablo Neruda
  \item Octavio Paz
  \item César Vallejo
\end{multicols}
\end{enumerate}
\vspace{0.15cm}
\item ¿Quién es el autor de El libro de la risa y el olvido?
\begin{enumerate}[label=\textbf{\Alph*)}]
\begin{multicols}{2}
  \item Milan Kundera
  \item Jorge Luis Borges
  \item Federico García Lorca
  \item Pablo Neruda
  \item Octavio Paz
  \item César Vallejo
\end{multicols}
\end{enumerate}
\vspace{0.15cm}
\item ¿Quién es el autor de La broma?
\begin{enumerate}[label=\textbf{\Alph*)}]
\begin{multicols}{2}
  \item Milan Kundera
  \item Jorge Luis Borges
  \item Federico García Lorca
  \item Pablo Neruda
  \item Octavio Paz
  \item César Vallejo
\end{multicols}
\end{enumerate}
\vspace{0.15cm}
\item ¿Quién es el autor de El tambor de hojalata?
\begin{enumerate}[label=\textbf{\Alph*)}]
\begin{multicols}{2}
  \item Günter Grass
  \item Jorge Luis Borges
  \item Federico García Lorca
  \item Pablo Neruda
  \item Octavio Paz
  \item César Vallejo
\end{multicols}
\end{enumerate}
\vspace{0.15cm}
\item ¿Quién es el autor de El gatopardo?
\begin{enumerate}[label=\textbf{\Alph*)}]
\begin{multicols}{2}
  \item Giuseppe Tomasi di Lampedusa
  \item Jorge Luis Borges
  \item Federico García Lorca
  \item Pablo Neruda
  \item Octavio Paz
  \item César Vallejo
\end{multicols}
\end{enumerate}
\vspace{0.15cm}
\item ¿Quién es el autor de La montaña mágica?
\begin{enumerate}[label=\textbf{\Alph*)}]
\begin{multicols}{2}
  \item Thomas Mann
  \item Jorge Luis Borges
  \item Federico García Lorca
  \item Pablo Neruda
  \item Octavio Paz
  \item César Vallejo
\end{multicols}
\end{enumerate}
\vspace{0.15cm}
\item ¿Quién es el autor de Muerte en Venecia?
\begin{enumerate}[label=\textbf{\Alph*)}]
\begin{multicols}{2}
  \item Thomas Mann
  \item Jorge Luis Borges
  \item Federico García Lorca
  \item Pablo Neruda
  \item Octavio Paz
  \item César Vallejo
\end{multicols}
\end{enumerate}
\vspace{0.15cm}
\end{enumerate}\n
\subsubsection*{SECCIÓN IV: NOVELA CONTEMPORÁNEA (141–190)}

\begin{enumerate}[label=\textbf{\arabic*.}]
\setcounter{enumi}{140}
\item ¿Quién es el autor de 1984?
\begin{enumerate}[label=\textbf{\Alph*)}]
\begin{multicols}{2}
  \item George Orwell
  \item Jorge Luis Borges
  \item Federico García Lorca
  \item Pablo Neruda
  \item Octavio Paz
  \item César Vallejo
\end{multicols}
\end{enumerate}
\vspace{0.15cm}
\item ¿Quién es el autor de Rebelión en la granja?
\begin{enumerate}[label=\textbf{\Alph*)}]
\begin{multicols}{2}
  \item George Orwell
  \item Jorge Luis Borges
  \item Federico García Lorca
  \item Pablo Neruda
  \item Octavio Paz
  \item César Vallejo
\end{multicols}
\end{enumerate}
\vspace{0.15cm}
\item ¿Quién es el autor de Un mundo feliz?
\begin{enumerate}[label=\textbf{\Alph*)}]
\begin{multicols}{2}
  \item Aldous Huxley
  \item Jorge Luis Borges
  \item Federico García Lorca
  \item Pablo Neruda
  \item Octavio Paz
  \item César Vallejo
\end{multicols}
\end{enumerate}
\vspace{0.15cm}
\item ¿Quién es el autor de Fahrenheit 451?
\begin{enumerate}[label=\textbf{\Alph*)}]
\begin{multicols}{2}
  \item Ray Bradbury
  \item Jorge Luis Borges
  \item Federico García Lorca
  \item Pablo Neruda
  \item Octavio Paz
  \item César Vallejo
\end{multicols}
\end{enumerate}
\vspace{0.15cm}
\item ¿Quién es el autor de El señor de las moscas?
\begin{enumerate}[label=\textbf{\Alph*)}]
\begin{multicols}{2}
  \item William Golding
  \item Jorge Luis Borges
  \item Federico García Lorca
  \item Pablo Neruda
  \item Octavio Paz
  \item César Vallejo
\end{multicols}
\end{enumerate}
\vspace{0.15cm}
\item ¿Quién es el autor de La naranja mecánica?
\begin{enumerate}[label=\textbf{\Alph*)}]
\begin{multicols}{2}
  \item Anthony Burgess
  \item Jorge Luis Borges
  \item Federico García Lorca
  \item Pablo Neruda
  \item Octavio Paz
  \item César Vallejo
\end{multicols}
\end{enumerate}
\vspace{0.15cm}
\item ¿Quién es el autor de El guardián entre el centeno?
\begin{enumerate}[label=\textbf{\Alph*)}]
\begin{multicols}{2}
  \item J.D. Salinger
  \item Jorge Luis Borges
  \item Federico García Lorca
  \item Pablo Neruda
  \item Octavio Paz
  \item César Vallejo
\end{multicols}
\end{enumerate}
\vspace{0.15cm}
\item ¿Quién es el autor de Trampa 22?
\begin{enumerate}[label=\textbf{\Alph*)}]
\begin{multicols}{2}
  \item Joseph Heller
  \item Jorge Luis Borges
  \item Federico García Lorca
  \item Pablo Neruda
  \item Octavio Paz
  \item César Vallejo
\end{multicols}
\end{enumerate}
\vspace{0.15cm}
\item ¿Quién es el autor de Matadero cinco?
\begin{enumerate}[label=\textbf{\Alph*)}]
\begin{multicols}{2}
  \item Kurt Vonnegut
  \item Jorge Luis Borges
  \item Federico García Lorca
  \item Pablo Neruda
  \item Octavio Paz
  \item César Vallejo
\end{multicols}
\end{enumerate}
\vspace{0.15cm}
\item ¿Quién es el autor de El arco iris de gravedad?
\begin{enumerate}[label=\textbf{\Alph*)}]
\begin{multicols}{2}
  \item Thomas Pynchon
  \item Jorge Luis Borges
  \item Federico García Lorca
  \item Pablo Neruda
  \item Octavio Paz
  \item César Vallejo
\end{multicols}
\end{enumerate}
\vspace{0.15cm}
\item ¿Quién es el autor de La conjura de los necios?
\begin{enumerate}[label=\textbf{\Alph*)}]
\begin{multicols}{2}
  \item John Kennedy Toole
  \item Jorge Luis Borges
  \item Federico García Lorca
  \item Pablo Neruda
  \item Octavio Paz
  \item César Vallejo
\end{multicols}
\end{enumerate}
\vspace{0.15cm}
\item ¿Quién es el autor de Beloved?
\begin{enumerate}[label=\textbf{\Alph*)}]
\begin{multicols}{2}
  \item Toni Morrison
  \item Jorge Luis Borges
  \item Federico García Lorca
  \item Pablo Neruda
  \item Octavio Paz
  \item César Vallejo
\end{multicols}
\end{enumerate}
\vspace{0.15cm}
\item ¿Quién es el autor de El ruido y la furia?
\begin{enumerate}[label=\textbf{\Alph*)}]
\begin{multicols}{2}
  \item William Faulkner
  \item Jorge Luis Borges
  \item Federico García Lorca
  \item Pablo Neruda
  \item Octavio Paz
  \item César Vallejo
\end{multicols}
\end{enumerate}
\vspace{0.15cm}
\item ¿Quién es el autor de Mientras agonizo?
\begin{enumerate}[label=\textbf{\Alph*)}]
\begin{multicols}{2}
  \item William Faulkner
  \item Jorge Luis Borges
  \item Federico García Lorca
  \item Pablo Neruda
  \item Octavio Paz
  \item César Vallejo
\end{multicols}
\end{enumerate}
\vspace{0.15cm}
\item ¿Quién es el autor de Absalón, Absalón?
\begin{enumerate}[label=\textbf{\Alph*)}]
\begin{multicols}{2}
  \item William Faulkner
  \item Jorge Luis Borges
  \item Federico García Lorca
  \item Pablo Neruda
  \item Octavio Paz
  \item César Vallejo
\end{multicols}
\end{enumerate}
\vspace{0.15cm}
\item ¿Quién es el autor de Las palmeras salvajes?
\begin{enumerate}[label=\textbf{\Alph*)}]
\begin{multicols}{2}
  \item William Faulkner
  \item Jorge Luis Borges
  \item Federico García Lorca
  \item Pablo Neruda
  \item Octavio Paz
  \item César Vallejo
\end{multicols}
\end{enumerate}
\vspace{0.15cm}
\item ¿Quién es el autor de Santuario?
\begin{enumerate}[label=\textbf{\Alph*)}]
\begin{multicols}{2}
  \item William Faulkner
  \item Jorge Luis Borges
  \item Federico García Lorca
  \item Pablo Neruda
  \item Octavio Paz
  \item César Vallejo
\end{multicols}
\end{enumerate}
\vspace{0.15cm}
\item ¿Quién es el autor de Por quién doblan las campanas?
\begin{enumerate}[label=\textbf{\Alph*)}]
\begin{multicols}{2}
  \item Ernest Hemingway
  \item Jorge Luis Borges
  \item Federico García Lorca
  \item Pablo Neruda
  \item Octavio Paz
  \item César Vallejo
\end{multicols}
\end{enumerate}
\vspace{0.15cm}
\item ¿Quién es el autor de El viejo y el mar?
\begin{enumerate}[label=\textbf{\Alph*)}]
\begin{multicols}{2}
  \item Ernest Hemingway
  \item Jorge Luis Borges
  \item Federico García Lorca
  \item Pablo Neruda
  \item Octavio Paz
  \item César Vallejo
\end{multicols}
\end{enumerate}
\vspace{0.15cm}
\item ¿Quién es el autor de Adiós a las armas?
\begin{enumerate}[label=\textbf{\Alph*)}]
\begin{multicols}{2}
  \item Ernest Hemingway
  \item Jorge Luis Borges
  \item Federico García Lorca
  \item Pablo Neruda
  \item Octavio Paz
  \item César Vallejo
\end{multicols}
\end{enumerate}
\vspace{0.15cm}
\item ¿Quién es el autor de Fiesta?
\begin{enumerate}[label=\textbf{\Alph*)}]
\begin{multicols}{2}
  \item Ernest Hemingway
  \item Jorge Luis Borges
  \item Federico García Lorca
  \item Pablo Neruda
  \item Octavio Paz
  \item César Vallejo
\end{multicols}
\end{enumerate}
\vspace{0.15cm}
\item ¿Quién es el autor de 1984?
\begin{enumerate}[label=\textbf{\Alph*)}]
\begin{multicols}{2}
  \item George Orwell
  \item Jorge Luis Borges
  \item Federico García Lorca
  \item Pablo Neruda
  \item Octavio Paz
  \item César Vallejo
\end{multicols}
\end{enumerate}
\vspace{0.15cm}
\item ¿Quién es el autor de Rebelión en la granja?
\begin{enumerate}[label=\textbf{\Alph*)}]
\begin{multicols}{2}
  \item George Orwell
  \item Jorge Luis Borges
  \item Federico García Lorca
  \item Pablo Neruda
  \item Octavio Paz
  \item César Vallejo
\end{multicols}
\end{enumerate}
\vspace{0.15cm}
\item ¿Quién es el autor de Un mundo feliz?
\begin{enumerate}[label=\textbf{\Alph*)}]
\begin{multicols}{2}
  \item Aldous Huxley
  \item Jorge Luis Borges
  \item Federico García Lorca
  \item Pablo Neruda
  \item Octavio Paz
  \item César Vallejo
\end{multicols}
\end{enumerate}
\vspace{0.15cm}
\item ¿Quién es el autor de Fahrenheit 451?
\begin{enumerate}[label=\textbf{\Alph*)}]
\begin{multicols}{2}
  \item Ray Bradbury
  \item Jorge Luis Borges
  \item Federico García Lorca
  \item Pablo Neruda
  \item Octavio Paz
  \item César Vallejo
\end{multicols}
\end{enumerate}
\vspace{0.15cm}
\item ¿Quién es el autor de El señor de las moscas?
\begin{enumerate}[label=\textbf{\Alph*)}]
\begin{multicols}{2}
  \item William Golding
  \item Jorge Luis Borges
  \item Federico García Lorca
  \item Pablo Neruda
  \item Octavio Paz
  \item César Vallejo
\end{multicols}
\end{enumerate}
\vspace{0.15cm}
\item ¿Quién es el autor de La naranja mecánica?
\begin{enumerate}[label=\textbf{\Alph*)}]
\begin{multicols}{2}
  \item Anthony Burgess
  \item Jorge Luis Borges
  \item Federico García Lorca
  \item Pablo Neruda
  \item Octavio Paz
  \item César Vallejo
\end{multicols}
\end{enumerate}
\vspace{0.15cm}
\item ¿Quién es el autor de El guardián entre el centeno?
\begin{enumerate}[label=\textbf{\Alph*)}]
\begin{multicols}{2}
  \item J.D. Salinger
  \item Jorge Luis Borges
  \item Federico García Lorca
  \item Pablo Neruda
  \item Octavio Paz
  \item César Vallejo
\end{multicols}
\end{enumerate}
\vspace{0.15cm}
\item ¿Quién es el autor de Trampa 22?
\begin{enumerate}[label=\textbf{\Alph*)}]
\begin{multicols}{2}
  \item Joseph Heller
  \item Jorge Luis Borges
  \item Federico García Lorca
  \item Pablo Neruda
  \item Octavio Paz
  \item César Vallejo
\end{multicols}
\end{enumerate}
\vspace{0.15cm}
\item ¿Quién es el autor de Matadero cinco?
\begin{enumerate}[label=\textbf{\Alph*)}]
\begin{multicols}{2}
  \item Kurt Vonnegut
  \item Jorge Luis Borges
  \item Federico García Lorca
  \item Pablo Neruda
  \item Octavio Paz
  \item César Vallejo
\end{multicols}
\end{enumerate}
\vspace{0.15cm}
\item ¿Quién es el autor de El arco iris de gravedad?
\begin{enumerate}[label=\textbf{\Alph*)}]
\begin{multicols}{2}
  \item Thomas Pynchon
  \item Jorge Luis Borges
  \item Federico García Lorca
  \item Pablo Neruda
  \item Octavio Paz
  \item César Vallejo
\end{multicols}
\end{enumerate}
\vspace{0.15cm}
\item ¿Quién es el autor de La conjura de los necios?
\begin{enumerate}[label=\textbf{\Alph*)}]
\begin{multicols}{2}
  \item John Kennedy Toole
  \item Jorge Luis Borges
  \item Federico García Lorca
  \item Pablo Neruda
  \item Octavio Paz
  \item César Vallejo
\end{multicols}
\end{enumerate}
\vspace{0.15cm}
\item ¿Quién es el autor de Beloved?
\begin{enumerate}[label=\textbf{\Alph*)}]
\begin{multicols}{2}
  \item Toni Morrison
  \item Jorge Luis Borges
  \item Federico García Lorca
  \item Pablo Neruda
  \item Octavio Paz
  \item César Vallejo
\end{multicols}
\end{enumerate}
\vspace{0.15cm}
\item ¿Quién es el autor de El ruido y la furia?
\begin{enumerate}[label=\textbf{\Alph*)}]
\begin{multicols}{2}
  \item William Faulkner
  \item Jorge Luis Borges
  \item Federico García Lorca
  \item Pablo Neruda
  \item Octavio Paz
  \item César Vallejo
\end{multicols}
\end{enumerate}
\vspace{0.15cm}
\item ¿Quién es el autor de Mientras agonizo?
\begin{enumerate}[label=\textbf{\Alph*)}]
\begin{multicols}{2}
  \item William Faulkner
  \item Jorge Luis Borges
  \item Federico García Lorca
  \item Pablo Neruda
  \item Octavio Paz
  \item César Vallejo
\end{multicols}
\end{enumerate}
\vspace{0.15cm}
\item ¿Quién es el autor de Absalón, Absalón?
\begin{enumerate}[label=\textbf{\Alph*)}]
\begin{multicols}{2}
  \item William Faulkner
  \item Jorge Luis Borges
  \item Federico García Lorca
  \item Pablo Neruda
  \item Octavio Paz
  \item César Vallejo
\end{multicols}
\end{enumerate}
\vspace{0.15cm}
\item ¿Quién es el autor de Las palmeras salvajes?
\begin{enumerate}[label=\textbf{\Alph*)}]
\begin{multicols}{2}
  \item William Faulkner
  \item Jorge Luis Borges
  \item Federico García Lorca
  \item Pablo Neruda
  \item Octavio Paz
  \item César Vallejo
\end{multicols}
\end{enumerate}
\vspace{0.15cm}
\item ¿Quién es el autor de Santuario?
\begin{enumerate}[label=\textbf{\Alph*)}]
\begin{multicols}{2}
  \item William Faulkner
  \item Jorge Luis Borges
  \item Federico García Lorca
  \item Pablo Neruda
  \item Octavio Paz
  \item César Vallejo
\end{multicols}
\end{enumerate}
\vspace{0.15cm}
\item ¿Quién es el autor de Por quién doblan las campanas?
\begin{enumerate}[label=\textbf{\Alph*)}]
\begin{multicols}{2}
  \item Ernest Hemingway
  \item Jorge Luis Borges
  \item Federico García Lorca
  \item Pablo Neruda
  \item Octavio Paz
  \item César Vallejo
\end{multicols}
\end{enumerate}
\vspace{0.15cm}
\item ¿Quién es el autor de El viejo y el mar?
\begin{enumerate}[label=\textbf{\Alph*)}]
\begin{multicols}{2}
  \item Ernest Hemingway
  \item Jorge Luis Borges
  \item Federico García Lorca
  \item Pablo Neruda
  \item Octavio Paz
  \item César Vallejo
\end{multicols}
\end{enumerate}
\vspace{0.15cm}
\item ¿Quién es el autor de Adiós a las armas?
\begin{enumerate}[label=\textbf{\Alph*)}]
\begin{multicols}{2}
  \item Ernest Hemingway
  \item Jorge Luis Borges
  \item Federico García Lorca
  \item Pablo Neruda
  \item Octavio Paz
  \item César Vallejo
\end{multicols}
\end{enumerate}
\vspace{0.15cm}
\item ¿Quién es el autor de Fiesta?
\begin{enumerate}[label=\textbf{\Alph*)}]
\begin{multicols}{2}
  \item Ernest Hemingway
  \item Jorge Luis Borges
  \item Federico García Lorca
  \item Pablo Neruda
  \item Octavio Paz
  \item César Vallejo
\end{multicols}
\end{enumerate}
\vspace{0.15cm}
\item ¿Quién es el autor de 1984?
\begin{enumerate}[label=\textbf{\Alph*)}]
\begin{multicols}{2}
  \item George Orwell
  \item Jorge Luis Borges
  \item Federico García Lorca
  \item Pablo Neruda
  \item Octavio Paz
  \item César Vallejo
\end{multicols}
\end{enumerate}
\vspace{0.15cm}
\item ¿Quién es el autor de Rebelión en la granja?
\begin{enumerate}[label=\textbf{\Alph*)}]
\begin{multicols}{2}
  \item George Orwell
  \item Jorge Luis Borges
  \item Federico García Lorca
  \item Pablo Neruda
  \item Octavio Paz
  \item César Vallejo
\end{multicols}
\end{enumerate}
\vspace{0.15cm}
\item ¿Quién es el autor de Un mundo feliz?
\begin{enumerate}[label=\textbf{\Alph*)}]
\begin{multicols}{2}
  \item Aldous Huxley
  \item Jorge Luis Borges
  \item Federico García Lorca
  \item Pablo Neruda
  \item Octavio Paz
  \item César Vallejo
\end{multicols}
\end{enumerate}
\vspace{0.15cm}
\item ¿Quién es el autor de Fahrenheit 451?
\begin{enumerate}[label=\textbf{\Alph*)}]
\begin{multicols}{2}
  \item Ray Bradbury
  \item Jorge Luis Borges
  \item Federico García Lorca
  \item Pablo Neruda
  \item Octavio Paz
  \item César Vallejo
\end{multicols}
\end{enumerate}
\vspace{0.15cm}
\item ¿Quién es el autor de El señor de las moscas?
\begin{enumerate}[label=\textbf{\Alph*)}]
\begin{multicols}{2}
  \item William Golding
  \item Jorge Luis Borges
  \item Federico García Lorca
  \item Pablo Neruda
  \item Octavio Paz
  \item César Vallejo
\end{multicols}
\end{enumerate}
\vspace{0.15cm}
\item ¿Quién es el autor de La naranja mecánica?
\begin{enumerate}[label=\textbf{\Alph*)}]
\begin{multicols}{2}
  \item Anthony Burgess
  \item Jorge Luis Borges
  \item Federico García Lorca
  \item Pablo Neruda
  \item Octavio Paz
  \item César Vallejo
\end{multicols}
\end{enumerate}
\vspace{0.15cm}
\item ¿Quién es el autor de El guardián entre el centeno?
\begin{enumerate}[label=\textbf{\Alph*)}]
\begin{multicols}{2}
  \item J.D. Salinger
  \item Jorge Luis Borges
  \item Federico García Lorca
  \item Pablo Neruda
  \item Octavio Paz
  \item César Vallejo
\end{multicols}
\end{enumerate}
\vspace{0.15cm}
\item ¿Quién es el autor de Trampa 22?
\begin{enumerate}[label=\textbf{\Alph*)}]
\begin{multicols}{2}
  \item Joseph Heller
  \item Jorge Luis Borges
  \item Federico García Lorca
  \item Pablo Neruda
  \item Octavio Paz
  \item César Vallejo
\end{multicols}
\end{enumerate}
\vspace{0.15cm}
\end{enumerate}\n
\subsubsection*{SECCIÓN V: TEATRO CONTEMPORÁNEO (191–230)}

\begin{enumerate}[label=\textbf{\arabic*.}]
\setcounter{enumi}{190}
\item ¿Quién es el autor de Historia de una escalera?
\begin{enumerate}[label=\textbf{\Alph*)}]
\begin{multicols}{2}
  \item Antonio Buero Vallejo
  \item Jorge Luis Borges
  \item Federico García Lorca
  \item Pablo Neruda
  \item Octavio Paz
  \item César Vallejo
\end{multicols}
\end{enumerate}
\vspace{0.15cm}
\item ¿Quién es el autor de En la ardiente oscuridad?
\begin{enumerate}[label=\textbf{\Alph*)}]
\begin{multicols}{2}
  \item Antonio Buero Vallejo
  \item Jorge Luis Borges
  \item Federico García Lorca
  \item Pablo Neruda
  \item Octavio Paz
  \item César Vallejo
\end{multicols}
\end{enumerate}
\vspace{0.15cm}
\item ¿Quién es el autor de El tragaluz?
\begin{enumerate}[label=\textbf{\Alph*)}]
\begin{multicols}{2}
  \item Antonio Buero Vallejo
  \item Jorge Luis Borges
  \item Federico García Lorca
  \item Pablo Neruda
  \item Octavio Paz
  \item César Vallejo
\end{multicols}
\end{enumerate}
\vspace{0.15cm}
\item ¿Quién es el autor de La fundación?
\begin{enumerate}[label=\textbf{\Alph*)}]
\begin{multicols}{2}
  \item Antonio Buero Vallejo
  \item Jorge Luis Borges
  \item Federico García Lorca
  \item Pablo Neruda
  \item Octavio Paz
  \item César Vallejo
\end{multicols}
\end{enumerate}
\vspace{0.15cm}
\item ¿Quién es el autor de Un soñador para un pueblo?
\begin{enumerate}[label=\textbf{\Alph*)}]
\begin{multicols}{2}
  \item Antonio Buero Vallejo
  \item Jorge Luis Borges
  \item Federico García Lorca
  \item Pablo Neruda
  \item Octavio Paz
  \item César Vallejo
\end{multicols}
\end{enumerate}
\vspace{0.15cm}
\item ¿Quién es el autor de Las Meninas?
\begin{enumerate}[label=\textbf{\Alph*)}]
\begin{multicols}{2}
  \item Antonio Buero Vallejo
  \item Jorge Luis Borges
  \item Federico García Lorca
  \item Pablo Neruda
  \item Octavio Paz
  \item César Vallejo
\end{multicols}
\end{enumerate}
\vspace{0.15cm}
\item ¿Quién es el autor de El concierto de San Ovidio?
\begin{enumerate}[label=\textbf{\Alph*)}]
\begin{multicols}{2}
  \item Antonio Buero Vallejo
  \item Jorge Luis Borges
  \item Federico García Lorca
  \item Pablo Neruda
  \item Octavio Paz
  \item César Vallejo
\end{multicols}
\end{enumerate}
\vspace{0.15cm}
\item ¿Quién es el autor de La doble historia del doctor Valmy?
\begin{enumerate}[label=\textbf{\Alph*)}]
\begin{multicols}{2}
  \item Antonio Buero Vallejo
  \item Jorge Luis Borges
  \item Federico García Lorca
  \item Pablo Neruda
  \item Octavio Paz
  \item César Vallejo
\end{multicols}
\end{enumerate}
\vspace{0.15cm}
\item ¿Quién es el autor de Muerte de un viajante?
\begin{enumerate}[label=\textbf{\Alph*)}]
\begin{multicols}{2}
  \item Arthur Miller
  \item Jorge Luis Borges
  \item Federico García Lorca
  \item Pablo Neruda
  \item Octavio Paz
  \item César Vallejo
\end{multicols}
\end{enumerate}
\vspace{0.15cm}
\item ¿Quién es el autor de Las brujas de Salem?
\begin{enumerate}[label=\textbf{\Alph*)}]
\begin{multicols}{2}
  \item Arthur Miller
  \item Jorge Luis Borges
  \item Federico García Lorca
  \item Pablo Neruda
  \item Octavio Paz
  \item César Vallejo
\end{multicols}
\end{enumerate}
\vspace{0.15cm}
\item ¿Quién es el autor de Un tranvía llamado Deseo?
\begin{enumerate}[label=\textbf{\Alph*)}]
\begin{multicols}{2}
  \item Tennessee Williams
  \item Jorge Luis Borges
  \item Federico García Lorca
  \item Pablo Neruda
  \item Octavio Paz
  \item César Vallejo
\end{multicols}
\end{enumerate}
\vspace{0.15cm}
\item ¿Quién es el autor de La gata sobre el tejado de zinc?
\begin{enumerate}[label=\textbf{\Alph*)}]
\begin{multicols}{2}
  \item Tennessee Williams
  \item Jorge Luis Borges
  \item Federico García Lorca
  \item Pablo Neruda
  \item Octavio Paz
  \item César Vallejo
\end{multicols}
\end{enumerate}
\vspace{0.15cm}
\item ¿Quién es el autor de El zoo de cristal?
\begin{enumerate}[label=\textbf{\Alph*)}]
\begin{multicols}{2}
  \item Tennessee Williams
  \item Jorge Luis Borges
  \item Federico García Lorca
  \item Pablo Neruda
  \item Octavio Paz
  \item César Vallejo
\end{multicols}
\end{enumerate}
\vspace{0.15cm}
\item ¿Quién es el autor de Deseo bajo los olmos?
\begin{enumerate}[label=\textbf{\Alph*)}]
\begin{multicols}{2}
  \item Eugene O'Neill
  \item Jorge Luis Borges
  \item Federico García Lorca
  \item Pablo Neruda
  \item Octavio Paz
  \item César Vallejo
\end{multicols}
\end{enumerate}
\vspace{0.15cm}
\item ¿Quién es el autor de Largo viaje hacia la noche?
\begin{enumerate}[label=\textbf{\Alph*)}]
\begin{multicols}{2}
  \item Eugene O'Neill
  \item Jorge Luis Borges
  \item Federico García Lorca
  \item Pablo Neruda
  \item Octavio Paz
  \item César Vallejo
\end{multicols}
\end{enumerate}
\vspace{0.15cm}
\item ¿Quién es el autor de El emperador Jones?
\begin{enumerate}[label=\textbf{\Alph*)}]
\begin{multicols}{2}
  \item Eugene O'Neill
  \item Jorge Luis Borges
  \item Federico García Lorca
  \item Pablo Neruda
  \item Octavio Paz
  \item César Vallejo
\end{multicols}
\end{enumerate}
\vspace{0.15cm}
\item ¿Quién es el autor de Calígula?
\begin{enumerate}[label=\textbf{\Alph*)}]
\begin{multicols}{2}
  \item Albert Camus
  \item Jorge Luis Borges
  \item Federico García Lorca
  \item Pablo Neruda
  \item Octavio Paz
  \item César Vallejo
\end{multicols}
\end{enumerate}
\vspace{0.15cm}
\item ¿Quién es el autor de Los justos?
\begin{enumerate}[label=\textbf{\Alph*)}]
\begin{multicols}{2}
  \item Albert Camus
  \item Jorge Luis Borges
  \item Federico García Lorca
  \item Pablo Neruda
  \item Octavio Paz
  \item César Vallejo
\end{multicols}
\end{enumerate}
\vspace{0.15cm}
\item ¿Quién es el autor de El diablo y Dios?
\begin{enumerate}[label=\textbf{\Alph*)}]
\begin{multicols}{2}
  \item Jean-Paul Sartre
  \item Jorge Luis Borges
  \item Federico García Lorca
  \item Pablo Neruda
  \item Octavio Paz
  \item César Vallejo
\end{multicols}
\end{enumerate}
\vspace{0.15cm}
\item ¿Quién es el autor de Las manos sucias?
\begin{enumerate}[label=\textbf{\Alph*)}]
\begin{multicols}{2}
  \item Jean-Paul Sartre
  \item Jorge Luis Borges
  \item Federico García Lorca
  \item Pablo Neruda
  \item Octavio Paz
  \item César Vallejo
\end{multicols}
\end{enumerate}
\vspace{0.15cm}
\item ¿Quién es el autor de Historia de una escalera?
\begin{enumerate}[label=\textbf{\Alph*)}]
\begin{multicols}{2}
  \item Antonio Buero Vallejo
  \item Jorge Luis Borges
  \item Federico García Lorca
  \item Pablo Neruda
  \item Octavio Paz
  \item César Vallejo
\end{multicols}
\end{enumerate}
\vspace{0.15cm}
\item ¿Quién es el autor de En la ardiente oscuridad?
\begin{enumerate}[label=\textbf{\Alph*)}]
\begin{multicols}{2}
  \item Antonio Buero Vallejo
  \item Jorge Luis Borges
  \item Federico García Lorca
  \item Pablo Neruda
  \item Octavio Paz
  \item César Vallejo
\end{multicols}
\end{enumerate}
\vspace{0.15cm}
\item ¿Quién es el autor de El tragaluz?
\begin{enumerate}[label=\textbf{\Alph*)}]
\begin{multicols}{2}
  \item Antonio Buero Vallejo
  \item Jorge Luis Borges
  \item Federico García Lorca
  \item Pablo Neruda
  \item Octavio Paz
  \item César Vallejo
\end{multicols}
\end{enumerate}
\vspace{0.15cm}
\item ¿Quién es el autor de La fundación?
\begin{enumerate}[label=\textbf{\Alph*)}]
\begin{multicols}{2}
  \item Antonio Buero Vallejo
  \item Jorge Luis Borges
  \item Federico García Lorca
  \item Pablo Neruda
  \item Octavio Paz
  \item César Vallejo
\end{multicols}
\end{enumerate}
\vspace{0.15cm}
\item ¿Quién es el autor de Un soñador para un pueblo?
\begin{enumerate}[label=\textbf{\Alph*)}]
\begin{multicols}{2}
  \item Antonio Buero Vallejo
  \item Jorge Luis Borges
  \item Federico García Lorca
  \item Pablo Neruda
  \item Octavio Paz
  \item César Vallejo
\end{multicols}
\end{enumerate}
\vspace{0.15cm}
\item ¿Quién es el autor de Las Meninas?
\begin{enumerate}[label=\textbf{\Alph*)}]
\begin{multicols}{2}
  \item Antonio Buero Vallejo
  \item Jorge Luis Borges
  \item Federico García Lorca
  \item Pablo Neruda
  \item Octavio Paz
  \item César Vallejo
\end{multicols}
\end{enumerate}
\vspace{0.15cm}
\item ¿Quién es el autor de El concierto de San Ovidio?
\begin{enumerate}[label=\textbf{\Alph*)}]
\begin{multicols}{2}
  \item Antonio Buero Vallejo
  \item Jorge Luis Borges
  \item Federico García Lorca
  \item Pablo Neruda
  \item Octavio Paz
  \item César Vallejo
\end{multicols}
\end{enumerate}
\vspace{0.15cm}
\item ¿Quién es el autor de La doble historia del doctor Valmy?
\begin{enumerate}[label=\textbf{\Alph*)}]
\begin{multicols}{2}
  \item Antonio Buero Vallejo
  \item Jorge Luis Borges
  \item Federico García Lorca
  \item Pablo Neruda
  \item Octavio Paz
  \item César Vallejo
\end{multicols}
\end{enumerate}
\vspace{0.15cm}
\item ¿Quién es el autor de Muerte de un viajante?
\begin{enumerate}[label=\textbf{\Alph*)}]
\begin{multicols}{2}
  \item Arthur Miller
  \item Jorge Luis Borges
  \item Federico García Lorca
  \item Pablo Neruda
  \item Octavio Paz
  \item César Vallejo
\end{multicols}
\end{enumerate}
\vspace{0.15cm}
\item ¿Quién es el autor de Las brujas de Salem?
\begin{enumerate}[label=\textbf{\Alph*)}]
\begin{multicols}{2}
  \item Arthur Miller
  \item Jorge Luis Borges
  \item Federico García Lorca
  \item Pablo Neruda
  \item Octavio Paz
  \item César Vallejo
\end{multicols}
\end{enumerate}
\vspace{0.15cm}
\item ¿Quién es el autor de Un tranvía llamado Deseo?
\begin{enumerate}[label=\textbf{\Alph*)}]
\begin{multicols}{2}
  \item Tennessee Williams
  \item Jorge Luis Borges
  \item Federico García Lorca
  \item Pablo Neruda
  \item Octavio Paz
  \item César Vallejo
\end{multicols}
\end{enumerate}
\vspace{0.15cm}
\item ¿Quién es el autor de La gata sobre el tejado de zinc?
\begin{enumerate}[label=\textbf{\Alph*)}]
\begin{multicols}{2}
  \item Tennessee Williams
  \item Jorge Luis Borges
  \item Federico García Lorca
  \item Pablo Neruda
  \item Octavio Paz
  \item César Vallejo
\end{multicols}
\end{enumerate}
\vspace{0.15cm}
\item ¿Quién es el autor de El zoo de cristal?
\begin{enumerate}[label=\textbf{\Alph*)}]
\begin{multicols}{2}
  \item Tennessee Williams
  \item Jorge Luis Borges
  \item Federico García Lorca
  \item Pablo Neruda
  \item Octavio Paz
  \item César Vallejo
\end{multicols}
\end{enumerate}
\vspace{0.15cm}
\item ¿Quién es el autor de Deseo bajo los olmos?
\begin{enumerate}[label=\textbf{\Alph*)}]
\begin{multicols}{2}
  \item Eugene O'Neill
  \item Jorge Luis Borges
  \item Federico García Lorca
  \item Pablo Neruda
  \item Octavio Paz
  \item César Vallejo
\end{multicols}
\end{enumerate}
\vspace{0.15cm}
\item ¿Quién es el autor de Largo viaje hacia la noche?
\begin{enumerate}[label=\textbf{\Alph*)}]
\begin{multicols}{2}
  \item Eugene O'Neill
  \item Jorge Luis Borges
  \item Federico García Lorca
  \item Pablo Neruda
  \item Octavio Paz
  \item César Vallejo
\end{multicols}
\end{enumerate}
\vspace{0.15cm}
\item ¿Quién es el autor de El emperador Jones?
\begin{enumerate}[label=\textbf{\Alph*)}]
\begin{multicols}{2}
  \item Eugene O'Neill
  \item Jorge Luis Borges
  \item Federico García Lorca
  \item Pablo Neruda
  \item Octavio Paz
  \item César Vallejo
\end{multicols}
\end{enumerate}
\vspace{0.15cm}
\item ¿Quién es el autor de Calígula?
\begin{enumerate}[label=\textbf{\Alph*)}]
\begin{multicols}{2}
  \item Albert Camus
  \item Jorge Luis Borges
  \item Federico García Lorca
  \item Pablo Neruda
  \item Octavio Paz
  \item César Vallejo
\end{multicols}
\end{enumerate}
\vspace{0.15cm}
\item ¿Quién es el autor de Los justos?
\begin{enumerate}[label=\textbf{\Alph*)}]
\begin{multicols}{2}
  \item Albert Camus
  \item Jorge Luis Borges
  \item Federico García Lorca
  \item Pablo Neruda
  \item Octavio Paz
  \item César Vallejo
\end{multicols}
\end{enumerate}
\vspace{0.15cm}
\item ¿Quién es el autor de El diablo y Dios?
\begin{enumerate}[label=\textbf{\Alph*)}]
\begin{multicols}{2}
  \item Jean-Paul Sartre
  \item Jorge Luis Borges
  \item Federico García Lorca
  \item Pablo Neruda
  \item Octavio Paz
  \item César Vallejo
\end{multicols}
\end{enumerate}
\vspace{0.15cm}
\item ¿Quién es el autor de Las manos sucias?
\begin{enumerate}[label=\textbf{\Alph*)}]
\begin{multicols}{2}
  \item Jean-Paul Sartre
  \item Jorge Luis Borges
  \item Federico García Lorca
  \item Pablo Neruda
  \item Octavio Paz
  \item César Vallejo
\end{multicols}
\end{enumerate}
\vspace{0.15cm}
\end{enumerate}\n
\subsubsection*{SECCIÓN VI: POESÍA CONTEMPORÁNEA (231–270)}

\begin{enumerate}[label=\textbf{\arabic*.}]
\setcounter{enumi}{230}
\item ¿Qué es la poesía contemporánea?
\begin{enumerate}[label=\textbf{\Alph*)}]
\begin{multicols}{2}
  \item La poesía del siglo XX que rompe con las formas tradicionales
  \item Jorge Luis Borges
  \item Federico García Lorca
  \item Pablo Neruda
  \item Octavio Paz
  \item César Vallejo
\end{multicols}
\end{enumerate}
\vspace{0.15cm}
\item ¿Qué es el verso libre?
\begin{enumerate}[label=\textbf{\Alph*)}]
\begin{multicols}{2}
  \item Un verso sin rima ni métrica fija
  \item Jorge Luis Borges
  \item Federico García Lorca
  \item Pablo Neruda
  \item Octavio Paz
  \item César Vallejo
\end{multicols}
\end{enumerate}
\vspace{0.15cm}
\item ¿Quién es el autor de Poeta en Nueva York?
\begin{enumerate}[label=\textbf{\Alph*)}]
\begin{multicols}{2}
  \item Federico García Lorca
  \item Jorge Luis Borges
  \item Pablo Neruda
  \item Octavio Paz
  \item César Vallejo
  \item Vicente Huidobro
\end{multicols}
\end{enumerate}
\vspace{0.15cm}
\item ¿Quién es el autor de Veinte poemas de amor y una canción desesperada?
\begin{enumerate}[label=\textbf{\Alph*)}]
\begin{multicols}{2}
  \item Pablo Neruda
  \item Jorge Luis Borges
  \item Federico García Lorca
  \item Octavio Paz
  \item César Vallejo
  \item Vicente Huidobro
\end{multicols}
\end{enumerate}
\vspace{0.15cm}
\item ¿Quién es el autor de Trilce?
\begin{enumerate}[label=\textbf{\Alph*)}]
\begin{multicols}{2}
  \item César Vallejo
  \item Jorge Luis Borges
  \item Federico García Lorca
  \item Pablo Neruda
  \item Octavio Paz
  \item Vicente Huidobro
\end{multicols}
\end{enumerate}
\vspace{0.15cm}
\item ¿Quién es el autor de Libertad bajo palabra?
\begin{enumerate}[label=\textbf{\Alph*)}]
\begin{multicols}{2}
  \item Octavio Paz
  \item Jorge Luis Borges
  \item Federico García Lorca
  \item Pablo Neruda
  \item César Vallejo
  \item Vicente Huidobro
\end{multicols}
\end{enumerate}
\vspace{0.15cm}
\item ¿Quién es el autor de Altazor?
\begin{enumerate}[label=\textbf{\Alph*)}]
\begin{multicols}{2}
  \item Vicente Huidobro
  \item Jorge Luis Borges
  \item Federico García Lorca
  \item Pablo Neruda
  \item Octavio Paz
  \item César Vallejo
\end{multicols}
\end{enumerate}
\vspace{0.15cm}
\item ¿Quién es el autor de Fervor de Buenos Aires?
\begin{enumerate}[label=\textbf{\Alph*)}]
\begin{multicols}{2}
  \item Jorge Luis Borges
  \item Federico García Lorca
  \item Pablo Neruda
  \item Octavio Paz
  \item César Vallejo
  \item Vicente Huidobro
\end{multicols}
\end{enumerate}
\vspace{0.15cm}
\item ¿Quién es el autor de Romancero gitano?
\begin{enumerate}[label=\textbf{\Alph*)}]
\begin{multicols}{2}
  \item Federico García Lorca
  \item Jorge Luis Borges
  \item Pablo Neruda
  \item Octavio Paz
  \item César Vallejo
  \item Vicente Huidobro
\end{multicols}
\end{enumerate}
\vspace{0.15cm}
\item ¿Quién es el autor de Residencia en la tierra?
\begin{enumerate}[label=\textbf{\Alph*)}]
\begin{multicols}{2}
  \item Pablo Neruda
  \item Jorge Luis Borges
  \item Federico García Lorca
  \item Octavio Paz
  \item César Vallejo
  \item Vicente Huidobro
\end{multicols}
\end{enumerate}
\vspace{0.15cm}
\item ¿Quién es el autor de España, aparta de mí este cáliz?
\begin{enumerate}[label=\textbf{\Alph*)}]
\begin{multicols}{2}
  \item César Vallejo
  \item Jorge Luis Borges
  \item Federico García Lorca
  \item Pablo Neruda
  \item Octavio Paz
  \item Vicente Huidobro
\end{multicols}
\end{enumerate}
\vspace{0.15cm}
\item ¿Quién es el autor de Piedra de sol?
\begin{enumerate}[label=\textbf{\Alph*)}]
\begin{multicols}{2}
  \item Octavio Paz
  \item Jorge Luis Borges
  \item Federico García Lorca
  \item Pablo Neruda
  \item César Vallejo
  \item Vicente Huidobro
\end{multicols}
\end{enumerate}
\vspace{0.15cm}
\item ¿Quién es el autor de El rayo que no cesa?
\begin{enumerate}[label=\textbf{\Alph*)}]
\begin{multicols}{2}
  \item Miguel Hernández
  \item Jorge Luis Borges
  \item Federico García Lorca
  \item Pablo Neruda
  \item Octavio Paz
  \item César Vallejo
\end{multicols}
\end{enumerate}
\vspace{0.15cm}
\item ¿Quién es el autor de Viento del pueblo?
\begin{enumerate}[label=\textbf{\Alph*)}]
\begin{multicols}{2}
  \item Miguel Hernández
  \item Jorge Luis Borges
  \item Federico García Lorca
  \item Pablo Neruda
  \item Octavio Paz
  \item César Vallejo
\end{multicols}
\end{enumerate}
\vspace{0.15cm}
\item ¿Quién es el autor de Cancionero y romancero de ausencias?
\begin{enumerate}[label=\textbf{\Alph*)}]
\begin{multicols}{2}
  \item Miguel Hernández
  \item Jorge Luis Borges
  \item Federico García Lorca
  \item Pablo Neruda
  \item Octavio Paz
  \item César Vallejo
\end{multicols}
\end{enumerate}
\vspace{0.15cm}
\item ¿Quién es el autor de Hijos de la ira?
\begin{enumerate}[label=\textbf{\Alph*)}]
\begin{multicols}{2}
  \item Dámaso Alonso
  \item Jorge Luis Borges
  \item Federico García Lorca
  \item Pablo Neruda
  \item Octavio Paz
  \item César Vallejo
\end{multicols}
\end{enumerate}
\vspace{0.15cm}
\item ¿Quién es el autor de Sombra del paraíso?
\begin{enumerate}[label=\textbf{\Alph*)}]
\begin{multicols}{2}
  \item Vicente Aleixandre
  \item Jorge Luis Borges
  \item Federico García Lorca
  \item Pablo Neruda
  \item Octavio Paz
  \item César Vallejo
\end{multicols}
\end{enumerate}
\vspace{0.15cm}
\item ¿Quién es el autor de La destrucción o el amor?
\begin{enumerate}[label=\textbf{\Alph*)}]
\begin{multicols}{2}
  \item Vicente Aleixandre
  \item Jorge Luis Borges
  \item Federico García Lorca
  \item Pablo Neruda
  \item Octavio Paz
  \item César Vallejo
\end{multicols}
\end{enumerate}
\vspace{0.15cm}
\item ¿Quién es el autor de Espadas como labios?
\begin{enumerate}[label=\textbf{\Alph*)}]
\begin{multicols}{2}
  \item Vicente Aleixandre
  \item Jorge Luis Borges
  \item Federico García Lorca
  \item Pablo Neruda
  \item Octavio Paz
  \item César Vallejo
\end{multicols}
\end{enumerate}
\vspace{0.15cm}
\item ¿Quién es el autor de Poemas humanos?
\begin{enumerate}[label=\textbf{\Alph*)}]
\begin{multicols}{2}
  \item César Vallejo
  \item Jorge Luis Borges
  \item Federico García Lorca
  \item Pablo Neruda
  \item Octavio Paz
  \item Vicente Huidobro
\end{multicols}
\end{enumerate}
\vspace{0.15cm}
\item ¿Quién es el autor de Los heraldos negros?
\begin{enumerate}[label=\textbf{\Alph*)}]
\begin{multicols}{2}
  \item César Vallejo
  \item Jorge Luis Borges
  \item Federico García Lorca
  \item Pablo Neruda
  \item Octavio Paz
  \item Vicente Huidobro
\end{multicols}
\end{enumerate}
\vspace{0.15cm}
\item ¿Quién es el autor de Canto general?
\begin{enumerate}[label=\textbf{\Alph*)}]
\begin{multicols}{2}
  \item Pablo Neruda
  \item Jorge Luis Borges
  \item Federico García Lorca
  \item Octavio Paz
  \item César Vallejo
  \item Vicente Huidobro
\end{multicols}
\end{enumerate}
\vspace{0.15cm}
\item ¿Quién es el autor de Odas elementales?
\begin{enumerate}[label=\textbf{\Alph*)}]
\begin{multicols}{2}
  \item Pablo Neruda
  \item Jorge Luis Borges
  \item Federico García Lorca
  \item Octavio Paz
  \item César Vallejo
  \item Vicente Huidobro
\end{multicols}
\end{enumerate}
\vspace{0.15cm}
\item ¿Quién es el autor de Confieso que he vivido?
\begin{enumerate}[label=\textbf{\Alph*)}]
\begin{multicols}{2}
  \item Pablo Neruda
  \item Jorge Luis Borges
  \item Federico García Lorca
  \item Octavio Paz
  \item César Vallejo
  \item Vicente Huidobro
\end{multicols}
\end{enumerate}
\vspace{0.15cm}
\item ¿Quién es el autor de El laberinto de la soledad?
\begin{enumerate}[label=\textbf{\Alph*)}]
\begin{multicols}{2}
  \item Octavio Paz
  \item Jorge Luis Borges
  \item Federico García Lorca
  \item Pablo Neruda
  \item César Vallejo
  \item Vicente Huidobro
\end{multicols}
\end{enumerate}
\vspace{0.15cm}
\item ¿Quién es el autor de Árbol adentro?
\begin{enumerate}[label=\textbf{\Alph*)}]
\begin{multicols}{2}
  \item Octavio Paz
  \item Jorge Luis Borges
  \item Federico García Lorca
  \item Pablo Neruda
  \item César Vallejo
  \item Vicente Huidobro
\end{multicols}
\end{enumerate}
\vspace{0.15cm}
\item ¿Quién es el autor de Blanco?
\begin{enumerate}[label=\textbf{\Alph*)}]
\begin{multicols}{2}
  \item Octavio Paz
  \item Jorge Luis Borges
  \item Federico García Lorca
  \item Pablo Neruda
  \item César Vallejo
  \item Vicente Huidobro
\end{multicols}
\end{enumerate}
\vspace{0.15cm}
\item ¿Quién es el autor de Ladera este?
\begin{enumerate}[label=\textbf{\Alph*)}]
\begin{multicols}{2}
  \item Octavio Paz
  \item Jorge Luis Borges
  \item Federico García Lorca
  \item Pablo Neruda
  \item César Vallejo
  \item Vicente Huidobro
\end{multicols}
\end{enumerate}
\vspace{0.15cm}
\item ¿Quién es el autor de Salamandra?
\begin{enumerate}[label=\textbf{\Alph*)}]
\begin{multicols}{2}
  \item Octavio Paz
  \item Jorge Luis Borges
  \item Federico García Lorca
  \item Pablo Neruda
  \item César Vallejo
  \item Vicente Huidobro
\end{multicols}
\end{enumerate}
\vspace{0.15cm}
\item ¿Quién es el autor de Pasado en claro?
\begin{enumerate}[label=\textbf{\Alph*)}]
\begin{multicols}{2}
  \item Octavio Paz
  \item Jorge Luis Borges
  \item Federico García Lorca
  \item Pablo Neruda
  \item César Vallejo
  \item Vicente Huidobro
\end{multicols}
\end{enumerate}
\vspace{0.15cm}
\item ¿Qué es la poesía contemporánea?
\begin{enumerate}[label=\textbf{\Alph*)}]
\begin{multicols}{2}
  \item La poesía del siglo XX que rompe con las formas tradicionales
  \item Jorge Luis Borges
  \item Federico García Lorca
  \item Pablo Neruda
  \item Octavio Paz
  \item César Vallejo
\end{multicols}
\end{enumerate}
\vspace{0.15cm}
\item ¿Qué es el verso libre?
\begin{enumerate}[label=\textbf{\Alph*)}]
\begin{multicols}{2}
  \item Un verso sin rima ni métrica fija
  \item Jorge Luis Borges
  \item Federico García Lorca
  \item Pablo Neruda
  \item Octavio Paz
  \item César Vallejo
\end{multicols}
\end{enumerate}
\vspace{0.15cm}
\item ¿Quién es el autor de Poeta en Nueva York?
\begin{enumerate}[label=\textbf{\Alph*)}]
\begin{multicols}{2}
  \item Federico García Lorca
  \item Jorge Luis Borges
  \item Pablo Neruda
  \item Octavio Paz
  \item César Vallejo
  \item Vicente Huidobro
\end{multicols}
\end{enumerate}
\vspace{0.15cm}
\item ¿Quién es el autor de Veinte poemas de amor y una canción desesperada?
\begin{enumerate}[label=\textbf{\Alph*)}]
\begin{multicols}{2}
  \item Pablo Neruda
  \item Jorge Luis Borges
  \item Federico García Lorca
  \item Octavio Paz
  \item César Vallejo
  \item Vicente Huidobro
\end{multicols}
\end{enumerate}
\vspace{0.15cm}
\item ¿Quién es el autor de Trilce?
\begin{enumerate}[label=\textbf{\Alph*)}]
\begin{multicols}{2}
  \item César Vallejo
  \item Jorge Luis Borges
  \item Federico García Lorca
  \item Pablo Neruda
  \item Octavio Paz
  \item Vicente Huidobro
\end{multicols}
\end{enumerate}
\vspace{0.15cm}
\item ¿Quién es el autor de Libertad bajo palabra?
\begin{enumerate}[label=\textbf{\Alph*)}]
\begin{multicols}{2}
  \item Octavio Paz
  \item Jorge Luis Borges
  \item Federico García Lorca
  \item Pablo Neruda
  \item César Vallejo
  \item Vicente Huidobro
\end{multicols}
\end{enumerate}
\vspace{0.15cm}
\item ¿Quién es el autor de Altazor?
\begin{enumerate}[label=\textbf{\Alph*)}]
\begin{multicols}{2}
  \item Vicente Huidobro
  \item Jorge Luis Borges
  \item Federico García Lorca
  \item Pablo Neruda
  \item Octavio Paz
  \item César Vallejo
\end{multicols}
\end{enumerate}
\vspace{0.15cm}
\item ¿Quién es el autor de Fervor de Buenos Aires?
\begin{enumerate}[label=\textbf{\Alph*)}]
\begin{multicols}{2}
  \item Jorge Luis Borges
  \item Federico García Lorca
  \item Pablo Neruda
  \item Octavio Paz
  \item César Vallejo
  \item Vicente Huidobro
\end{multicols}
\end{enumerate}
\vspace{0.15cm}
\item ¿Quién es el autor de Romancero gitano?
\begin{enumerate}[label=\textbf{\Alph*)}]
\begin{multicols}{2}
  \item Federico García Lorca
  \item Jorge Luis Borges
  \item Pablo Neruda
  \item Octavio Paz
  \item César Vallejo
  \item Vicente Huidobro
\end{multicols}
\end{enumerate}
\vspace{0.15cm}
\item ¿Quién es el autor de Residencia en la tierra?
\begin{enumerate}[label=\textbf{\Alph*)}]
\begin{multicols}{2}
  \item Pablo Neruda
  \item Jorge Luis Borges
  \item Federico García Lorca
  \item Octavio Paz
  \item César Vallejo
  \item Vicente Huidobro
\end{multicols}
\end{enumerate}
\vspace{0.15cm}
\end{enumerate}\n
\subsubsection*{SECCIÓN VII: TEMAS Y CARACTERÍSTICAS (271–300)}

\begin{enumerate}[label=\textbf{\arabic*.}]
\setcounter{enumi}{270}
\item ¿Cuál es el tema central de la literatura contemporánea?
\begin{enumerate}[label=\textbf{\Alph*)}]
\begin{multicols}{2}
  \item La crisis de la modernidad y la condición humana
  \item Jorge Luis Borges
  \item Federico García Lorca
  \item Pablo Neruda
  \item Octavio Paz
  \item César Vallejo
\end{multicols}
\end{enumerate}
\vspace{0.15cm}
\item ¿Qué es la alienación en la literatura contemporánea?
\begin{enumerate}[label=\textbf{\Alph*)}]
\begin{multicols}{2}
  \item La representación del individuo aislado y deshumanizado
  \item Jorge Luis Borges
  \item Federico García Lorca
  \item Pablo Neruda
  \item Octavio Paz
  \item César Vallejo
\end{multicols}
\end{enumerate}
\vspace{0.15cm}
\item ¿Qué es la incomunicación en la literatura contemporánea?
\begin{enumerate}[label=\textbf{\Alph*)}]
\begin{multicols}{2}
  \item La representación de la dificultad de comunicación entre las personas
  \item Jorge Luis Borges
  \item Federico García Lorca
  \item Pablo Neruda
  \item Octavio Paz
  \item César Vallejo
\end{multicols}
\end{enumerate}
\vspace{0.15cm}
\item ¿Qué es la angustia existencial?
\begin{enumerate}[label=\textbf{\Alph*)}]
\begin{multicols}{2}
  \item La representación de la preocupación por el sentido de la existencia
  \item Jorge Luis Borges
  \item Federico García Lorca
  \item Pablo Neruda
  \item Octavio Paz
  \item César Vallejo
\end{multicols}
\end{enumerate}
\vspace{0.15cm}
\item ¿Qué es la libertad en el existencialismo?
\begin{enumerate}[label=\textbf{\Alph*)}]
\begin{multicols}{2}
  \item La capacidad del individuo de elegir y responsabilizarse
  \item Jorge Luis Borges
  \item Federico García Lorca
  \item Pablo Neruda
  \item Octavio Paz
  \item César Vallejo
\end{multicols}
\end{enumerate}
\vspace{0.15cm}
\item ¿Qué es la responsabilidad en el existencialismo?
\begin{enumerate}[label=\textbf{\Alph*)}]
\begin{multicols}{2}
  \item La consecuencia de la libertad de elección
  \item Jorge Luis Borges
  \item Federico García Lorca
  \item Pablo Neruda
  \item Octavio Paz
  \item César Vallejo
\end{multicols}
\end{enumerate}
\vspace{0.15cm}
\item ¿Qué es el compromiso en la literatura contemporánea?
\begin{enumerate}[label=\textbf{\Alph*)}]
\begin{multicols}{2}
  \item La implicación del escritor en los problemas sociales y políticos
  \item Jorge Luis Borges
  \item Federico García Lorca
  \item Pablo Neruda
  \item Octavio Paz
  \item César Vallejo
\end{multicols}
\end{enumerate}
\vspace{0.15cm}
\item ¿Qué es la intertextualidad?
\begin{enumerate}[label=\textbf{\Alph*)}]
\begin{multicols}{2}
  \item La relación entre un texto y otros textos
  \item Jorge Luis Borges
  \item Federico García Lorca
  \item Pablo Neruda
  \item Octavio Paz
  \item César Vallejo
\end{multicols}
\end{enumerate}
\vspace{0.15cm}
\item ¿Qué es la metaficción?
\begin{enumerate}[label=\textbf{\Alph*)}]
\begin{multicols}{2}
  \item La reflexión sobre la propia narración dentro de la narración
  \item Jorge Luis Borges
  \item Federico García Lorca
  \item Pablo Neruda
  \item Octavio Paz
  \item César Vallejo
\end{multicols}
\end{enumerate}
\vspace{0.15cm}
\item ¿Qué es la fragmentación en la literatura contemporánea?
\begin{enumerate}[label=\textbf{\Alph*)}]
\begin{multicols}{2}
  \item La ruptura de la linealidad narrativa
  \item Jorge Luis Borges
  \item Federico García Lorca
  \item Pablo Neruda
  \item Octavio Paz
  \item César Vallejo
\end{multicols}
\end{enumerate}
\vspace{0.15cm}
\item ¿Qué es el perspectivismo en la literatura contemporánea?
\begin{enumerate}[label=\textbf{\Alph*)}]
\begin{multicols}{2}
  \item La multiplicidad de puntos de vista en la narración
  \item Jorge Luis Borges
  \item Federico García Lorca
  \item Pablo Neruda
  \item Octavio Paz
  \item César Vallejo
\end{multicols}
\end{enumerate}
\vspace{0.15cm}
\item ¿Qué es el tiempo circular en la literatura contemporánea?
\begin{enumerate}[label=\textbf{\Alph*)}]
\begin{multicols}{2}
  \item La repetición de los acontecimientos
  \item Jorge Luis Borges
  \item Federico García Lorca
  \item Pablo Neruda
  \item Octavio Paz
  \item César Vallejo
\end{multicols}
\end{enumerate}
\vspace{0.15cm}
\item ¿Qué es el realismo mágico?
\begin{enumerate}[label=\textbf{\Alph*)}]
\begin{multicols}{2}
  \item La mezcla de lo real y lo fantástico en la narración
  \item Jorge Luis Borges
  \item Federico García Lorca
  \item Pablo Neruda
  \item Octavio Paz
  \item César Vallejo
\end{multicols}
\end{enumerate}
\vspace{0.15cm}
\item ¿Qué es lo real maravilloso?
\begin{enumerate}[label=\textbf{\Alph*)}]
\begin{multicols}{2}
  \item La presencia de lo extraordinario en la realidad cotidiana
  \item Jorge Luis Borges
  \item Federico García Lorca
  \item Pablo Neruda
  \item Octavio Paz
  \item César Vallejo
\end{multicols}
\end{enumerate}
\vspace{0.15cm}
\item ¿Qué es la literatura del boom latinoamericano?
\begin{enumerate}[label=\textbf{\Alph*)}]
\begin{multicols}{2}
  \item El fenómeno editorial y literario de los años 60 y 70
  \item Jorge Luis Borges
  \item Federico García Lorca
  \item Pablo Neruda
  \item Octavio Paz
  \item César Vallejo
\end{multicols}
\end{enumerate}
\vspace{0.15cm}
\item ¿Quiénes son los principales representantes del boom?
\begin{enumerate}[label=\textbf{\Alph*)}]
\begin{multicols}{2}
  \item García Márquez, Cortázar, Vargas Llosa y Fuentes
  \item Jorge Luis Borges
  \item Federico García Lorca
  \item Pablo Neruda
  \item Octavio Paz
  \item César Vallejo
\end{multicols}
\end{enumerate}
\vspace{0.15cm}
\item ¿Qué es la novela total?
\begin{enumerate}[label=\textbf{\Alph*)}]
\begin{multicols}{2}
  \item Una novela que abarca múltiples aspectos de la realidad
  \item Jorge Luis Borges
  \item Federico García Lorca
  \item Pablo Neruda
  \item Octavio Paz
  \item César Vallejo
\end{multicols}
\end{enumerate}
\vspace{0.15cm}
\item ¿Qué es la literatura posmoderna?
\begin{enumerate}[label=\textbf{\Alph*)}]
\begin{multicols}{2}
  \item La literatura que cuestiona los grandes relatos y la verdad única
  \item Jorge Luis Borges
  \item Federico García Lorca
  \item Pablo Neruda
  \item Octavio Paz
  \item César Vallejo
\end{multicols}
\end{enumerate}
\vspace{0.15cm}
\item ¿Qué es la desconstrucción en la literatura?
\begin{enumerate}[label=\textbf{\Alph*)}]
\begin{multicols}{2}
  \item La crítica a las oposiciones jerárquicas del pensamiento occidental
  \item Jorge Luis Borges
  \item Federico García Lorca
  \item Pablo Neruda
  \item Octavio Paz
  \item César Vallejo
\end{multicols}
\end{enumerate}
\vspace{0.15cm}
\item ¿Qué es la hipertextualidad?
\begin{enumerate}[label=\textbf{\Alph*)}]
\begin{multicols}{2}
  \item La relación entre un texto y otro texto anterior
  \item Jorge Luis Borges
  \item Federico García Lorca
  \item Pablo Neruda
  \item Octavio Paz
  \item César Vallejo
\end{multicols}
\end{enumerate}
\vspace{0.15cm}
\item ¿Qué es la parodia en la literatura contemporánea?
\begin{enumerate}[label=\textbf{\Alph*)}]
\begin{multicols}{2}
  \item La imitación burlesca de un texto o género
  \item Jorge Luis Borges
  \item Federico García Lorca
  \item Pablo Neruda
  \item Octavio Paz
  \item César Vallejo
\end{multicols}
\end{enumerate}
\vspace{0.15cm}
\item ¿Qué es el pastiche en la literatura contemporánea?
\begin{enumerate}[label=\textbf{\Alph*)}]
\begin{multicols}{2}
  \item La mezcla de estilos y géneros diversos
  \item Jorge Luis Borges
  \item Federico García Lorca
  \item Pablo Neruda
  \item Octavio Paz
  \item César Vallejo
\end{multicols}
\end{enumerate}
\vspace{0.15cm}
\item ¿Qué es la ironía en la literatura contemporánea?
\begin{enumerate}[label=\textbf{\Alph*)}]
\begin{multicols}{2}
  \item La expresión de lo contrario de lo que se quiere decir
  \item Jorge Luis Borges
  \item Federico García Lorca
  \item Pablo Neruda
  \item Octavio Paz
  \item César Vallejo
\end{multicols}
\end{enumerate}
\vspace{0.15cm}
\item ¿Qué es el humor negro en la literatura contemporánea?
\begin{enumerate}[label=\textbf{\Alph*)}]
\begin{multicols}{2}
  \item El humor que trata temas tabúes o dolorosos
  \item Jorge Luis Borges
  \item Federico García Lorca
  \item Pablo Neruda
  \item Octavio Paz
  \item César Vallejo
\end{multicols}
\end{enumerate}
\vspace{0.15cm}
\item ¿Qué es la literatura del absurdo?
\begin{enumerate}[label=\textbf{\Alph*)}]
\begin{multicols}{2}
  \item La literatura que representa la falta de sentido de la existencia
  \item Jorge Luis Borges
  \item Federico García Lorca
  \item Pablo Neruda
  \item Octavio Paz
  \item César Vallejo
\end{multicols}
\end{enumerate}
\vspace{0.15cm}
\item ¿Qué es la literatura existencialista?
\begin{enumerate}[label=\textbf{\Alph*)}]
\begin{multicols}{2}
  \item La literatura que reflexiona sobre la existencia humana
  \item Jorge Luis Borges
  \item Federico García Lorca
  \item Pablo Neruda
  \item Octavio Paz
  \item César Vallejo
\end{multicols}
\end{enumerate}
\vspace{0.15cm}
\item ¿Qué es la literatura comprometida?
\begin{enumerate}[label=\textbf{\Alph*)}]
\begin{multicols}{2}
  \item La literatura que se implica en los problemas sociales y políticos
  \item Jorge Luis Borges
  \item Federico García Lorca
  \item Pablo Neruda
  \item Octavio Paz
  \item César Vallejo
\end{multicols}
\end{enumerate}
\vspace{0.15cm}
\item ¿Qué es la literatura experimental?
\begin{enumerate}[label=\textbf{\Alph*)}]
\begin{multicols}{2}
  \item La literatura que rompe con las formas tradicionales
  \item Jorge Luis Borges
  \item Federico García Lorca
  \item Pablo Neruda
  \item Octavio Paz
  \item César Vallejo
\end{multicols}
\end{enumerate}
\vspace{0.15cm}
\item ¿Qué es la literatura de vanguardia?
\begin{enumerate}[label=\textbf{\Alph*)}]
\begin{multicols}{2}
  \item La literatura que rompe con las normas y tradiciones
  \item Jorge Luis Borges
  \item Federico García Lorca
  \item Pablo Neruda
  \item Octavio Paz
  \item César Vallejo
\end{multicols}
\end{enumerate}
\vspace{0.15cm}
\item ¿Qué es la literatura contemporánea?
\begin{enumerate}[label=\textbf{\Alph*)}]
\begin{multicols}{2}
  \item La literatura del siglo XX y XXI que rompe con las formas tradicionales
  \item Jorge Luis Borges
  \item Federico García Lorca
  \item Pablo Neruda
  \item Octavio Paz
  \item César Vallejo
\end{multicols}
\end{enumerate}
\vspace{0.15cm}
\end{enumerate}\n
\vspace{0.4cm}
\begin{tcolorbox}[colback=gray!5!white,colframe=gray!60!black,title=\textbf{SOLUCIONARIO -- 3.3 Literatura del Siglo Contemporáneo}]
\begin{multicols}{4}
\begin{enumerate}[label=\textbf{\arabic*.}]
  \item \textbf{A}
  \item \textbf{A}
  \item \textbf{A}
  \item \textbf{A}
  \item \textbf{A}
  \item \textbf{A}
  \item \textbf{A}
  \item \textbf{A}
  \item \textbf{A}
  \item \textbf{A}
  \item \textbf{A}
  \item \textbf{A}
  \item \textbf{A}
  \item \textbf{A}
  \item \textbf{A}
  \item \textbf{A}
  \item \textbf{A}
  \item \textbf{A}
  \item \textbf{A}
  \item \textbf{A}
  \item \textbf{A}
  \item \textbf{A}
  \item \textbf{A}
  \item \textbf{A}
  \item \textbf{A}
  \item \textbf{A}
  \item \textbf{A}
  \item \textbf{A}
  \item \textbf{A}
  \item \textbf{A}
  \item \textbf{A}
  \item \textbf{A}
  \item \textbf{A}
  \item \textbf{A}
  \item \textbf{A}
  \item \textbf{A}
  \item \textbf{A}
  \item \textbf{A}
  \item \textbf{A}
  \item \textbf{A}
  \item \textbf{A}
  \item \textbf{A}
  \item \textbf{A}
  \item \textbf{A}
  \item \textbf{A}
  \item \textbf{A}
  \item \textbf{A}
  \item \textbf{A}
  \item \textbf{A}
  \item \textbf{A}
  \item \textbf{A}
  \item \textbf{A}
  \item \textbf{A}
  \item \textbf{A}
  \item \textbf{A}
  \item \textbf{A}
  \item \textbf{A}
  \item \textbf{A}
  \item \textbf{A}
  \item \textbf{A}
  \item \textbf{A}
  \item \textbf{A}
  \item \textbf{A}
  \item \textbf{A}
  \item \textbf{A}
  \item \textbf{A}
  \item \textbf{A}
  \item \textbf{A}
  \item \textbf{A}
  \item \textbf{A}
  \item \textbf{A}
  \item \textbf{A}
  \item \textbf{A}
  \item \textbf{A}
  \item \textbf{A}
  \item \textbf{A}
  \item \textbf{A}
  \item \textbf{A}
  \item \textbf{A}
  \item \textbf{A}
  \item \textbf{A}
  \item \textbf{A}
  \item \textbf{A}
  \item \textbf{A}
  \item \textbf{A}
  \item \textbf{A}
  \item \textbf{A}
  \item \textbf{A}
  \item \textbf{A}
  \item \textbf{A}
  \item \textbf{A}
  \item \textbf{A}
  \item \textbf{A}
  \item \textbf{A}
  \item \textbf{A}
  \item \textbf{A}
  \item \textbf{A}
  \item \textbf{A}
  \item \textbf{A}
  \item \textbf{A}
  \item \textbf{A}
  \item \textbf{A}
  \item \textbf{A}
  \item \textbf{A}
  \item \textbf{A}
  \item \textbf{A}
  \item \textbf{A}
  \item \textbf{A}
  \item \textbf{A}
  \item \textbf{A}
  \item \textbf{A}
  \item \textbf{A}
  \item \textbf{A}
  \item \textbf{A}
  \item \textbf{A}
  \item \textbf{A}
  \item \textbf{A}
  \item \textbf{A}
  \item \textbf{A}
  \item \textbf{A}
  \item \textbf{A}
  \item \textbf{A}
  \item \textbf{A}
  \item \textbf{A}
  \item \textbf{A}
  \item \textbf{A}
  \item \textbf{A}
  \item \textbf{A}
  \item \textbf{A}
  \item \textbf{A}
  \item \textbf{A}
  \item \textbf{A}
  \item \textbf{A}
  \item \textbf{A}
  \item \textbf{A}
  \item \textbf{A}
  \item \textbf{A}
  \item \textbf{A}
  \item \textbf{A}
  \item \textbf{A}
  \item \textbf{A}
  \item \textbf{A}
  \item \textbf{A}
  \item \textbf{A}
  \item \textbf{A}
  \item \textbf{A}
  \item \textbf{A}
  \item \textbf{A}
  \item \textbf{A}
  \item \textbf{A}
  \item \textbf{A}
  \item \textbf{A}
  \item \textbf{A}
  \item \textbf{A}
  \item \textbf{A}
  \item \textbf{A}
  \item \textbf{A}
  \item \textbf{A}
  \item \textbf{A}
  \item \textbf{A}
  \item \textbf{A}
  \item \textbf{A}
  \item \textbf{A}
  \item \textbf{A}
  \item \textbf{A}
  \item \textbf{A}
  \item \textbf{A}
  \item \textbf{A}
  \item \textbf{A}
  \item \textbf{A}
  \item \textbf{A}
  \item \textbf{A}
  \item \textbf{A}
  \item \textbf{A}
  \item \textbf{A}
  \item \textbf{A}
  \item \textbf{A}
  \item \textbf{A}
  \item \textbf{A}
  \item \textbf{A}
  \item \textbf{A}
  \item \textbf{A}
  \item \textbf{A}
  \item \textbf{A}
  \item \textbf{A}
  \item \textbf{A}
  \item \textbf{A}
  \item \textbf{A}
  \item \textbf{A}
  \item \textbf{A}
  \item \textbf{A}
  \item \textbf{A}
  \item \textbf{A}
  \item \textbf{A}
  \item \textbf{A}
  \item \textbf{A}
  \item \textbf{A}
  \item \textbf{A}
  \item \textbf{A}
  \item \textbf{A}
  \item \textbf{A}
  \item \textbf{A}
  \item \textbf{A}
  \item \textbf{A}
  \item \textbf{A}
  \item \textbf{A}
  \item \textbf{A}
  \item \textbf{A}
  \item \textbf{A}
  \item \textbf{A}
  \item \textbf{A}
  \item \textbf{A}
  \item \textbf{A}
  \item \textbf{A}
  \item \textbf{A}
  \item \textbf{A}
  \item \textbf{A}
  \item \textbf{A}
  \item \textbf{A}
  \item \textbf{A}
  \item \textbf{A}
  \item \textbf{A}
  \item \textbf{A}
  \item \textbf{A}
  \item \textbf{A}
  \item \textbf{A}
  \item \textbf{A}
  \item \textbf{A}
  \item \textbf{A}
  \item \textbf{A}
  \item \textbf{A}
  \item \textbf{A}
  \item \textbf{A}
  \item \textbf{A}
  \item \textbf{A}
  \item \textbf{A}
  \item \textbf{A}
  \item \textbf{A}
  \item \textbf{A}
  \item \textbf{A}
  \item \textbf{A}
  \item \textbf{A}
  \item \textbf{A}
  \item \textbf{A}
  \item \textbf{A}
  \item \textbf{A}
  \item \textbf{A}
  \item \textbf{A}
  \item \textbf{A}
  \item \textbf{A}
  \item \textbf{A}
  \item \textbf{A}
  \item \textbf{A}
  \item \textbf{A}
  \item \textbf{A}
  \item \textbf{A}
  \item \textbf{A}
  \item \textbf{A}
  \item \textbf{A}
  \item \textbf{A}
  \item \textbf{A}
  \item \textbf{A}
  \item \textbf{A}
  \item \textbf{A}
  \item \textbf{A}
  \item \textbf{A}
  \item \textbf{A}
  \item \textbf{A}
  \item \textbf{A}
  \item \textbf{A}
  \item \textbf{A}
  \item \textbf{A}
  \item \textbf{A}
  \item \textbf{A}
  \item \textbf{A}
  \item \textbf{A}
  \item \textbf{A}
  \item \textbf{A}
  \item \textbf{A}
  \item \textbf{A}
  \item \textbf{A}
  \item \textbf{A}
  \item \textbf{A}
  \item \textbf{A}
  \item \textbf{A}
  \item \textbf{A}
  \item \textbf{A}
  \item \textbf{A}
  \item \textbf{A}
  \item \textbf{A}
  \item \textbf{A}
  \item \textbf{A}
  \item \textbf{A}
  \item \textbf{A}
  \item \textbf{A}
  \item \textbf{A}
  \item \textbf{A}
  \item \textbf{A}
  \item \textbf{A}
  \item \textbf{A}
\end{enumerate}
\end{multicols}
\end{tcolorbox}
\newpage\n
\subsection{3.4 Literatura Latinoamericana}

\subsubsection*{SECCIÓN I: LITERATURA PRECOLOMBINA (1–40)}

\begin{enumerate}[label=\textbf{\arabic*.}]
\setcounter{enumi}{0}
\item ¿Qué es la literatura precolombina?
\begin{enumerate}[label=\textbf{\Alph*)}]
\begin{multicols}{2}
  \item La literatura producida antes de la llegada de Colón a América
  \item Gabriel García Márquez
  \item Julio Cortázar
  \item Mario Vargas Llosa
  \item Carlos Fuentes
  \item Juan Rulfo
\end{multicols}
\end{enumerate}
\vspace{0.15cm}
\item ¿Cuál es una característica de la literatura precolombina?
\begin{enumerate}[label=\textbf{\Alph*)}]
\begin{multicols}{2}
  \item Ser de tradición oral y anónima
  \item Gabriel García Márquez
  \item Julio Cortázar
  \item Mario Vargas Llosa
  \item Carlos Fuentes
  \item Juan Rulfo
\end{multicols}
\end{enumerate}
\vspace{0.15cm}
\item ¿Qué civilización produjo el Popol Vuh?
\begin{enumerate}[label=\textbf{\Alph*)}]
\begin{multicols}{2}
  \item Maya
  \item Gabriel García Márquez
  \item Julio Cortázar
  \item Mario Vargas Llosa
  \item Carlos Fuentes
  \item Juan Rulfo
\end{multicols}
\end{enumerate}
\vspace{0.15cm}
\item ¿Qué es el Popol Vuh?
\begin{enumerate}[label=\textbf{\Alph*)}]
\begin{multicols}{2}
  \item Un libro sagrado maya que narra la creación del mundo
  \item Gabriel García Márquez
  \item Julio Cortázar
  \item Mario Vargas Llosa
  \item Carlos Fuentes
  \item Juan Rulfo
\end{multicols}
\end{enumerate}
\vspace{0.15cm}
\item ¿Qué civilización produjo el Chilam Balam?
\begin{enumerate}[label=\textbf{\Alph*)}]
\begin{multicols}{2}
  \item Maya
  \item Gabriel García Márquez
  \item Julio Cortázar
  \item Mario Vargas Llosa
  \item Carlos Fuentes
  \item Juan Rulfo
\end{multicols}
\end{enumerate}
\vspace{0.15cm}
\item ¿Qué es el Chilam Balam?
\begin{enumerate}[label=\textbf{\Alph*)}]
\begin{multicols}{2}
  \item Un conjunto de libros mayas con profecías y conocimientos
  \item Gabriel García Márquez
  \item Julio Cortázar
  \item Mario Vargas Llosa
  \item Carlos Fuentes
  \item Juan Rulfo
\end{multicols}
\end{enumerate}
\vspace{0.15cm}
\item ¿Qué civilización produjo los Cantares mexicanos?
\begin{enumerate}[label=\textbf{\Alph*)}]
\begin{multicols}{2}
  \item Azteca
  \item Gabriel García Márquez
  \item Julio Cortázar
  \item Mario Vargas Llosa
  \item Carlos Fuentes
  \item Juan Rulfo
\end{multicols}
\end{enumerate}
\vspace{0.15cm}
\item ¿Qué es un "cuicatl" en la literatura azteca?
\begin{enumerate}[label=\textbf{\Alph*)}]
\begin{multicols}{2}
  \item Un canto o poema
  \item Gabriel García Márquez
  \item Julio Cortázar
  \item Mario Vargas Llosa
  \item Carlos Fuentes
  \item Juan Rulfo
\end{multicols}
\end{enumerate}
\vspace{0.15cm}
\item ¿Quién fue Nezahualcóyotl?
\begin{enumerate}[label=\textbf{\Alph*)}]
\begin{multicols}{2}
  \item Un poeta y rey azteca
  \item Gabriel García Márquez
  \item Julio Cortázar
  \item Mario Vargas Llosa
  \item Carlos Fuentes
  \item Juan Rulfo
\end{multicols}
\end{enumerate}
\vspace{0.15cm}
\item ¿Qué civilización produjo el Ollantay?
\begin{enumerate}[label=\textbf{\Alph*)}]
\begin{multicols}{2}
  \item Inca
  \item Gabriel García Márquez
  \item Julio Cortázar
  \item Mario Vargas Llosa
  \item Carlos Fuentes
  \item Juan Rulfo
\end{multicols}
\end{enumerate}
\vspace{0.15cm}
\item ¿Qué es el Ollantay?
\begin{enumerate}[label=\textbf{\Alph*)}]
\begin{multicols}{2}
  \item Un drama inca de tema amoroso y heroico
  \item Gabriel García Márquez
  \item Julio Cortázar
  \item Mario Vargas Llosa
  \item Carlos Fuentes
  \item Juan Rulfo
\end{multicols}
\end{enumerate}
\vspace{0.15cm}
\item ¿Qué es un "harawi" en la literatura inca?
\begin{enumerate}[label=\textbf{\Alph*)}]
\begin{multicols}{2}
  \item Un canto o poema lírico
  \item Gabriel García Márquez
  \item Julio Cortázar
  \item Mario Vargas Llosa
  \item Carlos Fuentes
  \item Juan Rulfo
\end{multicols}
\end{enumerate}
\vspace{0.15cm}
\item ¿Qué es un "haylli" en la literatura inca?
\begin{enumerate}[label=\textbf{\Alph*)}]
\begin{multicols}{2}
  \item Un canto de triunfo o guerra
  \item Gabriel García Márquez
  \item Julio Cortázar
  \item Mario Vargas Llosa
  \item Carlos Fuentes
  \item Juan Rulfo
\end{multicols}
\end{enumerate}
\vspace{0.15cm}
\item ¿Qué es un "aylli" en la literatura inca?
\begin{enumerate}[label=\textbf{\Alph*)}]
\begin{multicols}{2}
  \item Un canto de amor o alegría
  \item Gabriel García Márquez
  \item Julio Cortázar
  \item Mario Vargas Llosa
  \item Carlos Fuentes
  \item Juan Rulfo
\end{multicols}
\end{enumerate}
\vspace{0.15cm}
\item ¿Qué es un "umu" en la literatura inca?
\begin{enumerate}[label=\textbf{\Alph*)}]
\begin{multicols}{2}
  \item Un canto fúnebre
  \item Gabriel García Márquez
  \item Julio Cortázar
  \item Mario Vargas Llosa
  \item Carlos Fuentes
  \item Juan Rulfo
\end{multicols}
\end{enumerate}
\vspace{0.15cm}
\item ¿Qué es un "taqui" en la literatura inca?
\begin{enumerate}[label=\textbf{\Alph*)}]
\begin{multicols}{2}
  \item Un canto o danza
  \item Gabriel García Márquez
  \item Julio Cortázar
  \item Mario Vargas Llosa
  \item Carlos Fuentes
  \item Juan Rulfo
\end{multicols}
\end{enumerate}
\vspace{0.15cm}
\item ¿Qué civilización produjo el Rabinal Achí?
\begin{enumerate}[label=\textbf{\Alph*)}]
\begin{multicols}{2}
  \item Maya
  \item Gabriel García Márquez
  \item Julio Cortázar
  \item Mario Vargas Llosa
  \item Carlos Fuentes
  \item Juan Rulfo
\end{multicols}
\end{enumerate}
\vspace{0.15cm}
\item ¿Qué es el Rabinal Achí?
\begin{enumerate}[label=\textbf{\Alph*)}]
\begin{multicols}{2}
  \item Un drama maya de tema guerrero
  \item Gabriel García Márquez
  \item Julio Cortázar
  \item Mario Vargas Llosa
  \item Carlos Fuentes
  \item Juan Rulfo
\end{multicols}
\end{enumerate}
\vspace{0.15cm}
\item ¿Qué es la literatura precolombina?
\begin{enumerate}[label=\textbf{\Alph*)}]
\begin{multicols}{2}
  \item La literatura producida antes de la llegada de Colón a América
  \item Gabriel García Márquez
  \item Julio Cortázar
  \item Mario Vargas Llosa
  \item Carlos Fuentes
  \item Juan Rulfo
\end{multicols}
\end{enumerate}
\vspace{0.15cm}
\item ¿Cuál es una característica de la literatura precolombina?
\begin{enumerate}[label=\textbf{\Alph*)}]
\begin{multicols}{2}
  \item Ser de tradición oral y anónima
  \item Gabriel García Márquez
  \item Julio Cortázar
  \item Mario Vargas Llosa
  \item Carlos Fuentes
  \item Juan Rulfo
\end{multicols}
\end{enumerate}
\vspace{0.15cm}
\item ¿Qué civilización produjo el Popol Vuh?
\begin{enumerate}[label=\textbf{\Alph*)}]
\begin{multicols}{2}
  \item Maya
  \item Gabriel García Márquez
  \item Julio Cortázar
  \item Mario Vargas Llosa
  \item Carlos Fuentes
  \item Juan Rulfo
\end{multicols}
\end{enumerate}
\vspace{0.15cm}
\item ¿Qué es el Popol Vuh?
\begin{enumerate}[label=\textbf{\Alph*)}]
\begin{multicols}{2}
  \item Un libro sagrado maya que narra la creación del mundo
  \item Gabriel García Márquez
  \item Julio Cortázar
  \item Mario Vargas Llosa
  \item Carlos Fuentes
  \item Juan Rulfo
\end{multicols}
\end{enumerate}
\vspace{0.15cm}
\item ¿Qué civilización produjo el Chilam Balam?
\begin{enumerate}[label=\textbf{\Alph*)}]
\begin{multicols}{2}
  \item Maya
  \item Gabriel García Márquez
  \item Julio Cortázar
  \item Mario Vargas Llosa
  \item Carlos Fuentes
  \item Juan Rulfo
\end{multicols}
\end{enumerate}
\vspace{0.15cm}
\item ¿Qué es el Chilam Balam?
\begin{enumerate}[label=\textbf{\Alph*)}]
\begin{multicols}{2}
  \item Un conjunto de libros mayas con profecías y conocimientos
  \item Gabriel García Márquez
  \item Julio Cortázar
  \item Mario Vargas Llosa
  \item Carlos Fuentes
  \item Juan Rulfo
\end{multicols}
\end{enumerate}
\vspace{0.15cm}
\item ¿Qué civilización produjo los Cantares mexicanos?
\begin{enumerate}[label=\textbf{\Alph*)}]
\begin{multicols}{2}
  \item Azteca
  \item Gabriel García Márquez
  \item Julio Cortázar
  \item Mario Vargas Llosa
  \item Carlos Fuentes
  \item Juan Rulfo
\end{multicols}
\end{enumerate}
\vspace{0.15cm}
\item ¿Qué es un "cuicatl" en la literatura azteca?
\begin{enumerate}[label=\textbf{\Alph*)}]
\begin{multicols}{2}
  \item Un canto o poema
  \item Gabriel García Márquez
  \item Julio Cortázar
  \item Mario Vargas Llosa
  \item Carlos Fuentes
  \item Juan Rulfo
\end{multicols}
\end{enumerate}
\vspace{0.15cm}
\item ¿Quién fue Nezahualcóyotl?
\begin{enumerate}[label=\textbf{\Alph*)}]
\begin{multicols}{2}
  \item Un poeta y rey azteca
  \item Gabriel García Márquez
  \item Julio Cortázar
  \item Mario Vargas Llosa
  \item Carlos Fuentes
  \item Juan Rulfo
\end{multicols}
\end{enumerate}
\vspace{0.15cm}
\item ¿Qué civilización produjo el Ollantay?
\begin{enumerate}[label=\textbf{\Alph*)}]
\begin{multicols}{2}
  \item Inca
  \item Gabriel García Márquez
  \item Julio Cortázar
  \item Mario Vargas Llosa
  \item Carlos Fuentes
  \item Juan Rulfo
\end{multicols}
\end{enumerate}
\vspace{0.15cm}
\item ¿Qué es el Ollantay?
\begin{enumerate}[label=\textbf{\Alph*)}]
\begin{multicols}{2}
  \item Un drama inca de tema amoroso y heroico
  \item Gabriel García Márquez
  \item Julio Cortázar
  \item Mario Vargas Llosa
  \item Carlos Fuentes
  \item Juan Rulfo
\end{multicols}
\end{enumerate}
\vspace{0.15cm}
\item ¿Qué es un "harawi" en la literatura inca?
\begin{enumerate}[label=\textbf{\Alph*)}]
\begin{multicols}{2}
  \item Un canto o poema lírico
  \item Gabriel García Márquez
  \item Julio Cortázar
  \item Mario Vargas Llosa
  \item Carlos Fuentes
  \item Juan Rulfo
\end{multicols}
\end{enumerate}
\vspace{0.15cm}
\item ¿Qué es un "haylli" en la literatura inca?
\begin{enumerate}[label=\textbf{\Alph*)}]
\begin{multicols}{2}
  \item Un canto de triunfo o guerra
  \item Gabriel García Márquez
  \item Julio Cortázar
  \item Mario Vargas Llosa
  \item Carlos Fuentes
  \item Juan Rulfo
\end{multicols}
\end{enumerate}
\vspace{0.15cm}
\item ¿Qué es un "aylli" en la literatura inca?
\begin{enumerate}[label=\textbf{\Alph*)}]
\begin{multicols}{2}
  \item Un canto de amor o alegría
  \item Gabriel García Márquez
  \item Julio Cortázar
  \item Mario Vargas Llosa
  \item Carlos Fuentes
  \item Juan Rulfo
\end{multicols}
\end{enumerate}
\vspace{0.15cm}
\item ¿Qué es un "umu" en la literatura inca?
\begin{enumerate}[label=\textbf{\Alph*)}]
\begin{multicols}{2}
  \item Un canto fúnebre
  \item Gabriel García Márquez
  \item Julio Cortázar
  \item Mario Vargas Llosa
  \item Carlos Fuentes
  \item Juan Rulfo
\end{multicols}
\end{enumerate}
\vspace{0.15cm}
\item ¿Qué es un "taqui" en la literatura inca?
\begin{enumerate}[label=\textbf{\Alph*)}]
\begin{multicols}{2}
  \item Un canto o danza
  \item Gabriel García Márquez
  \item Julio Cortázar
  \item Mario Vargas Llosa
  \item Carlos Fuentes
  \item Juan Rulfo
\end{multicols}
\end{enumerate}
\vspace{0.15cm}
\item ¿Qué civilización produjo el Rabinal Achí?
\begin{enumerate}[label=\textbf{\Alph*)}]
\begin{multicols}{2}
  \item Maya
  \item Gabriel García Márquez
  \item Julio Cortázar
  \item Mario Vargas Llosa
  \item Carlos Fuentes
  \item Juan Rulfo
\end{multicols}
\end{enumerate}
\vspace{0.15cm}
\item ¿Qué es el Rabinal Achí?
\begin{enumerate}[label=\textbf{\Alph*)}]
\begin{multicols}{2}
  \item Un drama maya de tema guerrero
  \item Gabriel García Márquez
  \item Julio Cortázar
  \item Mario Vargas Llosa
  \item Carlos Fuentes
  \item Juan Rulfo
\end{multicols}
\end{enumerate}
\vspace{0.15cm}
\item ¿Qué es la literatura precolombina?
\begin{enumerate}[label=\textbf{\Alph*)}]
\begin{multicols}{2}
  \item La literatura producida antes de la llegada de Colón a América
  \item Gabriel García Márquez
  \item Julio Cortázar
  \item Mario Vargas Llosa
  \item Carlos Fuentes
  \item Juan Rulfo
\end{multicols}
\end{enumerate}
\vspace{0.15cm}
\item ¿Cuál es una característica de la literatura precolombina?
\begin{enumerate}[label=\textbf{\Alph*)}]
\begin{multicols}{2}
  \item Ser de tradición oral y anónima
  \item Gabriel García Márquez
  \item Julio Cortázar
  \item Mario Vargas Llosa
  \item Carlos Fuentes
  \item Juan Rulfo
\end{multicols}
\end{enumerate}
\vspace{0.15cm}
\item ¿Qué civilización produjo el Popol Vuh?
\begin{enumerate}[label=\textbf{\Alph*)}]
\begin{multicols}{2}
  \item Maya
  \item Gabriel García Márquez
  \item Julio Cortázar
  \item Mario Vargas Llosa
  \item Carlos Fuentes
  \item Juan Rulfo
\end{multicols}
\end{enumerate}
\vspace{0.15cm}
\item ¿Qué es el Popol Vuh?
\begin{enumerate}[label=\textbf{\Alph*)}]
\begin{multicols}{2}
  \item Un libro sagrado maya que narra la creación del mundo
  \item Gabriel García Márquez
  \item Julio Cortázar
  \item Mario Vargas Llosa
  \item Carlos Fuentes
  \item Juan Rulfo
\end{multicols}
\end{enumerate}
\vspace{0.15cm}
\end{enumerate}\n
\subsubsection*{SECCIÓN II: LITERATURA COLONIAL (41–80)}

\begin{enumerate}[label=\textbf{\arabic*.}]
\setcounter{enumi}{40}
\item ¿Qué es la literatura colonial?
\begin{enumerate}[label=\textbf{\Alph*)}]
\begin{multicols}{2}
  \item La literatura producida durante el período colonial en América
  \item Gabriel García Márquez
  \item Julio Cortázar
  \item Mario Vargas Llosa
  \item Carlos Fuentes
  \item Juan Rulfo
\end{multicols}
\end{enumerate}
\vspace{0.15cm}
\item ¿Quiénes fueron los principales cronistas de Indias?
\begin{enumerate}[label=\textbf{\Alph*)}]
\begin{multicols}{2}
  \item Colón, Cortés, Bernal Díaz del Castillo y Las Casas
  \item Gabriel García Márquez
  \item Julio Cortázar
  \item Mario Vargas Llosa
  \item Carlos Fuentes
  \item Juan Rulfo
\end{multicols}
\end{enumerate}
\vspace{0.15cm}
\item ¿Quién escribió las Cartas de relación?
\begin{enumerate}[label=\textbf{\Alph*)}]
\begin{multicols}{2}
  \item Hernán Cortés
  \item Gabriel García Márquez
  \item Julio Cortázar
  \item Mario Vargas Llosa
  \item Carlos Fuentes
  \item Juan Rulfo
\end{multicols}
\end{enumerate}
\vspace{0.15cm}
\item ¿Quién escribió la Historia verdadera de la conquista de la Nueva España?
\begin{enumerate}[label=\textbf{\Alph*)}]
\begin{multicols}{2}
  \item Bernal Díaz del Castillo
  \item Gabriel García Márquez
  \item Julio Cortázar
  \item Mario Vargas Llosa
  \item Carlos Fuentes
  \item Juan Rulfo
\end{multicols}
\end{enumerate}
\vspace{0.15cm}
\item ¿Quién escribió la Brevísima relación de la destrucción de las Indias?
\begin{enumerate}[label=\textbf{\Alph*)}]
\begin{multicols}{2}
  \item Bartolomé de las Casas
  \item Gabriel García Márquez
  \item Julio Cortázar
  \item Mario Vargas Llosa
  \item Carlos Fuentes
  \item Juan Rulfo
\end{multicols}
\end{enumerate}
\vspace{0.15cm}
\item ¿Quién escribió los Comentarios reales?
\begin{enumerate}[label=\textbf{\Alph*)}]
\begin{multicols}{2}
  \item Inca Garcilaso de la Vega
  \item Gabriel García Márquez
  \item Julio Cortázar
  \item Mario Vargas Llosa
  \item Carlos Fuentes
  \item Juan Rulfo
\end{multicols}
\end{enumerate}
\vspace{0.15cm}
\item ¿Quién fue el Inca Garcilaso de la Vega?
\begin{enumerate}[label=\textbf{\Alph*)}]
\begin{multicols}{2}
  \item Un cronista mestizo peruano
  \item Gabriel García Márquez
  \item Julio Cortázar
  \item Mario Vargas Llosa
  \item Carlos Fuentes
  \item Juan Rulfo
\end{multicols}
\end{enumerate}
\vspace{0.15cm}
\item ¿Quién escribió La Araucana?
\begin{enumerate}[label=\textbf{\Alph*)}]
\begin{multicols}{2}
  \item Alonso de Ercilla
  \item Gabriel García Márquez
  \item Julio Cortázar
  \item Mario Vargas Llosa
  \item Carlos Fuentes
  \item Juan Rulfo
\end{multicols}
\end{enumerate}
\vspace{0.15cm}
\item ¿Qué es La Araucana?
\begin{enumerate}[label=\textbf{\Alph*)}]
\begin{multicols}{2}
  \item Un poema épico sobre la conquista de Chile
  \item Gabriel García Márquez
  \item Julio Cortázar
  \item Mario Vargas Llosa
  \item Carlos Fuentes
  \item Juan Rulfo
\end{multicols}
\end{enumerate}
\vspace{0.15cm}
\item ¿Quién escribió el Diario de a bordo?
\begin{enumerate}[label=\textbf{\Alph*)}]
\begin{multicols}{2}
  \item Cristóbal Colón
  \item Gabriel García Márquez
  \item Julio Cortázar
  \item Mario Vargas Llosa
  \item Carlos Fuentes
  \item Juan Rulfo
\end{multicols}
\end{enumerate}
\vspace{0.15cm}
\item ¿Quién fue Sor Juana Inés de la Cruz?
\begin{enumerate}[label=\textbf{\Alph*)}]
\begin{multicols}{2}
  \item Una poetisa y monja mexicana del Barroco
  \item Gabriel García Márquez
  \item Julio Cortázar
  \item Mario Vargas Llosa
  \item Carlos Fuentes
  \item Juan Rulfo
\end{multicols}
\end{enumerate}
\vspace{0.15cm}
\item ¿Cuál es la obra más famosa de Sor Juana Inés de la Cruz?
\begin{enumerate}[label=\textbf{\Alph*)}]
\begin{multicols}{2}
  \item Primero sueño
  \item Gabriel García Márquez
  \item Julio Cortázar
  \item Mario Vargas Llosa
  \item Carlos Fuentes
  \item Juan Rulfo
\end{multicols}
\end{enumerate}
\vspace{0.15cm}
\item ¿Quién escribió Respuesta a sor Filotea?
\begin{enumerate}[label=\textbf{\Alph*)}]
\begin{multicols}{2}
  \item Sor Juana Inés de la Cruz
  \item Gabriel García Márquez
  \item Julio Cortázar
  \item Mario Vargas Llosa
  \item Carlos Fuentes
  \item Juan Rulfo
\end{multicols}
\end{enumerate}
\vspace{0.15cm}
\item ¿Qué es el Barroco americano?
\begin{enumerate}[label=\textbf{\Alph*)}]
\begin{multicols}{2}
  \item La adaptación del Barroco europeo a la realidad americana
  \item Gabriel García Márquez
  \item Julio Cortázar
  \item Mario Vargas Llosa
  \item Carlos Fuentes
  \item Juan Rulfo
\end{multicols}
\end{enumerate}
\vspace{0.15cm}
\item ¿Quién fue Juan Ruiz de Alarcón?
\begin{enumerate}[label=\textbf{\Alph*)}]
\begin{multicols}{2}
  \item Un dramaturgo mexicano del Siglo de Oro
  \item Gabriel García Márquez
  \item Julio Cortázar
  \item Mario Vargas Llosa
  \item Carlos Fuentes
  \item Juan Rulfo
\end{multicols}
\end{enumerate}
\vspace{0.15cm}
\item ¿Cuál es una obra de Juan Ruiz de Alarcón?
\begin{enumerate}[label=\textbf{\Alph*)}]
\begin{multicols}{2}
  \item La verdad sospechosa
  \item Gabriel García Márquez
  \item Julio Cortázar
  \item Mario Vargas Llosa
  \item Carlos Fuentes
  \item Juan Rulfo
\end{multicols}
\end{enumerate}
\vspace{0.15cm}
\item ¿Quién fue Carlos de Sigüenza y Góngora?
\begin{enumerate}[label=\textbf{\Alph*)}]
\begin{multicols}{2}
  \item Un escritor y científico mexicano del Barroco
  \item Gabriel García Márquez
  \item Julio Cortázar
  \item Mario Vargas Llosa
  \item Carlos Fuentes
  \item Juan Rulfo
\end{multicols}
\end{enumerate}
\vspace{0.15cm}
\item ¿Quién fue Pedro de Peralta Barnuevo?
\begin{enumerate}[label=\textbf{\Alph*)}]
\begin{multicols}{2}
  \item Un escritor peruano del Barroco
  \item Gabriel García Márquez
  \item Julio Cortázar
  \item Mario Vargas Llosa
  \item Carlos Fuentes
  \item Juan Rulfo
\end{multicols}
\end{enumerate}
\vspace{0.15cm}
\item ¿Qué es la literatura colonial?
\begin{enumerate}[label=\textbf{\Alph*)}]
\begin{multicols}{2}
  \item La literatura producida durante el período colonial en América
  \item Gabriel García Márquez
  \item Julio Cortázar
  \item Mario Vargas Llosa
  \item Carlos Fuentes
  \item Juan Rulfo
\end{multicols}
\end{enumerate}
\vspace{0.15cm}
\item ¿Quiénes fueron los principales cronistas de Indias?
\begin{enumerate}[label=\textbf{\Alph*)}]
\begin{multicols}{2}
  \item Colón, Cortés, Bernal Díaz del Castillo y Las Casas
  \item Gabriel García Márquez
  \item Julio Cortázar
  \item Mario Vargas Llosa
  \item Carlos Fuentes
  \item Juan Rulfo
\end{multicols}
\end{enumerate}
\vspace{0.15cm}
\item ¿Quién escribió las Cartas de relación?
\begin{enumerate}[label=\textbf{\Alph*)}]
\begin{multicols}{2}
  \item Hernán Cortés
  \item Gabriel García Márquez
  \item Julio Cortázar
  \item Mario Vargas Llosa
  \item Carlos Fuentes
  \item Juan Rulfo
\end{multicols}
\end{enumerate}
\vspace{0.15cm}
\item ¿Quién escribió la Historia verdadera de la conquista de la Nueva España?
\begin{enumerate}[label=\textbf{\Alph*)}]
\begin{multicols}{2}
  \item Bernal Díaz del Castillo
  \item Gabriel García Márquez
  \item Julio Cortázar
  \item Mario Vargas Llosa
  \item Carlos Fuentes
  \item Juan Rulfo
\end{multicols}
\end{enumerate}
\vspace{0.15cm}
\item ¿Quién escribió la Brevísima relación de la destrucción de las Indias?
\begin{enumerate}[label=\textbf{\Alph*)}]
\begin{multicols}{2}
  \item Bartolomé de las Casas
  \item Gabriel García Márquez
  \item Julio Cortázar
  \item Mario Vargas Llosa
  \item Carlos Fuentes
  \item Juan Rulfo
\end{multicols}
\end{enumerate}
\vspace{0.15cm}
\item ¿Quién escribió los Comentarios reales?
\begin{enumerate}[label=\textbf{\Alph*)}]
\begin{multicols}{2}
  \item Inca Garcilaso de la Vega
  \item Gabriel García Márquez
  \item Julio Cortázar
  \item Mario Vargas Llosa
  \item Carlos Fuentes
  \item Juan Rulfo
\end{multicols}
\end{enumerate}
\vspace{0.15cm}
\item ¿Quién fue el Inca Garcilaso de la Vega?
\begin{enumerate}[label=\textbf{\Alph*)}]
\begin{multicols}{2}
  \item Un cronista mestizo peruano
  \item Gabriel García Márquez
  \item Julio Cortázar
  \item Mario Vargas Llosa
  \item Carlos Fuentes
  \item Juan Rulfo
\end{multicols}
\end{enumerate}
\vspace{0.15cm}
\item ¿Quién escribió La Araucana?
\begin{enumerate}[label=\textbf{\Alph*)}]
\begin{multicols}{2}
  \item Alonso de Ercilla
  \item Gabriel García Márquez
  \item Julio Cortázar
  \item Mario Vargas Llosa
  \item Carlos Fuentes
  \item Juan Rulfo
\end{multicols}
\end{enumerate}
\vspace{0.15cm}
\item ¿Qué es La Araucana?
\begin{enumerate}[label=\textbf{\Alph*)}]
\begin{multicols}{2}
  \item Un poema épico sobre la conquista de Chile
  \item Gabriel García Márquez
  \item Julio Cortázar
  \item Mario Vargas Llosa
  \item Carlos Fuentes
  \item Juan Rulfo
\end{multicols}
\end{enumerate}
\vspace{0.15cm}
\item ¿Quién escribió el Diario de a bordo?
\begin{enumerate}[label=\textbf{\Alph*)}]
\begin{multicols}{2}
  \item Cristóbal Colón
  \item Gabriel García Márquez
  \item Julio Cortázar
  \item Mario Vargas Llosa
  \item Carlos Fuentes
  \item Juan Rulfo
\end{multicols}
\end{enumerate}
\vspace{0.15cm}
\item ¿Quién fue Sor Juana Inés de la Cruz?
\begin{enumerate}[label=\textbf{\Alph*)}]
\begin{multicols}{2}
  \item Una poetisa y monja mexicana del Barroco
  \item Gabriel García Márquez
  \item Julio Cortázar
  \item Mario Vargas Llosa
  \item Carlos Fuentes
  \item Juan Rulfo
\end{multicols}
\end{enumerate}
\vspace{0.15cm}
\item ¿Cuál es la obra más famosa de Sor Juana Inés de la Cruz?
\begin{enumerate}[label=\textbf{\Alph*)}]
\begin{multicols}{2}
  \item Primero sueño
  \item Gabriel García Márquez
  \item Julio Cortázar
  \item Mario Vargas Llosa
  \item Carlos Fuentes
  \item Juan Rulfo
\end{multicols}
\end{enumerate}
\vspace{0.15cm}
\item ¿Quién escribió Respuesta a sor Filotea?
\begin{enumerate}[label=\textbf{\Alph*)}]
\begin{multicols}{2}
  \item Sor Juana Inés de la Cruz
  \item Gabriel García Márquez
  \item Julio Cortázar
  \item Mario Vargas Llosa
  \item Carlos Fuentes
  \item Juan Rulfo
\end{multicols}
\end{enumerate}
\vspace{0.15cm}
\item ¿Qué es el Barroco americano?
\begin{enumerate}[label=\textbf{\Alph*)}]
\begin{multicols}{2}
  \item La adaptación del Barroco europeo a la realidad americana
  \item Gabriel García Márquez
  \item Julio Cortázar
  \item Mario Vargas Llosa
  \item Carlos Fuentes
  \item Juan Rulfo
\end{multicols}
\end{enumerate}
\vspace{0.15cm}
\item ¿Quién fue Juan Ruiz de Alarcón?
\begin{enumerate}[label=\textbf{\Alph*)}]
\begin{multicols}{2}
  \item Un dramaturgo mexicano del Siglo de Oro
  \item Gabriel García Márquez
  \item Julio Cortázar
  \item Mario Vargas Llosa
  \item Carlos Fuentes
  \item Juan Rulfo
\end{multicols}
\end{enumerate}
\vspace{0.15cm}
\item ¿Cuál es una obra de Juan Ruiz de Alarcón?
\begin{enumerate}[label=\textbf{\Alph*)}]
\begin{multicols}{2}
  \item La verdad sospechosa
  \item Gabriel García Márquez
  \item Julio Cortázar
  \item Mario Vargas Llosa
  \item Carlos Fuentes
  \item Juan Rulfo
\end{multicols}
\end{enumerate}
\vspace{0.15cm}
\item ¿Quién fue Carlos de Sigüenza y Góngora?
\begin{enumerate}[label=\textbf{\Alph*)}]
\begin{multicols}{2}
  \item Un escritor y científico mexicano del Barroco
  \item Gabriel García Márquez
  \item Julio Cortázar
  \item Mario Vargas Llosa
  \item Carlos Fuentes
  \item Juan Rulfo
\end{multicols}
\end{enumerate}
\vspace{0.15cm}
\item ¿Quién fue Pedro de Peralta Barnuevo?
\begin{enumerate}[label=\textbf{\Alph*)}]
\begin{multicols}{2}
  \item Un escritor peruano del Barroco
  \item Gabriel García Márquez
  \item Julio Cortázar
  \item Mario Vargas Llosa
  \item Carlos Fuentes
  \item Juan Rulfo
\end{multicols}
\end{enumerate}
\vspace{0.15cm}
\item ¿Qué es la literatura colonial?
\begin{enumerate}[label=\textbf{\Alph*)}]
\begin{multicols}{2}
  \item La literatura producida durante el período colonial en América
  \item Gabriel García Márquez
  \item Julio Cortázar
  \item Mario Vargas Llosa
  \item Carlos Fuentes
  \item Juan Rulfo
\end{multicols}
\end{enumerate}
\vspace{0.15cm}
\item ¿Quiénes fueron los principales cronistas de Indias?
\begin{enumerate}[label=\textbf{\Alph*)}]
\begin{multicols}{2}
  \item Colón, Cortés, Bernal Díaz del Castillo y Las Casas
  \item Gabriel García Márquez
  \item Julio Cortázar
  \item Mario Vargas Llosa
  \item Carlos Fuentes
  \item Juan Rulfo
\end{multicols}
\end{enumerate}
\vspace{0.15cm}
\item ¿Quién escribió las Cartas de relación?
\begin{enumerate}[label=\textbf{\Alph*)}]
\begin{multicols}{2}
  \item Hernán Cortés
  \item Gabriel García Márquez
  \item Julio Cortázar
  \item Mario Vargas Llosa
  \item Carlos Fuentes
  \item Juan Rulfo
\end{multicols}
\end{enumerate}
\vspace{0.15cm}
\item ¿Quién escribió la Historia verdadera de la conquista de la Nueva España?
\begin{enumerate}[label=\textbf{\Alph*)}]
\begin{multicols}{2}
  \item Bernal Díaz del Castillo
  \item Gabriel García Márquez
  \item Julio Cortázar
  \item Mario Vargas Llosa
  \item Carlos Fuentes
  \item Juan Rulfo
\end{multicols}
\end{enumerate}
\vspace{0.15cm}
\end{enumerate}\n
\subsubsection*{SECCIÓN III: ROMANTICISMO AMERICANO (81–120)}

\begin{enumerate}[label=\textbf{\arabic*.}]
\setcounter{enumi}{80}
\item ¿Qué es el Romanticismo americano?
\begin{enumerate}[label=\textbf{\Alph*)}]
\begin{multicols}{2}
  \item La adaptación del Romanticismo europeo a la realidad americana
  \item Gabriel García Márquez
  \item Julio Cortázar
  \item Mario Vargas Llosa
  \item Carlos Fuentes
  \item Juan Rulfo
\end{multicols}
\end{enumerate}
\vspace{0.15cm}
\item ¿Cuál es una característica del Romanticismo americano?
\begin{enumerate}[label=\textbf{\Alph*)}]
\begin{multicols}{2}
  \item La exaltación de la naturaleza y la identidad nacional
  \item Gabriel García Márquez
  \item Julio Cortázar
  \item Mario Vargas Llosa
  \item Carlos Fuentes
  \item Juan Rulfo
\end{multicols}
\end{enumerate}
\vspace{0.15cm}
\item ¿Quién es el principal representante del Romanticismo argentino?
\begin{enumerate}[label=\textbf{\Alph*)}]
\begin{multicols}{2}
  \item Esteban Echeverría
  \item Gabriel García Márquez
  \item Julio Cortázar
  \item Mario Vargas Llosa
  \item Carlos Fuentes
  \item Juan Rulfo
\end{multicols}
\end{enumerate}
\vspace{0.15cm}
\item ¿Quién escribió El matadero?
\begin{enumerate}[label=\textbf{\Alph*)}]
\begin{multicols}{2}
  \item Esteban Echeverría
  \item Gabriel García Márquez
  \item Julio Cortázar
  \item Mario Vargas Llosa
  \item Carlos Fuentes
  \item Juan Rulfo
\end{multicols}
\end{enumerate}
\vspace{0.15cm}
\item ¿Quién escribió Amalia?
\begin{enumerate}[label=\textbf{\Alph*)}]
\begin{multicols}{2}
  \item José Mármol
  \item Gabriel García Márquez
  \item Julio Cortázar
  \item Mario Vargas Llosa
  \item Carlos Fuentes
  \item Juan Rulfo
\end{multicols}
\end{enumerate}
\vspace{0.15cm}
\item ¿Quién escribió Facundo?
\begin{enumerate}[label=\textbf{\Alph*)}]
\begin{multicols}{2}
  \item Domingo Faustino Sarmiento
  \item Gabriel García Márquez
  \item Julio Cortázar
  \item Mario Vargas Llosa
  \item Carlos Fuentes
  \item Juan Rulfo
\end{multicols}
\end{enumerate}
\vspace{0.15cm}
\item ¿Quién escribió Martín Fierro?
\begin{enumerate}[label=\textbf{\Alph*)}]
\begin{multicols}{2}
  \item José Hernández
  \item Gabriel García Márquez
  \item Julio Cortázar
  \item Mario Vargas Llosa
  \item Carlos Fuentes
  \item Juan Rulfo
\end{multicols}
\end{enumerate}
\vspace{0.15cm}
\item ¿Qué es Martín Fierro?
\begin{enumerate}[label=\textbf{\Alph*)}]
\begin{multicols}{2}
  \item Un poema épico gauchesco argentino
  \item Gabriel García Márquez
  \item Julio Cortázar
  \item Mario Vargas Llosa
  \item Carlos Fuentes
  \item Juan Rulfo
\end{multicols}
\end{enumerate}
\vspace{0.15cm}
\item ¿Quién es el principal representante del Romanticismo cubano?
\begin{enumerate}[label=\textbf{\Alph*)}]
\begin{multicols}{2}
  \item José María Heredia
  \item Gabriel García Márquez
  \item Julio Cortázar
  \item Mario Vargas Llosa
  \item Carlos Fuentes
  \item Juan Rulfo
\end{multicols}
\end{enumerate}
\vspace{0.15cm}
\item ¿Quién escribió Cecilia Valdés?
\begin{enumerate}[label=\textbf{\Alph*)}]
\begin{multicols}{2}
  \item Cirilo Villaverde
  \item Gabriel García Márquez
  \item Julio Cortázar
  \item Mario Vargas Llosa
  \item Carlos Fuentes
  \item Juan Rulfo
\end{multicols}
\end{enumerate}
\vspace{0.15cm}
\item ¿Quién escribió Sab?
\begin{enumerate}[label=\textbf{\Alph*)}]
\begin{multicols}{2}
  \item Gertrudis Gómez de Avellaneda
  \item Gabriel García Márquez
  \item Julio Cortázar
  \item Mario Vargas Llosa
  \item Carlos Fuentes
  \item Juan Rulfo
\end{multicols}
\end{enumerate}
\vspace{0.15cm}
\item ¿Quién es el principal representante del Romanticismo colombiano?
\begin{enumerate}[label=\textbf{\Alph*)}]
\begin{multicols}{2}
  \item Jorge Isaacs
  \item Gabriel García Márquez
  \item Julio Cortázar
  \item Mario Vargas Llosa
  \item Carlos Fuentes
  \item Juan Rulfo
\end{multicols}
\end{enumerate}
\vspace{0.15cm}
\item ¿Quién escribió María?
\begin{enumerate}[label=\textbf{\Alph*)}]
\begin{multicols}{2}
  \item Jorge Isaacs
  \item Gabriel García Márquez
  \item Julio Cortázar
  \item Mario Vargas Llosa
  \item Carlos Fuentes
  \item Juan Rulfo
\end{multicols}
\end{enumerate}
\vspace{0.15cm}
\item ¿Qué es María?
\begin{enumerate}[label=\textbf{\Alph*)}]
\begin{multicols}{2}
  \item Una novela romántica colombiana
  \item Gabriel García Márquez
  \item Julio Cortázar
  \item Mario Vargas Llosa
  \item Carlos Fuentes
  \item Juan Rulfo
\end{multicols}
\end{enumerate}
\vspace{0.15cm}
\item ¿Quién es el principal representante del Romanticismo peruano?
\begin{enumerate}[label=\textbf{\Alph*)}]
\begin{multicols}{2}
  \item Ricardo Palma
  \item Gabriel García Márquez
  \item Julio Cortázar
  \item Mario Vargas Llosa
  \item Carlos Fuentes
  \item Juan Rulfo
\end{multicols}
\end{enumerate}
\vspace{0.15cm}
\item ¿Quién escribió las Tradiciones peruanas?
\begin{enumerate}[label=\textbf{\Alph*)}]
\begin{multicols}{2}
  \item Ricardo Palma
  \item Gabriel García Márquez
  \item Julio Cortázar
  \item Mario Vargas Llosa
  \item Carlos Fuentes
  \item Juan Rulfo
\end{multicols}
\end{enumerate}
\vspace{0.15cm}
\item ¿Qué son las Tradiciones peruanas?
\begin{enumerate}[label=\textbf{\Alph*)}]
\begin{multicols}{2}
  \item Una colección de relatos históricos y costumbristas
  \item Gabriel García Márquez
  \item Julio Cortázar
  \item Mario Vargas Llosa
  \item Carlos Fuentes
  \item Juan Rulfo
\end{multicols}
\end{enumerate}
\vspace{0.15cm}
\item ¿Quién es el principal representante del Romanticismo mexicano?
\begin{enumerate}[label=\textbf{\Alph*)}]
\begin{multicols}{2}
  \item Ignacio Manuel Altamirano
  \item Gabriel García Márquez
  \item Julio Cortázar
  \item Mario Vargas Llosa
  \item Carlos Fuentes
  \item Juan Rulfo
\end{multicols}
\end{enumerate}
\vspace{0.15cm}
\item ¿Quién escribió Clemencia?
\begin{enumerate}[label=\textbf{\Alph*)}]
\begin{multicols}{2}
  \item Ignacio Manuel Altamirano
  \item Gabriel García Márquez
  \item Julio Cortázar
  \item Mario Vargas Llosa
  \item Carlos Fuentes
  \item Juan Rulfo
\end{multicols}
\end{enumerate}
\vspace{0.15cm}
\item ¿Quién escribió El Zarco?
\begin{enumerate}[label=\textbf{\Alph*)}]
\begin{multicols}{2}
  \item Ignacio Manuel Altamirano
  \item Gabriel García Márquez
  \item Julio Cortázar
  \item Mario Vargas Llosa
  \item Carlos Fuentes
  \item Juan Rulfo
\end{multicols}
\end{enumerate}
\vspace{0.15cm}
\item ¿Qué es el Romanticismo americano?
\begin{enumerate}[label=\textbf{\Alph*)}]
\begin{multicols}{2}
  \item La adaptación del Romanticismo europeo a la realidad americana
  \item Gabriel García Márquez
  \item Julio Cortázar
  \item Mario Vargas Llosa
  \item Carlos Fuentes
  \item Juan Rulfo
\end{multicols}
\end{enumerate}
\vspace{0.15cm}
\item ¿Cuál es una característica del Romanticismo americano?
\begin{enumerate}[label=\textbf{\Alph*)}]
\begin{multicols}{2}
  \item La exaltación de la naturaleza y la identidad nacional
  \item Gabriel García Márquez
  \item Julio Cortázar
  \item Mario Vargas Llosa
  \item Carlos Fuentes
  \item Juan Rulfo
\end{multicols}
\end{enumerate}
\vspace{0.15cm}
\item ¿Quién es el principal representante del Romanticismo argentino?
\begin{enumerate}[label=\textbf{\Alph*)}]
\begin{multicols}{2}
  \item Esteban Echeverría
  \item Gabriel García Márquez
  \item Julio Cortázar
  \item Mario Vargas Llosa
  \item Carlos Fuentes
  \item Juan Rulfo
\end{multicols}
\end{enumerate}
\vspace{0.15cm}
\item ¿Quién escribió El matadero?
\begin{enumerate}[label=\textbf{\Alph*)}]
\begin{multicols}{2}
  \item Esteban Echeverría
  \item Gabriel García Márquez
  \item Julio Cortázar
  \item Mario Vargas Llosa
  \item Carlos Fuentes
  \item Juan Rulfo
\end{multicols}
\end{enumerate}
\vspace{0.15cm}
\item ¿Quién escribió Amalia?
\begin{enumerate}[label=\textbf{\Alph*)}]
\begin{multicols}{2}
  \item José Mármol
  \item Gabriel García Márquez
  \item Julio Cortázar
  \item Mario Vargas Llosa
  \item Carlos Fuentes
  \item Juan Rulfo
\end{multicols}
\end{enumerate}
\vspace{0.15cm}
\item ¿Quién escribió Facundo?
\begin{enumerate}[label=\textbf{\Alph*)}]
\begin{multicols}{2}
  \item Domingo Faustino Sarmiento
  \item Gabriel García Márquez
  \item Julio Cortázar
  \item Mario Vargas Llosa
  \item Carlos Fuentes
  \item Juan Rulfo
\end{multicols}
\end{enumerate}
\vspace{0.15cm}
\item ¿Quién escribió Martín Fierro?
\begin{enumerate}[label=\textbf{\Alph*)}]
\begin{multicols}{2}
  \item José Hernández
  \item Gabriel García Márquez
  \item Julio Cortázar
  \item Mario Vargas Llosa
  \item Carlos Fuentes
  \item Juan Rulfo
\end{multicols}
\end{enumerate}
\vspace{0.15cm}
\item ¿Qué es Martín Fierro?
\begin{enumerate}[label=\textbf{\Alph*)}]
\begin{multicols}{2}
  \item Un poema épico gauchesco argentino
  \item Gabriel García Márquez
  \item Julio Cortázar
  \item Mario Vargas Llosa
  \item Carlos Fuentes
  \item Juan Rulfo
\end{multicols}
\end{enumerate}
\vspace{0.15cm}
\item ¿Quién es el principal representante del Romanticismo cubano?
\begin{enumerate}[label=\textbf{\Alph*)}]
\begin{multicols}{2}
  \item José María Heredia
  \item Gabriel García Márquez
  \item Julio Cortázar
  \item Mario Vargas Llosa
  \item Carlos Fuentes
  \item Juan Rulfo
\end{multicols}
\end{enumerate}
\vspace{0.15cm}
\item ¿Quién escribió Cecilia Valdés?
\begin{enumerate}[label=\textbf{\Alph*)}]
\begin{multicols}{2}
  \item Cirilo Villaverde
  \item Gabriel García Márquez
  \item Julio Cortázar
  \item Mario Vargas Llosa
  \item Carlos Fuentes
  \item Juan Rulfo
\end{multicols}
\end{enumerate}
\vspace{0.15cm}
\item ¿Quién escribió Sab?
\begin{enumerate}[label=\textbf{\Alph*)}]
\begin{multicols}{2}
  \item Gertrudis Gómez de Avellaneda
  \item Gabriel García Márquez
  \item Julio Cortázar
  \item Mario Vargas Llosa
  \item Carlos Fuentes
  \item Juan Rulfo
\end{multicols}
\end{enumerate}
\vspace{0.15cm}
\item ¿Quién es el principal representante del Romanticismo colombiano?
\begin{enumerate}[label=\textbf{\Alph*)}]
\begin{multicols}{2}
  \item Jorge Isaacs
  \item Gabriel García Márquez
  \item Julio Cortázar
  \item Mario Vargas Llosa
  \item Carlos Fuentes
  \item Juan Rulfo
\end{multicols}
\end{enumerate}
\vspace{0.15cm}
\item ¿Quién escribió María?
\begin{enumerate}[label=\textbf{\Alph*)}]
\begin{multicols}{2}
  \item Jorge Isaacs
  \item Gabriel García Márquez
  \item Julio Cortázar
  \item Mario Vargas Llosa
  \item Carlos Fuentes
  \item Juan Rulfo
\end{multicols}
\end{enumerate}
\vspace{0.15cm}
\item ¿Qué es María?
\begin{enumerate}[label=\textbf{\Alph*)}]
\begin{multicols}{2}
  \item Una novela romántica colombiana
  \item Gabriel García Márquez
  \item Julio Cortázar
  \item Mario Vargas Llosa
  \item Carlos Fuentes
  \item Juan Rulfo
\end{multicols}
\end{enumerate}
\vspace{0.15cm}
\item ¿Quién es el principal representante del Romanticismo peruano?
\begin{enumerate}[label=\textbf{\Alph*)}]
\begin{multicols}{2}
  \item Ricardo Palma
  \item Gabriel García Márquez
  \item Julio Cortázar
  \item Mario Vargas Llosa
  \item Carlos Fuentes
  \item Juan Rulfo
\end{multicols}
\end{enumerate}
\vspace{0.15cm}
\item ¿Quién escribió las Tradiciones peruanas?
\begin{enumerate}[label=\textbf{\Alph*)}]
\begin{multicols}{2}
  \item Ricardo Palma
  \item Gabriel García Márquez
  \item Julio Cortázar
  \item Mario Vargas Llosa
  \item Carlos Fuentes
  \item Juan Rulfo
\end{multicols}
\end{enumerate}
\vspace{0.15cm}
\item ¿Qué son las Tradiciones peruanas?
\begin{enumerate}[label=\textbf{\Alph*)}]
\begin{multicols}{2}
  \item Una colección de relatos históricos y costumbristas
  \item Gabriel García Márquez
  \item Julio Cortázar
  \item Mario Vargas Llosa
  \item Carlos Fuentes
  \item Juan Rulfo
\end{multicols}
\end{enumerate}
\vspace{0.15cm}
\item ¿Quién es el principal representante del Romanticismo mexicano?
\begin{enumerate}[label=\textbf{\Alph*)}]
\begin{multicols}{2}
  \item Ignacio Manuel Altamirano
  \item Gabriel García Márquez
  \item Julio Cortázar
  \item Mario Vargas Llosa
  \item Carlos Fuentes
  \item Juan Rulfo
\end{multicols}
\end{enumerate}
\vspace{0.15cm}
\item ¿Quién escribió Clemencia?
\begin{enumerate}[label=\textbf{\Alph*)}]
\begin{multicols}{2}
  \item Ignacio Manuel Altamirano
  \item Gabriel García Márquez
  \item Julio Cortázar
  \item Mario Vargas Llosa
  \item Carlos Fuentes
  \item Juan Rulfo
\end{multicols}
\end{enumerate}
\vspace{0.15cm}
\item ¿Quién escribió El Zarco?
\begin{enumerate}[label=\textbf{\Alph*)}]
\begin{multicols}{2}
  \item Ignacio Manuel Altamirano
  \item Gabriel García Márquez
  \item Julio Cortázar
  \item Mario Vargas Llosa
  \item Carlos Fuentes
  \item Juan Rulfo
\end{multicols}
\end{enumerate}
\vspace{0.15cm}
\end{enumerate}\n
\subsubsection*{SECCIÓN IV: MODERNISMO – RUBÉN DARÍO (121–170)}

\begin{enumerate}[label=\textbf{\arabic*.}]
\setcounter{enumi}{120}
\item ¿Qué es el Modernismo?
\begin{enumerate}[label=\textbf{\Alph*)}]
\begin{multicols}{2}
  \item Un movimiento literario hispanoamericano de finales del siglo XIX
  \item Gabriel García Márquez
  \item Julio Cortázar
  \item Mario Vargas Llosa
  \item Carlos Fuentes
  \item Juan Rulfo
\end{multicols}
\end{enumerate}
\vspace{0.15cm}
\item ¿Quién es el principal representante del Modernismo?
\begin{enumerate}[label=\textbf{\Alph*)}]
\begin{multicols}{2}
  \item Rubén Darío
  \item Gabriel García Márquez
  \item Julio Cortázar
  \item Mario Vargas Llosa
  \item Carlos Fuentes
  \item Juan Rulfo
\end{multicols}
\end{enumerate}
\vspace{0.15cm}
\item ¿Quién es el autor de Azul?
\begin{enumerate}[label=\textbf{\Alph*)}]
\begin{multicols}{2}
  \item Rubén Darío
  \item Gabriel García Márquez
  \item Julio Cortázar
  \item Mario Vargas Llosa
  \item Carlos Fuentes
  \item Juan Rulfo
\end{multicols}
\end{enumerate}
\vspace{0.15cm}
\item ¿Quién es el autor de Prosas profanas?
\begin{enumerate}[label=\textbf{\Alph*)}]
\begin{multicols}{2}
  \item Rubén Darío
  \item Gabriel García Márquez
  \item Julio Cortázar
  \item Mario Vargas Llosa
  \item Carlos Fuentes
  \item Juan Rulfo
\end{multicols}
\end{enumerate}
\vspace{0.15cm}
\item ¿Quién es el autor de Cantos de vida y esperanza?
\begin{enumerate}[label=\textbf{\Alph*)}]
\begin{multicols}{2}
  \item Rubén Darío
  \item Gabriel García Márquez
  \item Julio Cortázar
  \item Mario Vargas Llosa
  \item Carlos Fuentes
  \item Juan Rulfo
\end{multicols}
\end{enumerate}
\vspace{0.15cm}
\item ¿Qué es el Modernismo?
\begin{enumerate}[label=\textbf{\Alph*)}]
\begin{multicols}{2}
  \item Un movimiento literario hispanoamericano de finales del siglo XIX
  \item Gabriel García Márquez
  \item Julio Cortázar
  \item Mario Vargas Llosa
  \item Carlos Fuentes
  \item Juan Rulfo
\end{multicols}
\end{enumerate}
\vspace{0.15cm}
\item ¿Quién es el principal representante del Modernismo?
\begin{enumerate}[label=\textbf{\Alph*)}]
\begin{multicols}{2}
  \item Rubén Darío
  \item Gabriel García Márquez
  \item Julio Cortázar
  \item Mario Vargas Llosa
  \item Carlos Fuentes
  \item Juan Rulfo
\end{multicols}
\end{enumerate}
\vspace{0.15cm}
\item ¿Quién es el autor de Azul?
\begin{enumerate}[label=\textbf{\Alph*)}]
\begin{multicols}{2}
  \item Rubén Darío
  \item Gabriel García Márquez
  \item Julio Cortázar
  \item Mario Vargas Llosa
  \item Carlos Fuentes
  \item Juan Rulfo
\end{multicols}
\end{enumerate}
\vspace{0.15cm}
\item ¿Quién es el autor de Prosas profanas?
\begin{enumerate}[label=\textbf{\Alph*)}]
\begin{multicols}{2}
  \item Rubén Darío
  \item Gabriel García Márquez
  \item Julio Cortázar
  \item Mario Vargas Llosa
  \item Carlos Fuentes
  \item Juan Rulfo
\end{multicols}
\end{enumerate}
\vspace{0.15cm}
\item ¿Quién es el autor de Cantos de vida y esperanza?
\begin{enumerate}[label=\textbf{\Alph*)}]
\begin{multicols}{2}
  \item Rubén Darío
  \item Gabriel García Márquez
  \item Julio Cortázar
  \item Mario Vargas Llosa
  \item Carlos Fuentes
  \item Juan Rulfo
\end{multicols}
\end{enumerate}
\vspace{0.15cm}
\item ¿Qué es el Modernismo?
\begin{enumerate}[label=\textbf{\Alph*)}]
\begin{multicols}{2}
  \item Un movimiento literario hispanoamericano de finales del siglo XIX
  \item Gabriel García Márquez
  \item Julio Cortázar
  \item Mario Vargas Llosa
  \item Carlos Fuentes
  \item Juan Rulfo
\end{multicols}
\end{enumerate}
\vspace{0.15cm}
\item ¿Quién es el principal representante del Modernismo?
\begin{enumerate}[label=\textbf{\Alph*)}]
\begin{multicols}{2}
  \item Rubén Darío
  \item Gabriel García Márquez
  \item Julio Cortázar
  \item Mario Vargas Llosa
  \item Carlos Fuentes
  \item Juan Rulfo
\end{multicols}
\end{enumerate}
\vspace{0.15cm}
\item ¿Quién es el autor de Azul?
\begin{enumerate}[label=\textbf{\Alph*)}]
\begin{multicols}{2}
  \item Rubén Darío
  \item Gabriel García Márquez
  \item Julio Cortázar
  \item Mario Vargas Llosa
  \item Carlos Fuentes
  \item Juan Rulfo
\end{multicols}
\end{enumerate}
\vspace{0.15cm}
\item ¿Quién es el autor de Prosas profanas?
\begin{enumerate}[label=\textbf{\Alph*)}]
\begin{multicols}{2}
  \item Rubén Darío
  \item Gabriel García Márquez
  \item Julio Cortázar
  \item Mario Vargas Llosa
  \item Carlos Fuentes
  \item Juan Rulfo
\end{multicols}
\end{enumerate}
\vspace{0.15cm}
\item ¿Quién es el autor de Cantos de vida y esperanza?
\begin{enumerate}[label=\textbf{\Alph*)}]
\begin{multicols}{2}
  \item Rubén Darío
  \item Gabriel García Márquez
  \item Julio Cortázar
  \item Mario Vargas Llosa
  \item Carlos Fuentes
  \item Juan Rulfo
\end{multicols}
\end{enumerate}
\vspace{0.15cm}
\item ¿Qué es el Modernismo?
\begin{enumerate}[label=\textbf{\Alph*)}]
\begin{multicols}{2}
  \item Un movimiento literario hispanoamericano de finales del siglo XIX
  \item Gabriel García Márquez
  \item Julio Cortázar
  \item Mario Vargas Llosa
  \item Carlos Fuentes
  \item Juan Rulfo
\end{multicols}
\end{enumerate}
\vspace{0.15cm}
\item ¿Quién es el principal representante del Modernismo?
\begin{enumerate}[label=\textbf{\Alph*)}]
\begin{multicols}{2}
  \item Rubén Darío
  \item Gabriel García Márquez
  \item Julio Cortázar
  \item Mario Vargas Llosa
  \item Carlos Fuentes
  \item Juan Rulfo
\end{multicols}
\end{enumerate}
\vspace{0.15cm}
\item ¿Quién es el autor de Azul?
\begin{enumerate}[label=\textbf{\Alph*)}]
\begin{multicols}{2}
  \item Rubén Darío
  \item Gabriel García Márquez
  \item Julio Cortázar
  \item Mario Vargas Llosa
  \item Carlos Fuentes
  \item Juan Rulfo
\end{multicols}
\end{enumerate}
\vspace{0.15cm}
\item ¿Quién es el autor de Prosas profanas?
\begin{enumerate}[label=\textbf{\Alph*)}]
\begin{multicols}{2}
  \item Rubén Darío
  \item Gabriel García Márquez
  \item Julio Cortázar
  \item Mario Vargas Llosa
  \item Carlos Fuentes
  \item Juan Rulfo
\end{multicols}
\end{enumerate}
\vspace{0.15cm}
\item ¿Quién es el autor de Cantos de vida y esperanza?
\begin{enumerate}[label=\textbf{\Alph*)}]
\begin{multicols}{2}
  \item Rubén Darío
  \item Gabriel García Márquez
  \item Julio Cortázar
  \item Mario Vargas Llosa
  \item Carlos Fuentes
  \item Juan Rulfo
\end{multicols}
\end{enumerate}
\vspace{0.15cm}
\item ¿Qué es el Modernismo?
\begin{enumerate}[label=\textbf{\Alph*)}]
\begin{multicols}{2}
  \item Un movimiento literario hispanoamericano de finales del siglo XIX
  \item Gabriel García Márquez
  \item Julio Cortázar
  \item Mario Vargas Llosa
  \item Carlos Fuentes
  \item Juan Rulfo
\end{multicols}
\end{enumerate}
\vspace{0.15cm}
\item ¿Quién es el principal representante del Modernismo?
\begin{enumerate}[label=\textbf{\Alph*)}]
\begin{multicols}{2}
  \item Rubén Darío
  \item Gabriel García Márquez
  \item Julio Cortázar
  \item Mario Vargas Llosa
  \item Carlos Fuentes
  \item Juan Rulfo
\end{multicols}
\end{enumerate}
\vspace{0.15cm}
\item ¿Quién es el autor de Azul?
\begin{enumerate}[label=\textbf{\Alph*)}]
\begin{multicols}{2}
  \item Rubén Darío
  \item Gabriel García Márquez
  \item Julio Cortázar
  \item Mario Vargas Llosa
  \item Carlos Fuentes
  \item Juan Rulfo
\end{multicols}
\end{enumerate}
\vspace{0.15cm}
\item ¿Quién es el autor de Prosas profanas?
\begin{enumerate}[label=\textbf{\Alph*)}]
\begin{multicols}{2}
  \item Rubén Darío
  \item Gabriel García Márquez
  \item Julio Cortázar
  \item Mario Vargas Llosa
  \item Carlos Fuentes
  \item Juan Rulfo
\end{multicols}
\end{enumerate}
\vspace{0.15cm}
\item ¿Quién es el autor de Cantos de vida y esperanza?
\begin{enumerate}[label=\textbf{\Alph*)}]
\begin{multicols}{2}
  \item Rubén Darío
  \item Gabriel García Márquez
  \item Julio Cortázar
  \item Mario Vargas Llosa
  \item Carlos Fuentes
  \item Juan Rulfo
\end{multicols}
\end{enumerate}
\vspace{0.15cm}
\item ¿Qué es el Modernismo?
\begin{enumerate}[label=\textbf{\Alph*)}]
\begin{multicols}{2}
  \item Un movimiento literario hispanoamericano de finales del siglo XIX
  \item Gabriel García Márquez
  \item Julio Cortázar
  \item Mario Vargas Llosa
  \item Carlos Fuentes
  \item Juan Rulfo
\end{multicols}
\end{enumerate}
\vspace{0.15cm}
\item ¿Quién es el principal representante del Modernismo?
\begin{enumerate}[label=\textbf{\Alph*)}]
\begin{multicols}{2}
  \item Rubén Darío
  \item Gabriel García Márquez
  \item Julio Cortázar
  \item Mario Vargas Llosa
  \item Carlos Fuentes
  \item Juan Rulfo
\end{multicols}
\end{enumerate}
\vspace{0.15cm}
\item ¿Quién es el autor de Azul?
\begin{enumerate}[label=\textbf{\Alph*)}]
\begin{multicols}{2}
  \item Rubén Darío
  \item Gabriel García Márquez
  \item Julio Cortázar
  \item Mario Vargas Llosa
  \item Carlos Fuentes
  \item Juan Rulfo
\end{multicols}
\end{enumerate}
\vspace{0.15cm}
\item ¿Quién es el autor de Prosas profanas?
\begin{enumerate}[label=\textbf{\Alph*)}]
\begin{multicols}{2}
  \item Rubén Darío
  \item Gabriel García Márquez
  \item Julio Cortázar
  \item Mario Vargas Llosa
  \item Carlos Fuentes
  \item Juan Rulfo
\end{multicols}
\end{enumerate}
\vspace{0.15cm}
\item ¿Quién es el autor de Cantos de vida y esperanza?
\begin{enumerate}[label=\textbf{\Alph*)}]
\begin{multicols}{2}
  \item Rubén Darío
  \item Gabriel García Márquez
  \item Julio Cortázar
  \item Mario Vargas Llosa
  \item Carlos Fuentes
  \item Juan Rulfo
\end{multicols}
\end{enumerate}
\vspace{0.15cm}
\item ¿Qué es el Modernismo?
\begin{enumerate}[label=\textbf{\Alph*)}]
\begin{multicols}{2}
  \item Un movimiento literario hispanoamericano de finales del siglo XIX
  \item Gabriel García Márquez
  \item Julio Cortázar
  \item Mario Vargas Llosa
  \item Carlos Fuentes
  \item Juan Rulfo
\end{multicols}
\end{enumerate}
\vspace{0.15cm}
\item ¿Quién es el principal representante del Modernismo?
\begin{enumerate}[label=\textbf{\Alph*)}]
\begin{multicols}{2}
  \item Rubén Darío
  \item Gabriel García Márquez
  \item Julio Cortázar
  \item Mario Vargas Llosa
  \item Carlos Fuentes
  \item Juan Rulfo
\end{multicols}
\end{enumerate}
\vspace{0.15cm}
\item ¿Quién es el autor de Azul?
\begin{enumerate}[label=\textbf{\Alph*)}]
\begin{multicols}{2}
  \item Rubén Darío
  \item Gabriel García Márquez
  \item Julio Cortázar
  \item Mario Vargas Llosa
  \item Carlos Fuentes
  \item Juan Rulfo
\end{multicols}
\end{enumerate}
\vspace{0.15cm}
\item ¿Quién es el autor de Prosas profanas?
\begin{enumerate}[label=\textbf{\Alph*)}]
\begin{multicols}{2}
  \item Rubén Darío
  \item Gabriel García Márquez
  \item Julio Cortázar
  \item Mario Vargas Llosa
  \item Carlos Fuentes
  \item Juan Rulfo
\end{multicols}
\end{enumerate}
\vspace{0.15cm}
\item ¿Quién es el autor de Cantos de vida y esperanza?
\begin{enumerate}[label=\textbf{\Alph*)}]
\begin{multicols}{2}
  \item Rubén Darío
  \item Gabriel García Márquez
  \item Julio Cortázar
  \item Mario Vargas Llosa
  \item Carlos Fuentes
  \item Juan Rulfo
\end{multicols}
\end{enumerate}
\vspace{0.15cm}
\item ¿Qué es el Modernismo?
\begin{enumerate}[label=\textbf{\Alph*)}]
\begin{multicols}{2}
  \item Un movimiento literario hispanoamericano de finales del siglo XIX
  \item Gabriel García Márquez
  \item Julio Cortázar
  \item Mario Vargas Llosa
  \item Carlos Fuentes
  \item Juan Rulfo
\end{multicols}
\end{enumerate}
\vspace{0.15cm}
\item ¿Quién es el principal representante del Modernismo?
\begin{enumerate}[label=\textbf{\Alph*)}]
\begin{multicols}{2}
  \item Rubén Darío
  \item Gabriel García Márquez
  \item Julio Cortázar
  \item Mario Vargas Llosa
  \item Carlos Fuentes
  \item Juan Rulfo
\end{multicols}
\end{enumerate}
\vspace{0.15cm}
\item ¿Quién es el autor de Azul?
\begin{enumerate}[label=\textbf{\Alph*)}]
\begin{multicols}{2}
  \item Rubén Darío
  \item Gabriel García Márquez
  \item Julio Cortázar
  \item Mario Vargas Llosa
  \item Carlos Fuentes
  \item Juan Rulfo
\end{multicols}
\end{enumerate}
\vspace{0.15cm}
\item ¿Quién es el autor de Prosas profanas?
\begin{enumerate}[label=\textbf{\Alph*)}]
\begin{multicols}{2}
  \item Rubén Darío
  \item Gabriel García Márquez
  \item Julio Cortázar
  \item Mario Vargas Llosa
  \item Carlos Fuentes
  \item Juan Rulfo
\end{multicols}
\end{enumerate}
\vspace{0.15cm}
\item ¿Quién es el autor de Cantos de vida y esperanza?
\begin{enumerate}[label=\textbf{\Alph*)}]
\begin{multicols}{2}
  \item Rubén Darío
  \item Gabriel García Márquez
  \item Julio Cortázar
  \item Mario Vargas Llosa
  \item Carlos Fuentes
  \item Juan Rulfo
\end{multicols}
\end{enumerate}
\vspace{0.15cm}
\item ¿Qué es el Modernismo?
\begin{enumerate}[label=\textbf{\Alph*)}]
\begin{multicols}{2}
  \item Un movimiento literario hispanoamericano de finales del siglo XIX
  \item Gabriel García Márquez
  \item Julio Cortázar
  \item Mario Vargas Llosa
  \item Carlos Fuentes
  \item Juan Rulfo
\end{multicols}
\end{enumerate}
\vspace{0.15cm}
\item ¿Quién es el principal representante del Modernismo?
\begin{enumerate}[label=\textbf{\Alph*)}]
\begin{multicols}{2}
  \item Rubén Darío
  \item Gabriel García Márquez
  \item Julio Cortázar
  \item Mario Vargas Llosa
  \item Carlos Fuentes
  \item Juan Rulfo
\end{multicols}
\end{enumerate}
\vspace{0.15cm}
\item ¿Quién es el autor de Azul?
\begin{enumerate}[label=\textbf{\Alph*)}]
\begin{multicols}{2}
  \item Rubén Darío
  \item Gabriel García Márquez
  \item Julio Cortázar
  \item Mario Vargas Llosa
  \item Carlos Fuentes
  \item Juan Rulfo
\end{multicols}
\end{enumerate}
\vspace{0.15cm}
\item ¿Quién es el autor de Prosas profanas?
\begin{enumerate}[label=\textbf{\Alph*)}]
\begin{multicols}{2}
  \item Rubén Darío
  \item Gabriel García Márquez
  \item Julio Cortázar
  \item Mario Vargas Llosa
  \item Carlos Fuentes
  \item Juan Rulfo
\end{multicols}
\end{enumerate}
\vspace{0.15cm}
\item ¿Quién es el autor de Cantos de vida y esperanza?
\begin{enumerate}[label=\textbf{\Alph*)}]
\begin{multicols}{2}
  \item Rubén Darío
  \item Gabriel García Márquez
  \item Julio Cortázar
  \item Mario Vargas Llosa
  \item Carlos Fuentes
  \item Juan Rulfo
\end{multicols}
\end{enumerate}
\vspace{0.15cm}
\item ¿Qué es el Modernismo?
\begin{enumerate}[label=\textbf{\Alph*)}]
\begin{multicols}{2}
  \item Un movimiento literario hispanoamericano de finales del siglo XIX
  \item Gabriel García Márquez
  \item Julio Cortázar
  \item Mario Vargas Llosa
  \item Carlos Fuentes
  \item Juan Rulfo
\end{multicols}
\end{enumerate}
\vspace{0.15cm}
\item ¿Quién es el principal representante del Modernismo?
\begin{enumerate}[label=\textbf{\Alph*)}]
\begin{multicols}{2}
  \item Rubén Darío
  \item Gabriel García Márquez
  \item Julio Cortázar
  \item Mario Vargas Llosa
  \item Carlos Fuentes
  \item Juan Rulfo
\end{multicols}
\end{enumerate}
\vspace{0.15cm}
\item ¿Quién es el autor de Azul?
\begin{enumerate}[label=\textbf{\Alph*)}]
\begin{multicols}{2}
  \item Rubén Darío
  \item Gabriel García Márquez
  \item Julio Cortázar
  \item Mario Vargas Llosa
  \item Carlos Fuentes
  \item Juan Rulfo
\end{multicols}
\end{enumerate}
\vspace{0.15cm}
\item ¿Quién es el autor de Prosas profanas?
\begin{enumerate}[label=\textbf{\Alph*)}]
\begin{multicols}{2}
  \item Rubén Darío
  \item Gabriel García Márquez
  \item Julio Cortázar
  \item Mario Vargas Llosa
  \item Carlos Fuentes
  \item Juan Rulfo
\end{multicols}
\end{enumerate}
\vspace{0.15cm}
\item ¿Quién es el autor de Cantos de vida y esperanza?
\begin{enumerate}[label=\textbf{\Alph*)}]
\begin{multicols}{2}
  \item Rubén Darío
  \item Gabriel García Márquez
  \item Julio Cortázar
  \item Mario Vargas Llosa
  \item Carlos Fuentes
  \item Juan Rulfo
\end{multicols}
\end{enumerate}
\vspace{0.15cm}
\end{enumerate}\n
\subsubsection*{SECCIÓN V: REALISMO CRÍTICO Y NOVELA DE LA TIERRA (171–210)}

\begin{enumerate}[label=\textbf{\arabic*.}]
\setcounter{enumi}{170}
\item ¿Qué es el realismo crítico en Latinoamérica?
\begin{enumerate}[label=\textbf{\Alph*)}]
\begin{multicols}{2}
  \item Una corriente que retrata los problemas sociales y políticos de la región
  \item Gabriel García Márquez
  \item Julio Cortázar
  \item Mario Vargas Llosa
  \item Carlos Fuentes
  \item Juan Rulfo
\end{multicols}
\end{enumerate}
\vspace{0.15cm}
\item ¿Qué es la novela de la tierra?
\begin{enumerate}[label=\textbf{\Alph*)}]
\begin{multicols}{2}
  \item Una novela que retrata la naturaleza y los conflictos del hombre con ella
  \item Gabriel García Márquez
  \item Julio Cortázar
  \item Mario Vargas Llosa
  \item Carlos Fuentes
  \item Juan Rulfo
\end{multicols}
\end{enumerate}
\vspace{0.15cm}
\item ¿Quién es el autor de La vorágine?
\begin{enumerate}[label=\textbf{\Alph*)}]
\begin{multicols}{2}
  \item José Eustasio Rivera
  \item Gabriel García Márquez
  \item Julio Cortázar
  \item Mario Vargas Llosa
  \item Carlos Fuentes
  \item Juan Rulfo
\end{multicols}
\end{enumerate}
\vspace{0.15cm}
\item ¿Qué es La vorágine?
\begin{enumerate}[label=\textbf{\Alph*)}]
\begin{multicols}{2}
  \item Una novela de la selva colombiana
  \item Gabriel García Márquez
  \item Julio Cortázar
  \item Mario Vargas Llosa
  \item Carlos Fuentes
  \item Juan Rulfo
\end{multicols}
\end{enumerate}
\vspace{0.15cm}
\item ¿Quién es el autor de Doña Bárbara?
\begin{enumerate}[label=\textbf{\Alph*)}]
\begin{multicols}{2}
  \item Rómulo Gallegos
  \item Gabriel García Márquez
  \item Julio Cortázar
  \item Mario Vargas Llosa
  \item Carlos Fuentes
  \item Juan Rulfo
\end{multicols}
\end{enumerate}
\vspace{0.15cm}
\item ¿Qué es Doña Bárbara?
\begin{enumerate}[label=\textbf{\Alph*)}]
\begin{multicols}{2}
  \item Una novela de la llanura venezolana
  \item Gabriel García Márquez
  \item Julio Cortázar
  \item Mario Vargas Llosa
  \item Carlos Fuentes
  \item Juan Rulfo
\end{multicols}
\end{enumerate}
\vspace{0.15cm}
\item ¿Quién es el autor de Don Segundo Sombra?
\begin{enumerate}[label=\textbf{\Alph*)}]
\begin{multicols}{2}
  \item Ricardo Güiraldes
  \item Gabriel García Márquez
  \item Julio Cortázar
  \item Mario Vargas Llosa
  \item Carlos Fuentes
  \item Juan Rulfo
\end{multicols}
\end{enumerate}
\vspace{0.15cm}
\item ¿Qué es Don Segundo Sombra?
\begin{enumerate}[label=\textbf{\Alph*)}]
\begin{multicols}{2}
  \item Una novela gauchesca argentina
  \item Gabriel García Márquez
  \item Julio Cortázar
  \item Mario Vargas Llosa
  \item Carlos Fuentes
  \item Juan Rulfo
\end{multicols}
\end{enumerate}
\vspace{0.15cm}
\item ¿Quién es el autor de Los de abajo?
\begin{enumerate}[label=\textbf{\Alph*)}]
\begin{multicols}{2}
  \item Mariano Azuela
  \item Gabriel García Márquez
  \item Julio Cortázar
  \item Mario Vargas Llosa
  \item Carlos Fuentes
  \item Juan Rulfo
\end{multicols}
\end{enumerate}
\vspace{0.15cm}
\item ¿Qué es Los de abajo?
\begin{enumerate}[label=\textbf{\Alph*)}]
\begin{multicols}{2}
  \item Una novela de la Revolución Mexicana
  \item Gabriel García Márquez
  \item Julio Cortázar
  \item Mario Vargas Llosa
  \item Carlos Fuentes
  \item Juan Rulfo
\end{multicols}
\end{enumerate}
\vspace{0.15cm}
\item ¿Quién es el autor de El águila y la serpiente?
\begin{enumerate}[label=\textbf{\Alph*)}]
\begin{multicols}{2}
  \item Martín Luis Guzmán
  \item Gabriel García Márquez
  \item Julio Cortázar
  \item Mario Vargas Llosa
  \item Carlos Fuentes
  \item Juan Rulfo
\end{multicols}
\end{enumerate}
\vspace{0.15cm}
\item ¿Quién es el autor de Ulises criollo?
\begin{enumerate}[label=\textbf{\Alph*)}]
\begin{multicols}{2}
  \item José Vasconcelos
  \item Gabriel García Márquez
  \item Julio Cortázar
  \item Mario Vargas Llosa
  \item Carlos Fuentes
  \item Juan Rulfo
\end{multicols}
\end{enumerate}
\vspace{0.15cm}
\item ¿Quién es el autor de Huasipungo?
\begin{enumerate}[label=\textbf{\Alph*)}]
\begin{multicols}{2}
  \item Jorge Icaza
  \item Gabriel García Márquez
  \item Julio Cortázar
  \item Mario Vargas Llosa
  \item Carlos Fuentes
  \item Juan Rulfo
\end{multicols}
\end{enumerate}
\vspace{0.15cm}
\item ¿Qué es Huasipungo?
\begin{enumerate}[label=\textbf{\Alph*)}]
\begin{multicols}{2}
  \item Una novela indigenista ecuatoriana
  \item Gabriel García Márquez
  \item Julio Cortázar
  \item Mario Vargas Llosa
  \item Carlos Fuentes
  \item Juan Rulfo
\end{multicols}
\end{enumerate}
\vspace{0.15cm}
\item ¿Quién es el autor de El mundo es ancho y ajeno?
\begin{enumerate}[label=\textbf{\Alph*)}]
\begin{multicols}{2}
  \item Ciro Alegría
  \item Gabriel García Márquez
  \item Julio Cortázar
  \item Mario Vargas Llosa
  \item Carlos Fuentes
  \item Juan Rulfo
\end{multicols}
\end{enumerate}
\vspace{0.15cm}
\item ¿Qué es El mundo es ancho y ajeno?
\begin{enumerate}[label=\textbf{\Alph*)}]
\begin{multicols}{2}
  \item Una novela indigenista peruana
  \item Gabriel García Márquez
  \item Julio Cortázar
  \item Mario Vargas Llosa
  \item Carlos Fuentes
  \item Juan Rulfo
\end{multicols}
\end{enumerate}
\vspace{0.15cm}
\item ¿Quién es el autor de Yawar Fiesta?
\begin{enumerate}[label=\textbf{\Alph*)}]
\begin{multicols}{2}
  \item José María Arguedas
  \item Gabriel García Márquez
  \item Julio Cortázar
  \item Mario Vargas Llosa
  \item Carlos Fuentes
  \item Juan Rulfo
\end{multicols}
\end{enumerate}
\vspace{0.15cm}
\item ¿Quién es el autor de Los ríos profundos?
\begin{enumerate}[label=\textbf{\Alph*)}]
\begin{multicols}{2}
  \item José María Arguedas
  \item Gabriel García Márquez
  \item Julio Cortázar
  \item Mario Vargas Llosa
  \item Carlos Fuentes
  \item Juan Rulfo
\end{multicols}
\end{enumerate}
\vspace{0.15cm}
\item ¿Quién es el autor de Todas las sangres?
\begin{enumerate}[label=\textbf{\Alph*)}]
\begin{multicols}{2}
  \item José María Arguedas
  \item Gabriel García Márquez
  \item Julio Cortázar
  \item Mario Vargas Llosa
  \item Carlos Fuentes
  \item Juan Rulfo
\end{multicols}
\end{enumerate}
\vspace{0.15cm}
\item ¿Qué es el indigenismo?
\begin{enumerate}[label=\textbf{\Alph*)}]
\begin{multicols}{2}
  \item Una corriente literaria que reivindica la cultura indígena
  \item Gabriel García Márquez
  \item Julio Cortázar
  \item Mario Vargas Llosa
  \item Carlos Fuentes
  \item Juan Rulfo
\end{multicols}
\end{enumerate}
\vspace{0.15cm}
\item ¿Qué es el realismo crítico en Latinoamérica?
\begin{enumerate}[label=\textbf{\Alph*)}]
\begin{multicols}{2}
  \item Una corriente que retrata los problemas sociales y políticos de la región
  \item Gabriel García Márquez
  \item Julio Cortázar
  \item Mario Vargas Llosa
  \item Carlos Fuentes
  \item Juan Rulfo
\end{multicols}
\end{enumerate}
\vspace{0.15cm}
\item ¿Qué es la novela de la tierra?
\begin{enumerate}[label=\textbf{\Alph*)}]
\begin{multicols}{2}
  \item Una novela que retrata la naturaleza y los conflictos del hombre con ella
  \item Gabriel García Márquez
  \item Julio Cortázar
  \item Mario Vargas Llosa
  \item Carlos Fuentes
  \item Juan Rulfo
\end{multicols}
\end{enumerate}
\vspace{0.15cm}
\item ¿Quién es el autor de La vorágine?
\begin{enumerate}[label=\textbf{\Alph*)}]
\begin{multicols}{2}
  \item José Eustasio Rivera
  \item Gabriel García Márquez
  \item Julio Cortázar
  \item Mario Vargas Llosa
  \item Carlos Fuentes
  \item Juan Rulfo
\end{multicols}
\end{enumerate}
\vspace{0.15cm}
\item ¿Qué es La vorágine?
\begin{enumerate}[label=\textbf{\Alph*)}]
\begin{multicols}{2}
  \item Una novela de la selva colombiana
  \item Gabriel García Márquez
  \item Julio Cortázar
  \item Mario Vargas Llosa
  \item Carlos Fuentes
  \item Juan Rulfo
\end{multicols}
\end{enumerate}
\vspace{0.15cm}
\item ¿Quién es el autor de Doña Bárbara?
\begin{enumerate}[label=\textbf{\Alph*)}]
\begin{multicols}{2}
  \item Rómulo Gallegos
  \item Gabriel García Márquez
  \item Julio Cortázar
  \item Mario Vargas Llosa
  \item Carlos Fuentes
  \item Juan Rulfo
\end{multicols}
\end{enumerate}
\vspace{0.15cm}
\item ¿Qué es Doña Bárbara?
\begin{enumerate}[label=\textbf{\Alph*)}]
\begin{multicols}{2}
  \item Una novela de la llanura venezolana
  \item Gabriel García Márquez
  \item Julio Cortázar
  \item Mario Vargas Llosa
  \item Carlos Fuentes
  \item Juan Rulfo
\end{multicols}
\end{enumerate}
\vspace{0.15cm}
\item ¿Quién es el autor de Don Segundo Sombra?
\begin{enumerate}[label=\textbf{\Alph*)}]
\begin{multicols}{2}
  \item Ricardo Güiraldes
  \item Gabriel García Márquez
  \item Julio Cortázar
  \item Mario Vargas Llosa
  \item Carlos Fuentes
  \item Juan Rulfo
\end{multicols}
\end{enumerate}
\vspace{0.15cm}
\item ¿Qué es Don Segundo Sombra?
\begin{enumerate}[label=\textbf{\Alph*)}]
\begin{multicols}{2}
  \item Una novela gauchesca argentina
  \item Gabriel García Márquez
  \item Julio Cortázar
  \item Mario Vargas Llosa
  \item Carlos Fuentes
  \item Juan Rulfo
\end{multicols}
\end{enumerate}
\vspace{0.15cm}
\item ¿Quién es el autor de Los de abajo?
\begin{enumerate}[label=\textbf{\Alph*)}]
\begin{multicols}{2}
  \item Mariano Azuela
  \item Gabriel García Márquez
  \item Julio Cortázar
  \item Mario Vargas Llosa
  \item Carlos Fuentes
  \item Juan Rulfo
\end{multicols}
\end{enumerate}
\vspace{0.15cm}
\item ¿Qué es Los de abajo?
\begin{enumerate}[label=\textbf{\Alph*)}]
\begin{multicols}{2}
  \item Una novela de la Revolución Mexicana
  \item Gabriel García Márquez
  \item Julio Cortázar
  \item Mario Vargas Llosa
  \item Carlos Fuentes
  \item Juan Rulfo
\end{multicols}
\end{enumerate}
\vspace{0.15cm}
\item ¿Quién es el autor de El águila y la serpiente?
\begin{enumerate}[label=\textbf{\Alph*)}]
\begin{multicols}{2}
  \item Martín Luis Guzmán
  \item Gabriel García Márquez
  \item Julio Cortázar
  \item Mario Vargas Llosa
  \item Carlos Fuentes
  \item Juan Rulfo
\end{multicols}
\end{enumerate}
\vspace{0.15cm}
\item ¿Quién es el autor de Ulises criollo?
\begin{enumerate}[label=\textbf{\Alph*)}]
\begin{multicols}{2}
  \item José Vasconcelos
  \item Gabriel García Márquez
  \item Julio Cortázar
  \item Mario Vargas Llosa
  \item Carlos Fuentes
  \item Juan Rulfo
\end{multicols}
\end{enumerate}
\vspace{0.15cm}
\item ¿Quién es el autor de Huasipungo?
\begin{enumerate}[label=\textbf{\Alph*)}]
\begin{multicols}{2}
  \item Jorge Icaza
  \item Gabriel García Márquez
  \item Julio Cortázar
  \item Mario Vargas Llosa
  \item Carlos Fuentes
  \item Juan Rulfo
\end{multicols}
\end{enumerate}
\vspace{0.15cm}
\item ¿Qué es Huasipungo?
\begin{enumerate}[label=\textbf{\Alph*)}]
\begin{multicols}{2}
  \item Una novela indigenista ecuatoriana
  \item Gabriel García Márquez
  \item Julio Cortázar
  \item Mario Vargas Llosa
  \item Carlos Fuentes
  \item Juan Rulfo
\end{multicols}
\end{enumerate}
\vspace{0.15cm}
\item ¿Quién es el autor de El mundo es ancho y ajeno?
\begin{enumerate}[label=\textbf{\Alph*)}]
\begin{multicols}{2}
  \item Ciro Alegría
  \item Gabriel García Márquez
  \item Julio Cortázar
  \item Mario Vargas Llosa
  \item Carlos Fuentes
  \item Juan Rulfo
\end{multicols}
\end{enumerate}
\vspace{0.15cm}
\item ¿Qué es El mundo es ancho y ajeno?
\begin{enumerate}[label=\textbf{\Alph*)}]
\begin{multicols}{2}
  \item Una novela indigenista peruana
  \item Gabriel García Márquez
  \item Julio Cortázar
  \item Mario Vargas Llosa
  \item Carlos Fuentes
  \item Juan Rulfo
\end{multicols}
\end{enumerate}
\vspace{0.15cm}
\item ¿Quién es el autor de Yawar Fiesta?
\begin{enumerate}[label=\textbf{\Alph*)}]
\begin{multicols}{2}
  \item José María Arguedas
  \item Gabriel García Márquez
  \item Julio Cortázar
  \item Mario Vargas Llosa
  \item Carlos Fuentes
  \item Juan Rulfo
\end{multicols}
\end{enumerate}
\vspace{0.15cm}
\item ¿Quién es el autor de Los ríos profundos?
\begin{enumerate}[label=\textbf{\Alph*)}]
\begin{multicols}{2}
  \item José María Arguedas
  \item Gabriel García Márquez
  \item Julio Cortázar
  \item Mario Vargas Llosa
  \item Carlos Fuentes
  \item Juan Rulfo
\end{multicols}
\end{enumerate}
\vspace{0.15cm}
\item ¿Quién es el autor de Todas las sangres?
\begin{enumerate}[label=\textbf{\Alph*)}]
\begin{multicols}{2}
  \item José María Arguedas
  \item Gabriel García Márquez
  \item Julio Cortázar
  \item Mario Vargas Llosa
  \item Carlos Fuentes
  \item Juan Rulfo
\end{multicols}
\end{enumerate}
\vspace{0.15cm}
\item ¿Qué es el indigenismo?
\begin{enumerate}[label=\textbf{\Alph*)}]
\begin{multicols}{2}
  \item Una corriente literaria que reivindica la cultura indígena
  \item Gabriel García Márquez
  \item Julio Cortázar
  \item Mario Vargas Llosa
  \item Carlos Fuentes
  \item Juan Rulfo
\end{multicols}
\end{enumerate}
\vspace{0.15cm}
\end{enumerate}\n
\subsubsection*{SECCIÓN VI: BOOM LATINOAMERICANO Y REALISMO MÁGICO (211–260)}

\begin{enumerate}[label=\textbf{\arabic*.}]
\setcounter{enumi}{210}
\item ¿Qué es el Boom latinoamericano?
\begin{enumerate}[label=\textbf{\Alph*)}]
\begin{multicols}{2}
  \item Un fenómeno literario y editorial de los años 60 y 70
  \item Gabriel García Márquez
  \item Julio Cortázar
  \item Mario Vargas Llosa
  \item Carlos Fuentes
  \item Juan Rulfo
\end{multicols}
\end{enumerate}
\vspace{0.15cm}
\item ¿Quiénes son los principales representantes del Boom?
\begin{enumerate}[label=\textbf{\Alph*)}]
\begin{multicols}{2}
  \item García Márquez, Cortázar, Vargas Llosa y Fuentes
  \item Gabriel García Márquez
  \item Julio Cortázar
  \item Mario Vargas Llosa
  \item Carlos Fuentes
  \item Juan Rulfo
\end{multicols}
\end{enumerate}
\vspace{0.15cm}
\item ¿Qué es el realismo mágico?
\begin{enumerate}[label=\textbf{\Alph*)}]
\begin{multicols}{2}
  \item La mezcla de lo real y lo fantástico en la narración
  \item Gabriel García Márquez
  \item Julio Cortázar
  \item Mario Vargas Llosa
  \item Carlos Fuentes
  \item Juan Rulfo
\end{multicols}
\end{enumerate}
\vspace{0.15cm}
\item ¿Quién es el principal representante del realismo mágico?
\begin{enumerate}[label=\textbf{\Alph*)}]
\begin{multicols}{2}
  \item Gabriel García Márquez
  \item Julio Cortázar
  \item Mario Vargas Llosa
  \item Carlos Fuentes
  \item Juan Rulfo
  \item Jorge Luis Borges
\end{multicols}
\end{enumerate}
\vspace{0.15cm}
\item ¿Quién es el autor de Cien años de soledad?
\begin{enumerate}[label=\textbf{\Alph*)}]
\begin{multicols}{2}
  \item Gabriel García Márquez
  \item Julio Cortázar
  \item Mario Vargas Llosa
  \item Carlos Fuentes
  \item Juan Rulfo
  \item Jorge Luis Borges
\end{multicols}
\end{enumerate}
\vspace{0.15cm}
\item ¿Qué es Cien años de soledad?
\begin{enumerate}[label=\textbf{\Alph*)}]
\begin{multicols}{2}
  \item Una novela que narra la historia de la familia Buendía en Macondo
  \item Gabriel García Márquez
  \item Julio Cortázar
  \item Mario Vargas Llosa
  \item Carlos Fuentes
  \item Juan Rulfo
\end{multicols}
\end{enumerate}
\vspace{0.15cm}
\item ¿Quién es el autor de Rayuela?
\begin{enumerate}[label=\textbf{\Alph*)}]
\begin{multicols}{2}
  \item Julio Cortázar
  \item Gabriel García Márquez
  \item Mario Vargas Llosa
  \item Carlos Fuentes
  \item Juan Rulfo
  \item Jorge Luis Borges
\end{multicols}
\end{enumerate}
\vspace{0.15cm}
\item ¿Qué es Rayuela?
\begin{enumerate}[label=\textbf{\Alph*)}]
\begin{multicols}{2}
  \item Una novela experimental que propone múltiples lecturas
  \item Gabriel García Márquez
  \item Julio Cortázar
  \item Mario Vargas Llosa
  \item Carlos Fuentes
  \item Juan Rulfo
\end{multicols}
\end{enumerate}
\vspace{0.15cm}
\item ¿Quién es el autor de La ciudad y los perros?
\begin{enumerate}[label=\textbf{\Alph*)}]
\begin{multicols}{2}
  \item Mario Vargas Llosa
  \item Gabriel García Márquez
  \item Julio Cortázar
  \item Carlos Fuentes
  \item Juan Rulfo
  \item Jorge Luis Borges
\end{multicols}
\end{enumerate}
\vspace{0.15cm}
\item ¿Quién es el autor de La muerte de Artemio Cruz?
\begin{enumerate}[label=\textbf{\Alph*)}]
\begin{multicols}{2}
  \item Carlos Fuentes
  \item Gabriel García Márquez
  \item Julio Cortázar
  \item Mario Vargas Llosa
  \item Juan Rulfo
  \item Jorge Luis Borges
\end{multicols}
\end{enumerate}
\vspace{0.15cm}
\item ¿Quién es el autor de Pedro Páramo?
\begin{enumerate}[label=\textbf{\Alph*)}]
\begin{multicols}{2}
  \item Juan Rulfo
  \item Gabriel García Márquez
  \item Julio Cortázar
  \item Mario Vargas Llosa
  \item Carlos Fuentes
  \item Jorge Luis Borges
\end{multicols}
\end{enumerate}
\vspace{0.15cm}
\item ¿Qué es Pedro Páramo?
\begin{enumerate}[label=\textbf{\Alph*)}]
\begin{multicols}{2}
  \item Una novela que narra la historia de Comala y su cacique
  \item Gabriel García Márquez
  \item Julio Cortázar
  \item Mario Vargas Llosa
  \item Carlos Fuentes
  \item Juan Rulfo
\end{multicols}
\end{enumerate}
\vspace{0.15cm}
\item ¿Quién es el autor de Ficciones?
\begin{enumerate}[label=\textbf{\Alph*)}]
\begin{multicols}{2}
  \item Jorge Luis Borges
  \item Gabriel García Márquez
  \item Julio Cortázar
  \item Mario Vargas Llosa
  \item Carlos Fuentes
  \item Juan Rulfo
\end{multicols}
\end{enumerate}
\vspace{0.15cm}
\item ¿Quién es el autor de El túnel?
\begin{enumerate}[label=\textbf{\Alph*)}]
\begin{multicols}{2}
  \item Ernesto Sábato
  \item Gabriel García Márquez
  \item Julio Cortázar
  \item Mario Vargas Llosa
  \item Carlos Fuentes
  \item Juan Rulfo
\end{multicols}
\end{enumerate}
\vspace{0.15cm}
\item ¿Quién es el autor de La invención de Morel?
\begin{enumerate}[label=\textbf{\Alph*)}]
\begin{multicols}{2}
  \item Adolfo Bioy Casares
  \item Gabriel García Márquez
  \item Julio Cortázar
  \item Mario Vargas Llosa
  \item Carlos Fuentes
  \item Juan Rulfo
\end{multicols}
\end{enumerate}
\vspace{0.15cm}
\item ¿Quién es el autor de La tregua?
\begin{enumerate}[label=\textbf{\Alph*)}]
\begin{multicols}{2}
  \item Mario Benedetti
  \item Gabriel García Márquez
  \item Julio Cortázar
  \item Mario Vargas Llosa
  \item Carlos Fuentes
  \item Juan Rulfo
\end{multicols}
\end{enumerate}
\vspace{0.15cm}
\item ¿Quién es el autor de El astillero?
\begin{enumerate}[label=\textbf{\Alph*)}]
\begin{multicols}{2}
  \item Juan Carlos Onetti
  \item Gabriel García Márquez
  \item Julio Cortázar
  \item Mario Vargas Llosa
  \item Carlos Fuentes
  \item Juan Rulfo
\end{multicols}
\end{enumerate}
\vspace{0.15cm}
\item ¿Quién es el autor de Paradiso?
\begin{enumerate}[label=\textbf{\Alph*)}]
\begin{multicols}{2}
  \item José Lezama Lima
  \item Gabriel García Márquez
  \item Julio Cortázar
  \item Mario Vargas Llosa
  \item Carlos Fuentes
  \item Juan Rulfo
\end{multicols}
\end{enumerate}
\vspace{0.15cm}
\item ¿Quién es el autor de Tres tristes tigres?
\begin{enumerate}[label=\textbf{\Alph*)}]
\begin{multicols}{2}
  \item Guillermo Cabrera Infante
  \item Gabriel García Márquez
  \item Julio Cortázar
  \item Mario Vargas Llosa
  \item Carlos Fuentes
  \item Juan Rulfo
\end{multicols}
\end{enumerate}
\vspace{0.15cm}
\item ¿Quién es el autor de El nombre de la rosa?
\begin{enumerate}[label=\textbf{\Alph*)}]
\begin{multicols}{2}
  \item Umberto Eco
  \item Gabriel García Márquez
  \item Julio Cortázar
  \item Mario Vargas Llosa
  \item Carlos Fuentes
  \item Juan Rulfo
\end{multicols}
\end{enumerate}
\vspace{0.15cm}
\item ¿Qué es el Boom latinoamericano?
\begin{enumerate}[label=\textbf{\Alph*)}]
\begin{multicols}{2}
  \item Un fenómeno literario y editorial de los años 60 y 70
  \item Gabriel García Márquez
  \item Julio Cortázar
  \item Mario Vargas Llosa
  \item Carlos Fuentes
  \item Juan Rulfo
\end{multicols}
\end{enumerate}
\vspace{0.15cm}
\item ¿Quiénes son los principales representantes del Boom?
\begin{enumerate}[label=\textbf{\Alph*)}]
\begin{multicols}{2}
  \item García Márquez, Cortázar, Vargas Llosa y Fuentes
  \item Gabriel García Márquez
  \item Julio Cortázar
  \item Mario Vargas Llosa
  \item Carlos Fuentes
  \item Juan Rulfo
\end{multicols}
\end{enumerate}
\vspace{0.15cm}
\item ¿Qué es el realismo mágico?
\begin{enumerate}[label=\textbf{\Alph*)}]
\begin{multicols}{2}
  \item La mezcla de lo real y lo fantástico en la narración
  \item Gabriel García Márquez
  \item Julio Cortázar
  \item Mario Vargas Llosa
  \item Carlos Fuentes
  \item Juan Rulfo
\end{multicols}
\end{enumerate}
\vspace{0.15cm}
\item ¿Quién es el principal representante del realismo mágico?
\begin{enumerate}[label=\textbf{\Alph*)}]
\begin{multicols}{2}
  \item Gabriel García Márquez
  \item Julio Cortázar
  \item Mario Vargas Llosa
  \item Carlos Fuentes
  \item Juan Rulfo
  \item Jorge Luis Borges
\end{multicols}
\end{enumerate}
\vspace{0.15cm}
\item ¿Quién es el autor de Cien años de soledad?
\begin{enumerate}[label=\textbf{\Alph*)}]
\begin{multicols}{2}
  \item Gabriel García Márquez
  \item Julio Cortázar
  \item Mario Vargas Llosa
  \item Carlos Fuentes
  \item Juan Rulfo
  \item Jorge Luis Borges
\end{multicols}
\end{enumerate}
\vspace{0.15cm}
\item ¿Qué es Cien años de soledad?
\begin{enumerate}[label=\textbf{\Alph*)}]
\begin{multicols}{2}
  \item Una novela que narra la historia de la familia Buendía en Macondo
  \item Gabriel García Márquez
  \item Julio Cortázar
  \item Mario Vargas Llosa
  \item Carlos Fuentes
  \item Juan Rulfo
\end{multicols}
\end{enumerate}
\vspace{0.15cm}
\item ¿Quién es el autor de Rayuela?
\begin{enumerate}[label=\textbf{\Alph*)}]
\begin{multicols}{2}
  \item Julio Cortázar
  \item Gabriel García Márquez
  \item Mario Vargas Llosa
  \item Carlos Fuentes
  \item Juan Rulfo
  \item Jorge Luis Borges
\end{multicols}
\end{enumerate}
\vspace{0.15cm}
\item ¿Qué es Rayuela?
\begin{enumerate}[label=\textbf{\Alph*)}]
\begin{multicols}{2}
  \item Una novela experimental que propone múltiples lecturas
  \item Gabriel García Márquez
  \item Julio Cortázar
  \item Mario Vargas Llosa
  \item Carlos Fuentes
  \item Juan Rulfo
\end{multicols}
\end{enumerate}
\vspace{0.15cm}
\item ¿Quién es el autor de La ciudad y los perros?
\begin{enumerate}[label=\textbf{\Alph*)}]
\begin{multicols}{2}
  \item Mario Vargas Llosa
  \item Gabriel García Márquez
  \item Julio Cortázar
  \item Carlos Fuentes
  \item Juan Rulfo
  \item Jorge Luis Borges
\end{multicols}
\end{enumerate}
\vspace{0.15cm}
\item ¿Quién es el autor de La muerte de Artemio Cruz?
\begin{enumerate}[label=\textbf{\Alph*)}]
\begin{multicols}{2}
  \item Carlos Fuentes
  \item Gabriel García Márquez
  \item Julio Cortázar
  \item Mario Vargas Llosa
  \item Juan Rulfo
  \item Jorge Luis Borges
\end{multicols}
\end{enumerate}
\vspace{0.15cm}
\item ¿Quién es el autor de Pedro Páramo?
\begin{enumerate}[label=\textbf{\Alph*)}]
\begin{multicols}{2}
  \item Juan Rulfo
  \item Gabriel García Márquez
  \item Julio Cortázar
  \item Mario Vargas Llosa
  \item Carlos Fuentes
  \item Jorge Luis Borges
\end{multicols}
\end{enumerate}
\vspace{0.15cm}
\item ¿Qué es Pedro Páramo?
\begin{enumerate}[label=\textbf{\Alph*)}]
\begin{multicols}{2}
  \item Una novela que narra la historia de Comala y su cacique
  \item Gabriel García Márquez
  \item Julio Cortázar
  \item Mario Vargas Llosa
  \item Carlos Fuentes
  \item Juan Rulfo
\end{multicols}
\end{enumerate}
\vspace{0.15cm}
\item ¿Quién es el autor de Ficciones?
\begin{enumerate}[label=\textbf{\Alph*)}]
\begin{multicols}{2}
  \item Jorge Luis Borges
  \item Gabriel García Márquez
  \item Julio Cortázar
  \item Mario Vargas Llosa
  \item Carlos Fuentes
  \item Juan Rulfo
\end{multicols}
\end{enumerate}
\vspace{0.15cm}
\item ¿Quién es el autor de El túnel?
\begin{enumerate}[label=\textbf{\Alph*)}]
\begin{multicols}{2}
  \item Ernesto Sábato
  \item Gabriel García Márquez
  \item Julio Cortázar
  \item Mario Vargas Llosa
  \item Carlos Fuentes
  \item Juan Rulfo
\end{multicols}
\end{enumerate}
\vspace{0.15cm}
\item ¿Quién es el autor de La invención de Morel?
\begin{enumerate}[label=\textbf{\Alph*)}]
\begin{multicols}{2}
  \item Adolfo Bioy Casares
  \item Gabriel García Márquez
  \item Julio Cortázar
  \item Mario Vargas Llosa
  \item Carlos Fuentes
  \item Juan Rulfo
\end{multicols}
\end{enumerate}
\vspace{0.15cm}
\item ¿Quién es el autor de La tregua?
\begin{enumerate}[label=\textbf{\Alph*)}]
\begin{multicols}{2}
  \item Mario Benedetti
  \item Gabriel García Márquez
  \item Julio Cortázar
  \item Mario Vargas Llosa
  \item Carlos Fuentes
  \item Juan Rulfo
\end{multicols}
\end{enumerate}
\vspace{0.15cm}
\item ¿Quién es el autor de El astillero?
\begin{enumerate}[label=\textbf{\Alph*)}]
\begin{multicols}{2}
  \item Juan Carlos Onetti
  \item Gabriel García Márquez
  \item Julio Cortázar
  \item Mario Vargas Llosa
  \item Carlos Fuentes
  \item Juan Rulfo
\end{multicols}
\end{enumerate}
\vspace{0.15cm}
\item ¿Quién es el autor de Paradiso?
\begin{enumerate}[label=\textbf{\Alph*)}]
\begin{multicols}{2}
  \item José Lezama Lima
  \item Gabriel García Márquez
  \item Julio Cortázar
  \item Mario Vargas Llosa
  \item Carlos Fuentes
  \item Juan Rulfo
\end{multicols}
\end{enumerate}
\vspace{0.15cm}
\item ¿Quién es el autor de Tres tristes tigres?
\begin{enumerate}[label=\textbf{\Alph*)}]
\begin{multicols}{2}
  \item Guillermo Cabrera Infante
  \item Gabriel García Márquez
  \item Julio Cortázar
  \item Mario Vargas Llosa
  \item Carlos Fuentes
  \item Juan Rulfo
\end{multicols}
\end{enumerate}
\vspace{0.15cm}
\item ¿Quién es el autor de El nombre de la rosa?
\begin{enumerate}[label=\textbf{\Alph*)}]
\begin{multicols}{2}
  \item Umberto Eco
  \item Gabriel García Márquez
  \item Julio Cortázar
  \item Mario Vargas Llosa
  \item Carlos Fuentes
  \item Juan Rulfo
\end{multicols}
\end{enumerate}
\vspace{0.15cm}
\item ¿Qué es el Boom latinoamericano?
\begin{enumerate}[label=\textbf{\Alph*)}]
\begin{multicols}{2}
  \item Un fenómeno literario y editorial de los años 60 y 70
  \item Gabriel García Márquez
  \item Julio Cortázar
  \item Mario Vargas Llosa
  \item Carlos Fuentes
  \item Juan Rulfo
\end{multicols}
\end{enumerate}
\vspace{0.15cm}
\item ¿Quiénes son los principales representantes del Boom?
\begin{enumerate}[label=\textbf{\Alph*)}]
\begin{multicols}{2}
  \item García Márquez, Cortázar, Vargas Llosa y Fuentes
  \item Gabriel García Márquez
  \item Julio Cortázar
  \item Mario Vargas Llosa
  \item Carlos Fuentes
  \item Juan Rulfo
\end{multicols}
\end{enumerate}
\vspace{0.15cm}
\item ¿Qué es el realismo mágico?
\begin{enumerate}[label=\textbf{\Alph*)}]
\begin{multicols}{2}
  \item La mezcla de lo real y lo fantástico en la narración
  \item Gabriel García Márquez
  \item Julio Cortázar
  \item Mario Vargas Llosa
  \item Carlos Fuentes
  \item Juan Rulfo
\end{multicols}
\end{enumerate}
\vspace{0.15cm}
\item ¿Quién es el principal representante del realismo mágico?
\begin{enumerate}[label=\textbf{\Alph*)}]
\begin{multicols}{2}
  \item Gabriel García Márquez
  \item Julio Cortázar
  \item Mario Vargas Llosa
  \item Carlos Fuentes
  \item Juan Rulfo
  \item Jorge Luis Borges
\end{multicols}
\end{enumerate}
\vspace{0.15cm}
\item ¿Quién es el autor de Cien años de soledad?
\begin{enumerate}[label=\textbf{\Alph*)}]
\begin{multicols}{2}
  \item Gabriel García Márquez
  \item Julio Cortázar
  \item Mario Vargas Llosa
  \item Carlos Fuentes
  \item Juan Rulfo
  \item Jorge Luis Borges
\end{multicols}
\end{enumerate}
\vspace{0.15cm}
\item ¿Qué es Cien años de soledad?
\begin{enumerate}[label=\textbf{\Alph*)}]
\begin{multicols}{2}
  \item Una novela que narra la historia de la familia Buendía en Macondo
  \item Gabriel García Márquez
  \item Julio Cortázar
  \item Mario Vargas Llosa
  \item Carlos Fuentes
  \item Juan Rulfo
\end{multicols}
\end{enumerate}
\vspace{0.15cm}
\item ¿Quién es el autor de Rayuela?
\begin{enumerate}[label=\textbf{\Alph*)}]
\begin{multicols}{2}
  \item Julio Cortázar
  \item Gabriel García Márquez
  \item Mario Vargas Llosa
  \item Carlos Fuentes
  \item Juan Rulfo
  \item Jorge Luis Borges
\end{multicols}
\end{enumerate}
\vspace{0.15cm}
\item ¿Qué es Rayuela?
\begin{enumerate}[label=\textbf{\Alph*)}]
\begin{multicols}{2}
  \item Una novela experimental que propone múltiples lecturas
  \item Gabriel García Márquez
  \item Julio Cortázar
  \item Mario Vargas Llosa
  \item Carlos Fuentes
  \item Juan Rulfo
\end{multicols}
\end{enumerate}
\vspace{0.15cm}
\item ¿Quién es el autor de La ciudad y los perros?
\begin{enumerate}[label=\textbf{\Alph*)}]
\begin{multicols}{2}
  \item Mario Vargas Llosa
  \item Gabriel García Márquez
  \item Julio Cortázar
  \item Carlos Fuentes
  \item Juan Rulfo
  \item Jorge Luis Borges
\end{multicols}
\end{enumerate}
\vspace{0.15cm}
\item ¿Quién es el autor de La muerte de Artemio Cruz?
\begin{enumerate}[label=\textbf{\Alph*)}]
\begin{multicols}{2}
  \item Carlos Fuentes
  \item Gabriel García Márquez
  \item Julio Cortázar
  \item Mario Vargas Llosa
  \item Juan Rulfo
  \item Jorge Luis Borges
\end{multicols}
\end{enumerate}
\vspace{0.15cm}
\end{enumerate}\n
\subsubsection*{SECCIÓN VII: LITERATURA CONTEMPORÁNEA Y TEMAS (261–300)}

\begin{enumerate}[label=\textbf{\arabic*.}]
\setcounter{enumi}{260}
\item ¿Qué es la literatura contemporánea latinoamericana?
\begin{enumerate}[label=\textbf{\Alph*)}]
\begin{multicols}{2}
  \item La literatura producida desde mediados del siglo XX hasta la actualidad
  \item Gabriel García Márquez
  \item Julio Cortázar
  \item Mario Vargas Llosa
  \item Carlos Fuentes
  \item Juan Rulfo
\end{multicols}
\end{enumerate}
\vspace{0.15cm}
\item ¿Qué es la literatura del posboom?
\begin{enumerate}[label=\textbf{\Alph*)}]
\begin{multicols}{2}
  \item La literatura posterior al Boom que explora nuevas formas
  \item Gabriel García Márquez
  \item Julio Cortázar
  \item Mario Vargas Llosa
  \item Carlos Fuentes
  \item Juan Rulfo
\end{multicols}
\end{enumerate}
\vspace{0.15cm}
\item ¿Quiénes son representantes del posboom?
\begin{enumerate}[label=\textbf{\Alph*)}]
\begin{multicols}{2}
  \item Isabel Allende, Laura Esquivel y Roberto Bolaño
  \item Gabriel García Márquez
  \item Julio Cortázar
  \item Mario Vargas Llosa
  \item Carlos Fuentes
  \item Juan Rulfo
\end{multicols}
\end{enumerate}
\vspace{0.15cm}
\item ¿Quién es la autora de La casa de los espíritus?
\begin{enumerate}[label=\textbf{\Alph*)}]
\begin{multicols}{2}
  \item Isabel Allende
  \item Gabriel García Márquez
  \item Julio Cortázar
  \item Mario Vargas Llosa
  \item Carlos Fuentes
  \item Juan Rulfo
\end{multicols}
\end{enumerate}
\vspace{0.15cm}
\item ¿Quién es la autora de Como agua para chocolate?
\begin{enumerate}[label=\textbf{\Alph*)}]
\begin{multicols}{2}
  \item Laura Esquivel
  \item Gabriel García Márquez
  \item Julio Cortázar
  \item Mario Vargas Llosa
  \item Carlos Fuentes
  \item Juan Rulfo
\end{multicols}
\end{enumerate}
\vspace{0.15cm}
\item ¿Quién es el autor de Los detectives salvajes?
\begin{enumerate}[label=\textbf{\Alph*)}]
\begin{multicols}{2}
  \item Roberto Bolaño
  \item Gabriel García Márquez
  \item Julio Cortázar
  \item Mario Vargas Llosa
  \item Carlos Fuentes
  \item Juan Rulfo
\end{multicols}
\end{enumerate}
\vspace{0.15cm}
\item ¿Quién es el autor de 2666?
\begin{enumerate}[label=\textbf{\Alph*)}]
\begin{multicols}{2}
  \item Roberto Bolaño
  \item Gabriel García Márquez
  \item Julio Cortázar
  \item Mario Vargas Llosa
  \item Carlos Fuentes
  \item Juan Rulfo
\end{multicols}
\end{enumerate}
\vspace{0.15cm}
\item ¿Quién es el autor de La noche de Tlatelolco?
\begin{enumerate}[label=\textbf{\Alph*)}]
\begin{multicols}{2}
  \item Elena Poniatowska
  \item Gabriel García Márquez
  \item Julio Cortázar
  \item Mario Vargas Llosa
  \item Carlos Fuentes
  \item Juan Rulfo
\end{multicols}
\end{enumerate}
\vspace{0.15cm}
\item ¿Quién es el autor de Arráncame la vida?
\begin{enumerate}[label=\textbf{\Alph*)}]
\begin{multicols}{2}
  \item Ángeles Mastretta
  \item Gabriel García Márquez
  \item Julio Cortázar
  \item Mario Vargas Llosa
  \item Carlos Fuentes
  \item Juan Rulfo
\end{multicols}
\end{enumerate}
\vspace{0.15cm}
\item ¿Quién es el autor de El beso de la mujer araña?
\begin{enumerate}[label=\textbf{\Alph*)}]
\begin{multicols}{2}
  \item Manuel Puig
  \item Gabriel García Márquez
  \item Julio Cortázar
  \item Mario Vargas Llosa
  \item Carlos Fuentes
  \item Juan Rulfo
\end{multicols}
\end{enumerate}
\vspace{0.15cm}
\item ¿Qué es la literatura contemporánea latinoamericana?
\begin{enumerate}[label=\textbf{\Alph*)}]
\begin{multicols}{2}
  \item La literatura producida desde mediados del siglo XX hasta la actualidad
  \item Gabriel García Márquez
  \item Julio Cortázar
  \item Mario Vargas Llosa
  \item Carlos Fuentes
  \item Juan Rulfo
\end{multicols}
\end{enumerate}
\vspace{0.15cm}
\item ¿Qué es la literatura del posboom?
\begin{enumerate}[label=\textbf{\Alph*)}]
\begin{multicols}{2}
  \item La literatura posterior al Boom que explora nuevas formas
  \item Gabriel García Márquez
  \item Julio Cortázar
  \item Mario Vargas Llosa
  \item Carlos Fuentes
  \item Juan Rulfo
\end{multicols}
\end{enumerate}
\vspace{0.15cm}
\item ¿Quiénes son representantes del posboom?
\begin{enumerate}[label=\textbf{\Alph*)}]
\begin{multicols}{2}
  \item Isabel Allende, Laura Esquivel y Roberto Bolaño
  \item Gabriel García Márquez
  \item Julio Cortázar
  \item Mario Vargas Llosa
  \item Carlos Fuentes
  \item Juan Rulfo
\end{multicols}
\end{enumerate}
\vspace{0.15cm}
\item ¿Quién es la autora de La casa de los espíritus?
\begin{enumerate}[label=\textbf{\Alph*)}]
\begin{multicols}{2}
  \item Isabel Allende
  \item Gabriel García Márquez
  \item Julio Cortázar
  \item Mario Vargas Llosa
  \item Carlos Fuentes
  \item Juan Rulfo
\end{multicols}
\end{enumerate}
\vspace{0.15cm}
\item ¿Quién es la autora de Como agua para chocolate?
\begin{enumerate}[label=\textbf{\Alph*)}]
\begin{multicols}{2}
  \item Laura Esquivel
  \item Gabriel García Márquez
  \item Julio Cortázar
  \item Mario Vargas Llosa
  \item Carlos Fuentes
  \item Juan Rulfo
\end{multicols}
\end{enumerate}
\vspace{0.15cm}
\item ¿Quién es el autor de Los detectives salvajes?
\begin{enumerate}[label=\textbf{\Alph*)}]
\begin{multicols}{2}
  \item Roberto Bolaño
  \item Gabriel García Márquez
  \item Julio Cortázar
  \item Mario Vargas Llosa
  \item Carlos Fuentes
  \item Juan Rulfo
\end{multicols}
\end{enumerate}
\vspace{0.15cm}
\item ¿Quién es el autor de 2666?
\begin{enumerate}[label=\textbf{\Alph*)}]
\begin{multicols}{2}
  \item Roberto Bolaño
  \item Gabriel García Márquez
  \item Julio Cortázar
  \item Mario Vargas Llosa
  \item Carlos Fuentes
  \item Juan Rulfo
\end{multicols}
\end{enumerate}
\vspace{0.15cm}
\item ¿Quién es el autor de La noche de Tlatelolco?
\begin{enumerate}[label=\textbf{\Alph*)}]
\begin{multicols}{2}
  \item Elena Poniatowska
  \item Gabriel García Márquez
  \item Julio Cortázar
  \item Mario Vargas Llosa
  \item Carlos Fuentes
  \item Juan Rulfo
\end{multicols}
\end{enumerate}
\vspace{0.15cm}
\item ¿Quién es el autor de Arráncame la vida?
\begin{enumerate}[label=\textbf{\Alph*)}]
\begin{multicols}{2}
  \item Ángeles Mastretta
  \item Gabriel García Márquez
  \item Julio Cortázar
  \item Mario Vargas Llosa
  \item Carlos Fuentes
  \item Juan Rulfo
\end{multicols}
\end{enumerate}
\vspace{0.15cm}
\item ¿Quién es el autor de El beso de la mujer araña?
\begin{enumerate}[label=\textbf{\Alph*)}]
\begin{multicols}{2}
  \item Manuel Puig
  \item Gabriel García Márquez
  \item Julio Cortázar
  \item Mario Vargas Llosa
  \item Carlos Fuentes
  \item Juan Rulfo
\end{multicols}
\end{enumerate}
\vspace{0.15cm}
\item ¿Qué es la literatura contemporánea latinoamericana?
\begin{enumerate}[label=\textbf{\Alph*)}]
\begin{multicols}{2}
  \item La literatura producida desde mediados del siglo XX hasta la actualidad
  \item Gabriel García Márquez
  \item Julio Cortázar
  \item Mario Vargas Llosa
  \item Carlos Fuentes
  \item Juan Rulfo
\end{multicols}
\end{enumerate}
\vspace{0.15cm}
\item ¿Qué es la literatura del posboom?
\begin{enumerate}[label=\textbf{\Alph*)}]
\begin{multicols}{2}
  \item La literatura posterior al Boom que explora nuevas formas
  \item Gabriel García Márquez
  \item Julio Cortázar
  \item Mario Vargas Llosa
  \item Carlos Fuentes
  \item Juan Rulfo
\end{multicols}
\end{enumerate}
\vspace{0.15cm}
\item ¿Quiénes son representantes del posboom?
\begin{enumerate}[label=\textbf{\Alph*)}]
\begin{multicols}{2}
  \item Isabel Allende, Laura Esquivel y Roberto Bolaño
  \item Gabriel García Márquez
  \item Julio Cortázar
  \item Mario Vargas Llosa
  \item Carlos Fuentes
  \item Juan Rulfo
\end{multicols}
\end{enumerate}
\vspace{0.15cm}
\item ¿Quién es la autora de La casa de los espíritus?
\begin{enumerate}[label=\textbf{\Alph*)}]
\begin{multicols}{2}
  \item Isabel Allende
  \item Gabriel García Márquez
  \item Julio Cortázar
  \item Mario Vargas Llosa
  \item Carlos Fuentes
  \item Juan Rulfo
\end{multicols}
\end{enumerate}
\vspace{0.15cm}
\item ¿Quién es la autora de Como agua para chocolate?
\begin{enumerate}[label=\textbf{\Alph*)}]
\begin{multicols}{2}
  \item Laura Esquivel
  \item Gabriel García Márquez
  \item Julio Cortázar
  \item Mario Vargas Llosa
  \item Carlos Fuentes
  \item Juan Rulfo
\end{multicols}
\end{enumerate}
\vspace{0.15cm}
\item ¿Quién es el autor de Los detectives salvajes?
\begin{enumerate}[label=\textbf{\Alph*)}]
\begin{multicols}{2}
  \item Roberto Bolaño
  \item Gabriel García Márquez
  \item Julio Cortázar
  \item Mario Vargas Llosa
  \item Carlos Fuentes
  \item Juan Rulfo
\end{multicols}
\end{enumerate}
\vspace{0.15cm}
\item ¿Quién es el autor de 2666?
\begin{enumerate}[label=\textbf{\Alph*)}]
\begin{multicols}{2}
  \item Roberto Bolaño
  \item Gabriel García Márquez
  \item Julio Cortázar
  \item Mario Vargas Llosa
  \item Carlos Fuentes
  \item Juan Rulfo
\end{multicols}
\end{enumerate}
\vspace{0.15cm}
\item ¿Quién es el autor de La noche de Tlatelolco?
\begin{enumerate}[label=\textbf{\Alph*)}]
\begin{multicols}{2}
  \item Elena Poniatowska
  \item Gabriel García Márquez
  \item Julio Cortázar
  \item Mario Vargas Llosa
  \item Carlos Fuentes
  \item Juan Rulfo
\end{multicols}
\end{enumerate}
\vspace{0.15cm}
\item ¿Quién es el autor de Arráncame la vida?
\begin{enumerate}[label=\textbf{\Alph*)}]
\begin{multicols}{2}
  \item Ángeles Mastretta
  \item Gabriel García Márquez
  \item Julio Cortázar
  \item Mario Vargas Llosa
  \item Carlos Fuentes
  \item Juan Rulfo
\end{multicols}
\end{enumerate}
\vspace{0.15cm}
\item ¿Quién es el autor de El beso de la mujer araña?
\begin{enumerate}[label=\textbf{\Alph*)}]
\begin{multicols}{2}
  \item Manuel Puig
  \item Gabriel García Márquez
  \item Julio Cortázar
  \item Mario Vargas Llosa
  \item Carlos Fuentes
  \item Juan Rulfo
\end{multicols}
\end{enumerate}
\vspace{0.15cm}
\item ¿Qué es la literatura contemporánea latinoamericana?
\begin{enumerate}[label=\textbf{\Alph*)}]
\begin{multicols}{2}
  \item La literatura producida desde mediados del siglo XX hasta la actualidad
  \item Gabriel García Márquez
  \item Julio Cortázar
  \item Mario Vargas Llosa
  \item Carlos Fuentes
  \item Juan Rulfo
\end{multicols}
\end{enumerate}
\vspace{0.15cm}
\item ¿Qué es la literatura del posboom?
\begin{enumerate}[label=\textbf{\Alph*)}]
\begin{multicols}{2}
  \item La literatura posterior al Boom que explora nuevas formas
  \item Gabriel García Márquez
  \item Julio Cortázar
  \item Mario Vargas Llosa
  \item Carlos Fuentes
  \item Juan Rulfo
\end{multicols}
\end{enumerate}
\vspace{0.15cm}
\item ¿Quiénes son representantes del posboom?
\begin{enumerate}[label=\textbf{\Alph*)}]
\begin{multicols}{2}
  \item Isabel Allende, Laura Esquivel y Roberto Bolaño
  \item Gabriel García Márquez
  \item Julio Cortázar
  \item Mario Vargas Llosa
  \item Carlos Fuentes
  \item Juan Rulfo
\end{multicols}
\end{enumerate}
\vspace{0.15cm}
\item ¿Quién es la autora de La casa de los espíritus?
\begin{enumerate}[label=\textbf{\Alph*)}]
\begin{multicols}{2}
  \item Isabel Allende
  \item Gabriel García Márquez
  \item Julio Cortázar
  \item Mario Vargas Llosa
  \item Carlos Fuentes
  \item Juan Rulfo
\end{multicols}
\end{enumerate}
\vspace{0.15cm}
\item ¿Quién es la autora de Como agua para chocolate?
\begin{enumerate}[label=\textbf{\Alph*)}]
\begin{multicols}{2}
  \item Laura Esquivel
  \item Gabriel García Márquez
  \item Julio Cortázar
  \item Mario Vargas Llosa
  \item Carlos Fuentes
  \item Juan Rulfo
\end{multicols}
\end{enumerate}
\vspace{0.15cm}
\item ¿Quién es el autor de Los detectives salvajes?
\begin{enumerate}[label=\textbf{\Alph*)}]
\begin{multicols}{2}
  \item Roberto Bolaño
  \item Gabriel García Márquez
  \item Julio Cortázar
  \item Mario Vargas Llosa
  \item Carlos Fuentes
  \item Juan Rulfo
\end{multicols}
\end{enumerate}
\vspace{0.15cm}
\item ¿Quién es el autor de 2666?
\begin{enumerate}[label=\textbf{\Alph*)}]
\begin{multicols}{2}
  \item Roberto Bolaño
  \item Gabriel García Márquez
  \item Julio Cortázar
  \item Mario Vargas Llosa
  \item Carlos Fuentes
  \item Juan Rulfo
\end{multicols}
\end{enumerate}
\vspace{0.15cm}
\item ¿Quién es el autor de La noche de Tlatelolco?
\begin{enumerate}[label=\textbf{\Alph*)}]
\begin{multicols}{2}
  \item Elena Poniatowska
  \item Gabriel García Márquez
  \item Julio Cortázar
  \item Mario Vargas Llosa
  \item Carlos Fuentes
  \item Juan Rulfo
\end{multicols}
\end{enumerate}
\vspace{0.15cm}
\item ¿Quién es el autor de Arráncame la vida?
\begin{enumerate}[label=\textbf{\Alph*)}]
\begin{multicols}{2}
  \item Ángeles Mastretta
  \item Gabriel García Márquez
  \item Julio Cortázar
  \item Mario Vargas Llosa
  \item Carlos Fuentes
  \item Juan Rulfo
\end{multicols}
\end{enumerate}
\vspace{0.15cm}
\item ¿Quién es el autor de El beso de la mujer araña?
\begin{enumerate}[label=\textbf{\Alph*)}]
\begin{multicols}{2}
  \item Manuel Puig
  \item Gabriel García Márquez
  \item Julio Cortázar
  \item Mario Vargas Llosa
  \item Carlos Fuentes
  \item Juan Rulfo
\end{multicols}
\end{enumerate}
\vspace{0.15cm}
\end{enumerate}\n
\vspace{0.4cm}
\begin{tcolorbox}[colback=gray!5!white,colframe=gray!60!black,title=\textbf{SOLUCIONARIO -- 3.4 Literatura Latinoamericana}]
\begin{multicols}{4}
\begin{enumerate}[label=\textbf{\arabic*.}]
  \item \textbf{A}
  \item \textbf{A}
  \item \textbf{A}
  \item \textbf{A}
  \item \textbf{A}
  \item \textbf{A}
  \item \textbf{A}
  \item \textbf{A}
  \item \textbf{A}
  \item \textbf{A}
  \item \textbf{A}
  \item \textbf{A}
  \item \textbf{A}
  \item \textbf{A}
  \item \textbf{A}
  \item \textbf{A}
  \item \textbf{A}
  \item \textbf{A}
  \item \textbf{A}
  \item \textbf{A}
  \item \textbf{A}
  \item \textbf{A}
  \item \textbf{A}
  \item \textbf{A}
  \item \textbf{A}
  \item \textbf{A}
  \item \textbf{A}
  \item \textbf{A}
  \item \textbf{A}
  \item \textbf{A}
  \item \textbf{A}
  \item \textbf{A}
  \item \textbf{A}
  \item \textbf{A}
  \item \textbf{A}
  \item \textbf{A}
  \item \textbf{A}
  \item \textbf{A}
  \item \textbf{A}
  \item \textbf{A}
  \item \textbf{A}
  \item \textbf{A}
  \item \textbf{A}
  \item \textbf{A}
  \item \textbf{A}
  \item \textbf{A}
  \item \textbf{A}
  \item \textbf{A}
  \item \textbf{A}
  \item \textbf{A}
  \item \textbf{A}
  \item \textbf{A}
  \item \textbf{A}
  \item \textbf{A}
  \item \textbf{A}
  \item \textbf{A}
  \item \textbf{A}
  \item \textbf{A}
  \item \textbf{A}
  \item \textbf{A}
  \item \textbf{A}
  \item \textbf{A}
  \item \textbf{A}
  \item \textbf{A}
  \item \textbf{A}
  \item \textbf{A}
  \item \textbf{A}
  \item \textbf{A}
  \item \textbf{A}
  \item \textbf{A}
  \item \textbf{A}
  \item \textbf{A}
  \item \textbf{A}
  \item \textbf{A}
  \item \textbf{A}
  \item \textbf{A}
  \item \textbf{A}
  \item \textbf{A}
  \item \textbf{A}
  \item \textbf{A}
  \item \textbf{A}
  \item \textbf{A}
  \item \textbf{A}
  \item \textbf{A}
  \item \textbf{A}
  \item \textbf{A}
  \item \textbf{A}
  \item \textbf{A}
  \item \textbf{A}
  \item \textbf{A}
  \item \textbf{A}
  \item \textbf{A}
  \item \textbf{A}
  \item \textbf{A}
  \item \textbf{A}
  \item \textbf{A}
  \item \textbf{A}
  \item \textbf{A}
  \item \textbf{A}
  \item \textbf{A}
  \item \textbf{A}
  \item \textbf{A}
  \item \textbf{A}
  \item \textbf{A}
  \item \textbf{A}
  \item \textbf{A}
  \item \textbf{A}
  \item \textbf{A}
  \item \textbf{A}
  \item \textbf{A}
  \item \textbf{A}
  \item \textbf{A}
  \item \textbf{A}
  \item \textbf{A}
  \item \textbf{A}
  \item \textbf{A}
  \item \textbf{A}
  \item \textbf{A}
  \item \textbf{A}
  \item \textbf{A}
  \item \textbf{A}
  \item \textbf{A}
  \item \textbf{A}
  \item \textbf{A}
  \item \textbf{A}
  \item \textbf{A}
  \item \textbf{A}
  \item \textbf{A}
  \item \textbf{A}
  \item \textbf{A}
  \item \textbf{A}
  \item \textbf{A}
  \item \textbf{A}
  \item \textbf{A}
  \item \textbf{A}
  \item \textbf{A}
  \item \textbf{A}
  \item \textbf{A}
  \item \textbf{A}
  \item \textbf{A}
  \item \textbf{A}
  \item \textbf{A}
  \item \textbf{A}
  \item \textbf{A}
  \item \textbf{A}
  \item \textbf{A}
  \item \textbf{A}
  \item \textbf{A}
  \item \textbf{A}
  \item \textbf{A}
  \item \textbf{A}
  \item \textbf{A}
  \item \textbf{A}
  \item \textbf{A}
  \item \textbf{A}
  \item \textbf{A}
  \item \textbf{A}
  \item \textbf{A}
  \item \textbf{A}
  \item \textbf{A}
  \item \textbf{A}
  \item \textbf{A}
  \item \textbf{A}
  \item \textbf{A}
  \item \textbf{A}
  \item \textbf{A}
  \item \textbf{A}
  \item \textbf{A}
  \item \textbf{A}
  \item \textbf{A}
  \item \textbf{A}
  \item \textbf{A}
  \item \textbf{A}
  \item \textbf{A}
  \item \textbf{A}
  \item \textbf{A}
  \item \textbf{A}
  \item \textbf{A}
  \item \textbf{A}
  \item \textbf{A}
  \item \textbf{A}
  \item \textbf{A}
  \item \textbf{A}
  \item \textbf{A}
  \item \textbf{A}
  \item \textbf{A}
  \item \textbf{A}
  \item \textbf{A}
  \item \textbf{A}
  \item \textbf{A}
  \item \textbf{A}
  \item \textbf{A}
  \item \textbf{A}
  \item \textbf{A}
  \item \textbf{A}
  \item \textbf{A}
  \item \textbf{A}
  \item \textbf{A}
  \item \textbf{A}
  \item \textbf{A}
  \item \textbf{A}
  \item \textbf{A}
  \item \textbf{A}
  \item \textbf{A}
  \item \textbf{A}
  \item \textbf{A}
  \item \textbf{A}
  \item \textbf{A}
  \item \textbf{A}
  \item \textbf{A}
  \item \textbf{A}
  \item \textbf{A}
  \item \textbf{A}
  \item \textbf{A}
  \item \textbf{A}
  \item \textbf{A}
  \item \textbf{A}
  \item \textbf{A}
  \item \textbf{A}
  \item \textbf{A}
  \item \textbf{A}
  \item \textbf{A}
  \item \textbf{A}
  \item \textbf{A}
  \item \textbf{A}
  \item \textbf{A}
  \item \textbf{A}
  \item \textbf{A}
  \item \textbf{A}
  \item \textbf{A}
  \item \textbf{A}
  \item \textbf{A}
  \item \textbf{A}
  \item \textbf{A}
  \item \textbf{A}
  \item \textbf{A}
  \item \textbf{A}
  \item \textbf{A}
  \item \textbf{A}
  \item \textbf{A}
  \item \textbf{A}
  \item \textbf{A}
  \item \textbf{A}
  \item \textbf{A}
  \item \textbf{A}
  \item \textbf{A}
  \item \textbf{A}
  \item \textbf{A}
  \item \textbf{A}
  \item \textbf{A}
  \item \textbf{A}
  \item \textbf{A}
  \item \textbf{A}
  \item \textbf{A}
  \item \textbf{A}
  \item \textbf{A}
  \item \textbf{A}
  \item \textbf{A}
  \item \textbf{A}
  \item \textbf{A}
  \item \textbf{A}
  \item \textbf{A}
  \item \textbf{A}
  \item \textbf{A}
  \item \textbf{A}
  \item \textbf{A}
  \item \textbf{A}
  \item \textbf{A}
  \item \textbf{A}
  \item \textbf{A}
  \item \textbf{A}
  \item \textbf{A}
  \item \textbf{A}
  \item \textbf{A}
  \item \textbf{A}
  \item \textbf{A}
  \item \textbf{A}
  \item \textbf{A}
  \item \textbf{A}
  \item \textbf{A}
  \item \textbf{A}
  \item \textbf{A}
  \item \textbf{A}
  \item \textbf{A}
  \item \textbf{A}
  \item \textbf{A}
  \item \textbf{A}
  \item \textbf{A}
  \item \textbf{A}
  \item \textbf{A}
  \item \textbf{A}
  \item \textbf{A}
  \item \textbf{A}
  \item \textbf{A}
  \item \textbf{A}
  \item \textbf{A}
  \item \textbf{A}
  \item \textbf{A}
  \item \textbf{A}
  \item \textbf{A}
\end{enumerate}
\end{multicols}
\end{tcolorbox}
\newpage\n
\subsection{3.5 Literatura Salvadoreña}

\subsubsection*{SECCIÓN I: CONTEXTO HISTÓRICO Y LITERATURA PRECOLOMBINA (1–40)}

\begin{enumerate}[label=\textbf{\arabic*.}]
\setcounter{enumi}{0}
\item ¿Qué civilizaciones precolombinas habitaron el actual territorio salvadoreño?
\begin{enumerate}[label=\textbf{\Alph*)}]
\begin{multicols}{2}
  \item Mayas, lencas, pipiles y chortíes
  \item Francisco Gavidia
  \item Juan Ramón Molina
  \item Alberto Masferrer
  \item Arturo Ambrogi
  \item Claudia Lars
\end{multicols}
\end{enumerate}
\vspace{0.15cm}
\item ¿Cuál fue la civilización más desarrollada en el actual El Salvador?
\begin{enumerate}[label=\textbf{\Alph*)}]
\begin{multicols}{2}
  \item Los mayas
  \item Francisco Gavidia
  \item Juan Ramón Molina
  \item Alberto Masferrer
  \item Arturo Ambrogi
  \item Claudia Lars
\end{multicols}
\end{enumerate}
\vspace{0.15cm}
\item ¿Qué sitio arqueológico salvadoreño es Patrimonio de la Humanidad?
\begin{enumerate}[label=\textbf{\Alph*)}]
\begin{multicols}{2}
  \item Joya de Cerén
  \item Francisco Gavidia
  \item Juan Ramón Molina
  \item Alberto Masferrer
  \item Arturo Ambrogi
  \item Claudia Lars
\end{multicols}
\end{enumerate}
\vspace{0.15cm}
\item ¿Qué es Joya de Cerén?
\begin{enumerate}[label=\textbf{\Alph*)}]
\begin{multicols}{2}
  \item Un sitio arqueológico maya conocido como la Pompeya de América
  \item Francisco Gavidia
  \item Juan Ramón Molina
  \item Alberto Masferrer
  \item Arturo Ambrogi
  \item Claudia Lars
\end{multicols}
\end{enumerate}
\vspace{0.15cm}
\item ¿Qué lengua precolombina se hablaba en el actual El Salvador?
\begin{enumerate}[label=\textbf{\Alph*)}]
\begin{multicols}{2}
  \item Náhuat
  \item Francisco Gavidia
  \item Juan Ramón Molina
  \item Alberto Masferrer
  \item Arturo Ambrogi
  \item Claudia Lars
\end{multicols}
\end{enumerate}
\vspace{0.15cm}
\item ¿Qué es el náhuat?
\begin{enumerate}[label=\textbf{\Alph*)}]
\begin{multicols}{2}
  \item La lengua de los pipiles en El Salvador
  \item Francisco Gavidia
  \item Juan Ramón Molina
  \item Alberto Masferrer
  \item Arturo Ambrogi
  \item Claudia Lars
\end{multicols}
\end{enumerate}
\vspace{0.15cm}
\item ¿Qué literatura precolombina se produjo en El Salvador?
\begin{enumerate}[label=\textbf{\Alph*)}]
\begin{multicols}{2}
  \item Poesía náhuat, mitos y leyendas orales
  \item Francisco Gavidia
  \item Juan Ramón Molina
  \item Alberto Masferrer
  \item Arturo Ambrogi
  \item Claudia Lars
\end{multicols}
\end{enumerate}
\vspace{0.15cm}
\item ¿Qué son los "cantares" en la tradición náhuat?
\begin{enumerate}[label=\textbf{\Alph*)}]
\begin{multicols}{2}
  \item Poemas y cantos rituales
  \item Francisco Gavidia
  \item Juan Ramón Molina
  \item Alberto Masferrer
  \item Arturo Ambrogi
  \item Claudia Lars
\end{multicols}
\end{enumerate}
\vspace{0.15cm}
\item ¿Qué es el "tacón" en la literatura náhuat?
\begin{enumerate}[label=\textbf{\Alph*)}]
\begin{multicols}{2}
  \item Un canto o poema
  \item Francisco Gavidia
  \item Juan Ramón Molina
  \item Alberto Masferrer
  \item Arturo Ambrogi
  \item Claudia Lars
\end{multicols}
\end{enumerate}
\vspace{0.15cm}
\item ¿Qué es el "huipil" en la cultura náhuat?
\begin{enumerate}[label=\textbf{\Alph*)}]
\begin{multicols}{2}
  \item Una prenda de vestir tradicional
  \item Francisco Gavidia
  \item Juan Ramón Molina
  \item Alberto Masferrer
  \item Arturo Ambrogi
  \item Claudia Lars
\end{multicols}
\end{enumerate}
\vspace{0.15cm}
\item ¿Qué mito precolombino salvadoreño es conocido?
\begin{enumerate}[label=\textbf{\Alph*)}]
\begin{multicols}{2}
  \item El mito de los gemelos mayas
  \item Francisco Gavidia
  \item Juan Ramón Molina
  \item Alberto Masferrer
  \item Arturo Ambrogi
  \item Claudia Lars
\end{multicols}
\end{enumerate}
\vspace{0.15cm}
\item ¿Qué es el Popol Vuh?
\begin{enumerate}[label=\textbf{\Alph*)}]
\begin{multicols}{2}
  \item Un libro sagrado maya con mitos de creación
  \item Francisco Gavidia
  \item Juan Ramón Molina
  \item Alberto Masferrer
  \item Arturo Ambrogi
  \item Claudia Lars
\end{multicols}
\end{enumerate}
\vspace{0.15cm}
\item ¿Qué relación tiene el Popol Vuh con El Salvador?
\begin{enumerate}[label=\textbf{\Alph*)}]
\begin{multicols}{2}
  \item Los mayas que habitaron El Salvador compartían esta tradición
  \item Francisco Gavidia
  \item Juan Ramón Molina
  \item Alberto Masferrer
  \item Arturo Ambrogi
  \item Claudia Lars
\end{multicols}
\end{enumerate}
\vspace{0.15cm}
\item ¿Qué es el Chilam Balam?
\begin{enumerate}[label=\textbf{\Alph*)}]
\begin{multicols}{2}
  \item Un conjunto de libros mayas con profecías
  \item Francisco Gavidia
  \item Juan Ramón Molina
  \item Alberto Masferrer
  \item Arturo Ambrogi
  \item Claudia Lars
\end{multicols}
\end{enumerate}
\vspace{0.15cm}
\item ¿Qué relación tiene el Chilam Balam con El Salvador?
\begin{enumerate}[label=\textbf{\Alph*)}]
\begin{multicols}{2}
  \item Los mayas que habitaron El Salvador compartían esta tradición
  \item Francisco Gavidia
  \item Juan Ramón Molina
  \item Alberto Masferrer
  \item Arturo Ambrogi
  \item Claudia Lars
\end{multicols}
\end{enumerate}
\vspace{0.15cm}
\item ¿Qué es la literatura precolombina?
\begin{enumerate}[label=\textbf{\Alph*)}]
\begin{multicols}{2}
  \item La producida antes de la llegada de los españoles
  \item Francisco Gavidia
  \item Juan Ramón Molina
  \item Alberto Masferrer
  \item Arturo Ambrogi
  \item Claudia Lars
\end{multicols}
\end{enumerate}
\vspace{0.15cm}
\item ¿Cuál es una característica de la literatura precolombina salvadoreña?
\begin{enumerate}[label=\textbf{\Alph*)}]
\begin{multicols}{2}
  \item Ser de tradición oral y anónima
  \item Francisco Gavidia
  \item Juan Ramón Molina
  \item Alberto Masferrer
  \item Arturo Ambrogi
  \item Claudia Lars
\end{multicols}
\end{enumerate}
\vspace{0.15cm}
\item ¿Qué es el "calpulli" en la organización social náhuat?
\begin{enumerate}[label=\textbf{\Alph*)}]
\begin{multicols}{2}
  \item Un clan o comunidad
  \item Francisco Gavidia
  \item Juan Ramón Molina
  \item Alberto Masferrer
  \item Arturo Ambrogi
  \item Claudia Lars
\end{multicols}
\end{enumerate}
\vspace{0.15cm}
\item ¿Qué es el "tianguis" en la cultura náhuat?
\begin{enumerate}[label=\textbf{\Alph*)}]
\begin{multicols}{2}
  \item Un mercado
  \item Francisco Gavidia
  \item Juan Ramón Molina
  \item Alberto Masferrer
  \item Arturo Ambrogi
  \item Claudia Lars
\end{multicols}
\end{enumerate}
\vspace{0.15cm}
\item ¿Qué es el "teocalli" en la cultura náhuat?
\begin{enumerate}[label=\textbf{\Alph*)}]
\begin{multicols}{2}
  \item Un templo
  \item Francisco Gavidia
  \item Juan Ramón Molina
  \item Alberto Masferrer
  \item Arturo Ambrogi
  \item Claudia Lars
\end{multicols}
\end{enumerate}
\vspace{0.15cm}
\item ¿Qué es el "tlahtoani" en la cultura náhuat?
\begin{enumerate}[label=\textbf{\Alph*)}]
\begin{multicols}{2}
  \item Un gobernante o rey
  \item Francisco Gavidia
  \item Juan Ramón Molina
  \item Alberto Masferrer
  \item Arturo Ambrogi
  \item Claudia Lars
\end{multicols}
\end{enumerate}
\vspace{0.15cm}
\item ¿Qué es el "macehual" en la cultura náhuat?
\begin{enumerate}[label=\textbf{\Alph*)}]
\begin{multicols}{2}
  \item Un plebeyo o campesino
  \item Francisco Gavidia
  \item Juan Ramón Molina
  \item Alberto Masferrer
  \item Arturo Ambrogi
  \item Claudia Lars
\end{multicols}
\end{enumerate}
\vspace{0.15cm}
\item ¿Qué es el "pipil"?
\begin{enumerate}[label=\textbf{\Alph*)}]
\begin{multicols}{2}
  \item El pueblo náhuat que habitó El Salvador
  \item Francisco Gavidia
  \item Juan Ramón Molina
  \item Alberto Masferrer
  \item Arturo Ambrogi
  \item Claudia Lars
\end{multicols}
\end{enumerate}
\vspace{0.15cm}
\item ¿Qué es el "lenca"?
\begin{enumerate}[label=\textbf{\Alph*)}]
\begin{multicols}{2}
  \item Un pueblo indígena que habitó El Salvador
  \item Francisco Gavidia
  \item Juan Ramón Molina
  \item Alberto Masferrer
  \item Arturo Ambrogi
  \item Claudia Lars
\end{multicols}
\end{enumerate}
\vspace{0.15cm}
\item ¿Qué es el "chortí"?
\begin{enumerate}[label=\textbf{\Alph*)}]
\begin{multicols}{2}
  \item Un pueblo maya que habitó El Salvador
  \item Francisco Gavidia
  \item Juan Ramón Molina
  \item Alberto Masferrer
  \item Arturo Ambrogi
  \item Claudia Lars
\end{multicols}
\end{enumerate}
\vspace{0.15cm}
\item ¿Qué es el "xinca"?
\begin{enumerate}[label=\textbf{\Alph*)}]
\begin{multicols}{2}
  \item Un pueblo indígena que habitó El Salvador
  \item Francisco Gavidia
  \item Juan Ramón Molina
  \item Alberto Masferrer
  \item Arturo Ambrogi
  \item Claudia Lars
\end{multicols}
\end{enumerate}
\vspace{0.15cm}
\item ¿Qué es el "cacao" en la cultura precolombina salvadoreña?
\begin{enumerate}[label=\textbf{\Alph*)}]
\begin{multicols}{2}
  \item Un producto importante para la economía y rituales
  \item Francisco Gavidia
  \item Juan Ramón Molina
  \item Alberto Masferrer
  \item Arturo Ambrogi
  \item Claudia Lars
\end{multicols}
\end{enumerate}
\vspace{0.15cm}
\item ¿Qué es el "maíz" en la cultura precolombina salvadoreña?
\begin{enumerate}[label=\textbf{\Alph*)}]
\begin{multicols}{2}
  \item Un producto fundamental para la alimentación y rituales
  \item Francisco Gavidia
  \item Juan Ramón Molina
  \item Alberto Masferrer
  \item Arturo Ambrogi
  \item Claudia Lars
\end{multicols}
\end{enumerate}
\vspace{0.15cm}
\item ¿Qué es el "jade" en la cultura precolombina salvadoreña?
\begin{enumerate}[label=\textbf{\Alph*)}]
\begin{multicols}{2}
  \item Una piedra preciosa usada en rituales
  \item Francisco Gavidia
  \item Juan Ramón Molina
  \item Alberto Masferrer
  \item Arturo Ambrogi
  \item Claudia Lars
\end{multicols}
\end{enumerate}
\vspace{0.15cm}
\item ¿Qué es el "obsidiana" en la cultura precolombina salvadoreña?
\begin{enumerate}[label=\textbf{\Alph*)}]
\begin{multicols}{2}
  \item Un mineral usado para herramientas y rituales
  \item Francisco Gavidia
  \item Juan Ramón Molina
  \item Alberto Masferrer
  \item Arturo Ambrogi
  \item Claudia Lars
\end{multicols}
\end{enumerate}
\vspace{0.15cm}
\item ¿Qué es el "cerámica" en la cultura precolombina salvadoreña?
\begin{enumerate}[label=\textbf{\Alph*)}]
\begin{multicols}{2}
  \item Un arte que refleja la vida cotidiana y ritual
  \item Francisco Gavidia
  \item Juan Ramón Molina
  \item Alberto Masferrer
  \item Arturo Ambrogi
  \item Claudia Lars
\end{multicols}
\end{enumerate}
\vspace{0.15cm}
\item ¿Qué es el "petrograbado" en la cultura precolombina salvadoreña?
\begin{enumerate}[label=\textbf{\Alph*)}]
\begin{multicols}{2}
  \item Un arte rupestre con significados rituales
  \item Francisco Gavidia
  \item Juan Ramón Molina
  \item Alberto Masferrer
  \item Arturo Ambrogi
  \item Claudia Lars
\end{multicols}
\end{enumerate}
\vspace{0.15cm}
\item ¿Qué es el "metate" en la cultura precolombina salvadoreña?
\begin{enumerate}[label=\textbf{\Alph*)}]
\begin{multicols}{2}
  \item Una piedra para moler maíz
  \item Francisco Gavidia
  \item Juan Ramón Molina
  \item Alberto Masferrer
  \item Arturo Ambrogi
  \item Claudia Lars
\end{multicols}
\end{enumerate}
\vspace{0.15cm}
\item ¿Qué es el "comal" en la cultura precolombina salvadoreña?
\begin{enumerate}[label=\textbf{\Alph*)}]
\begin{multicols}{2}
  \item Un utensilio para cocinar tortillas
  \item Francisco Gavidia
  \item Juan Ramón Molina
  \item Alberto Masferrer
  \item Arturo Ambrogi
  \item Claudia Lars
\end{multicols}
\end{enumerate}
\vspace{0.15cm}
\item ¿Qué es el "molcajete" en la cultura precolombina salvadoreña?
\begin{enumerate}[label=\textbf{\Alph*)}]
\begin{multicols}{2}
  \item Un utensilio para moler
  \item Francisco Gavidia
  \item Juan Ramón Molina
  \item Alberto Masferrer
  \item Arturo Ambrogi
  \item Claudia Lars
\end{multicols}
\end{enumerate}
\vspace{0.15cm}
\item ¿Qué es el "huacal" en la cultura precolombina salvadoreña?
\begin{enumerate}[label=\textbf{\Alph*)}]
\begin{multicols}{2}
  \item Un recipiente
  \item Francisco Gavidia
  \item Juan Ramón Molina
  \item Alberto Masferrer
  \item Arturo Ambrogi
  \item Claudia Lars
\end{multicols}
\end{enumerate}
\vspace{0.15cm}
\item ¿Qué es el "tamal" en la cultura precolombina salvadoreña?
\begin{enumerate}[label=\textbf{\Alph*)}]
\begin{multicols}{2}
  \item Un alimento tradicional
  \item Francisco Gavidia
  \item Juan Ramón Molina
  \item Alberto Masferrer
  \item Arturo Ambrogi
  \item Claudia Lars
\end{multicols}
\end{enumerate}
\vspace{0.15cm}
\item ¿Qué es el "pupusa" en la cultura salvadoreña?
\begin{enumerate}[label=\textbf{\Alph*)}]
\begin{multicols}{2}
  \item Un alimento tradicional de origen precolombino
  \item Francisco Gavidia
  \item Juan Ramón Molina
  \item Alberto Masferrer
  \item Arturo Ambrogi
  \item Claudia Lars
\end{multicols}
\end{enumerate}
\vspace{0.15cm}
\item ¿Qué es el "atol" en la cultura precolombina salvadoreña?
\begin{enumerate}[label=\textbf{\Alph*)}]
\begin{multicols}{2}
  \item Una bebida tradicional
  \item Francisco Gavidia
  \item Juan Ramón Molina
  \item Alberto Masferrer
  \item Arturo Ambrogi
  \item Claudia Lars
\end{multicols}
\end{enumerate}
\vspace{0.15cm}
\item ¿Qué es el "chicha" en la cultura precolombina salvadoreña?
\begin{enumerate}[label=\textbf{\Alph*)}]
\begin{multicols}{2}
  \item Una bebida fermentada
  \item Francisco Gavidia
  \item Juan Ramón Molina
  \item Alberto Masferrer
  \item Arturo Ambrogi
  \item Claudia Lars
\end{multicols}
\end{enumerate}
\vspace{0.15cm}
\end{enumerate}\n
\subsubsection*{SECCIÓN II: LITERATURA COLONIAL SALVADOREÑA (41–80)}

\begin{enumerate}[label=\textbf{\arabic*.}]
\setcounter{enumi}{40}
\item ¿Qué es la literatura colonial salvadoreña?
\begin{enumerate}[label=\textbf{\Alph*)}]
\begin{multicols}{2}
  \item La producida durante el período colonial en El Salvador
  \item Francisco Gavidia
  \item Juan Ramón Molina
  \item Alberto Masferrer
  \item Arturo Ambrogi
  \item Claudia Lars
\end{multicols}
\end{enumerate}
\vspace{0.15cm}
\item ¿Quiénes fueron los principales cronistas coloniales en El Salvador?
\begin{enumerate}[label=\textbf{\Alph*)}]
\begin{multicols}{2}
  \item Pedro de Alvarado, Gonzalo de Alvarado y Diego de Holguín
  \item Francisco Gavidia
  \item Juan Ramón Molina
  \item Alberto Masferrer
  \item Arturo Ambrogi
  \item Claudia Lars
\end{multicols}
\end{enumerate}
\vspace{0.15cm}
\item ¿Quién fue Pedro de Alvarado?
\begin{enumerate}[label=\textbf{\Alph*)}]
\begin{multicols}{2}
  \item El conquistador de El Salvador
  \item Francisco Gavidia
  \item Juan Ramón Molina
  \item Alberto Masferrer
  \item Arturo Ambrogi
  \item Claudia Lars
\end{multicols}
\end{enumerate}
\vspace{0.15cm}
\item ¿Qué escribió Pedro de Alvarado?
\begin{enumerate}[label=\textbf{\Alph*)}]
\begin{multicols}{2}
  \item Cartas de relación sobre la conquista
  \item Francisco Gavidia
  \item Juan Ramón Molina
  \item Alberto Masferrer
  \item Arturo Ambrogi
  \item Claudia Lars
\end{multicols}
\end{enumerate}
\vspace{0.15cm}
\item ¿Quién fue Diego de Holguín?
\begin{enumerate}[label=\textbf{\Alph*)}]
\begin{multicols}{2}
  \item Un conquistador y cronista de El Salvador
  \item Francisco Gavidia
  \item Juan Ramón Molina
  \item Alberto Masferrer
  \item Arturo Ambrogi
  \item Claudia Lars
\end{multicols}
\end{enumerate}
\vspace{0.15cm}
\item ¿Qué es la "Relación de la provincia de Guatemala"?
\begin{enumerate}[label=\textbf{\Alph*)}]
\begin{multicols}{2}
  \item Una crónica colonial sobre la región
  \item Francisco Gavidia
  \item Juan Ramón Molina
  \item Alberto Masferrer
  \item Arturo Ambrogi
  \item Claudia Lars
\end{multicols}
\end{enumerate}
\vspace{0.15cm}
\item ¿Quién fue Fray Bartolomé de las Casas?
\begin{enumerate}[label=\textbf{\Alph*)}]
\begin{multicols}{2}
  \item Un fraile defensor de los indígenas
  \item Francisco Gavidia
  \item Juan Ramón Molina
  \item Alberto Masferrer
  \item Arturo Ambrogi
  \item Claudia Lars
\end{multicols}
\end{enumerate}
\vspace{0.15cm}
\item ¿Qué escribió Fray Bartolomé de las Casas?
\begin{enumerate}[label=\textbf{\Alph*)}]
\begin{multicols}{2}
  \item Brevísima relación de la destrucción de las Indias
  \item Francisco Gavidia
  \item Juan Ramón Molina
  \item Alberto Masferrer
  \item Arturo Ambrogi
  \item Claudia Lars
\end{multicols}
\end{enumerate}
\vspace{0.15cm}
\item ¿Qué es la "Historia de la provincia de San Vicente de Chiapa y Guatemala"?
\begin{enumerate}[label=\textbf{\Alph*)}]
\begin{multicols}{2}
  \item Una crónica colonial
  \item Francisco Gavidia
  \item Juan Ramón Molina
  \item Alberto Masferrer
  \item Arturo Ambrogi
  \item Claudia Lars
\end{multicols}
\end{enumerate}
\vspace{0.15cm}
\item ¿Quién fue Fray Antonio de Remesal?
\begin{enumerate}[label=\textbf{\Alph*)}]
\begin{multicols}{2}
  \item Un cronista colonial
  \item Francisco Gavidia
  \item Juan Ramón Molina
  \item Alberto Masferrer
  \item Arturo Ambrogi
  \item Claudia Lars
\end{multicols}
\end{enumerate}
\vspace{0.15cm}
\item ¿Qué escribió Fray Antonio de Remesal?
\begin{enumerate}[label=\textbf{\Alph*)}]
\begin{multicols}{2}
  \item Historia de la provincia de San Vicente de Chiapa y Guatemala
  \item Francisco Gavidia
  \item Juan Ramón Molina
  \item Alberto Masferrer
  \item Arturo Ambrogi
  \item Claudia Lars
\end{multicols}
\end{enumerate}
\vspace{0.15cm}
\item ¿Qué es la "Crónica de la provincia de Guatemala"?
\begin{enumerate}[label=\textbf{\Alph*)}]
\begin{multicols}{2}
  \item Una crónica colonial
  \item Francisco Gavidia
  \item Juan Ramón Molina
  \item Alberto Masferrer
  \item Arturo Ambrogi
  \item Claudia Lars
\end{multicols}
\end{enumerate}
\vspace{0.15cm}
\item ¿Quién fue Fray Francisco Ximénez?
\begin{enumerate}[label=\textbf{\Alph*)}]
\begin{multicols}{2}
  \item Un cronista colonial
  \item Francisco Gavidia
  \item Juan Ramón Molina
  \item Alberto Masferrer
  \item Arturo Ambrogi
  \item Claudia Lars
\end{multicols}
\end{enumerate}
\vspace{0.15cm}
\item ¿Qué escribió Fray Francisco Ximénez?
\begin{enumerate}[label=\textbf{\Alph*)}]
\begin{multicols}{2}
  \item Historia de la provincia de Guatemala
  \item Francisco Gavidia
  \item Juan Ramón Molina
  \item Alberto Masferrer
  \item Arturo Ambrogi
  \item Claudia Lars
\end{multicols}
\end{enumerate}
\vspace{0.15cm}
\item ¿Qué es el "Popol Vuh" en relación con El Salvador?
\begin{enumerate}[label=\textbf{\Alph*)}]
\begin{multicols}{2}
  \item Un texto maya compartido por los pueblos de la región
  \item Francisco Gavidia
  \item Juan Ramón Molina
  \item Alberto Masferrer
  \item Arturo Ambrogi
  \item Claudia Lars
\end{multicols}
\end{enumerate}
\vspace{0.15cm}
\item ¿Qué es la "literatura de viajes" en la colonia?
\begin{enumerate}[label=\textbf{\Alph*)}]
\begin{multicols}{2}
  \item Relatos de exploradores y cronistas
  \item Francisco Gavidia
  \item Juan Ramón Molina
  \item Alberto Masferrer
  \item Arturo Ambrogi
  \item Claudia Lars
\end{multicols}
\end{enumerate}
\vspace{0.15cm}
\item ¿Quién fue Fray Pedro de Betancur?
\begin{enumerate}[label=\textbf{\Alph*)}]
\begin{multicols}{2}
  \item Un religioso y escritor colonial
  \item Francisco Gavidia
  \item Juan Ramón Molina
  \item Alberto Masferrer
  \item Arturo Ambrogi
  \item Claudia Lars
\end{multicols}
\end{enumerate}
\vspace{0.15cm}
\item ¿Qué es la "poesía colonial salvadoreña"?
\begin{enumerate}[label=\textbf{\Alph*)}]
\begin{multicols}{2}
  \item La producida durante la colonia, principalmente religiosa
  \item Francisco Gavidia
  \item Juan Ramón Molina
  \item Alberto Masferrer
  \item Arturo Ambrogi
  \item Claudia Lars
\end{multicols}
\end{enumerate}
\vspace{0.15cm}
\item ¿Qué es el "teatro colonial salvadoreño"?
\begin{enumerate}[label=\textbf{\Alph*)}]
\begin{multicols}{2}
  \item El teatro religioso y didáctico de la colonia
  \item Francisco Gavidia
  \item Juan Ramón Molina
  \item Alberto Masferrer
  \item Arturo Ambrogi
  \item Claudia Lars
\end{multicols}
\end{enumerate}
\vspace{0.15cm}
\item ¿Qué es el "auto sacramental" en la colonia?
\begin{enumerate}[label=\textbf{\Alph*)}]
\begin{multicols}{2}
  \item Una obra teatral de tema religioso
  \item Francisco Gavidia
  \item Juan Ramón Molina
  \item Alberto Masferrer
  \item Arturo Ambrogi
  \item Claudia Lars
\end{multicols}
\end{enumerate}
\vspace{0.15cm}
\item ¿Qué es la "loa" en el teatro colonial?
\begin{enumerate}[label=\textbf{\Alph*)}]
\begin{multicols}{2}
  \item Una obra breve de alabanza
  \item Francisco Gavidia
  \item Juan Ramón Molina
  \item Alberto Masferrer
  \item Arturo Ambrogi
  \item Claudia Lars
\end{multicols}
\end{enumerate}
\vspace{0.15cm}
\item ¿Qué es el "entremés" en el teatro colonial?
\begin{enumerate}[label=\textbf{\Alph*)}]
\begin{multicols}{2}
  \item Una obra breve cómica
  \item Francisco Gavidia
  \item Juan Ramón Molina
  \item Alberto Masferrer
  \item Arturo Ambrogi
  \item Claudia Lars
\end{multicols}
\end{enumerate}
\vspace{0.15cm}
\item ¿Qué es la "comedia" en el teatro colonial?
\begin{enumerate}[label=\textbf{\Alph*)}]
\begin{multicols}{2}
  \item Una obra teatral de tema profano
  \item Francisco Gavidia
  \item Juan Ramón Molina
  \item Alberto Masferrer
  \item Arturo Ambrogi
  \item Claudia Lars
\end{multicols}
\end{enumerate}
\vspace{0.15cm}
\item ¿Qué es la "tragedia" en el teatro colonial?
\begin{enumerate}[label=\textbf{\Alph*)}]
\begin{multicols}{2}
  \item Una obra teatral de tema serio
  \item Francisco Gavidia
  \item Juan Ramón Molina
  \item Alberto Masferrer
  \item Arturo Ambrogi
  \item Claudia Lars
\end{multicols}
\end{enumerate}
\vspace{0.15cm}
\item ¿Qué es el "coloquio" en el teatro colonial?
\begin{enumerate}[label=\textbf{\Alph*)}]
\begin{multicols}{2}
  \item Una obra teatral dialogada
  \item Francisco Gavidia
  \item Juan Ramón Molina
  \item Alberto Masferrer
  \item Arturo Ambrogi
  \item Claudia Lars
\end{multicols}
\end{enumerate}
\vspace{0.15cm}
\item ¿Qué es la "pastorela" en el teatro colonial?
\begin{enumerate}[label=\textbf{\Alph*)}]
\begin{multicols}{2}
  \item Una obra teatral de tema navideño
  \item Francisco Gavidia
  \item Juan Ramón Molina
  \item Alberto Masferrer
  \item Arturo Ambrogi
  \item Claudia Lars
\end{multicols}
\end{enumerate}
\vspace{0.15cm}
\item ¿Qué es la "posada" en la cultura colonial?
\begin{enumerate}[label=\textbf{\Alph*)}]
\begin{multicols}{2}
  \item Una celebración religiosa
  \item Francisco Gavidia
  \item Juan Ramón Molina
  \item Alberto Masferrer
  \item Arturo Ambrogi
  \item Claudia Lars
\end{multicols}
\end{enumerate}
\vspace{0.15cm}
\item ¿Qué es la "cofradía" en la colonia?
\begin{enumerate}[label=\textbf{\Alph*)}]
\begin{multicols}{2}
  \item Una organización religiosa
  \item Francisco Gavidia
  \item Juan Ramón Molina
  \item Alberto Masferrer
  \item Arturo Ambrogi
  \item Claudia Lars
\end{multicols}
\end{enumerate}
\vspace{0.15cm}
\item ¿Qué es el "cabildo" en la colonia?
\begin{enumerate}[label=\textbf{\Alph*)}]
\begin{multicols}{2}
  \item Una institución de gobierno municipal
  \item Francisco Gavidia
  \item Juan Ramón Molina
  \item Alberto Masferrer
  \item Arturo Ambrogi
  \item Claudia Lars
\end{multicols}
\end{enumerate}
\vspace{0.15cm}
\item ¿Qué es la "audiencia" en la colonia?
\begin{enumerate}[label=\textbf{\Alph*)}]
\begin{multicols}{2}
  \item Una institución de gobierno colonial
  \item Francisco Gavidia
  \item Juan Ramón Molina
  \item Alberto Masferrer
  \item Arturo Ambrogi
  \item Claudia Lars
\end{multicols}
\end{enumerate}
\vspace{0.15cm}
\item ¿Qué es el "virreinato" en la colonia?
\begin{enumerate}[label=\textbf{\Alph*)}]
\begin{multicols}{2}
  \item Una división administrativa colonial
  \item Francisco Gavidia
  \item Juan Ramón Molina
  \item Alberto Masferrer
  \item Arturo Ambrogi
  \item Claudia Lars
\end{multicols}
\end{enumerate}
\vspace{0.15cm}
\item ¿Qué es la "capitanía general" en la colonia?
\begin{enumerate}[label=\textbf{\Alph*)}]
\begin{multicols}{2}
  \item Una división administrativa colonial
  \item Francisco Gavidia
  \item Juan Ramón Molina
  \item Alberto Masferrer
  \item Arturo Ambrogi
  \item Claudia Lars
\end{multicols}
\end{enumerate}
\vspace{0.15cm}
\item ¿Qué es la "encomienda" en la colonia?
\begin{enumerate}[label=\textbf{\Alph*)}]
\begin{multicols}{2}
  \item Una institución de explotación indígena
  \item Francisco Gavidia
  \item Juan Ramón Molina
  \item Alberto Masferrer
  \item Arturo Ambrogi
  \item Claudia Lars
\end{multicols}
\end{enumerate}
\vspace{0.15cm}
\item ¿Qué es la "mita" en la colonia?
\begin{enumerate}[label=\textbf{\Alph*)}]
\begin{multicols}{2}
  \item Un sistema de trabajo forzado
  \item Francisco Gavidia
  \item Juan Ramón Molina
  \item Alberto Masferrer
  \item Arturo Ambrogi
  \item Claudia Lars
\end{multicols}
\end{enumerate}
\vspace{0.15cm}
\item ¿Qué es el "obraje" en la colonia?
\begin{enumerate}[label=\textbf{\Alph*)}]
\begin{multicols}{2}
  \item Un taller textil colonial
  \item Francisco Gavidia
  \item Juan Ramón Molina
  \item Alberto Masferrer
  \item Arturo Ambrogi
  \item Claudia Lars
\end{multicols}
\end{enumerate}
\vspace{0.15cm}
\item ¿Qué es la "hacienda" en la colonia?
\begin{enumerate}[label=\textbf{\Alph*)}]
\begin{multicols}{2}
  \item Una gran propiedad agrícola
  \item Francisco Gavidia
  \item Juan Ramón Molina
  \item Alberto Masferrer
  \item Arturo Ambrogi
  \item Claudia Lars
\end{multicols}
\end{enumerate}
\vspace{0.15cm}
\item ¿Qué es el "trapiche" en la colonia?
\begin{enumerate}[label=\textbf{\Alph*)}]
\begin{multicols}{2}
  \item Un molino para caña de azúcar
  \item Francisco Gavidia
  \item Juan Ramón Molina
  \item Alberto Masferrer
  \item Arturo Ambrogi
  \item Claudia Lars
\end{multicols}
\end{enumerate}
\vspace{0.15cm}
\item ¿Qué es el "añil" en la colonia salvadoreña?
\begin{enumerate}[label=\textbf{\Alph*)}]
\begin{multicols}{2}
  \item Un cultivo de tinte importante
  \item Francisco Gavidia
  \item Juan Ramón Molina
  \item Alberto Masferrer
  \item Arturo Ambrogi
  \item Claudia Lars
\end{multicols}
\end{enumerate}
\vspace{0.15cm}
\item ¿Qué es el "cacao" en la colonia salvadoreña?
\begin{enumerate}[label=\textbf{\Alph*)}]
\begin{multicols}{2}
  \item Un cultivo importante para la economía
  \item Francisco Gavidia
  \item Juan Ramón Molina
  \item Alberto Masferrer
  \item Arturo Ambrogi
  \item Claudia Lars
\end{multicols}
\end{enumerate}
\vspace{0.15cm}
\item ¿Qué es el "bálsamo" en la colonia salvadoreña?
\begin{enumerate}[label=\textbf{\Alph*)}]
\begin{multicols}{2}
  \item Un producto medicinal y aromático
  \item Francisco Gavidia
  \item Juan Ramón Molina
  \item Alberto Masferrer
  \item Arturo Ambrogi
  \item Claudia Lars
\end{multicols}
\end{enumerate}
\vspace{0.15cm}
\end{enumerate}\n
\subsubsection*{SECCIÓN III: ROMANTICISMO Y COSTUMBRISMO SALVADOREÑO (81–120)}

\begin{enumerate}[label=\textbf{\arabic*.}]
\setcounter{enumi}{80}
\item ¿Qué es el Romanticismo salvadoreño?
\begin{enumerate}[label=\textbf{\Alph*)}]
\begin{multicols}{2}
  \item La adaptación del Romanticismo europeo a la realidad salvadoreña
  \item Francisco Gavidia
  \item Juan Ramón Molina
  \item Alberto Masferrer
  \item Arturo Ambrogi
  \item Claudia Lars
\end{multicols}
\end{enumerate}
\vspace{0.15cm}
\item ¿Cuál es una característica del Romanticismo salvadoreño?
\begin{enumerate}[label=\textbf{\Alph*)}]
\begin{multicols}{2}
  \item La exaltación de la naturaleza y la identidad nacional
  \item Francisco Gavidia
  \item Juan Ramón Molina
  \item Alberto Masferrer
  \item Arturo Ambrogi
  \item Claudia Lars
\end{multicols}
\end{enumerate}
\vspace{0.15cm}
\item ¿Quién es el principal representante del Romanticismo salvadoreño?
\begin{enumerate}[label=\textbf{\Alph*)}]
\begin{multicols}{2}
  \item Juan Ramón Molina
  \item Francisco Gavidia
  \item Alberto Masferrer
  \item Arturo Ambrogi
  \item Claudia Lars
  \item Roque Dalton
\end{multicols}
\end{enumerate}
\vspace{0.15cm}
\item ¿Quién escribió "Tierra de infancia"?
\begin{enumerate}[label=\textbf{\Alph*)}]
\begin{multicols}{2}
  \item Claudia Lars
  \item Francisco Gavidia
  \item Juan Ramón Molina
  \item Alberto Masferrer
  \item Arturo Ambrogi
  \item Roque Dalton
\end{multicols}
\end{enumerate}
\vspace{0.15cm}
\item ¿Quién escribió "El libro de la vida"?
\begin{enumerate}[label=\textbf{\Alph*)}]
\begin{multicols}{2}
  \item Alberto Masferrer
  \item Francisco Gavidia
  \item Juan Ramón Molina
  \item Arturo Ambrogi
  \item Claudia Lars
  \item Roque Dalton
\end{multicols}
\end{enumerate}
\vspace{0.15cm}
\item ¿Quién escribió "El dinero maldito"?
\begin{enumerate}[label=\textbf{\Alph*)}]
\begin{multicols}{2}
  \item Alberto Masferrer
  \item Francisco Gavidia
  \item Juan Ramón Molina
  \item Arturo Ambrogi
  \item Claudia Lars
  \item Roque Dalton
\end{multicols}
\end{enumerate}
\vspace{0.15cm}
\item ¿Quién escribió "Ensayo sobre el dinero"?
\begin{enumerate}[label=\textbf{\Alph*)}]
\begin{multicols}{2}
  \item Alberto Masferrer
  \item Francisco Gavidia
  \item Juan Ramón Molina
  \item Arturo Ambrogi
  \item Claudia Lars
  \item Roque Dalton
\end{multicols}
\end{enumerate}
\vspace{0.15cm}
\item ¿Quién escribió "El mínimum vital"?
\begin{enumerate}[label=\textbf{\Alph*)}]
\begin{multicols}{2}
  \item Alberto Masferrer
  \item Francisco Gavidia
  \item Juan Ramón Molina
  \item Arturo Ambrogi
  \item Claudia Lars
  \item Roque Dalton
\end{multicols}
\end{enumerate}
\vspace{0.15cm}
\item ¿Quién es el autor de "Cuentos de barro"?
\begin{enumerate}[label=\textbf{\Alph*)}]
\begin{multicols}{2}
  \item Salarrué
  \item Francisco Gavidia
  \item Juan Ramón Molina
  \item Alberto Masferrer
  \item Arturo Ambrogi
  \item Claudia Lars
\end{multicols}
\end{enumerate}
\vspace{0.15cm}
\item ¿Quién es el autor de "Cuentos de cipotes"?
\begin{enumerate}[label=\textbf{\Alph*)}]
\begin{multicols}{2}
  \item Salarrué
  \item Francisco Gavidia
  \item Juan Ramón Molina
  \item Alberto Masferrer
  \item Arturo Ambrogi
  \item Claudia Lars
\end{multicols}
\end{enumerate}
\vspace{0.15cm}
\item ¿Quién es el autor de "El Cristo negro"?
\begin{enumerate}[label=\textbf{\Alph*)}]
\begin{multicols}{2}
  \item Salarrué
  \item Francisco Gavidia
  \item Juan Ramón Molina
  \item Alberto Masferrer
  \item Arturo Ambrogi
  \item Claudia Lars
\end{multicols}
\end{enumerate}
\vspace{0.15cm}
\item ¿Quién es el autor de "Trasmallo"?
\begin{enumerate}[label=\textbf{\Alph*)}]
\begin{multicols}{2}
  \item Salarrué
  \item Francisco Gavidia
  \item Juan Ramón Molina
  \item Alberto Masferrer
  \item Arturo Ambrogi
  \item Claudia Lars
\end{multicols}
\end{enumerate}
\vspace{0.15cm}
\item ¿Quién es el autor de "La espada y otras narraciones"?
\begin{enumerate}[label=\textbf{\Alph*)}]
\begin{multicols}{2}
  \item Salarrué
  \item Francisco Gavidia
  \item Juan Ramón Molina
  \item Alberto Masferrer
  \item Arturo Ambrogi
  \item Claudia Lars
\end{multicols}
\end{enumerate}
\vspace{0.15cm}
\item ¿Quién es el autor de "Sombra y otros poemas"?
\begin{enumerate}[label=\textbf{\Alph*)}]
\begin{multicols}{2}
  \item Claudia Lars
  \item Francisco Gavidia
  \item Juan Ramón Molina
  \item Alberto Masferrer
  \item Arturo Ambrogi
  \item Roque Dalton
\end{multicols}
\end{enumerate}
\vspace{0.15cm}
\item ¿Quién es el autor de "Donde comienza el aire"?
\begin{enumerate}[label=\textbf{\Alph*)}]
\begin{multicols}{2}
  \item Claudia Lars
  \item Francisco Gavidia
  \item Juan Ramón Molina
  \item Alberto Masferrer
  \item Arturo Ambrogi
  \item Roque Dalton
\end{multicols}
\end{enumerate}
\vspace{0.15cm}
\item ¿Quién es el autor de "Poesía última"?
\begin{enumerate}[label=\textbf{\Alph*)}]
\begin{multicols}{2}
  \item Claudia Lars
  \item Francisco Gavidia
  \item Juan Ramón Molina
  \item Alberto Masferrer
  \item Arturo Ambrogi
  \item Roque Dalton
\end{multicols}
\end{enumerate}
\vspace{0.15cm}
\item ¿Quién es el autor de "Canción redonda"?
\begin{enumerate}[label=\textbf{\Alph*)}]
\begin{multicols}{2}
  \item Claudia Lars
  \item Francisco Gavidia
  \item Juan Ramón Molina
  \item Alberto Masferrer
  \item Arturo Ambrogi
  \item Roque Dalton
\end{multicols}
\end{enumerate}
\vspace{0.15cm}
\item ¿Quién es el autor de "El árbol de fuego"?
\begin{enumerate}[label=\textbf{\Alph*)}]
\begin{multicols}{2}
  \item Claudia Lars
  \item Francisco Gavidia
  \item Juan Ramón Molina
  \item Alberto Masferrer
  \item Arturo Ambrogi
  \item Roque Dalton
\end{multicols}
\end{enumerate}
\vspace{0.15cm}
\item ¿Quién es el autor de "Romancero del puerto"?
\begin{enumerate}[label=\textbf{\Alph*)}]
\begin{multicols}{2}
  \item Claudia Lars
  \item Francisco Gavidia
  \item Juan Ramón Molina
  \item Alberto Masferrer
  \item Arturo Ambrogi
  \item Roque Dalton
\end{multicols}
\end{enumerate}
\vspace{0.15cm}
\item ¿Qué es el costumbrismo salvadoreño?
\begin{enumerate}[label=\textbf{\Alph*)}]
\begin{multicols}{2}
  \item La descripción de las costumbres y tipos salvadoreños
  \item Francisco Gavidia
  \item Juan Ramón Molina
  \item Alberto Masferrer
  \item Arturo Ambrogi
  \item Claudia Lars
\end{multicols}
\end{enumerate}
\vspace{0.15cm}
\item ¿Quién es el principal representante del costumbrismo salvadoreño?
\begin{enumerate}[label=\textbf{\Alph*)}]
\begin{multicols}{2}
  \item Arturo Ambrogi
  \item Francisco Gavidia
  \item Juan Ramón Molina
  \item Alberto Masferrer
  \item Claudia Lars
  \item Roque Dalton
\end{multicols}
\end{enumerate}
\vspace{0.15cm}
\item ¿Quién escribió "El libro de los espejos"?
\begin{enumerate}[label=\textbf{\Alph*)}]
\begin{multicols}{2}
  \item Arturo Ambrogi
  \item Francisco Gavidia
  \item Juan Ramón Molina
  \item Alberto Masferrer
  \item Claudia Lars
  \item Roque Dalton
\end{multicols}
\end{enumerate}
\vspace{0.15cm}
\item ¿Quién escribió "Máscaras, retablos y fantasmas"?
\begin{enumerate}[label=\textbf{\Alph*)}]
\begin{multicols}{2}
  \item Arturo Ambrogi
  \item Francisco Gavidia
  \item Juan Ramón Molina
  \item Alberto Masferrer
  \item Claudia Lars
  \item Roque Dalton
\end{multicols}
\end{enumerate}
\vspace{0.15cm}
\item ¿Quién escribió "Sensaciones de viaje"?
\begin{enumerate}[label=\textbf{\Alph*)}]
\begin{multicols}{2}
  \item Arturo Ambrogi
  \item Francisco Gavidia
  \item Juan Ramón Molina
  \item Alberto Masferrer
  \item Claudia Lars
  \item Roque Dalton
\end{multicols}
\end{enumerate}
\vspace{0.15cm}
\item ¿Qué es el "cuento salvadoreño"?
\begin{enumerate}[label=\textbf{\Alph*)}]
\begin{multicols}{2}
  \item Un género narrativo breve que refleja la realidad salvadoreña
  \item Francisco Gavidia
  \item Juan Ramón Molina
  \item Alberto Masferrer
  \item Arturo Ambrogi
  \item Claudia Lars
\end{multicols}
\end{enumerate}
\vspace{0.15cm}
\item ¿Quién es el principal representante del cuento salvadoreño?
\begin{enumerate}[label=\textbf{\Alph*)}]
\begin{multicols}{2}
  \item Salarrué
  \item Francisco Gavidia
  \item Juan Ramón Molina
  \item Alberto Masferrer
  \item Arturo Ambrogi
  \item Claudia Lars
\end{multicols}
\end{enumerate}
\vspace{0.15cm}
\item ¿Qué es "Cuentos de barro"?
\begin{enumerate}[label=\textbf{\Alph*)}]
\begin{multicols}{2}
  \item Una colección de cuentos que retratan la vida rural salvadoreña
  \item Francisco Gavidia
  \item Juan Ramón Molina
  \item Alberto Masferrer
  \item Arturo Ambrogi
  \item Claudia Lars
\end{multicols}
\end{enumerate}
\vspace{0.15cm}
\item ¿Qué es "Cuentos de cipotes"?
\begin{enumerate}[label=\textbf{\Alph*)}]
\begin{multicols}{2}
  \item Una colección de cuentos infantiles salvadoreños
  \item Francisco Gavidia
  \item Juan Ramón Molina
  \item Alberto Masferrer
  \item Arturo Ambrogi
  \item Claudia Lars
\end{multicols}
\end{enumerate}
\vspace{0.15cm}
\item ¿Qué es "El Cristo negro"?
\begin{enumerate}[label=\textbf{\Alph*)}]
\begin{multicols}{2}
  \item Un cuento que retrata la religiosidad popular salvadoreña
  \item Francisco Gavidia
  \item Juan Ramón Molina
  \item Alberto Masferrer
  \item Arturo Ambrogi
  \item Claudia Lars
\end{multicols}
\end{enumerate}
\vspace{0.15cm}
\item ¿Qué es "Trasmallo"?
\begin{enumerate}[label=\textbf{\Alph*)}]
\begin{multicols}{2}
  \item Una colección de cuentos que retrata la vida de los pescadores
  \item Francisco Gavidia
  \item Juan Ramón Molina
  \item Alberto Masferrer
  \item Arturo Ambrogi
  \item Claudia Lars
\end{multicols}
\end{enumerate}
\vspace{0.15cm}
\item ¿Qué es "La espada y otras narraciones"?
\begin{enumerate}[label=\textbf{\Alph*)}]
\begin{multicols}{2}
  \item Una colección de cuentos salvadoreños
  \item Francisco Gavidia
  \item Juan Ramón Molina
  \item Alberto Masferrer
  \item Arturo Ambrogi
  \item Claudia Lars
\end{multicols}
\end{enumerate}
\vspace{0.15cm}
\item ¿Qué es el "realismo salvadoreño"?
\begin{enumerate}[label=\textbf{\Alph*)}]
\begin{multicols}{2}
  \item La representación objetiva de la realidad salvadoreña
  \item Francisco Gavidia
  \item Juan Ramón Molina
  \item Alberto Masferrer
  \item Arturo Ambrogi
  \item Claudia Lars
\end{multicols}
\end{enumerate}
\vspace{0.15cm}
\item ¿Quién es el principal representante del realismo salvadoreño?
\begin{enumerate}[label=\textbf{\Alph*)}]
\begin{multicols}{2}
  \item Salarrué
  \item Francisco Gavidia
  \item Juan Ramón Molina
  \item Alberto Masferrer
  \item Arturo Ambrogi
  \item Claudia Lars
\end{multicols}
\end{enumerate}
\vspace{0.15cm}
\item ¿Qué es la "novela salvadoreña"?
\begin{enumerate}[label=\textbf{\Alph*)}]
\begin{multicols}{2}
  \item Un género narrativo extenso que refleja la realidad salvadoreña
  \item Francisco Gavidia
  \item Juan Ramón Molina
  \item Alberto Masferrer
  \item Arturo Ambrogi
  \item Claudia Lars
\end{multicols}
\end{enumerate}
\vspace{0.15cm}
\item ¿Quién es el principal representante de la novela salvadoreña?
\begin{enumerate}[label=\textbf{\Alph*)}]
\begin{multicols}{2}
  \item Manlio Argueta
  \item Francisco Gavidia
  \item Juan Ramón Molina
  \item Alberto Masferrer
  \item Arturo Ambrogi
  \item Claudia Lars
\end{multicols}
\end{enumerate}
\vspace{0.15cm}
\item ¿Quién escribió "Un día en la vida"?
\begin{enumerate}[label=\textbf{\Alph*)}]
\begin{multicols}{2}
  \item Manlio Argueta
  \item Francisco Gavidia
  \item Juan Ramón Molina
  \item Alberto Masferrer
  \item Arturo Ambrogi
  \item Claudia Lars
\end{multicols}
\end{enumerate}
\vspace{0.15cm}
\item ¿Quién escribió "Cuzcatlán donde bate la mar del sur"?
\begin{enumerate}[label=\textbf{\Alph*)}]
\begin{multicols}{2}
  \item Manlio Argueta
  \item Francisco Gavidia
  \item Juan Ramón Molina
  \item Alberto Masferrer
  \item Arturo Ambrogi
  \item Claudia Lars
\end{multicols}
\end{enumerate}
\vspace{0.15cm}
\item ¿Quién escribió "Milagro de la paz"?
\begin{enumerate}[label=\textbf{\Alph*)}]
\begin{multicols}{2}
  \item Manlio Argueta
  \item Francisco Gavidia
  \item Juan Ramón Molina
  \item Alberto Masferrer
  \item Arturo Ambrogi
  \item Claudia Lars
\end{multicols}
\end{enumerate}
\vspace{0.15cm}
\item ¿Qué es "Un día en la vida"?
\begin{enumerate}[label=\textbf{\Alph*)}]
\begin{multicols}{2}
  \item Una novela que retrata la vida de una campesina salvadoreña
  \item Francisco Gavidia
  \item Juan Ramón Molina
  \item Alberto Masferrer
  \item Arturo Ambrogi
  \item Claudia Lars
\end{multicols}
\end{enumerate}
\vspace{0.15cm}
\item ¿Qué es el Romanticismo salvadoreño?
\begin{enumerate}[label=\textbf{\Alph*)}]
\begin{multicols}{2}
  \item La adaptación del Romanticismo europeo a la realidad salvadoreña
  \item Francisco Gavidia
  \item Juan Ramón Molina
  \item Alberto Masferrer
  \item Arturo Ambrogi
  \item Claudia Lars
\end{multicols}
\end{enumerate}
\vspace{0.15cm}
\end{enumerate}\n
\subsubsection*{SECCIÓN IV: MODERNISMO SALVADOREÑO (121–160)}

\begin{enumerate}[label=\textbf{\arabic*.}]
\setcounter{enumi}{120}
\item ¿Qué es el Modernismo salvadoreño?
\begin{enumerate}[label=\textbf{\Alph*)}]
\begin{multicols}{2}
  \item La adaptación del Modernismo hispanoamericano a El Salvador
  \item Francisco Gavidia
  \item Juan Ramón Molina
  \item Alberto Masferrer
  \item Arturo Ambrogi
  \item Claudia Lars
\end{multicols}
\end{enumerate}
\vspace{0.15cm}
\item ¿Quién es el principal representante del Modernismo salvadoreño?
\begin{enumerate}[label=\textbf{\Alph*)}]
\begin{multicols}{2}
  \item Francisco Gavidia
  \item Juan Ramón Molina
  \item Alberto Masferrer
  \item Arturo Ambrogi
  \item Claudia Lars
  \item Roque Dalton
\end{multicols}
\end{enumerate}
\vspace{0.15cm}
\item ¿Quién es el autor de "Versos"?
\begin{enumerate}[label=\textbf{\Alph*)}]
\begin{multicols}{2}
  \item Francisco Gavidia
  \item Juan Ramón Molina
  \item Alberto Masferrer
  \item Arturo Ambrogi
  \item Claudia Lars
  \item Roque Dalton
\end{multicols}
\end{enumerate}
\vspace{0.15cm}
\item ¿Quién es el autor de "Sólo así"?
\begin{enumerate}[label=\textbf{\Alph*)}]
\begin{multicols}{2}
  \item Francisco Gavidia
  \item Juan Ramón Molina
  \item Alberto Masferrer
  \item Arturo Ambrogi
  \item Claudia Lars
  \item Roque Dalton
\end{multicols}
\end{enumerate}
\vspace{0.15cm}
\item ¿Quién es el autor de "Júbilos"?
\begin{enumerate}[label=\textbf{\Alph*)}]
\begin{multicols}{2}
  \item Francisco Gavidia
  \item Juan Ramón Molina
  \item Alberto Masferrer
  \item Arturo Ambrogi
  \item Claudia Lars
  \item Roque Dalton
\end{multicols}
\end{enumerate}
\vspace{0.15cm}
\item ¿Quién es el autor de "Tierras de sol"?
\begin{enumerate}[label=\textbf{\Alph*)}]
\begin{multicols}{2}
  \item Juan Ramón Molina
  \item Francisco Gavidia
  \item Alberto Masferrer
  \item Arturo Ambrogi
  \item Claudia Lars
  \item Roque Dalton
\end{multicols}
\end{enumerate}
\vspace{0.15cm}
\item ¿Quién es el autor de "Campo"?
\begin{enumerate}[label=\textbf{\Alph*)}]
\begin{multicols}{2}
  \item Juan Ramón Molina
  \item Francisco Gavidia
  \item Alberto Masferrer
  \item Arturo Ambrogi
  \item Claudia Lars
  \item Roque Dalton
\end{multicols}
\end{enumerate}
\vspace{0.15cm}
\item ¿Quién es el autor de "Pesca"?
\begin{enumerate}[label=\textbf{\Alph*)}]
\begin{multicols}{2}
  \item Juan Ramón Molina
  \item Francisco Gavidia
  \item Alberto Masferrer
  \item Arturo Ambrogi
  \item Claudia Lars
  \item Roque Dalton
\end{multicols}
\end{enumerate}
\vspace{0.15cm}
\item ¿Quién es el autor de "La sombra de la palmera"?
\begin{enumerate}[label=\textbf{\Alph*)}]
\begin{multicols}{2}
  \item Juan Ramón Molina
  \item Francisco Gavidia
  \item Alberto Masferrer
  \item Arturo Ambrogi
  \item Claudia Lars
  \item Roque Dalton
\end{multicols}
\end{enumerate}
\vspace{0.15cm}
\item ¿Quién es el autor de "El canto del cisne"?
\begin{enumerate}[label=\textbf{\Alph*)}]
\begin{multicols}{2}
  \item Juan Ramón Molina
  \item Francisco Gavidia
  \item Alberto Masferrer
  \item Arturo Ambrogi
  \item Claudia Lars
  \item Roque Dalton
\end{multicols}
\end{enumerate}
\vspace{0.15cm}
\item ¿Qué es el Modernismo salvadoreño?
\begin{enumerate}[label=\textbf{\Alph*)}]
\begin{multicols}{2}
  \item La adaptación del Modernismo hispanoamericano a El Salvador
  \item Francisco Gavidia
  \item Juan Ramón Molina
  \item Alberto Masferrer
  \item Arturo Ambrogi
  \item Claudia Lars
\end{multicols}
\end{enumerate}
\vspace{0.15cm}
\item ¿Quién es el principal representante del Modernismo salvadoreño?
\begin{enumerate}[label=\textbf{\Alph*)}]
\begin{multicols}{2}
  \item Francisco Gavidia
  \item Juan Ramón Molina
  \item Alberto Masferrer
  \item Arturo Ambrogi
  \item Claudia Lars
  \item Roque Dalton
\end{multicols}
\end{enumerate}
\vspace{0.15cm}
\item ¿Quién es el autor de "Versos"?
\begin{enumerate}[label=\textbf{\Alph*)}]
\begin{multicols}{2}
  \item Francisco Gavidia
  \item Juan Ramón Molina
  \item Alberto Masferrer
  \item Arturo Ambrogi
  \item Claudia Lars
  \item Roque Dalton
\end{multicols}
\end{enumerate}
\vspace{0.15cm}
\item ¿Quién es el autor de "Sólo así"?
\begin{enumerate}[label=\textbf{\Alph*)}]
\begin{multicols}{2}
  \item Francisco Gavidia
  \item Juan Ramón Molina
  \item Alberto Masferrer
  \item Arturo Ambrogi
  \item Claudia Lars
  \item Roque Dalton
\end{multicols}
\end{enumerate}
\vspace{0.15cm}
\item ¿Quién es el autor de "Júbilos"?
\begin{enumerate}[label=\textbf{\Alph*)}]
\begin{multicols}{2}
  \item Francisco Gavidia
  \item Juan Ramón Molina
  \item Alberto Masferrer
  \item Arturo Ambrogi
  \item Claudia Lars
  \item Roque Dalton
\end{multicols}
\end{enumerate}
\vspace{0.15cm}
\item ¿Quién es el autor de "Tierras de sol"?
\begin{enumerate}[label=\textbf{\Alph*)}]
\begin{multicols}{2}
  \item Juan Ramón Molina
  \item Francisco Gavidia
  \item Alberto Masferrer
  \item Arturo Ambrogi
  \item Claudia Lars
  \item Roque Dalton
\end{multicols}
\end{enumerate}
\vspace{0.15cm}
\item ¿Quién es el autor de "Campo"?
\begin{enumerate}[label=\textbf{\Alph*)}]
\begin{multicols}{2}
  \item Juan Ramón Molina
  \item Francisco Gavidia
  \item Alberto Masferrer
  \item Arturo Ambrogi
  \item Claudia Lars
  \item Roque Dalton
\end{multicols}
\end{enumerate}
\vspace{0.15cm}
\item ¿Quién es el autor de "Pesca"?
\begin{enumerate}[label=\textbf{\Alph*)}]
\begin{multicols}{2}
  \item Juan Ramón Molina
  \item Francisco Gavidia
  \item Alberto Masferrer
  \item Arturo Ambrogi
  \item Claudia Lars
  \item Roque Dalton
\end{multicols}
\end{enumerate}
\vspace{0.15cm}
\item ¿Quién es el autor de "La sombra de la palmera"?
\begin{enumerate}[label=\textbf{\Alph*)}]
\begin{multicols}{2}
  \item Juan Ramón Molina
  \item Francisco Gavidia
  \item Alberto Masferrer
  \item Arturo Ambrogi
  \item Claudia Lars
  \item Roque Dalton
\end{multicols}
\end{enumerate}
\vspace{0.15cm}
\item ¿Quién es el autor de "El canto del cisne"?
\begin{enumerate}[label=\textbf{\Alph*)}]
\begin{multicols}{2}
  \item Juan Ramón Molina
  \item Francisco Gavidia
  \item Alberto Masferrer
  \item Arturo Ambrogi
  \item Claudia Lars
  \item Roque Dalton
\end{multicols}
\end{enumerate}
\vspace{0.15cm}
\item ¿Qué es el Modernismo salvadoreño?
\begin{enumerate}[label=\textbf{\Alph*)}]
\begin{multicols}{2}
  \item La adaptación del Modernismo hispanoamericano a El Salvador
  \item Francisco Gavidia
  \item Juan Ramón Molina
  \item Alberto Masferrer
  \item Arturo Ambrogi
  \item Claudia Lars
\end{multicols}
\end{enumerate}
\vspace{0.15cm}
\item ¿Quién es el principal representante del Modernismo salvadoreño?
\begin{enumerate}[label=\textbf{\Alph*)}]
\begin{multicols}{2}
  \item Francisco Gavidia
  \item Juan Ramón Molina
  \item Alberto Masferrer
  \item Arturo Ambrogi
  \item Claudia Lars
  \item Roque Dalton
\end{multicols}
\end{enumerate}
\vspace{0.15cm}
\item ¿Quién es el autor de "Versos"?
\begin{enumerate}[label=\textbf{\Alph*)}]
\begin{multicols}{2}
  \item Francisco Gavidia
  \item Juan Ramón Molina
  \item Alberto Masferrer
  \item Arturo Ambrogi
  \item Claudia Lars
  \item Roque Dalton
\end{multicols}
\end{enumerate}
\vspace{0.15cm}
\item ¿Quién es el autor de "Sólo así"?
\begin{enumerate}[label=\textbf{\Alph*)}]
\begin{multicols}{2}
  \item Francisco Gavidia
  \item Juan Ramón Molina
  \item Alberto Masferrer
  \item Arturo Ambrogi
  \item Claudia Lars
  \item Roque Dalton
\end{multicols}
\end{enumerate}
\vspace{0.15cm}
\item ¿Quién es el autor de "Júbilos"?
\begin{enumerate}[label=\textbf{\Alph*)}]
\begin{multicols}{2}
  \item Francisco Gavidia
  \item Juan Ramón Molina
  \item Alberto Masferrer
  \item Arturo Ambrogi
  \item Claudia Lars
  \item Roque Dalton
\end{multicols}
\end{enumerate}
\vspace{0.15cm}
\item ¿Quién es el autor de "Tierras de sol"?
\begin{enumerate}[label=\textbf{\Alph*)}]
\begin{multicols}{2}
  \item Juan Ramón Molina
  \item Francisco Gavidia
  \item Alberto Masferrer
  \item Arturo Ambrogi
  \item Claudia Lars
  \item Roque Dalton
\end{multicols}
\end{enumerate}
\vspace{0.15cm}
\item ¿Quién es el autor de "Campo"?
\begin{enumerate}[label=\textbf{\Alph*)}]
\begin{multicols}{2}
  \item Juan Ramón Molina
  \item Francisco Gavidia
  \item Alberto Masferrer
  \item Arturo Ambrogi
  \item Claudia Lars
  \item Roque Dalton
\end{multicols}
\end{enumerate}
\vspace{0.15cm}
\item ¿Quién es el autor de "Pesca"?
\begin{enumerate}[label=\textbf{\Alph*)}]
\begin{multicols}{2}
  \item Juan Ramón Molina
  \item Francisco Gavidia
  \item Alberto Masferrer
  \item Arturo Ambrogi
  \item Claudia Lars
  \item Roque Dalton
\end{multicols}
\end{enumerate}
\vspace{0.15cm}
\item ¿Quién es el autor de "La sombra de la palmera"?
\begin{enumerate}[label=\textbf{\Alph*)}]
\begin{multicols}{2}
  \item Juan Ramón Molina
  \item Francisco Gavidia
  \item Alberto Masferrer
  \item Arturo Ambrogi
  \item Claudia Lars
  \item Roque Dalton
\end{multicols}
\end{enumerate}
\vspace{0.15cm}
\item ¿Quién es el autor de "El canto del cisne"?
\begin{enumerate}[label=\textbf{\Alph*)}]
\begin{multicols}{2}
  \item Juan Ramón Molina
  \item Francisco Gavidia
  \item Alberto Masferrer
  \item Arturo Ambrogi
  \item Claudia Lars
  \item Roque Dalton
\end{multicols}
\end{enumerate}
\vspace{0.15cm}
\item ¿Qué es el Modernismo salvadoreño?
\begin{enumerate}[label=\textbf{\Alph*)}]
\begin{multicols}{2}
  \item La adaptación del Modernismo hispanoamericano a El Salvador
  \item Francisco Gavidia
  \item Juan Ramón Molina
  \item Alberto Masferrer
  \item Arturo Ambrogi
  \item Claudia Lars
\end{multicols}
\end{enumerate}
\vspace{0.15cm}
\item ¿Quién es el principal representante del Modernismo salvadoreño?
\begin{enumerate}[label=\textbf{\Alph*)}]
\begin{multicols}{2}
  \item Francisco Gavidia
  \item Juan Ramón Molina
  \item Alberto Masferrer
  \item Arturo Ambrogi
  \item Claudia Lars
  \item Roque Dalton
\end{multicols}
\end{enumerate}
\vspace{0.15cm}
\item ¿Quién es el autor de "Versos"?
\begin{enumerate}[label=\textbf{\Alph*)}]
\begin{multicols}{2}
  \item Francisco Gavidia
  \item Juan Ramón Molina
  \item Alberto Masferrer
  \item Arturo Ambrogi
  \item Claudia Lars
  \item Roque Dalton
\end{multicols}
\end{enumerate}
\vspace{0.15cm}
\item ¿Quién es el autor de "Sólo así"?
\begin{enumerate}[label=\textbf{\Alph*)}]
\begin{multicols}{2}
  \item Francisco Gavidia
  \item Juan Ramón Molina
  \item Alberto Masferrer
  \item Arturo Ambrogi
  \item Claudia Lars
  \item Roque Dalton
\end{multicols}
\end{enumerate}
\vspace{0.15cm}
\item ¿Quién es el autor de "Júbilos"?
\begin{enumerate}[label=\textbf{\Alph*)}]
\begin{multicols}{2}
  \item Francisco Gavidia
  \item Juan Ramón Molina
  \item Alberto Masferrer
  \item Arturo Ambrogi
  \item Claudia Lars
  \item Roque Dalton
\end{multicols}
\end{enumerate}
\vspace{0.15cm}
\item ¿Quién es el autor de "Tierras de sol"?
\begin{enumerate}[label=\textbf{\Alph*)}]
\begin{multicols}{2}
  \item Juan Ramón Molina
  \item Francisco Gavidia
  \item Alberto Masferrer
  \item Arturo Ambrogi
  \item Claudia Lars
  \item Roque Dalton
\end{multicols}
\end{enumerate}
\vspace{0.15cm}
\item ¿Quién es el autor de "Campo"?
\begin{enumerate}[label=\textbf{\Alph*)}]
\begin{multicols}{2}
  \item Juan Ramón Molina
  \item Francisco Gavidia
  \item Alberto Masferrer
  \item Arturo Ambrogi
  \item Claudia Lars
  \item Roque Dalton
\end{multicols}
\end{enumerate}
\vspace{0.15cm}
\item ¿Quién es el autor de "Pesca"?
\begin{enumerate}[label=\textbf{\Alph*)}]
\begin{multicols}{2}
  \item Juan Ramón Molina
  \item Francisco Gavidia
  \item Alberto Masferrer
  \item Arturo Ambrogi
  \item Claudia Lars
  \item Roque Dalton
\end{multicols}
\end{enumerate}
\vspace{0.15cm}
\item ¿Quién es el autor de "La sombra de la palmera"?
\begin{enumerate}[label=\textbf{\Alph*)}]
\begin{multicols}{2}
  \item Juan Ramón Molina
  \item Francisco Gavidia
  \item Alberto Masferrer
  \item Arturo Ambrogi
  \item Claudia Lars
  \item Roque Dalton
\end{multicols}
\end{enumerate}
\vspace{0.15cm}
\item ¿Quién es el autor de "El canto del cisne"?
\begin{enumerate}[label=\textbf{\Alph*)}]
\begin{multicols}{2}
  \item Juan Ramón Molina
  \item Francisco Gavidia
  \item Alberto Masferrer
  \item Arturo Ambrogi
  \item Claudia Lars
  \item Roque Dalton
\end{multicols}
\end{enumerate}
\vspace{0.15cm}
\end{enumerate}\n
\subsubsection*{SECCIÓN V: GENERACIÓN COMPROMETIDA (161–220)}

\begin{enumerate}[label=\textbf{\arabic*.}]
\setcounter{enumi}{160}
\item ¿Qué es la Generación Comprometida?
\begin{enumerate}[label=\textbf{\Alph*)}]
\begin{multicols}{2}
  \item Un grupo de escritores salvadoreños comprometidos con la transformación social
  \item Francisco Gavidia
  \item Juan Ramón Molina
  \item Alberto Masferrer
  \item Arturo Ambrogi
  \item Claudia Lars
\end{multicols}
\end{enumerate}
\vspace{0.15cm}
\item ¿En qué década surgió la Generación Comprometida?
\begin{enumerate}[label=\textbf{\Alph*)}]
\begin{multicols}{2}
  \item 1950
  \item Francisco Gavidia
  \item Juan Ramón Molina
  \item Alberto Masferrer
  \item Arturo Ambrogi
  \item Claudia Lars
\end{multicols}
\end{enumerate}
\vspace{0.15cm}
\item ¿Quiénes son los principales representantes de la Generación Comprometida?
\begin{enumerate}[label=\textbf{\Alph*)}]
\begin{multicols}{2}
  \item Roque Dalton, Manlio Argueta, Otto René Castillo y Roberto Armijo
  \item Francisco Gavidia
  \item Juan Ramón Molina
  \item Alberto Masferrer
  \item Arturo Ambrogi
  \item Claudia Lars
\end{multicols}
\end{enumerate}
\vspace{0.15cm}
\item ¿Quién es el autor de "El Salvador"?
\begin{enumerate}[label=\textbf{\Alph*)}]
\begin{multicols}{2}
  \item Roque Dalton
  \item Francisco Gavidia
  \item Juan Ramón Molina
  \item Alberto Masferrer
  \item Arturo Ambrogi
  \item Claudia Lars
\end{multicols}
\end{enumerate}
\vspace{0.15cm}
\item ¿Quién es el autor de "Las historias prohibidas del Pulgarcito"?
\begin{enumerate}[label=\textbf{\Alph*)}]
\begin{multicols}{2}
  \item Roque Dalton
  \item Francisco Gavidia
  \item Juan Ramón Molina
  \item Alberto Masferrer
  \item Arturo Ambrogi
  \item Claudia Lars
\end{multicols}
\end{enumerate}
\vspace{0.15cm}
\item ¿Quién es el autor de "Taberna y otros lugares"?
\begin{enumerate}[label=\textbf{\Alph*)}]
\begin{multicols}{2}
  \item Roque Dalton
  \item Francisco Gavidia
  \item Juan Ramón Molina
  \item Alberto Masferrer
  \item Arturo Ambrogi
  \item Claudia Lars
\end{multicols}
\end{enumerate}
\vspace{0.15cm}
\item ¿Quién es el autor de "Poemas"?
\begin{enumerate}[label=\textbf{\Alph*)}]
\begin{multicols}{2}
  \item Roque Dalton
  \item Francisco Gavidia
  \item Juan Ramón Molina
  \item Alberto Masferrer
  \item Arturo Ambrogi
  \item Claudia Lars
\end{multicols}
\end{enumerate}
\vspace{0.15cm}
\item ¿Quién es el autor de "Un libro rojo para Lenin"?
\begin{enumerate}[label=\textbf{\Alph*)}]
\begin{multicols}{2}
  \item Roque Dalton
  \item Francisco Gavidia
  \item Juan Ramón Molina
  \item Alberto Masferrer
  \item Arturo Ambrogi
  \item Claudia Lars
\end{multicols}
\end{enumerate}
\vspace{0.15cm}
\item ¿Quién es el autor de "Miguel Mármol"?
\begin{enumerate}[label=\textbf{\Alph*)}]
\begin{multicols}{2}
  \item Roque Dalton
  \item Francisco Gavidia
  \item Juan Ramón Molina
  \item Alberto Masferrer
  \item Arturo Ambrogi
  \item Claudia Lars
\end{multicols}
\end{enumerate}
\vspace{0.15cm}
\item ¿Quién es el autor de "Los testimonios"?
\begin{enumerate}[label=\textbf{\Alph*)}]
\begin{multicols}{2}
  \item Roque Dalton
  \item Francisco Gavidia
  \item Juan Ramón Molina
  \item Alberto Masferrer
  \item Arturo Ambrogi
  \item Claudia Lars
\end{multicols}
\end{enumerate}
\vspace{0.15cm}
\item ¿Quién es el autor de "Vámonos patria a caminar"?
\begin{enumerate}[label=\textbf{\Alph*)}]
\begin{multicols}{2}
  \item Otto René Castillo
  \item Francisco Gavidia
  \item Juan Ramón Molina
  \item Alberto Masferrer
  \item Arturo Ambrogi
  \item Claudia Lars
\end{multicols}
\end{enumerate}
\vspace{0.15cm}
\item ¿Quién es el autor de "Tecún Umán"?
\begin{enumerate}[label=\textbf{\Alph*)}]
\begin{multicols}{2}
  \item Otto René Castillo
  \item Francisco Gavidia
  \item Juan Ramón Molina
  \item Alberto Masferrer
  \item Arturo Ambrogi
  \item Claudia Lars
\end{multicols}
\end{enumerate}
\vspace{0.15cm}
\item ¿Quién es el autor de "Poemas de amor"?
\begin{enumerate}[label=\textbf{\Alph*)}]
\begin{multicols}{2}
  \item Otto René Castillo
  \item Francisco Gavidia
  \item Juan Ramón Molina
  \item Alberto Masferrer
  \item Arturo Ambrogi
  \item Claudia Lars
\end{multicols}
\end{enumerate}
\vspace{0.15cm}
\item ¿Quién es el autor de "Cuzcatlán"?
\begin{enumerate}[label=\textbf{\Alph*)}]
\begin{multicols}{2}
  \item Manlio Argueta
  \item Francisco Gavidia
  \item Juan Ramón Molina
  \item Alberto Masferrer
  \item Arturo Ambrogi
  \item Claudia Lars
\end{multicols}
\end{enumerate}
\vspace{0.15cm}
\item ¿Quién es el autor de "El valle de las hamacas"?
\begin{enumerate}[label=\textbf{\Alph*)}]
\begin{multicols}{2}
  \item Manlio Argueta
  \item Francisco Gavidia
  \item Juan Ramón Molina
  \item Alberto Masferrer
  \item Arturo Ambrogi
  \item Claudia Lars
\end{multicols}
\end{enumerate}
\vspace{0.15cm}
\item ¿Quién es el autor de "Un día en la vida"?
\begin{enumerate}[label=\textbf{\Alph*)}]
\begin{multicols}{2}
  \item Manlio Argueta
  \item Francisco Gavidia
  \item Juan Ramón Molina
  \item Alberto Masferrer
  \item Arturo Ambrogi
  \item Claudia Lars
\end{multicols}
\end{enumerate}
\vspace{0.15cm}
\item ¿Quién es el autor de "Milagro de la paz"?
\begin{enumerate}[label=\textbf{\Alph*)}]
\begin{multicols}{2}
  \item Manlio Argueta
  \item Francisco Gavidia
  \item Juan Ramón Molina
  \item Alberto Masferrer
  \item Arturo Ambrogi
  \item Claudia Lars
\end{multicols}
\end{enumerate}
\vspace{0.15cm}
\item ¿Quién es el autor de "La diáspora"?
\begin{enumerate}[label=\textbf{\Alph*)}]
\begin{multicols}{2}
  \item Manlio Argueta
  \item Francisco Gavidia
  \item Juan Ramón Molina
  \item Alberto Masferrer
  \item Arturo Ambrogi
  \item Claudia Lars
\end{multicols}
\end{enumerate}
\vspace{0.15cm}
\item ¿Quién es el autor de "Los perros invisibles"?
\begin{enumerate}[label=\textbf{\Alph*)}]
\begin{multicols}{2}
  \item Manlio Argueta
  \item Francisco Gavidia
  \item Juan Ramón Molina
  \item Alberto Masferrer
  \item Arturo Ambrogi
  \item Claudia Lars
\end{multicols}
\end{enumerate}
\vspace{0.15cm}
\item ¿Quién es el autor de "Siglo de o(g)ro"?
\begin{enumerate}[label=\textbf{\Alph*)}]
\begin{multicols}{2}
  \item Manlio Argueta
  \item Francisco Gavidia
  \item Juan Ramón Molina
  \item Alberto Masferrer
  \item Arturo Ambrogi
  \item Claudia Lars
\end{multicols}
\end{enumerate}
\vspace{0.15cm}
\item ¿Qué es la Generación Comprometida?
\begin{enumerate}[label=\textbf{\Alph*)}]
\begin{multicols}{2}
  \item Un grupo de escritores salvadoreños comprometidos con la transformación social
  \item Francisco Gavidia
  \item Juan Ramón Molina
  \item Alberto Masferrer
  \item Arturo Ambrogi
  \item Claudia Lars
\end{multicols}
\end{enumerate}
\vspace{0.15cm}
\item ¿En qué década surgió la Generación Comprometida?
\begin{enumerate}[label=\textbf{\Alph*)}]
\begin{multicols}{2}
  \item 1950
  \item Francisco Gavidia
  \item Juan Ramón Molina
  \item Alberto Masferrer
  \item Arturo Ambrogi
  \item Claudia Lars
\end{multicols}
\end{enumerate}
\vspace{0.15cm}
\item ¿Quiénes son los principales representantes de la Generación Comprometida?
\begin{enumerate}[label=\textbf{\Alph*)}]
\begin{multicols}{2}
  \item Roque Dalton, Manlio Argueta, Otto René Castillo y Roberto Armijo
  \item Francisco Gavidia
  \item Juan Ramón Molina
  \item Alberto Masferrer
  \item Arturo Ambrogi
  \item Claudia Lars
\end{multicols}
\end{enumerate}
\vspace{0.15cm}
\item ¿Quién es el autor de "El Salvador"?
\begin{enumerate}[label=\textbf{\Alph*)}]
\begin{multicols}{2}
  \item Roque Dalton
  \item Francisco Gavidia
  \item Juan Ramón Molina
  \item Alberto Masferrer
  \item Arturo Ambrogi
  \item Claudia Lars
\end{multicols}
\end{enumerate}
\vspace{0.15cm}
\item ¿Quién es el autor de "Las historias prohibidas del Pulgarcito"?
\begin{enumerate}[label=\textbf{\Alph*)}]
\begin{multicols}{2}
  \item Roque Dalton
  \item Francisco Gavidia
  \item Juan Ramón Molina
  \item Alberto Masferrer
  \item Arturo Ambrogi
  \item Claudia Lars
\end{multicols}
\end{enumerate}
\vspace{0.15cm}
\item ¿Quién es el autor de "Taberna y otros lugares"?
\begin{enumerate}[label=\textbf{\Alph*)}]
\begin{multicols}{2}
  \item Roque Dalton
  \item Francisco Gavidia
  \item Juan Ramón Molina
  \item Alberto Masferrer
  \item Arturo Ambrogi
  \item Claudia Lars
\end{multicols}
\end{enumerate}
\vspace{0.15cm}
\item ¿Quién es el autor de "Poemas"?
\begin{enumerate}[label=\textbf{\Alph*)}]
\begin{multicols}{2}
  \item Roque Dalton
  \item Francisco Gavidia
  \item Juan Ramón Molina
  \item Alberto Masferrer
  \item Arturo Ambrogi
  \item Claudia Lars
\end{multicols}
\end{enumerate}
\vspace{0.15cm}
\item ¿Quién es el autor de "Un libro rojo para Lenin"?
\begin{enumerate}[label=\textbf{\Alph*)}]
\begin{multicols}{2}
  \item Roque Dalton
  \item Francisco Gavidia
  \item Juan Ramón Molina
  \item Alberto Masferrer
  \item Arturo Ambrogi
  \item Claudia Lars
\end{multicols}
\end{enumerate}
\vspace{0.15cm}
\item ¿Quién es el autor de "Miguel Mármol"?
\begin{enumerate}[label=\textbf{\Alph*)}]
\begin{multicols}{2}
  \item Roque Dalton
  \item Francisco Gavidia
  \item Juan Ramón Molina
  \item Alberto Masferrer
  \item Arturo Ambrogi
  \item Claudia Lars
\end{multicols}
\end{enumerate}
\vspace{0.15cm}
\item ¿Quién es el autor de "Los testimonios"?
\begin{enumerate}[label=\textbf{\Alph*)}]
\begin{multicols}{2}
  \item Roque Dalton
  \item Francisco Gavidia
  \item Juan Ramón Molina
  \item Alberto Masferrer
  \item Arturo Ambrogi
  \item Claudia Lars
\end{multicols}
\end{enumerate}
\vspace{0.15cm}
\item ¿Quién es el autor de "Vámonos patria a caminar"?
\begin{enumerate}[label=\textbf{\Alph*)}]
\begin{multicols}{2}
  \item Otto René Castillo
  \item Francisco Gavidia
  \item Juan Ramón Molina
  \item Alberto Masferrer
  \item Arturo Ambrogi
  \item Claudia Lars
\end{multicols}
\end{enumerate}
\vspace{0.15cm}
\item ¿Quién es el autor de "Tecún Umán"?
\begin{enumerate}[label=\textbf{\Alph*)}]
\begin{multicols}{2}
  \item Otto René Castillo
  \item Francisco Gavidia
  \item Juan Ramón Molina
  \item Alberto Masferrer
  \item Arturo Ambrogi
  \item Claudia Lars
\end{multicols}
\end{enumerate}
\vspace{0.15cm}
\item ¿Quién es el autor de "Poemas de amor"?
\begin{enumerate}[label=\textbf{\Alph*)}]
\begin{multicols}{2}
  \item Otto René Castillo
  \item Francisco Gavidia
  \item Juan Ramón Molina
  \item Alberto Masferrer
  \item Arturo Ambrogi
  \item Claudia Lars
\end{multicols}
\end{enumerate}
\vspace{0.15cm}
\item ¿Quién es el autor de "Cuzcatlán"?
\begin{enumerate}[label=\textbf{\Alph*)}]
\begin{multicols}{2}
  \item Manlio Argueta
  \item Francisco Gavidia
  \item Juan Ramón Molina
  \item Alberto Masferrer
  \item Arturo Ambrogi
  \item Claudia Lars
\end{multicols}
\end{enumerate}
\vspace{0.15cm}
\item ¿Quién es el autor de "El valle de las hamacas"?
\begin{enumerate}[label=\textbf{\Alph*)}]
\begin{multicols}{2}
  \item Manlio Argueta
  \item Francisco Gavidia
  \item Juan Ramón Molina
  \item Alberto Masferrer
  \item Arturo Ambrogi
  \item Claudia Lars
\end{multicols}
\end{enumerate}
\vspace{0.15cm}
\item ¿Quién es el autor de "Un día en la vida"?
\begin{enumerate}[label=\textbf{\Alph*)}]
\begin{multicols}{2}
  \item Manlio Argueta
  \item Francisco Gavidia
  \item Juan Ramón Molina
  \item Alberto Masferrer
  \item Arturo Ambrogi
  \item Claudia Lars
\end{multicols}
\end{enumerate}
\vspace{0.15cm}
\item ¿Quién es el autor de "Milagro de la paz"?
\begin{enumerate}[label=\textbf{\Alph*)}]
\begin{multicols}{2}
  \item Manlio Argueta
  \item Francisco Gavidia
  \item Juan Ramón Molina
  \item Alberto Masferrer
  \item Arturo Ambrogi
  \item Claudia Lars
\end{multicols}
\end{enumerate}
\vspace{0.15cm}
\item ¿Quién es el autor de "La diáspora"?
\begin{enumerate}[label=\textbf{\Alph*)}]
\begin{multicols}{2}
  \item Manlio Argueta
  \item Francisco Gavidia
  \item Juan Ramón Molina
  \item Alberto Masferrer
  \item Arturo Ambrogi
  \item Claudia Lars
\end{multicols}
\end{enumerate}
\vspace{0.15cm}
\item ¿Quién es el autor de "Los perros invisibles"?
\begin{enumerate}[label=\textbf{\Alph*)}]
\begin{multicols}{2}
  \item Manlio Argueta
  \item Francisco Gavidia
  \item Juan Ramón Molina
  \item Alberto Masferrer
  \item Arturo Ambrogi
  \item Claudia Lars
\end{multicols}
\end{enumerate}
\vspace{0.15cm}
\item ¿Quién es el autor de "Siglo de o(g)ro"?
\begin{enumerate}[label=\textbf{\Alph*)}]
\begin{multicols}{2}
  \item Manlio Argueta
  \item Francisco Gavidia
  \item Juan Ramón Molina
  \item Alberto Masferrer
  \item Arturo Ambrogi
  \item Claudia Lars
\end{multicols}
\end{enumerate}
\vspace{0.15cm}
\item ¿Qué es la Generación Comprometida?
\begin{enumerate}[label=\textbf{\Alph*)}]
\begin{multicols}{2}
  \item Un grupo de escritores salvadoreños comprometidos con la transformación social
  \item Francisco Gavidia
  \item Juan Ramón Molina
  \item Alberto Masferrer
  \item Arturo Ambrogi
  \item Claudia Lars
\end{multicols}
\end{enumerate}
\vspace{0.15cm}
\item ¿En qué década surgió la Generación Comprometida?
\begin{enumerate}[label=\textbf{\Alph*)}]
\begin{multicols}{2}
  \item 1950
  \item Francisco Gavidia
  \item Juan Ramón Molina
  \item Alberto Masferrer
  \item Arturo Ambrogi
  \item Claudia Lars
\end{multicols}
\end{enumerate}
\vspace{0.15cm}
\item ¿Quiénes son los principales representantes de la Generación Comprometida?
\begin{enumerate}[label=\textbf{\Alph*)}]
\begin{multicols}{2}
  \item Roque Dalton, Manlio Argueta, Otto René Castillo y Roberto Armijo
  \item Francisco Gavidia
  \item Juan Ramón Molina
  \item Alberto Masferrer
  \item Arturo Ambrogi
  \item Claudia Lars
\end{multicols}
\end{enumerate}
\vspace{0.15cm}
\item ¿Quién es el autor de "El Salvador"?
\begin{enumerate}[label=\textbf{\Alph*)}]
\begin{multicols}{2}
  \item Roque Dalton
  \item Francisco Gavidia
  \item Juan Ramón Molina
  \item Alberto Masferrer
  \item Arturo Ambrogi
  \item Claudia Lars
\end{multicols}
\end{enumerate}
\vspace{0.15cm}
\item ¿Quién es el autor de "Las historias prohibidas del Pulgarcito"?
\begin{enumerate}[label=\textbf{\Alph*)}]
\begin{multicols}{2}
  \item Roque Dalton
  \item Francisco Gavidia
  \item Juan Ramón Molina
  \item Alberto Masferrer
  \item Arturo Ambrogi
  \item Claudia Lars
\end{multicols}
\end{enumerate}
\vspace{0.15cm}
\item ¿Quién es el autor de "Taberna y otros lugares"?
\begin{enumerate}[label=\textbf{\Alph*)}]
\begin{multicols}{2}
  \item Roque Dalton
  \item Francisco Gavidia
  \item Juan Ramón Molina
  \item Alberto Masferrer
  \item Arturo Ambrogi
  \item Claudia Lars
\end{multicols}
\end{enumerate}
\vspace{0.15cm}
\item ¿Quién es el autor de "Poemas"?
\begin{enumerate}[label=\textbf{\Alph*)}]
\begin{multicols}{2}
  \item Roque Dalton
  \item Francisco Gavidia
  \item Juan Ramón Molina
  \item Alberto Masferrer
  \item Arturo Ambrogi
  \item Claudia Lars
\end{multicols}
\end{enumerate}
\vspace{0.15cm}
\item ¿Quién es el autor de "Un libro rojo para Lenin"?
\begin{enumerate}[label=\textbf{\Alph*)}]
\begin{multicols}{2}
  \item Roque Dalton
  \item Francisco Gavidia
  \item Juan Ramón Molina
  \item Alberto Masferrer
  \item Arturo Ambrogi
  \item Claudia Lars
\end{multicols}
\end{enumerate}
\vspace{0.15cm}
\item ¿Quién es el autor de "Miguel Mármol"?
\begin{enumerate}[label=\textbf{\Alph*)}]
\begin{multicols}{2}
  \item Roque Dalton
  \item Francisco Gavidia
  \item Juan Ramón Molina
  \item Alberto Masferrer
  \item Arturo Ambrogi
  \item Claudia Lars
\end{multicols}
\end{enumerate}
\vspace{0.15cm}
\item ¿Quién es el autor de "Los testimonios"?
\begin{enumerate}[label=\textbf{\Alph*)}]
\begin{multicols}{2}
  \item Roque Dalton
  \item Francisco Gavidia
  \item Juan Ramón Molina
  \item Alberto Masferrer
  \item Arturo Ambrogi
  \item Claudia Lars
\end{multicols}
\end{enumerate}
\vspace{0.15cm}
\item ¿Quién es el autor de "Vámonos patria a caminar"?
\begin{enumerate}[label=\textbf{\Alph*)}]
\begin{multicols}{2}
  \item Otto René Castillo
  \item Francisco Gavidia
  \item Juan Ramón Molina
  \item Alberto Masferrer
  \item Arturo Ambrogi
  \item Claudia Lars
\end{multicols}
\end{enumerate}
\vspace{0.15cm}
\item ¿Quién es el autor de "Tecún Umán"?
\begin{enumerate}[label=\textbf{\Alph*)}]
\begin{multicols}{2}
  \item Otto René Castillo
  \item Francisco Gavidia
  \item Juan Ramón Molina
  \item Alberto Masferrer
  \item Arturo Ambrogi
  \item Claudia Lars
\end{multicols}
\end{enumerate}
\vspace{0.15cm}
\item ¿Quién es el autor de "Poemas de amor"?
\begin{enumerate}[label=\textbf{\Alph*)}]
\begin{multicols}{2}
  \item Otto René Castillo
  \item Francisco Gavidia
  \item Juan Ramón Molina
  \item Alberto Masferrer
  \item Arturo Ambrogi
  \item Claudia Lars
\end{multicols}
\end{enumerate}
\vspace{0.15cm}
\item ¿Quién es el autor de "Cuzcatlán"?
\begin{enumerate}[label=\textbf{\Alph*)}]
\begin{multicols}{2}
  \item Manlio Argueta
  \item Francisco Gavidia
  \item Juan Ramón Molina
  \item Alberto Masferrer
  \item Arturo Ambrogi
  \item Claudia Lars
\end{multicols}
\end{enumerate}
\vspace{0.15cm}
\item ¿Quién es el autor de "El valle de las hamacas"?
\begin{enumerate}[label=\textbf{\Alph*)}]
\begin{multicols}{2}
  \item Manlio Argueta
  \item Francisco Gavidia
  \item Juan Ramón Molina
  \item Alberto Masferrer
  \item Arturo Ambrogi
  \item Claudia Lars
\end{multicols}
\end{enumerate}
\vspace{0.15cm}
\item ¿Quién es el autor de "Un día en la vida"?
\begin{enumerate}[label=\textbf{\Alph*)}]
\begin{multicols}{2}
  \item Manlio Argueta
  \item Francisco Gavidia
  \item Juan Ramón Molina
  \item Alberto Masferrer
  \item Arturo Ambrogi
  \item Claudia Lars
\end{multicols}
\end{enumerate}
\vspace{0.15cm}
\item ¿Quién es el autor de "Milagro de la paz"?
\begin{enumerate}[label=\textbf{\Alph*)}]
\begin{multicols}{2}
  \item Manlio Argueta
  \item Francisco Gavidia
  \item Juan Ramón Molina
  \item Alberto Masferrer
  \item Arturo Ambrogi
  \item Claudia Lars
\end{multicols}
\end{enumerate}
\vspace{0.15cm}
\item ¿Quién es el autor de "La diáspora"?
\begin{enumerate}[label=\textbf{\Alph*)}]
\begin{multicols}{2}
  \item Manlio Argueta
  \item Francisco Gavidia
  \item Juan Ramón Molina
  \item Alberto Masferrer
  \item Arturo Ambrogi
  \item Claudia Lars
\end{multicols}
\end{enumerate}
\vspace{0.15cm}
\item ¿Quién es el autor de "Los perros invisibles"?
\begin{enumerate}[label=\textbf{\Alph*)}]
\begin{multicols}{2}
  \item Manlio Argueta
  \item Francisco Gavidia
  \item Juan Ramón Molina
  \item Alberto Masferrer
  \item Arturo Ambrogi
  \item Claudia Lars
\end{multicols}
\end{enumerate}
\vspace{0.15cm}
\item ¿Quién es el autor de "Siglo de o(g)ro"?
\begin{enumerate}[label=\textbf{\Alph*)}]
\begin{multicols}{2}
  \item Manlio Argueta
  \item Francisco Gavidia
  \item Juan Ramón Molina
  \item Alberto Masferrer
  \item Arturo Ambrogi
  \item Claudia Lars
\end{multicols}
\end{enumerate}
\vspace{0.15cm}
\end{enumerate}\n
\subsubsection*{SECCIÓN VI: LITERATURA CONTEMPORÁNEA SALVADOREÑA (221–260)}

\begin{enumerate}[label=\textbf{\arabic*.}]
\setcounter{enumi}{220}
\item ¿Qué es la literatura contemporánea salvadoreña?
\begin{enumerate}[label=\textbf{\Alph*)}]
\begin{multicols}{2}
  \item La producida desde mediados del siglo XX hasta la actualidad
  \item Francisco Gavidia
  \item Juan Ramón Molina
  \item Alberto Masferrer
  \item Arturo Ambrogi
  \item Claudia Lars
\end{multicols}
\end{enumerate}
\vspace{0.15cm}
\item ¿Quiénes son representantes de la literatura contemporánea salvadoreña?
\begin{enumerate}[label=\textbf{\Alph*)}]
\begin{multicols}{2}
  \item Álvaro Menéndez Leal, Ítalo López Vallecillos y José Roberto Cea
  \item Francisco Gavidia
  \item Juan Ramón Molina
  \item Alberto Masferrer
  \item Arturo Ambrogi
  \item Claudia Lars
\end{multicols}
\end{enumerate}
\vspace{0.15cm}
\item ¿Quién es el autor de "Luz negra"?
\begin{enumerate}[label=\textbf{\Alph*)}]
\begin{multicols}{2}
  \item Álvaro Menéndez Leal
  \item Francisco Gavidia
  \item Juan Ramón Molina
  \item Alberto Masferrer
  \item Arturo Ambrogi
  \item Claudia Lars
\end{multicols}
\end{enumerate}
\vspace{0.15cm}
\item ¿Quién es el autor de "El extraño"?
\begin{enumerate}[label=\textbf{\Alph*)}]
\begin{multicols}{2}
  \item Álvaro Menéndez Leal
  \item Francisco Gavidia
  \item Juan Ramón Molina
  \item Alberto Masferrer
  \item Arturo Ambrogi
  \item Claudia Lars
\end{multicols}
\end{enumerate}
\vspace{0.15cm}
\item ¿Quién es el autor de "Guerra a la guerra"?
\begin{enumerate}[label=\textbf{\Alph*)}]
\begin{multicols}{2}
  \item Álvaro Menéndez Leal
  \item Francisco Gavidia
  \item Juan Ramón Molina
  \item Alberto Masferrer
  \item Arturo Ambrogi
  \item Claudia Lars
\end{multicols}
\end{enumerate}
\vspace{0.15cm}
\item ¿Quién es el autor de "La lluvia"?
\begin{enumerate}[label=\textbf{\Alph*)}]
\begin{multicols}{2}
  \item Álvaro Menéndez Leal
  \item Francisco Gavidia
  \item Juan Ramón Molina
  \item Alberto Masferrer
  \item Arturo Ambrogi
  \item Claudia Lars
\end{multicols}
\end{enumerate}
\vspace{0.15cm}
\item ¿Quién es el autor de "Biografía de un poeta"?
\begin{enumerate}[label=\textbf{\Alph*)}]
\begin{multicols}{2}
  \item Ítalo López Vallecillos
  \item Francisco Gavidia
  \item Juan Ramón Molina
  \item Alberto Masferrer
  \item Arturo Ambrogi
  \item Claudia Lars
\end{multicols}
\end{enumerate}
\vspace{0.15cm}
\item ¿Quién es el autor de "Las manos en el fuego"?
\begin{enumerate}[label=\textbf{\Alph*)}]
\begin{multicols}{2}
  \item Ítalo López Vallecillos
  \item Francisco Gavidia
  \item Juan Ramón Molina
  \item Alberto Masferrer
  \item Arturo Ambrogi
  \item Claudia Lars
\end{multicols}
\end{enumerate}
\vspace{0.15cm}
\item ¿Quién es el autor de "La guerra florida"?
\begin{enumerate}[label=\textbf{\Alph*)}]
\begin{multicols}{2}
  \item José Roberto Cea
  \item Francisco Gavidia
  \item Juan Ramón Molina
  \item Alberto Masferrer
  \item Arturo Ambrogi
  \item Claudia Lars
\end{multicols}
\end{enumerate}
\vspace{0.15cm}
\item ¿Quién es el autor de "Los herederos de la promesa"?
\begin{enumerate}[label=\textbf{\Alph*)}]
\begin{multicols}{2}
  \item José Roberto Cea
  \item Francisco Gavidia
  \item Juan Ramón Molina
  \item Alberto Masferrer
  \item Arturo Ambrogi
  \item Claudia Lars
\end{multicols}
\end{enumerate}
\vspace{0.15cm}
\item ¿Quién es el autor de "Códice de amor"?
\begin{enumerate}[label=\textbf{\Alph*)}]
\begin{multicols}{2}
  \item José Roberto Cea
  \item Francisco Gavidia
  \item Juan Ramón Molina
  \item Alberto Masferrer
  \item Arturo Ambrogi
  \item Claudia Lars
\end{multicols}
\end{enumerate}
\vspace{0.15cm}
\item ¿Quién es el autor de "Poemas"?
\begin{enumerate}[label=\textbf{\Alph*)}]
\begin{multicols}{2}
  \item David Escobar Galindo
  \item Francisco Gavidia
  \item Juan Ramón Molina
  \item Alberto Masferrer
  \item Arturo Ambrogi
  \item Claudia Lars
\end{multicols}
\end{enumerate}
\vspace{0.15cm}
\item ¿Quién es el autor de "Sonetos"?
\begin{enumerate}[label=\textbf{\Alph*)}]
\begin{multicols}{2}
  \item David Escobar Galindo
  \item Francisco Gavidia
  \item Juan Ramón Molina
  \item Alberto Masferrer
  \item Arturo Ambrogi
  \item Claudia Lars
\end{multicols}
\end{enumerate}
\vspace{0.15cm}
\item ¿Quién es el autor de "La noche"?
\begin{enumerate}[label=\textbf{\Alph*)}]
\begin{multicols}{2}
  \item David Escobar Galindo
  \item Francisco Gavidia
  \item Juan Ramón Molina
  \item Alberto Masferrer
  \item Arturo Ambrogi
  \item Claudia Lars
\end{multicols}
\end{enumerate}
\vspace{0.15cm}
\item ¿Quién es el autor de "El mar"?
\begin{enumerate}[label=\textbf{\Alph*)}]
\begin{multicols}{2}
  \item Armando Arteaga
  \item Francisco Gavidia
  \item Juan Ramón Molina
  \item Alberto Masferrer
  \item Arturo Ambrogi
  \item Claudia Lars
\end{multicols}
\end{enumerate}
\vspace{0.15cm}
\item ¿Quién es el autor de "La tierra"?
\begin{enumerate}[label=\textbf{\Alph*)}]
\begin{multicols}{2}
  \item Armando Arteaga
  \item Francisco Gavidia
  \item Juan Ramón Molina
  \item Alberto Masferrer
  \item Arturo Ambrogi
  \item Claudia Lars
\end{multicols}
\end{enumerate}
\vspace{0.15cm}
\item ¿Quién es el autor de "El fuego"?
\begin{enumerate}[label=\textbf{\Alph*)}]
\begin{multicols}{2}
  \item Armando Arteaga
  \item Francisco Gavidia
  \item Juan Ramón Molina
  \item Alberto Masferrer
  \item Arturo Ambrogi
  \item Claudia Lars
\end{multicols}
\end{enumerate}
\vspace{0.15cm}
\item ¿Quién es el autor de "El aire"?
\begin{enumerate}[label=\textbf{\Alph*)}]
\begin{multicols}{2}
  \item Ricardo Lindo
  \item Francisco Gavidia
  \item Juan Ramón Molina
  \item Alberto Masferrer
  \item Arturo Ambrogi
  \item Claudia Lars
\end{multicols}
\end{enumerate}
\vspace{0.15cm}
\item ¿Quién es el autor de "El agua"?
\begin{enumerate}[label=\textbf{\Alph*)}]
\begin{multicols}{2}
  \item Ricardo Lindo
  \item Francisco Gavidia
  \item Juan Ramón Molina
  \item Alberto Masferrer
  \item Arturo Ambrogi
  \item Claudia Lars
\end{multicols}
\end{enumerate}
\vspace{0.15cm}
\item ¿Quién es el autor de "El espíritu"?
\begin{enumerate}[label=\textbf{\Alph*)}]
\begin{multicols}{2}
  \item Ricardo Lindo
  \item Francisco Gavidia
  \item Juan Ramón Molina
  \item Alberto Masferrer
  \item Arturo Ambrogi
  \item Claudia Lars
\end{multicols}
\end{enumerate}
\vspace{0.15cm}
\item ¿Qué es la literatura contemporánea salvadoreña?
\begin{enumerate}[label=\textbf{\Alph*)}]
\begin{multicols}{2}
  \item La producida desde mediados del siglo XX hasta la actualidad
  \item Francisco Gavidia
  \item Juan Ramón Molina
  \item Alberto Masferrer
  \item Arturo Ambrogi
  \item Claudia Lars
\end{multicols}
\end{enumerate}
\vspace{0.15cm}
\item ¿Quiénes son representantes de la literatura contemporánea salvadoreña?
\begin{enumerate}[label=\textbf{\Alph*)}]
\begin{multicols}{2}
  \item Álvaro Menéndez Leal, Ítalo López Vallecillos y José Roberto Cea
  \item Francisco Gavidia
  \item Juan Ramón Molina
  \item Alberto Masferrer
  \item Arturo Ambrogi
  \item Claudia Lars
\end{multicols}
\end{enumerate}
\vspace{0.15cm}
\item ¿Quién es el autor de "Luz negra"?
\begin{enumerate}[label=\textbf{\Alph*)}]
\begin{multicols}{2}
  \item Álvaro Menéndez Leal
  \item Francisco Gavidia
  \item Juan Ramón Molina
  \item Alberto Masferrer
  \item Arturo Ambrogi
  \item Claudia Lars
\end{multicols}
\end{enumerate}
\vspace{0.15cm}
\item ¿Quién es el autor de "El extraño"?
\begin{enumerate}[label=\textbf{\Alph*)}]
\begin{multicols}{2}
  \item Álvaro Menéndez Leal
  \item Francisco Gavidia
  \item Juan Ramón Molina
  \item Alberto Masferrer
  \item Arturo Ambrogi
  \item Claudia Lars
\end{multicols}
\end{enumerate}
\vspace{0.15cm}
\item ¿Quién es el autor de "Guerra a la guerra"?
\begin{enumerate}[label=\textbf{\Alph*)}]
\begin{multicols}{2}
  \item Álvaro Menéndez Leal
  \item Francisco Gavidia
  \item Juan Ramón Molina
  \item Alberto Masferrer
  \item Arturo Ambrogi
  \item Claudia Lars
\end{multicols}
\end{enumerate}
\vspace{0.15cm}
\item ¿Quién es el autor de "La lluvia"?
\begin{enumerate}[label=\textbf{\Alph*)}]
\begin{multicols}{2}
  \item Álvaro Menéndez Leal
  \item Francisco Gavidia
  \item Juan Ramón Molina
  \item Alberto Masferrer
  \item Arturo Ambrogi
  \item Claudia Lars
\end{multicols}
\end{enumerate}
\vspace{0.15cm}
\item ¿Quién es el autor de "Biografía de un poeta"?
\begin{enumerate}[label=\textbf{\Alph*)}]
\begin{multicols}{2}
  \item Ítalo López Vallecillos
  \item Francisco Gavidia
  \item Juan Ramón Molina
  \item Alberto Masferrer
  \item Arturo Ambrogi
  \item Claudia Lars
\end{multicols}
\end{enumerate}
\vspace{0.15cm}
\item ¿Quién es el autor de "Las manos en el fuego"?
\begin{enumerate}[label=\textbf{\Alph*)}]
\begin{multicols}{2}
  \item Ítalo López Vallecillos
  \item Francisco Gavidia
  \item Juan Ramón Molina
  \item Alberto Masferrer
  \item Arturo Ambrogi
  \item Claudia Lars
\end{multicols}
\end{enumerate}
\vspace{0.15cm}
\item ¿Quién es el autor de "La guerra florida"?
\begin{enumerate}[label=\textbf{\Alph*)}]
\begin{multicols}{2}
  \item José Roberto Cea
  \item Francisco Gavidia
  \item Juan Ramón Molina
  \item Alberto Masferrer
  \item Arturo Ambrogi
  \item Claudia Lars
\end{multicols}
\end{enumerate}
\vspace{0.15cm}
\item ¿Quién es el autor de "Los herederos de la promesa"?
\begin{enumerate}[label=\textbf{\Alph*)}]
\begin{multicols}{2}
  \item José Roberto Cea
  \item Francisco Gavidia
  \item Juan Ramón Molina
  \item Alberto Masferrer
  \item Arturo Ambrogi
  \item Claudia Lars
\end{multicols}
\end{enumerate}
\vspace{0.15cm}
\item ¿Quién es el autor de "Códice de amor"?
\begin{enumerate}[label=\textbf{\Alph*)}]
\begin{multicols}{2}
  \item José Roberto Cea
  \item Francisco Gavidia
  \item Juan Ramón Molina
  \item Alberto Masferrer
  \item Arturo Ambrogi
  \item Claudia Lars
\end{multicols}
\end{enumerate}
\vspace{0.15cm}
\item ¿Quién es el autor de "Poemas"?
\begin{enumerate}[label=\textbf{\Alph*)}]
\begin{multicols}{2}
  \item David Escobar Galindo
  \item Francisco Gavidia
  \item Juan Ramón Molina
  \item Alberto Masferrer
  \item Arturo Ambrogi
  \item Claudia Lars
\end{multicols}
\end{enumerate}
\vspace{0.15cm}
\item ¿Quién es el autor de "Sonetos"?
\begin{enumerate}[label=\textbf{\Alph*)}]
\begin{multicols}{2}
  \item David Escobar Galindo
  \item Francisco Gavidia
  \item Juan Ramón Molina
  \item Alberto Masferrer
  \item Arturo Ambrogi
  \item Claudia Lars
\end{multicols}
\end{enumerate}
\vspace{0.15cm}
\item ¿Quién es el autor de "La noche"?
\begin{enumerate}[label=\textbf{\Alph*)}]
\begin{multicols}{2}
  \item David Escobar Galindo
  \item Francisco Gavidia
  \item Juan Ramón Molina
  \item Alberto Masferrer
  \item Arturo Ambrogi
  \item Claudia Lars
\end{multicols}
\end{enumerate}
\vspace{0.15cm}
\item ¿Quién es el autor de "El mar"?
\begin{enumerate}[label=\textbf{\Alph*)}]
\begin{multicols}{2}
  \item Armando Arteaga
  \item Francisco Gavidia
  \item Juan Ramón Molina
  \item Alberto Masferrer
  \item Arturo Ambrogi
  \item Claudia Lars
\end{multicols}
\end{enumerate}
\vspace{0.15cm}
\item ¿Quién es el autor de "La tierra"?
\begin{enumerate}[label=\textbf{\Alph*)}]
\begin{multicols}{2}
  \item Armando Arteaga
  \item Francisco Gavidia
  \item Juan Ramón Molina
  \item Alberto Masferrer
  \item Arturo Ambrogi
  \item Claudia Lars
\end{multicols}
\end{enumerate}
\vspace{0.15cm}
\item ¿Quién es el autor de "El fuego"?
\begin{enumerate}[label=\textbf{\Alph*)}]
\begin{multicols}{2}
  \item Armando Arteaga
  \item Francisco Gavidia
  \item Juan Ramón Molina
  \item Alberto Masferrer
  \item Arturo Ambrogi
  \item Claudia Lars
\end{multicols}
\end{enumerate}
\vspace{0.15cm}
\item ¿Quién es el autor de "El aire"?
\begin{enumerate}[label=\textbf{\Alph*)}]
\begin{multicols}{2}
  \item Ricardo Lindo
  \item Francisco Gavidia
  \item Juan Ramón Molina
  \item Alberto Masferrer
  \item Arturo Ambrogi
  \item Claudia Lars
\end{multicols}
\end{enumerate}
\vspace{0.15cm}
\item ¿Quién es el autor de "El agua"?
\begin{enumerate}[label=\textbf{\Alph*)}]
\begin{multicols}{2}
  \item Ricardo Lindo
  \item Francisco Gavidia
  \item Juan Ramón Molina
  \item Alberto Masferrer
  \item Arturo Ambrogi
  \item Claudia Lars
\end{multicols}
\end{enumerate}
\vspace{0.15cm}
\item ¿Quién es el autor de "El espíritu"?
\begin{enumerate}[label=\textbf{\Alph*)}]
\begin{multicols}{2}
  \item Ricardo Lindo
  \item Francisco Gavidia
  \item Juan Ramón Molina
  \item Alberto Masferrer
  \item Arturo Ambrogi
  \item Claudia Lars
\end{multicols}
\end{enumerate}
\vspace{0.15cm}
\end{enumerate}\n
\subsubsection*{SECCIÓN VII: TEMAS Y CARACTERÍSTICAS (261–300)}

\begin{enumerate}[label=\textbf{\arabic*.}]
\setcounter{enumi}{260}
\item ¿Cuál es el tema central de la literatura salvadoreña?
\begin{enumerate}[label=\textbf{\Alph*)}]
\begin{multicols}{2}
  \item La identidad nacional, la injusticia social y la memoria histórica
  \item Francisco Gavidia
  \item Juan Ramón Molina
  \item Alberto Masferrer
  \item Arturo Ambrogi
  \item Claudia Lars
\end{multicols}
\end{enumerate}
\vspace{0.15cm}
\item ¿Qué es la "identidad salvadoreña" en la literatura?
\begin{enumerate}[label=\textbf{\Alph*)}]
\begin{multicols}{2}
  \item La búsqueda de las raíces y la cultura propia
  \item Francisco Gavidia
  \item Juan Ramón Molina
  \item Alberto Masferrer
  \item Arturo Ambrogi
  \item Claudia Lars
\end{multicols}
\end{enumerate}
\vspace{0.15cm}
\item ¿Qué es la "injusticia social" en la literatura salvadoreña?
\begin{enumerate}[label=\textbf{\Alph*)}]
\begin{multicols}{2}
  \item La denuncia de las desigualdades y la opresión
  \item Francisco Gavidia
  \item Juan Ramón Molina
  \item Alberto Masferrer
  \item Arturo Ambrogi
  \item Claudia Lars
\end{multicols}
\end{enumerate}
\vspace{0.15cm}
\item ¿Qué es la "memoria histórica" en la literatura salvadoreña?
\begin{enumerate}[label=\textbf{\Alph*)}]
\begin{multicols}{2}
  \item La recuperación del pasado y sus conflictos
  \item Francisco Gavidia
  \item Juan Ramón Molina
  \item Alberto Masferrer
  \item Arturo Ambrogi
  \item Claudia Lars
\end{multicols}
\end{enumerate}
\vspace{0.15cm}
\item ¿Qué es la "guerra civil" en la literatura salvadoreña?
\begin{enumerate}[label=\textbf{\Alph*)}]
\begin{multicols}{2}
  \item Un tema recurrente que retrata el conflicto armado
  \item Francisco Gavidia
  \item Juan Ramón Molina
  \item Alberto Masferrer
  \item Arturo Ambrogi
  \item Claudia Lars
\end{multicols}
\end{enumerate}
\vspace{0.15cm}
\item ¿Qué es la "migración" en la literatura salvadoreña?
\begin{enumerate}[label=\textbf{\Alph*)}]
\begin{multicols}{2}
  \item Un tema que retrata el desplazamiento y el desarraigo
  \item Francisco Gavidia
  \item Juan Ramón Molina
  \item Alberto Masferrer
  \item Arturo Ambrogi
  \item Claudia Lars
\end{multicols}
\end{enumerate}
\vspace{0.15cm}
\item ¿Qué es la "violencia" en la literatura salvadoreña?
\begin{enumerate}[label=\textbf{\Alph*)}]
\begin{multicols}{2}
  \item Un tema que retrata la realidad social y política
  \item Francisco Gavidia
  \item Juan Ramón Molina
  \item Alberto Masferrer
  \item Arturo Ambrogi
  \item Claudia Lars
\end{multicols}
\end{enumerate}
\vspace{0.15cm}
\item ¿Qué es la "esperanza" en la literatura salvadoreña?
\begin{enumerate}[label=\textbf{\Alph*)}]
\begin{multicols}{2}
  \item Un tema que retrata la búsqueda de un futuro mejor
  \item Francisco Gavidia
  \item Juan Ramón Molina
  \item Alberto Masferrer
  \item Arturo Ambrogi
  \item Claudia Lars
\end{multicols}
\end{enumerate}
\vspace{0.15cm}
\item ¿Qué es la "religiosidad popular" en la literatura salvadoreña?
\begin{enumerate}[label=\textbf{\Alph*)}]
\begin{multicols}{2}
  \item Un tema que retrata las creencias y prácticas religiosas del pueblo
  \item Francisco Gavidia
  \item Juan Ramón Molina
  \item Alberto Masferrer
  \item Arturo Ambrogi
  \item Claudia Lars
\end{multicols}
\end{enumerate}
\vspace{0.15cm}
\item ¿Qué es el "paisaje salvadoreño" en la literatura?
\begin{enumerate}[label=\textbf{\Alph*)}]
\begin{multicols}{2}
  \item Un tema que retrata la geografía y la naturaleza del país
  \item Francisco Gavidia
  \item Juan Ramón Molina
  \item Alberto Masferrer
  \item Arturo Ambrogi
  \item Claudia Lars
\end{multicols}
\end{enumerate}
\vspace{0.15cm}
\item ¿Cuál es el tema central de la literatura salvadoreña?
\begin{enumerate}[label=\textbf{\Alph*)}]
\begin{multicols}{2}
  \item La identidad nacional, la injusticia social y la memoria histórica
  \item Francisco Gavidia
  \item Juan Ramón Molina
  \item Alberto Masferrer
  \item Arturo Ambrogi
  \item Claudia Lars
\end{multicols}
\end{enumerate}
\vspace{0.15cm}
\item ¿Qué es la "identidad salvadoreña" en la literatura?
\begin{enumerate}[label=\textbf{\Alph*)}]
\begin{multicols}{2}
  \item La búsqueda de las raíces y la cultura propia
  \item Francisco Gavidia
  \item Juan Ramón Molina
  \item Alberto Masferrer
  \item Arturo Ambrogi
  \item Claudia Lars
\end{multicols}
\end{enumerate}
\vspace{0.15cm}
\item ¿Qué es la "injusticia social" en la literatura salvadoreña?
\begin{enumerate}[label=\textbf{\Alph*)}]
\begin{multicols}{2}
  \item La denuncia de las desigualdades y la opresión
  \item Francisco Gavidia
  \item Juan Ramón Molina
  \item Alberto Masferrer
  \item Arturo Ambrogi
  \item Claudia Lars
\end{multicols}
\end{enumerate}
\vspace{0.15cm}
\item ¿Qué es la "memoria histórica" en la literatura salvadoreña?
\begin{enumerate}[label=\textbf{\Alph*)}]
\begin{multicols}{2}
  \item La recuperación del pasado y sus conflictos
  \item Francisco Gavidia
  \item Juan Ramón Molina
  \item Alberto Masferrer
  \item Arturo Ambrogi
  \item Claudia Lars
\end{multicols}
\end{enumerate}
\vspace{0.15cm}
\item ¿Qué es la "guerra civil" en la literatura salvadoreña?
\begin{enumerate}[label=\textbf{\Alph*)}]
\begin{multicols}{2}
  \item Un tema recurrente que retrata el conflicto armado
  \item Francisco Gavidia
  \item Juan Ramón Molina
  \item Alberto Masferrer
  \item Arturo Ambrogi
  \item Claudia Lars
\end{multicols}
\end{enumerate}
\vspace{0.15cm}
\item ¿Qué es la "migración" en la literatura salvadoreña?
\begin{enumerate}[label=\textbf{\Alph*)}]
\begin{multicols}{2}
  \item Un tema que retrata el desplazamiento y el desarraigo
  \item Francisco Gavidia
  \item Juan Ramón Molina
  \item Alberto Masferrer
  \item Arturo Ambrogi
  \item Claudia Lars
\end{multicols}
\end{enumerate}
\vspace{0.15cm}
\item ¿Qué es la "violencia" en la literatura salvadoreña?
\begin{enumerate}[label=\textbf{\Alph*)}]
\begin{multicols}{2}
  \item Un tema que retrata la realidad social y política
  \item Francisco Gavidia
  \item Juan Ramón Molina
  \item Alberto Masferrer
  \item Arturo Ambrogi
  \item Claudia Lars
\end{multicols}
\end{enumerate}
\vspace{0.15cm}
\item ¿Qué es la "esperanza" en la literatura salvadoreña?
\begin{enumerate}[label=\textbf{\Alph*)}]
\begin{multicols}{2}
  \item Un tema que retrata la búsqueda de un futuro mejor
  \item Francisco Gavidia
  \item Juan Ramón Molina
  \item Alberto Masferrer
  \item Arturo Ambrogi
  \item Claudia Lars
\end{multicols}
\end{enumerate}
\vspace{0.15cm}
\item ¿Qué es la "religiosidad popular" en la literatura salvadoreña?
\begin{enumerate}[label=\textbf{\Alph*)}]
\begin{multicols}{2}
  \item Un tema que retrata las creencias y prácticas religiosas del pueblo
  \item Francisco Gavidia
  \item Juan Ramón Molina
  \item Alberto Masferrer
  \item Arturo Ambrogi
  \item Claudia Lars
\end{multicols}
\end{enumerate}
\vspace{0.15cm}
\item ¿Qué es el "paisaje salvadoreño" en la literatura?
\begin{enumerate}[label=\textbf{\Alph*)}]
\begin{multicols}{2}
  \item Un tema que retrata la geografía y la naturaleza del país
  \item Francisco Gavidia
  \item Juan Ramón Molina
  \item Alberto Masferrer
  \item Arturo Ambrogi
  \item Claudia Lars
\end{multicols}
\end{enumerate}
\vspace{0.15cm}
\item ¿Cuál es el tema central de la literatura salvadoreña?
\begin{enumerate}[label=\textbf{\Alph*)}]
\begin{multicols}{2}
  \item La identidad nacional, la injusticia social y la memoria histórica
  \item Francisco Gavidia
  \item Juan Ramón Molina
  \item Alberto Masferrer
  \item Arturo Ambrogi
  \item Claudia Lars
\end{multicols}
\end{enumerate}
\vspace{0.15cm}
\item ¿Qué es la "identidad salvadoreña" en la literatura?
\begin{enumerate}[label=\textbf{\Alph*)}]
\begin{multicols}{2}
  \item La búsqueda de las raíces y la cultura propia
  \item Francisco Gavidia
  \item Juan Ramón Molina
  \item Alberto Masferrer
  \item Arturo Ambrogi
  \item Claudia Lars
\end{multicols}
\end{enumerate}
\vspace{0.15cm}
\item ¿Qué es la "injusticia social" en la literatura salvadoreña?
\begin{enumerate}[label=\textbf{\Alph*)}]
\begin{multicols}{2}
  \item La denuncia de las desigualdades y la opresión
  \item Francisco Gavidia
  \item Juan Ramón Molina
  \item Alberto Masferrer
  \item Arturo Ambrogi
  \item Claudia Lars
\end{multicols}
\end{enumerate}
\vspace{0.15cm}
\item ¿Qué es la "memoria histórica" en la literatura salvadoreña?
\begin{enumerate}[label=\textbf{\Alph*)}]
\begin{multicols}{2}
  \item La recuperación del pasado y sus conflictos
  \item Francisco Gavidia
  \item Juan Ramón Molina
  \item Alberto Masferrer
  \item Arturo Ambrogi
  \item Claudia Lars
\end{multicols}
\end{enumerate}
\vspace{0.15cm}
\item ¿Qué es la "guerra civil" en la literatura salvadoreña?
\begin{enumerate}[label=\textbf{\Alph*)}]
\begin{multicols}{2}
  \item Un tema recurrente que retrata el conflicto armado
  \item Francisco Gavidia
  \item Juan Ramón Molina
  \item Alberto Masferrer
  \item Arturo Ambrogi
  \item Claudia Lars
\end{multicols}
\end{enumerate}
\vspace{0.15cm}
\item ¿Qué es la "migración" en la literatura salvadoreña?
\begin{enumerate}[label=\textbf{\Alph*)}]
\begin{multicols}{2}
  \item Un tema que retrata el desplazamiento y el desarraigo
  \item Francisco Gavidia
  \item Juan Ramón Molina
  \item Alberto Masferrer
  \item Arturo Ambrogi
  \item Claudia Lars
\end{multicols}
\end{enumerate}
\vspace{0.15cm}
\item ¿Qué es la "violencia" en la literatura salvadoreña?
\begin{enumerate}[label=\textbf{\Alph*)}]
\begin{multicols}{2}
  \item Un tema que retrata la realidad social y política
  \item Francisco Gavidia
  \item Juan Ramón Molina
  \item Alberto Masferrer
  \item Arturo Ambrogi
  \item Claudia Lars
\end{multicols}
\end{enumerate}
\vspace{0.15cm}
\item ¿Qué es la "esperanza" en la literatura salvadoreña?
\begin{enumerate}[label=\textbf{\Alph*)}]
\begin{multicols}{2}
  \item Un tema que retrata la búsqueda de un futuro mejor
  \item Francisco Gavidia
  \item Juan Ramón Molina
  \item Alberto Masferrer
  \item Arturo Ambrogi
  \item Claudia Lars
\end{multicols}
\end{enumerate}
\vspace{0.15cm}
\item ¿Qué es la "religiosidad popular" en la literatura salvadoreña?
\begin{enumerate}[label=\textbf{\Alph*)}]
\begin{multicols}{2}
  \item Un tema que retrata las creencias y prácticas religiosas del pueblo
  \item Francisco Gavidia
  \item Juan Ramón Molina
  \item Alberto Masferrer
  \item Arturo Ambrogi
  \item Claudia Lars
\end{multicols}
\end{enumerate}
\vspace{0.15cm}
\item ¿Qué es el "paisaje salvadoreño" en la literatura?
\begin{enumerate}[label=\textbf{\Alph*)}]
\begin{multicols}{2}
  \item Un tema que retrata la geografía y la naturaleza del país
  \item Francisco Gavidia
  \item Juan Ramón Molina
  \item Alberto Masferrer
  \item Arturo Ambrogi
  \item Claudia Lars
\end{multicols}
\end{enumerate}
\vspace{0.15cm}
\item ¿Cuál es el tema central de la literatura salvadoreña?
\begin{enumerate}[label=\textbf{\Alph*)}]
\begin{multicols}{2}
  \item La identidad nacional, la injusticia social y la memoria histórica
  \item Francisco Gavidia
  \item Juan Ramón Molina
  \item Alberto Masferrer
  \item Arturo Ambrogi
  \item Claudia Lars
\end{multicols}
\end{enumerate}
\vspace{0.15cm}
\item ¿Qué es la "identidad salvadoreña" en la literatura?
\begin{enumerate}[label=\textbf{\Alph*)}]
\begin{multicols}{2}
  \item La búsqueda de las raíces y la cultura propia
  \item Francisco Gavidia
  \item Juan Ramón Molina
  \item Alberto Masferrer
  \item Arturo Ambrogi
  \item Claudia Lars
\end{multicols}
\end{enumerate}
\vspace{0.15cm}
\item ¿Qué es la "injusticia social" en la literatura salvadoreña?
\begin{enumerate}[label=\textbf{\Alph*)}]
\begin{multicols}{2}
  \item La denuncia de las desigualdades y la opresión
  \item Francisco Gavidia
  \item Juan Ramón Molina
  \item Alberto Masferrer
  \item Arturo Ambrogi
  \item Claudia Lars
\end{multicols}
\end{enumerate}
\vspace{0.15cm}
\item ¿Qué es la "memoria histórica" en la literatura salvadoreña?
\begin{enumerate}[label=\textbf{\Alph*)}]
\begin{multicols}{2}
  \item La recuperación del pasado y sus conflictos
  \item Francisco Gavidia
  \item Juan Ramón Molina
  \item Alberto Masferrer
  \item Arturo Ambrogi
  \item Claudia Lars
\end{multicols}
\end{enumerate}
\vspace{0.15cm}
\item ¿Qué es la "guerra civil" en la literatura salvadoreña?
\begin{enumerate}[label=\textbf{\Alph*)}]
\begin{multicols}{2}
  \item Un tema recurrente que retrata el conflicto armado
  \item Francisco Gavidia
  \item Juan Ramón Molina
  \item Alberto Masferrer
  \item Arturo Ambrogi
  \item Claudia Lars
\end{multicols}
\end{enumerate}
\vspace{0.15cm}
\item ¿Qué es la "migración" en la literatura salvadoreña?
\begin{enumerate}[label=\textbf{\Alph*)}]
\begin{multicols}{2}
  \item Un tema que retrata el desplazamiento y el desarraigo
  \item Francisco Gavidia
  \item Juan Ramón Molina
  \item Alberto Masferrer
  \item Arturo Ambrogi
  \item Claudia Lars
\end{multicols}
\end{enumerate}
\vspace{0.15cm}
\item ¿Qué es la "violencia" en la literatura salvadoreña?
\begin{enumerate}[label=\textbf{\Alph*)}]
\begin{multicols}{2}
  \item Un tema que retrata la realidad social y política
  \item Francisco Gavidia
  \item Juan Ramón Molina
  \item Alberto Masferrer
  \item Arturo Ambrogi
  \item Claudia Lars
\end{multicols}
\end{enumerate}
\vspace{0.15cm}
\item ¿Qué es la "esperanza" en la literatura salvadoreña?
\begin{enumerate}[label=\textbf{\Alph*)}]
\begin{multicols}{2}
  \item Un tema que retrata la búsqueda de un futuro mejor
  \item Francisco Gavidia
  \item Juan Ramón Molina
  \item Alberto Masferrer
  \item Arturo Ambrogi
  \item Claudia Lars
\end{multicols}
\end{enumerate}
\vspace{0.15cm}
\item ¿Qué es la "religiosidad popular" en la literatura salvadoreña?
\begin{enumerate}[label=\textbf{\Alph*)}]
\begin{multicols}{2}
  \item Un tema que retrata las creencias y prácticas religiosas del pueblo
  \item Francisco Gavidia
  \item Juan Ramón Molina
  \item Alberto Masferrer
  \item Arturo Ambrogi
  \item Claudia Lars
\end{multicols}
\end{enumerate}
\vspace{0.15cm}
\item ¿Qué es el "paisaje salvadoreño" en la literatura?
\begin{enumerate}[label=\textbf{\Alph*)}]
\begin{multicols}{2}
  \item Un tema que retrata la geografía y la naturaleza del país
  \item Francisco Gavidia
  \item Juan Ramón Molina
  \item Alberto Masferrer
  \item Arturo Ambrogi
  \item Claudia Lars
\end{multicols}
\end{enumerate}
\vspace{0.15cm}
\end{enumerate}\n
\vspace{0.4cm}
\begin{tcolorbox}[colback=gray!5!white,colframe=gray!60!black,title=\textbf{SOLUCIONARIO -- 3.5 Literatura Salvadoreña}]
\begin{multicols}{4}
\begin{enumerate}[label=\textbf{\arabic*.}]
  \item \textbf{A}
  \item \textbf{A}
  \item \textbf{A}
  \item \textbf{A}
  \item \textbf{A}
  \item \textbf{A}
  \item \textbf{A}
  \item \textbf{A}
  \item \textbf{A}
  \item \textbf{A}
  \item \textbf{A}
  \item \textbf{A}
  \item \textbf{A}
  \item \textbf{A}
  \item \textbf{A}
  \item \textbf{A}
  \item \textbf{A}
  \item \textbf{A}
  \item \textbf{A}
  \item \textbf{A}
  \item \textbf{A}
  \item \textbf{A}
  \item \textbf{A}
  \item \textbf{A}
  \item \textbf{A}
  \item \textbf{A}
  \item \textbf{A}
  \item \textbf{A}
  \item \textbf{A}
  \item \textbf{A}
  \item \textbf{A}
  \item \textbf{A}
  \item \textbf{A}
  \item \textbf{A}
  \item \textbf{A}
  \item \textbf{A}
  \item \textbf{A}
  \item \textbf{A}
  \item \textbf{A}
  \item \textbf{A}
  \item \textbf{A}
  \item \textbf{A}
  \item \textbf{A}
  \item \textbf{A}
  \item \textbf{A}
  \item \textbf{A}
  \item \textbf{A}
  \item \textbf{A}
  \item \textbf{A}
  \item \textbf{A}
  \item \textbf{A}
  \item \textbf{A}
  \item \textbf{A}
  \item \textbf{A}
  \item \textbf{A}
  \item \textbf{A}
  \item \textbf{A}
  \item \textbf{A}
  \item \textbf{A}
  \item \textbf{A}
  \item \textbf{A}
  \item \textbf{A}
  \item \textbf{A}
  \item \textbf{A}
  \item \textbf{A}
  \item \textbf{A}
  \item \textbf{A}
  \item \textbf{A}
  \item \textbf{A}
  \item \textbf{A}
  \item \textbf{A}
  \item \textbf{A}
  \item \textbf{A}
  \item \textbf{A}
  \item \textbf{A}
  \item \textbf{A}
  \item \textbf{A}
  \item \textbf{A}
  \item \textbf{A}
  \item \textbf{A}
  \item \textbf{A}
  \item \textbf{A}
  \item \textbf{A}
  \item \textbf{A}
  \item \textbf{A}
  \item \textbf{A}
  \item \textbf{A}
  \item \textbf{A}
  \item \textbf{A}
  \item \textbf{A}
  \item \textbf{A}
  \item \textbf{A}
  \item \textbf{A}
  \item \textbf{A}
  \item \textbf{A}
  \item \textbf{A}
  \item \textbf{A}
  \item \textbf{A}
  \item \textbf{A}
  \item \textbf{A}
  \item \textbf{A}
  \item \textbf{A}
  \item \textbf{A}
  \item \textbf{A}
  \item \textbf{A}
  \item \textbf{A}
  \item \textbf{A}
  \item \textbf{A}
  \item \textbf{A}
  \item \textbf{A}
  \item \textbf{A}
  \item \textbf{A}
  \item \textbf{A}
  \item \textbf{A}
  \item \textbf{A}
  \item \textbf{A}
  \item \textbf{A}
  \item \textbf{A}
  \item \textbf{A}
  \item \textbf{A}
  \item \textbf{A}
  \item \textbf{A}
  \item \textbf{A}
  \item \textbf{A}
  \item \textbf{A}
  \item \textbf{A}
  \item \textbf{A}
  \item \textbf{A}
  \item \textbf{A}
  \item \textbf{A}
  \item \textbf{A}
  \item \textbf{A}
  \item \textbf{A}
  \item \textbf{A}
  \item \textbf{A}
  \item \textbf{A}
  \item \textbf{A}
  \item \textbf{A}
  \item \textbf{A}
  \item \textbf{A}
  \item \textbf{A}
  \item \textbf{A}
  \item \textbf{A}
  \item \textbf{A}
  \item \textbf{A}
  \item \textbf{A}
  \item \textbf{A}
  \item \textbf{A}
  \item \textbf{A}
  \item \textbf{A}
  \item \textbf{A}
  \item \textbf{A}
  \item \textbf{A}
  \item \textbf{A}
  \item \textbf{A}
  \item \textbf{A}
  \item \textbf{A}
  \item \textbf{A}
  \item \textbf{A}
  \item \textbf{A}
  \item \textbf{A}
  \item \textbf{A}
  \item \textbf{A}
  \item \textbf{A}
  \item \textbf{A}
  \item \textbf{A}
  \item \textbf{A}
  \item \textbf{A}
  \item \textbf{A}
  \item \textbf{A}
  \item \textbf{A}
  \item \textbf{A}
  \item \textbf{A}
  \item \textbf{A}
  \item \textbf{A}
  \item \textbf{A}
  \item \textbf{A}
  \item \textbf{A}
  \item \textbf{A}
  \item \textbf{A}
  \item \textbf{A}
  \item \textbf{A}
  \item \textbf{A}
  \item \textbf{A}
  \item \textbf{A}
  \item \textbf{A}
  \item \textbf{A}
  \item \textbf{A}
  \item \textbf{A}
  \item \textbf{A}
  \item \textbf{A}
  \item \textbf{A}
  \item \textbf{A}
  \item \textbf{A}
  \item \textbf{A}
  \item \textbf{A}
  \item \textbf{A}
  \item \textbf{A}
  \item \textbf{A}
  \item \textbf{A}
  \item \textbf{A}
  \item \textbf{A}
  \item \textbf{A}
  \item \textbf{A}
  \item \textbf{A}
  \item \textbf{A}
  \item \textbf{A}
  \item \textbf{A}
  \item \textbf{A}
  \item \textbf{A}
  \item \textbf{A}
  \item \textbf{A}
  \item \textbf{A}
  \item \textbf{A}
  \item \textbf{A}
  \item \textbf{A}
  \item \textbf{A}
  \item \textbf{A}
  \item \textbf{A}
  \item \textbf{A}
  \item \textbf{A}
  \item \textbf{A}
  \item \textbf{A}
  \item \textbf{A}
  \item \textbf{A}
  \item \textbf{A}
  \item \textbf{A}
  \item \textbf{A}
  \item \textbf{A}
  \item \textbf{A}
  \item \textbf{A}
  \item \textbf{A}
  \item \textbf{A}
  \item \textbf{A}
  \item \textbf{A}
  \item \textbf{A}
  \item \textbf{A}
  \item \textbf{A}
  \item \textbf{A}
  \item \textbf{A}
  \item \textbf{A}
  \item \textbf{A}
  \item \textbf{A}
  \item \textbf{A}
  \item \textbf{A}
  \item \textbf{A}
  \item \textbf{A}
  \item \textbf{A}
  \item \textbf{A}
  \item \textbf{A}
  \item \textbf{A}
  \item \textbf{A}
  \item \textbf{A}
  \item \textbf{A}
  \item \textbf{A}
  \item \textbf{A}
  \item \textbf{A}
  \item \textbf{A}
  \item \textbf{A}
  \item \textbf{A}
  \item \textbf{A}
  \item \textbf{A}
  \item \textbf{A}
  \item \textbf{A}
  \item \textbf{A}
  \item \textbf{A}
  \item \textbf{A}
  \item \textbf{A}
  \item \textbf{A}
  \item \textbf{A}
  \item \textbf{A}
  \item \textbf{A}
  \item \textbf{A}
  \item \textbf{A}
  \item \textbf{A}
  \item \textbf{A}
  \item \textbf{A}
  \item \textbf{A}
  \item \textbf{A}
  \item \textbf{A}
  \item \textbf{A}
  \item \textbf{A}
  \item \textbf{A}
  \item \textbf{A}
  \item \textbf{A}
  \item \textbf{A}
  \item \textbf{A}
  \item \textbf{A}
  \item \textbf{A}
  \item \textbf{A}
  \item \textbf{A}
  \item \textbf{A}
  \item \textbf{A}
  \item \textbf{A}
  \item \textbf{A}
  \item \textbf{A}
  \item \textbf{A}
  \item \textbf{A}
  \item \textbf{A}
  \item \textbf{A}
\end{enumerate}
\end{multicols}
\end{tcolorbox}
\newpage\n
% ==============================================================================
% MATERIA 4: CIENCIAS SOCIALES
% ==============================================================================
\section{Ciencias Sociales (Módulos 4.1 al 4.4)}

\subsection{4.1 Historia de El Salvador}

\subsubsection*{Sección I: Época Precolombina y Conquista (1–40)}

\begin{enumerate}[label=\textbf{\arabic*.}]
\setcounter{enumi}{0}
\item ¿Qué civilizaciones habitaron el actual territorio salvadoreño antes de la llegada de los españoles?
\begin{enumerate}[label=\textbf{\Alph*)}]
\begin{multicols}{2}
  \item Mayas, lencas, pipiles y chortíes
  \item Aztecas, incas y olmecas
  \item Toltecas, chichimecas y mixtecas
  \item Zapotecas, purépechas y tarascos
  \item Quechuas, aymaras y guaraníes
  \item Caribes, taínos y arahuacos
\end{multicols}
\end{enumerate}
\vspace{0.15cm}
\item ¿Cuál fue el principal señorío indígena en el actual El Salvador?
\begin{enumerate}[label=\textbf{\Alph*)}]
\begin{multicols}{2}
  \item Cuzcatlán
  \item Tenochtitlán
  \item Cusco
  \item Chichén Itzá
  \item Teotihuacán
  \item Monte Albán
\end{multicols}
\end{enumerate}
\vspace{0.15cm}
\item ¿Quiénes eran los pipiles?
\begin{enumerate}[label=\textbf{\Alph*)}]
\begin{multicols}{2}
  \item Un pueblo náhuat que habitó el centro y occidente de El Salvador
  \item Un pueblo maya del oriente
  \item Un pueblo lenca del norte
  \item Un pueblo chortí del sur
  \item Un pueblo azteca del valle de México
  \item Un pueblo inca del Perú
\end{multicols}
\end{enumerate}
\vspace{0.15cm}
\item ¿Quiénes eran los lencas?
\begin{enumerate}[label=\textbf{\Alph*)}]
\begin{multicols}{2}
  \item Un pueblo indígena que habitó el oriente de El Salvador
  \item Un pueblo náhuat del centro
  \item Un pueblo maya del occidente
  \item Un pueblo chortí del sur
  \item Un pueblo azteca del norte
  \item Un pueblo inca del sur
\end{multicols}
\end{enumerate}
\vspace{0.15cm}
\item ¿Quiénes eran los chortíes?
\begin{enumerate}[label=\textbf{\Alph*)}]
\begin{multicols}{2}
  \item Un pueblo maya que habitó el oriente de El Salvador
  \item Un pueblo náhuat del centro
  \item Un pueblo lenca del norte
  \item Un pueblo pipil del occidente
  \item Un pueblo azteca del sur
  \item Un pueblo inca del norte
\end{multicols}
\end{enumerate}
\vspace{0.15cm}
\item ¿Cuál era la lengua principal de los pipiles?
\begin{enumerate}[label=\textbf{\Alph*)}]
\begin{multicols}{2}
  \item Náhuat
  \item Lenca
  \item Chortí
  \item Maya
  \item Quechua
  \item Guaraní
\end{multicols}
\end{enumerate}
\vspace{0.15cm}
\item ¿Qué sitio arqueológico salvadoreño es Patrimonio de la Humanidad?
\begin{enumerate}[label=\textbf{\Alph*)}]
\begin{multicols}{2}
  \item Joya de Cerén
  \item Tazumal
  \item San Andrés
  \item Cihuatán
  \item Quelepa
  \item Casa Blanca
\end{multicols}
\end{enumerate}
\vspace{0.15cm}
\item ¿Qué es Joya de Cerén?
\begin{enumerate}[label=\textbf{\Alph*)}]
\begin{multicols}{2}
  \item Un sitio arqueológico maya conocido como la Pompeya de América
  \item Una pirámide azteca
  \item Un templo inca
  \item Una ciudad olmeca
  \item Un centro ceremonial tolteca
  \item Una fortaleza chichimeca
\end{multicols}
\end{enumerate}
\vspace{0.15cm}
\item ¿Quién dirigió la conquista española del actual El Salvador?
\begin{enumerate}[label=\textbf{\Alph*)}]
\begin{multicols}{2}
  \item Pedro de Alvarado
  \item Hernán Cortés
  \item Francisco Pizarro
  \item Diego de Almagro
  \item Pedro de Valdivia
  \item Sebastián de Belalcázar
\end{multicols}
\end{enumerate}
\vspace{0.15cm}
\item ¿En qué año llegó Pedro de Alvarado a Cuzcatlán?
\begin{enumerate}[label=\textbf{\Alph*)}]
\begin{multicols}{2}
  \item 1524
  \item 1492
  \item 1521
  \item 1533
  \item 1540
  \item 1519
\end{multicols}
\end{enumerate}
\vspace{0.15cm}
\item ¿Qué batalla enfrentó a los españoles con los pipiles en 1524?
\begin{enumerate}[label=\textbf{\Alph*)}]
\begin{multicols}{2}
  \item Batalla de Acajutla
  \item Batalla de Cuzcatlán
  \item Batalla de Tecún Umán
  \item Batalla de la Noche Triste
  \item Batalla de Otumba
  \item Batalla de Cajamarca
\end{multicols}
\end{enumerate}
\vspace{0.15cm}
\item ¿Quién fue Atlacatl?
\begin{enumerate}[label=\textbf{\Alph*)}]
\begin{multicols}{2}
  \item Un líder pipil que resistió a los españoles
  \item Un conquistador español
  \item Un rey maya
  \item Un sacerdote lenca
  \item Un cacique chortí
  \item Un gobernante azteca
\end{multicols}
\end{enumerate}
\vspace{0.15cm}
\item ¿Qué significó la conquista para los pueblos indígenas?
\begin{enumerate}[label=\textbf{\Alph*)}]
\begin{multicols}{2}
  \item La pérdida de tierras, cultura y población
  \item La mejora de sus condiciones de vida
  \item La preservación de su independencia
  \item La alianza con los españoles
  \item La expansión de su territorio
  \item El aumento de su población
\end{multicols}
\end{enumerate}
\vspace{0.15cm}
\item ¿Qué fue la encomienda?
\begin{enumerate}[label=\textbf{\Alph*)}]
\begin{multicols}{2}
  \item Una institución que otorgaba tierras e indígenas a los conquistadores
  \item Un sistema de comercio
  \item Una forma de gobierno municipal
  \item Una organización religiosa
  \item Un tipo de tributo
  \item Una escuela para indígenas
\end{multicols}
\end{enumerate}
\vspace{0.15cm}
\item ¿Qué fue la mita?
\begin{enumerate}[label=\textbf{\Alph*)}]
\begin{multicols}{2}
  \item Un sistema de trabajo forzado indígena
  \item Un sistema de comercio
  \item Una forma de gobierno
  \item Una organización religiosa
  \item Un tipo de tributo
  \item Una escuela
\end{multicols}
\end{enumerate}
\vspace{0.15cm}
\item ¿Qué fue el tributo colonial?
\begin{enumerate}[label=\textbf{\Alph*)}]
\begin{multicols}{2}
  \item Un pago obligatorio de los indígenas a los españoles
  \item Un impuesto sobre la tierra
  \item Un diezmo religioso
  \item Un arancel comercial
  \item Una contribución voluntaria
  \item Un préstamo
\end{multicols}
\end{enumerate}
\vspace{0.15cm}
\item ¿Qué fue la reducción de indígenas?
\begin{enumerate}[label=\textbf{\Alph*)}]
\begin{multicols}{2}
  \item La concentración de indígenas en pueblos bajo control español
  \item La eliminación de los indígenas
  \item La liberación de los indígenas
  \item La educación de los indígenas
  \item La evangelización de los indígenas
  \item La migración de los indígenas
\end{multicols}
\end{enumerate}
\vspace{0.15cm}
\item ¿Qué fue la evangelización?
\begin{enumerate}[label=\textbf{\Alph*)}]
\begin{multicols}{2}
  \item La conversión de los indígenas al cristianismo
  \item La educación de los indígenas
  \item La liberación de los indígenas
  \item La eliminación de los indígenas
  \item La migración de los indígenas
  \item La reducción de los indígenas
\end{multicols}
\end{enumerate}
\vspace{0.15cm}
\item ¿Qué orden religiosa tuvo un papel importante en la colonización?
\begin{enumerate}[label=\textbf{\Alph*)}]
\begin{multicols}{2}
  \item Los dominicos, franciscanos y jesuitas
  \item Los benedictinos, agustinos y carmelitas
  \item Los templarios, hospitalarios y teutónicos
  \item Los salesianos, capuchinos y paúles
  \item Los maristas, lasallistas y jesuitas
  \item Los claretianos, redentoristas y pasionistas
\end{multicols}
\end{enumerate}
\vspace{0.15cm}
\item ¿Qué fue la Capitanía General de Guatemala?
\begin{enumerate}[label=\textbf{\Alph*)}]
\begin{multicols}{2}
  \item La división administrativa que incluía a El Salvador
  \item Una provincia independiente
  \item Un virreinato
  \item Una audiencia
  \item Un obispado
  \item Un cabildo
\end{multicols}
\end{enumerate}
\vspace{0.15cm}
\item ¿Qué ciudad fue la capital de la Capitanía General de Guatemala?
\begin{enumerate}[label=\textbf{\Alph*)}]
\begin{multicols}{2}
  \item Santiago de Guatemala
  \item San Salvador
  \item Ciudad de Guatemala
  \item Tegucigalpa
  \item León
  \item Cartago
\end{multicols}
\end{enumerate}
\vspace{0.15cm}
\item ¿Qué fue la Audiencia de Guatemala?
\begin{enumerate}[label=\textbf{\Alph*)}]
\begin{multicols}{2}
  \item El tribunal de justicia de la Capitanía
  \item El gobierno municipal
  \item El cabildo
  \item El obispado
  \item La encomienda
  \item La mita
\end{multicols}
\end{enumerate}
\vspace{0.15cm}
\item ¿Qué producto fue importante en la economía colonial salvadoreña?
\begin{enumerate}[label=\textbf{\Alph*)}]
\begin{multicols}{2}
  \item El añil, el cacao y el bálsamo
  \item El café, el azúcar y el algodón
  \item El maíz, el frijol y el arroz
  \item El trigo, la cebada y el centeno
  \item La plata, el oro y el cobre
  \item El tabaco, el ron y el azúcar
\end{multicols}
\end{enumerate}
\vspace{0.15cm}
\item ¿Qué fue el añil?
\begin{enumerate}[label=\textbf{\Alph*)}]
\begin{multicols}{2}
  \item Un cultivo de tinte muy importante en la colonia
  \item Un mineral precioso
  \item Una bebida tradicional
  \item Un tipo de ganado
  \item Una herramienta agrícola
  \item Un impuesto colonial
\end{multicols}
\end{enumerate}
\vspace{0.15cm}
\item ¿Qué fue el bálsamo?
\begin{enumerate}[label=\textbf{\Alph*)}]
\begin{multicols}{2}
  \item Un producto medicinal y aromático de El Salvador
  \item Un mineral precioso
  \item Una bebida tradicional
  \item Un tipo de ganado
  \item Una herramienta agrícola
  \item Un impuesto colonial
\end{multicols}
\end{enumerate}
\vspace{0.15cm}
\item ¿Qué fue el cacao?
\begin{enumerate}[label=\textbf{\Alph*)}]
\begin{multicols}{2}
  \item Un cultivo importante para la economía colonial
  \item Un mineral precioso
  \item Una bebida tradicional
  \item Un tipo de ganado
  \item Una herramienta agrícola
  \item Un impuesto colonial
\end{multicols}
\end{enumerate}
\vspace{0.15cm}
\item ¿Qué fue el obraje?
\begin{enumerate}[label=\textbf{\Alph*)}]
\begin{multicols}{2}
  \item Un taller textil colonial
  \item Una mina
  \item Una hacienda
  \item Un trapiche
  \item Un mercado
  \item Una escuela
\end{multicols}
\end{enumerate}
\vspace{0.15cm}
\item ¿Qué fue el trapiche?
\begin{enumerate}[label=\textbf{\Alph*)}]
\begin{multicols}{2}
  \item Un molino para caña de azúcar
  \item Una mina
  \item Una hacienda
  \item Un obraje
  \item Un mercado
  \item Una escuela
\end{multicols}
\end{enumerate}
\vspace{0.15cm}
\item ¿Qué fue la hacienda?
\begin{enumerate}[label=\textbf{\Alph*)}]
\begin{multicols}{2}
  \item Una gran propiedad agrícola colonial
  \item Una mina
  \item Un obraje
  \item Un trapiche
  \item Un mercado
  \item Una escuela
\end{multicols}
\end{enumerate}
\vspace{0.15cm}
\item ¿Qué fue el cabildo?
\begin{enumerate}[label=\textbf{\Alph*)}]
\begin{multicols}{2}
  \item Una institución de gobierno municipal colonial
  \item Una mina
  \item Una hacienda
  \item Un obraje
  \item Un mercado
  \item Una escuela
\end{multicols}
\end{enumerate}
\vspace{0.15cm}
\item ¿Qué fue la cofradía?
\begin{enumerate}[label=\textbf{\Alph*)}]
\begin{multicols}{2}
  \item Una organización religiosa colonial
  \item Una mina
  \item Una hacienda
  \item Un obraje
  \item Un mercado
  \item Una escuela
\end{multicols}
\end{enumerate}
\vspace{0.15cm}
\item ¿Qué fue la posada?
\begin{enumerate}[label=\textbf{\Alph*)}]
\begin{multicols}{2}
  \item Una celebración religiosa colonial
  \item Una mina
  \item Una hacienda
  \item Un obraje
  \item Un mercado
  \item Una escuela
\end{multicols}
\end{enumerate}
\vspace{0.15cm}
\item ¿Qué fue el auto sacramental?
\begin{enumerate}[label=\textbf{\Alph*)}]
\begin{multicols}{2}
  \item Una obra teatral religiosa colonial
  \item Una mina
  \item Una hacienda
  \item Un obraje
  \item Un mercado
  \item Una escuela
\end{multicols}
\end{enumerate}
\vspace{0.15cm}
\item ¿Qué fue la loa?
\begin{enumerate}[label=\textbf{\Alph*)}]
\begin{multicols}{2}
  \item Una obra breve de alabanza colonial
  \item Una mina
  \item Una hacienda
  \item Un obraje
  \item Un mercado
  \item Una escuela
\end{multicols}
\end{enumerate}
\vspace{0.15cm}
\item ¿Qué fue el entremés?
\begin{enumerate}[label=\textbf{\Alph*)}]
\begin{multicols}{2}
  \item Una obra breve cómica colonial
  \item Una mina
  \item Una hacienda
  \item Un obraje
  \item Un mercado
  \item Una escuela
\end{multicols}
\end{enumerate}
\vspace{0.15cm}
\item ¿Qué fue la comedia?
\begin{enumerate}[label=\textbf{\Alph*)}]
\begin{multicols}{2}
  \item Una obra teatral profana colonial
  \item Una mina
  \item Una hacienda
  \item Un obraje
  \item Un mercado
  \item Una escuela
\end{multicols}
\end{enumerate}
\vspace{0.15cm}
\item ¿Qué fue la tragedia?
\begin{enumerate}[label=\textbf{\Alph*)}]
\begin{multicols}{2}
  \item Una obra teatral seria colonial
  \item Una mina
  \item Una hacienda
  \item Un obraje
  \item Un mercado
  \item Una escuela
\end{multicols}
\end{enumerate}
\vspace{0.15cm}
\item ¿Qué fue el coloquio?
\begin{enumerate}[label=\textbf{\Alph*)}]
\begin{multicols}{2}
  \item Una obra teatral dialogada colonial
  \item Una mina
  \item Una hacienda
  \item Un obraje
  \item Un mercado
  \item Una escuela
\end{multicols}
\end{enumerate}
\vspace{0.15cm}
\item ¿Qué fue la pastorela?
\begin{enumerate}[label=\textbf{\Alph*)}]
\begin{multicols}{2}
  \item Una obra teatral navideña colonial
  \item Una mina
  \item Una hacienda
  \item Un obraje
  \item Un mercado
  \item Una escuela
\end{multicols}
\end{enumerate}
\vspace{0.15cm}
\item ¿Qué fue la reducción?
\begin{enumerate}[label=\textbf{\Alph*)}]
\begin{multicols}{2}
  \item La concentración de indígenas en pueblos coloniales
  \item Una mina
  \item Una hacienda
  \item Un obraje
  \item Un mercado
  \item Una escuela
\end{multicols}
\end{enumerate}
\vspace{0.15cm}
\end{enumerate}\n
\subsubsection*{Sección II: Independencia y Federación (41–80)}

\begin{enumerate}[label=\textbf{\arabic*.}]
\setcounter{enumi}{40}
\item ¿En qué año se proclamó la independencia de Centroamérica?
\begin{enumerate}[label=\textbf{\Alph*)}]
\begin{multicols}{2}
  \item 1821
  \item 1810
  \item 1811
  \item 1823
  \item 1838
  \item 1840
\end{multicols}
\end{enumerate}
\vspace{0.15cm}
\item ¿Quién fue el prócer salvadoreño que firmó el Acta de Independencia?
\begin{enumerate}[label=\textbf{\Alph*)}]
\begin{multicols}{2}
  \item José Matías Delgado
  \item Manuel José Arce
  \item Pedro Barriere
  \item Juan Manuel Rodríguez
  \item Mariano Prado
  \item José María Cornejo
\end{multicols}
\end{enumerate}
\vspace{0.15cm}
\item ¿Quién es considerado el "Padre de la Patria" en El Salvador?
\begin{enumerate}[label=\textbf{\Alph*)}]
\begin{multicols}{2}
  \item José Matías Delgado
  \item Manuel José Arce
  \item Pedro Barriere
  \item Juan Manuel Rodríguez
  \item Mariano Prado
  \item José María Cornejo
\end{multicols}
\end{enumerate}
\vspace{0.15cm}
\item ¿Qué fue el Primer Grito de Independencia en El Salvador?
\begin{enumerate}[label=\textbf{\Alph*)}]
\begin{multicols}{2}
  \item El movimiento de 1811
  \item El movimiento de 1821
  \item El movimiento de 1823
  \item El movimiento de 1838
  \item El movimiento de 1840
  \item El movimiento de 1810
\end{multicols}
\end{enumerate}
\vspace{0.15cm}
\item ¿Quién lideró el Primer Grito de Independencia de 1811?
\begin{enumerate}[label=\textbf{\Alph*)}]
\begin{multicols}{2}
  \item José Matías Delgado
  \item Manuel José Arce
  \item Pedro Barriere
  \item Juan Manuel Rodríguez
  \item Mariano Prado
  \item José María Cornejo
\end{multicols}
\end{enumerate}
\vspace{0.15cm}
\item ¿Qué fue la anexión a México?
\begin{enumerate}[label=\textbf{\Alph*)}]
\begin{multicols}{2}
  \item La incorporación de Centroamérica al Imperio Mexicano de Iturbide
  \item La independencia de Centroamérica
  \item La formación de la Federación
  \item La separación de El Salvador
  \item La invasión española
  \item La intervención inglesa
\end{multicols}
\end{enumerate}
\vspace{0.15cm}
\item ¿Quién se opuso a la anexión a México?
\begin{enumerate}[label=\textbf{\Alph*)}]
\begin{multicols}{2}
  \item José Matías Delgado
  \item Manuel José Arce
  \item Pedro Barriere
  \item Juan Manuel Rodríguez
  \item Mariano Prado
  \item José María Cornejo
\end{multicols}
\end{enumerate}
\vspace{0.15cm}
\item ¿En qué año se formó la Federación Centroamericana?
\begin{enumerate}[label=\textbf{\Alph*)}]
\begin{multicols}{2}
  \item 1823
  \item 1821
  \item 1811
  \item 1838
  \item 1840
  \item 1825
\end{multicols}
\end{enumerate}
\vspace{0.15cm}
\item ¿Quién fue el primer presidente de la Federación Centroamericana?
\begin{enumerate}[label=\textbf{\Alph*)}]
\begin{multicols}{2}
  \item Manuel José Arce
  \item José Matías Delgado
  \item Pedro Barriere
  \item Juan Manuel Rodríguez
  \item Mariano Prado
  \item José María Cornejo
\end{multicols}
\end{enumerate}
\vspace{0.15cm}
\item ¿Qué fue la Federación Centroamericana?
\begin{enumerate}[label=\textbf{\Alph*)}]
\begin{multicols}{2}
  \item Una unión política de los países centroamericanos
  \item Una colonia española
  \item Un imperio
  \item Una república independiente
  \item Una monarquía
  \item Una dictadura
\end{multicols}
\end{enumerate}
\vspace{0.15cm}
\item ¿Quién fue Francisco Morazán?
\begin{enumerate}[label=\textbf{\Alph*)}]
\begin{multicols}{2}
  \item Un líder liberal hondureño que presidió la Federación
  \item Un conservador salvadoreño
  \item Un rey español
  \item Un presidente mexicano
  \item Un caudillo guatemalteco
  \item Un obispo costarricense
\end{multicols}
\end{enumerate}
\vspace{0.15cm}
\item ¿Qué ideas defendía Francisco Morazán?
\begin{enumerate}[label=\textbf{\Alph*)}]
\begin{multicols}{2}
  \item El liberalismo y la unión centroamericana
  \item El conservadurismo y la fragmentación
  \item La monarquía y la tradición
  \item El colonialismo y la dependencia
  \item El socialismo y la revolución
  \item El anarquismo y la libertad
\end{multicols}
\end{enumerate}
\vspace{0.15cm}
\item ¿Qué fue la guerra entre liberales y conservadores?
\begin{enumerate}[label=\textbf{\Alph*)}]
\begin{multicols}{2}
  \item Un conflicto por el control político de Centroamérica
  \item Una guerra contra España
  \item Una guerra contra México
  \item Una guerra contra Inglaterra
  \item Una guerra contra Estados Unidos
  \item Una guerra civil en El Salvador
\end{multicols}
\end{enumerate}
\vspace{0.15cm}
\item ¿En qué año se disolvió la Federación Centroamericana?
\begin{enumerate}[label=\textbf{\Alph*)}]
\begin{multicols}{2}
  \item 1838
  \item 1823
  \item 1821
  \item 1840
  \item 1825
  \item 1850
\end{multicols}
\end{enumerate}
\vspace{0.15cm}
\item ¿Quién fue el último presidente de la Federación?
\begin{enumerate}[label=\textbf{\Alph*)}]
\begin{multicols}{2}
  \item Francisco Morazán
  \item Manuel José Arce
  \item José Matías Delgado
  \item Pedro Barriere
  \item Juan Manuel Rodríguez
  \item Mariano Prado
\end{multicols}
\end{enumerate}
\vspace{0.15cm}
\item ¿Qué fue la República de El Salvador?
\begin{enumerate}[label=\textbf{\Alph*)}]
\begin{multicols}{2}
  \item El Estado independiente que surgió tras la disolución de la Federación
  \item Una colonia española
  \item Un imperio
  \item Una federación
  \item Una monarquía
  \item Una dictadura
\end{multicols}
\end{enumerate}
\vspace{0.15cm}
\item ¿En qué año se declaró El Salvador como república independiente?
\begin{enumerate}[label=\textbf{\Alph*)}]
\begin{multicols}{2}
  \item 1841
  \item 1821
  \item 1823
  \item 1838
  \item 1840
  \item 1850
\end{multicols}
\end{enumerate}
\vspace{0.15cm}
\item ¿Quién fue el primer presidente de El Salvador como república independiente?
\begin{enumerate}[label=\textbf{\Alph*)}]
\begin{multicols}{2}
  \item Juan Lindo
  \item Manuel José Arce
  \item José Matías Delgado
  \item Francisco Morazán
  \item Juan Manuel Rodríguez
  \item Mariano Prado
\end{multicols}
\end{enumerate}
\vspace{0.15cm}
\item ¿Qué fue la Constitución de 1841?
\begin{enumerate}[label=\textbf{\Alph*)}]
\begin{multicols}{2}
  \item La primera constitución de El Salvador como república
  \item La constitución de la Federación
  \item La constitución colonial
  \item La constitución de 1824
  \item La constitución de 1886
  \item La constitución de 1950
\end{multicols}
\end{enumerate}
\vspace{0.15cm}
\item ¿Qué fue el caudillismo?
\begin{enumerate}[label=\textbf{\Alph*)}]
\begin{multicols}{2}
  \item El gobierno de caudillos militares en el siglo XIX
  \item El gobierno democrático
  \item El gobierno colonial
  \item El gobierno liberal
  \item El gobierno conservador
  \item El gobierno federal
\end{multicols}
\end{enumerate}
\vspace{0.15cm}
\item ¿Quién fue Gerardo Barrios?
\begin{enumerate}[label=\textbf{\Alph*)}]
\begin{multicols}{2}
  \item Un presidente salvadoreño que impulsó el liberalismo
  \item Un conservador salvadoreño
  \item Un rey español
  \item Un presidente mexicano
  \item Un caudillo guatemalteco
  \item Un obispo costarricense
\end{multicols}
\end{enumerate}
\vspace{0.15cm}
\item ¿Qué impulsó Gerardo Barrios?
\begin{enumerate}[label=\textbf{\Alph*)}]
\begin{multicols}{2}
  \item La educación, el comercio y la modernización
  \item La tradición y el conservadurismo
  \item La monarquía y la dependencia
  \item El colonialismo y la opresión
  \item El socialismo y la revolución
  \item El anarquismo y la libertad
\end{multicols}
\end{enumerate}
\vspace{0.15cm}
\item ¿Quién fue Francisco Dueñas?
\begin{enumerate}[label=\textbf{\Alph*)}]
\begin{multicols}{2}
  \item Un presidente conservador salvadoreño
  \item Un liberal salvadoreño
  \item Un rey español
  \item Un presidente mexicano
  \item Un caudillo guatemalteco
  \item Un obispo costarricense
\end{multicols}
\end{enumerate}
\vspace{0.15cm}
\item ¿Qué defendía Francisco Dueñas?
\begin{enumerate}[label=\textbf{\Alph*)}]
\begin{multicols}{2}
  \item El conservadurismo y la tradición
  \item El liberalismo y la modernización
  \item La monarquía y la dependencia
  \item El colonialismo y la opresión
  \item El socialismo y la revolución
  \item El anarquismo y la libertad
\end{multicols}
\end{enumerate}
\vspace{0.15cm}
\item ¿Qué fue la Revolución Liberal de 1871?
\begin{enumerate}[label=\textbf{\Alph*)}]
\begin{multicols}{2}
  \item Un movimiento que impulsó la modernización y el liberalismo
  \item Un movimiento conservador
  \item Una invasión extranjera
  \item Una guerra civil
  \item Una revolución socialista
  \item Una revolución comunista
\end{multicols}
\end{enumerate}
\vspace{0.15cm}
\item ¿Quién lideró la Revolución Liberal de 1871?
\begin{enumerate}[label=\textbf{\Alph*)}]
\begin{multicols}{2}
  \item Santiago González
  \item Gerardo Barrios
  \item Francisco Dueñas
  \item Manuel José Arce
  \item José Matías Delgado
  \item Francisco Morazán
\end{multicols}
\end{enumerate}
\vspace{0.15cm}
\item ¿Qué fue la Constitución de 1886?
\begin{enumerate}[label=\textbf{\Alph*)}]
\begin{multicols}{2}
  \item Una constitución liberal que rigió por mucho tiempo
  \item Una constitución conservadora
  \item Una constitución colonial
  \item Una constitución federal
  \item Una constitución socialista
  \item Una constitución comunista
\end{multicols}
\end{enumerate}
\vspace{0.15cm}
\item ¿Qué estableció la Constitución de 1886?
\begin{enumerate}[label=\textbf{\Alph*)}]
\begin{multicols}{2}
  \item Garantías individuales, educación laica y propiedad privada
  \item El poder absoluto del presidente
  \item La monarquía constitucional
  \item La dependencia colonial
  \item El socialismo de estado
  \item La propiedad comunal
\end{multicols}
\end{enumerate}
\vspace{0.15cm}
\item ¿Qué fue el café?
\begin{enumerate}[label=\textbf{\Alph*)}]
\begin{multicols}{2}
  \item El principal producto de exportación desde finales del siglo XIX
  \item Un producto secundario
  \item Un producto colonial
  \item Un producto industrial
  \item Un producto minero
  \item Un producto artesanal
\end{multicols}
\end{enumerate}
\vspace{0.15cm}
\item ¿Qué fue la "república cafetalera"?
\begin{enumerate}[label=\textbf{\Alph*)}]
\begin{multicols}{2}
  \item El período en que el café dominó la economía y la política
  \item El período colonial
  \item El período de la Federación
  \item El período de la independencia
  \item El período de la guerra civil
  \item El período post-guerra
\end{multicols}
\end{enumerate}
\vspace{0.15cm}
\item ¿Qué fue la oligarquía cafetalera?
\begin{enumerate}[label=\textbf{\Alph*)}]
\begin{multicols}{2}
  \item El grupo de familias que controlaban la tierra y el poder
  \item El grupo de campesinos
  \item El grupo de obreros
  \item El grupo de comerciantes
  \item El grupo de militares
  \item El grupo de intelectuales
\end{multicols}
\end{enumerate}
\vspace{0.15cm}
\item ¿Qué fue la ley de extinción de ejidos?
\begin{enumerate}[label=\textbf{\Alph*)}]
\begin{multicols}{2}
  \item Una ley que privatizó las tierras comunales
  \item Una ley que protegió las tierras comunales
  \item Una ley que distribuyó la tierra
  \item Una ley que expropió la tierra
  \item Una ley que nacionalizó la tierra
  \item Una ley que colectivizó la tierra
\end{multicols}
\end{enumerate}
\vspace{0.15cm}
\item ¿En qué año se aprobó la ley de extinción de ejidos?
\begin{enumerate}[label=\textbf{\Alph*)}]
\begin{multicols}{2}
  \item 1882
  \item 1871
  \item 1886
  \item 1890
  \item 1900
  \item 1850
\end{multicols}
\end{enumerate}
\vspace{0.15cm}
\item ¿Qué consecuencia tuvo la ley de extinción de ejidos?
\begin{enumerate}[label=\textbf{\Alph*)}]
\begin{multicols}{2}
  \item La concentración de la tierra en pocas manos
  \item La distribución equitativa de la tierra
  \item La desaparición del café
  \item El aumento de la producción agrícola
  \item La mejora de las condiciones campesinas
  \item La reducción de la pobreza
\end{multicols}
\end{enumerate}
\vspace{0.15cm}
\item ¿Qué fue el comunismo?
\begin{enumerate}[label=\textbf{\Alph*)}]
\begin{multicols}{2}
  \item Un sistema de propiedad comunal de la tierra
  \item Un sistema de propiedad privada
  \item Un sistema de propiedad estatal
  \item Un sistema de propiedad feudal
  \item Un sistema de propiedad capitalista
  \item Un sistema de propiedad mixta
\end{multicols}
\end{enumerate}
\vspace{0.15cm}
\item ¿Qué fue la matanza de 1932?
\begin{enumerate}[label=\textbf{\Alph*)}]
\begin{multicols}{2}
  \item La represión del levantamiento campesino e indígena
  \item Una masacre colonial
  \item Una guerra civil
  \item Una invasión extranjera
  \item Un terremoto
  \item Una epidemia
\end{multicols}
\end{enumerate}
\vspace{0.15cm}
\item ¿Quién ordenó la matanza de 1932?
\begin{enumerate}[label=\textbf{\Alph*)}]
\begin{multicols}{2}
  \item Maximiliano Hernández Martínez
  \item Gerardo Barrios
  \item Francisco Dueñas
  \item Manuel José Arce
  \item José Matías Delgado
  \item Francisco Morazán
\end{multicols}
\end{enumerate}
\vspace{0.15cm}
\item ¿Qué fue el levantamiento campesino de 1932?
\begin{enumerate}[label=\textbf{\Alph*)}]
\begin{multicols}{2}
  \item Una rebelión contra las condiciones de explotación
  \item Una rebelión colonial
  \item Una rebelión conservadora
  \item Una rebelión liberal
  \item Una rebelión militar
  \item Una rebelión estudiantil
\end{multicols}
\end{enumerate}
\vspace{0.15cm}
\item ¿Quién fue Agustín Farabundo Martí?
\begin{enumerate}[label=\textbf{\Alph*)}]
\begin{multicols}{2}
  \item Un líder comunista salvadoreño
  \item Un presidente conservador
  \item Un militar liberal
  \item Un obispo salvadoreño
  \item Un caudillo guatemalteco
  \item Un rey español
\end{multicols}
\end{enumerate}
\vspace{0.15cm}
\item ¿Qué pasó con Farabundo Martí?
\begin{enumerate}[label=\textbf{\Alph*)}]
\begin{multicols}{2}
  \item Fue capturado y fusilado tras el levantamiento de 1932
  \item Fue presidente de El Salvador
  \item Fue exiliado
  \item Fue encarcelado y liberado
  \item Fue nombrado obispo
  \item Fue nombrado general
\end{multicols}
\end{enumerate}
\vspace{0.15cm}
\end{enumerate}\n
\subsubsection*{Sección III: Regímenes Militares y Conflicto (81–140)}

\begin{enumerate}[label=\textbf{\arabic*.}]
\setcounter{enumi}{80}
\item ¿Quién fue Maximiliano Hernández Martínez?
\begin{enumerate}[label=\textbf{\Alph*)}]
\begin{multicols}{2}
  \item Un dictador militar que gobernó El Salvador
  \item Un presidente democrático
  \item Un líder comunista
  \item Un obispo salvadoreño
  \item Un caudillo guatemalteco
  \item Un rey español
\end{multicols}
\end{enumerate}
\vspace{0.15cm}
\item ¿En qué año llegó al poder Maximiliano Hernández Martínez?
\begin{enumerate}[label=\textbf{\Alph*)}]
\begin{multicols}{2}
  \item 1931
  \item 1932
  \item 1944
  \item 1950
  \item 1960
  \item 1920
\end{multicols}
\end{enumerate}
\vspace{0.15cm}
\item ¿Qué caracterizó el gobierno de Hernández Martínez?
\begin{enumerate}[label=\textbf{\Alph*)}]
\begin{multicols}{2}
  \item La represión, el autoritarismo y el anticomunismo
  \item La democracia y la libertad
  \item El socialismo y la revolución
  \item El liberalismo y la modernización
  \item El conservadurismo y la tradición
  \item El anarquismo y la libertad
\end{multicols}
\end{enumerate}
\vspace{0.15cm}
\item ¿Qué fue la huelga de brazos caídos de 1944?
\begin{enumerate}[label=\textbf{\Alph*)}]
\begin{multicols}{2}
  \item Una protesta que llevó a la caída de Hernández Martínez
  \item Una protesta colonial
  \item Una protesta conservadora
  \item Una protesta liberal
  \item Una protesta militar
  \item Una protesta estudiantil
\end{multicols}
\end{enumerate}
\vspace{0.15cm}
\item ¿Quién sucedió a Hernández Martínez?
\begin{enumerate}[label=\textbf{\Alph*)}]
\begin{multicols}{2}
  \item Osmín Aguirre y Salinas
  \item Salvador Castaneda Castro
  \item Óscar Osorio
  \item José María Lemus
  \item Julio Adalberto Rivera
  \item Fidel Sánchez Hernández
\end{multicols}
\end{enumerate}
\vspace{0.15cm}
\item ¿Qué fue el golpe de estado de 1948?
\begin{enumerate}[label=\textbf{\Alph*)}]
\begin{multicols}{2}
  \item Un golpe militar que llevó al poder a Óscar Osorio
  \item Un golpe colonial
  \item Un golpe conservador
  \item Un golpe liberal
  \item Un golpe comunista
  \item Un golpe estudiantil
\end{multicols}
\end{enumerate}
\vspace{0.15cm}
\item ¿Quién fue Óscar Osorio?
\begin{enumerate}[label=\textbf{\Alph*)}]
\begin{multicols}{2}
  \item Un presidente militar que impulsó reformas sociales
  \item Un dictador conservador
  \item Un líder comunista
  \item Un obispo salvadoreño
  \item Un caudillo guatemalteco
  \item Un rey español
\end{multicols}
\end{enumerate}
\vspace{0.15cm}
\item ¿Qué impulsó Óscar Osorio?
\begin{enumerate}[label=\textbf{\Alph*)}]
\begin{multicols}{2}
  \item La seguridad social, la educación y la infraestructura
  \item La represión y el autoritarismo
  \item El socialismo y la revolución
  \item El liberalismo y la modernización
  \item El conservadurismo y la tradición
  \item El anarquismo y la libertad
\end{multicols}
\end{enumerate}
\vspace{0.15cm}
\item ¿Quién fue José María Lemus?
\begin{enumerate}[label=\textbf{\Alph*)}]
\begin{multicols}{2}
  \item Un presidente militar que fue derrocado
  \item Un dictador conservador
  \item Un líder comunista
  \item Un obispo salvadoreño
  \item Un caudillo guatemalteco
  \item Un rey español
\end{multicols}
\end{enumerate}
\vspace{0.15cm}
\item ¿Qué fue el golpe de estado de 1960?
\begin{enumerate}[label=\textbf{\Alph*)}]
\begin{multicols}{2}
  \item Un golpe militar que derrocó a Lemus
  \item Un golpe colonial
  \item Un golpe conservador
  \item Un golpe liberal
  \item Un golpe comunista
  \item Un golpe estudiantil
\end{multicols}
\end{enumerate}
\vspace{0.15cm}
\item ¿Quién fue Julio Adalberto Rivera?
\begin{enumerate}[label=\textbf{\Alph*)}]
\begin{multicols}{2}
  \item Un presidente militar que impulsó la Alianza para el Progreso
  \item Un dictador conservador
  \item Un líder comunista
  \item Un obispo salvadoreño
  \item Un caudillo guatemalteco
  \item Un rey español
\end{multicols}
\end{enumerate}
\vspace{0.15cm}
\item ¿Qué fue la Alianza para el Progreso?
\begin{enumerate}[label=\textbf{\Alph*)}]
\begin{multicols}{2}
  \item Un programa de ayuda de Estados Unidos para América Latina
  \item Un programa de ayuda europea
  \item Un programa de ayuda soviética
  \item Un programa de ayuda china
  \item Un programa de ayuda japonesa
  \item Un programa de ayuda africana
\end{multicols}
\end{enumerate}
\vspace{0.15cm}
\item ¿Quién fue Fidel Sánchez Hernández?
\begin{enumerate}[label=\textbf{\Alph*)}]
\begin{multicols}{2}
  \item Un presidente militar que enfrentó la guerra con Honduras
  \item Un dictador conservador
  \item Un líder comunista
  \item Un obispo salvadoreño
  \item Un caudillo guatemalteco
  \item Un rey español
\end{multicols}
\end{enumerate}
\vspace{0.15cm}
\item ¿Qué fue la guerra de las 100 horas?
\begin{enumerate}[label=\textbf{\Alph*)}]
\begin{multicols}{2}
  \item El conflicto entre El Salvador y Honduras en 1969
  \item Una guerra civil
  \item Una invasión extranjera
  \item Una guerra colonial
  \item Una guerra mundial
  \item Una guerra fría
\end{multicols}
\end{enumerate}
\vspace{0.15cm}
\item ¿Cuál fue la causa de la guerra de las 100 horas?
\begin{enumerate}[label=\textbf{\Alph*)}]
\begin{multicols}{2}
  \item La migración salvadoreña a Honduras y la reforma agraria hondureña
  \item Una disputa territorial
  \item Una disputa comercial
  \item Una disputa política
  \item Una disputa religiosa
  \item Una disputa étnica
\end{multicols}
\end{enumerate}
\vspace{0.15cm}
\item ¿Qué fue la Organización Democrática Nacionalista (ORDEN)?
\begin{enumerate}[label=\textbf{\Alph*)}]
\begin{multicols}{2}
  \item Un grupo paramilitar de extrema derecha
  \item Un grupo guerrillero
  \item Un partido político
  \item Un sindicato
  \item Una organización religiosa
  \item Una organización estudiantil
\end{multicols}
\end{enumerate}
\vspace{0.15cm}
\item ¿Qué fue el Frente Farabundo Martí para la Liberación Nacional (FMLN)?
\begin{enumerate}[label=\textbf{\Alph*)}]
\begin{multicols}{2}
  \item Una organización guerrillera de izquierda
  \item Un partido político de derecha
  \item Un grupo paramilitar
  \item Un sindicato
  \item Una organización religiosa
  \item Una organización estudiantil
\end{multicols}
\end{enumerate}
\vspace{0.15cm}
\item ¿En qué año se fundó el FMLN?
\begin{enumerate}[label=\textbf{\Alph*)}]
\begin{multicols}{2}
  \item 1980
  \item 1970
  \item 1975
  \item 1985
  \item 1990
  \item 1992
\end{multicols}
\end{enumerate}
\vspace{0.15cm}
\item ¿Qué fue la guerra civil salvadoreña?
\begin{enumerate}[label=\textbf{\Alph*)}]
\begin{multicols}{2}
  \item Un conflicto armado entre el gobierno y la guerrilla
  \item Una guerra colonial
  \item Una guerra internacional
  \item Una guerra civil corta
  \item Una guerra mundial
  \item Una guerra fría
\end{multicols}
\end{enumerate}
\vspace{0.15cm}
\item ¿En qué año comenzó la guerra civil salvadoreña?
\begin{enumerate}[label=\textbf{\Alph*)}]
\begin{multicols}{2}
  \item 1980
  \item 1970
  \item 1975
  \item 1985
  \item 1990
  \item 1992
\end{multicols}
\end{enumerate}
\vspace{0.15cm}
\item ¿Qué fue el asesinato de Monseñor Óscar Arnulfo Romero?
\begin{enumerate}[label=\textbf{\Alph*)}]
\begin{multicols}{2}
  \item El asesinato de un obispo defensor de los derechos humanos
  \item El asesinato de un presidente
  \item El asesinato de un militar
  \item El asesinato de un guerrillero
  \item El asesinato de un empresario
  \item El asesinato de un estudiante
\end{multicols}
\end{enumerate}
\vspace{0.15cm}
\item ¿En qué año fue asesinado Monseñor Romero?
\begin{enumerate}[label=\textbf{\Alph*)}]
\begin{multicols}{2}
  \item 1980
  \item 1970
  \item 1975
  \item 1985
  \item 1990
  \item 1992
\end{multicols}
\end{enumerate}
\vspace{0.15cm}
\item ¿Qué fue la masacre de El Mozote?
\begin{enumerate}[label=\textbf{\Alph*)}]
\begin{multicols}{2}
  \item Una masacre de campesinos por el ejército en 1981
  \item Una masacre colonial
  \item Una masacre guerrillera
  \item Una masacre militar
  \item Una masacre estudiantil
  \item Una masacre religiosa
\end{multicols}
\end{enumerate}
\vspace{0.15cm}
\item ¿Qué fue la masacre de los Jesuitas?
\begin{enumerate}[label=\textbf{\Alph*)}]
\begin{multicols}{2}
  \item El asesinato de seis sacerdotes jesuitas en 1989
  \item Una masacre colonial
  \item Una masacre guerrillera
  \item Una masacre militar
  \item Una masacre estudiantil
  \item Una masacre religiosa
\end{multicols}
\end{enumerate}
\vspace{0.15cm}
\item ¿Quiénes fueron asesinados en la masacre de los Jesuitas?
\begin{enumerate}[label=\textbf{\Alph*)}]
\begin{multicols}{2}
  \item Ignacio Ellacuría, Ignacio Martín-Baró y otros
  \item Monseñor Romero
  \item Farabundo Martí
  \item Maximiliano Hernández Martínez
  \item Gerardo Barrios
  \item Francisco Dueñas
\end{multicols}
\end{enumerate}
\vspace{0.15cm}
\item ¿Qué fue la doctrina de seguridad nacional?
\begin{enumerate}[label=\textbf{\Alph*)}]
\begin{multicols}{2}
  \item Una doctrina anticomunista que justificó la represión
  \item Una doctrina democrática
  \item Una doctrina socialista
  \item Una doctrina liberal
  \item Una doctrina conservadora
  \item Una doctrina anarquista
\end{multicols}
\end{enumerate}
\vspace{0.15cm}
\item ¿Qué fue el batallón Atlacatl?
\begin{enumerate}[label=\textbf{\Alph*)}]
\begin{multicols}{2}
  \item Un batallón de élite del ejército salvadoreño
  \item Un batallón guerrillero
  \item Un batallón colonial
  \item Un batallón conservador
  \item Un batallón liberal
  \item Un batallón estudiantil
\end{multicols}
\end{enumerate}
\vspace{0.15cm}
\item ¿Qué fue la Fuerza Armada de El Salvador?
\begin{enumerate}[label=\textbf{\Alph*)}]
\begin{multicols}{2}
  \item La institución militar que combatió a la guerrilla
  \item Una organización guerrillera
  \item Un partido político
  \item Un sindicato
  \item Una organización religiosa
  \item Una organización estudiantil
\end{multicols}
\end{enumerate}
\vspace{0.15cm}
\item ¿Qué fue el Ejército Revolucionario del Pueblo (ERP)?
\begin{enumerate}[label=\textbf{\Alph*)}]
\begin{multicols}{2}
  \item Una organización guerrillera salvadoreña
  \item Un partido político
  \item Un grupo paramilitar
  \item Un sindicato
  \item Una organización religiosa
  \item Una organización estudiantil
\end{multicols}
\end{enumerate}
\vspace{0.15cm}
\item ¿Qué fue las Fuerzas Populares de Liberación (FPL)?
\begin{enumerate}[label=\textbf{\Alph*)}]
\begin{multicols}{2}
  \item Una organización guerrillera salvadoreña
  \item Un partido político
  \item Un grupo paramilitar
  \item Un sindicato
  \item Una organización religiosa
  \item Una organización estudiantil
\end{multicols}
\end{enumerate}
\vspace{0.15cm}
\item ¿Qué fue la Resistencia Nacional (RN)?
\begin{enumerate}[label=\textbf{\Alph*)}]
\begin{multicols}{2}
  \item Una organización guerrillera salvadoreña
  \item Un partido político
  \item Un grupo paramilitar
  \item Un sindicato
  \item Una organización religiosa
  \item Una organización estudiantil
\end{multicols}
\end{enumerate}
\vspace{0.15cm}
\item ¿Qué fue el Partido Comunista de El Salvador (PCS)?
\begin{enumerate}[label=\textbf{\Alph*)}]
\begin{multicols}{2}
  \item Una organización política de izquierda
  \item Un partido político de derecha
  \item Un grupo paramilitar
  \item Un sindicato
  \item Una organización religiosa
  \item Una organización estudiantil
\end{multicols}
\end{enumerate}
\vspace{0.15cm}
\item ¿Qué fue el Partido Demócrata Cristiano (PDC)?
\begin{enumerate}[label=\textbf{\Alph*)}]
\begin{multicols}{2}
  \item Un partido político de centro
  \item Un partido político de izquierda
  \item Un partido político de derecha
  \item Un grupo paramilitar
  \item Un sindicato
  \item Una organización religiosa
\end{multicols}
\end{enumerate}
\vspace{0.15cm}
\item ¿Quién fue José Napoleón Duarte?
\begin{enumerate}[label=\textbf{\Alph*)}]
\begin{multicols}{2}
  \item Un líder del PDC y presidente de El Salvador
  \item Un líder guerrillero
  \item Un militar
  \item Un obispo
  \item Un empresario
  \item Un estudiante
\end{multicols}
\end{enumerate}
\vspace{0.15cm}
\item ¿Qué fue la Alianza Republicana Nacionalista (ARENA)?
\begin{enumerate}[label=\textbf{\Alph*)}]
\begin{multicols}{2}
  \item Un partido político de derecha
  \item Un partido político de izquierda
  \item Un partido político de centro
  \item Un grupo paramilitar
  \item Un sindicato
  \item Una organización religiosa
\end{multicols}
\end{enumerate}
\vspace{0.15cm}
\item ¿Quién fundó ARENA?
\begin{enumerate}[label=\textbf{\Alph*)}]
\begin{multicols}{2}
  \item Roberto D'Aubuisson
  \item José Napoleón Duarte
  \item Alfredo Cristiani
  \item Armando Calderón Sol
  \item Francisco Flores
  \item Antonio Saca
\end{multicols}
\end{enumerate}
\vspace{0.15cm}
\item ¿Qué fue el gobierno de Alfredo Cristiani?
\begin{enumerate}[label=\textbf{\Alph*)}]
\begin{multicols}{2}
  \item El gobierno que firmó los Acuerdos de Paz
  \item El gobierno que inició la guerra
  \item El gobierno que continuó la guerra
  \item El gobierno que reprimió la protesta
  \item El gobierno que impulsó la reforma agraria
  \item El gobierno que nacionalizó la banca
\end{multicols}
\end{enumerate}
\vspace{0.15cm}
\item ¿En qué año se firmaron los Acuerdos de Paz?
\begin{enumerate}[label=\textbf{\Alph*)}]
\begin{multicols}{2}
  \item 1992
  \item 1980
  \item 1985
  \item 1990
  \item 1995
  \item 2000
\end{multicols}
\end{enumerate}
\vspace{0.15cm}
\item ¿Dónde se firmaron los Acuerdos de Paz?
\begin{enumerate}[label=\textbf{\Alph*)}]
\begin{multicols}{2}
  \item Chapultepec, México
  \item San Salvador, El Salvador
  \item Guatemala, Guatemala
  \item Tegucigalpa, Honduras
  \item Managua, Nicaragua
  \item San José, Costa Rica
\end{multicols}
\end{enumerate}
\vspace{0.15cm}
\item ¿Qué fue la Comisión de la Verdad?
\begin{enumerate}[label=\textbf{\Alph*)}]
\begin{multicols}{2}
  \item Una comisión que investigó las violaciones a los derechos humanos
  \item Una comisión guerrillera
  \item Una comisión militar
  \item Una comisión política
  \item Una comisión religiosa
  \item Una comisión estudiantil
\end{multicols}
\end{enumerate}
\vspace{0.15cm}
\item ¿Quién fue Maximiliano Hernández Martínez?
\begin{enumerate}[label=\textbf{\Alph*)}]
\begin{multicols}{2}
  \item Un dictador militar que gobernó El Salvador
  \item Un presidente democrático
  \item Un líder comunista
  \item Un obispo salvadoreño
  \item Un caudillo guatemalteco
  \item Un rey español
\end{multicols}
\end{enumerate}
\vspace{0.15cm}
\item ¿En qué año llegó al poder Maximiliano Hernández Martínez?
\begin{enumerate}[label=\textbf{\Alph*)}]
\begin{multicols}{2}
  \item 1931
  \item 1932
  \item 1944
  \item 1950
  \item 1960
  \item 1920
\end{multicols}
\end{enumerate}
\vspace{0.15cm}
\item ¿Qué caracterizó el gobierno de Hernández Martínez?
\begin{enumerate}[label=\textbf{\Alph*)}]
\begin{multicols}{2}
  \item La represión, el autoritarismo y el anticomunismo
  \item La democracia y la libertad
  \item El socialismo y la revolución
  \item El liberalismo y la modernización
  \item El conservadurismo y la tradición
  \item El anarquismo y la libertad
\end{multicols}
\end{enumerate}
\vspace{0.15cm}
\item ¿Qué fue la huelga de brazos caídos de 1944?
\begin{enumerate}[label=\textbf{\Alph*)}]
\begin{multicols}{2}
  \item Una protesta que llevó a la caída de Hernández Martínez
  \item Una protesta colonial
  \item Una protesta conservadora
  \item Una protesta liberal
  \item Una protesta militar
  \item Una protesta estudiantil
\end{multicols}
\end{enumerate}
\vspace{0.15cm}
\item ¿Quién sucedió a Hernández Martínez?
\begin{enumerate}[label=\textbf{\Alph*)}]
\begin{multicols}{2}
  \item Osmín Aguirre y Salinas
  \item Salvador Castaneda Castro
  \item Óscar Osorio
  \item José María Lemus
  \item Julio Adalberto Rivera
  \item Fidel Sánchez Hernández
\end{multicols}
\end{enumerate}
\vspace{0.15cm}
\item ¿Qué fue el golpe de estado de 1948?
\begin{enumerate}[label=\textbf{\Alph*)}]
\begin{multicols}{2}
  \item Un golpe militar que llevó al poder a Óscar Osorio
  \item Un golpe colonial
  \item Un golpe conservador
  \item Un golpe liberal
  \item Un golpe comunista
  \item Un golpe estudiantil
\end{multicols}
\end{enumerate}
\vspace{0.15cm}
\item ¿Quién fue Óscar Osorio?
\begin{enumerate}[label=\textbf{\Alph*)}]
\begin{multicols}{2}
  \item Un presidente militar que impulsó reformas sociales
  \item Un dictador conservador
  \item Un líder comunista
  \item Un obispo salvadoreño
  \item Un caudillo guatemalteco
  \item Un rey español
\end{multicols}
\end{enumerate}
\vspace{0.15cm}
\item ¿Qué impulsó Óscar Osorio?
\begin{enumerate}[label=\textbf{\Alph*)}]
\begin{multicols}{2}
  \item La seguridad social, la educación y la infraestructura
  \item La represión y el autoritarismo
  \item El socialismo y la revolución
  \item El liberalismo y la modernización
  \item El conservadurismo y la tradición
  \item El anarquismo y la libertad
\end{multicols}
\end{enumerate}
\vspace{0.15cm}
\item ¿Quién fue José María Lemus?
\begin{enumerate}[label=\textbf{\Alph*)}]
\begin{multicols}{2}
  \item Un presidente militar que fue derrocado
  \item Un dictador conservador
  \item Un líder comunista
  \item Un obispo salvadoreño
  \item Un caudillo guatemalteco
  \item Un rey español
\end{multicols}
\end{enumerate}
\vspace{0.15cm}
\item ¿Qué fue el golpe de estado de 1960?
\begin{enumerate}[label=\textbf{\Alph*)}]
\begin{multicols}{2}
  \item Un golpe militar que derrocó a Lemus
  \item Un golpe colonial
  \item Un golpe conservador
  \item Un golpe liberal
  \item Un golpe comunista
  \item Un golpe estudiantil
\end{multicols}
\end{enumerate}
\vspace{0.15cm}
\item ¿Quién fue Julio Adalberto Rivera?
\begin{enumerate}[label=\textbf{\Alph*)}]
\begin{multicols}{2}
  \item Un presidente militar que impulsó la Alianza para el Progreso
  \item Un dictador conservador
  \item Un líder comunista
  \item Un obispo salvadoreño
  \item Un caudillo guatemalteco
  \item Un rey español
\end{multicols}
\end{enumerate}
\vspace{0.15cm}
\item ¿Qué fue la Alianza para el Progreso?
\begin{enumerate}[label=\textbf{\Alph*)}]
\begin{multicols}{2}
  \item Un programa de ayuda de Estados Unidos para América Latina
  \item Un programa de ayuda europea
  \item Un programa de ayuda soviética
  \item Un programa de ayuda china
  \item Un programa de ayuda japonesa
  \item Un programa de ayuda africana
\end{multicols}
\end{enumerate}
\vspace{0.15cm}
\item ¿Quién fue Fidel Sánchez Hernández?
\begin{enumerate}[label=\textbf{\Alph*)}]
\begin{multicols}{2}
  \item Un presidente militar que enfrentó la guerra con Honduras
  \item Un dictador conservador
  \item Un líder comunista
  \item Un obispo salvadoreño
  \item Un caudillo guatemalteco
  \item Un rey español
\end{multicols}
\end{enumerate}
\vspace{0.15cm}
\item ¿Qué fue la guerra de las 100 horas?
\begin{enumerate}[label=\textbf{\Alph*)}]
\begin{multicols}{2}
  \item El conflicto entre El Salvador y Honduras en 1969
  \item Una guerra civil
  \item Una invasión extranjera
  \item Una guerra colonial
  \item Una guerra mundial
  \item Una guerra fría
\end{multicols}
\end{enumerate}
\vspace{0.15cm}
\item ¿Cuál fue la causa de la guerra de las 100 horas?
\begin{enumerate}[label=\textbf{\Alph*)}]
\begin{multicols}{2}
  \item La migración salvadoreña a Honduras y la reforma agraria hondureña
  \item Una disputa territorial
  \item Una disputa comercial
  \item Una disputa política
  \item Una disputa religiosa
  \item Una disputa étnica
\end{multicols}
\end{enumerate}
\vspace{0.15cm}
\item ¿Qué fue la Organización Democrática Nacionalista (ORDEN)?
\begin{enumerate}[label=\textbf{\Alph*)}]
\begin{multicols}{2}
  \item Un grupo paramilitar de extrema derecha
  \item Un grupo guerrillero
  \item Un partido político
  \item Un sindicato
  \item Una organización religiosa
  \item Una organización estudiantil
\end{multicols}
\end{enumerate}
\vspace{0.15cm}
\item ¿Qué fue el Frente Farabundo Martí para la Liberación Nacional (FMLN)?
\begin{enumerate}[label=\textbf{\Alph*)}]
\begin{multicols}{2}
  \item Una organización guerrillera de izquierda
  \item Un partido político de derecha
  \item Un grupo paramilitar
  \item Un sindicato
  \item Una organización religiosa
  \item Una organización estudiantil
\end{multicols}
\end{enumerate}
\vspace{0.15cm}
\item ¿En qué año se fundó el FMLN?
\begin{enumerate}[label=\textbf{\Alph*)}]
\begin{multicols}{2}
  \item 1980
  \item 1970
  \item 1975
  \item 1985
  \item 1990
  \item 1992
\end{multicols}
\end{enumerate}
\vspace{0.15cm}
\item ¿Qué fue la guerra civil salvadoreña?
\begin{enumerate}[label=\textbf{\Alph*)}]
\begin{multicols}{2}
  \item Un conflicto armado entre el gobierno y la guerrilla
  \item Una guerra colonial
  \item Una guerra internacional
  \item Una guerra civil corta
  \item Una guerra mundial
  \item Una guerra fría
\end{multicols}
\end{enumerate}
\vspace{0.15cm}
\item ¿En qué año comenzó la guerra civil salvadoreña?
\begin{enumerate}[label=\textbf{\Alph*)}]
\begin{multicols}{2}
  \item 1980
  \item 1970
  \item 1975
  \item 1985
  \item 1990
  \item 1992
\end{multicols}
\end{enumerate}
\vspace{0.15cm}
\end{enumerate}\n
\subsubsection*{Sección IV: Acuerdos de Paz y Posguerra (141–200)}

\begin{enumerate}[label=\textbf{\arabic*.}]
\setcounter{enumi}{140}
\item ¿Qué fue el asesinato de Monseñor Óscar Arnulfo Romero?
\begin{enumerate}[label=\textbf{\Alph*)}]
\begin{multicols}{2}
  \item El asesinato de un obispo defensor de los derechos humanos
  \item El asesinato de un presidente
  \item El asesinato de un militar
  \item El asesinato de un guerrillero
  \item El asesinato de un empresario
  \item El asesinato de un estudiante
\end{multicols}
\end{enumerate}
\vspace{0.15cm}
\item ¿En qué año fue asesinado Monseñor Romero?
\begin{enumerate}[label=\textbf{\Alph*)}]
\begin{multicols}{2}
  \item 1980
  \item 1970
  \item 1975
  \item 1985
  \item 1990
  \item 1992
\end{multicols}
\end{enumerate}
\vspace{0.15cm}
\item ¿Qué fue la masacre de El Mozote?
\begin{enumerate}[label=\textbf{\Alph*)}]
\begin{multicols}{2}
  \item Una masacre de campesinos por el ejército en 1981
  \item Una masacre colonial
  \item Una masacre guerrillera
  \item Una masacre militar
  \item Una masacre estudiantil
  \item Una masacre religiosa
\end{multicols}
\end{enumerate}
\vspace{0.15cm}
\item ¿Qué fue la masacre de los Jesuitas?
\begin{enumerate}[label=\textbf{\Alph*)}]
\begin{multicols}{2}
  \item El asesinato de seis sacerdotes jesuitas en 1989
  \item Una masacre colonial
  \item Una masacre guerrillera
  \item Una masacre militar
  \item Una masacre estudiantil
  \item Una masacre religiosa
\end{multicols}
\end{enumerate}
\vspace{0.15cm}
\item ¿Quiénes fueron asesinados en la masacre de los Jesuitas?
\begin{enumerate}[label=\textbf{\Alph*)}]
\begin{multicols}{2}
  \item Ignacio Ellacuría, Ignacio Martín-Baró y otros
  \item Monseñor Romero
  \item Farabundo Martí
  \item Maximiliano Hernández Martínez
  \item Gerardo Barrios
  \item Francisco Dueñas
\end{multicols}
\end{enumerate}
\vspace{0.15cm}
\item ¿Qué fue la doctrina de seguridad nacional?
\begin{enumerate}[label=\textbf{\Alph*)}]
\begin{multicols}{2}
  \item Una doctrina anticomunista que justificó la represión
  \item Una doctrina democrática
  \item Una doctrina socialista
  \item Una doctrina liberal
  \item Una doctrina conservadora
  \item Una doctrina anarquista
\end{multicols}
\end{enumerate}
\vspace{0.15cm}
\item ¿Qué fue el batallón Atlacatl?
\begin{enumerate}[label=\textbf{\Alph*)}]
\begin{multicols}{2}
  \item Un batallón de élite del ejército salvadoreño
  \item Un batallón guerrillero
  \item Un batallón colonial
  \item Un batallón conservador
  \item Un batallón liberal
  \item Un batallón estudiantil
\end{multicols}
\end{enumerate}
\vspace{0.15cm}
\item ¿Qué fue la Fuerza Armada de El Salvador?
\begin{enumerate}[label=\textbf{\Alph*)}]
\begin{multicols}{2}
  \item La institución militar que combatió a la guerrilla
  \item Una organización guerrillera
  \item Un partido político
  \item Un sindicato
  \item Una organización religiosa
  \item Una organización estudiantil
\end{multicols}
\end{enumerate}
\vspace{0.15cm}
\item ¿Qué fue el Ejército Revolucionario del Pueblo (ERP)?
\begin{enumerate}[label=\textbf{\Alph*)}]
\begin{multicols}{2}
  \item Una organización guerrillera salvadoreña
  \item Un partido político
  \item Un grupo paramilitar
  \item Un sindicato
  \item Una organización religiosa
  \item Una organización estudiantil
\end{multicols}
\end{enumerate}
\vspace{0.15cm}
\item ¿Qué fue las Fuerzas Populares de Liberación (FPL)?
\begin{enumerate}[label=\textbf{\Alph*)}]
\begin{multicols}{2}
  \item Una organización guerrillera salvadoreña
  \item Un partido político
  \item Un grupo paramilitar
  \item Un sindicato
  \item Una organización religiosa
  \item Una organización estudiantil
\end{multicols}
\end{enumerate}
\vspace{0.15cm}
\item ¿Qué fue la Resistencia Nacional (RN)?
\begin{enumerate}[label=\textbf{\Alph*)}]
\begin{multicols}{2}
  \item Una organización guerrillera salvadoreña
  \item Un partido político
  \item Un grupo paramilitar
  \item Un sindicato
  \item Una organización religiosa
  \item Una organización estudiantil
\end{multicols}
\end{enumerate}
\vspace{0.15cm}
\item ¿Qué fue el Partido Comunista de El Salvador (PCS)?
\begin{enumerate}[label=\textbf{\Alph*)}]
\begin{multicols}{2}
  \item Una organización política de izquierda
  \item Un partido político de derecha
  \item Un grupo paramilitar
  \item Un sindicato
  \item Una organización religiosa
  \item Una organización estudiantil
\end{multicols}
\end{enumerate}
\vspace{0.15cm}
\item ¿Qué fue el Partido Demócrata Cristiano (PDC)?
\begin{enumerate}[label=\textbf{\Alph*)}]
\begin{multicols}{2}
  \item Un partido político de centro
  \item Un partido político de izquierda
  \item Un partido político de derecha
  \item Un grupo paramilitar
  \item Un sindicato
  \item Una organización religiosa
\end{multicols}
\end{enumerate}
\vspace{0.15cm}
\item ¿Quién fue José Napoleón Duarte?
\begin{enumerate}[label=\textbf{\Alph*)}]
\begin{multicols}{2}
  \item Un líder del PDC y presidente de El Salvador
  \item Un líder guerrillero
  \item Un militar
  \item Un obispo
  \item Un empresario
  \item Un estudiante
\end{multicols}
\end{enumerate}
\vspace{0.15cm}
\item ¿Qué fue la Alianza Republicana Nacionalista (ARENA)?
\begin{enumerate}[label=\textbf{\Alph*)}]
\begin{multicols}{2}
  \item Un partido político de derecha
  \item Un partido político de izquierda
  \item Un partido político de centro
  \item Un grupo paramilitar
  \item Un sindicato
  \item Una organización religiosa
\end{multicols}
\end{enumerate}
\vspace{0.15cm}
\item ¿Quién fundó ARENA?
\begin{enumerate}[label=\textbf{\Alph*)}]
\begin{multicols}{2}
  \item Roberto D'Aubuisson
  \item José Napoleón Duarte
  \item Alfredo Cristiani
  \item Armando Calderón Sol
  \item Francisco Flores
  \item Antonio Saca
\end{multicols}
\end{enumerate}
\vspace{0.15cm}
\item ¿Qué fue el gobierno de Alfredo Cristiani?
\begin{enumerate}[label=\textbf{\Alph*)}]
\begin{multicols}{2}
  \item El gobierno que firmó los Acuerdos de Paz
  \item El gobierno que inició la guerra
  \item El gobierno que continuó la guerra
  \item El gobierno que reprimió la protesta
  \item El gobierno que impulsó la reforma agraria
  \item El gobierno que nacionalizó la banca
\end{multicols}
\end{enumerate}
\vspace{0.15cm}
\item ¿En qué año se firmaron los Acuerdos de Paz?
\begin{enumerate}[label=\textbf{\Alph*)}]
\begin{multicols}{2}
  \item 1992
  \item 1980
  \item 1985
  \item 1990
  \item 1995
  \item 2000
\end{multicols}
\end{enumerate}
\vspace{0.15cm}
\item ¿Dónde se firmaron los Acuerdos de Paz?
\begin{enumerate}[label=\textbf{\Alph*)}]
\begin{multicols}{2}
  \item Chapultepec, México
  \item San Salvador, El Salvador
  \item Guatemala, Guatemala
  \item Tegucigalpa, Honduras
  \item Managua, Nicaragua
  \item San José, Costa Rica
\end{multicols}
\end{enumerate}
\vspace{0.15cm}
\item ¿Qué fue la Comisión de la Verdad?
\begin{enumerate}[label=\textbf{\Alph*)}]
\begin{multicols}{2}
  \item Una comisión que investigó las violaciones a los derechos humanos
  \item Una comisión guerrillera
  \item Una comisión militar
  \item Una comisión política
  \item Una comisión religiosa
  \item Una comisión estudiantil
\end{multicols}
\end{enumerate}
\vspace{0.15cm}
\item ¿Qué establecieron los Acuerdos de Paz?
\begin{enumerate}[label=\textbf{\Alph*)}]
\begin{multicols}{2}
  \item El fin del conflicto armado y reformas políticas
  \item La continuación de la guerra
  \item La victoria de la guerrilla
  \item La victoria del ejército
  \item La intervención extranjera
  \item La dictadura militar
\end{multicols}
\end{enumerate}
\vspace{0.15cm}
\item ¿Qué fue la Procuraduría para la Defensa de los Derechos Humanos?
\begin{enumerate}[label=\textbf{\Alph*)}]
\begin{multicols}{2}
  \item Una institución creada por los Acuerdos de Paz
  \item Una institución colonial
  \item Una institución militar
  \item Una institución guerrillera
  \item Una institución religiosa
  \item Una institución estudiantil
\end{multicols}
\end{enumerate}
\vspace{0.15cm}
\item ¿Qué fue la Policía Nacional Civil (PNC)?
\begin{enumerate}[label=\textbf{\Alph*)}]
\begin{multicols}{2}
  \item Una institución creada por los Acuerdos de Paz
  \item Una institución colonial
  \item Una institución militar
  \item Una institución guerrillera
  \item Una institución religiosa
  \item Una institución estudiantil
\end{multicols}
\end{enumerate}
\vspace{0.15cm}
\item ¿Qué fue el Tribunal Supremo Electoral (TSE)?
\begin{enumerate}[label=\textbf{\Alph*)}]
\begin{multicols}{2}
  \item Una institución creada por los Acuerdos de Paz
  \item Una institución colonial
  \item Una institución militar
  \item Una institución guerrillera
  \item Una institución religiosa
  \item Una institución estudiantil
\end{multicols}
\end{enumerate}
\vspace{0.15cm}
\item ¿Qué fue la reforma electoral?
\begin{enumerate}[label=\textbf{\Alph*)}]
\begin{multicols}{2}
  \item Una reforma política establecida en los Acuerdos de Paz
  \item Una reforma colonial
  \item Una reforma militar
  \item Una reforma guerrillera
  \item Una reforma religiosa
  \item Una reforma estudiantil
\end{multicols}
\end{enumerate}
\vspace{0.15cm}
\item ¿Qué fue la reforma judicial?
\begin{enumerate}[label=\textbf{\Alph*)}]
\begin{multicols}{2}
  \item Una reforma establecida en los Acuerdos de Paz
  \item Una reforma colonial
  \item Una reforma militar
  \item Una reforma guerrillera
  \item Una reforma religiosa
  \item Una reforma estudiantil
\end{multicols}
\end{enumerate}
\vspace{0.15cm}
\item ¿Qué fue la reforma agraria?
\begin{enumerate}[label=\textbf{\Alph*)}]
\begin{multicols}{2}
  \item Una reforma establecida en los Acuerdos de Paz
  \item Una reforma colonial
  \item Una reforma militar
  \item Una reforma guerrillera
  \item Una reforma religiosa
  \item Una reforma estudiantil
\end{multicols}
\end{enumerate}
\vspace{0.15cm}
\item ¿Qué fue la desmilitarización?
\begin{enumerate}[label=\textbf{\Alph*)}]
\begin{multicols}{2}
  \item La reducción del poder del ejército tras los Acuerdos de Paz
  \item El aumento del poder del ejército
  \item La eliminación del ejército
  \item La creación del ejército
  \item La reforma del ejército
  \item La profesionalización del ejército
\end{multicols}
\end{enumerate}
\vspace{0.15cm}
\item ¿Qué fue la reintegración de los excombatientes?
\begin{enumerate}[label=\textbf{\Alph*)}]
\begin{multicols}{2}
  \item La incorporación de los exguerrilleros y exmilitares a la vida civil
  \item La expulsión de los excombatientes
  \item La eliminación de los excombatientes
  \item La educación de los excombatientes
  \item La migración de los excombatientes
  \item La reducción de los excombatientes
\end{multicols}
\end{enumerate}
\vspace{0.15cm}
\item ¿Qué fue el FMLN como partido político?
\begin{enumerate}[label=\textbf{\Alph*)}]
\begin{multicols}{2}
  \item La transformación de la guerrilla en partido político tras los Acuerdos de Paz
  \item La continuación de la guerrilla
  \item La eliminación de la guerrilla
  \item La militarización de la guerrilla
  \item La educación de la guerrilla
  \item La migración de la guerrilla
\end{multicols}
\end{enumerate}
\vspace{0.15cm}
\item ¿Qué fue el asesinato de Monseñor Óscar Arnulfo Romero?
\begin{enumerate}[label=\textbf{\Alph*)}]
\begin{multicols}{2}
  \item El asesinato de un obispo defensor de los derechos humanos
  \item El asesinato de un presidente
  \item El asesinato de un militar
  \item El asesinato de un guerrillero
  \item El asesinato de un empresario
  \item El asesinato de un estudiante
\end{multicols}
\end{enumerate}
\vspace{0.15cm}
\item ¿En qué año fue asesinado Monseñor Romero?
\begin{enumerate}[label=\textbf{\Alph*)}]
\begin{multicols}{2}
  \item 1980
  \item 1970
  \item 1975
  \item 1985
  \item 1990
  \item 1992
\end{multicols}
\end{enumerate}
\vspace{0.15cm}
\item ¿Qué fue la masacre de El Mozote?
\begin{enumerate}[label=\textbf{\Alph*)}]
\begin{multicols}{2}
  \item Una masacre de campesinos por el ejército en 1981
  \item Una masacre colonial
  \item Una masacre guerrillera
  \item Una masacre militar
  \item Una masacre estudiantil
  \item Una masacre religiosa
\end{multicols}
\end{enumerate}
\vspace{0.15cm}
\item ¿Qué fue la masacre de los Jesuitas?
\begin{enumerate}[label=\textbf{\Alph*)}]
\begin{multicols}{2}
  \item El asesinato de seis sacerdotes jesuitas en 1989
  \item Una masacre colonial
  \item Una masacre guerrillera
  \item Una masacre militar
  \item Una masacre estudiantil
  \item Una masacre religiosa
\end{multicols}
\end{enumerate}
\vspace{0.15cm}
\item ¿Quiénes fueron asesinados en la masacre de los Jesuitas?
\begin{enumerate}[label=\textbf{\Alph*)}]
\begin{multicols}{2}
  \item Ignacio Ellacuría, Ignacio Martín-Baró y otros
  \item Monseñor Romero
  \item Farabundo Martí
  \item Maximiliano Hernández Martínez
  \item Gerardo Barrios
  \item Francisco Dueñas
\end{multicols}
\end{enumerate}
\vspace{0.15cm}
\item ¿Qué fue la doctrina de seguridad nacional?
\begin{enumerate}[label=\textbf{\Alph*)}]
\begin{multicols}{2}
  \item Una doctrina anticomunista que justificó la represión
  \item Una doctrina democrática
  \item Una doctrina socialista
  \item Una doctrina liberal
  \item Una doctrina conservadora
  \item Una doctrina anarquista
\end{multicols}
\end{enumerate}
\vspace{0.15cm}
\item ¿Qué fue el batallón Atlacatl?
\begin{enumerate}[label=\textbf{\Alph*)}]
\begin{multicols}{2}
  \item Un batallón de élite del ejército salvadoreño
  \item Un batallón guerrillero
  \item Un batallón colonial
  \item Un batallón conservador
  \item Un batallón liberal
  \item Un batallón estudiantil
\end{multicols}
\end{enumerate}
\vspace{0.15cm}
\item ¿Qué fue la Fuerza Armada de El Salvador?
\begin{enumerate}[label=\textbf{\Alph*)}]
\begin{multicols}{2}
  \item La institución militar que combatió a la guerrilla
  \item Una organización guerrillera
  \item Un partido político
  \item Un sindicato
  \item Una organización religiosa
  \item Una organización estudiantil
\end{multicols}
\end{enumerate}
\vspace{0.15cm}
\item ¿Qué fue el Ejército Revolucionario del Pueblo (ERP)?
\begin{enumerate}[label=\textbf{\Alph*)}]
\begin{multicols}{2}
  \item Una organización guerrillera salvadoreña
  \item Un partido político
  \item Un grupo paramilitar
  \item Un sindicato
  \item Una organización religiosa
  \item Una organización estudiantil
\end{multicols}
\end{enumerate}
\vspace{0.15cm}
\item ¿Qué fue las Fuerzas Populares de Liberación (FPL)?
\begin{enumerate}[label=\textbf{\Alph*)}]
\begin{multicols}{2}
  \item Una organización guerrillera salvadoreña
  \item Un partido político
  \item Un grupo paramilitar
  \item Un sindicato
  \item Una organización religiosa
  \item Una organización estudiantil
\end{multicols}
\end{enumerate}
\vspace{0.15cm}
\item ¿Qué fue la Resistencia Nacional (RN)?
\begin{enumerate}[label=\textbf{\Alph*)}]
\begin{multicols}{2}
  \item Una organización guerrillera salvadoreña
  \item Un partido político
  \item Un grupo paramilitar
  \item Un sindicato
  \item Una organización religiosa
  \item Una organización estudiantil
\end{multicols}
\end{enumerate}
\vspace{0.15cm}
\item ¿Qué fue el Partido Comunista de El Salvador (PCS)?
\begin{enumerate}[label=\textbf{\Alph*)}]
\begin{multicols}{2}
  \item Una organización política de izquierda
  \item Un partido político de derecha
  \item Un grupo paramilitar
  \item Un sindicato
  \item Una organización religiosa
  \item Una organización estudiantil
\end{multicols}
\end{enumerate}
\vspace{0.15cm}
\item ¿Qué fue el Partido Demócrata Cristiano (PDC)?
\begin{enumerate}[label=\textbf{\Alph*)}]
\begin{multicols}{2}
  \item Un partido político de centro
  \item Un partido político de izquierda
  \item Un partido político de derecha
  \item Un grupo paramilitar
  \item Un sindicato
  \item Una organización religiosa
\end{multicols}
\end{enumerate}
\vspace{0.15cm}
\item ¿Quién fue José Napoleón Duarte?
\begin{enumerate}[label=\textbf{\Alph*)}]
\begin{multicols}{2}
  \item Un líder del PDC y presidente de El Salvador
  \item Un líder guerrillero
  \item Un militar
  \item Un obispo
  \item Un empresario
  \item Un estudiante
\end{multicols}
\end{enumerate}
\vspace{0.15cm}
\item ¿Qué fue la Alianza Republicana Nacionalista (ARENA)?
\begin{enumerate}[label=\textbf{\Alph*)}]
\begin{multicols}{2}
  \item Un partido político de derecha
  \item Un partido político de izquierda
  \item Un partido político de centro
  \item Un grupo paramilitar
  \item Un sindicato
  \item Una organización religiosa
\end{multicols}
\end{enumerate}
\vspace{0.15cm}
\item ¿Quién fundó ARENA?
\begin{enumerate}[label=\textbf{\Alph*)}]
\begin{multicols}{2}
  \item Roberto D'Aubuisson
  \item José Napoleón Duarte
  \item Alfredo Cristiani
  \item Armando Calderón Sol
  \item Francisco Flores
  \item Antonio Saca
\end{multicols}
\end{enumerate}
\vspace{0.15cm}
\item ¿Qué fue el gobierno de Alfredo Cristiani?
\begin{enumerate}[label=\textbf{\Alph*)}]
\begin{multicols}{2}
  \item El gobierno que firmó los Acuerdos de Paz
  \item El gobierno que inició la guerra
  \item El gobierno que continuó la guerra
  \item El gobierno que reprimió la protesta
  \item El gobierno que impulsó la reforma agraria
  \item El gobierno que nacionalizó la banca
\end{multicols}
\end{enumerate}
\vspace{0.15cm}
\item ¿En qué año se firmaron los Acuerdos de Paz?
\begin{enumerate}[label=\textbf{\Alph*)}]
\begin{multicols}{2}
  \item 1992
  \item 1980
  \item 1985
  \item 1990
  \item 1995
  \item 2000
\end{multicols}
\end{enumerate}
\vspace{0.15cm}
\item ¿Dónde se firmaron los Acuerdos de Paz?
\begin{enumerate}[label=\textbf{\Alph*)}]
\begin{multicols}{2}
  \item Chapultepec, México
  \item San Salvador, El Salvador
  \item Guatemala, Guatemala
  \item Tegucigalpa, Honduras
  \item Managua, Nicaragua
  \item San José, Costa Rica
\end{multicols}
\end{enumerate}
\vspace{0.15cm}
\item ¿Qué fue la Comisión de la Verdad?
\begin{enumerate}[label=\textbf{\Alph*)}]
\begin{multicols}{2}
  \item Una comisión que investigó las violaciones a los derechos humanos
  \item Una comisión guerrillera
  \item Una comisión militar
  \item Una comisión política
  \item Una comisión religiosa
  \item Una comisión estudiantil
\end{multicols}
\end{enumerate}
\vspace{0.15cm}
\item ¿Qué establecieron los Acuerdos de Paz?
\begin{enumerate}[label=\textbf{\Alph*)}]
\begin{multicols}{2}
  \item El fin del conflicto armado y reformas políticas
  \item La continuación de la guerra
  \item La victoria de la guerrilla
  \item La victoria del ejército
  \item La intervención extranjera
  \item La dictadura militar
\end{multicols}
\end{enumerate}
\vspace{0.15cm}
\item ¿Qué fue la Procuraduría para la Defensa de los Derechos Humanos?
\begin{enumerate}[label=\textbf{\Alph*)}]
\begin{multicols}{2}
  \item Una institución creada por los Acuerdos de Paz
  \item Una institución colonial
  \item Una institución militar
  \item Una institución guerrillera
  \item Una institución religiosa
  \item Una institución estudiantil
\end{multicols}
\end{enumerate}
\vspace{0.15cm}
\item ¿Qué fue la Policía Nacional Civil (PNC)?
\begin{enumerate}[label=\textbf{\Alph*)}]
\begin{multicols}{2}
  \item Una institución creada por los Acuerdos de Paz
  \item Una institución colonial
  \item Una institución militar
  \item Una institución guerrillera
  \item Una institución religiosa
  \item Una institución estudiantil
\end{multicols}
\end{enumerate}
\vspace{0.15cm}
\item ¿Qué fue el Tribunal Supremo Electoral (TSE)?
\begin{enumerate}[label=\textbf{\Alph*)}]
\begin{multicols}{2}
  \item Una institución creada por los Acuerdos de Paz
  \item Una institución colonial
  \item Una institución militar
  \item Una institución guerrillera
  \item Una institución religiosa
  \item Una institución estudiantil
\end{multicols}
\end{enumerate}
\vspace{0.15cm}
\item ¿Qué fue la reforma electoral?
\begin{enumerate}[label=\textbf{\Alph*)}]
\begin{multicols}{2}
  \item Una reforma política establecida en los Acuerdos de Paz
  \item Una reforma colonial
  \item Una reforma militar
  \item Una reforma guerrillera
  \item Una reforma religiosa
  \item Una reforma estudiantil
\end{multicols}
\end{enumerate}
\vspace{0.15cm}
\item ¿Qué fue la reforma judicial?
\begin{enumerate}[label=\textbf{\Alph*)}]
\begin{multicols}{2}
  \item Una reforma establecida en los Acuerdos de Paz
  \item Una reforma colonial
  \item Una reforma militar
  \item Una reforma guerrillera
  \item Una reforma religiosa
  \item Una reforma estudiantil
\end{multicols}
\end{enumerate}
\vspace{0.15cm}
\item ¿Qué fue la reforma agraria?
\begin{enumerate}[label=\textbf{\Alph*)}]
\begin{multicols}{2}
  \item Una reforma establecida en los Acuerdos de Paz
  \item Una reforma colonial
  \item Una reforma militar
  \item Una reforma guerrillera
  \item Una reforma religiosa
  \item Una reforma estudiantil
\end{multicols}
\end{enumerate}
\vspace{0.15cm}
\item ¿Qué fue la desmilitarización?
\begin{enumerate}[label=\textbf{\Alph*)}]
\begin{multicols}{2}
  \item La reducción del poder del ejército tras los Acuerdos de Paz
  \item El aumento del poder del ejército
  \item La eliminación del ejército
  \item La creación del ejército
  \item La reforma del ejército
  \item La profesionalización del ejército
\end{multicols}
\end{enumerate}
\vspace{0.15cm}
\item ¿Qué fue la reintegración de los excombatientes?
\begin{enumerate}[label=\textbf{\Alph*)}]
\begin{multicols}{2}
  \item La incorporación de los exguerrilleros y exmilitares a la vida civil
  \item La expulsión de los excombatientes
  \item La eliminación de los excombatientes
  \item La educación de los excombatientes
  \item La migración de los excombatientes
  \item La reducción de los excombatientes
\end{multicols}
\end{enumerate}
\vspace{0.15cm}
\item ¿Qué fue el FMLN como partido político?
\begin{enumerate}[label=\textbf{\Alph*)}]
\begin{multicols}{2}
  \item La transformación de la guerrilla en partido político tras los Acuerdos de Paz
  \item La continuación de la guerrilla
  \item La eliminación de la guerrilla
  \item La militarización de la guerrilla
  \item La educación de la guerrilla
  \item La migración de la guerrilla
\end{multicols}
\end{enumerate}
\vspace{0.15cm}
\end{enumerate}\n
\subsubsection*{Sección V: El Salvador Contemporáneo (201–300)}

\begin{enumerate}[label=\textbf{\arabic*.}]
\setcounter{enumi}{200}
\item ¿Quién fue Armando Calderón Sol?
\begin{enumerate}[label=\textbf{\Alph*)}]
\begin{multicols}{2}
  \item Un presidente de ARENA
  \item Un presidente del FMLN
  \item Un presidente del PDC
  \item Un presidente militar
  \item Un presidente independiente
  \item Un presidente del PCN
\end{multicols}
\end{enumerate}
\vspace{0.15cm}
\item ¿En qué período gobernó Armando Calderón Sol?
\begin{enumerate}[label=\textbf{\Alph*)}]
\begin{multicols}{2}
  \item 1994-1999
  \item 1989-1994
  \item 1999-2004
  \item 2004-2009
  \item 2009-2014
  \item 2014-2019
\end{multicols}
\end{enumerate}
\vspace{0.15cm}
\item ¿Qué impulsó Armando Calderón Sol?
\begin{enumerate}[label=\textbf{\Alph*)}]
\begin{multicols}{2}
  \item La privatización de empresas estatales
  \item La nacionalización de la banca
  \item La reforma agraria
  \item La desmilitarización
  \item La reforma electoral
  \item La reforma judicial
\end{multicols}
\end{enumerate}
\vspace{0.15cm}
\item ¿Quién fue Francisco Flores?
\begin{enumerate}[label=\textbf{\Alph*)}]
\begin{multicols}{2}
  \item Un presidente de ARENA
  \item Un presidente del FMLN
  \item Un presidente del PDC
  \item Un presidente militar
  \item Un presidente independiente
  \item Un presidente del PCN
\end{multicols}
\end{enumerate}
\vspace{0.15cm}
\item ¿En qué período gobernó Francisco Flores?
\begin{enumerate}[label=\textbf{\Alph*)}]
\begin{multicols}{2}
  \item 1999-2004
  \item 1994-1999
  \item 1989-1994
  \item 2004-2009
  \item 2009-2014
  \item 2014-2019
\end{multicols}
\end{enumerate}
\vspace{0.15cm}
\item ¿Qué impulsó Francisco Flores?
\begin{enumerate}[label=\textbf{\Alph*)}]
\begin{multicols}{2}
  \item La dolarización de la economía
  \item La privatización de empresas
  \item La nacionalización de la banca
  \item La reforma agraria
  \item La desmilitarización
  \item La reforma electoral
\end{multicols}
\end{enumerate}
\vspace{0.15cm}
\item ¿En qué año se dolarizó la economía salvadoreña?
\begin{enumerate}[label=\textbf{\Alph*)}]
\begin{multicols}{2}
  \item 2001
  \item 1999
  \item 2004
  \item 1994
  \item 2009
  \item 2014
\end{multicols}
\end{enumerate}
\vspace{0.15cm}
\item ¿Quién fue Antonio Saca?
\begin{enumerate}[label=\textbf{\Alph*)}]
\begin{multicols}{2}
  \item Un presidente de ARENA
  \item Un presidente del FMLN
  \item Un presidente del PDC
  \item Un presidente militar
  \item Un presidente independiente
  \item Un presidente del PCN
\end{multicols}
\end{enumerate}
\vspace{0.15cm}
\item ¿En qué período gobernó Antonio Saca?
\begin{enumerate}[label=\textbf{\Alph*)}]
\begin{multicols}{2}
  \item 2004-2009
  \item 1999-2004
  \item 1994-1999
  \item 1989-1994
  \item 2009-2014
  \item 2014-2019
\end{multicols}
\end{enumerate}
\vspace{0.15cm}
\item ¿Qué impulsó Antonio Saca?
\begin{enumerate}[label=\textbf{\Alph*)}]
\begin{multicols}{2}
  \item La firma del Tratado de Libre Comercio con Estados Unidos
  \item La dolarización
  \item La privatización
  \item La reforma agraria
  \item La desmilitarización
  \item La reforma electoral
\end{multicols}
\end{enumerate}
\vspace{0.15cm}
\item ¿Quién fue Mauricio Funes?
\begin{enumerate}[label=\textbf{\Alph*)}]
\begin{multicols}{2}
  \item El primer presidente del FMLN
  \item Un presidente de ARENA
  \item Un presidente del PDC
  \item Un presidente militar
  \item Un presidente independiente
  \item Un presidente del PCN
\end{multicols}
\end{enumerate}
\vspace{0.15cm}
\item ¿En qué período gobernó Mauricio Funes?
\begin{enumerate}[label=\textbf{\Alph*)}]
\begin{multicols}{2}
  \item 2009-2014
  \item 2004-2009
  \item 1999-2004
  \item 1994-1999
  \item 2014-2019
  \item 2019-2024
\end{multicols}
\end{enumerate}
\vspace{0.15cm}
\item ¿Qué impulsó Mauricio Funes?
\begin{enumerate}[label=\textbf{\Alph*)}]
\begin{multicols}{2}
  \item Programas sociales y la salida de tropas de Irak
  \item La dolarización
  \item La privatización
  \item La reforma agraria
  \item La desmilitarización
  \item La reforma electoral
\end{multicols}
\end{enumerate}
\vspace{0.15cm}
\item ¿Quién fue Salvador Sánchez Cerén?
\begin{enumerate}[label=\textbf{\Alph*)}]
\begin{multicols}{2}
  \item Un presidente del FMLN
  \item Un presidente de ARENA
  \item Un presidente del PDC
  \item Un presidente militar
  \item Un presidente independiente
  \item Un presidente del PCN
\end{multicols}
\end{enumerate}
\vspace{0.15cm}
\item ¿En qué período gobernó Salvador Sánchez Cerén?
\begin{enumerate}[label=\textbf{\Alph*)}]
\begin{multicols}{2}
  \item 2014-2019
  \item 2009-2014
  \item 2004-2009
  \item 1999-2004
  \item 1994-1999
  \item 2019-2024
\end{multicols}
\end{enumerate}
\vspace{0.15cm}
\item ¿Quién fue Nayib Bukele?
\begin{enumerate}[label=\textbf{\Alph*)}]
\begin{multicols}{2}
  \item Un presidente de Nuevas Ideas
  \item Un presidente de ARENA
  \item Un presidente del FMLN
  \item Un presidente del PDC
  \item Un presidente militar
  \item Un presidente del PCN
\end{multicols}
\end{enumerate}
\vspace{0.15cm}
\item ¿En qué año asumió Nayib Bukele la presidencia?
\begin{enumerate}[label=\textbf{\Alph*)}]
\begin{multicols}{2}
  \item 2019
  \item 2014
  \item 2009
  \item 2004
  \item 1999
  \item 2024
\end{multicols}
\end{enumerate}
\vspace{0.15cm}
\item ¿Qué impulsó Nayib Bukele?
\begin{enumerate}[label=\textbf{\Alph*)}]
\begin{multicols}{2}
  \item La seguridad, la tecnología y la imagen del país
  \item La dolarización
  \item La privatización
  \item La reforma agraria
  \item La desmilitarización
  \item La reforma electoral
\end{multicols}
\end{enumerate}
\vspace{0.15cm}
\item ¿Qué es el régimen de excepción?
\begin{enumerate}[label=\textbf{\Alph*)}]
\begin{multicols}{2}
  \item Una medida que suspende garantías constitucionales para combatir la violencia
  \item Una medida que amplía las garantías constitucionales
  \item Una medida que elimina las garantías constitucionales
  \item Una medida que reforma las garantías constitucionales
  \item Una medida que protege las garantías constitucionales
  \item Una medida que ignora las garantías constitucionales
\end{multicols}
\end{enumerate}
\vspace{0.15cm}
\item ¿Qué es la guerra contra las pandillas?
\begin{enumerate}[label=\textbf{\Alph*)}]
\begin{multicols}{2}
  \item Una estrategia de seguridad para combatir a las maras
  \item Una estrategia de diálogo con las maras
  \item Una estrategia de negociación con las maras
  \item Una estrategia de integración de las maras
  \item Una estrategia de educación de las maras
  \item Una estrategia de prevención de las maras
\end{multicols}
\end{enumerate}
\vspace{0.15cm}
\item ¿Quién fue Armando Calderón Sol?
\begin{enumerate}[label=\textbf{\Alph*)}]
\begin{multicols}{2}
  \item Un presidente de ARENA
  \item Un presidente del FMLN
  \item Un presidente del PDC
  \item Un presidente militar
  \item Un presidente independiente
  \item Un presidente del PCN
\end{multicols}
\end{enumerate}
\vspace{0.15cm}
\item ¿En qué período gobernó Armando Calderón Sol?
\begin{enumerate}[label=\textbf{\Alph*)}]
\begin{multicols}{2}
  \item 1994-1999
  \item 1989-1994
  \item 1999-2004
  \item 2004-2009
  \item 2009-2014
  \item 2014-2019
\end{multicols}
\end{enumerate}
\vspace{0.15cm}
\item ¿Qué impulsó Armando Calderón Sol?
\begin{enumerate}[label=\textbf{\Alph*)}]
\begin{multicols}{2}
  \item La privatización de empresas estatales
  \item La nacionalización de la banca
  \item La reforma agraria
  \item La desmilitarización
  \item La reforma electoral
  \item La reforma judicial
\end{multicols}
\end{enumerate}
\vspace{0.15cm}
\item ¿Quién fue Francisco Flores?
\begin{enumerate}[label=\textbf{\Alph*)}]
\begin{multicols}{2}
  \item Un presidente de ARENA
  \item Un presidente del FMLN
  \item Un presidente del PDC
  \item Un presidente militar
  \item Un presidente independiente
  \item Un presidente del PCN
\end{multicols}
\end{enumerate}
\vspace{0.15cm}
\item ¿En qué período gobernó Francisco Flores?
\begin{enumerate}[label=\textbf{\Alph*)}]
\begin{multicols}{2}
  \item 1999-2004
  \item 1994-1999
  \item 1989-1994
  \item 2004-2009
  \item 2009-2014
  \item 2014-2019
\end{multicols}
\end{enumerate}
\vspace{0.15cm}
\item ¿Qué impulsó Francisco Flores?
\begin{enumerate}[label=\textbf{\Alph*)}]
\begin{multicols}{2}
  \item La dolarización de la economía
  \item La privatización de empresas
  \item La nacionalización de la banca
  \item La reforma agraria
  \item La desmilitarización
  \item La reforma electoral
\end{multicols}
\end{enumerate}
\vspace{0.15cm}
\item ¿En qué año se dolarizó la economía salvadoreña?
\begin{enumerate}[label=\textbf{\Alph*)}]
\begin{multicols}{2}
  \item 2001
  \item 1999
  \item 2004
  \item 1994
  \item 2009
  \item 2014
\end{multicols}
\end{enumerate}
\vspace{0.15cm}
\item ¿Quién fue Antonio Saca?
\begin{enumerate}[label=\textbf{\Alph*)}]
\begin{multicols}{2}
  \item Un presidente de ARENA
  \item Un presidente del FMLN
  \item Un presidente del PDC
  \item Un presidente militar
  \item Un presidente independiente
  \item Un presidente del PCN
\end{multicols}
\end{enumerate}
\vspace{0.15cm}
\item ¿En qué período gobernó Antonio Saca?
\begin{enumerate}[label=\textbf{\Alph*)}]
\begin{multicols}{2}
  \item 2004-2009
  \item 1999-2004
  \item 1994-1999
  \item 1989-1994
  \item 2009-2014
  \item 2014-2019
\end{multicols}
\end{enumerate}
\vspace{0.15cm}
\item ¿Qué impulsó Antonio Saca?
\begin{enumerate}[label=\textbf{\Alph*)}]
\begin{multicols}{2}
  \item La firma del Tratado de Libre Comercio con Estados Unidos
  \item La dolarización
  \item La privatización
  \item La reforma agraria
  \item La desmilitarización
  \item La reforma electoral
\end{multicols}
\end{enumerate}
\vspace{0.15cm}
\item ¿Quién fue Mauricio Funes?
\begin{enumerate}[label=\textbf{\Alph*)}]
\begin{multicols}{2}
  \item El primer presidente del FMLN
  \item Un presidente de ARENA
  \item Un presidente del PDC
  \item Un presidente militar
  \item Un presidente independiente
  \item Un presidente del PCN
\end{multicols}
\end{enumerate}
\vspace{0.15cm}
\item ¿En qué período gobernó Mauricio Funes?
\begin{enumerate}[label=\textbf{\Alph*)}]
\begin{multicols}{2}
  \item 2009-2014
  \item 2004-2009
  \item 1999-2004
  \item 1994-1999
  \item 2014-2019
  \item 2019-2024
\end{multicols}
\end{enumerate}
\vspace{0.15cm}
\item ¿Qué impulsó Mauricio Funes?
\begin{enumerate}[label=\textbf{\Alph*)}]
\begin{multicols}{2}
  \item Programas sociales y la salida de tropas de Irak
  \item La dolarización
  \item La privatización
  \item La reforma agraria
  \item La desmilitarización
  \item La reforma electoral
\end{multicols}
\end{enumerate}
\vspace{0.15cm}
\item ¿Quién fue Salvador Sánchez Cerén?
\begin{enumerate}[label=\textbf{\Alph*)}]
\begin{multicols}{2}
  \item Un presidente del FMLN
  \item Un presidente de ARENA
  \item Un presidente del PDC
  \item Un presidente militar
  \item Un presidente independiente
  \item Un presidente del PCN
\end{multicols}
\end{enumerate}
\vspace{0.15cm}
\item ¿En qué período gobernó Salvador Sánchez Cerén?
\begin{enumerate}[label=\textbf{\Alph*)}]
\begin{multicols}{2}
  \item 2014-2019
  \item 2009-2014
  \item 2004-2009
  \item 1999-2004
  \item 1994-1999
  \item 2019-2024
\end{multicols}
\end{enumerate}
\vspace{0.15cm}
\item ¿Quién fue Nayib Bukele?
\begin{enumerate}[label=\textbf{\Alph*)}]
\begin{multicols}{2}
  \item Un presidente de Nuevas Ideas
  \item Un presidente de ARENA
  \item Un presidente del FMLN
  \item Un presidente del PDC
  \item Un presidente militar
  \item Un presidente del PCN
\end{multicols}
\end{enumerate}
\vspace{0.15cm}
\item ¿En qué año asumió Nayib Bukele la presidencia?
\begin{enumerate}[label=\textbf{\Alph*)}]
\begin{multicols}{2}
  \item 2019
  \item 2014
  \item 2009
  \item 2004
  \item 1999
  \item 2024
\end{multicols}
\end{enumerate}
\vspace{0.15cm}
\item ¿Qué impulsó Nayib Bukele?
\begin{enumerate}[label=\textbf{\Alph*)}]
\begin{multicols}{2}
  \item La seguridad, la tecnología y la imagen del país
  \item La dolarización
  \item La privatización
  \item La reforma agraria
  \item La desmilitarización
  \item La reforma electoral
\end{multicols}
\end{enumerate}
\vspace{0.15cm}
\item ¿Qué es el régimen de excepción?
\begin{enumerate}[label=\textbf{\Alph*)}]
\begin{multicols}{2}
  \item Una medida que suspende garantías constitucionales para combatir la violencia
  \item Una medida que amplía las garantías constitucionales
  \item Una medida que elimina las garantías constitucionales
  \item Una medida que reforma las garantías constitucionales
  \item Una medida que protege las garantías constitucionales
  \item Una medida que ignora las garantías constitucionales
\end{multicols}
\end{enumerate}
\vspace{0.15cm}
\item ¿Qué es la guerra contra las pandillas?
\begin{enumerate}[label=\textbf{\Alph*)}]
\begin{multicols}{2}
  \item Una estrategia de seguridad para combatir a las maras
  \item Una estrategia de diálogo con las maras
  \item Una estrategia de negociación con las maras
  \item Una estrategia de integración de las maras
  \item Una estrategia de educación de las maras
  \item Una estrategia de prevención de las maras
\end{multicols}
\end{enumerate}
\vspace{0.15cm}
\item ¿Quién fue Armando Calderón Sol?
\begin{enumerate}[label=\textbf{\Alph*)}]
\begin{multicols}{2}
  \item Un presidente de ARENA
  \item Un presidente del FMLN
  \item Un presidente del PDC
  \item Un presidente militar
  \item Un presidente independiente
  \item Un presidente del PCN
\end{multicols}
\end{enumerate}
\vspace{0.15cm}
\item ¿En qué período gobernó Armando Calderón Sol?
\begin{enumerate}[label=\textbf{\Alph*)}]
\begin{multicols}{2}
  \item 1994-1999
  \item 1989-1994
  \item 1999-2004
  \item 2004-2009
  \item 2009-2014
  \item 2014-2019
\end{multicols}
\end{enumerate}
\vspace{0.15cm}
\item ¿Qué impulsó Armando Calderón Sol?
\begin{enumerate}[label=\textbf{\Alph*)}]
\begin{multicols}{2}
  \item La privatización de empresas estatales
  \item La nacionalización de la banca
  \item La reforma agraria
  \item La desmilitarización
  \item La reforma electoral
  \item La reforma judicial
\end{multicols}
\end{enumerate}
\vspace{0.15cm}
\item ¿Quién fue Francisco Flores?
\begin{enumerate}[label=\textbf{\Alph*)}]
\begin{multicols}{2}
  \item Un presidente de ARENA
  \item Un presidente del FMLN
  \item Un presidente del PDC
  \item Un presidente militar
  \item Un presidente independiente
  \item Un presidente del PCN
\end{multicols}
\end{enumerate}
\vspace{0.15cm}
\item ¿En qué período gobernó Francisco Flores?
\begin{enumerate}[label=\textbf{\Alph*)}]
\begin{multicols}{2}
  \item 1999-2004
  \item 1994-1999
  \item 1989-1994
  \item 2004-2009
  \item 2009-2014
  \item 2014-2019
\end{multicols}
\end{enumerate}
\vspace{0.15cm}
\item ¿Qué impulsó Francisco Flores?
\begin{enumerate}[label=\textbf{\Alph*)}]
\begin{multicols}{2}
  \item La dolarización de la economía
  \item La privatización de empresas
  \item La nacionalización de la banca
  \item La reforma agraria
  \item La desmilitarización
  \item La reforma electoral
\end{multicols}
\end{enumerate}
\vspace{0.15cm}
\item ¿En qué año se dolarizó la economía salvadoreña?
\begin{enumerate}[label=\textbf{\Alph*)}]
\begin{multicols}{2}
  \item 2001
  \item 1999
  \item 2004
  \item 1994
  \item 2009
  \item 2014
\end{multicols}
\end{enumerate}
\vspace{0.15cm}
\item ¿Quién fue Antonio Saca?
\begin{enumerate}[label=\textbf{\Alph*)}]
\begin{multicols}{2}
  \item Un presidente de ARENA
  \item Un presidente del FMLN
  \item Un presidente del PDC
  \item Un presidente militar
  \item Un presidente independiente
  \item Un presidente del PCN
\end{multicols}
\end{enumerate}
\vspace{0.15cm}
\item ¿En qué período gobernó Antonio Saca?
\begin{enumerate}[label=\textbf{\Alph*)}]
\begin{multicols}{2}
  \item 2004-2009
  \item 1999-2004
  \item 1994-1999
  \item 1989-1994
  \item 2009-2014
  \item 2014-2019
\end{multicols}
\end{enumerate}
\vspace{0.15cm}
\item ¿Qué impulsó Antonio Saca?
\begin{enumerate}[label=\textbf{\Alph*)}]
\begin{multicols}{2}
  \item La firma del Tratado de Libre Comercio con Estados Unidos
  \item La dolarización
  \item La privatización
  \item La reforma agraria
  \item La desmilitarización
  \item La reforma electoral
\end{multicols}
\end{enumerate}
\vspace{0.15cm}
\item ¿Quién fue Mauricio Funes?
\begin{enumerate}[label=\textbf{\Alph*)}]
\begin{multicols}{2}
  \item El primer presidente del FMLN
  \item Un presidente de ARENA
  \item Un presidente del PDC
  \item Un presidente militar
  \item Un presidente independiente
  \item Un presidente del PCN
\end{multicols}
\end{enumerate}
\vspace{0.15cm}
\item ¿En qué período gobernó Mauricio Funes?
\begin{enumerate}[label=\textbf{\Alph*)}]
\begin{multicols}{2}
  \item 2009-2014
  \item 2004-2009
  \item 1999-2004
  \item 1994-1999
  \item 2014-2019
  \item 2019-2024
\end{multicols}
\end{enumerate}
\vspace{0.15cm}
\item ¿Qué impulsó Mauricio Funes?
\begin{enumerate}[label=\textbf{\Alph*)}]
\begin{multicols}{2}
  \item Programas sociales y la salida de tropas de Irak
  \item La dolarización
  \item La privatización
  \item La reforma agraria
  \item La desmilitarización
  \item La reforma electoral
\end{multicols}
\end{enumerate}
\vspace{0.15cm}
\item ¿Quién fue Salvador Sánchez Cerén?
\begin{enumerate}[label=\textbf{\Alph*)}]
\begin{multicols}{2}
  \item Un presidente del FMLN
  \item Un presidente de ARENA
  \item Un presidente del PDC
  \item Un presidente militar
  \item Un presidente independiente
  \item Un presidente del PCN
\end{multicols}
\end{enumerate}
\vspace{0.15cm}
\item ¿En qué período gobernó Salvador Sánchez Cerén?
\begin{enumerate}[label=\textbf{\Alph*)}]
\begin{multicols}{2}
  \item 2014-2019
  \item 2009-2014
  \item 2004-2009
  \item 1999-2004
  \item 1994-1999
  \item 2019-2024
\end{multicols}
\end{enumerate}
\vspace{0.15cm}
\item ¿Quién fue Nayib Bukele?
\begin{enumerate}[label=\textbf{\Alph*)}]
\begin{multicols}{2}
  \item Un presidente de Nuevas Ideas
  \item Un presidente de ARENA
  \item Un presidente del FMLN
  \item Un presidente del PDC
  \item Un presidente militar
  \item Un presidente del PCN
\end{multicols}
\end{enumerate}
\vspace{0.15cm}
\item ¿En qué año asumió Nayib Bukele la presidencia?
\begin{enumerate}[label=\textbf{\Alph*)}]
\begin{multicols}{2}
  \item 2019
  \item 2014
  \item 2009
  \item 2004
  \item 1999
  \item 2024
\end{multicols}
\end{enumerate}
\vspace{0.15cm}
\item ¿Qué impulsó Nayib Bukele?
\begin{enumerate}[label=\textbf{\Alph*)}]
\begin{multicols}{2}
  \item La seguridad, la tecnología y la imagen del país
  \item La dolarización
  \item La privatización
  \item La reforma agraria
  \item La desmilitarización
  \item La reforma electoral
\end{multicols}
\end{enumerate}
\vspace{0.15cm}
\item ¿Qué es el régimen de excepción?
\begin{enumerate}[label=\textbf{\Alph*)}]
\begin{multicols}{2}
  \item Una medida que suspende garantías constitucionales para combatir la violencia
  \item Una medida que amplía las garantías constitucionales
  \item Una medida que elimina las garantías constitucionales
  \item Una medida que reforma las garantías constitucionales
  \item Una medida que protege las garantías constitucionales
  \item Una medida que ignora las garantías constitucionales
\end{multicols}
\end{enumerate}
\vspace{0.15cm}
\item ¿Qué es la guerra contra las pandillas?
\begin{enumerate}[label=\textbf{\Alph*)}]
\begin{multicols}{2}
  \item Una estrategia de seguridad para combatir a las maras
  \item Una estrategia de diálogo con las maras
  \item Una estrategia de negociación con las maras
  \item Una estrategia de integración de las maras
  \item Una estrategia de educación de las maras
  \item Una estrategia de prevención de las maras
\end{multicols}
\end{enumerate}
\vspace{0.15cm}
\item ¿Quién fue Armando Calderón Sol?
\begin{enumerate}[label=\textbf{\Alph*)}]
\begin{multicols}{2}
  \item Un presidente de ARENA
  \item Un presidente del FMLN
  \item Un presidente del PDC
  \item Un presidente militar
  \item Un presidente independiente
  \item Un presidente del PCN
\end{multicols}
\end{enumerate}
\vspace{0.15cm}
\item ¿En qué período gobernó Armando Calderón Sol?
\begin{enumerate}[label=\textbf{\Alph*)}]
\begin{multicols}{2}
  \item 1994-1999
  \item 1989-1994
  \item 1999-2004
  \item 2004-2009
  \item 2009-2014
  \item 2014-2019
\end{multicols}
\end{enumerate}
\vspace{0.15cm}
\item ¿Qué impulsó Armando Calderón Sol?
\begin{enumerate}[label=\textbf{\Alph*)}]
\begin{multicols}{2}
  \item La privatización de empresas estatales
  \item La nacionalización de la banca
  \item La reforma agraria
  \item La desmilitarización
  \item La reforma electoral
  \item La reforma judicial
\end{multicols}
\end{enumerate}
\vspace{0.15cm}
\item ¿Quién fue Francisco Flores?
\begin{enumerate}[label=\textbf{\Alph*)}]
\begin{multicols}{2}
  \item Un presidente de ARENA
  \item Un presidente del FMLN
  \item Un presidente del PDC
  \item Un presidente militar
  \item Un presidente independiente
  \item Un presidente del PCN
\end{multicols}
\end{enumerate}
\vspace{0.15cm}
\item ¿En qué período gobernó Francisco Flores?
\begin{enumerate}[label=\textbf{\Alph*)}]
\begin{multicols}{2}
  \item 1999-2004
  \item 1994-1999
  \item 1989-1994
  \item 2004-2009
  \item 2009-2014
  \item 2014-2019
\end{multicols}
\end{enumerate}
\vspace{0.15cm}
\item ¿Qué impulsó Francisco Flores?
\begin{enumerate}[label=\textbf{\Alph*)}]
\begin{multicols}{2}
  \item La dolarización de la economía
  \item La privatización de empresas
  \item La nacionalización de la banca
  \item La reforma agraria
  \item La desmilitarización
  \item La reforma electoral
\end{multicols}
\end{enumerate}
\vspace{0.15cm}
\item ¿En qué año se dolarizó la economía salvadoreña?
\begin{enumerate}[label=\textbf{\Alph*)}]
\begin{multicols}{2}
  \item 2001
  \item 1999
  \item 2004
  \item 1994
  \item 2009
  \item 2014
\end{multicols}
\end{enumerate}
\vspace{0.15cm}
\item ¿Quién fue Antonio Saca?
\begin{enumerate}[label=\textbf{\Alph*)}]
\begin{multicols}{2}
  \item Un presidente de ARENA
  \item Un presidente del FMLN
  \item Un presidente del PDC
  \item Un presidente militar
  \item Un presidente independiente
  \item Un presidente del PCN
\end{multicols}
\end{enumerate}
\vspace{0.15cm}
\item ¿En qué período gobernó Antonio Saca?
\begin{enumerate}[label=\textbf{\Alph*)}]
\begin{multicols}{2}
  \item 2004-2009
  \item 1999-2004
  \item 1994-1999
  \item 1989-1994
  \item 2009-2014
  \item 2014-2019
\end{multicols}
\end{enumerate}
\vspace{0.15cm}
\item ¿Qué impulsó Antonio Saca?
\begin{enumerate}[label=\textbf{\Alph*)}]
\begin{multicols}{2}
  \item La firma del Tratado de Libre Comercio con Estados Unidos
  \item La dolarización
  \item La privatización
  \item La reforma agraria
  \item La desmilitarización
  \item La reforma electoral
\end{multicols}
\end{enumerate}
\vspace{0.15cm}
\item ¿Quién fue Mauricio Funes?
\begin{enumerate}[label=\textbf{\Alph*)}]
\begin{multicols}{2}
  \item El primer presidente del FMLN
  \item Un presidente de ARENA
  \item Un presidente del PDC
  \item Un presidente militar
  \item Un presidente independiente
  \item Un presidente del PCN
\end{multicols}
\end{enumerate}
\vspace{0.15cm}
\item ¿En qué período gobernó Mauricio Funes?
\begin{enumerate}[label=\textbf{\Alph*)}]
\begin{multicols}{2}
  \item 2009-2014
  \item 2004-2009
  \item 1999-2004
  \item 1994-1999
  \item 2014-2019
  \item 2019-2024
\end{multicols}
\end{enumerate}
\vspace{0.15cm}
\item ¿Qué impulsó Mauricio Funes?
\begin{enumerate}[label=\textbf{\Alph*)}]
\begin{multicols}{2}
  \item Programas sociales y la salida de tropas de Irak
  \item La dolarización
  \item La privatización
  \item La reforma agraria
  \item La desmilitarización
  \item La reforma electoral
\end{multicols}
\end{enumerate}
\vspace{0.15cm}
\item ¿Quién fue Salvador Sánchez Cerén?
\begin{enumerate}[label=\textbf{\Alph*)}]
\begin{multicols}{2}
  \item Un presidente del FMLN
  \item Un presidente de ARENA
  \item Un presidente del PDC
  \item Un presidente militar
  \item Un presidente independiente
  \item Un presidente del PCN
\end{multicols}
\end{enumerate}
\vspace{0.15cm}
\item ¿En qué período gobernó Salvador Sánchez Cerén?
\begin{enumerate}[label=\textbf{\Alph*)}]
\begin{multicols}{2}
  \item 2014-2019
  \item 2009-2014
  \item 2004-2009
  \item 1999-2004
  \item 1994-1999
  \item 2019-2024
\end{multicols}
\end{enumerate}
\vspace{0.15cm}
\item ¿Quién fue Nayib Bukele?
\begin{enumerate}[label=\textbf{\Alph*)}]
\begin{multicols}{2}
  \item Un presidente de Nuevas Ideas
  \item Un presidente de ARENA
  \item Un presidente del FMLN
  \item Un presidente del PDC
  \item Un presidente militar
  \item Un presidente del PCN
\end{multicols}
\end{enumerate}
\vspace{0.15cm}
\item ¿En qué año asumió Nayib Bukele la presidencia?
\begin{enumerate}[label=\textbf{\Alph*)}]
\begin{multicols}{2}
  \item 2019
  \item 2014
  \item 2009
  \item 2004
  \item 1999
  \item 2024
\end{multicols}
\end{enumerate}
\vspace{0.15cm}
\item ¿Qué impulsó Nayib Bukele?
\begin{enumerate}[label=\textbf{\Alph*)}]
\begin{multicols}{2}
  \item La seguridad, la tecnología y la imagen del país
  \item La dolarización
  \item La privatización
  \item La reforma agraria
  \item La desmilitarización
  \item La reforma electoral
\end{multicols}
\end{enumerate}
\vspace{0.15cm}
\item ¿Qué es el régimen de excepción?
\begin{enumerate}[label=\textbf{\Alph*)}]
\begin{multicols}{2}
  \item Una medida que suspende garantías constitucionales para combatir la violencia
  \item Una medida que amplía las garantías constitucionales
  \item Una medida que elimina las garantías constitucionales
  \item Una medida que reforma las garantías constitucionales
  \item Una medida que protege las garantías constitucionales
  \item Una medida que ignora las garantías constitucionales
\end{multicols}
\end{enumerate}
\vspace{0.15cm}
\item ¿Qué es la guerra contra las pandillas?
\begin{enumerate}[label=\textbf{\Alph*)}]
\begin{multicols}{2}
  \item Una estrategia de seguridad para combatir a las maras
  \item Una estrategia de diálogo con las maras
  \item Una estrategia de negociación con las maras
  \item Una estrategia de integración de las maras
  \item Una estrategia de educación de las maras
  \item Una estrategia de prevención de las maras
\end{multicols}
\end{enumerate}
\vspace{0.15cm}
\item ¿Quién fue Armando Calderón Sol?
\begin{enumerate}[label=\textbf{\Alph*)}]
\begin{multicols}{2}
  \item Un presidente de ARENA
  \item Un presidente del FMLN
  \item Un presidente del PDC
  \item Un presidente militar
  \item Un presidente independiente
  \item Un presidente del PCN
\end{multicols}
\end{enumerate}
\vspace{0.15cm}
\item ¿En qué período gobernó Armando Calderón Sol?
\begin{enumerate}[label=\textbf{\Alph*)}]
\begin{multicols}{2}
  \item 1994-1999
  \item 1989-1994
  \item 1999-2004
  \item 2004-2009
  \item 2009-2014
  \item 2014-2019
\end{multicols}
\end{enumerate}
\vspace{0.15cm}
\item ¿Qué impulsó Armando Calderón Sol?
\begin{enumerate}[label=\textbf{\Alph*)}]
\begin{multicols}{2}
  \item La privatización de empresas estatales
  \item La nacionalización de la banca
  \item La reforma agraria
  \item La desmilitarización
  \item La reforma electoral
  \item La reforma judicial
\end{multicols}
\end{enumerate}
\vspace{0.15cm}
\item ¿Quién fue Francisco Flores?
\begin{enumerate}[label=\textbf{\Alph*)}]
\begin{multicols}{2}
  \item Un presidente de ARENA
  \item Un presidente del FMLN
  \item Un presidente del PDC
  \item Un presidente militar
  \item Un presidente independiente
  \item Un presidente del PCN
\end{multicols}
\end{enumerate}
\vspace{0.15cm}
\item ¿En qué período gobernó Francisco Flores?
\begin{enumerate}[label=\textbf{\Alph*)}]
\begin{multicols}{2}
  \item 1999-2004
  \item 1994-1999
  \item 1989-1994
  \item 2004-2009
  \item 2009-2014
  \item 2014-2019
\end{multicols}
\end{enumerate}
\vspace{0.15cm}
\item ¿Qué impulsó Francisco Flores?
\begin{enumerate}[label=\textbf{\Alph*)}]
\begin{multicols}{2}
  \item La dolarización de la economía
  \item La privatización de empresas
  \item La nacionalización de la banca
  \item La reforma agraria
  \item La desmilitarización
  \item La reforma electoral
\end{multicols}
\end{enumerate}
\vspace{0.15cm}
\item ¿En qué año se dolarizó la economía salvadoreña?
\begin{enumerate}[label=\textbf{\Alph*)}]
\begin{multicols}{2}
  \item 2001
  \item 1999
  \item 2004
  \item 1994
  \item 2009
  \item 2014
\end{multicols}
\end{enumerate}
\vspace{0.15cm}
\item ¿Quién fue Antonio Saca?
\begin{enumerate}[label=\textbf{\Alph*)}]
\begin{multicols}{2}
  \item Un presidente de ARENA
  \item Un presidente del FMLN
  \item Un presidente del PDC
  \item Un presidente militar
  \item Un presidente independiente
  \item Un presidente del PCN
\end{multicols}
\end{enumerate}
\vspace{0.15cm}
\item ¿En qué período gobernó Antonio Saca?
\begin{enumerate}[label=\textbf{\Alph*)}]
\begin{multicols}{2}
  \item 2004-2009
  \item 1999-2004
  \item 1994-1999
  \item 1989-1994
  \item 2009-2014
  \item 2014-2019
\end{multicols}
\end{enumerate}
\vspace{0.15cm}
\item ¿Qué impulsó Antonio Saca?
\begin{enumerate}[label=\textbf{\Alph*)}]
\begin{multicols}{2}
  \item La firma del Tratado de Libre Comercio con Estados Unidos
  \item La dolarización
  \item La privatización
  \item La reforma agraria
  \item La desmilitarización
  \item La reforma electoral
\end{multicols}
\end{enumerate}
\vspace{0.15cm}
\item ¿Quién fue Mauricio Funes?
\begin{enumerate}[label=\textbf{\Alph*)}]
\begin{multicols}{2}
  \item El primer presidente del FMLN
  \item Un presidente de ARENA
  \item Un presidente del PDC
  \item Un presidente militar
  \item Un presidente independiente
  \item Un presidente del PCN
\end{multicols}
\end{enumerate}
\vspace{0.15cm}
\item ¿En qué período gobernó Mauricio Funes?
\begin{enumerate}[label=\textbf{\Alph*)}]
\begin{multicols}{2}
  \item 2009-2014
  \item 2004-2009
  \item 1999-2004
  \item 1994-1999
  \item 2014-2019
  \item 2019-2024
\end{multicols}
\end{enumerate}
\vspace{0.15cm}
\item ¿Qué impulsó Mauricio Funes?
\begin{enumerate}[label=\textbf{\Alph*)}]
\begin{multicols}{2}
  \item Programas sociales y la salida de tropas de Irak
  \item La dolarización
  \item La privatización
  \item La reforma agraria
  \item La desmilitarización
  \item La reforma electoral
\end{multicols}
\end{enumerate}
\vspace{0.15cm}
\item ¿Quién fue Salvador Sánchez Cerén?
\begin{enumerate}[label=\textbf{\Alph*)}]
\begin{multicols}{2}
  \item Un presidente del FMLN
  \item Un presidente de ARENA
  \item Un presidente del PDC
  \item Un presidente militar
  \item Un presidente independiente
  \item Un presidente del PCN
\end{multicols}
\end{enumerate}
\vspace{0.15cm}
\item ¿En qué período gobernó Salvador Sánchez Cerén?
\begin{enumerate}[label=\textbf{\Alph*)}]
\begin{multicols}{2}
  \item 2014-2019
  \item 2009-2014
  \item 2004-2009
  \item 1999-2004
  \item 1994-1999
  \item 2019-2024
\end{multicols}
\end{enumerate}
\vspace{0.15cm}
\item ¿Quién fue Nayib Bukele?
\begin{enumerate}[label=\textbf{\Alph*)}]
\begin{multicols}{2}
  \item Un presidente de Nuevas Ideas
  \item Un presidente de ARENA
  \item Un presidente del FMLN
  \item Un presidente del PDC
  \item Un presidente militar
  \item Un presidente del PCN
\end{multicols}
\end{enumerate}
\vspace{0.15cm}
\item ¿En qué año asumió Nayib Bukele la presidencia?
\begin{enumerate}[label=\textbf{\Alph*)}]
\begin{multicols}{2}
  \item 2019
  \item 2014
  \item 2009
  \item 2004
  \item 1999
  \item 2024
\end{multicols}
\end{enumerate}
\vspace{0.15cm}
\item ¿Qué impulsó Nayib Bukele?
\begin{enumerate}[label=\textbf{\Alph*)}]
\begin{multicols}{2}
  \item La seguridad, la tecnología y la imagen del país
  \item La dolarización
  \item La privatización
  \item La reforma agraria
  \item La desmilitarización
  \item La reforma electoral
\end{multicols}
\end{enumerate}
\vspace{0.15cm}
\item ¿Qué es el régimen de excepción?
\begin{enumerate}[label=\textbf{\Alph*)}]
\begin{multicols}{2}
  \item Una medida que suspende garantías constitucionales para combatir la violencia
  \item Una medida que amplía las garantías constitucionales
  \item Una medida que elimina las garantías constitucionales
  \item Una medida que reforma las garantías constitucionales
  \item Una medida que protege las garantías constitucionales
  \item Una medida que ignora las garantías constitucionales
\end{multicols}
\end{enumerate}
\vspace{0.15cm}
\item ¿Qué es la guerra contra las pandillas?
\begin{enumerate}[label=\textbf{\Alph*)}]
\begin{multicols}{2}
  \item Una estrategia de seguridad para combatir a las maras
  \item Una estrategia de diálogo con las maras
  \item Una estrategia de negociación con las maras
  \item Una estrategia de integración de las maras
  \item Una estrategia de educación de las maras
  \item Una estrategia de prevención de las maras
\end{multicols}
\end{enumerate}
\vspace{0.15cm}
\end{enumerate}\n
\vspace{0.4cm}
\begin{tcolorbox}[colback=gray!5!white,colframe=gray!60!black,title=\textbf{SOLUCIONARIO -- 4.1 Historia de El Salvador}]
\begin{multicols}{4}
\begin{enumerate}[label=\textbf{\arabic*.}]
  \item \textbf{A}
  \item \textbf{A}
  \item \textbf{A}
  \item \textbf{A}
  \item \textbf{A}
  \item \textbf{A}
  \item \textbf{A}
  \item \textbf{A}
  \item \textbf{A}
  \item \textbf{A}
  \item \textbf{A}
  \item \textbf{A}
  \item \textbf{A}
  \item \textbf{A}
  \item \textbf{A}
  \item \textbf{A}
  \item \textbf{A}
  \item \textbf{A}
  \item \textbf{A}
  \item \textbf{A}
  \item \textbf{A}
  \item \textbf{A}
  \item \textbf{A}
  \item \textbf{A}
  \item \textbf{A}
  \item \textbf{A}
  \item \textbf{A}
  \item \textbf{A}
  \item \textbf{A}
  \item \textbf{A}
  \item \textbf{A}
  \item \textbf{A}
  \item \textbf{A}
  \item \textbf{A}
  \item \textbf{A}
  \item \textbf{A}
  \item \textbf{A}
  \item \textbf{A}
  \item \textbf{A}
  \item \textbf{A}
  \item \textbf{A}
  \item \textbf{A}
  \item \textbf{A}
  \item \textbf{A}
  \item \textbf{A}
  \item \textbf{A}
  \item \textbf{A}
  \item \textbf{A}
  \item \textbf{A}
  \item \textbf{A}
  \item \textbf{A}
  \item \textbf{A}
  \item \textbf{A}
  \item \textbf{A}
  \item \textbf{A}
  \item \textbf{A}
  \item \textbf{A}
  \item \textbf{A}
  \item \textbf{A}
  \item \textbf{A}
  \item \textbf{A}
  \item \textbf{A}
  \item \textbf{A}
  \item \textbf{A}
  \item \textbf{A}
  \item \textbf{A}
  \item \textbf{A}
  \item \textbf{A}
  \item \textbf{A}
  \item \textbf{A}
  \item \textbf{A}
  \item \textbf{A}
  \item \textbf{A}
  \item \textbf{A}
  \item \textbf{A}
  \item \textbf{A}
  \item \textbf{A}
  \item \textbf{A}
  \item \textbf{A}
  \item \textbf{A}
  \item \textbf{A}
  \item \textbf{A}
  \item \textbf{A}
  \item \textbf{A}
  \item \textbf{A}
  \item \textbf{A}
  \item \textbf{A}
  \item \textbf{A}
  \item \textbf{A}
  \item \textbf{A}
  \item \textbf{A}
  \item \textbf{A}
  \item \textbf{A}
  \item \textbf{A}
  \item \textbf{A}
  \item \textbf{A}
  \item \textbf{A}
  \item \textbf{A}
  \item \textbf{A}
  \item \textbf{A}
  \item \textbf{A}
  \item \textbf{A}
  \item \textbf{A}
  \item \textbf{A}
  \item \textbf{A}
  \item \textbf{A}
  \item \textbf{A}
  \item \textbf{A}
  \item \textbf{A}
  \item \textbf{A}
  \item \textbf{A}
  \item \textbf{A}
  \item \textbf{A}
  \item \textbf{A}
  \item \textbf{A}
  \item \textbf{A}
  \item \textbf{A}
  \item \textbf{A}
  \item \textbf{A}
  \item \textbf{A}
  \item \textbf{A}
  \item \textbf{A}
  \item \textbf{A}
  \item \textbf{A}
  \item \textbf{A}
  \item \textbf{A}
  \item \textbf{A}
  \item \textbf{A}
  \item \textbf{A}
  \item \textbf{A}
  \item \textbf{A}
  \item \textbf{A}
  \item \textbf{A}
  \item \textbf{A}
  \item \textbf{A}
  \item \textbf{A}
  \item \textbf{A}
  \item \textbf{A}
  \item \textbf{A}
  \item \textbf{A}
  \item \textbf{A}
  \item \textbf{A}
  \item \textbf{A}
  \item \textbf{A}
  \item \textbf{A}
  \item \textbf{A}
  \item \textbf{A}
  \item \textbf{A}
  \item \textbf{A}
  \item \textbf{A}
  \item \textbf{A}
  \item \textbf{A}
  \item \textbf{A}
  \item \textbf{A}
  \item \textbf{A}
  \item \textbf{A}
  \item \textbf{A}
  \item \textbf{A}
  \item \textbf{A}
  \item \textbf{A}
  \item \textbf{A}
  \item \textbf{A}
  \item \textbf{A}
  \item \textbf{A}
  \item \textbf{A}
  \item \textbf{A}
  \item \textbf{A}
  \item \textbf{A}
  \item \textbf{A}
  \item \textbf{A}
  \item \textbf{A}
  \item \textbf{A}
  \item \textbf{A}
  \item \textbf{A}
  \item \textbf{A}
  \item \textbf{A}
  \item \textbf{A}
  \item \textbf{A}
  \item \textbf{A}
  \item \textbf{A}
  \item \textbf{A}
  \item \textbf{A}
  \item \textbf{A}
  \item \textbf{A}
  \item \textbf{A}
  \item \textbf{A}
  \item \textbf{A}
  \item \textbf{A}
  \item \textbf{A}
  \item \textbf{A}
  \item \textbf{A}
  \item \textbf{A}
  \item \textbf{A}
  \item \textbf{A}
  \item \textbf{A}
  \item \textbf{A}
  \item \textbf{A}
  \item \textbf{A}
  \item \textbf{A}
  \item \textbf{A}
  \item \textbf{A}
  \item \textbf{A}
  \item \textbf{A}
  \item \textbf{A}
  \item \textbf{A}
  \item \textbf{A}
  \item \textbf{A}
  \item \textbf{A}
  \item \textbf{A}
  \item \textbf{A}
  \item \textbf{A}
  \item \textbf{A}
  \item \textbf{A}
  \item \textbf{A}
  \item \textbf{A}
  \item \textbf{A}
  \item \textbf{A}
  \item \textbf{A}
  \item \textbf{A}
  \item \textbf{A}
  \item \textbf{A}
  \item \textbf{A}
  \item \textbf{A}
  \item \textbf{A}
  \item \textbf{A}
  \item \textbf{A}
  \item \textbf{A}
  \item \textbf{A}
  \item \textbf{A}
  \item \textbf{A}
  \item \textbf{A}
  \item \textbf{A}
  \item \textbf{A}
  \item \textbf{A}
  \item \textbf{A}
  \item \textbf{A}
  \item \textbf{A}
  \item \textbf{A}
  \item \textbf{A}
  \item \textbf{A}
  \item \textbf{A}
  \item \textbf{A}
  \item \textbf{A}
  \item \textbf{A}
  \item \textbf{A}
  \item \textbf{A}
  \item \textbf{A}
  \item \textbf{A}
  \item \textbf{A}
  \item \textbf{A}
  \item \textbf{A}
  \item \textbf{A}
  \item \textbf{A}
  \item \textbf{A}
  \item \textbf{A}
  \item \textbf{A}
  \item \textbf{A}
  \item \textbf{A}
  \item \textbf{A}
  \item \textbf{A}
  \item \textbf{A}
  \item \textbf{A}
  \item \textbf{A}
  \item \textbf{A}
  \item \textbf{A}
  \item \textbf{A}
  \item \textbf{A}
  \item \textbf{A}
  \item \textbf{A}
  \item \textbf{A}
  \item \textbf{A}
  \item \textbf{A}
  \item \textbf{A}
  \item \textbf{A}
  \item \textbf{A}
  \item \textbf{A}
  \item \textbf{A}
  \item \textbf{A}
  \item \textbf{A}
  \item \textbf{A}
  \item \textbf{A}
  \item \textbf{A}
  \item \textbf{A}
  \item \textbf{A}
  \item \textbf{A}
  \item \textbf{A}
  \item \textbf{A}
  \item \textbf{A}
  \item \textbf{A}
  \item \textbf{A}
  \item \textbf{A}
  \item \textbf{A}
  \item \textbf{A}
  \item \textbf{A}
  \item \textbf{A}
  \item \textbf{A}
  \item \textbf{A}
  \item \textbf{A}
  \item \textbf{A}
  \item \textbf{A}
\end{enumerate}
\end{multicols}
\end{tcolorbox}
\newpage\n
\subsection{4.2 Contexto Mundial}

\subsubsection*{Módulo 1: Reactivos 1 - 50}

\begin{enumerate}[label=\textbf{\arabic*.}]
\setcounter{enumi}{0}
\item ¿Qué conferencia celebrada en el año $1945$ determinó la división de Alemania en cuatro zonas de ocupación militar tras la derrota del Eje en la Segunda Guerra Mundial?
\begin{enumerate}[label=\textbf{\Alph*)}]
\begin{multicols}{2}
  \item Conferencia de Bretton Woods
  \item Conferencia de Yalta
  \item Conferencia de Bandung
  \item Conferencia de Potsdam
  \item Tratado de Versalles
  \item Conferencia de San Francisco
\end{multicols}
\end{enumerate}
\vspace{0.15cm}
\item El plan económico diseñado por Estados Unidos para la reconstrucción de Europa Occidental y la contención del comunismo, el cual inyectó aproximadamente $\$13 \times 10^{9}$ dólares, se denomina:
\begin{enumerate}[label=\textbf{\Alph*)}]
\begin{multicols}{2}
  \item Doctrina Truman
  \item Plan Dawes
  \item Plan Marshall
  \item Plan Molotov
  \item COMECON
  \item Nuevo Trato (New Deal)
\end{multicols}
\end{enumerate}
\vspace{0.15cm}
\item ¿Cuál fue la alianza militar intergubernamental creada en el año $1949$, fundamentada en el principio de defensa colectiva frente a la expansión de la Unión Soviética?
\begin{enumerate}[label=\textbf{\Alph*)}]
\begin{multicols}{2}
  \item Pacto de Varsovia
  \item Organización de las Naciones Unidas (ONU)
  \item Organización del Tratado del Atlántico Norte (OTAN)
  \item Sociedad de Naciones
  \item Organización de Estados Americanos (OEA)
  \item Movimiento de Países No Alineados
\end{multicols}
\end{enumerate}
\vspace{0.15cm}
\item En respuesta a la alianza militar de Occidente occidental, la Unión Soviética y sus Estados satélites de Europa del Este firmaron en $1955$ un tratado de asistencia mutua conocido como:
\begin{enumerate}[label=\textbf{\Alph*)}]
\begin{multicols}{2}
  \item Kominform
  \item Pacto de Varsovia
  \item COMECON
  \item Tratado de Roma
  \item Pacto de Acero
  \item Tratado de Moscú
\end{multicols}
\end{enumerate}
\vspace{0.15cm}
\item Durante la Guerra Fría, la península de Corea fue dividida en dos naciones bajo zonas de influencia opuestas. ¿Qué línea de latitud geográfica marcó esta frontera geopolítica?
\begin{enumerate}[label=\textbf{\Alph*)}]
\begin{multicols}{2}
  \item Paralelo $17^{\circ}$ Norte
  \item Paralelo $38^{\circ}$ Norte
  \item Meridiano de Greenwich ($0^{\circ}$)
  \item Paralelo $24^{\circ}$ Sur
  \item Trópico de Cáncer ($23.5^{\circ}$ Norte)
  \item Paralelo $50^{\circ}$ Norte
\end{multicols}
\end{enumerate}
\vspace{0.15cm}
\item Los Acuerdos de Bretton Woods ($1944$) rediseñaron la arquitectura financiera global, estableciendo la paridad del dólar con el oro en una tasa de $\$35 \text{ USD} = 1 \text{ onza}$. ¿Qué organismos internacionales nacieron de esta conferencia?
\begin{enumerate}[label=\textbf{\Alph*)}]
\begin{multicols}{2}
  \item La OMC y la OIT
  \item El Banco Mundial (BM) y el Fondo Monetario Internacional (FMI)
  \item La CEPAL y el BID
  \item La Liga de las Naciones y la Cruz Roja
  \item El COMECON y la OTAN
  \item El Tribunal de La Haya y la UNESCO
\end{multicols}
\end{enumerate}
\vspace{0.15cm}
\item El proceso de Descolonización de Asia y África ganó impulso tras la Segunda Guerra Mundial. ¿En qué año declaró la India su independencia del Imperio Británico bajo el liderazgo de Mahatma Gandhi y el Partido del Congreso?
\begin{enumerate}[label=\textbf{\Alph*)}]
\begin{multicols}{2}
  \item $1945$
  \item $1947$
  \item $1950$
  \item $1960$
  \item $1954$
  \item $1962$
\end{multicols}
\end{enumerate}
\vspace{0.15cm}
\item El demógrafo francés Alfred Sauvy acuñó en $1952$ un término para describir a los países subdesarrollados que no pertenecían ni al bloque capitalista ni al bloque comunista. ¿Cuál fue este término?
\begin{enumerate}[label=\textbf{\Alph*)}]
\begin{multicols}{2}
  \item Países Periféricos
  \item El Cuarto Mundo
  \item El Tercer Mundo
  \item Economías Emergentes
  \item Países No Alineados
  \item El Sur Global
\end{multicols}
\end{enumerate}
\vspace{0.15cm}
\item En la economía política del $\text{Siglo XX}$, la Comisión Económica para América Latina y el Caribe (CEPAL) desarrolló una teoría que explicaba el subdesarrollo como un problema estructural asimétrico entre naciones productoras de materias primas y naciones industrializadas. ¿Cómo se llama esta teoría?
\begin{enumerate}[label=\textbf{\Alph*)}]
\begin{multicols}{2}
  \item Teoría de la Ventaja Comparativa
  \item Teoría Neoliberal
  \item Teoría de la Dependencia
  \item Teoría del Valor-Trabajo
  \item Consenso de Washington
  \item Teoría Monetarista
\end{multicols}
\end{enumerate}
\vspace{0.15cm}
\item La Conferencia de Bandung ($1955$) reunió a líderes de $29$ Estados asiáticos y africanos de reciente independencia. ¿Cuál fue el principal movimiento político internacional que germinó a partir de este encuentro?
\begin{enumerate}[label=\textbf{\Alph*)}]
\begin{multicols}{2}
  \item El Movimiento Panafricano
  \item La Liga Árabe
  \item El Movimiento de Países No Alineados (MNOAL)
  \item La Organización de la Unidad Africana (OUA)
  \item El Bloque Sino-Soviético
  \item El G-77
\end{multicols}
\end{enumerate}
\vspace{0.15cm}
\item La "Carrera Espacial" fue un escenario de competencia tecnológica. ¿Qué satélite artificial de la historia, lanzado en $1957$, otorgó a la URSS la primera victoria en esta contienda?
\begin{enumerate}[label=\textbf{\Alph*)}]
\begin{multicols}{2}
  \item Apollo 11
  \item Explorer 1
  \item Vostok 1
  \item Sputnik 1
  \item Soyuz 2
  \item Telstar
\end{multicols}
\end{enumerate}
\vspace{0.15cm}
\item ¿En qué año se desencadenó la Crisis de los Misiles en Cuba, el punto de máxima tensión nuclear donde la probabilidad de una guerra atómica $P(\text{guerra}) \to 1$ estuvo a escasos días de materializarse?
\begin{enumerate}[label=\textbf{\Alph*)}]
\begin{multicols}{2}
  \item $1959$
  \item $1961$
  \item $1962$
  \item $1965$
  \item $1970$
  \item $1989$
\end{multicols}
\end{enumerate}
\vspace{0.15cm}
\item Durante la Guerra de Vietnam, la frontera entre Vietnam del Norte (comunista) y Vietnam del Sur (capitalista) se estableció provisionalmente en los Acuerdos de Ginebra de $1954$ sobre la siguiente línea geográfica:
\begin{enumerate}[label=\textbf{\Alph*)}]
\begin{multicols}{2}
  \item Paralelo $38^{\circ}$ Norte
  \item Paralelo $17^{\circ}$ Norte
  \item Río Mekong
  \item Golfo de Tonkín
  \item Paralelo $24^{\circ}$ Sur
  \item Trópico de Capricornio
\end{multicols}
\end{enumerate}
\vspace{0.15cm}
\item La política exterior estadounidense formulada por Dwight D. Eisenhower postulaba que si un Estado en una región caía bajo el comunismo, los países vecinos también caerían en secuencia. ¿Cómo se denominó este postulado analítico?
\begin{enumerate}[label=\textbf{\Alph*)}]
\begin{multicols}{2}
  \item Doctrina del Destino Manifiesto
  \item Teoría del Dominó
  \item Estrategia de Contención Estricta
  \item Doctrina de Represalia Masiva
  \item Plan de Coexistencia Pacífica
  \item Equilibrio del Terror
\end{multicols}
\end{enumerate}
\vspace{0.15cm}
\item El "Equilibrio del Terror" o Destrucción Mutua Asegurada (MAD, por sus siglas en inglés) funcionaba bajo la premisa estratégica militar de que:
\begin{enumerate}[label=\textbf{\Alph*)}]
\begin{multicols}{2}
  \item El uso de armas nucleares por dos bandos opuestos causaría la aniquilación total de ambos atacante y defensor.
  \item Solo la URSS tenía capacidad de un segundo ataque nuclear.
  \item Las guerras subsidiarias (proxy wars) no escalarían a conflictos nucleares si se desarrollaban en el Tercer Mundo.
  \item Estados Unidos desmantelaría su arsenal si la URSS se retiraba de Europa del Este.
  \item El despliegue de misiles intercontinentales (ICBM) sería supervisado por la ONU.
  \item Las armas químicas sustituirían a las nucleares en el $\text{Siglo XXI}$.
\end{multicols}
\end{enumerate}
\vspace{0.15cm}
\item En agosto de $1961$, el bloque comunista inició la construcción de una barrera física que se convirtió en el símbolo máximo del mundo bipolar. ¿Cuál fue esta estructura?
\begin{enumerate}[label=\textbf{\Alph*)}]
\begin{multicols}{2}
  \item La Cortina de Hierro
  \item El Muro de Berlín
  \item La Línea Maginot
  \item El Paralelo $38^{\circ}$
  \item El Puente Aéreo de Alemania
  \item La Frontera de Checkpoint Charlie
\end{multicols}
\end{enumerate}
\vspace{0.15cm}
\item ¿Cómo se llamó el proceso de reestructuración económica introducido en la URSS en $1985$ por Mijaíl Gorbachov, cuyo objetivo era descentralizar la economía socialista introduciendo dinámicas de mercado controlado?
\begin{enumerate}[label=\textbf{\Alph*)}]
\begin{multicols}{2}
  \item Glasnost
  \item Komintern
  \item Nomenklatura
  \item Perestroika
  \item Gosplan
  \item Koljós
\end{multicols}
\end{enumerate}
\vspace{0.15cm}
\item Acompañando a las reformas económicas, Mijaíl Gorbachov implementó una política de "apertura" y transparencia que permitió la libertad de expresión y prensa, debilitando el monopolio del Partido Comunista. ¿Cómo se denominó en ruso?
\begin{enumerate}[label=\textbf{\Alph*)}]
\begin{multicols}{2}
  \item Duma
  \item Pravda
  \item Glasnost
  \item Perestroika
  \item Bolchevique
  \item Soviet
\end{multicols}
\end{enumerate}
\vspace{0.15cm}
\item La crisis del petróleo de $1973$ cuadruplicó los precios del crudo, alterando radicalmente la curva de oferta agregada global. Este evento fue provocado por el embargo petrolero ejecutado por:
\begin{enumerate}[label=\textbf{\Alph*)}]
\begin{multicols}{2}
  \item La Organización de las Naciones Unidas (ONU)
  \item El bloque soviético (COMECON)
  \item La Organización de Países Exportadores de Petróleo (OPEP)
  \item El Fondo Monetario Internacional (FMI)
  \item La OTAN
  \item El Banco Central Europeo
\end{multicols}
\end{enumerate}
\vspace{0.15cm}
\item El embargo petrolero de $1973$ fue una represalia directa de las naciones árabes contra los países occidentales que apoyaron a Israel durante un conflicto armado específico. ¿Qué guerra desencadenó este embargo?
\begin{enumerate}[label=\textbf{\Alph*)}]
\begin{multicols}{2}
  \item La Guerra de los Seis Días
  \item La Guerra de Yom Kipur
  \item La Primera Guerra del Golfo
  \item La Guerra Irán-Irak
  \item La Crisis de Suez
  \item La Guerra de la Independencia de Israel
\end{multicols}
\end{enumerate}
\vspace{0.15cm}
\item A finales del $\text{Siglo XX}$, ¿en qué año ocurrió la caída del Muro de Berlín, evento catalizador que precipitó el colapso de los regímenes de Europa del Este?
\begin{enumerate}[label=\textbf{\Alph*)}]
\begin{multicols}{2}
  \item $1985$
  \item $1988$
  \item $1989$
  \item $1990$
  \item $1991$
  \item $1993$
\end{multicols}
\end{enumerate}
\vspace{0.15cm}
\item La Unión de Repúblicas Socialistas Soviéticas (URSS) se disolvió oficialmente, fragmentándose en $15$ repúblicas independientes, poniendo fin formal a la Guerra Fría el 25 de diciembre del año:
\begin{enumerate}[label=\textbf{\Alph*)}]
\begin{multicols}{2}
  \item $1989$
  \item $1990$
  \item $1991$
  \item $1992$
  \item $1993$
  \item $1995$
\end{multicols}
\end{enumerate}
\vspace{0.15cm}
\item ¿Cuál fue la alianza económica y comercial establecida por la URSS en $1949$ para integrar las economías de los países de Europa del Este bajo el modelo de planificación centralizada estatal?
\begin{enumerate}[label=\textbf{\Alph*)}]
\begin{multicols}{2}
  \item COMECON (CAME)
  \item Pacto de Varsovia
  \item Mercado Común Europeo
  \item Asociación Europea de Libre Comercio (EFTA)
  \item Tratado de Maastricht
  \item Organización de Cooperación de Shanghái
\end{multicols}
\end{enumerate}
\vspace{0.15cm}
\item En el contexto de América Latina, las dictaduras cívico-militares de las décadas de $1970$ y $1980$ ejecutaron una red de inteligencia transnacional coordinada para la persecución de opositores políticos de izquierda. Este plan se denominó:
\begin{enumerate}[label=\textbf{\Alph*)}]
\begin{multicols}{2}
  \item Operación Barbarroja
  \item Plan Marshall Latinoamericano
  \item Operación Cóndor
  \item Alianza para el Progreso
  \item Doctrina de Seguridad Nacional
  \item Operación Causa Justa
\end{multicols}
\end{enumerate}
\vspace{0.15cm}
\item El "Consenso de Washington" ($1989$) fue un decálogo de medidas de política económica promovidas por organismos como el FMI y el Banco Mundial, orientadas a países en desarrollo. ¿Qué modelo económico implementaban estas medidas?
\begin{enumerate}[label=\textbf{\Alph*)}]
\begin{multicols}{2}
  \item Estado de Bienestar (Keynesianismo)
  \item Modelo de Sustitución de Importaciones (ISI)
  \item Neoliberalismo y apertura de mercados
  \item Socialismo del Siglo XXI
  \item Economía Centralizada Planificada
  \item Corporativismo de Estado
\end{multicols}
\end{enumerate}
\vspace{0.15cm}
\item ¿Qué presidente estadounidense propuso en $1983$ la Iniciativa de Defensa Estratégica (SDI), apodada "Star Wars", un escudo espacial para interceptar misiles balísticos intercontinentales soviéticos con tecnología láser?
\begin{enumerate}[label=\textbf{\Alph*)}]
\begin{multicols}{2}
  \item Richard Nixon
  \item Jimmy Carter
  \item Ronald Reagan
  \item George H.W. Bush
  \item John F. Kennedy
  \item Lyndon B. Johnson
\end{multicols}
\end{enumerate}
\vspace{0.15cm}
\item A partir de la década de $1970$, Estados Unidos y la Unión Soviética entraron en una fase de "distensión" caracterizada por acuerdos para limitar el crecimiento exponencial de sus arsenales nucleares. ¿Qué siglas identificaron a estas negociaciones?
\begin{enumerate}[label=\textbf{\Alph*)}]
\begin{multicols}{2}
  \item SALT (I y II)
  \item NAFTA
  \item START (I y II)
  \item SEATO
  \item CENTO
  \item ANZUS
\end{multicols}
\end{enumerate}
\vspace{0.15cm}
\item La Guerra de Vietnam finalizó con la captura de Saigón por el Ejército de Vietnam del Norte y el Frente Nacional de Liberación (Viet Cong). ¿En qué año se consolidó la reunificación del país bajo un régimen comunista?
\begin{enumerate}[label=\textbf{\Alph*)}]
\begin{multicols}{2}
  \item $1968$
  \item $1973$
  \item $1975$
  \item $1979$
  \item $1982$
  \item $1985$
\end{multicols}
\end{enumerate}
\vspace{0.15cm}
\item Durante el año $1968$, un país de Europa del Este intentó implementar el "socialismo con rostro humano", un proceso de liberalización política y social que fue aplastado militarmente por los tanques del Pacto de Varsovia. ¿A qué evento corresponde esta descripción?
\begin{enumerate}[label=\textbf{\Alph*)}]
\begin{multicols}{2}
  \item Revolución Húngara
  \item Primavera de Praga (Checoslovaquia)
  \item Levantamiento de Varsovia (Polonia)
  \item Revolución de Terciopelo
  \item Revolución de los Claveles
  \item Secesión Yugoslava
\end{multicols}
\end{enumerate}
\vspace{0.15cm}
\item La Revolución Islámica de $1979$ derrocó al régimen prooccidental del Sha Mohammad Reza Pahlaví y estableció una república teocrática chiita, alterando el equilibrio geopolítico de Oriente Medio. ¿En qué país ocurrió?
\begin{enumerate}[label=\textbf{\Alph*)}]
\begin{multicols}{2}
  \item Irak
  \item Afganistán
  \item Irán
  \item Arabia Saudita
  \item Siria
  \item Líbano
\end{multicols}
\end{enumerate}
\vspace{0.15cm}
\item En el periodo posterior a la Guerra Fría, ¿qué organización supranacional se constituyó en $1995$ para regular, supervisar y liberalizar el comercio global, reemplazando al Acuerdo General sobre Aranceles Aduaneros y Comercio (GATT)?
\begin{enumerate}[label=\textbf{\Alph*)}]
\begin{multicols}{2}
  \item Organización Mundial de la Salud (OMS)
  \item Banco Mundial (BM)
  \item Organización Mundial del Comercio (OMC)
  \item Foro Económico Mundial (WEF)
  \item Comunidad Económica Europea (CEE)
  \item Asociación Transatlántica de Comercio e Inversión (TTIP)
\end{multicols}
\end{enumerate}
\vspace{0.15cm}
\item El continente africano vivió en $1960$ un hito en su geopolítica cuando $17$ naciones del África subsahariana declararon su emancipación política, logrando ingresar a la ONU. En la historiografía, el año $1960$ se conoce como:
\begin{enumerate}[label=\textbf{\Alph*)}]
\begin{multicols}{2}
  \item El Año de África
  \item La Primavera Árabe
  \item El Renacimiento Panafricano
  \item La Década Perdida
  \item La Era del Imperialismo
  \item La Gran Descolonización
\end{multicols}
\end{enumerate}
\vspace{0.15cm}
\item El periodo conocido como "Macartismo" o la Segunda "Cacería de Brujas" ($1950$-$1954$) en Estados Unidos fue una intensa campaña política impulsada por el senador Joseph McCarthy orientada a:
\begin{enumerate}[label=\textbf{\Alph*)}]
\begin{multicols}{2}
  \item Perseguir judicialmente a los miembros del Ku Klux Klan.
  \item Investigar y encarcelar a supuestos espías y simpatizantes comunistas en instituciones, Hollywood y el gobierno.
  \item Expulsar a inmigrantes indocumentados de origen latinoamericano.
  \item Bloquear las políticas del New Deal argumentando fraude electoral.
  \item Promover los derechos civiles de las minorías afroamericanas.
  \item Investigar la corrupción en las empresas petroleras de Texas.
\end{multicols}
\end{enumerate}
\vspace{0.15cm}
\item ¿Qué conflicto armado se desató en $1979$ cuando el ejército de la URSS intervino militarmente para sostener a un gobierno socialista, resultando en una guerra de guerrillas de una década contra los muyahidines financiados por EE. UU.?
\begin{enumerate}[label=\textbf{\Alph*)}]
\begin{multicols}{2}
  \item La Guerra Civil Siria
  \item La Guerra de Corea
  \item La Guerra de los Balcanes
  \item La Guerra Soviético-Afgana
  \item La Guerra Irán-Irak
  \item La Guerra de Chechenia
\end{multicols}
\end{enumerate}
\vspace{0.15cm}
\item La "Doctrina Truman", proclamada en $1947$, delineó la estrategia geopolítica fundamental de Estados Unidos frente al comunismo, la cual se basaba en el concepto de:
\begin{enumerate}[label=\textbf{\Alph*)}]
\begin{multicols}{2}
  \item Política de apaciguamiento
  \item Contención (Containment)
  \item Aislamiento total (Isolationism)
  \item Reversión armada (Rollback)
  \item Imperialismo formal
  \item No intervención recíproca
\end{multicols}
\end{enumerate}
\vspace{0.15cm}
\item En $1956$, el presidente egipcio Gamal Abdel Nasser nacionalizó una infraestructura logística clave para el comercio global de petróleo, lo que provocó la intervención armada de Reino Unido, Francia e Israel. ¿A qué crisis se hace referencia?
\begin{enumerate}[label=\textbf{\Alph*)}]
\begin{multicols}{2}
  \item Crisis del Estrecho de Ormuz
  \item Crisis del Canal de Panamá
  \item Crisis de Suez
  \item Crisis de los Dardanelos
  \item Guerra del Sinaí
  \item Crisis del Golfo Pérsico
\end{multicols}
\end{enumerate}
\vspace{0.15cm}
\item ¿Qué líder soviético, tras la muerte de Iósif Stalin en $1953$, inició el proceso de "Desestalinización" y postuló la doctrina de la "Coexistencia Pacífica" con el bloque capitalista?
\begin{enumerate}[label=\textbf{\Alph*)}]
\begin{multicols}{2}
  \item Vladimir Lenin
  \item León Trotski
  \item Nikita Jruschov
  \item Leonid Brézhnev
  \item Yuri Andrópov
  \item Mijaíl Gorbachov
\end{multicols}
\end{enumerate}
\vspace{0.15cm}
\item Durante el periodo de entreguerras y el comienzo de la Guerra Fría, John Maynard Keynes formuló un modelo económico en el cual la ecuación de demanda agregada $DA = C + I + G + (X - M)$ requiere que el Estado (G) intervenga mediante:
\begin{enumerate}[label=\textbf{\Alph*)}]
\begin{multicols}{2}
  \item La eliminación total de los impuestos sobre la renta.
  \item La privatización absoluta de las empresas estratégicas.
  \item El incremento del gasto público para estimular la economía durante las recesiones.
  \item La fijación exclusiva del precio del dinero en un patrón oro estricto.
  \item El cese de toda emisión monetaria (déficit cero).
  \item La abolición de las exportaciones e importaciones.
\end{multicols}
\end{enumerate}
\vspace{0.15cm}
\item El sistema de segregación racial institucionalizado legalmente en Sudáfrica que garantizaba el dominio político y económico de la minoría blanca ($20\%$ de la población) sobre la mayoría negra se denominó:
\begin{enumerate}[label=\textbf{\Alph*)}]
\begin{multicols}{2}
  \item Jim Crow
  \item Apartheid
  \item Bantustán
  \item Destino Manifiesto
  \item Supremacismo Afrikáner
  \item Etnocentrismo Institucional
\end{multicols}
\end{enumerate}
\vspace{0.15cm}
\item En $1994$, Sudáfrica celebró sus primeras elecciones democráticas multirraciales, poniendo fin oficial al Apartheid. ¿Quién fue el primer presidente negro de la nación, elegido tras pasar $27$ años en prisión?
\begin{enumerate}[label=\textbf{\Alph*)}]
\begin{multicols}{2}
  \item Desmond Tutu
  \item Steve Biko
  \item Nelson Mandela
  \item F.W. de Klerk
  \item Patrice Lumumba
  \item Kwame Nkrumah
\end{multicols}
\end{enumerate}
\vspace{0.15cm}
\item ¿Qué revolución latinoamericana triunfó en el año $1959$, instaurando posteriormente el primer gobierno de corte socialista y aliado de la Unión Soviética en el Hemisferio Occidental?
\begin{enumerate}[label=\textbf{\Alph*)}]
\begin{multicols}{2}
  \item Revolución Mexicana
  \item Revolución Sandinista en Nicaragua
  \item Revolución Bolivariana
  \item Revolución Cubana
  \item Revolución de Octubre
  \item Revolución de los Claveles
\end{multicols}
\end{enumerate}
\vspace{0.15cm}
\item A finales de la década de $1950$, surgió una profunda ruptura ideológica y geopolítica dentro del bloque comunista, originada por disputas doctrinales sobre la ortodoxia marxista y el liderazgo mundial. Esta división histórica se conoce como:
\begin{enumerate}[label=\textbf{\Alph*)}]
\begin{multicols}{2}
  \item Cisma de Yugoslavia
  \item Ruptura Sino-Soviética (China y URSS)
  \item Conflicto Ruso-Japonés
  \item Crisis de los Urales
  \item División Coreana
  \item Ruptura del Kominform
\end{multicols}
\end{enumerate}
\vspace{0.15cm}
\item Para la medición del desarrollo en el contexto mundial contemporáneo, el Programa de las Naciones Unidas para el Desarrollo (PNUD) creó en $1990$ un índice que integra esperanza de vida, nivel educativo y Producto Interno Bruto (PIB) per cápita. ¿Cómo se llama este indicador?
\begin{enumerate}[label=\textbf{\Alph*)}]
\begin{multicols}{2}
  \item Índice de Gini (IG)
  \item Índice de Pobreza Multidimensional (IPM)
  \item Índice de Desarrollo Humano (IDH)
  \item Producto Nacional Bruto (PNB)
  \item Índice de Precios al Consumidor (IPC)
  \item Coeficiente de Lorenz
\end{multicols}
\end{enumerate}
\vspace{0.15cm}
\item Durante el conflicto armado en El Salvador ($1980$-$1992$), el contexto mundial de la Guerra Fría polarizó la guerra civil. El FMLN recibió apoyo ideológico y logístico del bloque comunista a través de Cuba y Nicaragua, mientras que el gobierno militar/civil salvadoreño fue respaldado financiera y militarmente por:
\begin{enumerate}[label=\textbf{\Alph*)}]
\begin{multicols}{2}
  \item La Unión Europea
  \item La Liga Árabe
  \item La Organización de Estados Americanos
  \item Estados Unidos
  \item La República de China (Taiwán)
  \item La ONU a través de los Cascos Azules
\end{multicols}
\end{enumerate}
\vspace{0.15cm}
\item Tras el fin de la Segunda Guerra Mundial, Japón experimentó un crecimiento económico acelerado apoyado por Estados Unidos, logrando tasas de crecimiento interanual del PIB superiores al $10\%$ durante las décadas de $1950$ y $1960$. A este fenómeno se le conoce como:
\begin{enumerate}[label=\textbf{\Alph*)}]
\begin{multicols}{2}
  \item El Gran Salto Adelante
  \item El Milagro Económico Japonés
  \item La Restauración Meiji
  \item El Plan Dodge
  \item La Perestroika Asiática
  \item El Tigre Asiático Original
\end{multicols}
\end{enumerate}
\vspace{0.15cm}
\item En geopolítica del $\text{Siglo XX}$, ¿qué región asiática experimentó el surgimiento de economías altamente industrializadas orientadas a la exportación, conocidas como los "Cuatro Tigres Asiáticos" (Corea del Sur, Taiwán, Singapur y Hong Kong)?
\begin{enumerate}[label=\textbf{\Alph*)}]
\begin{multicols}{2}
  \item Asia Central
  \item Sudeste Asiático y Asia Oriental
  \item Subcontinente Indio
  \item Medio Oriente
  \item Indochina
  \item Polinesia
\end{multicols}
\end{enumerate}
\vspace{0.15cm}
\item La creación de la Comunidad Económica Europea (CEE) en $1957$ mediante el Tratado de Roma sentó las bases institucionales para la actual Unión Europea. Su principal objetivo original fue:
\begin{enumerate}[label=\textbf{\Alph*)}]
\begin{multicols}{2}
  \item Crear un bloque militar superior a la OTAN.
  \item Establecer un mercado común con libre circulación de personas, bienes, servicios y capitales.
  \item Eliminar la soberanía nacional de los estados miembros de forma inmediata.
  \item Adoptar el Euro como moneda única al día siguiente.
  \item Financiar proyectos espaciales conjuntos para competir con la URSS.
  \item Construir una barrera aduanera contra Estados Unidos.
\end{multicols}
\end{enumerate}
\vspace{0.15cm}
\item ¿Qué tratado internacional firmado en $1968$ buscó limitar la expansión de armas nucleares permitiendo su posesión únicamente a cinco naciones (EE. UU., URSS, Reino Unido, Francia y China) e impulsando el uso pacífico de la energía atómica?
\begin{enumerate}[label=\textbf{\Alph*)}]
\begin{multicols}{2}
  \item Tratado de Kioto
  \item Tratado de Tlatelolco
  \item Tratado de No Proliferación Nuclear (TNP)
  \item Acuerdo de París
  \item Convención de Ginebra
  \item Tratado de Versalles
\end{multicols}
\end{enumerate}
\vspace{0.15cm}
\item El fenómeno de la "Globalización" se aceleró vertiginosamente a finales del $\text{Siglo XX}$ ($1990$s) impulsado por dos vectores fundamentales:
\begin{enumerate}[label=\textbf{\Alph*)}]
\begin{multicols}{2}
  \item La expansión del socialismo y la agricultura a gran escala.
  \item La desregulación financiera (Neoliberalismo) y la revolución de las Tecnologías de la Información y Comunicación (TIC).
  \item El fortalecimiento del proteccionismo arancelario y la minería.
  \item El aumento de las tasas de interés y la construcción de infraestructura ferroviaria.
  \item La colonización de nuevos territorios en África y la energía de vapor.
  \item La Guerra Fría y el nacionalismo económico extremo.
\end{multicols}
\end{enumerate}
\vspace{0.15cm}
\item ¿Qué conflicto bélico internacional estalló en $1990$ tras la invasión iraquí a Kuwait, provocando la formación de una coalición internacional de $34$ países sancionada por la ONU y liderada por EE. UU.?
\begin{enumerate}[label=\textbf{\Alph*)}]
\begin{multicols}{2}
  \item La Guerra Irán-Irak
  \item La Guerra de Afganistán
  \item La Primera Guerra del Golfo (Operación Tormenta del Desierto)
  \item La Guerra de los Seis Días
  \item La Crisis de Suez
  \item La Guerra del Líbano
\end{multicols}
\end{enumerate}
\vspace{0.15cm}
\end{enumerate}\n
\subsubsection*{Módulo 2: Reactivos 51 - 100}

\begin{enumerate}[label=\textbf{\arabic*.}]
\setcounter{enumi}{50}
\item La Organización de las Naciones Unidas (ONU) fue fundada en octubre de 1945 con la firma de la Carta de San Francisco. Su órgano ejecutivo con capacidad de emitir resoluciones vinculantes y autorizar intervenciones militares es el Consejo de Seguridad, el cual otorga poder de veto absoluto a 5 miembros permanentes. ¿Cuáles son estas naciones?
\begin{enumerate}[label=\textbf{\Alph*)}]
\begin{multicols}{2}
  \item EE. UU., URSS (Rusia), Reino Unido, Alemania y Japón.
  \item EE. UU., URSS (Rusia), Reino Unido, Francia y China.
  \item EE. UU., URSS (Rusia), Canadá, Francia y China.
  \item EE. UU., Reino Unido, Francia, Alemania e Italia.
  \item EE. UU., URSS (Rusia), Brasil, India y China.
  \item EE. UU., Reino Unido, Francia, Japón y Australia.
\end{multicols}
\end{enumerate}
\vspace{0.15cm}
\item Tras una prolongada guerra civil, en el año 1949 el Ejército Popular de Liberación venció a las fuerzas nacionalistas del Kuomintang (que huyeron a Taiwán), proclamando la República Popular China. ¿Qué líder histórico encabezó esta revolución comunista?
\begin{enumerate}[label=\textbf{\Alph*)}]
\begin{multicols}{2}
  \item Sun Yat-sen
  \item Chiang Kai-shek
  \item Deng Xiaoping
  \item Mao Zedong
  \item Zhou Enlai
  \item Ho Chi Minh
\end{multicols}
\end{enumerate}
\vspace{0.15cm}
\item Entre 1958 y 1961, el gobierno chino implementó una campaña de medidas económicas, sociales y políticas para transformar rápidamente la base agraria del país en una sociedad industrializada a través de comunas agrícolas. Este plan, que resultó en una hambruna catastrófica, se denominó:
\begin{enumerate}[label=\textbf{\Alph*)}]
\begin{multicols}{2}
  \item La Revolución Cultural
  \item El Gran Salto Adelante
  \item El Plan Quinquenal de Acero
  \item La Nueva Política Económica (NEP)
  \item El Movimiento de Rectificación
  \item La Reforma de los Cien Flores
\end{multicols}
\end{enumerate}
\vspace{0.15cm}
\item En la cúspide de la Carrera Espacial, el 12 de abril de 1961, la Unión Soviética logró otra victoria tecnológica al enviar al primer ser humano al espacio exterior a bordo de la cápsula Vostok 1. ¿Quién fue este cosmonauta?
\begin{enumerate}[label=\textbf{\Alph*)}]
\begin{multicols}{2}
  \item Neil Armstrong
  \item Buzz Aldrin
  \item Yuri Gagarin
  \item Valentina Tereshkova
  \item Alexei Leonov
  \item Laika
\end{multicols}
\end{enumerate}
\vspace{0.15cm}
\item El proyecto tecnológico más ambicioso de Estados Unidos durante la Guerra Fría culminó el 20 de julio de 1969, logrando alunizar con éxito y ganar la carrera lunar soviético-estadounidense. ¿Qué misión logró este hito histórico?
\begin{enumerate}[label=\textbf{\Alph*)}]
\begin{multicols}{2}
  \item Apollo 8
  \item Apollo 11
  \item Gemini 4
  \item Mercury-Redstone 3
  \item Sputnik 2
  \item Explorer 1
\end{multicols}
\end{enumerate}
\vspace{0.15cm}
\item El proceso de descolonización en África tuvo episodios sumamente violentos. Entre 1954 y 1962, el Frente de Liberación Nacional (FLN) libró una guerra asimétrica contra una potencia europea que consideraba a este territorio magrebí como parte integral de su metrópoli. ¿A qué país africano corresponde este proceso independentista?
\begin{enumerate}[label=\textbf{\Alph*)}]
\begin{multicols}{2}
  \item Congo
  \item Angola
  \item Sudáfrica
  \item Argelia
  \item Mozambique
  \item Kenia
\end{multicols}
\end{enumerate}
\vspace{0.15cm}
\item En 1971, el presidente estadounidense Richard Nixon anunció la suspensión unilateral de la convertibilidad del dólar en oro, rompiendo la paridad fija de $\$35 \text{ USD} = 1 \text{ oz}$ de oro establecida en 1944. ¿Cómo se denomina el sistema monetario internacional que colapsó con esta medida?
\begin{enumerate}[label=\textbf{\Alph*)}]
\begin{multicols}{2}
  \item Patrón Plata de Londres
  \item Sistema del Acuerdo de Plaza
  \item Sistema de Bretton Woods
  \item Sistema Monetario Europeo
  \item Consenso de Washington
  \item Patrón Fiduciario Global
\end{multicols}
\end{enumerate}
\vspace{0.15cm}
\item Durante el año 1967, una coalición militar árabe se movilizó contra el Estado de Israel. En un ataque preventivo devastador, Israel destruyó la aviación de sus enemigos y capturó la península del Sinaí, Gaza, Cisjordania, Jerusalén Este y los Altos del Golán en menos de una semana. ¿Qué nombre recibe este conflicto?
\begin{enumerate}[label=\textbf{\Alph*)}]
\begin{multicols}{2}
  \item Guerra de Yom Kipur
  \item Guerra de los Seis Días
  \item Crisis de Suez
  \item Primera Intifada
  \item Guerra de Independencia de Israel
  \item Operación Plomo Fundido
\end{multicols}
\end{enumerate}
\vspace{0.15cm}
\item La resolución diplomática de los conflictos de Oriente Medio dio un giro trascendental en 1978 cuando el presidente de Egipto, Anwar el-Sadat, y el primer ministro de Israel, Menájem Beguin, firmaron un acuerdo de paz mediado por Jimmy Carter. Estos tratados se conocen como:
\begin{enumerate}[label=\textbf{\Alph*)}]
\begin{multicols}{2}
  \item Acuerdos de Oslo
  \item Acuerdos de Camp David
  \item Tratado de Versalles
  \item Protocolos de Madrid
  \item Acuerdos de Ginebra
  \item Resolución 242 de la ONU
\end{multicols}
\end{enumerate}
\vspace{0.15cm}
\item En 1947, la Asamblea General de la ONU aprobó la Resolución 181, la cual recomendaba una reestructuración geopolítica en Oriente Medio debido a las tensiones intercomunitarias tras el Holocausto y la migración masiva. ¿Qué proponía dicha resolución?
\begin{enumerate}[label=\textbf{\Alph*)}]
\begin{multicols}{2}
  \item La entrega de todo el territorio a la Liga Árabe.
  \item La partición del Mandato Británico de Palestina en un Estado judío, un Estado árabe y un régimen internacional para Jerusalén.
  \item La creación de un único Estado binacional secular bajo administración de la OTAN.
  \item La expulsión incondicional del Imperio Británico y la entrega territorial a Turquía.
  \item El desplazamiento de las poblaciones judías hacia la isla de Madagascar.
  \item La prohibición de asentamientos y la creación de un mandato exclusivo del Vaticano.
\end{multicols}
\end{enumerate}
\vspace{0.15cm}
\item Fundada en 1964, esta organización política y paramilitar agrupó diversas facciones palestinas con el objetivo original de liberar el territorio mediante la lucha armada, siendo liderada durante décadas por Yasser Arafat.
\begin{enumerate}[label=\textbf{\Alph*)}]
\begin{multicols}{2}
  \item Hamás
  \item Hezbolá
  \item Organización para la Liberación de Palestina (OLP)
  \item Hermanos Musulmanes
  \item Al-Fatah (como entidad única separada de la alianza)
  \item Yihad Islámica Palestina
\end{multicols}
\end{enumerate}
\vspace{0.15cm}
\item El 11 de septiembre de 1973, la Fuerza Aérea bombardeó el Palacio de La Moneda, concretando un golpe de Estado militar respaldado por la CIA que derrocó al presidente socialista democráticamente electo Salvador Allende. ¿En qué país latinoamericano ocurrió este suceso?
\begin{enumerate}[label=\textbf{\Alph*)}]
\begin{multicols}{2}
  \item Argentina
  \item Brasil
  \item Uruguay
  \item Chile
  \item Perú
  \item Bolivia
\end{multicols}
\end{enumerate}
\vspace{0.15cm}
\item En la década de 1980, la "Teoría de la Disuasión Nuclear" escaló con el concepto de la "Tríada Nuclear". En la lógica militar, poseer una tríada significa que un Estado tiene la capacidad de lanzar ojivas atómicas $W$ mediante tres vectores de entrega simultáneos. ¿Cuáles son estos tres vectores?
\begin{enumerate}[label=\textbf{\Alph*)}]
\begin{multicols}{2}
  \item Portaaviones, satélites militares y misiles de corto alcance.
  \item Misiles balísticos intercontinentales (ICBM) terrestres, bombarderos estratégicos y submarinos de misiles balísticos (SSBN).
  \item Cañones de artillería, drones no tripulados y cruceros acorazados.
  \item Fuerzas especiales, misiles de crucero convencionales y ojivas de neutrones.
  \item Tanques pesados, cazas tácticos y bases espaciales.
  \item Helicópteros artillados, sistemas de defensa antimisiles y minas navales.
\end{multicols}
\end{enumerate}
\vspace{0.15cm}
\item El 2 de abril de 1982, la dictadura militar argentina desembarcó tropas en un archipiélago del Atlántico Sur bajo dominio británico, desatando una guerra convencional que duró 74 días y culminó con la victoria del Reino Unido. ¿Cuál fue este archipiélago?
\begin{enumerate}[label=\textbf{\Alph*)}]
\begin{multicols}{2}
  \item Islas Georgias del Sur
  \item Islas Canarias
  \item Islas Malvinas (Falklands)
  \item Islas Azores
  \item Islas Galápagos
  \item Isla de Pascua
\end{multicols}
\end{enumerate}
\vspace{0.15cm}
\item En julio de 1979, una coalición guerrillera de ideología marxista-leninista y nacionalista derrocó a la dictadura militar de la familia Somoza en Centroamérica, estableciendo una Junta de Gobierno de Reconstrucción Nacional. ¿Cómo se conoce a esta revolución?
\begin{enumerate}[label=\textbf{\Alph*)}]
\begin{multicols}{2}
  \item Revolución Cubana
  \item Revolución Sandinista
  \item Revolución Zapatista
  \item Revolución Bolivariana
  \item Revolución de Octubre
  \item Revolución de los Claveles
\end{multicols}
\end{enumerate}
\vspace{0.15cm}
\item A nivel físico, la Carrera Armamentista de la Guerra Fría evolucionó de las bombas de fisión (bomba atómica) a las bombas termonucleares, las cuales generan una energía destructiva exponencial $E \gg mc^2$ mediante la unión de isótopos de hidrógeno bajo calor extremo. Esta arma se denomina comúnmente:
\begin{enumerate}[label=\textbf{\Alph*)}]
\begin{multicols}{2}
  \item Bomba de Neutrones
  \item Bomba de Hidrógeno (Bomba H)
  \item Bomba de Cobalto
  \item Arma Radiológica (Bomba Sucia)
  \item Proyectil de Uranio Empobrecido
  \item Ojiva Biológica
\end{multicols}
\end{enumerate}
\vspace{0.15cm}
\item Aprobado en 1992, este tratado transformó la Comunidad Económica Europea en la actual Unión Europea, estableciendo la ciudadanía europea y los pilares para la creación de una moneda única (euro) bajo parámetros estrictos de inflación y déficit fiscal. ¿Cuál es su nombre?
\begin{enumerate}[label=\textbf{\Alph*)}]
\begin{multicols}{2}
  \item Tratado de Roma
  \item Tratado de Lisboa
  \item Tratado de Ámsterdam
  \item Tratado de Maastricht
  \item Tratado de Schengen
  \item Tratado de Niza
\end{multicols}
\end{enumerate}
\vspace{0.15cm}
\item Durante la Revolución de Terciopelo de 1989, el gobierno comunista de un país de Europa Central fue derrocado de manera pacífica mediante masivas protestas civiles, permitiendo la elección del dramaturgo Václav Havel como presidente. ¿De qué país se trata?
\begin{enumerate}[label=\textbf{\Alph*)}]
\begin{multicols}{2}
  \item Polonia
  \item Rumania
  \item Checoslovaquia
  \item Hungría
  \item Bulgaria
  \item Alemania Oriental
\end{multicols}
\end{enumerate}
\vspace{0.15cm}
\item Tras la disolución de un Estado federal socialista en la Península de los Balcanes en 1991, se desataron las peores guerras interétnicas en Europa desde la SGM (1991-2001), caracterizadas por asedios, limpieza étnica y la intervención de la OTAN. ¿Qué país se desintegró en este conflicto?
\begin{enumerate}[label=\textbf{\Alph*)}]
\begin{multicols}{2}
  \item Checoslovaquia
  \item Unión Soviética
  \item Imperio Austrohúngaro
  \item Yugoslavia
  \item Rumania
  \item Albania
\end{multicols}
\end{enumerate}
\vspace{0.15cm}
\item En 1994, la pasividad de la comunidad internacional permitió que en 100 días se ejecutara uno de los genocidios más brutales del Siglo XX, donde la mayoría étnica hutu asesinó sistemáticamente a más de 800,000 personas de la minoría tutsi. ¿En qué país africano ocurrió esta masacre?
\begin{enumerate}[label=\textbf{\Alph*)}]
\begin{multicols}{2}
  \item Sudáfrica
  \item Sudán del Sur
  \item Ruanda
  \item Somalia
  \item República Democrática del Congo
  \item Nigeria
\end{multicols}
\end{enumerate}
\vspace{0.15cm}
\item Los Acuerdos de Oslo, firmados en 1993 en la Casa Blanca tras negociaciones secretas en Noruega, establecieron un hito en el conflicto árabe-israelí mediante:
\begin{enumerate}[label=\textbf{\Alph*)}]
\begin{multicols}{2}
  \item La devolución incondicional de toda la ciudad de Jerusalén a los palestinos.
  \item La expulsión del Estado de Israel de la ONU.
  \item El reconocimiento mutuo entre Israel y la OLP y la creación de la Autoridad Nacional Palestina para el autogobierno limitado.
  \item La creación militar conjunta de un Estado binacional armado.
  \item El establecimiento de fronteras abiertas y un mercado común en Medio Oriente.
  \item El fin inmediato y total de todos los asentamientos judíos en Cisjordania.
\end{multicols}
\end{enumerate}
\vspace{0.15cm}
\item Tras la muerte de Mao Zedong, el líder Deng Xiaoping inició en 1978 una serie de reformas estructurales bajo el pragmatismo resumido en su frase: "No importa si el gato es blanco o negro, siempre y cuando cace ratones". El modelo económico resultante se conoce oficialmente en China como:
\begin{enumerate}[label=\textbf{\Alph*)}]
\begin{multicols}{2}
  \item Capitalismo de Laissez-faire
  \item Economía de Mercado Socialista (Socialismo con características chinas)
  \item Comunismo Ortodoxo Soviético
  \item Neoliberalismo Autoritario
  \item Estado de Bienestar Nórdico
  \item Economía Feudal Agraria
\end{multicols}
\end{enumerate}
\vspace{0.15cm}
\item En 1966, Mao Zedong impulsó una masiva purga política y social mediante movilizaciones de los "Guardias Rojos", con el fin de eliminar elementos capitalistas, tradiciones antiguas y opositores burgueses infiltrados en el partido y la sociedad. ¿Cómo se conoce a este periodo histórico?
\begin{enumerate}[label=\textbf{\Alph*)}]
\begin{multicols}{2}
  \item El Gran Salto Adelante
  \item La Revolución Cultural Proletaria
  \item La Marcha Larga
  \item La Guerra del Opio
  \item La Campaña de las Cien Flores
  \item El Levantamiento de los Bóxers
\end{multicols}
\end{enumerate}
\vspace{0.15cm}
\item El 26 de abril de 1986, el reactor $N^{\circ}4$ de una central nuclear soviética explotó liberando radiación equivalente a $\approx 400$ bombas atómicas de Hiroshima. Este desastre ecológico desnudó las fallas del secretismo burocrático de la URSS e impulsó la política de Glasnost. ¿Cuál es la central accidentada?
\begin{enumerate}[label=\textbf{\Alph*)}]
\begin{multicols}{2}
  \item Fukushima Daiichi
  \item Three Mile Island
  \item Chernóbil
  \item Zaporiyia
  \item Kyshtym
  \item Sellafield
\end{multicols}
\end{enumerate}
\vspace{0.15cm}
\item Entre junio de 1948 y mayo de 1949, Iósif Stalin bloqueó todas las vías de acceso terrestre y fluvial al sector occidental de una capital europea enclavada dentro de territorio comunista. Para evitar la rendición por hambre y desabastecimiento, EE. UU. y el Reino Unido implementaron:
\begin{enumerate}[label=\textbf{\Alph*)}]
\begin{multicols}{2}
  \item Un ataque preventivo con misiles balísticos intercontinentales.
  \item Un enorme puente aéreo ininterrumpido (Luftbrücke) para suministrar $\approx 2.3 \times 10^{6}$ toneladas de insumos diarios.
  \item Una invasión terrestre con tanques Sherman y divisiones blindadas.
  \item La entrega formal de la capital a cambio de territorios en Polonia.
  \item Un embargo comercial global inmediato contra China y la URSS.
  \item La construcción del Muro de contención para sitiar a las fuerzas rusas.
\end{multicols}
\end{enumerate}
\vspace{0.15cm}
\item Como respuesta al Plan Marshall occidental, en 1947 la Unión Soviética creó una organización supranacional para coordinar ideológicamente a los distintos partidos comunistas de Europa y garantizar su obediencia a las directrices de Moscú. Esta entidad se llamó:
\begin{enumerate}[label=\textbf{\Alph*)}]
\begin{multicols}{2}
  \item COMECON
  \item Pacto de Varsovia
  \item Kominform (Oficina de Información de los Partidos Comunistas)
  \item KGB
  \item Stasi
  \item Politburó
\end{multicols}
\end{enumerate}
\vspace{0.15cm}
\item La intervención militar estadounidense en América Latina incluyó una política articulada en 1961 por John F. Kennedy. Esta iniciativa ofrecía ayuda económica, social y técnica a los países del hemisferio con el fin expreso de contrarrestar el atractivo de la Revolución Cubana y promover el reformismo democrático. Se denominó:
\begin{enumerate}[label=\textbf{\Alph*)}]
\begin{multicols}{2}
  \item Plan Cóndor
  \item Alianza para el Progreso
  \item Doctrina del Gran Garrote (Big Stick)
  \item Iniciativa de la Cuenca del Caribe
  \item Política del Buen Vecino
  \item Consenso de Washington
\end{multicols}
\end{enumerate}
\vspace{0.15cm}
\item Durante la Revolución de los Claveles, ocurrida el 25 de abril de 1974, un levantamiento militar de oficiales de izquierda derrocó sin derramamiento de sangre a la dictadura corporativista de casi 50 años (Estado Novo). ¿En qué país ocurrió, facilitando además la independencia de Angola y Mozambique?
\begin{enumerate}[label=\textbf{\Alph*)}]
\begin{multicols}{2}
  \item España
  \item Grecia
  \item Italia
  \item Portugal
  \item Brasil
  \item Yugoslavia
\end{multicols}
\end{enumerate}
\vspace{0.15cm}
\item Tras su división de la India británica en 1947, India (de mayoría hindú) y Pakistán (de mayoría musulmana) se han enfrentado en múltiples guerras por el control soberano de una región montañosa fronteriza altamente militarizada. ¿Cuál es esta región geopolíticamente estratégica?
\begin{enumerate}[label=\textbf{\Alph*)}]
\begin{multicols}{2}
  \item Tíbet
  \item Punyab
  \item Bengala
  \item Cachemira
  \item Baluchistán
  \item Ceilán (Sri Lanka)
\end{multicols}
\end{enumerate}
\vspace{0.15cm}
\item En la Crisis de los Misiles de Cuba (1962), EE. UU. descubrió mediante vuelos espías $U-2$ los misiles soviéticos R-12 en la isla. El acuerdo final para que la URSS retirara su arsenal de Cuba incluyó el compromiso secreto de Estados Unidos de retirar sus misiles balísticos (modelo Jupiter) estacionados geográficamente en:
\begin{enumerate}[label=\textbf{\Alph*)}]
\begin{multicols}{2}
  \item Alemania Occidental
  \item Italia
  \item Grecia
  \item Turquía
  \item Japón
  \item Corea del Sur
\end{multicols}
\end{enumerate}
\vspace{0.15cm}
\item La "Década Perdida" en América Latina corresponde a los años de 1980, donde una profunda recesión macroeconómica, hiperinflación y estancamiento productivo afectaron a la región, detonada principalmente por:
\begin{enumerate}[label=\textbf{\Alph*)}]
\begin{multicols}{2}
  \item La nacionalización masiva de todas las industrias mineras.
  \item La crisis de la deuda externa soberana, originada por el abrupto incremento en las tasas de interés impulsado por la Reserva Federal estadounidense (Shock Volcker).
  \item El embargo petrolero de la OPEP contra México y Venezuela.
  \item La destrucción de infraestructura debido al fenómeno climático del Niño.
  \item La negativa del FMI a otorgar préstamos a regímenes democráticos.
  \item El establecimiento del Euro y la pérdida de competitividad de las monedas locales.
\end{multicols}
\end{enumerate}
\vspace{0.15cm}
\item El Movimiento de Países No Alineados (MNOAL) fue articulado por líderes históricos que buscaban una "tercera vía" para proteger su soberanía nacional frente a Washington y Moscú. Identifique el grupo correcto que contiene a tres de sus fundadores principales:
\begin{enumerate}[label=\textbf{\Alph*)}]
\begin{multicols}{2}
  \item Winston Churchill, Franklin D. Roosevelt, Iósif Stalin
  \item Josip Broz "Tito", Jawaharlal Nehru, Gamal Abdel Nasser
  \item Fidel Castro, Ernesto "Che" Guevara, Salvador Allende
  \item Mao Zedong, Kim Il-sung, Ho Chi Minh
  \item Charles de Gaulle, Konrad Adenauer, Alcide De Gasperi
  \item Nelson Mandela, Patrice Lumumba, Kwame Nkrumah
\end{multicols}
\end{enumerate}
\vspace{0.15cm}
\item A partir del año 1975, la coordinación económica del bloque de países capitalistas e industrializados se formalizó mediante cumbres anuales para discutir políticas macroeconómicas y estabilidad de divisas. Este exclusivo grupo, que representa gran parte de la riqueza mundial (excluyendo a la URSS y China), se denomina:
\begin{enumerate}[label=\textbf{\Alph*)}]
\begin{multicols}{2}
  \item BRICS
  \item Grupo de los 77 (G-77)
  \item Grupo de los Siete (G-7)
  \item OPEP
  \item OCDE Ampliada
  \item Foro de Davos
\end{multicols}
\end{enumerate}
\vspace{0.15cm}
\item Durante la Guerra Fría, la Unión Soviética ejerció control doctrinario sobre el Pacto de Varsovia estableciendo en 1968 un axioma político que declaraba que "cuando las fuerzas hostiles al socialismo intentan desviar a un país hacia el capitalismo, esto no es solo problema de ese país, sino de todos los países socialistas". Esta justificación de intervención militar se conoce como:
\begin{enumerate}[label=\textbf{\Alph*)}]
\begin{multicols}{2}
  \item Doctrina Truman
  \item Doctrina Jrushchov
  \item Doctrina Brézhnev
  \item Doctrina Monroe
  \item Política de Glasnost
  \item Estrategia de Choque
\end{multicols}
\end{enumerate}
\vspace{0.15cm}
\item El 20 de diciembre de 1989, el gobierno de Estados Unidos lanzó la "Operación Causa Justa" desplegando unos 27,000 soldados de élite. El objetivo explícito de esta invasión fue desmantelar a las Fuerzas de Defensa y arrestar bajo cargos de narcotráfico al dictador general Manuel Antonio Noriega. ¿En qué país ocurrió?
\begin{enumerate}[label=\textbf{\Alph*)}]
\begin{multicols}{2}
  \item Nicaragua
  \item Cuba
  \item Granada
  \item Panamá
  \item República Dominicana
  \item Haití
\end{multicols}
\end{enumerate}
\vspace{0.15cm}
\item En el contexto de El Salvador (1980), el recrudecimiento de la violencia y el preludio de la guerra civil a gran escala acapararon la atención mundial tras un magnicidio ejecutado el 24 de marzo por los escuadrones de la muerte de la ultraderecha. ¿Qué líder eclesiástico y defensor de los derechos humanos fue asesinado mientras oficiaba misa?
\begin{enumerate}[label=\textbf{\Alph*)}]
\begin{multicols}{2}
  \item Ignacio Ellacuría
  \item Rutilio Grande
  \item Monseñor Óscar Arnulfo Romero
  \item Arturo Rivera y Damas
  \item Juan Gerardi
  \item Fernando Lugo
\end{multicols}
\end{enumerate}
\vspace{0.15cm}
\item Tras la finalización de la Segunda Guerra Mundial, la mayoría de los países de Europa Occidental adoptaron y expandieron un modelo sociopolítico sustentado en políticas fiscales redistributivas ($G \uparrow$), con el objetivo de garantizar educación, sanidad pública universal, pensiones y seguros de desempleo a la ciudadanía. Este modelo se denomina:
\begin{enumerate}[label=\textbf{\Alph*)}]
\begin{multicols}{2}
  \item Anarcocapitalismo
  \item Corporativismo Sindical Totalitario
  \item Estado de Bienestar (Welfare State)
  \item Capitalismo de Libre Mercado Desregulado (Laissez-faire)
  \item Comunismo de Guerra
  \item Socialismo Estructural Planificado
\end{multicols}
\end{enumerate}
\vspace{0.15cm}
\item En el sudeste asiático, un proceso de descolonización enfrentó al gobierno nacionalista indonesio (liderado por Sukarno) contra el intento de reconquista imperial de la antigua metrópoli tras el fin de la ocupación japonesa en la SGM. Tras cuatro años de guerra diplomática y armada (1945-1949), este enorme archipiélago se independizó de:
\begin{enumerate}[label=\textbf{\Alph*)}]
\begin{multicols}{2}
  \item Imperio Británico
  \item Francia
  \item Imperio Español
  \item Países Bajos (Holanda)
  \item Portugal
  \item Estados Unidos
\end{multicols}
\end{enumerate}
\vspace{0.15cm}
\item En agosto de 1991, un grupo de la línea dura del Partido Comunista, el KGB y el Ejército Rojo, temiendo la desintegración del país por la firma del Nuevo Tratado de la Unión, dio un golpe de Estado fallido (Putsch de Moscú) secuestrando temporalmente a:
\begin{enumerate}[label=\textbf{\Alph*)}]
\begin{multicols}{2}
  \item Vladímir Putin
  \item Borís Yeltsin
  \item Mijaíl Gorbachov
  \item Leonid Brézhnev
  \item Eduard Shevardnadze
  \item Yuri Andrópov
\end{multicols}
\end{enumerate}
\vspace{0.15cm}
\item Tras el golpe de Estado fallido de 1991, ¿qué figura política rusa lideró la resistencia popular subiéndose a un tanque frente a la Casa Blanca de Moscú, precipitando la proscripción del Partido Comunista y convirtiéndose en el primer presidente de la Federación de Rusia postsoviética?
\begin{enumerate}[label=\textbf{\Alph*)}]
\begin{multicols}{2}
  \item Mijaíl Gorbachov
  \item Borís Yeltsin
  \item Vladímir Putin
  \item Aleksandr Rutskói
  \item Dmitri Medvédev
  \item Alexéi Navalni
\end{multicols}
\end{enumerate}
\vspace{0.15cm}
\item La guerra Irán-Irak (1980-1988) fue uno de los conflictos convencionales más letales de Medio Oriente. Comenzó cuando el líder iraquí, buscando capitalizar la inestabilidad post-revolución en Irán y controlar zonas ricas en petróleo como Juzestán, invadió su país vecino. ¿Quién era el gobernante de Irak durante esta guerra?
\begin{enumerate}[label=\textbf{\Alph*)}]
\begin{multicols}{2}
  \item Ruhollah Jomeiní
  \item Muamar el Gadafi
  \item Saddam Hussein
  \item Bashar al-Ásad
  \item Yasir Arafat
  \item Hosni Mubarak
\end{multicols}
\end{enumerate}
\vspace{0.15cm}
\item A mediados de los años 70, Cuba proyectó su política exterior militar apoyando activamente a movimientos de liberación marxistas y regímenes pro-soviéticos en África. La mayor intervención cubana, enviando a decenas de miles de soldados (Operación Carlota), apoyó al Movimiento Popular de Liberación (MPLA) en la guerra civil de:
\begin{enumerate}[label=\textbf{\Alph*)}]
\begin{multicols}{2}
  \item Sudáfrica
  \item Zimbabue
  \item Ruanda
  \item Angola
  \item Egipto
  \item Marruecos
\end{multicols}
\end{enumerate}
\vspace{0.15cm}
\item El sistema del Apartheid sudafricano generó múltiples protestas reprimidas violentamente por el Estado. En 1976, una serie de manifestaciones lideradas por miles de estudiantes de secundaria negros protestando contra la imposición del idioma afrikáans en la enseñanza fue brutalmente ametrallada. Este evento clave de resistencia global se conoce como:
\begin{enumerate}[label=\textbf{\Alph*)}]
\begin{multicols}{2}
  \item Levantamiento de Soweto
  \item Masacre de Sharpeville
  \item Marcha de la Sal
  \item Noche de los Cristales Rotos
  \item Domingo Sangriento
  \item Masacre de Tlatelolco
\end{multicols}
\end{enumerate}
\vspace{0.15cm}
\item El 1 de enero de 1994 entró en vigencia un tratado económico que eliminó barreras arancelarias creando una de las zonas de libre comercio más grandes del mundo, transformando las cadenas de suministro $(X, M)$ en América del Norte e integrando fuertemente la economía mexicana. ¿Cuál fue este tratado?
\begin{enumerate}[label=\textbf{\Alph*)}]
\begin{multicols}{2}
  \item MERCOSUR
  \item Tratado de Libre Comercio de América del Norte (TLCAN o NAFTA)
  \item Alianza del Pacífico
  \item Mercado Común Centroamericano (MCCA)
  \item Acuerdo Transpacífico de Cooperación Económica (TPP)
  \item Área de Libre Comercio de las Américas (ALCA)
\end{multicols}
\end{enumerate}
\vspace{0.15cm}
\item En el marco del desarme nuclear de finales de la Guerra Fría, los presidentes Ronald Reagan y Mijaíl Gorbachov firmaron en 1987 un tratado sin precedentes que eliminó físicamente toda una clase de armamento: misiles balísticos y de crucero con rango de vuelo $R \in [500,5500]\text{ km}$. ¿A qué tratado corresponde?
\begin{enumerate}[label=\textbf{\Alph*)}]
\begin{multicols}{2}
  \item Tratado de Cielos Abiertos
  \item Tratado START II
  \item Tratado INF (Fuerzas Nucleares de Alcance Intermedio)
  \item Tratado ABM (Antimisiles Balísticos)
  \item Tratado de No Proliferación Nuclear (TNP)
  \item SALT I
\end{multicols}
\end{enumerate}
\vspace{0.15cm}
\item Las "Doctrinas de Seguridad Nacional", impartidas militarmente en la Escuela de las Américas durante la Guerra Fría, redefinieron la función de las fuerzas armadas en el Cono Sur latinoamericano. Su postulado central dictaba que el principal enemigo existencial de la nación no era un invasor extranjero tradicional, sino:
\begin{enumerate}[label=\textbf{\Alph*)}]
\begin{multicols}{2}
  \item Los cárteles del narcotráfico transnacional.
  \item El "enemigo interno", constituido por disidentes políticos, estudiantes, sindicatos e ideologías subversivas (comunismo).
  \item Las organizaciones no gubernamentales (ONG) extranjeras de corte ecologista.
  \item Las misiones diplomáticas del bloque europeo que buscaban re-colonizar la región.
  \item El capitalismo global financiero y los bancos de inversión de Wall Street.
  \item El avance del protestantismo evangélico en detrimento de la iglesia católica estatal.
\end{multicols}
\end{enumerate}
\vspace{0.15cm}
\item Las Conferencias sobre Seguridad y la Cooperación en Europa concluyeron en 1975 con la firma de un acta que reconoció formalmente las fronteras europeas establecidas tras la SGM (validando de facto la dominación soviética sobre Europa del Este), pero obligaba al bloque comunista a respetar derechos humanos elementales. ¿De qué acuerdos se trata?
\begin{enumerate}[label=\textbf{\Alph*)}]
\begin{multicols}{2}
  \item Acuerdos de Múnich
  \item Acuerdos de Helsinki
  \item Pacto de San José
  \item Acuerdos de Ginebra
  \item Tratado de Letrán
  \item Acuerdos de Schengen
\end{multicols}
\end{enumerate}
\vspace{0.15cm}
\item En 1986, el bloque soviético lanzó un proyecto espacial ambicioso colocando en órbita la primera estación de investigación modular habitada de manera permanente para misiones de microgravedad a largo plazo. Su funcionamiento simbolizó un éxito del programa civil aeroespacial ruso. ¿Cuál fue su nombre?
\begin{enumerate}[label=\textbf{\Alph*)}]
\begin{multicols}{2}
  \item Estación Espacial Internacional (ISS)
  \item Skylab
  \item Estación Mir
  \item Soyuz-T
  \item Vostok 6
  \item Proyecto Apollo-Soyuz
\end{multicols}
\end{enumerate}
\vspace{0.15cm}
\item En la década de 1990, con la popularización de las computadoras personales conectadas a redes satelitales y fibra óptica, un protocolo de red transformó radicalmente las telecomunicaciones internacionales al crear un espacio hipervinculado, siendo el motor de la "Tercera Revolución Industrial". Esta invención (creada por Tim Berners-Lee en el CERN) es:
\begin{enumerate}[label=\textbf{\Alph*)}]
\begin{multicols}{2}
  \item El Microprocesador de Silicio
  \item La World Wide Web (WWW)
  \item El Sistema de Posicionamiento Global (GPS)
  \item La Inteligencia Artificial
  \item El Teléfono Celular (GSM)
  \item La Máquina de Enigma
\end{multicols}
\end{enumerate}
\vspace{0.15cm}
\item El fin oficial del orden bipolar y la conclusión historiográfica del marco central de la Guerra Fría están datados el día que la bandera roja de la URSS fue arriada del Kremlin por última vez, oficializando su desintegración legal en 15 estados soberanos. ¿Qué fecha exacta marca este acontecimiento final?
\begin{enumerate}[label=\textbf{\Alph*)}]
\begin{multicols}{2}
  \item 9 de noviembre de 1989
  \item 25 de diciembre de 1991
  \item 1 de enero de 1990
  \item 3 de octubre de 1990
  \item 15 de agosto de 1991
  \item 31 de diciembre de 1999
\end{multicols}
\end{enumerate}
\vspace{0.15cm}
\end{enumerate}\n
\subsubsection*{Módulo 3: Reactivos 101 - 150}

\begin{enumerate}[label=\textbf{\arabic*.}]
\setcounter{enumi}{100}
\item Entre noviembre de $1945$ y octubre de $1946$, las potencias Aliadas organizaron una serie de tribunales militares internacionales para enjuiciar a $24$ altos jerarcas del Tercer Reich. Estos procesos establecieron el concepto jurídico de "Crímenes contra la humanidad". ¿Cómo se denominan estos juicios?
\begin{enumerate}[label=\textbf{\Alph*)}]
\begin{multicols}{2}
  \item Juicios de Tokio
  \item Tribunal de La Haya
  \item Juicios de Núremberg
  \item Tribunal de Múnich
  \item Tratado de Versalles
  \item Convención de Ginebra
\end{multicols}
\end{enumerate}
\vspace{0.15cm}
\item En el contexto de la contención inicial durante $1947$, el presidente estadounidense Harry Truman solicitó al Congreso la inyección urgente de $\approx \$400 \times 10^6$ en ayuda militar y económica para evitar que dos naciones del Mediterráneo cayeran bajo influencia de guerrillas comunistas y presiones soviéticas. ¿Cuáles fueron estos dos países?
\begin{enumerate}[label=\textbf{\Alph*)}]
\begin{multicols}{2}
  \item España y Portugal
  \item Italia y Yugoslavia
  \item Francia y Austria
  \item Grecia y Turquía
  \item Egipto e Israel
  \item Irán y Afganistán
\end{multicols}
\end{enumerate}
\vspace{0.15cm}
\item Durante la Guerra de Corea ($1950$-$1953$), las fuerzas surcoreanas y de la ONU estuvieron a punto de ser expulsadas hacia el mar en el Perímetro de Pusan. La guerra dio un giro táctico radical cuando el General Douglas MacArthur ejecutó un audaz desembarco anfibio detrás de las líneas norcoreanas en:
\begin{enumerate}[label=\textbf{\Alph*)}]
\begin{multicols}{2}
  \item Seúl
  \item Pionyang
  \item Inchon
  \item Wonsan
  \item Bahía de Cochinos
  \item Isla de Jeju
\end{multicols}
\end{enumerate}
\vspace{0.15cm}
\item El fin del dominio colonial francés en la península de Indochina se selló en $1954$, cuando las fuerzas de la Unión Francesa sufrieron una devastadora derrota militar (con la pérdida de $\approx 11,000$ hombres) frente al Viet Minh, liderado estratégicamente por el general Vo Nguyen Giap en la batalla de:
\begin{enumerate}[label=\textbf{\Alph*)}]
\begin{multicols}{2}
  \item Saigón
  \item Dien Bien Phu
  \item Hué
  \item Golfo de Tonkín
  \item Valle de Ia Drang
  \item Khe Sanh
\end{multicols}
\end{enumerate}
\vspace{0.15cm}
\item La "Guerra Fría" dependía en gran medida de operaciones de espionaje y acciones encubiertas. El organismo central de inteligencia, contraespionaje y policía secreta de la Unión Soviética, establecido en $1954$ y operante hasta $1991$, respondía a las siglas:
\begin{enumerate}[label=\textbf{\Alph*)}]
\begin{multicols}{2}
  \item CIA
  \item MI6
  \item Stasi
  \item Mossad
  \item KGB
  \item Gestapo
\end{multicols}
\end{enumerate}
\vspace{0.15cm}
\item En la década de $1980$, la ortodoxia económica del Estado de Bienestar keynesiano fue desplazada por una corriente conservadora y neoliberal (privatizaciones, desregulación financiera, reducción de la carga fiscal $T \downarrow$). Los principales arquitectos políticos de este cambio global fueron:
\begin{enumerate}[label=\textbf{\Alph*)}]
\begin{multicols}{2}
  \item Franklin D. Roosevelt y Winston Churchill
  \item John F. Kennedy y Harold Macmillan
  \item Ronald Reagan y Margaret Thatcher
  \item Bill Clinton y Tony Blair
  \item George H.W. Bush y François Mitterrand
  \item Jimmy Carter y James Callaghan
\end{multicols}
\end{enumerate}
\vspace{0.15cm}
\item En el año $1954$, la CIA orquestó la Operación PBSuccess en Centroamérica. Esta intervención encubierta derrocó a un presidente democrático debido a sus políticas de reforma agraria que expropiaban tierras ociosas de la United Fruit Company. ¿A qué líder guatemalteco se refiere?
\begin{enumerate}[label=\textbf{\Alph*)}]
\begin{multicols}{2}
  \item Jorge Ubico
  \item Carlos Castillo Armas
  \item Jacobo Árbenz Guzmán
  \item Juan José Arévalo
  \item Manuel Estrada Cabrera
  \item Efraín Ríos Montt
\end{multicols}
\end{enumerate}
\vspace{0.15cm}
\item El sistema económico global sufrió una segunda recesión inflacionaria profunda en $1979$, caracterizada por un choque en la oferta ($\Delta Q_s \ll 0$) que disparó exponencialmente el precio del barril de petróleo. Este evento geopolítico fue el resultado directo de:
\begin{enumerate}[label=\textbf{\Alph*)}]
\begin{multicols}{2}
  \item La invasión soviética de Afganistán
  \item La Revolución Islámica de Irán
  \item La Guerra del Golfo Pérsico
  \item La Nacionalización del Canal de Suez
  \item La Guerra de Yom Kipur
  \item El embargo de Venezuela
\end{multicols}
\end{enumerate}
\vspace{0.15cm}
\item En junio de $1963$, el programa espacial soviético demostró nuevamente su superioridad técnica (y un avance en equidad de género propagandístico) al enviar a la primera mujer en orbitar el planeta Tierra a bordo de la cápsula Vostok 6. Su nombre es:
\begin{enumerate}[label=\textbf{\Alph*)}]
\begin{multicols}{2}
  \item Sally Ride
  \item Svetlana Savítskaya
  \item Valentina Tereshkova
  \item Mae Jemison
  \item Christa McAuliffe
  \item Elena Kondakova
\end{multicols}
\end{enumerate}
\vspace{0.15cm}
\item El mes de mayo de $1968$ se convirtió en el epicentro de la contracultura europea contemporánea. Comenzando como una serie de protestas estudiantiles en la Universidad de Nanterre, evolucionó hacia la mayor huelga general ($\approx 9 \times 10^6$ de trabajadores) de la historia moderna de Europa, paralizando el país. ¿En qué nación ocurrió?
\begin{enumerate}[label=\textbf{\Alph*)}]
\begin{multicols}{2}
  \item Alemania Occidental (RFA)
  \item Reino Unido
  \item Italia
  \item Francia
  \item Checoslovaquia
  \item España
\end{multicols}
\end{enumerate}
\vspace{0.15cm}
\item El $4$ de abril de $1968$, Estados Unidos experimentó profundos disturbios raciales tras el asesinato de su principal líder del movimiento por los derechos civiles. Este líder, premio Nobel de la Paz y defensor de la desobediencia civil no violenta, fue:
\begin{enumerate}[label=\textbf{\Alph*)}]
\begin{multicols}{2}
  \item Malcolm X
  \item Rosa Parks
  \item John F. Kennedy
  \item Martin Luther King Jr.
  \item Medgar Evers
  \item Stokely Carmichael
\end{multicols}
\end{enumerate}
\vspace{0.15cm}
\item En abril de $1961$, un ejército compuesto por aproximadamente $1,400$ exiliados cubanos (Brigada 2506), entrenados y financiados por la CIA, intentó invadir Cuba para derrocar a Fidel Castro. La operación fue un fracaso militar y táctico estrepitoso conocido como:
\begin{enumerate}[label=\textbf{\Alph*)}]
\begin{multicols}{2}
  \item Crisis de Octubre
  \item Operación Mangosta
  \item Invasión de Bahía de Cochinos (Playa Girón)
  \item Batalla de Santa Clara
  \item Operación Causa Justa
  \item El Asalto al Cuartel Moncada
\end{multicols}
\end{enumerate}
\vspace{0.15cm}
\item En respuesta a las atrocidades de la Segunda Guerra Mundial, la Asamblea General de la ONU adoptó el $10$ de diciembre de $1948$ un documento histórico que establece los derechos inalienables fundamentales que corresponden a toda persona por su condición humana. Este documento es:
\begin{enumerate}[label=\textbf{\Alph*)}]
\begin{multicols}{2}
  \item La Carta de las Naciones Unidas
  \item La Declaración de los Derechos del Hombre y del Ciudadano
  \item Los Protocolos de Ginebra
  \item La Declaración Universal de los Derechos Humanos (DUDH)
  \item El Pacto de Varsovia
  \item La Carta Magna
\end{multicols}
\end{enumerate}
\vspace{0.15cm}
\item Entre $1975$ y $1979$, un régimen comunista fanático tomó el poder en Camboya (Kampuchea Democrática). Su intento de forzar un retorno inmediato a una sociedad agraria pura culminó en el exterminio de $\approx 25\%$ de su propia población ($> 1.5 \times 10^6$ de muertes). ¿Cómo se conocía a este grupo y a su líder?
\begin{enumerate}[label=\textbf{\Alph*)}]
\begin{multicols}{2}
  \item Viet Cong liderado por Ho Chi Minh
  \item Pathet Lao liderado por Souphanouvong
  \item Jemeres Rojos (Khmer Rouge) liderados por Pol Pot
  \item Kuomintang liderado por Chiang Kai-shek
  \item Sendero Luminoso liderado por Abimael Guzmán
  \item Ejército Rojo liderado por Mao Zedong
\end{multicols}
\end{enumerate}
\vspace{0.15cm}
\item En el año $1960$, la República Democrática del Congo declaró su independencia de Bélgica. Su primer ministro democráticamente electo, de tendencia anticolonialista y panafricanista, fue asesinado en $1961$ por fuerzas separatistas con la complicidad directa de los servicios de inteligencia belgas y la CIA. ¿De quién se trata?
\begin{enumerate}[label=\textbf{\Alph*)}]
\begin{multicols}{2}
  \item Nelson Mandela
  \item Idi Amin
  \item Patrice Lumumba
  \item Thomas Sankara
  \item Mobutu Sese Seko
  \item Jomo Kenyatta
\end{multicols}
\end{enumerate}
\vspace{0.15cm}
\item El primer hito histórico para la integración pacífica europea posterior a la SGM ocurrió en $1951$ bajo el Plan Schuman. Se basó en la premisa de que "quien no controla los materiales básicos para la guerra, no puede hacer la guerra". ¿Qué organización supranacional franco-alemana nació de este plan?
\begin{enumerate}[label=\textbf{\Alph*)}]
\begin{multicols}{2}
  \item Comunidad Económica Europea (CEE)
  \item Unión Europea (UE)
  \item Organización del Tratado del Atlántico Norte (OTAN)
  \item Comunidad Europea del Carbón y del Acero (CECA)
  \item Asociación Europea de Libre Comercio (EFTA)
  \item Espacio Schengen
\end{multicols}
\end{enumerate}
\vspace{0.15cm}
\item Durante la Guerra de las Malvinas de $1982$, el hundimiento de un crucero ligero por un submarino de propulsión nuclear británico (el HMS Conqueror) representó el golpe moral y humano más duro para la Armada Argentina, provocando $323$ bajas mortales. ¿Cuál era el nombre de la embarcación argentina hundida?
\begin{enumerate}[label=\textbf{\Alph*)}]
\begin{multicols}{2}
  \item ARA San Juan
  \item ARA Veinticinco de Mayo
  \item ARA General Belgrano
  \item ARA Santísima Trinidad
  \item SS Atlantic Conveyor
  \item ARA Almirante Brown
\end{multicols}
\end{enumerate}
\vspace{0.15cm}
\item En $1987$, tras un accidente de tráfico provocado por un camión militar israelí en Gaza, estalló un levantamiento popular y asimétrico palestino que se extendió por Cisjordania, caracterizado por jóvenes lanzando piedras contra tanques. Este fenómeno, que transformó el conflicto, se denomina:
\begin{enumerate}[label=\textbf{\Alph*)}]
\begin{multicols}{2}
  \item Operación Plomo Fundido
  \item La Guerra de los Seis Días
  \item La Primera Intifada
  \item La Segunda Intifada (Intifada de Al-Aqsa)
  \item El Septiembre Negro
  \item La Guerra de Yom Kipur
\end{multicols}
\end{enumerate}
\vspace{0.15cm}
\item En julio de $1995$, la Guerra de Bosnia presenció el mayor genocidio en suelo europeo desde la SGM. Las fuerzas serbobosnias, bajo el mando del general Ratko Mladić, ejecutaron a más de $8,000$ hombres y niños musulmanes en un enclave que, teóricamente, estaba bajo protección de cascos azules de la ONU. Este hecho ocurrió en:
\begin{enumerate}[label=\textbf{\Alph*)}]
\begin{multicols}{2}
  \item Sarajevo
  \item Mostar
  \item Kosovo
  \item Srebrenica
  \item Pristina
  \item Belgrado
\end{multicols}
\end{enumerate}
\vspace{0.15cm}
\item A principios de la década de $1980$, en un país de Europa del Este perteneciente al Pacto de Varsovia, surgió un sindicato obrero católico e independiente que organizó masivas huelgas en los astilleros de Gdansk. Liderado por Lech Wałęsa, este movimiento fue el principal catalizador de la caída del comunismo en la región. Se llamó:
\begin{enumerate}[label=\textbf{\Alph*)}]
\begin{multicols}{2}
  \item Solidaridad (Solidarność)
  \item Primavera de Praga
  \item Revolución de Terciopelo
  \item Komsomol
  \item Oposición Obrera
  \item Glasnost
\end{multicols}
\end{enumerate}
\vspace{0.15cm}
\item Tras la Crisis de los Misiles Cubanos de $1962$, las potencias nucleares reconocieron el gravísimo riesgo de un holocausto atómico desatado por fallas de comunicación. Para evitar un error de cálculo (donde el tiempo de respuesta $\Delta t \to 0$), EE. UU. y la URSS establecieron en $1963$:
\begin{enumerate}[label=\textbf{\Alph*)}]
\begin{multicols}{2}
  \item El desarme total supervisado por la ONU.
  \item El "Teléfono Rojo" (una línea directa de teletipo entre Washington y Moscú).
  \item El Muro de Berlín.
  \item La base conjunta de Guantánamo.
  \item El intercambio permanente de embajadores presidenciales.
  \item El satélite compartido de comunicaciones "Mir".
\end{multicols}
\end{enumerate}
\vspace{0.15cm}
\item Días antes de los Juegos Olímpicos de verano de $1968$, un mitin estudiantil en la Plaza de las Tres Culturas (Tlatelolco) fue brutalmente ametrallado por francotiradores gubernamentales y el Batallón Olimpia del ejército, sofocando de raíz el movimiento democrático en ese país. ¿Dónde sucedió este evento?
\begin{enumerate}[label=\textbf{\Alph*)}]
\begin{multicols}{2}
  \item Argentina
  \item España
  \item Colombia
  \item México
  \item Perú
  \item Brasil
\end{multicols}
\end{enumerate}
\vspace{0.15cm}
\item En el contexto de la Guerra Civil Salvadoreña, el FMLN ejecutó en noviembre de $1989$ la ofensiva táctica de mayor envergadura militar urbana, llevando la guerra a las zonas residenciales de la clase alta y la capital, demostrando que el conflicto era militarmente insoluble. ¿Cómo se denominó esta operación?
\begin{enumerate}[label=\textbf{\Alph*)}]
\begin{multicols}{2}
  \item Ofensiva Final de 1981
  \item Operación Rescate
  \item Ofensiva "Hasta el Tope"
  \item Campaña de Morazán
  \item Plan Guazapa
  \item Operación Fénix
\end{multicols}
\end{enumerate}
\vspace{0.15cm}
\item El $16$ de septiembre de $1987$, la comunidad global demostró capacidad de cooperación científica sin precedentes al firmar un tratado para eliminar gradualmente la producción de gases Clorofluorocarbonos (CFCs), los cuales destruían químicamente la capa de ozono estratosférico $O_3$. ¿Cuál es el nombre de este instrumento internacional?
\begin{enumerate}[label=\textbf{\Alph*)}]
\begin{multicols}{2}
  \item Protocolo de Kioto
  \item Acuerdo de París
  \item Protocolo de Montreal
  \item Tratado de Copenhague
  \item Convenio de Viena
  \item Carta de la Tierra
\end{multicols}
\end{enumerate}
\vspace{0.15cm}
\item El proceso de emancipación poscolonial africana estuvo plagado de tensiones interétnicas, provocadas por fronteras dibujadas artificialmente en Europa. Entre $1967$ y $1970$, Nigeria sufrió una devastadora guerra civil tras la secesión de su provincia sureste rica en petróleo. ¿Cómo se conoce esta guerra?
\begin{enumerate}[label=\textbf{\Alph*)}]
\begin{multicols}{2}
  \item Guerra de Secesión del Congo
  \item Guerra de Biafra
  \item Guerra del Darfur
  \item Revolución de Zanzíbar
  \item Guerra del Rif
  \item Guerra del Ogadén
\end{multicols}
\end{enumerate}
\vspace{0.15cm}
\item En el marco geopolítico de la Guerra Fría de $1957$, un presidente de EE. UU. promulgó una doctrina que garantizaba la asistencia económica y el uso de las Fuerzas Armadas estadounidenses para proteger a cualquier nación del Medio Oriente que fuese amenazada por la "agresión armada de cualquier país controlado por el comunismo internacional". Esta fue la:
\begin{enumerate}[label=\textbf{\Alph*)}]
\begin{multicols}{2}
  \item Doctrina Truman
  \item Doctrina Monroe
  \item Doctrina Reagan
  \item Doctrina Eisenhower
  \item Doctrina Nixon
  \item Doctrina Brézhnev
\end{multicols}
\end{enumerate}
\vspace{0.15cm}
\item ¿Qué líder soviético gobernó la URSS entre $1964$ y $1982$, periodo caracterizado por la paridad nuclear con EE. UU., un gigantesco complejo industrial-militar, pero también por una severa esclerosis burocrática y déficit de bienes de consumo interno denominado "La Era de Estancamiento"?
\begin{enumerate}[label=\textbf{\Alph*)}]
\begin{multicols}{2}
  \item Iósif Stalin
  \item Nikita Jruschov
  \item Mijaíl Gorbachov
  \item Yuri Andrópov
  \item Konstantín Chernenko
  \item Leonid Brézhnev
\end{multicols}
\end{enumerate}
\vspace{0.15cm}
\item La polarización sobre el proceso de paz en Medio Oriente tuvo su episodio más crítico en noviembre de $1995$, cuando el Primer Ministro israelí (firmante de los Acuerdos de Oslo) fue asesinado a tiros en Tel Aviv por Yigal Amir, un estudiante de derecho judío ultranacionalista. ¿De qué líder se trata?
\begin{enumerate}[label=\textbf{\Alph*)}]
\begin{multicols}{2}
  \item Ariel Sharón
  \item Menájem Beguin
  \item Golda Meir
  \item Ehud Barak
  \item Shimon Peres
  \item Yitzhak Rabin
\end{multicols}
\end{enumerate}
\vspace{0.15cm}
\item El descubrimiento de los "Archivos del Terror" en Paraguay ($1992$) aportó las pruebas documentales de un mecanismo transnacional de exterminio donde las dictaduras militares de América del Sur compartían inteligencia y escuadrones para secuestrar, torturar y desaparecer a opositores marxistas en la década de los $70$s. Su nombre operativo fue:
\begin{enumerate}[label=\textbf{\Alph*)}]
\begin{multicols}{2}
  \item Operación Barbarroja
  \item Plan Marshall
  \item Operación Cóndor
  \item Alianza para el Progreso
  \item Escuela de las Américas
  \item Plan Colombia
\end{multicols}
\end{enumerate}
\vspace{0.15cm}
\item ¿Qué resolución, aprobada en agosto de $1964$ tras un supuesto ataque de lanchas torpederas norvietnamitas contra el destructor *USS Maddox*, sirvió como cheque en blanco (autoridad legal sin declaración formal de guerra) para que el presidente Lyndon B. Johnson intensificara masivamente la presencia militar en Vietnam?
\begin{enumerate}[label=\textbf{\Alph*)}]
\begin{multicols}{2}
  \item Resolución de Ginebra
  \item Tratado de la SEATO
  \item Resolución de Poderes de Guerra
  \item Resolución del Golfo de Tonkín
  \item Acta de Seguridad Nacional
  \item Pacto de París
\end{multicols}
\end{enumerate}
\vspace{0.15cm}
\item Tras su independencia en $1975$, esta excolonia portuguesa en África se hundió en una guerra civil de $27$ años. El MPLA (marxista, apoyado por Cuba y la URSS) luchó contra la UNITA (financiada por EE. UU. y la Sudáfrica del Apartheid), convirtiendo el territorio en un tablero sangriento de la Guerra Fría. ¿Cuál es la nación descrita?
\begin{enumerate}[label=\textbf{\Alph*)}]
\begin{multicols}{2}
  \item Mozambique
  \item Guinea-Bisáu
  \item Angola
  \item Cabo Verde
  \item Santo Tomé y Príncipe
  \item Guinea Ecuatorial
\end{multicols}
\end{enumerate}
\vspace{0.15cm}
\item En octubre de $1961$, la Unión Soviética detonó la "Bomba del Zar" (*Tsar Bomba*) sobre el archipiélago de Nueva Zembla. Con un rendimiento de $\approx 50 \text{ Mt}$ (megatones), este artefacto se caracterizó por ser:
\begin{enumerate}[label=\textbf{\Alph*)}]
\begin{multicols}{2}
  \item La primera bomba nuclear de fisión de la historia.
  \item La bomba termonuclear más potente jamás detonada por la humanidad.
  \item El primer misil nuclear lanzado desde un submarino.
  \item La bomba biológica responsable del ántrax en la era soviética.
  \item Un arma fallida que no logró alcanzar la masa crítica.
  \item Una ojiva química lanzada desde el espacio.
\end{multicols}
\end{enumerate}
\vspace{0.15cm}
\item Un gravísimo escándalo constitucional estalló en el gobierno de Estados Unidos en $1986$ (Administración Reagan) cuando se descubrió que funcionarios del Consejo de Seguridad Nacional vendían armas clandestinamente a un país embargado, utilizando las ganancias secretas para financiar a la contraguerrilla nicaragüense. ¿Cómo se bautizó este escándalo?
\begin{enumerate}[label=\textbf{\Alph*)}]
\begin{multicols}{2}
  \item Watergate
  \item Irán-Contra (Irangate)
  \item Los Documentos del Pentágono
  \item Escándalo Lewinsky
  \item El Asunto XYZ
  \item Operación Rápido y Furioso
\end{multicols}
\end{enumerate}
\vspace{0.15cm}
\item En mayo de $1960$, la Unión Soviética logró derribar mediante un misil tierra-aire un avión espía estadounidense $U-2$ que sobrevolaba el espacio aéreo soviético tomando fotografías de alta resolución. El piloto sobrevivió, fue capturado por el KGB y enjuiciado, arruinando la cumbre de paz de París. El nombre del piloto es:
\begin{enumerate}[label=\textbf{\Alph*)}]
\begin{multicols}{2}
  \item Chuck Yeager
  \item John Glenn
  \item Neil Armstrong
  \item Francis Gary Powers
  \item Charles Lindbergh
  \item Yuri Gagarin
\end{multicols}
\end{enumerate}
\vspace{0.15cm}
\item El Tratado "Dos más Cuatro" (firmado en Moscú en $1990$) fue la resolución diplomática magistral que puso fin oficial a las divisiones territoriales de la Segunda Guerra Mundial. A través de este documento, EE. UU., URSS, Francia y Reino Unido concedían la soberanía plena para permitir:
\begin{enumerate}[label=\textbf{\Alph*)}]
\begin{multicols}{2}
  \item La independencia de Yugoslavia
  \item La desintegración pacífica de Checoslovaquia
  \item La reunificación de Alemania (RFA y RDA)
  \item La partición de Corea
  \item La autonomía de Taiwán
  \item El retiro de la ONU de Palestina
\end{multicols}
\end{enumerate}
\vspace{0.15cm}
\item Durante la Crisis de los Misiles en Cuba ($1962$), la decisión táctica tomada por el presidente John F. Kennedy para impedir que más barcos de carga soviéticos aprovisionaran armamento atómico a la isla, eludiendo una declaración de guerra formal inmediata, consistió en:
\begin{enumerate}[label=\textbf{\Alph*)}]
\begin{multicols}{2}
  \item Bombardeo quirúrgico aéreo de los puertos cubanos.
  \item Invasión anfibia con los Marines.
  \item Un cerco o "cuarentena" naval alrededor de la isla.
  \item El lanzamiento de un misil preventivo de advertencia sobre La Habana.
  \item La minería exhaustiva del Golfo de México.
  \item El asesinato selectivo de la cúpula revolucionaria.
\end{multicols}
\end{enumerate}
\vspace{0.15cm}
\item Firmados en la base militar de Wright-Patterson en Ohio y ratificados en París en $1995$, estos acuerdos forzaron el fin de las atrocidades en los Balcanes. Estructuraron al país en dos entidades autónomas (la Federación de Bosnia-Herzegovina y la República Srpska). ¿Cómo se conocen?
\begin{enumerate}[label=\textbf{\Alph*)}]
\begin{multicols}{2}
  \item Acuerdos de Múnich
  \item Acuerdos de Brioni
  \item Tratado de Versalles
  \item Acuerdos de Oslo
  \item Acuerdos de Dayton
  \item Tratado de Lisboa
\end{multicols}
\end{enumerate}
\vspace{0.15cm}
\item El siglo XX fue testigo de uno de los mayores triunfos de la salud pública mundial en el marco institucional de las Naciones Unidas. En el año $1980$, la Organización Mundial de la Salud (OMS) certificó oficialmente la erradicación global y definitiva ($R_0 \to 0$) de una enfermedad infecciosa devastadora. ¿Cuál fue?
\begin{enumerate}[label=\textbf{\Alph*)}]
\begin{multicols}{2}
  \item El Sarampión
  \item La Poliomielitis
  \item El Cólera
  \item La Tuberculosis
  \item La Viruela
  \item El VIH/SIDA
\end{multicols}
\end{enumerate}
\vspace{0.15cm}
\item En julio de $1997$, los inversores financieros comenzaron a retirar capitales masivamente de Tailandia forzando la devaluación del baht frente al dólar. Este efecto dominó contagió las monedas y bolsas de valores de Corea del Sur, Indonesia y Malasia, requiriendo rescates multimillonarios del FMI. Este evento se conoce como:
\begin{enumerate}[label=\textbf{\Alph*)}]
\begin{multicols}{2}
  \item Lunes Negro (Crash del 87)
  \item Efecto Tequila
  \item Crisis Subprime
  \item Crisis Financiera Asiática
  \item Burbuja de las Dot-com
  \item Gran Depresión
\end{multicols}
\end{enumerate}
\vspace{0.15cm}
\item En octubre de $1983$, la administración Reagan ordenó la Operación *Urgent Fury*, consistente en el despliegue militar paracaidista y anfibio para intervenir una isla caribeña tras el asesinato de su primer ministro marxista Maurice Bishop y la presencia de constructores cubanos en un aeropuerto. ¿Qué país fue invadido?
\begin{enumerate}[label=\textbf{\Alph*)}]
\begin{multicols}{2}
  \item Cuba
  \item Haití
  \item República Dominicana
  \item Jamaica
  \item Granada
  \item Puerto Rico
\end{multicols}
\end{enumerate}
\vspace{0.15cm}
\item En materia ambiental global, la CMNUCC aprobó en el año $1997$ un protocolo pionero y jurídicamente vinculante para las naciones industrializadas que establecía cuotas obligatorias de reducción de emisiones de Gases de Efecto Invernadero (GEI) buscando mitigar que el calentamiento global $\Delta T > 0$. Su nombre es:
\begin{enumerate}[label=\textbf{\Alph*)}]
\begin{multicols}{2}
  \item Protocolo de Montreal
  \item Protocolo de Kioto
  \item Acuerdo de París
  \item Convenio de Basilea
  \item Tratado de Estocolmo
  \item Cumbre de Río
\end{multicols}
\end{enumerate}
\vspace{0.15cm}
\item Tras una intensa guerra de guerrillas iniciada en los $1960$s contra el gobierno supremacista blanco de Ian Smith, una colonia africana (antiguamente Rodesia del Sur) logró en $1980$ su independencia bajo el mando del líder Robert Mugabe. El país fue rebautizado con el nombre de:
\begin{enumerate}[label=\textbf{\Alph*)}]
\begin{multicols}{2}
  \item Zambia
  \item Botsuana
  \item Zimbabue
  \item Namibia
  \item Malaui
  \item Tanzania
\end{multicols}
\end{enumerate}
\vspace{0.15cm}
\item La arquitectura monetaria de la posguerra (Bretton Woods, $1944$) fue un choque de visiones económicas. El plan estadounidense del dólar-oro de Harry Dexter White se impuso frente a la propuesta de crear una divisa global apolítica ("Bancor") ideada por el representante británico. ¿Quién fue este insigne economista?
\begin{enumerate}[label=\textbf{\Alph*)}]
\begin{multicols}{2}
  \item Adam Smith
  \item Milton Friedman
  \item Friedrich Hayek
  \item Karl Marx
  \item John Maynard Keynes
  \item Joseph Schumpeter
\end{multicols}
\end{enumerate}
\vspace{0.15cm}
\item El mes de octubre de $1956$ vio una revuelta antisoviética masiva que exigía democracia pluripartidista y la retirada del país del Pacto de Varsovia. El líder reformista Imre Nagy formó gobierno brevemente hasta que los tanques del Ejército Rojo entraron en la capital ahogando en sangre el movimiento. ¿En qué país de Europa del Este ocurrió?
\begin{enumerate}[label=\textbf{\Alph*)}]
\begin{multicols}{2}
  \item Checoslovaquia
  \item Polonia
  \item Alemania Oriental
  \item Hungría
  \item Rumania
  \item Yugoslavia
\end{multicols}
\end{enumerate}
\vspace{0.15cm}
\item Un economista belga-estadounidense predijo en $1960$ la inevitable caída del patrón oro de Bretton Woods ($\$35 \text{ USD} = 1 \text{ oz}$). Afirmó que la liquidez global dependía del déficit comercial de EE. UU., pero ese mismo déficit socavaba la confianza en el dólar al superar las reservas de oro estadounidenses. Este axioma se llama:
\begin{enumerate}[label=\textbf{\Alph*)}]
\begin{multicols}{2}
  \item Curva de Phillips
  \item Teorema de Pareto
  \item Dilema de Triffin
  \item Efecto Multiplicador
  \item Paradoja del Ahorro
  \item Ley de Say
\end{multicols}
\end{enumerate}
\vspace{0.15cm}
\item En la Carrera Espacial, no solo importaba la maquinaria orbital, sino la supervivencia biológica autónoma en el vacío ($P = 0 \text{ atm}$). El cosmonauta soviético que pasó a la historia en $1965$ por realizar la primera caminata espacial (Actividad Extravehicular, EVA) atado a su nave Vosjod 2 por un cordón umbilical fue:
\begin{enumerate}[label=\textbf{\Alph*)}]
\begin{multicols}{2}
  \item Yuri Gagarin
  \item Vladímir Komarov
  \item Guerman Titov
  \item Alexei Leonov
  \item Pavel Popovich
  \item Valeri Poliakov
\end{multicols}
\end{enumerate}
\vspace{0.15cm}
\item Entre los años $1965$ y $1966$, el Sudeste Asiático vivió uno de los crímenes de masa de la Guerra Fría. Tras un autogolpe de Estado militar apoyado disimuladamente por EE. UU., las fuerzas armadas dirigidas por el general Suharto ejecutaron y purgaron masivamente a unos $\approx 500,000$ civiles miembros o simpatizantes del Partido Comunista. ¿En qué archipiélago sucedió?
\begin{enumerate}[label=\textbf{\Alph*)}]
\begin{multicols}{2}
  \item Filipinas
  \item Malasia
  \item Indonesia
  \item Japón
  \item Sri Lanka
  \item Taiwán
\end{multicols}
\end{enumerate}
\vspace{0.15cm}
\item Como parte del tejido de la Guerra Fría en el Istmo, la Revolución Sandinista movilizó a $95,000$ jóvenes y estudiantes al campo en $1980$. Su objetivo pedagógico y político logró, en menos de un año, reducir la tasa de analfabetismo nacional del $50.35\%$ al $12.96\%$. Este hito histórico se denominó:
\begin{enumerate}[label=\textbf{\Alph*)}]
\begin{multicols}{2}
  \item Revolución Cultural Nicaragüense
  \item Cruzada Nacional de Alfabetización
  \item Campaña del Maestro Rural
  \item Misión Sucre
  \item Proyecto de Liberación Pedagógica
  \item Juventud Comunista Educadora
\end{multicols}
\end{enumerate}
\vspace{0.15cm}
\item Las revoluciones de Europa del Este de $1989$ fueron mayormente pacíficas, con una severa excepción. En este país, el estallido social y la insubordinación militar culminaron el día de Navidad con el juicio militar sumario y la ejecución por fusilamiento de su líder dictatorial y su esposa. ¿Cuál es el dictador y el país?
\begin{enumerate}[label=\textbf{\Alph*)}]
\begin{multicols}{2}
  \item Slobodan Milošević (Yugoslavia)
  \item Todor Zhivkov (Bulgaria)
  \item Erich Honecker (Alemania Oriental)
  \item Wojciech Jaruzelski (Polonia)
  \item Nicolae Ceaușescu (Rumania)
  \item Enver Hoxha (Albania)
\end{multicols}
\end{enumerate}
\vspace{0.15cm}
\item Creada en el año $1945$ en Quebec, es una agencia especializada de las Naciones Unidas cuyo mandato específico, codificado en el sistema global, es erradicar el hambre en el mundo, garantizando la seguridad alimentaria ($Q_{\text{alimentos}} \ge \text{demanda base}$) y el desarrollo agrícola rural. ¿A qué siglas nos referimos?
\begin{enumerate}[label=\textbf{\Alph*)}]
\begin{multicols}{2}
  \item OMS
  \item FMI
  \item FAO (Organización de las Naciones Unidas para la Alimentación y la Agricultura)
  \item ACNUR
  \item UNESCO
  \item OIT
\end{multicols}
\end{enumerate}
\vspace{0.15cm}
\end{enumerate}\n
\subsubsection*{Módulo 4: Reactivos 151 - 200}

\begin{enumerate}[label=\textbf{\arabic*.}]
\setcounter{enumi}{150}
\item El $14$ de mayo de $1948$, horas antes de que expirara el Mandato Británico sobre Palestina, se leyó la Declaración de Independencia del Estado de Israel. ¿Qué líder sionista histórico, quien se convirtió en el primer Primer Ministro del país, proclamó esta declaración?
\begin{enumerate}[label=\textbf{\Alph*)}]
\begin{multicols}{2}
  \item Chaim Weizmann
  \item Golda Meir
  \item Yitzhak Rabin
  \item David Ben-Gurión
  \item Menájem Beguin
  \item Ariel Sharón
\end{multicols}
\end{enumerate}
\vspace{0.15cm}
\item La Guerra de Corea culminó temporalmente en $1953$ con la firma de un armisticio (no un tratado de paz definitivo), estableciendo una Zona Desmilitarizada (ZDM) de $\approx 4 \text{ km}$ de ancho cerca del paralelo $\lambda = 38^{\circ}$. ¿En qué aldea fronteriza se firmó este documento militar?
\begin{enumerate}[label=\textbf{\Alph*)}]
\begin{multicols}{2}
  \item Inchon
  \item Seúl
  \item Pionyang
  \item Panmunjom
  \item Wonsan
  \item Kaesong
\end{multicols}
\end{enumerate}
\vspace{0.15cm}
\item En septiembre de $1960$, cinco naciones productoras de petróleo se reunieron en la Conferencia de Bagdad para fundar la OPEP, buscando recuperar el control sobre los precios de su crudo frente a las "Siete Hermanas" (grandes corporaciones occidentales). ¿Cuáles fueron estas $n=5$ naciones fundadoras?
\begin{enumerate}[label=\textbf{\Alph*)}]
\begin{multicols}{2}
  \item Arabia Saudita, Irán, Irak, Kuwait y Venezuela.
  \item Arabia Saudita, Egipto, Siria, Jordania e Irán.
  \item Venezuela, México, Ecuador, Brasil y Argentina.
  \item Irán, Irak, Qatar, Emiratos Árabes Unidos y Libia.
  \item Arabia Saudita, Nigeria, Argelia, Libia y Angola.
  \item Estados Unidos, Reino Unido, URSS, Irán y Arabia Saudita.
\end{multicols}
\end{enumerate}
\vspace{0.15cm}
\item El $5$ de marzo de $1946$, en un célebre discurso pronunciado en Fulton (Misuri), se acuñó la metáfora geopolítica definitiva que describiría la división ideológica y física de Europa durante la Guerra Fría: "Desde Stettin en el Báltico, hasta Trieste en el Adriático, ha caído sobre el continente una...".
\begin{enumerate}[label=\textbf{\Alph*)}]
\begin{multicols}{2}
  \item Muralla Soviética
  \item Cortina de Hierro (Telón de Acero)
  \item Barrera Antifascista
  \item Línea Maginot Ideológica
  \item Frontera Roja
  \item Sombra Bolchevique
\end{multicols}
\end{enumerate}
\vspace{0.15cm}
\item El apogeo de la "Distensión" (Détente) espacial se materializó en $1975$ con el acoplamiento orbital ($v \approx 28,000 \text{ km/h}$) de dos naves pertenecientes a superpotencias rivales, permitiendo un apretón de manos histórico en el espacio. ¿Cómo se denominó este proyecto conjunto?
\begin{enumerate}[label=\textbf{\Alph*)}]
\begin{multicols}{2}
  \item Proyecto Mercury-Vostok
  \item Estación Espacial Internacional (ISS)
  \item Misión Skylab-Salyut
  \item Proyecto de Pruebas Apolo-Soyuz
  \item Misión Gemini-Cosmos
  \item Telescopio Hubble Conjunto
\end{multicols}
\end{enumerate}
\vspace{0.15cm}
\item Durante la década de $1950$, el Imperio Británico enfrentó una feroz rebelión anticolonial en África Oriental protagonizada por guerrilleros de la etnia kikuyu, quienes luchaban contra la expropiación de sus tierras por colonos blancos. Este conflicto se conoce históricamente como:
\begin{enumerate}[label=\textbf{\Alph*)}]
\begin{multicols}{2}
  \item Rebelión de los Bóxers
  \item Guerra de los Bóeres
  \item Rebelión del Mau Mau (Kenia)
  \item Levantamiento del Mahdi (Sudán)
  \item Guerra Zulú
  \item Independencia del Tanganica
\end{multicols}
\end{enumerate}
\vspace{0.15cm}
\item El $14$ de agosto de $1947$, un día antes que la India, se fundó un nuevo Estado soberano de mayoría musulmana, concebido como un refugio religioso tras el colapso del Raj británico. ¿Qué líder político y abogado fue el fundador de Pakistán liderando la Liga Musulmana?
\begin{enumerate}[label=\textbf{\Alph*)}]
\begin{multicols}{2}
  \item Mahatma Gandhi
  \item Jawaharlal Nehru
  \item Subhas Chandra Bose
  \item Muhammad Ali Jinnah
  \item Liaquat Ali Khan
  \item Zulfikar Ali Bhutto
\end{multicols}
\end{enumerate}
\vspace{0.15cm}
\item Días antes de la masiva movilización ciudadana que culminaría con la caída del Muro de Berlín en noviembre de $1989$, el rígido líder histórico de la República Democrática Alemana (RDA) fue forzado a dimitir por su propio partido. ¿De quién se trata?
\begin{enumerate}[label=\textbf{\Alph*)}]
\begin{multicols}{2}
  \item Alexander Dubček
  \item Walter Ulbricht
  \item Erich Honecker
  \item Egon Krenz
  \item Willy Brandt
  \item Helmut Kohl
\end{multicols}
\end{enumerate}
\vspace{0.15cm}
\item En $1954$, Estados Unidos promovió la creación de una alianza militar regional para contener el comunismo en el Sudeste Asiático (una especie de "OTAN asiática"). Sus siglas en inglés corresponden a *Southeast Asia Treaty Organization*. ¿Qué acrónimo la identificaba?
\begin{enumerate}[label=\textbf{\Alph*)}]
\begin{multicols}{2}
  \item CENTO
  \item ANZUS
  \item SEATO
  \item ASEAN
  \item APEC
  \item Pacto de Varsovia Asiático
\end{multicols}
\end{enumerate}
\vspace{0.15cm}
\item En la jerga estratégica de la Carrera Armamentista, un misil capaz de ser lanzado desde un continente para impactar en otro con ojivas termonucleares, requiriendo un rango de vuelo balístico suborbital mínimo de $r \ge 5,500 \text{ km}$, recibe las siglas de:
\begin{enumerate}[label=\textbf{\Alph*)}]
\begin{multicols}{2}
  \item SLBM (Submarine-Launched Ballistic Missile)
  \item SAM (Surface-to-Air Missile)
  \item ICBM (Intercontinental Ballistic Missile)
  \item MRBM (Medium-Range Ballistic Missile)
  \item ALCM (Air-Launched Cruise Missile)
  \item ABM (Anti-Ballistic Missile)
\end{multicols}
\end{enumerate}
\vspace{0.15cm}
\item Entre $1934$ y $1935$, el Ejército Rojo del Partido Comunista de China realizó un monumental repliegue táctico militar, recorriendo $\approx 9,000 \text{ km}$ a pie a través de montañas y pantanos para escapar del cerco del Kuomintang. Este evento forjó el liderazgo absoluto de Mao Zedong.
\begin{enumerate}[label=\textbf{\Alph*)}]
\begin{multicols}{2}
  \item La Larga Marcha
  \item El Gran Salto Adelante
  \item La Revolución Cultural
  \item La Campaña de las Cien Flores
  \item El Levantamiento de Wuchang
  \item La Expedición del Norte
\end{multicols}
\end{enumerate}
\vspace{0.15cm}
\item En noviembre de $1979$, estudiantes islámicos radicalizados asaltaron la embajada de Estados Unidos en Teherán, secuestrando a $52$ diplomáticos y ciudadanos estadounidenses durante $444$ días. Este evento destruyó políticamente las opciones de reelección del presidente:
\begin{enumerate}[label=\textbf{\Alph*)}]
\begin{multicols}{2}
  \item Richard Nixon
  \item Gerald Ford
  \item Jimmy Carter
  \item Ronald Reagan
  \item George H.W. Bush
  \item Bill Clinton
\end{multicols}
\end{enumerate}
\vspace{0.15cm}
\item En el año $1952$, América del Sur vivió una revolución social profunda donde el Movimiento Nacionalista Revolucionario (MNR) nacionalizó las grandes minas de estaño, decretó el sufragio universal y ejecutó una masiva reforma agraria. ¿En qué país ocurrió?
\begin{enumerate}[label=\textbf{\Alph*)}]
\begin{multicols}{2}
  \item Perú
  \item Ecuador
  \item Bolivia
  \item Paraguay
  \item Chile
  \item Colombia
\end{multicols}
\end{enumerate}
\vspace{0.15cm}
\item Tras prolongadas y frustrantes negociaciones, en enero de $1973$ se firmó el tratado que ponía fin a la intervención militar directa de Estados Unidos en el conflicto de Vietnam, permitiendo el retiro de sus tropas de combate bajo la fórmula de "Paz con Honor" de Nixon.
\begin{enumerate}[label=\textbf{\Alph*)}]
\begin{multicols}{2}
  \item Acuerdos de Ginebra
  \item Tratado de Versalles
  \item Acuerdos de París
  \item Tratado de Washington
  \item Convención de La Haya
  \item Protocolo de Saigón
\end{multicols}
\end{enumerate}
\vspace{0.15cm}
\item La diplomacia multilateral creó en $1956$ un nuevo mecanismo de intervención internacional: la Fuerza de Emergencia de las Naciones Unidas (UNEF), coloquialmente conocida como "Cascos Azules". Su primer despliegue histórico fue para interponerse pacíficamente en:
\begin{enumerate}[label=\textbf{\Alph*)}]
\begin{multicols}{2}
  \item La Guerra de Corea
  \item La Crisis de Suez
  \item La Guerra Civil en el Líbano
  \item El conflicto de Cachemira
  \item La Guerra de los Balcanes
  \item La Revolución Húngara
\end{multicols}
\end{enumerate}
\vspace{0.15cm}
\item En $1972$, Richard Nixon se convirtió en el primer presidente de EE. UU. en visitar Moscú. Allí firmó con Leonid Brézhnev el Tratado sobre Misiles Antibalísticos (ABM) y el acuerdo interino que congelaba el número de lanzadores estratégicos intercontinentales. Este acuerdo se conoce como:
\begin{enumerate}[label=\textbf{\Alph*)}]
\begin{multicols}{2}
  \item SALT I
  \item START I
  \item Tratado de Tlatelolco
  \item Tratado INF
  \item Acuerdo de Helsinki
  \item Tratado de No Proliferación (TNP)
\end{multicols}
\end{enumerate}
\vspace{0.15cm}
\item El $1$ de enero de $1993$, las repúblicas checa y eslovaca disolvieron pacíficamente su Estado federal común, dividiendo la nación en la República Checa y la República Eslovaca. En la geopolítica y diplomacia europea contemporánea, este suceso cívico se denomina:
\begin{enumerate}[label=\textbf{\Alph*)}]
\begin{multicols}{2}
  \item El Divorcio de Terciopelo
  \item La Balcanización
  \item La Secesión de Praga
  \item El Pacto de Varsovia
  \item La Primavera de Bratislava
  \item El Cisma Bohemio
\end{multicols}
\end{enumerate}
\vspace{0.15cm}
\item En septiembre de $1970$, el Rey Hussein de Jordania ordenó a sus fuerzas armadas expulsar violentamente a la Organización para la Liberación de Palestina (OLP) de su territorio, debido a que esta operaba como un "Estado dentro de un Estado". Este sangriento conflicto intrarregional se conoce como:
\begin{enumerate}[label=\textbf{\Alph*)}]
\begin{multicols}{2}
  \item Septiembre Negro
  \item La Nakba
  \item La Guerra de los Seis Días
  \item La Primera Intifada
  \item La Revolución de los Cedros
  \item El Levantamiento Druso
\end{multicols}
\end{enumerate}
\vspace{0.15cm}
\item Durante los años $1980$s, en un esfuerzo por evitar la militarización total e intervención directa de EE. UU. en las guerras civiles de Centroamérica, Colombia, México, Panamá y Venezuela formaron en $1983$ una instancia diplomática conjunta conocida como:
\begin{enumerate}[label=\textbf{\Alph*)}]
\begin{multicols}{2}
  \item Organización de Estados Americanos (OEA)
  \item Grupo de Río
  \item Grupo de Contadora
  \item ALBA (Alianza Bolivariana)
  \item Mercado Común Centroamericano (MCCA)
  \item Sistema de la Integración Centroamericana (SICA)
\end{multicols}
\end{enumerate}
\vspace{0.15cm}
\item Tras la invasión soviética de Afganistán en $1979$, el presidente de EE. UU. declaró oficialmente ante el Congreso que cualquier intento de una fuerza extranjera de obtener el control del Golfo Pérsico sería considerado un ataque a los intereses vitales estadounidenses y repelido mediante fuerza militar. Esta política exterior fue la:
\begin{enumerate}[label=\textbf{\Alph*)}]
\begin{multicols}{2}
  \item Doctrina Monroe
  \item Doctrina Truman
  \item Doctrina Reagan
  \item Doctrina Carter
  \item Doctrina Nixon
  \item Doctrina del Destino Manifiesto
\end{multicols}
\end{enumerate}
\vspace{0.15cm}
\item En abril de $1971$, la URSS logró colocar en la Órbita Terrestre Baja (LEO) la primera estación espacial tripulada de la historia humana, sentando las bases para misiones de permanencia biológica prolongada. El nombre de este primer programa de estaciones orbitales fue:
\begin{enumerate}[label=\textbf{\Alph*)}]
\begin{multicols}{2}
  \item Skylab
  \item Mir
  \item Salyut
  \item Vostok
  \item Soyuz
  \item Sputnik
\end{multicols}
\end{enumerate}
\vspace{0.15cm}
\item Firmado inicialmente en $1985$ por cinco países de la CEE, este acuerdo suprimió gradualmente los controles en las fronteras interiores de Europa, creando un espacio de libre circulación transnacional donde un ciudadano puede cruzar de país sin mostrar el pasaporte.
\begin{enumerate}[label=\textbf{\Alph*)}]
\begin{multicols}{2}
  \item Tratado de Maastricht
  \item Acuerdo de Schengen
  \item Tratado de Roma
  \item Tratado de Lisboa
  \item Tratado de Ámsterdam
  \item Pacto de Varsovia
\end{multicols}
\end{enumerate}
\vspace{0.15cm}
\item En agosto de $1982$, el ministro de Hacienda de un país latinoamericano anunció en Washington que su nación no podía cumplir con el servicio de su deuda externa. Este default soberano desató la "Crisis de la Deuda" (o Década Perdida) en la región. El país fue:
\begin{enumerate}[label=\textbf{\Alph*)}]
\begin{multicols}{2}
  \item Brasil
  \item Argentina
  \item México
  \item Venezuela
  \item Chile
  \item Perú
\end{multicols}
\end{enumerate}
\vspace{0.15cm}
\item Creado en la conferencia de Bretton Woods ($1944$), el Banco Internacional de Reconstrucción y Fomento (BIRF), que hoy constituye el núcleo del Banco Mundial, tuvo como su primera y principal misión institucional original:
\begin{enumerate}[label=\textbf{\Alph*)}]
\begin{multicols}{2}
  \item Otorgar microcréditos a campesinos en América Latina.
  \item Financiar programas espaciales conjuntos entre EE. UU. y Europa.
  \item Proveer préstamos a largo plazo para la reconstrucción de la infraestructura física de Europa occidental destruida por la SGM.
  \item Comprar armas nucleares a la URSS para su desmantelamiento.
  \item Regular el mercado global del café y el azúcar.
  \item Financiar exclusivamente la transición económica de China.
\end{multicols}
\end{enumerate}
\vspace{0.15cm}
\item El $1$ de septiembre de $1983$, la tensión de la Guerra Fría llegó a niveles críticos cuando aviones cazas interceptores soviéticos Su-15 derribaron un avión comercial civil (con $269$ personas a bordo) alegando que era una aeronave espía adentrándose en espacio aéreo restringido.
\begin{enumerate}[label=\textbf{\Alph*)}]
\begin{multicols}{2}
  \item Vuelo Pan Am 103 (Lockerbie)
  \item Vuelo Korean Air Lines 007 (KAL 007)
  \item Vuelo Malaysia Airlines 17 (MH17)
  \item Vuelo Iran Air 655
  \item Vuelo Air France 447
  \item Vuelo TWA 800
\end{multicols}
\end{enumerate}
\vspace{0.15cm}
\item En Sudáfrica, la lucha cívica y armada contra el régimen racista del Apartheid fue liderada hegemónicamente por el partido político al que pertenecía Nelson Mandela. Creado en $1912$, este movimiento fue legalizado nuevamente en $1990$. Su nombre es:
\begin{enumerate}[label=\textbf{\Alph*)}]
\begin{multicols}{2}
  \item Congreso Nacional Africano (ANC)
  \item Partido Nacional (NP)
  \item Partido de la Libertad Inkatha
  \item Frente Patriótico Ruandés
  \item Movimiento Popular de Liberación (MPLA)
  \item Partido del Trabajo Afrikáner
\end{multicols}
\end{enumerate}
\vspace{0.15cm}
\item En $1989$, el gobierno soviético de Mijaíl Gorbachov declaró, con una metáfora humorística basada en una famosa canción estadounidense, que los países del Pacto de Varsovia eran libres de elegir su propio camino político interno sin temor a una invasión de tanques rusos. La prensa occidental la bautizó como:
\begin{enumerate}[label=\textbf{\Alph*)}]
\begin{multicols}{2}
  \item Doctrina Monroe
  \item Doctrina de Represalia Masiva
  \item Doctrina Sinatra ("My Way")
  \item Doctrina Brézhnev
  \item Política del Gran Garrote
  \item Doctrina de Contención
\end{multicols}
\end{enumerate}
\vspace{0.15cm}
\item En febrero de $1967$, los Estados de América Latina y el Caribe firmaron un tratado pionero a nivel mundial que prohibió el ensayo, uso, fabricación, producción o adquisición de armas nucleares en toda la región ($n_{\text{armas}} = 0$). Este acuerdo se conoce como:
\begin{enumerate}[label=\textbf{\Alph*)}]
\begin{multicols}{2}
  \item Tratado de Versalles
  \item Tratado de No Proliferación (TNP)
  \item Tratado de Tlatelolco
  \item Acuerdos de Helsinki
  \item Pacto de San José
  \item Protocolo de Kioto
\end{multicols}
\end{enumerate}
\vspace{0.15cm}
\item En $1978$, tras intensas negociaciones secretas de doce días en la residencia presidencial de retiro en Maryland, el presidente estadounidense Jimmy Carter logró que Menájem Beguin (Israel) y Anwar el-Sadat (Egipto) firmaran el acuerdo marco para la paz en Medio Oriente denominado:
\begin{enumerate}[label=\textbf{\Alph*)}]
\begin{multicols}{2}
  \item Acuerdos de Oslo
  \item Acuerdos de Camp David
  \item Resolución 242 de la ONU
  \item Acuerdos de Ginebra
  \item Tratado del Sinaí
  \item Pacto Tripartito
\end{multicols}
\end{enumerate}
\vspace{0.15cm}
\item A finales de enero de $1968$, durante el alto el fuego pactado por la festividad del Año Nuevo lunar, el Viet Cong y el Ejército de Vietnam del Norte lanzaron un masivo y sorpresivo ataque simultáneo contra más de $100$ ciudades survietnamitas. Aunque militarmente fracasaron, destruyeron el apoyo público estadounidense a la guerra. Fue la:
\begin{enumerate}[label=\textbf{\Alph*)}]
\begin{multicols}{2}
  \item Batalla de Dien Bien Phu
  \item Ofensiva del Tet
  \item Caída de Saigón
  \item Operación Rolling Thunder
  \item Masacre de My Lai
  \item Batalla de Hue
\end{multicols}
\end{enumerate}
\vspace{0.15cm}
\item Bajo la ocupación estadounidense (General MacArthur), Japón promulgó en $1947$ una nueva Constitución redactada con ideales pacifistas. En su famoso "Artículo $9$", el Estado japonés renuncia formal e indefinidamente a:
\begin{enumerate}[label=\textbf{\Alph*)}]
\begin{multicols}{2}
  \item Poseer un emperador.
  \item La guerra como derecho soberano y a la amenaza o uso de la fuerza armada para resolver disputas internacionales.
  \item Exportar bienes manufacturados a Estados Unidos.
  \item La energía nuclear civil.
  \item Participar en la Asamblea General de la ONU.
  \item Acoger bases militares extranjeras.
\end{multicols}
\end{enumerate}
\vspace{0.15cm}
\item El $16$ de julio de $1945$, a las $5:29\text{ a.m.}$, en el desierto de Jornada del Muerto en Nuevo México, la ciencia humana transformó la masa en energía destructiva según la equivalencia $E = mc^2$, detonando el primer artefacto nuclear de la historia bajo el Proyecto Manhattan. La prueba se codificó como:
\begin{enumerate}[label=\textbf{\Alph*)}]
\begin{multicols}{2}
  \item Proyecto Trinity (Trinidad)
  \item Operación Crossroads
  \item Little Boy
  \item Fat Man
  \item Castle Bravo
  \item Zar Bomba
\end{multicols}
\end{enumerate}
\vspace{0.15cm}
\item En $1987$, buscando erradicar la violencia patrocinada por la Guerra Fría en el istmo, los cinco presidentes centroamericanos firmaron en Guatemala un acuerdo de paz que promovía amnistías, cese al fuego y democratización. Su arquitecto, el presidente de Costa Rica Óscar Arias, recibió el Nobel de la Paz.
\begin{enumerate}[label=\textbf{\Alph*)}]
\begin{multicols}{2}
  \item Acuerdos de Chapultepec
  \item Acuerdo de Esquipulas II
  \item Tratado de Contadora
  \item Pacto de San José
  \item Protocolo de Tegucigalpa
  \item Tratado de Managua
\end{multicols}
\end{enumerate}
\vspace{0.15cm}
\item El $4$ de julio de $1946$, Estados Unidos cumplió la promesa hecha años antes y renunció formalmente a su soberanía sobre un vasto archipiélago asiático, reconociendo mediante el Tratado de Manila la independencia de:
\begin{enumerate}[label=\textbf{\Alph*)}]
\begin{multicols}{2}
  \item Indonesia
  \item Taiwán
  \item Filipinas
  \item Vietnam del Sur
  \item Malasia
  \item Hawái
\end{multicols}
\end{enumerate}
\vspace{0.15cm}
\item Durante la Guerra Fría, las agencias de espionaje occidentales sufrieron una profunda infiltración soviética. Una célebre red de espías ($n=5$) reclutados por el KGB durante sus años universitarios en la década de los $30$ logró pasar secretos atómicos y diplomáticos británicos desde dentro del MI6 y la Cancillería. Se les conoce como:
\begin{enumerate}[label=\textbf{\Alph*)}]
\begin{multicols}{2}
  \item Los Cinco de Cambridge
  \item El Anillo de Los Álamos
  \item La Red Enigma
  \item Los Infiltrados de Oxford
  \item El Grupo de Múnich
  \item La Conexión Francesa
\end{multicols}
\end{enumerate}
\vspace{0.15cm}
\item El XX Congreso del Partido Comunista de la Unión Soviética ($1956$) fue el escenario de un momento de ruptura, cuando Nikita Jruschov leyó el "Discurso Secreto" ante los delegados, inaugurando la política de Desestalinización. ¿De qué se trataba este discurso?
\begin{enumerate}[label=\textbf{\Alph*)}]
\begin{multicols}{2}
  \item Anunciaba la disolución del Pacto de Varsovia.
  \item Denunciaba el "culto a la personalidad" y los crímenes de Estado masivos (purgas) cometidos por Iósif Stalin.
  \item Declaraba la guerra nuclear inmediata a Estados Unidos.
  \item Proclamaba la independencia de Ucrania y Bielorrusia.
  \item Promovía la reinstauración de la monarquía del Zar.
  \item Adoptaba formalmente la economía de libre mercado capitalista.
\end{multicols}
\end{enumerate}
\vspace{0.15cm}
\item La Comisión Económica para América Latina (CEPAL), bajo el liderazgo intelectual del argentino Raúl Prebisch, impulsó un modelo macroeconómico entre las décadas de $1950$ y $1970$ basado en aranceles protectores y subsidios para crear industria local y dejar de depender de manufacturas extranjeras. Este modelo es:
\begin{enumerate}[label=\textbf{\Alph*)}]
\begin{multicols}{2}
  \item Laissez-faire absoluto
  \item Economía de Mercado Socialista
  \item Modelo Agroexportador Monopólico
  \item Industrialización por Sustitución de Importaciones (ISI)
  \item Patrón Oro Clásico
  \item Modelo Neoliberal del Consenso de Washington
\end{multicols}
\end{enumerate}
\vspace{0.15cm}
\item El prolongado y mortífero conflicto asimétrico y de trincheras entre la República Islámica de Irán y el Irak de Saddam Hussein finalizó en $1988$ (tras el agotamiento mutuo y el uso de gas mostaza por Irak) al aceptar ambos la resolución pacífica de la ONU. ¿Cuál fue el instrumento internacional que marcó el alto el fuego?
\begin{enumerate}[label=\textbf{\Alph*)}]
\begin{multicols}{2}
  \item Resolución 242 de la ONU
  \item Tratado de Versalles
  \item Resolución 598 del Consejo de Seguridad de la ONU
  \item Acuerdos de Camp David
  \item Tratado del Golfo
  \item Tratado de no agresión de Bagdad
\end{multicols}
\end{enumerate}
\vspace{0.15cm}
\item En febrero de $1989$, el gobierno de Venezuela (Carlos Andrés Pérez) implementó un severo paquete de medidas de austeridad recetadas por el FMI (aumento de la gasolina en un $100\%$ e incremento del transporte público). Esto detonó una ola de saqueos y protestas civiles masivas que fueron reprimidas violentamente por el ejército. A este suceso se le denomina:
\begin{enumerate}[label=\textbf{\Alph*)}]
\begin{multicols}{2}
  \item El Bogotazo
  \item La Noche de los Lápices
  \item El Caracazo
  \item El Cordobazo
  \item La Primavera Venezolana
  \item La Marcha de los Cuatro Suyos
\end{multicols}
\end{enumerate}
\vspace{0.15cm}
\item La balística nuclear evolucionó en los $1970$s hacia una tecnología donde un solo misil balístico intercontinental ya no contenía una única bomba, sino que su cabeza albergaba múltiples vehículos de reentrada (cada uno con una ojiva atómica) que podían dirigirse a diferentes ciudades simultáneamente. Se conoce con las siglas:
\begin{enumerate}[label=\textbf{\Alph*)}]
\begin{multicols}{2}
  \item ABM (Anti-Ballistic Missile)
  \item MAD (Mutual Assured Destruction)
  \item MIRV (Multiple Independently targetable Reentry Vehicle)
  \item EMP (Electromagnetic Pulse)
  \item SLBM (Submarine Launched Ballistic Missile)
  \item SDI (Strategic Defense Initiative)
\end{multicols}
\end{enumerate}
\vspace{0.15cm}
\item El $28$ de marzo de $1979$, la industria de la energía nuclear civil de EE. UU. sufrió su peor accidente (nivel $5$ en la escala INES) cuando un reactor PWR sufrió una fusión parcial del núcleo debido a fallos mecánicos y errores humanos, paralizando la construcción de nuevas plantas en el país por décadas. Ocurrió en:
\begin{enumerate}[label=\textbf{\Alph*)}]
\begin{multicols}{2}
  \item Chernóbil (Ucrania)
  \item Fukushima (Japón)
  \item Three Mile Island (Pensilvania)
  \item Sellafield (Reino Unido)
  \item Los Álamos (Nuevo México)
  \item Kyshtym (URSS)
\end{multicols}
\end{enumerate}
\vspace{0.15cm}
\item Tras la caída del Muro de Berlín, el politólogo estadounidense Francis Fukuyama publicó en $1989$ un ensayo sumamente influyente. En él, argumentaba que la derrota del fascismo y luego del comunismo demostraba que la democracia liberal occidental era el punto culminante de la evolución sociocultural y política de la humanidad. Su tesis se denomina:
\begin{enumerate}[label=\textbf{\Alph*)}]
\begin{multicols}{2}
  \item El Choque de Civilizaciones
  \item El Fin de la Historia
  \item La Tercera Vía
  \item El Destino Manifiesto Global
  \item La Sociedad Abierta
  \item El Consenso de Washington Definitivo
\end{multicols}
\end{enumerate}
\vspace{0.15cm}
\item En la Federación Rusa postsoviética, el presidente Borís Yeltsin ordenó en $1994$ la invasión militar de una pequeña república caucásica de mayoría musulmana que había declarado su independencia. La dura guerra urbana (especialmente en la capital, Grozni) resultó en una humillante retirada rusa en $1996$. ¿De qué guerra se trata?
\begin{enumerate}[label=\textbf{\Alph*)}]
\begin{multicols}{2}
  \item Primera Guerra Chechena
  \item Guerra de Osetia del Sur
  \item Conflicto del Alto Karabaj
  \item Guerra Ruso-Georgiana
  \item Intervención en Afganistán
  \item Guerra Civil Tayika
\end{multicols}
\end{enumerate}
\vspace{0.15cm}
\item En el Castillo de Chapultepec (Ciudad de México), bajo el auspicio de la ONU, el gobierno de un país centroamericano y la guerrilla marxista firmaron en enero de $1992$ un tratado de paz que desmovilizó a los insurgentes para convertirlos en partido político y reformó profundamente las fuerzas armadas locales. Este fue el fin de la guerra civil de:
\begin{enumerate}[label=\textbf{\Alph*)}]
\begin{multicols}{2}
  \item Nicaragua
  \item Guatemala
  \item El Salvador
  \item Honduras
  \item Colombia
  \item Panamá
\end{multicols}
\end{enumerate}
\vspace{0.15cm}
\item La Conferencia de Bandung ($1955$), cuna del Movimiento de Países No Alineados (MNOAL), fue impulsada por gigantes del liderazgo poscolonial asiático y africano, incluyendo a Sukarno (Indonesia), Gamal Abdel Nasser (Egipto) y el primer jefe de gobierno de la India independiente:
\begin{enumerate}[label=\textbf{\Alph*)}]
\begin{multicols}{2}
  \item Mahatma Gandhi
  \item Indira Gandhi
  \item Jawaharlal Nehru
  \item Muhammad Ali Jinnah
  \item Subhas Chandra Bose
  \item Rajiv Gandhi
\end{multicols}
\end{enumerate}
\vspace{0.15cm}
\item En marzo de $1957$, la región del África subsahariana vivió un hito descolonizador cuando la colonia británica de la Costa de Oro logró su independencia, cambiando su nombre. Su líder, Kwame Nkrumah, se convirtió en un pilar del panafricanismo global. ¿De qué país africano se trata?
\begin{enumerate}[label=\textbf{\Alph*)}]
\begin{multicols}{2}
  \item Sudáfrica
  \item Ghana
  \item Nigeria
  \item Kenia
  \item Etiopía
  \item República Democrática del Congo
\end{multicols}
\end{enumerate}
\vspace{0.15cm}
\item El $6$ de octubre de $1973$, coincidiendo deliberadamente con el día más sagrado del calendario judío (Día de la Expiación), las fuerzas armadas de Egipto y Siria lanzaron un ataque militar sorpresa coordinado contra las posiciones de Israel en el Sinaí y el Golán. ¿Cuál es el nombre historiográfico de este conflicto?
\begin{enumerate}[label=\textbf{\Alph*)}]
\begin{multicols}{2}
  \item Guerra de los Seis Días
  \item Primera Guerra del Golfo
  \item Guerra de la Independencia de Israel
  \item Guerra de Yom Kipur
  \item Crisis de Suez
  \item Operación Tormenta del Desierto
\end{multicols}
\end{enumerate}
\vspace{0.15cm}
\item Durante la guerra mencionada en el reactivo anterior ($1973$), el Estado de Israel estuvo al borde de la derrota en las primeras $48$ horas. La contraofensiva acorazada y el puente aéreo de suministros desde EE. UU. (Operación *Nickel Grass*) fueron coordinados por la "Dama de Hierro" israelí, quien era la Primera Ministra en ese momento:
\begin{enumerate}[label=\textbf{\Alph*)}]
\begin{multicols}{2}
  \item Indira Gandhi
  \item Margaret Thatcher
  \item Golda Meir
  \item Angela Merkel
  \item Yael Dayan
  \item Benazir Bhutto
\end{multicols}
\end{enumerate}
\vspace{0.15cm}
\item El Movimiento de Países No Alineados (MNOAL) celebró en el año $1979$ su Sexta Cumbre. Fue la primera vez que esta coalición del Tercer Mundo se reunió en el hemisferio occidental (América Latina), generando controversia por la sede elegida, dado su explícito alineamiento con el bloque soviético en la práctica. La cumbre fue en:
\begin{enumerate}[label=\textbf{\Alph*)}]
\begin{multicols}{2}
  \item Buenos Aires, Argentina
  \item La Habana, Cuba
  \item Ciudad de México, México
  \item Managua, Nicaragua
  \item Caracas, Venezuela
  \item Santiago, Chile
\end{multicols}
\end{enumerate}
\vspace{0.15cm}
\item En diciembre de $1991$, semanas antes de la disolución formal y total de la URSS, los líderes de Rusia, Ucrania y Bielorrusia firmaron el Tratado de Belavezha. En él, declaraban que la URSS "dejaba de existir como sujeto de derecho internacional" y acordaban sustituirla geopolíticamente por una organización laxa llamada:
\begin{enumerate}[label=\textbf{\Alph*)}]
\begin{multicols}{2}
  \item Pacto de Varsovia
  \item Organización del Tratado de la Seguridad Colectiva (OTSC)
  \item Unión Económica Euroasiática
  \item Comunidad de Estados Independientes (CEI)
  \item Federación de Repúblicas Eslavas
  \item Liga de Naciones del Este
\end{multicols}
\end{enumerate}
\vspace{0.15cm}
\end{enumerate}\n
\subsubsection*{Módulo 5: Reactivos 201 - 250}

\begin{enumerate}[label=\textbf{\arabic*.}]
\setcounter{enumi}{200}
\item En el año $1946$, un diplomático estadounidense destinado en Moscú envió el "Telegrama Largo" (de $\approx 8,000$ palabras) a Washington. En este documento argumentaba que la URSS era intrínsecamente expansionista y que Estados Unidos debía adoptar una política de "contención" ($C$) sistemática. ¿Quién fue el autor de este influyente texto geopolítico?
\begin{enumerate}[label=\textbf{\Alph*)}]
\begin{multicols}{2}
  \item Henry Kissinger
  \item John Foster Dulles
  \item George F. Kennan
  \item Dean Acheson
  \item George Marshall
  \item Robert McNamara
\end{multicols}
\end{enumerate}
\vspace{0.15cm}
\item Tras el colapso de la URSS, la Federación Rusa bajo el mandato de Borís Yeltsin aplicó una política macroeconómica diseñada por economistas occidentales. Consistió en la liberalización instantánea de los precios ($P$) y la privatización acelerada de las empresas estatales, generando una hiperinflación y pobreza masiva a corto plazo. Este modelo se denomina:
\begin{enumerate}[label=\textbf{\Alph*)}]
\begin{multicols}{2}
  \item Economía Social de Mercado
  \item Terapia de Choque (Shock Therapy)
  \item Nueva Política Económica (NEP)
  \item Perestroika Financiera
  \item Plan Quinquenal Capitalista
  \item Keynesianismo Estructural
\end{multicols}
\end{enumerate}
\vspace{0.15cm}
\item Entre abril y junio de $1989$, el gobierno de la República Popular China enfrentó masivas protestas lideradas por estudiantes universitarios que exigían reformas democráticas y libertad de expresión. Estas movilizaciones fueron brutalmente aplastadas por el Ejército Popular de Liberación en:
\begin{enumerate}[label=\textbf{\Alph*)}]
\begin{multicols}{2}
  \item La Plaza Roja
  \item La Plaza de Tiananmén
  \item La Plaza de las Tres Culturas
  \item El Zócalo
  \item La Plaza Tahrir
  \item El Palacio de Invierno
\end{multicols}
\end{enumerate}
\vspace{0.15cm}
\item Durante los Juegos Olímpicos de verano de $1972$, la organización terrorista palestina "Septiembre Negro" secuestró y asesinó a un grupo de atletas israelíes ($n=11$) en la villa olímpica, alterando para siempre los protocolos de seguridad de los eventos globales. ¿En qué ciudad ocurrieron estos hechos?
\begin{enumerate}[label=\textbf{\Alph*)}]
\begin{multicols}{2}
  \item Múnich
  \item Berlín
  \item Moscú
  \item Los Ángeles
  \item Montreal
  \item Ciudad de México
\end{multicols}
\end{enumerate}
\vspace{0.15cm}
\item Durante la Guerra de Vietnam, la Fuerza Aérea de Estados Unidos roció millones de galones de un herbicida defoliante táctico sobre la selva (Operación *Ranch Hand*) para destruir la cobertura vegetal del Viet Cong. Este químico, altamente tóxico y causante de mutaciones genéticas por su contenido de dioxina (TCDD), se conoce como:
\begin{enumerate}[label=\textbf{\Alph*)}]
\begin{multicols}{2}
  \item Napalm
  \item Gas Mostaza
  \item Agente Naranja
  \item Fósforo Blanco
  \item Gas Sarín
  \item Zyklon B
\end{multicols}
\end{enumerate}
\vspace{0.15cm}
\item Tras el golpe de Estado de $1973$ en Chile, la dictadura militar implementó un modelo neoliberal ortodoxo, privatizando el sistema de pensiones y reduciendo el gasto fiscal ($G \downarrow$). Estas reformas fueron diseñadas por un grupo de economistas chilenos educados bajo las teorías de Milton Friedman. Este grupo es conocido como:
\begin{enumerate}[label=\textbf{\Alph*)}]
\begin{multicols}{2}
  \item Los Keynesianos de los Andes
  \item El Grupo de Chicago (Chicago Boys)
  \item Los Tecnócratas de Harvard
  \item La Comisión Prebisch
  \item El Círculo de Viena
  \item Los Estructuralistas de la CEPAL
\end{multicols}
\end{enumerate}
\vspace{0.15cm}
\item Durante el genocidio de $1994$ en un país de la región de los Grandes Lagos en África, las masacres sistémicas fueron ejecutadas principalmente por milicias civiles extremistas de la etnia hutu armadas con machetes, apoyadas por propaganda radial. ¿Cuál era el nombre de esta milicia?
\begin{enumerate}[label=\textbf{\Alph*)}]
\begin{multicols}{2}
  \item Boko Haram
  \item Al-Shabaab
  \item Interahamwe
  \item Janjaweed
  \item Frente Polisario
  \item Lord's Resistance Army (LRA)
\end{multicols}
\end{enumerate}
\vspace{0.15cm}
\item En la década de $1950$, el Secretario de Estado de EE. UU., John Foster Dulles, formuló una doctrina de "Política al Borde del Abismo" (Brinkmanship). Según la Teoría de Juegos, esto implicaba amenazar con escalar un conflicto hasta la guerra nuclear total ($P(\text{Guerra}) \to 1$) para obligar al adversario soviético a retroceder. Esta estrategia dependía militarmente de:
\begin{enumerate}[label=\textbf{\Alph*)}]
\begin{multicols}{2}
  \item La infiltración guerrillera comunista.
  \item La "Represalia Masiva" (Massive Retaliation).
  \item El desarme unilateral.
  \item La Coexistencia Pacífica.
  \item El boicot comercial de Naciones Unidas.
  \item La política de Apaciguamiento.
\end{multicols}
\end{enumerate}
\vspace{0.15cm}
\item El proceso de descolonización de la India británica concluyó en $1947$ mediante un instrumento jurídico impulsado por el último virrey de la India, el cual decretó la partición del subcontinente en dos dominios basándose en criterios demográfico-religiosos. Este plan histórico se conoce como:
\begin{enumerate}[label=\textbf{\Alph*)}]
\begin{multicols}{2}
  \item Plan Marshall
  \item Acuerdos de Simla
  \item Plan Mountbatten
  \item Tratado de Delhi
  \item Declaración Balfour
  \item Línea Curzon
\end{multicols}
\end{enumerate}
\vspace{0.15cm}
\item En febrero de $1962$, en plena Carrera Espacial, Estados Unidos logró igualar el hito orbital soviético previo de Yuri Gagarin al enviar al primer astronauta estadounidense a orbitar el planeta a bordo de la cápsula *Friendship 7*. Su nombre era:
\begin{enumerate}[label=\textbf{\Alph*)}]
\begin{multicols}{2}
  \item Neil Armstrong
  \item Alan Shepard
  \item Buzz Aldrin
  \item John Glenn
  \item Michael Collins
  \item Chuck Yeager
\end{multicols}
\end{enumerate}
\vspace{0.15cm}
\item A finales de la década de $1960$, el Canciller de Alemania Occidental, Willy Brandt, inició una política de normalización de las relaciones diplomáticas y comerciales con la Alemania Oriental y el bloque soviético, aceptando el *statu quo* territorial. Esta política se denominó:
\begin{enumerate}[label=\textbf{\Alph*)}]
\begin{multicols}{2}
  \item Realpolitik
  \item Ostpolitik
  \item Glasnost
  \item Blitzkrieg
  \item Lebensraum
  \item Containment
\end{multicols}
\end{enumerate}
\vspace{0.15cm}
\item En el año $1972$, el "Club de Roma" encargó al MIT un estudio computacional sobre la viabilidad del crecimiento económico continuo frente a recursos planetarios finitos ($N(t) = N_0 e^{rt}$). Este influyente informe, que sentó las bases del ecologismo moderno, se tituló:
\begin{enumerate}[label=\textbf{\Alph*)}]
\begin{multicols}{2}
  \item Primavera Silenciosa (Silent Spring)
  \item El Fin de la Historia
  \item Los Límites del Crecimiento
  \item Nuestro Futuro Común
  \item La Riqueza de las Naciones
  \item El Manifiesto Comunista
\end{multicols}
\end{enumerate}
\vspace{0.15cm}
\item Durante el escándalo Irán-Contra en EE. UU., se demostró que el poder ejecutivo había evadido una ley explícita del Congreso que prohibía estrictamente destinar fondos federales ($\$ \Delta = 0$) para financiar a los rebeldes nicaragüenses (la "Contra"). ¿Cómo se llamaba esta enmienda legislativa?
\begin{enumerate}[label=\textbf{\Alph*)}]
\begin{multicols}{2}
  \item Enmienda Platt
  \item Enmienda Teller
  \item Enmienda Boland
  \item Doctrina Monroe
  \item Acta de Seguridad Nacional
  \item Resolución de Poderes de Guerra
\end{multicols}
\end{enumerate}
\vspace{0.15cm}
\item En la estructura social y política de la URSS, existía una élite burocrática del Partido Comunista que gozaba de privilegios exclusivos (acceso a tiendas especiales, hospitales y viviendas de lujo), contradiciendo el ideal marxista de una sociedad sin clases. A esta casta estatal se le denominaba:
\begin{enumerate}[label=\textbf{\Alph*)}]
\begin{multicols}{2}
  \item Proletariado
  \item Gulag
  \item Kulaks
  \item Nomenklatura
  \item Soviets Supremos
  \item Inteligencia (Intelligentsia)
\end{multicols}
\end{enumerate}
\vspace{0.15cm}
\item En mayo de $1974$, la dinámica geopolítica del sur de Asia cambió drásticamente cuando un país no alineado realizó su primera detonación nuclear subterránea exitosa, codificada irónicamente bajo el nombre de "Buda Sonriente" (Smiling Buddha). Este país fue:
\begin{enumerate}[label=\textbf{\Alph*)}]
\begin{multicols}{2}
  \item Pakistán
  \item Irán
  \item China
  \item India
  \item Israel
  \item Corea del Norte
\end{multicols}
\end{enumerate}
\vspace{0.15cm}
\item En julio de $1969$, buscando reducir el costo humano de la Guerra de Vietnam, el presidente de Estados Unidos anunció una nueva doctrina militar: EE. UU. proporcionaría apoyo económico y armamento a sus aliados asiáticos, pero exigiría que ellos aportaran las tropas terrestres para su propia defensa. Esta fue la:
\begin{enumerate}[label=\textbf{\Alph*)}]
\begin{multicols}{2}
  \item Doctrina Truman
  \item Doctrina Reagan
  \item Doctrina Nixon (Doctrina de Guam)
  \item Estrategia de Contención
  \item Respuesta Flexible
  \item Destrucción Mutua Asegurada
\end{multicols}
\end{enumerate}
\vspace{0.15cm}
\item En $1971$, un peculiar intercambio deportivo rompió el hielo diplomático de más de dos décadas entre Washington y Pekín, allanando el camino para la histórica visita presidencial del año siguiente. A este suceso se le bautizó periodísticamente como:
\begin{enumerate}[label=\textbf{\Alph*)}]
\begin{multicols}{2}
  \item Diplomacia del Béisbol
  \item Diplomacia del Ping-Pong
  \item El Milagro sobre el Hielo
  \item Diplomacia del Baloncesto
  \item La Guerra Fría del Ajedrez
  \item El Pacto Olímpico
\end{multicols}
\end{enumerate}
\vspace{0.15cm}
\item En abril de $1951$, en plena Guerra de Corea, el presidente Harry Truman relevó de sus funciones al General Douglas MacArthur debido a la insubordinación pública de este último. La fricción principal se debió a que MacArthur solicitaba escalar el conflicto ($\Delta E \to \infty$) mediante el uso de:
\begin{enumerate}[label=\textbf{\Alph*)}]
\begin{multicols}{2}
  \item Armas biológicas contra Pionyang.
  \item Bombardeos nucleares contra las tropas de la República Popular China.
  \item El bloqueo comercial total de Japón.
  \item La invasión terrestre directa de la URSS.
  \item Armas químicas (Gas Sarín) en Seúl.
  \item El abandono total de la península coreana.
\end{multicols}
\end{enumerate}
\vspace{0.15cm}
\item Durante el apogeo del terror anticomunista de la década de $1950$ en EE. UU. (Macartismo), un matrimonio estadounidense fue enjuiciado, condenado y ejecutado en la silla eléctrica en $1953$ bajo cargos de entregar los planos técnicos de la bomba atómica a la Unión Soviética. ¿Cuál era su apellido?
\begin{enumerate}[label=\textbf{\Alph*)}]
\begin{multicols}{2}
  \item Oppenheimer
  \item McCarthy
  \item Rosenberg
  \item Hiss
  \item Fuchs
  \item Truman
\end{multicols}
\end{enumerate}
\vspace{0.15cm}
\item En el contexto de la descolonización, el brutal conflicto librado en el norte de África entre independentistas y el ejército francés culminó en marzo de $1962$ con la firma de unos acuerdos que reconocieron el nacimiento de la República Argelina Democrática y Popular. Estos fueron los:
\begin{enumerate}[label=\textbf{\Alph*)}]
\begin{multicols}{2}
  \item Acuerdos de París
  \item Acuerdos de Ginebra
  \item Acuerdos de Évian
  \item Tratado de Versalles
  \item Declaración de Argel
  \item Acuerdos de Camp David
\end{multicols}
\end{enumerate}
\vspace{0.15cm}
\item Como pilar fundamental de la segregación racial del sistema de Apartheid en Sudáfrica, el gobierno blanco creó unos "territorios autónomos" empobrecidos, despojando de su ciudadanía sudafricana a la población negra y forzándola a vivir allí. Estas falsas naciones se denominaban:
\begin{enumerate}[label=\textbf{\Alph*)}]
\begin{multicols}{2}
  \item Gulags
  \item Guetos
  \item Kibutz
  \item Bantustanes
  \item Zonas Económicas Especiales
  \item Colonias Penales
\end{multicols}
\end{enumerate}
\vspace{0.15cm}
\item El sistema de alerta militar de las Fuerzas Armadas de EE. UU. mide la preparación para el combate nuclear en una escala inversa del $5$ al $1$. Durante la Crisis de los Misiles Cubanos ($1962$), el Comando Aéreo Estratégico (SAC) alcanzó su nivel histórico de máxima tensión (a un paso de la guerra termonuclear inminente), situándose en:
\begin{enumerate}[label=\textbf{\Alph*)}]
\begin{multicols}{2}
  \item DEFCON 5
  \item DEFCON 4
  \item DEFCON 3
  \item DEFCON 2
  \item DEFCON 1
  \item DEFCON 0
\end{multicols}
\end{enumerate}
\vspace{0.15cm}
\item Para garantizar la lealtad ideológica de los ejércitos latinoamericanos durante la Guerra Fría, Estados Unidos operó un centro de adoctrinamiento militar en Panamá (posteriormente trasladado a Georgia). Allí se capacitó a futuros dictadores y oficiales en tácticas de contrainsurgencia y tortura. Esta institución se denominaba:
\begin{enumerate}[label=\textbf{\Alph*)}]
\begin{multicols}{2}
  \item Academia Militar de West Point
  \item El Pentágono Sur
  \item Escuela de las Américas
  \item Centro de Inteligencia del Comando Sur
  \item Instituto del Hemisferio Occidental
  \item Base Naval de Guantánamo
\end{multicols}
\end{enumerate}
\vspace{0.15cm}
\item En el año $1947$, el Comité de Actividades Antiestadounidenses (HUAC) del Congreso citó a declarar a numerosos profesionales del cine. Un grupo de guionistas y directores se amparó en la Primera Enmienda negándose a testificar sobre sus afiliaciones políticas, por lo que fueron encarcelados y colocados en una "lista negra". Se les conoce como:
\begin{enumerate}[label=\textbf{\Alph*)}]
\begin{multicols}{2}
  \item Los Mártires de Hollywood
  \item Los Diez de Hollywood (The Hollywood Ten)
  \item El Grupo de California
  \item La Resistencia de Burbank
  \item El Círculo de Guionistas Libres
  \item Los Subversivos de Los Ángeles
\end{multicols}
\end{enumerate}
\vspace{0.15cm}
\item Antes de la Revolución Islámica de $1979$, el Sha de Irán gobernaba con el respaldo de Occidente, manteniendo un férreo control interno y aplastando cualquier disidencia izquierdista o religiosa. Su temida agencia de policía secreta e inteligencia nacional, entrenada por la CIA y el Mossad, se llamaba:
\begin{enumerate}[label=\textbf{\Alph*)}]
\begin{multicols}{2}
  \item KGB
  \item Gestapo
  \item SAVAK
  \item Stasi
  \item DINA
  \item Securitate
\end{multicols}
\end{enumerate}
\vspace{0.15cm}
\item En $1986$, la Comunidad Económica Europea firmó su primera gran revisión del Tratado de Roma. Este nuevo acuerdo estableció el objetivo cronológico de crear un mercado único genuino para finales de $1992$, eliminando barreras no arancelarias para bienes, servicios, personas y capitales. Este tratado es el:
\begin{enumerate}[label=\textbf{\Alph*)}]
\begin{multicols}{2}
  \item Tratado de Maastricht
  \item Acuerdo de Schengen
  \item Acta Única Europea
  \item Tratado de Lisboa
  \item Tratado de Ámsterdam
  \item Protocolo de Londres
\end{multicols}
\end{enumerate}
\vspace{0.15cm}
\item El Frente Farabundo Martí para la Liberación Nacional (FMLN), creado en $1980$ al unificar a $n=5$ organizaciones político-militares salvadoreñas, tomó su nombre de Farabundo Martí, quien fue un histórico líder comunista fusilado tras el fallido levantamiento indígena y campesino del año:
\begin{enumerate}[label=\textbf{\Alph*)}]
\begin{multicols}{2}
  \item 1920
  \item 1932
  \item 1944
  \item 1969
  \item 1979
  \item 1981
\end{multicols}
\end{enumerate}
\vspace{0.15cm}
\item A principios de la década de $1980$, la comunidad médica mundial detectó un síndrome inmunológico hasta entonces desconocido, el cual alteraría la demografía global y las políticas de salud pública del fin del $\text{Siglo XX}$. Causado por un lentivirus (VIH), esta pandemia fue oficialmente reconocida como:
\begin{enumerate}[label=\textbf{\Alph*)}]
\begin{multicols}{2}
  \item Ébola
  \item Gripe A (H1N1)
  \item Síndrome de Inmunodeficiencia Adquirida (SIDA)
  \item Tuberculosis Multirresistente
  \item SARS
  \item Cólera Asiático
\end{multicols}
\end{enumerate}
\vspace{0.15cm}
\item En mayo de $1997$, concluyó la Primera Guerra del Congo en África Central. Las fuerzas rebeldes de Laurent-Désiré Kabila, apoyadas por Ruanda y Uganda, lograron derrocar al veterano dictador cleptócrata apoyado por Occidente durante la Guerra Fría y rebautizaron Zaire como República Democrática del Congo. ¿Quién fue el dictador depuesto?
\begin{enumerate}[label=\textbf{\Alph*)}]
\begin{multicols}{2}
  \item Idi Amin
  \item Muamar el Gadafi
  \item Jean-Bédel Bokassa
  \item Mobutu Sese Seko
  \item Patrice Lumumba
  \item Robert Mugabe
\end{multicols}
\end{enumerate}
\vspace{0.15cm}
\item El colapso del bloque soviético culminó físicamente con un evento jurídico y político que modificó la cartografía de Europa Occidental. La absorción de la antigua República Democrática Alemana (RDA) por la República Federal de Alemania (RFA) se oficializó el día:
\begin{enumerate}[label=\textbf{\Alph*)}]
\begin{multicols}{2}
  \item 9 de noviembre de 1989
  \item 3 de octubre de 1990
  \item 25 de diciembre de 1991
  \item 1 de enero de 1993
  \item 15 de agosto de 1945
  \item 11 de septiembre de 2001
\end{multicols}
\end{enumerate}
\vspace{0.15cm}
\item Durante la Guerra de las Malvinas de $1982$, las fuerzas aéreas de la dictadura argentina lograron un impacto psicológico y táctico profundo al hundir el destructor británico de última generación *HMS Sheffield*, demostrando la vulnerabilidad de las flotas modernas, utilizando misiles antibuque de origen francés llamados:
\begin{enumerate}[label=\textbf{\Alph*)}]
\begin{multicols}{2}
  \item Tomahawk
  \item Scud
  \item Exocet
  \item Harpoon
  \item Patriot
  \item Sidewinder
\end{multicols}
\end{enumerate}
\vspace{0.15cm}
\item La arquitectura teórica de la Guerra Fría se basó fuertemente en la Teoría de Juegos. La estrategia de "Destrucción Mutua Asegurada" (MAD) es un escenario donde ningún actor tiene incentivos para desatar un ataque atómico por miedo a la aniquilación total de su propio país. En términos matemáticos, esta estabilidad del terror representa un:
\begin{enumerate}[label=\textbf{\Alph*)}]
\begin{multicols}{2}
  \item Teorema de Pitágoras
  \item Equilibrio de Nash
  \item Paradoja de Fermi
  \item Principio de Incertidumbre de Heisenberg
  \item Efecto Mariposa
  \item Movimiento Browniano
\end{multicols}
\end{enumerate}
\vspace{0.15cm}
\item En enero de $1967$, la Carrera Espacial de Estados Unidos sufrió una tragedia letal ($n_{\text{víctimas}} = 3$) cuando un incendio devastador consumió el módulo de comando presurizado con oxígeno puro ($100\%\ O_2$) durante un ensayo en tierra firme. ¿A qué misión pertenecía esta tripulación?
\begin{enumerate}[label=\textbf{\Alph*)}]
\begin{multicols}{2}
  \item Apolo 11
  \item Apolo 1
  \item Apolo 13
  \item Challenger
  \item Columbia
  \item Gemini 8
\end{multicols}
\end{enumerate}
\vspace{0.15cm}
\item La Guerra de Yom Kipur de $1973$ inició cuando el ejército egipcio, burlando la inteligencia de occidente, cruzó exitosamente el Canal de Suez empleando bombas de agua a alta presión para perforar las defensas de arena de la "Línea Bar Lev" israelí. Esta audaz maniobra anfibia se codificó como:
\begin{enumerate}[label=\textbf{\Alph*)}]
\begin{multicols}{2}
  \item Operación Barbarroja
  \item Operación Badr
  \item Operación Tormenta del Desierto
  \item Operación Entebbe
  \item Operación León Marino
  \item Operación Ciclón
\end{multicols}
\end{enumerate}
\vspace{0.15cm}
\item Uno de los movimientos anticoloniales pioneros en el África subsahariana fue liderado por Jomo Kenyatta. Tras años de cárcel durante la rebelión del Mau Mau, Kenyatta fue elegido en $1963$ como el primer presidente constitucional e independiente de:
\begin{enumerate}[label=\textbf{\Alph*)}]
\begin{multicols}{2}
  \item Ghana
  \item Zimbabue
  \item Sudáfrica
  \item Kenia
  \item Nigeria
  \item Angola
\end{multicols}
\end{enumerate}
\vspace{0.15cm}
\item Mientras que el Tribunal de Núremberg juzgó a los nazis en Europa, el Tribunal Penal Militar Internacional para el Lejano Oriente fue establecido en $1946$ para procesar a los líderes del Imperio Japonés por crímenes de guerra. Su figura central condenada a la horca, quien fuera Primer Ministro y general del ejército imperial, fue:
\begin{enumerate}[label=\textbf{\Alph*)}]
\begin{multicols}{2}
  \item Emperador Hirohito
  \item Almirante Isoroku Yamamoto
  \item General Hideki Tojo
  \item Príncipe Asaka
  \item General Tomoyuki Yamashita
  \item Almirante Chuichi Nagumo
\end{multicols}
\end{enumerate}
\vspace{0.15cm}
\item A finales de la década de $1960$, previendo la crisis de liquidez del dólar, el Fondo Monetario Internacional (FMI) creó un activo de reserva internacional sintético cuyo valor se basa en una cesta de las principales monedas globales, destinado a complementar las reservas oficiales de los países miembros. Se denominan:
\begin{enumerate}[label=\textbf{\Alph*)}]
\begin{multicols}{2}
  \item Dólares de Plata
  \item Bitcoins
  \item Derechos Especiales de Giro (DEG)
  \item Bancor
  \item Bonos Brady
  \item Criptoactivos Soberanos
\end{multicols}
\end{enumerate}
\vspace{0.15cm}
\item En el otoño de $1962$, coincidiendo temporalmente con la Crisis de los Misiles en Cuba, dos potencias demográficas asiáticas entraron en una breve pero cruenta guerra fronteriza por el control de la región himalaya de Aksai Chin y Arunachal Pradesh. El conflicto fue entre:
\begin{enumerate}[label=\textbf{\Alph*)}]
\begin{multicols}{2}
  \item India y Pakistán
  \item URSS y China
  \item China y Vietnam
  \item India y China
  \item Corea del Norte y Corea del Sur
  \item Japón y Corea del Sur
\end{multicols}
\end{enumerate}
\vspace{0.15cm}
\item Durante la Revolución Cultural Proletaria ($1966$-$1976$), los jóvenes de los "Guardias Rojos" utilizaban como guía dogmática e instrumento de adoctrinamiento una recopilación de discursos y escritos de su líder supremo, cuya lectura y posesión era un requisito de supervivencia política. El libro se popularizó en Occidente como:
\begin{enumerate}[label=\textbf{\Alph*)}]
\begin{multicols}{2}
  \item El Capital
  \item El Libro Rojo de Mao
  \item Mi Lucha
  \item El Manifiesto Comunista
  \item El Libro Verde de Gadafi
  \item El Archipiélago Gulag
\end{multicols}
\end{enumerate}
\vspace{0.15cm}
\item Durante la ocupación soviética de Afganistán ($1979$-$1989$), la balanza del conflicto asimétrico cambió a favor de los guerrilleros afganos (muyahidines) a mediados de la década, cuando la CIA comenzó a suministrarles una letal tecnología de misiles portátiles (MANPADS) que neutralizó la supremacía de los helicópteros soviéticos de asalto *Mi-24*. Estos misiles eran los:
\begin{enumerate}[label=\textbf{\Alph*)}]
\begin{multicols}{2}
  \item Misiles Scud
  \item Misiles Patriot
  \item Misiles FIM-92 Stinger
  \item Lanzacohetes RPG-7
  \item Misiles Javelin
  \item Misiles Tomahawk
\end{multicols}
\end{enumerate}
\vspace{0.15cm}
\item En la década de $1980$, la nación andina de Perú enfrentó a uno de los movimientos insurgentes marxista-leninista-maoístas más violentos del continente. Liderado por el profesor de filosofía Abimael Guzmán, este grupo utilizó tácticas de terrorismo extremo, causando decenas de miles de muertes. Su nombre es:
\begin{enumerate}[label=\textbf{\Alph*)}]
\begin{multicols}{2}
  \item Tupamaros
  \item FARC
  \item Sendero Luminoso (PCP-SL)
  \item ELN
  \item Montoneros
  \item Movimiento de Izquierda Revolucionaria (MIR)
\end{multicols}
\end{enumerate}
\vspace{0.15cm}
\item En enero de $1991$, la "Operación Tormenta del Desierto" expulsó a las tropas iraquíes de Kuwait, marcando un hito tecnológico de superioridad aérea y guerra acorazada occidental bajo el mando de un mediático general de EE. UU. (apodado *Stormin' Norman*). ¿Quién era este comandante?
\begin{enumerate}[label=\textbf{\Alph*)}]
\begin{multicols}{2}
  \item General George S. Patton
  \item General Colin Powell
  \item General H. Norman Schwarzkopf
  \item General Douglas MacArthur
  \item General William Westmoreland
  \item General David Petraeus
\end{multicols}
\end{enumerate}
\vspace{0.15cm}
\item Entre abril de $1992$ y febrero de $1996$, la capital de Bosnia-Herzegovina fue sometida por fuerzas serbias al asedio militar más prolongado de una capital europea en la historia de la guerra moderna ($\approx 1,425$ días), caracterizado por el terrorismo de francotiradores y fuego de mortero contra civiles. ¿A qué ciudad nos referimos?
\begin{enumerate}[label=\textbf{\Alph*)}]
\begin{multicols}{2}
  \item Belgrado
  \item Srebrenica
  \item Kosovo
  \item Sarajevo
  \item Pristina
  \item Zagreb
\end{multicols}
\end{enumerate}
\vspace{0.15cm}
\item En el Reino Unido, durante la Guerra Fría tardía, la primera mujer en ocupar el cargo de Primera Ministra impuso políticas conservadoras inflexibles, derrotó a los sindicatos mineros y recuperó las Islas Malvinas militarmente. Un periodista soviético la bautizó con el célebre apodo de:
\begin{enumerate}[label=\textbf{\Alph*)}]
\begin{multicols}{2}
  \item La Dama de Hierro
  \item La Reina Victoria Moderna
  \item La Canciller de Acero
  \item La Mujer Biónica
  \item El Halcón Británico
  \item La Valquiria Conservadora
\end{multicols}
\end{enumerate}
\vspace{0.15cm}
\item Como arquitecto clave de la política exterior soviética a finales de los $1980$s, el ministro de Asuntos Exteriores de Mijaíl Gorbachov fue fundamental para distender la Guerra Fría y promover la reunificación alemana de forma pacífica, siendo posteriormente el presidente independiente de Georgia. ¿De quién se trata?
\begin{enumerate}[label=\textbf{\Alph*)}]
\begin{multicols}{2}
  \item Yuri Andrópov
  \item Borís Yeltsin
  \item Vyacheslav Molotov
  \item Eduard Shevardnadze
  \item Andrei Gromyko
  \item Aleksandr Rutskói
\end{multicols}
\end{enumerate}
\vspace{0.15cm}
\item En el sur de África, entre $1987$ y $1988$, se libró la batalla convencional más grande del continente desde la SGM. Las fuerzas conjuntas de Angola (FAPLA) y las tropas expedicionarias de Cuba frenaron el avance del ejército de la Sudáfrica del Apartheid. ¿Cuál fue el nombre de esta histórica batalla?
\begin{enumerate}[label=\textbf{\Alph*)}]
\begin{multicols}{2}
  \item Batalla de El Alamein
  \item Batalla de Cuito Cuanavale
  \item Batalla de Rorke's Drift
  \item Sitio de Tobruk
  \item Batalla de Mogadiscio
  \item Batalla de Isandlwana
\end{multicols}
\end{enumerate}
\vspace{0.15cm}
\item El Tratado de Roma de $1957$ no solo instituyó la Comunidad Económica Europea, sino que en paralelo fundó otra entidad supranacional destinada a coordinar y monopolizar los programas pacíficos de investigación y energía nuclear civil entre los países miembros fundadores. Esta comunidad es conocida como:
\begin{enumerate}[label=\textbf{\Alph*)}]
\begin{multicols}{2}
  \item OIEA (Agencia Internacional de Energía Atómica)
  \item Euratom
  \item CERN (Organización Europea para la Investigación Nuclear)
  \item AELC (Asociación Europea de Libre Comercio)
  \item Pacto de Varsovia Nuclear
  \item Tratado de No Proliferación Europea
\end{multicols}
\end{enumerate}
\vspace{0.15cm}
\item En octubre de $1961$, el nivel de tensión nuclear alcanzó un pico durante una crisis fronteriza táctica, cuando tanques M48 Patton estadounidenses y tanques T-54 soviéticos se enfrentaron a menos de $100$ metros de distancia, apuntándose mutuamente durante $16$ horas en el paso fronterizo interurbano conocido como:
\begin{enumerate}[label=\textbf{\Alph*)}]
\begin{multicols}{2}
  \item Checkpoint Charlie
  \item Puente de Glienicke
  \item Muro de los Lamentos
  \item Zona Desmilitarizada de Corea
  \item El Paso de Calais
  \item Línea Maginot
\end{multicols}
\end{enumerate}
\vspace{0.15cm}
\item En el año $1977$, la NASA lanzó dos sondas espaciales diseñadas para explorar los planetas exteriores del Sistema Solar. Llevaban a bordo un "Disco de Oro" (*Golden Record*) con sonidos e imágenes de la Tierra, concebido como un mensaje pacífico para posibles civilizaciones extraterrestres. Estas misiones fueron:
\begin{enumerate}[label=\textbf{\Alph*)}]
\begin{multicols}{2}
  \item Viking 1 y 2
  \item Pioneer 10 y 11
  \item Voyager 1 y 2
  \item Apollo 15 y 16
  \item Mariner 4 y 9
  \item Sputnik y Vostok
\end{multicols}
\end{enumerate}
\vspace{0.15cm}
\item A principios de diciembre de $1989$, apenas unas semanas después de la caída del Muro de Berlín, George H. W. Bush y Mijaíl Gorbachov se reunieron en una cumbre naval a bordo de buques de guerra en el mar Mediterráneo. En este histórico encuentro, declararon oficial y públicamente el fin de la Guerra Fría. Esta fue la:
\begin{enumerate}[label=\textbf{\Alph*)}]
\begin{multicols}{2}
  \item Cumbre de Reikiavik
  \item Cumbre de Yalta
  \item Cumbre de Ginebra
  \item Cumbre de Malta
  \item Conferencia de Potsdam
  \item Conferencia de Teherán
\end{multicols}
\end{enumerate}
\vspace{0.15cm}
\end{enumerate}\n
\subsubsection*{Módulo 6 Final: Reactivos 251 - 300}

\begin{enumerate}[label=\textbf{\arabic*.}]
\setcounter{enumi}{250}
\item En el año $1959$, en medio de la escalada de la Guerra Fría y la Carrera Armamentista, $12$ naciones firmaron un tratado histórico que garantizaba que un continente entero sería utilizado exclusivamente para fines pacíficos y científicos, prohibiendo expresamente la militarización ($n_{\text{armas}} = 0$) al sur del paralelo $\lambda = 60^{\circ}$ Sur. ¿Cuál es este tratado?
\begin{enumerate}[label=\textbf{\Alph*)}]
\begin{multicols}{2}
  \item Tratado de Groenlandia
  \item Tratado Antártico
  \item Convención de Svalbard
  \item Tratado del Atlántico Norte
  \item Pacto Ártico
  \item Tratado de Tlatelolco
\end{multicols}
\end{enumerate}
\vspace{0.15cm}
\item En la geopolítica de la Guerra Fría, las superpotencias (EE. UU. y la URSS) evitaron el enfrentamiento militar directo (donde el conflicto bélico frontal $A \cap B = \emptyset$) debido al riesgo de aniquilación nuclear. En su lugar, financiaron y armaron a facciones opuestas en países del Tercer Mundo. Este fenómeno se define conceptualmente como:
\begin{enumerate}[label=\textbf{\Alph*)}]
\begin{multicols}{2}
  \item Guerra Híbrida
  \item Guerra Relámpago (Blitzkrieg)
  \item Guerra Subsidiaria o *Proxy War*
  \item Guerra de Trincheras
  \item Guerra Comercial Asimétrica
  \item Guerra Irregular de Baja Intensidad
\end{multicols}
\end{enumerate}
\vspace{0.15cm}
\item En julio de $1969$, El Salvador y Honduras protagonizaron un conflicto armado de aproximadamente $100\text{ horas}$. Aunque la prensa internacional la frivolizó llamándola "Guerra del Fútbol", sus causas estructurales radicaban en la presión demográfica ($\Delta\text{Población} > 0$), la migración campesina salvadoreña y las políticas de reforma agraria hondureñas. Este conflicto se conoce históricamente como:
\begin{enumerate}[label=\textbf{\Alph*)}]
\begin{multicols}{2}
  \item La Guerra de los Cien Días
  \item La Guerra de las $100$ Horas o Guerra de Legítima Defensa
  \item La Batalla de Centroamérica
  \item La Guerra del Golfo de Fonseca
  \item La Crisis de la Integración
  \item La Guerra de Morazán
\end{multicols}
\end{enumerate}
\vspace{0.15cm}
\item Durante la dictadura cívico-militar en Argentina ($1976$-$1983$), un grupo de mujeres comenzó a marchar semanalmente alrededor de la pirámide central frente a la Casa Rosada, exigiendo al Estado la aparición con vida de sus hijos secuestrados ($\Delta\text{Vidas} < 0$). Este emblemático movimiento de derechos humanos es:
\begin{enumerate}[label=\textbf{\Alph*)}]
\begin{multicols}{2}
  \item Mujeres de Negro
  \item Las Sufragistas de Buenos Aires
  \item Madres de Plaza de Mayo
  \item Abuelas de la Revolución
  \item Las Piqueteras
  \item Movimiento Evita
\end{multicols}
\end{enumerate}
\vspace{0.15cm}
\item Tras décadas de violencia sectaria y terrorismo paramilitar (*The Troubles*) que dejaron más de $3,500$ muertos en el Reino Unido e Irlanda, las facciones católicas (republicanas) y protestantes (unionistas) firmaron en $1998$ un acuerdo de paz mediado por Estados Unidos. ¿Cuál es el nombre de este tratado?
\begin{enumerate}[label=\textbf{\Alph*)}]
\begin{multicols}{2}
  \item Acuerdos de Oslo
  \item Acuerdo de Viernes Santo (Good Friday Agreement)
  \item Tratado de Londres
  \item Pacto de Dublín
  \item Acuerdos de Dayton
  \item Tratado de Edimburgo
\end{multicols}
\end{enumerate}
\vspace{0.15cm}
\item En diciembre de $1984$, la globalización industrial evidenció sus riesgos letales cuando una planta de pesticidas de la corporación estadounidense *Union Carbide* en India sufrió una fuga masiva de isocianato de metilo ($\text{C}_2\text{H}_3\text{NO}$), provocando la muerte de más de $15,000$ personas. Este desastre industrial se conoce como:
\begin{enumerate}[label=\textbf{\Alph*)}]
\begin{multicols}{2}
  \item Desastre de Chernóbil
  \item Tragedia de Minamata
  \item Desastre de Bhopal
  \item Fuga de Seveso
  \item Derrame del Exxon Valdez
  \item Incidente de Fukushima
\end{multicols}
\end{enumerate}
\vspace{0.15cm}
\item Entre $1986$ y $1994$, el comercio mundial se reestructuró bajo la "Ronda de Uruguay". Este complejo proceso de negociaciones arancelarias culminó con el Tratado de Marrakech ($1994$), el cual transformó el frágil Acuerdo General sobre Aranceles Aduaneros y Comercio (GATT) en una organización internacional vinculante llamada:
\begin{enumerate}[label=\textbf{\Alph*)}]
\begin{multicols}{2}
  \item Banco Mundial (BM)
  \item Cámara de Comercio Internacional (CCI)
  \item Organización para la Cooperación y el Desarrollo Económicos (OCDE)
  \item Organización Mundial del Comercio (OMC)
  \item Fondo Monetario Internacional (FMI)
  \item Foro Económico Mundial (WEF)
\end{multicols}
\end{enumerate}
\vspace{0.15cm}
\item El final de la Segunda Guerra Mundial marcó el inicio de un fenómeno demográfico en Estados Unidos y Europa Occidental. Entre $1946$ y $1964$, la prosperidad económica y la paz generaron un incremento exponencial en las tasas de fecundidad ($TFR \uparrow$). A esta cohorte generacional se le denomina:
\begin{enumerate}[label=\textbf{\Alph*)}]
\begin{multicols}{2}
  \item Generación Silenciosa
  \item Generación X
  \item Millennials
  \item Baby Boomers
  \item Generación Perdida
  \item Centennials
\end{multicols}
\end{enumerate}
\vspace{0.15cm}
\item En noviembre de $1957$, apenas un mes después del *Sputnik 1*, el programa espacial soviético lanzó la misión *Sputnik 2*, la cual marcó un hito biológico al poner en órbita ($v \approx 8\text{ km/s}$) al primer ser vivo de la historia, demostrando que un organismo podía sobrevivir al lanzamiento. ¿Quién fue este tripulante?
\begin{enumerate}[label=\textbf{\Alph*)}]
\begin{multicols}{2}
  \item El chimpancé Ham
  \item La perra Laika
  \item El mono Albert II
  \item La perra Belka
  \item El ratón Hector
  \item La tortuga Horsfield
\end{multicols}
\end{enumerate}
\vspace{0.15cm}
\item A mediados del siglo XX, un ingeniero agrónomo estadounidense introdujo variedades de alto rendimiento (VAR) de trigo y arroz, combinadas con fertilizantes y riego masivo, salvando a cientos de millones de personas de la inanición en México, India y Pakistán ($Q_{\text{agrícola}} \gg \text{demanda}$). Este proceso, liderado por Norman Borlaug (Nobel de Paz $1970$), se denomina:
\begin{enumerate}[label=\textbf{\Alph*)}]
\begin{multicols}{2}
  \item Revolución Agraria
  \item Gran Salto Adelante
  \item Revolución Verde
  \item Agricultura Orgánica
  \item Colectivización Forzosa
  \item Plan Marshall Agrícola
\end{multicols}
\end{enumerate}
\vspace{0.15cm}
\item Nelson Mandela fue el prisionero político más célebre de la Guerra Fría. De sus $27$ años de reclusión bajo el régimen del Apartheid, pasó $18$ años realizando trabajos forzados en una isla prisión de máxima seguridad frente a la costa de Ciudad del Cabo. ¿Cuál es el nombre de esta isla?
\begin{enumerate}[label=\textbf{\Alph*)}]
\begin{multicols}{2}
  \item Isla Santa Elena
  \item Isla Robben (Robben Island)
  \item Isla de Alcatraz
  \item Isla del Diablo
  \item Isla Gorée
  \item Madagascar
\end{multicols}
\end{enumerate}
\vspace{0.15cm}
\item En $1954$, la Corte Suprema de los Estados Unidos emitió un fallo unánime histórico que declaró inconstitucional el sistema de segregación racial en las escuelas públicas, dictaminando que la doctrina de "Separados pero Iguales" es inherentemente desigual ($S \neq I$). Este caso judicial paradigmático se titula:
\begin{enumerate}[label=\textbf{\Alph*)}]
\begin{multicols}{2}
  \item Roe v. Wade
  \item Plessy v. Ferguson
  \item Miranda v. Arizona
  \item Brown v. Board of Education (Consejo de Educación)
  \item Marbury v. Madison
  \item Dred Scott v. Sandford
\end{multicols}
\end{enumerate}
\vspace{0.15cm}
\item La intervención militar internacional bajo la bandera de la ONU en la Guerra de Corea ($1950$) fue jurídicamente posible en el Consejo de Seguridad debido a un factor geopolítico clave, donde los votos en contra ($V_{\text{contra}} = 0$) no incluyeron el veto soviético. Esto ocurrió porque:
\begin{enumerate}[label=\textbf{\Alph*)}]
\begin{multicols}{2}
  \item La URSS había sido expulsada temporalmente de la ONU.
  \item La URSS estaba boicoteando las sesiones de la ONU en protesta porque el asiento de China lo ocupaba Taiwán.
  \item Stalin apoyaba secretamente a Corea del Sur.
  \item Estados Unidos sobornó a la delegación soviética.
  \item La URSS carecía de poder de veto en $1950$.
  \item La resolución fue aprobada en secreto durante la madrugada.
\end{multicols}
\end{enumerate}
\vspace{0.15cm}
\item El régimen totalitario camboyano de los Jemeres Rojos impuso en $1975$ un reinicio del calendario, decretando la abolición de la historia previa, la moneda, la familia y la tecnología urbana. El dictador Pol Pot denominó a este inicio forzado hacia una utopía agraria primitiva como:
\begin{enumerate}[label=\textbf{\Alph*)}]
\begin{multicols}{2}
  \item Año Cero ($t = 0$)
  \item El Gran Salto
  \item El Amanecer Rojo
  \item El Nuevo Orden
  \item La Era del Campesino
  \item El Primer Plan Quinquenal
\end{multicols}
\end{enumerate}
\vspace{0.15cm}
\item Durante las fases finales de la Guerra Irán-Irak, el $16$ de marzo de $1988$, la fuerza aérea iraquí bombardeó una ciudad kurda utilizando un cóctel letal de gas mostaza y agentes nerviosos (Sarín, Tabún), matando asfixiados a $\approx 5,000$ civiles en un solo día. Esta masacre se conoce como:
\begin{enumerate}[label=\textbf{\Alph*)}]
\begin{multicols}{2}
  \item Masacre de Srebrenica
  \item Ataque químico de Guta
  \item Masacre de Halabja
  \item Sitio de Kobane
  \item Tragedia de Al-Anfal
  \item Masacre de Sabra y Chatila
\end{multicols}
\end{enumerate}
\vspace{0.15cm}
\item En $1989$, el economista británico John Williamson acuñó un término para describir el decálogo de políticas de ajuste macroeconómico (disciplina fiscal, privatización, liberalización del comercio $\Delta X \uparrow$) exigidas a los países latinoamericanos asfixiados por la deuda. El término es:
\begin{enumerate}[label=\textbf{\Alph*)}]
\begin{multicols}{2}
  \item Plan Marshall
  \item Consenso de Washington
  \item Modelo Keynesiano
  \item Teoría de la Dependencia
  \item Consenso de Ginebra
  \item Nuevo Trato (New Deal)
\end{multicols}
\end{enumerate}
\vspace{0.15cm}
\item La Organización de Países Exportadores de Petróleo (OPEP) es un cártel intergubernamental que controla la oferta global ($Q_s$) para influir en los precios del crudo. Curiosamente, su sede central no se encuentra en Medio Oriente, sino en una capital europea neutral y diplomática. ¿Cuál es la sede de la OPEP desde $1965$?
\begin{enumerate}[label=\textbf{\Alph*)}]
\begin{multicols}{2}
  \item Ginebra
  \item Bruselas
  \item París
  \item Viena
  \item Londres
  \item Zúrich
\end{multicols}
\end{enumerate}
\vspace{0.15cm}
\item El Artículo $5$ del Tratado del Atlántico Norte consagra la doctrina de "defensa colectiva" (un ataque contra uno es un ataque contra todos). Durante todo el $\text{Siglo XX}$ de la Guerra Fría jamás se activó. La única vez en la historia que la OTAN ha invocado este artículo fue en respuesta a:
\begin{enumerate}[label=\textbf{\Alph*)}]
\begin{multicols}{2}
  \item La invasión soviética de Afganistán ($1979$).
  \item Los ataques terroristas del $11$ de septiembre de $2001$ en EE. UU.
  \item La Guerra de las Malvinas ($1982$).
  \item La Crisis de los Misiles en Cuba ($1962$).
  \item La Guerra de Corea ($1950$).
  \item El derribo del vuelo KAL 007 ($1983$).
\end{multicols}
\end{enumerate}
\vspace{0.15cm}
\item El Pacto de Varsovia, la alianza militar del bloque comunista, no se fundó en $1949$ simultáneamente con la OTAN, sino en $1955$. El detonante geopolítico exacto que impulsó a Moscú a crear este pacto fue:
\begin{enumerate}[label=\textbf{\Alph*)}]
\begin{multicols}{2}
  \item El éxito del Plan Marshall.
  \item La muerte de Iósif Stalin.
  \item El ingreso de la República Federal de Alemania (RFA) a la OTAN y su rearme.
  \item La victoria de Mao Zedong en China.
  \item El estallido de la Guerra de Vietnam.
  \item La detonación de la primera bomba de hidrógeno estadounidense.
\end{multicols}
\end{enumerate}
\vspace{0.15cm}
\item La Crisis de los Misiles Cubanos de octubre de $1962$ es considerada el momento en que el Reloj del Juicio Final estuvo más cerca de la medianoche nuclear. En la historiografía, ¿cuántos días duró exactamente este punto de tensión máxima desde el descubrimiento de los misiles hasta el acuerdo de retirada?
\begin{enumerate}[label=\textbf{\Alph*)}]
\begin{multicols}{2}
  \item $7\text{ días}$
  \item $13\text{ días}$
  \item $30\text{ días}$
  \item $60\text{ días}$
  \item $100\text{ días}$
  \item $3\text{ días}$
\end{multicols}
\end{enumerate}
\vspace{0.15cm}
\item A partir de la década de $1960$, un grupo de países asiáticos experimentó un crecimiento económico asombroso ($\Delta\text{PIB} > 7\%$ anual) basado en la Industrialización Orientada a la Exportación, ganándose el apodo de "Los Cuatro Tigres Asiáticos". ¿Qué país de la siguiente lista **NO** pertenece a este grupo original?
\begin{enumerate}[label=\textbf{\Alph*)}]
\begin{multicols}{2}
  \item Corea del Sur
  \item Taiwán
  \item Singapur
  \item Hong Kong
  \item Filipinas
  \item (Todos los anteriores pertenecen)
\end{multicols}
\end{enumerate}
\vspace{0.15cm}
\item En abril de $1980$, el presidente estadounidense Jimmy Carter autorizó una operación militar de fuerzas especiales para rescatar a los rehenes de la embajada en Teherán. La misión fue abortada debido a tormentas de arena y colisiones de helicópteros que dejaron $8$ militares muertos, humillando al Pentágono. La operación se llamó:
\begin{enumerate}[label=\textbf{\Alph*)}]
\begin{multicols}{2}
  \item Tormenta del Desierto
  \item Furia Urgente
  \item Garra de Águila (Eagle Claw)
  \item Lanza de Neptuno
  \item Cóndor Negro
  \item Justicia Infinita
\end{multicols}
\end{enumerate}
\vspace{0.15cm}
\item Durante el escándalo Irán-Contra, la CIA y la administración Reagan canalizaron fondos encubiertos a la guerrilla antisandinista en Nicaragua. Esta milicia financiada por Washington (que utilizaba tácticas de sabotaje económico y guerra de guerrillas) fue bautizada mediáticamente como:
\begin{enumerate}[label=\textbf{\Alph*)}]
\begin{multicols}{2}
  \item Frente Sandinista
  \item Los Contras
  \item FARC
  \item Sendero Luminoso
  \item Escuadrones de la Muerte
  \item Batallón Atlácatl
\end{multicols}
\end{enumerate}
\vspace{0.15cm}
\item La Guerra de las Malvinas inició el $2$ de abril de $1982$, cuando comandos anfibios argentinos desembarcaron en Puerto Argentino (Port Stanley) y rindieron a la pequeña guarnición de *Royal Marines* británicos sin causar bajas enemigas. El nombre clave de la armada argentina para esta invasión fue:
\begin{enumerate}[label=\textbf{\Alph*)}]
\begin{multicols}{2}
  \item Operación Soberanía
  \item Operación Cóndor
  \item Operación Rosario
  \item Operación Albatros
  \item Operación Independencia
  \item Operación Atlántico Sur
\end{multicols}
\end{enumerate}
\vspace{0.15cm}
\item El Genocidio de Ruanda ($1994$) evidenció el poder devastador de los medios de comunicación en la guerra psicológica. Las masacres fueron coordinadas y azuzadas ($f \approx 100\%$ de incitación al odio) mediante transmisiones que llamaban a los tutsis "cucarachas". Esta infame estación de radio se denominaba:
\begin{enumerate}[label=\textbf{\Alph*)}]
\begin{multicols}{2}
  \item Radio Free Europe
  \item Voice of America
  \item RTLM (Radio Télévision Libre des Mille Collines)
  \item Radio Moscú
  \item BBC África
  \item Radio Venceremos
\end{multicols}
\end{enumerate}
\vspace{0.15cm}
\item En las horas finales de la Guerra del Golfo ($1991$), el ejército iraquí en retirada desde Kuwait fue inmovilizado y bombardeado implacablemente por la aviación de coalición en la Autopista 80, dejando miles de vehículos calcinados. Las imágenes de este evento forzaron a EE. UU. a declarar el alto el fuego. A este tramo se le conoce como:
\begin{enumerate}[label=\textbf{\Alph*)}]
\begin{multicols}{2}
  \item La Ruta de la Seda
  \item La Vía Láctea
  \item La Autopista de la Muerte
  \item El Sendero de Lágrimas
  \item El Valle de las Sombras
  \item La Carretera del Infierno
\end{multicols}
\end{enumerate}
\vspace{0.15cm}
\item En $1968$, el líder del Partido Comunista de Checoslovaquia, Alexander Dubček, abolió la censura, permitió la libertad de prensa y promovió la democratización interna sin salirse del Pacto de Varsovia. Su eslogan filosófico para este experimento sociopolítico reprimido por los tanques soviéticos fue:
\begin{enumerate}[label=\textbf{\Alph*)}]
\begin{multicols}{2}
  \item "Todo el poder a los Soviets"
  \item "La Dictadura del Proletariado"
  \item "Socialismo con rostro humano"
  \item "Perestroika y Glasnost"
  \item "Comunismo en un solo país"
  \item "Revolución Permanente"
\end{multicols}
\end{enumerate}
\vspace{0.15cm}
\item La noche del $9$ de noviembre de $1989$, el Muro de Berlín cayó debido a un error de comunicación burocrática. ¿Qué funcionario de la RDA leyó precipitadamente en una rueda de prensa internacional que las restricciones de viaje se levantaban "de inmediato, sin demora", provocando la avalancha cívica a los puestos de control?
\begin{enumerate}[label=\textbf{\Alph*)}]
\begin{multicols}{2}
  \item Erich Honecker
  \item Mijaíl Gorbachov
  \item Günter Schabowski
  \item Egon Krenz
  \item Markus Wolf
  \item Helmut Kohl
\end{multicols}
\end{enumerate}
\vspace{0.15cm}
\item La descolonización pacífica de la India ($1947$) fue antecedida por la "Marcha de la Sal" en $1930$. En ella, Mahatma Gandhi lideró una caminata de $\approx 390\text{ km}$ hacia el mar para producir su propia sal, desafiando el monopolio y los impuestos británicos mediante la filosofía de resistencia civil no violenta denominada:
\begin{enumerate}[label=\textbf{\Alph*)}]
\begin{multicols}{2}
  \item Ahimsa
  \item Satyagraha (Fuerza de la Verdad)
  \item Yihad
  \item Apartheid
  \item Intifada
  \item Karma
\end{multicols}
\end{enumerate}
\vspace{0.15cm}
\item El $26$ de julio de $1953$, un grupo de rebeldes liderado por el joven abogado Fidel Castro atacó la segunda fortaleza militar más grande de Cuba con el objetivo de iniciar una insurrección contra la dictadura de Fulgencio Batista. Aunque fracasó tácticamente, este evento ($n_{\text{atacantes}} \approx 135$) dio nombre al movimiento revolucionario. Fue el asalto al:
\begin{enumerate}[label=\textbf{\Alph*)}]
\begin{multicols}{2}
  \item Cuartel Moncada
  \item Palacio Presidencial
  \item Fuerte de San Carlos
  \item Castillo del Morro
  \item Base de Guantánamo
  \item Cuartel de la Cabaña
\end{multicols}
\end{enumerate}
\vspace{0.15cm}
\item Durante la Primera Guerra de Indochina ($1946$-$1954$), la amalgama de fuerzas nacionalistas y comunistas que lucharon contra el imperio colonial francés por la independencia de Vietnam fue fundada en $1941$ por Ho Chi Minh. Esta liga guerrillera anticolonial se conocía como:
\begin{enumerate}[label=\textbf{\Alph*)}]
\begin{multicols}{2}
  \item Viet Cong
  \item Jemeres Rojos
  \item Viet Minh
  \item Pathet Lao
  \item Kuomintang
  \item Ejército Popular de Liberación
\end{multicols}
\end{enumerate}
\vspace{0.15cm}
\item El $21$ de diciembre de $1991$, los representantes de once ex repúblicas soviéticas firmaron una declaración que ratificaba el fin de la URSS y el nacimiento jurídico de la Comunidad de Estados Independientes (CEI), sellando la disolución del bloque del Este. Este documento es la:
\begin{enumerate}[label=\textbf{\Alph*)}]
\begin{multicols}{2}
  \item Declaración de Moscú
  \item Declaración de Alma-Ata (Almaty)
  \item Acuerdos de Minsk
  \item Tratado de Kiev
  \item Pacto de San Petersburgo
  \item Tratado de Yalta
\end{multicols}
\end{enumerate}
\vspace{0.15cm}
\item Como cumplimiento imperativo de los Acuerdos de Paz de Chapultepec ($1992$), la Organización de las Naciones Unidas (ONU) instauró una comisión en El Salvador encargada de investigar los graves actos de violencia y crímenes de lesa humanidad ocurridos durante la guerra civil. Su informe ($1993$) se tituló "De la locura a la esperanza". Esta entidad fue:
\begin{enumerate}[label=\textbf{\Alph*)}]
\begin{multicols}{2}
  \item La Comisión de la Verdad
  \item La Corte Penal Internacional
  \item La Fiscalía Especial contra la Impunidad
  \item El Tribunal Supremo Electoral
  \item La Comisión Internacional contra la Impunidad
  \item El Tribunal de Justicia Transicional
\end{multicols}
\end{enumerate}
\vspace{0.15cm}
\item Durante la década de $1980$, el conflicto afgano se convirtió en la guerra *proxy* más cara del bloque capitalista. La Agencia Central de Inteligencia (CIA) ejecutó una operación de miles de millones de dólares ($\$ \approx 2 \times 10^9$) para financiar, armar y entrenar a los insurgentes islámicos contra los soviéticos. El nombre en clave fue:
\begin{enumerate}[label=\textbf{\Alph*)}]
\begin{multicols}{2}
  \item Operación Cóndor
  \item Operación Gladio
  \item Operación Ciclón
  \item Operación Fénix
  \item Operación Barbarroja
  \item Operación Tormenta del Desierto
\end{multicols}
\end{enumerate}
\vspace{0.15cm}
\item El año $1968$ fue testigo de una ola de insurgencia global. En octubre de ese año, en el marco de la Masacre de Tlatelolco en México, el ejército y paramilitares abrieron fuego contra estudiantes pacíficos en la Plaza de las Tres Culturas. El batallón paramilitar encubierto que inició los disparos para generar el caos se denominaba:
\begin{enumerate}[label=\textbf{\Alph*)}]
\begin{multicols}{2}
  \item Batallón Atlácatl
  \item Escuadrón de la Muerte
  \item Batallón Olimpia
  \item Guardias Rojos
  \item Sección de Asalto
  \item Kaibiles
\end{multicols}
\end{enumerate}
\vspace{0.15cm}
\item La crisis petrolera de $1973$ rompió la Curva de Phillips tradicional de la macroeconomía ($Inflación \propto \frac{1}{Desempleo}$). Occidente sufrió un fenómeno donde la alta inflación coexistía paradójicamente con el estancamiento del PIB ($\Delta \text{PIB} \to 0$) y alto desempleo. Este escenario económico se denomina:
\begin{enumerate}[label=\textbf{\Alph*)}]
\begin{multicols}{2}
  \item Deflación
  \item Hiperinflación
  \item Estanflación (Stagflation)
  \item Depresión Económica
  \item Recesión de Balance
  \item Sobrecalentamiento Keynesiano
\end{multicols}
\end{enumerate}
\vspace{0.15cm}
\item En $1968$, el Tratado de No Proliferación Nuclear (TNP) estableció que las potencias atómicas no transferirían armas y los demás Estados no las adquirirían. Para auditar esto, el TNP designó a un organismo técnico de la ONU encargado de inspeccionar las instalaciones de uranio ($U-235$) del mundo. ¿Cuál es este organismo?
\begin{enumerate}[label=\textbf{\Alph*)}]
\begin{multicols}{2}
  \item Consejo de Seguridad (CSNU)
  \item Cascos Azules (DPKO)
  \item Corte Internacional de Justicia (CIJ)
  \item Organismo Internacional de Energía Atómica (OIEA o IAEA)
  \item Interpol
  \item Programa de las Naciones Unidas para el Medio Ambiente (PNUMA)
\end{multicols}
\end{enumerate}
\vspace{0.15cm}
\item En febrero de $1945$, con la derrota de la Alemania nazi inminente, Roosevelt, Churchill y Stalin se reunieron en una ciudad balneario soviética junto al Mar Negro. Allí acordaron en secreto dividir a Europa en "zonas de influencia" de posguerra y exigir la rendición incondicional del Eje. Fue la:
\begin{enumerate}[label=\textbf{\Alph*)}]
\begin{multicols}{2}
  \item Conferencia de Potsdam
  \item Conferencia de Teherán
  \item Conferencia de Yalta
  \item Conferencia de Múnich
  \item Conferencia de Bretton Woods
  \item Conferencia de San Francisco
\end{multicols}
\end{enumerate}
\vspace{0.15cm}
\item El $31$ de julio de $1991$, semanas antes del colapso de la URSS, Bush y Gorbachov firmaron un tratado bilateral que, por primera vez en la Guerra Fría, no solo "congelaba" la Carrera Armamentista, sino que obligaba a una reducción masiva ($\approx 30\%$) de los arsenales estratégicos ($N_{EEUU} \downarrow, N_{URSS} \downarrow$). Este tratado fue el:
\begin{enumerate}[label=\textbf{\Alph*)}]
\begin{multicols}{2}
  \item SALT I
  \item SALT II
  \item Tratado INF
  \item Tratado START I (Tratado de Reducción de Armas Estratégicas)
  \item Tratado ABM
  \item Tratado de Tlatelolco
\end{multicols}
\end{enumerate}
\vspace{0.15cm}
\item A partir de $1965$, la administración de Lyndon B. Johnson ordenó una campaña de bombardeo aéreo masivo, continuo y estratégico sobre Vietnam del Norte. Su objetivo (que fracasó) era destruir la voluntad de combate comunista y su capacidad industrial arrojando $\approx 864,000$ toneladas de bombas. Esta operación fue:
\begin{enumerate}[label=\textbf{\Alph*)}]
\begin{multicols}{2}
  \item Operación Ciclón
  \item Operación Rolling Thunder (Trueno Retumbante)
  \item Operación Tormenta del Desierto
  \item Operación Garra de Águila
  \item Operación Barbarroja
  \item Operación Ranch Hand
\end{multicols}
\end{enumerate}
\vspace{0.15cm}
\item La Guerra Civil China finalizó en el continente en $1949$ con la victoria de Mao Zedong. El gobierno nacionalista del Kuomintang, liderado por Chiang Kai-shek, evacuó sus tropas, reservas de oro y a unos $\approx 2 \times 10^6$ de civiles cruzando el estrecho hacia una isla, fundando la "República de China" en el exilio. Esta isla (antiguamente Formosa) es:
\begin{enumerate}[label=\textbf{\Alph*)}]
\begin{multicols}{2}
  \item Hainan
  \item Hong Kong
  \item Macao
  \item Taiwán
  \item Okinawa
  \item Jeju
\end{multicols}
\end{enumerate}
\vspace{0.15cm}
\item La intervención de Estados Unidos en la geopolítica del $\text{Siglo XX}$ tiene profundas raíces ideológicas en el $\text{Siglo XIX}$. La creencia cultural y religiosa de que la nación estadounidense estaba predestinada por Dios ($P(\text{expansión}) = 1$) para expandir su dominio y sus valores democráticos por el mundo se conoce como:
\begin{enumerate}[label=\textbf{\Alph*)}]
\begin{multicols}{2}
  \item Doctrina Truman
  \item Doctrina del Destino Manifiesto
  \item Política del Buen Vecino
  \item Plan Marshall
  \item Corolario Roosevelt
  \item Aislacionismo
\end{multicols}
\end{enumerate}
\vspace{0.15cm}
\item Al finalizar la descolonización en la década de los $1960$s, las potencias europeas retiraron sus ejércitos y administradores políticos. Sin embargo, los nuevos países del Tercer Mundo denunciaron que la subordinación continuaba mediante deudas asimétricas, control corporativo y fijación de precios de materias primas. A este fenómeno se le denomina:
\begin{enumerate}[label=\textbf{\Alph*)}]
\begin{multicols}{2}
  \item Imperialismo Clásico
  \item Mercantilismo
  \item Neocolonialismo
  \item Federalismo
  \item Panafricanismo
  \item Estado de Bienestar
\end{multicols}
\end{enumerate}
\vspace{0.15cm}
\item El $26$ de junio de $1963$, frente al Muro de Berlín, el presidente John F. Kennedy pronunció un discurso histórico subrayando el apoyo de Estados Unidos a la Alemania Democrática Occidental frente a la tiranía soviética. Su frase culminante, pronunciada en idioma alemán, fue:
\begin{enumerate}[label=\textbf{\Alph*)}]
\begin{multicols}{2}
  \item "¡Señor Gorbachov, derribe este muro!"
  \item "¡Viva la libertad!"
  \item "¡Ich bin ein Berliner!" (Yo soy un berlinés)
  \item "El mundo libre prevalecerá"
  \item "No pasaran"
  \item "Berlín es la capital del mundo"
\end{multicols}
\end{enumerate}
\vspace{0.15cm}
\item El shock psicológico que supuso el lanzamiento del satélite soviético *Sputnik 1* convenció al presidente Dwight D. Eisenhower y al Congreso de EE. UU. de que estaban perdiendo la supremacía tecnológica. En $1958$, como respuesta directa, el gobierno estadounidense creó la agencia civil:
\begin{enumerate}[label=\textbf{\Alph*)}]
\begin{multicols}{2}
  \item NSA (Agencia de Seguridad Nacional)
  \item CIA (Agencia Central de Inteligencia)
  \item NASA (Administración Nacional de Aeronáutica y el Espacio)
  \item DARPA
  \item Fuerza Espacial
  \item FBI
\end{multicols}
\end{enumerate}
\vspace{0.15cm}
\item El programa del transbordador espacial estadounidense sufrió su primer gran revés en enero de $1986$. Apenas $73\text{ segundos}$ después del despegue, una falla en los anillos de sellado (O-rings) debido a temperaturas gélidas ($T < 0^{\circ}\text{C}$) provocó la desintegración del vehículo, matando a sus $7$ tripulantes (incluyendo una maestra civil). Fue la tragedia del transbordador:
\begin{enumerate}[label=\textbf{\Alph*)}]
\begin{multicols}{2}
  \item Columbia
  \item Discovery
  \item Challenger
  \item Atlantis
  \item Endeavour
  \item Enterprise
\end{multicols}
\end{enumerate}
\vspace{0.15cm}
\item En febrero de $1990$, ante la inminente bancarrota económica por las sanciones internacionales y las huelgas internas, el último presidente blanco de Sudáfrica sorprendió al mundo legalizando el Congreso Nacional Africano (ANC) y liberando a Nelson Mandela. ¿De qué mandatario se trata?
\begin{enumerate}[label=\textbf{\Alph*)}]
\begin{multicols}{2}
  \item P. W. Botha
  \item Hendrik Verwoerd
  \item F. W. de Klerk
  \item Ian Smith
  \item Cecil Rhodes
  \item John Vorster
\end{multicols}
\end{enumerate}
\vspace{0.15cm}
\item El $1$ de enero de $1994$, el mismo día que entraba en vigor el Tratado de Libre Comercio de América del Norte (TLCAN/NAFTA), un grupo guerrillero indigenista se levantó en armas en el estado de Chiapas (México), declarando la guerra al gobierno y denunciando el neoliberalismo como "una sentencia de muerte". Este grupo es el:
\begin{enumerate}[label=\textbf{\Alph*)}]
\begin{multicols}{2}
  \item Sendero Luminoso
  \item FARC-EP
  \item Ejército Zapatista de Liberación Nacional (EZLN)
  \item Frente Sandinista
  \item Movimiento 19 de Abril (M-19)
  \item Montoneros
\end{multicols}
\end{enumerate}
\vspace{0.15cm}
\item A finales del $\text{Siglo XX}$, el motor principal de la globalización económica ($>50\%$ del comercio mundial) dejó de ser el intercambio entre Estados soberanos, pasando a ser liderado por gigantescas estructuras corporativas con cadenas de suministro fragmentadas por todo el globo. Estas entidades se denominan:
\begin{enumerate}[label=\textbf{\Alph*)}]
\begin{multicols}{2}
  \item Organizaciones No Gubernamentales (ONG)
  \item Empresas Transnacionales (ETN) o Multinacionales
  \item Gremios Sindicales Globales
  \item Bancos Centrales Nacionales
  \item Cooperativas Agrícolas
  \item Fondos Soberanos de Inversión
\end{multicols}
\end{enumerate}
\vspace{0.15cm}
\item Tras la desintegración de la Unión Soviética en $1991$ y antes de la guerra contra el terrorismo (post $2001$), el mundo experimentó una década ($1990$s) bajo un sistema internacional unipolar. Para describir la hegemonía militar, económica y cultural absoluta de Estados Unidos, el exministro francés Hubert Védrine popularizó el término:
\begin{enumerate}[label=\textbf{\Alph*)}]
\begin{multicols}{2}
  \item El Nuevo Imperio Romano
  \item La Hiperpotencia (Hyperpower)
  \item El Leviatán Democrático
  \item El Hegemón Solitario
  \item La República Imperial
  \item El Orden Neoliberal Absoluto
\end{multicols}
\end{enumerate}
\vspace{0.15cm}
\end{enumerate}\n
\vspace{0.4cm}
\begin{tcolorbox}[colback=gray!5!white,colframe=gray!60!black,title=\textbf{SOLUCIONARIO -- 4.2 Contexto Mundial}]
\begin{multicols}{4}
\begin{enumerate}[label=\textbf{\arabic*.}]
  \item \textbf{D}
  \item \textbf{C}
  \item \textbf{C}
  \item \textbf{B}
  \item \textbf{B}
  \item \textbf{B}
  \item \textbf{B}
  \item \textbf{C}
  \item \textbf{C}
  \item \textbf{C}
  \item \textbf{D}
  \item \textbf{C}
  \item \textbf{B}
  \item \textbf{B}
  \item \textbf{A}
  \item \textbf{B}
  \item \textbf{D}
  \item \textbf{C}
  \item \textbf{C}
  \item \textbf{B}
  \item \textbf{C}
  \item \textbf{C}
  \item \textbf{A}
  \item \textbf{C}
  \item \textbf{C}
  \item \textbf{C}
  \item \textbf{A}
  \item \textbf{C}
  \item \textbf{B}
  \item \textbf{C}
  \item \textbf{C}
  \item \textbf{A}
  \item \textbf{B}
  \item \textbf{D}
  \item \textbf{B}
  \item \textbf{C}
  \item \textbf{C}
  \item \textbf{C}
  \item \textbf{B}
  \item \textbf{C}
  \item \textbf{D}
  \item \textbf{B}
  \item \textbf{C}
  \item \textbf{D}
  \item \textbf{B}
  \item \textbf{B}
  \item \textbf{B}
  \item \textbf{C}
  \item \textbf{B}
  \item \textbf{C}
  \item \textbf{B}
  \item \textbf{D}
  \item \textbf{B}
  \item \textbf{C}
  \item \textbf{B}
  \item \textbf{D}
  \item \textbf{C}
  \item \textbf{B}
  \item \textbf{B}
  \item \textbf{B}
  \item \textbf{C}
  \item \textbf{D}
  \item \textbf{B}
  \item \textbf{C}
  \item \textbf{B}
  \item \textbf{B}
  \item \textbf{D}
  \item \textbf{C}
  \item \textbf{D}
  \item \textbf{C}
  \item \textbf{C}
  \item \textbf{B}
  \item \textbf{B}
  \item \textbf{C}
  \item \textbf{B}
  \item \textbf{C}
  \item \textbf{B}
  \item \textbf{D}
  \item \textbf{D}
  \item \textbf{D}
  \item \textbf{B}
  \item \textbf{B}
  \item \textbf{C}
  \item \textbf{C}
  \item \textbf{D}
  \item \textbf{C}
  \item \textbf{C}
  \item \textbf{D}
  \item \textbf{C}
  \item \textbf{B}
  \item \textbf{C}
  \item \textbf{D}
  \item \textbf{A}
  \item \textbf{B}
  \item \textbf{C}
  \item \textbf{B}
  \item \textbf{B}
  \item \textbf{C}
  \item \textbf{B}
  \item \textbf{B}
  \item \textbf{C}
  \item \textbf{D}
  \item \textbf{C}
  \item \textbf{B}
  \item \textbf{E}
  \item \textbf{C}
  \item \textbf{C}
  \item \textbf{B}
  \item \textbf{C}
  \item \textbf{D}
  \item \textbf{D}
  \item \textbf{C}
  \item \textbf{D}
  \item \textbf{C}
  \item \textbf{C}
  \item \textbf{D}
  \item \textbf{C}
  \item \textbf{C}
  \item \textbf{D}
  \item \textbf{A}
  \item \textbf{B}
  \item \textbf{D}
  \item \textbf{C}
  \item \textbf{C}
  \item \textbf{B}
  \item \textbf{D}
  \item \textbf{F}
  \item \textbf{F}
  \item \textbf{C}
  \item \textbf{D}
  \item \textbf{C}
  \item \textbf{B}
  \item \textbf{B}
  \item \textbf{D}
  \item \textbf{C}
  \item \textbf{C}
  \item \textbf{E}
  \item \textbf{E}
  \item \textbf{D}
  \item \textbf{E}
  \item \textbf{B}
  \item \textbf{C}
  \item \textbf{E}
  \item \textbf{D}
  \item \textbf{C}
  \item \textbf{D}
  \item \textbf{C}
  \item \textbf{B}
  \item \textbf{E}
  \item \textbf{C}
  \item \textbf{D}
  \item \textbf{D}
  \item \textbf{A}
  \item \textbf{B}
  \item \textbf{D}
  \item \textbf{C}
  \item \textbf{D}
  \item \textbf{C}
  \item \textbf{C}
  \item \textbf{C}
  \item \textbf{A}
  \item \textbf{C}
  \item \textbf{C}
  \item \textbf{C}
  \item \textbf{B}
  \item \textbf{A}
  \item \textbf{A}
  \item \textbf{A}
  \item \textbf{C}
  \item \textbf{D}
  \item \textbf{C}
  \item \textbf{B}
  \item \textbf{C}
  \item \textbf{C}
  \item \textbf{B}
  \item \textbf{A}
  \item \textbf{C}
  \item \textbf{C}
  \item \textbf{B}
  \item \textbf{B}
  \item \textbf{B}
  \item \textbf{A}
  \item \textbf{B}
  \item \textbf{C}
  \item \textbf{A}
  \item \textbf{B}
  \item \textbf{D}
  \item \textbf{C}
  \item \textbf{C}
  \item \textbf{C}
  \item \textbf{C}
  \item \textbf{B}
  \item \textbf{A}
  \item \textbf{C}
  \item \textbf{C}
  \item \textbf{B}
  \item \textbf{D}
  \item \textbf{C}
  \item \textbf{B}
  \item \textbf{D}
  \item \textbf{C}
  \item \textbf{B}
  \item \textbf{B}
  \item \textbf{A}
  \item \textbf{C}
  \item \textbf{B}
  \item \textbf{C}
  \item \textbf{B}
  \item \textbf{C}
  \item \textbf{D}
  \item \textbf{B}
  \item \textbf{C}
  \item \textbf{C}
  \item \textbf{D}
  \item \textbf{D}
  \item \textbf{C}
  \item \textbf{B}
  \item \textbf{B}
  \item \textbf{C}
  \item \textbf{C}
  \item \textbf{D}
  \item \textbf{D}
  \item \textbf{C}
  \item \textbf{B}
  \item \textbf{C}
  \item \textbf{C}
  \item \textbf{B}
  \item \textbf{C}
  \item \textbf{D}
  \item \textbf{B}
  \item \textbf{C}
  \item \textbf{B}
  \item \textbf{B}
  \item \textbf{B}
  \item \textbf{D}
  \item \textbf{C}
  \item \textbf{C}
  \item \textbf{D}
  \item \textbf{B}
  \item \textbf{C}
  \item \textbf{C}
  \item \textbf{C}
  \item \textbf{D}
  \item \textbf{A}
  \item \textbf{D}
  \item \textbf{B}
  \item \textbf{B}
  \item \textbf{A}
  \item \textbf{C}
  \item \textbf{D}
  \item \textbf{B}
  \item \textbf{C}
  \item \textbf{B}
  \item \textbf{C}
  \item \textbf{B}
  \item \textbf{C}
  \item \textbf{D}
  \item \textbf{D}
  \item \textbf{B}
  \item \textbf{C}
  \item \textbf{B}
  \item \textbf{D}
  \item \textbf{B}
  \item \textbf{A}
  \item \textbf{C}
  \item \textbf{B}
  \item \textbf{D}
  \item \textbf{B}
  \item \textbf{C}
  \item \textbf{B}
  \item \textbf{E}
  \item \textbf{C}
  \item \textbf{B}
  \item \textbf{C}
  \item \textbf{C}
  \item \textbf{C}
  \item \textbf{C}
  \item \textbf{C}
  \item \textbf{B}
  \item \textbf{A}
  \item \textbf{C}
  \item \textbf{B}
  \item \textbf{A}
  \item \textbf{C}
  \item \textbf{C}
  \item \textbf{C}
  \item \textbf{D}
  \item \textbf{C}
  \item \textbf{D}
  \item \textbf{B}
  \item \textbf{D}
  \item \textbf{B}
  \item \textbf{C}
  \item \textbf{C}
  \item \textbf{C}
  \item \textbf{C}
  \item \textbf{C}
  \item \textbf{C}
  \item \textbf{B}
  \item \textbf{B}
\end{enumerate}
\end{multicols}
\end{tcolorbox}
\newpage\n
\subsection{4.3 Sociedad, Cultura y Población}

\subsubsection*{Sección I: Sociedad y Cultura – Conceptos Básicos (1–40)}

\begin{enumerate}[label=\textbf{\arabic*.}]
\setcounter{enumi}{0}
\item ¿Qué es la sociedad?
\begin{enumerate}[label=\textbf{\Alph*)}]
\begin{multicols}{2}
  \item Un conjunto de individuos que interactúan y comparten normas y valores
  \item Un conjunto de edificios
  \item Un conjunto de leyes
  \item Un conjunto de animales
  \item Un conjunto de ríos
  \item Un conjunto de montañas
\end{multicols}
\end{enumerate}
\vspace{0.15cm}
\item ¿Qué es la cultura?
\begin{enumerate}[label=\textbf{\Alph*)}]
\begin{multicols}{2}
  \item El conjunto de conocimientos, creencias, arte, moral, derecho, costumbres y hábitos adquiridos por el ser humano en sociedad
  \item El conjunto de bienes materiales
  \item El conjunto de leyes jurídicas
  \item El conjunto de fenómenos naturales
  \item El conjunto de recursos económicos
  \item El conjunto de instituciones políticas
\end{multicols}
\end{enumerate}
\vspace{0.15cm}
\item ¿Qué es un fenómeno sociocultural?
\begin{enumerate}[label=\textbf{\Alph*)}]
\begin{multicols}{2}
  \item Un hecho o proceso que involucra la interacción social y cultural
  \item Un fenómeno natural
  \item Un fenómeno físico
  \item Un fenómeno químico
  \item Un fenómeno biológico
  \item Un fenómeno matemático
\end{multicols}
\end{enumerate}
\vspace{0.15cm}
\item ¿Qué es la socialización?
\begin{enumerate}[label=\textbf{\Alph*)}]
\begin{multicols}{2}
  \item El proceso mediante el cual el individuo aprende e interioriza la cultura de su sociedad
  \item El proceso de aislamiento social
  \item El proceso de migración
  \item El proceso de urbanización
  \item El proceso de industrialización
  \item El proceso de globalización
\end{multicols}
\end{enumerate}
\vspace{0.15cm}
\item ¿Qué es una norma social?
\begin{enumerate}[label=\textbf{\Alph*)}]
\begin{multicols}{2}
  \item Una regla de conducta compartida por un grupo social
  \item Una ley jurídica
  \item Un fenómeno natural
  \item Un recurso económico
  \item Un bien material
  \item Un fenómeno físico
\end{multicols}
\end{enumerate}
\vspace{0.15cm}
\item ¿Qué es un valor social?
\begin{enumerate}[label=\textbf{\Alph*)}]
\begin{multicols}{2}
  \item Un principio o creencia que guía la conducta de un grupo social
  \item Una norma jurídica
  \item Un fenómeno natural
  \item Un recurso económico
  \item Un bien material
  \item Un fenómeno físico
\end{multicols}
\end{enumerate}
\vspace{0.15cm}
\item ¿Qué es una institución social?
\begin{enumerate}[label=\textbf{\Alph*)}]
\begin{multicols}{2}
  \item Una organización establecida que satisface necesidades sociales
  \item Una norma jurídica
  \item Un fenómeno natural
  \item Un recurso económico
  \item Un bien material
  \item Un fenómeno físico
\end{multicols}
\end{enumerate}
\vspace{0.15cm}
\item ¿Cuál de las siguientes es una institución social?
\begin{enumerate}[label=\textbf{\Alph*)}]
\begin{multicols}{2}
  \item La familia
  \item Un río
  \item Una montaña
  \item Un mineral
  \item Un planeta
  \item Una estrella
\end{multicols}
\end{enumerate}
\vspace{0.15cm}
\item ¿Cuál de las siguientes es una institución social?
\begin{enumerate}[label=\textbf{\Alph*)}]
\begin{multicols}{2}
  \item La escuela
  \item Un árbol
  \item Un animal
  \item Una nube
  \item Una roca
  \item Un océano
\end{multicols}
\end{enumerate}
\vspace{0.15cm}
\item ¿Cuál de las siguientes es una institución social?
\begin{enumerate}[label=\textbf{\Alph*)}]
\begin{multicols}{2}
  \item El Estado
  \item Un volcán
  \item Un lago
  \item Un desierto
  \item Una selva
  \item Una playa
\end{multicols}
\end{enumerate}
\vspace{0.15cm}
\item ¿Qué es la familia?
\begin{enumerate}[label=\textbf{\Alph*)}]
\begin{multicols}{2}
  \item La institución social básica formada por personas unidas por parentesco
  \item Una organización política
  \item Una organización económica
  \item Una organización religiosa
  \item Una organización deportiva
  \item Una organización cultural
\end{multicols}
\end{enumerate}
\vspace{0.15cm}
\item ¿Qué es la comunidad?
\begin{enumerate}[label=\textbf{\Alph*)}]
\begin{multicols}{2}
  \item Un grupo de personas que comparten un territorio e intereses comunes
  \item Un grupo de personas sin relación
  \item Un grupo de edificios
  \item Un grupo de leyes
  \item Un grupo de animales
  \item Un grupo de plantas
\end{multicols}
\end{enumerate}
\vspace{0.15cm}
\item ¿Qué es un grupo social?
\begin{enumerate}[label=\textbf{\Alph*)}]
\begin{multicols}{2}
  \item Un conjunto de personas que interactúan y comparten normas
  \item Un conjunto de edificios
  \item Un conjunto de leyes
  \item Un conjunto de animales
  \item Un conjunto de plantas
  \item Un conjunto de minerales
\end{multicols}
\end{enumerate}
\vspace{0.15cm}
\item ¿Qué es el rol social?
\begin{enumerate}[label=\textbf{\Alph*)}]
\begin{multicols}{2}
  \item El papel o función que desempeña un individuo en la sociedad
  \item Una norma jurídica
  \item Un fenómeno natural
  \item Un recurso económico
  \item Un bien material
  \item Un fenómeno físico
\end{multicols}
\end{enumerate}
\vspace{0.15cm}
\item ¿Qué es el estatus social?
\begin{enumerate}[label=\textbf{\Alph*)}]
\begin{multicols}{2}
  \item La posición o rango que ocupa un individuo en la sociedad
  \item Una norma jurídica
  \item Un fenómeno natural
  \item Un recurso económico
  \item Un bien material
  \item Un fenómeno físico
\end{multicols}
\end{enumerate}
\vspace{0.15cm}
\item ¿Qué es la estratificación social?
\begin{enumerate}[label=\textbf{\Alph*)}]
\begin{multicols}{2}
  \item La división de la sociedad en estratos o clases jerárquicas
  \item La unión de la sociedad
  \item La igualdad social
  \item La movilidad social
  \item La cohesión social
  \item La integración social
\end{multicols}
\end{enumerate}
\vspace{0.15cm}
\item ¿Qué es la clase social?
\begin{enumerate}[label=\textbf{\Alph*)}]
\begin{multicols}{2}
  \item Un grupo de personas con características socioeconómicas similares
  \item Un grupo de personas con la misma edad
  \item Un grupo de personas con el mismo género
  \item Un grupo de personas con la misma religión
  \item Un grupo de personas con la misma nacionalidad
  \item Un grupo de personas con la misma profesión
\end{multicols}
\end{enumerate}
\vspace{0.15cm}
\item ¿Qué es la movilidad social?
\begin{enumerate}[label=\textbf{\Alph*)}]
\begin{multicols}{2}
  \item El cambio de posición social de un individuo o grupo
  \item La falta de cambio social
  \item La igualdad social
  \item La estratificación social
  \item La cohesión social
  \item La integración social
\end{multicols}
\end{enumerate}
\vspace{0.15cm}
\item ¿Qué es el cambio social?
\begin{enumerate}[label=\textbf{\Alph*)}]
\begin{multicols}{2}
  \item La transformación de las estructuras y relaciones sociales
  \item La permanencia de las estructuras sociales
  \item La igualdad social
  \item La estratificación social
  \item La cohesión social
  \item La integración social
\end{multicols}
\end{enumerate}
\vspace{0.15cm}
\item ¿Qué es el conflicto social?
\begin{enumerate}[label=\textbf{\Alph*)}]
\begin{multicols}{2}
  \item La oposición entre grupos con intereses divergentes
  \item La cooperación entre grupos
  \item La igualdad social
  \item La estratificación social
  \item La cohesión social
  \item La integración social
\end{multicols}
\end{enumerate}
\vspace{0.15cm}
\item ¿Qué es la cohesión social?
\begin{enumerate}[label=\textbf{\Alph*)}]
\begin{multicols}{2}
  \item La unión y solidaridad entre los miembros de una sociedad
  \item La división social
  \item El conflicto social
  \item La estratificación social
  \item La movilidad social
  \item El cambio social
\end{multicols}
\end{enumerate}
\vspace{0.15cm}
\item ¿Qué es la integración social?
\begin{enumerate}[label=\textbf{\Alph*)}]
\begin{multicols}{2}
  \item La incorporación de individuos o grupos a la sociedad
  \item La exclusión social
  \item El conflicto social
  \item La estratificación social
  \item La movilidad social
  \item El cambio social
\end{multicols}
\end{enumerate}
\vspace{0.15cm}
\item ¿Qué es la exclusión social?
\begin{enumerate}[label=\textbf{\Alph*)}]
\begin{multicols}{2}
  \item La marginación de individuos o grupos de la participación social
  \item La integración social
  \item La cohesión social
  \item La movilidad social
  \item El cambio social
  \item La igualdad social
\end{multicols}
\end{enumerate}
\vspace{0.15cm}
\item ¿Qué es la desigualdad social?
\begin{enumerate}[label=\textbf{\Alph*)}]
\begin{multicols}{2}
  \item La distribución inequitativa de recursos y oportunidades
  \item La igualdad social
  \item La cohesión social
  \item La movilidad social
  \item El cambio social
  \item La integración social
\end{multicols}
\end{enumerate}
\vspace{0.15cm}
\item ¿Qué es la pobreza?
\begin{enumerate}[label=\textbf{\Alph*)}]
\begin{multicols}{2}
  \item La carencia de recursos para satisfacer necesidades básicas
  \item La abundancia de recursos
  \item La igualdad social
  \item La cohesión social
  \item La movilidad social
  \item El cambio social
\end{multicols}
\end{enumerate}
\vspace{0.15cm}
\item ¿Qué es la pobreza extrema?
\begin{enumerate}[label=\textbf{\Alph*)}]
\begin{multicols}{2}
  \item La carencia de recursos para satisfacer necesidades alimentarias básicas
  \item La carencia de recursos para satisfacer necesidades no básicas
  \item La abundancia de recursos
  \item La igualdad social
  \item La cohesión social
  \item La movilidad social
\end{multicols}
\end{enumerate}
\vspace{0.15cm}
\item ¿Qué es el desarrollo humano?
\begin{enumerate}[label=\textbf{\Alph*)}]
\begin{multicols}{2}
  \item El proceso de ampliación de las oportunidades y capacidades de las personas
  \item El crecimiento económico
  \item La industrialización
  \item La urbanización
  \item La globalización
  \item La modernización
\end{multicols}
\end{enumerate}
\vspace{0.15cm}
\item ¿Qué es el Índice de Desarrollo Humano (IDH)?
\begin{enumerate}[label=\textbf{\Alph*)}]
\begin{multicols}{2}
  \item Un indicador que mide el desarrollo humano de un país
  \item Un indicador económico
  \item Un indicador político
  \item Un indicador ambiental
  \item Un indicador cultural
  \item Un indicador religioso
\end{multicols}
\end{enumerate}
\vspace{0.15cm}
\item ¿Qué dimensiones mide el IDH?
\begin{enumerate}[label=\textbf{\Alph*)}]
\begin{multicols}{2}
  \item Salud, educación e ingresos
  \item Salud, educación y política
  \item Salud, educación y religión
  \item Salud, educación y cultura
  \item Salud, educación y medio ambiente
  \item Salud, educación y tecnología
\end{multicols}
\end{enumerate}
\vspace{0.15cm}
\item ¿Qué es la calidad de vida?
\begin{enumerate}[label=\textbf{\Alph*)}]
\begin{multicols}{2}
  \item El grado de bienestar y satisfacción de las necesidades humanas
  \item La cantidad de bienes materiales
  \item El nivel de ingresos
  \item El nivel educativo
  \item El nivel de salud
  \item El nivel de consumo
\end{multicols}
\end{enumerate}
\vspace{0.15cm}
\item ¿Qué es el bienestar social?
\begin{enumerate}[label=\textbf{\Alph*)}]
\begin{multicols}{2}
  \item El conjunto de condiciones que permiten la satisfacción de las necesidades humanas
  \item La cantidad de bienes materiales
  \item El nivel de ingresos
  \item El nivel educativo
  \item El nivel de salud
  \item El nivel de consumo
\end{multicols}
\end{enumerate}
\vspace{0.15cm}
\item ¿Qué es la política social?
\begin{enumerate}[label=\textbf{\Alph*)}]
\begin{multicols}{2}
  \item El conjunto de acciones del Estado para satisfacer necesidades sociales
  \item La política económica
  \item La política exterior
  \item La política militar
  \item La política cultural
  \item La política ambiental
\end{multicols}
\end{enumerate}
\vspace{0.15cm}
\item ¿Qué es la seguridad social?
\begin{enumerate}[label=\textbf{\Alph*)}]
\begin{multicols}{2}
  \item El sistema de protección social para los ciudadanos
  \item La seguridad nacional
  \item La seguridad pública
  \item La seguridad alimentaria
  \item La seguridad energética
  \item La seguridad ambiental
\end{multicols}
\end{enumerate}
\vspace{0.15cm}
\item ¿Qué es la educación?
\begin{enumerate}[label=\textbf{\Alph*)}]
\begin{multicols}{2}
  \item El proceso de transmisión de conocimientos, valores y habilidades
  \item El proceso de producción económica
  \item El proceso de urbanización
  \item El proceso de industrialización
  \item El proceso de globalización
  \item El proceso de modernización
\end{multicols}
\end{enumerate}
\vspace{0.15cm}
\item ¿Qué es la salud?
\begin{enumerate}[label=\textbf{\Alph*)}]
\begin{multicols}{2}
  \item El estado de bienestar físico, mental y social
  \item La ausencia de enfermedad
  \item La presencia de enfermedad
  \item El estado de bienestar económico
  \item El estado de bienestar político
  \item El estado de bienestar cultural
\end{multicols}
\end{enumerate}
\vspace{0.15cm}
\item ¿Qué es la vivienda?
\begin{enumerate}[label=\textbf{\Alph*)}]
\begin{multicols}{2}
  \item El espacio físico donde habita una persona o familia
  \item El espacio de trabajo
  \item El espacio de recreación
  \item El espacio de educación
  \item El espacio de salud
  \item El espacio de comercio
\end{multicols}
\end{enumerate}
\vspace{0.15cm}
\item ¿Qué es el empleo?
\begin{enumerate}[label=\textbf{\Alph*)}]
\begin{multicols}{2}
  \item La actividad remunerada que realiza una persona
  \item La actividad no remunerada
  \item La actividad de ocio
  \item La actividad de estudio
  \item La actividad de voluntariado
  \item La actividad doméstica
\end{multicols}
\end{enumerate}
\vspace{0.15cm}
\item ¿Qué es el desempleo?
\begin{enumerate}[label=\textbf{\Alph*)}]
\begin{multicols}{2}
  \item La situación de una persona que busca trabajo y no lo encuentra
  \item La situación de una persona que tiene trabajo
  \item La situación de una persona que estudia
  \item La situación de una persona que se jubila
  \item La situación de una persona que trabaja
  \item La situación de una persona que emprende
\end{multicols}
\end{enumerate}
\vspace{0.15cm}
\item ¿Qué es el subempleo?
\begin{enumerate}[label=\textbf{\Alph*)}]
\begin{multicols}{2}
  \item La situación de una persona que trabaja menos horas o en empleos de menor calificación
  \item La situación de una persona que trabaja a tiempo completo
  \item La situación de una persona que estudia
  \item La situación de una persona que se jubila
  \item La situación de una persona que emprende
  \item La situación de una persona que invierte
\end{multicols}
\end{enumerate}
\vspace{0.15cm}
\item ¿Qué es la economía informal?
\begin{enumerate}[label=\textbf{\Alph*)}]
\begin{multicols}{2}
  \item El conjunto de actividades económicas no reguladas ni registradas
  \item El conjunto de actividades económicas reguladas
  \item El conjunto de actividades económicas formales
  \item El conjunto de actividades económicas estatales
  \item El conjunto de actividades económicas privadas
  \item El conjunto de actividades económicas mixtas
\end{multicols}
\end{enumerate}
\vspace{0.15cm}
\end{enumerate}\n
\subsubsection*{Sección II: Identidad y Diversidad Cultural (41–80)}

\begin{enumerate}[label=\textbf{\arabic*.}]
\setcounter{enumi}{40}
\item ¿Qué es la identidad cultural?
\begin{enumerate}[label=\textbf{\Alph*)}]
\begin{multicols}{2}
  \item El conjunto de rasgos culturales que definen a un grupo social
  \item El conjunto de rasgos biológicos
  \item El conjunto de rasgos económicos
  \item El conjunto de rasgos políticos
  \item El conjunto de rasgos religiosos
  \item El conjunto de rasgos geográficos
\end{multicols}
\end{enumerate}
\vspace{0.15cm}
\item ¿Qué es la diversidad cultural?
\begin{enumerate}[label=\textbf{\Alph*)}]
\begin{multicols}{2}
  \item La existencia de múltiples culturas en un mismo espacio
  \item La existencia de una sola cultura
  \item La existencia de múltiples culturas en espacios separados
  \item La existencia de una cultura dominante
  \item La existencia de una cultura homogénea
  \item La existencia de una cultura global
\end{multicols}
\end{enumerate}
\vspace{0.15cm}
\item ¿Qué es el multiculturalismo?
\begin{enumerate}[label=\textbf{\Alph*)}]
\begin{multicols}{2}
  \item La coexistencia de múltiples culturas en una sociedad
  \item La existencia de una sola cultura
  \item La existencia de una cultura dominante
  \item La existencia de una cultura homogénea
  \item La existencia de una cultura global
  \item La existencia de una cultura nacional
\end{multicols}
\end{enumerate}
\vspace{0.15cm}
\item ¿Qué es la interculturalidad?
\begin{enumerate}[label=\textbf{\Alph*)}]
\begin{multicols}{2}
  \item La interacción y diálogo entre culturas diferentes
  \item La existencia de una sola cultura
  \item La existencia de culturas separadas
  \item La existencia de una cultura dominante
  \item La existencia de una cultura homogénea
  \item La existencia de una cultura global
\end{multicols}
\end{enumerate}
\vspace{0.15cm}
\item ¿Qué es el etnocentrismo?
\begin{enumerate}[label=\textbf{\Alph*)}]
\begin{multicols}{2}
  \item La tendencia a juzgar otras culturas desde los parámetros de la propia
  \item La valoración de todas las culturas por igual
  \item La mezcla de culturas
  \item La separación de culturas
  \item La dominación de una cultura sobre otra
  \item La desaparición de una cultura
\end{multicols}
\end{enumerate}
\vspace{0.15cm}
\item ¿Qué es el relativismo cultural?
\begin{enumerate}[label=\textbf{\Alph*)}]
\begin{multicols}{2}
  \item La idea de que cada cultura debe ser entendida desde sus propios parámetros
  \item La idea de que todas las culturas son iguales
  \item La idea de que una cultura es superior a otras
  \item La idea de que las culturas deben desaparecer
  \item La idea de que las culturas deben mezclarse
  \item La idea de que las culturas deben separarse
\end{multicols}
\end{enumerate}
\vspace{0.15cm}
\item ¿Qué es la aculturación?
\begin{enumerate}[label=\textbf{\Alph*)}]
\begin{multicols}{2}
  \item El proceso de adopción de elementos de otra cultura
  \item El proceso de rechazo de otra cultura
  \item El proceso de aislamiento cultural
  \item El proceso de dominación cultural
  \item El proceso de desaparición cultural
  \item El proceso de mezcla cultural
\end{multicols}
\end{enumerate}
\vspace{0.15cm}
\item ¿Qué es la transculturación?
\begin{enumerate}[label=\textbf{\Alph*)}]
\begin{multicols}{2}
  \item El proceso de fusión e intercambio de elementos culturales
  \item El proceso de rechazo cultural
  \item El proceso de aislamiento cultural
  \item El proceso de dominación cultural
  \item El proceso de desaparición cultural
  \item El proceso de adopción cultural
\end{multicols}
\end{enumerate}
\vspace{0.15cm}
\item ¿Qué es la globalización cultural?
\begin{enumerate}[label=\textbf{\Alph*)}]
\begin{multicols}{2}
  \item La difusión de elementos culturales a nivel mundial
  \item La preservación de culturas locales
  \item La separación de culturas
  \item La dominación de una cultura sobre otras
  \item La desaparición de culturas
  \item La homogeneización cultural
\end{multicols}
\end{enumerate}
\vspace{0.15cm}
\item ¿Qué es la identidad nacional?
\begin{enumerate}[label=\textbf{\Alph*)}]
\begin{multicols}{2}
  \item El conjunto de rasgos que definen a una nación
  \item El conjunto de rasgos que definen a una región
  \item El conjunto de rasgos que definen a una etnia
  \item El conjunto de rasgos que definen a una familia
  \item El conjunto de rasgos que definen a un individuo
  \item El conjunto de rasgos que definen a una comunidad
\end{multicols}
\end{enumerate}
\vspace{0.15cm}
\item ¿Qué es el patrimonio cultural?
\begin{enumerate}[label=\textbf{\Alph*)}]
\begin{multicols}{2}
  \item El conjunto de bienes culturales materiales e inmateriales de un pueblo
  \item El conjunto de bienes económicos
  \item El conjunto de bienes naturales
  \item El conjunto de bienes políticos
  \item El conjunto de bienes religiosos
  \item El conjunto de bienes militares
\end{multicols}
\end{enumerate}
\vspace{0.15cm}
\item ¿Qué es el patrimonio cultural material?
\begin{enumerate}[label=\textbf{\Alph*)}]
\begin{multicols}{2}
  \item El conjunto de bienes culturales tangibles
  \item El conjunto de bienes culturales intangibles
  \item El conjunto de bienes naturales
  \item El conjunto de bienes económicos
  \item El conjunto de bienes políticos
  \item El conjunto de bienes religiosos
\end{multicols}
\end{enumerate}
\vspace{0.15cm}
\item ¿Qué es el patrimonio cultural inmaterial?
\begin{enumerate}[label=\textbf{\Alph*)}]
\begin{multicols}{2}
  \item El conjunto de bienes culturales intangibles
  \item El conjunto de bienes culturales tangibles
  \item El conjunto de bienes naturales
  \item El conjunto de bienes económicos
  \item El conjunto de bienes políticos
  \item El conjunto de bienes religiosos
\end{multicols}
\end{enumerate}
\vspace{0.15cm}
\item ¿Cuál es un ejemplo de patrimonio cultural material?
\begin{enumerate}[label=\textbf{\Alph*)}]
\begin{multicols}{2}
  \item Un sitio arqueológico
  \item Una canción
  \item Una danza
  \item Una leyenda
  \item Una costumbre
  \item Una lengua
\end{multicols}
\end{enumerate}
\vspace{0.15cm}
\item ¿Cuál es un ejemplo de patrimonio cultural inmaterial?
\begin{enumerate}[label=\textbf{\Alph*)}]
\begin{multicols}{2}
  \item Una danza tradicional
  \item Un edificio histórico
  \item Una escultura
  \item Una pintura
  \item Un monumento
  \item Un sitio arqueológico
\end{multicols}
\end{enumerate}
\vspace{0.15cm}
\item ¿Qué es la lengua?
\begin{enumerate}[label=\textbf{\Alph*)}]
\begin{multicols}{2}
  \item Un sistema de comunicación verbal propio de una comunidad
  \item Un sistema de comunicación no verbal
  \item Un sistema de comunicación animal
  \item Un sistema de comunicación tecnológico
  \item Un sistema de comunicación escrito
  \item Un sistema de comunicación simbólico
\end{multicols}
\end{enumerate}
\vspace{0.15cm}
\item ¿Qué es la religión?
\begin{enumerate}[label=\textbf{\Alph*)}]
\begin{multicols}{2}
  \item Un sistema de creencias y prácticas relacionadas con lo sagrado
  \item Un sistema de creencias políticas
  \item Un sistema de creencias económicas
  \item Un sistema de creencias científicas
  \item Un sistema de creencias artísticas
  \item Un sistema de creencias deportivas
\end{multicols}
\end{enumerate}
\vspace{0.15cm}
\item ¿Qué es la cosmovisión?
\begin{enumerate}[label=\textbf{\Alph*)}]
\begin{multicols}{2}
  \item La forma en que un grupo social interpreta el mundo
  \item La forma en que un individuo interpreta el mundo
  \item La forma en que una nación interpreta el mundo
  \item La forma en que una familia interpreta el mundo
  \item La forma en que una comunidad interpreta el mundo
  \item La forma en que una etnia interpreta el mundo
\end{multicols}
\end{enumerate}
\vspace{0.15cm}
\item ¿Qué es el sincretismo cultural?
\begin{enumerate}[label=\textbf{\Alph*)}]
\begin{multicols}{2}
  \item La fusión de elementos de diferentes culturas
  \item La separación de culturas
  \item La dominación de una cultura
  \item La desaparición de una cultura
  \item La preservación de una cultura
  \item La homogeneización cultural
\end{multicols}
\end{enumerate}
\vspace{0.15cm}
\item ¿Qué es la hibridación cultural?
\begin{enumerate}[label=\textbf{\Alph*)}]
\begin{multicols}{2}
  \item La mezcla de elementos culturales diversos
  \item La separación de culturas
  \item La dominación de una cultura
  \item La desaparición de una cultura
  \item La preservación de una cultura
  \item La homogeneización cultural
\end{multicols}
\end{enumerate}
\vspace{0.15cm}
\item ¿Qué es un estereotipo?
\begin{enumerate}[label=\textbf{\Alph*)}]
\begin{multicols}{2}
  \item Una imagen simplificada y generalizada de un grupo social
  \item Una imagen compleja de un grupo social
  \item Una imagen objetiva de un grupo social
  \item Una imagen científica de un grupo social
  \item Una imagen artística de un grupo social
  \item Una imagen histórica de un grupo social
\end{multicols}
\end{enumerate}
\vspace{0.15cm}
\item ¿Qué es un prejuicio?
\begin{enumerate}[label=\textbf{\Alph*)}]
\begin{multicols}{2}
  \item Una opinión preconcebida sobre un grupo social
  \item Una opinión fundamentada sobre un grupo social
  \item Una opinión objetiva sobre un grupo social
  \item Una opinión científica sobre un grupo social
  \item Una opinión artística sobre un grupo social
  \item Una opinión histórica sobre un grupo social
\end{multicols}
\end{enumerate}
\vspace{0.15cm}
\item ¿Qué es la discriminación?
\begin{enumerate}[label=\textbf{\Alph*)}]
\begin{multicols}{2}
  \item El trato desigual e injusto hacia un grupo social
  \item El trato igualitario hacia un grupo social
  \item El trato justo hacia un grupo social
  \item El trato objetivo hacia un grupo social
  \item El trato científico hacia un grupo social
  \item El trato artístico hacia un grupo social
\end{multicols}
\end{enumerate}
\vspace{0.15cm}
\item ¿Qué es el racismo?
\begin{enumerate}[label=\textbf{\Alph*)}]
\begin{multicols}{2}
  \item La discriminación basada en la raza o etnia
  \item La discriminación basada en el género
  \item La discriminación basada en la edad
  \item La discriminación basada en la religión
  \item La discriminación basada en la clase social
  \item La discriminación basada en la nacionalidad
\end{multicols}
\end{enumerate}
\vspace{0.15cm}
\item ¿Qué es el sexismo?
\begin{enumerate}[label=\textbf{\Alph*)}]
\begin{multicols}{2}
  \item La discriminación basada en el sexo o género
  \item La discriminación basada en la raza
  \item La discriminación basada en la edad
  \item La discriminación basada en la religión
  \item La discriminación basada en la clase social
  \item La discriminación basada en la nacionalidad
\end{multicols}
\end{enumerate}
\vspace{0.15cm}
\item ¿Qué es la xenofobia?
\begin{enumerate}[label=\textbf{\Alph*)}]
\begin{multicols}{2}
  \item El rechazo hacia los extranjeros
  \item El rechazo hacia los nacionales
  \item El rechazo hacia los ancianos
  \item El rechazo hacia los niños
  \item El rechazo hacia los discapacitados
  \item El rechazo hacia los pobres
\end{multicols}
\end{enumerate}
\vspace{0.15cm}
\item ¿Qué es la homofobia?
\begin{enumerate}[label=\textbf{\Alph*)}]
\begin{multicols}{2}
  \item El rechazo hacia las personas homosexuales
  \item El rechazo hacia las personas heterosexuales
  \item El rechazo hacia las personas bisexuales
  \item El rechazo hacia las personas transgénero
  \item El rechazo hacia las personas intersexuales
  \item El rechazo hacia las personas asexuales
\end{multicols}
\end{enumerate}
\vspace{0.15cm}
\item ¿Qué es la inclusión social?
\begin{enumerate}[label=\textbf{\Alph*)}]
\begin{multicols}{2}
  \item La incorporación de todos los grupos sociales a la vida social
  \item La exclusión de algunos grupos sociales
  \item La marginación de algunos grupos sociales
  \item La discriminación de algunos grupos sociales
  \item La segregación de algunos grupos sociales
  \item La separación de algunos grupos sociales
\end{multicols}
\end{enumerate}
\vspace{0.15cm}
\item ¿Qué es la equidad de género?
\begin{enumerate}[label=\textbf{\Alph*)}]
\begin{multicols}{2}
  \item La igualdad de derechos y oportunidades entre hombres y mujeres
  \item La superioridad de un género sobre otro
  \item La inferioridad de un género sobre otro
  \item La diferencia de derechos entre géneros
  \item La diferencia de oportunidades entre géneros
  \item La separación de géneros
\end{multicols}
\end{enumerate}
\vspace{0.15cm}
\item ¿Qué es el empoderamiento?
\begin{enumerate}[label=\textbf{\Alph*)}]
\begin{multicols}{2}
  \item El proceso de adquisición de poder y autonomía por parte de grupos vulnerables
  \item El proceso de pérdida de poder
  \item El proceso de dominación
  \item El proceso de exclusión
  \item El proceso de marginación
  \item El proceso de discriminación
\end{multicols}
\end{enumerate}
\vspace{0.15cm}
\item ¿Qué es la ciudadanía?
\begin{enumerate}[label=\textbf{\Alph*)}]
\begin{multicols}{2}
  \item La condición de pertenencia a un Estado con derechos y deberes
  \item La condición de pertenencia a una familia
  \item La condición de pertenencia a una comunidad
  \item La condición de pertenencia a una etnia
  \item La condición de pertenencia a una religión
  \item La condición de pertenencia a un partido
\end{multicols}
\end{enumerate}
\vspace{0.15cm}
\item ¿Qué es la participación ciudadana?
\begin{enumerate}[label=\textbf{\Alph*)}]
\begin{multicols}{2}
  \item La intervención de los ciudadanos en los asuntos públicos
  \item La abstención de los ciudadanos
  \item La indiferencia de los ciudadanos
  \item La pasividad de los ciudadanos
  \item La exclusión de los ciudadanos
  \item La marginación de los ciudadanos
\end{multicols}
\end{enumerate}
\vspace{0.15cm}
\item ¿Qué es el capital social?
\begin{enumerate}[label=\textbf{\Alph*)}]
\begin{multicols}{2}
  \item El conjunto de redes y relaciones sociales basadas en la confianza
  \item El conjunto de bienes materiales
  \item El conjunto de bienes económicos
  \item El conjunto de bienes políticos
  \item El conjunto de bienes culturales
  \item El conjunto de bienes naturales
\end{multicols}
\end{enumerate}
\vspace{0.15cm}
\item ¿Qué es la sociedad civil?
\begin{enumerate}[label=\textbf{\Alph*)}]
\begin{multicols}{2}
  \item El conjunto de organizaciones sociales independientes del Estado
  \item El conjunto de instituciones estatales
  \item El conjunto de empresas privadas
  \item El conjunto de partidos políticos
  \item El conjunto de sindicatos
  \item El conjunto de organizaciones religiosas
\end{multicols}
\end{enumerate}
\vspace{0.15cm}
\item ¿Qué es una ONG?
\begin{enumerate}[label=\textbf{\Alph*)}]
\begin{multicols}{2}
  \item Una organización no gubernamental
  \item Una organización gubernamental
  \item Una organización empresarial
  \item Una organización política
  \item Una organización religiosa
  \item Una organización deportiva
\end{multicols}
\end{enumerate}
\vspace{0.15cm}
\item ¿Qué es el voluntariado?
\begin{enumerate}[label=\textbf{\Alph*)}]
\begin{multicols}{2}
  \item El trabajo no remunerado realizado libremente para ayudar a otros
  \item El trabajo remunerado
  \item El trabajo forzado
  \item El trabajo informal
  \item El trabajo formal
  \item El trabajo doméstico
\end{multicols}
\end{enumerate}
\vspace{0.15cm}
\item ¿Qué es la solidaridad?
\begin{enumerate}[label=\textbf{\Alph*)}]
\begin{multicols}{2}
  \item La ayuda mutua entre los miembros de una sociedad
  \item La indiferencia entre los miembros de una sociedad
  \item La competencia entre los miembros de una sociedad
  \item La exclusión entre los miembros de una sociedad
  \item La marginación entre los miembros de una sociedad
  \item La discriminación entre los miembros de una sociedad
\end{multicols}
\end{enumerate}
\vspace{0.15cm}
\item ¿Qué es la tolerancia?
\begin{enumerate}[label=\textbf{\Alph*)}]
\begin{multicols}{2}
  \item El respeto hacia las ideas y prácticas diferentes
  \item El rechazo hacia las ideas diferentes
  \item La indiferencia hacia las ideas diferentes
  \item La imposición de las ideas propias
  \item La discriminación de las ideas diferentes
  \item La exclusión de las ideas diferentes
\end{multicols}
\end{enumerate}
\vspace{0.15cm}
\item ¿Qué es el diálogo intercultural?
\begin{enumerate}[label=\textbf{\Alph*)}]
\begin{multicols}{2}
  \item La comunicación y entendimiento entre culturas diferentes
  \item La imposición de una cultura sobre otra
  \item La separación de culturas
  \item La dominación de una cultura
  \item La desaparición de una cultura
  \item La homogeneización cultural
\end{multicols}
\end{enumerate}
\vspace{0.15cm}
\item ¿Qué es la paz social?
\begin{enumerate}[label=\textbf{\Alph*)}]
\begin{multicols}{2}
  \item La convivencia armónica y sin conflictos violentos en la sociedad
  \item La ausencia de conflictos
  \item La presencia de conflictos
  \item La violencia social
  \item La guerra civil
  \item La revolución social
\end{multicols}
\end{enumerate}
\vspace{0.15cm}
\end{enumerate}\n
\subsubsection*{Sección III: Globalización (81–120)}

\begin{enumerate}[label=\textbf{\arabic*.}]
\setcounter{enumi}{80}
\item ¿Qué es la globalización?
\begin{enumerate}[label=\textbf{\Alph*)}]
\begin{multicols}{2}
  \item El proceso de integración económica, política, social y cultural a nivel mundial
  \item El proceso de aislamiento de los países
  \item El proceso de separación de las culturas
  \item El proceso de fragmentación política
  \item El proceso de localización económica
  \item El proceso de regionalización
\end{multicols}
\end{enumerate}
\vspace{0.15cm}
\item ¿Cuál es una dimensión de la globalización?
\begin{enumerate}[label=\textbf{\Alph*)}]
\begin{multicols}{2}
  \item La económica
  \item La biológica
  \item La geológica
  \item La astronómica
  \item La química
  \item La física
\end{multicols}
\end{enumerate}
\vspace{0.15cm}
\item ¿Cuál es una dimensión de la globalización?
\begin{enumerate}[label=\textbf{\Alph*)}]
\begin{multicols}{2}
  \item La política
  \item La biológica
  \item La geológica
  \item La astronómica
  \item La química
  \item La física
\end{multicols}
\end{enumerate}
\vspace{0.15cm}
\item ¿Cuál es una dimensión de la globalización?
\begin{enumerate}[label=\textbf{\Alph*)}]
\begin{multicols}{2}
  \item La cultural
  \item La biológica
  \item La geológica
  \item La astronómica
  \item La química
  \item La física
\end{multicols}
\end{enumerate}
\vspace{0.15cm}
\item ¿Cuál es una dimensión de la globalización?
\begin{enumerate}[label=\textbf{\Alph*)}]
\begin{multicols}{2}
  \item La social
  \item La biológica
  \item La geológica
  \item La astronómica
  \item La química
  \item La física
\end{multicols}
\end{enumerate}
\vspace{0.15cm}
\item ¿Qué es la globalización económica?
\begin{enumerate}[label=\textbf{\Alph*)}]
\begin{multicols}{2}
  \item La integración de los mercados y economías a nivel mundial
  \item La separación de los mercados
  \item La localización de los mercados
  \item La regionalización de los mercados
  \item La fragmentación de los mercados
  \item El aislamiento de los mercados
\end{multicols}
\end{enumerate}
\vspace{0.15cm}
\item ¿Qué es el comercio internacional?
\begin{enumerate}[label=\textbf{\Alph*)}]
\begin{multicols}{2}
  \item El intercambio de bienes y servicios entre países
  \item El intercambio de bienes y servicios dentro de un país
  \item El intercambio de bienes y servicios en una región
  \item El intercambio de bienes y servicios en una ciudad
  \item El intercambio de bienes y servicios en una comunidad
  \item El intercambio de bienes y servicios en una familia
\end{multicols}
\end{enumerate}
\vspace{0.15cm}
\item ¿Qué es la inversión extranjera?
\begin{enumerate}[label=\textbf{\Alph*)}]
\begin{multicols}{2}
  \item La inversión de capital de un país en otro
  \item La inversión de capital dentro del mismo país
  \item La inversión de capital en una región
  \item La inversión de capital en una ciudad
  \item La inversión de capital en una comunidad
  \item La inversión de capital en una familia
\end{multicols}
\end{enumerate}
\vspace{0.15cm}
\item ¿Qué es una empresa transnacional?
\begin{enumerate}[label=\textbf{\Alph*)}]
\begin{multicols}{2}
  \item Una empresa que opera en varios países
  \item Una empresa que opera en un solo país
  \item Una empresa que opera en una región
  \item Una empresa que opera en una ciudad
  \item Una empresa que opera en una comunidad
  \item Una empresa que opera en una familia
\end{multicols}
\end{enumerate}
\vspace{0.15cm}
\item ¿Qué es la deslocalización?
\begin{enumerate}[label=\textbf{\Alph*)}]
\begin{multicols}{2}
  \item El traslado de la producción a países con menores costos
  \item El traslado de la producción al país de origen
  \item El traslado de la producción a una región
  \item El traslado de la producción a una ciudad
  \item El traslado de la producción a una comunidad
  \item El traslado de la producción a una familia
\end{multicols}
\end{enumerate}
\vspace{0.15cm}
\item ¿Qué es la globalización política?
\begin{enumerate}[label=\textbf{\Alph*)}]
\begin{multicols}{2}
  \item La creciente interconexión de los Estados y organizaciones internacionales
  \item La separación de los Estados
  \item El aislamiento de los Estados
  \item La fragmentación de los Estados
  \item La localización de los Estados
  \item La regionalización de los Estados
\end{multicols}
\end{enumerate}
\vspace{0.15cm}
\item ¿Qué es un organismo internacional?
\begin{enumerate}[label=\textbf{\Alph*)}]
\begin{multicols}{2}
  \item Una organización formada por varios Estados para cooperar en temas comunes
  \item Una organización formada por un solo Estado
  \item Una organización formada por empresas
  \item Una organización formada por individuos
  \item Una organización formada por comunidades
  \item Una organización formada por familias
\end{multicols}
\end{enumerate}
\vspace{0.15cm}
\item ¿Cuál es un organismo internacional?
\begin{enumerate}[label=\textbf{\Alph*)}]
\begin{multicols}{2}
  \item La ONU
  \item Una alcaldía
  \item Un ministerio
  \item Un congreso
  \item Un tribunal
  \item Una gobernación
\end{multicols}
\end{enumerate}
\vspace{0.15cm}
\item ¿Cuál es un organismo internacional?
\begin{enumerate}[label=\textbf{\Alph*)}]
\begin{multicols}{2}
  \item La OEA
  \item Una alcaldía
  \item Un ministerio
  \item Un congreso
  \item Un tribunal
  \item Una gobernación
\end{multicols}
\end{enumerate}
\vspace{0.15cm}
\item ¿Cuál es un organismo internacional?
\begin{enumerate}[label=\textbf{\Alph*)}]
\begin{multicols}{2}
  \item El FMI
  \item Una alcaldía
  \item Un ministerio
  \item Un congreso
  \item Un tribunal
  \item Una gobernación
\end{multicols}
\end{enumerate}
\vspace{0.15cm}
\item ¿Cuál es un organismo internacional?
\begin{enumerate}[label=\textbf{\Alph*)}]
\begin{multicols}{2}
  \item El Banco Mundial
  \item Una alcaldía
  \item Un ministerio
  \item Un congreso
  \item Un tribunal
  \item Una gobernación
\end{multicols}
\end{enumerate}
\vspace{0.15cm}
\item ¿Cuál es un organismo internacional?
\begin{enumerate}[label=\textbf{\Alph*)}]
\begin{multicols}{2}
  \item La OMC
  \item Una alcaldía
  \item Un ministerio
  \item Un congreso
  \item Un tribunal
  \item Una gobernación
\end{multicols}
\end{enumerate}
\vspace{0.15cm}
\item ¿Qué es la globalización cultural?
\begin{enumerate}[label=\textbf{\Alph*)}]
\begin{multicols}{2}
  \item La difusión de ideas, valores y prácticas culturales a nivel mundial
  \item La preservación de culturas locales
  \item La separación de culturas
  \item La dominación de una cultura
  \item La desaparición de culturas
  \item La homogeneización cultural
\end{multicols}
\end{enumerate}
\vspace{0.15cm}
\item ¿Qué es la aldea global?
\begin{enumerate}[label=\textbf{\Alph*)}]
\begin{multicols}{2}
  \item La idea de que el mundo se ha interconectado como una aldea
  \item La idea de que el mundo está aislado
  \item La idea de que el mundo está fragmentado
  \item La idea de que el mundo está separado
  \item La idea de que el mundo está localizado
  \item La idea de que el mundo está regionalizado
\end{multicols}
\end{enumerate}
\vspace{0.15cm}
\item ¿Qué es la sociedad de la información?
\begin{enumerate}[label=\textbf{\Alph*)}]
\begin{multicols}{2}
  \item Una sociedad en la que la información y el conocimiento son centrales
  \item Una sociedad en la que la agricultura es central
  \item Una sociedad en la que la industria es central
  \item Una sociedad en la que el comercio es central
  \item Una sociedad en la que la minería es central
  \item Una sociedad en la que la pesca es central
\end{multicols}
\end{enumerate}
\vspace{0.15cm}
\item ¿Qué es la brecha digital?
\begin{enumerate}[label=\textbf{\Alph*)}]
\begin{multicols}{2}
  \item La desigualdad en el acceso y uso de las tecnologías de la información
  \item La igualdad en el acceso a la tecnología
  \item La abundancia de tecnología
  \item La escasez de tecnología
  \item La distribución equitativa de tecnología
  \item La democratización de la tecnología
\end{multicols}
\end{enumerate}
\vspace{0.15cm}
\item ¿Qué es la sociedad del conocimiento?
\begin{enumerate}[label=\textbf{\Alph*)}]
\begin{multicols}{2}
  \item Una sociedad en la que el conocimiento es el principal recurso
  \item Una sociedad en la que la agricultura es el principal recurso
  \item Una sociedad en la que la industria es el principal recurso
  \item Una sociedad en la que el comercio es el principal recurso
  \item Una sociedad en la que la minería es el principal recurso
  \item Una sociedad en la que la pesca es el principal recurso
\end{multicols}
\end{enumerate}
\vspace{0.15cm}
\item ¿Qué es la globalización social?
\begin{enumerate}[label=\textbf{\Alph*)}]
\begin{multicols}{2}
  \item La creciente interconexión de las sociedades y personas a nivel mundial
  \item La separación de las sociedades
  \item El aislamiento de las sociedades
  \item La fragmentación de las sociedades
  \item La localización de las sociedades
  \item La regionalización de las sociedades
\end{multicols}
\end{enumerate}
\vspace{0.15cm}
\item ¿Qué es la migración internacional?
\begin{enumerate}[label=\textbf{\Alph*)}]
\begin{multicols}{2}
  \item El desplazamiento de personas de un país a otro
  \item El desplazamiento de personas dentro de un país
  \item El desplazamiento de personas dentro de una región
  \item El desplazamiento de personas dentro de una ciudad
  \item El desplazamiento de personas dentro de una comunidad
  \item El desplazamiento de personas dentro de una familia
\end{multicols}
\end{enumerate}
\vspace{0.15cm}
\item ¿Qué es la diáspora?
\begin{enumerate}[label=\textbf{\Alph*)}]
\begin{multicols}{2}
  \item La dispersión de un pueblo o comunidad por el mundo
  \item La concentración de un pueblo en un lugar
  \item El aislamiento de un pueblo
  \item La separación de un pueblo
  \item La fragmentación de un pueblo
  \item La localización de un pueblo
\end{multicols}
\end{enumerate}
\vspace{0.15cm}
\item ¿Qué es la ciudadanía global?
\begin{enumerate}[label=\textbf{\Alph*)}]
\begin{multicols}{2}
  \item La conciencia de pertenencia a una comunidad mundial
  \item La conciencia de pertenencia a una nación
  \item La conciencia de pertenencia a una región
  \item La conciencia de pertenencia a una ciudad
  \item La conciencia de pertenencia a una comunidad
  \item La conciencia de pertenencia a una familia
\end{multicols}
\end{enumerate}
\vspace{0.15cm}
\item ¿Qué es el consumo responsable?
\begin{enumerate}[label=\textbf{\Alph*)}]
\begin{multicols}{2}
  \item El consumo que considera el impacto social y ambiental
  \item El consumo sin considerar el impacto
  \item El consumo excesivo
  \item El consumo irracional
  \item El consumo compulsivo
  \item El consumo desmedido
\end{multicols}
\end{enumerate}
\vspace{0.15cm}
\item ¿Qué es el desarrollo sostenible?
\begin{enumerate}[label=\textbf{\Alph*)}]
\begin{multicols}{2}
  \item El desarrollo que satisface las necesidades actuales sin comprometer las futuras
  \item El desarrollo que satisface solo las necesidades actuales
  \item El desarrollo que satisface solo las necesidades futuras
  \item El desarrollo que no satisface necesidades
  \item El desarrollo que ignora las necesidades
  \item El desarrollo que prioriza lo económico
\end{multicols}
\end{enumerate}
\vspace{0.15cm}
\item ¿Qué es la huella ecológica?
\begin{enumerate}[label=\textbf{\Alph*)}]
\begin{multicols}{2}
  \item La medida del impacto de las actividades humanas en el medio ambiente
  \item La medida del impacto de las actividades naturales
  \item La medida del impacto de las actividades económicas
  \item La medida del impacto de las actividades políticas
  \item La medida del impacto de las actividades culturales
  \item La medida del impacto de las actividades religiosas
\end{multicols}
\end{enumerate}
\vspace{0.15cm}
\item ¿Qué es la globalización alternativa?
\begin{enumerate}[label=\textbf{\Alph*)}]
\begin{multicols}{2}
  \item Una globalización basada en la solidaridad y la justicia social
  \item Una globalización basada en el libre mercado
  \item Una globalización basada en la competencia
  \item Una globalización basada en la exclusión
  \item Una globalización basada en la desigualdad
  \item Una globalización basada en la dominación
\end{multicols}
\end{enumerate}
\vspace{0.15cm}
\item ¿Qué es la antiglobalización?
\begin{enumerate}[label=\textbf{\Alph*)}]
\begin{multicols}{2}
  \item Un movimiento que critica los efectos negativos de la globalización
  \item Un movimiento que apoya la globalización
  \item Un movimiento que ignora la globalización
  \item Un movimiento que promueve la globalización
  \item Un movimiento que estudia la globalización
  \item Un movimiento que celebra la globalización
\end{multicols}
\end{enumerate}
\vspace{0.15cm}
\item ¿Qué es el Foro Social Mundial?
\begin{enumerate}[label=\textbf{\Alph*)}]
\begin{multicols}{2}
  \item Un espacio de encuentro de movimientos sociales críticos de la globalización
  \item Un espacio de encuentro de empresas transnacionales
  \item Un espacio de encuentro de gobiernos
  \item Un espacio de encuentro de partidos políticos
  \item Un espacio de encuentro de organizaciones religiosas
  \item Un espacio de encuentro de organizaciones deportivas
\end{multicols}
\end{enumerate}
\vspace{0.15cm}
\item ¿Qué es el Foro Económico Mundial?
\begin{enumerate}[label=\textbf{\Alph*)}]
\begin{multicols}{2}
  \item Un espacio de encuentro de líderes económicos y políticos
  \item Un espacio de encuentro de movimientos sociales
  \item Un espacio de encuentro de organizaciones no gubernamentales
  \item Un espacio de encuentro de sindicatos
  \item Un espacio de encuentro de estudiantes
  \item Un espacio de encuentro de artistas
\end{multicols}
\end{enumerate}
\vspace{0.15cm}
\item ¿Qué es el G20?
\begin{enumerate}[label=\textbf{\Alph*)}]
\begin{multicols}{2}
  \item Un grupo de países con las economías más grandes del mundo
  \item Un grupo de países pequeños
  \item Un grupo de países en desarrollo
  \item Un grupo de países subdesarrollados
  \item Un grupo de países sin economía
  \item Un grupo de países sin influencia
\end{multicols}
\end{enumerate}
\vspace{0.15cm}
\item ¿Qué es el BRICS?
\begin{enumerate}[label=\textbf{\Alph*)}]
\begin{multicols}{2}
  \item Un grupo de países emergentes: Brasil, Rusia, India, China y Sudáfrica
  \item Un grupo de países desarrollados
  \item Un grupo de países subdesarrollados
  \item Un grupo de países sin economía
  \item Un grupo de países sin influencia
  \item Un grupo de países pequeños
\end{multicols}
\end{enumerate}
\vspace{0.15cm}
\item ¿Qué es la Unión Europea?
\begin{enumerate}[label=\textbf{\Alph*)}]
\begin{multicols}{2}
  \item Una organización de integración política y económica de países europeos
  \item Una organización de países americanos
  \item Una organización de países asiáticos
  \item Una organización de países africanos
  \item Una organización de países oceánicos
  \item Una organización de países antárticos
\end{multicols}
\end{enumerate}
\vspace{0.15cm}
\item ¿Qué es el Mercosur?
\begin{enumerate}[label=\textbf{\Alph*)}]
\begin{multicols}{2}
  \item Una organización de integración económica de países sudamericanos
  \item Una organización de países centroamericanos
  \item Una organización de países norteamericanos
  \item Una organización de países europeos
  \item Una organización de países asiáticos
  \item Una organización de países africanos
\end{multicols}
\end{enumerate}
\vspace{0.15cm}
\item ¿Qué es el SICA?
\begin{enumerate}[label=\textbf{\Alph*)}]
\begin{multicols}{2}
  \item El Sistema de la Integración Centroamericana
  \item El Sistema de la Integración Sudamericana
  \item El Sistema de la Integración Norteamericana
  \item El Sistema de la Integración Europea
  \item El Sistema de la Integración Asiática
  \item El Sistema de la Integración Africana
\end{multicols}
\end{enumerate}
\vspace{0.15cm}
\item ¿Qué es el TLC?
\begin{enumerate}[label=\textbf{\Alph*)}]
\begin{multicols}{2}
  \item Un Tratado de Libre Comercio
  \item Un Tratado de Libre Tránsito
  \item Un Tratado de Libre Trabajo
  \item Un Tratado de Libre Territorio
  \item Un Tratado de Libre Tecnología
  \item Un Tratado de Libre Turismo
\end{multicols}
\end{enumerate}
\vspace{0.15cm}
\item ¿Qué es la cooperación internacional?
\begin{enumerate}[label=\textbf{\Alph*)}]
\begin{multicols}{2}
  \item La colaboración entre países para alcanzar objetivos comunes
  \item La competencia entre países
  \item El conflicto entre países
  \item La separación entre países
  \item El aislamiento entre países
  \item La dominación entre países
\end{multicols}
\end{enumerate}
\vspace{0.15cm}
\end{enumerate}\n
\subsubsection*{Sección IV: Demografía (121–160)}

\begin{enumerate}[label=\textbf{\arabic*.}]
\setcounter{enumi}{120}
\item ¿Qué es la demografía?
\begin{enumerate}[label=\textbf{\Alph*)}]
\begin{multicols}{2}
  \item La ciencia que estudia la población humana
  \item La ciencia que estudia los animales
  \item La ciencia que estudia las plantas
  \item La ciencia que estudia los minerales
  \item La ciencia que estudia los astros
  \item La ciencia que estudia los fenómenos físicos
\end{multicols}
\end{enumerate}
\vspace{0.15cm}
\item ¿Qué es la población?
\begin{enumerate}[label=\textbf{\Alph*)}]
\begin{multicols}{2}
  \item El conjunto de personas que habitan un territorio
  \item El conjunto de animales que habitan un territorio
  \item El conjunto de plantas que habitan un territorio
  \item El conjunto de minerales que habitan un territorio
  \item El conjunto de astros que habitan un territorio
  \item El conjunto de fenómenos que habitan un territorio
\end{multicols}
\end{enumerate}
\vspace{0.15cm}
\item ¿Qué es el crecimiento poblacional?
\begin{enumerate}[label=\textbf{\Alph*)}]
\begin{multicols}{2}
  \item El aumento del número de habitantes de una población
  \item La disminución del número de habitantes
  \item La estabilidad del número de habitantes
  \item La migración de habitantes
  \item La distribución de habitantes
  \item La composición de habitantes
\end{multicols}
\end{enumerate}
\vspace{0.15cm}
\item ¿Qué es la tasa de natalidad?
\begin{enumerate}[label=\textbf{\Alph*)}]
\begin{multicols}{2}
  \item El número de nacimientos por cada mil habitantes en un año
  \item El número de muertes por cada mil habitantes en un año
  \item El número de migrantes por cada mil habitantes en un año
  \item El número de habitantes por kilómetro cuadrado
  \item El número de familias por cada mil habitantes
  \item El número de matrimonios por cada mil habitantes
\end{multicols}
\end{enumerate}
\vspace{0.15cm}
\item ¿Qué es la tasa de mortalidad?
\begin{enumerate}[label=\textbf{\Alph*)}]
\begin{multicols}{2}
  \item El número de muertes por cada mil habitantes en un año
  \item El número de nacimientos por cada mil habitantes en un año
  \item El número de migrantes por cada mil habitantes en un año
  \item El número de habitantes por kilómetro cuadrado
  \item El número de familias por cada mil habitantes
  \item El número de matrimonios por cada mil habitantes
\end{multicols}
\end{enumerate}
\vspace{0.15cm}
\item ¿Qué es la tasa de mortalidad infantil?
\begin{enumerate}[label=\textbf{\Alph*)}]
\begin{multicols}{2}
  \item El número de muertes de niños menores de un año por cada mil nacidos vivos
  \item El número de muertes de niños menores de cinco años
  \item El número de muertes de niños menores de diez años
  \item El número de muertes de niños menores de quince años
  \item El número de muertes de niños menores de dieciocho años
  \item El número de muertes de niños menores de veintiún años
\end{multicols}
\end{enumerate}
\vspace{0.15cm}
\item ¿Qué es la esperanza de vida?
\begin{enumerate}[label=\textbf{\Alph*)}]
\begin{multicols}{2}
  \item El número promedio de años que se espera que viva una persona
  \item El número máximo de años que puede vivir una persona
  \item El número mínimo de años que puede vivir una persona
  \item El número de años que ha vivido una persona
  \item El número de años que le quedan a una persona
  \item El número de años que vive una población
\end{multicols}
\end{enumerate}
\vspace{0.15cm}
\item ¿Qué es la tasa de fecundidad?
\begin{enumerate}[label=\textbf{\Alph*)}]
\begin{multicols}{2}
  \item El número promedio de hijos por mujer en edad reproductiva
  \item El número promedio de hijos por familia
  \item El número promedio de hijos por hombre
  \item El número promedio de hijos por pareja
  \item El número promedio de hijos por comunidad
  \item El número promedio de hijos por país
\end{multicols}
\end{enumerate}
\vspace{0.15cm}
\item ¿Qué es la transición demográfica?
\begin{enumerate}[label=\textbf{\Alph*)}]
\begin{multicols}{2}
  \item El paso de altas tasas de natalidad y mortalidad a bajas tasas
  \item El paso de bajas tasas de natalidad y mortalidad a altas tasas
  \item La estabilidad de las tasas de natalidad y mortalidad
  \item El aumento de la tasa de natalidad
  \item El aumento de la tasa de mortalidad
  \item La disminución de la tasa de natalidad
\end{multicols}
\end{enumerate}
\vspace{0.15cm}
\item ¿Qué es la pirámide poblacional?
\begin{enumerate}[label=\textbf{\Alph*)}]
\begin{multicols}{2}
  \item Una representación gráfica de la distribución de la población por edad y sexo
  \item Una representación gráfica de la distribución de la población por ingresos
  \item Una representación gráfica de la distribución de la población por educación
  \item Una representación gráfica de la distribución de la población por religión
  \item Una representación gráfica de la distribución de la población por etnia
  \item Una representación gráfica de la distribución de la población por ocupación
\end{multicols}
\end{enumerate}
\vspace{0.15cm}
\item ¿Qué es la densidad de población?
\begin{enumerate}[label=\textbf{\Alph*)}]
\begin{multicols}{2}
  \item El número de habitantes por unidad de superficie
  \item El número de habitantes por familia
  \item El número de habitantes por vivienda
  \item El número de habitantes por ciudad
  \item El número de habitantes por región
  \item El número de habitantes por país
\end{multicols}
\end{enumerate}
\vspace{0.15cm}
\item ¿Qué es la población urbana?
\begin{enumerate}[label=\textbf{\Alph*)}]
\begin{multicols}{2}
  \item La población que vive en ciudades
  \item La población que vive en el campo
  \item La población que vive en pueblos
  \item La población que vive en aldeas
  \item La población que vive en caseríos
  \item La población que vive en comunidades rurales
\end{multicols}
\end{enumerate}
\vspace{0.15cm}
\item ¿Qué es la población rural?
\begin{enumerate}[label=\textbf{\Alph*)}]
\begin{multicols}{2}
  \item La población que vive en el campo
  \item La población que vive en ciudades
  \item La población que vive en metrópolis
  \item La población que vive en capitales
  \item La población que vive en zonas urbanas
  \item La población que vive en zonas metropolitanas
\end{multicols}
\end{enumerate}
\vspace{0.15cm}
\item ¿Qué es la urbanización?
\begin{enumerate}[label=\textbf{\Alph*)}]
\begin{multicols}{2}
  \item El proceso de concentración de la población en las ciudades
  \item El proceso de dispersión de la población en el campo
  \item El proceso de migración de la ciudad al campo
  \item El proceso de migración del campo a la ciudad
  \item El proceso de crecimiento de la población rural
  \item El proceso de disminución de la población urbana
\end{multicols}
\end{enumerate}
\vspace{0.15cm}
\item ¿Qué es la migración?
\begin{enumerate}[label=\textbf{\Alph*)}]
\begin{multicols}{2}
  \item El desplazamiento de personas de un lugar a otro
  \item La permanencia de personas en un lugar
  \item El crecimiento de la población
  \item La disminución de la población
  \item La distribución de la población
  \item La composición de la población
\end{multicols}
\end{enumerate}
\vspace{0.15cm}
\item ¿Qué es la emigración?
\begin{enumerate}[label=\textbf{\Alph*)}]
\begin{multicols}{2}
  \item La salida de personas de un lugar para establecerse en otro
  \item La llegada de personas a un lugar
  \item La permanencia de personas en un lugar
  \item El crecimiento de la población
  \item La disminución de la población
  \item La distribución de la población
\end{multicols}
\end{enumerate}
\vspace{0.15cm}
\item ¿Qué es la inmigración?
\begin{enumerate}[label=\textbf{\Alph*)}]
\begin{multicols}{2}
  \item La llegada de personas a un lugar para establecerse
  \item La salida de personas de un lugar
  \item La permanencia de personas en un lugar
  \item El crecimiento de la población
  \item La disminución de la población
  \item La distribución de la población
\end{multicols}
\end{enumerate}
\vspace{0.15cm}
\item ¿Qué es el saldo migratorio?
\begin{enumerate}[label=\textbf{\Alph*)}]
\begin{multicols}{2}
  \item La diferencia entre inmigrantes y emigrantes
  \item La suma de inmigrantes y emigrantes
  \item El producto de inmigrantes y emigrantes
  \item El cociente de inmigrantes y emigrantes
  \item El promedio de inmigrantes y emigrantes
  \item La tasa de inmigrantes y emigrantes
\end{multicols}
\end{enumerate}
\vspace{0.15cm}
\item ¿Qué es la migración interna?
\begin{enumerate}[label=\textbf{\Alph*)}]
\begin{multicols}{2}
  \item El desplazamiento de personas dentro de un mismo país
  \item El desplazamiento de personas entre países
  \item El desplazamiento de personas dentro de una ciudad
  \item El desplazamiento de personas dentro de una región
  \item El desplazamiento de personas dentro de una comunidad
  \item El desplazamiento de personas dentro de una familia
\end{multicols}
\end{enumerate}
\vspace{0.15cm}
\item ¿Qué es la migración internacional?
\begin{enumerate}[label=\textbf{\Alph*)}]
\begin{multicols}{2}
  \item El desplazamiento de personas de un país a otro
  \item El desplazamiento de personas dentro de un país
  \item El desplazamiento de personas dentro de una ciudad
  \item El desplazamiento de personas dentro de una región
  \item El desplazamiento de personas dentro de una comunidad
  \item El desplazamiento de personas dentro de una familia
\end{multicols}
\end{enumerate}
\vspace{0.15cm}
\item ¿Qué es la migración forzada?
\begin{enumerate}[label=\textbf{\Alph*)}]
\begin{multicols}{2}
  \item El desplazamiento de personas obligadas a abandonar su lugar de origen
  \item El desplazamiento de personas por voluntad propia
  \item El desplazamiento de personas por turismo
  \item El desplazamiento de personas por negocios
  \item El desplazamiento de personas por estudio
  \item El desplazamiento de personas por vacaciones
\end{multicols}
\end{enumerate}
\vspace{0.15cm}
\item ¿Qué es el desplazamiento interno?
\begin{enumerate}[label=\textbf{\Alph*)}]
\begin{multicols}{2}
  \item El desplazamiento de personas dentro de su propio país por conflictos o desastres
  \item El desplazamiento de personas fuera de su país
  \item El desplazamiento de personas por turismo
  \item El desplazamiento de personas por negocios
  \item El desplazamiento de personas por estudio
  \item El desplazamiento de personas por vacaciones
\end{multicols}
\end{enumerate}
\vspace{0.15cm}
\item ¿Qué es un refugiado?
\begin{enumerate}[label=\textbf{\Alph*)}]
\begin{multicols}{2}
  \item Una persona que huye de su país por persecución o conflicto
  \item Una persona que viaja por turismo
  \item Una persona que viaja por negocios
  \item Una persona que viaja por estudio
  \item Una persona que viaja por vacaciones
  \item Una persona que viaja por placer
\end{multicols}
\end{enumerate}
\vspace{0.15cm}
\item ¿Qué es un desplazado interno?
\begin{enumerate}[label=\textbf{\Alph*)}]
\begin{multicols}{2}
  \item Una persona que ha huido de su hogar pero permanece en su país
  \item Una persona que ha huido de su país
  \item Una persona que viaja por turismo
  \item Una persona que viaja por negocios
  \item Una persona que viaja por estudio
  \item Una persona que viaja por vacaciones
\end{multicols}
\end{enumerate}
\vspace{0.15cm}
\item ¿Qué es la trata de personas?
\begin{enumerate}[label=\textbf{\Alph*)}]
\begin{multicols}{2}
  \item El comercio ilegal de seres humanos con fines de explotación
  \item El comercio legal de seres humanos
  \item El comercio de bienes
  \item El comercio de servicios
  \item El comercio de tecnología
  \item El comercio de alimentos
\end{multicols}
\end{enumerate}
\vspace{0.15cm}
\item ¿Qué es el tráfico de personas?
\begin{enumerate}[label=\textbf{\Alph*)}]
\begin{multicols}{2}
  \item El traslado ilegal de personas a través de fronteras
  \item El traslado legal de personas
  \item El traslado de bienes
  \item El traslado de servicios
  \item El traslado de tecnología
  \item El traslado de alimentos
\end{multicols}
\end{enumerate}
\vspace{0.15cm}
\item ¿Qué es la diáspora?
\begin{enumerate}[label=\textbf{\Alph*)}]
\begin{multicols}{2}
  \item La dispersión de un pueblo o comunidad por el mundo
  \item La concentración de un pueblo en un lugar
  \item El aislamiento de un pueblo
  \item La separación de un pueblo
  \item La fragmentación de un pueblo
  \item La localización de un pueblo
\end{multicols}
\end{enumerate}
\vspace{0.15cm}
\item ¿Qué son las remesas?
\begin{enumerate}[label=\textbf{\Alph*)}]
\begin{multicols}{2}
  \item El dinero que los migrantes envían a sus familias en el país de origen
  \item El dinero que los gobiernos envían a otros países
  \item El dinero que las empresas envían a otros países
  \item El dinero que las ONG envían a otros países
  \item El dinero que los bancos envían a otros países
  \item El dinero que los individuos envían a otros individuos
\end{multicols}
\end{enumerate}
\vspace{0.15cm}
\item ¿Qué es la fuga de cerebros?
\begin{enumerate}[label=\textbf{\Alph*)}]
\begin{multicols}{2}
  \item La emigración de profesionales altamente calificados
  \item La inmigración de profesionales altamente calificados
  \item La emigración de trabajadores no calificados
  \item La inmigración de trabajadores no calificados
  \item La emigración de estudiantes
  \item La inmigración de estudiantes
\end{multicols}
\end{enumerate}
\vspace{0.15cm}
\item ¿Qué es la migración de retorno?
\begin{enumerate}[label=\textbf{\Alph*)}]
\begin{multicols}{2}
  \item El regreso de los migrantes a su país de origen
  \item La salida de los migrantes de su país de origen
  \item La llegada de los migrantes a un nuevo país
  \item La permanencia de los migrantes en un país
  \item La circulación de los migrantes entre países
  \item La dispersión de los migrantes por el mundo
\end{multicols}
\end{enumerate}
\vspace{0.15cm}
\item ¿Qué es la población económicamente activa?
\begin{enumerate}[label=\textbf{\Alph*)}]
\begin{multicols}{2}
  \item El conjunto de personas en edad de trabajar que trabajan o buscan trabajo
  \item El conjunto de personas que no trabajan
  \item El conjunto de personas jubiladas
  \item El conjunto de personas estudiantes
  \item El conjunto de personas menores de edad
  \item El conjunto de personas mayores de edad
\end{multicols}
\end{enumerate}
\vspace{0.15cm}
\item ¿Qué es la tasa de dependencia?
\begin{enumerate}[label=\textbf{\Alph*)}]
\begin{multicols}{2}
  \item La relación entre la población dependiente y la población activa
  \item La relación entre la población activa y la población dependiente
  \item La relación entre la población urbana y la rural
  \item La relación entre la población joven y la vieja
  \item La relación entre la población masculina y la femenina
  \item La relación entre la población nacional y la extranjera
\end{multicols}
\end{enumerate}
\vspace{0.15cm}
\item ¿Qué es el bono demográfico?
\begin{enumerate}[label=\textbf{\Alph*)}]
\begin{multicols}{2}
  \item El período en que la población activa es mayor que la dependiente
  \item El período en que la población dependiente es mayor que la activa
  \item El período en que la población urbana es mayor que la rural
  \item El período en que la población joven es mayor que la vieja
  \item El período en que la población masculina es mayor que la femenina
  \item El período en que la población nacional es mayor que la extranjera
\end{multicols}
\end{enumerate}
\vspace{0.15cm}
\item ¿Qué es el envejecimiento poblacional?
\begin{enumerate}[label=\textbf{\Alph*)}]
\begin{multicols}{2}
  \item El aumento de la proporción de personas mayores en la población
  \item La disminución de la proporción de personas mayores
  \item El aumento de la proporción de personas jóvenes
  \item La disminución de la proporción de personas jóvenes
  \item El aumento de la proporción de niños
  \item La disminución de la proporción de niños
\end{multicols}
\end{enumerate}
\vspace{0.15cm}
\item ¿Qué es la mortalidad materna?
\begin{enumerate}[label=\textbf{\Alph*)}]
\begin{multicols}{2}
  \item La muerte de mujeres durante el embarazo, parto o puerperio
  \item La muerte de mujeres por causas naturales
  \item La muerte de mujeres por accidentes
  \item La muerte de mujeres por enfermedades crónicas
  \item La muerte de mujeres por violencia
  \item La muerte de mujeres por suicidio
\end{multicols}
\end{enumerate}
\vspace{0.15cm}
\item ¿Qué es la planificación familiar?
\begin{enumerate}[label=\textbf{\Alph*)}]
\begin{multicols}{2}
  \item El conjunto de prácticas para decidir el número y espaciamiento de los hijos
  \item El conjunto de prácticas para tener más hijos
  \item El conjunto de prácticas para tener menos hijos
  \item El conjunto de prácticas para no tener hijos
  \item El conjunto de prácticas para adoptar hijos
  \item El conjunto de prácticas para educar hijos
\end{multicols}
\end{enumerate}
\vspace{0.15cm}
\item ¿Qué es la salud reproductiva?
\begin{enumerate}[label=\textbf{\Alph*)}]
\begin{multicols}{2}
  \item El estado de bienestar en relación con el sistema reproductivo
  \item El estado de bienestar en relación con el sistema digestivo
  \item El estado de bienestar en relación con el sistema respiratorio
  \item El estado de bienestar en relación con el sistema circulatorio
  \item El estado de bienestar en relación con el sistema nervioso
  \item El estado de bienestar en relación con el sistema óseo
\end{multicols}
\end{enumerate}
\vspace{0.15cm}
\item ¿Qué es el embarazo adolescente?
\begin{enumerate}[label=\textbf{\Alph*)}]
\begin{multicols}{2}
  \item El embarazo que ocurre en mujeres menores de 19 años
  \item El embarazo que ocurre en mujeres mayores de 19 años
  \item El embarazo que ocurre en mujeres mayores de 30 años
  \item El embarazo que ocurre en mujeres mayores de 40 años
  \item El embarazo que ocurre en mujeres mayores de 50 años
  \item El embarazo que ocurre en mujeres mayores de 60 años
\end{multicols}
\end{enumerate}
\vspace{0.15cm}
\item ¿Qué es la violencia de género?
\begin{enumerate}[label=\textbf{\Alph*)}]
\begin{multicols}{2}
  \item La violencia ejercida contra una persona por su género
  \item La violencia ejercida contra una persona por su edad
  \item La violencia ejercida contra una persona por su raza
  \item La violencia ejercida contra una persona por su religión
  \item La violencia ejercida contra una persona por su clase social
  \item La violencia ejercida contra una persona por su nacionalidad
\end{multicols}
\end{enumerate}
\vspace{0.15cm}
\item ¿Qué es la equidad de género?
\begin{enumerate}[label=\textbf{\Alph*)}]
\begin{multicols}{2}
  \item La igualdad de derechos y oportunidades entre hombres y mujeres
  \item La superioridad de un género sobre otro
  \item La inferioridad de un género sobre otro
  \item La diferencia de derechos entre géneros
  \item La diferencia de oportunidades entre géneros
  \item La separación de géneros
\end{multicols}
\end{enumerate}
\vspace{0.15cm}
\end{enumerate}\n
\subsubsection*{Sección V: Medio Ambiente y Recursos (161–200)}

\begin{enumerate}[label=\textbf{\arabic*.}]
\setcounter{enumi}{160}
\item ¿Qué es el medio ambiente?
\begin{enumerate}[label=\textbf{\Alph*)}]
\begin{multicols}{2}
  \item El conjunto de elementos naturales y sociales que rodean al ser humano
  \item El conjunto de elementos naturales
  \item El conjunto de elementos sociales
  \item El conjunto de elementos económicos
  \item El conjunto de elementos políticos
  \item El conjunto de elementos culturales
\end{multicols}
\end{enumerate}
\vspace{0.15cm}
\item ¿Qué es la contaminación ambiental?
\begin{enumerate}[label=\textbf{\Alph*)}]
\begin{multicols}{2}
  \item La introducción de sustancias nocivas en el medio ambiente
  \item La introducción de sustancias beneficiosas
  \item La introducción de sustancias neutras
  \item La introducción de sustancias naturales
  \item La introducción de sustancias artificiales
  \item La introducción de sustancias orgánicas
\end{multicols}
\end{enumerate}
\vspace{0.15cm}
\item ¿Qué es el cambio climático?
\begin{enumerate}[label=\textbf{\Alph*)}]
\begin{multicols}{2}
  \item La alteración del clima global por causas naturales y humanas
  \item La alteración del clima local
  \item La alteración del clima regional
  \item La alteración del clima nacional
  \item La alteración del clima continental
  \item La alteración del clima oceánico
\end{multicols}
\end{enumerate}
\vspace{0.15cm}
\item ¿Qué es el efecto invernadero?
\begin{enumerate}[label=\textbf{\Alph*)}]
\begin{multicols}{2}
  \item El fenómeno por el cual ciertos gases retienen el calor en la atmósfera
  \item El fenómeno por el cual ciertos gases liberan el calor
  \item El fenómeno por el cual ciertos gases reflejan el calor
  \item El fenómeno por el cual ciertos gases absorben el calor
  \item El fenómeno por el cual ciertos gases emiten el calor
  \item El fenómeno por el cual ciertos gases transmiten el calor
\end{multicols}
\end{enumerate}
\vspace{0.15cm}
\item ¿Qué es el calentamiento global?
\begin{enumerate}[label=\textbf{\Alph*)}]
\begin{multicols}{2}
  \item El aumento de la temperatura promedio de la Tierra
  \item La disminución de la temperatura promedio de la Tierra
  \item La estabilidad de la temperatura promedio de la Tierra
  \item La variación de la temperatura promedio de la Tierra
  \item La fluctuación de la temperatura promedio de la Tierra
  \item La oscilación de la temperatura promedio de la Tierra
\end{multicols}
\end{enumerate}
\vspace{0.15cm}
\item ¿Qué es la deforestación?
\begin{enumerate}[label=\textbf{\Alph*)}]
\begin{multicols}{2}
  \item La pérdida de bosques por causas naturales y humanas
  \item La ganancia de bosques
  \item La preservación de bosques
  \item La reforestación
  \item La forestación
  \item La conservación de bosques
\end{multicols}
\end{enumerate}
\vspace{0.15cm}
\item ¿Qué es la desertificación?
\begin{enumerate}[label=\textbf{\Alph*)}]
\begin{multicols}{2}
  \item La degradación de tierras fértiles en desiertos
  \item La transformación de desiertos en tierras fértiles
  \item La preservación de tierras fértiles
  \item La conservación de desiertos
  \item La reforestación de desiertos
  \item La forestación de desiertos
\end{multicols}
\end{enumerate}
\vspace{0.15cm}
\item ¿Qué es la pérdida de biodiversidad?
\begin{enumerate}[label=\textbf{\Alph*)}]
\begin{multicols}{2}
  \item La disminución de la variedad de especies y ecosistemas
  \item El aumento de la variedad de especies
  \item La preservación de la variedad de especies
  \item La conservación de la variedad de especies
  \item La restauración de la variedad de especies
  \item La protección de la variedad de especies
\end{multicols}
\end{enumerate}
\vspace{0.15cm}
\item ¿Qué es la sobreexplotación de recursos?
\begin{enumerate}[label=\textbf{\Alph*)}]
\begin{multicols}{2}
  \item El uso excesivo de recursos naturales más allá de su capacidad de renovación
  \item El uso moderado de recursos naturales
  \item El uso eficiente de recursos naturales
  \item El uso sostenible de recursos naturales
  \item El uso racional de recursos naturales
  \item El uso responsable de recursos naturales
\end{multicols}
\end{enumerate}
\vspace{0.15cm}
\item ¿Qué es la huella de carbono?
\begin{enumerate}[label=\textbf{\Alph*)}]
\begin{multicols}{2}
  \item La medida de las emisiones de gases de efecto invernadero
  \item La medida de las emisiones de gases contaminantes
  \item La medida de las emisiones de gases nobles
  \item La medida de las emisiones de gases inertes
  \item La medida de las emisiones de gases tóxicos
  \item La medida de las emisiones de gases radiactivos
\end{multicols}
\end{enumerate}
\vspace{0.15cm}
\item ¿Qué es la energía renovable?
\begin{enumerate}[label=\textbf{\Alph*)}]
\begin{multicols}{2}
  \item La energía que proviene de fuentes naturales inagotables
  \item La energía que proviene de fuentes fósiles
  \item La energía que proviene de fuentes nucleares
  \item La energía que proviene de fuentes minerales
  \item La energía que proviene de fuentes metálicas
  \item La energía que proviene de fuentes sintéticas
\end{multicols}
\end{enumerate}
\vspace{0.15cm}
\item ¿Cuál es una fuente de energía renovable?
\begin{enumerate}[label=\textbf{\Alph*)}]
\begin{multicols}{2}
  \item La solar
  \item El petróleo
  \item El carbón
  \item El gas natural
  \item La gasolina
  \item El diésel
\end{multicols}
\end{enumerate}
\vspace{0.15cm}
\item ¿Cuál es una fuente de energía renovable?
\begin{enumerate}[label=\textbf{\Alph*)}]
\begin{multicols}{2}
  \item La eólica
  \item El petróleo
  \item El carbón
  \item El gas natural
  \item La gasolina
  \item El diésel
\end{multicols}
\end{enumerate}
\vspace{0.15cm}
\item ¿Cuál es una fuente de energía renovable?
\begin{enumerate}[label=\textbf{\Alph*)}]
\begin{multicols}{2}
  \item La hidroeléctrica
  \item El petróleo
  \item El carbón
  \item El gas natural
  \item La gasolina
  \item El diésel
\end{multicols}
\end{enumerate}
\vspace{0.15cm}
\item ¿Cuál es una fuente de energía renovable?
\begin{enumerate}[label=\textbf{\Alph*)}]
\begin{multicols}{2}
  \item La geotérmica
  \item El petróleo
  \item El carbón
  \item El gas natural
  \item La gasolina
  \item El diésel
\end{multicols}
\end{enumerate}
\vspace{0.15cm}
\item ¿Cuál es una fuente de energía renovable?
\begin{enumerate}[label=\textbf{\Alph*)}]
\begin{multicols}{2}
  \item La biomasa
  \item El petróleo
  \item El carbón
  \item El gas natural
  \item La gasolina
  \item El diésel
\end{multicols}
\end{enumerate}
\vspace{0.15cm}
\item ¿Qué es la energía no renovable?
\begin{enumerate}[label=\textbf{\Alph*)}]
\begin{multicols}{2}
  \item La energía que proviene de fuentes finitas
  \item La energía que proviene de fuentes inagotables
  \item La energía que proviene de fuentes naturales
  \item La energía que proviene de fuentes limpias
  \item La energía que proviene de fuentes verdes
  \item La energía que proviene de fuentes sostenibles
\end{multicols}
\end{enumerate}
\vspace{0.15cm}
\item ¿Cuál es una fuente de energía no renovable?
\begin{enumerate}[label=\textbf{\Alph*)}]
\begin{multicols}{2}
  \item El petróleo
  \item La solar
  \item La eólica
  \item La hidroeléctrica
  \item La geotérmica
  \item La biomasa
\end{multicols}
\end{enumerate}
\vspace{0.15cm}
\item ¿Cuál es una fuente de energía no renovable?
\begin{enumerate}[label=\textbf{\Alph*)}]
\begin{multicols}{2}
  \item El carbón
  \item La solar
  \item La eólica
  \item La hidroeléctrica
  \item La geotérmica
  \item La biomasa
\end{multicols}
\end{enumerate}
\vspace{0.15cm}
\item ¿Cuál es una fuente de energía no renovable?
\begin{enumerate}[label=\textbf{\Alph*)}]
\begin{multicols}{2}
  \item El gas natural
  \item La solar
  \item La eólica
  \item La hidroeléctrica
  \item La geotérmica
  \item La biomasa
\end{multicols}
\end{enumerate}
\vspace{0.15cm}
\item ¿Qué es el desarrollo sostenible?
\begin{enumerate}[label=\textbf{\Alph*)}]
\begin{multicols}{2}
  \item El desarrollo que satisface las necesidades actuales sin comprometer las futuras
  \item El desarrollo que satisface solo las necesidades actuales
  \item El desarrollo que satisface solo las necesidades futuras
  \item El desarrollo que no satisface necesidades
  \item El desarrollo que ignora las necesidades
  \item El desarrollo que prioriza lo económico
\end{multicols}
\end{enumerate}
\vspace{0.15cm}
\item ¿Qué es la economía circular?
\begin{enumerate}[label=\textbf{\Alph*)}]
\begin{multicols}{2}
  \item Un modelo económico que busca reducir, reutilizar y reciclar
  \item Un modelo económico lineal
  \item Un modelo económico extractivista
  \item Un modelo económico consumista
  \item Un modelo económico depredador
  \item Un modelo económico insostenible
\end{multicols}
\end{enumerate}
\vspace{0.15cm}
\item ¿Qué es el reciclaje?
\begin{enumerate}[label=\textbf{\Alph*)}]
\begin{multicols}{2}
  \item El proceso de transformar residuos en nuevos productos
  \item El proceso de eliminar residuos
  \item El proceso de acumular residuos
  \item El proceso de quemar residuos
  \item El proceso de enterrar residuos
  \item El proceso de arrojar residuos
\end{multicols}
\end{enumerate}
\vspace{0.15cm}
\item ¿Qué es la gestión de residuos?
\begin{enumerate}[label=\textbf{\Alph*)}]
\begin{multicols}{2}
  \item El conjunto de acciones para manejar adecuadamente los residuos
  \item El conjunto de acciones para eliminar residuos
  \item El conjunto de acciones para acumular residuos
  \item El conjunto de acciones para quemar residuos
  \item El conjunto de acciones para enterrar residuos
  \item El conjunto de acciones para arrojar residuos
\end{multicols}
\end{enumerate}
\vspace{0.15cm}
\item ¿Qué es la contaminación del agua?
\begin{enumerate}[label=\textbf{\Alph*)}]
\begin{multicols}{2}
  \item La introducción de sustancias nocivas en cuerpos de agua
  \item La introducción de sustancias beneficiosas
  \item La introducción de sustancias neutras
  \item La introducción de sustancias naturales
  \item La introducción de sustancias artificiales
  \item La introducción de sustancias orgánicas
\end{multicols}
\end{enumerate}
\vspace{0.15cm}
\item ¿Qué es la contaminación del aire?
\begin{enumerate}[label=\textbf{\Alph*)}]
\begin{multicols}{2}
  \item La introducción de sustancias nocivas en la atmósfera
  \item La introducción de sustancias beneficiosas
  \item La introducción de sustancias neutras
  \item La introducción de sustancias naturales
  \item La introducción de sustancias artificiales
  \item La introducción de sustancias orgánicas
\end{multicols}
\end{enumerate}
\vspace{0.15cm}
\item ¿Qué es la contaminación del suelo?
\begin{enumerate}[label=\textbf{\Alph*)}]
\begin{multicols}{2}
  \item La introducción de sustancias nocivas en el suelo
  \item La introducción de sustancias beneficiosas
  \item La introducción de sustancias neutras
  \item La introducción de sustancias naturales
  \item La introducción de sustancias artificiales
  \item La introducción de sustancias orgánicas
\end{multicols}
\end{enumerate}
\vspace{0.15cm}
\item ¿Qué es la lluvia ácida?
\begin{enumerate}[label=\textbf{\Alph*)}]
\begin{multicols}{2}
  \item La precipitación con altos niveles de ácidos por contaminación
  \item La precipitación con altos niveles de bases
  \item La precipitación con altos niveles de sales
  \item La precipitación con altos niveles de minerales
  \item La precipitación con altos niveles de nutrientes
  \item La precipitación con altos niveles de oxígeno
\end{multicols}
\end{enumerate}
\vspace{0.15cm}
\item ¿Qué es la capa de ozono?
\begin{enumerate}[label=\textbf{\Alph*)}]
\begin{multicols}{2}
  \item Una capa de la atmósfera que protege de la radiación ultravioleta
  \item Una capa de la atmósfera que protege de la radiación infrarroja
  \item Una capa de la atmósfera que protege de la radiación visible
  \item Una capa de la atmósfera que protege de la radiación gamma
  \item Una capa de la atmósfera que protege de la radiación X
  \item Una capa de la atmósfera que protege de la radiación cósmica
\end{multicols}
\end{enumerate}
\vspace{0.15cm}
\item ¿Qué es el agujero de la capa de ozono?
\begin{enumerate}[label=\textbf{\Alph*)}]
\begin{multicols}{2}
  \item La disminución de la capa de ozono por contaminantes
  \item El aumento de la capa de ozono
  \item La estabilidad de la capa de ozono
  \item La variación de la capa de ozono
  \item La fluctuación de la capa de ozono
  \item La oscilación de la capa de ozono
\end{multicols}
\end{enumerate}
\vspace{0.15cm}
\item ¿Qué es un desastre natural?
\begin{enumerate}[label=\textbf{\Alph*)}]
\begin{multicols}{2}
  \item Un fenómeno natural que causa daños y víctimas
  \item Un fenómeno natural beneficioso
  \item Un fenómeno natural neutro
  \item Un fenómeno natural controlado
  \item Un fenómeno natural predecible
  \item Un fenómeno natural evitable
\end{multicols}
\end{enumerate}
\vspace{0.15cm}
\item ¿Cuál es un desastre natural?
\begin{enumerate}[label=\textbf{\Alph*)}]
\begin{multicols}{2}
  \item Un terremoto
  \item Una lluvia
  \item Un viento
  \item Una nube
  \item Un amanecer
  \item Un atardecer
\end{multicols}
\end{enumerate}
\vspace{0.15cm}
\item ¿Cuál es un desastre natural?
\begin{enumerate}[label=\textbf{\Alph*)}]
\begin{multicols}{2}
  \item Un huracán
  \item Una lluvia
  \item Un viento
  \item Una nube
  \item Un amanecer
  \item Un atardecer
\end{multicols}
\end{enumerate}
\vspace{0.15cm}
\item ¿Cuál es un desastre natural?
\begin{enumerate}[label=\textbf{\Alph*)}]
\begin{multicols}{2}
  \item Una inundación
  \item Una lluvia
  \item Un viento
  \item Una nube
  \item Un amanecer
  \item Un atardecer
\end{multicols}
\end{enumerate}
\vspace{0.15cm}
\item ¿Cuál es un desastre natural?
\begin{enumerate}[label=\textbf{\Alph*)}]
\begin{multicols}{2}
  \item Una sequía
  \item Una lluvia
  \item Un viento
  \item Una nube
  \item Un amanecer
  \item Un atardecer
\end{multicols}
\end{enumerate}
\vspace{0.15cm}
\item ¿Cuál es un desastre natural?
\begin{enumerate}[label=\textbf{\Alph*)}]
\begin{multicols}{2}
  \item Un deslizamiento
  \item Una lluvia
  \item Un viento
  \item Una nube
  \item Un amanecer
  \item Un atardecer
\end{multicols}
\end{enumerate}
\vspace{0.15cm}
\item ¿Qué es la gestión de riesgos?
\begin{enumerate}[label=\textbf{\Alph*)}]
\begin{multicols}{2}
  \item El conjunto de acciones para prevenir y mitigar desastres
  \item El conjunto de acciones para ignorar desastres
  \item El conjunto de acciones para aumentar desastres
  \item El conjunto de acciones para provocar desastres
  \item El conjunto de acciones para negar desastres
  \item El conjunto de acciones para olvidar desastres
\end{multicols}
\end{enumerate}
\vspace{0.15cm}
\item ¿Qué es la adaptación al cambio climático?
\begin{enumerate}[label=\textbf{\Alph*)}]
\begin{multicols}{2}
  \item El conjunto de acciones para reducir la vulnerabilidad al cambio climático
  \item El conjunto de acciones para aumentar la vulnerabilidad
  \item El conjunto de acciones para ignorar el cambio climático
  \item El conjunto de acciones para negar el cambio climático
  \item El conjunto de acciones para provocar el cambio climático
  \item El conjunto de acciones para olvidar el cambio climático
\end{multicols}
\end{enumerate}
\vspace{0.15cm}
\item ¿Qué es la mitigación del cambio climático?
\begin{enumerate}[label=\textbf{\Alph*)}]
\begin{multicols}{2}
  \item El conjunto de acciones para reducir las emisiones de gases de efecto invernadero
  \item El conjunto de acciones para aumentar las emisiones
  \item El conjunto de acciones para ignorar las emisiones
  \item El conjunto de acciones para negar las emisiones
  \item El conjunto de acciones para provocar las emisiones
  \item El conjunto de acciones para olvidar las emisiones
\end{multicols}
\end{enumerate}
\vspace{0.15cm}
\item ¿Qué es la justicia ambiental?
\begin{enumerate}[label=\textbf{\Alph*)}]
\begin{multicols}{2}
  \item La distribución equitativa de los beneficios y cargas ambientales
  \item La distribución inequitativa de los beneficios ambientales
  \item La concentración de los beneficios ambientales
  \item La concentración de las cargas ambientales
  \item La exclusión de los beneficios ambientales
  \item La exclusión de las cargas ambientales
\end{multicols}
\end{enumerate}
\vspace{0.15cm}
\end{enumerate}\n
\subsubsection*{Sección VI: Problemas Sociales Contemporáneos (201–300)}

\begin{enumerate}[label=\textbf{\arabic*.}]
\setcounter{enumi}{200}
\item ¿Qué es la violencia social?
\begin{enumerate}[label=\textbf{\Alph*)}]
\begin{multicols}{2}
  \item El uso de la fuerza para causar daño en la sociedad
  \item La convivencia pacífica
  \item La resolución pacífica de conflictos
  \item La mediación de conflictos
  \item La negociación de conflictos
  \item La reconciliación social
\end{multicols}
\end{enumerate}
\vspace{0.15cm}
\item ¿Qué es la delincuencia?
\begin{enumerate}[label=\textbf{\Alph*)}]
\begin{multicols}{2}
  \item El conjunto de actos ilegales cometidos por personas
  \item El conjunto de actos legales
  \item El conjunto de actos morales
  \item El conjunto de actos éticos
  \item El conjunto de actos solidarios
  \item El conjunto de actos cooperativos
\end{multicols}
\end{enumerate}
\vspace{0.15cm}
\item ¿Qué es una pandilla?
\begin{enumerate}[label=\textbf{\Alph*)}]
\begin{multicols}{2}
  \item Un grupo organizado que comete actos delictivos
  \item Un grupo organizado que comete actos benéficos
  \item Un grupo organizado que comete actos culturales
  \item Un grupo organizado que comete actos deportivos
  \item Un grupo organizado que comete actos religiosos
  \item Un grupo organizado que comete actos políticos
\end{multicols}
\end{enumerate}
\vspace{0.15cm}
\item ¿Qué es el narcotráfico?
\begin{enumerate}[label=\textbf{\Alph*)}]
\begin{multicols}{2}
  \item El comercio ilegal de drogas
  \item El comercio legal de drogas
  \item El comercio de medicamentos
  \item El comercio de alimentos
  \item El comercio de tecnología
  \item El comercio de textiles
\end{multicols}
\end{enumerate}
\vspace{0.15cm}
\item ¿Qué es la corrupción?
\begin{enumerate}[label=\textbf{\Alph*)}]
\begin{multicols}{2}
  \item El abuso del poder público para obtener beneficios privados
  \item El uso adecuado del poder público
  \item El uso transparente del poder público
  \item El uso honesto del poder público
  \item El uso ético del poder público
  \item El uso responsable del poder público
\end{multicols}
\end{enumerate}
\vspace{0.15cm}
\item ¿Qué es la impunidad?
\begin{enumerate}[label=\textbf{\Alph*)}]
\begin{multicols}{2}
  \item La falta de castigo para los delitos cometidos
  \item El castigo para los delitos cometidos
  \item La prevención de los delitos
  \item La investigación de los delitos
  \item La sanción de los delitos
  \item La reparación de los delitos
\end{multicols}
\end{enumerate}
\vspace{0.15cm}
\item ¿Qué es la trata de personas?
\begin{enumerate}[label=\textbf{\Alph*)}]
\begin{multicols}{2}
  \item El comercio ilegal de seres humanos con fines de explotación
  \item El comercio legal de seres humanos
  \item El comercio de bienes
  \item El comercio de servicios
  \item El comercio de tecnología
  \item El comercio de alimentos
\end{multicols}
\end{enumerate}
\vspace{0.15cm}
\item ¿Qué es la explotación laboral?
\begin{enumerate}[label=\textbf{\Alph*)}]
\begin{multicols}{2}
  \item El abuso en las condiciones de trabajo
  \item El respeto a las condiciones de trabajo
  \item La mejora de las condiciones de trabajo
  \item La protección de las condiciones de trabajo
  \item La garantía de las condiciones de trabajo
  \item La promoción de las condiciones de trabajo
\end{multicols}
\end{enumerate}
\vspace{0.15cm}
\item ¿Qué es el trabajo infantil?
\begin{enumerate}[label=\textbf{\Alph*)}]
\begin{multicols}{2}
  \item La explotación de niños en actividades laborales
  \item La educación de niños
  \item La recreación de niños
  \item La protección de niños
  \item La alimentación de niños
  \item La salud de niños
\end{multicols}
\end{enumerate}
\vspace{0.15cm}
\item ¿Qué es la desigualdad económica?
\begin{enumerate}[label=\textbf{\Alph*)}]
\begin{multicols}{2}
  \item La distribución inequitativa de la riqueza
  \item La distribución equitativa de la riqueza
  \item La concentración de la riqueza
  \item La dispersión de la riqueza
  \item La redistribución de la riqueza
  \item La generación de riqueza
\end{multicols}
\end{enumerate}
\vspace{0.15cm}
\item ¿Qué es la pobreza multidimensional?
\begin{enumerate}[label=\textbf{\Alph*)}]
\begin{multicols}{2}
  \item La carencia en múltiples dimensiones del bienestar
  \item La carencia en una sola dimensión
  \item La abundancia en múltiples dimensiones
  \item La abundancia en una sola dimensión
  \item La distribución equitativa del bienestar
  \item La distribución inequitativa del bienestar
\end{multicols}
\end{enumerate}
\vspace{0.15cm}
\item ¿Qué es la exclusión social?
\begin{enumerate}[label=\textbf{\Alph*)}]
\begin{multicols}{2}
  \item La marginación de individuos o grupos de la participación social
  \item La integración de individuos o grupos
  \item La inclusión de individuos o grupos
  \item La participación de individuos o grupos
  \item La incorporación de individuos o grupos
  \item La aceptación de individuos o grupos
\end{multicols}
\end{enumerate}
\vspace{0.15cm}
\item ¿Qué es la discriminación?
\begin{enumerate}[label=\textbf{\Alph*)}]
\begin{multicols}{2}
  \item El trato desigual e injusto hacia un grupo social
  \item El trato igualitario hacia un grupo social
  \item El trato justo hacia un grupo social
  \item El trato objetivo hacia un grupo social
  \item El trato científico hacia un grupo social
  \item El trato artístico hacia un grupo social
\end{multicols}
\end{enumerate}
\vspace{0.15cm}
\item ¿Qué es el racismo?
\begin{enumerate}[label=\textbf{\Alph*)}]
\begin{multicols}{2}
  \item La discriminación basada en la raza o etnia
  \item La discriminación basada en el género
  \item La discriminación basada en la edad
  \item La discriminación basada en la religión
  \item La discriminación basada en la clase social
  \item La discriminación basada en la nacionalidad
\end{multicols}
\end{enumerate}
\vspace{0.15cm}
\item ¿Qué es el sexismo?
\begin{enumerate}[label=\textbf{\Alph*)}]
\begin{multicols}{2}
  \item La discriminación basada en el sexo o género
  \item La discriminación basada en la raza
  \item La discriminación basada en la edad
  \item La discriminación basada en la religión
  \item La discriminación basada en la clase social
  \item La discriminación basada en la nacionalidad
\end{multicols}
\end{enumerate}
\vspace{0.15cm}
\item ¿Qué es la xenofobia?
\begin{enumerate}[label=\textbf{\Alph*)}]
\begin{multicols}{2}
  \item El rechazo hacia los extranjeros
  \item El rechazo hacia los nacionales
  \item El rechazo hacia los ancianos
  \item El rechazo hacia los niños
  \item El rechazo hacia los discapacitados
  \item El rechazo hacia los pobres
\end{multicols}
\end{enumerate}
\vspace{0.15cm}
\item ¿Qué es la homofobia?
\begin{enumerate}[label=\textbf{\Alph*)}]
\begin{multicols}{2}
  \item El rechazo hacia las personas homosexuales
  \item El rechazo hacia las personas heterosexuales
  \item El rechazo hacia las personas bisexuales
  \item El rechazo hacia las personas transgénero
  \item El rechazo hacia las personas intersexuales
  \item El rechazo hacia las personas asexuales
\end{multicols}
\end{enumerate}
\vspace{0.15cm}
\item ¿Qué es la violencia de género?
\begin{enumerate}[label=\textbf{\Alph*)}]
\begin{multicols}{2}
  \item La violencia ejercida contra una persona por su género
  \item La violencia ejercida contra una persona por su edad
  \item La violencia ejercida contra una persona por su raza
  \item La violencia ejercida contra una persona por su religión
  \item La violencia ejercida contra una persona por su clase social
  \item La violencia ejercida contra una persona por su nacionalidad
\end{multicols}
\end{enumerate}
\vspace{0.15cm}
\item ¿Qué es el feminicidio?
\begin{enumerate}[label=\textbf{\Alph*)}]
\begin{multicols}{2}
  \item El asesinato de mujeres por razones de género
  \item El asesinato de hombres por razones de género
  \item El asesinato de niños por razones de género
  \item El asesinato de ancianos por razones de género
  \item El asesinato de extranjeros por razones de género
  \item El asesinato de pobres por razones de género
\end{multicols}
\end{enumerate}
\vspace{0.15cm}
\item ¿Qué es la trata de personas?
\begin{enumerate}[label=\textbf{\Alph*)}]
\begin{multicols}{2}
  \item El comercio ilegal de seres humanos con fines de explotación
  \item El comercio legal de seres humanos
  \item El comercio de bienes
  \item El comercio de servicios
  \item El comercio de tecnología
  \item El comercio de alimentos
\end{multicols}
\end{enumerate}
\vspace{0.15cm}
\item ¿Qué es la migración irregular?
\begin{enumerate}[label=\textbf{\Alph*)}]
\begin{multicols}{2}
  \item El desplazamiento de personas sin documentación legal
  \item El desplazamiento de personas con documentación legal
  \item El desplazamiento de personas por turismo
  \item El desplazamiento de personas por negocios
  \item El desplazamiento de personas por estudio
  \item El desplazamiento de personas por vacaciones
\end{multicols}
\end{enumerate}
\vspace{0.15cm}
\item ¿Qué es la deportación?
\begin{enumerate}[label=\textbf{\Alph*)}]
\begin{multicols}{2}
  \item La expulsión de un extranjero de un país
  \item La admisión de un extranjero en un país
  \item La naturalización de un extranjero
  \item La residencia de un extranjero
  \item La nacionalización de un extranjero
  \item La integración de un extranjero
\end{multicols}
\end{enumerate}
\vspace{0.15cm}
\item ¿Qué es el asilo político?
\begin{enumerate}[label=\textbf{\Alph*)}]
\begin{multicols}{2}
  \item La protección que un país brinda a una persona perseguida por razones políticas
  \item La protección que un país brinda a una persona por razones económicas
  \item La protección que un país brinda a una persona por razones sociales
  \item La protección que un país brinda a una persona por razones culturales
  \item La protección que un país brinda a una persona por razones religiosas
  \item La protección que un país brinda a una persona por razones ambientales
\end{multicols}
\end{enumerate}
\vspace{0.15cm}
\item ¿Qué es la reunificación familiar?
\begin{enumerate}[label=\textbf{\Alph*)}]
\begin{multicols}{2}
  \item El proceso de reunir a los miembros de una familia separada por la migración
  \item El proceso de separar a los miembros de una familia
  \item El proceso de dispersar a los miembros de una familia
  \item El proceso de dividir a los miembros de una familia
  \item El proceso de aislar a los miembros de una familia
  \item El proceso de excluir a los miembros de una familia
\end{multicols}
\end{enumerate}
\vspace{0.15cm}
\item ¿Qué es la integración de los migrantes?
\begin{enumerate}[label=\textbf{\Alph*)}]
\begin{multicols}{2}
  \item El proceso de incorporación de los migrantes a la sociedad de acogida
  \item El proceso de exclusión de los migrantes
  \item El proceso de marginación de los migrantes
  \item El proceso de discriminación de los migrantes
  \item El proceso de segregación de los migrantes
  \item El proceso de aislamiento de los migrantes
\end{multicols}
\end{enumerate}
\vspace{0.15cm}
\item ¿Qué es la fuga de cerebros?
\begin{enumerate}[label=\textbf{\Alph*)}]
\begin{multicols}{2}
  \item La emigración de profesionales altamente calificados
  \item La inmigración de profesionales altamente calificados
  \item La emigración de trabajadores no calificados
  \item La inmigración de trabajadores no calificados
  \item La emigración de estudiantes
  \item La inmigración de estudiantes
\end{multicols}
\end{enumerate}
\vspace{0.15cm}
\item ¿Qué son las remesas?
\begin{enumerate}[label=\textbf{\Alph*)}]
\begin{multicols}{2}
  \item El dinero que los migrantes envían a sus familias en el país de origen
  \item El dinero que los gobiernos envían a otros países
  \item El dinero que las empresas envían a otros países
  \item El dinero que las ONG envían a otros países
  \item El dinero que los bancos envían a otros países
  \item El dinero que los individuos envían a otros individuos
\end{multicols}
\end{enumerate}
\vspace{0.15cm}
\item ¿Qué es la diáspora?
\begin{enumerate}[label=\textbf{\Alph*)}]
\begin{multicols}{2}
  \item La dispersión de un pueblo o comunidad por el mundo
  \item La concentración de un pueblo en un lugar
  \item El aislamiento de un pueblo
  \item La separación de un pueblo
  \item La fragmentación de un pueblo
  \item La localización de un pueblo
\end{multicols}
\end{enumerate}
\vspace{0.15cm}
\item ¿Qué es el transnacionalismo?
\begin{enumerate}[label=\textbf{\Alph*)}]
\begin{multicols}{2}
  \item La creación de vínculos sociales, culturales y económicos entre países
  \item La separación de los países
  \item El aislamiento de los países
  \item La fragmentación de los países
  \item La localización de los países
  \item La regionalización de los países
\end{multicols}
\end{enumerate}
\vspace{0.15cm}
\item ¿Qué es la ciudadanía global?
\begin{enumerate}[label=\textbf{\Alph*)}]
\begin{multicols}{2}
  \item La conciencia de pertenencia a una comunidad mundial
  \item La conciencia de pertenencia a una nación
  \item La conciencia de pertenencia a una región
  \item La conciencia de pertenencia a una ciudad
  \item La conciencia de pertenencia a una comunidad
  \item La conciencia de pertenencia a una familia
\end{multicols}
\end{enumerate}
\vspace{0.15cm}
\item ¿Qué es la sociedad del conocimiento?
\begin{enumerate}[label=\textbf{\Alph*)}]
\begin{multicols}{2}
  \item Una sociedad en la que el conocimiento es el principal recurso
  \item Una sociedad en la que la agricultura es el principal recurso
  \item Una sociedad en la que la industria es el principal recurso
  \item Una sociedad en la que el comercio es el principal recurso
  \item Una sociedad en la que la minería es el principal recurso
  \item Una sociedad en la que la pesca es el principal recurso
\end{multicols}
\end{enumerate}
\vspace{0.15cm}
\item ¿Qué es la brecha digital?
\begin{enumerate}[label=\textbf{\Alph*)}]
\begin{multicols}{2}
  \item La desigualdad en el acceso y uso de las tecnologías de la información
  \item La igualdad en el acceso a la tecnología
  \item La abundancia de tecnología
  \item La escasez de tecnología
  \item La distribución equitativa de tecnología
  \item La democratización de la tecnología
\end{multicols}
\end{enumerate}
\vspace{0.15cm}
\item ¿Qué es la inteligencia artificial?
\begin{enumerate}[label=\textbf{\Alph*)}]
\begin{multicols}{2}
  \item La capacidad de las máquinas para realizar tareas que requieren inteligencia humana
  \item La capacidad de las máquinas para realizar tareas mecánicas
  \item La capacidad de las máquinas para realizar tareas repetitivas
  \item La capacidad de las máquinas para realizar tareas simples
  \item La capacidad de las máquinas para realizar tareas básicas
  \item La capacidad de las máquinas para realizar tareas rutinarias
\end{multicols}
\end{enumerate}
\vspace{0.15cm}
\item ¿Qué es la robótica?
\begin{enumerate}[label=\textbf{\Alph*)}]
\begin{multicols}{2}
  \item La ciencia que estudia el diseño y uso de robots
  \item La ciencia que estudia los animales
  \item La ciencia que estudia las plantas
  \item La ciencia que estudia los minerales
  \item La ciencia que estudia los astros
  \item La ciencia que estudia los fenómenos físicos
\end{multicols}
\end{enumerate}
\vspace{0.15cm}
\item ¿Qué es la automatización?
\begin{enumerate}[label=\textbf{\Alph*)}]
\begin{multicols}{2}
  \item El uso de tecnología para realizar tareas sin intervención humana
  \item El uso de tecnología para realizar tareas con intervención humana
  \item El uso de tecnología para realizar tareas manuales
  \item El uso de tecnología para realizar tareas artesanales
  \item El uso de tecnología para realizar tareas tradicionales
  \item El uso de tecnología para realizar tareas rudimentarias
\end{multicols}
\end{enumerate}
\vspace{0.15cm}
\item ¿Qué es el teletrabajo?
\begin{enumerate}[label=\textbf{\Alph*)}]
\begin{multicols}{2}
  \item El trabajo realizado a distancia mediante tecnologías de la información
  \item El trabajo realizado en una oficina
  \item El trabajo realizado en una fábrica
  \item El trabajo realizado en un campo
  \item El trabajo realizado en una tienda
  \item El trabajo realizado en un taller
\end{multicols}
\end{enumerate}
\vspace{0.15cm}
\item ¿Qué es la educación virtual?
\begin{enumerate}[label=\textbf{\Alph*)}]
\begin{multicols}{2}
  \item La educación impartida a través de plataformas digitales
  \item La educación impartida en aulas presenciales
  \item La educación impartida en bibliotecas
  \item La educación impartida en laboratorios
  \item La educación impartida en talleres
  \item La educación impartida en seminarios
\end{multicols}
\end{enumerate}
\vspace{0.15cm}
\item ¿Qué es la sociedad de la información?
\begin{enumerate}[label=\textbf{\Alph*)}]
\begin{multicols}{2}
  \item Una sociedad en la que la información y el conocimiento son centrales
  \item Una sociedad en la que la agricultura es central
  \item Una sociedad en la que la industria es central
  \item Una sociedad en la que el comercio es central
  \item Una sociedad en la que la minería es central
  \item Una sociedad en la que la pesca es central
\end{multicols}
\end{enumerate}
\vspace{0.15cm}
\item ¿Qué es la ciberseguridad?
\begin{enumerate}[label=\textbf{\Alph*)}]
\begin{multicols}{2}
  \item La protección de sistemas informáticos y datos contra ataques
  \item La protección de sistemas agrícolas
  \item La protección de sistemas industriales
  \item La protección de sistemas comerciales
  \item La protección de sistemas mineros
  \item La protección de sistemas pesqueros
\end{multicols}
\end{enumerate}
\vspace{0.15cm}
\item ¿Qué es la privacidad digital?
\begin{enumerate}[label=\textbf{\Alph*)}]
\begin{multicols}{2}
  \item El derecho a controlar la información personal en el entorno digital
  \item El derecho a controlar la información pública
  \item El derecho a controlar la información gubernamental
  \item El derecho a controlar la información empresarial
  \item El derecho a controlar la información académica
  \item El derecho a controlar la información religiosa
\end{multicols}
\end{enumerate}
\vspace{0.15cm}
\item ¿Qué es la violencia social?
\begin{enumerate}[label=\textbf{\Alph*)}]
\begin{multicols}{2}
  \item El uso de la fuerza para causar daño en la sociedad
  \item La convivencia pacífica
  \item La resolución pacífica de conflictos
  \item La mediación de conflictos
  \item La negociación de conflictos
  \item La reconciliación social
\end{multicols}
\end{enumerate}
\vspace{0.15cm}
\item ¿Qué es la delincuencia?
\begin{enumerate}[label=\textbf{\Alph*)}]
\begin{multicols}{2}
  \item El conjunto de actos ilegales cometidos por personas
  \item El conjunto de actos legales
  \item El conjunto de actos morales
  \item El conjunto de actos éticos
  \item El conjunto de actos solidarios
  \item El conjunto de actos cooperativos
\end{multicols}
\end{enumerate}
\vspace{0.15cm}
\item ¿Qué es una pandilla?
\begin{enumerate}[label=\textbf{\Alph*)}]
\begin{multicols}{2}
  \item Un grupo organizado que comete actos delictivos
  \item Un grupo organizado que comete actos benéficos
  \item Un grupo organizado que comete actos culturales
  \item Un grupo organizado que comete actos deportivos
  \item Un grupo organizado que comete actos religiosos
  \item Un grupo organizado que comete actos políticos
\end{multicols}
\end{enumerate}
\vspace{0.15cm}
\item ¿Qué es el narcotráfico?
\begin{enumerate}[label=\textbf{\Alph*)}]
\begin{multicols}{2}
  \item El comercio ilegal de drogas
  \item El comercio legal de drogas
  \item El comercio de medicamentos
  \item El comercio de alimentos
  \item El comercio de tecnología
  \item El comercio de textiles
\end{multicols}
\end{enumerate}
\vspace{0.15cm}
\item ¿Qué es la corrupción?
\begin{enumerate}[label=\textbf{\Alph*)}]
\begin{multicols}{2}
  \item El abuso del poder público para obtener beneficios privados
  \item El uso adecuado del poder público
  \item El uso transparente del poder público
  \item El uso honesto del poder público
  \item El uso ético del poder público
  \item El uso responsable del poder público
\end{multicols}
\end{enumerate}
\vspace{0.15cm}
\item ¿Qué es la impunidad?
\begin{enumerate}[label=\textbf{\Alph*)}]
\begin{multicols}{2}
  \item La falta de castigo para los delitos cometidos
  \item El castigo para los delitos cometidos
  \item La prevención de los delitos
  \item La investigación de los delitos
  \item La sanción de los delitos
  \item La reparación de los delitos
\end{multicols}
\end{enumerate}
\vspace{0.15cm}
\item ¿Qué es la trata de personas?
\begin{enumerate}[label=\textbf{\Alph*)}]
\begin{multicols}{2}
  \item El comercio ilegal de seres humanos con fines de explotación
  \item El comercio legal de seres humanos
  \item El comercio de bienes
  \item El comercio de servicios
  \item El comercio de tecnología
  \item El comercio de alimentos
\end{multicols}
\end{enumerate}
\vspace{0.15cm}
\item ¿Qué es la explotación laboral?
\begin{enumerate}[label=\textbf{\Alph*)}]
\begin{multicols}{2}
  \item El abuso en las condiciones de trabajo
  \item El respeto a las condiciones de trabajo
  \item La mejora de las condiciones de trabajo
  \item La protección de las condiciones de trabajo
  \item La garantía de las condiciones de trabajo
  \item La promoción de las condiciones de trabajo
\end{multicols}
\end{enumerate}
\vspace{0.15cm}
\item ¿Qué es el trabajo infantil?
\begin{enumerate}[label=\textbf{\Alph*)}]
\begin{multicols}{2}
  \item La explotación de niños en actividades laborales
  \item La educación de niños
  \item La recreación de niños
  \item La protección de niños
  \item La alimentación de niños
  \item La salud de niños
\end{multicols}
\end{enumerate}
\vspace{0.15cm}
\item ¿Qué es la desigualdad económica?
\begin{enumerate}[label=\textbf{\Alph*)}]
\begin{multicols}{2}
  \item La distribución inequitativa de la riqueza
  \item La distribución equitativa de la riqueza
  \item La concentración de la riqueza
  \item La dispersión de la riqueza
  \item La redistribución de la riqueza
  \item La generación de riqueza
\end{multicols}
\end{enumerate}
\vspace{0.15cm}
\item ¿Qué es la pobreza multidimensional?
\begin{enumerate}[label=\textbf{\Alph*)}]
\begin{multicols}{2}
  \item La carencia en múltiples dimensiones del bienestar
  \item La carencia en una sola dimensión
  \item La abundancia en múltiples dimensiones
  \item La abundancia en una sola dimensión
  \item La distribución equitativa del bienestar
  \item La distribución inequitativa del bienestar
\end{multicols}
\end{enumerate}
\vspace{0.15cm}
\item ¿Qué es la exclusión social?
\begin{enumerate}[label=\textbf{\Alph*)}]
\begin{multicols}{2}
  \item La marginación de individuos o grupos de la participación social
  \item La integración de individuos o grupos
  \item La inclusión de individuos o grupos
  \item La participación de individuos o grupos
  \item La incorporación de individuos o grupos
  \item La aceptación de individuos o grupos
\end{multicols}
\end{enumerate}
\vspace{0.15cm}
\item ¿Qué es la discriminación?
\begin{enumerate}[label=\textbf{\Alph*)}]
\begin{multicols}{2}
  \item El trato desigual e injusto hacia un grupo social
  \item El trato igualitario hacia un grupo social
  \item El trato justo hacia un grupo social
  \item El trato objetivo hacia un grupo social
  \item El trato científico hacia un grupo social
  \item El trato artístico hacia un grupo social
\end{multicols}
\end{enumerate}
\vspace{0.15cm}
\item ¿Qué es el racismo?
\begin{enumerate}[label=\textbf{\Alph*)}]
\begin{multicols}{2}
  \item La discriminación basada en la raza o etnia
  \item La discriminación basada en el género
  \item La discriminación basada en la edad
  \item La discriminación basada en la religión
  \item La discriminación basada en la clase social
  \item La discriminación basada en la nacionalidad
\end{multicols}
\end{enumerate}
\vspace{0.15cm}
\item ¿Qué es el sexismo?
\begin{enumerate}[label=\textbf{\Alph*)}]
\begin{multicols}{2}
  \item La discriminación basada en el sexo o género
  \item La discriminación basada en la raza
  \item La discriminación basada en la edad
  \item La discriminación basada en la religión
  \item La discriminación basada en la clase social
  \item La discriminación basada en la nacionalidad
\end{multicols}
\end{enumerate}
\vspace{0.15cm}
\item ¿Qué es la xenofobia?
\begin{enumerate}[label=\textbf{\Alph*)}]
\begin{multicols}{2}
  \item El rechazo hacia los extranjeros
  \item El rechazo hacia los nacionales
  \item El rechazo hacia los ancianos
  \item El rechazo hacia los niños
  \item El rechazo hacia los discapacitados
  \item El rechazo hacia los pobres
\end{multicols}
\end{enumerate}
\vspace{0.15cm}
\item ¿Qué es la homofobia?
\begin{enumerate}[label=\textbf{\Alph*)}]
\begin{multicols}{2}
  \item El rechazo hacia las personas homosexuales
  \item El rechazo hacia las personas heterosexuales
  \item El rechazo hacia las personas bisexuales
  \item El rechazo hacia las personas transgénero
  \item El rechazo hacia las personas intersexuales
  \item El rechazo hacia las personas asexuales
\end{multicols}
\end{enumerate}
\vspace{0.15cm}
\item ¿Qué es la violencia de género?
\begin{enumerate}[label=\textbf{\Alph*)}]
\begin{multicols}{2}
  \item La violencia ejercida contra una persona por su género
  \item La violencia ejercida contra una persona por su edad
  \item La violencia ejercida contra una persona por su raza
  \item La violencia ejercida contra una persona por su religión
  \item La violencia ejercida contra una persona por su clase social
  \item La violencia ejercida contra una persona por su nacionalidad
\end{multicols}
\end{enumerate}
\vspace{0.15cm}
\item ¿Qué es el feminicidio?
\begin{enumerate}[label=\textbf{\Alph*)}]
\begin{multicols}{2}
  \item El asesinato de mujeres por razones de género
  \item El asesinato de hombres por razones de género
  \item El asesinato de niños por razones de género
  \item El asesinato de ancianos por razones de género
  \item El asesinato de extranjeros por razones de género
  \item El asesinato de pobres por razones de género
\end{multicols}
\end{enumerate}
\vspace{0.15cm}
\item ¿Qué es la trata de personas?
\begin{enumerate}[label=\textbf{\Alph*)}]
\begin{multicols}{2}
  \item El comercio ilegal de seres humanos con fines de explotación
  \item El comercio legal de seres humanos
  \item El comercio de bienes
  \item El comercio de servicios
  \item El comercio de tecnología
  \item El comercio de alimentos
\end{multicols}
\end{enumerate}
\vspace{0.15cm}
\item ¿Qué es la migración irregular?
\begin{enumerate}[label=\textbf{\Alph*)}]
\begin{multicols}{2}
  \item El desplazamiento de personas sin documentación legal
  \item El desplazamiento de personas con documentación legal
  \item El desplazamiento de personas por turismo
  \item El desplazamiento de personas por negocios
  \item El desplazamiento de personas por estudio
  \item El desplazamiento de personas por vacaciones
\end{multicols}
\end{enumerate}
\vspace{0.15cm}
\item ¿Qué es la deportación?
\begin{enumerate}[label=\textbf{\Alph*)}]
\begin{multicols}{2}
  \item La expulsión de un extranjero de un país
  \item La admisión de un extranjero en un país
  \item La naturalización de un extranjero
  \item La residencia de un extranjero
  \item La nacionalización de un extranjero
  \item La integración de un extranjero
\end{multicols}
\end{enumerate}
\vspace{0.15cm}
\item ¿Qué es el asilo político?
\begin{enumerate}[label=\textbf{\Alph*)}]
\begin{multicols}{2}
  \item La protección que un país brinda a una persona perseguida por razones políticas
  \item La protección que un país brinda a una persona por razones económicas
  \item La protección que un país brinda a una persona por razones sociales
  \item La protección que un país brinda a una persona por razones culturales
  \item La protección que un país brinda a una persona por razones religiosas
  \item La protección que un país brinda a una persona por razones ambientales
\end{multicols}
\end{enumerate}
\vspace{0.15cm}
\item ¿Qué es la reunificación familiar?
\begin{enumerate}[label=\textbf{\Alph*)}]
\begin{multicols}{2}
  \item El proceso de reunir a los miembros de una familia separada por la migración
  \item El proceso de separar a los miembros de una familia
  \item El proceso de dispersar a los miembros de una familia
  \item El proceso de dividir a los miembros de una familia
  \item El proceso de aislar a los miembros de una familia
  \item El proceso de excluir a los miembros de una familia
\end{multicols}
\end{enumerate}
\vspace{0.15cm}
\item ¿Qué es la integración de los migrantes?
\begin{enumerate}[label=\textbf{\Alph*)}]
\begin{multicols}{2}
  \item El proceso de incorporación de los migrantes a la sociedad de acogida
  \item El proceso de exclusión de los migrantes
  \item El proceso de marginación de los migrantes
  \item El proceso de discriminación de los migrantes
  \item El proceso de segregación de los migrantes
  \item El proceso de aislamiento de los migrantes
\end{multicols}
\end{enumerate}
\vspace{0.15cm}
\item ¿Qué es la fuga de cerebros?
\begin{enumerate}[label=\textbf{\Alph*)}]
\begin{multicols}{2}
  \item La emigración de profesionales altamente calificados
  \item La inmigración de profesionales altamente calificados
  \item La emigración de trabajadores no calificados
  \item La inmigración de trabajadores no calificados
  \item La emigración de estudiantes
  \item La inmigración de estudiantes
\end{multicols}
\end{enumerate}
\vspace{0.15cm}
\item ¿Qué son las remesas?
\begin{enumerate}[label=\textbf{\Alph*)}]
\begin{multicols}{2}
  \item El dinero que los migrantes envían a sus familias en el país de origen
  \item El dinero que los gobiernos envían a otros países
  \item El dinero que las empresas envían a otros países
  \item El dinero que las ONG envían a otros países
  \item El dinero que los bancos envían a otros países
  \item El dinero que los individuos envían a otros individuos
\end{multicols}
\end{enumerate}
\vspace{0.15cm}
\item ¿Qué es la diáspora?
\begin{enumerate}[label=\textbf{\Alph*)}]
\begin{multicols}{2}
  \item La dispersión de un pueblo o comunidad por el mundo
  \item La concentración de un pueblo en un lugar
  \item El aislamiento de un pueblo
  \item La separación de un pueblo
  \item La fragmentación de un pueblo
  \item La localización de un pueblo
\end{multicols}
\end{enumerate}
\vspace{0.15cm}
\item ¿Qué es el transnacionalismo?
\begin{enumerate}[label=\textbf{\Alph*)}]
\begin{multicols}{2}
  \item La creación de vínculos sociales, culturales y económicos entre países
  \item La separación de los países
  \item El aislamiento de los países
  \item La fragmentación de los países
  \item La localización de los países
  \item La regionalización de los países
\end{multicols}
\end{enumerate}
\vspace{0.15cm}
\item ¿Qué es la ciudadanía global?
\begin{enumerate}[label=\textbf{\Alph*)}]
\begin{multicols}{2}
  \item La conciencia de pertenencia a una comunidad mundial
  \item La conciencia de pertenencia a una nación
  \item La conciencia de pertenencia a una región
  \item La conciencia de pertenencia a una ciudad
  \item La conciencia de pertenencia a una comunidad
  \item La conciencia de pertenencia a una familia
\end{multicols}
\end{enumerate}
\vspace{0.15cm}
\item ¿Qué es la sociedad del conocimiento?
\begin{enumerate}[label=\textbf{\Alph*)}]
\begin{multicols}{2}
  \item Una sociedad en la que el conocimiento es el principal recurso
  \item Una sociedad en la que la agricultura es el principal recurso
  \item Una sociedad en la que la industria es el principal recurso
  \item Una sociedad en la que el comercio es el principal recurso
  \item Una sociedad en la que la minería es el principal recurso
  \item Una sociedad en la que la pesca es el principal recurso
\end{multicols}
\end{enumerate}
\vspace{0.15cm}
\item ¿Qué es la brecha digital?
\begin{enumerate}[label=\textbf{\Alph*)}]
\begin{multicols}{2}
  \item La desigualdad en el acceso y uso de las tecnologías de la información
  \item La igualdad en el acceso a la tecnología
  \item La abundancia de tecnología
  \item La escasez de tecnología
  \item La distribución equitativa de tecnología
  \item La democratización de la tecnología
\end{multicols}
\end{enumerate}
\vspace{0.15cm}
\item ¿Qué es la inteligencia artificial?
\begin{enumerate}[label=\textbf{\Alph*)}]
\begin{multicols}{2}
  \item La capacidad de las máquinas para realizar tareas que requieren inteligencia humana
  \item La capacidad de las máquinas para realizar tareas mecánicas
  \item La capacidad de las máquinas para realizar tareas repetitivas
  \item La capacidad de las máquinas para realizar tareas simples
  \item La capacidad de las máquinas para realizar tareas básicas
  \item La capacidad de las máquinas para realizar tareas rutinarias
\end{multicols}
\end{enumerate}
\vspace{0.15cm}
\item ¿Qué es la robótica?
\begin{enumerate}[label=\textbf{\Alph*)}]
\begin{multicols}{2}
  \item La ciencia que estudia el diseño y uso de robots
  \item La ciencia que estudia los animales
  \item La ciencia que estudia las plantas
  \item La ciencia que estudia los minerales
  \item La ciencia que estudia los astros
  \item La ciencia que estudia los fenómenos físicos
\end{multicols}
\end{enumerate}
\vspace{0.15cm}
\item ¿Qué es la automatización?
\begin{enumerate}[label=\textbf{\Alph*)}]
\begin{multicols}{2}
  \item El uso de tecnología para realizar tareas sin intervención humana
  \item El uso de tecnología para realizar tareas con intervención humana
  \item El uso de tecnología para realizar tareas manuales
  \item El uso de tecnología para realizar tareas artesanales
  \item El uso de tecnología para realizar tareas tradicionales
  \item El uso de tecnología para realizar tareas rudimentarias
\end{multicols}
\end{enumerate}
\vspace{0.15cm}
\item ¿Qué es el teletrabajo?
\begin{enumerate}[label=\textbf{\Alph*)}]
\begin{multicols}{2}
  \item El trabajo realizado a distancia mediante tecnologías de la información
  \item El trabajo realizado en una oficina
  \item El trabajo realizado en una fábrica
  \item El trabajo realizado en un campo
  \item El trabajo realizado en una tienda
  \item El trabajo realizado en un taller
\end{multicols}
\end{enumerate}
\vspace{0.15cm}
\item ¿Qué es la educación virtual?
\begin{enumerate}[label=\textbf{\Alph*)}]
\begin{multicols}{2}
  \item La educación impartida a través de plataformas digitales
  \item La educación impartida en aulas presenciales
  \item La educación impartida en bibliotecas
  \item La educación impartida en laboratorios
  \item La educación impartida en talleres
  \item La educación impartida en seminarios
\end{multicols}
\end{enumerate}
\vspace{0.15cm}
\item ¿Qué es la sociedad de la información?
\begin{enumerate}[label=\textbf{\Alph*)}]
\begin{multicols}{2}
  \item Una sociedad en la que la información y el conocimiento son centrales
  \item Una sociedad en la que la agricultura es central
  \item Una sociedad en la que la industria es central
  \item Una sociedad en la que el comercio es central
  \item Una sociedad en la que la minería es central
  \item Una sociedad en la que la pesca es central
\end{multicols}
\end{enumerate}
\vspace{0.15cm}
\item ¿Qué es la ciberseguridad?
\begin{enumerate}[label=\textbf{\Alph*)}]
\begin{multicols}{2}
  \item La protección de sistemas informáticos y datos contra ataques
  \item La protección de sistemas agrícolas
  \item La protección de sistemas industriales
  \item La protección de sistemas comerciales
  \item La protección de sistemas mineros
  \item La protección de sistemas pesqueros
\end{multicols}
\end{enumerate}
\vspace{0.15cm}
\item ¿Qué es la privacidad digital?
\begin{enumerate}[label=\textbf{\Alph*)}]
\begin{multicols}{2}
  \item El derecho a controlar la información personal en el entorno digital
  \item El derecho a controlar la información pública
  \item El derecho a controlar la información gubernamental
  \item El derecho a controlar la información empresarial
  \item El derecho a controlar la información académica
  \item El derecho a controlar la información religiosa
\end{multicols}
\end{enumerate}
\vspace{0.15cm}
\item ¿Qué es la violencia social?
\begin{enumerate}[label=\textbf{\Alph*)}]
\begin{multicols}{2}
  \item El uso de la fuerza para causar daño en la sociedad
  \item La convivencia pacífica
  \item La resolución pacífica de conflictos
  \item La mediación de conflictos
  \item La negociación de conflictos
  \item La reconciliación social
\end{multicols}
\end{enumerate}
\vspace{0.15cm}
\item ¿Qué es la delincuencia?
\begin{enumerate}[label=\textbf{\Alph*)}]
\begin{multicols}{2}
  \item El conjunto de actos ilegales cometidos por personas
  \item El conjunto de actos legales
  \item El conjunto de actos morales
  \item El conjunto de actos éticos
  \item El conjunto de actos solidarios
  \item El conjunto de actos cooperativos
\end{multicols}
\end{enumerate}
\vspace{0.15cm}
\item ¿Qué es una pandilla?
\begin{enumerate}[label=\textbf{\Alph*)}]
\begin{multicols}{2}
  \item Un grupo organizado que comete actos delictivos
  \item Un grupo organizado que comete actos benéficos
  \item Un grupo organizado que comete actos culturales
  \item Un grupo organizado que comete actos deportivos
  \item Un grupo organizado que comete actos religiosos
  \item Un grupo organizado que comete actos políticos
\end{multicols}
\end{enumerate}
\vspace{0.15cm}
\item ¿Qué es el narcotráfico?
\begin{enumerate}[label=\textbf{\Alph*)}]
\begin{multicols}{2}
  \item El comercio ilegal de drogas
  \item El comercio legal de drogas
  \item El comercio de medicamentos
  \item El comercio de alimentos
  \item El comercio de tecnología
  \item El comercio de textiles
\end{multicols}
\end{enumerate}
\vspace{0.15cm}
\item ¿Qué es la corrupción?
\begin{enumerate}[label=\textbf{\Alph*)}]
\begin{multicols}{2}
  \item El abuso del poder público para obtener beneficios privados
  \item El uso adecuado del poder público
  \item El uso transparente del poder público
  \item El uso honesto del poder público
  \item El uso ético del poder público
  \item El uso responsable del poder público
\end{multicols}
\end{enumerate}
\vspace{0.15cm}
\item ¿Qué es la impunidad?
\begin{enumerate}[label=\textbf{\Alph*)}]
\begin{multicols}{2}
  \item La falta de castigo para los delitos cometidos
  \item El castigo para los delitos cometidos
  \item La prevención de los delitos
  \item La investigación de los delitos
  \item La sanción de los delitos
  \item La reparación de los delitos
\end{multicols}
\end{enumerate}
\vspace{0.15cm}
\item ¿Qué es la trata de personas?
\begin{enumerate}[label=\textbf{\Alph*)}]
\begin{multicols}{2}
  \item El comercio ilegal de seres humanos con fines de explotación
  \item El comercio legal de seres humanos
  \item El comercio de bienes
  \item El comercio de servicios
  \item El comercio de tecnología
  \item El comercio de alimentos
\end{multicols}
\end{enumerate}
\vspace{0.15cm}
\item ¿Qué es la explotación laboral?
\begin{enumerate}[label=\textbf{\Alph*)}]
\begin{multicols}{2}
  \item El abuso en las condiciones de trabajo
  \item El respeto a las condiciones de trabajo
  \item La mejora de las condiciones de trabajo
  \item La protección de las condiciones de trabajo
  \item La garantía de las condiciones de trabajo
  \item La promoción de las condiciones de trabajo
\end{multicols}
\end{enumerate}
\vspace{0.15cm}
\item ¿Qué es el trabajo infantil?
\begin{enumerate}[label=\textbf{\Alph*)}]
\begin{multicols}{2}
  \item La explotación de niños en actividades laborales
  \item La educación de niños
  \item La recreación de niños
  \item La protección de niños
  \item La alimentación de niños
  \item La salud de niños
\end{multicols}
\end{enumerate}
\vspace{0.15cm}
\item ¿Qué es la desigualdad económica?
\begin{enumerate}[label=\textbf{\Alph*)}]
\begin{multicols}{2}
  \item La distribución inequitativa de la riqueza
  \item La distribución equitativa de la riqueza
  \item La concentración de la riqueza
  \item La dispersión de la riqueza
  \item La redistribución de la riqueza
  \item La generación de riqueza
\end{multicols}
\end{enumerate}
\vspace{0.15cm}
\item ¿Qué es la pobreza multidimensional?
\begin{enumerate}[label=\textbf{\Alph*)}]
\begin{multicols}{2}
  \item La carencia en múltiples dimensiones del bienestar
  \item La carencia en una sola dimensión
  \item La abundancia en múltiples dimensiones
  \item La abundancia en una sola dimensión
  \item La distribución equitativa del bienestar
  \item La distribución inequitativa del bienestar
\end{multicols}
\end{enumerate}
\vspace{0.15cm}
\item ¿Qué es la exclusión social?
\begin{enumerate}[label=\textbf{\Alph*)}]
\begin{multicols}{2}
  \item La marginación de individuos o grupos de la participación social
  \item La integración de individuos o grupos
  \item La inclusión de individuos o grupos
  \item La participación de individuos o grupos
  \item La incorporación de individuos o grupos
  \item La aceptación de individuos o grupos
\end{multicols}
\end{enumerate}
\vspace{0.15cm}
\item ¿Qué es la discriminación?
\begin{enumerate}[label=\textbf{\Alph*)}]
\begin{multicols}{2}
  \item El trato desigual e injusto hacia un grupo social
  \item El trato igualitario hacia un grupo social
  \item El trato justo hacia un grupo social
  \item El trato objetivo hacia un grupo social
  \item El trato científico hacia un grupo social
  \item El trato artístico hacia un grupo social
\end{multicols}
\end{enumerate}
\vspace{0.15cm}
\item ¿Qué es el racismo?
\begin{enumerate}[label=\textbf{\Alph*)}]
\begin{multicols}{2}
  \item La discriminación basada en la raza o etnia
  \item La discriminación basada en el género
  \item La discriminación basada en la edad
  \item La discriminación basada en la religión
  \item La discriminación basada en la clase social
  \item La discriminación basada en la nacionalidad
\end{multicols}
\end{enumerate}
\vspace{0.15cm}
\item ¿Qué es el sexismo?
\begin{enumerate}[label=\textbf{\Alph*)}]
\begin{multicols}{2}
  \item La discriminación basada en el sexo o género
  \item La discriminación basada en la raza
  \item La discriminación basada en la edad
  \item La discriminación basada en la religión
  \item La discriminación basada en la clase social
  \item La discriminación basada en la nacionalidad
\end{multicols}
\end{enumerate}
\vspace{0.15cm}
\item ¿Qué es la xenofobia?
\begin{enumerate}[label=\textbf{\Alph*)}]
\begin{multicols}{2}
  \item El rechazo hacia los extranjeros
  \item El rechazo hacia los nacionales
  \item El rechazo hacia los ancianos
  \item El rechazo hacia los niños
  \item El rechazo hacia los discapacitados
  \item El rechazo hacia los pobres
\end{multicols}
\end{enumerate}
\vspace{0.15cm}
\item ¿Qué es la homofobia?
\begin{enumerate}[label=\textbf{\Alph*)}]
\begin{multicols}{2}
  \item El rechazo hacia las personas homosexuales
  \item El rechazo hacia las personas heterosexuales
  \item El rechazo hacia las personas bisexuales
  \item El rechazo hacia las personas transgénero
  \item El rechazo hacia las personas intersexuales
  \item El rechazo hacia las personas asexuales
\end{multicols}
\end{enumerate}
\vspace{0.15cm}
\item ¿Qué es la violencia de género?
\begin{enumerate}[label=\textbf{\Alph*)}]
\begin{multicols}{2}
  \item La violencia ejercida contra una persona por su género
  \item La violencia ejercida contra una persona por su edad
  \item La violencia ejercida contra una persona por su raza
  \item La violencia ejercida contra una persona por su religión
  \item La violencia ejercida contra una persona por su clase social
  \item La violencia ejercida contra una persona por su nacionalidad
\end{multicols}
\end{enumerate}
\vspace{0.15cm}
\item ¿Qué es el feminicidio?
\begin{enumerate}[label=\textbf{\Alph*)}]
\begin{multicols}{2}
  \item El asesinato de mujeres por razones de género
  \item El asesinato de hombres por razones de género
  \item El asesinato de niños por razones de género
  \item El asesinato de ancianos por razones de género
  \item El asesinato de extranjeros por razones de género
  \item El asesinato de pobres por razones de género
\end{multicols}
\end{enumerate}
\vspace{0.15cm}
\item ¿Qué es la trata de personas?
\begin{enumerate}[label=\textbf{\Alph*)}]
\begin{multicols}{2}
  \item El comercio ilegal de seres humanos con fines de explotación
  \item El comercio legal de seres humanos
  \item El comercio de bienes
  \item El comercio de servicios
  \item El comercio de tecnología
  \item El comercio de alimentos
\end{multicols}
\end{enumerate}
\vspace{0.15cm}
\end{enumerate}\n
\vspace{0.4cm}
\begin{tcolorbox}[colback=gray!5!white,colframe=gray!60!black,title=\textbf{SOLUCIONARIO -- 4.3 Sociedad, Cultura y Población}]
\begin{multicols}{4}
\begin{enumerate}[label=\textbf{\arabic*.}]
  \item \textbf{A}
  \item \textbf{A}
  \item \textbf{A}
  \item \textbf{A}
  \item \textbf{A}
  \item \textbf{A}
  \item \textbf{A}
  \item \textbf{A}
  \item \textbf{A}
  \item \textbf{A}
  \item \textbf{A}
  \item \textbf{A}
  \item \textbf{A}
  \item \textbf{A}
  \item \textbf{A}
  \item \textbf{A}
  \item \textbf{A}
  \item \textbf{A}
  \item \textbf{A}
  \item \textbf{A}
  \item \textbf{A}
  \item \textbf{A}
  \item \textbf{A}
  \item \textbf{A}
  \item \textbf{A}
  \item \textbf{A}
  \item \textbf{A}
  \item \textbf{A}
  \item \textbf{A}
  \item \textbf{A}
  \item \textbf{A}
  \item \textbf{A}
  \item \textbf{A}
  \item \textbf{A}
  \item \textbf{A}
  \item \textbf{A}
  \item \textbf{A}
  \item \textbf{A}
  \item \textbf{A}
  \item \textbf{A}
  \item \textbf{A}
  \item \textbf{A}
  \item \textbf{A}
  \item \textbf{A}
  \item \textbf{A}
  \item \textbf{A}
  \item \textbf{A}
  \item \textbf{A}
  \item \textbf{A}
  \item \textbf{A}
  \item \textbf{A}
  \item \textbf{A}
  \item \textbf{A}
  \item \textbf{A}
  \item \textbf{A}
  \item \textbf{A}
  \item \textbf{A}
  \item \textbf{A}
  \item \textbf{A}
  \item \textbf{A}
  \item \textbf{A}
  \item \textbf{A}
  \item \textbf{A}
  \item \textbf{A}
  \item \textbf{A}
  \item \textbf{A}
  \item \textbf{A}
  \item \textbf{A}
  \item \textbf{A}
  \item \textbf{A}
  \item \textbf{A}
  \item \textbf{A}
  \item \textbf{A}
  \item \textbf{A}
  \item \textbf{A}
  \item \textbf{A}
  \item \textbf{A}
  \item \textbf{A}
  \item \textbf{A}
  \item \textbf{A}
  \item \textbf{A}
  \item \textbf{A}
  \item \textbf{A}
  \item \textbf{A}
  \item \textbf{A}
  \item \textbf{A}
  \item \textbf{A}
  \item \textbf{A}
  \item \textbf{A}
  \item \textbf{A}
  \item \textbf{A}
  \item \textbf{A}
  \item \textbf{A}
  \item \textbf{A}
  \item \textbf{A}
  \item \textbf{A}
  \item \textbf{A}
  \item \textbf{A}
  \item \textbf{A}
  \item \textbf{A}
  \item \textbf{A}
  \item \textbf{A}
  \item \textbf{A}
  \item \textbf{A}
  \item \textbf{A}
  \item \textbf{A}
  \item \textbf{A}
  \item \textbf{A}
  \item \textbf{A}
  \item \textbf{A}
  \item \textbf{A}
  \item \textbf{A}
  \item \textbf{A}
  \item \textbf{A}
  \item \textbf{A}
  \item \textbf{A}
  \item \textbf{A}
  \item \textbf{A}
  \item \textbf{A}
  \item \textbf{A}
  \item \textbf{A}
  \item \textbf{A}
  \item \textbf{A}
  \item \textbf{A}
  \item \textbf{A}
  \item \textbf{A}
  \item \textbf{A}
  \item \textbf{A}
  \item \textbf{A}
  \item \textbf{A}
  \item \textbf{A}
  \item \textbf{A}
  \item \textbf{A}
  \item \textbf{A}
  \item \textbf{A}
  \item \textbf{A}
  \item \textbf{A}
  \item \textbf{A}
  \item \textbf{A}
  \item \textbf{A}
  \item \textbf{A}
  \item \textbf{A}
  \item \textbf{A}
  \item \textbf{A}
  \item \textbf{A}
  \item \textbf{A}
  \item \textbf{A}
  \item \textbf{A}
  \item \textbf{A}
  \item \textbf{A}
  \item \textbf{A}
  \item \textbf{A}
  \item \textbf{A}
  \item \textbf{A}
  \item \textbf{A}
  \item \textbf{A}
  \item \textbf{A}
  \item \textbf{A}
  \item \textbf{A}
  \item \textbf{A}
  \item \textbf{A}
  \item \textbf{A}
  \item \textbf{A}
  \item \textbf{A}
  \item \textbf{A}
  \item \textbf{A}
  \item \textbf{A}
  \item \textbf{A}
  \item \textbf{A}
  \item \textbf{A}
  \item \textbf{A}
  \item \textbf{A}
  \item \textbf{A}
  \item \textbf{A}
  \item \textbf{A}
  \item \textbf{A}
  \item \textbf{A}
  \item \textbf{A}
  \item \textbf{A}
  \item \textbf{A}
  \item \textbf{A}
  \item \textbf{A}
  \item \textbf{A}
  \item \textbf{A}
  \item \textbf{A}
  \item \textbf{A}
  \item \textbf{A}
  \item \textbf{A}
  \item \textbf{A}
  \item \textbf{A}
  \item \textbf{A}
  \item \textbf{A}
  \item \textbf{A}
  \item \textbf{A}
  \item \textbf{A}
  \item \textbf{A}
  \item \textbf{A}
  \item \textbf{A}
  \item \textbf{A}
  \item \textbf{A}
  \item \textbf{A}
  \item \textbf{A}
  \item \textbf{A}
  \item \textbf{A}
  \item \textbf{A}
  \item \textbf{A}
  \item \textbf{A}
  \item \textbf{A}
  \item \textbf{A}
  \item \textbf{A}
  \item \textbf{A}
  \item \textbf{A}
  \item \textbf{A}
  \item \textbf{A}
  \item \textbf{A}
  \item \textbf{A}
  \item \textbf{A}
  \item \textbf{A}
  \item \textbf{A}
  \item \textbf{A}
  \item \textbf{A}
  \item \textbf{A}
  \item \textbf{A}
  \item \textbf{A}
  \item \textbf{A}
  \item \textbf{A}
  \item \textbf{A}
  \item \textbf{A}
  \item \textbf{A}
  \item \textbf{A}
  \item \textbf{A}
  \item \textbf{A}
  \item \textbf{A}
  \item \textbf{A}
  \item \textbf{A}
  \item \textbf{A}
  \item \textbf{A}
  \item \textbf{A}
  \item \textbf{A}
  \item \textbf{A}
  \item \textbf{A}
  \item \textbf{A}
  \item \textbf{A}
  \item \textbf{A}
  \item \textbf{A}
  \item \textbf{A}
  \item \textbf{A}
  \item \textbf{A}
  \item \textbf{A}
  \item \textbf{A}
  \item \textbf{A}
  \item \textbf{A}
  \item \textbf{A}
  \item \textbf{A}
  \item \textbf{A}
  \item \textbf{A}
  \item \textbf{A}
  \item \textbf{A}
  \item \textbf{A}
  \item \textbf{A}
  \item \textbf{A}
  \item \textbf{A}
  \item \textbf{A}
  \item \textbf{A}
  \item \textbf{A}
  \item \textbf{A}
  \item \textbf{A}
  \item \textbf{A}
  \item \textbf{A}
  \item \textbf{A}
  \item \textbf{A}
  \item \textbf{A}
  \item \textbf{A}
  \item \textbf{A}
  \item \textbf{A}
  \item \textbf{A}
  \item \textbf{A}
  \item \textbf{A}
  \item \textbf{A}
  \item \textbf{A}
  \item \textbf{A}
  \item \textbf{A}
  \item \textbf{A}
  \item \textbf{A}
  \item \textbf{A}
  \item \textbf{A}
  \item \textbf{A}
  \item \textbf{A}
  \item \textbf{A}
  \item \textbf{A}
  \item \textbf{A}
  \item \textbf{A}
  \item \textbf{A}
  \item \textbf{A}
  \item \textbf{A}
  \item \textbf{A}
  \item \textbf{A}
  \item \textbf{A}
  \item \textbf{A}
  \item \textbf{A}
\end{enumerate}
\end{multicols}
\end{tcolorbox}
\newpage\n
\subsection{4.4 Derechos y Sociedad Actual}

\subsubsection*{Sección I: Derechos Humanos – Conceptos Básicos (1–50)}

\begin{enumerate}[label=\textbf{\arabic*.}]
\setcounter{enumi}{0}
\item ¿Qué son los derechos humanos?
\begin{enumerate}[label=\textbf{\Alph*)}]
\begin{multicols}{2}
  \item Facultades y libertades inherentes a todas las personas por el hecho de ser humanas
  \item Privilegios otorgados por el Estado
  \item Beneficios económicos de los ciudadanos
  \item Concesiones de las empresas privadas
  \item Regalos de los gobiernos
  \item Favores de las autoridades
\end{multicols}
\end{enumerate}
\vspace{0.15cm}
\item ¿Cuál es la característica principal de los derechos humanos?
\begin{enumerate}[label=\textbf{\Alph*)}]
\begin{multicols}{2}
  \item Son universales, inalienables e indivisibles
  \item Son locales y temporales
  \item Son negociables y revocables
  \item Son opcionales y condicionales
  \item Son particulares y divisibles
  \item Son transitorios y limitados
\end{multicols}
\end{enumerate}
\vspace{0.15cm}
\item ¿Qué significa que los derechos humanos sean universales?
\begin{enumerate}[label=\textbf{\Alph*)}]
\begin{multicols}{2}
  \item Que pertenecen a todas las personas sin excepción
  \item Que pertenecen solo a algunas personas
  \item Que pertenecen solo a los ciudadanos
  \item Que pertenecen solo a los adultos
  \item Que pertenecen solo a los nacionales
  \item Que pertenecen solo a los ricos
\end{multicols}
\end{enumerate}
\vspace{0.15cm}
\item ¿Qué significa que los derechos humanos sean inalienables?
\begin{enumerate}[label=\textbf{\Alph*)}]
\begin{multicols}{2}
  \item Que no se pueden renunciar a ellos ni transferirlos
  \item Que se pueden renunciar a ellos
  \item Que se pueden transferir
  \item Que se pueden vender
  \item Que se pueden negociar
  \item Que se pueden ceder
\end{multicols}
\end{enumerate}
\vspace{0.15cm}
\item ¿Qué significa que los derechos humanos sean indivisibles?
\begin{enumerate}[label=\textbf{\Alph*)}]
\begin{multicols}{2}
  \item Que no se pueden separar ni jerarquizar
  \item Que se pueden separar
  \item Que se pueden jerarquizar
  \item Que se pueden dividir
  \item Que se pueden fragmentar
  \item Que se pueden clasificar
\end{multicols}
\end{enumerate}
\vspace{0.15cm}
\item ¿Cuál es la primera generación de derechos humanos?
\begin{enumerate}[label=\textbf{\Alph*)}]
\begin{multicols}{2}
  \item Los derechos civiles y políticos
  \item Los derechos económicos, sociales y culturales
  \item Los derechos de solidaridad
  \item Los derechos ambientales
  \item Los derechos digitales
  \item Los derechos colectivos
\end{multicols}
\end{enumerate}
\vspace{0.15cm}
\item ¿Cuál es la segunda generación de derechos humanos?
\begin{enumerate}[label=\textbf{\Alph*)}]
\begin{multicols}{2}
  \item Los derechos económicos, sociales y culturales
  \item Los derechos civiles y políticos
  \item Los derechos de solidaridad
  \item Los derechos ambientales
  \item Los derechos digitales
  \item Los derechos colectivos
\end{multicols}
\end{enumerate}
\vspace{0.15cm}
\item ¿Cuál es la tercera generación de derechos humanos?
\begin{enumerate}[label=\textbf{\Alph*)}]
\begin{multicols}{2}
  \item Los derechos de solidaridad o colectivos
  \item Los derechos civiles y políticos
  \item Los derechos económicos, sociales y culturales
  \item Los derechos ambientales
  \item Los derechos digitales
  \item Los derechos individuales
\end{multicols}
\end{enumerate}
\vspace{0.15cm}
\item ¿Cuál es la cuarta generación de derechos humanos?
\begin{enumerate}[label=\textbf{\Alph*)}]
\begin{multicols}{2}
  \item Los derechos digitales y tecnológicos
  \item Los derechos civiles y políticos
  \item Los derechos económicos, sociales y culturales
  \item Los derechos de solidaridad
  \item Los derechos ambientales
  \item Los derechos individuales
\end{multicols}
\end{enumerate}
\vspace{0.15cm}
\item ¿Qué es la dignidad humana?
\begin{enumerate}[label=\textbf{\Alph*)}]
\begin{multicols}{2}
  \item El valor inherente de toda persona por el hecho de ser humana
  \item El valor económico de una persona
  \item El valor social de una persona
  \item El valor político de una persona
  \item El valor cultural de una persona
  \item El valor religioso de una persona
\end{multicols}
\end{enumerate}
\vspace{0.15cm}
\item ¿Qué es la libertad?
\begin{enumerate}[label=\textbf{\Alph*)}]
\begin{multicols}{2}
  \item La facultad de actuar según la propia voluntad sin coacción
  \item La facultad de actuar según la voluntad de otros
  \item La facultad de actuar sin responsabilidad
  \item La facultad de actuar sin límites
  \item La facultad de actuar sin consecuencias
  \item La facultad de actuar sin normas
\end{multicols}
\end{enumerate}
\vspace{0.15cm}
\item ¿Qué es la igualdad?
\begin{enumerate}[label=\textbf{\Alph*)}]
\begin{multicols}{2}
  \item El principio de que todas las personas tienen los mismos derechos
  \item El principio de que algunas personas tienen más derechos
  \item El principio de que algunas personas tienen menos derechos
  \item El principio de que los derechos dependen del mérito
  \item El principio de que los derechos dependen de la clase social
  \item El principio de que los derechos dependen de la nacionalidad
\end{multicols}
\end{enumerate}
\vspace{0.15cm}
\item ¿Qué es la no discriminación?
\begin{enumerate}[label=\textbf{\Alph*)}]
\begin{multicols}{2}
  \item El principio de que nadie puede ser tratado de manera desigual por razones arbitrarias
  \item El principio de que se puede discriminar por razones arbitrarias
  \item El principio de que se puede discriminar por razones justificadas
  \item El principio de que se puede discriminar por razones legales
  \item El principio de que se puede discriminar por razones culturales
  \item El principio de que se puede discriminar por razones religiosas
\end{multicols}
\end{enumerate}
\vspace{0.15cm}
\item ¿Qué es la solidaridad?
\begin{enumerate}[label=\textbf{\Alph*)}]
\begin{multicols}{2}
  \item La ayuda mutua entre los miembros de una sociedad
  \item La indiferencia entre los miembros de una sociedad
  \item La competencia entre los miembros de una sociedad
  \item La exclusión entre los miembros de una sociedad
  \item La marginación entre los miembros de una sociedad
  \item La discriminación entre los miembros de una sociedad
\end{multicols}
\end{enumerate}
\vspace{0.15cm}
\item ¿Qué es la justicia social?
\begin{enumerate}[label=\textbf{\Alph*)}]
\begin{multicols}{2}
  \item La distribución equitativa de los bienes y oportunidades
  \item La distribución inequitativa de los bienes
  \item La concentración de los bienes
  \item La acumulación de los bienes
  \item La privatización de los bienes
  \item La estatización de los bienes
\end{multicols}
\end{enumerate}
\vspace{0.15cm}
\item ¿Qué es el Estado de derecho?
\begin{enumerate}[label=\textbf{\Alph*)}]
\begin{multicols}{2}
  \item El principio de que todos, incluidos los gobernantes, están sujetos a la ley
  \item El principio de que los gobernantes están por encima de la ley
  \item El principio de que la ley no se aplica a los gobernantes
  \item El principio de que la ley es opcional
  \item El principio de que la ley es relativa
  \item El principio de que la ley es arbitraria
\end{multicols}
\end{enumerate}
\vspace{0.15cm}
\item ¿Qué es la democracia?
\begin{enumerate}[label=\textbf{\Alph*)}]
\begin{multicols}{2}
  \item Un sistema de gobierno en el que el poder reside en el pueblo
  \item Un sistema de gobierno en el que el poder reside en un rey
  \item Un sistema de gobierno en el que el poder reside en un dictador
  \item Un sistema de gobierno en el que el poder reside en los militares
  \item Un sistema de gobierno en el que el poder reside en los jueces
  \item Un sistema de gobierno en el que el poder reside en los empresarios
\end{multicols}
\end{enumerate}
\vspace{0.15cm}
\item ¿Qué es la ciudadanía?
\begin{enumerate}[label=\textbf{\Alph*)}]
\begin{multicols}{2}
  \item La condición de pertenencia a un Estado con derechos y deberes
  \item La condición de pertenencia a una familia
  \item La condición de pertenencia a una comunidad
  \item La condición de pertenencia a una etnia
  \item La condición de pertenencia a una religión
  \item La condición de pertenencia a un partido
\end{multicols}
\end{enumerate}
\vspace{0.15cm}
\item ¿Qué es la participación ciudadana?
\begin{enumerate}[label=\textbf{\Alph*)}]
\begin{multicols}{2}
  \item La intervención de los ciudadanos en los asuntos públicos
  \item La abstención de los ciudadanos
  \item La indiferencia de los ciudadanos
  \item La pasividad de los ciudadanos
  \item La exclusión de los ciudadanos
  \item La marginación de los ciudadanos
\end{multicols}
\end{enumerate}
\vspace{0.15cm}
\item ¿Qué es la rendición de cuentas?
\begin{enumerate}[label=\textbf{\Alph*)}]
\begin{multicols}{2}
  \item La obligación de los gobernantes de informar y justificar sus acciones
  \item La obligación de los ciudadanos de informar
  \item La obligación de las empresas de informar
  \item La obligación de las ONG de informar
  \item La obligación de los medios de informar
  \item La obligación de los jueces de informar
\end{multicols}
\end{enumerate}
\vspace{0.15cm}
\item ¿Qué es la transparencia?
\begin{enumerate}[label=\textbf{\Alph*)}]
\begin{multicols}{2}
  \item La obligación de los gobiernos de actuar de manera pública y verificable
  \item La obligación de los ciudadanos de actuar públicamente
  \item La obligación de las empresas de actuar públicamente
  \item La obligación de las ONG de actuar públicamente
  \item La obligación de los medios de actuar públicamente
  \item La obligación de los jueces de actuar públicamente
\end{multicols}
\end{enumerate}
\vspace{0.15cm}
\item ¿Qué es la corrupción?
\begin{enumerate}[label=\textbf{\Alph*)}]
\begin{multicols}{2}
  \item El abuso del poder público para obtener beneficios privados
  \item El uso adecuado del poder público
  \item El uso transparente del poder público
  \item El uso honesto del poder público
  \item El uso ético del poder público
  \item El uso responsable del poder público
\end{multicols}
\end{enumerate}
\vspace{0.15cm}
\item ¿Qué es la impunidad?
\begin{enumerate}[label=\textbf{\Alph*)}]
\begin{multicols}{2}
  \item La falta de castigo para los delitos cometidos
  \item El castigo para los delitos cometidos
  \item La prevención de los delitos
  \item La investigación de los delitos
  \item La sanción de los delitos
  \item La reparación de los delitos
\end{multicols}
\end{enumerate}
\vspace{0.15cm}
\item ¿Qué es la violencia?
\begin{enumerate}[label=\textbf{\Alph*)}]
\begin{multicols}{2}
  \item El uso de la fuerza para causar daño
  \item El uso de la fuerza para prevenir daño
  \item El uso de la fuerza para proteger
  \item El uso de la fuerza para defender
  \item El uso de la fuerza para auxiliar
  \item El uso de la fuerza para ayudar
\end{multicols}
\end{enumerate}
\vspace{0.15cm}
\item ¿Qué es la paz?
\begin{enumerate}[label=\textbf{\Alph*)}]
\begin{multicols}{2}
  \item La ausencia de violencia y la presencia de justicia
  \item La ausencia de violencia solamente
  \item La presencia de violencia
  \item La ausencia de justicia
  \item La presencia de conflictos
  \item La ausencia de conflictos
\end{multicols}
\end{enumerate}
\vspace{0.15cm}
\item ¿Qué es la tolerancia?
\begin{enumerate}[label=\textbf{\Alph*)}]
\begin{multicols}{2}
  \item El respeto hacia las ideas y prácticas diferentes
  \item El rechazo hacia las ideas diferentes
  \item La indiferencia hacia las ideas diferentes
  \item La imposición de las ideas propias
  \item La discriminación de las ideas diferentes
  \item La exclusión de las ideas diferentes
\end{multicols}
\end{enumerate}
\vspace{0.15cm}
\item ¿Qué es la libertad de expresión?
\begin{enumerate}[label=\textbf{\Alph*)}]
\begin{multicols}{2}
  \item El derecho a expresar opiniones sin censura
  \item El derecho a expresar solo opiniones oficiales
  \item El derecho a expresar solo opiniones mayoritarias
  \item El derecho a expresar solo opiniones religiosas
  \item El derecho a expresar solo opiniones políticas
  \item El derecho a expresar solo opiniones culturales
\end{multicols}
\end{enumerate}
\vspace{0.15cm}
\item ¿Qué es la libertad de prensa?
\begin{enumerate}[label=\textbf{\Alph*)}]
\begin{multicols}{2}
  \item El derecho de los medios a informar sin censura
  \item El derecho de los medios a informar con censura
  \item El derecho de los medios a informar solo lo oficial
  \item El derecho de los medios a informar solo lo comercial
  \item El derecho de los medios a informar solo lo religioso
  \item El derecho de los medios a informar solo lo político
\end{multicols}
\end{enumerate}
\vspace{0.15cm}
\item ¿Qué es la libertad de asociación?
\begin{enumerate}[label=\textbf{\Alph*)}]
\begin{multicols}{2}
  \item El derecho a formar parte de organizaciones sin coacción
  \item El derecho a formar parte solo de organizaciones oficiales
  \item El derecho a formar parte solo de organizaciones religiosas
  \item El derecho a formar parte solo de organizaciones políticas
  \item El derecho a formar parte solo de organizaciones culturales
  \item El derecho a formar parte solo de organizaciones económicas
\end{multicols}
\end{enumerate}
\vspace{0.15cm}
\item ¿Qué es la libertad de reunión?
\begin{enumerate}[label=\textbf{\Alph*)}]
\begin{multicols}{2}
  \item El derecho a reunirse pacíficamente sin coacción
  \item El derecho a reunirse solo con permiso oficial
  \item El derecho a reunirse solo con permiso religioso
  \item El derecho a reunirse solo con permiso político
  \item El derecho a reunirse solo con permiso cultural
  \item El derecho a reunirse solo con permiso económico
\end{multicols}
\end{enumerate}
\vspace{0.15cm}
\item ¿Qué es la libertad de religión?
\begin{enumerate}[label=\textbf{\Alph*)}]
\begin{multicols}{2}
  \item El derecho a profesar y practicar la religión de elección
  \item El derecho a profesar solo la religión oficial
  \item El derecho a profesar solo la religión mayoritaria
  \item El derecho a profesar solo la religión tradicional
  \item El derecho a profesar solo la religión estatal
  \item El derecho a profesar solo la religión nacional
\end{multicols}
\end{enumerate}
\vspace{0.15cm}
\item ¿Qué es la libertad de conciencia?
\begin{enumerate}[label=\textbf{\Alph*)}]
\begin{multicols}{2}
  \item El derecho a tener y cambiar las propias creencias
  \item El derecho a tener solo creencias oficiales
  \item El derecho a tener solo creencias mayoritarias
  \item El derecho a tener solo creencias tradicionales
  \item El derecho a tener solo creencias religiosas
  \item El derecho a tener solo creencias políticas
\end{multicols}
\end{enumerate}
\vspace{0.15cm}
\item ¿Qué es el derecho a la vida?
\begin{enumerate}[label=\textbf{\Alph*)}]
\begin{multicols}{2}
  \item El derecho fundamental de toda persona a existir
  \item El derecho de algunas personas a existir
  \item El derecho condicional a existir
  \item El derecho opcional a existir
  \item El derecho temporal a existir
  \item El derecho relativo a existir
\end{multicols}
\end{enumerate}
\vspace{0.15cm}
\item ¿Qué es el derecho a la integridad personal?
\begin{enumerate}[label=\textbf{\Alph*)}]
\begin{multicols}{2}
  \item El derecho a no ser sometido a torturas ni tratos crueles
  \item El derecho a ser sometido a torturas
  \item El derecho a ser sometido a tratos crueles
  \item El derecho a ser sometido a penas degradantes
  \item El derecho a ser sometido a violencia
  \item El derecho a ser sometido a abuso
\end{multicols}
\end{enumerate}
\vspace{0.15cm}
\item ¿Qué es el derecho a la libertad personal?
\begin{enumerate}[label=\textbf{\Alph*)}]
\begin{multicols}{2}
  \item El derecho a no ser detenido arbitrariamente
  \item El derecho a ser detenido arbitrariamente
  \item El derecho a ser encarcelado sin causa
  \item El derecho a ser privado de libertad sin juicio
  \item El derecho a ser secuestrado
  \item El derecho a ser arrestado sin orden
\end{multicols}
\end{enumerate}
\vspace{0.15cm}
\item ¿Qué es el derecho al debido proceso?
\begin{enumerate}[label=\textbf{\Alph*)}]
\begin{multicols}{2}
  \item El derecho a un juicio justo con garantías
  \item El derecho a un juicio sin garantías
  \item El derecho a un juicio arbitrario
  \item El derecho a un juicio secreto
  \item El derecho a un juicio sin abogado
  \item El derecho a un juicio sin pruebas
\end{multicols}
\end{enumerate}
\vspace{0.15cm}
\item ¿Qué es el derecho a la presunción de inocencia?
\begin{enumerate}[label=\textbf{\Alph*)}]
\begin{multicols}{2}
  \item El derecho a ser considerado inocente hasta que se pruebe lo contrario
  \item El derecho a ser considerado culpable hasta que se pruebe lo contrario
  \item El derecho a ser considerado culpable sin pruebas
  \item El derecho a ser considerado culpable sin juicio
  \item El derecho a ser considerado culpable sin defensa
  \item El derecho a ser considerado culpable sin apelación
\end{multicols}
\end{enumerate}
\vspace{0.15cm}
\item ¿Qué es el derecho a la privacidad?
\begin{enumerate}[label=\textbf{\Alph*)}]
\begin{multicols}{2}
  \item El derecho a proteger la vida privada y los datos personales
  \item El derecho a invadir la vida privada
  \item El derecho a publicar datos personales
  \item El derecho a difundir información privada
  \item El derecho a espiar a otros
  \item El derecho a vigilar a otros
\end{multicols}
\end{enumerate}
\vspace{0.15cm}
\item ¿Qué es el derecho a la educación?
\begin{enumerate}[label=\textbf{\Alph*)}]
\begin{multicols}{2}
  \item El derecho de toda persona a recibir educación
  \item El derecho de algunas personas a recibir educación
  \item El derecho condicional a la educación
  \item El derecho opcional a la educación
  \item El derecho temporal a la educación
  \item El derecho relativo a la educación
\end{multicols}
\end{enumerate}
\vspace{0.15cm}
\item ¿Qué es el derecho a la salud?
\begin{enumerate}[label=\textbf{\Alph*)}]
\begin{multicols}{2}
  \item El derecho de toda persona a disfrutar del más alto nivel de salud
  \item El derecho de algunas personas a la salud
  \item El derecho condicional a la salud
  \item El derecho opcional a la salud
  \item El derecho temporal a la salud
  \item El derecho relativo a la salud
\end{multicols}
\end{enumerate}
\vspace{0.15cm}
\item ¿Qué es el derecho al trabajo?
\begin{enumerate}[label=\textbf{\Alph*)}]
\begin{multicols}{2}
  \item El derecho de toda persona a trabajar en condiciones dignas
  \item El derecho de algunas personas a trabajar
  \item El derecho condicional al trabajo
  \item El derecho opcional al trabajo
  \item El derecho temporal al trabajo
  \item El derecho relativo al trabajo
\end{multicols}
\end{enumerate}
\vspace{0.15cm}
\item ¿Qué es el derecho a la vivienda?
\begin{enumerate}[label=\textbf{\Alph*)}]
\begin{multicols}{2}
  \item El derecho de toda persona a una vivienda adecuada
  \item El derecho de algunas personas a una vivienda
  \item El derecho condicional a la vivienda
  \item El derecho opcional a la vivienda
  \item El derecho temporal a la vivienda
  \item El derecho relativo a la vivienda
\end{multicols}
\end{enumerate}
\vspace{0.15cm}
\item ¿Qué es el derecho a la alimentación?
\begin{enumerate}[label=\textbf{\Alph*)}]
\begin{multicols}{2}
  \item El derecho de toda persona a una alimentación adecuada
  \item El derecho de algunas personas a la alimentación
  \item El derecho condicional a la alimentación
  \item El derecho opcional a la alimentación
  \item El derecho temporal a la alimentación
  \item El derecho relativo a la alimentación
\end{multicols}
\end{enumerate}
\vspace{0.15cm}
\item ¿Qué es el derecho al agua?
\begin{enumerate}[label=\textbf{\Alph*)}]
\begin{multicols}{2}
  \item El derecho de toda persona a acceder al agua potable
  \item El derecho de algunas personas al agua
  \item El derecho condicional al agua
  \item El derecho opcional al agua
  \item El derecho temporal al agua
  \item El derecho relativo al agua
\end{multicols}
\end{enumerate}
\vspace{0.15cm}
\item ¿Qué es el derecho al desarrollo?
\begin{enumerate}[label=\textbf{\Alph*)}]
\begin{multicols}{2}
  \item El derecho de toda persona a participar y beneficiarse del desarrollo
  \item El derecho de algunas personas al desarrollo
  \item El derecho condicional al desarrollo
  \item El derecho opcional al desarrollo
  \item El derecho temporal al desarrollo
  \item El derecho relativo al desarrollo
\end{multicols}
\end{enumerate}
\vspace{0.15cm}
\item ¿Qué es el derecho a la paz?
\begin{enumerate}[label=\textbf{\Alph*)}]
\begin{multicols}{2}
  \item El derecho de toda persona a vivir en un mundo sin violencia
  \item El derecho de algunas personas a la paz
  \item El derecho condicional a la paz
  \item El derecho opcional a la paz
  \item El derecho temporal a la paz
  \item El derecho relativo a la paz
\end{multicols}
\end{enumerate}
\vspace{0.15cm}
\item ¿Qué es el derecho a un medio ambiente sano?
\begin{enumerate}[label=\textbf{\Alph*)}]
\begin{multicols}{2}
  \item El derecho de toda persona a vivir en un ambiente equilibrado
  \item El derecho de algunas personas a un ambiente sano
  \item El derecho condicional a un ambiente sano
  \item El derecho opcional a un ambiente sano
  \item El derecho temporal a un ambiente sano
  \item El derecho relativo a un ambiente sano
\end{multicols}
\end{enumerate}
\vspace{0.15cm}
\item ¿Qué es el derecho a la verdad?
\begin{enumerate}[label=\textbf{\Alph*)}]
\begin{multicols}{2}
  \item El derecho de las víctimas a conocer la verdad sobre las violaciones
  \item El derecho de algunas personas a la verdad
  \item El derecho condicional a la verdad
  \item El derecho opcional a la verdad
  \item El derecho temporal a la verdad
  \item El derecho relativo a la verdad
\end{multicols}
\end{enumerate}
\vspace{0.15cm}
\item ¿Qué es el derecho a la justicia?
\begin{enumerate}[label=\textbf{\Alph*)}]
\begin{multicols}{2}
  \item El derecho de las víctimas a obtener justicia
  \item El derecho de algunas personas a la justicia
  \item El derecho condicional a la justicia
  \item El derecho opcional a la justicia
  \item El derecho temporal a la justicia
  \item El derecho relativo a la justicia
\end{multicols}
\end{enumerate}
\vspace{0.15cm}
\item ¿Qué es el derecho a la reparación?
\begin{enumerate}[label=\textbf{\Alph*)}]
\begin{multicols}{2}
  \item El derecho de las víctimas a ser reparadas por los daños sufridos
  \item El derecho de algunas personas a la reparación
  \item El derecho condicional a la reparación
  \item El derecho opcional a la reparación
  \item El derecho temporal a la reparación
  \item El derecho relativo a la reparación
\end{multicols}
\end{enumerate}
\vspace{0.15cm}
\end{enumerate}\n
\subsubsection*{Sección II: Evolución Histórica de los Derechos Humanos (51–100)}

\begin{enumerate}[label=\textbf{\arabic*.}]
\setcounter{enumi}{50}
\item ¿Qué documento antiguo es considerado un antecedente de los derechos humanos?
\begin{enumerate}[label=\textbf{\Alph*)}]
\begin{multicols}{2}
  \item El Cilindro de Ciro
  \item La Ilíada
  \item La Odisea
  \item La Eneida
  \item El Ramayana
  \item El Mahabharata
\end{multicols}
\end{enumerate}
\vspace{0.15cm}
\item ¿Qué documento inglés de 1215 limitó el poder del rey?
\begin{enumerate}[label=\textbf{\Alph*)}]
\begin{multicols}{2}
  \item La Carta Magna
  \item La Petición de Derechos
  \item La Declaración de Derechos
  \item El Acta de Hábeas Corpus
  \item La Carta de Derechos
  \item El Bill of Rights
\end{multicols}
\end{enumerate}
\vspace{0.15cm}
\item ¿Qué documento inglés de 1679 estableció el hábeas corpus?
\begin{enumerate}[label=\textbf{\Alph*)}]
\begin{multicols}{2}
  \item El Acta de Hábeas Corpus
  \item La Carta Magna
  \item La Petición de Derechos
  \item La Declaración de Derechos
  \item La Carta de Derechos
  \item El Bill of Rights
\end{multicols}
\end{enumerate}
\vspace{0.15cm}
\item ¿Qué documento inglés de 1689 limitó el poder de la monarquía?
\begin{enumerate}[label=\textbf{\Alph*)}]
\begin{multicols}{2}
  \item El Bill of Rights
  \item La Carta Magna
  \item La Petición de Derechos
  \item La Declaración de Derechos
  \item El Acta de Hábeas Corpus
  \item La Carta de Derechos
\end{multicols}
\end{enumerate}
\vspace{0.15cm}
\item ¿Qué documento estadounidense de 1776 proclamó la independencia y los derechos?
\begin{enumerate}[label=\textbf{\Alph*)}]
\begin{multicols}{2}
  \item La Declaración de Independencia
  \item La Constitución
  \item La Declaración de Derechos
  \item El Bill of Rights
  \item La Carta de Derechos
  \item La Petición de Derechos
\end{multicols}
\end{enumerate}
\vspace{0.15cm}
\item ¿Qué documento francés de 1789 proclamó los derechos del hombre?
\begin{enumerate}[label=\textbf{\Alph*)}]
\begin{multicols}{2}
  \item La Declaración de los Derechos del Hombre y del Ciudadano
  \item La Constitución
  \item La Declaración de Independencia
  \item El Bill of Rights
  \item La Carta Magna
  \item El Acta de Hábeas Corpus
\end{multicols}
\end{enumerate}
\vspace{0.15cm}
\item ¿Qué documento de 1948 proclamó los derechos humanos a nivel internacional?
\begin{enumerate}[label=\textbf{\Alph*)}]
\begin{multicols}{2}
  \item La Declaración Universal de los Derechos Humanos
  \item La Carta de las Naciones Unidas
  \item La Declaración de los Derechos del Hombre
  \item El Pacto Internacional de Derechos Civiles
  \item El Pacto Internacional de Derechos Económicos
  \item La Convención Americana de Derechos Humanos
\end{multicols}
\end{enumerate}
\vspace{0.15cm}
\item ¿Qué organización aprobó la Declaración Universal de los Derechos Humanos?
\begin{enumerate}[label=\textbf{\Alph*)}]
\begin{multicols}{2}
  \item La ONU
  \item La OEA
  \item La Unión Europea
  \item La Liga Árabe
  \item La Unión Africana
  \item La ASEAN
\end{multicols}
\end{enumerate}
\vspace{0.15cm}
\item ¿En qué año se aprobó la Declaración Universal de los Derechos Humanos?
\begin{enumerate}[label=\textbf{\Alph*)}]
\begin{multicols}{2}
  \item 1948
  \item 1945
  \item 1950
  \item 1960
  \item 1970
  \item 1980
\end{multicols}
\end{enumerate}
\vspace{0.15cm}
\item ¿Cuántos artículos tiene la Declaración Universal de los Derechos Humanos?
\begin{enumerate}[label=\textbf{\Alph*)}]
\begin{multicols}{2}
  \item 30
  \item 20
  \item 40
  \item 50
  \item 10
  \item 60
\end{multicols}
\end{enumerate}
\vspace{0.15cm}
\item ¿Qué documento de 1966 desarrolló los derechos civiles y políticos?
\begin{enumerate}[label=\textbf{\Alph*)}]
\begin{multicols}{2}
  \item El Pacto Internacional de Derechos Civiles y Políticos
  \item El Pacto Internacional de Derechos Económicos, Sociales y Culturales
  \item La Declaración Universal de los Derechos Humanos
  \item La Convención Americana de Derechos Humanos
  \item La Carta de las Naciones Unidas
  \item La Declaración de los Derechos del Hombre
\end{multicols}
\end{enumerate}
\vspace{0.15cm}
\item ¿Qué documento de 1966 desarrolló los derechos económicos, sociales y culturales?
\begin{enumerate}[label=\textbf{\Alph*)}]
\begin{multicols}{2}
  \item El Pacto Internacional de Derechos Económicos, Sociales y Culturales
  \item El Pacto Internacional de Derechos Civiles y Políticos
  \item La Declaración Universal de los Derechos Humanos
  \item La Convención Americana de Derechos Humanos
  \item La Carta de las Naciones Unidas
  \item La Declaración de los Derechos del Hombre
\end{multicols}
\end{enumerate}
\vspace{0.15cm}
\item ¿Qué documento de 1969 protege los derechos humanos en América?
\begin{enumerate}[label=\textbf{\Alph*)}]
\begin{multicols}{2}
  \item La Convención Americana de Derechos Humanos
  \item La Declaración Universal de los Derechos Humanos
  \item El Pacto Internacional de Derechos Civiles
  \item El Pacto Internacional de Derechos Económicos
  \item La Carta de las Naciones Unidas
  \item La Declaración de los Derechos del Hombre
\end{multicols}
\end{enumerate}
\vspace{0.15cm}
\item ¿Qué organismo protege los derechos humanos en América?
\begin{enumerate}[label=\textbf{\Alph*)}]
\begin{multicols}{2}
  \item La Comisión Interamericana de Derechos Humanos
  \item La Corte Internacional de Justicia
  \item El Consejo de Seguridad
  \item La Asamblea General
  \item El Consejo Económico y Social
  \item La Secretaría General
\end{multicols}
\end{enumerate}
\vspace{0.15cm}
\item ¿Qué organismo juzga los casos de violaciones a los derechos humanos en América?
\begin{enumerate}[label=\textbf{\Alph*)}]
\begin{multicols}{2}
  \item La Corte Interamericana de Derechos Humanos
  \item La Comisión Interamericana de Derechos Humanos
  \item La Corte Internacional de Justicia
  \item El Consejo de Seguridad
  \item La Asamblea General
  \item El Consejo Económico y Social
\end{multicols}
\end{enumerate}
\vspace{0.15cm}
\item ¿Qué documento de 1979 protege los derechos de las mujeres?
\begin{enumerate}[label=\textbf{\Alph*)}]
\begin{multicols}{2}
  \item La Convención sobre la Eliminación de Todas las Formas de Discriminación contra la Mujer
  \item La Declaración Universal de los Derechos Humanos
  \item El Pacto Internacional de Derechos Civiles
  \item El Pacto Internacional de Derechos Económicos
  \item La Convención Americana de Derechos Humanos
  \item La Carta de las Naciones Unidas
\end{multicols}
\end{enumerate}
\vspace{0.15cm}
\item ¿Qué documento de 1989 protege los derechos de los niños?
\begin{enumerate}[label=\textbf{\Alph*)}]
\begin{multicols}{2}
  \item La Convención sobre los Derechos del Niño
  \item La Declaración Universal de los Derechos Humanos
  \item El Pacto Internacional de Derechos Civiles
  \item El Pacto Internacional de Derechos Económicos
  \item La Convención Americana de Derechos Humanos
  \item La Carta de las Naciones Unidas
\end{multicols}
\end{enumerate}
\vspace{0.15cm}
\item ¿Qué documento de 2006 protege los derechos de las personas con discapacidad?
\begin{enumerate}[label=\textbf{\Alph*)}]
\begin{multicols}{2}
  \item La Convención sobre los Derechos de las Personas con Discapacidad
  \item La Declaración Universal de los Derechos Humanos
  \item El Pacto Internacional de Derechos Civiles
  \item El Pacto Internacional de Derechos Económicos
  \item La Convención Americana de Derechos Humanos
  \item La Carta de las Naciones Unidas
\end{multicols}
\end{enumerate}
\vspace{0.15cm}
\item ¿Qué documento de 2007 protege los derechos de los pueblos indígenas?
\begin{enumerate}[label=\textbf{\Alph*)}]
\begin{multicols}{2}
  \item La Declaración de las Naciones Unidas sobre los Derechos de los Pueblos Indígenas
  \item La Declaración Universal de los Derechos Humanos
  \item El Pacto Internacional de Derechos Civiles
  \item El Pacto Internacional de Derechos Económicos
  \item La Convención Americana de Derechos Humanos
  \item La Carta de las Naciones Unidas
\end{multicols}
\end{enumerate}
\vspace{0.15cm}
\item ¿Qué es la Carta de las Naciones Unidas?
\begin{enumerate}[label=\textbf{\Alph*)}]
\begin{multicols}{2}
  \item El documento fundacional de la ONU que promueve los derechos humanos
  \item Un tratado comercial
  \item Un acuerdo militar
  \item Un pacto económico
  \item Una alianza política
  \item Un convenio cultural
\end{multicols}
\end{enumerate}
\vspace{0.15cm}
\item ¿Qué es el Consejo de Derechos Humanos de la ONU?
\begin{enumerate}[label=\textbf{\Alph*)}]
\begin{multicols}{2}
  \item El organismo que promueve y protege los derechos humanos
  \item El organismo que promueve el comercio
  \item El organismo que promueve la seguridad
  \item El organismo que promueve la economía
  \item El organismo que promueve la cultura
  \item El organismo que promueve la educación
\end{multicols}
\end{enumerate}
\vspace{0.15cm}
\item ¿Qué es el Alto Comisionado de las Naciones Unidas para los Derechos Humanos?
\begin{enumerate}[label=\textbf{\Alph*)}]
\begin{multicols}{2}
  \item La oficina que coordina las actividades de derechos humanos de la ONU
  \item La oficina que coordina el comercio
  \item La oficina que coordina la seguridad
  \item La oficina que coordina la economía
  \item La oficina que coordina la cultura
  \item La oficina que coordina la educación
\end{multicols}
\end{enumerate}
\vspace{0.15cm}
\item ¿Qué es la Corte Penal Internacional?
\begin{enumerate}[label=\textbf{\Alph*)}]
\begin{multicols}{2}
  \item El tribunal que juzga crímenes de guerra y lesa humanidad
  \item El tribunal que juzga delitos comunes
  \item El tribunal que juzga delitos comerciales
  \item El tribunal que juzga delitos políticos
  \item El tribunal que juzga delitos religiosos
  \item El tribunal que juzga delitos culturales
\end{multicols}
\end{enumerate}
\vspace{0.15cm}
\item ¿Qué es el Estatuto de Roma?
\begin{enumerate}[label=\textbf{\Alph*)}]
\begin{multicols}{2}
  \item El tratado que creó la Corte Penal Internacional
  \item El tratado que creó la ONU
  \item El tratado que creó la OEA
  \item El tratado que creó la Unión Europea
  \item El tratado que creó la Unión Africana
  \item El tratado que creó la ASEAN
\end{multicols}
\end{enumerate}
\vspace{0.15cm}
\item ¿Qué es el genocidio?
\begin{enumerate}[label=\textbf{\Alph*)}]
\begin{multicols}{2}
  \item La eliminación sistemática de un grupo humano
  \item La protección de un grupo humano
  \item La integración de un grupo humano
  \item La educación de un grupo humano
  \item La salud de un grupo humano
  \item La vivienda de un grupo humano
\end{multicols}
\end{enumerate}
\vspace{0.15cm}
\item ¿Qué es un crimen de lesa humanidad?
\begin{enumerate}[label=\textbf{\Alph*)}]
\begin{multicols}{2}
  \item Un acto inhumano cometido contra una población civil
  \item Un acto humanitario
  \item Un acto de solidaridad
  \item Un acto de cooperación
  \item Un acto de ayuda
  \item Un acto de apoyo
\end{multicols}
\end{enumerate}
\vspace{0.15cm}
\item ¿Qué es un crimen de guerra?
\begin{enumerate}[label=\textbf{\Alph*)}]
\begin{multicols}{2}
  \item Una violación grave del derecho internacional humanitario
  \item Una violación leve del derecho internacional
  \item Una violación comercial
  \item Una violación política
  \item Una violación religiosa
  \item Una violación cultural
\end{multicols}
\end{enumerate}
\vspace{0.15cm}
\item ¿Qué es el derecho internacional humanitario?
\begin{enumerate}[label=\textbf{\Alph*)}]
\begin{multicols}{2}
  \item El conjunto de normas que protegen a las personas en conflictos armados
  \item El conjunto de normas que protegen el comercio
  \item El conjunto de normas que protegen la economía
  \item El conjunto de normas que protegen la política
  \item El conjunto de normas que protegen la religión
  \item El conjunto de normas que protegen la cultura
\end{multicols}
\end{enumerate}
\vspace{0.15cm}
\item ¿Qué son los Convenios de Ginebra?
\begin{enumerate}[label=\textbf{\Alph*)}]
\begin{multicols}{2}
  \item Los tratados que protegen a las víctimas de conflictos armados
  \item Los tratados que protegen el comercio
  \item Los tratados que protegen la economía
  \item Los tratados que protegen la política
  \item Los tratados que protegen la religión
  \item Los tratados que protegen la cultura
\end{multicols}
\end{enumerate}
\vspace{0.15cm}
\item ¿Qué es el Tribunal de Núremberg?
\begin{enumerate}[label=\textbf{\Alph*)}]
\begin{multicols}{2}
  \item El tribunal que juzgó a los criminales nazis tras la Segunda Guerra Mundial
  \item El tribunal que juzgó a los criminales de guerra japoneses
  \item El tribunal que juzgó a los criminales de guerra italianos
  \item El tribunal que juzgó a los criminales de guerra alemanes
  \item El tribunal que juzgó a los criminales de guerra rusos
  \item El tribunal que juzgó a los criminales de guerra chinos
\end{multicols}
\end{enumerate}
\vspace{0.15cm}
\item ¿Qué es el Tribunal de Tokio?
\begin{enumerate}[label=\textbf{\Alph*)}]
\begin{multicols}{2}
  \item El tribunal que juzgó a los criminales de guerra japoneses
  \item El tribunal que juzgó a los criminales nazis
  \item El tribunal que juzgó a los criminales italianos
  \item El tribunal que juzgó a los criminales alemanes
  \item El tribunal que juzgó a los criminales rusos
  \item El tribunal que juzgó a los criminales chinos
\end{multicols}
\end{enumerate}
\vspace{0.15cm}
\item ¿Qué es la Comisión de la Verdad?
\begin{enumerate}[label=\textbf{\Alph*)}]
\begin{multicols}{2}
  \item El organismo que investiga violaciones a los derechos humanos
  \item El organismo que investiga delitos comunes
  \item El organismo que investiga delitos comerciales
  \item El organismo que investiga delitos políticos
  \item El organismo que investiga delitos religiosos
  \item El organismo que investiga delitos culturales
\end{multicols}
\end{enumerate}
\vspace{0.15cm}
\item ¿Qué es la justicia transicional?
\begin{enumerate}[label=\textbf{\Alph*)}]
\begin{multicols}{2}
  \item El conjunto de medidas para enfrentar violaciones a los derechos humanos
  \item El conjunto de medidas para olvidar violaciones
  \item El conjunto de medidas para negar violaciones
  \item El conjunto de medidas para ocultar violaciones
  \item El conjunto de medidas para justificar violaciones
  \item El conjunto de medidas para perpetuar violaciones
\end{multicols}
\end{enumerate}
\vspace{0.15cm}
\item ¿Qué es la reparación integral?
\begin{enumerate}[label=\textbf{\Alph*)}]
\begin{multicols}{2}
  \item El conjunto de medidas para reparar a las víctimas
  \item El conjunto de medidas para ignorar a las víctimas
  \item El conjunto de medidas para olvidar a las víctimas
  \item El conjunto de medidas para negar a las víctimas
  \item El conjunto de medidas para excluir a las víctimas
  \item El conjunto de medidas para marginar a las víctimas
\end{multicols}
\end{enumerate}
\vspace{0.15cm}
\item ¿Qué es la memoria histórica?
\begin{enumerate}[label=\textbf{\Alph*)}]
\begin{multicols}{2}
  \item La recuperación del pasado para evitar repetir violaciones
  \item La negación del pasado
  \item El olvido del pasado
  \item La ocultación del pasado
  \item La justificación del pasado
  \item La perpetuación del pasado
\end{multicols}
\end{enumerate}
\vspace{0.15cm}
\item ¿Qué es la educación en derechos humanos?
\begin{enumerate}[label=\textbf{\Alph*)}]
\begin{multicols}{2}
  \item El proceso de formar conciencia sobre los derechos humanos
  \item El proceso de formar conciencia sobre el comercio
  \item El proceso de formar conciencia sobre la economía
  \item El proceso de formar conciencia sobre la política
  \item El proceso de formar conciencia sobre la religión
  \item El proceso de formar conciencia sobre la cultura
\end{multicols}
\end{enumerate}
\vspace{0.15cm}
\item ¿Qué es la cultura de paz?
\begin{enumerate}[label=\textbf{\Alph*)}]
\begin{multicols}{2}
  \item El conjunto de valores y actitudes que rechazan la violencia
  \item El conjunto de valores que promueven la violencia
  \item El conjunto de valores que ignoran la violencia
  \item El conjunto de valores que justifican la violencia
  \item El conjunto de valores que perpetúan la violencia
  \item El conjunto de valores que ocultan la violencia
\end{multicols}
\end{enumerate}
\vspace{0.15cm}
\item ¿Qué es la no violencia?
\begin{enumerate}[label=\textbf{\Alph*)}]
\begin{multicols}{2}
  \item La estrategia de lucha que rechaza el uso de la fuerza
  \item La estrategia de lucha que usa la fuerza
  \item La estrategia de lucha que usa la violencia
  \item La estrategia de lucha que usa la agresión
  \item La estrategia de lucha que usa la coerción
  \item La estrategia de lucha que usa la represión
\end{multicols}
\end{enumerate}
\vspace{0.15cm}
\item ¿Qué es la desobediencia civil?
\begin{enumerate}[label=\textbf{\Alph*)}]
\begin{multicols}{2}
  \item La resistencia pacífica a leyes injustas
  \item La obediencia a leyes injustas
  \item La aceptación de leyes injustas
  \item La sumisión a leyes injustas
  \item La conformidad con leyes injustas
  \item La resignación ante leyes injustas
\end{multicols}
\end{enumerate}
\vspace{0.15cm}
\item ¿Qué es la objeción de conciencia?
\begin{enumerate}[label=\textbf{\Alph*)}]
\begin{multicols}{2}
  \item El derecho a negarse a cumplir una obligación por razones de conciencia
  \item El derecho a cumplir una obligación
  \item El derecho a aceptar una obligación
  \item El derecho a asumir una obligación
  \item El derecho a realizar una obligación
  \item El derecho a ejecutar una obligación
\end{multicols}
\end{enumerate}
\vspace{0.15cm}
\item ¿Qué es el hábeas corpus?
\begin{enumerate}[label=\textbf{\Alph*)}]
\begin{multicols}{2}
  \item El derecho a comparecer ante un juez para verificar la legalidad de una detención
  \item El derecho a ser detenido sin causa
  \item El derecho a ser encarcelado sin juicio
  \item El derecho a ser privado de libertad
  \item El derecho a ser arrestado sin orden
  \item El derecho a ser secuestrado
\end{multicols}
\end{enumerate}
\vspace{0.15cm}
\item ¿Qué es el hábeas data?
\begin{enumerate}[label=\textbf{\Alph*)}]
\begin{multicols}{2}
  \item El derecho a acceder y actualizar la información personal
  \item El derecho a acceder a la información pública
  \item El derecho a acceder a la información gubernamental
  \item El derecho a acceder a la información empresarial
  \item El derecho a acceder a la información académica
  \item El derecho a acceder a la información religiosa
\end{multicols}
\end{enumerate}
\vspace{0.15cm}
\item ¿Qué es el amparo?
\begin{enumerate}[label=\textbf{\Alph*)}]
\begin{multicols}{2}
  \item La acción judicial para proteger derechos fundamentales
  \item La acción judicial para proteger intereses comerciales
  \item La acción judicial para proteger intereses económicos
  \item La acción judicial para proteger intereses políticos
  \item La acción judicial para proteger intereses religiosos
  \item La acción judicial para proteger intereses culturales
\end{multicols}
\end{enumerate}
\vspace{0.15cm}
\item ¿Qué es el debido proceso?
\begin{enumerate}[label=\textbf{\Alph*)}]
\begin{multicols}{2}
  \item El conjunto de garantías para un juicio justo
  \item El conjunto de garantías para un juicio arbitrario
  \item El conjunto de garantías para un juicio secreto
  \item El conjunto de garantías para un juicio sin defensa
  \item El conjunto de garantías para un juicio sin pruebas
  \item El conjunto de garantías para un juicio sin juez
\end{multicols}
\end{enumerate}
\vspace{0.15cm}
\item ¿Qué es la tutela judicial efectiva?
\begin{enumerate}[label=\textbf{\Alph*)}]
\begin{multicols}{2}
  \item El derecho a acceder a la justicia y obtener una respuesta
  \item El derecho a acceder a la justicia sin respuesta
  \item El derecho a acceder a la justicia sin juez
  \item El derecho a acceder a la justicia sin abogado
  \item El derecho a acceder a la justicia sin pruebas
  \item El derecho a acceder a la justicia sin ley
\end{multicols}
\end{enumerate}
\vspace{0.15cm}
\item ¿Qué es la independencia judicial?
\begin{enumerate}[label=\textbf{\Alph*)}]
\begin{multicols}{2}
  \item La autonomía de los jueces para decidir sin presiones
  \item La dependencia de los jueces del poder político
  \item La dependencia de los jueces del poder económico
  \item La dependencia de los jueces del poder militar
  \item La dependencia de los jueces del poder religioso
  \item La dependencia de los jueces del poder cultural
\end{multicols}
\end{enumerate}
\vspace{0.15cm}
\item ¿Qué es la imparcialidad judicial?
\begin{enumerate}[label=\textbf{\Alph*)}]
\begin{multicols}{2}
  \item La obligación de los jueces de decidir sin favoritismos
  \item La obligación de los jueces de decidir con favoritismos
  \item La obligación de los jueces de decidir con prejuicios
  \item La obligación de los jueces de decidir con sesgos
  \item La obligación de los jueces de decidir con intereses
  \item La obligación de los jueces de decidir con presiones
\end{multicols}
\end{enumerate}
\vspace{0.15cm}
\item ¿Qué es la cosa juzgada?
\begin{enumerate}[label=\textbf{\Alph*)}]
\begin{multicols}{2}
  \item La imposibilidad de volver a juzgar un caso ya resuelto
  \item La posibilidad de volver a juzgar un caso
  \item La posibilidad de revisar un caso indefinidamente
  \item La posibilidad de apelar indefinidamente
  \item La posibilidad de reabrir un caso cerrado
  \item La posibilidad de anular una sentencia
\end{multicols}
\end{enumerate}
\vspace{0.15cm}
\item ¿Qué es la prescripción?
\begin{enumerate}[label=\textbf{\Alph*)}]
\begin{multicols}{2}
  \item La extinción de la acción penal por el paso del tiempo
  \item La extinción de la acción civil por el paso del tiempo
  \item La extinción de la acción administrativa por el paso del tiempo
  \item La extinción de la acción laboral por el paso del tiempo
  \item La extinción de la acción comercial por el paso del tiempo
  \item La extinción de la acción fiscal por el paso del tiempo
\end{multicols}
\end{enumerate}
\vspace{0.15cm}
\item ¿Qué es la imprescriptibilidad?
\begin{enumerate}[label=\textbf{\Alph*)}]
\begin{multicols}{2}
  \item La imposibilidad de que ciertos delitos prescriban
  \item La posibilidad de que ciertos delitos prescriban
  \item La posibilidad de que ciertos delitos se archiven
  \item La posibilidad de que ciertos delitos se olviden
  \item La posibilidad de que ciertos delitos se perdonen
  \item La posibilidad de que ciertos delitos se amnistíen
\end{multicols}
\end{enumerate}
\vspace{0.15cm}
\end{enumerate}\n
\subsubsection*{Sección III: Tratados Internacionales de Derechos Humanos (101–150)}

\begin{enumerate}[label=\textbf{\arabic*.}]
\setcounter{enumi}{100}
\item ¿Qué es un tratado internacional?
\begin{enumerate}[label=\textbf{\Alph*)}]
\begin{multicols}{2}
  \item Un acuerdo escrito entre Estados regido por el derecho internacional
  \item Un acuerdo verbal entre Estados
  \item Un acuerdo informal entre Estados
  \item Un acuerdo comercial entre empresas
  \item Un acuerdo político entre partidos
  \item Un acuerdo religioso entre iglesias
\end{multicols}
\end{enumerate}
\vspace{0.15cm}
\item ¿Qué es la ratificación de un tratado?
\begin{enumerate}[label=\textbf{\Alph*)}]
\begin{multicols}{2}
  \item El acto por el cual un Estado se obliga a cumplir el tratado
  \item El acto por el cual un Estado rechaza el tratado
  \item El acto por el cual un Estado ignora el tratado
  \item El acto por el cual un Estado olvida el tratado
  \item El acto por el cual un Estado niega el tratado
  \item El acto por el cual un Estado oculta el tratado
\end{multicols}
\end{enumerate}
\vspace{0.15cm}
\item ¿Qué es la adhesión a un tratado?
\begin{enumerate}[label=\textbf{\Alph*)}]
\begin{multicols}{2}
  \item El acto por el cual un Estado se suma a un tratado ya existente
  \item El acto por el cual un Estado se retira de un tratado
  \item El acto por el cual un Estado rechaza un tratado
  \item El acto por el cual un Estado ignora un tratado
  \item El acto por el cual un Estado olvida un tratado
  \item El acto por el cual un Estado niega un tratado
\end{multicols}
\end{enumerate}
\vspace{0.15cm}
\item ¿Qué es la reserva en un tratado?
\begin{enumerate}[label=\textbf{\Alph*)}]
\begin{multicols}{2}
  \item La declaración por la cual un Estado excluye o modifica una disposición
  \item La declaración por la cual un Estado acepta todo el tratado
  \item La declaración por la cual un Estado rechaza todo el tratado
  \item La declaración por la cual un Estado ignora todo el tratado
  \item La declaración por la cual un Estado olvida todo el tratado
  \item La declaración por la cual un Estado niega todo el tratado
\end{multicols}
\end{enumerate}
\vspace{0.15cm}
\item ¿Qué es la denuncia de un tratado?
\begin{enumerate}[label=\textbf{\Alph*)}]
\begin{multicols}{2}
  \item El acto por el cual un Estado se retira de un tratado
  \item El acto por el cual un Estado se suma a un tratado
  \item El acto por el cual un Estado acepta un tratado
  \item El acto por el cual un Estado ratifica un tratado
  \item El acto por el cual un Estado firma un tratado
  \item El acto por el cual un Estado negocia un tratado
\end{multicols}
\end{enumerate}
\vspace{0.15cm}
\item ¿Qué es el Sistema Universal de Protección de los Derechos Humanos?
\begin{enumerate}[label=\textbf{\Alph*)}]
\begin{multicols}{2}
  \item El conjunto de mecanismos de la ONU para proteger los derechos humanos
  \item El conjunto de mecanismos de la OEA
  \item El conjunto de mecanismos de la Unión Europea
  \item El conjunto de mecanismos de la Unión Africana
  \item El conjunto de mecanismos de la Liga Árabe
  \item El conjunto de mecanismos de la ASEAN
\end{multicols}
\end{enumerate}
\vspace{0.15cm}
\item ¿Qué es el Sistema Interamericano de Protección de los Derechos Humanos?
\begin{enumerate}[label=\textbf{\Alph*)}]
\begin{multicols}{2}
  \item El conjunto de mecanismos de la OEA para proteger los derechos humanos
  \item El conjunto de mecanismos de la ONU
  \item El conjunto de mecanismos de la Unión Europea
  \item El conjunto de mecanismos de la Unión Africana
  \item El conjunto de mecanismos de la Liga Árabe
  \item El conjunto de mecanismos de la ASEAN
\end{multicols}
\end{enumerate}
\vspace{0.15cm}
\item ¿Qué es la Comisión Interamericana de Derechos Humanos?
\begin{enumerate}[label=\textbf{\Alph*)}]
\begin{multicols}{2}
  \item El órgano de la OEA que promueve y protege los derechos humanos
  \item El órgano de la ONU que promueve los derechos humanos
  \item El órgano de la Unión Europea que promueve los derechos humanos
  \item El órgano de la Unión Africana que promueve los derechos humanos
  \item El órgano de la Liga Árabe que promueve los derechos humanos
  \item El órgano de la ASEAN que promueve los derechos humanos
\end{multicols}
\end{enumerate}
\vspace{0.15cm}
\item ¿Qué es la Corte Interamericana de Derechos Humanos?
\begin{enumerate}[label=\textbf{\Alph*)}]
\begin{multicols}{2}
  \item El tribunal de la OEA que juzga violaciones a los derechos humanos
  \item El tribunal de la ONU que juzga violaciones a los derechos humanos
  \item El tribunal de la Unión Europea que juzga violaciones
  \item El tribunal de la Unión Africana que juzga violaciones
  \item El tribunal de la Liga Árabe que juzga violaciones
  \item El tribunal de la ASEAN que juzga violaciones
\end{multicols}
\end{enumerate}
\vspace{0.15cm}
\item ¿Dónde tiene su sede la Corte Interamericana de Derechos Humanos?
\begin{enumerate}[label=\textbf{\Alph*)}]
\begin{multicols}{2}
  \item San José, Costa Rica
  \item Washington, Estados Unidos
  \item Ginebra, Suiza
  \item Estrasburgo, Francia
  \item La Haya, Países Bajos
  \item Bruselas, Bélgica
\end{multicols}
\end{enumerate}
\vspace{0.15cm}
\item ¿Qué es el Pacto de San José de Costa Rica?
\begin{enumerate}[label=\textbf{\Alph*)}]
\begin{multicols}{2}
  \item La Convención Americana sobre Derechos Humanos
  \item La Declaración Universal de los Derechos Humanos
  \item El Pacto Internacional de Derechos Civiles
  \item El Pacto Internacional de Derechos Económicos
  \item La Carta de las Naciones Unidas
  \item La Declaración de los Derechos del Hombre
\end{multicols}
\end{enumerate}
\vspace{0.15cm}
\item ¿Qué es el Protocolo de San Salvador?
\begin{enumerate}[label=\textbf{\Alph*)}]
\begin{multicols}{2}
  \item El protocolo adicional a la Convención Americana sobre derechos económicos, sociales y culturales
  \item El protocolo adicional a la Convención Americana sobre derechos civiles
  \item El protocolo adicional a la Convención Americana sobre derechos políticos
  \item El protocolo adicional a la Convención Americana sobre derechos ambientales
  \item El protocolo adicional a la Convención Americana sobre derechos digitales
  \item El protocolo adicional a la Convención Americana sobre derechos colectivos
\end{multicols}
\end{enumerate}
\vspace{0.15cm}
\item ¿Qué es la Comisión de la Verdad para El Salvador?
\begin{enumerate}[label=\textbf{\Alph*)}]
\begin{multicols}{2}
  \item El organismo que investigó las violaciones a los derechos humanos durante la guerra civil
  \item El organismo que investigó las violaciones durante la colonia
  \item El organismo que investigó las violaciones durante la independencia
  \item El organismo que investigó las violaciones durante la federación
  \item El organismo que investigó las violaciones durante el siglo XIX
  \item El organismo que investigó las violaciones durante el siglo XX
\end{multicols}
\end{enumerate}
\vspace{0.15cm}
\item ¿Qué documento surgió de la Comisión de la Verdad para El Salvador?
\begin{enumerate}[label=\textbf{\Alph*)}]
\begin{multicols}{2}
  \item El informe "De la locura a la esperanza"
  \item El informe "Nunca más"
  \item El informe "Basta ya"
  \item El informe "Guatemala: memoria del silencio"
  \item El informe "Chile: la memoria obstinada"
  \item El informe "Argentina: la verdad"
\end{multicols}
\end{enumerate}
\vspace{0.15cm}
\item ¿Qué es la Oficina de Derechos Humanos de la ONU en El Salvador?
\begin{enumerate}[label=\textbf{\Alph*)}]
\begin{multicols}{2}
  \item La oficina que monitorea la situación de los derechos humanos
  \item La oficina que monitorea la situación económica
  \item La oficina que monitorea la situación política
  \item La oficina que monitorea la situación social
  \item La oficina que monitorea la situación cultural
  \item La oficina que monitorea la situación ambiental
\end{multicols}
\end{enumerate}
\vspace{0.15cm}
\item ¿Qué es la Procuraduría para la Defensa de los Derechos Humanos de El Salvador?
\begin{enumerate}[label=\textbf{\Alph*)}]
\begin{multicols}{2}
  \item La institución que protege y promueve los derechos humanos
  \item La institución que protege los derechos comerciales
  \item La institución que protege los derechos económicos
  \item La institución que protege los derechos políticos
  \item La institución que protege los derechos religiosos
  \item La institución que protege los derechos culturales
\end{multicols}
\end{enumerate}
\vspace{0.15cm}
\item ¿Qué es la Convención contra la Tortura?
\begin{enumerate}[label=\textbf{\Alph*)}]
\begin{multicols}{2}
  \item El tratado que prohíbe la tortura y otros tratos crueles
  \item El tratado que permite la tortura
  \item El tratado que regula la tortura
  \item El tratado que limita la tortura
  \item El tratado que condiciona la tortura
  \item El tratado que autoriza la tortura
\end{multicols}
\end{enumerate}
\vspace{0.15cm}
\item ¿Qué es la Convención sobre los Desaparecidos?
\begin{enumerate}[label=\textbf{\Alph*)}]
\begin{multicols}{2}
  \item El tratado que protege a las personas contra las desapariciones forzadas
  \item El tratado que permite las desapariciones forzadas
  \item El tratado que regula las desapariciones forzadas
  \item El tratado que limita las desapariciones forzadas
  \item El tratado que condiciona las desapariciones forzadas
  \item El tratado que autoriza las desapariciones forzadas
\end{multicols}
\end{enumerate}
\vspace{0.15cm}
\item ¿Qué es el Estatuto de Roma?
\begin{enumerate}[label=\textbf{\Alph*)}]
\begin{multicols}{2}
  \item El tratado que creó la Corte Penal Internacional
  \item El tratado que creó la ONU
  \item El tratado que creó la OEA
  \item El tratado que creó la Unión Europea
  \item El tratado que creó la Unión Africana
  \item El tratado que creó la ASEAN
\end{multicols}
\end{enumerate}
\vspace{0.15cm}
\item ¿Qué es la Corte Penal Internacional?
\begin{enumerate}[label=\textbf{\Alph*)}]
\begin{multicols}{2}
  \item El tribunal que juzga genocidio, crímenes de guerra y lesa humanidad
  \item El tribunal que juzga delitos comunes
  \item El tribunal que juzga delitos comerciales
  \item El tribunal que juzga delitos políticos
  \item El tribunal que juzga delitos religiosos
  \item El tribunal que juzga delitos culturales
\end{multicols}
\end{enumerate}
\vspace{0.15cm}
\item ¿Qué es el genocidio según la Convención de 1948?
\begin{enumerate}[label=\textbf{\Alph*)}]
\begin{multicols}{2}
  \item La matanza o eliminación de un grupo nacional, étnico, racial o religioso
  \item La matanza de un grupo político
  \item La matanza de un grupo económico
  \item La matanza de un grupo social
  \item La matanza de un grupo cultural
  \item La matanza de un grupo profesional
\end{multicols}
\end{enumerate}
\vspace{0.15cm}
\item ¿Qué es el apartheid?
\begin{enumerate}[label=\textbf{\Alph*)}]
\begin{multicols}{2}
  \item La política de segregación racial institucionalizada
  \item La política de integración racial
  \item La política de igualdad racial
  \item La política de inclusión racial
  \item La política de diversidad racial
  \item La política de convivencia racial
\end{multicols}
\end{enumerate}
\vspace{0.15cm}
\item ¿Qué es la discriminación racial?
\begin{enumerate}[label=\textbf{\Alph*)}]
\begin{multicols}{2}
  \item El trato desigual por motivos de raza o etnia
  \item El trato igualitario por motivos de raza
  \item El trato justo por motivos de raza
  \item El trato objetivo por motivos de raza
  \item El trato científico por motivos de raza
  \item El trato artístico por motivos de raza
\end{multicols}
\end{enumerate}
\vspace{0.15cm}
\item ¿Qué es la Convención contra la Discriminación Racial?
\begin{enumerate}[label=\textbf{\Alph*)}]
\begin{multicols}{2}
  \item El tratado que prohíbe la discriminación racial
  \item El tratado que permite la discriminación racial
  \item El tratado que regula la discriminación racial
  \item El tratado que limita la discriminación racial
  \item El tratado que condiciona la discriminación racial
  \item El tratado que autoriza la discriminación racial
\end{multicols}
\end{enumerate}
\vspace{0.15cm}
\item ¿Qué es la Convención contra la Discriminación contra la Mujer?
\begin{enumerate}[label=\textbf{\Alph*)}]
\begin{multicols}{2}
  \item El tratado que prohíbe la discriminación contra la mujer
  \item El tratado que permite la discriminación contra la mujer
  \item El tratado que regula la discriminación contra la mujer
  \item El tratado que limita la discriminación contra la mujer
  \item El tratado que condiciona la discriminación contra la mujer
  \item El tratado que autoriza la discriminación contra la mujer
\end{multicols}
\end{enumerate}
\vspace{0.15cm}
\item ¿Qué es la Convención sobre los Derechos del Niño?
\begin{enumerate}[label=\textbf{\Alph*)}]
\begin{multicols}{2}
  \item El tratado que protege los derechos de los menores de 18 años
  \item El tratado que protege los derechos de los mayores de 18 años
  \item El tratado que protege los derechos de los adultos
  \item El tratado que protege los derechos de los ancianos
  \item El tratado que protege los derechos de los discapacitados
  \item El tratado que protege los derechos de los indígenas
\end{multicols}
\end{enumerate}
\vspace{0.15cm}
\item ¿Qué es el interés superior del niño?
\begin{enumerate}[label=\textbf{\Alph*)}]
\begin{multicols}{2}
  \item El principio que orienta las decisiones sobre los niños
  \item El principio que orienta las decisiones sobre los adultos
  \item El principio que orienta las decisiones sobre los ancianos
  \item El principio que orienta las decisiones sobre los discapacitados
  \item El principio que orienta las decisiones sobre los indígenas
  \item El principio que orienta las decisiones sobre los migrantes
\end{multicols}
\end{enumerate}
\vspace{0.15cm}
\item ¿Qué es la Convención sobre los Derechos de las Personas con Discapacidad?
\begin{enumerate}[label=\textbf{\Alph*)}]
\begin{multicols}{2}
  \item El tratado que protege los derechos de las personas con discapacidad
  \item El tratado que protege los derechos de las personas sin discapacidad
  \item El tratado que protege los derechos de los niños
  \item El tratado que protege los derechos de las mujeres
  \item El tratado que protege los derechos de los indígenas
  \item El tratado que protege los derechos de los migrantes
\end{multicols}
\end{enumerate}
\vspace{0.15cm}
\item ¿Qué es la accesibilidad?
\begin{enumerate}[label=\textbf{\Alph*)}]
\begin{multicols}{2}
  \item La eliminación de barreras para la participación de las personas con discapacidad
  \item La creación de barreras para las personas con discapacidad
  \item La regulación de barreras para las personas con discapacidad
  \item La limitación de barreras para las personas con discapacidad
  \item La condición de barreras para las personas con discapacidad
  \item La autorización de barreras para las personas con discapacidad
\end{multicols}
\end{enumerate}
\vspace{0.15cm}
\item ¿Qué es la Declaración de las Naciones Unidas sobre los Derechos de los Pueblos Indígenas?
\begin{enumerate}[label=\textbf{\Alph*)}]
\begin{multicols}{2}
  \item El documento que protege los derechos colectivos de los pueblos indígenas
  \item El documento que protege los derechos individuales de los indígenas
  \item El documento que protege los derechos comerciales de los indígenas
  \item El documento que protege los derechos económicos de los indígenas
  \item El documento que protege los derechos políticos de los indígenas
  \item El documento que protege los derechos religiosos de los indígenas
\end{multicols}
\end{enumerate}
\vspace{0.15cm}
\item ¿Qué es el Convenio 169 de la OIT?
\begin{enumerate}[label=\textbf{\Alph*)}]
\begin{multicols}{2}
  \item El tratado que protege los derechos de los pueblos indígenas y tribales
  \item El tratado que protege los derechos de los trabajadores migrantes
  \item El tratado que protege los derechos de las mujeres trabajadoras
  \item El tratado que protege los derechos de los niños trabajadores
  \item El tratado que protege los derechos de los trabajadores discapacitados
  \item El tratado que protege los derechos de los trabajadores informales
\end{multicols}
\end{enumerate}
\vspace{0.15cm}
\item ¿Qué es la consulta previa?
\begin{enumerate}[label=\textbf{\Alph*)}]
\begin{multicols}{2}
  \item El derecho de los pueblos indígenas a ser consultados antes de decisiones que les afecten
  \item El derecho de los pueblos indígenas a ser informados después de las decisiones
  \item El derecho de los pueblos indígenas a ser ignorados
  \item El derecho de los pueblos indígenas a ser excluidos
  \item El derecho de los pueblos indígenas a ser marginados
  \item El derecho de los pueblos indígenas a ser discriminados
\end{multicols}
\end{enumerate}
\vspace{0.15cm}
\item ¿Qué es el consentimiento libre, previo e informado?
\begin{enumerate}[label=\textbf{\Alph*)}]
\begin{multicols}{2}
  \item El derecho de los pueblos indígenas a dar o negar su consentimiento
  \item El derecho de los pueblos indígenas a ser ignorados
  \item El derecho de los pueblos indígenas a ser excluidos
  \item El derecho de los pueblos indígenas a ser marginados
  \item El derecho de los pueblos indígenas a ser discriminados
  \item El derecho de los pueblos indígenas a ser olvidados
\end{multicols}
\end{enumerate}
\vspace{0.15cm}
\item ¿Qué es la Convención sobre los Trabajadores Migrantes?
\begin{enumerate}[label=\textbf{\Alph*)}]
\begin{multicols}{2}
  \item El tratado que protege los derechos de los trabajadores migrantes y sus familias
  \item El tratado que protege los derechos de los trabajadores locales
  \item El tratado que protege los derechos de los trabajadores informales
  \item El tratado que protege los derechos de los trabajadores formales
  \item El tratado que protege los derechos de los trabajadores domésticos
  \item El tratado que protege los derechos de los trabajadores agrícolas
\end{multicols}
\end{enumerate}
\vspace{0.15cm}
\item ¿Qué es la Convención contra la Desaparición Forzada?
\begin{enumerate}[label=\textbf{\Alph*)}]
\begin{multicols}{2}
  \item El tratado que protege contra las desapariciones forzadas
  \item El tratado que permite las desapariciones forzadas
  \item El tratado que regula las desapariciones forzadas
  \item El tratado que limita las desapariciones forzadas
  \item El tratado que condiciona las desapariciones forzadas
  \item El tratado que autoriza las desapariciones forzadas
\end{multicols}
\end{enumerate}
\vspace{0.15cm}
\item ¿Qué es la Corte Africana de Derechos Humanos?
\begin{enumerate}[label=\textbf{\Alph*)}]
\begin{multicols}{2}
  \item El tribunal regional que protege los derechos humanos en África
  \item El tribunal regional que protege los derechos humanos en América
  \item El tribunal regional que protege los derechos humanos en Europa
  \item El tribunal regional que protege los derechos humanos en Asia
  \item El tribunal regional que protege los derechos humanos en Oceanía
  \item El tribunal regional que protege los derechos humanos en la Antártida
\end{multicols}
\end{enumerate}
\vspace{0.15cm}
\item ¿Qué es el Tribunal Europeo de Derechos Humanos?
\begin{enumerate}[label=\textbf{\Alph*)}]
\begin{multicols}{2}
  \item El tribunal regional que protege los derechos humanos en Europa
  \item El tribunal regional que protege los derechos humanos en América
  \item El tribunal regional que protege los derechos humanos en África
  \item El tribunal regional que protege los derechos humanos en Asia
  \item El tribunal regional que protege los derechos humanos en Oceanía
  \item El tribunal regional que protege los derechos humanos en la Antártida
\end{multicols}
\end{enumerate}
\vspace{0.15cm}
\item ¿Qué es la Carta Africana de Derechos Humanos?
\begin{enumerate}[label=\textbf{\Alph*)}]
\begin{multicols}{2}
  \item El tratado regional que protege los derechos humanos en África
  \item El tratado regional que protege los derechos humanos en América
  \item El tratado regional que protege los derechos humanos en Europa
  \item El tratado regional que protege los derechos humanos en Asia
  \item El tratado regional que protege los derechos humanos en Oceanía
  \item El tratado regional que protege los derechos humanos en la Antártida
\end{multicols}
\end{enumerate}
\vspace{0.15cm}
\item ¿Qué es la Carta Árabe de Derechos Humanos?
\begin{enumerate}[label=\textbf{\Alph*)}]
\begin{multicols}{2}
  \item El tratado regional que protege los derechos humanos en el mundo árabe
  \item El tratado regional que protege los derechos humanos en América
  \item El tratado regional que protege los derechos humanos en Europa
  \item El tratado regional que protege los derechos humanos en África
  \item El tratado regional que protege los derechos humanos en Asia
  \item El tratado regional que protege los derechos humanos en Oceanía
\end{multicols}
\end{enumerate}
\vspace{0.15cm}
\item ¿Qué es la Declaración de los Derechos Humanos en el Islam?
\begin{enumerate}[label=\textbf{\Alph*)}]
\begin{multicols}{2}
  \item El documento que protege los derechos humanos según la tradición islámica
  \item El documento que protege los derechos humanos según la tradición cristiana
  \item El documento que protege los derechos humanos según la tradición judía
  \item El documento que protege los derechos humanos según la tradición hindú
  \item El documento que protege los derechos humanos según la tradición budista
  \item El documento que protege los derechos humanos según la tradición laica
\end{multicols}
\end{enumerate}
\vspace{0.15cm}
\item ¿Qué es el Sistema Europeo de Protección de los Derechos Humanos?
\begin{enumerate}[label=\textbf{\Alph*)}]
\begin{multicols}{2}
  \item El conjunto de mecanismos del Consejo de Europa
  \item El conjunto de mecanismos de la ONU
  \item El conjunto de mecanismos de la OEA
  \item El conjunto de mecanismos de la Unión Africana
  \item El conjunto de mecanismos de la Liga Árabe
  \item El conjunto de mecanismos de la ASEAN
\end{multicols}
\end{enumerate}
\vspace{0.15cm}
\item ¿Qué es el Sistema Africano de Protección de los Derechos Humanos?
\begin{enumerate}[label=\textbf{\Alph*)}]
\begin{multicols}{2}
  \item El conjunto de mecanismos de la Unión Africana
  \item El conjunto de mecanismos de la ONU
  \item El conjunto de mecanismos de la OEA
  \item El conjunto de mecanismos del Consejo de Europa
  \item El conjunto de mecanismos de la Liga Árabe
  \item El conjunto de mecanismos de la ASEAN
\end{multicols}
\end{enumerate}
\vspace{0.15cm}
\item ¿Qué es el Sistema Asiático de Protección de los Derechos Humanos?
\begin{enumerate}[label=\textbf{\Alph*)}]
\begin{multicols}{2}
  \item El conjunto de mecanismos de la ASEAN
  \item El conjunto de mecanismos de la ONU
  \item El conjunto de mecanismos de la OEA
  \item El conjunto de mecanismos del Consejo de Europa
  \item El conjunto de mecanismos de la Unión Africana
  \item El conjunto de mecanismos de la Liga Árabe
\end{multicols}
\end{enumerate}
\vspace{0.15cm}
\item ¿Qué es la ASEAN?
\begin{enumerate}[label=\textbf{\Alph*)}]
\begin{multicols}{2}
  \item La Asociación de Naciones del Sudeste Asiático
  \item La Asociación de Naciones del Sur de Asia
  \item La Asociación de Naciones del Este de Asia
  \item La Asociación de Naciones del Norte de Asia
  \item La Asociación de Naciones del Oeste de Asia
  \item La Asociación de Naciones del Centro de Asia
\end{multicols}
\end{enumerate}
\vspace{0.15cm}
\item ¿Qué es la Liga Árabe?
\begin{enumerate}[label=\textbf{\Alph*)}]
\begin{multicols}{2}
  \item La organización de países árabes
  \item La organización de países africanos
  \item La organización de países asiáticos
  \item La organización de países europeos
  \item La organización de países americanos
  \item La organización de países oceánicos
\end{multicols}
\end{enumerate}
\vspace{0.15cm}
\item ¿Qué es la Unión Africana?
\begin{enumerate}[label=\textbf{\Alph*)}]
\begin{multicols}{2}
  \item La organización de países africanos
  \item La organización de países árabes
  \item La organización de países asiáticos
  \item La organización de países europeos
  \item La organización de países americanos
  \item La organización de países oceánicos
\end{multicols}
\end{enumerate}
\vspace{0.15cm}
\item ¿Qué es el Consejo de Europa?
\begin{enumerate}[label=\textbf{\Alph*)}]
\begin{multicols}{2}
  \item La organización de países europeos para la cooperación y los derechos humanos
  \item La organización de países americanos
  \item La organización de países africanos
  \item La organización de países asiáticos
  \item La organización de países árabes
  \item La organización de países oceánicos
\end{multicols}
\end{enumerate}
\vspace{0.15cm}
\item ¿Qué es la OSCE?
\begin{enumerate}[label=\textbf{\Alph*)}]
\begin{multicols}{2}
  \item La Organización para la Seguridad y la Cooperación en Europa
  \item La Organización para la Seguridad y la Cooperación en América
  \item La Organización para la Seguridad y la Cooperación en África
  \item La Organización para la Seguridad y la Cooperación en Asia
  \item La Organización para la Seguridad y la Cooperación en el mundo árabe
  \item La Organización para la Seguridad y la Cooperación en Oceanía
\end{multicols}
\end{enumerate}
\vspace{0.15cm}
\item ¿Qué es el Consejo de Derechos Humanos de la ONU?
\begin{enumerate}[label=\textbf{\Alph*)}]
\begin{multicols}{2}
  \item El organismo intergubernamental que promueve los derechos humanos
  \item El organismo intergubernamental que promueve el comercio
  \item El organismo intergubernamental que promueve la seguridad
  \item El organismo intergubernamental que promueve la economía
  \item El organismo intergubernamental que promueve la cultura
  \item El organismo intergubernamental que promueve la educación
\end{multicols}
\end{enumerate}
\vspace{0.15cm}
\item ¿Qué es el Examen Periódico Universal?
\begin{enumerate}[label=\textbf{\Alph*)}]
\begin{multicols}{2}
  \item El mecanismo que evalúa la situación de derechos humanos de todos los países
  \item El mecanismo que evalúa la situación económica de todos los países
  \item El mecanismo que evalúa la situación política de todos los países
  \item El mecanismo que evalúa la situación social de todos los países
  \item El mecanismo que evalúa la situación cultural de todos los países
  \item El mecanismo que evalúa la situación ambiental de todos los países
\end{multicols}
\end{enumerate}
\vspace{0.15cm}
\end{enumerate}\n
\subsubsection*{Sección IV: Sociedad de la Información y Tecnologías Digitales (151–210)}

\begin{enumerate}[label=\textbf{\arabic*.}]
\setcounter{enumi}{150}
\item ¿Qué es la sociedad de la información?
\begin{enumerate}[label=\textbf{\Alph*)}]
\begin{multicols}{2}
  \item Una sociedad en la que la información y el conocimiento son centrales
  \item Una sociedad en la que la agricultura es central
  \item Una sociedad en la que la industria es central
  \item Una sociedad en la que el comercio es central
  \item Una sociedad en la que la minería es central
  \item Una sociedad en la que la pesca es central
\end{multicols}
\end{enumerate}
\vspace{0.15cm}
\item ¿Qué son las tecnologías de la información y la comunicación (TIC)?
\begin{enumerate}[label=\textbf{\Alph*)}]
\begin{multicols}{2}
  \item El conjunto de tecnologías para procesar y transmitir información
  \item El conjunto de tecnologías para la agricultura
  \item El conjunto de tecnologías para la industria
  \item El conjunto de tecnologías para el comercio
  \item El conjunto de tecnologías para la minería
  \item El conjunto de tecnologías para la pesca
\end{multicols}
\end{enumerate}
\vspace{0.15cm}
\item ¿Qué es internet?
\begin{enumerate}[label=\textbf{\Alph*)}]
\begin{multicols}{2}
  \item Una red global de computadoras interconectadas
  \item Una red local de computadoras
  \item Una red regional de computadoras
  \item Una red nacional de computadoras
  \item Una red continental de computadoras
  \item Una red empresarial de computadoras
\end{multicols}
\end{enumerate}
\vspace{0.15cm}
\item ¿Qué es la World Wide Web?
\begin{enumerate}[label=\textbf{\Alph*)}]
\begin{multicols}{2}
  \item El sistema de documentos interconectados accesibles por internet
  \item El sistema de correo electrónico
  \item El sistema de mensajería instantánea
  \item El sistema de videoconferencia
  \item El sistema de transferencia de archivos
  \item El sistema de bases de datos
\end{multicols}
\end{enumerate}
\vspace{0.15cm}
\item ¿Qué es una red social?
\begin{enumerate}[label=\textbf{\Alph*)}]
\begin{multicols}{2}
  \item Una plataforma digital para la interacción social
  \item Una plataforma digital para el comercio
  \item Una plataforma digital para la educación
  \item Una plataforma digital para la salud
  \item Una plataforma digital para el gobierno
  \item Una plataforma digital para la banca
\end{multicols}
\end{enumerate}
\vspace{0.15cm}
\item ¿Qué es la brecha digital?
\begin{enumerate}[label=\textbf{\Alph*)}]
\begin{multicols}{2}
  \item La desigualdad en el acceso y uso de las tecnologías
  \item La igualdad en el acceso a la tecnología
  \item La abundancia de tecnología
  \item La escasez de tecnología
  \item La distribución equitativa de tecnología
  \item La democratización de la tecnología
\end{multicols}
\end{enumerate}
\vspace{0.15cm}
\item ¿Qué es la alfabetización digital?
\begin{enumerate}[label=\textbf{\Alph*)}]
\begin{multicols}{2}
  \item La capacidad de usar las tecnologías de la información
  \item La capacidad de usar la agricultura
  \item La capacidad de usar la industria
  \item La capacidad de usar el comercio
  \item La capacidad de usar la minería
  \item La capacidad de usar la pesca
\end{multicols}
\end{enumerate}
\vspace{0.15cm}
\item ¿Qué es la ciudadanía digital?
\begin{enumerate}[label=\textbf{\Alph*)}]
\begin{multicols}{2}
  \item El ejercicio de derechos y deberes en el entorno digital
  \item El ejercicio de derechos y deberes en el entorno físico
  \item El ejercicio de derechos y deberes en el entorno rural
  \item El ejercicio de derechos y deberes en el entorno urbano
  \item El ejercicio de derechos y deberes en el entorno laboral
  \item El ejercicio de derechos y deberes en el entorno escolar
\end{multicols}
\end{enumerate}
\vspace{0.15cm}
\item ¿Qué es la identidad digital?
\begin{enumerate}[label=\textbf{\Alph*)}]
\begin{multicols}{2}
  \item La representación de una persona en el entorno digital
  \item La representación de una persona en el entorno físico
  \item La representación de una persona en el entorno rural
  \item La representación de una persona en el entorno urbano
  \item La representación de una persona en el entorno laboral
  \item La representación de una persona en el entorno escolar
\end{multicols}
\end{enumerate}
\vspace{0.15cm}
\item ¿Qué es la huella digital?
\begin{enumerate}[label=\textbf{\Alph*)}]
\begin{multicols}{2}
  \item El rastro que deja una persona al usar internet
  \item El rastro que deja una persona al caminar
  \item El rastro que deja una persona al escribir
  \item El rastro que deja una persona al hablar
  \item El rastro que deja una persona al comer
  \item El rastro que deja una persona al dormir
\end{multicols}
\end{enumerate}
\vspace{0.15cm}
\item ¿Qué es la privacidad digital?
\begin{enumerate}[label=\textbf{\Alph*)}]
\begin{multicols}{2}
  \item El derecho a controlar la información personal en el entorno digital
  \item El derecho a controlar la información pública
  \item El derecho a controlar la información gubernamental
  \item El derecho a controlar la información empresarial
  \item El derecho a controlar la información académica
  \item El derecho a controlar la información religiosa
\end{multicols}
\end{enumerate}
\vspace{0.15cm}
\item ¿Qué es la protección de datos personales?
\begin{enumerate}[label=\textbf{\Alph*)}]
\begin{multicols}{2}
  \item El derecho a proteger la información personal
  \item El derecho a proteger la información pública
  \item El derecho a proteger la información gubernamental
  \item El derecho a proteger la información empresarial
  \item El derecho a proteger la información académica
  \item El derecho a proteger la información religiosa
\end{multicols}
\end{enumerate}
\vspace{0.15cm}
\item ¿Qué es el consentimiento informado en el entorno digital?
\begin{enumerate}[label=\textbf{\Alph*)}]
\begin{multicols}{2}
  \item La autorización del usuario para el tratamiento de sus datos
  \item La autorización del gobierno para el tratamiento de datos
  \item La autorización de la empresa para el tratamiento de datos
  \item La autorización de la escuela para el tratamiento de datos
  \item La autorización de la iglesia para el tratamiento de datos
  \item La autorización de la familia para el tratamiento de datos
\end{multicols}
\end{enumerate}
\vspace{0.15cm}
\item ¿Qué es la ciberseguridad?
\begin{enumerate}[label=\textbf{\Alph*)}]
\begin{multicols}{2}
  \item La protección de sistemas informáticos y datos contra ataques
  \item La protección de sistemas agrícolas
  \item La protección de sistemas industriales
  \item La protección de sistemas comerciales
  \item La protección de sistemas mineros
  \item La protección de sistemas pesqueros
\end{multicols}
\end{enumerate}
\vspace{0.15cm}
\item ¿Qué es el ciberacoso?
\begin{enumerate}[label=\textbf{\Alph*)}]
\begin{multicols}{2}
  \item El acoso realizado a través de medios digitales
  \item El acoso realizado en persona
  \item El acoso realizado por teléfono
  \item El acoso realizado por carta
  \item El acoso realizado por radio
  \item El acoso realizado por televisión
\end{multicols}
\end{enumerate}
\vspace{0.15cm}
\item ¿Qué es la violencia digital de género?
\begin{enumerate}[label=\textbf{\Alph*)}]
\begin{multicols}{2}
  \item La violencia contra las mujeres en el entorno digital
  \item La violencia contra los hombres en el entorno digital
  \item La violencia contra los niños en el entorno digital
  \item La violencia contra los ancianos en el entorno digital
  \item La violencia contra los discapacitados en el entorno digital
  \item La violencia contra los indígenas en el entorno digital
\end{multicols}
\end{enumerate}
\vspace{0.15cm}
\item ¿Qué es la desinformación?
\begin{enumerate}[label=\textbf{\Alph*)}]
\begin{multicols}{2}
  \item La difusión de información falsa o engañosa
  \item La difusión de información verdadera
  \item La difusión de información verificada
  \item La difusión de información objetiva
  \item La difusión de información científica
  \item La difusión de información académica
\end{multicols}
\end{enumerate}
\vspace{0.15cm}
\item ¿Qué son las noticias falsas (fake news)?
\begin{enumerate}[label=\textbf{\Alph*)}]
\begin{multicols}{2}
  \item Informaciones falsas difundidas como si fueran verdaderas
  \item Informaciones verdaderas difundidas como falsas
  \item Informaciones verificadas difundidas como falsas
  \item Informaciones objetivas difundidas como falsas
  \item Informaciones científicas difundidas como falsas
  \item Informaciones académicas difundidas como falsas
\end{multicols}
\end{enumerate}
\vspace{0.15cm}
\item ¿Qué es la verificación de hechos (fact-checking)?
\begin{enumerate}[label=\textbf{\Alph*)}]
\begin{multicols}{2}
  \item El proceso de comprobar la veracidad de la información
  \item El proceso de crear información falsa
  \item El proceso de difundir información falsa
  \item El proceso de ocultar información verdadera
  \item El proceso de manipular información
  \item El proceso de distorsionar información
\end{multicols}
\end{enumerate}
\vspace{0.15cm}
\item ¿Qué es el periodismo digital?
\begin{enumerate}[label=\textbf{\Alph*)}]
\begin{multicols}{2}
  \item El periodismo realizado a través de plataformas digitales
  \item El periodismo realizado en papel
  \item El periodismo realizado en radio
  \item El periodismo realizado en televisión
  \item El periodismo realizado en cine
  \item El periodismo realizado en teatro
\end{multicols}
\end{enumerate}
\vspace{0.15cm}
\item ¿Qué es el gobierno electrónico?
\begin{enumerate}[label=\textbf{\Alph*)}]
\begin{multicols}{2}
  \item El uso de tecnologías digitales para prestar servicios públicos
  \item El uso de tecnologías digitales para el comercio
  \item El uso de tecnologías digitales para la educación
  \item El uso de tecnologías digitales para la salud
  \item El uso de tecnologías digitales para la banca
  \item El uso de tecnologías digitales para el transporte
\end{multicols}
\end{enumerate}
\vspace{0.15cm}
\item ¿Qué es la participación electrónica?
\begin{enumerate}[label=\textbf{\Alph*)}]
\begin{multicols}{2}
  \item La participación ciudadana a través de medios digitales
  \item La participación ciudadana en persona
  \item La participación ciudadana por teléfono
  \item La participación ciudadana por carta
  \item La participación ciudadana por radio
  \item La participación ciudadana por televisión
\end{multicols}
\end{enumerate}
\vspace{0.15cm}
\item ¿Qué es el voto electrónico?
\begin{enumerate}[label=\textbf{\Alph*)}]
\begin{multicols}{2}
  \item El voto realizado a través de medios digitales
  \item El voto realizado en papel
  \item El voto realizado en urna
  \item El voto realizado por correo
  \item El voto realizado por teléfono
  \item El voto realizado por fax
\end{multicols}
\end{enumerate}
\vspace{0.15cm}
\item ¿Qué es la inteligencia artificial?
\begin{enumerate}[label=\textbf{\Alph*)}]
\begin{multicols}{2}
  \item La capacidad de las máquinas para realizar tareas que requieren inteligencia humana
  \item La capacidad de las máquinas para realizar tareas mecánicas
  \item La capacidad de las máquinas para realizar tareas repetitivas
  \item La capacidad de las máquinas para realizar tareas simples
  \item La capacidad de las máquinas para realizar tareas básicas
  \item La capacidad de las máquinas para realizar tareas rutinarias
\end{multicols}
\end{enumerate}
\vspace{0.15cm}
\item ¿Qué es el big data?
\begin{enumerate}[label=\textbf{\Alph*)}]
\begin{multicols}{2}
  \item El conjunto de datos masivos que requieren tecnologías especiales para su análisis
  \item El conjunto de datos pequeños
  \item El conjunto de datos medianos
  \item El conjunto de datos simples
  \item El conjunto de datos básicos
  \item El conjunto de datos rutinarios
\end{multicols}
\end{enumerate}
\vspace{0.15cm}
\item ¿Qué es el internet de las cosas?
\begin{enumerate}[label=\textbf{\Alph*)}]
\begin{multicols}{2}
  \item La interconexión de dispositivos cotidianos a través de internet
  \item La interconexión de personas a través de internet
  \item La interconexión de empresas a través de internet
  \item La interconexión de gobiernos a través de internet
  \item La interconexión de escuelas a través de internet
  \item La interconexión de hospitales a través de internet
\end{multicols}
\end{enumerate}
\vspace{0.15cm}
\item ¿Qué es la computación en la nube?
\begin{enumerate}[label=\textbf{\Alph*)}]
\begin{multicols}{2}
  \item El uso de servicios informáticos a través de internet
  \item El uso de servicios informáticos locales
  \item El uso de servicios informáticos físicos
  \item El uso de servicios informáticos tradicionales
  \item El uso de servicios informáticos antiguos
  \item El uso de servicios informáticos obsoletos
\end{multicols}
\end{enumerate}
\vspace{0.15cm}
\item ¿Qué es el teletrabajo?
\begin{enumerate}[label=\textbf{\Alph*)}]
\begin{multicols}{2}
  \item El trabajo realizado a distancia mediante tecnologías de la información
  \item El trabajo realizado en una oficina
  \item El trabajo realizado en una fábrica
  \item El trabajo realizado en un campo
  \item El trabajo realizado en una tienda
  \item El trabajo realizado en un taller
\end{multicols}
\end{enumerate}
\vspace{0.15cm}
\item ¿Qué es la educación virtual?
\begin{enumerate}[label=\textbf{\Alph*)}]
\begin{multicols}{2}
  \item La educación impartida a través de plataformas digitales
  \item La educación impartida en aulas presenciales
  \item La educación impartida en bibliotecas
  \item La educación impartida en laboratorios
  \item La educación impartida en talleres
  \item La educación impartida en seminarios
\end{multicols}
\end{enumerate}
\vspace{0.15cm}
\item ¿Qué es la telemedicina?
\begin{enumerate}[label=\textbf{\Alph*)}]
\begin{multicols}{2}
  \item La prestación de servicios médicos a distancia mediante tecnologías
  \item La prestación de servicios médicos presenciales
  \item La prestación de servicios médicos domiciliarios
  \item La prestación de servicios médicos hospitalarios
  \item La prestación de servicios médicos ambulatorios
  \item La prestación de servicios médicos de emergencia
\end{multicols}
\end{enumerate}
\vspace{0.15cm}
\item ¿Qué es el comercio electrónico?
\begin{enumerate}[label=\textbf{\Alph*)}]
\begin{multicols}{2}
  \item La compra y venta de bienes y servicios a través de internet
  \item La compra y venta de bienes en tiendas físicas
  \item La compra y venta de bienes en mercados
  \item La compra y venta de bienes en ferias
  \item La compra y venta de bienes en subastas
  \item La compra y venta de bienes en remates
\end{multicols}
\end{enumerate}
\vspace{0.15cm}
\item ¿Qué es la banca electrónica?
\begin{enumerate}[label=\textbf{\Alph*)}]
\begin{multicols}{2}
  \item La prestación de servicios bancarios a través de internet
  \item La prestación de servicios bancarios en sucursales
  \item La prestación de servicios bancarios en cajeros
  \item La prestación de servicios bancarios por teléfono
  \item La prestación de servicios bancarios por correo
  \item La prestación de servicios bancarios en persona
\end{multicols}
\end{enumerate}
\vspace{0.15cm}
\item ¿Qué es la firma electrónica?
\begin{enumerate}[label=\textbf{\Alph*)}]
\begin{multicols}{2}
  \item El mecanismo digital para autenticar la identidad en documentos electrónicos
  \item El mecanismo manual para firmar documentos
  \item El mecanismo físico para firmar documentos
  \item El mecanismo tradicional para firmar documentos
  \item El mecanismo antiguo para firmar documentos
  \item El mecanismo obsoleto para firmar documentos
\end{multicols}
\end{enumerate}
\vspace{0.15cm}
\item ¿Qué es el documento electrónico?
\begin{enumerate}[label=\textbf{\Alph*)}]
\begin{multicols}{2}
  \item El documento almacenado y transmitido en formato digital
  \item El documento almacenado en papel
  \item El documento almacenado en microfilm
  \item El documento almacenado en cinta
  \item El documento almacenado en disco
  \item El documento almacenado en casete
\end{multicols}
\end{enumerate}
\vspace{0.15cm}
\item ¿Qué es la administración electrónica?
\begin{enumerate}[label=\textbf{\Alph*)}]
\begin{multicols}{2}
  \item El uso de tecnologías digitales en la administración pública
  \item El uso de tecnologías digitales en la empresa privada
  \item El uso de tecnologías digitales en la sociedad civil
  \item El uso de tecnologías digitales en las ONG
  \item El uso de tecnologías digitales en los sindicatos
  \item El uso de tecnologías digitales en los partidos
\end{multicols}
\end{enumerate}
\vspace{0.15cm}
\item ¿Qué es la inclusión digital?
\begin{enumerate}[label=\textbf{\Alph*)}]
\begin{multicols}{2}
  \item La incorporación de todos los grupos sociales a la sociedad de la información
  \item La exclusión de algunos grupos de la sociedad de la información
  \item La marginación de algunos grupos de la sociedad de la información
  \item La discriminación de algunos grupos de la sociedad de la información
  \item La segregación de algunos grupos de la sociedad de la información
  \item La separación de algunos grupos de la sociedad de la información
\end{multicols}
\end{enumerate}
\vspace{0.15cm}
\item ¿Qué es la accesibilidad web?
\begin{enumerate}[label=\textbf{\Alph*)}]
\begin{multicols}{2}
  \item La eliminación de barreras para que todos puedan usar internet
  \item La creación de barreras para el uso de internet
  \item La regulación de barreras para el uso de internet
  \item La limitación de barreras para el uso de internet
  \item La condición de barreras para el uso de internet
  \item La autorización de barreras para el uso de internet
\end{multicols}
\end{enumerate}
\vspace{0.15cm}
\item ¿Qué es el software libre?
\begin{enumerate}[label=\textbf{\Alph*)}]
\begin{multicols}{2}
  \item El software que puede ser usado, modificado y distribuido libremente
  \item El software que no puede ser modificado
  \item El software que no puede ser distribuido
  \item El software que no puede ser usado
  \item El software que no puede ser copiado
  \item El software que no puede ser compartido
\end{multicols}
\end{enumerate}
\vspace{0.15cm}
\item ¿Qué es el código abierto?
\begin{enumerate}[label=\textbf{\Alph*)}]
\begin{multicols}{2}
  \item El software cuyo código fuente es accesible para su modificación
  \item El software cuyo código fuente es secreto
  \item El software cuyo código fuente es privado
  \item El software cuyo código fuente es restringido
  \item El software cuyo código fuente es limitado
  \item El software cuyo código fuente es cerrado
\end{multicols}
\end{enumerate}
\vspace{0.15cm}
\item ¿Qué es la propiedad intelectual?
\begin{enumerate}[label=\textbf{\Alph*)}]
\begin{multicols}{2}
  \item El conjunto de derechos sobre las creaciones del intelecto
  \item El conjunto de derechos sobre los bienes materiales
  \item El conjunto de derechos sobre los bienes inmuebles
  \item El conjunto de derechos sobre los bienes muebles
  \item El conjunto de derechos sobre los bienes naturales
  \item El conjunto de derechos sobre los bienes económicos
\end{multicols}
\end{enumerate}
\vspace{0.15cm}
\item ¿Qué es el derecho de autor?
\begin{enumerate}[label=\textbf{\Alph*)}]
\begin{multicols}{2}
  \item El derecho que protege las obras literarias y artísticas
  \item El derecho que protege las invenciones
  \item El derecho que protege las marcas
  \item El derecho que protege los diseños industriales
  \item El derecho que protege los secretos comerciales
  \item El derecho que protege las variedades vegetales
\end{multicols}
\end{enumerate}
\vspace{0.15cm}
\item ¿Qué es la piratería?
\begin{enumerate}[label=\textbf{\Alph*)}]
\begin{multicols}{2}
  \item La reproducción y distribución no autorizada de obras protegidas
  \item La reproducción autorizada de obras
  \item La distribución autorizada de obras
  \item La reproducción legal de obras
  \item La distribución legal de obras
  \item La reproducción permitida de obras
\end{multicols}
\end{enumerate}
\vspace{0.15cm}
\item ¿Qué es el plagio?
\begin{enumerate}[label=\textbf{\Alph*)}]
\begin{multicols}{2}
  \item La presentación de ideas o palabras de otro como propias
  \item La presentación de ideas propias
  \item La presentación de ideas originales
  \item La presentación de ideas novedosas
  \item La presentación de ideas innovadoras
  \item La presentación de ideas creativas
\end{multicols}
\end{enumerate}
\vspace{0.15cm}
\item ¿Qué es la cita?
\begin{enumerate}[label=\textbf{\Alph*)}]
\begin{multicols}{2}
  \item La mención de la fuente de una idea o texto
  \item La omisión de la fuente
  \item La ocultación de la fuente
  \item La negación de la fuente
  \item La distorsión de la fuente
  \item La manipulación de la fuente
\end{multicols}
\end{enumerate}
\vspace{0.15cm}
\item ¿Qué es la referencia bibliográfica?
\begin{enumerate}[label=\textbf{\Alph*)}]
\begin{multicols}{2}
  \item La lista de fuentes citadas en un trabajo
  \item La lista de fuentes no citadas
  \item La lista de fuentes inventadas
  \item La lista de fuentes falsas
  \item La lista de fuentes ocultas
  \item La lista de fuentes negadas
\end{multicols}
\end{enumerate}
\vspace{0.15cm}
\item ¿Qué es el acceso abierto?
\begin{enumerate}[label=\textbf{\Alph*)}]
\begin{multicols}{2}
  \item El acceso libre y gratuito a la información científica
  \item El acceso restringido a la información científica
  \item El acceso privado a la información científica
  \item El acceso limitado a la información científica
  \item El acceso condicionado a la información científica
  \item El acceso autorizado a la información científica
\end{multicols}
\end{enumerate}
\vspace{0.15cm}
\item ¿Qué es la neutralidad de la red?
\begin{enumerate}[label=\textbf{\Alph*)}]
\begin{multicols}{2}
  \item El principio de que todo el tráfico de internet debe ser tratado por igual
  \item El principio de que el tráfico debe ser discriminado
  \item El principio de que el tráfico debe ser priorizado
  \item El principio de que el tráfico debe ser filtrado
  \item El principio de que el tráfico debe ser bloqueado
  \item El principio de que el tráfico debe ser restringido
\end{multicols}
\end{enumerate}
\vspace{0.15cm}
\item ¿Qué es la gobernanza de internet?
\begin{enumerate}[label=\textbf{\Alph*)}]
\begin{multicols}{2}
  \item El conjunto de normas y procesos que rigen el funcionamiento de internet
  \item El conjunto de normas que rigen el comercio
  \item El conjunto de normas que rigen la economía
  \item El conjunto de normas que rigen la política
  \item El conjunto de normas que rigen la religión
  \item El conjunto de normas que rigen la cultura
\end{multicols}
\end{enumerate}
\vspace{0.15cm}
\item ¿Qué es la sociedad del conocimiento?
\begin{enumerate}[label=\textbf{\Alph*)}]
\begin{multicols}{2}
  \item Una sociedad en la que el conocimiento es el principal recurso
  \item Una sociedad en la que la agricultura es el principal recurso
  \item Una sociedad en la que la industria es el principal recurso
  \item Una sociedad en la que el comercio es el principal recurso
  \item Una sociedad en la que la minería es el principal recurso
  \item Una sociedad en la que la pesca es el principal recurso
\end{multicols}
\end{enumerate}
\vspace{0.15cm}
\item ¿Qué es la economía del conocimiento?
\begin{enumerate}[label=\textbf{\Alph*)}]
\begin{multicols}{2}
  \item La economía basada en la producción y uso del conocimiento
  \item La economía basada en la agricultura
  \item La economía basada en la industria
  \item La economía basada en el comercio
  \item La economía basada en la minería
  \item La economía basada en la pesca
\end{multicols}
\end{enumerate}
\vspace{0.15cm}
\item ¿Qué es la innovación?
\begin{enumerate}[label=\textbf{\Alph*)}]
\begin{multicols}{2}
  \item La introducción de nuevas ideas, productos o procesos
  \item La repetición de ideas existentes
  \item La copia de productos existentes
  \item La imitación de procesos existentes
  \item La reproducción de ideas existentes
  \item La duplicación de productos existentes
\end{multicols}
\end{enumerate}
\vspace{0.15cm}
\item ¿Qué es la investigación y desarrollo (I+D)?
\begin{enumerate}[label=\textbf{\Alph*)}]
\begin{multicols}{2}
  \item El conjunto de actividades para generar nuevo conocimiento
  \item El conjunto de actividades para repetir conocimiento
  \item El conjunto de actividades para copiar conocimiento
  \item El conjunto de actividades para imitar conocimiento
  \item El conjunto de actividades para duplicar conocimiento
  \item El conjunto de actividades para reproducir conocimiento
\end{multicols}
\end{enumerate}
\vspace{0.15cm}
\item ¿Qué es la transferencia de tecnología?
\begin{enumerate}[label=\textbf{\Alph*)}]
\begin{multicols}{2}
  \item El proceso de compartir tecnología entre países o empresas
  \item El proceso de ocultar tecnología
  \item El proceso de negar tecnología
  \item El proceso de restringir tecnología
  \item El proceso de limitar tecnología
  \item El proceso de condicionar tecnología
\end{multicols}
\end{enumerate}
\vspace{0.15cm}
\item ¿Qué es la cooperación tecnológica?
\begin{enumerate}[label=\textbf{\Alph*)}]
\begin{multicols}{2}
  \item La colaboración entre países o empresas para desarrollar tecnología
  \item La competencia entre países o empresas
  \item El conflicto entre países o empresas
  \item La separación entre países o empresas
  \item El aislamiento entre países o empresas
  \item La fragmentación entre países o empresas
\end{multicols}
\end{enumerate}
\vspace{0.15cm}
\item ¿Qué es la soberanía tecnológica?
\begin{enumerate}[label=\textbf{\Alph*)}]
\begin{multicols}{2}
  \item La capacidad de un país para desarrollar y controlar su propia tecnología
  \item La dependencia de un país de la tecnología extranjera
  \item La subordinación de un país a la tecnología extranjera
  \item La sumisión de un país a la tecnología extranjera
  \item La dependencia de un país de la tecnología nacional
  \item La subordinación de un país a la tecnología nacional
\end{multicols}
\end{enumerate}
\vspace{0.15cm}
\item ¿Qué es la autonomía tecnológica?
\begin{enumerate}[label=\textbf{\Alph*)}]
\begin{multicols}{2}
  \item La capacidad de un país para tomar decisiones sobre su tecnología
  \item La dependencia de un país de decisiones externas
  \item La subordinación de un país a decisiones externas
  \item La sumisión de un país a decisiones externas
  \item La dependencia de un país de decisiones internas
  \item La subordinación de un país a decisiones internas
\end{multicols}
\end{enumerate}
\vspace{0.15cm}
\item ¿Qué es la ética digital?
\begin{enumerate}[label=\textbf{\Alph*)}]
\begin{multicols}{2}
  \item El conjunto de principios que guían el comportamiento en el entorno digital
  \item El conjunto de principios que guían el comportamiento en el entorno físico
  \item El conjunto de principios que guían el comportamiento en el entorno rural
  \item El conjunto de principios que guían el comportamiento en el entorno urbano
  \item El conjunto de principios que guían el comportamiento en el entorno laboral
  \item El conjunto de principios que guían el comportamiento en el entorno escolar
\end{multicols}
\end{enumerate}
\vspace{0.15cm}
\item ¿Qué es la responsabilidad digital?
\begin{enumerate}[label=\textbf{\Alph*)}]
\begin{multicols}{2}
  \item La obligación de actuar éticamente en el entorno digital
  \item La obligación de actuar éticamente en el entorno físico
  \item La obligación de actuar éticamente en el entorno rural
  \item La obligación de actuar éticamente en el entorno urbano
  \item La obligación de actuar éticamente en el entorno laboral
  \item La obligación de actuar éticamente en el entorno escolar
\end{multicols}
\end{enumerate}
\vspace{0.15cm}
\item ¿Qué es la solidaridad digital?
\begin{enumerate}[label=\textbf{\Alph*)}]
\begin{multicols}{2}
  \item La ayuda mutua en el entorno digital
  \item La indiferencia en el entorno digital
  \item La competencia en el entorno digital
  \item La exclusión en el entorno digital
  \item La marginación en el entorno digital
  \item La discriminación en el entorno digital
\end{multicols}
\end{enumerate}
\vspace{0.15cm}
\item ¿Qué es la justicia digital?
\begin{enumerate}[label=\textbf{\Alph*)}]
\begin{multicols}{2}
  \item La distribución equitativa de los beneficios y cargas de la tecnología
  \item La distribución inequitativa de los beneficios de la tecnología
  \item La concentración de los beneficios de la tecnología
  \item La concentración de las cargas de la tecnología
  \item La exclusión de los beneficios de la tecnología
  \item La exclusión de las cargas de la tecnología
\end{multicols}
\end{enumerate}
\vspace{0.15cm}
\end{enumerate}\n
\subsubsection*{Sección V: Migración Contemporánea (211–260)}

\begin{enumerate}[label=\textbf{\arabic*.}]
\setcounter{enumi}{210}
\item ¿Qué es la migración?
\begin{enumerate}[label=\textbf{\Alph*)}]
\begin{multicols}{2}
  \item El desplazamiento de personas de un lugar a otro
  \item La permanencia de personas en un lugar
  \item El crecimiento de la población
  \item La disminución de la población
  \item La distribución de la población
  \item La composición de la población
\end{multicols}
\end{enumerate}
\vspace{0.15cm}
\item ¿Qué es la migración internacional?
\begin{enumerate}[label=\textbf{\Alph*)}]
\begin{multicols}{2}
  \item El desplazamiento de personas de un país a otro
  \item El desplazamiento de personas dentro de un país
  \item El desplazamiento de personas dentro de una ciudad
  \item El desplazamiento de personas dentro de una región
  \item El desplazamiento de personas dentro de una comunidad
  \item El desplazamiento de personas dentro de una familia
\end{multicols}
\end{enumerate}
\vspace{0.15cm}
\item ¿Qué es la migración interna?
\begin{enumerate}[label=\textbf{\Alph*)}]
\begin{multicols}{2}
  \item El desplazamiento de personas dentro de un mismo país
  \item El desplazamiento de personas entre países
  \item El desplazamiento de personas dentro de una ciudad
  \item El desplazamiento de personas dentro de una región
  \item El desplazamiento de personas dentro de una comunidad
  \item El desplazamiento de personas dentro de una familia
\end{multicols}
\end{enumerate}
\vspace{0.15cm}
\item ¿Qué es la emigración?
\begin{enumerate}[label=\textbf{\Alph*)}]
\begin{multicols}{2}
  \item La salida de personas de un lugar para establecerse en otro
  \item La llegada de personas a un lugar
  \item La permanencia de personas en un lugar
  \item El crecimiento de la población
  \item La disminución de la población
  \item La distribución de la población
\end{multicols}
\end{enumerate}
\vspace{0.15cm}
\item ¿Qué es la inmigración?
\begin{enumerate}[label=\textbf{\Alph*)}]
\begin{multicols}{2}
  \item La llegada de personas a un lugar para establecerse
  \item La salida de personas de un lugar
  \item La permanencia de personas en un lugar
  \item El crecimiento de la población
  \item La disminución de la población
  \item La distribución de la población
\end{multicols}
\end{enumerate}
\vspace{0.15cm}
\item ¿Qué es el saldo migratorio?
\begin{enumerate}[label=\textbf{\Alph*)}]
\begin{multicols}{2}
  \item La diferencia entre inmigrantes y emigrantes
  \item La suma de inmigrantes y emigrantes
  \item El producto de inmigrantes y emigrantes
  \item El cociente de inmigrantes y emigrantes
  \item El promedio de inmigrantes y emigrantes
  \item La tasa de inmigrantes y emigrantes
\end{multicols}
\end{enumerate}
\vspace{0.15cm}
\item ¿Qué es la migración forzada?
\begin{enumerate}[label=\textbf{\Alph*)}]
\begin{multicols}{2}
  \item El desplazamiento de personas obligadas a abandonar su lugar de origen
  \item El desplazamiento de personas por voluntad propia
  \item El desplazamiento de personas por turismo
  \item El desplazamiento de personas por negocios
  \item El desplazamiento de personas por estudio
  \item El desplazamiento de personas por vacaciones
\end{multicols}
\end{enumerate}
\vspace{0.15cm}
\item ¿Qué es la migración voluntaria?
\begin{enumerate}[label=\textbf{\Alph*)}]
\begin{multicols}{2}
  \item El desplazamiento de personas por voluntad propia
  \item El desplazamiento de personas obligadas
  \item El desplazamiento de personas por conflicto
  \item El desplazamiento de personas por desastre
  \item El desplazamiento de personas por persecución
  \item El desplazamiento de personas por violencia
\end{multicols}
\end{enumerate}
\vspace{0.15cm}
\item ¿Qué es la migración irregular?
\begin{enumerate}[label=\textbf{\Alph*)}]
\begin{multicols}{2}
  \item El desplazamiento de personas sin documentación legal
  \item El desplazamiento de personas con documentación legal
  \item El desplazamiento de personas por turismo
  \item El desplazamiento de personas por negocios
  \item El desplazamiento de personas por estudio
  \item El desplazamiento de personas por vacaciones
\end{multicols}
\end{enumerate}
\vspace{0.15cm}
\item ¿Qué es la migración regular?
\begin{enumerate}[label=\textbf{\Alph*)}]
\begin{multicols}{2}
  \item El desplazamiento de personas con documentación legal
  \item El desplazamiento de personas sin documentación legal
  \item El desplazamiento de personas por conflicto
  \item El desplazamiento de personas por desastre
  \item El desplazamiento de personas por persecución
  \item El desplazamiento de personas por violencia
\end{multicols}
\end{enumerate}
\vspace{0.15cm}
\item ¿Qué es un refugiado?
\begin{enumerate}[label=\textbf{\Alph*)}]
\begin{multicols}{2}
  \item Una persona que huye de su país por persecución o conflicto
  \item Una persona que viaja por turismo
  \item Una persona que viaja por negocios
  \item Una persona que viaja por estudio
  \item Una persona que viaja por vacaciones
  \item Una persona que viaja por placer
\end{multicols}
\end{enumerate}
\vspace{0.15cm}
\item ¿Qué es un desplazado interno?
\begin{enumerate}[label=\textbf{\Alph*)}]
\begin{multicols}{2}
  \item Una persona que ha huido de su hogar pero permanece en su país
  \item Una persona que ha huido de su país
  \item Una persona que viaja por turismo
  \item Una persona que viaja por negocios
  \item Una persona que viaja por estudio
  \item Una persona que viaja por vacaciones
\end{multicols}
\end{enumerate}
\vspace{0.15cm}
\item ¿Qué es un solicitante de asilo?
\begin{enumerate}[label=\textbf{\Alph*)}]
\begin{multicols}{2}
  \item Una persona que pide protección internacional
  \item Una persona que rechaza la protección internacional
  \item Una persona que ignora la protección internacional
  \item Una persona que olvida la protección internacional
  \item Una persona que niega la protección internacional
  \item Una persona que oculta la protección internacional
\end{multicols}
\end{enumerate}
\vspace{0.15cm}
\item ¿Qué es el asilo político?
\begin{enumerate}[label=\textbf{\Alph*)}]
\begin{multicols}{2}
  \item La protección que un país brinda a una persona perseguida por razones políticas
  \item La protección que un país brinda a una persona por razones económicas
  \item La protección que un país brinda a una persona por razones sociales
  \item La protección que un país brinda a una persona por razones culturales
  \item La protección que un país brinda a una persona por razones religiosas
  \item La protección que un país brinda a una persona por razones ambientales
\end{multicols}
\end{enumerate}
\vspace{0.15cm}
\item ¿Qué es el refugio?
\begin{enumerate}[label=\textbf{\Alph*)}]
\begin{multicols}{2}
  \item La protección que un país brinda a una persona que huye de conflictos
  \item La protección que un país brinda a una persona por razones económicas
  \item La protección que un país brinda a una persona por razones sociales
  \item La protección que un país brinda a una persona por razones culturales
  \item La protección que un país brinda a una persona por razones religiosas
  \item La protección que un país brinda a una persona por razones ambientales
\end{multicols}
\end{enumerate}
\vspace{0.15cm}
\item ¿Qué es el principio de no devolución?
\begin{enumerate}[label=\textbf{\Alph*)}]
\begin{multicols}{2}
  \item La prohibición de devolver a una persona a un país donde corre peligro
  \item La obligación de devolver a una persona a su país
  \item La posibilidad de devolver a una persona a su país
  \item La opción de devolver a una persona a su país
  \item La alternativa de devolver a una persona a su país
  \item La decisión de devolver a una persona a su país
\end{multicols}
\end{enumerate}
\vspace{0.15cm}
\item ¿Qué es la trata de personas?
\begin{enumerate}[label=\textbf{\Alph*)}]
\begin{multicols}{2}
  \item El comercio ilegal de seres humanos con fines de explotación
  \item El comercio legal de seres humanos
  \item El comercio de bienes
  \item El comercio de servicios
  \item El comercio de tecnología
  \item El comercio de alimentos
\end{multicols}
\end{enumerate}
\vspace{0.15cm}
\item ¿Qué es el tráfico de personas?
\begin{enumerate}[label=\textbf{\Alph*)}]
\begin{multicols}{2}
  \item El traslado ilegal de personas a través de fronteras
  \item El traslado legal de personas
  \item El traslado de bienes
  \item El traslado de servicios
  \item El traslado de tecnología
  \item El traslado de alimentos
\end{multicols}
\end{enumerate}
\vspace{0.15cm}
\item ¿Qué es el contrabando de personas?
\begin{enumerate}[label=\textbf{\Alph*)}]
\begin{multicols}{2}
  \item El traslado ilegal de personas con su consentimiento
  \item El traslado legal de personas
  \item El traslado de bienes
  \item El traslado de servicios
  \item El traslado de tecnología
  \item El traslado de alimentos
\end{multicols}
\end{enumerate}
\vspace{0.15cm}
\item ¿Qué es la explotación laboral de migrantes?
\begin{enumerate}[label=\textbf{\Alph*)}]
\begin{multicols}{2}
  \item El abuso en las condiciones de trabajo de los migrantes
  \item El respeto a las condiciones de trabajo de los migrantes
  \item La mejora de las condiciones de trabajo de los migrantes
  \item La protección de las condiciones de trabajo de los migrantes
  \item La garantía de las condiciones de trabajo de los migrantes
  \item La promoción de las condiciones de trabajo de los migrantes
\end{multicols}
\end{enumerate}
\vspace{0.15cm}
\item ¿Qué es la xenofobia?
\begin{enumerate}[label=\textbf{\Alph*)}]
\begin{multicols}{2}
  \item El rechazo hacia los extranjeros
  \item El rechazo hacia los nacionales
  \item El rechazo hacia los ancianos
  \item El rechazo hacia los niños
  \item El rechazo hacia los discapacitados
  \item El rechazo hacia los pobres
\end{multicols}
\end{enumerate}
\vspace{0.15cm}
\item ¿Qué es la discriminación hacia los migrantes?
\begin{enumerate}[label=\textbf{\Alph*)}]
\begin{multicols}{2}
  \item El trato desigual e injusto hacia los migrantes
  \item El trato igualitario hacia los migrantes
  \item El trato justo hacia los migrantes
  \item El trato objetivo hacia los migrantes
  \item El trato científico hacia los migrantes
  \item El trato artístico hacia los migrantes
\end{multicols}
\end{enumerate}
\vspace{0.15cm}
\item ¿Qué es la integración de los migrantes?
\begin{enumerate}[label=\textbf{\Alph*)}]
\begin{multicols}{2}
  \item El proceso de incorporación de los migrantes a la sociedad de acogida
  \item El proceso de exclusión de los migrantes
  \item El proceso de marginación de los migrantes
  \item El proceso de discriminación de los migrantes
  \item El proceso de segregación de los migrantes
  \item El proceso de aislamiento de los migrantes
\end{multicols}
\end{enumerate}
\vspace{0.15cm}
\item ¿Qué es la inclusión de los migrantes?
\begin{enumerate}[label=\textbf{\Alph*)}]
\begin{multicols}{2}
  \item La incorporación de los migrantes a la vida social, económica y cultural
  \item La exclusión de los migrantes
  \item La marginación de los migrantes
  \item La discriminación de los migrantes
  \item La segregación de los migrantes
  \item El aislamiento de los migrantes
\end{multicols}
\end{enumerate}
\vspace{0.15cm}
\item ¿Qué es la multiculturalidad?
\begin{enumerate}[label=\textbf{\Alph*)}]
\begin{multicols}{2}
  \item La coexistencia de múltiples culturas en una sociedad
  \item La existencia de una sola cultura
  \item La existencia de una cultura dominante
  \item La existencia de una cultura homogénea
  \item La existencia de una cultura global
  \item La existencia de una cultura nacional
\end{multicols}
\end{enumerate}
\vspace{0.15cm}
\item ¿Qué es la interculturalidad?
\begin{enumerate}[label=\textbf{\Alph*)}]
\begin{multicols}{2}
  \item La interacción y diálogo entre culturas diferentes
  \item La existencia de una sola cultura
  \item La existencia de culturas separadas
  \item La existencia de una cultura dominante
  \item La existencia de una cultura homogénea
  \item La existencia de una cultura global
\end{multicols}
\end{enumerate}
\vspace{0.15cm}
\item ¿Qué es la ciudadanía intercultural?
\begin{enumerate}[label=\textbf{\Alph*)}]
\begin{multicols}{2}
  \item El ejercicio de la ciudadanía en contextos de diversidad cultural
  \item El ejercicio de la ciudadanía en contextos homogéneos
  \item El ejercicio de la ciudadanía en contextos monoculturales
  \item El ejercicio de la ciudadanía en contextos cerrados
  \item El ejercicio de la ciudadanía en contextos aislados
  \item El ejercicio de la ciudadanía en contextos separados
\end{multicols}
\end{enumerate}
\vspace{0.15cm}
\item ¿Qué es la fuga de cerebros?
\begin{enumerate}[label=\textbf{\Alph*)}]
\begin{multicols}{2}
  \item La emigración de profesionales altamente calificados
  \item La inmigración de profesionales altamente calificados
  \item La emigración de trabajadores no calificados
  \item La inmigración de trabajadores no calificados
  \item La emigración de estudiantes
  \item La inmigración de estudiantes
\end{multicols}
\end{enumerate}
\vspace{0.15cm}
\item ¿Qué son las remesas?
\begin{enumerate}[label=\textbf{\Alph*)}]
\begin{multicols}{2}
  \item El dinero que los migrantes envían a sus familias en el país de origen
  \item El dinero que los gobiernos envían a otros países
  \item El dinero que las empresas envían a otros países
  \item El dinero que las ONG envían a otros países
  \item El dinero que los bancos envían a otros países
  \item El dinero que los individuos envían a otros individuos
\end{multicols}
\end{enumerate}
\vspace{0.15cm}
\item ¿Qué es la diáspora?
\begin{enumerate}[label=\textbf{\Alph*)}]
\begin{multicols}{2}
  \item La dispersión de un pueblo o comunidad por el mundo
  \item La concentración de un pueblo en un lugar
  \item El aislamiento de un pueblo
  \item La separación de un pueblo
  \item La fragmentación de un pueblo
  \item La localización de un pueblo
\end{multicols}
\end{enumerate}
\vspace{0.15cm}
\item ¿Qué es el transnacionalismo?
\begin{enumerate}[label=\textbf{\Alph*)}]
\begin{multicols}{2}
  \item La creación de vínculos sociales, culturales y económicos entre países
  \item La separación de los países
  \item El aislamiento de los países
  \item La fragmentación de los países
  \item La localización de los países
  \item La regionalización de los países
\end{multicols}
\end{enumerate}
\vspace{0.15cm}
\item ¿Qué es la ciudadanía global?
\begin{enumerate}[label=\textbf{\Alph*)}]
\begin{multicols}{2}
  \item La conciencia de pertenencia a una comunidad mundial
  \item La conciencia de pertenencia a una nación
  \item La conciencia de pertenencia a una región
  \item La conciencia de pertenencia a una ciudad
  \item La conciencia de pertenencia a una comunidad
  \item La conciencia de pertenencia a una familia
\end{multicols}
\end{enumerate}
\vspace{0.15cm}
\item ¿Qué es la reunificación familiar?
\begin{enumerate}[label=\textbf{\Alph*)}]
\begin{multicols}{2}
  \item El proceso de reunir a los miembros de una familia separada por la migración
  \item El proceso de separar a los miembros de una familia
  \item El proceso de dispersar a los miembros de una familia
  \item El proceso de dividir a los miembros de una familia
  \item El proceso de aislar a los miembros de una familia
  \item El proceso de excluir a los miembros de una familia
\end{multicols}
\end{enumerate}
\vspace{0.15cm}
\item ¿Qué es la deportación?
\begin{enumerate}[label=\textbf{\Alph*)}]
\begin{multicols}{2}
  \item La expulsión de un extranjero de un país
  \item La admisión de un extranjero en un país
  \item La naturalización de un extranjero
  \item La residencia de un extranjero
  \item La nacionalización de un extranjero
  \item La integración de un extranjero
\end{multicols}
\end{enumerate}
\vspace{0.15cm}
\item ¿Qué es la expulsión?
\begin{enumerate}[label=\textbf{\Alph*)}]
\begin{multicols}{2}
  \item La obligación de abandonar un país por decisión administrativa
  \item La autorización para permanecer en un país
  \item La admisión en un país
  \item La naturalización en un país
  \item La residencia en un país
  \item La integración en un país
\end{multicols}
\end{enumerate}
\vspace{0.15cm}
\item ¿Qué es la naturalización?
\begin{enumerate}[label=\textbf{\Alph*)}]
\begin{multicols}{2}
  \item El proceso por el cual un extranjero adquiere la nacionalidad
  \item El proceso por el cual un nacional pierde la nacionalidad
  \item El proceso por el cual un extranjero es expulsado
  \item El proceso por el cual un extranjero es deportado
  \item El proceso por el cual un extranjero es rechazado
  \item El proceso por el cual un extranjero es excluido
\end{multicols}
\end{enumerate}
\vspace{0.15cm}
\item ¿Qué es la doble nacionalidad?
\begin{enumerate}[label=\textbf{\Alph*)}]
\begin{multicols}{2}
  \item La posesión de la nacionalidad de dos países
  \item La posesión de la nacionalidad de un solo país
  \item La pérdida de la nacionalidad
  \item La renuncia a la nacionalidad
  \item La suspensión de la nacionalidad
  \item La restricción de la nacionalidad
\end{multicols}
\end{enumerate}
\vspace{0.15cm}
\item ¿Qué es el derecho internacional de los refugiados?
\begin{enumerate}[label=\textbf{\Alph*)}]
\begin{multicols}{2}
  \item El conjunto de normas que protegen a los refugiados
  \item El conjunto de normas que protegen a los turistas
  \item El conjunto de normas que protegen a los comerciantes
  \item El conjunto de normas que protegen a los estudiantes
  \item El conjunto de normas que protegen a los trabajadores
  \item El conjunto de normas que protegen a los inversionistas
\end{multicols}
\end{enumerate}
\vspace{0.15cm}
\item ¿Qué es el ACNUR?
\begin{enumerate}[label=\textbf{\Alph*)}]
\begin{multicols}{2}
  \item El Alto Comisionado de las Naciones Unidas para los Refugiados
  \item El Alto Comisionado de las Naciones Unidas para los Derechos Humanos
  \item El Alto Comisionado de las Naciones Unidas para el Desarrollo
  \item El Alto Comisionado de las Naciones Unidas para la Educación
  \item El Alto Comisionado de las Naciones Unidas para la Salud
  \item El Alto Comisionado de las Naciones Unidas para el Medio Ambiente
\end{multicols}
\end{enumerate}
\vspace{0.15cm}
\item ¿Qué es la OIM?
\begin{enumerate}[label=\textbf{\Alph*)}]
\begin{multicols}{2}
  \item La Organización Internacional para las Migraciones
  \item La Organización Internacional para los Derechos Humanos
  \item La Organización Internacional para el Desarrollo
  \item La Organización Internacional para la Educación
  \item La Organización Internacional para la Salud
  \item La Organización Internacional para el Medio Ambiente
\end{multicols}
\end{enumerate}
\vspace{0.15cm}
\item ¿Qué es la migración climática?
\begin{enumerate}[label=\textbf{\Alph*)}]
\begin{multicols}{2}
  \item El desplazamiento de personas por causas ambientales
  \item El desplazamiento de personas por causas económicas
  \item El desplazamiento de personas por causas políticas
  \item El desplazamiento de personas por causas sociales
  \item El desplazamiento de personas por causas culturales
  \item El desplazamiento de personas por causas religiosas
\end{multicols}
\end{enumerate}
\vspace{0.15cm}
\item ¿Qué es la migración económica?
\begin{enumerate}[label=\textbf{\Alph*)}]
\begin{multicols}{2}
  \item El desplazamiento de personas por causas económicas
  \item El desplazamiento de personas por causas ambientales
  \item El desplazamiento de personas por causas políticas
  \item El desplazamiento de personas por causas sociales
  \item El desplazamiento de personas por causas culturales
  \item El desplazamiento de personas por causas religiosas
\end{multicols}
\end{enumerate}
\vspace{0.15cm}
\item ¿Qué es la migración política?
\begin{enumerate}[label=\textbf{\Alph*)}]
\begin{multicols}{2}
  \item El desplazamiento de personas por causas políticas
  \item El desplazamiento de personas por causas económicas
  \item El desplazamiento de personas por causas ambientales
  \item El desplazamiento de personas por causas sociales
  \item El desplazamiento de personas por causas culturales
  \item El desplazamiento de personas por causas religiosas
\end{multicols}
\end{enumerate}
\vspace{0.15cm}
\item ¿Qué es la migración social?
\begin{enumerate}[label=\textbf{\Alph*)}]
\begin{multicols}{2}
  \item El desplazamiento de personas por causas sociales
  \item El desplazamiento de personas por causas económicas
  \item El desplazamiento de personas por causas ambientales
  \item El desplazamiento de personas por causas políticas
  \item El desplazamiento de personas por causas culturales
  \item El desplazamiento de personas por causas religiosas
\end{multicols}
\end{enumerate}
\vspace{0.15cm}
\item ¿Qué es la migración cultural?
\begin{enumerate}[label=\textbf{\Alph*)}]
\begin{multicols}{2}
  \item El desplazamiento de personas por causas culturales
  \item El desplazamiento de personas por causas económicas
  \item El desplazamiento de personas por causas ambientales
  \item El desplazamiento de personas por causas políticas
  \item El desplazamiento de personas por causas sociales
  \item El desplazamiento de personas por causas religiosas
\end{multicols}
\end{enumerate}
\vspace{0.15cm}
\item ¿Qué es la migración religiosa?
\begin{enumerate}[label=\textbf{\Alph*)}]
\begin{multicols}{2}
  \item El desplazamiento de personas por causas religiosas
  \item El desplazamiento de personas por causas económicas
  \item El desplazamiento de personas por causas ambientales
  \item El desplazamiento de personas por causas políticas
  \item El desplazamiento de personas por causas sociales
  \item El desplazamiento de personas por causas culturales
\end{multicols}
\end{enumerate}
\vspace{0.15cm}
\item ¿Qué es la migración de tránsito?
\begin{enumerate}[label=\textbf{\Alph*)}]
\begin{multicols}{2}
  \item El desplazamiento de personas que cruzan un país para llegar a otro
  \item El desplazamiento de personas que se quedan en un país
  \item El desplazamiento de personas que regresan a su país
  \item El desplazamiento de personas que salen de su país
  \item El desplazamiento de personas que llegan a su país
  \item El desplazamiento de personas que viven en su país
\end{multicols}
\end{enumerate}
\vspace{0.15cm}
\item ¿Qué es la migración de retorno?
\begin{enumerate}[label=\textbf{\Alph*)}]
\begin{multicols}{2}
  \item El regreso de los migrantes a su país de origen
  \item La salida de los migrantes de su país de origen
  \item La llegada de los migrantes a un nuevo país
  \item La permanencia de los migrantes en un país
  \item La circulación de los migrantes entre países
  \item La dispersión de los migrantes por el mundo
\end{multicols}
\end{enumerate}
\vspace{0.15cm}
\item ¿Qué es la migración circular?
\begin{enumerate}[label=\textbf{\Alph*)}]
\begin{multicols}{2}
  \item El desplazamiento repetido de personas entre dos o más países
  \item El desplazamiento único de personas
  \item El desplazamiento permanente de personas
  \item El desplazamiento definitivo de personas
  \item El desplazamiento temporal de personas
  \item El desplazamiento estacional de personas
\end{multicols}
\end{enumerate}
\vspace{0.15cm}
\item ¿Qué es la migración pendular?
\begin{enumerate}[label=\textbf{\Alph*)}]
\begin{multicols}{2}
  \item El desplazamiento diario o semanal de personas entre su lugar de residencia y trabajo
  \item El desplazamiento anual de personas
  \item El desplazamiento mensual de personas
  \item El desplazamiento trimestral de personas
  \item El desplazamiento semestral de personas
  \item El desplazamiento decenal de personas
\end{multicols}
\end{enumerate}
\vspace{0.15cm}
\end{enumerate}\n
\subsubsection*{Sección VI: Derechos y Sociedad Actual – Temas Contemporáneos (261–300)}

\begin{enumerate}[label=\textbf{\arabic*.}]
\setcounter{enumi}{260}
\item ¿Qué es la globalización?
\begin{enumerate}[label=\textbf{\Alph*)}]
\begin{multicols}{2}
  \item El proceso de integración económica, política, social y cultural a nivel mundial
  \item El proceso de aislamiento de los países
  \item El proceso de separación de las culturas
  \item El proceso de fragmentación política
  \item El proceso de localización económica
  \item El proceso de regionalización
\end{multicols}
\end{enumerate}
\vspace{0.15cm}
\item ¿Qué es la sociedad civil?
\begin{enumerate}[label=\textbf{\Alph*)}]
\begin{multicols}{2}
  \item El conjunto de organizaciones sociales independientes del Estado
  \item El conjunto de instituciones estatales
  \item El conjunto de empresas privadas
  \item El conjunto de partidos políticos
  \item El conjunto de sindicatos
  \item El conjunto de organizaciones religiosas
\end{multicols}
\end{enumerate}
\vspace{0.15cm}
\item ¿Qué es una ONG?
\begin{enumerate}[label=\textbf{\Alph*)}]
\begin{multicols}{2}
  \item Una organización no gubernamental
  \item Una organización gubernamental
  \item Una organización empresarial
  \item Una organización política
  \item Una organización religiosa
  \item Una organización deportiva
\end{multicols}
\end{enumerate}
\vspace{0.15cm}
\item ¿Qué es el voluntariado?
\begin{enumerate}[label=\textbf{\Alph*)}]
\begin{multicols}{2}
  \item El trabajo no remunerado realizado libremente para ayudar a otros
  \item El trabajo remunerado
  \item El trabajo forzado
  \item El trabajo informal
  \item El trabajo formal
  \item El trabajo doméstico
\end{multicols}
\end{enumerate}
\vspace{0.15cm}
\item ¿Qué es la solidaridad?
\begin{enumerate}[label=\textbf{\Alph*)}]
\begin{multicols}{2}
  \item La ayuda mutua entre los miembros de una sociedad
  \item La indiferencia entre los miembros de una sociedad
  \item La competencia entre los miembros de una sociedad
  \item La exclusión entre los miembros de una sociedad
  \item La marginación entre los miembros de una sociedad
  \item La discriminación entre los miembros de una sociedad
\end{multicols}
\end{enumerate}
\vspace{0.15cm}
\item ¿Qué es la tolerancia?
\begin{enumerate}[label=\textbf{\Alph*)}]
\begin{multicols}{2}
  \item El respeto hacia las ideas y prácticas diferentes
  \item El rechazo hacia las ideas diferentes
  \item La indiferencia hacia las ideas diferentes
  \item La imposición de las ideas propias
  \item La discriminación de las ideas diferentes
  \item La exclusión de las ideas diferentes
\end{multicols}
\end{enumerate}
\vspace{0.15cm}
\item ¿Qué es el diálogo intercultural?
\begin{enumerate}[label=\textbf{\Alph*)}]
\begin{multicols}{2}
  \item La comunicación y entendimiento entre culturas diferentes
  \item La imposición de una cultura sobre otra
  \item La separación de culturas
  \item La dominación de una cultura
  \item La desaparición de una cultura
  \item La homogeneización cultural
\end{multicols}
\end{enumerate}
\vspace{0.15cm}
\item ¿Qué es la paz social?
\begin{enumerate}[label=\textbf{\Alph*)}]
\begin{multicols}{2}
  \item La convivencia armónica y sin conflictos violentos en la sociedad
  \item La ausencia de conflictos
  \item La presencia de conflictos
  \item La violencia social
  \item La guerra civil
  \item La revolución social
\end{multicols}
\end{enumerate}
\vspace{0.15cm}
\item ¿Qué es la justicia social?
\begin{enumerate}[label=\textbf{\Alph*)}]
\begin{multicols}{2}
  \item La distribución equitativa de los bienes y oportunidades
  \item La distribución inequitativa de los bienes
  \item La concentración de los bienes
  \item La acumulación de los bienes
  \item La privatización de los bienes
  \item La estatización de los bienes
\end{multicols}
\end{enumerate}
\vspace{0.15cm}
\item ¿Qué es la equidad?
\begin{enumerate}[label=\textbf{\Alph*)}]
\begin{multicols}{2}
  \item La justicia basada en el reconocimiento de las diferencias
  \item La igualdad absoluta
  \item La uniformidad
  \item La homogeneidad
  \item La identidad
  \item La similitud
\end{multicols}
\end{enumerate}
\vspace{0.15cm}
\item ¿Qué es la inclusión social?
\begin{enumerate}[label=\textbf{\Alph*)}]
\begin{multicols}{2}
  \item La incorporación de todos los grupos sociales a la vida social
  \item La exclusión de algunos grupos sociales
  \item La marginación de algunos grupos sociales
  \item La discriminación de algunos grupos sociales
  \item La segregación de algunos grupos sociales
  \item La separación de algunos grupos sociales
\end{multicols}
\end{enumerate}
\vspace{0.15cm}
\item ¿Qué es la cohesión social?
\begin{enumerate}[label=\textbf{\Alph*)}]
\begin{multicols}{2}
  \item La unión y solidaridad entre los miembros de una sociedad
  \item La división social
  \item El conflicto social
  \item La estratificación social
  \item La movilidad social
  \item El cambio social
\end{multicols}
\end{enumerate}
\vspace{0.15cm}
\item ¿Qué es el capital social?
\begin{enumerate}[label=\textbf{\Alph*)}]
\begin{multicols}{2}
  \item El conjunto de redes y relaciones sociales basadas en la confianza
  \item El conjunto de bienes materiales
  \item El conjunto de bienes económicos
  \item El conjunto de bienes políticos
  \item El conjunto de bienes culturales
  \item El conjunto de bienes naturales
\end{multicols}
\end{enumerate}
\vspace{0.15cm}
\item ¿Qué es el desarrollo humano?
\begin{enumerate}[label=\textbf{\Alph*)}]
\begin{multicols}{2}
  \item El proceso de ampliación de las oportunidades y capacidades de las personas
  \item El crecimiento económico
  \item La industrialización
  \item La urbanización
  \item La globalización
  \item La modernización
\end{multicols}
\end{enumerate}
\vspace{0.15cm}
\item ¿Qué es el Índice de Desarrollo Humano (IDH)?
\begin{enumerate}[label=\textbf{\Alph*)}]
\begin{multicols}{2}
  \item Un indicador que mide el desarrollo humano de un país
  \item Un indicador económico
  \item Un indicador político
  \item Un indicador ambiental
  \item Un indicador cultural
  \item Un indicador religioso
\end{multicols}
\end{enumerate}
\vspace{0.15cm}
\item ¿Qué es la calidad de vida?
\begin{enumerate}[label=\textbf{\Alph*)}]
\begin{multicols}{2}
  \item El grado de bienestar y satisfacción de las necesidades humanas
  \item La cantidad de bienes materiales
  \item El nivel de ingresos
  \item El nivel educativo
  \item El nivel de salud
  \item El nivel de consumo
\end{multicols}
\end{enumerate}
\vspace{0.15cm}
\item ¿Qué es el bienestar social?
\begin{enumerate}[label=\textbf{\Alph*)}]
\begin{multicols}{2}
  \item El conjunto de condiciones que permiten la satisfacción de las necesidades humanas
  \item La cantidad de bienes materiales
  \item El nivel de ingresos
  \item El nivel educativo
  \item El nivel de salud
  \item El nivel de consumo
\end{multicols}
\end{enumerate}
\vspace{0.15cm}
\item ¿Qué es la política social?
\begin{enumerate}[label=\textbf{\Alph*)}]
\begin{multicols}{2}
  \item El conjunto de acciones del Estado para satisfacer necesidades sociales
  \item La política económica
  \item La política exterior
  \item La política militar
  \item La política cultural
  \item La política ambiental
\end{multicols}
\end{enumerate}
\vspace{0.15cm}
\item ¿Qué es la seguridad social?
\begin{enumerate}[label=\textbf{\Alph*)}]
\begin{multicols}{2}
  \item El sistema de protección social para los ciudadanos
  \item La seguridad nacional
  \item La seguridad pública
  \item La seguridad alimentaria
  \item La seguridad energética
  \item La seguridad ambiental
\end{multicols}
\end{enumerate}
\vspace{0.15cm}
\item ¿Qué es la educación?
\begin{enumerate}[label=\textbf{\Alph*)}]
\begin{multicols}{2}
  \item El proceso de transmisión de conocimientos, valores y habilidades
  \item El proceso de producción económica
  \item El proceso de urbanización
  \item El proceso de industrialización
  \item El proceso de globalización
  \item El proceso de modernización
\end{multicols}
\end{enumerate}
\vspace{0.15cm}
\item ¿Qué es la salud?
\begin{enumerate}[label=\textbf{\Alph*)}]
\begin{multicols}{2}
  \item El estado de bienestar físico, mental y social
  \item La ausencia de enfermedad
  \item La presencia de enfermedad
  \item El estado de bienestar económico
  \item El estado de bienestar político
  \item El estado de bienestar cultural
\end{multicols}
\end{enumerate}
\vspace{0.15cm}
\item ¿Qué es la vivienda?
\begin{enumerate}[label=\textbf{\Alph*)}]
\begin{multicols}{2}
  \item El espacio físico donde habita una persona o familia
  \item El espacio de trabajo
  \item El espacio de recreación
  \item El espacio de educación
  \item El espacio de salud
  \item El espacio de comercio
\end{multicols}
\end{enumerate}
\vspace{0.15cm}
\item ¿Qué es el empleo?
\begin{enumerate}[label=\textbf{\Alph*)}]
\begin{multicols}{2}
  \item La actividad remunerada que realiza una persona
  \item La actividad no remunerada
  \item La actividad de ocio
  \item La actividad de estudio
  \item La actividad de voluntariado
  \item La actividad doméstica
\end{multicols}
\end{enumerate}
\vspace{0.15cm}
\item ¿Qué es el desempleo?
\begin{enumerate}[label=\textbf{\Alph*)}]
\begin{multicols}{2}
  \item La situación de una persona que busca trabajo y no lo encuentra
  \item La situación de una persona que tiene trabajo
  \item La situación de una persona que estudia
  \item La situación de una persona que se jubila
  \item La situación de una persona que trabaja
  \item La situación de una persona que emprende
\end{multicols}
\end{enumerate}
\vspace{0.15cm}
\item ¿Qué es el subempleo?
\begin{enumerate}[label=\textbf{\Alph*)}]
\begin{multicols}{2}
  \item La situación de una persona que trabaja menos horas o en empleos de menor calificación
  \item La situación de una persona que trabaja a tiempo completo
  \item La situación de una persona que estudia
  \item La situación de una persona que se jubila
  \item La situación de una persona que emprende
  \item La situación de una persona que invierte
\end{multicols}
\end{enumerate}
\vspace{0.15cm}
\item ¿Qué es la economía informal?
\begin{enumerate}[label=\textbf{\Alph*)}]
\begin{multicols}{2}
  \item El conjunto de actividades económicas no reguladas ni registradas
  \item El conjunto de actividades económicas reguladas
  \item El conjunto de actividades económicas formales
  \item El conjunto de actividades económicas estatales
  \item El conjunto de actividades económicas privadas
  \item El conjunto de actividades económicas mixtas
\end{multicols}
\end{enumerate}
\vspace{0.15cm}
\item ¿Qué es la desigualdad social?
\begin{enumerate}[label=\textbf{\Alph*)}]
\begin{multicols}{2}
  \item La distribución inequitativa de recursos y oportunidades
  \item La igualdad social
  \item La cohesión social
  \item La movilidad social
  \item El cambio social
  \item La integración social
\end{multicols}
\end{enumerate}
\vspace{0.15cm}
\item ¿Qué es la pobreza?
\begin{enumerate}[label=\textbf{\Alph*)}]
\begin{multicols}{2}
  \item La carencia de recursos para satisfacer necesidades básicas
  \item La abundancia de recursos
  \item La igualdad social
  \item La cohesión social
  \item La movilidad social
  \item El cambio social
\end{multicols}
\end{enumerate}
\vspace{0.15cm}
\item ¿Qué es la pobreza extrema?
\begin{enumerate}[label=\textbf{\Alph*)}]
\begin{multicols}{2}
  \item La carencia de recursos para satisfacer necesidades alimentarias básicas
  \item La carencia de recursos para satisfacer necesidades no básicas
  \item La abundancia de recursos
  \item La igualdad social
  \item La cohesión social
  \item La movilidad social
\end{multicols}
\end{enumerate}
\vspace{0.15cm}
\item ¿Qué es la exclusión social?
\begin{enumerate}[label=\textbf{\Alph*)}]
\begin{multicols}{2}
  \item La marginación de individuos o grupos de la participación social
  \item La integración de individuos o grupos
  \item La inclusión de individuos o grupos
  \item La participación de individuos o grupos
  \item La incorporación de individuos o grupos
  \item La aceptación de individuos o grupos
\end{multicols}
\end{enumerate}
\vspace{0.15cm}
\item ¿Qué es la discriminación?
\begin{enumerate}[label=\textbf{\Alph*)}]
\begin{multicols}{2}
  \item El trato desigual e injusto hacia un grupo social
  \item El trato igualitario hacia un grupo social
  \item El trato justo hacia un grupo social
  \item El trato objetivo hacia un grupo social
  \item El trato científico hacia un grupo social
  \item El trato artístico hacia un grupo social
\end{multicols}
\end{enumerate}
\vspace{0.15cm}
\item ¿Qué es el racismo?
\begin{enumerate}[label=\textbf{\Alph*)}]
\begin{multicols}{2}
  \item La discriminación basada en la raza o etnia
  \item La discriminación basada en el género
  \item La discriminación basada en la edad
  \item La discriminación basada en la religión
  \item La discriminación basada en la clase social
  \item La discriminación basada en la nacionalidad
\end{multicols}
\end{enumerate}
\vspace{0.15cm}
\item ¿Qué es el sexismo?
\begin{enumerate}[label=\textbf{\Alph*)}]
\begin{multicols}{2}
  \item La discriminación basada en el sexo o género
  \item La discriminación basada en la raza
  \item La discriminación basada en la edad
  \item La discriminación basada en la religión
  \item La discriminación basada en la clase social
  \item La discriminación basada en la nacionalidad
\end{multicols}
\end{enumerate}
\vspace{0.15cm}
\item ¿Qué es la xenofobia?
\begin{enumerate}[label=\textbf{\Alph*)}]
\begin{multicols}{2}
  \item El rechazo hacia los extranjeros
  \item El rechazo hacia los nacionales
  \item El rechazo hacia los ancianos
  \item El rechazo hacia los niños
  \item El rechazo hacia los discapacitados
  \item El rechazo hacia los pobres
\end{multicols}
\end{enumerate}
\vspace{0.15cm}
\item ¿Qué es la homofobia?
\begin{enumerate}[label=\textbf{\Alph*)}]
\begin{multicols}{2}
  \item El rechazo hacia las personas homosexuales
  \item El rechazo hacia las personas heterosexuales
  \item El rechazo hacia las personas bisexuales
  \item El rechazo hacia las personas transgénero
  \item El rechazo hacia las personas intersexuales
  \item El rechazo hacia las personas asexuales
\end{multicols}
\end{enumerate}
\vspace{0.15cm}
\item ¿Qué es la violencia de género?
\begin{enumerate}[label=\textbf{\Alph*)}]
\begin{multicols}{2}
  \item La violencia ejercida contra una persona por su género
  \item La violencia ejercida contra una persona por su edad
  \item La violencia ejercida contra una persona por su raza
  \item La violencia ejercida contra una persona por su religión
  \item La violencia ejercida contra una persona por su clase social
  \item La violencia ejercida contra una persona por su nacionalidad
\end{multicols}
\end{enumerate}
\vspace{0.15cm}
\item ¿Qué es el feminicidio?
\begin{enumerate}[label=\textbf{\Alph*)}]
\begin{multicols}{2}
  \item El asesinato de mujeres por razones de género
  \item El asesinato de hombres por razones de género
  \item El asesinato de niños por razones de género
  \item El asesinato de ancianos por razones de género
  \item El asesinato de extranjeros por razones de género
  \item El asesinato de pobres por razones de género
\end{multicols}
\end{enumerate}
\vspace{0.15cm}
\item ¿Qué es la trata de personas?
\begin{enumerate}[label=\textbf{\Alph*)}]
\begin{multicols}{2}
  \item El comercio ilegal de seres humanos con fines de explotación
  \item El comercio legal de seres humanos
  \item El comercio de bienes
  \item El comercio de servicios
  \item El comercio de tecnología
  \item El comercio de alimentos
\end{multicols}
\end{enumerate}
\vspace{0.15cm}
\item ¿Qué es la explotación laboral?
\begin{enumerate}[label=\textbf{\Alph*)}]
\begin{multicols}{2}
  \item El abuso en las condiciones de trabajo
  \item El respeto a las condiciones de trabajo
  \item La mejora de las condiciones de trabajo
  \item La protección de las condiciones de trabajo
  \item La garantía de las condiciones de trabajo
  \item La promoción de las condiciones de trabajo
\end{multicols}
\end{enumerate}
\vspace{0.15cm}
\item ¿Qué es el trabajo infantil?
\begin{enumerate}[label=\textbf{\Alph*)}]
\begin{multicols}{2}
  \item La explotación de niños en actividades laborales
  \item La educación de niños
  \item La recreación de niños
  \item La protección de niños
  \item La alimentación de niños
  \item La salud de niños
\end{multicols}
\end{enumerate}
\vspace{0.15cm}
\end{enumerate}\n
\vspace{0.4cm}
\begin{tcolorbox}[colback=gray!5!white,colframe=gray!60!black,title=\textbf{SOLUCIONARIO -- 4.4 Derechos y Sociedad Actual}]
\begin{multicols}{4}
\begin{enumerate}[label=\textbf{\arabic*.}]
  \item \textbf{A}
  \item \textbf{A}
  \item \textbf{A}
  \item \textbf{A}
  \item \textbf{A}
  \item \textbf{A}
  \item \textbf{A}
  \item \textbf{A}
  \item \textbf{A}
  \item \textbf{A}
  \item \textbf{A}
  \item \textbf{A}
  \item \textbf{A}
  \item \textbf{A}
  \item \textbf{A}
  \item \textbf{A}
  \item \textbf{A}
  \item \textbf{A}
  \item \textbf{A}
  \item \textbf{A}
  \item \textbf{A}
  \item \textbf{A}
  \item \textbf{A}
  \item \textbf{A}
  \item \textbf{A}
  \item \textbf{A}
  \item \textbf{A}
  \item \textbf{A}
  \item \textbf{A}
  \item \textbf{A}
  \item \textbf{A}
  \item \textbf{A}
  \item \textbf{A}
  \item \textbf{A}
  \item \textbf{A}
  \item \textbf{A}
  \item \textbf{A}
  \item \textbf{A}
  \item \textbf{A}
  \item \textbf{A}
  \item \textbf{A}
  \item \textbf{A}
  \item \textbf{A}
  \item \textbf{A}
  \item \textbf{A}
  \item \textbf{A}
  \item \textbf{A}
  \item \textbf{A}
  \item \textbf{A}
  \item \textbf{A}
  \item \textbf{A}
  \item \textbf{A}
  \item \textbf{A}
  \item \textbf{A}
  \item \textbf{A}
  \item \textbf{A}
  \item \textbf{A}
  \item \textbf{A}
  \item \textbf{A}
  \item \textbf{A}
  \item \textbf{A}
  \item \textbf{A}
  \item \textbf{A}
  \item \textbf{A}
  \item \textbf{A}
  \item \textbf{A}
  \item \textbf{A}
  \item \textbf{A}
  \item \textbf{A}
  \item \textbf{A}
  \item \textbf{A}
  \item \textbf{A}
  \item \textbf{A}
  \item \textbf{A}
  \item \textbf{A}
  \item \textbf{A}
  \item \textbf{A}
  \item \textbf{A}
  \item \textbf{A}
  \item \textbf{A}
  \item \textbf{A}
  \item \textbf{A}
  \item \textbf{A}
  \item \textbf{A}
  \item \textbf{A}
  \item \textbf{A}
  \item \textbf{A}
  \item \textbf{A}
  \item \textbf{A}
  \item \textbf{A}
  \item \textbf{A}
  \item \textbf{A}
  \item \textbf{A}
  \item \textbf{A}
  \item \textbf{A}
  \item \textbf{A}
  \item \textbf{A}
  \item \textbf{A}
  \item \textbf{A}
  \item \textbf{A}
  \item \textbf{A}
  \item \textbf{A}
  \item \textbf{A}
  \item \textbf{A}
  \item \textbf{A}
  \item \textbf{A}
  \item \textbf{A}
  \item \textbf{A}
  \item \textbf{A}
  \item \textbf{A}
  \item \textbf{A}
  \item \textbf{A}
  \item \textbf{A}
  \item \textbf{A}
  \item \textbf{A}
  \item \textbf{A}
  \item \textbf{A}
  \item \textbf{A}
  \item \textbf{A}
  \item \textbf{A}
  \item \textbf{A}
  \item \textbf{A}
  \item \textbf{A}
  \item \textbf{A}
  \item \textbf{A}
  \item \textbf{A}
  \item \textbf{A}
  \item \textbf{A}
  \item \textbf{A}
  \item \textbf{A}
  \item \textbf{A}
  \item \textbf{A}
  \item \textbf{A}
  \item \textbf{A}
  \item \textbf{A}
  \item \textbf{A}
  \item \textbf{A}
  \item \textbf{A}
  \item \textbf{A}
  \item \textbf{A}
  \item \textbf{A}
  \item \textbf{A}
  \item \textbf{A}
  \item \textbf{A}
  \item \textbf{A}
  \item \textbf{A}
  \item \textbf{A}
  \item \textbf{A}
  \item \textbf{A}
  \item \textbf{A}
  \item \textbf{A}
  \item \textbf{A}
  \item \textbf{A}
  \item \textbf{A}
  \item \textbf{A}
  \item \textbf{A}
  \item \textbf{A}
  \item \textbf{A}
  \item \textbf{A}
  \item \textbf{A}
  \item \textbf{A}
  \item \textbf{A}
  \item \textbf{A}
  \item \textbf{A}
  \item \textbf{A}
  \item \textbf{A}
  \item \textbf{A}
  \item \textbf{A}
  \item \textbf{A}
  \item \textbf{A}
  \item \textbf{A}
  \item \textbf{A}
  \item \textbf{A}
  \item \textbf{A}
  \item \textbf{A}
  \item \textbf{A}
  \item \textbf{A}
  \item \textbf{A}
  \item \textbf{A}
  \item \textbf{A}
  \item \textbf{A}
  \item \textbf{A}
  \item \textbf{A}
  \item \textbf{A}
  \item \textbf{A}
  \item \textbf{A}
  \item \textbf{A}
  \item \textbf{A}
  \item \textbf{A}
  \item \textbf{A}
  \item \textbf{A}
  \item \textbf{A}
  \item \textbf{A}
  \item \textbf{A}
  \item \textbf{A}
  \item \textbf{A}
  \item \textbf{A}
  \item \textbf{A}
  \item \textbf{A}
  \item \textbf{A}
  \item \textbf{A}
  \item \textbf{A}
  \item \textbf{A}
  \item \textbf{A}
  \item \textbf{A}
  \item \textbf{A}
  \item \textbf{A}
  \item \textbf{A}
  \item \textbf{A}
  \item \textbf{A}
  \item \textbf{A}
  \item \textbf{A}
  \item \textbf{A}
  \item \textbf{A}
  \item \textbf{A}
  \item \textbf{A}
  \item \textbf{A}
  \item \textbf{A}
  \item \textbf{A}
  \item \textbf{A}
  \item \textbf{A}
  \item \textbf{A}
  \item \textbf{A}
  \item \textbf{A}
  \item \textbf{A}
  \item \textbf{A}
  \item \textbf{A}
  \item \textbf{A}
  \item \textbf{A}
  \item \textbf{A}
  \item \textbf{A}
  \item \textbf{A}
  \item \textbf{A}
  \item \textbf{A}
  \item \textbf{A}
  \item \textbf{A}
  \item \textbf{A}
  \item \textbf{A}
  \item \textbf{A}
  \item \textbf{A}
  \item \textbf{A}
  \item \textbf{A}
  \item \textbf{A}
  \item \textbf{A}
  \item \textbf{A}
  \item \textbf{A}
  \item \textbf{A}
  \item \textbf{A}
  \item \textbf{A}
  \item \textbf{A}
  \item \textbf{A}
  \item \textbf{A}
  \item \textbf{A}
  \item \textbf{A}
  \item \textbf{A}
  \item \textbf{A}
  \item \textbf{A}
  \item \textbf{A}
  \item \textbf{A}
  \item \textbf{A}
  \item \textbf{A}
  \item \textbf{A}
  \item \textbf{A}
  \item \textbf{A}
  \item \textbf{A}
  \item \textbf{A}
  \item \textbf{A}
  \item \textbf{A}
  \item \textbf{A}
  \item \textbf{A}
  \item \textbf{A}
  \item \textbf{A}
  \item \textbf{A}
  \item \textbf{A}
  \item \textbf{A}
  \item \textbf{A}
  \item \textbf{A}
  \item \textbf{A}
  \item \textbf{A}
  \item \textbf{A}
  \item \textbf{A}
  \item \textbf{A}
  \item \textbf{A}
  \item \textbf{A}
  \item \textbf{A}
  \item \textbf{A}
  \item \textbf{A}
  \item \textbf{A}
  \item \textbf{A}
  \item \textbf{A}
  \item \textbf{A}
  \item \textbf{A}
  \item \textbf{A}
  \item \textbf{A}
  \item \textbf{A}
  \item \textbf{A}
  \item \textbf{A}
  \item \textbf{A}
  \item \textbf{A}
  \item \textbf{A}
\end{enumerate}
\end{multicols}
\end{tcolorbox}
\newpage\n
% ==============================================================================
% MATERIA 5: CIENCIAS NATURALES
% ==============================================================================
\section{Ciencias Naturales (Módulos 5.1 al 5.4)}

\subsection{5.1 Biología y Seres Vivos}

\subsubsection*{SECCIÓN I: PROPIEDADES DE LA MATERIA VIVA (1–50)}

\begin{enumerate}[label=\textbf{\arabic*.}]
\setcounter{enumi}{0}
\item ¿Qué es la biología?
\begin{enumerate}[label=\textbf{\Alph*)}]
\begin{multicols}{2}
  \item La ciencia que estudia los seres vivos
  \item La ciencia que estudia los minerales
  \item La ciencia que estudia los astros
  \item La ciencia que estudia los fenómenos físicos
  \item La ciencia que estudia las reacciones químicas
  \item La ciencia que estudia los números
\end{multicols}
\end{enumerate}
\vspace{0.15cm}
\item ¿Cuál es la característica principal de los seres vivos?
\begin{enumerate}[label=\textbf{\Alph*)}]
\begin{multicols}{2}
  \item Están formados por células
  \item Están formados por minerales
  \item Están formados por átomos metálicos
  \item Están formados por gases
  \item Están formados por líquidos
  \item Están formados por plásticos
\end{multicols}
\end{enumerate}
\vspace{0.15cm}
\item ¿Qué es la célula?
\begin{enumerate}[label=\textbf{\Alph*)}]
\begin{multicols}{2}
  \item La unidad estructural y funcional de los seres vivos
  \item La unidad estructural de los minerales
  \item La unidad funcional de los astros
  \item La unidad estructural del agua
  \item La unidad funcional del aire
  \item La unidad estructural del suelo
\end{multicols}
\end{enumerate}
\vspace{0.15cm}
\item ¿Quién descubrió la célula?
\begin{enumerate}[label=\textbf{\Alph*)}]
\begin{multicols}{2}
  \item Robert Hooke
  \item Charles Darwin
  \item Gregor Mendel
  \item Louis Pasteur
  \item Alexander Fleming
  \item Isaac Newton
\end{multicols}
\end{enumerate}
\vspace{0.15cm}
\item ¿Qué es la teoría celular?
\begin{enumerate}[label=\textbf{\Alph*)}]
\begin{multicols}{2}
  \item Todos los seres vivos están formados por células
  \item Todos los seres vivos están formados por minerales
  \item Todos los seres vivos están formados por átomos
  \item Todos los seres vivos están formados por moléculas
  \item Todos los seres vivos están formados por gases
  \item Todos los seres vivos están formados por líquidos
\end{multicols}
\end{enumerate}
\vspace{0.15cm}
\item ¿Cuál de los siguientes postulados corresponde a la teoría celular?
\begin{enumerate}[label=\textbf{\Alph*)}]
\begin{multicols}{2}
  \item La célula es la unidad básica de la vida
  \item La célula es la unidad básica de los minerales
  \item La célula es la unidad básica del agua
  \item La célula es la unidad básica del aire
  \item La célula es la unidad básica del suelo
  \item La célula es la unidad básica del fuego
\end{multicols}
\end{enumerate}
\vspace{0.15cm}
\item ¿Qué es el metabolismo?
\begin{enumerate}[label=\textbf{\Alph*)}]
\begin{multicols}{2}
  \item El conjunto de reacciones químicas que ocurren en los seres vivos
  \item El conjunto de reacciones físicas que ocurren en los minerales
  \item El conjunto de reacciones químicas que ocurren en el agua
  \item El conjunto de reacciones físicas que ocurren en el aire
  \item El conjunto de reacciones químicas que ocurren en el suelo
  \item El conjunto de reacciones físicas que ocurren en los astros
\end{multicols}
\end{enumerate}
\vspace{0.15cm}
\item ¿Qué es el anabolismo?
\begin{enumerate}[label=\textbf{\Alph*)}]
\begin{multicols}{2}
  \item El conjunto de reacciones de síntesis que construyen moléculas
  \item El conjunto de reacciones de degradación
  \item El conjunto de reacciones de oxidación
  \item El conjunto de reacciones de reducción
  \item El conjunto de reacciones de hidrólisis
  \item El conjunto de reacciones de combustión
\end{multicols}
\end{enumerate}
\vspace{0.15cm}
\item ¿Qué es el catabolismo?
\begin{enumerate}[label=\textbf{\Alph*)}]
\begin{multicols}{2}
  \item El conjunto de reacciones de degradación que liberan energía
  \item El conjunto de reacciones de síntesis
  \item El conjunto de reacciones de construcción
  \item El conjunto de reacciones de formación
  \item El conjunto de reacciones de producción
  \item El conjunto de reacciones de elaboración
\end{multicols}
\end{enumerate}
\vspace{0.15cm}
\item ¿Qué es la homeostasis?
\begin{enumerate}[label=\textbf{\Alph*)}]
\begin{multicols}{2}
  \item La capacidad de mantener el equilibrio interno
  \item La capacidad de mantener el desequilibrio interno
  \item La capacidad de mantener el equilibrio externo
  \item La capacidad de mantener el desequilibrio externo
  \item La capacidad de cambiar el equilibrio interno
  \item La capacidad de cambiar el equilibrio externo
\end{multicols}
\end{enumerate}
\vspace{0.15cm}
\item ¿Cuál es un ejemplo de homeostasis?
\begin{enumerate}[label=\textbf{\Alph*)}]
\begin{multicols}{2}
  \item La regulación de la temperatura corporal
  \item La regulación de la temperatura ambiental
  \item La regulación de la temperatura del agua
  \item La regulación de la temperatura del aire
  \item La regulación de la temperatura del suelo
  \item La regulación de la temperatura de los astros
\end{multicols}
\end{enumerate}
\vspace{0.15cm}
\item ¿Qué es la irritabilidad?
\begin{enumerate}[label=\textbf{\Alph*)}]
\begin{multicols}{2}
  \item La capacidad de responder a estímulos
  \item La capacidad de no responder a estímulos
  \item La capacidad de ignorar estímulos
  \item La capacidad de olvidar estímulos
  \item La capacidad de negar estímulos
  \item La capacidad de ocultar estímulos
\end{multicols}
\end{enumerate}
\vspace{0.15cm}
\item ¿Qué es la reproducción?
\begin{enumerate}[label=\textbf{\Alph*)}]
\begin{multicols}{2}
  \item La capacidad de generar nuevos individuos
  \item La capacidad de generar minerales
  \item La capacidad de generar agua
  \item La capacidad de generar aire
  \item La capacidad de generar suelo
  \item La capacidad de generar astros
\end{multicols}
\end{enumerate}
\vspace{0.15cm}
\item ¿Qué es la reproducción asexual?
\begin{enumerate}[label=\textbf{\Alph*)}]
\begin{multicols}{2}
  \item La que produce descendientes genéticamente idénticos al progenitor
  \item La que produce descendientes genéticamente diferentes
  \item La que requiere dos progenitores
  \item La que requiere tres progenitores
  \item La que requiere cuatro progenitores
  \item La que requiere cinco progenitores
\end{multicols}
\end{enumerate}
\vspace{0.15cm}
\item ¿Qué es la reproducción sexual?
\begin{enumerate}[label=\textbf{\Alph*)}]
\begin{multicols}{2}
  \item La que requiere dos progenitores y produce descendientes diferentes
  \item La que requiere un progenitor
  \item La que produce descendientes idénticos
  \item La que no requiere progenitores
  \item La que requiere tres progenitores
  \item La que requiere cuatro progenitores
\end{multicols}
\end{enumerate}
\vspace{0.15cm}
\item ¿Qué es la adaptación?
\begin{enumerate}[label=\textbf{\Alph*)}]
\begin{multicols}{2}
  \item La capacidad de los seres vivos de ajustarse al ambiente
  \item La capacidad de los seres vivos de ignorar el ambiente
  \item La capacidad de los seres vivos de cambiar el ambiente
  \item La capacidad de los seres vivos de destruir el ambiente
  \item La capacidad de los seres vivos de olvidar el ambiente
  \item La capacidad de los seres vivos de negar el ambiente
\end{multicols}
\end{enumerate}
\vspace{0.15cm}
\item ¿Qué es la evolución?
\begin{enumerate}[label=\textbf{\Alph*)}]
\begin{multicols}{2}
  \item El cambio de las especies a lo largo del tiempo
  \item La permanencia de las especies
  \item La desaparición de las especies
  \item La creación de las especies
  \item La transformación de las especies en minerales
  \item La transformación de las especies en agua
\end{multicols}
\end{enumerate}
\vspace{0.15cm}
\item ¿Quién propuso la teoría de la evolución por selección natural?
\begin{enumerate}[label=\textbf{\Alph*)}]
\begin{multicols}{2}
  \item Charles Darwin
  \item Gregor Mendel
  \item Louis Pasteur
  \item Robert Hooke
  \item Alexander Fleming
  \item Isaac Newton
\end{multicols}
\end{enumerate}
\vspace{0.15cm}
\item ¿Qué es la selección natural?
\begin{enumerate}[label=\textbf{\Alph*)}]
\begin{multicols}{2}
  \item El proceso por el cual los individuos mejor adaptados sobreviven y se reproducen
  \item El proceso por el cual los individuos peor adaptados sobreviven
  \item El proceso por el cual los individuos no se reproducen
  \item El proceso por el cual los individuos no se adaptan
  \item El proceso por el cual los individuos no evolucionan
  \item El proceso por el cual los individuos no cambian
\end{multicols}
\end{enumerate}
\vspace{0.15cm}
\item ¿Qué es el ADN?
\begin{enumerate}[label=\textbf{\Alph*)}]
\begin{multicols}{2}
  \item El ácido desoxirribonucleico que contiene la información genética
  \item El ácido ribonucleico que contiene la información genética
  \item El ácido desoxirribonucleico que contiene la información proteica
  \item El ácido ribonucleico que contiene la información proteica
  \item El ácido desoxirribonucleico que contiene la información lipídica
  \item El ácido ribonucleico que contiene la información lipídica
\end{multicols}
\end{enumerate}
\vspace{0.15cm}
\item ¿Dónde se encuentra el ADN en las células eucariotas?
\begin{enumerate}[label=\textbf{\Alph*)}]
\begin{multicols}{2}
  \item En el núcleo
  \item En el citoplasma
  \item En la membrana
  \item En la pared celular
  \item En los ribosomas
  \item En las mitocondrias
\end{multicols}
\end{enumerate}
\vspace{0.15cm}
\item ¿Qué es el ARN?
\begin{enumerate}[label=\textbf{\Alph*)}]
\begin{multicols}{2}
  \item El ácido ribonucleico que participa en la síntesis de proteínas
  \item El ácido desoxirribonucleico que participa en la síntesis de proteínas
  \item El ácido ribonucleico que participa en la síntesis de lípidos
  \item El ácido desoxirribonucleico que participa en la síntesis de lípidos
  \item El ácido ribonucleico que participa en la síntesis de carbohidratos
  \item El ácido desoxirribonucleico que participa en la síntesis de carbohidratos
\end{multicols}
\end{enumerate}
\vspace{0.15cm}
\item ¿Qué es un gen?
\begin{enumerate}[label=\textbf{\Alph*)}]
\begin{multicols}{2}
  \item Un segmento de ADN que contiene información para un rasgo
  \item Un segmento de ARN que contiene información para un rasgo
  \item Un segmento de proteína que contiene información para un rasgo
  \item Un segmento de lípido que contiene información para un rasgo
  \item Un segmento de carbohidrato que contiene información para un rasgo
  \item Un segmento de mineral que contiene información para un rasgo
\end{multicols}
\end{enumerate}
\vspace{0.15cm}
\item ¿Qué es el genoma?
\begin{enumerate}[label=\textbf{\Alph*)}]
\begin{multicols}{2}
  \item El conjunto completo de genes de un organismo
  \item El conjunto completo de proteínas de un organismo
  \item El conjunto completo de lípidos de un organismo
  \item El conjunto completo de carbohidratos de un organismo
  \item El conjunto completo de minerales de un organismo
  \item El conjunto completo de células de un organismo
\end{multicols}
\end{enumerate}
\vspace{0.15cm}
\item ¿Qué es la nutrición?
\begin{enumerate}[label=\textbf{\Alph*)}]
\begin{multicols}{2}
  \item El proceso de obtener y utilizar nutrientes
  \item El proceso de obtener y utilizar minerales
  \item El proceso de obtener y utilizar gases
  \item El proceso de obtener y utilizar líquidos
  \item El proceso de obtener y utilizar astros
  \item El proceso de obtener y utilizar rocas
\end{multicols}
\end{enumerate}
\vspace{0.15cm}
\item ¿Qué es la nutrición autótrofa?
\begin{enumerate}[label=\textbf{\Alph*)}]
\begin{multicols}{2}
  \item La que produce su propio alimento
  \item La que consume alimento de otros
  \item La que no necesita alimento
  \item La que no produce alimento
  \item La que no consume alimento
  \item La que no utiliza alimento
\end{multicols}
\end{enumerate}
\vspace{0.15cm}
\item ¿Qué es la nutrición heterótrofa?
\begin{enumerate}[label=\textbf{\Alph*)}]
\begin{multicols}{2}
  \item La que consume alimento producido por otros
  \item La que produce su propio alimento
  \item La que no necesita alimento
  \item La que no produce alimento
  \item La que no consume alimento
  \item La que no utiliza alimento
\end{multicols}
\end{enumerate}
\vspace{0.15cm}
\item ¿Qué es la fotosíntesis?
\begin{enumerate}[label=\textbf{\Alph*)}]
\begin{multicols}{2}
  \item El proceso por el cual las plantas producen alimento usando luz solar
  \item El proceso por el cual las plantas producen alimento usando luz lunar
  \item El proceso por el cual las plantas producen alimento usando luz artificial
  \item El proceso por el cual las plantas producen alimento usando luz ultravioleta
  \item El proceso por el cual las plantas producen alimento usando luz infrarroja
  \item El proceso por el cual las plantas producen alimento usando luz gamma
\end{multicols}
\end{enumerate}
\vspace{0.15cm}
\item ¿Qué es la respiración celular?
\begin{enumerate}[label=\textbf{\Alph*)}]
\begin{multicols}{2}
  \item El proceso por el cual las células obtienen energía de los nutrientes
  \item El proceso por el cual las células obtienen energía de la luz
  \item El proceso por el cual las células obtienen energía del agua
  \item El proceso por el cual las células obtienen energía del aire
  \item El proceso por el cual las células obtienen energía del suelo
  \item El proceso por el cual las células obtienen energía de los minerales
\end{multicols}
\end{enumerate}
\vspace{0.15cm}
\item ¿Qué es la mitosis?
\begin{enumerate}[label=\textbf{\Alph*)}]
\begin{multicols}{2}
  \item El proceso de división celular que produce dos células idénticas
  \item El proceso de división celular que produce cuatro células diferentes
  \item El proceso de división celular que produce dos células diferentes
  \item El proceso de división celular que produce cuatro células idénticas
  \item El proceso de división celular que produce una célula
  \item El proceso de división celular que produce tres células
\end{multicols}
\end{enumerate}
\vspace{0.15cm}
\item ¿Qué es la meiosis?
\begin{enumerate}[label=\textbf{\Alph*)}]
\begin{multicols}{2}
  \item El proceso de división celular que produce cuatro células diferentes
  \item El proceso de división celular que produce dos células idénticas
  \item El proceso de división celular que produce dos células diferentes
  \item El proceso de división celular que produce cuatro células idénticas
  \item El proceso de división celular que produce una célula
  \item El proceso de división celular que produce tres células
\end{multicols}
\end{enumerate}
\vspace{0.15cm}
\item ¿Cuál es la función de la mitosis?
\begin{enumerate}[label=\textbf{\Alph*)}]
\begin{multicols}{2}
  \item El crecimiento y reparación de tejidos
  \item La producción de gametos
  \item La producción de células diferentes
  \item La producción de células sexuales
  \item La producción de células reproductoras
  \item La producción de células germinales
\end{multicols}
\end{enumerate}
\vspace{0.15cm}
\item ¿Cuál es la función de la meiosis?
\begin{enumerate}[label=\textbf{\Alph*)}]
\begin{multicols}{2}
  \item La producción de gametos
  \item El crecimiento y reparación de tejidos
  \item La producción de células idénticas
  \item La producción de células somáticas
  \item La producción de células corporales
  \item La producción de células diploides
\end{multicols}
\end{enumerate}
\vspace{0.15cm}
\item ¿Qué es la célula procariota?
\begin{enumerate}[label=\textbf{\Alph*)}]
\begin{multicols}{2}
  \item La que no tiene núcleo definido
  \item La que tiene núcleo definido
  \item La que tiene dos núcleos
  \item La que tiene tres núcleos
  \item La que tiene cuatro núcleos
  \item La que no tiene membrana
\end{multicols}
\end{enumerate}
\vspace{0.15cm}
\item ¿Qué es la célula eucariota?
\begin{enumerate}[label=\textbf{\Alph*)}]
\begin{multicols}{2}
  \item La que tiene núcleo definido
  \item La que no tiene núcleo definido
  \item La que tiene dos núcleos
  \item La que tiene tres núcleos
  \item La que tiene cuatro núcleos
  \item La que no tiene membrana
\end{multicols}
\end{enumerate}
\vspace{0.15cm}
\item ¿Cuál es un ejemplo de célula procariota?
\begin{enumerate}[label=\textbf{\Alph*)}]
\begin{multicols}{2}
  \item La bacteria
  \item La célula animal
  \item La célula vegetal
  \item La célula fúngica
  \item La célula protista
  \item La célula nerviosa
\end{multicols}
\end{enumerate}
\vspace{0.15cm}
\item ¿Cuál es un ejemplo de célula eucariota?
\begin{enumerate}[label=\textbf{\Alph*)}]
\begin{multicols}{2}
  \item La célula animal
  \item La bacteria
  \item La cianobacteria
  \item La arquea
  \item La bacteria intestinal
  \item La bacteria del suelo
\end{multicols}
\end{enumerate}
\vspace{0.15cm}
\item ¿Qué es la membrana celular?
\begin{enumerate}[label=\textbf{\Alph*)}]
\begin{multicols}{2}
  \item La estructura que delimita la célula y regula el paso de sustancias
  \item La estructura que no delimita la célula
  \item La estructura que no regula el paso de sustancias
  \item La estructura que delimita el núcleo
  \item La estructura que delimita la mitocondria
  \item La estructura que delimita el ribosoma
\end{multicols}
\end{enumerate}
\vspace{0.15cm}
\item ¿Qué es el citoplasma?
\begin{enumerate}[label=\textbf{\Alph*)}]
\begin{multicols}{2}
  \item El medio interno de la célula donde ocurren las reacciones
  \item El medio externo de la célula
  \item El medio interno del núcleo
  \item El medio externo del núcleo
  \item El medio interno de la mitocondria
  \item El medio externo de la mitocondria
\end{multicols}
\end{enumerate}
\vspace{0.15cm}
\item ¿Qué es el núcleo celular?
\begin{enumerate}[label=\textbf{\Alph*)}]
\begin{multicols}{2}
  \item La estructura que contiene el material genético
  \item La estructura que contiene el material proteico
  \item La estructura que contiene el material lipídico
  \item La estructura que contiene el material carbohidrático
  \item La estructura que contiene el material mineral
  \item La estructura que contiene el material gaseoso
\end{multicols}
\end{enumerate}
\vspace{0.15cm}
\item ¿Qué es la mitocondria?
\begin{enumerate}[label=\textbf{\Alph*)}]
\begin{multicols}{2}
  \item El orgánulo que produce energía
  \item El orgánulo que produce proteínas
  \item El orgánulo que produce lípidos
  \item El orgánulo que produce carbohidratos
  \item El orgánulo que produce minerales
  \item El orgánulo que produce gases
\end{multicols}
\end{enumerate}
\vspace{0.15cm}
\item ¿Qué es el ribosoma?
\begin{enumerate}[label=\textbf{\Alph*)}]
\begin{multicols}{2}
  \item El orgánulo que sintetiza proteínas
  \item El orgánulo que sintetiza lípidos
  \item El orgánulo que sintetiza carbohidratos
  \item El orgánulo que sintetiza minerales
  \item El orgánulo que sintetiza gases
  \item El orgánulo que sintetiza agua
\end{multicols}
\end{enumerate}
\vspace{0.15cm}
\item ¿Qué es el cloroplasto?
\begin{enumerate}[label=\textbf{\Alph*)}]
\begin{multicols}{2}
  \item El orgánulo que realiza la fotosíntesis
  \item El orgánulo que realiza la respiración
  \item El orgánulo que realiza la digestión
  \item El orgánulo que realiza la excreción
  \item El orgánulo que realiza la reproducción
  \item El orgánulo que realiza la síntesis
\end{multicols}
\end{enumerate}
\vspace{0.15cm}
\item ¿Qué es la pared celular?
\begin{enumerate}[label=\textbf{\Alph*)}]
\begin{multicols}{2}
  \item La estructura que rodea la membrana de células vegetales
  \item La estructura que rodea la membrana de células animales
  \item La estructura que rodea el núcleo
  \item La estructura que rodea la mitocondria
  \item La estructura que rodea el ribosoma
  \item La estructura que rodea el cloroplasto
\end{multicols}
\end{enumerate}
\vspace{0.15cm}
\item ¿Qué es el aparato de Golgi?
\begin{enumerate}[label=\textbf{\Alph*)}]
\begin{multicols}{2}
  \item El orgánulo que modifica y empaqueta proteínas
  \item El orgánulo que modifica y empaqueta lípidos
  \item El orgánulo que modifica y empaqueta carbohidratos
  \item El orgánulo que modifica y empaqueta minerales
  \item El orgánulo que modifica y empaqueta gases
  \item El orgánulo que modifica y empaqueta agua
\end{multicols}
\end{enumerate}
\vspace{0.15cm}
\item ¿Qué es el retículo endoplasmático?
\begin{enumerate}[label=\textbf{\Alph*)}]
\begin{multicols}{2}
  \item La red de membranas que transporta sustancias
  \item La red de membranas que no transporta sustancias
  \item La red de membranas que delimita el núcleo
  \item La red de membranas que delimita la mitocondria
  \item La red de membranas que delimita el ribosoma
  \item La red de membranas que delimita el cloroplasto
\end{multicols}
\end{enumerate}
\vspace{0.15cm}
\item ¿Qué es el lisosoma?
\begin{enumerate}[label=\textbf{\Alph*)}]
\begin{multicols}{2}
  \item El orgánulo que digiere sustancias
  \item El orgánulo que sintetiza sustancias
  \item El orgánulo que transporta sustancias
  \item El orgánulo que almacena sustancias
  \item El orgánulo que excreta sustancias
  \item El orgánulo que absorbe sustancias
\end{multicols}
\end{enumerate}
\vspace{0.15cm}
\item ¿Qué es el citoesqueleto?
\begin{enumerate}[label=\textbf{\Alph*)}]
\begin{multicols}{2}
  \item La red de proteínas que da forma y movimiento a la célula
  \item La red de lípidos que da forma a la célula
  \item La red de carbohidratos que da forma a la célula
  \item La red de minerales que da forma a la célula
  \item La red de gases que da forma a la célula
  \item La red de agua que da forma a la célula
\end{multicols}
\end{enumerate}
\vspace{0.15cm}
\item ¿Qué es la vacuola?
\begin{enumerate}[label=\textbf{\Alph*)}]
\begin{multicols}{2}
  \item El orgánulo que almacena sustancias
  \item El orgánulo que sintetiza sustancias
  \item El orgánulo que transporta sustancias
  \item El orgánulo que digiere sustancias
  \item El orgánulo que excreta sustancias
  \item El orgánulo que absorbe sustancias
\end{multicols}
\end{enumerate}
\vspace{0.15cm}
\item ¿Qué es el centriolo?
\begin{enumerate}[label=\textbf{\Alph*)}]
\begin{multicols}{2}
  \item El orgánulo que participa en la división celular
  \item El orgánulo que participa en la síntesis de proteínas
  \item El orgánulo que participa en la síntesis de lípidos
  \item El orgánulo que participa en la fotosíntesis
  \item El orgánulo que participa en la respiración
  \item El orgánulo que participa en la digestión ---
\end{multicols}
\end{enumerate}
\vspace{0.15cm}
\end{enumerate}\n
\subsubsection*{SECCIÓN II: CICLO CELULAR – MITOSIS Y MEIOSIS (51–120)}

\begin{enumerate}[label=\textbf{\arabic*.}]
\setcounter{enumi}{50}
\item ¿Qué es el ciclo celular?
\begin{enumerate}[label=\textbf{\Alph*)}]
\begin{multicols}{2}
  \item La secuencia de etapas por las que pasa una célula desde que nace hasta que se divide
  \item La secuencia de etapas por las que pasa una célula desde que muere hasta que nace
  \item La secuencia de etapas por las que pasa una célula sin dividirse
  \item La secuencia de etapas por las que pasa una célula sin crecer
  \item La secuencia de etapas por las que pasa una célula sin reproducirse
  \item La secuencia de etapas por las que pasa una célula sin cambiar
\end{multicols}
\end{enumerate}
\vspace{0.15cm}
\item ¿Cuáles son las fases del ciclo celular?
\begin{enumerate}[label=\textbf{\Alph*)}]
\begin{multicols}{2}
  \item Interfase y fase mitótica
  \item Interfase y fase meiótica
  \item Profase y metafase
  \item Anafase y telofase
  \item G1 y G2
  \item S y M
\end{multicols}
\end{enumerate}
\vspace{0.15cm}
\item ¿Qué es la interfase?
\begin{enumerate}[label=\textbf{\Alph*)}]
\begin{multicols}{2}
  \item La fase de crecimiento y preparación para la división
  \item La fase de división celular
  \item La fase de muerte celular
  \item La fase de reposo celular
  \item La fase de diferenciación celular
  \item La fase de especialización celular
\end{multicols}
\end{enumerate}
\vspace{0.15cm}
\item ¿Cuáles son las subfases de la interfase?
\begin{enumerate}[label=\textbf{\Alph*)}]
\begin{multicols}{2}
  \item G1, S y G2
  \item Profase, metafase y anafase
  \item Telofase, citocinesis y G0
  \item G0, G1 y G2
  \item S, G2 y M
  \item M, G1 y S
\end{multicols}
\end{enumerate}
\vspace{0.15cm}
\item ¿Qué ocurre en la fase G1?
\begin{enumerate}[label=\textbf{\Alph*)}]
\begin{multicols}{2}
  \item La célula crece y sintetiza proteínas
  \item La célula duplica su ADN
  \item La célula se divide
  \item La célula muere
  \item La célula se diferencia
  \item La célula se especializa
\end{multicols}
\end{enumerate}
\vspace{0.15cm}
\item ¿Qué ocurre en la fase S?
\begin{enumerate}[label=\textbf{\Alph*)}]
\begin{multicols}{2}
  \item La célula duplica su ADN
  \item La célula crece
  \item La célula se divide
  \item La célula muere
  \item La célula se diferencia
  \item La célula se especializa
\end{multicols}
\end{enumerate}
\vspace{0.15cm}
\item ¿Qué ocurre en la fase G2?
\begin{enumerate}[label=\textbf{\Alph*)}]
\begin{multicols}{2}
  \item La célula se prepara para la división
  \item La célula duplica su ADN
  \item La célula crece
  \item La célula muere
  \item La célula se diferencia
  \item La célula se especializa
\end{multicols}
\end{enumerate}
\vspace{0.15cm}
\item ¿Qué es la fase M?
\begin{enumerate}[label=\textbf{\Alph*)}]
\begin{multicols}{2}
  \item La fase de división celular (mitosis y citocinesis)
  \item La fase de crecimiento celular
  \item La fase de duplicación del ADN
  \item La fase de muerte celular
  \item La fase de reposo celular
  \item La fase de diferenciación celular
\end{multicols}
\end{enumerate}
\vspace{0.15cm}
\item ¿Cuáles son las etapas de la mitosis?
\begin{enumerate}[label=\textbf{\Alph*)}]
\begin{multicols}{2}
  \item Profase, metafase, anafase y telofase
  \item Profase, metafase y anafase
  \item Profase, anafase y telofase
  \item Metafase, anafase y telofase
  \item Profase, metafase, anafase, telofase y citocinesis
  \item Interfase, profase, metafase y anafase
\end{multicols}
\end{enumerate}
\vspace{0.15cm}
\item ¿Qué ocurre en la profase?
\begin{enumerate}[label=\textbf{\Alph*)}]
\begin{multicols}{2}
  \item La cromatina se condensa en cromosomas y desaparece el nucléolo
  \item Los cromosomas se alinean en el ecuador
  \item Las cromátidas se separan
  \item Se forman dos núcleos
  \item Se divide el citoplasma
  \item Se duplica el ADN
\end{multicols}
\end{enumerate}
\vspace{0.15cm}
\item ¿Qué ocurre en la metafase?
\begin{enumerate}[label=\textbf{\Alph*)}]
\begin{multicols}{2}
  \item Los cromosomas se alinean en el plano ecuatorial
  \item La cromatina se condensa
  \item Las cromátidas se separan
  \item Se forman dos núcleos
  \item Se divide el citoplasma
  \item Se duplica el ADN
\end{multicols}
\end{enumerate}
\vspace{0.15cm}
\item ¿Qué ocurre en la anafase?
\begin{enumerate}[label=\textbf{\Alph*)}]
\begin{multicols}{2}
  \item Las cromátidas hermanas se separan y migran a los polos
  \item Los cromosomas se alinean
  \item La cromatina se condensa
  \item Se forman dos núcleos
  \item Se divide el citoplasma
  \item Se duplica el ADN
\end{multicols}
\end{enumerate}
\vspace{0.15cm}
\item ¿Qué ocurre en la telofase?
\begin{enumerate}[label=\textbf{\Alph*)}]
\begin{multicols}{2}
  \item Se forman dos núcleos hijos y se descondensa la cromatina
  \item Los cromosomas se alinean
  \item Las cromátidas se separan
  \item La cromatina se condensa
  \item Se divide el citoplasma
  \item Se duplica el ADN
\end{multicols}
\end{enumerate}
\vspace{0.15cm}
\item ¿Qué es la citocinesis?
\begin{enumerate}[label=\textbf{\Alph*)}]
\begin{multicols}{2}
  \item La división del citoplasma
  \item La división del núcleo
  \item La división del ADN
  \item La división de los cromosomas
  \item La división de los ribosomas
  \item La división de las mitocondrias
\end{multicols}
\end{enumerate}
\vspace{0.15cm}
\item ¿Cuántas células resultan de la mitosis?
\begin{enumerate}[label=\textbf{\Alph*)}]
\begin{multicols}{2}
  \item Dos células diploides idénticas
  \item Cuatro células haploides
  \item Dos células haploides
  \item Cuatro células diploides
  \item Una célula diploide
  \item Tres células haploides
\end{multicols}
\end{enumerate}
\vspace{0.15cm}
\item ¿Cuántas células resultan de la meiosis?
\begin{enumerate}[label=\textbf{\Alph*)}]
\begin{multicols}{2}
  \item Cuatro células haploides diferentes
  \item Dos células diploides idénticas
  \item Dos células haploides
  \item Cuatro células diploides
  \item Una célula haploide
  \item Tres células diploides
\end{multicols}
\end{enumerate}
\vspace{0.15cm}
\item ¿Qué es la meiosis I?
\begin{enumerate}[label=\textbf{\Alph*)}]
\begin{multicols}{2}
  \item La primera división meiótica que separa cromosomas homólogos
  \item La segunda división meiótica que separa cromátidas
  \item La división mitótica
  \item La división del citoplasma
  \item La duplicación del ADN
  \item La condensación de la cromatina
\end{multicols}
\end{enumerate}
\vspace{0.15cm}
\item ¿Qué es la meiosis II?
\begin{enumerate}[label=\textbf{\Alph*)}]
\begin{multicols}{2}
  \item La segunda división meiótica que separa cromátidas hermanas
  \item La primera división meiótica
  \item La división mitótica
  \item La división del citoplasma
  \item La duplicación del ADN
  \item La condensación de la cromatina
\end{multicols}
\end{enumerate}
\vspace{0.15cm}
\item ¿Qué es el entrecruzamiento o crossing over?
\begin{enumerate}[label=\textbf{\Alph*)}]
\begin{multicols}{2}
  \item El intercambio de segmentos entre cromosomas homólogos
  \item La separación de cromátidas hermanas
  \item La duplicación del ADN
  \item La condensación de la cromatina
  \item La formación del huso mitótico
  \item La desaparición del nucléolo
\end{multicols}
\end{enumerate}
\vspace{0.15cm}
\item ¿En qué fase ocurre el entrecruzamiento?
\begin{enumerate}[label=\textbf{\Alph*)}]
\begin{multicols}{2}
  \item En la profase I de la meiosis
  \item En la metafase I de la meiosis
  \item En la anafase I de la meiosis
  \item En la telofase I de la meiosis
  \item En la profase de la mitosis
  \item En la metafase de la mitosis
\end{multicols}
\end{enumerate}
\vspace{0.15cm}
\item ¿Cuál es la importancia del entrecruzamiento?
\begin{enumerate}[label=\textbf{\Alph*)}]
\begin{multicols}{2}
  \item Genera variabilidad genética
  \item Genera células idénticas
  \item Genera células diploides
  \item Genera células sin variación
  \item Genera células sin ADN
  \item Genera células sin núcleo
\end{multicols}
\end{enumerate}
\vspace{0.15cm}
\item ¿Qué es un cromosoma homólogo?
\begin{enumerate}[label=\textbf{\Alph*)}]
\begin{multicols}{2}
  \item Un par de cromosomas con la misma información genética
  \item Un par de cromosomas con información diferente
  \item Un cromosoma sin par
  \item Un cromosoma sexual
  \item Un cromosoma duplicado
  \item Un cromosoma condensado
\end{multicols}
\end{enumerate}
\vspace{0.15cm}
\item ¿Qué es la cromátida?
\begin{enumerate}[label=\textbf{\Alph*)}]
\begin{multicols}{2}
  \item Cada una de las dos copias de un cromosoma duplicado
  \item Un cromosoma sin duplicar
  \item Un cromosoma condensado
  \item Un cromosoma descondensado
  \item Un cromosoma sexual
  \item Un cromosoma homólogo
\end{multicols}
\end{enumerate}
\vspace{0.15cm}
\item ¿Qué es el centrómero?
\begin{enumerate}[label=\textbf{\Alph*)}]
\begin{multicols}{2}
  \item La región que une las cromátidas hermanas
  \item La región que separa las cromátidas
  \item La región que duplica el ADN
  \item La región que condensa la cromatina
  \item La región que forma el huso
  \item La región que desaparece el nucléolo
\end{multicols}
\end{enumerate}
\vspace{0.15cm}
\item ¿Qué es el huso mitótico?
\begin{enumerate}[label=\textbf{\Alph*)}]
\begin{multicols}{2}
  \item La estructura que separa los cromosomas durante la división
  \item La estructura que une los cromosomas
  \item La estructura que duplica los cromosomas
  \item La estructura que condensa los cromosomas
  \item La estructura que descondensa los cromosomas
  \item La estructura que forma el núcleo
\end{multicols}
\end{enumerate}
\vspace{0.15cm}
\item ¿Qué es la citocinesis en células animales?
\begin{enumerate}[label=\textbf{\Alph*)}]
\begin{multicols}{2}
  \item La división del citoplasma mediante un anillo contráctil
  \item La división del citoplasma mediante una placa celular
  \item La división del núcleo
  \item La división del ADN
  \item La división de los cromosomas
  \item La división de los ribosomas
\end{multicols}
\end{enumerate}
\vspace{0.15cm}
\item ¿Qué es la citocinesis en células vegetales?
\begin{enumerate}[label=\textbf{\Alph*)}]
\begin{multicols}{2}
  \item La división del citoplasma mediante una placa celular
  \item La división del citoplasma mediante un anillo contráctil
  \item La división del núcleo
  \item La división del ADN
  \item La división de los cromosomas
  \item La división de los ribosomas
\end{multicols}
\end{enumerate}
\vspace{0.15cm}
\item ¿Qué es la apoptosis?
\begin{enumerate}[label=\textbf{\Alph*)}]
\begin{multicols}{2}
  \item La muerte celular programada
  \item La división celular
  \item El crecimiento celular
  \item La duplicación celular
  \item La diferenciación celular
  \item La especialización celular
\end{multicols}
\end{enumerate}
\vspace{0.15cm}
\item ¿Qué es el cáncer?
\begin{enumerate}[label=\textbf{\Alph*)}]
\begin{multicols}{2}
  \item El crecimiento descontrolado de células
  \item El crecimiento controlado de células
  \item La muerte controlada de células
  \item La muerte descontrolada de células
  \item La división controlada de células
  \item La división descontrolada de células
\end{multicols}
\end{enumerate}
\vspace{0.15cm}
\item ¿Qué es un tumor benigno?
\begin{enumerate}[label=\textbf{\Alph*)}]
\begin{multicols}{2}
  \item Un crecimiento celular que no se disemina
  \item Un crecimiento celular que se disemina
  \item Un crecimiento celular que invade tejidos
  \item Un crecimiento celular que metastatiza
  \item Un crecimiento celular que destruye órganos
  \item Un crecimiento celular que causa muerte
\end{multicols}
\end{enumerate}
\vspace{0.15cm}
\item ¿Qué es un tumor maligno?
\begin{enumerate}[label=\textbf{\Alph*)}]
\begin{multicols}{2}
  \item Un crecimiento celular que se disemina e invade tejidos
  \item Un crecimiento celular que no se disemina
  \item Un crecimiento celular que no invade tejidos
  \item Un crecimiento celular que no metastatiza
  \item Un crecimiento celular que no destruye órganos
  \item Un crecimiento celular que no causa muerte
\end{multicols}
\end{enumerate}
\vspace{0.15cm}
\item ¿Qué es la metástasis?
\begin{enumerate}[label=\textbf{\Alph*)}]
\begin{multicols}{2}
  \item La diseminación de células cancerosas a otros órganos
  \item La diseminación de células normales
  \item La diseminación de células benignas
  \item La diseminación de células sanas
  \item La diseminación de células musculares
  \item La diseminación de células nerviosas
\end{multicols}
\end{enumerate}
\vspace{0.15cm}
\item ¿Qué es un oncogén?
\begin{enumerate}[label=\textbf{\Alph*)}]
\begin{multicols}{2}
  \item Un gen que puede causar cáncer
  \item Un gen que previene el cáncer
  \item Un gen que repara el ADN
  \item Un gen que controla el ciclo celular
  \item Un gen que induce la apoptosis
  \item Un gen que inhibe el crecimiento
\end{multicols}
\end{enumerate}
\vspace{0.15cm}
\item ¿Qué es un gen supresor de tumores?
\begin{enumerate}[label=\textbf{\Alph*)}]
\begin{multicols}{2}
  \item Un gen que previene el cáncer
  \item Un gen que causa cáncer
  \item Un gen que repara el ADN
  \item Un gen que controla el ciclo celular
  \item Un gen que induce la apoptosis
  \item Un gen que inhibe el crecimiento
\end{multicols}
\end{enumerate}
\vspace{0.15cm}
\item ¿Qué es la mutación?
\begin{enumerate}[label=\textbf{\Alph*)}]
\begin{multicols}{2}
  \item Un cambio en la secuencia del ADN
  \item Un cambio en la secuencia del ARN
  \item Un cambio en la secuencia de proteínas
  \item Un cambio en la secuencia de lípidos
  \item Un cambio en la secuencia de carbohidratos
  \item Un cambio en la secuencia de minerales
\end{multicols}
\end{enumerate}
\vspace{0.15cm}
\item ¿Qué es una mutación puntual?
\begin{enumerate}[label=\textbf{\Alph*)}]
\begin{multicols}{2}
  \item Un cambio en un solo nucleótido
  \item Un cambio en varios nucleótidos
  \item Un cambio en un cromosoma completo
  \item Un cambio en varios cromosomas
  \item Un cambio en un gen completo
  \item Un cambio en varios genes
\end{multicols}
\end{enumerate}
\vspace{0.15cm}
\item ¿Qué es una mutación cromosómica?
\begin{enumerate}[label=\textbf{\Alph*)}]
\begin{multicols}{2}
  \item Un cambio en la estructura o número de cromosomas
  \item Un cambio en un solo nucleótido
  \item Un cambio en varios nucleótidos
  \item Un cambio en un gen completo
  \item Un cambio en varios genes
  \item Un cambio en una proteína
\end{multicols}
\end{enumerate}
\vspace{0.15cm}
\item ¿Qué es la aneuploidía?
\begin{enumerate}[label=\textbf{\Alph*)}]
\begin{multicols}{2}
  \item La presencia de un número anormal de cromosomas
  \item La presencia de un número normal de cromosomas
  \item La presencia de un número par de cromosomas
  \item La presencia de un número impar de cromosomas
  \item La presencia de un número múltiplo de cromosomas
  \item La presencia de un número fraccionario de cromosomas
\end{multicols}
\end{enumerate}
\vspace{0.15cm}
\item ¿Qué es el síndrome de Down?
\begin{enumerate}[label=\textbf{\Alph*)}]
\begin{multicols}{2}
  \item Una trisomía del cromosoma 21
  \item Una monosomía del cromosoma 21
  \item Una trisomía del cromosoma 18
  \item Una monosomía del cromosoma 18
  \item Una trisomía del cromosoma 13
  \item Una monosomía del cromosoma 13
\end{multicols}
\end{enumerate}
\vspace{0.15cm}
\item ¿Qué es el síndrome de Turner?
\begin{enumerate}[label=\textbf{\Alph*)}]
\begin{multicols}{2}
  \item Una monosomía del cromosoma X
  \item Una trisomía del cromosoma X
  \item Una monosomía del cromosoma Y
  \item Una trisomía del cromosoma Y
  \item Una monosomía del cromosoma 21
  \item Una trisomía del cromosoma 21
\end{multicols}
\end{enumerate}
\vspace{0.15cm}
\item ¿Qué es el síndrome de Klinefelter?
\begin{enumerate}[label=\textbf{\Alph*)}]
\begin{multicols}{2}
  \item Una trisomía de los cromosomas sexuales (XXY)
  \item Una monosomía de los cromosomas sexuales
  \item Una trisomía del cromosoma 21
  \item Una monosomía del cromosoma 21
  \item Una trisomía del cromosoma 18
  \item Una monosomía del cromosoma 18
\end{multicols}
\end{enumerate}
\vspace{0.15cm}
\item ¿Qué es un cariotipo?
\begin{enumerate}[label=\textbf{\Alph*)}]
\begin{multicols}{2}
  \item La representación ordenada de los cromosomas de una célula
  \item La representación ordenada de los genes de una célula
  \item La representación ordenada de las proteínas de una célula
  \item La representación ordenada de los lípidos de una célula
  \item La representación ordenada de los carbohidratos de una célula
  \item La representación ordenada de los minerales de una célula
\end{multicols}
\end{enumerate}
\vspace{0.15cm}
\item ¿Cuántos cromosomas tiene la especie humana?
\begin{enumerate}[label=\textbf{\Alph*)}]
\begin{multicols}{2}
  \item 46 (23 pares)
  \item 44 (22 pares)
  \item 48 (24 pares)
  \item 42 (21 pares)
  \item 50 (25 pares)
  \item 40 (20 pares)
\end{multicols}
\end{enumerate}
\vspace{0.15cm}
\item ¿Cuántos cromosomas sexuales tiene la mujer?
\begin{enumerate}[label=\textbf{\Alph*)}]
\begin{multicols}{2}
  \item XX
  \item XY
  \item XO
  \item XXY
  \item XYY
  \item YY
\end{multicols}
\end{enumerate}
\vspace{0.15cm}
\item ¿Cuántos cromosomas sexuales tiene el hombre?
\begin{enumerate}[label=\textbf{\Alph*)}]
\begin{multicols}{2}
  \item XY
  \item XX
  \item XO
  \item XXY
  \item XYY
  \item YY
\end{multicols}
\end{enumerate}
\vspace{0.15cm}
\item ¿Qué es la autosoma?
\begin{enumerate}[label=\textbf{\Alph*)}]
\begin{multicols}{2}
  \item Un cromosoma no sexual
  \item Un cromosoma sexual
  \item Un cromosoma X
  \item Un cromosoma Y
  \item Un cromosoma duplicado
  \item Un cromosoma homólogo
\end{multicols}
\end{enumerate}
\vspace{0.15cm}
\item ¿Qué es el genotipo?
\begin{enumerate}[label=\textbf{\Alph*)}]
\begin{multicols}{2}
  \item La composición genética de un organismo
  \item La apariencia externa de un organismo
  \item El comportamiento de un organismo
  \item El ambiente de un organismo
  \item La alimentación de un organismo
  \item La reproducción de un organismo
\end{multicols}
\end{enumerate}
\vspace{0.15cm}
\item ¿Qué es el fenotipo?
\begin{enumerate}[label=\textbf{\Alph*)}]
\begin{multicols}{2}
  \item La manifestación observable del genotipo
  \item La composición genética de un organismo
  \item El comportamiento de un organismo
  \item El ambiente de un organismo
  \item La alimentación de un organismo
  \item La reproducción de un organismo
\end{multicols}
\end{enumerate}
\vspace{0.15cm}
\item ¿Qué es un alelo?
\begin{enumerate}[label=\textbf{\Alph*)}]
\begin{multicols}{2}
  \item Cada una de las variantes de un gen
  \item Cada una de las variantes de un cromosoma
  \item Cada una de las variantes de una proteína
  \item Cada una de las variantes de un lípido
  \item Cada una de las variantes de un carbohidrato
  \item Cada una de las variantes de un mineral
\end{multicols}
\end{enumerate}
\vspace{0.15cm}
\item ¿Qué es un organismo homocigoto?
\begin{enumerate}[label=\textbf{\Alph*)}]
\begin{multicols}{2}
  \item El que tiene dos alelos iguales para un gen
  \item El que tiene dos alelos diferentes para un gen
  \item El que tiene un solo alelo para un gen
  \item El que tiene tres alelos para un gen
  \item El que tiene cuatro alelos para un gen
  \item El que no tiene alelos para un gen
\end{multicols}
\end{enumerate}
\vspace{0.15cm}
\item ¿Qué es un organismo heterocigoto?
\begin{enumerate}[label=\textbf{\Alph*)}]
\begin{multicols}{2}
  \item El que tiene dos alelos diferentes para un gen
  \item El que tiene dos alelos iguales para un gen
  \item El que tiene un solo alelo para un gen
  \item El que tiene tres alelos para un gen
  \item El que tiene cuatro alelos para un gen
  \item El que no tiene alelos para un gen
\end{multicols}
\end{enumerate}
\vspace{0.15cm}
\item ¿Qué es un alelo dominante?
\begin{enumerate}[label=\textbf{\Alph*)}]
\begin{multicols}{2}
  \item El que se manifiesta en el fenotipo aunque haya otro alelo diferente
  \item El que no se manifiesta en el fenotipo
  \item El que se manifiesta solo si hay dos copias
  \item El que se manifiesta solo si hay tres copias
  \item El que se manifiesta solo si hay cuatro copias
  \item El que no se manifiesta nunca
\end{multicols}
\end{enumerate}
\vspace{0.15cm}
\item ¿Qué es un alelo recesivo?
\begin{enumerate}[label=\textbf{\Alph*)}]
\begin{multicols}{2}
  \item El que se manifiesta solo si hay dos copias iguales
  \item El que se manifiesta siempre
  \item El que se manifiesta aunque haya otro alelo diferente
  \item El que se manifiesta solo si hay tres copias
  \item El que se manifiesta solo si hay cuatro copias
  \item El que se manifiesta solo si hay una copia
\end{multicols}
\end{enumerate}
\vspace{0.15cm}
\item ¿Qué es la mitosis?
\begin{enumerate}[label=\textbf{\Alph*)}]
\begin{multicols}{2}
  \item La división celular que produce dos células idénticas
  \item La división celular que produce cuatro células diferentes
  \item La división celular que produce dos células diferentes
  \item La división celular que produce cuatro células idénticas
  \item La división celular que produce una célula
  \item La división celular que produce tres células
\end{multicols}
\end{enumerate}
\vspace{0.15cm}
\item ¿Qué es la meiosis?
\begin{enumerate}[label=\textbf{\Alph*)}]
\begin{multicols}{2}
  \item La división celular que produce cuatro células diferentes
  \item La división celular que produce dos células idénticas
  \item La división celular que produce dos células diferentes
  \item La división celular que produce cuatro células idénticas
  \item La división celular que produce una célula
  \item La división celular que produce tres células
\end{multicols}
\end{enumerate}
\vspace{0.15cm}
\item ¿Cuál es la función de la mitosis?
\begin{enumerate}[label=\textbf{\Alph*)}]
\begin{multicols}{2}
  \item El crecimiento y reparación de tejidos
  \item La producción de gametos
  \item La producción de células diferentes
  \item La producción de células sexuales
  \item La producción de células reproductoras
  \item La producción de células germinales
\end{multicols}
\end{enumerate}
\vspace{0.15cm}
\item ¿Cuál es la función de la meiosis?
\begin{enumerate}[label=\textbf{\Alph*)}]
\begin{multicols}{2}
  \item La producción de gametos
  \item El crecimiento y reparación de tejidos
  \item La producción de células idénticas
  \item La producción de células somáticas
  \item La producción de células corporales
  \item La producción de células diploides
\end{multicols}
\end{enumerate}
\vspace{0.15cm}
\item ¿Qué es la célula procariota?
\begin{enumerate}[label=\textbf{\Alph*)}]
\begin{multicols}{2}
  \item La que no tiene núcleo definido
  \item La que tiene núcleo definido
  \item La que tiene dos núcleos
  \item La que tiene tres núcleos
  \item La que tiene cuatro núcleos
  \item La que no tiene membrana
\end{multicols}
\end{enumerate}
\vspace{0.15cm}
\item ¿Qué es la célula eucariota?
\begin{enumerate}[label=\textbf{\Alph*)}]
\begin{multicols}{2}
  \item La que tiene núcleo definido
  \item La que no tiene núcleo definido
  \item La que tiene dos núcleos
  \item La que tiene tres núcleos
  \item La que tiene cuatro núcleos
  \item La que no tiene membrana
\end{multicols}
\end{enumerate}
\vspace{0.15cm}
\item ¿Cuál es un ejemplo de célula procariota?
\begin{enumerate}[label=\textbf{\Alph*)}]
\begin{multicols}{2}
  \item La bacteria
  \item La célula animal
  \item La célula vegetal
  \item La célula fúngica
  \item La célula protista
  \item La célula nerviosa
\end{multicols}
\end{enumerate}
\vspace{0.15cm}
\item ¿Cuál es un ejemplo de célula eucariota?
\begin{enumerate}[label=\textbf{\Alph*)}]
\begin{multicols}{2}
  \item La célula animal
  \item La bacteria
  \item La cianobacteria
  \item La arquea
  \item La bacteria intestinal
  \item La bacteria del suelo
\end{multicols}
\end{enumerate}
\vspace{0.15cm}
\item ¿Qué es la membrana celular?
\begin{enumerate}[label=\textbf{\Alph*)}]
\begin{multicols}{2}
  \item La estructura que delimita la célula y regula el paso de sustancias
  \item La estructura que no delimita la célula
  \item La estructura que no regula el paso de sustancias
  \item La estructura que delimita el núcleo
  \item La estructura que delimita la mitocondria
  \item La estructura que delimita el ribosoma
\end{multicols}
\end{enumerate}
\vspace{0.15cm}
\item ¿Qué es el citoplasma?
\begin{enumerate}[label=\textbf{\Alph*)}]
\begin{multicols}{2}
  \item El medio interno de la célula donde ocurren las reacciones
  \item El medio externo de la célula
  \item El medio interno del núcleo
  \item El medio externo del núcleo
  \item El medio interno de la mitocondria
  \item El medio externo de la mitocondria
\end{multicols}
\end{enumerate}
\vspace{0.15cm}
\item ¿Qué es el núcleo celular?
\begin{enumerate}[label=\textbf{\Alph*)}]
\begin{multicols}{2}
  \item La estructura que contiene el material genético
  \item La estructura que contiene el material proteico
  \item La estructura que contiene el material lipídico
  \item La estructura que contiene el material carbohidrático
  \item La estructura que contiene el material mineral
  \item La estructura que contiene el material gaseoso
\end{multicols}
\end{enumerate}
\vspace{0.15cm}
\item ¿Qué es la mitocondria?
\begin{enumerate}[label=\textbf{\Alph*)}]
\begin{multicols}{2}
  \item El orgánulo que produce energía
  \item El orgánulo que produce proteínas
  \item El orgánulo que produce lípidos
  \item El orgánulo que produce carbohidratos
  \item El orgánulo que produce minerales
  \item El orgánulo que produce gases
\end{multicols}
\end{enumerate}
\vspace{0.15cm}
\item ¿Qué es el ribosoma?
\begin{enumerate}[label=\textbf{\Alph*)}]
\begin{multicols}{2}
  \item El orgánulo que sintetiza proteínas
  \item El orgánulo que sintetiza lípidos
  \item El orgánulo que sintetiza carbohidratos
  \item El orgánulo que sintetiza minerales
  \item El orgánulo que sintetiza gases
  \item El orgánulo que sintetiza agua
\end{multicols}
\end{enumerate}
\vspace{0.15cm}
\item ¿Qué es el cloroplasto?
\begin{enumerate}[label=\textbf{\Alph*)}]
\begin{multicols}{2}
  \item El orgánulo que realiza la fotosíntesis
  \item El orgánulo que realiza la respiración
  \item El orgánulo que realiza la digestión
  \item El orgánulo que realiza la excreción
  \item El orgánulo que realiza la reproducción
  \item El orgánulo que realiza la síntesis
\end{multicols}
\end{enumerate}
\vspace{0.15cm}
\item ¿Qué es la pared celular?
\begin{enumerate}[label=\textbf{\Alph*)}]
\begin{multicols}{2}
  \item La estructura que rodea la membrana de células vegetales
  \item La estructura que rodea la membrana de células animales
  \item La estructura que rodea el núcleo
  \item La estructura que rodea la mitocondria
  \item La estructura que rodea el ribosoma
  \item La estructura que rodea el cloroplasto
\end{multicols}
\end{enumerate}
\vspace{0.15cm}
\item ¿Qué es el aparato de Golgi?
\begin{enumerate}[label=\textbf{\Alph*)}]
\begin{multicols}{2}
  \item El orgánulo que modifica y empaqueta proteínas
  \item El orgánulo que modifica y empaqueta lípidos
  \item El orgánulo que modifica y empaqueta carbohidratos
  \item El orgánulo que modifica y empaqueta minerales
  \item El orgánulo que modifica y empaqueta gases
  \item El orgánulo que modifica y empaqueta agua
\end{multicols}
\end{enumerate}
\vspace{0.15cm}
\item ¿Qué es el retículo endoplasmático?
\begin{enumerate}[label=\textbf{\Alph*)}]
\begin{multicols}{2}
  \item La red de membranas que transporta sustancias
  \item La red de membranas que no transporta sustancias
  \item La red de membranas que delimita el núcleo
  \item La red de membranas que delimita la mitocondria
  \item La red de membranas que delimita el ribosoma
  \item La red de membranas que delimita el cloroplasto ---
\end{multicols}
\end{enumerate}
\vspace{0.15cm}
\end{enumerate}\n
\subsubsection*{SECCIÓN III: GENÉTICA MENDELIANA (121–180)}

\begin{enumerate}[label=\textbf{\arabic*.}]
\setcounter{enumi}{120}
\item ¿Quién es considerado el padre de la genética?
\begin{enumerate}[label=\textbf{\Alph*)}]
\begin{multicols}{2}
  \item Gregor Mendel
  \item Charles Darwin
  \item Louis Pasteur
  \item Robert Hooke
  \item Alexander Fleming
  \item Isaac Newton
\end{multicols}
\end{enumerate}
\vspace{0.15cm}
\item ¿Qué estudió Gregor Mendel?
\begin{enumerate}[label=\textbf{\Alph*)}]
\begin{multicols}{2}
  \item La herencia de los caracteres en las plantas de guisante
  \item La herencia de los caracteres en los animales
  \item La herencia de los caracteres en los hongos
  \item La herencia de los caracteres en las bacterias
  \item La herencia de los caracteres en los virus
  \item La herencia de los caracteres en los minerales
\end{multicols}
\end{enumerate}
\vspace{0.15cm}
\item ¿Qué es un carácter hereditario?
\begin{enumerate}[label=\textbf{\Alph*)}]
\begin{multicols}{2}
  \item Una característica que se transmite de padres a hijos
  \item Una característica que no se transmite
  \item Una característica adquirida
  \item Una característica ambiental
  \item Una característica cultural
  \item Una característica social
\end{multicols}
\end{enumerate}
\vspace{0.15cm}
\item ¿Qué es la primera ley de Mendel?
\begin{enumerate}[label=\textbf{\Alph*)}]
\begin{multicols}{2}
  \item La ley de la uniformidad de los híbridos de la primera generación
  \item La ley de la segregación de los caracteres
  \item La ley de la distribución independiente
  \item La ley de la herencia ligada al sexo
  \item La ley de la herencia citoplasmática
  \item La ley de la herencia poligénica
\end{multicols}
\end{enumerate}
\vspace{0.15cm}
\item ¿Qué es la segunda ley de Mendel?
\begin{enumerate}[label=\textbf{\Alph*)}]
\begin{multicols}{2}
  \item La ley de la segregación de los caracteres
  \item La ley de la uniformidad
  \item La ley de la distribución independiente
  \item La ley de la herencia ligada al sexo
  \item La ley de la herencia citoplasmática
  \item La ley de la herencia poligénica
\end{multicols}
\end{enumerate}
\vspace{0.15cm}
\item ¿Qué es la tercera ley de Mendel?
\begin{enumerate}[label=\textbf{\Alph*)}]
\begin{multicols}{2}
  \item La ley de la distribución independiente
  \item La ley de la uniformidad
  \item La ley de la segregación
  \item La ley de la herencia ligada al sexo
  \item La ley de la herencia citoplasmática
  \item La ley de la herencia poligénica
\end{multicols}
\end{enumerate}
\vspace{0.15cm}
\item ¿Qué es un cruce monohíbrido?
\begin{enumerate}[label=\textbf{\Alph*)}]
\begin{multicols}{2}
  \item El cruce que estudia un solo carácter
  \item El cruce que estudia dos caracteres
  \item El cruce que estudia tres caracteres
  \item El cruce que estudia cuatro caracteres
  \item El cruce que estudia cinco caracteres
  \item El cruce que estudia seis caracteres
\end{multicols}
\end{enumerate}
\vspace{0.15cm}
\item ¿Qué es un cruce dihíbrido?
\begin{enumerate}[label=\textbf{\Alph*)}]
\begin{multicols}{2}
  \item El cruce que estudia dos caracteres
  \item El cruce que estudia un carácter
  \item El cruce que estudia tres caracteres
  \item El cruce que estudia cuatro caracteres
  \item El cruce que estudia cinco caracteres
  \item El cruce que estudia seis caracteres
\end{multicols}
\end{enumerate}
\vspace{0.15cm}
\item ¿Qué es el cuadro de Punnett?
\begin{enumerate}[label=\textbf{\Alph*)}]
\begin{multicols}{2}
  \item Una herramienta para predecir las combinaciones genéticas
  \item Una herramienta para predecir las combinaciones ambientales
  \item Una herramienta para predecir las combinaciones culturales
  \item Una herramienta para predecir las combinaciones sociales
  \item Una herramienta para predecir las combinaciones económicas
  \item Una herramienta para predecir las combinaciones políticas
\end{multicols}
\end{enumerate}
\vspace{0.15cm}
\item ¿Qué proporción fenotípica se obtiene en la F2 de un cruce monohíbrido?
\begin{enumerate}[label=\textbf{\Alph*)}]
\begin{multicols}{2}
  \item 3:1
  \item 1:1
  \item 1:2:1
  \item 9:3:3:1
  \item 2:1
  \item 1:3
\end{multicols}
\end{enumerate}
\vspace{0.15cm}
\item ¿Qué proporción genotípica se obtiene en la F2 de un cruce monohíbrido?
\begin{enumerate}[label=\textbf{\Alph*)}]
\begin{multicols}{2}
  \item 1:2:1
  \item 3:1
  \item 1:1
  \item 9:3:3:1
  \item 2:1
  \item 1:3
\end{multicols}
\end{enumerate}
\vspace{0.15cm}
\item ¿Qué proporción fenotípica se obtiene en la F2 de un cruce dihíbrido?
\begin{enumerate}[label=\textbf{\Alph*)}]
\begin{multicols}{2}
  \item 9:3:3:1
  \item 3:1
  \item 1:1
  \item 1:2:1
  \item 2:1
  \item 1:3
\end{multicols}
\end{enumerate}
\vspace{0.15cm}
\item ¿Qué es un gen?
\begin{enumerate}[label=\textbf{\Alph*)}]
\begin{multicols}{2}
  \item Un segmento de ADN que contiene información para un rasgo
  \item Un segmento de ARN que contiene información para un rasgo
  \item Un segmento de proteína que contiene información para un rasgo
  \item Un segmento de lípido que contiene información para un rasgo
  \item Un segmento de carbohidrato que contiene información para un rasgo
  \item Un segmento de mineral que contiene información para un rasgo
\end{multicols}
\end{enumerate}
\vspace{0.15cm}
\item ¿Qué es el genotipo?
\begin{enumerate}[label=\textbf{\Alph*)}]
\begin{multicols}{2}
  \item La composición genética de un organismo
  \item La apariencia externa de un organismo
  \item El comportamiento de un organismo
  \item El ambiente de un organismo
  \item La alimentación de un organismo
  \item La reproducción de un organismo
\end{multicols}
\end{enumerate}
\vspace{0.15cm}
\item ¿Qué es el fenotipo?
\begin{enumerate}[label=\textbf{\Alph*)}]
\begin{multicols}{2}
  \item La manifestación observable del genotipo
  \item La composición genética de un organismo
  \item El comportamiento de un organismo
  \item El ambiente de un organismo
  \item La alimentación de un organismo
  \item La reproducción de un organismo
\end{multicols}
\end{enumerate}
\vspace{0.15cm}
\item ¿Qué es un alelo?
\begin{enumerate}[label=\textbf{\Alph*)}]
\begin{multicols}{2}
  \item Cada una de las variantes de un gen
  \item Cada una de las variantes de un cromosoma
  \item Cada una de las variantes de una proteína
  \item Cada una de las variantes de un lípido
  \item Cada una de las variantes de un carbohidrato
  \item Cada una de las variantes de un mineral
\end{multicols}
\end{enumerate}
\vspace{0.15cm}
\item ¿Qué es un organismo homocigoto?
\begin{enumerate}[label=\textbf{\Alph*)}]
\begin{multicols}{2}
  \item El que tiene dos alelos iguales para un gen
  \item El que tiene dos alelos diferentes para un gen
  \item El que tiene un solo alelo para un gen
  \item El que tiene tres alelos para un gen
  \item El que tiene cuatro alelos para un gen
  \item El que no tiene alelos para un gen
\end{multicols}
\end{enumerate}
\vspace{0.15cm}
\item ¿Qué es un organismo heterocigoto?
\begin{enumerate}[label=\textbf{\Alph*)}]
\begin{multicols}{2}
  \item El que tiene dos alelos diferentes para un gen
  \item El que tiene dos alelos iguales para un gen
  \item El que tiene un solo alelo para un gen
  \item El que tiene tres alelos para un gen
  \item El que tiene cuatro alelos para un gen
  \item El que no tiene alelos para un gen
\end{multicols}
\end{enumerate}
\vspace{0.15cm}
\item ¿Qué es un alelo dominante?
\begin{enumerate}[label=\textbf{\Alph*)}]
\begin{multicols}{2}
  \item El que se manifiesta en el fenotipo aunque haya otro alelo diferente
  \item El que no se manifiesta en el fenotipo
  \item El que se manifiesta solo si hay dos copias
  \item El que se manifiesta solo si hay tres copias
  \item El que se manifiesta solo si hay cuatro copias
  \item El que no se manifiesta nunca
\end{multicols}
\end{enumerate}
\vspace{0.15cm}
\item ¿Qué es un alelo recesivo?
\begin{enumerate}[label=\textbf{\Alph*)}]
\begin{multicols}{2}
  \item El que se manifiesta solo si hay dos copias iguales
  \item El que se manifiesta siempre
  \item El que se manifiesta aunque haya otro alelo diferente
  \item El que se manifiesta solo si hay tres copias
  \item El que se manifiesta solo si hay cuatro copias
  \item El que se manifiesta solo si hay una copia
\end{multicols}
\end{enumerate}
\vspace{0.15cm}
\item ¿Qué es un cruce monohíbrido?
\begin{enumerate}[label=\textbf{\Alph*)}]
\begin{multicols}{2}
  \item El cruce que estudia un solo carácter
  \item El cruce que estudia dos caracteres
  \item El cruce que estudia tres caracteres
  \item El cruce que estudia cuatro caracteres
  \item El cruce que estudia cinco caracteres
  \item El cruce que estudia seis caracteres
\end{multicols}
\end{enumerate}
\vspace{0.15cm}
\item ¿Qué es un cruce dihíbrido?
\begin{enumerate}[label=\textbf{\Alph*)}]
\begin{multicols}{2}
  \item El cruce que estudia dos caracteres
  \item El cruce que estudia un carácter
  \item El cruce que estudia tres caracteres
  \item El cruce que estudia cuatro caracteres
  \item El cruce que estudia cinco caracteres
  \item El cruce que estudia seis caracteres
\end{multicols}
\end{enumerate}
\vspace{0.15cm}
\item ¿Qué es el cuadro de Punnett?
\begin{enumerate}[label=\textbf{\Alph*)}]
\begin{multicols}{2}
  \item Una herramienta para predecir las combinaciones genéticas
  \item Una herramienta para predecir las combinaciones ambientales
  \item Una herramienta para predecir las combinaciones culturales
  \item Una herramienta para predecir las combinaciones sociales
  \item Una herramienta para predecir las combinaciones económicas
  \item Una herramienta para predecir las combinaciones políticas
\end{multicols}
\end{enumerate}
\vspace{0.15cm}
\item ¿Qué proporción fenotípica se obtiene en la F2 de un cruce monohíbrido?
\begin{enumerate}[label=\textbf{\Alph*)}]
\begin{multicols}{2}
  \item 3:1
  \item 1:1
  \item 1:2:1
  \item 9:3:3:1
  \item 2:1
  \item 1:3
\end{multicols}
\end{enumerate}
\vspace{0.15cm}
\item ¿Qué proporción genotípica se obtiene en la F2 de un cruce monohíbrido?
\begin{enumerate}[label=\textbf{\Alph*)}]
\begin{multicols}{2}
  \item 1:2:1
  \item 3:1
  \item 1:1
  \item 9:3:3:1
  \item 2:1
  \item 1:3
\end{multicols}
\end{enumerate}
\vspace{0.15cm}
\item ¿Qué proporción fenotípica se obtiene en la F2 de un cruce dihíbrido?
\begin{enumerate}[label=\textbf{\Alph*)}]
\begin{multicols}{2}
  \item 9:3:3:1
  \item 3:1
  \item 1:1
  \item 1:2:1
  \item 2:1
  \item 1:3
\end{multicols}
\end{enumerate}
\vspace{0.15cm}
\item ¿Qué es la herencia ligada al sexo?
\begin{enumerate}[label=\textbf{\Alph*)}]
\begin{multicols}{2}
  \item La herencia de caracteres determinados por genes en los cromosomas sexuales
  \item La herencia de caracteres determinados por genes en los autosomas
  \item La herencia de caracteres determinados por genes en las mitocondrias
  \item La herencia de caracteres determinados por genes en los cloroplastos
  \item La herencia de caracteres determinados por genes en los ribosomas
  \item La herencia de caracteres determinados por genes en el citoplasma
\end{multicols}
\end{enumerate}
\vspace{0.15cm}
\item ¿Cuál es un ejemplo de herencia ligada al sexo?
\begin{enumerate}[label=\textbf{\Alph*)}]
\begin{multicols}{2}
  \item El daltonismo
  \item La fenilcetonuria
  \item La fibrosis quística
  \item La anemia falciforme
  \item La enfermedad de Huntington
  \item La corea de Huntington
\end{multicols}
\end{enumerate}
\vspace{0.15cm}
\item ¿Qué es un árbol genealógico?
\begin{enumerate}[label=\textbf{\Alph*)}]
\begin{multicols}{2}
  \item Una representación gráfica de la herencia de una familia
  \item Una representación gráfica de la alimentación de una familia
  \item Una representación gráfica de la economía de una familia
  \item Una representación gráfica de la cultura de una familia
  \item Una representación gráfica de la sociedad de una familia
  \item Una representación gráfica de la política de una familia
\end{multicols}
\end{enumerate}
\vspace{0.15cm}
\item ¿Qué es un portador?
\begin{enumerate}[label=\textbf{\Alph*)}]
\begin{multicols}{2}
  \item Un individuo heterocigoto que no manifiesta la enfermedad pero puede transmitirla
  \item Un individuo homocigoto que manifiesta la enfermedad
  \item Un individuo homocigoto que no manifiesta la enfermedad
  \item Un individuo heterocigoto que manifiesta la enfermedad
  \item Un individuo que no tiene el gen
  \item Un individuo que tiene tres alelos
\end{multicols}
\end{enumerate}
\vspace{0.15cm}
\item ¿Qué es una enfermedad autosómica recesiva?
\begin{enumerate}[label=\textbf{\Alph*)}]
\begin{multicols}{2}
  \item La que se manifiesta solo si hay dos copias del alelo recesivo
  \item La que se manifiesta con una sola copia del alelo
  \item La que se manifiesta solo en hombres
  \item La que se manifiesta solo en mujeres
  \item La que se manifiesta solo en heterocigotos
  \item La que se manifiesta solo en homocigotos dominantes
\end{multicols}
\end{enumerate}
\vspace{0.15cm}
\item ¿Qué es una enfermedad autosómica dominante?
\begin{enumerate}[label=\textbf{\Alph*)}]
\begin{multicols}{2}
  \item La que se manifiesta con una sola copia del alelo dominante
  \item La que se manifiesta solo con dos copias
  \item La que se manifiesta solo en hombres
  \item La que se manifiesta solo en mujeres
  \item La que se manifiesta solo en heterocigotos recesivos
  \item La que se manifiesta solo en homocigotos recesivos
\end{multicols}
\end{enumerate}
\vspace{0.15cm}
\item ¿Cuál es un ejemplo de enfermedad autosómica recesiva?
\begin{enumerate}[label=\textbf{\Alph*)}]
\begin{multicols}{2}
  \item La fibrosis quística
  \item La enfermedad de Huntington
  \item La acondroplasia
  \item El síndrome de Marfan
  \item La neurofibromatosis
  \item La polidactilia
\end{multicols}
\end{enumerate}
\vspace{0.15cm}
\item ¿Cuál es un ejemplo de enfermedad autosómica dominante?
\begin{enumerate}[label=\textbf{\Alph*)}]
\begin{multicols}{2}
  \item La enfermedad de Huntington
  \item La fibrosis quística
  \item La fenilcetonuria
  \item La anemia falciforme
  \item La talasemia
  \item La galactosemia
\end{multicols}
\end{enumerate}
\vspace{0.15cm}
\item ¿Qué es la codominancia?
\begin{enumerate}[label=\textbf{\Alph*)}]
\begin{multicols}{2}
  \item La expresión simultánea de dos alelos diferentes
  \item La expresión de un solo alelo
  \item La expresión de ningún alelo
  \item La expresión de tres alelos
  \item La expresión de cuatro alelos
  \item La expresión de cinco alelos
\end{multicols}
\end{enumerate}
\vspace{0.15cm}
\item ¿Cuál es un ejemplo de codominancia?
\begin{enumerate}[label=\textbf{\Alph*)}]
\begin{multicols}{2}
  \item El grupo sanguíneo AB
  \item El grupo sanguíneo A
  \item El grupo sanguíneo B
  \item El grupo sanguíneo O
  \item El color de ojos
  \item El color de cabello
\end{multicols}
\end{enumerate}
\vspace{0.15cm}
\item ¿Qué es la herencia poligénica?
\begin{enumerate}[label=\textbf{\Alph*)}]
\begin{multicols}{2}
  \item La determinada por múltiples genes
  \item La determinada por un solo gen
  \item La determinada por dos genes
  \item La determinada por tres genes
  \item La determinada por cuatro genes
  \item La determinada por cinco genes
\end{multicols}
\end{enumerate}
\vspace{0.15cm}
\item ¿Cuál es un ejemplo de herencia poligénica?
\begin{enumerate}[label=\textbf{\Alph*)}]
\begin{multicols}{2}
  \item La estatura
  \item El grupo sanguíneo
  \item El daltonismo
  \item La fibrosis quística
  \item La anemia falciforme
  \item La enfermedad de Huntington
\end{multicols}
\end{enumerate}
\vspace{0.15cm}
\item ¿Qué es la pleiotropía?
\begin{enumerate}[label=\textbf{\Alph*)}]
\begin{multicols}{2}
  \item La capacidad de un gen de afectar múltiples caracteres
  \item La capacidad de múltiples genes de afectar un carácter
  \item La capacidad de un gen de afectar un carácter
  \item La capacidad de múltiples genes de afectar múltiples caracteres
  \item La capacidad de un gen de no afectar caracteres
  \item La capacidad de múltiples genes de no afectar caracteres
\end{multicols}
\end{enumerate}
\vspace{0.15cm}
\item ¿Qué es la penetrancia?
\begin{enumerate}[label=\textbf{\Alph*)}]
\begin{multicols}{2}
  \item La proporción de individuos con un genotipo que manifiestan el fenotipo
  \item La proporción de individuos con un fenotipo que tienen el genotipo
  \item La proporción de individuos con un genotipo que no manifiestan el fenotipo
  \item La proporción de individuos con un fenotipo que no tienen el genotipo
  \item La proporción de individuos con un genotipo que tienen otro genotipo
  \item La proporción de individuos con un fenotipo que tienen otro fenotipo
\end{multicols}
\end{enumerate}
\vspace{0.15cm}
\item ¿Qué es la expresividad variable?
\begin{enumerate}[label=\textbf{\Alph*)}]
\begin{multicols}{2}
  \item El grado en que un genotipo se manifiesta en el fenotipo
  \item La proporción de individuos que manifiestan el fenotipo
  \item La proporción de individuos que no manifiestan el fenotipo
  \item La proporción de individuos que tienen el genotipo
  \item La proporción de individuos que no tienen el genotipo
  \item La proporción de individuos que tienen otro genotipo
\end{multicols}
\end{enumerate}
\vspace{0.15cm}
\item ¿Qué es la herencia citoplasmática?
\begin{enumerate}[label=\textbf{\Alph*)}]
\begin{multicols}{2}
  \item La herencia de genes ubicados en el citoplasma
  \item La herencia de genes ubicados en el núcleo
  \item La herencia de genes ubicados en los cromosomas
  \item La herencia de genes ubicados en los autosomas
  \item La herencia de genes ubicados en los cromosomas sexuales
  \item La herencia de genes ubicados en los ribosomas
\end{multicols}
\end{enumerate}
\vspace{0.15cm}
\item ¿Qué es la herencia mitocondrial?
\begin{enumerate}[label=\textbf{\Alph*)}]
\begin{multicols}{2}
  \item La herencia de genes ubicados en las mitocondrias
  \item La herencia de genes ubicados en el núcleo
  \item La herencia de genes ubicados en los cromosomas
  \item La herencia de genes ubicados en los autosomas
  \item La herencia de genes ubicados en los cromosomas sexuales
  \item La herencia de genes ubicados en los ribosomas
\end{multicols}
\end{enumerate}
\vspace{0.15cm}
\item ¿Qué es la impronta genómica?
\begin{enumerate}[label=\textbf{\Alph*)}]
\begin{multicols}{2}
  \item La expresión diferencial de genes según su origen parental
  \item La expresión igual de genes sin importar su origen
  \item La expresión de genes solo maternos
  \item La expresión de genes solo paternos
  \item La expresión de genes solo en hombres
  \item La expresión de genes solo en mujeres
\end{multicols}
\end{enumerate}
\vspace{0.15cm}
\item ¿Qué es la mutación?
\begin{enumerate}[label=\textbf{\Alph*)}]
\begin{multicols}{2}
  \item Un cambio en la secuencia del ADN
  \item Un cambio en la secuencia del ARN
  \item Un cambio en la secuencia de proteínas
  \item Un cambio en la secuencia de lípidos
  \item Un cambio en la secuencia de carbohidratos
  \item Un cambio en la secuencia de minerales
\end{multicols}
\end{enumerate}
\vspace{0.15cm}
\item ¿Qué es una mutación puntual?
\begin{enumerate}[label=\textbf{\Alph*)}]
\begin{multicols}{2}
  \item Un cambio en un solo nucleótido
  \item Un cambio en varios nucleótidos
  \item Un cambio en un cromosoma completo
  \item Un cambio en varios cromosomas
  \item Un cambio en un gen completo
  \item Un cambio en varios genes
\end{multicols}
\end{enumerate}
\vspace{0.15cm}
\item ¿Qué es una mutación cromosómica?
\begin{enumerate}[label=\textbf{\Alph*)}]
\begin{multicols}{2}
  \item Un cambio en la estructura o número de cromosomas
  \item Un cambio en un solo nucleótido
  \item Un cambio en varios nucleótidos
  \item Un cambio en un gen completo
  \item Un cambio en varios genes
  \item Un cambio en una proteína
\end{multicols}
\end{enumerate}
\vspace{0.15cm}
\item ¿Qué es la aneuploidía?
\begin{enumerate}[label=\textbf{\Alph*)}]
\begin{multicols}{2}
  \item La presencia de un número anormal de cromosomas
  \item La presencia de un número normal de cromosomas
  \item La presencia de un número par de cromosomas
  \item La presencia de un número impar de cromosomas
  \item La presencia de un número múltiplo de cromosomas
  \item La presencia de un número fraccionario de cromosomas
\end{multicols}
\end{enumerate}
\vspace{0.15cm}
\item ¿Qué es el síndrome de Down?
\begin{enumerate}[label=\textbf{\Alph*)}]
\begin{multicols}{2}
  \item Una trisomía del cromosoma 21
  \item Una monosomía del cromosoma 21
  \item Una trisomía del cromosoma 18
  \item Una monosomía del cromosoma 18
  \item Una trisomía del cromosoma 13
  \item Una monosomía del cromosoma 13
\end{multicols}
\end{enumerate}
\vspace{0.15cm}
\item ¿Qué es el síndrome de Turner?
\begin{enumerate}[label=\textbf{\Alph*)}]
\begin{multicols}{2}
  \item Una monosomía del cromosoma X
  \item Una trisomía del cromosoma X
  \item Una monosomía del cromosoma Y
  \item Una trisomía del cromosoma Y
  \item Una monosomía del cromosoma 21
  \item Una trisomía del cromosoma 21
\end{multicols}
\end{enumerate}
\vspace{0.15cm}
\item ¿Qué es el síndrome de Klinefelter?
\begin{enumerate}[label=\textbf{\Alph*)}]
\begin{multicols}{2}
  \item Una trisomía de los cromosomas sexuales (XXY)
  \item Una monosomía de los cromosomas sexuales
  \item Una trisomía del cromosoma 21
  \item Una monosomía del cromosoma 21
  \item Una trisomía del cromosoma 18
  \item Una monosomía del cromosoma 18
\end{multicols}
\end{enumerate}
\vspace{0.15cm}
\item ¿Qué es un cariotipo?
\begin{enumerate}[label=\textbf{\Alph*)}]
\begin{multicols}{2}
  \item La representación ordenada de los cromosomas de una célula
  \item La representación ordenada de los genes de una célula
  \item La representación ordenada de las proteínas de una célula
  \item La representación ordenada de los lípidos de una célula
  \item La representación ordenada de los carbohidratos de una célula
  \item La representación ordenada de los minerales de una célula
\end{multicols}
\end{enumerate}
\vspace{0.15cm}
\item ¿Cuántos cromosomas tiene la especie humana?
\begin{enumerate}[label=\textbf{\Alph*)}]
\begin{multicols}{2}
  \item 46 (23 pares)
  \item 44 (22 pares)
  \item 48 (24 pares)
  \item 42 (21 pares)
  \item 50 (25 pares)
  \item 40 (20 pares)
\end{multicols}
\end{enumerate}
\vspace{0.15cm}
\item ¿Cuántos cromosomas sexuales tiene la mujer?
\begin{enumerate}[label=\textbf{\Alph*)}]
\begin{multicols}{2}
  \item XX
  \item XY
  \item XO
  \item XXY
  \item XYY
  \item YY
\end{multicols}
\end{enumerate}
\vspace{0.15cm}
\item ¿Cuántos cromosomas sexuales tiene el hombre?
\begin{enumerate}[label=\textbf{\Alph*)}]
\begin{multicols}{2}
  \item XY
  \item XX
  \item XO
  \item XXY
  \item XYY
  \item YY
\end{multicols}
\end{enumerate}
\vspace{0.15cm}
\item ¿Qué es la autosoma?
\begin{enumerate}[label=\textbf{\Alph*)}]
\begin{multicols}{2}
  \item Un cromosoma no sexual
  \item Un cromosoma sexual
  \item Un cromosoma X
  \item Un cromosoma Y
  \item Un cromosoma duplicado
  \item Un cromosoma homólogo
\end{multicols}
\end{enumerate}
\vspace{0.15cm}
\item ¿Qué es el genotipo?
\begin{enumerate}[label=\textbf{\Alph*)}]
\begin{multicols}{2}
  \item La composición genética de un organismo
  \item La apariencia externa de un organismo
  \item El comportamiento de un organismo
  \item El ambiente de un organismo
  \item La alimentación de un organismo
  \item La reproducción de un organismo
\end{multicols}
\end{enumerate}
\vspace{0.15cm}
\item ¿Qué es el fenotipo?
\begin{enumerate}[label=\textbf{\Alph*)}]
\begin{multicols}{2}
  \item La manifestación observable del genotipo
  \item La composición genética de un organismo
  \item El comportamiento de un organismo
  \item El ambiente de un organismo
  \item La alimentación de un organismo
  \item La reproducción de un organismo
\end{multicols}
\end{enumerate}
\vspace{0.15cm}
\item ¿Qué es un alelo?
\begin{enumerate}[label=\textbf{\Alph*)}]
\begin{multicols}{2}
  \item Cada una de las variantes de un gen
  \item Cada una de las variantes de un cromosoma
  \item Cada una de las variantes de una proteína
  \item Cada una de las variantes de un lípido
  \item Cada una de las variantes de un carbohidrato
  \item Cada una de las variantes de un mineral
\end{multicols}
\end{enumerate}
\vspace{0.15cm}
\item ¿Qué es un organismo homocigoto?
\begin{enumerate}[label=\textbf{\Alph*)}]
\begin{multicols}{2}
  \item El que tiene dos alelos iguales para un gen
  \item El que tiene dos alelos diferentes para un gen
  \item El que tiene un solo alelo para un gen
  \item El que tiene tres alelos para un gen
  \item El que tiene cuatro alelos para un gen
  \item El que no tiene alelos para un gen ---
\end{multicols}
\end{enumerate}
\vspace{0.15cm}
\end{enumerate}\n
\subsubsection*{SECCIÓN IV: GENÉTICA MOLECULAR – ADN Y HERENCIA (181–240)}

\begin{enumerate}[label=\textbf{\arabic*.}]
\setcounter{enumi}{180}
\item ¿Qué es el ADN?
\begin{enumerate}[label=\textbf{\Alph*)}]
\begin{multicols}{2}
  \item El ácido desoxirribonucleico que contiene la información genética
  \item El ácido ribonucleico que contiene la información genética
  \item El ácido desoxirribonucleico que contiene la información proteica
  \item El ácido ribonucleico que contiene la información proteica
  \item El ácido desoxirribonucleico que contiene la información lipídica
  \item El ácido ribonucleico que contiene la información lipídica
\end{multicols}
\end{enumerate}
\vspace{0.15cm}
\item ¿Quiénes descubrieron la estructura de doble hélice del ADN?
\begin{enumerate}[label=\textbf{\Alph*)}]
\begin{multicols}{2}
  \item Watson y Crick
  \item Mendel y Darwin
  \item Pasteur y Koch
  \item Hooke y Leeuwenhoek
  \item Fleming y Florey
  \item Newton y Einstein
\end{multicols}
\end{enumerate}
\vspace{0.15cm}
\item ¿Cuáles son las bases nitrogenadas del ADN?
\begin{enumerate}[label=\textbf{\Alph*)}]
\begin{multicols}{2}
  \item Adenina, timina, citosina y guanina
  \item Adenina, uracilo, citosina y guanina
  \item Adenina, timina, citosina y uracilo
  \item Adenina, uracilo, citosina y timina
  \item Adenina, guanina, uracilo y timina
  \item Adenina, citosina, uracilo y guanina
\end{multicols}
\end{enumerate}
\vspace{0.15cm}
\item ¿Cuáles son las bases nitrogenadas del ARN?
\begin{enumerate}[label=\textbf{\Alph*)}]
\begin{multicols}{2}
  \item Adenina, uracilo, citosina y guanina
  \item Adenina, timina, citosina y guanina
  \item Adenina, timina, citosina y uracilo
  \item Adenina, uracilo, citosina y timina
  \item Adenina, guanina, uracilo y timina
  \item Adenina, citosina, uracilo y guanina
\end{multicols}
\end{enumerate}
\vspace{0.15cm}
\item ¿Con qué base se empareja la adenina en el ADN?
\begin{enumerate}[label=\textbf{\Alph*)}]
\begin{multicols}{2}
  \item Con timina
  \item Con uracilo
  \item Con citosina
  \item Con guanina
  \item Con adenina
  \item Con otra base
\end{multicols}
\end{enumerate}
\vspace{0.15cm}
\item ¿Con qué base se empareja la guanina en el ADN?
\begin{enumerate}[label=\textbf{\Alph*)}]
\begin{multicols}{2}
  \item Con citosina
  \item Con timina
  \item Con uracilo
  \item Con adenina
  \item Con guanina
  \item Con otra base
\end{multicols}
\end{enumerate}
\vspace{0.15cm}
\item ¿Con qué base se empareja la adenina en el ARN?
\begin{enumerate}[label=\textbf{\Alph*)}]
\begin{multicols}{2}
  \item Con uracilo
  \item Con timina
  \item Con citosina
  \item Con guanina
  \item Con adenina
  \item Con otra base
\end{multicols}
\end{enumerate}
\vspace{0.15cm}
\item ¿Qué es la replicación del ADN?
\begin{enumerate}[label=\textbf{\Alph*)}]
\begin{multicols}{2}
  \item El proceso de copiar el ADN antes de la división celular
  \item El proceso de copiar el ARN
  \item El proceso de copiar proteínas
  \item El proceso de copiar lípidos
  \item El proceso de copiar carbohidratos
  \item El proceso de copiar minerales
\end{multicols}
\end{enumerate}
\vspace{0.15cm}
\item ¿Qué es la transcripción?
\begin{enumerate}[label=\textbf{\Alph*)}]
\begin{multicols}{2}
  \item El proceso de copiar el ADN en ARN
  \item El proceso de copiar el ARN en ADN
  \item El proceso de copiar el ARN en proteínas
  \item El proceso de copiar proteínas en ADN
  \item El proceso de copiar lípidos en ADN
  \item El proceso de copiar carbohidratos en ADN
\end{multicols}
\end{enumerate}
\vspace{0.15cm}
\item ¿Qué es la traducción?
\begin{enumerate}[label=\textbf{\Alph*)}]
\begin{multicols}{2}
  \item El proceso de sintetizar proteínas a partir del ARN
  \item El proceso de sintetizar ARN a partir del ADN
  \item El proceso de sintetizar ADN a partir del ARN
  \item El proceso de sintetizar lípidos a partir del ARN
  \item El proceso de sintetizar carbohidratos a partir del ARN
  \item El proceso de sintetizar minerales a partir del ARN
\end{multicols}
\end{enumerate}
\vspace{0.15cm}
\item ¿Qué es el código genético?
\begin{enumerate}[label=\textbf{\Alph*)}]
\begin{multicols}{2}
  \item El conjunto de reglas que relacionan los codones con los aminoácidos
  \item El conjunto de reglas que relacionan los genes con los cromosomas
  \item El conjunto de reglas que relacionan las proteínas con los lípidos
  \item El conjunto de reglas que relacionan los carbohidratos con los minerales
  \item El conjunto de reglas que relacionan las células con los tejidos
  \item El conjunto de reglas que relacionan los órganos con los sistemas
\end{multicols}
\end{enumerate}
\vspace{0.15cm}
\item ¿Qué es un codón?
\begin{enumerate}[label=\textbf{\Alph*)}]
\begin{multicols}{2}
  \item Un triplete de nucleótidos que codifica un aminoácido
  \item Un par de nucleótidos que codifica un aminoácido
  \item Un cuarteto de nucleótidos que codifica un aminoácido
  \item Un quinteto de nucleótidos que codifica un aminoácido
  \item Un sexteto de nucleótidos que codifica un aminoácido
  \item Un septeto de nucleótidos que codifica un aminoácido
\end{multicols}
\end{enumerate}
\vspace{0.15cm}
\item ¿Cuántos codones existen?
\begin{enumerate}[label=\textbf{\Alph*)}]
\begin{multicols}{2}
  \item 64
  \item 20
  \item 32
  \item 16
  \item 48
  \item 60
\end{multicols}
\end{enumerate}
\vspace{0.15cm}
\item ¿Cuántos aminoácidos existen?
\begin{enumerate}[label=\textbf{\Alph*)}]
\begin{multicols}{2}
  \item 20
  \item 64
  \item 32
  \item 16
  \item 48
  \item 60
\end{multicols}
\end{enumerate}
\vspace{0.15cm}
\item ¿Qué es un anticodón?
\begin{enumerate}[label=\textbf{\Alph*)}]
\begin{multicols}{2}
  \item Un triplete de nucleótidos del ARNt que se empareja con el codón
  \item Un par de nucleótidos del ARNt
  \item Un cuarteto de nucleótidos del ARNt
  \item Un quinteto de nucleótidos del ARNt
  \item Un sexteto de nucleótidos del ARNt
  \item Un septeto de nucleótidos del ARNt
\end{multicols}
\end{enumerate}
\vspace{0.15cm}
\item ¿Qué es el ARN mensajero?
\begin{enumerate}[label=\textbf{\Alph*)}]
\begin{multicols}{2}
  \item El que lleva la información genética del ADN a los ribosomas
  \item El que lleva los aminoácidos a los ribosomas
  \item El que forma parte de los ribosomas
  \item El que regula la expresión génica
  \item El que cataliza reacciones químicas
  \item El que almacena energía
\end{multicols}
\end{enumerate}
\vspace{0.15cm}
\item ¿Qué es el ARN de transferencia?
\begin{enumerate}[label=\textbf{\Alph*)}]
\begin{multicols}{2}
  \item El que lleva los aminoácidos a los ribosomas
  \item El que lleva la información genética del ADN
  \item El que forma parte de los ribosomas
  \item El que regula la expresión génica
  \item El que cataliza reacciones químicas
  \item El que almacena energía
\end{multicols}
\end{enumerate}
\vspace{0.15cm}
\item ¿Qué es el ARN ribosómico?
\begin{enumerate}[label=\textbf{\Alph*)}]
\begin{multicols}{2}
  \item El que forma parte de los ribosomas
  \item El que lleva la información genética del ADN
  \item El que lleva los aminoácidos a los ribosomas
  \item El que regula la expresión génica
  \item El que cataliza reacciones químicas
  \item El que almacena energía
\end{multicols}
\end{enumerate}
\vspace{0.15cm}
\item ¿Qué es la síntesis de proteínas?
\begin{enumerate}[label=\textbf{\Alph*)}]
\begin{multicols}{2}
  \item El proceso de formación de proteínas a partir de aminoácidos
  \item El proceso de formación de aminoácidos a partir de proteínas
  \item El proceso de formación de ADN a partir de proteínas
  \item El proceso de formación de ARN a partir de proteínas
  \item El proceso de formación de lípidos a partir de proteínas
  \item El proceso de formación de carbohidratos a partir de proteínas
\end{multicols}
\end{enumerate}
\vspace{0.15cm}
\item ¿Qué es una enzima?
\begin{enumerate}[label=\textbf{\Alph*)}]
\begin{multicols}{2}
  \item Una proteína que cataliza reacciones químicas
  \item Un lípido que cataliza reacciones químicas
  \item Un carbohidrato que cataliza reacciones químicas
  \item Un ácido nucleico que cataliza reacciones químicas
  \item Un mineral que cataliza reacciones químicas
  \item Un gas que cataliza reacciones químicas
\end{multicols}
\end{enumerate}
\vspace{0.15cm}
\item ¿Qué es el ADN?
\begin{enumerate}[label=\textbf{\Alph*)}]
\begin{multicols}{2}
  \item El ácido desoxirribonucleico que contiene la información genética
  \item El ácido ribonucleico que contiene la información genética
  \item El ácido desoxirribonucleico que contiene la información proteica
  \item El ácido ribonucleico que contiene la información proteica
  \item El ácido desoxirribonucleico que contiene la información lipídica
  \item El ácido ribonucleico que contiene la información lipídica
\end{multicols}
\end{enumerate}
\vspace{0.15cm}
\item ¿Quiénes descubrieron la estructura de doble hélice del ADN?
\begin{enumerate}[label=\textbf{\Alph*)}]
\begin{multicols}{2}
  \item Watson y Crick
  \item Mendel y Darwin
  \item Pasteur y Koch
  \item Hooke y Leeuwenhoek
  \item Fleming y Florey
  \item Newton y Einstein
\end{multicols}
\end{enumerate}
\vspace{0.15cm}
\item ¿Cuáles son las bases nitrogenadas del ADN?
\begin{enumerate}[label=\textbf{\Alph*)}]
\begin{multicols}{2}
  \item Adenina, timina, citosina y guanina
  \item Adenina, uracilo, citosina y guanina
  \item Adenina, timina, citosina y uracilo
  \item Adenina, uracilo, citosina y timina
  \item Adenina, guanina, uracilo y timina
  \item Adenina, citosina, uracilo y guanina
\end{multicols}
\end{enumerate}
\vspace{0.15cm}
\item ¿Cuáles son las bases nitrogenadas del ARN?
\begin{enumerate}[label=\textbf{\Alph*)}]
\begin{multicols}{2}
  \item Adenina, uracilo, citosina y guanina
  \item Adenina, timina, citosina y guanina
  \item Adenina, timina, citosina y uracilo
  \item Adenina, uracilo, citosina y timina
  \item Adenina, guanina, uracilo y timina
  \item Adenina, citosina, uracilo y guanina
\end{multicols}
\end{enumerate}
\vspace{0.15cm}
\item ¿Con qué base se empareja la adenina en el ADN?
\begin{enumerate}[label=\textbf{\Alph*)}]
\begin{multicols}{2}
  \item Con timina
  \item Con uracilo
  \item Con citosina
  \item Con guanina
  \item Con adenina
  \item Con otra base
\end{multicols}
\end{enumerate}
\vspace{0.15cm}
\item ¿Con qué base se empareja la guanina en el ADN?
\begin{enumerate}[label=\textbf{\Alph*)}]
\begin{multicols}{2}
  \item Con citosina
  \item Con timina
  \item Con uracilo
  \item Con adenina
  \item Con guanina
  \item Con otra base
\end{multicols}
\end{enumerate}
\vspace{0.15cm}
\item ¿Con qué base se empareja la adenina en el ARN?
\begin{enumerate}[label=\textbf{\Alph*)}]
\begin{multicols}{2}
  \item Con uracilo
  \item Con timina
  \item Con citosina
  \item Con guanina
  \item Con adenina
  \item Con otra base
\end{multicols}
\end{enumerate}
\vspace{0.15cm}
\item ¿Qué es la replicación del ADN?
\begin{enumerate}[label=\textbf{\Alph*)}]
\begin{multicols}{2}
  \item El proceso de copiar el ADN antes de la división celular
  \item El proceso de copiar el ARN
  \item El proceso de copiar proteínas
  \item El proceso de copiar lípidos
  \item El proceso de copiar carbohidratos
  \item El proceso de copiar minerales
\end{multicols}
\end{enumerate}
\vspace{0.15cm}
\item ¿Qué es la transcripción?
\begin{enumerate}[label=\textbf{\Alph*)}]
\begin{multicols}{2}
  \item El proceso de copiar el ADN en ARN
  \item El proceso de copiar el ARN en ADN
  \item El proceso de copiar el ARN en proteínas
  \item El proceso de copiar proteínas en ADN
  \item El proceso de copiar lípidos en ADN
  \item El proceso de copiar carbohidratos en ADN
\end{multicols}
\end{enumerate}
\vspace{0.15cm}
\item ¿Qué es la traducción?
\begin{enumerate}[label=\textbf{\Alph*)}]
\begin{multicols}{2}
  \item El proceso de sintetizar proteínas a partir del ARN
  \item El proceso de sintetizar ARN a partir del ADN
  \item El proceso de sintetizar ADN a partir del ARN
  \item El proceso de sintetizar lípidos a partir del ARN
  \item El proceso de sintetizar carbohidratos a partir del ARN
  \item El proceso de sintetizar minerales a partir del ARN
\end{multicols}
\end{enumerate}
\vspace{0.15cm}
\item ¿Qué es el código genético?
\begin{enumerate}[label=\textbf{\Alph*)}]
\begin{multicols}{2}
  \item El conjunto de reglas que relacionan los codones con los aminoácidos
  \item El conjunto de reglas que relacionan los genes con los cromosomas
  \item El conjunto de reglas que relacionan las proteínas con los lípidos
  \item El conjunto de reglas que relacionan los carbohidratos con los minerales
  \item El conjunto de reglas que relacionan las células con los tejidos
  \item El conjunto de reglas que relacionan los órganos con los sistemas
\end{multicols}
\end{enumerate}
\vspace{0.15cm}
\item ¿Qué es un codón?
\begin{enumerate}[label=\textbf{\Alph*)}]
\begin{multicols}{2}
  \item Un triplete de nucleótidos que codifica un aminoácido
  \item Un par de nucleótidos que codifica un aminoácido
  \item Un cuarteto de nucleótidos que codifica un aminoácido
  \item Un quinteto de nucleótidos que codifica un aminoácido
  \item Un sexteto de nucleótidos que codifica un aminoácido
  \item Un septeto de nucleótidos que codifica un aminoácido
\end{multicols}
\end{enumerate}
\vspace{0.15cm}
\item ¿Cuántos codones existen?
\begin{enumerate}[label=\textbf{\Alph*)}]
\begin{multicols}{2}
  \item 64
  \item 20
  \item 32
  \item 16
  \item 48
  \item 60
\end{multicols}
\end{enumerate}
\vspace{0.15cm}
\item ¿Cuántos aminoácidos existen?
\begin{enumerate}[label=\textbf{\Alph*)}]
\begin{multicols}{2}
  \item 20
  \item 64
  \item 32
  \item 16
  \item 48
  \item 60
\end{multicols}
\end{enumerate}
\vspace{0.15cm}
\item ¿Qué es un anticodón?
\begin{enumerate}[label=\textbf{\Alph*)}]
\begin{multicols}{2}
  \item Un triplete de nucleótidos del ARNt que se empareja con el codón
  \item Un par de nucleótidos del ARNt
  \item Un cuarteto de nucleótidos del ARNt
  \item Un quinteto de nucleótidos del ARNt
  \item Un sexteto de nucleótidos del ARNt
  \item Un septeto de nucleótidos del ARNt
\end{multicols}
\end{enumerate}
\vspace{0.15cm}
\item ¿Qué es el ARN mensajero?
\begin{enumerate}[label=\textbf{\Alph*)}]
\begin{multicols}{2}
  \item El que lleva la información genética del ADN a los ribosomas
  \item El que lleva los aminoácidos a los ribosomas
  \item El que forma parte de los ribosomas
  \item El que regula la expresión génica
  \item El que cataliza reacciones químicas
  \item El que almacena energía
\end{multicols}
\end{enumerate}
\vspace{0.15cm}
\item ¿Qué es el ARN de transferencia?
\begin{enumerate}[label=\textbf{\Alph*)}]
\begin{multicols}{2}
  \item El que lleva los aminoácidos a los ribosomas
  \item El que lleva la información genética del ADN
  \item El que forma parte de los ribosomas
  \item El que regula la expresión génica
  \item El que cataliza reacciones químicas
  \item El que almacena energía
\end{multicols}
\end{enumerate}
\vspace{0.15cm}
\item ¿Qué es el ARN ribosómico?
\begin{enumerate}[label=\textbf{\Alph*)}]
\begin{multicols}{2}
  \item El que forma parte de los ribosomas
  \item El que lleva la información genética del ADN
  \item El que lleva los aminoácidos a los ribosomas
  \item El que regula la expresión génica
  \item El que cataliza reacciones químicas
  \item El que almacena energía
\end{multicols}
\end{enumerate}
\vspace{0.15cm}
\item ¿Qué es la síntesis de proteínas?
\begin{enumerate}[label=\textbf{\Alph*)}]
\begin{multicols}{2}
  \item El proceso de formación de proteínas a partir de aminoácidos
  \item El proceso de formación de aminoácidos a partir de proteínas
  \item El proceso de formación de ADN a partir de proteínas
  \item El proceso de formación de ARN a partir de proteínas
  \item El proceso de formación de lípidos a partir de proteínas
  \item El proceso de formación de carbohidratos a partir de proteínas
\end{multicols}
\end{enumerate}
\vspace{0.15cm}
\item ¿Qué es una enzima?
\begin{enumerate}[label=\textbf{\Alph*)}]
\begin{multicols}{2}
  \item Una proteína que cataliza reacciones químicas
  \item Un lípido que cataliza reacciones químicas
  \item Un carbohidrato que cataliza reacciones químicas
  \item Un ácido nucleico que cataliza reacciones químicas
  \item Un mineral que cataliza reacciones químicas
  \item Un gas que cataliza reacciones químicas
\end{multicols}
\end{enumerate}
\vspace{0.15cm}
\item ¿Qué es la mutación?
\begin{enumerate}[label=\textbf{\Alph*)}]
\begin{multicols}{2}
  \item Un cambio en la secuencia del ADN
  \item Un cambio en la secuencia del ARN
  \item Un cambio en la secuencia de proteínas
  \item Un cambio en la secuencia de lípidos
  \item Un cambio en la secuencia de carbohidratos
  \item Un cambio en la secuencia de minerales
\end{multicols}
\end{enumerate}
\vspace{0.15cm}
\item ¿Qué es una mutación puntual?
\begin{enumerate}[label=\textbf{\Alph*)}]
\begin{multicols}{2}
  \item Un cambio en un solo nucleótido
  \item Un cambio en varios nucleótidos
  \item Un cambio en un cromosoma completo
  \item Un cambio en varios cromosomas
  \item Un cambio en un gen completo
  \item Un cambio en varios genes
\end{multicols}
\end{enumerate}
\vspace{0.15cm}
\item ¿Qué es una mutación cromosómica?
\begin{enumerate}[label=\textbf{\Alph*)}]
\begin{multicols}{2}
  \item Un cambio en la estructura o número de cromosomas
  \item Un cambio en un solo nucleótido
  \item Un cambio en varios nucleótidos
  \item Un cambio en un gen completo
  \item Un cambio en varios genes
  \item Un cambio en una proteína
\end{multicols}
\end{enumerate}
\vspace{0.15cm}
\item ¿Qué es la aneuploidía?
\begin{enumerate}[label=\textbf{\Alph*)}]
\begin{multicols}{2}
  \item La presencia de un número anormal de cromosomas
  \item La presencia de un número normal de cromosomas
  \item La presencia de un número par de cromosomas
  \item La presencia de un número impar de cromosomas
  \item La presencia de un número múltiplo de cromosomas
  \item La presencia de un número fraccionario de cromosomas
\end{multicols}
\end{enumerate}
\vspace{0.15cm}
\item ¿Qué es el síndrome de Down?
\begin{enumerate}[label=\textbf{\Alph*)}]
\begin{multicols}{2}
  \item Una trisomía del cromosoma 21
  \item Una monosomía del cromosoma 21
  \item Una trisomía del cromosoma 18
  \item Una monosomía del cromosoma 18
  \item Una trisomía del cromosoma 13
  \item Una monosomía del cromosoma 13
\end{multicols}
\end{enumerate}
\vspace{0.15cm}
\item ¿Qué es el síndrome de Turner?
\begin{enumerate}[label=\textbf{\Alph*)}]
\begin{multicols}{2}
  \item Una monosomía del cromosoma X
  \item Una trisomía del cromosoma X
  \item Una monosomía del cromosoma Y
  \item Una trisomía del cromosoma Y
  \item Una monosomía del cromosoma 21
  \item Una trisomía del cromosoma 21
\end{multicols}
\end{enumerate}
\vspace{0.15cm}
\item ¿Qué es el síndrome de Klinefelter?
\begin{enumerate}[label=\textbf{\Alph*)}]
\begin{multicols}{2}
  \item Una trisomía de los cromosomas sexuales (XXY)
  \item Una monosomía de los cromosomas sexuales
  \item Una trisomía del cromosoma 21
  \item Una monosomía del cromosoma 21
  \item Una trisomía del cromosoma 18
  \item Una monosomía del cromosoma 18
\end{multicols}
\end{enumerate}
\vspace{0.15cm}
\item ¿Qué es un cariotipo?
\begin{enumerate}[label=\textbf{\Alph*)}]
\begin{multicols}{2}
  \item La representación ordenada de los cromosomas de una célula
  \item La representación ordenada de los genes de una célula
  \item La representación ordenada de las proteínas de una célula
  \item La representación ordenada de los lípidos de una célula
  \item La representación ordenada de los carbohidratos de una célula
  \item La representación ordenada de los minerales de una célula
\end{multicols}
\end{enumerate}
\vspace{0.15cm}
\item ¿Cuántos cromosomas tiene la especie humana?
\begin{enumerate}[label=\textbf{\Alph*)}]
\begin{multicols}{2}
  \item 46 (23 pares)
  \item 44 (22 pares)
  \item 48 (24 pares)
  \item 42 (21 pares)
  \item 50 (25 pares)
  \item 40 (20 pares)
\end{multicols}
\end{enumerate}
\vspace{0.15cm}
\item ¿Cuántos cromosomas sexuales tiene la mujer?
\begin{enumerate}[label=\textbf{\Alph*)}]
\begin{multicols}{2}
  \item XX
  \item XY
  \item XO
  \item XXY
  \item XYY
  \item YY
\end{multicols}
\end{enumerate}
\vspace{0.15cm}
\item ¿Cuántos cromosomas sexuales tiene el hombre?
\begin{enumerate}[label=\textbf{\Alph*)}]
\begin{multicols}{2}
  \item XY
  \item XX
  \item XO
  \item XXY
  \item XYY
  \item YY
\end{multicols}
\end{enumerate}
\vspace{0.15cm}
\item ¿Qué es la autosoma?
\begin{enumerate}[label=\textbf{\Alph*)}]
\begin{multicols}{2}
  \item Un cromosoma no sexual
  \item Un cromosoma sexual
  \item Un cromosoma X
  \item Un cromosoma Y
  \item Un cromosoma duplicado
  \item Un cromosoma homólogo
\end{multicols}
\end{enumerate}
\vspace{0.15cm}
\item ¿Qué es el genotipo?
\begin{enumerate}[label=\textbf{\Alph*)}]
\begin{multicols}{2}
  \item La composición genética de un organismo
  \item La apariencia externa de un organismo
  \item El comportamiento de un organismo
  \item El ambiente de un organismo
  \item La alimentación de un organismo
  \item La reproducción de un organismo
\end{multicols}
\end{enumerate}
\vspace{0.15cm}
\item ¿Qué es el fenotipo?
\begin{enumerate}[label=\textbf{\Alph*)}]
\begin{multicols}{2}
  \item La manifestación observable del genotipo
  \item La composición genética de un organismo
  \item El comportamiento de un organismo
  \item El ambiente de un organismo
  \item La alimentación de un organismo
  \item La reproducción de un organismo
\end{multicols}
\end{enumerate}
\vspace{0.15cm}
\item ¿Qué es un alelo?
\begin{enumerate}[label=\textbf{\Alph*)}]
\begin{multicols}{2}
  \item Cada una de las variantes de un gen
  \item Cada una de las variantes de un cromosoma
  \item Cada una de las variantes de una proteína
  \item Cada una de las variantes de un lípido
  \item Cada una de las variantes de un carbohidrato
  \item Cada una de las variantes de un mineral
\end{multicols}
\end{enumerate}
\vspace{0.15cm}
\item ¿Qué es un organismo homocigoto?
\begin{enumerate}[label=\textbf{\Alph*)}]
\begin{multicols}{2}
  \item El que tiene dos alelos iguales para un gen
  \item El que tiene dos alelos diferentes para un gen
  \item El que tiene un solo alelo para un gen
  \item El que tiene tres alelos para un gen
  \item El que tiene cuatro alelos para un gen
  \item El que no tiene alelos para un gen
\end{multicols}
\end{enumerate}
\vspace{0.15cm}
\item ¿Qué es un organismo heterocigoto?
\begin{enumerate}[label=\textbf{\Alph*)}]
\begin{multicols}{2}
  \item El que tiene dos alelos diferentes para un gen
  \item El que tiene dos alelos iguales para un gen
  \item El que tiene un solo alelo para un gen
  \item El que tiene tres alelos para un gen
  \item El que tiene cuatro alelos para un gen
  \item El que no tiene alelos para un gen
\end{multicols}
\end{enumerate}
\vspace{0.15cm}
\item ¿Qué es un alelo dominante?
\begin{enumerate}[label=\textbf{\Alph*)}]
\begin{multicols}{2}
  \item El que se manifiesta en el fenotipo aunque haya otro alelo diferente
  \item El que no se manifiesta en el fenotipo
  \item El que se manifiesta solo si hay dos copias
  \item El que se manifiesta solo si hay tres copias
  \item El que se manifiesta solo si hay cuatro copias
  \item El que no se manifiesta nunca
\end{multicols}
\end{enumerate}
\vspace{0.15cm}
\item ¿Qué es un alelo recesivo?
\begin{enumerate}[label=\textbf{\Alph*)}]
\begin{multicols}{2}
  \item El que se manifiesta solo si hay dos copias iguales
  \item El que se manifiesta siempre
  \item El que se manifiesta aunque haya otro alelo diferente
  \item El que se manifiesta solo si hay tres copias
  \item El que se manifiesta solo si hay cuatro copias
  \item El que se manifiesta solo si hay una copia
\end{multicols}
\end{enumerate}
\vspace{0.15cm}
\item ¿Qué es un cruce monohíbrido?
\begin{enumerate}[label=\textbf{\Alph*)}]
\begin{multicols}{2}
  \item El cruce que estudia un solo carácter
  \item El cruce que estudia dos caracteres
  \item El cruce que estudia tres caracteres
  \item El cruce que estudia cuatro caracteres
  \item El cruce que estudia cinco caracteres
  \item El cruce que estudia seis caracteres ---
\end{multicols}
\end{enumerate}
\vspace{0.15cm}
\end{enumerate}\n
\subsubsection*{SECCIÓN V: BIOLOGÍA CELULAR Y HERENCIA – APLICACIONES (241–300)}

\begin{enumerate}[label=\textbf{\arabic*.}]
\setcounter{enumi}{240}
\item ¿Qué es la biotecnología?
\begin{enumerate}[label=\textbf{\Alph*)}]
\begin{multicols}{2}
  \item El uso de organismos vivos o sus componentes para crear productos
  \item El uso de minerales para crear productos
  \item El uso de gases para crear productos
  \item El uso de líquidos para crear productos
  \item El uso de astros para crear productos
  \item El uso de rocas para crear productos
\end{multicols}
\end{enumerate}
\vspace{0.15cm}
\item ¿Qué es la ingeniería genética?
\begin{enumerate}[label=\textbf{\Alph*)}]
\begin{multicols}{2}
  \item La manipulación directa del ADN de un organismo
  \item La manipulación directa del ARN de un organismo
  \item La manipulación directa de las proteínas de un organismo
  \item La manipulación directa de los lípidos de un organismo
  \item La manipulación directa de los carbohidratos de un organismo
  \item La manipulación directa de los minerales de un organismo
\end{multicols}
\end{enumerate}
\vspace{0.15cm}
\item ¿Qué es un organismo transgénico?
\begin{enumerate}[label=\textbf{\Alph*)}]
\begin{multicols}{2}
  \item Un organismo que contiene genes de otra especie
  \item Un organismo que contiene genes de la misma especie
  \item Un organismo que no contiene genes
  \item Un organismo que contiene genes artificiales
  \item Un organismo que contiene genes naturales
  \item Un organismo que contiene genes sintéticos
\end{multicols}
\end{enumerate}
\vspace{0.15cm}
\item ¿Qué es la clonación?
\begin{enumerate}[label=\textbf{\Alph*)}]
\begin{multicols}{2}
  \item La producción de copias genéticamente idénticas de un organismo
  \item La producción de copias genéticamente diferentes
  \item La producción de copias artificiales
  \item La producción de copias naturales
  \item La producción de copias sintéticas
  \item La producción de copias híbridas
\end{multicols}
\end{enumerate}
\vspace{0.15cm}
\item ¿Qué es la terapia génica?
\begin{enumerate}[label=\textbf{\Alph*)}]
\begin{multicols}{2}
  \item La introducción de genes en células para tratar enfermedades
  \item La introducción de proteínas en células
  \item La introducción de lípidos en células
  \item La introducción de carbohidratos en células
  \item La introducción de minerales en células
  \item La introducción de gases en células
\end{multicols}
\end{enumerate}
\vspace{0.15cm}
\item ¿Qué es la farmacogenómica?
\begin{enumerate}[label=\textbf{\Alph*)}]
\begin{multicols}{2}
  \item El estudio de cómo los genes afectan la respuesta a los fármacos
  \item El estudio de cómo los fármacos afectan los genes
  \item El estudio de cómo los genes afectan las proteínas
  \item El estudio de cómo las proteínas afectan los genes
  \item El estudio de cómo los genes afectan los lípidos
  \item El estudio de cómo los lípidos afectan los genes
\end{multicols}
\end{enumerate}
\vspace{0.15cm}
\item ¿Qué es la genómica?
\begin{enumerate}[label=\textbf{\Alph*)}]
\begin{multicols}{2}
  \item El estudio del genoma completo de un organismo
  \item El estudio de un solo gen
  \item El estudio de un solo cromosoma
  \item El estudio de una sola proteína
  \item El estudio de un solo lípido
  \item El estudio de un solo carbohidrato
\end{multicols}
\end{enumerate}
\vspace{0.15cm}
\item ¿Qué es la proteómica?
\begin{enumerate}[label=\textbf{\Alph*)}]
\begin{multicols}{2}
  \item El estudio del conjunto de proteínas de un organismo
  \item El estudio del conjunto de genes de un organismo
  \item El estudio del conjunto de lípidos de un organismo
  \item El estudio del conjunto de carbohidratos de un organismo
  \item El estudio del conjunto de minerales de un organismo
  \item El estudio del conjunto de células de un organismo
\end{multicols}
\end{enumerate}
\vspace{0.15cm}
\item ¿Qué es la bioinformática?
\begin{enumerate}[label=\textbf{\Alph*)}]
\begin{multicols}{2}
  \item El uso de herramientas computacionales para analizar datos biológicos
  \item El uso de herramientas computacionales para analizar datos químicos
  \item El uso de herramientas computacionales para analizar datos físicos
  \item El uso de herramientas computacionales para analizar datos matemáticos
  \item El uso de herramientas computacionales para analizar datos geológicos
  \item El uso de herramientas computacionales para analizar datos astronómicos
\end{multicols}
\end{enumerate}
\vspace{0.15cm}
\item ¿Qué es la medicina personalizada?
\begin{enumerate}[label=\textbf{\Alph*)}]
\begin{multicols}{2}
  \item La medicina que adapta el tratamiento al perfil genético del paciente
  \item La medicina que trata a todos los pacientes igual
  \item La medicina que trata a todos los pacientes con el mismo fármaco
  \item La medicina que trata a todos los pacientes con la misma dosis
  \item La medicina que trata a todos los pacientes con la misma terapia
  \item La medicina que trata a todos los pacientes con el mismo diagnóstico
\end{multicols}
\end{enumerate}
\vspace{0.15cm}
\item ¿Qué es la biología?
\begin{enumerate}[label=\textbf{\Alph*)}]
\begin{multicols}{2}
  \item La ciencia que estudia los seres vivos
  \item La ciencia que estudia los minerales
  \item La ciencia que estudia los astros
  \item La ciencia que estudia los fenómenos físicos
  \item La ciencia que estudia las reacciones químicas
  \item La ciencia que estudia los números
\end{multicols}
\end{enumerate}
\vspace{0.15cm}
\item ¿Cuál es la característica principal de los seres vivos?
\begin{enumerate}[label=\textbf{\Alph*)}]
\begin{multicols}{2}
  \item Están formados por células
  \item Están formados por minerales
  \item Están formados por átomos metálicos
  \item Están formados por gases
  \item Están formados por líquidos
  \item Están formados por plásticos
\end{multicols}
\end{enumerate}
\vspace{0.15cm}
\item ¿Qué es la célula?
\begin{enumerate}[label=\textbf{\Alph*)}]
\begin{multicols}{2}
  \item La unidad estructural y funcional de los seres vivos
  \item La unidad estructural de los minerales
  \item La unidad funcional de los astros
  \item La unidad estructural del agua
  \item La unidad funcional del aire
  \item La unidad estructural del suelo
\end{multicols}
\end{enumerate}
\vspace{0.15cm}
\item ¿Quién descubrió la célula?
\begin{enumerate}[label=\textbf{\Alph*)}]
\begin{multicols}{2}
  \item Robert Hooke
  \item Charles Darwin
  \item Gregor Mendel
  \item Louis Pasteur
  \item Alexander Fleming
  \item Isaac Newton
\end{multicols}
\end{enumerate}
\vspace{0.15cm}
\item ¿Qué es la teoría celular?
\begin{enumerate}[label=\textbf{\Alph*)}]
\begin{multicols}{2}
  \item Todos los seres vivos están formados por células
  \item Todos los seres vivos están formados por minerales
  \item Todos los seres vivos están formados por átomos
  \item Todos los seres vivos están formados por moléculas
  \item Todos los seres vivos están formados por gases
  \item Todos los seres vivos están formados por líquidos
\end{multicols}
\end{enumerate}
\vspace{0.15cm}
\item ¿Cuál de los siguientes postulados corresponde a la teoría celular?
\begin{enumerate}[label=\textbf{\Alph*)}]
\begin{multicols}{2}
  \item La célula es la unidad básica de la vida
  \item La célula es la unidad básica de los minerales
  \item La célula es la unidad básica del agua
  \item La célula es la unidad básica del aire
  \item La célula es la unidad básica del suelo
  \item La célula es la unidad básica del fuego
\end{multicols}
\end{enumerate}
\vspace{0.15cm}
\item ¿Qué es el metabolismo?
\begin{enumerate}[label=\textbf{\Alph*)}]
\begin{multicols}{2}
  \item El conjunto de reacciones químicas que ocurren en los seres vivos
  \item El conjunto de reacciones físicas que ocurren en los minerales
  \item El conjunto de reacciones químicas que ocurren en el agua
  \item El conjunto de reacciones físicas que ocurren en el aire
  \item El conjunto de reacciones químicas que ocurren en el suelo
  \item El conjunto de reacciones físicas que ocurren en los astros
\end{multicols}
\end{enumerate}
\vspace{0.15cm}
\item ¿Qué es el anabolismo?
\begin{enumerate}[label=\textbf{\Alph*)}]
\begin{multicols}{2}
  \item El conjunto de reacciones de síntesis que construyen moléculas
  \item El conjunto de reacciones de degradación
  \item El conjunto de reacciones de oxidación
  \item El conjunto de reacciones de reducción
  \item El conjunto de reacciones de hidrólisis
  \item El conjunto de reacciones de combustión
\end{multicols}
\end{enumerate}
\vspace{0.15cm}
\item ¿Qué es el catabolismo?
\begin{enumerate}[label=\textbf{\Alph*)}]
\begin{multicols}{2}
  \item El conjunto de reacciones de degradación que liberan energía
  \item El conjunto de reacciones de síntesis
  \item El conjunto de reacciones de construcción
  \item El conjunto de reacciones de formación
  \item El conjunto de reacciones de producción
  \item El conjunto de reacciones de elaboración
\end{multicols}
\end{enumerate}
\vspace{0.15cm}
\item ¿Qué es la homeostasis?
\begin{enumerate}[label=\textbf{\Alph*)}]
\begin{multicols}{2}
  \item La capacidad de mantener el equilibrio interno
  \item La capacidad de mantener el desequilibrio interno
  \item La capacidad de mantener el equilibrio externo
  \item La capacidad de mantener el desequilibrio externo
  \item La capacidad de cambiar el equilibrio interno
  \item La capacidad de cambiar el equilibrio externo
\end{multicols}
\end{enumerate}
\vspace{0.15cm}
\item ¿Cuál es un ejemplo de homeostasis?
\begin{enumerate}[label=\textbf{\Alph*)}]
\begin{multicols}{2}
  \item La regulación de la temperatura corporal
  \item La regulación de la temperatura ambiental
  \item La regulación de la temperatura del agua
  \item La regulación de la temperatura del aire
  \item La regulación de la temperatura del suelo
  \item La regulación de la temperatura de los astros
\end{multicols}
\end{enumerate}
\vspace{0.15cm}
\item ¿Qué es la irritabilidad?
\begin{enumerate}[label=\textbf{\Alph*)}]
\begin{multicols}{2}
  \item La capacidad de responder a estímulos
  \item La capacidad de no responder a estímulos
  \item La capacidad de ignorar estímulos
  \item La capacidad de olvidar estímulos
  \item La capacidad de negar estímulos
  \item La capacidad de ocultar estímulos
\end{multicols}
\end{enumerate}
\vspace{0.15cm}
\item ¿Qué es la reproducción?
\begin{enumerate}[label=\textbf{\Alph*)}]
\begin{multicols}{2}
  \item La capacidad de generar nuevos individuos
  \item La capacidad de generar minerales
  \item La capacidad de generar agua
  \item La capacidad de generar aire
  \item La capacidad de generar suelo
  \item La capacidad de generar astros
\end{multicols}
\end{enumerate}
\vspace{0.15cm}
\item ¿Qué es la reproducción asexual?
\begin{enumerate}[label=\textbf{\Alph*)}]
\begin{multicols}{2}
  \item La que produce descendientes genéticamente idénticos al progenitor
  \item La que produce descendientes genéticamente diferentes
  \item La que requiere dos progenitores
  \item La que requiere tres progenitores
  \item La que requiere cuatro progenitores
  \item La que requiere cinco progenitores
\end{multicols}
\end{enumerate}
\vspace{0.15cm}
\item ¿Qué es la reproducción sexual?
\begin{enumerate}[label=\textbf{\Alph*)}]
\begin{multicols}{2}
  \item La que requiere dos progenitores y produce descendientes diferentes
  \item La que requiere un progenitor
  \item La que produce descendientes idénticos
  \item La que no requiere progenitores
  \item La que requiere tres progenitores
  \item La que requiere cuatro progenitores
\end{multicols}
\end{enumerate}
\vspace{0.15cm}
\item ¿Qué es la adaptación?
\begin{enumerate}[label=\textbf{\Alph*)}]
\begin{multicols}{2}
  \item La capacidad de los seres vivos de ajustarse al ambiente
  \item La capacidad de los seres vivos de ignorar el ambiente
  \item La capacidad de los seres vivos de cambiar el ambiente
  \item La capacidad de los seres vivos de destruir el ambiente
  \item La capacidad de los seres vivos de olvidar el ambiente
  \item La capacidad de los seres vivos de negar el ambiente
\end{multicols}
\end{enumerate}
\vspace{0.15cm}
\item ¿Qué es la evolución?
\begin{enumerate}[label=\textbf{\Alph*)}]
\begin{multicols}{2}
  \item El cambio de las especies a lo largo del tiempo
  \item La permanencia de las especies
  \item La desaparición de las especies
  \item La creación de las especies
  \item La transformación de las especies en minerales
  \item La transformación de las especies en agua
\end{multicols}
\end{enumerate}
\vspace{0.15cm}
\item ¿Quién propuso la teoría de la evolución por selección natural?
\begin{enumerate}[label=\textbf{\Alph*)}]
\begin{multicols}{2}
  \item Charles Darwin
  \item Gregor Mendel
  \item Louis Pasteur
  \item Robert Hooke
  \item Alexander Fleming
  \item Isaac Newton
\end{multicols}
\end{enumerate}
\vspace{0.15cm}
\item ¿Qué es la selección natural?
\begin{enumerate}[label=\textbf{\Alph*)}]
\begin{multicols}{2}
  \item El proceso por el cual los individuos mejor adaptados sobreviven y se reproducen
  \item El proceso por el cual los individuos peor adaptados sobreviven
  \item El proceso por el cual los individuos no se reproducen
  \item El proceso por el cual los individuos no se adaptan
  \item El proceso por el cual los individuos no evolucionan
  \item El proceso por el cual los individuos no cambian
\end{multicols}
\end{enumerate}
\vspace{0.15cm}
\item ¿Qué es el ADN?
\begin{enumerate}[label=\textbf{\Alph*)}]
\begin{multicols}{2}
  \item El ácido desoxirribonucleico que contiene la información genética
  \item El ácido ribonucleico que contiene la información genética
  \item El ácido desoxirribonucleico que contiene la información proteica
  \item El ácido ribonucleico que contiene la información proteica
  \item El ácido desoxirribonucleico que contiene la información lipídica
  \item El ácido ribonucleico que contiene la información lipídica
\end{multicols}
\end{enumerate}
\vspace{0.15cm}
\item ¿Dónde se encuentra el ADN en las células eucariotas?
\begin{enumerate}[label=\textbf{\Alph*)}]
\begin{multicols}{2}
  \item En el núcleo
  \item En el citoplasma
  \item En la membrana
  \item En la pared celular
  \item En los ribosomas
  \item En las mitocondrias
\end{multicols}
\end{enumerate}
\vspace{0.15cm}
\item ¿Qué es el ARN?
\begin{enumerate}[label=\textbf{\Alph*)}]
\begin{multicols}{2}
  \item El ácido ribonucleico que participa en la síntesis de proteínas
  \item El ácido desoxirribonucleico que participa en la síntesis de proteínas
  \item El ácido ribonucleico que participa en la síntesis de lípidos
  \item El ácido desoxirribonucleico que participa en la síntesis de lípidos
  \item El ácido ribonucleico que participa en la síntesis de carbohidratos
  \item El ácido desoxirribonucleico que participa en la síntesis de carbohidratos
\end{multicols}
\end{enumerate}
\vspace{0.15cm}
\item ¿Qué es un gen?
\begin{enumerate}[label=\textbf{\Alph*)}]
\begin{multicols}{2}
  \item Un segmento de ADN que contiene información para un rasgo
  \item Un segmento de ARN que contiene información para un rasgo
  \item Un segmento de proteína que contiene información para un rasgo
  \item Un segmento de lípido que contiene información para un rasgo
  \item Un segmento de carbohidrato que contiene información para un rasgo
  \item Un segmento de mineral que contiene información para un rasgo
\end{multicols}
\end{enumerate}
\vspace{0.15cm}
\item ¿Qué es el genoma?
\begin{enumerate}[label=\textbf{\Alph*)}]
\begin{multicols}{2}
  \item El conjunto completo de genes de un organismo
  \item El conjunto completo de proteínas de un organismo
  \item El conjunto completo de lípidos de un organismo
  \item El conjunto completo de carbohidratos de un organismo
  \item El conjunto completo de minerales de un organismo
  \item El conjunto completo de células de un organismo
\end{multicols}
\end{enumerate}
\vspace{0.15cm}
\item ¿Qué es la nutrición?
\begin{enumerate}[label=\textbf{\Alph*)}]
\begin{multicols}{2}
  \item El proceso de obtener y utilizar nutrientes
  \item El proceso de obtener y utilizar minerales
  \item El proceso de obtener y utilizar gases
  \item El proceso de obtener y utilizar líquidos
  \item El proceso de obtener y utilizar astros
  \item El proceso de obtener y utilizar rocas
\end{multicols}
\end{enumerate}
\vspace{0.15cm}
\item ¿Qué es la nutrición autótrofa?
\begin{enumerate}[label=\textbf{\Alph*)}]
\begin{multicols}{2}
  \item La que produce su propio alimento
  \item La que consume alimento de otros
  \item La que no necesita alimento
  \item La que no produce alimento
  \item La que no consume alimento
  \item La que no utiliza alimento
\end{multicols}
\end{enumerate}
\vspace{0.15cm}
\item ¿Qué es la nutrición heterótrofa?
\begin{enumerate}[label=\textbf{\Alph*)}]
\begin{multicols}{2}
  \item La que consume alimento producido por otros
  \item La que produce su propio alimento
  \item La que no necesita alimento
  \item La que no produce alimento
  \item La que no consume alimento
  \item La que no utiliza alimento
\end{multicols}
\end{enumerate}
\vspace{0.15cm}
\item ¿Qué es la fotosíntesis?
\begin{enumerate}[label=\textbf{\Alph*)}]
\begin{multicols}{2}
  \item El proceso por el cual las plantas producen alimento usando luz solar
  \item El proceso por el cual las plantas producen alimento usando luz lunar
  \item El proceso por el cual las plantas producen alimento usando luz artificial
  \item El proceso por el cual las plantas producen alimento usando luz ultravioleta
  \item El proceso por el cual las plantas producen alimento usando luz infrarroja
  \item El proceso por el cual las plantas producen alimento usando luz gamma
\end{multicols}
\end{enumerate}
\vspace{0.15cm}
\item ¿Qué es la respiración celular?
\begin{enumerate}[label=\textbf{\Alph*)}]
\begin{multicols}{2}
  \item El proceso por el cual las células obtienen energía de los nutrientes
  \item El proceso por el cual las células obtienen energía de la luz
  \item El proceso por el cual las células obtienen energía del agua
  \item El proceso por el cual las células obtienen energía del aire
  \item El proceso por el cual las células obtienen energía del suelo
  \item El proceso por el cual las células obtienen energía de los minerales
\end{multicols}
\end{enumerate}
\vspace{0.15cm}
\item ¿Qué es la mitosis?
\begin{enumerate}[label=\textbf{\Alph*)}]
\begin{multicols}{2}
  \item El proceso de división celular que produce dos células idénticas
  \item El proceso de división celular que produce cuatro células diferentes
  \item El proceso de división celular que produce dos células diferentes
  \item El proceso de división celular que produce cuatro células idénticas
  \item El proceso de división celular que produce una célula
  \item El proceso de división celular que produce tres células
\end{multicols}
\end{enumerate}
\vspace{0.15cm}
\item ¿Qué es la meiosis?
\begin{enumerate}[label=\textbf{\Alph*)}]
\begin{multicols}{2}
  \item El proceso de división celular que produce cuatro células diferentes
  \item El proceso de división celular que produce dos células idénticas
  \item El proceso de división celular que produce dos células diferentes
  \item El proceso de división celular que produce cuatro células idénticas
  \item El proceso de división celular que produce una célula
  \item El proceso de división celular que produce tres células
\end{multicols}
\end{enumerate}
\vspace{0.15cm}
\item ¿Cuál es la función de la mitosis?
\begin{enumerate}[label=\textbf{\Alph*)}]
\begin{multicols}{2}
  \item El crecimiento y reparación de tejidos
  \item La producción de gametos
  \item La producción de células diferentes
  \item La producción de células sexuales
  \item La producción de células reproductoras
  \item La producción de células germinales
\end{multicols}
\end{enumerate}
\vspace{0.15cm}
\item ¿Cuál es la función de la meiosis?
\begin{enumerate}[label=\textbf{\Alph*)}]
\begin{multicols}{2}
  \item La producción de gametos
  \item El crecimiento y reparación de tejidos
  \item La producción de células idénticas
  \item La producción de células somáticas
  \item La producción de células corporales
  \item La producción de células diploides
\end{multicols}
\end{enumerate}
\vspace{0.15cm}
\item ¿Qué es la célula procariota?
\begin{enumerate}[label=\textbf{\Alph*)}]
\begin{multicols}{2}
  \item La que no tiene núcleo definido
  \item La que tiene núcleo definido
  \item La que tiene dos núcleos
  \item La que tiene tres núcleos
  \item La que tiene cuatro núcleos
  \item La que no tiene membrana
\end{multicols}
\end{enumerate}
\vspace{0.15cm}
\item ¿Qué es la célula eucariota?
\begin{enumerate}[label=\textbf{\Alph*)}]
\begin{multicols}{2}
  \item La que tiene núcleo definido
  \item La que no tiene núcleo definido
  \item La que tiene dos núcleos
  \item La que tiene tres núcleos
  \item La que tiene cuatro núcleos
  \item La que no tiene membrana
\end{multicols}
\end{enumerate}
\vspace{0.15cm}
\item ¿Cuál es un ejemplo de célula procariota?
\begin{enumerate}[label=\textbf{\Alph*)}]
\begin{multicols}{2}
  \item La bacteria
  \item La célula animal
  \item La célula vegetal
  \item La célula fúngica
  \item La célula protista
  \item La célula nerviosa
\end{multicols}
\end{enumerate}
\vspace{0.15cm}
\item ¿Cuál es un ejemplo de célula eucariota?
\begin{enumerate}[label=\textbf{\Alph*)}]
\begin{multicols}{2}
  \item La célula animal
  \item La bacteria
  \item La cianobacteria
  \item La arquea
  \item La bacteria intestinal
  \item La bacteria del suelo
\end{multicols}
\end{enumerate}
\vspace{0.15cm}
\item ¿Qué es la membrana celular?
\begin{enumerate}[label=\textbf{\Alph*)}]
\begin{multicols}{2}
  \item La estructura que delimita la célula y regula el paso de sustancias
  \item La estructura que no delimita la célula
  \item La estructura que no regula el paso de sustancias
  \item La estructura que delimita el núcleo
  \item La estructura que delimita la mitocondria
  \item La estructura que delimita el ribosoma
\end{multicols}
\end{enumerate}
\vspace{0.15cm}
\item ¿Qué es el citoplasma?
\begin{enumerate}[label=\textbf{\Alph*)}]
\begin{multicols}{2}
  \item El medio interno de la célula donde ocurren las reacciones
  \item El medio externo de la célula
  \item El medio interno del núcleo
  \item El medio externo del núcleo
  \item El medio interno de la mitocondria
  \item El medio externo de la mitocondria
\end{multicols}
\end{enumerate}
\vspace{0.15cm}
\item ¿Qué es el núcleo celular?
\begin{enumerate}[label=\textbf{\Alph*)}]
\begin{multicols}{2}
  \item La estructura que contiene el material genético
  \item La estructura que contiene el material proteico
  \item La estructura que contiene el material lipídico
  \item La estructura que contiene el material carbohidrático
  \item La estructura que contiene el material mineral
  \item La estructura que contiene el material gaseoso
\end{multicols}
\end{enumerate}
\vspace{0.15cm}
\item ¿Qué es la mitocondria?
\begin{enumerate}[label=\textbf{\Alph*)}]
\begin{multicols}{2}
  \item El orgánulo que produce energía
  \item El orgánulo que produce proteínas
  \item El orgánulo que produce lípidos
  \item El orgánulo que produce carbohidratos
  \item El orgánulo que produce minerales
  \item El orgánulo que produce gases
\end{multicols}
\end{enumerate}
\vspace{0.15cm}
\item ¿Qué es el ribosoma?
\begin{enumerate}[label=\textbf{\Alph*)}]
\begin{multicols}{2}
  \item El orgánulo que sintetiza proteínas
  \item El orgánulo que sintetiza lípidos
  \item El orgánulo que sintetiza carbohidratos
  \item El orgánulo que sintetiza minerales
  \item El orgánulo que sintetiza gases
  \item El orgánulo que sintetiza agua
\end{multicols}
\end{enumerate}
\vspace{0.15cm}
\item ¿Qué es el cloroplasto?
\begin{enumerate}[label=\textbf{\Alph*)}]
\begin{multicols}{2}
  \item El orgánulo que realiza la fotosíntesis
  \item El orgánulo que realiza la respiración
  \item El orgánulo que realiza la digestión
  \item El orgánulo que realiza la excreción
  \item El orgánulo que realiza la reproducción
  \item El orgánulo que realiza la síntesis
\end{multicols}
\end{enumerate}
\vspace{0.15cm}
\item ¿Qué es la pared celular?
\begin{enumerate}[label=\textbf{\Alph*)}]
\begin{multicols}{2}
  \item La estructura que rodea la membrana de células vegetales
  \item La estructura que rodea la membrana de células animales
  \item La estructura que rodea el núcleo
  \item La estructura que rodea la mitocondria
  \item La estructura que rodea el ribosoma
  \item La estructura que rodea el cloroplasto
\end{multicols}
\end{enumerate}
\vspace{0.15cm}
\item ¿Qué es el aparato de Golgi?
\begin{enumerate}[label=\textbf{\Alph*)}]
\begin{multicols}{2}
  \item El orgánulo que modifica y empaqueta proteínas
  \item El orgánulo que modifica y empaqueta lípidos
  \item El orgánulo que modifica y empaqueta carbohidratos
  \item El orgánulo que modifica y empaqueta minerales
  \item El orgánulo que modifica y empaqueta gases
  \item El orgánulo que modifica y empaqueta agua
\end{multicols}
\end{enumerate}
\vspace{0.15cm}
\item ¿Qué es el retículo endoplasmático?
\begin{enumerate}[label=\textbf{\Alph*)}]
\begin{multicols}{2}
  \item La red de membranas que transporta sustancias
  \item La red de membranas que no transporta sustancias
  \item La red de membranas que delimita el núcleo
  \item La red de membranas que delimita la mitocondria
  \item La red de membranas que delimita el ribosoma
  \item La red de membranas que delimita el cloroplasto
\end{multicols}
\end{enumerate}
\vspace{0.15cm}
\item ¿Qué es el lisosoma?
\begin{enumerate}[label=\textbf{\Alph*)}]
\begin{multicols}{2}
  \item El orgánulo que digiere sustancias
  \item El orgánulo que sintetiza sustancias
  \item El orgánulo que transporta sustancias
  \item El orgánulo que almacena sustancias
  \item El orgánulo que excreta sustancias
  \item El orgánulo que absorbe sustancias
\end{multicols}
\end{enumerate}
\vspace{0.15cm}
\item ¿Qué es el citoesqueleto?
\begin{enumerate}[label=\textbf{\Alph*)}]
\begin{multicols}{2}
  \item La red de proteínas que da forma y movimiento a la célula
  \item La red de lípidos que da forma a la célula
  \item La red de carbohidratos que da forma a la célula
  \item La red de minerales que da forma a la célula
  \item La red de gases que da forma a la célula
  \item La red de agua que da forma a la célula
\end{multicols}
\end{enumerate}
\vspace{0.15cm}
\item ¿Qué es la vacuola?
\begin{enumerate}[label=\textbf{\Alph*)}]
\begin{multicols}{2}
  \item El orgánulo que almacena sustancias
  \item El orgánulo que sintetiza sustancias
  \item El orgánulo que transporta sustancias
  \item El orgánulo que digiere sustancias
  \item El orgánulo que excreta sustancias
  \item El orgánulo que absorbe sustancias
\end{multicols}
\end{enumerate}
\vspace{0.15cm}
\item ¿Qué es el centriolo?
\begin{enumerate}[label=\textbf{\Alph*)}]
\begin{multicols}{2}
  \item El orgánulo que participa en la división celular
  \item El orgánulo que participa en la síntesis de proteínas
  \item El orgánulo que participa en la síntesis de lípidos
  \item El orgánulo que participa en la fotosíntesis
  \item El orgánulo que participa en la respiración
  \item El orgánulo que participa en la digestión ---
\end{multicols}
\end{enumerate}
\vspace{0.15cm}
\end{enumerate}\n
\vspace{0.4cm}
\begin{tcolorbox}[colback=gray!5!white,colframe=gray!60!black,title=\textbf{SOLUCIONARIO -- 5.1 Biología y Seres Vivos}]
\begin{multicols}{4}
\begin{enumerate}[label=\textbf{\arabic*.}]
  \item \textbf{A}
  \item \textbf{A}
  \item \textbf{A}
  \item \textbf{A}
  \item \textbf{A}
  \item \textbf{A}
  \item \textbf{A}
  \item \textbf{A}
  \item \textbf{A}
  \item \textbf{A}
  \item \textbf{A}
  \item \textbf{A}
  \item \textbf{A}
  \item \textbf{A}
  \item \textbf{A}
  \item \textbf{A}
  \item \textbf{A}
  \item \textbf{A}
  \item \textbf{A}
  \item \textbf{A}
  \item \textbf{A}
  \item \textbf{A}
  \item \textbf{A}
  \item \textbf{A}
  \item \textbf{A}
  \item \textbf{A}
  \item \textbf{A}
  \item \textbf{A}
  \item \textbf{A}
  \item \textbf{A}
  \item \textbf{A}
  \item \textbf{A}
  \item \textbf{A}
  \item \textbf{A}
  \item \textbf{A}
  \item \textbf{A}
  \item \textbf{A}
  \item \textbf{A}
  \item \textbf{A}
  \item \textbf{A}
  \item \textbf{A}
  \item \textbf{A}
  \item \textbf{A}
  \item \textbf{A}
  \item \textbf{A}
  \item \textbf{A}
  \item \textbf{A}
  \item \textbf{A}
  \item \textbf{A}
  \item \textbf{A}
  \item \textbf{A}
  \item \textbf{A}
  \item \textbf{A}
  \item \textbf{A}
  \item \textbf{A}
  \item \textbf{A}
  \item \textbf{A}
  \item \textbf{A}
  \item \textbf{A}
  \item \textbf{A}
  \item \textbf{A}
  \item \textbf{A}
  \item \textbf{A}
  \item \textbf{A}
  \item \textbf{A}
  \item \textbf{A}
  \item \textbf{A}
  \item \textbf{A}
  \item \textbf{A}
  \item \textbf{A}
  \item \textbf{A}
  \item \textbf{A}
  \item \textbf{A}
  \item \textbf{A}
  \item \textbf{A}
  \item \textbf{A}
  \item \textbf{A}
  \item \textbf{A}
  \item \textbf{A}
  \item \textbf{A}
  \item \textbf{A}
  \item \textbf{A}
  \item \textbf{A}
  \item \textbf{A}
  \item \textbf{A}
  \item \textbf{A}
  \item \textbf{A}
  \item \textbf{A}
  \item \textbf{A}
  \item \textbf{A}
  \item \textbf{A}
  \item \textbf{A}
  \item \textbf{A}
  \item \textbf{A}
  \item \textbf{A}
  \item \textbf{A}
  \item \textbf{A}
  \item \textbf{A}
  \item \textbf{A}
  \item \textbf{A}
  \item \textbf{A}
  \item \textbf{A}
  \item \textbf{A}
  \item \textbf{A}
  \item \textbf{A}
  \item \textbf{A}
  \item \textbf{A}
  \item \textbf{A}
  \item \textbf{A}
  \item \textbf{A}
  \item \textbf{A}
  \item \textbf{A}
  \item \textbf{A}
  \item \textbf{A}
  \item \textbf{A}
  \item \textbf{A}
  \item \textbf{A}
  \item \textbf{A}
  \item \textbf{A}
  \item \textbf{A}
  \item \textbf{A}
  \item \textbf{A}
  \item \textbf{A}
  \item \textbf{A}
  \item \textbf{A}
  \item \textbf{A}
  \item \textbf{A}
  \item \textbf{A}
  \item \textbf{A}
  \item \textbf{A}
  \item \textbf{A}
  \item \textbf{A}
  \item \textbf{A}
  \item \textbf{A}
  \item \textbf{A}
  \item \textbf{A}
  \item \textbf{A}
  \item \textbf{A}
  \item \textbf{A}
  \item \textbf{A}
  \item \textbf{A}
  \item \textbf{A}
  \item \textbf{A}
  \item \textbf{A}
  \item \textbf{A}
  \item \textbf{A}
  \item \textbf{A}
  \item \textbf{A}
  \item \textbf{A}
  \item \textbf{A}
  \item \textbf{A}
  \item \textbf{A}
  \item \textbf{A}
  \item \textbf{A}
  \item \textbf{A}
  \item \textbf{A}
  \item \textbf{A}
  \item \textbf{A}
  \item \textbf{A}
  \item \textbf{A}
  \item \textbf{A}
  \item \textbf{A}
  \item \textbf{A}
  \item \textbf{A}
  \item \textbf{A}
  \item \textbf{A}
  \item \textbf{A}
  \item \textbf{A}
  \item \textbf{A}
  \item \textbf{A}
  \item \textbf{A}
  \item \textbf{A}
  \item \textbf{A}
  \item \textbf{A}
  \item \textbf{A}
  \item \textbf{A}
  \item \textbf{A}
  \item \textbf{A}
  \item \textbf{A}
  \item \textbf{A}
  \item \textbf{A}
  \item \textbf{A}
  \item \textbf{A}
  \item \textbf{A}
  \item \textbf{A}
  \item \textbf{A}
  \item \textbf{A}
  \item \textbf{A}
  \item \textbf{A}
  \item \textbf{A}
  \item \textbf{A}
  \item \textbf{A}
  \item \textbf{A}
  \item \textbf{A}
  \item \textbf{A}
  \item \textbf{A}
  \item \textbf{A}
  \item \textbf{A}
  \item \textbf{A}
  \item \textbf{A}
  \item \textbf{A}
  \item \textbf{A}
  \item \textbf{A}
  \item \textbf{A}
  \item \textbf{A}
  \item \textbf{A}
  \item \textbf{A}
  \item \textbf{A}
  \item \textbf{A}
  \item \textbf{A}
  \item \textbf{A}
  \item \textbf{A}
  \item \textbf{A}
  \item \textbf{A}
  \item \textbf{A}
  \item \textbf{A}
  \item \textbf{A}
  \item \textbf{A}
  \item \textbf{A}
  \item \textbf{A}
  \item \textbf{A}
  \item \textbf{A}
  \item \textbf{A}
  \item \textbf{A}
  \item \textbf{A}
  \item \textbf{A}
  \item \textbf{A}
  \item \textbf{A}
  \item \textbf{A}
  \item \textbf{A}
  \item \textbf{A}
  \item \textbf{A}
  \item \textbf{A}
  \item \textbf{A}
  \item \textbf{A}
  \item \textbf{A}
  \item \textbf{A}
  \item \textbf{A}
  \item \textbf{A}
  \item \textbf{A}
  \item \textbf{A}
  \item \textbf{A}
  \item \textbf{A}
  \item \textbf{A}
  \item \textbf{A}
  \item \textbf{A}
  \item \textbf{A}
  \item \textbf{A}
  \item \textbf{A}
  \item \textbf{A}
  \item \textbf{A}
  \item \textbf{A}
  \item \textbf{A}
  \item \textbf{A}
  \item \textbf{A}
  \item \textbf{A}
  \item \textbf{A}
  \item \textbf{A}
  \item \textbf{A}
  \item \textbf{A}
  \item \textbf{A}
  \item \textbf{A}
  \item \textbf{A}
  \item \textbf{A}
  \item \textbf{A}
  \item \textbf{A}
  \item \textbf{A}
  \item \textbf{A}
  \item \textbf{A}
  \item \textbf{A}
  \item \textbf{A}
  \item \textbf{A}
  \item \textbf{A}
  \item \textbf{A}
  \item \textbf{A}
  \item \textbf{A}
  \item \textbf{A}
  \item \textbf{A}
  \item \textbf{A}
  \item \textbf{A}
  \item \textbf{A}
  \item \textbf{A}
  \item \textbf{A}
  \item \textbf{A}
  \item \textbf{A}
  \item \textbf{A}
  \item \textbf{A}
  \item \textbf{A}
  \item \textbf{A}
  \item \textbf{A}
  \item \textbf{A}
  \item \textbf{A}
  \item \textbf{A}
  \item \textbf{A}
  \item \textbf{A}
  \item \textbf{A}
  \item \textbf{A}
  \item \textbf{A}
  \item \textbf{A}
  \item \textbf{A}
\end{enumerate}
\end{multicols}
\end{tcolorbox}
\newpage\n
\subsection{5.2 Química General y Reacciones}

\subsubsection*{SECCIÓN I: PROPIEDADES DE LA MATERIA (1–50)}

\begin{enumerate}[label=\textbf{\arabic*.}]
\setcounter{enumi}{0}
\item ¿Qué es la materia?
\begin{enumerate}[label=\textbf{\Alph*)}]
\begin{multicols}{2}
  \item Todo aquello que tiene masa y ocupa un volumen
  \item Todo aquello que tiene energía y ocupa un espacio
  \item Todo aquello que tiene peso y ocupa un lugar
  \item Todo aquello que tiene color y ocupa un área
  \item Todo aquello que tiene temperatura y ocupa una región
  \item Todo aquello que tiene densidad y ocupa una superficie
\end{multicols}
\end{enumerate}
\vspace{0.15cm}
\item ¿Cuál es una propiedad general de la materia?
\begin{enumerate}[label=\textbf{\Alph*)}]
\begin{multicols}{2}
  \item La masa
  \item El color
  \item El sabor
  \item El olor
  \item La textura
  \item La forma
\end{multicols}
\end{enumerate}
\vspace{0.15cm}
\item ¿Cuál es una propiedad específica de la materia?
\begin{enumerate}[label=\textbf{\Alph*)}]
\begin{multicols}{2}
  \item La densidad
  \item La masa
  \item El volumen
  \item La longitud
  \item El peso
  \item La inercia
\end{multicols}
\end{enumerate}
\vspace{0.15cm}
\item ¿Qué es una propiedad física?
\begin{enumerate}[label=\textbf{\Alph*)}]
\begin{multicols}{2}
  \item La que se puede observar sin cambiar la composición de la sustancia
  \item La que cambia la composición de la sustancia
  \item La que se observa solo con reacciones químicas
  \item La que se observa solo con electricidad
  \item La que se observa solo con calor
  \item La que se observa solo con luz
\end{multicols}
\end{enumerate}
\vspace{0.15cm}
\item ¿Qué es una propiedad química?
\begin{enumerate}[label=\textbf{\Alph*)}]
\begin{multicols}{2}
  \item La que describe la capacidad de una sustancia para cambiar su composición
  \item La que se observa sin cambiar la composición
  \item La que se observa solo con cambios físicos
  \item La que se observa solo con cambios de estado
  \item La que se observa solo con cambios de forma
  \item La que se observa solo con cambios de tamaño
\end{multicols}
\end{enumerate}
\vspace{0.15cm}
\item ¿Cuál es un ejemplo de propiedad física?
\begin{enumerate}[label=\textbf{\Alph*)}]
\begin{multicols}{2}
  \item El punto de fusión
  \item La combustión
  \item La oxidación
  \item La fermentación
  \item La corrosión
  \item La descomposición
\end{multicols}
\end{enumerate}
\vspace{0.15cm}
\item ¿Cuál es un ejemplo de propiedad química?
\begin{enumerate}[label=\textbf{\Alph*)}]
\begin{multicols}{2}
  \item La combustión
  \item El punto de ebullición
  \item La densidad
  \item La solubilidad
  \item La conductividad
  \item La ductilidad
\end{multicols}
\end{enumerate}
\vspace{0.15cm}
\item ¿Qué es un cambio físico?
\begin{enumerate}[label=\textbf{\Alph*)}]
\begin{multicols}{2}
  \item El que no altera la composición de la sustancia
  \item El que altera la composición de la sustancia
  \item El que forma nuevas sustancias
  \item El que produce reacciones químicas
  \item El que descompone la materia
  \item El que sintetiza compuestos
\end{multicols}
\end{enumerate}
\vspace{0.15cm}
\item ¿Qué es un cambio químico?
\begin{enumerate}[label=\textbf{\Alph*)}]
\begin{multicols}{2}
  \item El que altera la composición de la sustancia
  \item El que no altera la composición
  \item El que solo cambia de estado
  \item El que solo cambia de forma
  \item El que solo cambia de tamaño
  \item El que solo cambia de color
\end{multicols}
\end{enumerate}
\vspace{0.15cm}
\item ¿Cuál es un ejemplo de cambio físico?
\begin{enumerate}[label=\textbf{\Alph*)}]
\begin{multicols}{2}
  \item La fusión del hielo
  \item La combustión del papel
  \item La oxidación del hierro
  \item La fermentación del azúcar
  \item La corrosión del metal
  \item La descomposición del agua
\end{multicols}
\end{enumerate}
\vspace{0.15cm}
\item ¿Cuál es un ejemplo de cambio químico?
\begin{enumerate}[label=\textbf{\Alph*)}]
\begin{multicols}{2}
  \item La oxidación del hierro
  \item La fusión del hielo
  \item La evaporación del agua
  \item La sublimación del yodo
  \item La condensación del vapor
  \item La solidificación del agua
\end{multicols}
\end{enumerate}
\vspace{0.15cm}
\item ¿Qué es la masa?
\begin{enumerate}[label=\textbf{\Alph*)}]
\begin{multicols}{2}
  \item La cantidad de materia de un cuerpo
  \item La cantidad de energía de un cuerpo
  \item La cantidad de volumen de un cuerpo
  \item La cantidad de peso de un cuerpo
  \item La cantidad de densidad de un cuerpo
  \item La cantidad de temperatura de un cuerpo
\end{multicols}
\end{enumerate}
\vspace{0.15cm}
\item ¿Qué es el volumen?
\begin{enumerate}[label=\textbf{\Alph*)}]
\begin{multicols}{2}
  \item El espacio que ocupa un cuerpo
  \item La cantidad de materia de un cuerpo
  \item La cantidad de energía de un cuerpo
  \item La cantidad de peso de un cuerpo
  \item La cantidad de densidad de un cuerpo
  \item La cantidad de temperatura de un cuerpo
\end{multicols}
\end{enumerate}
\vspace{0.15cm}
\item ¿Qué es la densidad?
\begin{enumerate}[label=\textbf{\Alph*)}]
\begin{multicols}{2}
  \item La relación entre la masa y el volumen
  \item La relación entre el volumen y la masa
  \item La relación entre el peso y la masa
  \item La relación entre la masa y el peso
  \item La relación entre el volumen y el peso
  \item La relación entre el peso y el volumen
\end{multicols}
\end{enumerate}
\vspace{0.15cm}
\item ¿Cuál es la unidad de masa en el Sistema Internacional?
\begin{enumerate}[label=\textbf{\Alph*)}]
\begin{multicols}{2}
  \item El kilogramo (kg)
  \item El gramo (g)
  \item La libra (lb)
  \item La onza (oz)
  \item La tonelada (t)
  \item El quintal (q)
\end{multicols}
\end{enumerate}
\vspace{0.15cm}
\item ¿Cuál es la unidad de volumen en el Sistema Internacional?
\begin{enumerate}[label=\textbf{\Alph*)}]
\begin{multicols}{2}
  \item El metro cúbico (m³)
  \item El litro (L)
  \item El centímetro cúbico (cm³)
  \item El mililitro (mL)
  \item El galón (gal)
  \item La pinta (pt)
\end{multicols}
\end{enumerate}
\vspace{0.15cm}
\item ¿Cuál es la unidad de densidad en el Sistema Internacional?
\begin{enumerate}[label=\textbf{\Alph*)}]
\begin{multicols}{2}
  \item kg/m³
  \item g/cm³
  \item kg/L
  \item g/mL
  \item lb/ft³
  \item oz/in³
\end{multicols}
\end{enumerate}
\vspace{0.15cm}
\item ¿Qué es la temperatura?
\begin{enumerate}[label=\textbf{\Alph*)}]
\begin{multicols}{2}
  \item La medida de la energía térmica de un cuerpo
  \item La medida de la energía mecánica
  \item La medida de la energía eléctrica
  \item La medida de la energía química
  \item La medida de la energía nuclear
  \item La medida de la energía radiante
\end{multicols}
\end{enumerate}
\vspace{0.15cm}
\item ¿Cuál es la unidad de temperatura en el Sistema Internacional?
\begin{enumerate}[label=\textbf{\Alph*)}]
\begin{multicols}{2}
  \item El kelvin (K)
  \item El grado Celsius (°C)
  \item El grado Fahrenheit (°F)
  \item El grado Rankine (°R)
  \item El grado Réaumur (°Ré)
  \item El grado Newton (°N)
\end{multicols}
\end{enumerate}
\vspace{0.15cm}
\item ¿Qué es la presión?
\begin{enumerate}[label=\textbf{\Alph*)}]
\begin{multicols}{2}
  \item La fuerza por unidad de área
  \item La masa por unidad de volumen
  \item El volumen por unidad de masa
  \item La energía por unidad de tiempo
  \item El trabajo por unidad de tiempo
  \item La fuerza por unidad de masa
\end{multicols}
\end{enumerate}
\vspace{0.15cm}
\item ¿Cuál es la unidad de presión en el Sistema Internacional?
\begin{enumerate}[label=\textbf{\Alph*)}]
\begin{multicols}{2}
  \item El pascal (Pa)
  \item La atmósfera (atm)
  \item El bar (bar)
  \item El milímetro de mercurio (mmHg)
  \item La libra por pulgada cuadrada (psi)
  \item El torr (Torr)
\end{multicols}
\end{enumerate}
\vspace{0.15cm}
\item ¿Qué es un elemento químico?
\begin{enumerate}[label=\textbf{\Alph*)}]
\begin{multicols}{2}
  \item Una sustancia pura formada por átomos del mismo tipo
  \item Una sustancia pura formada por átomos de diferente tipo
  \item Una mezcla de sustancias
  \item Una combinación de compuestos
  \item Una disolución de sustancias
  \item Una suspensión de partículas
\end{multicols}
\end{enumerate}
\vspace{0.15cm}
\item ¿Qué es un compuesto químico?
\begin{enumerate}[label=\textbf{\Alph*)}]
\begin{multicols}{2}
  \item Una sustancia pura formada por átomos de diferente tipo
  \item Una sustancia pura formada por átomos del mismo tipo
  \item Una mezcla de sustancias
  \item Una combinación de elementos
  \item Una disolución de sustancias
  \item Una suspensión de partículas
\end{multicols}
\end{enumerate}
\vspace{0.15cm}
\item ¿Qué es una mezcla?
\begin{enumerate}[label=\textbf{\Alph*)}]
\begin{multicols}{2}
  \item La unión de dos o más sustancias sin combinación química
  \item La unión de dos o más sustancias con combinación química
  \item La unión de dos o más elementos
  \item La unión de dos o más compuestos
  \item La unión de dos o más átomos
  \item La unión de dos o más moléculas
\end{multicols}
\end{enumerate}
\vspace{0.15cm}
\item ¿Qué es una mezcla homogénea?
\begin{enumerate}[label=\textbf{\Alph*)}]
\begin{multicols}{2}
  \item La que tiene composición uniforme
  \item La que tiene composición no uniforme
  \item La que tiene composición variable
  \item La que tiene composición definida
  \item La que tiene composición indefinida
  \item La que tiene composición constante
\end{multicols}
\end{enumerate}
\vspace{0.15cm}
\item ¿Qué es una mezcla heterogénea?
\begin{enumerate}[label=\textbf{\Alph*)}]
\begin{multicols}{2}
  \item La que tiene composición no uniforme
  \item La que tiene composición uniforme
  \item La que tiene composición variable
  \item La que tiene composición definida
  \item La que tiene composición indefinida
  \item La que tiene composición constante
\end{multicols}
\end{enumerate}
\vspace{0.15cm}
\item ¿Cuál es un ejemplo de mezcla homogénea?
\begin{enumerate}[label=\textbf{\Alph*)}]
\begin{multicols}{2}
  \item El agua salada
  \item El agua con aceite
  \item La arena con hierro
  \item El granito
  \item La madera
  \item El hormigón
\end{multicols}
\end{enumerate}
\vspace{0.15cm}
\item ¿Cuál es un ejemplo de mezcla heterogénea?
\begin{enumerate}[label=\textbf{\Alph*)}]
\begin{multicols}{2}
  \item El agua con aceite
  \item El agua salada
  \item El aire
  \item El acero
  \item El bronce
  \item El alcohol
\end{multicols}
\end{enumerate}
\vspace{0.15cm}
\item ¿Qué es una disolución?
\begin{enumerate}[label=\textbf{\Alph*)}]
\begin{multicols}{2}
  \item Una mezcla homogénea de soluto y disolvente
  \item Una mezcla heterogénea de soluto y disolvente
  \item Una combinación química de soluto y disolvente
  \item Una suspensión de soluto y disolvente
  \item Una emulsión de soluto y disolvente
  \item Una coloide de soluto y disolvente
\end{multicols}
\end{enumerate}
\vspace{0.15cm}
\item ¿Qué es el soluto?
\begin{enumerate}[label=\textbf{\Alph*)}]
\begin{multicols}{2}
  \item La sustancia que se disuelve
  \item La sustancia que disuelve
  \item La sustancia que se evapora
  \item La sustancia que se condensa
  \item La sustancia que se solidifica
  \item La sustancia que se sublima
\end{multicols}
\end{enumerate}
\vspace{0.15cm}
\item ¿Qué es el disolvente?
\begin{enumerate}[label=\textbf{\Alph*)}]
\begin{multicols}{2}
  \item La sustancia que disuelve
  \item La sustancia que se disuelve
  \item La sustancia que se evapora
  \item La sustancia que se condensa
  \item La sustancia que se solidifica
  \item La sustancia que se sublima
\end{multicols}
\end{enumerate}
\vspace{0.15cm}
\item ¿Qué es la solubilidad?
\begin{enumerate}[label=\textbf{\Alph*)}]
\begin{multicols}{2}
  \item La máxima cantidad de soluto que se puede disolver en un disolvente
  \item La mínima cantidad de soluto que se puede disolver
  \item La cantidad de disolvente necesaria
  \item La cantidad de soluto necesaria
  \item La cantidad de disolución obtenida
  \item La cantidad de precipitado formado
\end{multicols}
\end{enumerate}
\vspace{0.15cm}
\item ¿Qué es una disolución saturada?
\begin{enumerate}[label=\textbf{\Alph*)}]
\begin{multicols}{2}
  \item La que contiene la máxima cantidad de soluto disuelto
  \item La que contiene menos soluto del que puede disolver
  \item La que contiene más soluto del que puede disolver
  \item La que no contiene soluto
  \item La que no contiene disolvente
  \item La que contiene solo disolvente
\end{multicols}
\end{enumerate}
\vspace{0.15cm}
\item ¿Qué es una disolución insaturada?
\begin{enumerate}[label=\textbf{\Alph*)}]
\begin{multicols}{2}
  \item La que contiene menos soluto del que puede disolver
  \item La que contiene la máxima cantidad de soluto
  \item La que contiene más soluto del que puede disolver
  \item La que no contiene soluto
  \item La que no contiene disolvente
  \item La que contiene solo disolvente
\end{multicols}
\end{enumerate}
\vspace{0.15cm}
\item ¿Qué es una disolución sobresaturada?
\begin{enumerate}[label=\textbf{\Alph*)}]
\begin{multicols}{2}
  \item La que contiene más soluto del que puede disolver
  \item La que contiene menos soluto del que puede disolver
  \item La que contiene la máxima cantidad de soluto
  \item La que no contiene soluto
  \item La que no contiene disolvente
  \item La que contiene solo disolvente
\end{multicols}
\end{enumerate}
\vspace{0.15cm}
\item ¿Qué es la concentración de una disolución?
\begin{enumerate}[label=\textbf{\Alph*)}]
\begin{multicols}{2}
  \item La cantidad de soluto en una cantidad de disolución
  \item La cantidad de disolvente en una cantidad de disolución
  \item La cantidad de disolución en una cantidad de soluto
  \item La cantidad de disolución en una cantidad de disolvente
  \item La cantidad de soluto en una cantidad de disolvente
  \item La cantidad de disolvente en una cantidad de soluto
\end{multicols}
\end{enumerate}
\vspace{0.15cm}
\item ¿Cuál es una unidad de concentración?
\begin{enumerate}[label=\textbf{\Alph*)}]
\begin{multicols}{2}
  \item Molaridad (M)
  \item Kilogramo (kg)
  \item Metro (m)
  \item Segundo (s)
  \item Kelvin (K)
  \item Pascal (Pa)
\end{multicols}
\end{enumerate}
\vspace{0.15cm}
\item ¿Qué es la molaridad?
\begin{enumerate}[label=\textbf{\Alph*)}]
\begin{multicols}{2}
  \item Moles de soluto por litro de disolución
  \item Moles de soluto por kilogramo de disolvente
  \item Gramos de soluto por litro de disolución
  \item Gramos de soluto por kilogramo de disolvente
  \item Moles de disolvente por litro de disolución
  \item Gramos de disolvente por litro de disolución
\end{multicols}
\end{enumerate}
\vspace{0.15cm}
\item ¿Qué es la molalidad?
\begin{enumerate}[label=\textbf{\Alph*)}]
\begin{multicols}{2}
  \item Moles de soluto por kilogramo de disolvente
  \item Moles de soluto por litro de disolución
  \item Gramos de soluto por litro de disolución
  \item Gramos de soluto por kilogramo de disolvente
  \item Moles de disolvente por litro de disolución
  \item Gramos de disolvente por litro de disolución
\end{multicols}
\end{enumerate}
\vspace{0.15cm}
\item ¿Qué es el porcentaje en masa?
\begin{enumerate}[label=\textbf{\Alph*)}]
\begin{multicols}{2}
  \item Gramos de soluto por 100 gramos de disolución
  \item Gramos de soluto por 100 mL de disolución
  \item Moles de soluto por 100 gramos de disolución
  \item Moles de soluto por 100 mL de disolución
  \item Gramos de disolvente por 100 gramos de disolución
  \item Moles de disolvente por 100 mL de disolución
\end{multicols}
\end{enumerate}
\vspace{0.15cm}
\item ¿Qué es el porcentaje en volumen?
\begin{enumerate}[label=\textbf{\Alph*)}]
\begin{multicols}{2}
  \item mL de soluto por 100 mL de disolución
  \item mL de soluto por 100 gramos de disolución
  \item Gramos de soluto por 100 mL de disolución
  \item Gramos de soluto por 100 gramos de disolución
  \item mL de disolvente por 100 mL de disolución
  \item Gramos de disolvente por 100 mL de disolución
\end{multicols}
\end{enumerate}
\vspace{0.15cm}
\item ¿Qué es la parte por millón (ppm)?
\begin{enumerate}[label=\textbf{\Alph*)}]
\begin{multicols}{2}
  \item mg de soluto por kg de disolución
  \item g de soluto por kg de disolución
  \item mg de soluto por L de disolución
  \item g de soluto por L de disolución
  \item kg de soluto por kg de disolución
  \item L de soluto por L de disolución
\end{multicols}
\end{enumerate}
\vspace{0.15cm}
\item ¿Qué es un coloide?
\begin{enumerate}[label=\textbf{\Alph*)}]
\begin{multicols}{2}
  \item Una mezcla con partículas dispersas de tamaño intermedio
  \item Una mezcla homogénea verdadera
  \item Una mezcla heterogénea con partículas grandes
  \item Una combinación química
  \item Una sustancia pura
  \item Un elemento químico
\end{multicols}
\end{enumerate}
\vspace{0.15cm}
\item ¿Cuál es un ejemplo de coloide?
\begin{enumerate}[label=\textbf{\Alph*)}]
\begin{multicols}{2}
  \item La leche
  \item El agua salada
  \item El aire
  \item El acero
  \item El bronce
  \item El alcohol
\end{multicols}
\end{enumerate}
\vspace{0.15cm}
\item ¿Qué es una suspensión?
\begin{enumerate}[label=\textbf{\Alph*)}]
\begin{multicols}{2}
  \item Una mezcla heterogénea con partículas grandes que sedimentan
  \item Una mezcla homogénea verdadera
  \item Una mezcla con partículas dispersas de tamaño intermedio
  \item Una combinación química
  \item Una sustancia pura
  \item Un elemento químico
\end{multicols}
\end{enumerate}
\vspace{0.15cm}
\item ¿Cuál es un ejemplo de suspensión?
\begin{enumerate}[label=\textbf{\Alph*)}]
\begin{multicols}{2}
  \item El agua con arena
  \item El agua salada
  \item La leche
  \item El aire
  \item El acero
  \item El alcohol
\end{multicols}
\end{enumerate}
\vspace{0.15cm}
\item ¿Qué es la filtración?
\begin{enumerate}[label=\textbf{\Alph*)}]
\begin{multicols}{2}
  \item Un método de separación de mezclas heterogéneas
  \item Un método de separación de mezclas homogéneas
  \item Un método de separación de compuestos
  \item Un método de separación de elementos
  \item Un método de separación de átomos
  \item Un método de separación de moléculas
\end{multicols}
\end{enumerate}
\vspace{0.15cm}
\item ¿Qué es la destilación?
\begin{enumerate}[label=\textbf{\Alph*)}]
\begin{multicols}{2}
  \item Un método de separación basado en puntos de ebullición
  \item Un método de separación basado en densidades
  \item Un método de separación basado en solubilidades
  \item Un método de separación basado en tamaños
  \item Un método de separación basado en magnetismo
  \item Un método de separación basado en filtración
\end{multicols}
\end{enumerate}
\vspace{0.15cm}
\item ¿Qué es la cromatografía?
\begin{enumerate}[label=\textbf{\Alph*)}]
\begin{multicols}{2}
  \item Un método de separación basado en la velocidad de migración
  \item Un método de separación basado en puntos de ebullición
  \item Un método de separación basado en densidades
  \item Un método de separación basado en solubilidades
  \item Un método de separación basado en tamaños
  \item Un método de separación basado en magnetismo
\end{multicols}
\end{enumerate}
\vspace{0.15cm}
\item ¿Qué es la decantación?
\begin{enumerate}[label=\textbf{\Alph*)}]
\begin{multicols}{2}
  \item Un método de separación de líquidos inmiscibles o sólidos sedimentados
  \item Un método de separación de líquidos miscibles
  \item Un método de separación de gases
  \item Un método de separación de sólidos disueltos
  \item Un método de separación de coloides
  \item Un método de separación de suspensiones ---
\end{multicols}
\end{enumerate}
\vspace{0.15cm}
\end{enumerate}\n
\subsubsection*{SECCIÓN II: ESTRUCTURA ATÓMICA Y ENLACES (51–120)}

\begin{enumerate}[label=\textbf{\arabic*.}]
\setcounter{enumi}{50}
\item ¿Qué es un átomo?
\begin{enumerate}[label=\textbf{\Alph*)}]
\begin{multicols}{2}
  \item La unidad básica de la materia
  \item La unidad básica de la energía
  \item La unidad básica del volumen
  \item La unidad básica de la densidad
  \item La unidad básica de la temperatura
  \item La unidad básica de la presión
\end{multicols}
\end{enumerate}
\vspace{0.15cm}
\item ¿Cuáles son las partículas subatómicas principales?
\begin{enumerate}[label=\textbf{\Alph*)}]
\begin{multicols}{2}
  \item Protones, neutrones y electrones
  \item Protones, neutrones y fotones
  \item Protones, electrones y quarks
  \item Neutrones, electrones y leptones
  \item Protones, neutrones y bosones
  \item Electrones, fotones y gluones
\end{multicols}
\end{enumerate}
\vspace{0.15cm}
\item ¿Cuál es la carga del protón?
\begin{enumerate}[label=\textbf{\Alph*)}]
\begin{multicols}{2}
  \item Positiva
  \item Negativa
  \item Neutra
  \item Variable
  \item Nula
  \item Indefinida
\end{multicols}
\end{enumerate}
\vspace{0.15cm}
\item ¿Cuál es la carga del neutrón?
\begin{enumerate}[label=\textbf{\Alph*)}]
\begin{multicols}{2}
  \item Neutra
  \item Positiva
  \item Negativa
  \item Variable
  \item Nula
  \item Indefinida
\end{multicols}
\end{enumerate}
\vspace{0.15cm}
\item ¿Cuál es la carga del electrón?
\begin{enumerate}[label=\textbf{\Alph*)}]
\begin{multicols}{2}
  \item Negativa
  \item Positiva
  \item Neutra
  \item Variable
  \item Nula
  \item Indefinida
\end{multicols}
\end{enumerate}
\vspace{0.15cm}
\item ¿Dónde se encuentran los protones y neutrones?
\begin{enumerate}[label=\textbf{\Alph*)}]
\begin{multicols}{2}
  \item En el núcleo
  \item En la corteza
  \item En los orbitales
  \item En los enlaces
  \item En las moléculas
  \item En los iones
\end{multicols}
\end{enumerate}
\vspace{0.15cm}
\item ¿Dónde se encuentran los electrones?
\begin{enumerate}[label=\textbf{\Alph*)}]
\begin{multicols}{2}
  \item En la corteza, en orbitales
  \item En el núcleo
  \item En los protones
  \item En los neutrones
  \item En las moléculas
  \item En los enlaces
\end{multicols}
\end{enumerate}
\vspace{0.15cm}
\item ¿Qué es el número atómico?
\begin{enumerate}[label=\textbf{\Alph*)}]
\begin{multicols}{2}
  \item El número de protones en el núcleo
  \item El número de neutrones en el núcleo
  \item El número de electrones en la corteza
  \item El número de protones más neutrones
  \item El número de protones menos neutrones
  \item El número de electrones más protones
\end{multicols}
\end{enumerate}
\vspace{0.15cm}
\item ¿Qué es el número de masa?
\begin{enumerate}[label=\textbf{\Alph*)}]
\begin{multicols}{2}
  \item La suma de protones y neutrones
  \item La suma de protones y electrones
  \item La suma de neutrones y electrones
  \item La diferencia entre protones y neutrones
  \item La diferencia entre protones y electrones
  \item La diferencia entre neutrones y electrones
\end{multicols}
\end{enumerate}
\vspace{0.15cm}
\item ¿Qué es un isótopo?
\begin{enumerate}[label=\textbf{\Alph*)}]
\begin{multicols}{2}
  \item Un átomo con el mismo número atómico pero diferente número de masa
  \item Un átomo con diferente número atómico pero mismo número de masa
  \item Un átomo con mismo número de protones y electrones
  \item Un átomo con diferente número de protones y electrones
  \item Un átomo con mismo número de neutrones y electrones
  \item Un átomo con diferente número de neutrones y electrones
\end{multicols}
\end{enumerate}
\vspace{0.15cm}
\item ¿Qué es un ion?
\begin{enumerate}[label=\textbf{\Alph*)}]
\begin{multicols}{2}
  \item Un átomo con carga eléctrica
  \item Un átomo sin carga eléctrica
  \item Un átomo con igual número de protones y electrones
  \item Un átomo con igual número de protones y neutrones
  \item Un átomo con igual número de neutrones y electrones
  \item Un átomo con diferente número de protones y neutrones
\end{multicols}
\end{enumerate}
\vspace{0.15cm}
\item ¿Qué es un catión?
\begin{enumerate}[label=\textbf{\Alph*)}]
\begin{multicols}{2}
  \item Un ion con carga positiva
  \item Un ion con carga negativa
  \item Un ion sin carga
  \item Un átomo neutro
  \item Un isótopo
  \item Un isóbaro
\end{multicols}
\end{enumerate}
\vspace{0.15cm}
\item ¿Qué es un anión?
\begin{enumerate}[label=\textbf{\Alph*)}]
\begin{multicols}{2}
  \item Un ion con carga negativa
  \item Un ion con carga positiva
  \item Un ion sin carga
  \item Un átomo neutro
  \item Un isótopo
  \item Un isóbaro
\end{multicols}
\end{enumerate}
\vspace{0.15cm}
\item ¿Qué es un enlace químico?
\begin{enumerate}[label=\textbf{\Alph*)}]
\begin{multicols}{2}
  \item La unión entre átomos para formar compuestos
  \item La separación entre átomos
  \item La repulsión entre átomos
  \item La fusión entre núcleos
  \item La fisión entre núcleos
  \item La emisión de electrones
\end{multicols}
\end{enumerate}
\vspace{0.15cm}
\item ¿Qué es un enlace iónico?
\begin{enumerate}[label=\textbf{\Alph*)}]
\begin{multicols}{2}
  \item La transferencia de electrones entre átomos
  \item La compartición de electrones entre átomos
  \item La compartición de protones entre átomos
  \item La transferencia de neutrones entre átomos
  \item La compartición de neutrones entre átomos
  \item La transferencia de protones entre átomos
\end{multicols}
\end{enumerate}
\vspace{0.15cm}
\item ¿Qué es un enlace covalente?
\begin{enumerate}[label=\textbf{\Alph*)}]
\begin{multicols}{2}
  \item La compartición de electrones entre átomos
  \item La transferencia de electrones entre átomos
  \item La compartición de protones entre átomos
  \item La transferencia de neutrones entre átomos
  \item La compartición de neutrones entre átomos
  \item La transferencia de protones entre átomos
\end{multicols}
\end{enumerate}
\vspace{0.15cm}
\item ¿Qué es un enlace metálico?
\begin{enumerate}[label=\textbf{\Alph*)}]
\begin{multicols}{2}
  \item La unión entre átomos metálicos con electrones deslocalizados
  \item La unión entre átomos no metálicos
  \item La unión entre átomos metálicos y no metálicos
  \item La unión entre átomos con transferencia de electrones
  \item La unión entre átomos con compartición de electrones
  \item La unión entre átomos con compartición de protones
\end{multicols}
\end{enumerate}
\vspace{0.15cm}
\item ¿Cuál es un ejemplo de enlace iónico?
\begin{enumerate}[label=\textbf{\Alph*)}]
\begin{multicols}{2}
  \item El cloruro de sodio (NaCl)
  \item El agua (H₂O)
  \item El metano (CH₄)
  \item El hierro (Fe)
  \item El cobre (Cu)
  \item El oxígeno (O₂)
\end{multicols}
\end{enumerate}
\vspace{0.15cm}
\item ¿Cuál es un ejemplo de enlace covalente?
\begin{enumerate}[label=\textbf{\Alph*)}]
\begin{multicols}{2}
  \item El agua (H₂O)
  \item El cloruro de sodio (NaCl)
  \item El hierro (Fe)
  \item El cobre (Cu)
  \item El aluminio (Al)
  \item El zinc (Zn)
\end{multicols}
\end{enumerate}
\vspace{0.15cm}
\item ¿Cuál es un ejemplo de enlace metálico?
\begin{enumerate}[label=\textbf{\Alph*)}]
\begin{multicols}{2}
  \item El hierro (Fe)
  \item El agua (H₂O)
  \item El cloruro de sodio (NaCl)
  \item El metano (CH₄)
  \item El dióxido de carbono (CO₂)
  \item El amoníaco (NH₃)
\end{multicols}
\end{enumerate}
\vspace{0.15cm}
\item ¿Qué es la electronegatividad?
\begin{enumerate}[label=\textbf{\Alph*)}]
\begin{multicols}{2}
  \item La tendencia de un átomo a atraer electrones
  \item La tendencia de un átomo a ceder electrones
  \item La tendencia de un átomo a compartir protones
  \item La tendencia de un átomo a ceder protones
  \item La tendencia de un átomo a compartir neutrones
  \item La tendencia de un átomo a ceder neutrones
\end{multicols}
\end{enumerate}
\vspace{0.15cm}
\item ¿Qué es la valencia?
\begin{enumerate}[label=\textbf{\Alph*)}]
\begin{multicols}{2}
  \item La capacidad de combinación de un elemento
  \item La capacidad de ionización de un elemento
  \item La capacidad de oxidación de un elemento
  \item La capacidad de reducción de un elemento
  \item La capacidad de fusión de un elemento
  \item La capacidad de ebullición de un elemento
\end{multicols}
\end{enumerate}
\vspace{0.15cm}
\item ¿Qué es la regla del octeto?
\begin{enumerate}[label=\textbf{\Alph*)}]
\begin{multicols}{2}
  \item Los átomos tienden a tener 8 electrones en su capa externa
  \item Los átomos tienden a tener 8 protones en su núcleo
  \item Los átomos tienden a tener 8 neutrones en su núcleo
  \item Los átomos tienden a tener 8 electrones en su núcleo
  \item Los átomos tienden a tener 8 protones en su capa externa
  \item Los átomos tienden a tener 8 neutrones en su capa externa
\end{multicols}
\end{enumerate}
\vspace{0.15cm}
\item ¿Qué es la fórmula química?
\begin{enumerate}[label=\textbf{\Alph*)}]
\begin{multicols}{2}
  \item La representación de los elementos y sus proporciones en un compuesto
  \item La representación de los electrones de un átomo
  \item La representación de los protones de un átomo
  \item La representación de los neutrones de un átomo
  \item La representación de los isótopos de un elemento
  \item La representación de los iones de un elemento
\end{multicols}
\end{enumerate}
\vspace{0.15cm}
\item ¿Qué es una fórmula molecular?
\begin{enumerate}[label=\textbf{\Alph*)}]
\begin{multicols}{2}
  \item La que indica el número exacto de átomos de cada elemento
  \item La que indica la proporción más simple de átomos
  \item La que indica la estructura de los átomos
  \item La que indica los enlaces entre átomos
  \item La que indica la carga de los iones
  \item La que indica la masa de los átomos
\end{multicols}
\end{enumerate}
\vspace{0.15cm}
\item ¿Qué es una fórmula empírica?
\begin{enumerate}[label=\textbf{\Alph*)}]
\begin{multicols}{2}
  \item La que indica la proporción más simple de átomos
  \item La que indica el número exacto de átomos
  \item La que indica la estructura de los átomos
  \item La que indica los enlaces entre átomos
  \item La que indica la carga de los iones
  \item La que indica la masa de los átomos
\end{multicols}
\end{enumerate}
\vspace{0.15cm}
\item ¿Qué es una fórmula estructural?
\begin{enumerate}[label=\textbf{\Alph*)}]
\begin{multicols}{2}
  \item La que indica cómo están unidos los átomos
  \item La que indica el número exacto de átomos
  \item La que indica la proporción más simple de átomos
  \item La que indica la carga de los iones
  \item La que indica la masa de los átomos
  \item La que indica los isótopos
\end{multicols}
\end{enumerate}
\vspace{0.15cm}
\item ¿Qué es un mol?
\begin{enumerate}[label=\textbf{\Alph*)}]
\begin{multicols}{2}
  \item La cantidad de sustancia que contiene 6.022×10²³ partículas
  \item La cantidad de sustancia que contiene 6.022×10²² partículas
  \item La cantidad de sustancia que contiene 6.022×10²⁴ partículas
  \item La cantidad de sustancia que contiene 6.022×10²¹ partículas
  \item La cantidad de sustancia que contiene 6.022×10²⁰ partículas
  \item La cantidad de sustancia que contiene 6.022×10²⁵ partículas
\end{multicols}
\end{enumerate}
\vspace{0.15cm}
\item ¿Qué es el número de Avogadro?
\begin{enumerate}[label=\textbf{\Alph*)}]
\begin{multicols}{2}
  \item 6.022×10²³
  \item 6.022×10²²
  \item 6.022×10²⁴
  \item 6.022×10²¹
  \item 6.022×10²⁰
  \item 6.022×10²⁵
\end{multicols}
\end{enumerate}
\vspace{0.15cm}
\item ¿Qué es la masa molar?
\begin{enumerate}[label=\textbf{\Alph*)}]
\begin{multicols}{2}
  \item La masa de un mol de sustancia
  \item La masa de un átomo de sustancia
  \item La masa de una molécula de sustancia
  \item La masa de un ion de sustancia
  \item La masa de un isótopo de sustancia
  \item La masa de un electrón de sustancia
\end{multicols}
\end{enumerate}
\vspace{0.15cm}
\item ¿Cuál es la masa molar del agua (H₂O)?
\begin{enumerate}[label=\textbf{\Alph*)}]
\begin{multicols}{2}
  \item 18 g/mol
  \item 16 g/mol
  \item 20 g/mol
  \item 22 g/mol
  \item 24 g/mol
  \item 26 g/mol
\end{multicols}
\end{enumerate}
\vspace{0.15cm}
\item ¿Cuál es la masa molar del dióxido de carbono (CO₂)?
\begin{enumerate}[label=\textbf{\Alph*)}]
\begin{multicols}{2}
  \item 44 g/mol
  \item 28 g/mol
  \item 32 g/mol
  \item 40 g/mol
  \item 48 g/mol
  \item 52 g/mol
\end{multicols}
\end{enumerate}
\vspace{0.15cm}
\item ¿Cuál es la masa molar del cloruro de sodio (NaCl)?
\begin{enumerate}[label=\textbf{\Alph*)}]
\begin{multicols}{2}
  \item 58.5 g/mol
  \item 23 g/mol
  \item 35.5 g/mol
  \item 50 g/mol
  \item 60 g/mol
  \item 65 g/mol
\end{multicols}
\end{enumerate}
\vspace{0.15cm}
\item ¿Qué es la estequiometría?
\begin{enumerate}[label=\textbf{\Alph*)}]
\begin{multicols}{2}
  \item El cálculo de las relaciones cuantitativas en las reacciones químicas
  \item El cálculo de las relaciones cualitativas
  \item El cálculo de las propiedades físicas
  \item El cálculo de las propiedades químicas
  \item El cálculo de las estructuras moleculares
  \item El cálculo de las energías de enlace
\end{multicols}
\end{enumerate}
\vspace{0.15cm}
\item ¿Qué es un reactivo limitante?
\begin{enumerate}[label=\textbf{\Alph*)}]
\begin{multicols}{2}
  \item El que se consume primero en una reacción química
  \item El que se consume último
  \item El que no se consume
  \item El que se produce en la reacción
  \item El que cataliza la reacción
  \item El que inhibe la reacción
\end{multicols}
\end{enumerate}
\vspace{0.15cm}
\item ¿Qué es un reactivo en exceso?
\begin{enumerate}[label=\textbf{\Alph*)}]
\begin{multicols}{2}
  \item El que está en mayor cantidad de la necesaria
  \item El que está en menor cantidad de la necesaria
  \item El que se consume primero
  \item El que no se consume
  \item El que se produce en la reacción
  \item El que cataliza la reacción
\end{multicols}
\end{enumerate}
\vspace{0.15cm}
\item ¿Qué es el rendimiento de una reacción?
\begin{enumerate}[label=\textbf{\Alph*)}]
\begin{multicols}{2}
  \item La relación entre el producto obtenido y el teórico
  \item La relación entre el reactivo y el producto
  \item La relación entre el producto y el reactivo
  \item La relación entre los reactivos
  \item La relación entre los productos
  \item La relación entre el catalizador y el producto
\end{multicols}
\end{enumerate}
\vspace{0.15cm}
\item ¿Qué es un catalizador?
\begin{enumerate}[label=\textbf{\Alph*)}]
\begin{multicols}{2}
  \item Una sustancia que acelera una reacción sin consumirse
  \item Una sustancia que retarda una reacción
  \item Una sustancia que se consume en la reacción
  \item Una sustancia que se produce en la reacción
  \item Una sustancia que inhibe la reacción
  \item Una sustancia que detiene la reacción
\end{multicols}
\end{enumerate}
\vspace{0.15cm}
\item ¿Qué es un inhibidor?
\begin{enumerate}[label=\textbf{\Alph*)}]
\begin{multicols}{2}
  \item Una sustancia que retarda o detiene una reacción
  \item Una sustancia que acelera una reacción
  \item Una sustancia que se consume en la reacción
  \item Una sustancia que se produce en la reacción
  \item Una sustancia que cataliza la reacción
  \item Una sustancia que inicia la reacción
\end{multicols}
\end{enumerate}
\vspace{0.15cm}
\item ¿Qué es la energía de activación?
\begin{enumerate}[label=\textbf{\Alph*)}]
\begin{multicols}{2}
  \item La energía mínima necesaria para iniciar una reacción
  \item La energía liberada en una reacción
  \item La energía absorbida en una reacción
  \item La energía total de una reacción
  \item La energía de los reactivos
  \item La energía de los productos
\end{multicols}
\end{enumerate}
\vspace{0.15cm}
\item ¿Qué es una reacción exotérmica?
\begin{enumerate}[label=\textbf{\Alph*)}]
\begin{multicols}{2}
  \item La que libera energía en forma de calor
  \item La que absorbe energía en forma de calor
  \item La que no intercambia energía
  \item La que libera energía en forma de luz
  \item La que absorbe energía en forma de luz
  \item La que libera energía en forma de sonido
\end{multicols}
\end{enumerate}
\vspace{0.15cm}
\item ¿Qué es una reacción endotérmica?
\begin{enumerate}[label=\textbf{\Alph*)}]
\begin{multicols}{2}
  \item La que absorbe energía en forma de calor
  \item La que libera energía en forma de calor
  \item La que no intercambia energía
  \item La que libera energía en forma de luz
  \item La que absorbe energía en forma de luz
  \item La que libera energía en forma de sonido
\end{multicols}
\end{enumerate}
\vspace{0.15cm}
\item ¿Cuál es un ejemplo de reacción exotérmica?
\begin{enumerate}[label=\textbf{\Alph*)}]
\begin{multicols}{2}
  \item La combustión
  \item La fotosíntesis
  \item La electrólisis
  \item La fusión del hielo
  \item La evaporación del agua
  \item La sublimación del yodo
\end{multicols}
\end{enumerate}
\vspace{0.15cm}
\item ¿Cuál es un ejemplo de reacción endotérmica?
\begin{enumerate}[label=\textbf{\Alph*)}]
\begin{multicols}{2}
  \item La fotosíntesis
  \item La combustión
  \item La respiración celular
  \item La oxidación del hierro
  \item La fermentación
  \item La corrosión
\end{multicols}
\end{enumerate}
\vspace{0.15cm}
\item ¿Qué es el equilibrio químico?
\begin{enumerate}[label=\textbf{\Alph*)}]
\begin{multicols}{2}
  \item El estado en que las velocidades de reacción directa e inversa son iguales
  \item El estado en que la reacción directa es más rápida
  \item El estado en que la reacción inversa es más rápida
  \item El estado en que la reacción se detiene
  \item El estado en que no hay reacción
  \item El estado en que solo hay reactivos
\end{multicols}
\end{enumerate}
\vspace{0.15cm}
\item ¿Qué es el principio de Le Chatelier?
\begin{enumerate}[label=\textbf{\Alph*)}]
\begin{multicols}{2}
  \item Si se altera un sistema en equilibrio, este se desplaza para contrarrestar el cambio
  \item Si se altera un sistema en equilibrio, este no cambia
  \item Si se altera un sistema en equilibrio, este se destruye
  \item Si se altera un sistema en equilibrio, este se acelera
  \item Si se altera un sistema en equilibrio, este se detiene
  \item Si se altera un sistema en equilibrio, este se invierte
\end{multicols}
\end{enumerate}
\vspace{0.15cm}
\item ¿Qué es un ácido según Arrhenius?
\begin{enumerate}[label=\textbf{\Alph*)}]
\begin{multicols}{2}
  \item Una sustancia que libera iones H⁺ en agua
  \item Una sustancia que libera iones OH⁻ en agua
  \item Una sustancia que acepta protones
  \item Una sustancia que dona electrones
  \item Una sustancia que acepta electrones
  \item Una sustancia que dona neutrones
\end{multicols}
\end{enumerate}
\vspace{0.15cm}
\item ¿Qué es una base según Arrhenius?
\begin{enumerate}[label=\textbf{\Alph*)}]
\begin{multicols}{2}
  \item Una sustancia que libera iones OH⁻ en agua
  \item Una sustancia que libera iones H⁺ en agua
  \item Una sustancia que dona protones
  \item Una sustancia que acepta electrones
  \item Una sustancia que dona electrones
  \item Una sustancia que acepta neutrones
\end{multicols}
\end{enumerate}
\vspace{0.15cm}
\item ¿Qué es un ácido según Brønsted-Lowry?
\begin{enumerate}[label=\textbf{\Alph*)}]
\begin{multicols}{2}
  \item Una sustancia que dona protones
  \item Una sustancia que acepta protones
  \item Una sustancia que libera iones OH⁻
  \item Una sustancia que libera iones H⁺
  \item Una sustancia que acepta electrones
  \item Una sustancia que dona electrones
\end{multicols}
\end{enumerate}
\vspace{0.15cm}
\item ¿Qué es una base según Brønsted-Lowry?
\begin{enumerate}[label=\textbf{\Alph*)}]
\begin{multicols}{2}
  \item Una sustancia que acepta protones
  \item Una sustancia que dona protones
  \item Una sustancia que libera iones OH⁻
  \item Una sustancia que libera iones H⁺
  \item Una sustancia que acepta electrones
  \item Una sustancia que dona electrones
\end{multicols}
\end{enumerate}
\vspace{0.15cm}
\item ¿Qué es el pH?
\begin{enumerate}[label=\textbf{\Alph*)}]
\begin{multicols}{2}
  \item La medida de la acidez o basicidad de una disolución
  \item La medida de la temperatura de una disolución
  \item La medida de la densidad de una disolución
  \item La medida de la presión de una disolución
  \item La medida del volumen de una disolución
  \item La medida de la masa de una disolución
\end{multicols}
\end{enumerate}
\vspace{0.15cm}
\item ¿Cuál es el pH de una disolución neutra?
\begin{enumerate}[label=\textbf{\Alph*)}]
\begin{multicols}{2}
  \item 7
  \item 0
  \item 14
  \item 1
  \item 13
  \item 10
\end{multicols}
\end{enumerate}
\vspace{0.15cm}
\item ¿Cuál es el pH de una disolución ácida?
\begin{enumerate}[label=\textbf{\Alph*)}]
\begin{multicols}{2}
  \item Menor que 7
  \item Mayor que 7
  \item Igual a 7
  \item Igual a 0
  \item Igual a 14
  \item Mayor que 14
\end{multicols}
\end{enumerate}
\vspace{0.15cm}
\item ¿Cuál es el pH de una disolución básica?
\begin{enumerate}[label=\textbf{\Alph*)}]
\begin{multicols}{2}
  \item Mayor que 7
  \item Menor que 7
  \item Igual a 7
  \item Igual a 0
  \item Igual a 14
  \item Menor que 0
\end{multicols}
\end{enumerate}
\vspace{0.15cm}
\item ¿Qué es una reacción de neutralización?
\begin{enumerate}[label=\textbf{\Alph*)}]
\begin{multicols}{2}
  \item La reacción entre un ácido y una base para formar sal y agua
  \item La reacción entre dos ácidos
  \item La reacción entre dos bases
  \item La reacción entre un ácido y un metal
  \item La reacción entre una base y un no metal
  \item La reacción entre dos sales
\end{multicols}
\end{enumerate}
\vspace{0.15cm}
\item ¿Cuál es el producto de la reacción entre HCl y NaOH?
\begin{enumerate}[label=\textbf{\Alph*)}]
\begin{multicols}{2}
  \item NaCl y H₂O
  \item NaCl y H₂
  \item NaCl y O₂
  \item NaH y Cl₂
  \item Na y HCl
  \item NaOH y HCl
\end{multicols}
\end{enumerate}
\vspace{0.15cm}
\item ¿Qué es un indicador ácido-base?
\begin{enumerate}[label=\textbf{\Alph*)}]
\begin{multicols}{2}
  \item Una sustancia que cambia de color según el pH
  \item Una sustancia que cambia de temperatura
  \item Una sustancia que cambia de densidad
  \item Una sustancia que cambia de presión
  \item Una sustancia que cambia de volumen
  \item Una sustancia que cambia de masa
\end{multicols}
\end{enumerate}
\vspace{0.15cm}
\item ¿Cuál es un ejemplo de indicador ácido-base?
\begin{enumerate}[label=\textbf{\Alph*)}]
\begin{multicols}{2}
  \item La fenolftaleína
  \item El agua
  \item El alcohol
  \item La sal
  \item El azúcar
  \item El aceite
\end{multicols}
\end{enumerate}
\vspace{0.15cm}
\item ¿Qué es una sal?
\begin{enumerate}[label=\textbf{\Alph*)}]
\begin{multicols}{2}
  \item Un compuesto formado por un catión y un anión
  \item Un compuesto formado por dos cationes
  \item Un compuesto formado por dos aniones
  \item Un compuesto formado por un ácido y una base
  \item Un compuesto formado por dos ácidos
  \item Un compuesto formado por dos bases
\end{multicols}
\end{enumerate}
\vspace{0.15cm}
\item ¿Qué es la electrólisis?
\begin{enumerate}[label=\textbf{\Alph*)}]
\begin{multicols}{2}
  \item El proceso de descomposición de una sustancia mediante electricidad
  \item El proceso de descomposición de una sustancia mediante calor
  \item El proceso de descomposición de una sustancia mediante luz
  \item El proceso de descomposición de una sustancia mediante presión
  \item El proceso de descomposición de una sustancia mediante sonido
  \item El proceso de descomposición de una sustancia mediante radiación
\end{multicols}
\end{enumerate}
\vspace{0.15cm}
\item ¿Qué es una pila?
\begin{enumerate}[label=\textbf{\Alph*)}]
\begin{multicols}{2}
  \item Un dispositivo que convierte energía química en eléctrica
  \item Un dispositivo que convierte energía eléctrica en química
  \item Un dispositivo que convierte energía mecánica en eléctrica
  \item Un dispositivo que convierte energía térmica en eléctrica
  \item Un dispositivo que convierte energía solar en eléctrica
  \item Un dispositivo que convierte energía nuclear en eléctrica
\end{multicols}
\end{enumerate}
\vspace{0.15cm}
\item ¿Qué es la corrosión?
\begin{enumerate}[label=\textbf{\Alph*)}]
\begin{multicols}{2}
  \item La degradación de un metal por reacción química con el ambiente
  \item La degradación de un metal por reacción física
  \item La degradación de un metal por reacción mecánica
  \item La degradación de un metal por reacción térmica
  \item La degradación de un metal por reacción eléctrica
  \item La degradación de un metal por reacción nuclear
\end{multicols}
\end{enumerate}
\vspace{0.15cm}
\item ¿Qué es la oxidación?
\begin{enumerate}[label=\textbf{\Alph*)}]
\begin{multicols}{2}
  \item La pérdida de electrones por parte de una sustancia
  \item La ganancia de electrones por parte de una sustancia
  \item La pérdida de protones por parte de una sustancia
  \item La ganancia de protones por parte de una sustancia
  \item La pérdida de neutrones por parte de una sustancia
  \item La ganancia de neutrones por parte de una sustancia
\end{multicols}
\end{enumerate}
\vspace{0.15cm}
\item ¿Qué es la reducción?
\begin{enumerate}[label=\textbf{\Alph*)}]
\begin{multicols}{2}
  \item La ganancia de electrones por parte de una sustancia
  \item La pérdida de electrones por parte de una sustancia
  \item La ganancia de protones por parte de una sustancia
  \item La pérdida de protones por parte de una sustancia
  \item La ganancia de neutrones por parte de una sustancia
  \item La pérdida de neutrones por parte de una sustancia
\end{multicols}
\end{enumerate}
\vspace{0.15cm}
\item ¿Qué es una reacción redox?
\begin{enumerate}[label=\textbf{\Alph*)}]
\begin{multicols}{2}
  \item La que implica transferencia de electrones
  \item La que implica transferencia de protones
  \item La que implica transferencia de neutrones
  \item La que implica transferencia de átomos
  \item La que implica transferencia de moléculas
  \item La que implica transferencia de iones
\end{multicols}
\end{enumerate}
\vspace{0.15cm}
\item ¿Qué es el agente oxidante?
\begin{enumerate}[label=\textbf{\Alph*)}]
\begin{multicols}{2}
  \item La sustancia que gana electrones y se reduce
  \item La sustancia que pierde electrones y se oxida
  \item La sustancia que gana protones
  \item La sustancia que pierde protones
  \item La sustancia que gana neutrones
  \item La sustancia que pierde neutrones
\end{multicols}
\end{enumerate}
\vspace{0.15cm}
\item ¿Qué es el agente reductor?
\begin{enumerate}[label=\textbf{\Alph*)}]
\begin{multicols}{2}
  \item La sustancia que pierde electrones y se oxida
  \item La sustancia que gana electrones y se reduce
  \item La sustancia que gana protones
  \item La sustancia que pierde protones
  \item La sustancia que gana neutrones
  \item La sustancia que pierde neutrones
\end{multicols}
\end{enumerate}
\vspace{0.15cm}
\item ¿Qué es el número de oxidación?
\begin{enumerate}[label=\textbf{\Alph*)}]
\begin{multicols}{2}
  \item La carga que tendría un átomo si los electrones se asignaran al átomo más electronegativo
  \item La carga que tendría un átomo si los electrones se asignaran al átomo menos electronegativo
  \item La carga real de un átomo en un compuesto
  \item La carga de un ion en disolución
  \item La carga de un electrón
  \item La carga de un protón
\end{multicols}
\end{enumerate}
\vspace{0.15cm}
\item ¿Qué es una reacción de síntesis o combinación?
\begin{enumerate}[label=\textbf{\Alph*)}]
\begin{multicols}{2}
  \item La que forma un compuesto a partir de dos o más sustancias
  \item La que descompone un compuesto en dos o más sustancias
  \item La que sustituye un elemento por otro
  \item La que intercambia iones entre compuestos
  \item La que neutraliza un ácido y una base
  \item La que combustiona un compuesto
\end{multicols}
\end{enumerate}
\vspace{0.15cm}
\item ¿Qué es una reacción de descomposición?
\begin{enumerate}[label=\textbf{\Alph*)}]
\begin{multicols}{2}
  \item La que descompone un compuesto en dos o más sustancias
  \item La que forma un compuesto a partir de dos o más sustancias
  \item La que sustituye un elemento por otro
  \item La que intercambia iones entre compuestos
  \item La que neutraliza un ácido y una base
  \item La que combustiona un compuesto ---
\end{multicols}
\end{enumerate}
\vspace{0.15cm}
\end{enumerate}\n
\subsubsection*{SECCIÓN III: TIPOS DE REACCIONES QUÍMICAS (121–200)}

\begin{enumerate}[label=\textbf{\arabic*.}]
\setcounter{enumi}{120}
\item ¿Qué es una reacción de combinación?
\begin{enumerate}[label=\textbf{\Alph*)}]
\begin{multicols}{2}
  \item La que forma un compuesto a partir de dos o más sustancias
  \item La que descompone un compuesto
  \item La que sustituye un elemento
  \item La que intercambia iones
  \item La que neutraliza un ácido y una base
  \item La que combustiona un compuesto
\end{multicols}
\end{enumerate}
\vspace{0.15cm}
\item ¿Cuál es un ejemplo de reacción de combinación?
\begin{enumerate}[label=\textbf{\Alph*)}]
\begin{multicols}{2}
  \item 2H₂ + O₂ → 2H₂O
  \item 2H₂O → 2H₂ + O₂
  \item Zn + CuSO₄ → ZnSO₄ + Cu
  \item AgNO₃ + NaCl → AgCl + NaNO₃
  \item HCl + NaOH → NaCl + H₂O
  \item CH₄ + 2O₂ → CO₂ + 2H₂O
\end{multicols}
\end{enumerate}
\vspace{0.15cm}
\item ¿Qué es una reacción de descomposición?
\begin{enumerate}[label=\textbf{\Alph*)}]
\begin{multicols}{2}
  \item La que descompone un compuesto en dos o más sustancias
  \item La que forma un compuesto
  \item La que sustituye un elemento
  \item La que intercambia iones
  \item La que neutraliza un ácido y una base
  \item La que combustiona un compuesto
\end{multicols}
\end{enumerate}
\vspace{0.15cm}
\item ¿Cuál es un ejemplo de reacción de descomposición?
\begin{enumerate}[label=\textbf{\Alph*)}]
\begin{multicols}{2}
  \item 2H₂O → 2H₂ + O₂
  \item 2H₂ + O₂ → 2H₂O
  \item Zn + CuSO₄ → ZnSO₄ + Cu
  \item AgNO₃ + NaCl → AgCl + NaNO₃
  \item HCl + NaOH → NaCl + H₂O
  \item CH₄ + 2O₂ → CO₂ + 2H₂O
\end{multicols}
\end{enumerate}
\vspace{0.15cm}
\item ¿Qué es una reacción de sustitución o desplazamiento?
\begin{enumerate}[label=\textbf{\Alph*)}]
\begin{multicols}{2}
  \item La que sustituye un elemento por otro en un compuesto
  \item La que forma un compuesto
  \item La que descompone un compuesto
  \item La que intercambia iones
  \item La que neutraliza un ácido y una base
  \item La que combustiona un compuesto
\end{multicols}
\end{enumerate}
\vspace{0.15cm}
\item ¿Cuál es un ejemplo de reacción de sustitución?
\begin{enumerate}[label=\textbf{\Alph*)}]
\begin{multicols}{2}
  \item Zn + CuSO₄ → ZnSO₄ + Cu
  \item 2H₂ + O₂ → 2H₂O
  \item 2H₂O → 2H₂ + O₂
  \item AgNO₃ + NaCl → AgCl + NaNO₃
  \item HCl + NaOH → NaCl + H₂O
  \item CH₄ + 2O₂ → CO₂ + 2H₂O
\end{multicols}
\end{enumerate}
\vspace{0.15cm}
\item ¿Qué es una reacción de doble sustitución o metátesis?
\begin{enumerate}[label=\textbf{\Alph*)}]
\begin{multicols}{2}
  \item La que intercambia iones entre dos compuestos
  \item La que forma un compuesto
  \item La que descompone un compuesto
  \item La que sustituye un elemento
  \item La que neutraliza un ácido y una base
  \item La que combustiona un compuesto
\end{multicols}
\end{enumerate}
\vspace{0.15cm}
\item ¿Cuál es un ejemplo de reacción de doble sustitución?
\begin{enumerate}[label=\textbf{\Alph*)}]
\begin{multicols}{2}
  \item AgNO₃ + NaCl → AgCl + NaNO₃
  \item 2H₂ + O₂ → 2H₂O
  \item 2H₂O → 2H₂ + O₂
  \item Zn + CuSO₄ → ZnSO₄ + Cu
  \item HCl + NaOH → NaCl + H₂O
  \item CH₄ + 2O₂ → CO₂ + 2H₂O
\end{multicols}
\end{enumerate}
\vspace{0.15cm}
\item ¿Qué es una reacción de combustión?
\begin{enumerate}[label=\textbf{\Alph*)}]
\begin{multicols}{2}
  \item La que reacciona con oxígeno liberando energía
  \item La que reacciona con hidrógeno
  \item La que reacciona con nitrógeno
  \item La que reacciona con cloro
  \item La que reacciona con azufre
  \item La que reacciona con fósforo
\end{multicols}
\end{enumerate}
\vspace{0.15cm}
\item ¿Cuál es un ejemplo de reacción de combustión?
\begin{enumerate}[label=\textbf{\Alph*)}]
\begin{multicols}{2}
  \item CH₄ + 2O₂ → CO₂ + 2H₂O
  \item 2H₂ + O₂ → 2H₂O
  \item 2H₂O → 2H₂ + O₂
  \item Zn + CuSO₄ → ZnSO₄ + Cu
  \item AgNO₃ + NaCl → AgCl + NaNO₃
  \item HCl + NaOH → NaCl + H₂O
\end{multicols}
\end{enumerate}
\vspace{0.15cm}
\item ¿Qué es una reacción de neutralización?
\begin{enumerate}[label=\textbf{\Alph*)}]
\begin{multicols}{2}
  \item La que reacciona un ácido con una base para formar sal y agua
  \item La que reacciona un ácido con un metal
  \item La que reacciona una base con un no metal
  \item La que reacciona dos ácidos
  \item La que reacciona dos bases
  \item La que reacciona dos sales
\end{multicols}
\end{enumerate}
\vspace{0.15cm}
\item ¿Cuál es un ejemplo de reacción de neutralización?
\begin{enumerate}[label=\textbf{\Alph*)}]
\begin{multicols}{2}
  \item HCl + NaOH → NaCl + H₂O
  \item 2H₂ + O₂ → 2H₂O
  \item 2H₂O → 2H₂ + O₂
  \item Zn + CuSO₄ → ZnSO₄ + Cu
  \item AgNO₃ + NaCl → AgCl + NaNO₃
  \item CH₄ + 2O₂ → CO₂ + 2H₂O
\end{multicols}
\end{enumerate}
\vspace{0.15cm}
\item ¿Qué es una reacción exotérmica?
\begin{enumerate}[label=\textbf{\Alph*)}]
\begin{multicols}{2}
  \item La que libera energía en forma de calor
  \item La que absorbe energía en forma de calor
  \item La que no intercambia energía
  \item La que libera energía en forma de luz
  \item La que absorbe energía en forma de luz
  \item La que libera energía en forma de sonido
\end{multicols}
\end{enumerate}
\vspace{0.15cm}
\item ¿Qué es una reacción endotérmica?
\begin{enumerate}[label=\textbf{\Alph*)}]
\begin{multicols}{2}
  \item La que absorbe energía en forma de calor
  \item La que libera energía en forma de calor
  \item La que no intercambia energía
  \item La que libera energía en forma de luz
  \item La que absorbe energía en forma de luz
  \item La que libera energía en forma de sonido
\end{multicols}
\end{enumerate}
\vspace{0.15cm}
\item ¿Qué es una reacción reversible?
\begin{enumerate}[label=\textbf{\Alph*)}]
\begin{multicols}{2}
  \item La que puede ocurrir en ambos sentidos
  \item La que ocurre en un solo sentido
  \item La que no ocurre
  \item La que ocurre solo en presencia de catalizador
  \item La que ocurre solo a altas temperaturas
  \item La que ocurre solo a bajas temperaturas
\end{multicols}
\end{enumerate}
\vspace{0.15cm}
\item ¿Qué es una reacción irreversible?
\begin{enumerate}[label=\textbf{\Alph*)}]
\begin{multicols}{2}
  \item La que ocurre en un solo sentido
  \item La que ocurre en ambos sentidos
  \item La que no ocurre
  \item La que ocurre solo en presencia de catalizador
  \item La que ocurre solo a altas temperaturas
  \item La que ocurre solo a bajas temperaturas
\end{multicols}
\end{enumerate}
\vspace{0.15cm}
\item ¿Qué es una reacción de precipitación?
\begin{enumerate}[label=\textbf{\Alph*)}]
\begin{multicols}{2}
  \item La que forma un sólido insoluble
  \item La que forma un gas
  \item La que forma un líquido
  \item La que forma un sólido soluble
  \item La que forma un líquido inmiscible
  \item La que forma un gas soluble
\end{multicols}
\end{enumerate}
\vspace{0.15cm}
\item ¿Cuál es un ejemplo de reacción de precipitación?
\begin{enumerate}[label=\textbf{\Alph*)}]
\begin{multicols}{2}
  \item AgNO₃ + NaCl → AgCl↓ + NaNO₃
  \item 2H₂ + O₂ → 2H₂O
  \item 2H₂O → 2H₂ + O₂
  \item Zn + CuSO₄ → ZnSO₄ + Cu
  \item HCl + NaOH → NaCl + H₂O
  \item CH₄ + 2O₂ → CO₂ + 2H₂O
\end{multicols}
\end{enumerate}
\vspace{0.15cm}
\item ¿Qué es una reacción de formación de gas?
\begin{enumerate}[label=\textbf{\Alph*)}]
\begin{multicols}{2}
  \item La que produce un gas como producto
  \item La que produce un sólido
  \item La que produce un líquido
  \item La que produce un sólido soluble
  \item La que produce un líquido inmiscible
  \item La que produce un gas soluble
\end{multicols}
\end{enumerate}
\vspace{0.15cm}
\item ¿Cuál es un ejemplo de reacción de formación de gas?
\begin{enumerate}[label=\textbf{\Alph*)}]
\begin{multicols}{2}
  \item CaCO₃ + 2HCl → CaCl₂ + H₂O + CO₂↑
  \item 2H₂ + O₂ → 2H₂O
  \item 2H₂O → 2H₂ + O₂
  \item Zn + CuSO₄ → ZnSO₄ + Cu
  \item AgNO₃ + NaCl → AgCl↓ + NaNO₃
  \item HCl + NaOH → NaCl + H₂O
\end{multicols}
\end{enumerate}
\vspace{0.15cm}
\item ¿Qué es el balanceo de ecuaciones químicas?
\begin{enumerate}[label=\textbf{\Alph*)}]
\begin{multicols}{2}
  \item El ajuste de los coeficientes para cumplir la ley de conservación de la masa
  \item El ajuste de los subíndices
  \item El ajuste de los estados de oxidación
  \item El ajuste de las temperaturas
  \item El ajuste de las presiones
  \item El ajuste de los volúmenes
\end{multicols}
\end{enumerate}
\vspace{0.15cm}
\item ¿Qué establece la ley de conservación de la masa?
\begin{enumerate}[label=\textbf{\Alph*)}]
\begin{multicols}{2}
  \item La masa de los reactivos es igual a la masa de los productos
  \item La masa de los reactivos es mayor que la de los productos
  \item La masa de los reactivos es menor que la de los productos
  \item La masa no se conserva en las reacciones
  \item La masa se destruye en las reacciones
  \item La masa se crea en las reacciones
\end{multicols}
\end{enumerate}
\vspace{0.15cm}
\item ¿Cuál es el coeficiente correcto para balancear: __H₂ + O₂ → __H₂O?
\begin{enumerate}[label=\textbf{\Alph*)}]
\begin{multicols}{2}
  \item 2, 1, 2
  \item 1, 1, 1
  \item 2, 2, 2
  \item 1, 2, 1
  \item 2, 1, 1
  \item 1, 1, 2
\end{multicols}
\end{enumerate}
\vspace{0.15cm}
\item ¿Cuál es el coeficiente correcto para balancear: __N₂ + __H₂ → __NH₃?
\begin{enumerate}[label=\textbf{\Alph*)}]
\begin{multicols}{2}
  \item 1, 3, 2
  \item 1, 2, 3
  \item 2, 3, 1
  \item 3, 1, 2
  \item 2, 1, 3
  \item 1, 1, 1
\end{multicols}
\end{enumerate}
\vspace{0.15cm}
\item ¿Cuál es el coeficiente correcto para balancear: __CH₄ + __O₂ → __CO₂ + __H₂O?
\begin{enumerate}[label=\textbf{\Alph*)}]
\begin{multicols}{2}
  \item 1, 2, 1, 2
  \item 1, 1, 1, 2
  \item 1, 2, 2, 1
  \item 2, 1, 1, 2
  \item 1, 3, 1, 2
  \item 2, 2, 1, 1
\end{multicols}
\end{enumerate}
\vspace{0.15cm}
\item ¿Cuál es el coeficiente correcto para balancear: __Fe + __O₂ → __Fe₂O₃?
\begin{enumerate}[label=\textbf{\Alph*)}]
\begin{multicols}{2}
  \item 4, 3, 2
  \item 2, 3, 1
  \item 3, 4, 2
  \item 2, 1, 1
  \item 4, 2, 2
  \item 1, 1, 1
\end{multicols}
\end{enumerate}
\vspace{0.15cm}
\item ¿Cuál es el coeficiente correcto para balancear: __Al + __HCl → __AlCl₃ + __H₂?
\begin{enumerate}[label=\textbf{\Alph*)}]
\begin{multicols}{2}
  \item 2, 6, 2, 3
  \item 1, 3, 1, 1
  \item 2, 3, 2, 3
  \item 1, 6, 1, 3
  \item 2, 6, 3, 2
  \item 3, 6, 2, 2
\end{multicols}
\end{enumerate}
\vspace{0.15cm}
\item ¿Qué es una ecuación iónica?
\begin{enumerate}[label=\textbf{\Alph*)}]
\begin{multicols}{2}
  \item La que muestra los iones que participan en la reacción
  \item La que muestra las moléculas
  \item La que muestra los átomos
  \item La que muestra los electrones
  \item La que muestra los protones
  \item La que muestra los neutrones
\end{multicols}
\end{enumerate}
\vspace{0.15cm}
\item ¿Qué es una ecuación iónica neta?
\begin{enumerate}[label=\textbf{\Alph*)}]
\begin{multicols}{2}
  \item La que muestra solo los iones que reaccionan
  \item La que muestra todos los iones
  \item La que muestra las moléculas
  \item La que muestra los átomos
  \item La que muestra los electrones
  \item La que muestra los protones
\end{multicols}
\end{enumerate}
\vspace{0.15cm}
\item ¿Qué es un ion espectador?
\begin{enumerate}[label=\textbf{\Alph*)}]
\begin{multicols}{2}
  \item El que no participa en la reacción
  \item El que participa en la reacción
  \item El que se forma en la reacción
  \item El que se consume en la reacción
  \item El que cataliza la reacción
  \item El que inhibe la reacción
\end{multicols}
\end{enumerate}
\vspace{0.15cm}
\item ¿Qué es una reacción de combinación?
\begin{enumerate}[label=\textbf{\Alph*)}]
\begin{multicols}{2}
  \item La que forma un compuesto a partir de dos o más sustancias
  \item La que descompone un compuesto
  \item La que sustituye un elemento
  \item La que intercambia iones
  \item La que neutraliza un ácido y una base
  \item La que combustiona un compuesto
\end{multicols}
\end{enumerate}
\vspace{0.15cm}
\item ¿Cuál es un ejemplo de reacción de combinación?
\begin{enumerate}[label=\textbf{\Alph*)}]
\begin{multicols}{2}
  \item 2H₂ + O₂ → 2H₂O
  \item 2H₂O → 2H₂ + O₂
  \item Zn + CuSO₄ → ZnSO₄ + Cu
  \item AgNO₃ + NaCl → AgCl + NaNO₃
  \item HCl + NaOH → NaCl + H₂O
  \item CH₄ + 2O₂ → CO₂ + 2H₂O
\end{multicols}
\end{enumerate}
\vspace{0.15cm}
\item ¿Qué es una reacción de descomposición?
\begin{enumerate}[label=\textbf{\Alph*)}]
\begin{multicols}{2}
  \item La que descompone un compuesto en dos o más sustancias
  \item La que forma un compuesto
  \item La que sustituye un elemento
  \item La que intercambia iones
  \item La que neutraliza un ácido y una base
  \item La que combustiona un compuesto
\end{multicols}
\end{enumerate}
\vspace{0.15cm}
\item ¿Cuál es un ejemplo de reacción de descomposición?
\begin{enumerate}[label=\textbf{\Alph*)}]
\begin{multicols}{2}
  \item 2H₂O → 2H₂ + O₂
  \item 2H₂ + O₂ → 2H₂O
  \item Zn + CuSO₄ → ZnSO₄ + Cu
  \item AgNO₃ + NaCl → AgCl + NaNO₃
  \item HCl + NaOH → NaCl + H₂O
  \item CH₄ + 2O₂ → CO₂ + 2H₂O
\end{multicols}
\end{enumerate}
\vspace{0.15cm}
\item ¿Qué es una reacción de sustitución?
\begin{enumerate}[label=\textbf{\Alph*)}]
\begin{multicols}{2}
  \item La que sustituye un elemento por otro en un compuesto
  \item La que forma un compuesto
  \item La que descompone un compuesto
  \item La que intercambia iones
  \item La que neutraliza un ácido y una base
  \item La que combustiona un compuesto
\end{multicols}
\end{enumerate}
\vspace{0.15cm}
\item ¿Cuál es un ejemplo de reacción de sustitución?
\begin{enumerate}[label=\textbf{\Alph*)}]
\begin{multicols}{2}
  \item Zn + CuSO₄ → ZnSO₄ + Cu
  \item 2H₂ + O₂ → 2H₂O
  \item 2H₂O → 2H₂ + O₂
  \item AgNO₃ + NaCl → AgCl + NaNO₃
  \item HCl + NaOH → NaCl + H₂O
  \item CH₄ + 2O₂ → CO₂ + 2H₂O
\end{multicols}
\end{enumerate}
\vspace{0.15cm}
\item ¿Qué es una reacción de doble sustitución?
\begin{enumerate}[label=\textbf{\Alph*)}]
\begin{multicols}{2}
  \item La que intercambia iones entre dos compuestos
  \item La que forma un compuesto
  \item La que descompone un compuesto
  \item La que sustituye un elemento
  \item La que neutraliza un ácido y una base
  \item La que combustiona un compuesto
\end{multicols}
\end{enumerate}
\vspace{0.15cm}
\item ¿Cuál es un ejemplo de reacción de doble sustitución?
\begin{enumerate}[label=\textbf{\Alph*)}]
\begin{multicols}{2}
  \item AgNO₃ + NaCl → AgCl + NaNO₃
  \item 2H₂ + O₂ → 2H₂O
  \item 2H₂O → 2H₂ + O₂
  \item Zn + CuSO₄ → ZnSO₄ + Cu
  \item HCl + NaOH → NaCl + H₂O
  \item CH₄ + 2O₂ → CO₂ + 2H₂O
\end{multicols}
\end{enumerate}
\vspace{0.15cm}
\item ¿Qué es una reacción de combustión?
\begin{enumerate}[label=\textbf{\Alph*)}]
\begin{multicols}{2}
  \item La que reacciona con oxígeno liberando energía
  \item La que reacciona con hidrógeno
  \item La que reacciona con nitrógeno
  \item La que reacciona con cloro
  \item La que reacciona con azufre
  \item La que reacciona con fósforo
\end{multicols}
\end{enumerate}
\vspace{0.15cm}
\item ¿Cuál es un ejemplo de reacción de combustión?
\begin{enumerate}[label=\textbf{\Alph*)}]
\begin{multicols}{2}
  \item CH₄ + 2O₂ → CO₂ + 2H₂O
  \item 2H₂ + O₂ → 2H₂O
  \item 2H₂O → 2H₂ + O₂
  \item Zn + CuSO₄ → ZnSO₄ + Cu
  \item AgNO₃ + NaCl → AgCl + NaNO₃
  \item HCl + NaOH → NaCl + H₂O
\end{multicols}
\end{enumerate}
\vspace{0.15cm}
\item ¿Qué es una reacción de neutralización?
\begin{enumerate}[label=\textbf{\Alph*)}]
\begin{multicols}{2}
  \item La que reacciona un ácido con una base para formar sal y agua
  \item La que reacciona un ácido con un metal
  \item La que reacciona una base con un no metal
  \item La que reacciona dos ácidos
  \item La que reacciona dos bases
  \item La que reacciona dos sales
\end{multicols}
\end{enumerate}
\vspace{0.15cm}
\item ¿Cuál es un ejemplo de reacción de neutralización?
\begin{enumerate}[label=\textbf{\Alph*)}]
\begin{multicols}{2}
  \item HCl + NaOH → NaCl + H₂O
  \item 2H₂ + O₂ → 2H₂O
  \item 2H₂O → 2H₂ + O₂
  \item Zn + CuSO₄ → ZnSO₄ + Cu
  \item AgNO₃ + NaCl → AgCl + NaNO₃
  \item CH₄ + 2O₂ → CO₂ + 2H₂O
\end{multicols}
\end{enumerate}
\vspace{0.15cm}
\item ¿Qué es una reacción exotérmica?
\begin{enumerate}[label=\textbf{\Alph*)}]
\begin{multicols}{2}
  \item La que libera energía en forma de calor
  \item La que absorbe energía en forma de calor
  \item La que no intercambia energía
  \item La que libera energía en forma de luz
  \item La que absorbe energía en forma de luz
  \item La que libera energía en forma de sonido
\end{multicols}
\end{enumerate}
\vspace{0.15cm}
\item ¿Qué es una reacción endotérmica?
\begin{enumerate}[label=\textbf{\Alph*)}]
\begin{multicols}{2}
  \item La que absorbe energía en forma de calor
  \item La que libera energía en forma de calor
  \item La que no intercambia energía
  \item La que libera energía en forma de luz
  \item La que absorbe energía en forma de luz
  \item La que libera energía en forma de sonido
\end{multicols}
\end{enumerate}
\vspace{0.15cm}
\item ¿Qué es una reacción reversible?
\begin{enumerate}[label=\textbf{\Alph*)}]
\begin{multicols}{2}
  \item La que puede ocurrir en ambos sentidos
  \item La que ocurre en un solo sentido
  \item La que no ocurre
  \item La que ocurre solo en presencia de catalizador
  \item La que ocurre solo a altas temperaturas
  \item La que ocurre solo a bajas temperaturas
\end{multicols}
\end{enumerate}
\vspace{0.15cm}
\item ¿Qué es una reacción irreversible?
\begin{enumerate}[label=\textbf{\Alph*)}]
\begin{multicols}{2}
  \item La que ocurre en un solo sentido
  \item La que ocurre en ambos sentidos
  \item La que no ocurre
  \item La que ocurre solo en presencia de catalizador
  \item La que ocurre solo a altas temperaturas
  \item La que ocurre solo a bajas temperaturas
\end{multicols}
\end{enumerate}
\vspace{0.15cm}
\item ¿Qué es una reacción de precipitación?
\begin{enumerate}[label=\textbf{\Alph*)}]
\begin{multicols}{2}
  \item La que forma un sólido insoluble
  \item La que forma un gas
  \item La que forma un líquido
  \item La que forma un sólido soluble
  \item La que forma un líquido inmiscible
  \item La que forma un gas soluble
\end{multicols}
\end{enumerate}
\vspace{0.15cm}
\item ¿Cuál es un ejemplo de reacción de precipitación?
\begin{enumerate}[label=\textbf{\Alph*)}]
\begin{multicols}{2}
  \item AgNO₃ + NaCl → AgCl↓ + NaNO₃
  \item 2H₂ + O₂ → 2H₂O
  \item 2H₂O → 2H₂ + O₂
  \item Zn + CuSO₄ → ZnSO₄ + Cu
  \item HCl + NaOH → NaCl + H₂O
  \item CH₄ + 2O₂ → CO₂ + 2H₂O
\end{multicols}
\end{enumerate}
\vspace{0.15cm}
\item ¿Qué es una reacción de formación de gas?
\begin{enumerate}[label=\textbf{\Alph*)}]
\begin{multicols}{2}
  \item La que produce un gas como producto
  \item La que produce un sólido
  \item La que produce un líquido
  \item La que produce un sólido soluble
  \item La que produce un líquido inmiscible
  \item La que produce un gas soluble
\end{multicols}
\end{enumerate}
\vspace{0.15cm}
\item ¿Cuál es un ejemplo de reacción de formación de gas?
\begin{enumerate}[label=\textbf{\Alph*)}]
\begin{multicols}{2}
  \item CaCO₃ + 2HCl → CaCl₂ + H₂O + CO₂↑
  \item 2H₂ + O₂ → 2H₂O
  \item 2H₂O → 2H₂ + O₂
  \item Zn + CuSO₄ → ZnSO₄ + Cu
  \item AgNO₃ + NaCl → AgCl↓ + NaNO₃
  \item HCl + NaOH → NaCl + H₂O
\end{multicols}
\end{enumerate}
\vspace{0.15cm}
\item ¿Qué es el balanceo de ecuaciones químicas?
\begin{enumerate}[label=\textbf{\Alph*)}]
\begin{multicols}{2}
  \item El ajuste de los coeficientes para cumplir la ley de conservación de la masa
  \item El ajuste de los subíndices
  \item El ajuste de los estados de oxidación
  \item El ajuste de las temperaturas
  \item El ajuste de las presiones
  \item El ajuste de los volúmenes
\end{multicols}
\end{enumerate}
\vspace{0.15cm}
\item ¿Qué establece la ley de conservación de la masa?
\begin{enumerate}[label=\textbf{\Alph*)}]
\begin{multicols}{2}
  \item La masa de los reactivos es igual a la masa de los productos
  \item La masa de los reactivos es mayor que la de los productos
  \item La masa de los reactivos es menor que la de los productos
  \item La masa no se conserva en las reacciones
  \item La masa se destruye en las reacciones
  \item La masa se crea en las reacciones
\end{multicols}
\end{enumerate}
\vspace{0.15cm}
\item ¿Cuál es el coeficiente correcto para balancear: __H₂ + O₂ → __H₂O?
\begin{enumerate}[label=\textbf{\Alph*)}]
\begin{multicols}{2}
  \item 2, 1, 2
  \item 1, 1, 1
  \item 2, 2, 2
  \item 1, 2, 1
  \item 2, 1, 1
  \item 1, 1, 2
\end{multicols}
\end{enumerate}
\vspace{0.15cm}
\item ¿Cuál es el coeficiente correcto para balancear: __N₂ + __H₂ → __NH₃?
\begin{enumerate}[label=\textbf{\Alph*)}]
\begin{multicols}{2}
  \item 1, 3, 2
  \item 1, 2, 3
  \item 2, 3, 1
  \item 3, 1, 2
  \item 2, 1, 3
  \item 1, 1, 1
\end{multicols}
\end{enumerate}
\vspace{0.15cm}
\item ¿Cuál es el coeficiente correcto para balancear: __CH₄ + __O₂ → __CO₂ + __H₂O?
\begin{enumerate}[label=\textbf{\Alph*)}]
\begin{multicols}{2}
  \item 1, 2, 1, 2
  \item 1, 1, 1, 2
  \item 1, 2, 2, 1
  \item 2, 1, 1, 2
  \item 1, 3, 1, 2
  \item 2, 2, 1, 1
\end{multicols}
\end{enumerate}
\vspace{0.15cm}
\item ¿Cuál es el coeficiente correcto para balancear: __Fe + __O₂ → __Fe₂O₃?
\begin{enumerate}[label=\textbf{\Alph*)}]
\begin{multicols}{2}
  \item 4, 3, 2
  \item 2, 3, 1
  \item 3, 4, 2
  \item 2, 1, 1
  \item 4, 2, 2
  \item 1, 1, 1
\end{multicols}
\end{enumerate}
\vspace{0.15cm}
\item ¿Cuál es el coeficiente correcto para balancear: __Al + __HCl → __AlCl₃ + __H₂?
\begin{enumerate}[label=\textbf{\Alph*)}]
\begin{multicols}{2}
  \item 2, 6, 2, 3
  \item 1, 3, 1, 1
  \item 2, 3, 2, 3
  \item 1, 6, 1, 3
  \item 2, 6, 3, 2
  \item 3, 6, 2, 2
\end{multicols}
\end{enumerate}
\vspace{0.15cm}
\item ¿Qué es una ecuación iónica?
\begin{enumerate}[label=\textbf{\Alph*)}]
\begin{multicols}{2}
  \item La que muestra los iones que participan en la reacción
  \item La que muestra las moléculas
  \item La que muestra los átomos
  \item La que muestra los electrones
  \item La que muestra los protones
  \item La que muestra los neutrones
\end{multicols}
\end{enumerate}
\vspace{0.15cm}
\item ¿Qué es una ecuación iónica neta?
\begin{enumerate}[label=\textbf{\Alph*)}]
\begin{multicols}{2}
  \item La que muestra solo los iones que reaccionan
  \item La que muestra todos los iones
  \item La que muestra las moléculas
  \item La que muestra los átomos
  \item La que muestra los electrones
  \item La que muestra los protones
\end{multicols}
\end{enumerate}
\vspace{0.15cm}
\item ¿Qué es un ion espectador?
\begin{enumerate}[label=\textbf{\Alph*)}]
\begin{multicols}{2}
  \item El que no participa en la reacción
  \item El que participa en la reacción
  \item El que se forma en la reacción
  \item El que se consume en la reacción
  \item El que cataliza la reacción
  \item El que inhibe la reacción
\end{multicols}
\end{enumerate}
\vspace{0.15cm}
\item ¿Qué es una reacción de combinación?
\begin{enumerate}[label=\textbf{\Alph*)}]
\begin{multicols}{2}
  \item La que forma un compuesto a partir de dos o más sustancias
  \item La que descompone un compuesto
  \item La que sustituye un elemento
  \item La que intercambia iones
  \item La que neutraliza un ácido y una base
  \item La que combustiona un compuesto
\end{multicols}
\end{enumerate}
\vspace{0.15cm}
\item ¿Cuál es un ejemplo de reacción de combinación?
\begin{enumerate}[label=\textbf{\Alph*)}]
\begin{multicols}{2}
  \item 2H₂ + O₂ → 2H₂O
  \item 2H₂O → 2H₂ + O₂
  \item Zn + CuSO₄ → ZnSO₄ + Cu
  \item AgNO₃ + NaCl → AgCl + NaNO₃
  \item HCl + NaOH → NaCl + H₂O
  \item CH₄ + 2O₂ → CO₂ + 2H₂O
\end{multicols}
\end{enumerate}
\vspace{0.15cm}
\item ¿Qué es una reacción de descomposición?
\begin{enumerate}[label=\textbf{\Alph*)}]
\begin{multicols}{2}
  \item La que descompone un compuesto en dos o más sustancias
  \item La que forma un compuesto
  \item La que sustituye un elemento
  \item La que intercambia iones
  \item La que neutraliza un ácido y una base
  \item La que combustiona un compuesto
\end{multicols}
\end{enumerate}
\vspace{0.15cm}
\item ¿Cuál es un ejemplo de reacción de descomposición?
\begin{enumerate}[label=\textbf{\Alph*)}]
\begin{multicols}{2}
  \item 2H₂O → 2H₂ + O₂
  \item 2H₂ + O₂ → 2H₂O
  \item Zn + CuSO₄ → ZnSO₄ + Cu
  \item AgNO₃ + NaCl → AgCl + NaNO₃
  \item HCl + NaOH → NaCl + H₂O
  \item CH₄ + 2O₂ → CO₂ + 2H₂O
\end{multicols}
\end{enumerate}
\vspace{0.15cm}
\item ¿Qué es una reacción de sustitución?
\begin{enumerate}[label=\textbf{\Alph*)}]
\begin{multicols}{2}
  \item La que sustituye un elemento por otro en un compuesto
  \item La que forma un compuesto
  \item La que descompone un compuesto
  \item La que intercambia iones
  \item La que neutraliza un ácido y una base
  \item La que combustiona un compuesto
\end{multicols}
\end{enumerate}
\vspace{0.15cm}
\item ¿Cuál es un ejemplo de reacción de sustitución?
\begin{enumerate}[label=\textbf{\Alph*)}]
\begin{multicols}{2}
  \item Zn + CuSO₄ → ZnSO₄ + Cu
  \item 2H₂ + O₂ → 2H₂O
  \item 2H₂O → 2H₂ + O₂
  \item AgNO₃ + NaCl → AgCl + NaNO₃
  \item HCl + NaOH → NaCl + H₂O
  \item CH₄ + 2O₂ → CO₂ + 2H₂O
\end{multicols}
\end{enumerate}
\vspace{0.15cm}
\item ¿Qué es una reacción de doble sustitución?
\begin{enumerate}[label=\textbf{\Alph*)}]
\begin{multicols}{2}
  \item La que intercambia iones entre dos compuestos
  \item La que forma un compuesto
  \item La que descompone un compuesto
  \item La que sustituye un elemento
  \item La que neutraliza un ácido y una base
  \item La que combustiona un compuesto
\end{multicols}
\end{enumerate}
\vspace{0.15cm}
\item ¿Cuál es un ejemplo de reacción de doble sustitución?
\begin{enumerate}[label=\textbf{\Alph*)}]
\begin{multicols}{2}
  \item AgNO₃ + NaCl → AgCl + NaNO₃
  \item 2H₂ + O₂ → 2H₂O
  \item 2H₂O → 2H₂ + O₂
  \item Zn + CuSO₄ → ZnSO₄ + Cu
  \item HCl + NaOH → NaCl + H₂O
  \item CH₄ + 2O₂ → CO₂ + 2H₂O
\end{multicols}
\end{enumerate}
\vspace{0.15cm}
\item ¿Qué es una reacción de combustión?
\begin{enumerate}[label=\textbf{\Alph*)}]
\begin{multicols}{2}
  \item La que reacciona con oxígeno liberando energía
  \item La que reacciona con hidrógeno
  \item La que reacciona con nitrógeno
  \item La que reacciona con cloro
  \item La que reacciona con azufre
  \item La que reacciona con fósforo
\end{multicols}
\end{enumerate}
\vspace{0.15cm}
\item ¿Cuál es un ejemplo de reacción de combustión?
\begin{enumerate}[label=\textbf{\Alph*)}]
\begin{multicols}{2}
  \item CH₄ + 2O₂ → CO₂ + 2H₂O
  \item 2H₂ + O₂ → 2H₂O
  \item 2H₂O → 2H₂ + O₂
  \item Zn + CuSO₄ → ZnSO₄ + Cu
  \item AgNO₃ + NaCl → AgCl + NaNO₃
  \item HCl + NaOH → NaCl + H₂O
\end{multicols}
\end{enumerate}
\vspace{0.15cm}
\item ¿Qué es una reacción de neutralización?
\begin{enumerate}[label=\textbf{\Alph*)}]
\begin{multicols}{2}
  \item La que reacciona un ácido con una base para formar sal y agua
  \item La que reacciona un ácido con un metal
  \item La que reacciona una base con un no metal
  \item La que reacciona dos ácidos
  \item La que reacciona dos bases
  \item La que reacciona dos sales
\end{multicols}
\end{enumerate}
\vspace{0.15cm}
\item ¿Cuál es un ejemplo de reacción de neutralización?
\begin{enumerate}[label=\textbf{\Alph*)}]
\begin{multicols}{2}
  \item HCl + NaOH → NaCl + H₂O
  \item 2H₂ + O₂ → 2H₂O
  \item 2H₂O → 2H₂ + O₂
  \item Zn + CuSO₄ → ZnSO₄ + Cu
  \item AgNO₃ + NaCl → AgCl + NaNO₃
  \item CH₄ + 2O₂ → CO₂ + 2H₂O
\end{multicols}
\end{enumerate}
\vspace{0.15cm}
\item ¿Qué es una reacción exotérmica?
\begin{enumerate}[label=\textbf{\Alph*)}]
\begin{multicols}{2}
  \item La que libera energía en forma de calor
  \item La que absorbe energía en forma de calor
  \item La que no intercambia energía
  \item La que libera energía en forma de luz
  \item La que absorbe energía en forma de luz
  \item La que libera energía en forma de sonido
\end{multicols}
\end{enumerate}
\vspace{0.15cm}
\item ¿Qué es una reacción endotérmica?
\begin{enumerate}[label=\textbf{\Alph*)}]
\begin{multicols}{2}
  \item La que absorbe energía en forma de calor
  \item La que libera energía en forma de calor
  \item La que no intercambia energía
  \item La que libera energía en forma de luz
  \item La que absorbe energía en forma de luz
  \item La que libera energía en forma de sonido
\end{multicols}
\end{enumerate}
\vspace{0.15cm}
\item ¿Qué es una reacción reversible?
\begin{enumerate}[label=\textbf{\Alph*)}]
\begin{multicols}{2}
  \item La que puede ocurrir en ambos sentidos
  \item La que ocurre en un solo sentido
  \item La que no ocurre
  \item La que ocurre solo en presencia de catalizador
  \item La que ocurre solo a altas temperaturas
  \item La que ocurre solo a bajas temperaturas
\end{multicols}
\end{enumerate}
\vspace{0.15cm}
\item ¿Qué es una reacción irreversible?
\begin{enumerate}[label=\textbf{\Alph*)}]
\begin{multicols}{2}
  \item La que ocurre en un solo sentido
  \item La que ocurre en ambos sentidos
  \item La que no ocurre
  \item La que ocurre solo en presencia de catalizador
  \item La que ocurre solo a altas temperaturas
  \item La que ocurre solo a bajas temperaturas
\end{multicols}
\end{enumerate}
\vspace{0.15cm}
\item ¿Qué es una reacción de precipitación?
\begin{enumerate}[label=\textbf{\Alph*)}]
\begin{multicols}{2}
  \item La que forma un sólido insoluble
  \item La que forma un gas
  \item La que forma un líquido
  \item La que forma un sólido soluble
  \item La que forma un líquido inmiscible
  \item La que forma un gas soluble
\end{multicols}
\end{enumerate}
\vspace{0.15cm}
\item ¿Cuál es un ejemplo de reacción de precipitación?
\begin{enumerate}[label=\textbf{\Alph*)}]
\begin{multicols}{2}
  \item AgNO₃ + NaCl → AgCl↓ + NaNO₃
  \item 2H₂ + O₂ → 2H₂O
  \item 2H₂O → 2H₂ + O₂
  \item Zn + CuSO₄ → ZnSO₄ + Cu
  \item HCl + NaOH → NaCl + H₂O
  \item CH₄ + 2O₂ → CO₂ + 2H₂O
\end{multicols}
\end{enumerate}
\vspace{0.15cm}
\item ¿Qué es una reacción de formación de gas?
\begin{enumerate}[label=\textbf{\Alph*)}]
\begin{multicols}{2}
  \item La que produce un gas como producto
  \item La que produce un sólido
  \item La que produce un líquido
  \item La que produce un sólido soluble
  \item La que produce un líquido inmiscible
  \item La que produce un gas soluble
\end{multicols}
\end{enumerate}
\vspace{0.15cm}
\item ¿Cuál es un ejemplo de reacción de formación de gas?
\begin{enumerate}[label=\textbf{\Alph*)}]
\begin{multicols}{2}
  \item CaCO₃ + 2HCl → CaCl₂ + H₂O + CO₂↑
  \item 2H₂ + O₂ → 2H₂O
  \item 2H₂O → 2H₂ + O₂
  \item Zn + CuSO₄ → ZnSO₄ + Cu
  \item AgNO₃ + NaCl → AgCl↓ + NaNO₃
  \item HCl + NaOH → NaCl + H₂O ---
\end{multicols}
\end{enumerate}
\vspace{0.15cm}
\end{enumerate}\n
\subsubsection*{SECCIÓN IV: ESTEQUIOMETRÍA Y CÁLCULOS QUÍMICOS (201–250)}

\begin{enumerate}[label=\textbf{\arabic*.}]
\setcounter{enumi}{200}
\item ¿Qué es la estequiometría?
\begin{enumerate}[label=\textbf{\Alph*)}]
\begin{multicols}{2}
  \item El cálculo de las relaciones cuantitativas en las reacciones químicas
  \item El cálculo de las relaciones cualitativas
  \item El cálculo de las propiedades físicas
  \item El cálculo de las propiedades químicas
  \item El cálculo de las estructuras moleculares
  \item El cálculo de las energías de enlace
\end{multicols}
\end{enumerate}
\vspace{0.15cm}
\item ¿Qué es un mol?
\begin{enumerate}[label=\textbf{\Alph*)}]
\begin{multicols}{2}
  \item La cantidad de sustancia que contiene 6.022×10²³ partículas
  \item La cantidad de sustancia que contiene 6.022×10²² partículas
  \item La cantidad de sustancia que contiene 6.022×10²⁴ partículas
  \item La cantidad de sustancia que contiene 6.022×10²¹ partículas
  \item La cantidad de sustancia que contiene 6.022×10²⁰ partículas
  \item La cantidad de sustancia que contiene 6.022×10²⁵ partículas
\end{multicols}
\end{enumerate}
\vspace{0.15cm}
\item ¿Qué es el número de Avogadro?
\begin{enumerate}[label=\textbf{\Alph*)}]
\begin{multicols}{2}
  \item 6.022×10²³
  \item 6.022×10²²
  \item 6.022×10²⁴
  \item 6.022×10²¹
  \item 6.022×10²⁰
  \item 6.022×10²⁵
\end{multicols}
\end{enumerate}
\vspace{0.15cm}
\item ¿Qué es la masa molar?
\begin{enumerate}[label=\textbf{\Alph*)}]
\begin{multicols}{2}
  \item La masa de un mol de sustancia
  \item La masa de un átomo de sustancia
  \item La masa de una molécula de sustancia
  \item La masa de un ion de sustancia
  \item La masa de un isótopo de sustancia
  \item La masa de un electrón de sustancia
\end{multicols}
\end{enumerate}
\vspace{0.15cm}
\item ¿Cuál es la masa molar del agua (H₂O)?
\begin{enumerate}[label=\textbf{\Alph*)}]
\begin{multicols}{2}
  \item 18 g/mol
  \item 16 g/mol
  \item 20 g/mol
  \item 22 g/mol
  \item 24 g/mol
  \item 26 g/mol
\end{multicols}
\end{enumerate}
\vspace{0.15cm}
\item ¿Cuál es la masa molar del dióxido de carbono (CO₂)?
\begin{enumerate}[label=\textbf{\Alph*)}]
\begin{multicols}{2}
  \item 44 g/mol
  \item 28 g/mol
  \item 32 g/mol
  \item 40 g/mol
  \item 48 g/mol
  \item 52 g/mol
\end{multicols}
\end{enumerate}
\vspace{0.15cm}
\item ¿Cuál es la masa molar del cloruro de sodio (NaCl)?
\begin{enumerate}[label=\textbf{\Alph*)}]
\begin{multicols}{2}
  \item 58.5 g/mol
  \item 23 g/mol
  \item 35.5 g/mol
  \item 50 g/mol
  \item 60 g/mol
  \item 65 g/mol
\end{multicols}
\end{enumerate}
\vspace{0.15cm}
\item ¿Qué es un reactivo limitante?
\begin{enumerate}[label=\textbf{\Alph*)}]
\begin{multicols}{2}
  \item El que se consume primero en una reacción química
  \item El que se consume último
  \item El que no se consume
  \item El que se produce en la reacción
  \item El que cataliza la reacción
  \item El que inhibe la reacción
\end{multicols}
\end{enumerate}
\vspace{0.15cm}
\item ¿Qué es un reactivo en exceso?
\begin{enumerate}[label=\textbf{\Alph*)}]
\begin{multicols}{2}
  \item El que está en mayor cantidad de la necesaria
  \item El que está en menor cantidad de la necesaria
  \item El que se consume primero
  \item El que no se consume
  \item El que se produce en la reacción
  \item El que cataliza la reacción
\end{multicols}
\end{enumerate}
\vspace{0.15cm}
\item ¿Qué es el rendimiento de una reacción?
\begin{enumerate}[label=\textbf{\Alph*)}]
\begin{multicols}{2}
  \item La relación entre el producto obtenido y el teórico
  \item La relación entre el reactivo y el producto
  \item La relación entre el producto y el reactivo
  \item La relación entre los reactivos
  \item La relación entre los productos
  \item La relación entre el catalizador y el producto
\end{multicols}
\end{enumerate}
\vspace{0.15cm}
\item ¿Qué es un catalizador?
\begin{enumerate}[label=\textbf{\Alph*)}]
\begin{multicols}{2}
  \item Una sustancia que acelera una reacción sin consumirse
  \item Una sustancia que retarda una reacción
  \item Una sustancia que se consume en la reacción
  \item Una sustancia que se produce en la reacción
  \item Una sustancia que inhibe la reacción
  \item Una sustancia que detiene la reacción
\end{multicols}
\end{enumerate}
\vspace{0.15cm}
\item ¿Qué es un inhibidor?
\begin{enumerate}[label=\textbf{\Alph*)}]
\begin{multicols}{2}
  \item Una sustancia que retarda o detiene una reacción
  \item Una sustancia que acelera una reacción
  \item Una sustancia que se consume en la reacción
  \item Una sustancia que se produce en la reacción
  \item Una sustancia que cataliza la reacción
  \item Una sustancia que inicia la reacción
\end{multicols}
\end{enumerate}
\vspace{0.15cm}
\item ¿Qué es la energía de activación?
\begin{enumerate}[label=\textbf{\Alph*)}]
\begin{multicols}{2}
  \item La energía mínima necesaria para iniciar una reacción
  \item La energía liberada en una reacción
  \item La energía absorbida en una reacción
  \item La energía total de una reacción
  \item La energía de los reactivos
  \item La energía de los productos
\end{multicols}
\end{enumerate}
\vspace{0.15cm}
\item ¿Qué es una reacción exotérmica?
\begin{enumerate}[label=\textbf{\Alph*)}]
\begin{multicols}{2}
  \item La que libera energía en forma de calor
  \item La que absorbe energía en forma de calor
  \item La que no intercambia energía
  \item La que libera energía en forma de luz
  \item La que absorbe energía en forma de luz
  \item La que libera energía en forma de sonido
\end{multicols}
\end{enumerate}
\vspace{0.15cm}
\item ¿Qué es una reacción endotérmica?
\begin{enumerate}[label=\textbf{\Alph*)}]
\begin{multicols}{2}
  \item La que absorbe energía en forma de calor
  \item La que libera energía en forma de calor
  \item La que no intercambia energía
  \item La que libera energía en forma de luz
  \item La que absorbe energía en forma de luz
  \item La que libera energía en forma de sonido
\end{multicols}
\end{enumerate}
\vspace{0.15cm}
\item ¿Cuál es un ejemplo de reacción exotérmica?
\begin{enumerate}[label=\textbf{\Alph*)}]
\begin{multicols}{2}
  \item La combustión
  \item La fotosíntesis
  \item La electrólisis
  \item La fusión del hielo
  \item La evaporación del agua
  \item La sublimación del yodo
\end{multicols}
\end{enumerate}
\vspace{0.15cm}
\item ¿Cuál es un ejemplo de reacción endotérmica?
\begin{enumerate}[label=\textbf{\Alph*)}]
\begin{multicols}{2}
  \item La fotosíntesis
  \item La combustión
  \item La respiración celular
  \item La oxidación del hierro
  \item La fermentación
  \item La corrosión
\end{multicols}
\end{enumerate}
\vspace{0.15cm}
\item ¿Qué es el equilibrio químico?
\begin{enumerate}[label=\textbf{\Alph*)}]
\begin{multicols}{2}
  \item El estado en que las velocidades de reacción directa e inversa son iguales
  \item El estado en que la reacción directa es más rápida
  \item El estado en que la reacción inversa es más rápida
  \item El estado en que la reacción se detiene
  \item El estado en que no hay reacción
  \item El estado en que solo hay reactivos
\end{multicols}
\end{enumerate}
\vspace{0.15cm}
\item ¿Qué es el principio de Le Chatelier?
\begin{enumerate}[label=\textbf{\Alph*)}]
\begin{multicols}{2}
  \item Si se altera un sistema en equilibrio, este se desplaza para contrarrestar el cambio
  \item Si se altera un sistema en equilibrio, este no cambia
  \item Si se altera un sistema en equilibrio, este se destruye
  \item Si se altera un sistema en equilibrio, este se acelera
  \item Si se altera un sistema en equilibrio, este se detiene
  \item Si se altera un sistema en equilibrio, este se invierte
\end{multicols}
\end{enumerate}
\vspace{0.15cm}
\item ¿Qué es un ácido según Arrhenius?
\begin{enumerate}[label=\textbf{\Alph*)}]
\begin{multicols}{2}
  \item Una sustancia que libera iones H⁺ en agua
  \item Una sustancia que libera iones OH⁻ en agua
  \item Una sustancia que acepta protones
  \item Una sustancia que dona electrones
  \item Una sustancia que acepta electrones
  \item Una sustancia que dona neutrones
\end{multicols}
\end{enumerate}
\vspace{0.15cm}
\item ¿Qué es una base según Arrhenius?
\begin{enumerate}[label=\textbf{\Alph*)}]
\begin{multicols}{2}
  \item Una sustancia que libera iones OH⁻ en agua
  \item Una sustancia que libera iones H⁺ en agua
  \item Una sustancia que dona protones
  \item Una sustancia que acepta electrones
  \item Una sustancia que dona electrones
  \item Una sustancia que acepta neutrones
\end{multicols}
\end{enumerate}
\vspace{0.15cm}
\item ¿Qué es un ácido según Brønsted-Lowry?
\begin{enumerate}[label=\textbf{\Alph*)}]
\begin{multicols}{2}
  \item Una sustancia que dona protones
  \item Una sustancia que acepta protones
  \item Una sustancia que libera iones OH⁻
  \item Una sustancia que libera iones H⁺
  \item Una sustancia que acepta electrones
  \item Una sustancia que dona electrones
\end{multicols}
\end{enumerate}
\vspace{0.15cm}
\item ¿Qué es una base según Brønsted-Lowry?
\begin{enumerate}[label=\textbf{\Alph*)}]
\begin{multicols}{2}
  \item Una sustancia que acepta protones
  \item Una sustancia que dona protones
  \item Una sustancia que libera iones OH⁻
  \item Una sustancia que libera iones H⁺
  \item Una sustancia que acepta electrones
  \item Una sustancia que dona electrones
\end{multicols}
\end{enumerate}
\vspace{0.15cm}
\item ¿Qué es el pH?
\begin{enumerate}[label=\textbf{\Alph*)}]
\begin{multicols}{2}
  \item La medida de la acidez o basicidad de una disolución
  \item La medida de la temperatura de una disolución
  \item La medida de la densidad de una disolución
  \item La medida de la presión de una disolución
  \item La medida del volumen de una disolución
  \item La medida de la masa de una disolución
\end{multicols}
\end{enumerate}
\vspace{0.15cm}
\item ¿Cuál es el pH de una disolución neutra?
\begin{enumerate}[label=\textbf{\Alph*)}]
\begin{multicols}{2}
  \item 7
  \item 0
  \item 14
  \item 1
  \item 13
  \item 10
\end{multicols}
\end{enumerate}
\vspace{0.15cm}
\item ¿Cuál es el pH de una disolución ácida?
\begin{enumerate}[label=\textbf{\Alph*)}]
\begin{multicols}{2}
  \item Menor que 7
  \item Mayor que 7
  \item Igual a 7
  \item Igual a 0
  \item Igual a 14
  \item Mayor que 14
\end{multicols}
\end{enumerate}
\vspace{0.15cm}
\item ¿Cuál es el pH de una disolución básica?
\begin{enumerate}[label=\textbf{\Alph*)}]
\begin{multicols}{2}
  \item Mayor que 7
  \item Menor que 7
  \item Igual a 7
  \item Igual a 0
  \item Igual a 14
  \item Menor que 0
\end{multicols}
\end{enumerate}
\vspace{0.15cm}
\item ¿Qué es una reacción de neutralización?
\begin{enumerate}[label=\textbf{\Alph*)}]
\begin{multicols}{2}
  \item La reacción entre un ácido y una base para formar sal y agua
  \item La reacción entre dos ácidos
  \item La reacción entre dos bases
  \item La reacción entre un ácido y un metal
  \item La reacción entre una base y un no metal
  \item La reacción entre dos sales
\end{multicols}
\end{enumerate}
\vspace{0.15cm}
\item ¿Cuál es el producto de la reacción entre HCl y NaOH?
\begin{enumerate}[label=\textbf{\Alph*)}]
\begin{multicols}{2}
  \item NaCl y H₂O
  \item NaCl y H₂
  \item NaCl y O₂
  \item NaH y Cl₂
  \item Na y HCl
  \item NaOH y HCl
\end{multicols}
\end{enumerate}
\vspace{0.15cm}
\item ¿Qué es un indicador ácido-base?
\begin{enumerate}[label=\textbf{\Alph*)}]
\begin{multicols}{2}
  \item Una sustancia que cambia de color según el pH
  \item Una sustancia que cambia de temperatura
  \item Una sustancia que cambia de densidad
  \item Una sustancia que cambia de presión
  \item Una sustancia que cambia de volumen
  \item Una sustancia que cambia de masa
\end{multicols}
\end{enumerate}
\vspace{0.15cm}
\item ¿Cuál es un ejemplo de indicador ácido-base?
\begin{enumerate}[label=\textbf{\Alph*)}]
\begin{multicols}{2}
  \item La fenolftaleína
  \item El agua
  \item El alcohol
  \item La sal
  \item El azúcar
  \item El aceite
\end{multicols}
\end{enumerate}
\vspace{0.15cm}
\item ¿Qué es una sal?
\begin{enumerate}[label=\textbf{\Alph*)}]
\begin{multicols}{2}
  \item Un compuesto formado por un catión y un anión
  \item Un compuesto formado por dos cationes
  \item Un compuesto formado por dos aniones
  \item Un compuesto formado por un ácido y una base
  \item Un compuesto formado por dos ácidos
  \item Un compuesto formado por dos bases
\end{multicols}
\end{enumerate}
\vspace{0.15cm}
\item ¿Qué es la electrólisis?
\begin{enumerate}[label=\textbf{\Alph*)}]
\begin{multicols}{2}
  \item El proceso de descomposición de una sustancia mediante electricidad
  \item El proceso de descomposición de una sustancia mediante calor
  \item El proceso de descomposición de una sustancia mediante luz
  \item El proceso de descomposición de una sustancia mediante presión
  \item El proceso de descomposición de una sustancia mediante sonido
  \item El proceso de descomposición de una sustancia mediante radiación
\end{multicols}
\end{enumerate}
\vspace{0.15cm}
\item ¿Qué es una pila?
\begin{enumerate}[label=\textbf{\Alph*)}]
\begin{multicols}{2}
  \item Un dispositivo que convierte energía química en eléctrica
  \item Un dispositivo que convierte energía eléctrica en química
  \item Un dispositivo que convierte energía mecánica en eléctrica
  \item Un dispositivo que convierte energía térmica en eléctrica
  \item Un dispositivo que convierte energía solar en eléctrica
  \item Un dispositivo que convierte energía nuclear en eléctrica
\end{multicols}
\end{enumerate}
\vspace{0.15cm}
\item ¿Qué es la corrosión?
\begin{enumerate}[label=\textbf{\Alph*)}]
\begin{multicols}{2}
  \item La degradación de un metal por reacción química con el ambiente
  \item La degradación de un metal por reacción física
  \item La degradación de un metal por reacción mecánica
  \item La degradación de un metal por reacción térmica
  \item La degradación de un metal por reacción eléctrica
  \item La degradación de un metal por reacción nuclear
\end{multicols}
\end{enumerate}
\vspace{0.15cm}
\item ¿Qué es la oxidación?
\begin{enumerate}[label=\textbf{\Alph*)}]
\begin{multicols}{2}
  \item La pérdida de electrones por parte de una sustancia
  \item La ganancia de electrones por parte de una sustancia
  \item La pérdida de protones por parte de una sustancia
  \item La ganancia de protones por parte de una sustancia
  \item La pérdida de neutrones por parte de una sustancia
  \item La ganancia de neutrones por parte de una sustancia
\end{multicols}
\end{enumerate}
\vspace{0.15cm}
\item ¿Qué es la reducción?
\begin{enumerate}[label=\textbf{\Alph*)}]
\begin{multicols}{2}
  \item La ganancia de electrones por parte de una sustancia
  \item La pérdida de electrones por parte de una sustancia
  \item La ganancia de protones por parte de una sustancia
  \item La pérdida de protones por parte de una sustancia
  \item La ganancia de neutrones por parte de una sustancia
  \item La pérdida de neutrones por parte de una sustancia
\end{multicols}
\end{enumerate}
\vspace{0.15cm}
\item ¿Qué es una reacción redox?
\begin{enumerate}[label=\textbf{\Alph*)}]
\begin{multicols}{2}
  \item La que implica transferencia de electrones
  \item La que implica transferencia de protones
  \item La que implica transferencia de neutrones
  \item La que implica transferencia de átomos
  \item La que implica transferencia de moléculas
  \item La que implica transferencia de iones
\end{multicols}
\end{enumerate}
\vspace{0.15cm}
\item ¿Qué es el agente oxidante?
\begin{enumerate}[label=\textbf{\Alph*)}]
\begin{multicols}{2}
  \item La sustancia que gana electrones y se reduce
  \item La sustancia que pierde electrones y se oxida
  \item La sustancia que gana protones
  \item La sustancia que pierde protones
  \item La sustancia que gana neutrones
  \item La sustancia que pierde neutrones
\end{multicols}
\end{enumerate}
\vspace{0.15cm}
\item ¿Qué es el agente reductor?
\begin{enumerate}[label=\textbf{\Alph*)}]
\begin{multicols}{2}
  \item La sustancia que pierde electrones y se oxida
  \item La sustancia que gana electrones y se reduce
  \item La sustancia que gana protones
  \item La sustancia que pierde protones
  \item La sustancia que gana neutrones
  \item La sustancia que pierde neutrones
\end{multicols}
\end{enumerate}
\vspace{0.15cm}
\item ¿Qué es el número de oxidación?
\begin{enumerate}[label=\textbf{\Alph*)}]
\begin{multicols}{2}
  \item La carga que tendría un átomo si los electrones se asignaran al átomo más electronegativo
  \item La carga que tendría un átomo si los electrones se asignaran al átomo menos electronegativo
  \item La carga real de un átomo en un compuesto
  \item La carga de un ion en disolución
  \item La carga de un electrón
  \item La carga de un protón
\end{multicols}
\end{enumerate}
\vspace{0.15cm}
\item ¿Qué es una reacción de síntesis o combinación?
\begin{enumerate}[label=\textbf{\Alph*)}]
\begin{multicols}{2}
  \item La que forma un compuesto a partir de dos o más sustancias
  \item La que descompone un compuesto en dos o más sustancias
  \item La que sustituye un elemento por otro
  \item La que intercambia iones entre compuestos
  \item La que neutraliza un ácido y una base
  \item La que combustiona un compuesto
\end{multicols}
\end{enumerate}
\vspace{0.15cm}
\item ¿Qué es una reacción de descomposición?
\begin{enumerate}[label=\textbf{\Alph*)}]
\begin{multicols}{2}
  \item La que descompone un compuesto en dos o más sustancias
  \item La que forma un compuesto a partir de dos o más sustancias
  \item La que sustituye un elemento por otro
  \item La que intercambia iones entre compuestos
  \item La que neutraliza un ácido y una base
  \item La que combustiona un compuesto
\end{multicols}
\end{enumerate}
\vspace{0.15cm}
\item ¿Cuál es el coeficiente correcto para balancear: __H₂ + O₂ → __H₂O?
\begin{enumerate}[label=\textbf{\Alph*)}]
\begin{multicols}{2}
  \item 2, 1, 2
  \item 1, 1, 1
  \item 2, 2, 2
  \item 1, 2, 1
  \item 2, 1, 1
  \item 1, 1, 2
\end{multicols}
\end{enumerate}
\vspace{0.15cm}
\item ¿Cuál es el coeficiente correcto para balancear: __N₂ + __H₂ → __NH₃?
\begin{enumerate}[label=\textbf{\Alph*)}]
\begin{multicols}{2}
  \item 1, 3, 2
  \item 1, 2, 3
  \item 2, 3, 1
  \item 3, 1, 2
  \item 2, 1, 3
  \item 1, 1, 1
\end{multicols}
\end{enumerate}
\vspace{0.15cm}
\item ¿Cuál es el coeficiente correcto para balancear: __CH₄ + __O₂ → __CO₂ + __H₂O?
\begin{enumerate}[label=\textbf{\Alph*)}]
\begin{multicols}{2}
  \item 1, 2, 1, 2
  \item 1, 1, 1, 2
  \item 1, 2, 2, 1
  \item 2, 1, 1, 2
  \item 1, 3, 1, 2
  \item 2, 2, 1, 1
\end{multicols}
\end{enumerate}
\vspace{0.15cm}
\item ¿Cuál es el coeficiente correcto para balancear: __Fe + __O₂ → __Fe₂O₃?
\begin{enumerate}[label=\textbf{\Alph*)}]
\begin{multicols}{2}
  \item 4, 3, 2
  \item 2, 3, 1
  \item 3, 4, 2
  \item 2, 1, 1
  \item 4, 2, 2
  \item 1, 1, 1
\end{multicols}
\end{enumerate}
\vspace{0.15cm}
\item ¿Cuál es el coeficiente correcto para balancear: __Al + __HCl → __AlCl₃ + __H₂?
\begin{enumerate}[label=\textbf{\Alph*)}]
\begin{multicols}{2}
  \item 2, 6, 2, 3
  \item 1, 3, 1, 1
  \item 2, 3, 2, 3
  \item 1, 6, 1, 3
  \item 2, 6, 3, 2
  \item 3, 6, 2, 2
\end{multicols}
\end{enumerate}
\vspace{0.15cm}
\item ¿Qué es una ecuación iónica?
\begin{enumerate}[label=\textbf{\Alph*)}]
\begin{multicols}{2}
  \item La que muestra los iones que participan en la reacción
  \item La que muestra las moléculas
  \item La que muestra los átomos
  \item La que muestra los electrones
  \item La que muestra los protones
  \item La que muestra los neutrones
\end{multicols}
\end{enumerate}
\vspace{0.15cm}
\item ¿Qué es una ecuación iónica neta?
\begin{enumerate}[label=\textbf{\Alph*)}]
\begin{multicols}{2}
  \item La que muestra solo los iones que reaccionan
  \item La que muestra todos los iones
  \item La que muestra las moléculas
  \item La que muestra los átomos
  \item La que muestra los electrones
  \item La que muestra los protones ---
\end{multicols}
\end{enumerate}
\vspace{0.15cm}
\end{enumerate}\n
\subsubsection*{SECCIÓN V: QUÍMICA ORGÁNICA E INORGÁNICA (251–300)}

\begin{enumerate}[label=\textbf{\arabic*.}]
\setcounter{enumi}{250}
\item ¿Qué es la química orgánica?
\begin{enumerate}[label=\textbf{\Alph*)}]
\begin{multicols}{2}
  \item La que estudia los compuestos del carbono
  \item La que estudia los compuestos inorgánicos
  \item La que estudia los metales
  \item La que estudia los no metales
  \item La que estudia los gases nobles
  \item La que estudia los minerales
\end{multicols}
\end{enumerate}
\vspace{0.15cm}
\item ¿Qué es la química inorgánica?
\begin{enumerate}[label=\textbf{\Alph*)}]
\begin{multicols}{2}
  \item La que estudia los compuestos sin carbono (excepto algunos)
  \item La que estudia los compuestos del carbono
  \item La que estudia los seres vivos
  \item La que estudia los minerales
  \item La que estudia los metales
  \item La que estudia los gases
\end{multicols}
\end{enumerate}
\vspace{0.15cm}
\item ¿Cuál es el elemento central de la química orgánica?
\begin{enumerate}[label=\textbf{\Alph*)}]
\begin{multicols}{2}
  \item El carbono
  \item El hidrógeno
  \item El oxígeno
  \item El nitrógeno
  \item El azufre
  \item El fósforo
\end{multicols}
\end{enumerate}
\vspace{0.15cm}
\item ¿Cuántos enlaces puede formar el carbono?
\begin{enumerate}[label=\textbf{\Alph*)}]
\begin{multicols}{2}
  \item 4
  \item 2
  \item 3
  \item 5
  \item 6
  \item 1
\end{multicols}
\end{enumerate}
\vspace{0.15cm}
\item ¿Qué es un hidrocarburo?
\begin{enumerate}[label=\textbf{\Alph*)}]
\begin{multicols}{2}
  \item Un compuesto formado por carbono e hidrógeno
  \item Un compuesto formado por carbono y oxígeno
  \item Un compuesto formado por carbono y nitrógeno
  \item Un compuesto formado por carbono y azufre
  \item Un compuesto formado por carbono y fósforo
  \item Un compuesto formado por carbono y cloro
\end{multicols}
\end{enumerate}
\vspace{0.15cm}
\item ¿Qué es un alcano?
\begin{enumerate}[label=\textbf{\Alph*)}]
\begin{multicols}{2}
  \item Un hidrocarburo saturado con enlaces simples
  \item Un hidrocarburo insaturado con dobles enlaces
  \item Un hidrocarburo insaturado con triples enlaces
  \item Un hidrocarburo aromático
  \item Un hidrocarburo cíclico
  \item Un hidrocarburo ramificado
\end{multicols}
\end{enumerate}
\vspace{0.15cm}
\item ¿Qué es un alqueno?
\begin{enumerate}[label=\textbf{\Alph*)}]
\begin{multicols}{2}
  \item Un hidrocarburo insaturado con dobles enlaces
  \item Un hidrocarburo saturado con enlaces simples
  \item Un hidrocarburo insaturado con triples enlaces
  \item Un hidrocarburo aromático
  \item Un hidrocarburo cíclico
  \item Un hidrocarburo ramificado
\end{multicols}
\end{enumerate}
\vspace{0.15cm}
\item ¿Qué es un alquino?
\begin{enumerate}[label=\textbf{\Alph*)}]
\begin{multicols}{2}
  \item Un hidrocarburo insaturado con triples enlaces
  \item Un hidrocarburo saturado con enlaces simples
  \item Un hidrocarburo insaturado con dobles enlaces
  \item Un hidrocarburo aromático
  \item Un hidrocarburo cíclico
  \item Un hidrocarburo ramificado
\end{multicols}
\end{enumerate}
\vspace{0.15cm}
\item ¿Cuál es la fórmula general de los alcanos?
\begin{enumerate}[label=\textbf{\Alph*)}]
\begin{multicols}{2}
  \item CₙH₂ₙ₊₂
  \item CₙH₂ₙ
  \item CₙH₂ₙ₋₂
  \item CₙHₙ
  \item CₙH₂ₙ₊₁
  \item CₙH₂ₙ₋₁
\end{multicols}
\end{enumerate}
\vspace{0.15cm}
\item ¿Cuál es la fórmula general de los alquenos?
\begin{enumerate}[label=\textbf{\Alph*)}]
\begin{multicols}{2}
  \item CₙH₂ₙ
  \item CₙH₂ₙ₊₂
  \item CₙH₂ₙ₋₂
  \item CₙHₙ
  \item CₙH₂ₙ₊₁
  \item CₙH₂ₙ₋₁
\end{multicols}
\end{enumerate}
\vspace{0.15cm}
\item ¿Cuál es la fórmula general de los alquinos?
\begin{enumerate}[label=\textbf{\Alph*)}]
\begin{multicols}{2}
  \item CₙH₂ₙ₋₂
  \item CₙH₂ₙ₊₂
  \item CₙH₂ₙ
  \item CₙHₙ
  \item CₙH₂ₙ₊₁
  \item CₙH₂ₙ₋₁
\end{multicols}
\end{enumerate}
\vspace{0.15cm}
\item ¿Qué es un grupo funcional?
\begin{enumerate}[label=\textbf{\Alph*)}]
\begin{multicols}{2}
  \item Un átomo o grupo de átomos que determina las propiedades químicas
  \item Un átomo o grupo de átomos que determina las propiedades físicas
  \item Un átomo o grupo de átomos que determina la masa
  \item Un átomo o grupo de átomos que determina el volumen
  \item Un átomo o grupo de átomos que determina la densidad
  \item Un átomo o grupo de átomos que determina la temperatura
\end{multicols}
\end{enumerate}
\vspace{0.15cm}
\item ¿Cuál es el grupo funcional de los alcoholes?
\begin{enumerate}[label=\textbf{\Alph*)}]
\begin{multicols}{2}
  \item -OH (hidroxilo)
  \item -COOH (carboxilo)
  \item -CHO (aldehído)
  \item -CO- (cetona)
  \item -NH₂ (amino)
  \item -O- (éter)
\end{multicols}
\end{enumerate}
\vspace{0.15cm}
\item ¿Cuál es el grupo funcional de los ácidos carboxílicos?
\begin{enumerate}[label=\textbf{\Alph*)}]
\begin{multicols}{2}
  \item -COOH (carboxilo)
  \item -OH (hidroxilo)
  \item -CHO (aldehído)
  \item -CO- (cetona)
  \item -NH₂ (amino)
  \item -O- (éter)
\end{multicols}
\end{enumerate}
\vspace{0.15cm}
\item ¿Cuál es el grupo funcional de los aldehídos?
\begin{enumerate}[label=\textbf{\Alph*)}]
\begin{multicols}{2}
  \item -CHO (aldehído)
  \item -OH (hidroxilo)
  \item -COOH (carboxilo)
  \item -CO- (cetona)
  \item -NH₂ (amino)
  \item -O- (éter)
\end{multicols}
\end{enumerate}
\vspace{0.15cm}
\item ¿Cuál es el grupo funcional de las cetonas?
\begin{enumerate}[label=\textbf{\Alph*)}]
\begin{multicols}{2}
  \item -CO- (cetona)
  \item -OH (hidroxilo)
  \item -COOH (carboxilo)
  \item -CHO (aldehído)
  \item -NH₂ (amino)
  \item -O- (éter)
\end{multicols}
\end{enumerate}
\vspace{0.15cm}
\item ¿Cuál es el grupo funcional de las aminas?
\begin{enumerate}[label=\textbf{\Alph*)}]
\begin{multicols}{2}
  \item -NH₂ (amino)
  \item -OH (hidroxilo)
  \item -COOH (carboxilo)
  \item -CHO (aldehído)
  \item -CO- (cetona)
  \item -O- (éter)
\end{multicols}
\end{enumerate}
\vspace{0.15cm}
\item ¿Cuál es el grupo funcional de los éteres?
\begin{enumerate}[label=\textbf{\Alph*)}]
\begin{multicols}{2}
  \item -O- (éter)
  \item -OH (hidroxilo)
  \item -COOH (carboxilo)
  \item -CHO (aldehído)
  \item -CO- (cetona)
  \item -NH₂ (amino)
\end{multicols}
\end{enumerate}
\vspace{0.15cm}
\item ¿Qué es un isómero?
\begin{enumerate}[label=\textbf{\Alph*)}]
\begin{multicols}{2}
  \item Un compuesto con la misma fórmula molecular pero diferente estructura
  \item Un compuesto con diferente fórmula molecular
  \item Un compuesto con la misma estructura
  \item Un compuesto con diferentes átomos
  \item Un compuesto con diferentes enlaces
  \item Un compuesto con diferentes propiedades
\end{multicols}
\end{enumerate}
\vspace{0.15cm}
\item ¿Qué es un polímero?
\begin{enumerate}[label=\textbf{\Alph*)}]
\begin{multicols}{2}
  \item Una molécula grande formada por la unión de monómeros
  \item Una molécula pequeña
  \item Un átomo
  \item Un ion
  \item Un isótopo
  \item Un isómero
\end{multicols}
\end{enumerate}
\vspace{0.15cm}
\item ¿Cuál es un ejemplo de polímero natural?
\begin{enumerate}[label=\textbf{\Alph*)}]
\begin{multicols}{2}
  \item El almidón
  \item El polietileno
  \item El PVC
  \item El nylon
  \item El poliestireno
  \item El teflón
\end{multicols}
\end{enumerate}
\vspace{0.15cm}
\item ¿Cuál es un ejemplo de polímero sintético?
\begin{enumerate}[label=\textbf{\Alph*)}]
\begin{multicols}{2}
  \item El polietileno
  \item El almidón
  \item La celulosa
  \item La proteína
  \item El ADN
  \item El caucho natural
\end{multicols}
\end{enumerate}
\vspace{0.15cm}
\item ¿Qué es la fermentación?
\begin{enumerate}[label=\textbf{\Alph*)}]
\begin{multicols}{2}
  \item Un proceso anaerobio de transformación de azúcares
  \item Un proceso aerobio de transformación de azúcares
  \item Un proceso de oxidación de grasas
  \item Un proceso de descomposición de proteínas
  \item Un proceso de síntesis de lípidos
  \item Un proceso de síntesis de carbohidratos
\end{multicols}
\end{enumerate}
\vspace{0.15cm}
\item ¿Qué es la saponificación?
\begin{enumerate}[label=\textbf{\Alph*)}]
\begin{multicols}{2}
  \item La reacción de un ácido graso con una base para formar jabón
  \item La reacción de un ácido con un alcohol
  \item La reacción de un ácido con un metal
  \item La reacción de una base con un no metal
  \item La reacción de dos ácidos
  \item La reacción de dos bases
\end{multicols}
\end{enumerate}
\vspace{0.15cm}
\item ¿Qué es un jabón?
\begin{enumerate}[label=\textbf{\Alph*)}]
\begin{multicols}{2}
  \item Una sal de ácido graso
  \item Un ácido graso
  \item Un alcohol
  \item Un éster
  \item Una amida
  \item Una amina
\end{multicols}
\end{enumerate}
\vspace{0.15cm}
\item ¿Qué es un aminoácido?
\begin{enumerate}[label=\textbf{\Alph*)}]
\begin{multicols}{2}
  \item Un compuesto con grupo amino y grupo carboxilo
  \item Un compuesto con grupo hidroxilo y grupo carbonilo
  \item Un compuesto con grupo éter y grupo éster
  \item Un compuesto con grupo aldehído y grupo cetona
  \item Un compuesto con grupo nitro y grupo sulfónico
  \item Un compuesto con grupo fosfato y grupo sulfato
\end{multicols}
\end{enumerate}
\vspace{0.15cm}
\item ¿Qué es un enlace peptídico?
\begin{enumerate}[label=\textbf{\Alph*)}]
\begin{multicols}{2}
  \item El enlace entre dos aminoácidos
  \item El enlace entre dos monosacáridos
  \item El enlace entre dos ácidos grasos
  \item El enlace entre dos nucleótidos
  \item El enlace entre dos lípidos
  \item El enlace entre dos carbohidratos
\end{multicols}
\end{enumerate}
\vspace{0.15cm}
\item ¿Qué es una proteína?
\begin{enumerate}[label=\textbf{\Alph*)}]
\begin{multicols}{2}
  \item Un polímero de aminoácidos
  \item Un polímero de monosacáridos
  \item Un polímero de ácidos grasos
  \item Un polímero de nucleótidos
  \item Un polímero de lípidos
  \item Un polímero de carbohidratos
\end{multicols}
\end{enumerate}
\vspace{0.15cm}
\item ¿Qué es un carbohidrato?
\begin{enumerate}[label=\textbf{\Alph*)}]
\begin{multicols}{2}
  \item Un compuesto formado por carbono, hidrógeno y oxígeno
  \item Un compuesto formado por carbono, hidrógeno y nitrógeno
  \item Un compuesto formado por carbono, hidrógeno y azufre
  \item Un compuesto formado por carbono, hidrógeno y fósforo
  \item Un compuesto formado por carbono, hidrógeno y cloro
  \item Un compuesto formado por carbono, hidrógeno y flúor
\end{multicols}
\end{enumerate}
\vspace{0.15cm}
\item ¿Qué es un lípido?
\begin{enumerate}[label=\textbf{\Alph*)}]
\begin{multicols}{2}
  \item Un compuesto insoluble en agua y soluble en disolventes orgánicos
  \item Un compuesto soluble en agua
  \item Un compuesto soluble en agua y en disolventes orgánicos
  \item Un compuesto insoluble en agua y en disolventes orgánicos
  \item Un compuesto soluble en agua e insoluble en disolventes orgánicos
  \item Un compuesto insoluble en agua y soluble en agua
\end{multicols}
\end{enumerate}
\vspace{0.15cm}
\item ¿Qué es un nucleótido?
\begin{enumerate}[label=\textbf{\Alph*)}]
\begin{multicols}{2}
  \item Un compuesto formado por base nitrogenada, azúcar y fosfato
  \item Un compuesto formado por aminoácidos
  \item Un compuesto formado por monosacáridos
  \item Un compuesto formado por ácidos grasos
  \item Un compuesto formado por lípidos
  \item Un compuesto formado por carbohidratos
\end{multicols}
\end{enumerate}
\vspace{0.15cm}
\item ¿Qué es el ADN?
\begin{enumerate}[label=\textbf{\Alph*)}]
\begin{multicols}{2}
  \item Un polímero de nucleótidos
  \item Un polímero de aminoácidos
  \item Un polímero de monosacáridos
  \item Un polímero de ácidos grasos
  \item Un polímero de lípidos
  \item Un polímero de carbohidratos
\end{multicols}
\end{enumerate}
\vspace{0.15cm}
\item ¿Qué es el ARN?
\begin{enumerate}[label=\textbf{\Alph*)}]
\begin{multicols}{2}
  \item Un polímero de nucleótidos
  \item Un polímero de aminoácidos
  \item Un polímero de monosacáridos
  \item Un polímero de ácidos grasos
  \item Un polímero de lípidos
  \item Un polímero de carbohidratos
\end{multicols}
\end{enumerate}
\vspace{0.15cm}
\item ¿Qué es una enzima?
\begin{enumerate}[label=\textbf{\Alph*)}]
\begin{multicols}{2}
  \item Una proteína que cataliza reacciones químicas
  \item Un lípido que cataliza reacciones químicas
  \item Un carbohidrato que cataliza reacciones químicas
  \item Un ácido nucleico que cataliza reacciones químicas
  \item Un mineral que cataliza reacciones químicas
  \item Un gas que cataliza reacciones químicas
\end{multicols}
\end{enumerate}
\vspace{0.15cm}
\item ¿Qué es la catálisis?
\begin{enumerate}[label=\textbf{\Alph*)}]
\begin{multicols}{2}
  \item El proceso de acelerar una reacción mediante un catalizador
  \item El proceso de retardar una reacción
  \item El proceso de detener una reacción
  \item El proceso de invertir una reacción
  \item El proceso de iniciar una reacción
  \item El proceso de finalizar una reacción
\end{multicols}
\end{enumerate}
\vspace{0.15cm}
\item ¿Qué es un catalizador biológico?
\begin{enumerate}[label=\textbf{\Alph*)}]
\begin{multicols}{2}
  \item Una enzima
  \item Un ácido
  \item Una base
  \item Una sal
  \item Un óxido
  \item Un metal
\end{multicols}
\end{enumerate}
\vspace{0.15cm}
\item ¿Qué es la química analítica?
\begin{enumerate}[label=\textbf{\Alph*)}]
\begin{multicols}{2}
  \item La que estudia la composición de las sustancias
  \item La que estudia las reacciones químicas
  \item La que estudia las propiedades físicas
  \item La que estudia las propiedades químicas
  \item La que estudia las estructuras moleculares
  \item La que estudia las energías de enlace
\end{multicols}
\end{enumerate}
\vspace{0.15cm}
\item ¿Qué es la química física?
\begin{enumerate}[label=\textbf{\Alph*)}]
\begin{multicols}{2}
  \item La que estudia los principios físicos de los sistemas químicos
  \item La que estudia la composición de las sustancias
  \item La que estudia las reacciones químicas
  \item La que estudia las propiedades químicas
  \item La que estudia las estructuras moleculares
  \item La que estudia las energías de enlace
\end{multicols}
\end{enumerate}
\vspace{0.15cm}
\item ¿Qué es la bioquímica?
\begin{enumerate}[label=\textbf{\Alph*)}]
\begin{multicols}{2}
  \item La que estudia los procesos químicos en los seres vivos
  \item La que estudia la composición de las sustancias
  \item La que estudia las reacciones químicas
  \item La que estudia las propiedades físicas
  \item La que estudia las propiedades químicas
  \item La que estudia las estructuras moleculares
\end{multicols}
\end{enumerate}
\vspace{0.15cm}
\item ¿Qué es la química industrial?
\begin{enumerate}[label=\textbf{\Alph*)}]
\begin{multicols}{2}
  \item La que estudia la producción de sustancias a gran escala
  \item La que estudia la composición de las sustancias
  \item La que estudia las reacciones químicas
  \item La que estudia las propiedades físicas
  \item La que estudia las propiedades químicas
  \item La que estudia las estructuras moleculares
\end{multicols}
\end{enumerate}
\vspace{0.15cm}
\item ¿Qué es la química ambiental?
\begin{enumerate}[label=\textbf{\Alph*)}]
\begin{multicols}{2}
  \item La que estudia los procesos químicos en el medio ambiente
  \item La que estudia la composición de las sustancias
  \item La que estudia las reacciones químicas
  \item La que estudia las propiedades físicas
  \item La que estudia las propiedades químicas
  \item La que estudia las estructuras moleculares
\end{multicols}
\end{enumerate}
\vspace{0.15cm}
\item ¿Qué es la química de polímeros?
\begin{enumerate}[label=\textbf{\Alph*)}]
\begin{multicols}{2}
  \item La que estudia la síntesis y propiedades de los polímeros
  \item La que estudia la composición de las sustancias
  \item La que estudia las reacciones químicas
  \item La que estudia las propiedades físicas
  \item La que estudia las propiedades químicas
  \item La que estudia las estructuras moleculares
\end{multicols}
\end{enumerate}
\vspace{0.15cm}
\item ¿Qué es la química medicinal?
\begin{enumerate}[label=\textbf{\Alph*)}]
\begin{multicols}{2}
  \item La que estudia el diseño y síntesis de fármacos
  \item La que estudia la composición de las sustancias
  \item La que estudia las reacciones químicas
  \item La que estudia las propiedades físicas
  \item La que estudia las propiedades químicas
  \item La que estudia las estructuras moleculares
\end{multicols}
\end{enumerate}
\vspace{0.15cm}
\item ¿Qué es la química de materiales?
\begin{enumerate}[label=\textbf{\Alph*)}]
\begin{multicols}{2}
  \item La que estudia el diseño y síntesis de materiales
  \item La que estudia la composición de las sustancias
  \item La que estudia las reacciones químicas
  \item La que estudia las propiedades físicas
  \item La que estudia las propiedades químicas
  \item La que estudia las estructuras moleculares
\end{multicols}
\end{enumerate}
\vspace{0.15cm}
\item ¿Qué es la química computacional?
\begin{enumerate}[label=\textbf{\Alph*)}]
\begin{multicols}{2}
  \item La que usa computadoras para simular sistemas químicos
  \item La que estudia la composición de las sustancias
  \item La que estudia las reacciones químicas
  \item La que estudia las propiedades físicas
  \item La que estudia las propiedades químicas
  \item La que estudia las estructuras moleculares
\end{multicols}
\end{enumerate}
\vspace{0.15cm}
\item ¿Qué es la química verde?
\begin{enumerate}[label=\textbf{\Alph*)}]
\begin{multicols}{2}
  \item La que busca reducir el impacto ambiental de los procesos químicos
  \item La que estudia la composición de las sustancias
  \item La que estudia las reacciones químicas
  \item La que estudia las propiedades físicas
  \item La que estudia las propiedades químicas
  \item La que estudia las estructuras moleculares
\end{multicols}
\end{enumerate}
\vspace{0.15cm}
\item ¿Qué es la química de alimentos?
\begin{enumerate}[label=\textbf{\Alph*)}]
\begin{multicols}{2}
  \item La que estudia la composición y transformación de los alimentos
  \item La que estudia la composición de las sustancias
  \item La que estudia las reacciones químicas
  \item La que estudia las propiedades físicas
  \item La que estudia las propiedades químicas
  \item La que estudia las estructuras moleculares
\end{multicols}
\end{enumerate}
\vspace{0.15cm}
\item ¿Qué es la química agrícola?
\begin{enumerate}[label=\textbf{\Alph*)}]
\begin{multicols}{2}
  \item La que estudia los procesos químicos en la agricultura
  \item La que estudia la composición de las sustancias
  \item La que estudia las reacciones químicas
  \item La que estudia las propiedades físicas
  \item La que estudia las propiedades químicas
  \item La que estudia las estructuras moleculares
\end{multicols}
\end{enumerate}
\vspace{0.15cm}
\item ¿Qué es la química de los materiales?
\begin{enumerate}[label=\textbf{\Alph*)}]
\begin{multicols}{2}
  \item La que estudia la síntesis y propiedades de los materiales
  \item La que estudia la composición de las sustancias
  \item La que estudia las reacciones químicas
  \item La que estudia las propiedades físicas
  \item La que estudia las propiedades químicas
  \item La que estudia las estructuras moleculares
\end{multicols}
\end{enumerate}
\vspace{0.15cm}
\item ¿Qué es la química teórica?
\begin{enumerate}[label=\textbf{\Alph*)}]
\begin{multicols}{2}
  \item La que usa modelos matemáticos para explicar sistemas químicos
  \item La que estudia la composición de las sustancias
  \item La que estudia las reacciones químicas
  \item La que estudia las propiedades físicas
  \item La que estudia las propiedades químicas
  \item La que estudia las estructuras moleculares ---
\end{multicols}
\end{enumerate}
\vspace{0.15cm}
\end{enumerate}\n
\vspace{0.4cm}
\begin{tcolorbox}[colback=gray!5!white,colframe=gray!60!black,title=\textbf{SOLUCIONARIO -- 5.2 Química General y Reacciones}]
\begin{multicols}{4}
\begin{enumerate}[label=\textbf{\arabic*.}]
  \item \textbf{A}
  \item \textbf{A}
  \item \textbf{A}
  \item \textbf{A}
  \item \textbf{A}
  \item \textbf{A}
  \item \textbf{A}
  \item \textbf{A}
  \item \textbf{A}
  \item \textbf{A}
  \item \textbf{A}
  \item \textbf{A}
  \item \textbf{A}
  \item \textbf{A}
  \item \textbf{A}
  \item \textbf{A}
  \item \textbf{A}
  \item \textbf{A}
  \item \textbf{A}
  \item \textbf{A}
  \item \textbf{A}
  \item \textbf{A}
  \item \textbf{A}
  \item \textbf{A}
  \item \textbf{A}
  \item \textbf{A}
  \item \textbf{A}
  \item \textbf{A}
  \item \textbf{A}
  \item \textbf{A}
  \item \textbf{A}
  \item \textbf{A}
  \item \textbf{A}
  \item \textbf{A}
  \item \textbf{A}
  \item \textbf{A}
  \item \textbf{A}
  \item \textbf{A}
  \item \textbf{A}
  \item \textbf{A}
  \item \textbf{A}
  \item \textbf{A}
  \item \textbf{A}
  \item \textbf{A}
  \item \textbf{A}
  \item \textbf{A}
  \item \textbf{A}
  \item \textbf{A}
  \item \textbf{A}
  \item \textbf{A}
  \item \textbf{A}
  \item \textbf{A}
  \item \textbf{A}
  \item \textbf{A}
  \item \textbf{A}
  \item \textbf{A}
  \item \textbf{A}
  \item \textbf{A}
  \item \textbf{A}
  \item \textbf{A}
  \item \textbf{A}
  \item \textbf{A}
  \item \textbf{A}
  \item \textbf{A}
  \item \textbf{A}
  \item \textbf{A}
  \item \textbf{A}
  \item \textbf{A}
  \item \textbf{A}
  \item \textbf{A}
  \item \textbf{A}
  \item \textbf{A}
  \item \textbf{A}
  \item \textbf{A}
  \item \textbf{A}
  \item \textbf{A}
  \item \textbf{A}
  \item \textbf{A}
  \item \textbf{A}
  \item \textbf{A}
  \item \textbf{A}
  \item \textbf{A}
  \item \textbf{A}
  \item \textbf{A}
  \item \textbf{A}
  \item \textbf{A}
  \item \textbf{A}
  \item \textbf{A}
  \item \textbf{A}
  \item \textbf{A}
  \item \textbf{A}
  \item \textbf{A}
  \item \textbf{A}
  \item \textbf{A}
  \item \textbf{A}
  \item \textbf{A}
  \item \textbf{A}
  \item \textbf{A}
  \item \textbf{A}
  \item \textbf{A}
  \item \textbf{A}
  \item \textbf{A}
  \item \textbf{A}
  \item \textbf{A}
  \item \textbf{A}
  \item \textbf{A}
  \item \textbf{A}
  \item \textbf{A}
  \item \textbf{A}
  \item \textbf{A}
  \item \textbf{A}
  \item \textbf{A}
  \item \textbf{A}
  \item \textbf{A}
  \item \textbf{A}
  \item \textbf{A}
  \item \textbf{A}
  \item \textbf{A}
  \item \textbf{A}
  \item \textbf{A}
  \item \textbf{A}
  \item \textbf{A}
  \item \textbf{A}
  \item \textbf{A}
  \item \textbf{A}
  \item \textbf{A}
  \item \textbf{A}
  \item \textbf{A}
  \item \textbf{A}
  \item \textbf{A}
  \item \textbf{A}
  \item \textbf{A}
  \item \textbf{A}
  \item \textbf{A}
  \item \textbf{A}
  \item \textbf{A}
  \item \textbf{A}
  \item \textbf{A}
  \item \textbf{A}
  \item \textbf{A}
  \item \textbf{A}
  \item \textbf{A}
  \item \textbf{A}
  \item \textbf{A}
  \item \textbf{A}
  \item \textbf{A}
  \item \textbf{A}
  \item \textbf{A}
  \item \textbf{A}
  \item \textbf{A}
  \item \textbf{A}
  \item \textbf{A}
  \item \textbf{A}
  \item \textbf{A}
  \item \textbf{A}
  \item \textbf{A}
  \item \textbf{A}
  \item \textbf{A}
  \item \textbf{A}
  \item \textbf{A}
  \item \textbf{A}
  \item \textbf{A}
  \item \textbf{A}
  \item \textbf{A}
  \item \textbf{A}
  \item \textbf{A}
  \item \textbf{A}
  \item \textbf{A}
  \item \textbf{A}
  \item \textbf{A}
  \item \textbf{A}
  \item \textbf{A}
  \item \textbf{A}
  \item \textbf{A}
  \item \textbf{A}
  \item \textbf{A}
  \item \textbf{A}
  \item \textbf{A}
  \item \textbf{A}
  \item \textbf{A}
  \item \textbf{A}
  \item \textbf{A}
  \item \textbf{A}
  \item \textbf{A}
  \item \textbf{A}
  \item \textbf{A}
  \item \textbf{A}
  \item \textbf{A}
  \item \textbf{A}
  \item \textbf{A}
  \item \textbf{A}
  \item \textbf{A}
  \item \textbf{A}
  \item \textbf{A}
  \item \textbf{A}
  \item \textbf{A}
  \item \textbf{A}
  \item \textbf{A}
  \item \textbf{A}
  \item \textbf{A}
  \item \textbf{A}
  \item \textbf{A}
  \item \textbf{A}
  \item \textbf{A}
  \item \textbf{A}
  \item \textbf{A}
  \item \textbf{A}
  \item \textbf{A}
  \item \textbf{A}
  \item \textbf{A}
  \item \textbf{A}
  \item \textbf{A}
  \item \textbf{A}
  \item \textbf{A}
  \item \textbf{A}
  \item \textbf{A}
  \item \textbf{A}
  \item \textbf{A}
  \item \textbf{A}
  \item \textbf{A}
  \item \textbf{A}
  \item \textbf{A}
  \item \textbf{A}
  \item \textbf{A}
  \item \textbf{A}
  \item \textbf{A}
  \item \textbf{A}
  \item \textbf{A}
  \item \textbf{A}
  \item \textbf{A}
  \item \textbf{A}
  \item \textbf{A}
  \item \textbf{A}
  \item \textbf{A}
  \item \textbf{A}
  \item \textbf{A}
  \item \textbf{A}
  \item \textbf{A}
  \item \textbf{A}
  \item \textbf{A}
  \item \textbf{A}
  \item \textbf{A}
  \item \textbf{A}
  \item \textbf{A}
  \item \textbf{A}
  \item \textbf{A}
  \item \textbf{A}
  \item \textbf{A}
  \item \textbf{A}
  \item \textbf{A}
  \item \textbf{A}
  \item \textbf{A}
  \item \textbf{A}
  \item \textbf{A}
  \item \textbf{A}
  \item \textbf{A}
  \item \textbf{A}
  \item \textbf{A}
  \item \textbf{A}
  \item \textbf{A}
  \item \textbf{A}
  \item \textbf{A}
  \item \textbf{A}
  \item \textbf{A}
  \item \textbf{A}
  \item \textbf{A}
  \item \textbf{A}
  \item \textbf{A}
  \item \textbf{A}
  \item \textbf{A}
  \item \textbf{A}
  \item \textbf{A}
  \item \textbf{A}
  \item \textbf{A}
  \item \textbf{A}
  \item \textbf{A}
  \item \textbf{A}
  \item \textbf{A}
  \item \textbf{A}
  \item \textbf{A}
  \item \textbf{A}
  \item \textbf{A}
  \item \textbf{A}
  \item \textbf{A}
  \item \textbf{A}
  \item \textbf{A}
  \item \textbf{A}
  \item \textbf{A}
  \item \textbf{A}
  \item \textbf{A}
  \item \textbf{A}
  \item \textbf{A}
  \item \textbf{A}
  \item \textbf{A}
  \item \textbf{A}
  \item \textbf{A}
  \item \textbf{A}
  \item \textbf{A}
  \item \textbf{A}
  \item \textbf{A}
\end{enumerate}
\end{multicols}
\end{tcolorbox}
\newpage\n
\subsection{5.3 Física General y Mecánica}

\subsubsection*{SECCIÓN I: CONCEPTOS BÁSICOS Y MAGNITUDES (1–50)}

\begin{enumerate}[label=\textbf{\arabic*.}]
\setcounter{enumi}{0}
\item ¿Qué es la física?
\begin{enumerate}[label=\textbf{\Alph*)}]
\begin{multicols}{2}
  \item La ciencia que estudia la materia, la energía y sus interacciones
  \item La ciencia que estudia los seres vivos
  \item La ciencia que estudia los minerales
  \item La ciencia que estudia los astros
  \item La ciencia que estudia las reacciones químicas
  \item La ciencia que estudia los números
\end{multicols}
\end{enumerate}
\vspace{0.15cm}
\item ¿Qué es una magnitud física?
\begin{enumerate}[label=\textbf{\Alph*)}]
\begin{multicols}{2}
  \item Todo aquello que se puede medir y expresar con un número y una unidad
  \item Todo aquello que no se puede medir
  \item Todo aquello que se puede ver
  \item Todo aquello que se puede tocar
  \item Todo aquello que se puede oler
  \item Todo aquello que se puede saborear
\end{multicols}
\end{enumerate}
\vspace{0.15cm}
\item ¿Qué es una magnitud escalar?
\begin{enumerate}[label=\textbf{\Alph*)}]
\begin{multicols}{2}
  \item La que queda completamente definida por un número y una unidad
  \item La que necesita dirección y sentido
  \item La que necesita módulo, dirección y sentido
  \item La que no se puede medir
  \item La que no tiene unidades
  \item La que solo tiene dirección
\end{multicols}
\end{enumerate}
\vspace{0.15cm}
\item ¿Qué es una magnitud vectorial?
\begin{enumerate}[label=\textbf{\Alph*)}]
\begin{multicols}{2}
  \item La que necesita módulo, dirección y sentido
  \item La que solo necesita un número
  \item La que solo necesita una unidad
  \item La que no se puede medir
  \item La que no tiene unidades
  \item La que solo tiene sentido
\end{multicols}
\end{enumerate}
\vspace{0.15cm}
\item ¿Cuál es una magnitud escalar?
\begin{enumerate}[label=\textbf{\Alph*)}]
\begin{multicols}{2}
  \item La masa
  \item La velocidad
  \item La fuerza
  \item El desplazamiento
  \item La aceleración
  \item El peso
\end{multicols}
\end{enumerate}
\vspace{0.15cm}
\item ¿Cuál es una magnitud vectorial?
\begin{enumerate}[label=\textbf{\Alph*)}]
\begin{multicols}{2}
  \item La velocidad
  \item La masa
  \item El tiempo
  \item La temperatura
  \item La energía
  \item La densidad
\end{multicols}
\end{enumerate}
\vspace{0.15cm}
\item ¿Cuál es una magnitud escalar?
\begin{enumerate}[label=\textbf{\Alph*)}]
\begin{multicols}{2}
  \item La temperatura
  \item La velocidad
  \item La fuerza
  \item La aceleración
  \item El desplazamiento
  \item El peso
\end{multicols}
\end{enumerate}
\vspace{0.15cm}
\item ¿Cuál es una magnitud vectorial?
\begin{enumerate}[label=\textbf{\Alph*)}]
\begin{multicols}{2}
  \item La fuerza
  \item La masa
  \item El tiempo
  \item La temperatura
  \item La energía
  \item La densidad
\end{multicols}
\end{enumerate}
\vspace{0.15cm}
\item ¿Cuál es la unidad de longitud en el Sistema Internacional?
\begin{enumerate}[label=\textbf{\Alph*)}]
\begin{multicols}{2}
  \item El metro (m)
  \item El centímetro (cm)
  \item El kilómetro (km)
  \item La pulgada (in)
  \item El pie (ft)
  \item La milla (mi)
\end{multicols}
\end{enumerate}
\vspace{0.15cm}
\item ¿Cuál es la unidad de masa en el Sistema Internacional?
\begin{enumerate}[label=\textbf{\Alph*)}]
\begin{multicols}{2}
  \item El kilogramo (kg)
  \item El gramo (g)
  \item La libra (lb)
  \item La onza (oz)
  \item La tonelada (t)
  \item El quintal (q)
\end{multicols}
\end{enumerate}
\vspace{0.15cm}
\item ¿Cuál es la unidad de tiempo en el Sistema Internacional?
\begin{enumerate}[label=\textbf{\Alph*)}]
\begin{multicols}{2}
  \item El segundo (s)
  \item El minuto (min)
  \item La hora (h)
  \item El día (d)
  \item La semana (sem)
  \item El año (a)
\end{multicols}
\end{enumerate}
\vspace{0.15cm}
\item ¿Cuál es la unidad de velocidad en el Sistema Internacional?
\begin{enumerate}[label=\textbf{\Alph*)}]
\begin{multicols}{2}
  \item m/s
  \item km/h
  \item cm/s
  \item mi/h
  \item ft/s
  \item nudos
\end{multicols}
\end{enumerate}
\vspace{0.15cm}
\item ¿Cuál es la unidad de aceleración en el Sistema Internacional?
\begin{enumerate}[label=\textbf{\Alph*)}]
\begin{multicols}{2}
  \item m/s²
  \item km/h²
  \item cm/s²
  \item mi/h²
  \item ft/s²
  \item nudos/s
\end{multicols}
\end{enumerate}
\vspace{0.15cm}
\item ¿Cuál es la unidad de fuerza en el Sistema Internacional?
\begin{enumerate}[label=\textbf{\Alph*)}]
\begin{multicols}{2}
  \item El newton (N)
  \item La dina (dyn)
  \item La libra-fuerza (lbf)
  \item El kilogramo-fuerza (kgf)
  \item El poundal (pdl)
  \item El kilopondio (kp)
\end{multicols}
\end{enumerate}
\vspace{0.15cm}
\item ¿Qué es la notación científica?
\begin{enumerate}[label=\textbf{\Alph*)}]
\begin{multicols}{2}
  \item La forma de escribir números muy grandes o muy pequeños usando potencias de 10
  \item La forma de escribir números usando letras
  \item La forma de escribir números usando símbolos
  \item La forma de escribir números usando dibujos
  \item La forma de escribir números usando colores
  \item La forma de escribir números usando sonidos
\end{multicols}
\end{enumerate}
\vspace{0.15cm}
\item ¿Cómo se escribe 0.00045 en notación científica?
\begin{enumerate}[label=\textbf{\Alph*)}]
\begin{multicols}{2}
  \item \(4.5 	imes 10^{-4}\)
  \item \(4.5 	imes 10^{4}\)
  \item \(45 	imes 10^{-5}\)
  \item \(0.45 	imes 10^{-3}\)
  \item \(4.5 	imes 10^{-5}\)
  \item \(45 	imes 10^{-4}\)
\end{multicols}
\end{enumerate}
\vspace{0.15cm}
\item ¿Cómo se escribe 2500000 en notación científica?
\begin{enumerate}[label=\textbf{\Alph*)}]
\begin{multicols}{2}
  \item \(2.5 	imes 10^{6}\)
  \item \(2.5 	imes 10^{-6}\)
  \item \(25 	imes 10^{5}\)
  \item \(0.25 	imes 10^{7}\)
  \item \(2.5 	imes 10^{5}\)
  \item \(25 	imes 10^{6}\)
\end{multicols}
\end{enumerate}
\vspace{0.15cm}
\item ¿Qué es la precisión de una medida?
\begin{enumerate}[label=\textbf{\Alph*)}]
\begin{multicols}{2}
  \item El grado de concordancia entre medidas repetidas
  \item El grado de cercanía al valor verdadero
  \item El número de cifras significativas
  \item El error absoluto
  \item El error relativo
  \item La incertidumbre
\end{multicols}
\end{enumerate}
\vspace{0.15cm}
\item ¿Qué es la exactitud de una medida?
\begin{enumerate}[label=\textbf{\Alph*)}]
\begin{multicols}{2}
  \item El grado de cercanía al valor verdadero
  \item El grado de concordancia entre medidas repetidas
  \item El número de cifras significativas
  \item El error absoluto
  \item El error relativo
  \item La incertidumbre
\end{multicols}
\end{enumerate}
\vspace{0.15cm}
\item ¿Qué es el error absoluto?
\begin{enumerate}[label=\textbf{\Alph*)}]
\begin{multicols}{2}
  \item La diferencia entre el valor medido y el valor verdadero
  \item La diferencia entre dos medidas
  \item El cociente entre el error absoluto y el valor verdadero
  \item El número de cifras significativas
  \item La incertidumbre de la medida
  \item La precisión del instrumento
\end{multicols}
\end{enumerate}
\vspace{0.15cm}
\item ¿Qué es el error relativo?
\begin{enumerate}[label=\textbf{\Alph*)}]
\begin{multicols}{2}
  \item El cociente entre el error absoluto y el valor verdadero
  \item La diferencia entre el valor medido y el verdadero
  \item La diferencia entre dos medidas
  \item El número de cifras significativas
  \item La incertidumbre de la medida
  \item La precisión del instrumento
\end{multicols}
\end{enumerate}
\vspace{0.15cm}
\item ¿Qué es una cifra significativa?
\begin{enumerate}[label=\textbf{\Alph*)}]
\begin{multicols}{2}
  \item Cada uno de los dígitos que aportan información sobre la precisión de una medida
  \item Cada uno de los dígitos de un número
  \item Cada uno de los decimales de un número
  \item Cada uno de los enteros de un número
  \item Cada uno de los ceros de un número
  \item Cada uno de los unos de un número
\end{multicols}
\end{enumerate}
\vspace{0.15cm}
\item ¿Cuántas cifras significativas tiene 0.00320?
\begin{enumerate}[label=\textbf{\Alph*)}]
\begin{multicols}{2}
  \item 3
  \item 5
  \item 2
  \item 4
  \item 6
  \item 1
\end{multicols}
\end{enumerate}
\vspace{0.15cm}
\item ¿Cuántas cifras significativas tiene 1500?
\begin{enumerate}[label=\textbf{\Alph*)}]
\begin{multicols}{2}
  \item 2 (si no hay punto decimal)
  \item 4
  \item 3
  \item 1
  \item 5
  \item 6
\end{multicols}
\end{enumerate}
\vspace{0.15cm}
\item ¿Qué es la densidad?
\begin{enumerate}[label=\textbf{\Alph*)}]
\begin{multicols}{2}
  \item La masa por unidad de volumen
  \item El volumen por unidad de masa
  \item La masa por unidad de longitud
  \item La longitud por unidad de masa
  \item El peso por unidad de volumen
  \item El volumen por unidad de peso
\end{multicols}
\end{enumerate}
\vspace{0.15cm}
\item ¿Cuál es la densidad del agua?
\begin{enumerate}[label=\textbf{\Alph*)}]
\begin{multicols}{2}
  \item \(1000 \, 	ext{kg/m}^3\)
  \item \(100 \, 	ext{kg/m}^3\)
  \item \(10 \, 	ext{kg/m}^3\)
  \item \(1 \, 	ext{kg/m}^3\)
  \item \(10000 \, 	ext{kg/m}^3\)
  \item \(0.1 \, 	ext{kg/m}^3\)
\end{multicols}
\end{enumerate}
\vspace{0.15cm}
\item ¿Qué es el peso?
\begin{enumerate}[label=\textbf{\Alph*)}]
\begin{multicols}{2}
  \item La fuerza con que la Tierra atrae a un cuerpo
  \item La masa de un cuerpo
  \item El volumen de un cuerpo
  \item La densidad de un cuerpo
  \item La energía de un cuerpo
  \item El trabajo de un cuerpo
\end{multicols}
\end{enumerate}
\vspace{0.15cm}
\item ¿Cuál es la fórmula del peso?
\begin{enumerate}[label=\textbf{\Alph*)}]
\begin{multicols}{2}
  \item \(P = m \cdot g\)
  \item \(P = m / g\)
  \item \(P = g / m\)
  \item \(P = m + g\)
  \item \(P = m - g\)
  \item \(P = m \cdot g^2\)
\end{multicols}
\end{enumerate}
\vspace{0.15cm}
\item ¿Cuál es el valor aproximado de la gravedad en la superficie terrestre?
\begin{enumerate}[label=\textbf{\Alph*)}]
\begin{multicols}{2}
  \item \(9.8 \, 	ext{m/s}^2\)
  \item \(8.9 \, 	ext{m/s}^2\)
  \item \(10.8 \, 	ext{m/s}^2\)
  \item \(9.8 \, 	ext{m/s}\)
  \item \(9.8 \, 	ext{N}\)
  \item \(9.8 \, 	ext{kg}\)
\end{multicols}
\end{enumerate}
\vspace{0.15cm}
\item ¿Qué es la energía?
\begin{enumerate}[label=\textbf{\Alph*)}]
\begin{multicols}{2}
  \item La capacidad de realizar un trabajo
  \item La capacidad de ejercer una fuerza
  \item La capacidad de moverse
  \item La capacidad de cambiar de forma
  \item La capacidad de cambiar de temperatura
  \item La capacidad de cambiar de estado
\end{multicols}
\end{enumerate}
\vspace{0.15cm}
\item ¿Cuál es la unidad de energía en el Sistema Internacional?
\begin{enumerate}[label=\textbf{\Alph*)}]
\begin{multicols}{2}
  \item El julio (J)
  \item La caloría (cal)
  \item El kilovatio-hora (kWh)
  \item El electrón-voltio (eV)
  \item La BTU
  \item El ergio (erg)
\end{multicols}
\end{enumerate}
\vspace{0.15cm}
\item ¿Qué es el trabajo?
\begin{enumerate}[label=\textbf{\Alph*)}]
\begin{multicols}{2}
  \item El producto de la fuerza por el desplazamiento en la dirección de la fuerza
  \item El producto de la masa por la aceleración
  \item El producto de la masa por la velocidad
  \item El producto de la fuerza por el tiempo
  \item El producto de la potencia por el tiempo
  \item El producto de la energía por el tiempo
\end{multicols}
\end{enumerate}
\vspace{0.15cm}
\item ¿Cuál es la fórmula del trabajo?
\begin{enumerate}[label=\textbf{\Alph*)}]
\begin{multicols}{2}
  \item \(W = F \cdot d \cdot \cos	heta\)
  \item \(W = F \cdot d \cdot \sen	heta\)
  \item \(W = F \cdot d \cdot 	an	heta\)
  \item \(W = F / d\)
  \item \(W = d / F\)
  \item \(W = F + d\)
\end{multicols}
\end{enumerate}
\vspace{0.15cm}
\item ¿Qué es la potencia?
\begin{enumerate}[label=\textbf{\Alph*)}]
\begin{multicols}{2}
  \item La rapidez con que se realiza un trabajo
  \item La cantidad de trabajo realizado
  \item La cantidad de energía consumida
  \item La cantidad de fuerza aplicada
  \item La cantidad de desplazamiento
  \item La cantidad de tiempo empleado
\end{multicols}
\end{enumerate}
\vspace{0.15cm}
\item ¿Cuál es la unidad de potencia en el Sistema Internacional?
\begin{enumerate}[label=\textbf{\Alph*)}]
\begin{multicols}{2}
  \item El vatio (W)
  \item El julio (J)
  \item La caloría (cal)
  \item El kilovatio-hora (kWh)
  \item El caballo de vapor (CV)
  \item El ergio (erg)
\end{multicols}
\end{enumerate}
\vspace{0.15cm}
\item ¿Cuál es la fórmula de la potencia?
\begin{enumerate}[label=\textbf{\Alph*)}]
\begin{multicols}{2}
  \item \(P = W / t\)
  \item \(P = W \cdot t\)
  \item \(P = t / W\)
  \item \(P = W + t\)
  \item \(P = W - t\)
  \item \(P = W \cdot t^2\)
\end{multicols}
\end{enumerate}
\vspace{0.15cm}
\item ¿Qué es la energía cinética?
\begin{enumerate}[label=\textbf{\Alph*)}]
\begin{multicols}{2}
  \item La energía asociada al movimiento de un cuerpo
  \item La energía asociada a la posición de un cuerpo
  \item La energía asociada a la temperatura de un cuerpo
  \item La energía asociada a la deformación de un cuerpo
  \item La energía asociada a la carga de un cuerpo
  \item La energía asociada a la masa de un cuerpo
\end{multicols}
\end{enumerate}
\vspace{0.15cm}
\item ¿Cuál es la fórmula de la energía cinética?
\begin{enumerate}[label=\textbf{\Alph*)}]
\begin{multicols}{2}
  \item \(E_c = rac{1}{2} m v^2\)
  \item \(E_c = m v^2\)
  \item \(E_c = rac{1}{2} m v\)
  \item \(E_c = m v\)
  \item \(E_c = rac{1}{2} m^2 v\)
  \item \(E_c = m^2 v^2\)
\end{multicols}
\end{enumerate}
\vspace{0.15cm}
\item ¿Qué es la energía potencial gravitatoria?
\begin{enumerate}[label=\textbf{\Alph*)}]
\begin{multicols}{2}
  \item La energía asociada a la posición de un cuerpo en un campo gravitatorio
  \item La energía asociada al movimiento de un cuerpo
  \item La energía asociada a la temperatura de un cuerpo
  \item La energía asociada a la deformación de un cuerpo
  \item La energía asociada a la carga de un cuerpo
  \item La energía asociada a la masa de un cuerpo
\end{multicols}
\end{enumerate}
\vspace{0.15cm}
\item ¿Cuál es la fórmula de la energía potencial gravitatoria?
\begin{enumerate}[label=\textbf{\Alph*)}]
\begin{multicols}{2}
  \item \(E_p = m g h\)
  \item \(E_p = m g\)
  \item \(E_p = m h\)
  \item \(E_p = g h\)
  \item \(E_p = m g / h\)
  \item \(E_p = m h / g\)
\end{multicols}
\end{enumerate}
\vspace{0.15cm}
\item ¿Qué es la ley de conservación de la energía?
\begin{enumerate}[label=\textbf{\Alph*)}]
\begin{multicols}{2}
  \item La energía no se crea ni se destruye, solo se transforma
  \item La energía se crea y se destruye
  \item La energía solo se crea
  \item La energía solo se destruye
  \item La energía no se transforma
  \item La energía se destruye completamente
\end{multicols}
\end{enumerate}
\vspace{0.15cm}
\item ¿Qué es la cantidad de movimiento o momento lineal?
\begin{enumerate}[label=\textbf{\Alph*)}]
\begin{multicols}{2}
  \item El producto de la masa por la velocidad
  \item El producto de la masa por la aceleración
  \item El producto de la fuerza por el tiempo
  \item El producto de la masa por la distancia
  \item El producto de la velocidad por el tiempo
  \item El producto de la fuerza por la distancia
\end{multicols}
\end{enumerate}
\vspace{0.15cm}
\item ¿Cuál es la fórmula de la cantidad de movimiento?
\begin{enumerate}[label=\textbf{\Alph*)}]
\begin{multicols}{2}
  \item \(p = m v\)
  \item \(p = m a\)
  \item \(p = F t\)
  \item \(p = m d\)
  \item \(p = v t\)
  \item \(p = F d\)
\end{multicols}
\end{enumerate}
\vspace{0.15cm}
\item ¿Qué es el impulso?
\begin{enumerate}[label=\textbf{\Alph*)}]
\begin{multicols}{2}
  \item El producto de la fuerza por el tiempo de aplicación
  \item El producto de la masa por la velocidad
  \item El producto de la masa por la aceleración
  \item El producto de la fuerza por la distancia
  \item El producto de la velocidad por el tiempo
  \item El producto de la masa por la distancia
\end{multicols}
\end{enumerate}
\vspace{0.15cm}
\item ¿Cuál es la fórmula del impulso?
\begin{enumerate}[label=\textbf{\Alph*)}]
\begin{multicols}{2}
  \item \(I = F \cdot t\)
  \item \(I = m \cdot v\)
  \item \(I = m \cdot a\)
  \item \(I = F \cdot d\)
  \item \(I = v \cdot t\)
  \item \(I = m \cdot d\)
\end{multicols}
\end{enumerate}
\vspace{0.15cm}
\item ¿Qué establece el teorema del impulso?
\begin{enumerate}[label=\textbf{\Alph*)}]
\begin{multicols}{2}
  \item El impulso es igual al cambio de la cantidad de movimiento
  \item El impulso es igual a la fuerza
  \item El impulso es igual a la masa
  \item El impulso es igual a la velocidad
  \item El impulso es igual a la aceleración
  \item El impulso es igual a la distancia
\end{multicols}
\end{enumerate}
\vspace{0.15cm}
\item ¿Qué es una fuerza?
\begin{enumerate}[label=\textbf{\Alph*)}]
\begin{multicols}{2}
  \item Una interacción capaz de deformar un cuerpo o modificar su estado de movimiento
  \item Una magnitud escalar
  \item Una medida de la masa
  \item Una medida del volumen
  \item Una medida de la densidad
  \item Una medida de la temperatura
\end{multicols}
\end{enumerate}
\vspace{0.15cm}
\item ¿Qué establece la primera ley de Newton?
\begin{enumerate}[label=\textbf{\Alph*)}]
\begin{multicols}{2}
  \item Un cuerpo permanece en reposo o en movimiento rectilíneo uniforme si la fuerza neta es cero
  \item La fuerza es igual a la masa por la aceleración
  \item A toda acción corresponde una reacción igual y opuesta
  \item La energía se conserva
  \item La cantidad de movimiento se conserva
  \item La gravedad atrae a los cuerpos
\end{multicols}
\end{enumerate}
\vspace{0.15cm}
\item ¿Qué establece la segunda ley de Newton?
\begin{enumerate}[label=\textbf{\Alph*)}]
\begin{multicols}{2}
  \item La fuerza neta es igual a la masa por la aceleración
  \item Un cuerpo permanece en reposo si la fuerza neta es cero
  \item A toda acción corresponde una reacción igual y opuesta
  \item La energía se conserva
  \item La cantidad de movimiento se conserva
  \item La gravedad atrae a los cuerpos
\end{multicols}
\end{enumerate}
\vspace{0.15cm}
\item ¿Qué establece la tercera ley de Newton?
\begin{enumerate}[label=\textbf{\Alph*)}]
\begin{multicols}{2}
  \item A toda acción corresponde una reacción igual y opuesta
  \item La fuerza neta es igual a la masa por la aceleración
  \item Un cuerpo permanece en reposo si la fuerza neta es cero
  \item La energía se conserva
  \item La cantidad de movimiento se conserva
  \item La gravedad atrae a los cuerpos ---
\end{multicols}
\end{enumerate}
\vspace{0.15cm}
\end{enumerate}\n
\subsubsection*{SECCIÓN II: VECTORES – SUMA Y DESCOMPOSICIÓN (51–120)}

\begin{enumerate}[label=\textbf{\arabic*.}]
\setcounter{enumi}{50}
\item ¿Qué es un vector?
\begin{enumerate}[label=\textbf{\Alph*)}]
\begin{multicols}{2}
  \item Una magnitud con módulo, dirección y sentido
  \item Una magnitud con solo módulo
  \item Una magnitud con solo dirección
  \item Una magnitud con solo sentido
  \item Una magnitud sin módulo
  \item Una magnitud sin dirección
\end{multicols}
\end{enumerate}
\vspace{0.15cm}
\item ¿Cómo se representa gráficamente un vector?
\begin{enumerate}[label=\textbf{\Alph*)}]
\begin{multicols}{2}
  \item Con una flecha
  \item Con un punto
  \item Con una línea
  \item Con un círculo
  \item Con un cuadrado
  \item Con un triángulo
\end{multicols}
\end{enumerate}
\vspace{0.15cm}
\item ¿Qué es el módulo de un vector?
\begin{enumerate}[label=\textbf{\Alph*)}]
\begin{multicols}{2}
  \item La longitud del vector
  \item La dirección del vector
  \item El sentido del vector
  \item El origen del vector
  \item El extremo del vector
  \item La componente del vector
\end{multicols}
\end{enumerate}
\vspace{0.15cm}
\item ¿Qué es la dirección de un vector?
\begin{enumerate}[label=\textbf{\Alph*)}]
\begin{multicols}{2}
  \item El ángulo que forma con un eje de referencia
  \item La longitud del vector
  \item El sentido del vector
  \item El origen del vector
  \item El extremo del vector
  \item La componente del vector
\end{multicols}
\end{enumerate}
\vspace{0.15cm}
\item ¿Qué es el sentido de un vector?
\begin{enumerate}[label=\textbf{\Alph*)}]
\begin{multicols}{2}
  \item La orientación del vector a lo largo de su dirección
  \item La longitud del vector
  \item El ángulo del vector
  \item El origen del vector
  \item El extremo del vector
  \item La componente del vector
\end{multicols}
\end{enumerate}
\vspace{0.15cm}
\item ¿Qué son vectores equipolentes?
\begin{enumerate}[label=\textbf{\Alph*)}]
\begin{multicols}{2}
  \item Vectores con igual módulo, dirección y sentido
  \item Vectores con igual módulo pero diferente dirección
  \item Vectores con diferente módulo pero igual dirección
  \item Vectores con igual dirección pero diferente sentido
  \item Vectores con diferente módulo, dirección y sentido
  \item Vectores con igual módulo y dirección pero diferente sentido
\end{multicols}
\end{enumerate}
\vspace{0.15cm}
\item ¿Qué son vectores opuestos?
\begin{enumerate}[label=\textbf{\Alph*)}]
\begin{multicols}{2}
  \item Vectores con igual módulo y dirección pero sentido contrario
  \item Vectores con igual módulo, dirección y sentido
  \item Vectores con diferente módulo pero igual dirección
  \item Vectores con diferente módulo y dirección
  \item Vectores con igual módulo pero diferente dirección
  \item Vectores con diferente módulo y sentido
\end{multicols}
\end{enumerate}
\vspace{0.15cm}
\item ¿Qué es un vector unitario?
\begin{enumerate}[label=\textbf{\Alph*)}]
\begin{multicols}{2}
  \item Un vector de módulo 1
  \item Un vector de módulo 0
  \item Un vector de módulo 2
  \item Un vector de módulo 10
  \item Un vector de módulo infinito
  \item Un vector de módulo negativo
\end{multicols}
\end{enumerate}
\vspace{0.15cm}
\item ¿Cómo se calcula el vector unitario?
\begin{enumerate}[label=\textbf{\Alph*)}]
\begin{multicols}{2}
  \item Dividiendo el vector entre su módulo
  \item Multiplicando el vector por su módulo
  \item Sumando el vector a su módulo
  \item Restando el vector a su módulo
  \item Elevando el vector al cuadrado
  \item Sacando la raíz del vector
\end{multicols}
\end{enumerate}
\vspace{0.15cm}
\item ¿Qué son las componentes de un vector?
\begin{enumerate}[label=\textbf{\Alph*)}]
\begin{multicols}{2}
  \item Las proyecciones del vector sobre los ejes coordenados
  \item Las sumas del vector con los ejes
  \item Las restas del vector con los ejes
  \item Los productos del vector con los ejes
  \item Los cocientes del vector con los ejes
  \item Las potencias del vector con los ejes
\end{multicols}
\end{enumerate}
\vspace{0.15cm}
\item ¿Cuál es la componente \(x\) de un vector de módulo \(10\) y ángulo \(30^\circ\)?
\begin{enumerate}[label=\textbf{\Alph*)}]
\begin{multicols}{2}
  \item \(10\cos 30^\circ = 8.66\)
  \item \(10\sen 30^\circ = 5\)
  \item \(10	an 30^\circ = 5.77\)
  \item \(10/\cos 30^\circ = 11.55\)
  \item \(10/\sen 30^\circ = 20\)
  \item \(10	an 30^\circ = 5.77\)
\end{multicols}
\end{enumerate}
\vspace{0.15cm}
\item ¿Cuál es la componente \(y\) de un vector de módulo \(10\) y ángulo \(30^\circ\)?
\begin{enumerate}[label=\textbf{\Alph*)}]
\begin{multicols}{2}
  \item \(10\sen 30^\circ = 5\)
  \item \(10\cos 30^\circ = 8.66\)
  \item \(10	an 30^\circ = 5.77\)
  \item \(10/\cos 30^\circ = 11.55\)
  \item \(10/\sen 30^\circ = 20\)
  \item \(10	an 30^\circ = 5.77\)
\end{multicols}
\end{enumerate}
\vspace{0.15cm}
\item ¿Cuál es la componente \(x\) de un vector de módulo \(8\) y ángulo \(60^\circ\)?
\begin{enumerate}[label=\textbf{\Alph*)}]
\begin{multicols}{2}
  \item \(8\cos 60^\circ = 4\)
  \item \(8\sen 60^\circ = 6.93\)
  \item \(8	an 60^\circ = 13.86\)
  \item \(8/\cos 60^\circ = 16\)
  \item \(8/\sen 60^\circ = 9.24\)
  \item \(8	an 60^\circ = 13.86\)
\end{multicols}
\end{enumerate}
\vspace{0.15cm}
\item ¿Cuál es la componente \(y\) de un vector de módulo \(8\) y ángulo \(60^\circ\)?
\begin{enumerate}[label=\textbf{\Alph*)}]
\begin{multicols}{2}
  \item \(8\sen 60^\circ = 6.93\)
  \item \(8\cos 60^\circ = 4\)
  \item \(8	an 60^\circ = 13.86\)
  \item \(8/\cos 60^\circ = 16\)
  \item \(8/\sen 60^\circ = 9.24\)
  \item \(8	an 60^\circ = 13.86\)
\end{multicols}
\end{enumerate}
\vspace{0.15cm}
\item ¿Cuál es el módulo de un vector con componentes \(x = 3\) e \(y = 4\)?
\begin{enumerate}[label=\textbf{\Alph*)}]
\begin{multicols}{2}
  \item \(5\)
  \item \(7\)
  \item \(1\)
  \item \(12\)
  \item \(25\)
  \item \(\sqrt{7}\)
\end{multicols}
\end{enumerate}
\vspace{0.15cm}
\item ¿Cuál es el módulo de un vector con componentes \(x = 6\) e \(y = 8\)?
\begin{enumerate}[label=\textbf{\Alph*)}]
\begin{multicols}{2}
  \item \(10\)
  \item \(14\)
  \item \(2\)
  \item \(48\)
  \item \(100\)
  \item \(\sqrt{14}\)
\end{multicols}
\end{enumerate}
\vspace{0.15cm}
\item ¿Cuál es el módulo de un vector con componentes \(x = 5\) e \(y = 12\)?
\begin{enumerate}[label=\textbf{\Alph*)}]
\begin{multicols}{2}
  \item \(13\)
  \item \(17\)
  \item \(7\)
  \item \(60\)
  \item \(169\)
  \item \(\sqrt{17}\)
\end{multicols}
\end{enumerate}
\vspace{0.15cm}
\item ¿Cuál es la dirección de un vector con componentes \(x = 3\) e \(y = 4\)?
\begin{enumerate}[label=\textbf{\Alph*)}]
\begin{multicols}{2}
  \item \(	heta = 	an^{-1}(4/3) pprox 53.13^\circ\)
  \item \(	heta = 	an^{-1}(3/4) pprox 36.87^\circ\)
  \item \(	heta = 	an^{-1}(4/3) pprox 36.87^\circ\)
  \item \(	heta = 	an^{-1}(3/4) pprox 53.13^\circ\)
  \item \(	heta = 45^\circ\)
  \item \(	heta = 60^\circ\)
\end{multicols}
\end{enumerate}
\vspace{0.15cm}
\item ¿Cuál es la dirección de un vector con componentes \(x = 6\) e \(y = 8\)?
\begin{enumerate}[label=\textbf{\Alph*)}]
\begin{multicols}{2}
  \item \(	heta = 	an^{-1}(8/6) pprox 53.13^\circ\)
  \item \(	heta = 	an^{-1}(6/8) pprox 36.87^\circ\)
  \item \(	heta = 	an^{-1}(8/6) pprox 36.87^\circ\)
  \item \(	heta = 	an^{-1}(6/8) pprox 53.13^\circ\)
  \item \(	heta = 45^\circ\)
  \item \(	heta = 60^\circ\)
\end{multicols}
\end{enumerate}
\vspace{0.15cm}
\item ¿Cuál es la suma de los vectores \(ec{A} = (2, 3)\) y \(ec{B} = (4, 1)\)?
\begin{enumerate}[label=\textbf{\Alph*)}]
\begin{multicols}{2}
  \item \((6, 4)\)
  \item \((2, 2)\)
  \item \((8, 3)\)
  \item \((6, 2)\)
  \item \((4, 4)\)
  \item \((2, 4)\)
\end{multicols}
\end{enumerate}
\vspace{0.15cm}
\item ¿Cuál es la suma de los vectores \(ec{A} = (1, 5)\) y \(ec{B} = (3, -2)\)?
\begin{enumerate}[label=\textbf{\Alph*)}]
\begin{multicols}{2}
  \item \((4, 3)\)
  \item \((2, 7)\)
  \item \((4, 7)\)
  \item \((2, 3)\)
  \item \((3, 3)\)
  \item \((4, -3)\)
\end{multicols}
\end{enumerate}
\vspace{0.15cm}
\item ¿Cuál es la suma de los vectores \(ec{A} = (-2, 4)\) y \(ec{B} = (5, -1)\)?
\begin{enumerate}[label=\textbf{\Alph*)}]
\begin{multicols}{2}
  \item \((3, 3)\)
  \item \((7, 5)\)
  \item \((-7, 5)\)
  \item \((3, -5)\)
  \item \((-3, 3)\)
  \item \((7, -5)\)
\end{multicols}
\end{enumerate}
\vspace{0.15cm}
\item ¿Cuál es la resta de los vectores \(ec{A} = (5, 3)\) y \(ec{B} = (2, 1)\)?
\begin{enumerate}[label=\textbf{\Alph*)}]
\begin{multicols}{2}
  \item \((3, 2)\)
  \item \((7, 4)\)
  \item \((3, 4)\)
  \item \((7, 2)\)
  \item \((-3, -2)\)
  \item \((2, 3)\)
\end{multicols}
\end{enumerate}
\vspace{0.15cm}
\item ¿Cuál es la resta de los vectores \(ec{A} = (4, -1)\) y \(ec{B} = (1, 3)\)?
\begin{enumerate}[label=\textbf{\Alph*)}]
\begin{multicols}{2}
  \item \((3, -4)\)
  \item \((5, 2)\)
  \item \((3, 4)\)
  \item \((5, -4)\)
  \item \((-3, -4)\)
  \item \((4, -3)\)
\end{multicols}
\end{enumerate}
\vspace{0.15cm}
\item ¿Cuál es el producto de un vector \(ec{A} = (2, 3)\) por el escalar \(2\)?
\begin{enumerate}[label=\textbf{\Alph*)}]
\begin{multicols}{2}
  \item \((4, 6)\)
  \item \((2, 6)\)
  \item \((4, 3)\)
  \item \((1, 1.5)\)
  \item \((4, 9)\)
  \item \((6, 4)\)
\end{multicols}
\end{enumerate}
\vspace{0.15cm}
\item ¿Cuál es el producto de un vector \(ec{A} = (3, -2)\) por el escalar \(-1\)?
\begin{enumerate}[label=\textbf{\Alph*)}]
\begin{multicols}{2}
  \item \((-3, 2)\)
  \item \((3, 2)\)
  \item \((-3, -2)\)
  \item \((3, -2)\)
  \item \((-3, 0)\)
  \item \((0, -2)\)
\end{multicols}
\end{enumerate}
\vspace{0.15cm}
\item ¿Cuál es el producto escalar de \(ec{A} = (2, 3)\) y \(ec{B} = (4, 1)\)?
\begin{enumerate}[label=\textbf{\Alph*)}]
\begin{multicols}{2}
  \item \(11\)
  \item \(8\)
  \item \(3\)
  \item \(14\)
  \item \(10\)
  \item \(7\)
\end{multicols}
\end{enumerate}
\vspace{0.15cm}
\item ¿Cuál es el producto escalar de \(ec{A} = (1, 5)\) y \(ec{B} = (3, -2)\)?
\begin{enumerate}[label=\textbf{\Alph*)}]
\begin{multicols}{2}
  \item \(-7\)
  \item \(7\)
  \item \(13\)
  \item \(-13\)
  \item \(3\)
  \item \(-3\)
\end{multicols}
\end{enumerate}
\vspace{0.15cm}
\item ¿Cuál es el producto escalar de \(ec{A} = (2, 3)\) y \(ec{B} = (-1, 4)\)?
\begin{enumerate}[label=\textbf{\Alph*)}]
\begin{multicols}{2}
  \item \(10\)
  \item \(-10\)
  \item \(14\)
  \item \(-14\)
  \item \(2\)
  \item \(-2\)
\end{multicols}
\end{enumerate}
\vspace{0.15cm}
\item ¿Cuál es el producto vectorial de \(ec{A} = (1, 0, 0)\) y \(ec{B} = (0, 1, 0)\)?
\begin{enumerate}[label=\textbf{\Alph*)}]
\begin{multicols}{2}
  \item \((0, 0, 1)\)
  \item \((1, 1, 0)\)
  \item \((0, 1, 1)\)
  \item \((1, 0, 1)\)
  \item \((1, 1, 1)\)
  \item \((0, 0, 0)\)
\end{multicols}
\end{enumerate}
\vspace{0.15cm}
\item ¿Cuál es el módulo del vector \(ec{A} = (3, 4)\)?
\begin{enumerate}[label=\textbf{\Alph*)}]
\begin{multicols}{2}
  \item \(5\)
  \item \(7\)
  \item \(1\)
  \item \(12\)
  \item \(25\)
  \item \(\sqrt{7}\)
\end{multicols}
\end{enumerate}
\vspace{0.15cm}
\item ¿Cuál es el módulo del vector \(ec{A} = (5, 12)\)?
\begin{enumerate}[label=\textbf{\Alph*)}]
\begin{multicols}{2}
  \item \(13\)
  \item \(17\)
  \item \(7\)
  \item \(60\)
  \item \(169\)
  \item \(\sqrt{17}\)
\end{multicols}
\end{enumerate}
\vspace{0.15cm}
\item ¿Cuál es el módulo del vector \(ec{A} = (8, 15)\)?
\begin{enumerate}[label=\textbf{\Alph*)}]
\begin{multicols}{2}
  \item \(17\)
  \item \(23\)
  \item \(7\)
  \item \(120\)
  \item \(289\)
  \item \(\sqrt{23}\)
\end{multicols}
\end{enumerate}
\vspace{0.15cm}
\item ¿Cuál es el módulo del vector \(ec{A} = (1, 1)\)?
\begin{enumerate}[label=\textbf{\Alph*)}]
\begin{multicols}{2}
  \item \(\sqrt{2}\)
  \item \(2\)
  \item \(1\)
  \item \(\sqrt{3}\)
  \item \(0\)
  \item \(4\)
\end{multicols}
\end{enumerate}
\vspace{0.15cm}
\item ¿Cuál es el módulo del vector \(ec{A} = (1, \sqrt{3})\)?
\begin{enumerate}[label=\textbf{\Alph*)}]
\begin{multicols}{2}
  \item \(2\)
  \item \(\sqrt{3}\)
  \item \(1\)
  \item \(4\)
  \item \(3\)
  \item \(\sqrt{5}\)
\end{multicols}
\end{enumerate}
\vspace{0.15cm}
\item ¿Cuál es el vector unitario en la dirección de \(ec{A} = (3, 4)\)?
\begin{enumerate}[label=\textbf{\Alph*)}]
\begin{multicols}{2}
  \item \((3/5, 4/5)\)
  \item \((3/4, 4/3)\)
  \item \((3/7, 4/7)\)
  \item \((3/25, 4/25)\)
  \item \((3/12, 4/12)\)
  \item \((3/1, 4/1)\)
\end{multicols}
\end{enumerate}
\vspace{0.15cm}
\item ¿Cuál es el vector unitario en la dirección de \(ec{A} = (5, 12)\)?
\begin{enumerate}[label=\textbf{\Alph*)}]
\begin{multicols}{2}
  \item \((5/13, 12/13)\)
  \item \((5/12, 12/5)\)
  \item \((5/17, 12/17)\)
  \item \((5/169, 12/169)\)
  \item \((5/60, 12/60)\)
  \item \((5/1, 12/1)\)
\end{multicols}
\end{enumerate}
\vspace{0.15cm}
\item ¿Cuál es el ángulo entre dos vectores si su producto escalar es cero?
\begin{enumerate}[label=\textbf{\Alph*)}]
\begin{multicols}{2}
  \item \(90^\circ\)
  \item \(0^\circ\)
  \item \(180^\circ\)
  \item \(45^\circ\)
  \item \(60^\circ\)
  \item \(30^\circ\)
\end{multicols}
\end{enumerate}
\vspace{0.15cm}
\item ¿Cuál es el ángulo entre dos vectores si su producto escalar es igual al producto de sus módulos?
\begin{enumerate}[label=\textbf{\Alph*)}]
\begin{multicols}{2}
  \item \(0^\circ\)
  \item \(90^\circ\)
  \item \(180^\circ\)
  \item \(45^\circ\)
  \item \(60^\circ\)
  \item \(30^\circ\)
\end{multicols}
\end{enumerate}
\vspace{0.15cm}
\item ¿Cuál es el ángulo entre dos vectores si su producto escalar es igual al negativo del producto de sus módulos?
\begin{enumerate}[label=\textbf{\Alph*)}]
\begin{multicols}{2}
  \item \(180^\circ\)
  \item \(0^\circ\)
  \item \(90^\circ\)
  \item \(45^\circ\)
  \item \(60^\circ\)
  \item \(30^\circ\)
\end{multicols}
\end{enumerate}
\vspace{0.15cm}
\item ¿Qué es la descomposición vectorial?
\begin{enumerate}[label=\textbf{\Alph*)}]
\begin{multicols}{2}
  \item El proceso de expresar un vector como suma de sus componentes
  \item El proceso de sumar vectores
  \item El proceso de restar vectores
  \item El proceso de multiplicar vectores
  \item El proceso de dividir vectores
  \item El proceso de elevar vectores
\end{multicols}
\end{enumerate}
\vspace{0.15cm}
\item ¿Qué es un sistema de coordenadas cartesianas?
\begin{enumerate}[label=\textbf{\Alph*)}]
\begin{multicols}{2}
  \item Un sistema de ejes perpendiculares que permiten ubicar puntos en el plano
  \item Un sistema de ejes paralelos
  \item Un sistema de ejes oblicuos
  \item Un sistema de ejes circulares
  \item Un sistema de ejes esféricos
  \item Un sistema de ejes cilíndricos
\end{multicols}
\end{enumerate}
\vspace{0.15cm}
\item ¿Cuántos ejes tiene el sistema de coordenadas cartesianas en el plano?
\begin{enumerate}[label=\textbf{\Alph*)}]
\begin{multicols}{2}
  \item Dos: eje \(x\) y eje \(y\)
  \item Uno: eje \(x\)
  \item Tres: ejes \(x\), \(y\) y \(z\)
  \item Cuatro: ejes \(x\), \(y\), \(z\) y \(w\)
  \item Cinco: ejes \(x\), \(y\), \(z\), \(w\) y \(v\)
  \item Seis: ejes \(x\), \(y\), \(z\), \(w\), \(v\) y \(u\)
\end{multicols}
\end{enumerate}
\vspace{0.15cm}
\item ¿Cuántos ejes tiene el sistema de coordenadas cartesianas en el espacio?
\begin{enumerate}[label=\textbf{\Alph*)}]
\begin{multicols}{2}
  \item Tres: ejes \(x\), \(y\) y \(z\)
  \item Dos: ejes \(x\) y \(y\)
  \item Uno: eje \(x\)
  \item Cuatro: ejes \(x\), \(y\), \(z\) y \(w\)
  \item Cinco: ejes \(x\), \(y\), \(z\), \(w\) y \(v\)
  \item Seis: ejes \(x\), \(y\), \(z\), \(w\), \(v\) y \(u\)
\end{multicols}
\end{enumerate}
\vspace{0.15cm}
\item ¿Qué es un vector en el plano?
\begin{enumerate}[label=\textbf{\Alph*)}]
\begin{multicols}{2}
  \item Un vector con dos componentes \((x, y)\)
  \item Un vector con una componente \((x)\)
  \item Un vector con tres componentes \((x, y, z)\)
  \item Un vector con cuatro componentes \((x, y, z, w)\)
  \item Un vector con cinco componentes
  \item Un vector con seis componentes
\end{multicols}
\end{enumerate}
\vspace{0.15cm}
\item ¿Qué es un vector en el espacio?
\begin{enumerate}[label=\textbf{\Alph*)}]
\begin{multicols}{2}
  \item Un vector con tres componentes \((x, y, z)\)
  \item Un vector con dos componentes \((x, y)\)
  \item Un vector con una componente \((x)\)
  \item Un vector con cuatro componentes \((x, y, z, w)\)
  \item Un vector con cinco componentes
  \item Un vector con seis componentes
\end{multicols}
\end{enumerate}
\vspace{0.15cm}
\item ¿Cuál es la componente \(x\) de un vector de módulo \(20\) y ángulo \(45^\circ\)?
\begin{enumerate}[label=\textbf{\Alph*)}]
\begin{multicols}{2}
  \item \(20\cos 45^\circ = 14.14\)
  \item \(20\sen 45^\circ = 14.14\)
  \item \(20	an 45^\circ = 20\)
  \item \(20/\cos 45^\circ = 28.28\)
  \item \(20/\sen 45^\circ = 28.28\)
  \item \(20	an 45^\circ = 20\)
\end{multicols}
\end{enumerate}
\vspace{0.15cm}
\item ¿Cuál es la componente \(y\) de un vector de módulo \(20\) y ángulo \(45^\circ\)?
\begin{enumerate}[label=\textbf{\Alph*)}]
\begin{multicols}{2}
  \item \(20\sen 45^\circ = 14.14\)
  \item \(20\cos 45^\circ = 14.14\)
  \item \(20	an 45^\circ = 20\)
  \item \(20/\cos 45^\circ = 28.28\)
  \item \(20/\sen 45^\circ = 28.28\)
  \item \(20	an 45^\circ = 20\)
\end{multicols}
\end{enumerate}
\vspace{0.15cm}
\item ¿Cuál es la componente \(x\) de un vector de módulo \(15\) y ángulo \(37^\circ\)?
\begin{enumerate}[label=\textbf{\Alph*)}]
\begin{multicols}{2}
  \item \(15\cos 37^\circ pprox 11.98\)
  \item \(15\sen 37^\circ pprox 9.03\)
  \item \(15	an 37^\circ pprox 11.30\)
  \item \(15/\cos 37^\circ pprox 18.78\)
  \item \(15/\sen 37^\circ pprox 24.92\)
  \item \(15	an 37^\circ pprox 11.30\)
\end{multicols}
\end{enumerate}
\vspace{0.15cm}
\item ¿Cuál es la componente \(y\) de un vector de módulo \(15\) y ángulo \(37^\circ\)?
\begin{enumerate}[label=\textbf{\Alph*)}]
\begin{multicols}{2}
  \item \(15\sen 37^\circ pprox 9.03\)
  \item \(15\cos 37^\circ pprox 11.98\)
  \item \(15	an 37^\circ pprox 11.30\)
  \item \(15/\cos 37^\circ pprox 18.78\)
  \item \(15/\sen 37^\circ pprox 24.92\)
  \item \(15	an 37^\circ pprox 11.30\)
\end{multicols}
\end{enumerate}
\vspace{0.15cm}
\item ¿Cuál es la resultante de dos fuerzas perpendiculares de \(6 \, 	ext{N}\) y \(8 \, 	ext{N}\)?
\begin{enumerate}[label=\textbf{\Alph*)}]
\begin{multicols}{2}
  \item \(10 \, 	ext{N}\)
  \item \(14 \, 	ext{N}\)
  \item \(2 \, 	ext{N}\)
  \item \(48 \, 	ext{N}\)
  \item \(100 \, 	ext{N}\)
  \item \(\sqrt{14} \, 	ext{N}\)
\end{multicols}
\end{enumerate}
\vspace{0.15cm}
\item ¿Cuál es la resultante de dos fuerzas perpendiculares de \(3 \, 	ext{N}\) y \(4 \, 	ext{N}\)?
\begin{enumerate}[label=\textbf{\Alph*)}]
\begin{multicols}{2}
  \item \(5 \, 	ext{N}\)
  \item \(7 \, 	ext{N}\)
  \item \(1 \, 	ext{N}\)
  \item \(12 \, 	ext{N}\)
  \item \(25 \, 	ext{N}\)
  \item \(\sqrt{7} \, 	ext{N}\)
\end{multicols}
\end{enumerate}
\vspace{0.15cm}
\item ¿Cuál es la resultante de dos fuerzas perpendiculares de \(5 \, 	ext{N}\) y \(12 \, 	ext{N}\)?
\begin{enumerate}[label=\textbf{\Alph*)}]
\begin{multicols}{2}
  \item \(13 \, 	ext{N}\)
  \item \(17 \, 	ext{N}\)
  \item \(7 \, 	ext{N}\)
  \item \(60 \, 	ext{N}\)
  \item \(169 \, 	ext{N}\)
  \item \(\sqrt{17} \, 	ext{N}\)
\end{multicols}
\end{enumerate}
\vspace{0.15cm}
\item ¿Cuál es la resultante de dos fuerzas de \(10 \, 	ext{N}\) cada una que forman un ángulo de \(60^\circ\)?
\begin{enumerate}[label=\textbf{\Alph*)}]
\begin{multicols}{2}
  \item \(10\sqrt{3} \, 	ext{N} pprox 17.32 \, 	ext{N}\)
  \item \(20 \, 	ext{N}\)
  \item \(10 \, 	ext{N}\)
  \item \(5 \, 	ext{N}\)
  \item \(15 \, 	ext{N}\)
  \item \(10\sqrt{2} \, 	ext{N} pprox 14.14 \, 	ext{N}\)
\end{multicols}
\end{enumerate}
\vspace{0.15cm}
\item ¿Cuál es la resultante de dos fuerzas de \(10 \, 	ext{N}\) cada una que forman un ángulo de \(90^\circ\)?
\begin{enumerate}[label=\textbf{\Alph*)}]
\begin{multicols}{2}
  \item \(10\sqrt{2} \, 	ext{N} pprox 14.14 \, 	ext{N}\)
  \item \(20 \, 	ext{N}\)
  \item \(10 \, 	ext{N}\)
  \item \(5 \, 	ext{N}\)
  \item \(15 \, 	ext{N}\)
  \item \(10\sqrt{3} \, 	ext{N} pprox 17.32 \, 	ext{N}\)
\end{multicols}
\end{enumerate}
\vspace{0.15cm}
\item ¿Cuál es la resultante de dos fuerzas de \(10 \, 	ext{N}\) cada una que forman un ángulo de \(120^\circ\)?
\begin{enumerate}[label=\textbf{\Alph*)}]
\begin{multicols}{2}
  \item \(10 \, 	ext{N}\)
  \item \(20 \, 	ext{N}\)
  \item \(10\sqrt{3} \, 	ext{N} pprox 17.32 \, 	ext{N}\)
  \item \(10\sqrt{2} \, 	ext{N} pprox 14.14 \, 	ext{N}\)
  \item \(5 \, 	ext{N}\)
  \item \(15 \, 	ext{N}\)
\end{multicols}
\end{enumerate}
\vspace{0.15cm}
\item ¿Cuál es la resultante de dos fuerzas de \(10 \, 	ext{N}\) cada una que forman un ángulo de \(0^\circ\)?
\begin{enumerate}[label=\textbf{\Alph*)}]
\begin{multicols}{2}
  \item \(20 \, 	ext{N}\)
  \item \(10 \, 	ext{N}\)
  \item \(0 \, 	ext{N}\)
  \item \(10\sqrt{2} \, 	ext{N} pprox 14.14 \, 	ext{N}\)
  \item \(10\sqrt{3} \, 	ext{N} pprox 17.32 \, 	ext{N}\)
  \item \(5 \, 	ext{N}\)
\end{multicols}
\end{enumerate}
\vspace{0.15cm}
\item ¿Cuál es la resultante de dos fuerzas de \(10 \, 	ext{N}\) cada una que forman un ángulo de \(180^\circ\)?
\begin{enumerate}[label=\textbf{\Alph*)}]
\begin{multicols}{2}
  \item \(0 \, 	ext{N}\)
  \item \(20 \, 	ext{N}\)
  \item \(10 \, 	ext{N}\)
  \item \(10\sqrt{2} \, 	ext{N} pprox 14.14 \, 	ext{N}\)
  \item \(10\sqrt{3} \, 	ext{N} pprox 17.32 \, 	ext{N}\)
  \item \(5 \, 	ext{N}\)
\end{multicols}
\end{enumerate}
\vspace{0.15cm}
\item ¿Qué es la suma de vectores por el método del paralelogramo?
\begin{enumerate}[label=\textbf{\Alph*)}]
\begin{multicols}{2}
  \item Un método gráfico para sumar dos vectores
  \item Un método analítico para sumar dos vectores
  \item Un método gráfico para restar dos vectores
  \item Un método analítico para restar dos vectores
  \item Un método para multiplicar vectores
  \item Un método para dividir vectores
\end{multicols}
\end{enumerate}
\vspace{0.15cm}
\item ¿Qué es la suma de vectores por el método del polígono?
\begin{enumerate}[label=\textbf{\Alph*)}]
\begin{multicols}{2}
  \item Un método gráfico para sumar varios vectores
  \item Un método analítico para sumar varios vectores
  \item Un método gráfico para restar varios vectores
  \item Un método analítico para restar varios vectores
  \item Un método para multiplicar vectores
  \item Un método para dividir vectores
\end{multicols}
\end{enumerate}
\vspace{0.15cm}
\item ¿Qué es la suma de vectores por el método analítico?
\begin{enumerate}[label=\textbf{\Alph*)}]
\begin{multicols}{2}
  \item Un método que usa las componentes de los vectores
  \item Un método que usa dibujos
  \item Un método que usa reglas
  \item Un método que usa compases
  \item Un método que usa transportadores
  \item Un método que usa escuadras
\end{multicols}
\end{enumerate}
\vspace{0.15cm}
\item ¿Cuál es el módulo de la resultante de dos vectores de \(6 \, 	ext{N}\) y \(8 \, 	ext{N}\) que forman un ángulo de \(90^\circ\)?
\begin{enumerate}[label=\textbf{\Alph*)}]
\begin{multicols}{2}
  \item \(10 \, 	ext{N}\)
  \item \(14 \, 	ext{N}\)
  \item \(2 \, 	ext{N}\)
  \item \(48 \, 	ext{N}\)
  \item \(100 \, 	ext{N}\)
  \item \(\sqrt{14} \, 	ext{N}\)
\end{multicols}
\end{enumerate}
\vspace{0.15cm}
\item ¿Cuál es el módulo de la resultante de dos vectores de \(3 \, 	ext{N}\) y \(4 \, 	ext{N}\) que forman un ángulo de \(90^\circ\)?
\begin{enumerate}[label=\textbf{\Alph*)}]
\begin{multicols}{2}
  \item \(5 \, 	ext{N}\)
  \item \(7 \, 	ext{N}\)
  \item \(1 \, 	ext{N}\)
  \item \(12 \, 	ext{N}\)
  \item \(25 \, 	ext{N}\)
  \item \(\sqrt{7} \, 	ext{N}\)
\end{multicols}
\end{enumerate}
\vspace{0.15cm}
\item ¿Cuál es el módulo de la resultante de dos vectores de \(5 \, 	ext{N}\) y \(12 \, 	ext{N}\) que forman un ángulo de \(90^\circ\)?
\begin{enumerate}[label=\textbf{\Alph*)}]
\begin{multicols}{2}
  \item \(13 \, 	ext{N}\)
  \item \(17 \, 	ext{N}\)
  \item \(7 \, 	ext{N}\)
  \item \(60 \, 	ext{N}\)
  \item \(169 \, 	ext{N}\)
  \item \(\sqrt{17} \, 	ext{N}\)
\end{multicols}
\end{enumerate}
\vspace{0.15cm}
\item ¿Cuál es la dirección de la resultante de dos vectores de \(6 \, 	ext{N}\) y \(8 \, 	ext{N}\) que forman un ángulo de \(90^\circ\)?
\begin{enumerate}[label=\textbf{\Alph*)}]
\begin{multicols}{2}
  \item \(	heta = 	an^{-1}(8/6) pprox 53.13^\circ\)
  \item \(	heta = 	an^{-1}(6/8) pprox 36.87^\circ\)
  \item \(	heta = 45^\circ\)
  \item \(	heta = 30^\circ\)
  \item \(	heta = 60^\circ\)
  \item \(	heta = 90^\circ\)
\end{multicols}
\end{enumerate}
\vspace{0.15cm}
\item ¿Cuál es la dirección de la resultante de dos vectores de \(3 \, 	ext{N}\) y \(4 \, 	ext{N}\) que forman un ángulo de \(90^\circ\)?
\begin{enumerate}[label=\textbf{\Alph*)}]
\begin{multicols}{2}
  \item \(	heta = 	an^{-1}(4/3) pprox 53.13^\circ\)
  \item \(	heta = 	an^{-1}(3/4) pprox 36.87^\circ\)
  \item \(	heta = 45^\circ\)
  \item \(	heta = 30^\circ\)
  \item \(	heta = 60^\circ\)
  \item \(	heta = 90^\circ\)
\end{multicols}
\end{enumerate}
\vspace{0.15cm}
\item ¿Cuál es la dirección de la resultante de dos vectores de \(5 \, 	ext{N}\) y \(12 \, 	ext{N}\) que forman un ángulo de \(90^\circ\)?
\begin{enumerate}[label=\textbf{\Alph*)}]
\begin{multicols}{2}
  \item \(	heta = 	an^{-1}(12/5) pprox 67.38^\circ\)
  \item \(	heta = 	an^{-1}(5/12) pprox 22.62^\circ\)
  \item \(	heta = 45^\circ\)
  \item \(	heta = 30^\circ\)
  \item \(	heta = 60^\circ\)
  \item \(	heta = 90^\circ\)
\end{multicols}
\end{enumerate}
\vspace{0.15cm}
\item ¿Qué es un vector resultante?
\begin{enumerate}[label=\textbf{\Alph*)}]
\begin{multicols}{2}
  \item El vector que equivale a la suma de dos o más vectores
  \item El vector que equivale a la resta de dos o más vectores
  \item El vector que equivale al producto de dos o más vectores
  \item El vector que equivale a la división de dos o más vectores
  \item El vector que equivale a la potencia de dos o más vectores
  \item El vector que equivale a la raíz de dos o más vectores
\end{multicols}
\end{enumerate}
\vspace{0.15cm}
\item ¿Qué es un vector equilibrante?
\begin{enumerate}[label=\textbf{\Alph*)}]
\begin{multicols}{2}
  \item El vector que tiene igual módulo y dirección que la resultante pero sentido opuesto
  \item El vector que tiene igual módulo, dirección y sentido que la resultante
  \item El vector que tiene diferente módulo y dirección que la resultante
  \item El vector que tiene igual módulo pero diferente dirección
  \item El vector que tiene diferente módulo pero igual dirección
  \item El vector que tiene igual dirección pero diferente sentido
\end{multicols}
\end{enumerate}
\vspace{0.15cm}
\item ¿Cuál es la condición de equilibrio para un sistema de fuerzas?
\begin{enumerate}[label=\textbf{\Alph*)}]
\begin{multicols}{2}
  \item La suma vectorial de todas las fuerzas debe ser cero
  \item La suma vectorial de todas las fuerzas debe ser máxima
  \item La suma vectorial de todas las fuerzas debe ser positiva
  \item La suma vectorial de todas las fuerzas debe ser negativa
  \item La suma vectorial de todas las fuerzas debe ser infinita
  \item La suma vectorial de todas las fuerzas debe ser unitaria ---
\end{multicols}
\end{enumerate}
\vspace{0.15cm}
\end{enumerate}\n
\subsubsection*{SECCIÓN III: CINEMÁTICA – MRU (121–180)}

\begin{enumerate}[label=\textbf{\arabic*.}]
\setcounter{enumi}{120}
\item ¿Qué es la cinemática?
\begin{enumerate}[label=\textbf{\Alph*)}]
\begin{multicols}{2}
  \item La parte de la mecánica que estudia el movimiento sin considerar sus causas
  \item La parte de la mecánica que estudia el movimiento considerando sus causas
  \item La parte de la mecánica que estudia el equilibrio
  \item La parte de la mecánica que estudia las fuerzas
  \item La parte de la mecánica que estudia la energía
  \item La parte de la mecánica que estudia el trabajo
\end{multicols}
\end{enumerate}
\vspace{0.15cm}
\item ¿Qué es el movimiento?
\begin{enumerate}[label=\textbf{\Alph*)}]
\begin{multicols}{2}
  \item El cambio de posición de un cuerpo respecto a un sistema de referencia
  \item La permanencia de un cuerpo en su posición
  \item La deformación de un cuerpo
  \item La rotura de un cuerpo
  \item La fusión de un cuerpo
  \item La evaporación de un cuerpo
\end{multicols}
\end{enumerate}
\vspace{0.15cm}
\item ¿Qué es un sistema de referencia?
\begin{enumerate}[label=\textbf{\Alph*)}]
\begin{multicols}{2}
  \item Un conjunto de ejes y un origen que permiten describir el movimiento
  \item Un conjunto de fuerzas
  \item Un conjunto de energías
  \item Un conjunto de masas
  \item Un conjunto de volúmenes
  \item Un conjunto de densidades
\end{multicols}
\end{enumerate}
\vspace{0.15cm}
\item ¿Qué es la trayectoria?
\begin{enumerate}[label=\textbf{\Alph*)}]
\begin{multicols}{2}
  \item La línea que describe un cuerpo en su movimiento
  \item El punto de partida de un cuerpo
  \item El punto de llegada de un cuerpo
  \item La velocidad de un cuerpo
  \item La aceleración de un cuerpo
  \item El tiempo de un cuerpo
\end{multicols}
\end{enumerate}
\vspace{0.15cm}
\item ¿Qué es la distancia recorrida?
\begin{enumerate}[label=\textbf{\Alph*)}]
\begin{multicols}{2}
  \item La longitud total de la trayectoria
  \item El cambio de posición de un cuerpo
  \item La velocidad de un cuerpo
  \item La aceleración de un cuerpo
  \item El tiempo de un cuerpo
  \item La masa de un cuerpo
\end{multicols}
\end{enumerate}
\vspace{0.15cm}
\item ¿Qué es el desplazamiento?
\begin{enumerate}[label=\textbf{\Alph*)}]
\begin{multicols}{2}
  \item El cambio de posición de un cuerpo
  \item La longitud total de la trayectoria
  \item La velocidad de un cuerpo
  \item La aceleración de un cuerpo
  \item El tiempo de un cuerpo
  \item La masa de un cuerpo
\end{multicols}
\end{enumerate}
\vspace{0.15cm}
\item ¿Qué es la velocidad media?
\begin{enumerate}[label=\textbf{\Alph*)}]
\begin{multicols}{2}
  \item El desplazamiento dividido entre el tiempo empleado
  \item La distancia dividida entre el tiempo empleado
  \item La aceleración dividida entre el tiempo empleado
  \item La masa dividida entre el tiempo empleado
  \item La fuerza dividida entre el tiempo empleado
  \item La energía dividida entre el tiempo empleado
\end{multicols}
\end{enumerate}
\vspace{0.15cm}
\item ¿Qué es la rapidez media?
\begin{enumerate}[label=\textbf{\Alph*)}]
\begin{multicols}{2}
  \item La distancia recorrida dividida entre el tiempo empleado
  \item El desplazamiento dividido entre el tiempo empleado
  \item La aceleración dividida entre el tiempo empleado
  \item La masa dividida entre el tiempo empleado
  \item La fuerza dividida entre el tiempo empleado
  \item La energía dividida entre el tiempo empleado
\end{multicols}
\end{enumerate}
\vspace{0.15cm}
\item ¿Qué es la velocidad instantánea?
\begin{enumerate}[label=\textbf{\Alph*)}]
\begin{multicols}{2}
  \item La velocidad en un instante determinado
  \item La velocidad promedio en un intervalo
  \item La velocidad máxima de un cuerpo
  \item La velocidad mínima de un cuerpo
  \item La velocidad inicial de un cuerpo
  \item La velocidad final de un cuerpo
\end{multicols}
\end{enumerate}
\vspace{0.15cm}
\item ¿Qué es el movimiento rectilíneo uniforme (MRU)?
\begin{enumerate}[label=\textbf{\Alph*)}]
\begin{multicols}{2}
  \item El movimiento con velocidad constante en línea recta
  \item El movimiento con aceleración constante en línea recta
  \item El movimiento con velocidad variable en línea recta
  \item El movimiento con aceleración variable en línea recta
  \item El movimiento circular con velocidad constante
  \item El movimiento circular con aceleración constante
\end{multicols}
\end{enumerate}
\vspace{0.15cm}
\item ¿Cuál es la fórmula de la velocidad en el MRU?
\begin{enumerate}[label=\textbf{\Alph*)}]
\begin{multicols}{2}
  \item \(v = d / t\)
  \item \(v = d \cdot t\)
  \item \(v = t / d\)
  \item \(v = d + t\)
  \item \(v = d - t\)
  \item \(v = d \cdot t^2\)
\end{multicols}
\end{enumerate}
\vspace{0.15cm}
\item ¿Cuál es la fórmula de la distancia en el MRU?
\begin{enumerate}[label=\textbf{\Alph*)}]
\begin{multicols}{2}
  \item \(d = v \cdot t\)
  \item \(d = v / t\)
  \item \(d = t / v\)
  \item \(d = v + t\)
  \item \(d = v - t\)
  \item \(d = v \cdot t^2\)
\end{multicols}
\end{enumerate}
\vspace{0.15cm}
\item ¿Cuál es la fórmula del tiempo en el MRU?
\begin{enumerate}[label=\textbf{\Alph*)}]
\begin{multicols}{2}
  \item \(t = d / v\)
  \item \(t = d \cdot v\)
  \item \(t = v / d\)
  \item \(t = d + v\)
  \item \(t = d - v\)
  \item \(t = d \cdot v^2\)
\end{multicols}
\end{enumerate}
\vspace{0.15cm}
\item Un auto viaja a \(60 \, 	ext{km/h}\) durante \(2 \, 	ext{h}\). ¿Qué distancia recorre?
\begin{enumerate}[label=\textbf{\Alph*)}]
\begin{multicols}{2}
  \item \(120 \, 	ext{km}\)
  \item \(30 \, 	ext{km}\)
  \item \(60 \, 	ext{km}\)
  \item \(180 \, 	ext{km}\)
  \item \(90 \, 	ext{km}\)
  \item \(240 \, 	ext{km}\)
\end{multicols}
\end{enumerate}
\vspace{0.15cm}
\item Un auto viaja a \(80 \, 	ext{km/h}\) durante \(3 \, 	ext{h}\). ¿Qué distancia recorre?
\begin{enumerate}[label=\textbf{\Alph*)}]
\begin{multicols}{2}
  \item \(240 \, 	ext{km}\)
  \item \(160 \, 	ext{km}\)
  \item \(80 \, 	ext{km}\)
  \item \(320 \, 	ext{km}\)
  \item \(120 \, 	ext{km}\)
  \item \(400 \, 	ext{km}\)
\end{multicols}
\end{enumerate}
\vspace{0.15cm}
\item Un ciclista viaja a \(20 \, 	ext{km/h}\) durante \(4 \, 	ext{h}\). ¿Qué distancia recorre?
\begin{enumerate}[label=\textbf{\Alph*)}]
\begin{multicols}{2}
  \item \(80 \, 	ext{km}\)
  \item \(40 \, 	ext{km}\)
  \item \(20 \, 	ext{km}\)
  \item \(100 \, 	ext{km}\)
  \item \(60 \, 	ext{km}\)
  \item \(120 \, 	ext{km}\)
\end{multicols}
\end{enumerate}
\vspace{0.15cm}
\item Un tren viaja a \(120 \, 	ext{km/h}\) durante \(1.5 \, 	ext{h}\). ¿Qué distancia recorre?
\begin{enumerate}[label=\textbf{\Alph*)}]
\begin{multicols}{2}
  \item \(180 \, 	ext{km}\)
  \item \(120 \, 	ext{km}\)
  \item \(240 \, 	ext{km}\)
  \item \(60 \, 	ext{km}\)
  \item \(300 \, 	ext{km}\)
  \item \(150 \, 	ext{km}\)
\end{multicols}
\end{enumerate}
\vspace{0.15cm}
\item Un avión viaja a \(600 \, 	ext{km/h}\) durante \(4 \, 	ext{h}\). ¿Qué distancia recorre?
\begin{enumerate}[label=\textbf{\Alph*)}]
\begin{multicols}{2}
  \item \(2400 \, 	ext{km}\)
  \item \(1200 \, 	ext{km}\)
  \item \(600 \, 	ext{km}\)
  \item \(3000 \, 	ext{km}\)
  \item \(1800 \, 	ext{km}\)
  \item \(3600 \, 	ext{km}\)
\end{multicols}
\end{enumerate}
\vspace{0.15cm}
\item Un auto recorre \(300 \, 	ext{km}\) en \(5 \, 	ext{h}\). ¿Cuál es su velocidad media?
\begin{enumerate}[label=\textbf{\Alph*)}]
\begin{multicols}{2}
  \item \(60 \, 	ext{km/h}\)
  \item \(50 \, 	ext{km/h}\)
  \item \(70 \, 	ext{km/h}\)
  \item \(80 \, 	ext{km/h}\)
  \item \(100 \, 	ext{km/h}\)
  \item \(150 \, 	ext{km/h}\)
\end{multicols}
\end{enumerate}
\vspace{0.15cm}
\item Un ciclista recorre \(120 \, 	ext{km}\) en \(6 \, 	ext{h}\). ¿Cuál es su velocidad media?
\begin{enumerate}[label=\textbf{\Alph*)}]
\begin{multicols}{2}
  \item \(20 \, 	ext{km/h}\)
  \item \(30 \, 	ext{km/h}\)
  \item \(40 \, 	ext{km/h}\)
  \item \(10 \, 	ext{km/h}\)
  \item \(50 \, 	ext{km/h}\)
  \item \(60 \, 	ext{km/h}\)
\end{multicols}
\end{enumerate}
\vspace{0.15cm}
\item Un tren recorre \(450 \, 	ext{km}\) en \(3 \, 	ext{h}\). ¿Cuál es su velocidad media?
\begin{enumerate}[label=\textbf{\Alph*)}]
\begin{multicols}{2}
  \item \(150 \, 	ext{km/h}\)
  \item \(100 \, 	ext{km/h}\)
  \item \(200 \, 	ext{km/h}\)
  \item \(50 \, 	ext{km/h}\)
  \item \(250 \, 	ext{km/h}\)
  \item \(300 \, 	ext{km/h}\)
\end{multicols}
\end{enumerate}
\vspace{0.15cm}
\item Un avión recorre \(1800 \, 	ext{km}\) en \(3 \, 	ext{h}\). ¿Cuál es su velocidad media?
\begin{enumerate}[label=\textbf{\Alph*)}]
\begin{multicols}{2}
  \item \(600 \, 	ext{km/h}\)
  \item \(900 \, 	ext{km/h}\)
  \item \(300 \, 	ext{km/h}\)
  \item \(1200 \, 	ext{km/h}\)
  \item \(450 \, 	ext{km/h}\)
  \item \(150 \, 	ext{km/h}\)
\end{multicols}
\end{enumerate}
\vspace{0.15cm}
\item Un auto viaja a \(90 \, 	ext{km/h}\) y recorre \(270 \, 	ext{km}\). ¿Cuánto tiempo tarda?
\begin{enumerate}[label=\textbf{\Alph*)}]
\begin{multicols}{2}
  \item \(3 \, 	ext{h}\)
  \item \(2 \, 	ext{h}\)
  \item \(4 \, 	ext{h}\)
  \item \(5 \, 	ext{h}\)
  \item \(6 \, 	ext{h}\)
  \item \(1.5 \, 	ext{h}\)
\end{multicols}
\end{enumerate}
\vspace{0.15cm}
\item Un ciclista viaja a \(25 \, 	ext{km/h}\) y recorre \(100 \, 	ext{km}\). ¿Cuánto tiempo tarda?
\begin{enumerate}[label=\textbf{\Alph*)}]
\begin{multicols}{2}
  \item \(4 \, 	ext{h}\)
  \item \(5 \, 	ext{h}\)
  \item \(3 \, 	ext{h}\
  \item \(2 \, 	ext{h}\)
  \item \(6 \, 	ext{h}\)
  \item \(1 \, 	ext{h}\)
\end{multicols}
\end{enumerate}
\vspace{0.15cm}
\item Un tren viaja a \(150 \, 	ext{km/h}\) y recorre \(600 \, 	ext{km}\). ¿Cuánto tiempo tarda?
\begin{enumerate}[label=\textbf{\Alph*)}]
\begin{multicols}{2}
  \item \(4 \, 	ext{h}\)
  \item \(3 \, 	ext{h}\)
  \item \(5 \, 	ext{h}\)
  \item \(2 \, 	ext{h}\)
  \item \(6 \, 	ext{h}\)
  \item \(1.5 \, 	ext{h}\)
\end{multicols}
\end{enumerate}
\vspace{0.15cm}
\item Un avión viaja a \(800 \, 	ext{km/h}\) y recorre \(2400 \, 	ext{km}\). ¿Cuánto tiempo tarda?
\begin{enumerate}[label=\textbf{\Alph*)}]
\begin{multicols}{2}
  \item \(3 \, 	ext{h}\)
  \item \(4 \, 	ext{h}\)
  \item \(2 \, 	ext{h}\)
  \item \(5 \, 	ext{h}\)
  \item \(6 \, 	ext{h}\)
  \item \(1.5 \, 	ext{h}\)
\end{multicols}
\end{enumerate}
\vspace{0.15cm}
\item ¿Qué representa la pendiente en una gráfica de posición vs. tiempo en el MRU?
\begin{enumerate}[label=\textbf{\Alph*)}]
\begin{multicols}{2}
  \item La velocidad
  \item La aceleración
  \item La distancia
  \item El tiempo
  \item La masa
  \item La fuerza
\end{multicols}
\end{enumerate}
\vspace{0.15cm}
\item ¿Qué representa el área bajo la curva en una gráfica de velocidad vs. tiempo en el MRU?
\begin{enumerate}[label=\textbf{\Alph*)}]
\begin{multicols}{2}
  \item La distancia recorrida
  \item La aceleración
  \item La velocidad
  \item El tiempo
  \item La masa
  \item La fuerza
\end{multicols}
\end{enumerate}
\vspace{0.15cm}
\item ¿Cómo es la gráfica de posición vs. tiempo en el MRU?
\begin{enumerate}[label=\textbf{\Alph*)}]
\begin{multicols}{2}
  \item Una línea recta con pendiente constante
  \item Una parábola
  \item Una hipérbola
  \item Una circunferencia
  \item Una elipse
  \item Una línea curva
\end{multicols}
\end{enumerate}
\vspace{0.15cm}
\item ¿Cómo es la gráfica de velocidad vs. tiempo en el MRU?
\begin{enumerate}[label=\textbf{\Alph*)}]
\begin{multicols}{2}
  \item Una línea horizontal
  \item Una línea vertical
  \item Una parábola
  \item Una hipérbola
  \item Una circunferencia
  \item Una línea oblicua
\end{multicols}
\end{enumerate}
\vspace{0.15cm}
\item ¿Qué es la velocidad relativa?
\begin{enumerate}[label=\textbf{\Alph*)}]
\begin{multicols}{2}
  \item La velocidad de un cuerpo respecto a otro sistema de referencia
  \item La velocidad de un cuerpo respecto al suelo
  \item La velocidad de un cuerpo respecto al aire
  \item La velocidad de un cuerpo respecto al agua
  \item La velocidad de un cuerpo respecto al espacio
  \item La velocidad de un cuerpo respecto al tiempo
\end{multicols}
\end{enumerate}
\vspace{0.15cm}
\item Dos autos viajan en la misma dirección a \(60 \, 	ext{km/h}\) y \(80 \, 	ext{km/h}\). ¿Cuál es la velocidad relativa del segundo respecto al primero?
\begin{enumerate}[label=\textbf{\Alph*)}]
\begin{multicols}{2}
  \item \(20 \, 	ext{km/h}\)
  \item \(140 \, 	ext{km/h}\)
  \item \(60 \, 	ext{km/h}\)
  \item \(80 \, 	ext{km/h}\)
  \item \(100 \, 	ext{km/h}\)
  \item \(40 \, 	ext{km/h}\)
\end{multicols}
\end{enumerate}
\vspace{0.15cm}
\item Dos autos viajan en direcciones opuestas a \(60 \, 	ext{km/h}\) y \(80 \, 	ext{km/h}\). ¿Cuál es la velocidad relativa entre ellos?
\begin{enumerate}[label=\textbf{\Alph*)}]
\begin{multicols}{2}
  \item \(140 \, 	ext{km/h}\)
  \item \(20 \, 	ext{km/h}\)
  \item \(60 \, 	ext{km/h}\)
  \item \(80 \, 	ext{km/h}\)
  \item \(100 \, 	ext{km/h}\)
  \item \(40 \, 	ext{km/h}\)
\end{multicols}
\end{enumerate}
\vspace{0.15cm}
\item Un barco navega a \(20 \, 	ext{km/h}\) y la corriente es de \(5 \, 	ext{km/h}\) en la misma dirección. ¿Cuál es la velocidad del barco respecto a la orilla?
\begin{enumerate}[label=\textbf{\Alph*)}]
\begin{multicols}{2}
  \item \(25 \, 	ext{km/h}\)
  \item \(15 \, 	ext{km/h}\)
  \item \(20 \, 	ext{km/h}\)
  \item \(5 \, 	ext{km/h}\)
  \item \(100 \, 	ext{km/h}\)
  \item \(10 \, 	ext{km/h}\)
\end{multicols}
\end{enumerate}
\vspace{0.15cm}
\item Un barco navega a \(20 \, 	ext{km/h}\) y la corriente es de \(5 \, 	ext{km/h}\) en dirección opuesta. ¿Cuál es la velocidad del barco respecto a la orilla?
\begin{enumerate}[label=\textbf{\Alph*)}]
\begin{multicols}{2}
  \item \(15 \, 	ext{km/h}\)
  \item \(25 \, 	ext{km/h}\)
  \item \(20 \, 	ext{km/h}\)
  \item \(5 \, 	ext{km/h}\)
  \item \(100 \, 	ext{km/h}\)
  \item \(10 \, 	ext{km/h}\)
\end{multicols}
\end{enumerate}
\vspace{0.15cm}
\item ¿Qué es la aceleración?
\begin{enumerate}[label=\textbf{\Alph*)}]
\begin{multicols}{2}
  \item La variación de la velocidad por unidad de tiempo
  \item La variación de la posición por unidad de tiempo
  \item La variación de la masa por unidad de tiempo
  \item La variación de la fuerza por unidad de tiempo
  \item La variación de la energía por unidad de tiempo
  \item La variación de la potencia por unidad de tiempo
\end{multicols}
\end{enumerate}
\vspace{0.15cm}
\item ¿Cuál es la fórmula de la aceleración?
\begin{enumerate}[label=\textbf{\Alph*)}]
\begin{multicols}{2}
  \item \(a = rac{v_f - v_i}{t}\)
  \item \(a = rac{v_f + v_i}{t}\)
  \item \(a = rac{v_f \cdot v_i}{t}\)
  \item \(a = rac{v_f}{v_i \cdot t}\)
  \item \(a = rac{v_i}{v_f \cdot t}\)
  \item \(a = rac{t}{v_f - v_i}\)
\end{multicols}
\end{enumerate}
\vspace{0.15cm}
\item ¿Cuál es la unidad de aceleración en el Sistema Internacional?
\begin{enumerate}[label=\textbf{\Alph*)}]
\begin{multicols}{2}
  \item \( 	ext{m/s}^2\)
  \item \( 	ext{m/s}\)
  \item \( 	ext{m}\)
  \item \( 	ext{s}\)
  \item \( 	ext{m}^2/	ext{s}\)
  \item \( 	ext{m} \cdot 	ext{s}\)
\end{multicols}
\end{enumerate}
\vspace{0.15cm}
\item Un auto pasa de \(0 \, 	ext{m/s}\) a \(20 \, 	ext{m/s}\) en \(5 \, 	ext{s}\). ¿Cuál es su aceleración?
\begin{enumerate}[label=\textbf{\Alph*)}]
\begin{multicols}{2}
  \item \(4 \, 	ext{m/s}^2\)
  \item \(5 \, 	ext{m/s}^2\)
  \item \(2 \, 	ext{m/s}^2\)
  \item \(10 \, 	ext{m/s}^2\)
  \item \(20 \, 	ext{m/s}^2\)
  \item \(1 \, 	ext{m/s}^2\)
\end{multicols}
\end{enumerate}
\vspace{0.15cm}
\item Un auto pasa de \(10 \, 	ext{m/s}\) a \(30 \, 	ext{m/s}\) en \(4 \, 	ext{s}\). ¿Cuál es su aceleración?
\begin{enumerate}[label=\textbf{\Alph*)}]
\begin{multicols}{2}
  \item \(5 \, 	ext{m/s}^2\)
  \item \(10 \, 	ext{m/s}^2\)
  \item \(4 \, 	ext{m/s}^2\)
  \item \(2 \, 	ext{m/s}^2\)
  \item \(20 \, 	ext{m/s}^2\)
  \item \(1 \, 	ext{m/s}^2\)
\end{multicols}
\end{enumerate}
\vspace{0.15cm}
\item Un auto pasa de \(20 \, 	ext{m/s}\) a \(0 \, 	ext{m/s}\) en \(8 \, 	ext{s}\). ¿Cuál es su aceleración?
\begin{enumerate}[label=\textbf{\Alph*)}]
\begin{multicols}{2}
  \item \(-2.5 \, 	ext{m/s}^2\)
  \item \(2.5 \, 	ext{m/s}^2\)
  \item \(-5 \, 	ext{m/s}^2\)
  \item \(5 \, 	ext{m/s}^2\)
  \item \(-10 \, 	ext{m/s}^2\)
  \item \(10 \, 	ext{m/s}^2\)
\end{multicols}
\end{enumerate}
\vspace{0.15cm}
\item Un ciclista pasa de \(5 \, 	ext{m/s}\) a \(15 \, 	ext{m/s}\) en \(2 \, 	ext{s}\). ¿Cuál es su aceleración?
\begin{enumerate}[label=\textbf{\Alph*)}]
\begin{multicols}{2}
  \item \(5 \, 	ext{m/s}^2\)
  \item \(10 \, 	ext{m/s}^2\)
  \item \(2 \, 	ext{m/s}^2\)
  \item \(4 \, 	ext{m/s}^2\)
  \item \(20 \, 	ext{m/s}^2\)
  \item \(1 \, 	ext{m/s}^2\)
\end{multicols}
\end{enumerate}
\vspace{0.15cm}
\item Un tren pasa de \(30 \, 	ext{m/s}\) a \(50 \, 	ext{m/s}\) en \(10 \, 	ext{s}\). ¿Cuál es su aceleración?
\begin{enumerate}[label=\textbf{\Alph*)}]
\begin{multicols}{2}
  \item \(2 \, 	ext{m/s}^2\)
  \item \(5 \, 	ext{m/s}^2\)
  \item \(10 \, 	ext{m/s}^2\)
  \item \(4 \, 	ext{m/s}^2\)
  \item \(20 \, 	ext{m/s}^2\)
  \item \(1 \, 	ext{m/s}^2\)
\end{multicols}
\end{enumerate}
\vspace{0.15cm}
\item ¿Qué es el movimiento rectilíneo uniformemente variado (MRUV)?
\begin{enumerate}[label=\textbf{\Alph*)}]
\begin{multicols}{2}
  \item El movimiento con aceleración constante en línea recta
  \item El movimiento con velocidad constante en línea recta
  \item El movimiento con aceleración variable en línea recta
  \item El movimiento con velocidad variable en línea recta
  \item El movimiento circular con velocidad constante
  \item El movimiento circular con aceleración constante
\end{multicols}
\end{enumerate}
\vspace{0.15cm}
\item ¿Cuál es la fórmula de la velocidad final en el MRUV?
\begin{enumerate}[label=\textbf{\Alph*)}]
\begin{multicols}{2}
  \item \(v_f = v_i + a \cdot t\)
  \item \(v_f = v_i - a \cdot t\)
  \item \(v_f = v_i \cdot a \cdot t\)
  \item \(v_f = v_i / (a \cdot t)\)
  \item \(v_f = a \cdot t / v_i\)
  \item \(v_f = v_i + a / t\)
\end{multicols}
\end{enumerate}
\vspace{0.15cm}
\item ¿Cuál es la fórmula de la distancia en el MRUV?
\begin{enumerate}[label=\textbf{\Alph*)}]
\begin{multicols}{2}
  \item \(d = v_i \cdot t + rac{1}{2} a \cdot t^2\)
  \item \(d = v_i \cdot t - rac{1}{2} a \cdot t^2\)
  \item \(d = v_i \cdot t \cdot rac{1}{2} a \cdot t^2\)
  \item \(d = v_i \cdot t / (rac{1}{2} a \cdot t^2)\)
  \item \(d = rac{1}{2} a \cdot t^2 / (v_i \cdot t)\)
  \item \(d = v_i \cdot t + a \cdot t^2\)
\end{multicols}
\end{enumerate}
\vspace{0.15cm}
\item ¿Cuál es la fórmula de la velocidad final al cuadrado en el MRUV?
\begin{enumerate}[label=\textbf{\Alph*)}]
\begin{multicols}{2}
  \item \(v_f^2 = v_i^2 + 2 \cdot a \cdot d\)
  \item \(v_f^2 = v_i^2 - 2 \cdot a \cdot d\)
  \item \(v_f^2 = v_i^2 \cdot 2 \cdot a \cdot d\)
  \item \(v_f^2 = v_i^2 / (2 \cdot a \cdot d)\)
  \item \(v_f^2 = 2 \cdot a \cdot d / v_i^2\)
  \item \(v_f^2 = v_i^2 + a \cdot d\)
\end{multicols}
\end{enumerate}
\vspace{0.15cm}
\item Un auto parte del reposo con aceleración de \(2 \, 	ext{m/s}^2\) durante \(5 \, 	ext{s}\). ¿Cuál es su velocidad final?
\begin{enumerate}[label=\textbf{\Alph*)}]
\begin{multicols}{2}
  \item \(10 \, 	ext{m/s}\)
  \item \(5 \, 	ext{m/s}\)
  \item \(20 \, 	ext{m/s}\)
  \item \(2 \, 	ext{m/s}\)
  \item \(25 \, 	ext{m/s}\)
  \item \(15 \, 	ext{m/s}\)
\end{multicols}
\end{enumerate}
\vspace{0.15cm}
\item Un auto parte del reposo con aceleración de \(3 \, 	ext{m/s}^2\) durante \(4 \, 	ext{s}\). ¿Cuál es su velocidad final?
\begin{enumerate}[label=\textbf{\Alph*)}]
\begin{multicols}{2}
  \item \(12 \, 	ext{m/s}\)
  \item \(6 \, 	ext{m/s}\)
  \item \(24 \, 	ext{m/s}\)
  \item \(3 \, 	ext{m/s}\)
  \item \(36 \, 	ext{m/s}\)
  \item \(9 \, 	ext{m/s}\)
\end{multicols}
\end{enumerate}
\vspace{0.15cm}
\item Un auto parte del reposo con aceleración de \(4 \, 	ext{m/s}^2\) durante \(6 \, 	ext{s}\). ¿Cuál es su velocidad final?
\begin{enumerate}[label=\textbf{\Alph*)}]
\begin{multicols}{2}
  \item \(24 \, 	ext{m/s}\)
  \item \(12 \, 	ext{m/s}\)
  \item \(48 \, 	ext{m/s}\)
  \item \(4 \, 	ext{m/s}\)
  \item \(36 \, 	ext{m/s}\)
  \item \(18 \, 	ext{m/s}\)
\end{multicols}
\end{enumerate}
\vspace{0.15cm}
\item Un auto parte del reposo con aceleración de \(5 \, 	ext{m/s}^2\) durante \(8 \, 	ext{s}\). ¿Cuál es su velocidad final?
\begin{enumerate}[label=\textbf{\Alph*)}]
\begin{multicols}{2}
  \item \(40 \, 	ext{m/s}\)
  \item \(20 \, 	ext{m/s}\)
  \item \(80 \, 	ext{m/s}\)
  \item \(5 \, 	ext{m/s}\)
  \item \(60 \, 	ext{m/s}\)
  \item \(30 \, 	ext{m/s}\)
\end{multicols}
\end{enumerate}
\vspace{0.15cm}
\item Un auto parte del reposo con aceleración de \(2 \, 	ext{m/s}^2\) durante \(5 \, 	ext{s}\). ¿Qué distancia recorre?
\begin{enumerate}[label=\textbf{\Alph*)}]
\begin{multicols}{2}
  \item \(25 \, 	ext{m}\)
  \item \(10 \, 	ext{m}\)
  \item \(50 \, 	ext{m}\)
  \item \(5 \, 	ext{m}\)
  \item \(20 \, 	ext{m}\)
  \item \(15 \, 	ext{m}\)
\end{multicols}
\end{enumerate}
\vspace{0.15cm}
\item Un auto parte del reposo con aceleración de \(3 \, 	ext{m/s}^2\) durante \(4 \, 	ext{s}\). ¿Qué distancia recorre?
\begin{enumerate}[label=\textbf{\Alph*)}]
\begin{multicols}{2}
  \item \(24 \, 	ext{m}\)
  \item \(12 \, 	ext{m}\)
  \item \(48 \, 	ext{m}\)
  \item \(6 \, 	ext{m}\)
  \item \(36 \, 	ext{m}\)
  \item \(18 \, 	ext{m}\)
\end{multicols}
\end{enumerate}
\vspace{0.15cm}
\item Un auto parte del reposo con aceleración de \(4 \, 	ext{m/s}^2\) durante \(6 \, 	ext{s}\). ¿Qué distancia recorre?
\begin{enumerate}[label=\textbf{\Alph*)}]
\begin{multicols}{2}
  \item \(72 \, 	ext{m}\)
  \item \(36 \, 	ext{m}\)
  \item \(144 \, 	ext{m}\)
  \item \(24 \, 	ext{m}\)
  \item \(108 \, 	ext{m}\)
  \item \(48 \, 	ext{m}\)
\end{multicols}
\end{enumerate}
\vspace{0.15cm}
\item Un auto parte del reposo con aceleración de \(5 \, 	ext{m/s}^2\) durante \(8 \, 	ext{s}\). ¿Qué distancia recorre?
\begin{enumerate}[label=\textbf{\Alph*)}]
\begin{multicols}{2}
  \item \(160 \, 	ext{m}\)
  \item \(80 \, 	ext{m}\)
  \item \(320 \, 	ext{m}\)
  \item \(40 \, 	ext{m}\)
  \item \(240 \, 	ext{m}\)
  \item \(120 \, 	ext{m}\)
\end{multicols}
\end{enumerate}
\vspace{0.15cm}
\item Un auto viaja a \(20 \, 	ext{m/s}\) y acelera a \(3 \, 	ext{m/s}^2\) durante \(4 \, 	ext{s}\). ¿Cuál es su velocidad final?
\begin{enumerate}[label=\textbf{\Alph*)}]
\begin{multicols}{2}
  \item \(32 \, 	ext{m/s}\)
  \item \(12 \, 	ext{m/s}\)
  \item \(44 \, 	ext{m/s}\)
  \item \(8 \, 	ext{m/s}\)
  \item \(26 \, 	ext{m/s}\)
  \item \(36 \, 	ext{m/s}\)
\end{multicols}
\end{enumerate}
\vspace{0.15cm}
\item Un auto viaja a \(15 \, 	ext{m/s}\) y acelera a \(2 \, 	ext{m/s}^2\) durante \(5 \, 	ext{s}\). ¿Cuál es su velocidad final?
\begin{enumerate}[label=\textbf{\Alph*)}]
\begin{multicols}{2}
  \item \(25 \, 	ext{m/s}\)
  \item \(10 \, 	ext{m/s}\)
  \item \(35 \, 	ext{m/s}\)
  \item \(5 \, 	ext{m/s}\)
  \item \(20 \, 	ext{m/s}\)
  \item \(30 \, 	ext{m/s}\)
\end{multicols}
\end{enumerate}
\vspace{0.15cm}
\item Un auto viaja a \(10 \, 	ext{m/s}\) y acelera a \(4 \, 	ext{m/s}^2\) durante \(3 \, 	ext{s}\). ¿Cuál es su velocidad final?
\begin{enumerate}[label=\textbf{\Alph*)}]
\begin{multicols}{2}
  \item \(22 \, 	ext{m/s}\)
  \item \(12 \, 	ext{m/s}\)
  \item \(34 \, 	ext{m/s}\)
  \item \(6 \, 	ext{m/s}\)
  \item \(14 \, 	ext{m/s}\)
  \item \(26 \, 	ext{m/s}\)
\end{multicols}
\end{enumerate}
\vspace{0.15cm}
\item Un auto viaja a \(30 \, 	ext{m/s}\) y frena a \(-5 \, 	ext{m/s}^2\) durante \(4 \, 	ext{s}\). ¿Cuál es su velocidad final?
\begin{enumerate}[label=\textbf{\Alph*)}]
\begin{multicols}{2}
  \item \(10 \, 	ext{m/s}\)
  \item \(20 \, 	ext{m/s}\)
  \item \(0 \, 	ext{m/s}\)
  \item \(5 \, 	ext{m/s}\)
  \item \(15 \, 	ext{m/s}\)
  \item \(25 \, 	ext{m/s}\)
\end{multicols}
\end{enumerate}
\vspace{0.15cm}
\item Un auto viaja a \(40 \, 	ext{m/s}\) y frena a \(-8 \, 	ext{m/s}^2\) durante \(5 \, 	ext{s}\). ¿Cuál es su velocidad final?
\begin{enumerate}[label=\textbf{\Alph*)}]
\begin{multicols}{2}
  \item \(0 \, 	ext{m/s}\)
  \item \(10 \, 	ext{m/s}\)
  \item \(20 \, 	ext{m/s}\)
  \item \(5 \, 	ext{m/s}\)
  \item \(15 \, 	ext{m/s}\)
  \item \(25 \, 	ext{m/s}\) ---
\end{multicols}
\end{enumerate}
\vspace{0.15cm}
\end{enumerate}\n
\subsubsection*{SECCIÓN IV: CAÍDA LIBRE Y TIRO VERTICAL (181–240)}

\begin{enumerate}[label=\textbf{\arabic*.}]
\setcounter{enumi}{180}
\item ¿Qué es la caída libre?
\begin{enumerate}[label=\textbf{\Alph*)}]
\begin{multicols}{2}
  \item El movimiento de un cuerpo bajo la acción de la gravedad, sin resistencia del aire
  \item El movimiento de un cuerpo con velocidad constante
  \item El movimiento de un cuerpo con aceleración variable
  \item El movimiento de un cuerpo en círculos
  \item El movimiento de un cuerpo en el espacio
  \item El movimiento de un cuerpo en el agua
\end{multicols}
\end{enumerate}
\vspace{0.15cm}
\item ¿Cuál es la aceleración en la caída libre?
\begin{enumerate}[label=\textbf{\Alph*)}]
\begin{multicols}{2}
  \item \(9.8 \, 	ext{m/s}^2\) hacia abajo
  \item \(9.8 \, 	ext{m/s}^2\) hacia arriba
  \item \(0 \, 	ext{m/s}^2\)
  \item \(10 \, 	ext{m/s}^2\) hacia arriba
  \item \(5 \, 	ext{m/s}^2\) hacia abajo
  \item \(20 \, 	ext{m/s}^2\) hacia abajo
\end{multicols}
\end{enumerate}
\vspace{0.15cm}
\item ¿Cuál es la fórmula de la velocidad en la caída libre?
\begin{enumerate}[label=\textbf{\Alph*)}]
\begin{multicols}{2}
  \item \(v = g \cdot t\)
  \item \(v = g / t\)
  \item \(v = t / g\)
  \item \(v = g + t\)
  \item \(v = g - t\)
  \item \(v = g \cdot t^2\)
\end{multicols}
\end{enumerate}
\vspace{0.15cm}
\item ¿Cuál es la fórmula de la distancia en la caída libre?
\begin{enumerate}[label=\textbf{\Alph*)}]
\begin{multicols}{2}
  \item \(h = rac{1}{2} g \cdot t^2\)
  \item \(h = g \cdot t^2\)
  \item \(h = rac{1}{2} g \cdot t\)
  \item \(h = g \cdot t\)
  \item \(h = rac{1}{2} g^2 \cdot t\)
  \item \(h = g^2 \cdot t^2\)
\end{multicols}
\end{enumerate}
\vspace{0.15cm}
\item ¿Cuál es la fórmula de la velocidad final en la caída libre?
\begin{enumerate}[label=\textbf{\Alph*)}]
\begin{multicols}{2}
  \item \(v = \sqrt{2 g h}\)
  \item \(v = 2 g h\)
  \item \(v = \sqrt{g h}\)
  \item \(v = g h\)
  \item \(v = \sqrt{2 g / h}\)
  \item \(v = 2 g / h\)
\end{multicols}
\end{enumerate}
\vspace{0.15cm}
\item Un objeto cae libremente desde el reposo durante \(2 \, 	ext{s}\). ¿Cuál es su velocidad final?
\begin{enumerate}[label=\textbf{\Alph*)}]
\begin{multicols}{2}
  \item \(19.6 \, 	ext{m/s}\)
  \item \(9.8 \, 	ext{m/s}\)
  \item \(4.9 \, 	ext{m/s}\)
  \item \(39.2 \, 	ext{m/s}\)
  \item \(29.4 \, 	ext{m/s}\)
  \item \(14.7 \, 	ext{m/s}\)
\end{multicols}
\end{enumerate}
\vspace{0.15cm}
\item Un objeto cae libremente desde el reposo durante \(3 \, 	ext{s}\). ¿Cuál es su velocidad final?
\begin{enumerate}[label=\textbf{\Alph*)}]
\begin{multicols}{2}
  \item \(29.4 \, 	ext{m/s}\)
  \item \(9.8 \, 	ext{m/s}\)
  \item \(19.6 \, 	ext{m/s}\)
  \item \(39.2 \, 	ext{m/s}\)
  \item \(49.0 \, 	ext{m/s}\)
  \item \(14.7 \, 	ext{m/s}\)
\end{multicols}
\end{enumerate}
\vspace{0.15cm}
\item Un objeto cae libremente desde el reposo durante \(4 \, 	ext{s}\). ¿Cuál es su velocidad final?
\begin{enumerate}[label=\textbf{\Alph*)}]
\begin{multicols}{2}
  \item \(39.2 \, 	ext{m/s}\)
  \item \(9.8 \, 	ext{m/s}\)
  \item \(19.6 \, 	ext{m/s}\)
  \item \(29.4 \, 	ext{m/s}\)
  \item \(49.0 \, 	ext{m/s}\)
  \item \(14.7 \, 	ext{m/s}\)
\end{multicols}
\end{enumerate}
\vspace{0.15cm}
\item Un objeto cae libremente desde el reposo durante \(5 \, 	ext{s}\). ¿Cuál es su velocidad final?
\begin{enumerate}[label=\textbf{\Alph*)}]
\begin{multicols}{2}
  \item \(49.0 \, 	ext{m/s}\)
  \item \(9.8 \, 	ext{m/s}\)
  \item \(19.6 \, 	ext{m/s}\)
  \item \(29.4 \, 	ext{m/s}\)
  \item \(39.2 \, 	ext{m/s}\)
  \item \(14.7 \, 	ext{m/s}\)
\end{multicols}
\end{enumerate}
\vspace{0.15cm}
\item Un objeto cae libremente desde el reposo durante \(2 \, 	ext{s}\). ¿Qué distancia recorre?
\begin{enumerate}[label=\textbf{\Alph*)}]
\begin{multicols}{2}
  \item \(19.6 \, 	ext{m}\)
  \item \(9.8 \, 	ext{m}\)
  \item \(4.9 \, 	ext{m}\)
  \item \(39.2 \, 	ext{m}\)
  \item \(29.4 \, 	ext{m}\)
  \item \(14.7 \, 	ext{m}\)
\end{multicols}
\end{enumerate}
\vspace{0.15cm}
\item Un objeto cae libremente desde el reposo durante \(3 \, 	ext{s}\). ¿Qué distancia recorre?
\begin{enumerate}[label=\textbf{\Alph*)}]
\begin{multicols}{2}
  \item \(44.1 \, 	ext{m}\)
  \item \(9.8 \, 	ext{m}\)
  \item \(19.6 \, 	ext{m}\)
  \item \(29.4 \, 	ext{m}\)
  \item \(39.2 \, 	ext{m}\)
  \item \(14.7 \, 	ext{m}\)
\end{multicols}
\end{enumerate}
\vspace{0.15cm}
\item Un objeto cae libremente desde el reposo durante \(4 \, 	ext{s}\). ¿Qué distancia recorre?
\begin{enumerate}[label=\textbf{\Alph*)}]
\begin{multicols}{2}
  \item \(78.4 \, 	ext{m}\)
  \item \(9.8 \, 	ext{m}\)
  \item \(19.6 \, 	ext{m}\)
  \item \(29.4 \, 	ext{m}\)
  \item \(39.2 \, 	ext{m}\)
  \item \(14.7 \, 	ext{m}\)
\end{multicols}
\end{enumerate}
\vspace{0.15cm}
\item Un objeto cae libremente desde el reposo durante \(5 \, 	ext{s}\). ¿Qué distancia recorre?
\begin{enumerate}[label=\textbf{\Alph*)}]
\begin{multicols}{2}
  \item \(122.5 \, 	ext{m}\)
  \item \(9.8 \, 	ext{m}\)
  \item \(19.6 \, 	ext{m}\)
  \item \(29.4 \, 	ext{m}\)
  \item \(39.2 \, 	ext{m}\)
  \item \(14.7 \, 	ext{m}\)
\end{multicols}
\end{enumerate}
\vspace{0.15cm}
\item Un objeto se lanza verticalmente hacia arriba con una velocidad inicial de \(20 \, 	ext{m/s}\). ¿Cuál es su velocidad en el punto más alto?
\begin{enumerate}[label=\textbf{\Alph*)}]
\begin{multicols}{2}
  \item \(0 \, 	ext{m/s}\)
  \item \(20 \, 	ext{m/s}\)
  \item \(10 \, 	ext{m/s}\)
  \item \(-20 \, 	ext{m/s}\)
  \item \(9.8 \, 	ext{m/s}\)
  \item \(-9.8 \, 	ext{m/s}\)
\end{multicols}
\end{enumerate}
\vspace{0.15cm}
\item Un objeto se lanza verticalmente hacia arriba con una velocidad inicial de \(20 \, 	ext{m/s}\). ¿Cuánto tiempo tarda en alcanzar el punto más alto?
\begin{enumerate}[label=\textbf{\Alph*)}]
\begin{multicols}{2}
  \item \(2.04 \, 	ext{s}\)
  \item \(1.02 \, 	ext{s}\)
  \item \(4.08 \, 	ext{s}\)
  \item \(3.06 \, 	ext{s}\)
  \item \(5.10 \, 	ext{s}\)
  \item \(0.51 \, 	ext{s}\)
\end{multicols}
\end{enumerate}
\vspace{0.15cm}
\item Un objeto se lanza verticalmente hacia arriba con una velocidad inicial de \(30 \, 	ext{m/s}\). ¿Cuánto tiempo tarda en alcanzar el punto más alto?
\begin{enumerate}[label=\textbf{\Alph*)}]
\begin{multicols}{2}
  \item \(3.06 \, 	ext{s}\)
  \item \(1.53 \, 	ext{s}\)
  \item \(6.12 \, 	ext{s}\)
  \item \(4.59 \, 	ext{s}\)
  \item \(7.65 \, 	ext{s}\)
  \item \(0.77 \, 	ext{s}\)
\end{multicols}
\end{enumerate}
\vspace{0.15cm}
\item Un objeto se lanza verticalmente hacia arriba con una velocidad inicial de \(40 \, 	ext{m/s}\). ¿Cuánto tiempo tarda en alcanzar el punto más alto?
\begin{enumerate}[label=\textbf{\Alph*)}]
\begin{multicols}{2}
  \item \(4.08 \, 	ext{s}\)
  \item \(2.04 \, 	ext{s}\)
  \item \(8.16 \, 	ext{s}\)
  \item \(6.12 \, 	ext{s}\)
  \item \(10.20 \, 	ext{s}\)
  \item \(1.02 \, 	ext{s}\)
\end{multicols}
\end{enumerate}
\vspace{0.15cm}
\item Un objeto se lanza verticalmente hacia arriba con una velocidad inicial de \(20 \, 	ext{m/s}\). ¿Cuál es la altura máxima alcanzada?
\begin{enumerate}[label=\textbf{\Alph*)}]
\begin{multicols}{2}
  \item \(20.4 \, 	ext{m}\)
  \item \(10.2 \, 	ext{m}\)
  \item \(40.8 \, 	ext{m}\)
  \item \(30.6 \, 	ext{m}\)
  \item \(50.0 \, 	ext{m}\)
  \item \(15.0 \, 	ext{m}\)
\end{multicols}
\end{enumerate}
\vspace{0.15cm}
\item Un objeto se lanza verticalmente hacia arriba con una velocidad inicial de \(30 \, 	ext{m/s}\). ¿Cuál es la altura máxima alcanzada?
\begin{enumerate}[label=\textbf{\Alph*)}]
\begin{multicols}{2}
  \item \(45.9 \, 	ext{m}\)
  \item \(22.9 \, 	ext{m}\)
  \item \(91.8 \, 	ext{m}\)
  \item \(68.9 \, 	ext{m}\)
  \item \(100.0 \, 	ext{m}\)
  \item \(30.0 \, 	ext{m}\)
\end{multicols}
\end{enumerate}
\vspace{0.15cm}
\item Un objeto se lanza verticalmente hacia arriba con una velocidad inicial de \(40 \, 	ext{m/s}\). ¿Cuál es la altura máxima alcanzada?
\begin{enumerate}[label=\textbf{\Alph*)}]
\begin{multicols}{2}
  \item \(81.6 \, 	ext{m}\)
  \item \(40.8 \, 	ext{m}\)
  \item \(163.2 \, 	ext{m}\)
  \item \(122.4 \, 	ext{m}\)
  \item \(200.0 \, 	ext{m}\)
  \item \(60.0 \, 	ext{m}\)
\end{multicols}
\end{enumerate}
\vspace{0.15cm}
\item ¿Qué es el tiro parabólico?
\begin{enumerate}[label=\textbf{\Alph*)}]
\begin{multicols}{2}
  \item El movimiento de un cuerpo lanzado con una velocidad inicial que forma un ángulo con la horizontal
  \item El movimiento de un cuerpo lanzado verticalmente
  \item El movimiento de un cuerpo lanzado horizontalmente
  \item El movimiento de un cuerpo en círculos
  \item El movimiento de un cuerpo en el espacio
  \item El movimiento de un cuerpo en el agua
\end{multicols}
\end{enumerate}
\vspace{0.15cm}
\item ¿Cuál es la componente horizontal de la velocidad en el tiro parabólico?
\begin{enumerate}[label=\textbf{\Alph*)}]
\begin{multicols}{2}
  \item \(v_x = v_0 \cos	heta\) (constante)
  \item \(v_x = v_0 \sen	heta\) (constante)
  \item \(v_x = v_0 	an	heta\) (constante)
  \item \(v_x = v_0 \cos	heta\) (variable)
  \item \(v_x = v_0 \sen	heta\) (variable)
  \item \(v_x = v_0 	an	heta\) (variable)
\end{multicols}
\end{enumerate}
\vspace{0.15cm}
\item ¿Cuál es la componente vertical de la velocidad en el tiro parabólico?
\begin{enumerate}[label=\textbf{\Alph*)}]
\begin{multicols}{2}
  \item \(v_y = v_0 \sen	heta - g t\) (variable)
  \item \(v_y = v_0 \cos	heta - g t\) (variable)
  \item \(v_y = v_0 	an	heta - g t\) (variable)
  \item \(v_y = v_0 \sen	heta\) (constante)
  \item \(v_y = v_0 \cos	heta\) (constante)
  \item \(v_y = v_0 	an	heta\) (constante)
\end{multicols}
\end{enumerate}
\vspace{0.15cm}
\item ¿Cuál es el alcance horizontal máximo en el tiro parabólico?
\begin{enumerate}[label=\textbf{\Alph*)}]
\begin{multicols}{2}
  \item \(R = rac{v_0^2 \sen 2	heta}{g}\)
  \item \(R = rac{v_0^2 \cos 2	heta}{g}\)
  \item \(R = rac{v_0^2 	an 2	heta}{g}\)
  \item \(R = rac{v_0 \sen 2	heta}{g}\)
  \item \(R = rac{v_0 \cos 2	heta}{g}\)
  \item \(R = rac{v_0 	an 2	heta}{g}\)
\end{multicols}
\end{enumerate}
\vspace{0.15cm}
\item ¿Cuál es la altura máxima en el tiro parabólico?
\begin{enumerate}[label=\textbf{\Alph*)}]
\begin{multicols}{2}
  \item \(H = rac{v_0^2 \sen^2	heta}{2g}\)
  \item \(H = rac{v_0^2 \cos^2	heta}{2g}\)
  \item \(H = rac{v_0^2 	an^2	heta}{2g}\)
  \item \(H = rac{v_0 \sen^2	heta}{2g}\)
  \item \(H = rac{v_0 \cos^2	heta}{2g}\)
  \item \(H = rac{v_0 	an^2	heta}{2g}\)
\end{multicols}
\end{enumerate}
\vspace{0.15cm}
\item ¿Cuál es el tiempo de vuelo en el tiro parabólico?
\begin{enumerate}[label=\textbf{\Alph*)}]
\begin{multicols}{2}
  \item \(t_v = rac{2 v_0 \sen	heta}{g}\)
  \item \(t_v = rac{2 v_0 \cos	heta}{g}\)
  \item \(t_v = rac{2 v_0 	an	heta}{g}\)
  \item \(t_v = rac{v_0 \sen	heta}{g}\)
  \item \(t_v = rac{v_0 \cos	heta}{g}\)
  \item \(t_v = rac{v_0 	an	heta}{g}\)
\end{multicols}
\end{enumerate}
\vspace{0.15cm}
\item Un proyectil se lanza con \(v_0 = 20 \, 	ext{m/s}\) y \(	heta = 30^\circ\). ¿Cuál es su alcance máximo?
\begin{enumerate}[label=\textbf{\Alph*)}]
\begin{multicols}{2}
  \item \(35.35 \, 	ext{m}\)
  \item \(17.68 \, 	ext{m}\)
  \item \(70.70 \, 	ext{m}\)
  \item \(53.03 \, 	ext{m}\)
  \item \(88.38 \, 	ext{m}\)
  \item \(26.52 \, 	ext{m}\)
\end{multicols}
\end{enumerate}
\vspace{0.15cm}
\item Un proyectil se lanza con \(v_0 = 30 \, 	ext{m/s}\) y \(	heta = 45^\circ\). ¿Cuál es su alcance máximo?
\begin{enumerate}[label=\textbf{\Alph*)}]
\begin{multicols}{2}
  \item \(91.84 \, 	ext{m}\)
  \item \(45.92 \, 	ext{m}\)
  \item \(183.68 \, 	ext{m}\)
  \item \(137.76 \, 	ext{m}\)
  \item \(229.60 \, 	ext{m}\)
  \item \(68.88 \, 	ext{m}\)
\end{multicols}
\end{enumerate}
\vspace{0.15cm}
\item Un proyectil se lanza con \(v_0 = 40 \, 	ext{m/s}\) y \(	heta = 60^\circ\). ¿Cuál es su alcance máximo?
\begin{enumerate}[label=\textbf{\Alph*)}]
\begin{multicols}{2}
  \item \(141.39 \, 	ext{m}\)
  \item \(70.70 \, 	ext{m}\)
  \item \(282.78 \, 	ext{m}\)
  \item \(212.09 \, 	ext{m}\)
  \item \(353.48 \, 	ext{m}\)
  \item \(106.04 \, 	ext{m}\)
\end{multicols}
\end{enumerate}
\vspace{0.15cm}
\item Un proyectil se lanza con \(v_0 = 50 \, 	ext{m/s}\) y \(	heta = 30^\circ\). ¿Cuál es su altura máxima?
\begin{enumerate}[label=\textbf{\Alph*)}]
\begin{multicols}{2}
  \item \(31.89 \, 	ext{m}\)
  \item \(15.94 \, 	ext{m}\)
  \item \(63.78 \, 	ext{m}\)
  \item \(47.83 \, 	ext{m}\)
  \item \(79.72 \, 	ext{m}\)
  \item \(23.92 \, 	ext{m}\)
\end{multicols}
\end{enumerate}
\vspace{0.15cm}
\item Un proyectil se lanza con \(v_0 = 30 \, 	ext{m/s}\) y \(	heta = 45^\circ\). ¿Cuál es su altura máxima?
\begin{enumerate}[label=\textbf{\Alph*)}]
\begin{multicols}{2}
  \item \(22.96 \, 	ext{m}\)
  \item \(11.48 \, 	ext{m}\)
  \item \(45.92 \, 	ext{m}\)
  \item \(34.44 \, 	ext{m}\)
  \item \(57.40 \, 	ext{m}\)
  \item \(17.22 \, 	ext{m}\)
\end{multicols}
\end{enumerate}
\vspace{0.15cm}
\item Un proyectil se lanza con \(v_0 = 40 \, 	ext{m/s}\) y \(	heta = 60^\circ\). ¿Cuál es su altura máxima?
\begin{enumerate}[label=\textbf{\Alph*)}]
\begin{multicols}{2}
  \item \(61.22 \, 	ext{m}\)
  \item \(30.61 \, 	ext{m}\)
  \item \(122.44 \, 	ext{m}\)
  \item \(91.83 \, 	ext{m}\)
  \item \(153.05 \, 	ext{m}\)
  \item \(45.92 \, 	ext{m}\)
\end{multicols}
\end{enumerate}
\vspace{0.15cm}
\item Un proyectil se lanza con \(v_0 = 20 \, 	ext{m/s}\) y \(	heta = 30^\circ\). ¿Cuál es su tiempo de vuelo?
\begin{enumerate}[label=\textbf{\Alph*)}]
\begin{multicols}{2}
  \item \(2.04 \, 	ext{s}\)
  \item \(1.02 \, 	ext{s}\)
  \item \(4.08 \, 	ext{s}\)
  \item \(3.06 \, 	ext{s}\)
  \item \(5.10 \, 	ext{s}\)
  \item \(0.51 \, 	ext{s}\)
\end{multicols}
\end{enumerate}
\vspace{0.15cm}
\item Un proyectil se lanza con \(v_0 = 30 \, 	ext{m/s}\) y \(	heta = 45^\circ\). ¿Cuál es su tiempo de vuelo?
\begin{enumerate}[label=\textbf{\Alph*)}]
\begin{multicols}{2}
  \item \(4.33 \, 	ext{s}\)
  \item \(2.16 \, 	ext{s}\)
  \item \(8.66 \, 	ext{s}\)
  \item \(6.49 \, 	ext{s}\)
  \item \(10.82 \, 	ext{s}\)
  \item \(1.08 \, 	ext{s}\)
\end{multicols}
\end{enumerate}
\vspace{0.15cm}
\item Un proyectil se lanza con \(v_0 = 40 \, 	ext{m/s}\) y \(	heta = 60^\circ\). ¿Cuál es su tiempo de vuelo?
\begin{enumerate}[label=\textbf{\Alph*)}]
\begin{multicols}{2}
  \item \(7.07 \, 	ext{s}\)
  \item \(3.54 \, 	ext{s}\)
  \item \(14.14 \, 	ext{s}\)
  \item \(10.61 \, 	ext{s}\)
  \item \(17.68 \, 	ext{s}\)
  \item \(1.77 \, 	ext{s}\)
\end{multicols}
\end{enumerate}
\vspace{0.15cm}
\item ¿Qué es la velocidad angular?
\begin{enumerate}[label=\textbf{\Alph*)}]
\begin{multicols}{2}
  \item La variación del ángulo por unidad de tiempo
  \item La variación de la posición por unidad de tiempo
  \item La variación de la velocidad por unidad de tiempo
  \item La variación de la masa por unidad de tiempo
  \item La variación de la fuerza por unidad de tiempo
  \item La variación de la energía por unidad de tiempo
\end{multicols}
\end{enumerate}
\vspace{0.15cm}
\item ¿Cuál es la unidad de velocidad angular?
\begin{enumerate}[label=\textbf{\Alph*)}]
\begin{multicols}{2}
  \item rad/s
  \item m/s
  \item m/s²
  \item rad
  \item m
  \item s
\end{multicols}
\end{enumerate}
\vspace{0.15cm}
\item ¿Qué es el movimiento circular uniforme (MCU)?
\begin{enumerate}[label=\textbf{\Alph*)}]
\begin{multicols}{2}
  \item El movimiento con velocidad angular constante en una trayectoria circular
  \item El movimiento con velocidad angular variable
  \item El movimiento con aceleración angular constante
  \item El movimiento en línea recta
  \item El movimiento en el espacio
  \item El movimiento en el agua
\end{multicols}
\end{enumerate}
\vspace{0.15cm}
\item ¿Cuál es la fórmula de la velocidad tangencial en el MCU?
\begin{enumerate}[label=\textbf{\Alph*)}]
\begin{multicols}{2}
  \item \(v = \omega \cdot r\)
  \item \(v = \omega / r\)
  \item \(v = r / \omega\)
  \item \(v = \omega + r\)
  \item \(v = \omega - r\)
  \item \(v = \omega \cdot r^2\)
\end{multicols}
\end{enumerate}
\vspace{0.15cm}
\item ¿Cuál es la fórmula de la aceleración centrípeta en el MCU?
\begin{enumerate}[label=\textbf{\Alph*)}]
\begin{multicols}{2}
  \item \(a_c = rac{v^2}{r}\)
  \item \(a_c = rac{r}{v^2}\)
  \item \(a_c = v^2 \cdot r\)
  \item \(a_c = rac{v}{r}\)
  \item \(a_c = rac{r}{v}\)
  \item \(a_c = v \cdot r\)
\end{multicols}
\end{enumerate}
\vspace{0.15cm}
\item ¿Qué es la caída libre?
\begin{enumerate}[label=\textbf{\Alph*)}]
\begin{multicols}{2}
  \item El movimiento de un cuerpo bajo la acción de la gravedad, sin resistencia del aire
  \item El movimiento de un cuerpo con velocidad constante
  \item El movimiento de un cuerpo con aceleración variable
  \item El movimiento de un cuerpo en círculos
  \item El movimiento de un cuerpo en el espacio
  \item El movimiento de un cuerpo en el agua
\end{multicols}
\end{enumerate}
\vspace{0.15cm}
\item ¿Cuál es la aceleración en la caída libre?
\begin{enumerate}[label=\textbf{\Alph*)}]
\begin{multicols}{2}
  \item \(9.8 \, 	ext{m/s}^2\) hacia abajo
  \item \(9.8 \, 	ext{m/s}^2\) hacia arriba
  \item \(0 \, 	ext{m/s}^2\)
  \item \(10 \, 	ext{m/s}^2\) hacia arriba
  \item \(5 \, 	ext{m/s}^2\) hacia abajo
  \item \(20 \, 	ext{m/s}^2\) hacia abajo
\end{multicols}
\end{enumerate}
\vspace{0.15cm}
\item ¿Cuál es la fórmula de la velocidad en la caída libre?
\begin{enumerate}[label=\textbf{\Alph*)}]
\begin{multicols}{2}
  \item \(v = g \cdot t\)
  \item \(v = g / t\)
  \item \(v = t / g\)
  \item \(v = g + t\)
  \item \(v = g - t\)
  \item \(v = g \cdot t^2\)
\end{multicols}
\end{enumerate}
\vspace{0.15cm}
\item ¿Cuál es la fórmula de la distancia en la caída libre?
\begin{enumerate}[label=\textbf{\Alph*)}]
\begin{multicols}{2}
  \item \(h = rac{1}{2} g \cdot t^2\)
  \item \(h = g \cdot t^2\)
  \item \(h = rac{1}{2} g \cdot t\)
  \item \(h = g \cdot t\)
  \item \(h = rac{1}{2} g^2 \cdot t\)
  \item \(h = g^2 \cdot t^2\)
\end{multicols}
\end{enumerate}
\vspace{0.15cm}
\item ¿Cuál es la fórmula de la velocidad final en la caída libre?
\begin{enumerate}[label=\textbf{\Alph*)}]
\begin{multicols}{2}
  \item \(v = \sqrt{2 g h}\)
  \item \(v = 2 g h\)
  \item \(v = \sqrt{g h}\)
  \item \(v = g h\)
  \item \(v = \sqrt{2 g / h}\)
  \item \(v = 2 g / h\)
\end{multicols}
\end{enumerate}
\vspace{0.15cm}
\item Un objeto cae libremente desde el reposo durante \(2 \, 	ext{s}\). ¿Cuál es su velocidad final?
\begin{enumerate}[label=\textbf{\Alph*)}]
\begin{multicols}{2}
  \item \(19.6 \, 	ext{m/s}\)
  \item \(9.8 \, 	ext{m/s}\)
  \item \(4.9 \, 	ext{m/s}\)
  \item \(39.2 \, 	ext{m/s}\)
  \item \(29.4 \, 	ext{m/s}\)
  \item \(14.7 \, 	ext{m/s}\)
\end{multicols}
\end{enumerate}
\vspace{0.15cm}
\item Un objeto cae libremente desde el reposo durante \(3 \, 	ext{s}\). ¿Cuál es su velocidad final?
\begin{enumerate}[label=\textbf{\Alph*)}]
\begin{multicols}{2}
  \item \(29.4 \, 	ext{m/s}\)
  \item \(9.8 \, 	ext{m/s}\)
  \item \(19.6 \, 	ext{m/s}\)
  \item \(39.2 \, 	ext{m/s}\)
  \item \(49.0 \, 	ext{m/s}\)
  \item \(14.7 \, 	ext{m/s}\)
\end{multicols}
\end{enumerate}
\vspace{0.15cm}
\item Un objeto cae libremente desde el reposo durante \(4 \, 	ext{s}\). ¿Cuál es su velocidad final?
\begin{enumerate}[label=\textbf{\Alph*)}]
\begin{multicols}{2}
  \item \(39.2 \, 	ext{m/s}\)
  \item \(9.8 \, 	ext{m/s}\)
  \item \(19.6 \, 	ext{m/s}\)
  \item \(29.4 \, 	ext{m/s}\)
  \item \(49.0 \, 	ext{m/s}\)
  \item \(14.7 \, 	ext{m/s}\)
\end{multicols}
\end{enumerate}
\vspace{0.15cm}
\item Un objeto cae libremente desde el reposo durante \(5 \, 	ext{s}\). ¿Cuál es su velocidad final?
\begin{enumerate}[label=\textbf{\Alph*)}]
\begin{multicols}{2}
  \item \(49.0 \, 	ext{m/s}\)
  \item \(9.8 \, 	ext{m/s}\)
  \item \(19.6 \, 	ext{m/s}\)
  \item \(29.4 \, 	ext{m/s}\)
  \item \(39.2 \, 	ext{m/s}\)
  \item \(14.7 \, 	ext{m/s}\)
\end{multicols}
\end{enumerate}
\vspace{0.15cm}
\item Un objeto cae libremente desde el reposo durante \(2 \, 	ext{s}\). ¿Qué distancia recorre?
\begin{enumerate}[label=\textbf{\Alph*)}]
\begin{multicols}{2}
  \item \(19.6 \, 	ext{m}\)
  \item \(9.8 \, 	ext{m}\)
  \item \(4.9 \, 	ext{m}\)
  \item \(39.2 \, 	ext{m}\)
  \item \(29.4 \, 	ext{m}\)
  \item \(14.7 \, 	ext{m}\)
\end{multicols}
\end{enumerate}
\vspace{0.15cm}
\item Un objeto cae libremente desde el reposo durante \(3 \, 	ext{s}\). ¿Qué distancia recorre?
\begin{enumerate}[label=\textbf{\Alph*)}]
\begin{multicols}{2}
  \item \(44.1 \, 	ext{m}\)
  \item \(9.8 \, 	ext{m}\)
  \item \(19.6 \, 	ext{m}\)
  \item \(29.4 \, 	ext{m}\)
  \item \(39.2 \, 	ext{m}\)
  \item \(14.7 \, 	ext{m}\)
\end{multicols}
\end{enumerate}
\vspace{0.15cm}
\item Un objeto cae libremente desde el reposo durante \(4 \, 	ext{s}\). ¿Qué distancia recorre?
\begin{enumerate}[label=\textbf{\Alph*)}]
\begin{multicols}{2}
  \item \(78.4 \, 	ext{m}\)
  \item \(9.8 \, 	ext{m}\)
  \item \(19.6 \, 	ext{m}\)
  \item \(29.4 \, 	ext{m}\)
  \item \(39.2 \, 	ext{m}\)
  \item \(14.7 \, 	ext{m}\)
\end{multicols}
\end{enumerate}
\vspace{0.15cm}
\item Un objeto cae libremente desde el reposo durante \(5 \, 	ext{s}\). ¿Qué distancia recorre?
\begin{enumerate}[label=\textbf{\Alph*)}]
\begin{multicols}{2}
  \item \(122.5 \, 	ext{m}\)
  \item \(9.8 \, 	ext{m}\)
  \item \(19.6 \, 	ext{m}\)
  \item \(29.4 \, 	ext{m}\)
  \item \(39.2 \, 	ext{m}\)
  \item \(14.7 \, 	ext{m}\)
\end{multicols}
\end{enumerate}
\vspace{0.15cm}
\item Un objeto se lanza verticalmente hacia arriba con una velocidad inicial de \(20 \, 	ext{m/s}\). ¿Cuál es su velocidad en el punto más alto?
\begin{enumerate}[label=\textbf{\Alph*)}]
\begin{multicols}{2}
  \item \(0 \, 	ext{m/s}\)
  \item \(20 \, 	ext{m/s}\)
  \item \(10 \, 	ext{m/s}\)
  \item \(-20 \, 	ext{m/s}\)
  \item \(9.8 \, 	ext{m/s}\)
  \item \(-9.8 \, 	ext{m/s}\)
\end{multicols}
\end{enumerate}
\vspace{0.15cm}
\item Un objeto se lanza verticalmente hacia arriba con una velocidad inicial de \(20 \, 	ext{m/s}\). ¿Cuánto tiempo tarda en alcanzar el punto más alto?
\begin{enumerate}[label=\textbf{\Alph*)}]
\begin{multicols}{2}
  \item \(2.04 \, 	ext{s}\)
  \item \(1.02 \, 	ext{s}\)
  \item \(4.08 \, 	ext{s}\)
  \item \(3.06 \, 	ext{s}\)
  \item \(5.10 \, 	ext{s}\)
  \item \(0.51 \, 	ext{s}\)
\end{multicols}
\end{enumerate}
\vspace{0.15cm}
\item Un objeto se lanza verticalmente hacia arriba con una velocidad inicial de \(30 \, 	ext{m/s}\). ¿Cuánto tiempo tarda en alcanzar el punto más alto?
\begin{enumerate}[label=\textbf{\Alph*)}]
\begin{multicols}{2}
  \item \(3.06 \, 	ext{s}\)
  \item \(1.53 \, 	ext{s}\)
  \item \(6.12 \, 	ext{s}\)
  \item \(4.59 \, 	ext{s}\)
  \item \(7.65 \, 	ext{s}\)
  \item \(0.77 \, 	ext{s}\)
\end{multicols}
\end{enumerate}
\vspace{0.15cm}
\item Un objeto se lanza verticalmente hacia arriba con una velocidad inicial de \(40 \, 	ext{m/s}\). ¿Cuánto tiempo tarda en alcanzar el punto más alto?
\begin{enumerate}[label=\textbf{\Alph*)}]
\begin{multicols}{2}
  \item \(4.08 \, 	ext{s}\)
  \item \(2.04 \, 	ext{s}\)
  \item \(8.16 \, 	ext{s}\)
  \item \(6.12 \, 	ext{s}\)
  \item \(10.20 \, 	ext{s}\)
  \item \(1.02 \, 	ext{s}\)
\end{multicols}
\end{enumerate}
\vspace{0.15cm}
\item Un objeto se lanza verticalmente hacia arriba con una velocidad inicial de \(20 \, 	ext{m/s}\). ¿Cuál es la altura máxima alcanzada?
\begin{enumerate}[label=\textbf{\Alph*)}]
\begin{multicols}{2}
  \item \(20.4 \, 	ext{m}\)
  \item \(10.2 \, 	ext{m}\)
  \item \(40.8 \, 	ext{m}\)
  \item \(30.6 \, 	ext{m}\)
  \item \(50.0 \, 	ext{m}\)
  \item \(15.0 \, 	ext{m}\)
\end{multicols}
\end{enumerate}
\vspace{0.15cm}
\item Un objeto se lanza verticalmente hacia arriba con una velocidad inicial de \(30 \, 	ext{m/s}\). ¿Cuál es la altura máxima alcanzada?
\begin{enumerate}[label=\textbf{\Alph*)}]
\begin{multicols}{2}
  \item \(45.9 \, 	ext{m}\)
  \item \(22.9 \, 	ext{m}\)
  \item \(91.8 \, 	ext{m}\)
  \item \(68.9 \, 	ext{m}\)
  \item \(100.0 \, 	ext{m}\)
  \item \(30.0 \, 	ext{m}\)
\end{multicols}
\end{enumerate}
\vspace{0.15cm}
\item Un objeto se lanza verticalmente hacia arriba con una velocidad inicial de \(40 \, 	ext{m/s}\). ¿Cuál es la altura máxima alcanzada?
\begin{enumerate}[label=\textbf{\Alph*)}]
\begin{multicols}{2}
  \item \(81.6 \, 	ext{m}\)
  \item \(40.8 \, 	ext{m}\)
  \item \(163.2 \, 	ext{m}\)
  \item \(122.4 \, 	ext{m}\)
  \item \(200.0 \, 	ext{m}\)
  \item \(60.0 \, 	ext{m}\) ---
\end{multicols}
\end{enumerate}
\vspace{0.15cm}
\end{enumerate}\n
\subsubsection*{SECCIÓN V: PROBLEMAS DE APLICACIÓN (241–300)}

\begin{enumerate}[label=\textbf{\arabic*.}]
\setcounter{enumi}{240}
\item Un auto viaja a \(72 \, 	ext{km/h}\). ¿Cuál es su velocidad en m/s?
\begin{enumerate}[label=\textbf{\Alph*)}]
\begin{multicols}{2}
  \item \(20 \, 	ext{m/s}\)
  \item \(72 \, 	ext{m/s}\)
  \item \(7.2 \, 	ext{m/s}\)
  \item \(200 \, 	ext{m/s}\)
  \item \(36 \, 	ext{m/s}\)
  \item \(18 \, 	ext{m/s}\)
\end{multicols}
\end{enumerate}
\vspace{0.15cm}
\item Un auto viaja a \(108 \, 	ext{km/h}\). ¿Cuál es su velocidad en m/s?
\begin{enumerate}[label=\textbf{\Alph*)}]
\begin{multicols}{2}
  \item \(30 \, 	ext{m/s}\)
  \item \(108 \, 	ext{m/s}\)
  \item \(10.8 \, 	ext{m/s}\)
  \item \(300 \, 	ext{m/s}\)
  \item \(54 \, 	ext{m/s}\)
  \item \(27 \, 	ext{m/s}\)
\end{multicols}
\end{enumerate}
\vspace{0.15cm}
\item Un auto viaja a \(36 \, 	ext{km/h}\). ¿Cuál es su velocidad en m/s?
\begin{enumerate}[label=\textbf{\Alph*)}]
\begin{multicols}{2}
  \item \(10 \, 	ext{m/s}\)
  \item \(36 \, 	ext{m/s}\)
  \item \(3.6 \, 	ext{m/s}\)
  \item \(100 \, 	ext{m/s}\)
  \item \(18 \, 	ext{m/s}\)
  \item \(9 \, 	ext{m/s}\)
\end{multicols}
\end{enumerate}
\vspace{0.15cm}
\item Un auto viaja a \(144 \, 	ext{km/h}\). ¿Cuál es su velocidad en m/s?
\begin{enumerate}[label=\textbf{\Alph*)}]
\begin{multicols}{2}
  \item \(40 \, 	ext{m/s}\)
  \item \(144 \, 	ext{m/s}\)
  \item \(14.4 \, 	ext{m/s}\)
  \item \(400 \, 	ext{m/s}\)
  \item \(72 \, 	ext{m/s}\)
  \item \(36 \, 	ext{m/s}\)
\end{multicols}
\end{enumerate}
\vspace{0.15cm}
\item Un auto viaja a \(54 \, 	ext{km/h}\). ¿Cuál es su velocidad en m/s?
\begin{enumerate}[label=\textbf{\Alph*)}]
\begin{multicols}{2}
  \item \(15 \, 	ext{m/s}\)
  \item \(54 \, 	ext{m/s}\)
  \item \(5.4 \, 	ext{m/s}\)
  \item \(150 \, 	ext{m/s}\)
  \item \(27 \, 	ext{m/s}\)
  \item \(13.5 \, 	ext{m/s}\)
\end{multicols}
\end{enumerate}
\vspace{0.15cm}
\item Un auto viaja a \(25 \, 	ext{m/s}\). ¿Cuál es su velocidad en km/h?
\begin{enumerate}[label=\textbf{\Alph*)}]
\begin{multicols}{2}
  \item \(90 \, 	ext{km/h}\)
  \item \(25 \, 	ext{km/h}\)
  \item \(250 \, 	ext{km/h}\)
  \item \(60 \, 	ext{km/h}\)
  \item \(45 \, 	ext{km/h}\)
  \item \(75 \, 	ext{km/h}\)
\end{multicols}
\end{enumerate}
\vspace{0.15cm}
\item Un auto viaja a \(40 \, 	ext{m/s}\). ¿Cuál es su velocidad en km/h?
\begin{enumerate}[label=\textbf{\Alph*)}]
\begin{multicols}{2}
  \item \(144 \, 	ext{km/h}\)
  \item \(40 \, 	ext{km/h}\)
  \item \(400 \, 	ext{km/h}\)
  \item \(100 \, 	ext{km/h}\)
  \item \(72 \, 	ext{km/h}\)
  \item \(120 \, 	ext{km/h}\)
\end{multicols}
\end{enumerate}
\vspace{0.15cm}
\item Un auto viaja a \(15 \, 	ext{m/s}\). ¿Cuál es su velocidad en km/h?
\begin{enumerate}[label=\textbf{\Alph*)}]
\begin{multicols}{2}
  \item \(54 \, 	ext{km/h}\)
  \item \(15 \, 	ext{km/h}\)
  \item \(150 \, 	ext{km/h}\)
  \item \(36 \, 	ext{km/h}\)
  \item \(27 \, 	ext{km/h}\)
  \item \(45 \, 	ext{km/h}\)
\end{multicols}
\end{enumerate}
\vspace{0.15cm}
\item Un auto viaja a \(30 \, 	ext{m/s}\). ¿Cuál es su velocidad en km/h?
\begin{enumerate}[label=\textbf{\Alph*)}]
\begin{multicols}{2}
  \item \(108 \, 	ext{km/h}\)
  \item \(30 \, 	ext{km/h}\)
  \item \(300 \, 	ext{km/h}\)
  \item \(90 \, 	ext{km/h}\)
  \item \(54 \, 	ext{km/h}\)
  \item \(72 \, 	ext{km/h}\)
\end{multicols}
\end{enumerate}
\vspace{0.15cm}
\item Un auto viaja a \(10 \, 	ext{m/s}\). ¿Cuál es su velocidad en km/h?
\begin{enumerate}[label=\textbf{\Alph*)}]
\begin{multicols}{2}
  \item \(36 \, 	ext{km/h}\)
  \item \(10 \, 	ext{km/h}\)
  \item \(100 \, 	ext{km/h}\)
  \item \(18 \, 	ext{km/h}\)
  \item \(20 \, 	ext{km/h}\)
  \item \(30 \, 	ext{km/h}\)
\end{multicols}
\end{enumerate}
\vspace{0.15cm}
\item Un auto recorre \(150 \, 	ext{km}\) en \(2 \, 	ext{h}\). ¿Cuál es su velocidad media en m/s?
\begin{enumerate}[label=\textbf{\Alph*)}]
\begin{multicols}{2}
  \item \(20.83 \, 	ext{m/s}\)
  \item \(75 \, 	ext{m/s}\)
  \item \(41.67 \, 	ext{m/s}\)
  \item \(10.42 \, 	ext{m/s}\)
  \item \(30 \, 	ext{m/s}\)
  \item \(50 \, 	ext{m/s}\)
\end{multicols}
\end{enumerate}
\vspace{0.15cm}
\item Un ciclista recorre \(60 \, 	ext{km}\) en \(3 \, 	ext{h}\). ¿Cuál es su velocidad media en m/s?
\begin{enumerate}[label=\textbf{\Alph*)}]
\begin{multicols}{2}
  \item \(5.56 \, 	ext{m/s}\)
  \item \(20 \, 	ext{m/s}\)
  \item \(16.67 \, 	ext{m/s}\)
  \item \(10 \, 	ext{m/s}\)
  \item \(2.78 \, 	ext{m/s}\)
  \item \(15 \, 	ext{m/s}\)
\end{multicols}
\end{enumerate}
\vspace{0.15cm}
\item Un tren recorre \(240 \, 	ext{km}\) en \(2 \, 	ext{h}\). ¿Cuál es su velocidad media en m/s?
\begin{enumerate}[label=\textbf{\Alph*)}]
\begin{multicols}{2}
  \item \(33.33 \, 	ext{m/s}\)
  \item \(120 \, 	ext{m/s}\)
  \item \(66.67 \, 	ext{m/s}\)
  \item \(16.67 \, 	ext{m/s}\)
  \item \(50 \, 	ext{m/s}\)
  \item \(25 \, 	ext{m/s}\)
\end{multicols}
\end{enumerate}
\vspace{0.15cm}
\item Un avión recorre \(900 \, 	ext{km}\) en \(1.5 \, 	ext{h}\). ¿Cuál es su velocidad media en m/s?
\begin{enumerate}[label=\textbf{\Alph*)}]
\begin{multicols}{2}
  \item \(166.67 \, 	ext{m/s}\)
  \item \(600 \, 	ext{m/s}\)
  \item \(333.33 \, 	ext{m/s}\)
  \item \(83.33 \, 	ext{m/s}\)
  \item \(200 \, 	ext{m/s}\)
  \item \(150 \, 	ext{m/s}\)
\end{multicols}
\end{enumerate}
\vspace{0.15cm}
\item Un barco recorre \(120 \, 	ext{km}\) en \(4 \, 	ext{h}\). ¿Cuál es su velocidad media en m/s?
\begin{enumerate}[label=\textbf{\Alph*)}]
\begin{multicols}{2}
  \item \(8.33 \, 	ext{m/s}\)
  \item \(30 \, 	ext{m/s}\)
  \item \(16.67 \, 	ext{m/s}\)
  \item \(4.17 \, 	ext{m/s}\)
  \item \(12 \, 	ext{m/s}\)
  \item \(20 \, 	ext{m/s}\)
\end{multicols}
\end{enumerate}
\vspace{0.15cm}
\item Un auto viaja a \(20 \, 	ext{m/s}\) durante \(10 \, 	ext{s}\). ¿Qué distancia recorre?
\begin{enumerate}[label=\textbf{\Alph*)}]
\begin{multicols}{2}
  \item \(200 \, 	ext{m}\)
  \item \(20 \, 	ext{m}\)
  \item \(2 \, 	ext{m}\)
  \item \(100 \, 	ext{m}\)
  \item \(300 \, 	ext{m}\)
  \item \(50 \, 	ext{m}\)
\end{multicols}
\end{enumerate}
\vspace{0.15cm}
\item Un auto viaja a \(30 \, 	ext{m/s}\) durante \(15 \, 	ext{s}\). ¿Qué distancia recorre?
\begin{enumerate}[label=\textbf{\Alph*)}]
\begin{multicols}{2}
  \item \(450 \, 	ext{m}\)
  \item \(30 \, 	ext{m}\)
  \item \(2 \, 	ext{m}\)
  \item \(150 \, 	ext{m}\)
  \item \(600 \, 	ext{m}\)
  \item \(75 \, 	ext{m}\)
\end{multicols}
\end{enumerate}
\vspace{0.15cm}
\item Un auto viaja a \(25 \, 	ext{m/s}\) durante \(8 \, 	ext{s}\). ¿Qué distancia recorre?
\begin{enumerate}[label=\textbf{\Alph*)}]
\begin{multicols}{2}
  \item \(200 \, 	ext{m}\)
  \item \(25 \, 	ext{m}\)
  \item \(3.125 \, 	ext{m}\)
  \item \(100 \, 	ext{m}\)
  \item \(300 \, 	ext{m}\)
  \item \(50 \, 	ext{m}\)
\end{multicols}
\end{enumerate}
\vspace{0.15cm}
\item Un auto viaja a \(40 \, 	ext{m/s}\) durante \(12 \, 	ext{s}\). ¿Qué distancia recorre?
\begin{enumerate}[label=\textbf{\Alph*)}]
\begin{multicols}{2}
  \item \(480 \, 	ext{m}\)
  \item \(40 \, 	ext{m}\)
  \item \(3.33 \, 	ext{m}\)
  \item \(240 \, 	ext{m}\)
  \item \(600 \, 	ext{m}\)
  \item \(120 \, 	ext{m}\)
\end{multicols}
\end{enumerate}
\vspace{0.15cm}
\item Un auto viaja a \(15 \, 	ext{m/s}\) durante \(20 \, 	ext{s}\). ¿Qué distancia recorre?
\begin{enumerate}[label=\textbf{\Alph*)}]
\begin{multicols}{2}
  \item \(300 \, 	ext{m}\)
  \item \(15 \, 	ext{m}\)
  \item \(0.75 \, 	ext{m}\)
  \item \(150 \, 	ext{m}\)
  \item \(450 \, 	ext{m}\)
  \item \(75 \, 	ext{m}\)
\end{multicols}
\end{enumerate}
\vspace{0.15cm}
\item Un auto parte del reposo y acelera a \(2 \, 	ext{m/s}^2\) durante \(10 \, 	ext{s}\). ¿Qué distancia recorre?
\begin{enumerate}[label=\textbf{\Alph*)}]
\begin{multicols}{2}
  \item \(100 \, 	ext{m}\)
  \item \(20 \, 	ext{m}\)
  \item \(200 \, 	ext{m}\)
  \item \(50 \, 	ext{m}\)
  \item \(150 \, 	ext{m}\)
  \item \(25 \, 	ext{m}\)
\end{multicols}
\end{enumerate}
\vspace{0.15cm}
\item Un auto parte del reposo y acelera a \(3 \, 	ext{m/s}^2\) durante \(8 \, 	ext{s}\). ¿Qué distancia recorre?
\begin{enumerate}[label=\textbf{\Alph*)}]
\begin{multicols}{2}
  \item \(96 \, 	ext{m}\)
  \item \(24 \, 	ext{m}\)
  \item \(192 \, 	ext{m}\)
  \item \(48 \, 	ext{m}\)
  \item \(144 \, 	ext{m}\)
  \item \(12 \, 	ext{m}\)
\end{multicols}
\end{enumerate}
\vspace{0.15cm}
\item Un auto parte del reposo y acelera a \(4 \, 	ext{m/s}^2\) durante \(5 \, 	ext{s}\). ¿Qué distancia recorre?
\begin{enumerate}[label=\textbf{\Alph*)}]
\begin{multicols}{2}
  \item \(50 \, 	ext{m}\)
  \item \(20 \, 	ext{m}\)
  \item \(100 \, 	ext{m}\)
  \item \(25 \, 	ext{m}\)
  \item \(75 \, 	ext{m}\)
  \item \(10 \, 	ext{m}\)
\end{multicols}
\end{enumerate}
\vspace{0.15cm}
\item Un auto parte del reposo y acelera a \(5 \, 	ext{m/s}^2\) durante \(6 \, 	ext{s}\). ¿Qué distancia recorre?
\begin{enumerate}[label=\textbf{\Alph*)}]
\begin{multicols}{2}
  \item \(90 \, 	ext{m}\)
  \item \(30 \, 	ext{m}\)
  \item \(180 \, 	ext{m}\)
  \item \(45 \, 	ext{m}\)
  \item \(135 \, 	ext{m}\)
  \item \(15 \, 	ext{m}\)
\end{multicols}
\end{enumerate}
\vspace{0.15cm}
\item Un auto parte del reposo y acelera a \(6 \, 	ext{m/s}^2\) durante \(4 \, 	ext{s}\). ¿Qué distancia recorre?
\begin{enumerate}[label=\textbf{\Alph*)}]
\begin{multicols}{2}
  \item \(48 \, 	ext{m}\)
  \item \(24 \, 	ext{m}\)
  \item \(96 \, 	ext{m}\)
  \item \(12 \, 	ext{m}\)
  \item \(72 \, 	ext{m}\)
  \item \(6 \, 	ext{m}\)
\end{multicols}
\end{enumerate}
\vspace{0.15cm}
\item Un auto viaja a \(20 \, 	ext{m/s}\) y frena a \(-4 \, 	ext{m/s}^2\). ¿Cuánto tiempo tarda en detenerse?
\begin{enumerate}[label=\textbf{\Alph*)}]
\begin{multicols}{2}
  \item \(5 \, 	ext{s}\)
  \item \(4 \, 	ext{s}\)
  \item \(10 \, 	ext{s}\)
  \item \(2 \, 	ext{s}\)
  \item \(8 \, 	ext{s}\)
  \item \(1 \, 	ext{s}\)
\end{multicols}
\end{enumerate}
\vspace{0.15cm}
\item Un auto viaja a \(30 \, 	ext{m/s}\) y frena a \(-5 \, 	ext{m/s}^2\). ¿Cuánto tiempo tarda en detenerse?
\begin{enumerate}[label=\textbf{\Alph*)}]
\begin{multicols}{2}
  \item \(6 \, 	ext{s}\)
  \item \(5 \, 	ext{s}\)
  \item \(10 \, 	ext{s}\)
  \item \(3 \, 	ext{s}\)
  \item \(8 \, 	ext{s}\)
  \item \(2 \, 	ext{s}\)
\end{multicols}
\end{enumerate}
\vspace{0.15cm}
\item Un auto viaja a \(40 \, 	ext{m/s}\) y frena a \(-8 \, 	ext{m/s}^2\). ¿Cuánto tiempo tarda en detenerse?
\begin{enumerate}[label=\textbf{\Alph*)}]
\begin{multicols}{2}
  \item \(5 \, 	ext{s}\)
  \item \(4 \, 	ext{s}\)
  \item \(10 \, 	ext{s}\)
  \item \(2 \, 	ext{s}\)
  \item \(8 \, 	ext{s}\)
  \item \(1 \, 	ext{s}\)
\end{multicols}
\end{enumerate}
\vspace{0.15cm}
\item Un auto viaja a \(25 \, 	ext{m/s}\) y frena a \(-5 \, 	ext{m/s}^2\). ¿Cuánto tiempo tarda en detenerse?
\begin{enumerate}[label=\textbf{\Alph*)}]
\begin{multicols}{2}
  \item \(5 \, 	ext{s}\)
  \item \(4 \, 	ext{s}\)
  \item \(10 \, 	ext{s}\)
  \item \(2 \, 	ext{s}\)
  \item \(8 \, 	ext{s}\)
  \item \(1 \, 	ext{s}\)
\end{multicols}
\end{enumerate}
\vspace{0.15cm}
\item Un auto viaja a \(15 \, 	ext{m/s}\) y frena a \(-3 \, 	ext{m/s}^2\). ¿Cuánto tiempo tarda en detenerse?
\begin{enumerate}[label=\textbf{\Alph*)}]
\begin{multicols}{2}
  \item \(5 \, 	ext{s}\)
  \item \(4 \, 	ext{s}\)
  \item \(10 \, 	ext{s}\)
  \item \(2 \, 	ext{s}\)
  \item \(8 \, 	ext{s}\)
  \item \(1 \, 	ext{s}\)
\end{multicols}
\end{enumerate}
\vspace{0.15cm}
\item Un auto viaja a \(20 \, 	ext{m/s}\) y frena a \(-4 \, 	ext{m/s}^2\). ¿Qué distancia recorre hasta detenerse?
\begin{enumerate}[label=\textbf{\Alph*)}]
\begin{multicols}{2}
  \item \(50 \, 	ext{m}\)
  \item \(20 \, 	ext{m}\)
  \item \(100 \, 	ext{m}\)
  \item \(25 \, 	ext{m}\)
  \item \(75 \, 	ext{m}\)
  \item \(10 \, 	ext{m}\)
\end{multicols}
\end{enumerate}
\vspace{0.15cm}
\item Un auto viaja a \(30 \, 	ext{m/s}\) y frena a \(-5 \, 	ext{m/s}^2\). ¿Qué distancia recorre hasta detenerse?
\begin{enumerate}[label=\textbf{\Alph*)}]
\begin{multicols}{2}
  \item \(90 \, 	ext{m}\)
  \item \(30 \, 	ext{m}\)
  \item \(180 \, 	ext{m}\)
  \item \(45 \, 	ext{m}\)
  \item \(135 \, 	ext{m}\)
  \item \(15 \, 	ext{m}\)
\end{multicols}
\end{enumerate}
\vspace{0.15cm}
\item Un auto viaja a \(40 \, 	ext{m/s}\) y frena a \(-8 \, 	ext{m/s}^2\). ¿Qué distancia recorre hasta detenerse?
\begin{enumerate}[label=\textbf{\Alph*)}]
\begin{multicols}{2}
  \item \(100 \, 	ext{m}\)
  \item \(40 \, 	ext{m}\)
  \item \(200 \, 	ext{m}\)
  \item \(50 \, 	ext{m}\)
  \item \(150 \, 	ext{m}\)
  \item \(20 \, 	ext{m}\)
\end{multicols}
\end{enumerate}
\vspace{0.15cm}
\item Un auto viaja a \(25 \, 	ext{m/s}\) y frena a \(-5 \, 	ext{m/s}^2\). ¿Qué distancia recorre hasta detenerse?
\begin{enumerate}[label=\textbf{\Alph*)}]
\begin{multicols}{2}
  \item \(62.5 \, 	ext{m}\)
  \item \(25 \, 	ext{m}\)
  \item \(125 \, 	ext{m}\)
  \item \(31.25 \, 	ext{m}\)
  \item \(93.75 \, 	ext{m}\)
  \item \(12.5 \, 	ext{m}\)
\end{multicols}
\end{enumerate}
\vspace{0.15cm}
\item Un auto viaja a \(15 \, 	ext{m/s}\) y frena a \(-3 \, 	ext{m/s}^2\). ¿Qué distancia recorre hasta detenerse?
\begin{enumerate}[label=\textbf{\Alph*)}]
\begin{multicols}{2}
  \item \(37.5 \, 	ext{m}\)
  \item \(15 \, 	ext{m}\)
  \item \(75 \, 	ext{m}\)
  \item \(18.75 \, 	ext{m}\)
  \item \(56.25 \, 	ext{m}\)
  \item \(7.5 \, 	ext{m}\)
\end{multicols}
\end{enumerate}
\vspace{0.15cm}
\item Un objeto se deja caer desde una altura de \(45 \, 	ext{m}\). ¿Cuánto tiempo tarda en llegar al suelo?
\begin{enumerate}[label=\textbf{\Alph*)}]
\begin{multicols}{2}
  \item \(3.03 \, 	ext{s}\)
  \item \(1.52 \, 	ext{s}\)
  \item \(6.06 \, 	ext{s}\)
  \item \(4.55 \, 	ext{s}\)
  \item \(7.58 \, 	ext{s}\)
  \item \(0.76 \, 	ext{s}\)
\end{multicols}
\end{enumerate}
\vspace{0.15cm}
\item Un objeto se deja caer desde una altura de \(80 \, 	ext{m}\). ¿Cuánto tiempo tarda en llegar al suelo?
\begin{enumerate}[label=\textbf{\Alph*)}]
\begin{multicols}{2}
  \item \(4.04 \, 	ext{s}\)
  \item \(2.02 \, 	ext{s}\)
  \item \(8.08 \, 	ext{s}\)
  \item \(6.06 \, 	ext{s}\)
  \item \(10.10 \, 	ext{s}\)
  \item \(1.01 \, 	ext{s}\)
\end{multicols}
\end{enumerate}
\vspace{0.15cm}
\item Un objeto se deja caer desde una altura de \(20 \, 	ext{m}\). ¿Cuánto tiempo tarda en llegar al suelo?
\begin{enumerate}[label=\textbf{\Alph*)}]
\begin{multicols}{2}
  \item \(2.02 \, 	ext{s}\)
  \item \(1.01 \, 	ext{s}\)
  \item \(4.04 \, 	ext{s}\)
  \item \(3.03 \, 	ext{s}\)
  \item \(5.05 \, 	ext{s}\)
  \item \(0.51 \, 	ext{s}\)
\end{multicols}
\end{enumerate}
\vspace{0.15cm}
\item Un objeto se deja caer desde una altura de \(125 \, 	ext{m}\). ¿Cuánto tiempo tarda en llegar al suelo?
\begin{enumerate}[label=\textbf{\Alph*)}]
\begin{multicols}{2}
  \item \(5.05 \, 	ext{s}\)
  \item \(2.53 \, 	ext{s}\)
  \item \(10.10 \, 	ext{s}\)
  \item \(7.58 \, 	ext{s}\)
  \item \(12.63 \, 	ext{s}\)
  \item \(1.26 \, 	ext{s}\)
\end{multicols}
\end{enumerate}
\vspace{0.15cm}
\item Un objeto se deja caer desde una altura de \(180 \, 	ext{m}\). ¿Cuánto tiempo tarda en llegar al suelo?
\begin{enumerate}[label=\textbf{\Alph*)}]
\begin{multicols}{2}
  \item \(6.06 \, 	ext{s}\)
  \item \(3.03 \, 	ext{s}\)
  \item \(12.12 \, 	ext{s}\)
  \item \(9.09 \, 	ext{s}\)
  \item \(15.15 \, 	ext{s}\)
  \item \(1.52 \, 	ext{s}\)
\end{multicols}
\end{enumerate}
\vspace{0.15cm}
\item Un objeto se deja caer desde una altura de \(45 \, 	ext{m}\). ¿Con qué velocidad llega al suelo?
\begin{enumerate}[label=\textbf{\Alph*)}]
\begin{multicols}{2}
  \item \(29.7 \, 	ext{m/s}\)
  \item \(14.85 \, 	ext{m/s}\)
  \item \(59.4 \, 	ext{m/s}\)
  \item \(44.55 \, 	ext{m/s}\)
  \item \(74.25 \, 	ext{m/s}\)
  \item \(7.43 \, 	ext{m/s}\)
\end{multicols}
\end{enumerate}
\vspace{0.15cm}
\item Un objeto se deja caer desde una altura de \(80 \, 	ext{m}\). ¿Con qué velocidad llega al suelo?
\begin{enumerate}[label=\textbf{\Alph*)}]
\begin{multicols}{2}
  \item \(39.6 \, 	ext{m/s}\)
  \item \(19.8 \, 	ext{m/s}\)
  \item \(79.2 \, 	ext{m/s}\)
  \item \(59.4 \, 	ext{m/s}\)
  \item \(99.0 \, 	ext{m/s}\)
  \item \(9.9 \, 	ext{m/s}\)
\end{multicols}
\end{enumerate}
\vspace{0.15cm}
\item Un objeto se deja caer desde una altura de \(20 \, 	ext{m}\). ¿Con qué velocidad llega al suelo?
\begin{enumerate}[label=\textbf{\Alph*)}]
\begin{multicols}{2}
  \item \(19.8 \, 	ext{m/s}\)
  \item \(9.9 \, 	ext{m/s}\)
  \item \(39.6 \, 	ext{m/s}\)
  \item \(29.7 \, 	ext{m/s}\)
  \item \(49.5 \, 	ext{m/s}\)
  \item \(4.95 \, 	ext{m/s}\)
\end{multicols}
\end{enumerate}
\vspace{0.15cm}
\item Un objeto se deja caer desde una altura de \(125 \, 	ext{m}\). ¿Con qué velocidad llega al suelo?
\begin{enumerate}[label=\textbf{\Alph*)}]
\begin{multicols}{2}
  \item \(49.5 \, 	ext{m/s}\)
  \item \(24.75 \, 	ext{m/s}\)
  \item \(99.0 \, 	ext{m/s}\)
  \item \(74.25 \, 	ext{m/s}\)
  \item \(123.75 \, 	ext{m/s}\)
  \item \(12.38 \, 	ext{m/s}\)
\end{multicols}
\end{enumerate}
\vspace{0.15cm}
\item Un objeto se deja caer desde una altura de \(180 \, 	ext{m}\). ¿Con qué velocidad llega al suelo?
\begin{enumerate}[label=\textbf{\Alph*)}]
\begin{multicols}{2}
  \item \(59.4 \, 	ext{m/s}\)
  \item \(29.7 \, 	ext{m/s}\)
  \item \(118.8 \, 	ext{m/s}\)
  \item \(89.1 \, 	ext{m/s}\)
  \item \(148.5 \, 	ext{m/s}\)
  \item \(14.85 \, 	ext{m/s}\)
\end{multicols}
\end{enumerate}
\vspace{0.15cm}
\item Un proyectil se lanza con \(v_0 = 50 \, 	ext{m/s}\) y \(	heta = 30^\circ\). ¿Cuál es su alcance máximo?
\begin{enumerate}[label=\textbf{\Alph*)}]
\begin{multicols}{2}
  \item \(220.92 \, 	ext{m}\)
  \item \(110.46 \, 	ext{m}\)
  \item \(441.84 \, 	ext{m}\)
  \item \(331.38 \, 	ext{m}\)
  \item \(552.30 \, 	ext{m}\)
  \item \(165.69 \, 	ext{m}\)
\end{multicols}
\end{enumerate}
\vspace{0.15cm}
\item Un proyectil se lanza con \(v_0 = 60 \, 	ext{m/s}\) y \(	heta = 45^\circ\). ¿Cuál es su alcance máximo?
\begin{enumerate}[label=\textbf{\Alph*)}]
\begin{multicols}{2}
  \item \(367.35 \, 	ext{m}\)
  \item \(183.67 \, 	ext{m}\)
  \item \(734.70 \, 	ext{m}\)
  \item \(551.03 \, 	ext{m}\)
  \item \(918.38 \, 	ext{m}\)
  \item \(275.51 \, 	ext{m}\)
\end{multicols}
\end{enumerate}
\vspace{0.15cm}
\item Un proyectil se lanza con \(v_0 = 70 \, 	ext{m/s}\) y \(	heta = 60^\circ\). ¿Cuál es su alcance máximo?
\begin{enumerate}[label=\textbf{\Alph*)}]
\begin{multicols}{2}
  \item \(433.01 \, 	ext{m}\)
  \item \(216.51 \, 	ext{m}\)
  \item \(866.02 \, 	ext{m}\)
  \item \(649.52 \, 	ext{m}\)
  \item \(1082.53 \, 	ext{m}\)
  \item \(324.76 \, 	ext{m}\)
\end{multicols}
\end{enumerate}
\vspace{0.15cm}
\item Un proyectil se lanza con \(v_0 = 80 \, 	ext{m/s}\) y \(	heta = 30^\circ\). ¿Cuál es su altura máxima?
\begin{enumerate}[label=\textbf{\Alph*)}]
\begin{multicols}{2}
  \item \(81.63 \, 	ext{m}\)
  \item \(40.82 \, 	ext{m}\)
  \item \(163.26 \, 	ext{m}\)
  \item \(122.45 \, 	ext{m}\)
  \item \(204.08 \, 	ext{m}\)
  \item \(61.22 \, 	ext{m}\)
\end{multicols}
\end{enumerate}
\vspace{0.15cm}
\item Un proyectil se lanza con \(v_0 = 90 \, 	ext{m/s}\) y \(	heta = 45^\circ\). ¿Cuál es su altura máxima?
\begin{enumerate}[label=\textbf{\Alph*)}]
\begin{multicols}{2}
  \item \(206.63 \, 	ext{m}\)
  \item \(103.32 \, 	ext{m}\)
  \item \(413.26 \, 	ext{m}\)
  \item \(309.95 \, 	ext{m}\)
  \item \(516.58 \, 	ext{m}\)
  \item \(154.97 \, 	ext{m}\)
\end{multicols}
\end{enumerate}
\vspace{0.15cm}
\item Un proyectil se lanza con \(v_0 = 100 \, 	ext{m/s}\) y \(	heta = 60^\circ\). ¿Cuál es su altura máxima?
\begin{enumerate}[label=\textbf{\Alph*)}]
\begin{multicols}{2}
  \item \(382.65 \, 	ext{m}\)
  \item \(191.33 \, 	ext{m}\)
  \item \(765.30 \, 	ext{m}\)
  \item \(573.98 \, 	ext{m}\)
  \item \(956.63 \, 	ext{m}\)
  \item \(286.99 \, 	ext{m}\)
\end{multicols}
\end{enumerate}
\vspace{0.15cm}
\item Un proyectil se lanza con \(v_0 = 50 \, 	ext{m/s}\) y \(	heta = 30^\circ\). ¿Cuál es su tiempo de vuelo?
\begin{enumerate}[label=\textbf{\Alph*)}]
\begin{multicols}{2}
  \item \(5.10 \, 	ext{s}\)
  \item \(2.55 \, 	ext{s}\)
  \item \(10.20 \, 	ext{s}\)
  \item \(7.65 \, 	ext{s}\)
  \item \(12.75 \, 	ext{s}\)
  \item \(1.28 \, 	ext{s}\)
\end{multicols}
\end{enumerate}
\vspace{0.15cm}
\item Un proyectil se lanza con \(v_0 = 60 \, 	ext{m/s}\) y \(	heta = 45^\circ\). ¿Cuál es su tiempo de vuelo?
\begin{enumerate}[label=\textbf{\Alph*)}]
\begin{multicols}{2}
  \item \(8.66 \, 	ext{s}\)
  \item \(4.33 \, 	ext{s}\)
  \item \(17.32 \, 	ext{s}\)
  \item \(12.99 \, 	ext{s}\)
  \item \(21.65 \, 	ext{s}\)
  \item \(2.17 \, 	ext{s}\)
\end{multicols}
\end{enumerate}
\vspace{0.15cm}
\item Un proyectil se lanza con \(v_0 = 70 \, 	ext{m/s}\) y \(	heta = 60^\circ\). ¿Cuál es su tiempo de vuelo?
\begin{enumerate}[label=\textbf{\Alph*)}]
\begin{multicols}{2}
  \item \(12.37 \, 	ext{s}\)
  \item \(6.19 \, 	ext{s}\)
  \item \(24.74 \, 	ext{s}\)
  \item \(18.56 \, 	ext{s}\)
  \item \(30.93 \, 	ext{s}\)
  \item \(3.09 \, 	ext{s}\)
\end{multicols}
\end{enumerate}
\vspace{0.15cm}
\item Un auto viaja a \(72 \, 	ext{km/h}\) y frena a \(-4 \, 	ext{m/s}^2\). ¿Cuánto tiempo tarda en detenerse?
\begin{enumerate}[label=\textbf{\Alph*)}]
\begin{multicols}{2}
  \item \(5 \, 	ext{s}\)
  \item \(4 \, 	ext{s}\)
  \item \(10 \, 	ext{s}\)
  \item \(2 \, 	ext{s}\)
  \item \(8 \, 	ext{s}\)
  \item \(1 \, 	ext{s}\)
\end{multicols}
\end{enumerate}
\vspace{0.15cm}
\item Un auto viaja a \(108 \, 	ext{km/h}\) y frena a \(-6 \, 	ext{m/s}^2\). ¿Cuánto tiempo tarda en detenerse?
\begin{enumerate}[label=\textbf{\Alph*)}]
\begin{multicols}{2}
  \item \(5 \, 	ext{s}\)
  \item \(4 \, 	ext{s}\)
  \item \(10 \, 	ext{s}\)
  \item \(2 \, 	ext{s}\)
  \item \(8 \, 	ext{s}\)
  \item \(1 \, 	ext{s}\)
\end{multicols}
\end{enumerate}
\vspace{0.15cm}
\item Un auto viaja a \(144 \, 	ext{km/h}\) y frena a \(-8 \, 	ext{m/s}^2\). ¿Cuánto tiempo tarda en detenerse?
\begin{enumerate}[label=\textbf{\Alph*)}]
\begin{multicols}{2}
  \item \(5 \, 	ext{s}\)
  \item \(4 \, 	ext{s}\)
  \item \(10 \, 	ext{s}\)
  \item \(2 \, 	ext{s}\)
  \item \(8 \, 	ext{s}\)
  \item \(1 \, 	ext{s}\)
\end{multicols}
\end{enumerate}
\vspace{0.15cm}
\item Un auto viaja a \(72 \, 	ext{km/h}\) y frena a \(-4 \, 	ext{m/s}^2\). ¿Qué distancia recorre hasta detenerse?
\begin{enumerate}[label=\textbf{\Alph*)}]
\begin{multicols}{2}
  \item \(50 \, 	ext{m}\)
  \item \(20 \, 	ext{m}\)
  \item \(100 \, 	ext{m}\)
  \item \(25 \, 	ext{m}\)
  \item \(75 \, 	ext{m}\)
  \item \(10 \, 	ext{m}\)
\end{multicols}
\end{enumerate}
\vspace{0.15cm}
\item Un auto viaja a \(108 \, 	ext{km/h}\) y frena a \(-6 \, 	ext{m/s}^2\). ¿Qué distancia recorre hasta detenerse?
\begin{enumerate}[label=\textbf{\Alph*)}]
\begin{multicols}{2}
  \item \(75 \, 	ext{m}\)
  \item \(30 \, 	ext{m}\)
  \item \(150 \, 	ext{m}\)
  \item \(37.5 \, 	ext{m}\)
  \item \(112.5 \, 	ext{m}\)
  \item \(15 \, 	ext{m}\)
\end{multicols}
\end{enumerate}
\vspace{0.15cm}
\item Un auto viaja a \(144 \, 	ext{km/h}\) y frena a \(-8 \, 	ext{m/s}^2\). ¿Qué distancia recorre hasta detenerse?
\begin{enumerate}[label=\textbf{\Alph*)}]
\begin{multicols}{2}
  \item \(100 \, 	ext{m}\)
  \item \(40 \, 	ext{m}\)
  \item \(200 \, 	ext{m}\)
  \item \(50 \, 	ext{m}\)
  \item \(150 \, 	ext{m}\)
  \item \(20 \, 	ext{m}\) ---
\end{multicols}
\end{enumerate}
\vspace{0.15cm}
\end{enumerate}\n
\vspace{0.4cm}
\begin{tcolorbox}[colback=gray!5!white,colframe=gray!60!black,title=\textbf{SOLUCIONARIO -- 5.3 Física General y Mecánica}]
\begin{multicols}{4}
\begin{enumerate}[label=\textbf{\arabic*.}]
  \item \textbf{A}
  \item \textbf{A}
  \item \textbf{A}
  \item \textbf{A}
  \item \textbf{A}
  \item \textbf{A}
  \item \textbf{A}
  \item \textbf{A}
  \item \textbf{A}
  \item \textbf{A}
  \item \textbf{A}
  \item \textbf{A}
  \item \textbf{A}
  \item \textbf{A}
  \item \textbf{A}
  \item \textbf{A}
  \item \textbf{A}
  \item \textbf{A}
  \item \textbf{A}
  \item \textbf{A}
  \item \textbf{A}
  \item \textbf{A}
  \item \textbf{A}
  \item \textbf{A}
  \item \textbf{A}
  \item \textbf{A}
  \item \textbf{A}
  \item \textbf{A}
  \item \textbf{A}
  \item \textbf{A}
  \item \textbf{A}
  \item \textbf{A}
  \item \textbf{A}
  \item \textbf{A}
  \item \textbf{A}
  \item \textbf{A}
  \item \textbf{A}
  \item \textbf{A}
  \item \textbf{A}
  \item \textbf{A}
  \item \textbf{A}
  \item \textbf{A}
  \item \textbf{A}
  \item \textbf{A}
  \item \textbf{A}
  \item \textbf{A}
  \item \textbf{A}
  \item \textbf{A}
  \item \textbf{A}
  \item \textbf{A}
  \item \textbf{A}
  \item \textbf{A}
  \item \textbf{A}
  \item \textbf{A}
  \item \textbf{A}
  \item \textbf{A}
  \item \textbf{A}
  \item \textbf{A}
  \item \textbf{A}
  \item \textbf{A}
  \item \textbf{A}
  \item \textbf{A}
  \item \textbf{A}
  \item \textbf{A}
  \item \textbf{A}
  \item \textbf{A}
  \item \textbf{A}
  \item \textbf{A}
  \item \textbf{A}
  \item \textbf{A}
  \item \textbf{A}
  \item \textbf{A}
  \item \textbf{A}
  \item \textbf{A}
  \item \textbf{A}
  \item \textbf{A}
  \item \textbf{A}
  \item \textbf{A}
  \item \textbf{A}
  \item \textbf{A}
  \item \textbf{A}
  \item \textbf{A}
  \item \textbf{A}
  \item \textbf{A}
  \item \textbf{A}
  \item \textbf{A}
  \item \textbf{A}
  \item \textbf{A}
  \item \textbf{A}
  \item \textbf{A}
  \item \textbf{A}
  \item \textbf{A}
  \item \textbf{A}
  \item \textbf{A}
  \item \textbf{A}
  \item \textbf{A}
  \item \textbf{A}
  \item \textbf{A}
  \item \textbf{A}
  \item \textbf{A}
  \item \textbf{A}
  \item \textbf{A}
  \item \textbf{A}
  \item \textbf{A}
  \item \textbf{A}
  \item \textbf{A}
  \item \textbf{A}
  \item \textbf{A}
  \item \textbf{A}
  \item \textbf{A}
  \item \textbf{A}
  \item \textbf{A}
  \item \textbf{A}
  \item \textbf{A}
  \item \textbf{A}
  \item \textbf{A}
  \item \textbf{A}
  \item \textbf{A}
  \item \textbf{A}
  \item \textbf{A}
  \item \textbf{A}
  \item \textbf{A}
  \item \textbf{A}
  \item \textbf{A}
  \item \textbf{A}
  \item \textbf{A}
  \item \textbf{A}
  \item \textbf{A}
  \item \textbf{A}
  \item \textbf{A}
  \item \textbf{A}
  \item \textbf{A}
  \item \textbf{A}
  \item \textbf{A}
  \item \textbf{A}
  \item \textbf{A}
  \item \textbf{A}
  \item \textbf{A}
  \item \textbf{A}
  \item \textbf{A}
  \item \textbf{A}
  \item \textbf{A}
  \item \textbf{A}
  \item \textbf{A}
  \item \textbf{A}
  \item \textbf{A}
  \item \textbf{A}
  \item \textbf{A}
  \item \textbf{A}
  \item \textbf{A}
  \item \textbf{A}
  \item \textbf{A}
  \item \textbf{A}
  \item \textbf{A}
  \item \textbf{A}
  \item \textbf{A}
  \item \textbf{A}
  \item \textbf{A}
  \item \textbf{A}
  \item \textbf{A}
  \item \textbf{A}
  \item \textbf{A}
  \item \textbf{A}
  \item \textbf{A}
  \item \textbf{A}
  \item \textbf{A}
  \item \textbf{A}
  \item \textbf{A}
  \item \textbf{A}
  \item \textbf{A}
  \item \textbf{A}
  \item \textbf{A}
  \item \textbf{A}
  \item \textbf{A}
  \item \textbf{A}
  \item \textbf{A}
  \item \textbf{A}
  \item \textbf{A}
  \item \textbf{A}
  \item \textbf{A}
  \item \textbf{A}
  \item \textbf{A}
  \item \textbf{A}
  \item \textbf{A}
  \item \textbf{A}
  \item \textbf{A}
  \item \textbf{A}
  \item \textbf{A}
  \item \textbf{A}
  \item \textbf{A}
  \item \textbf{A}
  \item \textbf{A}
  \item \textbf{A}
  \item \textbf{A}
  \item \textbf{A}
  \item \textbf{A}
  \item \textbf{A}
  \item \textbf{A}
  \item \textbf{A}
  \item \textbf{A}
  \item \textbf{A}
  \item \textbf{A}
  \item \textbf{A}
  \item \textbf{A}
  \item \textbf{A}
  \item \textbf{A}
  \item \textbf{A}
  \item \textbf{A}
  \item \textbf{A}
  \item \textbf{A}
  \item \textbf{A}
  \item \textbf{A}
  \item \textbf{A}
  \item \textbf{A}
  \item \textbf{A}
  \item \textbf{A}
  \item \textbf{A}
  \item \textbf{A}
  \item \textbf{A}
  \item \textbf{A}
  \item \textbf{A}
  \item \textbf{A}
  \item \textbf{A}
  \item \textbf{A}
  \item \textbf{A}
  \item \textbf{A}
  \item \textbf{A}
  \item \textbf{A}
  \item \textbf{A}
  \item \textbf{A}
  \item \textbf{A}
  \item \textbf{A}
  \item \textbf{A}
  \item \textbf{A}
  \item \textbf{A}
  \item \textbf{A}
  \item \textbf{A}
  \item \textbf{A}
  \item \textbf{A}
  \item \textbf{A}
  \item \textbf{A}
  \item \textbf{A}
  \item \textbf{A}
  \item \textbf{A}
  \item \textbf{A}
  \item \textbf{A}
  \item \textbf{A}
  \item \textbf{A}
  \item \textbf{A}
  \item \textbf{A}
  \item \textbf{A}
  \item \textbf{A}
  \item \textbf{A}
  \item \textbf{A}
  \item \textbf{A}
  \item \textbf{A}
  \item \textbf{A}
  \item \textbf{A}
  \item \textbf{A}
  \item \textbf{A}
  \item \textbf{A}
  \item \textbf{A}
  \item \textbf{A}
  \item \textbf{A}
  \item \textbf{A}
  \item \textbf{A}
  \item \textbf{A}
  \item \textbf{A}
  \item \textbf{A}
  \item \textbf{A}
  \item \textbf{A}
  \item \textbf{A}
  \item \textbf{A}
  \item \textbf{A}
  \item \textbf{A}
  \item \textbf{A}
  \item \textbf{A}
  \item \textbf{A}
  \item \textbf{A}
  \item \textbf{A}
  \item \textbf{A}
  \item \textbf{A}
  \item \textbf{A}
  \item \textbf{A}
  \item \textbf{A}
  \item \textbf{A}
  \item \textbf{A}
  \item \textbf{A}
  \item \textbf{A}
  \item \textbf{A}
  \item \textbf{A}
  \item \textbf{A}
  \item \textbf{A}
  \item \textbf{A}
  \item \textbf{A}
  \item \textbf{A}
  \item \textbf{A}
  \item \textbf{A}
  \item \textbf{A}
  \item \textbf{A}
\end{enumerate}
\end{multicols}
\end{tcolorbox}
\newpage\n
\subsection{5.4 Taxonomía y Ecosistemas}

\subsubsection*{SECCIÓN I: FUNDAMENTOS Y CATEGORÍAS TAXONÓMICAS (1–50)}

\begin{enumerate}[label=\textbf{\arabic*.}]
\setcounter{enumi}{0}
\item ¿Qué es la taxonomía?
\begin{enumerate}[label=\textbf{\Alph*)}]
\begin{multicols}{2}
  \item La ciencia que clasifica y nombra a los seres vivos
  \item La ciencia que estudia los minerales
  \item La ciencia que estudia los astros
  \item La ciencia que estudia las reacciones químicas
  \item La ciencia que estudia los números
  \item La ciencia que estudia las fuerzas
\end{multicols}
\end{enumerate}
\vspace{0.15cm}
\item ¿Quién es considerado el padre de la taxonomía moderna?
\begin{enumerate}[label=\textbf{\Alph*)}]
\begin{multicols}{2}
  \item Carlos Linneo
  \item Charles Darwin
  \item Gregor Mendel
  \item Louis Pasteur
  \item Robert Hooke
  \item Alexander Fleming
\end{multicols}
\end{enumerate}
\vspace{0.15cm}
\item ¿Qué sistema de nomenclatura propuso Linneo?
\begin{enumerate}[label=\textbf{\Alph*)}]
\begin{multicols}{2}
  \item El sistema binomial
  \item El sistema trinomial
  \item El sistema uninominal
  \item El sistema polinomial
  \item El sistema decimal
  \item El sistema binario
\end{multicols}
\end{enumerate}
\vspace{0.15cm}
\item ¿Cuál es la categoría taxonómica básica?
\begin{enumerate}[label=\textbf{\Alph*)}]
\begin{multicols}{2}
  \item La especie
  \item El género
  \item La familia
  \item El orden
  \item La clase
  \item El filo
\end{multicols}
\end{enumerate}
\vspace{0.15cm}
\item ¿Qué es una especie?
\begin{enumerate}[label=\textbf{\Alph*)}]
\begin{multicols}{2}
  \item Un conjunto de poblaciones que se reproducen entre sí y originan descendencia fértil
  \item Un conjunto de géneros
  \item Un conjunto de familias
  \item Un conjunto de órdenes
  \item Un conjunto de clases
  \item Un conjunto de filos
\end{multicols}
\end{enumerate}
\vspace{0.15cm}
\item ¿Qué categoría agrupa especies similares?
\begin{enumerate}[label=\textbf{\Alph*)}]
\begin{multicols}{2}
  \item El género
  \item La familia
  \item El orden
  \item La clase
  \item El filo
  \item El reino
\end{multicols}
\end{enumerate}
\vspace{0.15cm}
\item ¿Qué categoría agrupa géneros similares?
\begin{enumerate}[label=\textbf{\Alph*)}]
\begin{multicols}{2}
  \item La familia
  \item El orden
  \item La clase
  \item El filo
  \item El reino
  \item El dominio
\end{multicols}
\end{enumerate}
\vspace{0.15cm}
\item ¿Qué categoría agrupa familias similares?
\begin{enumerate}[label=\textbf{\Alph*)}]
\begin{multicols}{2}
  \item El orden
  \item La clase
  \item El filo
  \item El reino
  \item El dominio
  \item La especie
\end{multicols}
\end{enumerate}
\vspace{0.15cm}
\item ¿Qué categoría agrupa órdenes similares?
\begin{enumerate}[label=\textbf{\Alph*)}]
\begin{multicols}{2}
  \item La clase
  \item El filo
  \item El reino
  \item El dominio
  \item La especie
  \item El género
\end{multicols}
\end{enumerate}
\vspace{0.15cm}
\item ¿Qué categoría agrupa clases similares?
\begin{enumerate}[label=\textbf{\Alph*)}]
\begin{multicols}{2}
  \item El filo
  \item El reino
  \item El dominio
  \item La especie
  \item El género
  \item La familia
\end{multicols}
\end{enumerate}
\vspace{0.15cm}
\item ¿Qué categoría agrupa filos similares?
\begin{enumerate}[label=\textbf{\Alph*)}]
\begin{multicols}{2}
  \item El reino
  \item El dominio
  \item La especie
  \item El género
  \item La familia
  \item El orden
\end{multicols}
\end{enumerate}
\vspace{0.15cm}
\item ¿Qué categoría agrupa reinos similares?
\begin{enumerate}[label=\textbf{\Alph*)}]
\begin{multicols}{2}
  \item El dominio
  \item La especie
  \item El género
  \item La familia
  \item El orden
  \item La clase
\end{multicols}
\end{enumerate}
\vspace{0.15cm}
\item ¿Qué es un taxón?
\begin{enumerate}[label=\textbf{\Alph*)}]
\begin{multicols}{2}
  \item Cualquier grupo de organismos en cualquier nivel taxonómico
  \item Solo la especie
  \item Solo el género
  \item Solo la familia
  \item Solo el orden
  \item Solo la clase
\end{multicols}
\end{enumerate}
\vspace{0.15cm}
\item ¿Qué es la clasificación artificial?
\begin{enumerate}[label=\textbf{\Alph*)}]
\begin{multicols}{2}
  \item La basada en pocos caracteres arbitrarios
  \item La basada en la evolución
  \item La basada en la genética
  \item La basada en la filogenia
  \item La basada en la cladística
  \item La basada en la molecular
\end{multicols}
\end{enumerate}
\vspace{0.15cm}
\item ¿Qué es la clasificación natural?
\begin{enumerate}[label=\textbf{\Alph*)}]
\begin{multicols}{2}
  \item La basada en las relaciones evolutivas
  \item La basada en caracteres arbitrarios
  \item La basada en el color
  \item La basada en el tamaño
  \item La basada en la forma
  \item La basada en el hábitat
\end{multicols}
\end{enumerate}
\vspace{0.15cm}
\item ¿Qué es la filogenia?
\begin{enumerate}[label=\textbf{\Alph*)}]
\begin{multicols}{2}
  \item La historia evolutiva de un grupo de organismos
  \item La clasificación arbitraria
  \item La nomenclatura binomial
  \item La descripción de especies
  \item La colección de especímenes
  \item La clave dicotómica
\end{multicols}
\end{enumerate}
\vspace{0.15cm}
\item ¿Qué es la cladística?
\begin{enumerate}[label=\textbf{\Alph*)}]
\begin{multicols}{2}
  \item El método que clasifica según caracteres derivados compartidos
  \item El método que clasifica según el color
  \item El método que clasifica según el tamaño
  \item El método que clasifica según la forma
  \item El método que clasifica según el hábitat
  \item El método que clasifica según la utilidad
\end{multicols}
\end{enumerate}
\vspace{0.15cm}
\item ¿Qué es una homología?
\begin{enumerate}[label=\textbf{\Alph*)}]
\begin{multicols}{2}
  \item Semejanza por ancestro común
  \item Semejanza por función
  \item Semejanza por ambiente
  \item Semejanza por convergencia
  \item Semejanza por casualidad
  \item Semejanza por imitación
\end{multicols}
\end{enumerate}
\vspace{0.15cm}
\item ¿Qué es una analogía?
\begin{enumerate}[label=\textbf{\Alph*)}]
\begin{multicols}{2}
  \item Semejanza por función, no por ancestro común
  \item Semejanza por ancestro común
  \item Semejanza por evolución
  \item Semejanza por genética
  \item Semejanza por filogenia
  \item Semejanza por cladística
\end{multicols}
\end{enumerate}
\vspace{0.15cm}
\item ¿Qué es un carácter ancestral?
\begin{enumerate}[label=\textbf{\Alph*)}]
\begin{multicols}{2}
  \item El que está presente en el ancestro común
  \item El que está presente solo en un grupo
  \item El que está ausente en el ancestro
  \item El que es nuevo
  \item El que es derivado
  \item El que es exclusivo
\end{multicols}
\end{enumerate}
\vspace{0.15cm}
\item ¿Qué es un carácter derivado?
\begin{enumerate}[label=\textbf{\Alph*)}]
\begin{multicols}{2}
  \item El que aparece en un linaje posterior
  \item El que está presente en el ancestro
  \item El que es antiguo
  \item El que es primitivo
  \item El que es común
  \item El que es general
\end{multicols}
\end{enumerate}
\vspace{0.15cm}
\item ¿Qué es un árbol filogenético?
\begin{enumerate}[label=\textbf{\Alph*)}]
\begin{multicols}{2}
  \item Un diagrama que representa relaciones evolutivas
  \item Una lista de especies
  \item Una clave dicotómica
  \item Una colección de especímenes
  \item Un catálogo de nombres
  \item Un herbario
\end{multicols}
\end{enumerate}
\vspace{0.15cm}
\item ¿Qué es un ancestro común?
\begin{enumerate}[label=\textbf{\Alph*)}]
\begin{multicols}{2}
  \item El organismo del que descienden dos o más linajes
  \item El organismo más antiguo
  \item El organismo más grande
  \item El organismo más pequeño
  \item El organismo más común
  \item El organismo más raro
\end{multicols}
\end{enumerate}
\vspace{0.15cm}
\item ¿Qué es la especiación?
\begin{enumerate}[label=\textbf{\Alph*)}]
\begin{multicols}{2}
  \item El proceso de formación de nuevas especies
  \item El proceso de extinción
  \item El proceso de migración
  \item El proceso de hibridación
  \item El proceso de mutación
  \item El proceso de selección
\end{multicols}
\end{enumerate}
\vspace{0.15cm}
\item ¿Qué es el aislamiento reproductivo?
\begin{enumerate}[label=\textbf{\Alph*)}]
\begin{multicols}{2}
  \item La incapacidad de dos poblaciones de reproducirse entre sí
  \item La capacidad de reproducirse
  \item La capacidad de hibridar
  \item La capacidad de migrar
  \item La capacidad de adaptarse
  \item La capacidad de evolucionar
\end{multicols}
\end{enumerate}
\vspace{0.15cm}
\item ¿Qué es una población?
\begin{enumerate}[label=\textbf{\Alph*)}]
\begin{multicols}{2}
  \item Un conjunto de individuos de la misma especie que viven en un área
  \item Un conjunto de especies
  \item Un conjunto de géneros
  \item Un conjunto de familias
  \item Un conjunto de órdenes
  \item Un conjunto de clases
\end{multicols}
\end{enumerate}
\vspace{0.15cm}
\item ¿Qué es la biodiversidad?
\begin{enumerate}[label=\textbf{\Alph*)}]
\begin{multicols}{2}
  \item La variedad de seres vivos en un lugar
  \item La cantidad de individuos
  \item La cantidad de especies
  \item La cantidad de ecosistemas
  \item La cantidad de genes
  \item La cantidad de poblaciones
\end{multicols}
\end{enumerate}
\vspace{0.15cm}
\item ¿Qué diferencia hay entre taxonomía y sistemática?
\begin{enumerate}[label=\textbf{\Alph*)}]
\begin{multicols}{2}
  \item La taxonomía clasifica y nombra; la sistemática estudia relaciones evolutivas
  \item Son sinónimos exactos
  \item La taxonomía estudia evolución; la sistemática nombra
  \item La taxonomía es molecular; la sistemática es morfológica
  \item La taxonomía es antigua; la sistemática es moderna
  \item La taxonomía es botánica; la sistemática es zoológica
\end{multicols}
\end{enumerate}
\vspace{0.15cm}
\item ¿Qué es la nomenclatura?
\begin{enumerate}[label=\textbf{\Alph*)}]
\begin{multicols}{2}
  \item El conjunto de reglas para nombrar a los organismos
  \item El conjunto de reglas para clasificar
  \item El conjunto de reglas para describir
  \item El conjunto de reglas para coleccionar
  \item El conjunto de reglas para conservar
  \item El conjunto de reglas para estudiar
\end{multicols}
\end{enumerate}
\vspace{0.15cm}
\item ¿Qué código regula la nomenclatura zoológica?
\begin{enumerate}[label=\textbf{\Alph*)}]
\begin{multicols}{2}
  \item El Código Internacional de Nomenclatura Zoológica
  \item El Código Internacional de Nomenclatura Botánica
  \item El Código Internacional de Nomenclatura Bacteriológica
  \item El Código Internacional de Nomenclatura Viral
  \item El Código Internacional de Nomenclatura Fúngica
  \item El Código Internacional de Nomenclatura Algal
\end{multicols}
\end{enumerate}
\vspace{0.15cm}
\item ¿Qué código regula la nomenclatura botánica?
\begin{enumerate}[label=\textbf{\Alph*)}]
\begin{multicols}{2}
  \item El Código Internacional de Nomenclatura para algas, hongos y plantas
  \item El Código Internacional de Nomenclatura Zoológica
  \item El Código Internacional de Nomenclatura Bacteriológica
  \item El Código Internacional de Nomenclatura Viral
  \item El Código Internacional de Nomenclatura Médica
  \item El Código Internacional de Nomenclatura Agrícola
\end{multicols}
\end{enumerate}
\vspace{0.15cm}
\item ¿Qué es el principio de prioridad?
\begin{enumerate}[label=\textbf{\Alph*)}]
\begin{multicols}{2}
  \item El nombre válido más antiguo tiene prioridad
  \item El nombre más nuevo tiene prioridad
  \item El nombre más usado tiene prioridad
  \item El nombre más corto tiene prioridad
  \item El nombre más descriptivo tiene prioridad
  \item El nombre más común tiene prioridad
\end{multicols}
\end{enumerate}
\vspace{0.15cm}
\item ¿Qué es un tipo nomenclatural?
\begin{enumerate}[label=\textbf{\Alph*)}]
\begin{multicols}{2}
  \item El espécimen de referencia que define un nombre
  \item El nombre más antiguo
  \item El nombre más usado
  \item El nombre más descriptivo
  \item El nombre más común
  \item El nombre más científico
\end{multicols}
\end{enumerate}
\vspace{0.15cm}
\item ¿Qué es un holotipo?
\begin{enumerate}[label=\textbf{\Alph*)}]
\begin{multicols}{2}
  \item El espécimen único designado por el autor original
  \item Un espécimen duplicado
  \item Un espécimen perdido
  \item Un espécimen ilustrado
  \item Un espécimen descrito después
  \item Un espécimen sin nombre
\end{multicols}
\end{enumerate}
\vspace{0.15cm}
\item ¿Qué es un paratipo?
\begin{enumerate}[label=\textbf{\Alph*)}]
\begin{multicols}{2}
  \item Un espécimen adicional citado por el autor original
  \item El espécimen único
  \item El espécimen perdido
  \item El espécimen ilustrado
  \item El espécimen descrito después
  \item El espécimen sin nombre
\end{multicols}
\end{enumerate}
\vspace{0.15cm}
\item ¿Qué es un lectotipo?
\begin{enumerate}[label=\textbf{\Alph*)}]
\begin{multicols}{2}
  \item Un espécimen seleccionado posteriormente como tipo
  \item El espécimen único original
  \item Un espécimen duplicado
  \item Un espécimen perdido
  \item Un espécimen ilustrado
  \item Un espécimen sin nombre
\end{multicols}
\end{enumerate}
\vspace{0.15cm}
\item ¿Qué es un neotipo?
\begin{enumerate}[label=\textbf{\Alph*)}]
\begin{multicols}{2}
  \item Un nuevo espécimen tipo cuando el original se pierde
  \item El espécimen único original
  \item Un espécimen duplicado
  \item Un espécimen ilustrado
  \item Un espécimen descrito después
  \item Un espécimen sin nombre
\end{multicols}
\end{enumerate}
\vspace{0.15cm}
\item ¿Qué es la sinonimia?
\begin{enumerate}[label=\textbf{\Alph*)}]
\begin{multicols}{2}
  \item Dos nombres diferentes para el mismo taxón
  \item Dos taxones con el mismo nombre
  \item Dos nombres para dos taxones
  \item Un nombre para dos taxones
  \item Un nombre para un taxón
  \item Dos nombres para dos taxones
\end{multicols}
\end{enumerate}
\vspace{0.15cm}
\item ¿Qué es la homonimia?
\begin{enumerate}[label=\textbf{\Alph*)}]
\begin{multicols}{2}
  \item El mismo nombre para dos taxones diferentes
  \item Dos nombres para el mismo taxón
  \item Un nombre para un taxón
  \item Dos nombres para dos taxones
  \item Un nombre para dos taxones
  \item Dos nombres para un taxón
\end{multicols}
\end{enumerate}
\vspace{0.15cm}
\item ¿Qué es la autoridad en nomenclatura?
\begin{enumerate}[label=\textbf{\Alph*)}]
\begin{multicols}{2}
  \item El autor que describe por primera vez un taxón
  \item El autor que clasifica
  \item El autor que colecciona
  \item El autor que ilustra
  \item El autor que traduce
  \item El autor que publica
\end{multicols}
\end{enumerate}
\vspace{0.15cm}
\item ¿Qué es un nombre válido?
\begin{enumerate}[label=\textbf{\Alph*)}]
\begin{multicols}{2}
  \item El que cumple las reglas del código
  \item El que no cumple las reglas
  \item El que es más usado
  \item El que es más antiguo
  \item El que es más descriptivo
  \item El que es más común
\end{multicols}
\end{enumerate}
\vspace{0.15cm}
\item ¿Qué es un nombre inválido?
\begin{enumerate}[label=\textbf{\Alph*)}]
\begin{multicols}{2}
  \item El que no cumple las reglas del código
  \item El que cumple las reglas
  \item El que es más usado
  \item El que es más antiguo
  \item El que es más descriptivo
  \item El que es más común
\end{multicols}
\end{enumerate}
\vspace{0.15cm}
\item ¿Qué es un diagnóstico taxonómico?
\begin{enumerate}[label=\textbf{\Alph*)}]
\begin{multicols}{2}
  \item La descripción de los caracteres que definen un taxón
  \item La lista de especies
  \item La clave dicotómica
  \item La colección de especímenes
  \item El catálogo de nombres
  \item El herbario
\end{multicols}
\end{enumerate}
\vspace{0.15cm}
\item ¿Qué es una clave dicotómica?
\begin{enumerate}[label=\textbf{\Alph*)}]
\begin{multicols}{2}
  \item Una herramienta con pares de opciones para identificar organismos
  \item Una lista de especies
  \item Una colección de especímenes
  \item Un catálogo de nombres
  \item Un herbario
  \item Un museo
\end{multicols}
\end{enumerate}
\vspace{0.15cm}
\item ¿Qué es una colección biológica?
\begin{enumerate}[label=\textbf{\Alph*)}]
\begin{multicols}{2}
  \item Un conjunto de especímenes preservados para estudio
  \item Una lista de especies
  \item Una clave dicotómica
  \item Un catálogo de nombres
  \item Un herbario
  \item Un museo
\end{multicols}
\end{enumerate}
\vspace{0.15cm}
\item ¿Qué es un herbario?
\begin{enumerate}[label=\textbf{\Alph*)}]
\begin{multicols}{2}
  \item Una colección de plantas secas y prensadas
  \item Una colección de animales
  \item Una colección de hongos
  \item Una colección de bacterias
  \item Una colección de fósiles
  \item Una colección de minerales
\end{multicols}
\end{enumerate}
\vspace{0.15cm}
\item ¿Qué es un museo de historia natural?
\begin{enumerate}[label=\textbf{\Alph*)}]
\begin{multicols}{2}
  \item Una institución que conserva y estudia especímenes
  \item Una institución que conserva plantas
  \item Una institución que conserva animales
  \item Una institución que conserva hongos
  \item Una institución que conserva bacterias
  \item Una institución que conserva fósiles
\end{multicols}
\end{enumerate}
\vspace{0.15cm}
\item ¿Qué es una especie tipo?
\begin{enumerate}[label=\textbf{\Alph*)}]
\begin{multicols}{2}
  \item La especie que define un género
  \item El espécimen que define una especie
  \item El género que define una familia
  \item La familia que define un orden
  \item El orden que define una clase
  \item La clase que define un filo
\end{multicols}
\end{enumerate}
\vspace{0.15cm}
\item ¿Qué es un género tipo?
\begin{enumerate}[label=\textbf{\Alph*)}]
\begin{multicols}{2}
  \item El género que define una familia
  \item La especie que define un género
  \item El espécimen que define una especie
  \item La familia que define un orden
  \item El orden que define una clase
  \item La clase que define un filo
\end{multicols}
\end{enumerate}
\vspace{0.15cm}
\item ¿Qué es una familia tipo?
\begin{enumerate}[label=\textbf{\Alph*)}]
\begin{multicols}{2}
  \item La familia que define un orden
  \item El género que define una familia
  \item La especie que define un género
  \item El espécimen que define una especie
  \item El orden que define una clase
  \item La clase que define un filo ---
\end{multicols}
\end{enumerate}
\vspace{0.15cm}
\end{enumerate}\n
\subsubsection*{SECCIÓN II: NOMENCLATURA BINOMIAL (51–100)}

\begin{enumerate}[label=\textbf{\arabic*.}]
\setcounter{enumi}{50}
\item ¿Cómo se escribe correctamente un nombre científico?
\begin{enumerate}[label=\textbf{\Alph*)}]
\begin{multicols}{2}
  \item Género con mayúscula y especie con minúscula, en cursiva
  \item Género con minúscula y especie con mayúscula
  \item Todo con mayúscula
  \item Todo con minúscula
  \item Género y especie con mayúscula
  \item Género y especie con minúscula
\end{multicols}
\end{enumerate}
\vspace{0.15cm}
\item ¿Cuál es la forma correcta de *Homo sapiens*?
\begin{enumerate}[label=\textbf{\Alph*)}]
\begin{multicols}{2}
  \item *Homo sapiens*
  \item *homo sapiens*
  \item *Homo Sapiens*
  \item HOMO SAPIENS
  \item homo Sapiens
  \item Homo Sapiens
\end{multicols}
\end{enumerate}
\vspace{0.15cm}
\item ¿Qué indica el primer nombre en la nomenclatura binomial?
\begin{enumerate}[label=\textbf{\Alph*)}]
\begin{multicols}{2}
  \item El género
  \item La especie
  \item La familia
  \item El orden
  \item La clase
  \item El filo
\end{multicols}
\end{enumerate}
\vspace{0.15cm}
\item ¿Qué indica el segundo nombre en la nomenclatura binomial?
\begin{enumerate}[label=\textbf{\Alph*)}]
\begin{multicols}{2}
  \item El epíteto específico
  \item El género
  \item La familia
  \item El orden
  \item La clase
  \item El filo
\end{multicols}
\end{enumerate}
\vspace{0.15cm}
\item ¿Cómo se abrevia *Escherichia coli*?
\begin{enumerate}[label=\textbf{\Alph*)}]
\begin{multicols}{2}
  \item *E. coli*
  \item *E. Coli*
  \item *e. coli*
  \item *E. Coli*
  \item *E.Coli*
  \item *E-coli*
\end{multicols}
\end{enumerate}
\vspace{0.15cm}
\item ¿Qué significa “sp.” en taxonomía?
\begin{enumerate}[label=\textbf{\Alph*)}]
\begin{multicols}{2}
  \item Especie no identificada
  \item Varias especies
  \item Subespecie
  \item Híbrido
  \item Variedad
  \item Forma
\end{multicols}
\end{enumerate}
\vspace{0.15cm}
\item ¿Qué significa “spp.” en taxonomía?
\begin{enumerate}[label=\textbf{\Alph*)}]
\begin{multicols}{2}
  \item Varias especies
  \item Especie no identificada
  \item Subespecie
  \item Híbrido
  \item Variedad
  \item Forma
\end{multicols}
\end{enumerate}
\vspace{0.15cm}
\item ¿Qué significa “cf.” en taxonomía?
\begin{enumerate}[label=\textbf{\Alph*)}]
\begin{multicols}{2}
  \item Confrontar, probablemente esa especie
  \item Varias especies
  \item Subespecie
  \item Híbrido
  \item Variedad
  \item Forma
\end{multicols}
\end{enumerate}
\vspace{0.15cm}
\item ¿Qué significa “aff.” en taxonomía?
\begin{enumerate}[label=\textbf{\Alph*)}]
\begin{multicols}{2}
  \item Afín, relacionado pero no confirmado
  \item Varias especies
  \item Subespecie
  \item Híbrido
  \item Variedad
  \item Forma
\end{multicols}
\end{enumerate}
\vspace{0.15cm}
\item ¿Qué es un trinomio?
\begin{enumerate}[label=\textbf{\Alph*)}]
\begin{multicols}{2}
  \item Nombre con género, especie y subespecie
  \item Nombre con género y especie
  \item Nombre con familia y género
  \item Nombre con orden y familia
  \item Nombre con clase y orden
  \item Nombre con filo y clase
\end{multicols}
\end{enumerate}
\vspace{0.15cm}
\item ¿Cómo se escribe una subespecie?
\begin{enumerate}[label=\textbf{\Alph*)}]
\begin{multicols}{2}
  \item Género especie subespecie, todo en cursiva
  \item Género Especie Subespecie
  \item género especie subespecie
  \item Género especie subespecie
  \item GENERO ESPECIE SUBESPECIE
  \item género Especie subespecie
\end{multicols}
\end{enumerate}
\vspace{0.15cm}
\item ¿Qué es una variedad en botánica?
\begin{enumerate}[label=\textbf{\Alph*)}]
\begin{multicols}{2}
  \item Un rango infraespecífico
  \item Un género
  \item Una familia
  \item Un orden
  \item Una clase
  \item Un filo
\end{multicols}
\end{enumerate}
\vspace{0.15cm}
\item ¿Qué es una forma en botánica?
\begin{enumerate}[label=\textbf{\Alph*)}]
\begin{multicols}{2}
  \item Un rango infraespecífico menor
  \item Un género
  \item Una familia
  \item Un orden
  \item Una clase
  \item Un filo
\end{multicols}
\end{enumerate}
\vspace{0.15cm}
\item ¿Qué es un cultivar?
\begin{enumerate}[label=\textbf{\Alph*)}]
\begin{multicols}{2}
  \item Una variedad cultivada por el ser humano
  \item Una especie silvestre
  \item Un género
  \item Una familia
  \item Un orden
  \item Una clase
\end{multicols}
\end{enumerate}
\vspace{0.15cm}
\item ¿Cómo se escribe un híbrido?
\begin{enumerate}[label=\textbf{\Alph*)}]
\begin{multicols}{2}
  \item Con el signo × entre los nombres
  \item Con un guion
  \item Con un punto
  \item Con una coma
  \item Con dos puntos
  \item Con punto y coma
\end{multicols}
\end{enumerate}
\vspace{0.15cm}
\item ¿Qué es un nombre común?
\begin{enumerate}[label=\textbf{\Alph*)}]
\begin{multicols}{2}
  \item El nombre popular de un organismo
  \item El nombre científico
  \item El nombre binomial
  \item El nombre trinomial
  \item El nombre taxonómico
  \item El nombre nomenclatural
\end{multicols}
\end{enumerate}
\vspace{0.15cm}
\item ¿Por qué se usa el nombre científico?
\begin{enumerate}[label=\textbf{\Alph*)}]
\begin{multicols}{2}
  \item Porque es universal y evita confusiones
  \item Porque es más corto
  \item Porque es más fácil
  \item Porque es más bonito
  \item Porque es más antiguo
  \item Porque es más común
\end{multicols}
\end{enumerate}
\vspace{0.15cm}
\item ¿Cuál es el nombre científico del ser humano?
\begin{enumerate}[label=\textbf{\Alph*)}]
\begin{multicols}{2}
  \item *Homo sapiens*
  \item *Homo erectus*
  \item *Homo habilis*
  \item *Homo neanderthalensis*
  \item *Pan troglodytes*
  \item *Gorilla gorilla*
\end{multicols}
\end{enumerate}
\vspace{0.15cm}
\item ¿Cuál es el nombre científico del león?
\begin{enumerate}[label=\textbf{\Alph*)}]
\begin{multicols}{2}
  \item *Panthera leo*
  \item *Panthera tigris*
  \item *Panthera pardus*
  \item *Panthera onca*
  \item *Felis catus*
  \item *Canis lupus*
\end{multicols}
\end{enumerate}
\vspace{0.15cm}
\item ¿Cuál es el nombre científico del lobo?
\begin{enumerate}[label=\textbf{\Alph*)}]
\begin{multicols}{2}
  \item *Canis lupus*
  \item *Canis familiaris*
  \item *Canis latrans*
  \item *Canis aureus*
  \item *Vulpes vulpes*
  \item *Lycaon pictus*
\end{multicols}
\end{enumerate}
\vspace{0.15cm}
\item ¿Cuál es el nombre científico del maíz?
\begin{enumerate}[label=\textbf{\Alph*)}]
\begin{multicols}{2}
  \item *Zea mays*
  \item *Triticum aestivum*
  \item *Oryza sativa*
  \item *Hordeum vulgare*
  \item *Sorghum bicolor*
  \item *Avena sativa*
\end{multicols}
\end{enumerate}
\vspace{0.15cm}
\item ¿Cuál es el nombre científico del trigo?
\begin{enumerate}[label=\textbf{\Alph*)}]
\begin{multicols}{2}
  \item *Triticum aestivum*
  \item *Zea mays*
  \item *Oryza sativa*
  \item *Hordeum vulgare*
  \item *Sorghum bicolor*
  \item *Avena sativa*
\end{multicols}
\end{enumerate}
\vspace{0.15cm}
\item ¿Cuál es el nombre científico del arroz?
\begin{enumerate}[label=\textbf{\Alph*)}]
\begin{multicols}{2}
  \item *Oryza sativa*
  \item *Zea mays*
  \item *Triticum aestivum*
  \item *Hordeum vulgare*
  \item *Sorghum bicolor*
  \item *Avena sativa*
\end{multicols}
\end{enumerate}
\vspace{0.15cm}
\item ¿Cuál es el nombre científico del perro?
\begin{enumerate}[label=\textbf{\Alph*)}]
\begin{multicols}{2}
  \item *Canis lupus familiaris*
  \item *Canis lupus*
  \item *Felis catus*
  \item *Vulpes vulpes*
  \item *Canis latrans*
  \item *Canis aureus*
\end{multicols}
\end{enumerate}
\vspace{0.15cm}
\item ¿Cuál es el nombre científico del gato?
\begin{enumerate}[label=\textbf{\Alph*)}]
\begin{multicols}{2}
  \item *Felis catus*
  \item *Canis lupus*
  \item *Panthera leo*
  \item *Panthera tigris*
  \item *Lynx lynx*
  \item *Puma concolor*
\end{multicols}
\end{enumerate}
\vspace{0.15cm}
\item ¿Cuál es el nombre científico del tigre?
\begin{enumerate}[label=\textbf{\Alph*)}]
\begin{multicols}{2}
  \item *Panthera tigris*
  \item *Panthera leo*
  \item *Panthera pardus*
  \item *Panthera onca*
  \item *Felis catus*
  \item *Canis lupus*
\end{multicols}
\end{enumerate}
\vspace{0.15cm}
\item ¿Cuál es el nombre científico del elefante africano?
\begin{enumerate}[label=\textbf{\Alph*)}]
\begin{multicols}{2}
  \item *Loxodonta africana*
  \item *Elephas maximus*
  \item *Elephas africanus*
  \item *Loxodonta asiatica*
  \item *Mammuthus primigenius*
  \item *Elephas indicus*
\end{multicols}
\end{enumerate}
\vspace{0.15cm}
\item ¿Cuál es el nombre científico de la ballena azul?
\begin{enumerate}[label=\textbf{\Alph*)}]
\begin{multicols}{2}
  \item *Balaenoptera musculus*
  \item *Physeter macrocephalus*
  \item *Megaptera novaeangliae*
  \item *Eschrichtius robustus*
  \item *Balaena mysticetus*
  \item *Orcinus orca*
\end{multicols}
\end{enumerate}
\vspace{0.15cm}
\item ¿Cuál es el nombre científico del murciélago común?
\begin{enumerate}[label=\textbf{\Alph*)}]
\begin{multicols}{2}
  \item *Myotis myotis*
  \item *Rousettus aegyptiacus*
  \item *Desmodus rotundus*
  \item *Pteropus vampyrus*
  \item *Rhinolophus ferrumequinum*
  \item *Eptesicus serotinus*
\end{multicols}
\end{enumerate}
\vspace{0.15cm}
\item ¿Cuál es el nombre científico del gorila?
\begin{enumerate}[label=\textbf{\Alph*)}]
\begin{multicols}{2}
  \item *Gorilla gorilla*
  \item *Pan troglodytes*
  \item *Pongo pygmaeus*
  \item *Homo sapiens*
  \item *Macaca mulatta*
  \item *Papio hamadryas*
\end{multicols}
\end{enumerate}
\vspace{0.15cm}
\item ¿Cuál es el nombre científico del chimpancé?
\begin{enumerate}[label=\textbf{\Alph*)}]
\begin{multicols}{2}
  \item *Pan troglodytes*
  \item *Gorilla gorilla*
  \item *Pongo pygmaeus*
  \item *Homo sapiens*
  \item *Macaca mulatta*
  \item *Papio hamadryas*
\end{multicols}
\end{enumerate}
\vspace{0.15cm}
\item ¿Cuál es el nombre científico del orangután?
\begin{enumerate}[label=\textbf{\Alph*)}]
\begin{multicols}{2}
  \item *Pongo pygmaeus*
  \item *Gorilla gorilla*
  \item *Pan troglodytes*
  \item *Homo sapiens*
  \item *Macaca mulatta*
  \item *Papio hamadryas*
\end{multicols}
\end{enumerate}
\vspace{0.15cm}
\item ¿Cuál es el nombre científico del caballo?
\begin{enumerate}[label=\textbf{\Alph*)}]
\begin{multicols}{2}
  \item *Equus caballus*
  \item *Equus asinus*
  \item *Equus zebra*
  \item *Equus quagga*
  \item *Equus ferus*
  \item *Equus hemionus*
\end{multicols}
\end{enumerate}
\vspace{0.15cm}
\item ¿Cuál es el nombre científico del burro?
\begin{enumerate}[label=\textbf{\Alph*)}]
\begin{multicols}{2}
  \item *Equus asinus*
  \item *Equus caballus*
  \item *Equus zebra*
  \item *Equus quagga*
  \item *Equus ferus*
  \item *Equus hemionus*
\end{multicols}
\end{enumerate}
\vspace{0.15cm}
\item ¿Cuál es el nombre científico de la cebra?
\begin{enumerate}[label=\textbf{\Alph*)}]
\begin{multicols}{2}
  \item *Equus zebra*
  \item *Equus asinus*
  \item *Equus caballus*
  \item *Equus quagga*
  \item *Equus ferus*
  \item *Equus hemionus*
\end{multicols}
\end{enumerate}
\vspace{0.15cm}
\item ¿Cuál es el nombre científico de la abeja melífera?
\begin{enumerate}[label=\textbf{\Alph*)}]
\begin{multicols}{2}
  \item *Apis mellifera*
  \item *Apis cerana*
  \item *Apis dorsata*
  \item *Apis florea*
  \item *Bombus terrestris*
  \item *Vespula vulgaris*
\end{multicols}
\end{enumerate}
\vspace{0.15cm}
\item ¿Cuál es el nombre científico de la mosca de la fruta?
\begin{enumerate}[label=\textbf{\Alph*)}]
\begin{multicols}{2}
  \item *Drosophila melanogaster*
  \item *Musca domestica*
  \item *Anopheles gambiae*
  \item *Aedes aegypti*
  \item *Culex pipiens*
  \item *Glossina morsitans*
\end{multicols}
\end{enumerate}
\vspace{0.15cm}
\item ¿Cuál es el nombre científico de la bacteria intestinal?
\begin{enumerate}[label=\textbf{\Alph*)}]
\begin{multicols}{2}
  \item *Escherichia coli*
  \item *Salmonella typhi*
  \item *Staphylococcus aureus*
  \item *Streptococcus pyogenes*
  \item *Bacillus subtilis*
  \item *Pseudomonas aeruginosa*
\end{multicols}
\end{enumerate}
\vspace{0.15cm}
\item ¿Cuál es el nombre científico de la levadura de pan?
\begin{enumerate}[label=\textbf{\Alph*)}]
\begin{multicols}{2}
  \item *Saccharomyces cerevisiae*
  \item *Candida albicans*
  \item *Aspergillus niger*
  \item *Penicillium chrysogenum*
  \item *Rhizopus stolonifer*
  \item *Agaricus bisporus*
\end{multicols}
\end{enumerate}
\vspace{0.15cm}
\item ¿Cuál es el nombre científico del champiñón común?
\begin{enumerate}[label=\textbf{\Alph*)}]
\begin{multicols}{2}
  \item *Agaricus bisporus*
  \item *Amanita muscaria*
  \item *Pleurotus ostreatus*
  \item *Lentinula edodes*
  \item *Boletus edulis*
  \item *Tuber melanosporum*
\end{multicols}
\end{enumerate}
\vspace{0.15cm}
\item ¿Cuál es el nombre científico del pino silvestre?
\begin{enumerate}[label=\textbf{\Alph*)}]
\begin{multicols}{2}
  \item *Pinus sylvestris*
  \item *Pinus pinaster*
  \item *Pinus halepensis*
  \item *Pinus nigra*
  \item *Pinus pinea*
  \item *Pinus radiata*
\end{multicols}
\end{enumerate}
\vspace{0.15cm}
\item ¿Cuál es el nombre científico del ginkgo?
\begin{enumerate}[label=\textbf{\Alph*)}]
\begin{multicols}{2}
  \item *Ginkgo biloba*
  \item *Ginkgo triloba*
  \item *Ginkgo quadrifolia*
  \item *Ginkgo sempervirens*
  \item *Ginkgo chinensis*
  \item *Ginkgo japonica*
\end{multicols}
\end{enumerate}
\vspace{0.15cm}
\item ¿Cuál es el nombre científico de la rosa?
\begin{enumerate}[label=\textbf{\Alph*)}]
\begin{multicols}{2}
  \item *Rosa canina*
  \item *Rosa gallica*
  \item *Rosa damascena*
  \item *Rosa chinensis*
  \item *Rosa multiflora*
  \item *Rosa rubiginosa*
\end{multicols}
\end{enumerate}
\vspace{0.15cm}
\item ¿Cuál es el nombre científico del girasol?
\begin{enumerate}[label=\textbf{\Alph*)}]
\begin{multicols}{2}
  \item *Helianthus annuus*
  \item *Helianthus tuberosus*
  \item *Helianthus maximiliani*
  \item *Helianthus petiolaris*
  \item *Helianthus argophyllus*
  \item *Helianthus debilis*
\end{multicols}
\end{enumerate}
\vspace{0.15cm}
\item ¿Cuál es el nombre científico de la papa?
\begin{enumerate}[label=\textbf{\Alph*)}]
\begin{multicols}{2}
  \item *Solanum tuberosum*
  \item *Solanum lycopersicum*
  \item *Solanum melongena*
  \item *Solanum nigrum*
  \item *Solanum dulcamara*
  \item *Solanum pseudocapsicum*
\end{multicols}
\end{enumerate}
\vspace{0.15cm}
\item ¿Cuál es el nombre científico del tomate?
\begin{enumerate}[label=\textbf{\Alph*)}]
\begin{multicols}{2}
  \item *Solanum lycopersicum*
  \item *Solanum tuberosum*
  \item *Solanum melongena*
  \item *Solanum nigrum*
  \item *Solanum dulcamara*
  \item *Solanum pseudocapsicum*
\end{multicols}
\end{enumerate}
\vspace{0.15cm}
\item ¿Cuál es el nombre científico de la berenjena?
\begin{enumerate}[label=\textbf{\Alph*)}]
\begin{multicols}{2}
  \item *Solanum melongena*
  \item *Solanum tuberosum*
  \item *Solanum lycopersicum*
  \item *Solanum nigrum*
  \item *Solanum dulcamara*
  \item *Solanum pseudocapsicum*
\end{multicols}
\end{enumerate}
\vspace{0.15cm}
\item ¿Cuál es el nombre científico del arroz?
\begin{enumerate}[label=\textbf{\Alph*)}]
\begin{multicols}{2}
  \item *Oryza sativa*
  \item *Oryza rufipogon*
  \item *Oryza glaberrima*
  \item *Oryza barthii*
  \item *Oryza longistaminata*
  \item *Oryza punctata*
\end{multicols}
\end{enumerate}
\vspace{0.15cm}
\item ¿Cuál es el nombre científico del cafeto?
\begin{enumerate}[label=\textbf{\Alph*)}]
\begin{multicols}{2}
  \item *Coffea arabica*
  \item *Coffea canephora*
  \item *Coffea liberica*
  \item *Coffea excelsa*
  \item *Coffea stenophylla*
  \item *Coffea racemosa*
\end{multicols}
\end{enumerate}
\vspace{0.15cm}
\item ¿Cuál es el nombre científico del cacao?
\begin{enumerate}[label=\textbf{\Alph*)}]
\begin{multicols}{2}
  \item *Theobroma cacao*
  \item *Theobroma bicolor*
  \item *Theobroma grandiflorum*
  \item *Theobroma speciosum*
  \item *Theobroma subincanum*
  \item *Theobroma microcarpum* ---
\end{multicols}
\end{enumerate}
\vspace{0.15cm}
\end{enumerate}\n
\subsubsection*{SECCIÓN III: DOMINIOS Y REINOS (101–150)}

\begin{enumerate}[label=\textbf{\arabic*.}]
\setcounter{enumi}{100}
\item ¿Cuántos dominios se reconocen actualmente?
\begin{enumerate}[label=\textbf{\Alph*)}]
\begin{multicols}{2}
  \item Tres
  \item Dos
  \item Cuatro
  \item Cinco
  \item Seis
  \item Siete
\end{multicols}
\end{enumerate}
\vspace{0.15cm}
\item ¿Cuáles son los tres dominios?
\begin{enumerate}[label=\textbf{\Alph*)}]
\begin{multicols}{2}
  \item Bacteria, Archaea y Eukarya
  \item Monera, Protista y Fungi
  \item Plantae, Animalia y Fungi
  \item Procariota, Eucariota y Virus
  \item Archaea, Protista y Plantae
  \item Bacteria, Fungi y Animalia
\end{multicols}
\end{enumerate}
\vspace{0.15cm}
\item ¿Quién propuso el sistema de tres dominios?
\begin{enumerate}[label=\textbf{\Alph*)}]
\begin{multicols}{2}
  \item Carl Woese
  \item Carlos Linneo
  \item Robert Whittaker
  \item Charles Darwin
  \item Gregor Mendel
  \item Louis Pasteur
\end{multicols}
\end{enumerate}
\vspace{0.15cm}
\item ¿Qué dominio agrupa a las bacterias?
\begin{enumerate}[label=\textbf{\Alph*)}]
\begin{multicols}{2}
  \item Bacteria
  \item Archaea
  \item Eukarya
  \item Monera
  \item Protista
  \item Fungi
\end{multicols}
\end{enumerate}
\vspace{0.15cm}
\item ¿Qué dominio agrupa a las arqueas?
\begin{enumerate}[label=\textbf{\Alph*)}]
\begin{multicols}{2}
  \item Archaea
  \item Bacteria
  \item Eukarya
  \item Monera
  \item Protista
  \item Fungi
\end{multicols}
\end{enumerate}
\vspace{0.15cm}
\item ¿Qué dominio agrupa a los eucariotas?
\begin{enumerate}[label=\textbf{\Alph*)}]
\begin{multicols}{2}
  \item Eukarya
  \item Bacteria
  \item Archaea
  \item Monera
  \item Protista
  \item Fungi
\end{multicols}
\end{enumerate}
\vspace{0.15cm}
\item ¿Cuántos reinos propuso Whittaker?
\begin{enumerate}[label=\textbf{\Alph*)}]
\begin{multicols}{2}
  \item Cinco
  \item Tres
  \item Cuatro
  \item Seis
  \item Dos
  \item Siete
\end{multicols}
\end{enumerate}
\vspace{0.15cm}
\item ¿Cuáles son los cinco reinos de Whittaker?
\begin{enumerate}[label=\textbf{\Alph*)}]
\begin{multicols}{2}
  \item Monera, Protista, Fungi, Plantae y Animalia
  \item Bacteria, Archaea, Fungi, Plantae y Animalia
  \item Monera, Protista, Plantae, Animalia y Virus
  \item Bacteria, Protista, Fungi, Plantae y Animalia
  \item Archaea, Protista, Fungi, Plantae y Animalia
  \item Monera, Archaea, Fungi, Plantae y Animalia
\end{multicols}
\end{enumerate}
\vspace{0.15cm}
\item ¿Qué reino agrupaba a los procariotas?
\begin{enumerate}[label=\textbf{\Alph*)}]
\begin{multicols}{2}
  \item Monera
  \item Protista
  \item Fungi
  \item Plantae
  \item Animalia
  \item Archaea
\end{multicols}
\end{enumerate}
\vspace{0.15cm}
\item ¿Qué reino agrupa a los protozoos y algas?
\begin{enumerate}[label=\textbf{\Alph*)}]
\begin{multicols}{2}
  \item Protista
  \item Monera
  \item Fungi
  \item Plantae
  \item Animalia
  \item Archaea
\end{multicols}
\end{enumerate}
\vspace{0.15cm}
\item ¿Qué reino agrupa a los hongos?
\begin{enumerate}[label=\textbf{\Alph*)}]
\begin{multicols}{2}
  \item Fungi
  \item Protista
  \item Monera
  \item Plantae
  \item Animalia
  \item Archaea
\end{multicols}
\end{enumerate}
\vspace{0.15cm}
\item ¿Qué reino agrupa a las plantas?
\begin{enumerate}[label=\textbf{\Alph*)}]
\begin{multicols}{2}
  \item Plantae
  \item Fungi
  \item Protista
  \item Monera
  \item Animalia
  \item Archaea
\end{multicols}
\end{enumerate}
\vspace{0.15cm}
\item ¿Qué reino agrupa a los animales?
\begin{enumerate}[label=\textbf{\Alph*)}]
\begin{multicols}{2}
  \item Animalia
  \item Plantae
  \item Fungi
  \item Protista
  \item Monera
  \item Archaea
\end{multicols}
\end{enumerate}
\vspace{0.15cm}
\item ¿Qué tipo de células tienen las bacterias?
\begin{enumerate}[label=\textbf{\Alph*)}]
\begin{multicols}{2}
  \item Procariotas
  \item Eucariotas
  \item Sin células
  \item Pluricelulares
  \item Con núcleo
  \item Con mitocondrias
\end{multicols}
\end{enumerate}
\vspace{0.15cm}
\item ¿Qué tipo de células tienen las arqueas?
\begin{enumerate}[label=\textbf{\Alph*)}]
\begin{multicols}{2}
  \item Procariotas
  \item Eucariotas
  \item Sin células
  \item Pluricelulares
  \item Con núcleo
  \item Con mitocondrias
\end{multicols}
\end{enumerate}
\vspace{0.15cm}
\item ¿Qué tipo de células tienen los eucariotas?
\begin{enumerate}[label=\textbf{\Alph*)}]
\begin{multicols}{2}
  \item Eucariotas
  \item Procariotas
  \item Sin células
  \item Pluricelulares
  \item Sin núcleo
  \item Sin mitocondrias
\end{multicols}
\end{enumerate}
\vspace{0.15cm}
\item ¿Las bacterias son unicelulares o pluricelulares?
\begin{enumerate}[label=\textbf{\Alph*)}]
\begin{multicols}{2}
  \item Unicelulares
  \item Pluricelulares
  \item Coloniales
  \item Filamentosas
  \item Ambas
  \item Ninguna
\end{multicols}
\end{enumerate}
\vspace{0.15cm}
\item ¿Las arqueas son unicelulares o pluricelulares?
\begin{enumerate}[label=\textbf{\Alph*)}]
\begin{multicols}{2}
  \item Unicelulares
  \item Pluricelulares
  \item Coloniales
  \item Filamentosas
  \item Ambas
  \item Ninguna
\end{multicols}
\end{enumerate}
\vspace{0.15cm}
\item ¿Los protozoos son autótrofos o heterótrofos?
\begin{enumerate}[label=\textbf{\Alph*)}]
\begin{multicols}{2}
  \item Heterótrofos
  \item Autótrofos
  \item Mixótrofos
  \item Fotosintéticos
  \item Quimiosintéticos
  \item Saprófitos
\end{multicols}
\end{enumerate}
\vspace{0.15cm}
\item ¿Las algas son autótrofas o heterótrofas?
\begin{enumerate}[label=\textbf{\Alph*)}]
\begin{multicols}{2}
  \item Autótrofas
  \item Heterótrofas
  \item Mixótrofas
  \item Saprófitas
  \item Parásitas
  \item Simbiontes
\end{multicols}
\end{enumerate}
\vspace{0.15cm}
\item ¿Los hongos son autótrofos o heterótrofos?
\begin{enumerate}[label=\textbf{\Alph*)}]
\begin{multicols}{2}
  \item Heterótrofos
  \item Autótrofos
  \item Mixótrofos
  \item Fotosintéticos
  \item Quimiosintéticos
  \item Saprófitos
\end{multicols}
\end{enumerate}
\vspace{0.15cm}
\item ¿Las plantas son autótrofas o heterótrofas?
\begin{enumerate}[label=\textbf{\Alph*)}]
\begin{multicols}{2}
  \item Autótrofas
  \item Heterótrofas
  \item Mixótrofas
  \item Saprófitas
  \item Parásitas
  \item Simbiontes
\end{multicols}
\end{enumerate}
\vspace{0.15cm}
\item ¿Los animales son autótrofos o heterótrofos?
\begin{enumerate}[label=\textbf{\Alph*)}]
\begin{multicols}{2}
  \item Heterótrofos
  \item Autótrofos
  \item Mixótrofos
  \item Fotosintéticos
  \item Quimiosintéticos
  \item Saprófitos
\end{multicols}
\end{enumerate}
\vspace{0.15cm}
\item ¿Qué tienen las células vegetales que no tienen las animales?
\begin{enumerate}[label=\textbf{\Alph*)}]
\begin{multicols}{2}
  \item Pared celular de celulosa y cloroplastos
  \item Pared celular de quitina
  \item Centriolos
  \item Lisosomas
  \item Flagelos
  \item Cilios
\end{multicols}
\end{enumerate}
\vspace{0.15cm}
\item ¿Qué tienen las células fúngicas que no tienen las animales?
\begin{enumerate}[label=\textbf{\Alph*)}]
\begin{multicols}{2}
  \item Pared celular de quitina
  \item Pared celular de celulosa
  \item Cloroplastos
  \item Centriolos
  \item Lisosomas
  \item Flagelos
\end{multicols}
\end{enumerate}
\vspace{0.15cm}
\item ¿Qué tipo de nutrición tienen las bacterias?
\begin{enumerate}[label=\textbf{\Alph*)}]
\begin{multicols}{2}
  \item Autótrofa o heterótrofa
  \item Solo autótrofa
  \item Solo heterótrofa
  \item Solo fotosintética
  \item Solo quimiosintética
  \item Solo saprófita
\end{multicols}
\end{enumerate}
\vspace{0.15cm}
\item ¿Qué tipo de reproducción tienen las bacterias?
\begin{enumerate}[label=\textbf{\Alph*)}]
\begin{multicols}{2}
  \item Asexual por fisión binaria
  \item Sexual por gametos
  \item Asexual por esporas
  \item Sexual por conjugación
  \item Asexual por gemación
  \item Sexual por meiosis
\end{multicols}
\end{enumerate}
\vspace{0.15cm}
\item ¿Qué tipo de reproducción tienen las arqueas?
\begin{enumerate}[label=\textbf{\Alph*)}]
\begin{multicols}{2}
  \item Asexual por fisión binaria
  \item Sexual por gametos
  \item Asexual por esporas
  \item Sexual por conjugación
  \item Asexual por gemación
  \item Sexual por meiosis
\end{multicols}
\end{enumerate}
\vspace{0.15cm}
\item ¿Dónde viven las arqueas extremófilas?
\begin{enumerate}[label=\textbf{\Alph*)}]
\begin{multicols}{2}
  \item En ambientes extremos
  \item En ambientes moderados
  \item En el cuerpo humano
  \item En las plantas
  \item En los animales
  \item En el suelo
\end{multicols}
\end{enumerate}
\vspace{0.15cm}
\item ¿Qué son los protozoos?
\begin{enumerate}[label=\textbf{\Alph*)}]
\begin{multicols}{2}
  \item Eucariotas unicelulares heterótrofos
  \item Procariotas unicelulares
  \item Eucariotas pluricelulares
  \item Hongos unicelulares
  \item Algas unicelulares
  \item Bacterias unicelulares
\end{multicols}
\end{enumerate}
\vspace{0.15cm}
\item ¿Qué son las algas?
\begin{enumerate}[label=\textbf{\Alph*)}]
\begin{multicols}{2}
  \item Eucariotas autótrofas, unicelulares o pluricelulares
  \item Procariotas autótrofas
  \item Eucariotas heterótrofas
  \item Hongos autótrofos
  \item Protozoos autótrofos
  \item Bacterias autótrofas
\end{multicols}
\end{enumerate}
\vspace{0.15cm}
\item ¿Qué son los hongos?
\begin{enumerate}[label=\textbf{\Alph*)}]
\begin{multicols}{2}
  \item Eucariotas heterótrofos con pared de quitina
  \item Eucariotas autótrofos con pared de celulosa
  \item Procariotas heterótrofos
  \item Eucariotas heterótrofos sin pared
  \item Eucariotas autótrofos sin pared
  \item Procariotas autótrofos
\end{multicols}
\end{enumerate}
\vspace{0.15cm}
\item ¿Qué son las levaduras?
\begin{enumerate}[label=\textbf{\Alph*)}]
\begin{multicols}{2}
  \item Hongos unicelulares
  \item Hongos pluricelulares
  \item Algas unicelulares
  \item Protozoos unicelulares
  \item Bacterias unicelulares
  \item Arqueas unicelulares
\end{multicols}
\end{enumerate}
\vspace{0.15cm}
\item ¿Qué son los mohos?
\begin{enumerate}[label=\textbf{\Alph*)}]
\begin{multicols}{2}
  \item Hongos pluricelulares filamentosos
  \item Hongos unicelulares
  \item Algas pluricelulares
  \item Protozoos pluricelulares
  \item Bacterias pluricelulares
  \item Arqueas pluricelulares
\end{multicols}
\end{enumerate}
\vspace{0.15cm}
\item ¿Qué son las briofitas?
\begin{enumerate}[label=\textbf{\Alph*)}]
\begin{multicols}{2}
  \item Plantas no vasculares
  \item Plantas vasculares sin semilla
  \item Plantas con semilla desnuda
  \item Plantas con flor
  \item Algas pluricelulares
  \item Hongos pluricelulares
\end{multicols}
\end{enumerate}
\vspace{0.15cm}
\item ¿Qué son las pteridofitas?
\begin{enumerate}[label=\textbf{\Alph*)}]
\begin{multicols}{2}
  \item Plantas vasculares sin semilla
  \item Plantas no vasculares
  \item Plantas con semilla desnuda
  \item Plantas con flor
  \item Algas pluricelulares
  \item Hongos pluricelulares
\end{multicols}
\end{enumerate}
\vspace{0.15cm}
\item ¿Qué son las gimnospermas?
\begin{enumerate}[label=\textbf{\Alph*)}]
\begin{multicols}{2}
  \item Plantas con semilla desnuda
  \item Plantas sin semilla
  \item Plantas con flor
  \item Plantas no vasculares
  \item Algas pluricelulares
  \item Hongos pluricelulares
\end{multicols}
\end{enumerate}
\vspace{0.15cm}
\item ¿Qué son las angiospermas?
\begin{enumerate}[label=\textbf{\Alph*)}]
\begin{multicols}{2}
  \item Plantas con flor y semilla en fruto
  \item Plantas sin flor
  \item Plantas sin semilla
  \item Plantas no vasculares
  \item Algas pluricelulares
  \item Hongos pluricelulares
\end{multicols}
\end{enumerate}
\vspace{0.15cm}
\item ¿Qué son los invertebrados?
\begin{enumerate}[label=\textbf{\Alph*)}]
\begin{multicols}{2}
  \item Animales sin columna vertebral
  \item Animales con columna vertebral
  \item Animales sin células
  \item Animales autótrofos
  \item Animales unicelulares
  \item Animales procariotas
\end{multicols}
\end{enumerate}
\vspace{0.15cm}
\item ¿Qué son los vertebrados?
\begin{enumerate}[label=\textbf{\Alph*)}]
\begin{multicols}{2}
  \item Animales con columna vertebral
  \item Animales sin columna vertebral
  \item Animales sin células
  \item Animales autótrofos
  \item Animales unicelulares
  \item Animales procariotas
\end{multicols}
\end{enumerate}
\vspace{0.15cm}
\item ¿Qué son los poríferos?
\begin{enumerate}[label=\textbf{\Alph*)}]
\begin{multicols}{2}
  \item Esponjas
  \item Medusas
  \item Gusanos planos
  \item Gusanos cilíndricos
  \item Moluscos
  \item Artrópodos
\end{multicols}
\end{enumerate}
\vspace{0.15cm}
\item ¿Qué son los cnidarios?
\begin{enumerate}[label=\textbf{\Alph*)}]
\begin{multicols}{2}
  \item Medusas, corales y anémonas
  \item Esponjas
  \item Gusanos planos
  \item Gusanos cilíndricos
  \item Moluscos
  \item Artrópodos
\end{multicols}
\end{enumerate}
\vspace{0.15cm}
\item ¿Qué son los platelmintos?
\begin{enumerate}[label=\textbf{\Alph*)}]
\begin{multicols}{2}
  \item Gusanos planos
  \item Gusanos cilíndricos
  \item Gusanos segmentados
  \item Moluscos
  \item Artrópodos
  \item Equinodermos
\end{multicols}
\end{enumerate}
\vspace{0.15cm}
\item ¿Qué son los nematodos?
\begin{enumerate}[label=\textbf{\Alph*)}]
\begin{multicols}{2}
  \item Gusanos cilíndricos
  \item Gusanos planos
  \item Gusanos segmentados
  \item Moluscos
  \item Artrópodos
  \item Equinodermos
\end{multicols}
\end{enumerate}
\vspace{0.15cm}
\item ¿Qué son los anélidos?
\begin{enumerate}[label=\textbf{\Alph*)}]
\begin{multicols}{2}
  \item Gusanos segmentados
  \item Gusanos planos
  \item Gusanos cilíndricos
  \item Moluscos
  \item Artrópodos
  \item Equinodermos
\end{multicols}
\end{enumerate}
\vspace{0.15cm}
\item ¿Qué son los moluscos?
\begin{enumerate}[label=\textbf{\Alph*)}]
\begin{multicols}{2}
  \item Animales de cuerpo blando, muchos con concha
  \item Animales de cuerpo duro
  \item Animales segmentados
  \item Animales con exoesqueleto
  \item Animales con espinas
  \item Animales con columna
\end{multicols}
\end{enumerate}
\vspace{0.15cm}
\item ¿Qué son los artrópodos?
\begin{enumerate}[label=\textbf{\Alph*)}]
\begin{multicols}{2}
  \item Animales con exoesqueleto y patas articuladas
  \item Animales de cuerpo blando
  \item Animales segmentados sin patas
  \item Animales con espinas
  \item Animales con columna
  \item Animales con concha
\end{multicols}
\end{enumerate}
\vspace{0.15cm}
\item ¿Qué son los equinodermos?
\begin{enumerate}[label=\textbf{\Alph*)}]
\begin{multicols}{2}
  \item Animales marinos con simetría radial y espinas
  \item Animales de cuerpo blando
  \item Animales segmentados
  \item Animales con exoesqueleto
  \item Animales con columna
  \item Animales con concha
\end{multicols}
\end{enumerate}
\vspace{0.15cm}
\item ¿Cuáles son los grupos de vertebrados?
\begin{enumerate}[label=\textbf{\Alph*)}]
\begin{multicols}{2}
  \item Peces, anfibios, reptiles, aves y mamíferos
  \item Peces, insectos, reptiles, aves y mamíferos
  \item Anfibios, reptiles, aves, mamíferos y crustáceos
  \item Peces, anfibios, reptiles, aves y arácnidos
  \item Peces, anfibios, reptiles, aves y moluscos
  \item Peces, anfibios, reptiles, aves y anélidos
\end{multicols}
\end{enumerate}
\vspace{0.15cm}
\item ¿Cuáles son los grupos de invertebrados?
\begin{enumerate}[label=\textbf{\Alph*)}]
\begin{multicols}{2}
  \item Poríferos, cnidarios, platelmintos, nematodos, anélidos, moluscos, artrópodos y equinodermos
  \item Peces, anfibios, reptiles, aves y mamíferos
  \item Solo artrópodos y moluscos
  \item Solo poríferos y cnidarios
  \item Solo platelmintos y nematodos
  \item Solo anélidos y equinodermos ---
\end{multicols}
\end{enumerate}
\vspace{0.15cm}
\end{enumerate}\n
\subsubsection*{SECCIÓN IV: CINCO REINOS – DETALLE (151–200)}

\begin{enumerate}[label=\textbf{\arabic*.}]
\setcounter{enumi}{150}
\item ¿El reino Monera es procariota o eucariota?
\begin{enumerate}[label=\textbf{\Alph*)}]
\begin{multicols}{2}
  \item Procariota
  \item Eucariota
  \item Ambos
  \item Ninguno
  \item Solo eucariota
  \item Solo procariota
\end{multicols}
\end{enumerate}
\vspace{0.15cm}
\item ¿El reino Monera es unicelular o pluricelular?
\begin{enumerate}[label=\textbf{\Alph*)}]
\begin{multicols}{2}
  \item Unicelular
  \item Pluricelular
  \item Colonial
  \item Filamentoso
  \item Ambos
  \item Ninguno
\end{multicols}
\end{enumerate}
\vspace{0.15cm}
\item ¿La nutrición en Monera puede ser?
\begin{enumerate}[label=\textbf{\Alph*)}]
\begin{multicols}{2}
  \item Autótrofa o heterótrofa
  \item Solo autótrofa
  \item Solo heterótrofa
  \item Solo fotosintética
  \item Solo quimiosintética
  \item Solo saprófita
\end{multicols}
\end{enumerate}
\vspace{0.15cm}
\item ¿El reino Protista es procariota o eucariota?
\begin{enumerate}[label=\textbf{\Alph*)}]
\begin{multicols}{2}
  \item Eucariota
  \item Procariota
  \item Ambos
  \item Ninguno
  \item Solo procariota
  \item Solo eucariota
\end{multicols}
\end{enumerate}
\vspace{0.15cm}
\item ¿El reino Protista es unicelular o pluricelular?
\begin{enumerate}[label=\textbf{\Alph*)}]
\begin{multicols}{2}
  \item Unicelular o colonial, algunos pluricelulares simples
  \item Solo pluricelular
  \item Solo unicelular
  \item Solo colonial
  \item Solo filamentoso
  \item Solo macroscópico
\end{multicols}
\end{enumerate}
\vspace{0.15cm}
\item ¿La nutrición en Protista puede ser?
\begin{enumerate}[label=\textbf{\Alph*)}]
\begin{multicols}{2}
  \item Autótrofa, heterótrofa o mixótrofa
  \item Solo autótrofa
  \item Solo heterótrofa
  \item Solo mixótrofa
  \item Solo fotosintética
  \item Solo saprófita
\end{multicols}
\end{enumerate}
\vspace{0.15cm}
\item ¿El reino Fungi es procariota o eucariota?
\begin{enumerate}[label=\textbf{\Alph*)}]
\begin{multicols}{2}
  \item Eucariota
  \item Procariota
  \item Ambos
  \item Ninguno
  \item Solo procariota
  \item Solo eucariota
\end{multicols}
\end{enumerate}
\vspace{0.15cm}
\item ¿La nutrición de Fungi es?
\begin{enumerate}[label=\textbf{\Alph*)}]
\begin{multicols}{2}
  \item Heterótrofa
  \item Autótrofa
  \item Mixótrofa
  \item Fotosintética
  \item Quimiosintética
  \item Saprófita
\end{multicols}
\end{enumerate}
\vspace{0.15cm}
\item ¿La pared celular de Fungi está formada por?
\begin{enumerate}[label=\textbf{\Alph*)}]
\begin{multicols}{2}
  \item Quitina
  \item Celulosa
  \item Peptidoglicano
  \item Lignina
  \item Pectina
  \item Silicato
\end{multicols}
\end{enumerate}
\vspace{0.15cm}
\item ¿El reino Plantae es procariota o eucariota?
\begin{enumerate}[label=\textbf{\Alph*)}]
\begin{multicols}{2}
  \item Eucariota
  \item Procariota
  \item Ambos
  \item Ninguno
  \item Solo procariota
  \item Solo eucariota
\end{multicols}
\end{enumerate}
\vspace{0.15cm}
\item ¿La nutrición de Plantae es?
\begin{enumerate}[label=\textbf{\Alph*)}]
\begin{multicols}{2}
  \item Autótrofa
  \item Heterótrofa
  \item Mixótrofa
  \item Saprófita
  \item Parásita
  \item Simbionte
\end{multicols}
\end{enumerate}
\vspace{0.15cm}
\item ¿La pared celular de Plantae está formada por?
\begin{enumerate}[label=\textbf{\Alph*)}]
\begin{multicols}{2}
  \item Celulosa
  \item Quitina
  \item Peptidoglicano
  \item Lignina
  \item Pectina
  \item Silicato
\end{multicols}
\end{enumerate}
\vspace{0.15cm}
\item ¿El reino Animalia es procariota o eucariota?
\begin{enumerate}[label=\textbf{\Alph*)}]
\begin{multicols}{2}
  \item Eucariota
  \item Procariota
  \item Ambos
  \item Ninguno
  \item Solo procariota
  \item Solo eucariota
\end{multicols}
\end{enumerate}
\vspace{0.15cm}
\item ¿La nutrición de Animalia es?
\begin{enumerate}[label=\textbf{\Alph*)}]
\begin{multicols}{2}
  \item Heterótrofa
  \item Autótrofa
  \item Mixótrofa
  \item Fotosintética
  \item Quimiosintética
  \item Saprófita
\end{multicols}
\end{enumerate}
\vspace{0.15cm}
\item ¿Las células animales tienen pared celular?
\begin{enumerate}[label=\textbf{\Alph*)}]
\begin{multicols}{2}
  \item No
  \item Sí, de celulosa
  \item Sí, de quitina
  \item Sí, de peptidoglicano
  \item Sí, de lignina
  \item Sí, de pectina
\end{multicols}
\end{enumerate}
\vspace{0.15cm}
\item ¿Cuáles son las clases de vertebrados?
\begin{enumerate}[label=\textbf{\Alph*)}]
\begin{multicols}{2}
  \item Mamíferos, aves, reptiles, anfibios y peces
  \item Mamíferos, aves, reptiles, anfibios e insectos
  \item Mamíferos, aves, reptiles, anfibios y crustáceos
  \item Mamíferos, aves, reptiles, anfibios y arácnidos
  \item Mamíferos, aves, reptiles, anfibios y moluscos
  \item Mamíferos, aves, reptiles, anfibios y anélidos
\end{multicols}
\end{enumerate}
\vspace{0.15cm}
\item ¿Qué caracteriza a los mamíferos?
\begin{enumerate}[label=\textbf{\Alph*)}]
\begin{multicols}{2}
  \item Tienen glándulas mamarias y pelo
  \item Tienen plumas
  \item Tienen escamas
  \item Tienen branquias
  \item Tienen aletas
  \item Tienen exoesqueleto
\end{multicols}
\end{enumerate}
\vspace{0.15cm}
\item ¿Qué caracteriza a las aves?
\begin{enumerate}[label=\textbf{\Alph*)}]
\begin{multicols}{2}
  \item Tienen plumas y pico
  \item Tienen pelo
  \item Tienen escamas
  \item Tienen branquias
  \item Tienen aletas
  \item Tienen exoesqueleto
\end{multicols}
\end{enumerate}
\vspace{0.15cm}
\item ¿Qué caracteriza a los reptiles?
\begin{enumerate}[label=\textbf{\Alph*)}]
\begin{multicols}{2}
  \item Tienen escamas y son de sangre fría
  \item Tienen plumas
  \item Tienen pelo
  \item Tienen branquias
  \item Tienen aletas
  \item Tienen exoesqueleto
\end{multicols}
\end{enumerate}
\vspace{0.15cm}
\item ¿Qué caracteriza a los anfibios?
\begin{enumerate}[label=\textbf{\Alph*)}]
\begin{multicols}{2}
  \item Dependen del agua para reproducirse
  \item Tienen plumas
  \item Tienen pelo
  \item Tienen escamas
  \item Tienen aletas
  \item Tienen exoesqueleto
\end{multicols}
\end{enumerate}
\vspace{0.15cm}
\item ¿Qué caracteriza a los peces?
\begin{enumerate}[label=\textbf{\Alph*)}]
\begin{multicols}{2}
  \item Tienen branquias y aletas
  \item Tienen pulmones
  \item Tienen plumas
  \item Tienen pelo
  \item Tienen escamas
  \item Tienen exoesqueleto
\end{multicols}
\end{enumerate}
\vspace{0.15cm}
\item ¿Qué caracteriza a los invertebrados?
\begin{enumerate}[label=\textbf{\Alph*)}]
\begin{multicols}{2}
  \item No tienen columna vertebral
  \item Tienen columna vertebral
  \item Tienen pelo
  \item Tienen plumas
  \item Tienen escamas
  \item Tienen branquias
\end{multicols}
\end{enumerate}
\vspace{0.15cm}
\item ¿Cuáles son las clases de artrópodos?
\begin{enumerate}[label=\textbf{\Alph*)}]
\begin{multicols}{2}
  \item Insectos, arácnidos, crustáceos y miriápodos
  \item Insectos, arácnidos, crustáceos y moluscos
  \item Insectos, arácnidos, crustáceos y anélidos
  \item Insectos, arácnidos, crustáceos y equinodermos
  \item Insectos, arácnidos, crustáceos y cnidarios
  \item Insectos, arácnidos, crustáceos y poríferos
\end{multicols}
\end{enumerate}
\vspace{0.15cm}
\item ¿Qué caracteriza a los insectos?
\begin{enumerate}[label=\textbf{\Alph*)}]
\begin{multicols}{2}
  \item Tienen seis patas y tres pares de patas
  \item Tienen ocho patas
  \item Tienen diez patas
  \item Tienen cuatro patas
  \item Tienen dos patas
  \item Tienen muchas patas
\end{multicols}
\end{enumerate}
\vspace{0.15cm}
\item ¿Qué caracteriza a los arácnidos?
\begin{enumerate}[label=\textbf{\Alph*)}]
\begin{multicols}{2}
  \item Tienen ocho patas
  \item Tienen seis patas
  \item Tienen diez patas
  \item Tienen cuatro patas
  \item Tienen dos patas
  \item Tienen muchas patas
\end{multicols}
\end{enumerate}
\vspace{0.15cm}
\item ¿Qué caracteriza a los crustáceos?
\begin{enumerate}[label=\textbf{\Alph*)}]
\begin{multicols}{2}
  \item Tienen diez patas y antenas
  \item Tienen seis patas
  \item Tienen ocho patas
  \item Tienen cuatro patas
  \item Tienen dos patas
  \item Tienen muchas patas
\end{multicols}
\end{enumerate}
\vspace{0.15cm}
\item ¿Qué caracteriza a los miriápodos?
\begin{enumerate}[label=\textbf{\Alph*)}]
\begin{multicols}{2}
  \item Tienen muchas patas
  \item Tienen seis patas
  \item Tienen ocho patas
  \item Tienen diez patas
  \item Tienen cuatro patas
  \item Tienen dos patas
\end{multicols}
\end{enumerate}
\vspace{0.15cm}
\item ¿Cuáles son las clases de moluscos?
\begin{enumerate}[label=\textbf{\Alph*)}]
\begin{multicols}{2}
  \item Gasterópodos, bivalvos y cefalópodos
  \item Gasterópodos, bivalvos y crustáceos
  \item Gasterópodos, bivalvos e insectos
  \item Gasterópodos, bivalvos y arácnidos
  \item Gasterópodos, bivalvos y anélidos
  \item Gasterópodos, bivalvos y equinodermos
\end{multicols}
\end{enumerate}
\vspace{0.15cm}
\item ¿Qué son los gasterópodos?
\begin{enumerate}[label=\textbf{\Alph*)}]
\begin{multicols}{2}
  \item Caracoles y babosas
  \item Almejas y ostras
  \item Pulpos y calamares
  \item Cangrejos y langostas
  \item Arañas y escorpiones
  \item Lombrices y sanguijuelas
\end{multicols}
\end{enumerate}
\vspace{0.15cm}
\item ¿Qué son los bivalvos?
\begin{enumerate}[label=\textbf{\Alph*)}]
\begin{multicols}{2}
  \item Almejas, ostras y mejillones
  \item Caracoles y babosas
  \item Pulpos y calamares
  \item Cangrejos y langostas
  \item Arañas y escorpiones
  \item Lombrices y sanguijuelas
\end{multicols}
\end{enumerate}
\vspace{0.15cm}
\item ¿Qué son los cefalópodos?
\begin{enumerate}[label=\textbf{\Alph*)}]
\begin{multicols}{2}
  \item Pulpos, calamares y sepias
  \item Caracoles y babosas
  \item Almejas y ostras
  \item Cangrejos y langostas
  \item Arañas y escorpiones
  \item Lombrices y sanguijuelas
\end{multicols}
\end{enumerate}
\vspace{0.15cm}
\item ¿Cuáles son las clases de anélidos?
\begin{enumerate}[label=\textbf{\Alph*)}]
\begin{multicols}{2}
  \item Poliquetos, oligoquetos e hirudíneos
  \item Poliquetos, oligoquetos e insectos
  \item Poliquetos, oligoquetos y crustáceos
  \item Poliquetos, oligoquetos y arácnidos
  \item Poliquetos, oligoquetos y moluscos
  \item Poliquetos, oligoquetos y equinodermos
\end{multicols}
\end{enumerate}
\vspace{0.15cm}
\item ¿Qué son los poliquetos?
\begin{enumerate}[label=\textbf{\Alph*)}]
\begin{multicols}{2}
  \item Gusanos marinos con muchas cerdas
  \item Gusanos de tierra
  \item Sanguijuelas
  \item Caracoles
  \item Almejas
  \item Pulpos
\end{multicols}
\end{enumerate}
\vspace{0.15cm}
\item ¿Qué son los oligoquetos?
\begin{enumerate}[label=\textbf{\Alph*)}]
\begin{multicols}{2}
  \item Lombrices de tierra
  \item Gusanos marinos
  \item Sanguijuelas
  \item Caracoles
  \item Almejas
  \item Pulpos
\end{multicols}
\end{enumerate}
\vspace{0.15cm}
\item ¿Qué son los hirudíneos?
\begin{enumerate}[label=\textbf{\Alph*)}]
\begin{multicols}{2}
  \item Sanguijuelas
  \item Lombrices de tierra
  \item Gusanos marinos
  \item Caracoles
  \item Almejas
  \item Pulpos
\end{multicols}
\end{enumerate}
\vspace{0.15cm}
\item ¿Cuáles son las clases de equinodermos?
\begin{enumerate}[label=\textbf{\Alph*)}]
\begin{multicols}{2}
  \item Asteroideos, equinoideos y holoturoideos
  \item Asteroideos, equinoideos e insectos
  \item Asteroideos, equinoideos y crustáceos
  \item Asteroideos, equinoideos y arácnidos
  \item Asteroideos, equinoideos y moluscos
  \item Asteroideos, equinoideos y anélidos
\end{multicols}
\end{enumerate}
\vspace{0.15cm}
\item ¿Qué son los asteroideos?
\begin{enumerate}[label=\textbf{\Alph*)}]
\begin{multicols}{2}
  \item Estrellas de mar
  \item Erizos de mar
  \item Pepinos de mar
  \item Caracoles
  \item Almejas
  \item Pulpos
\end{multicols}
\end{enumerate}
\vspace{0.15cm}
\item ¿Qué son los equinoideos?
\begin{enumerate}[label=\textbf{\Alph*)}]
\begin{multicols}{2}
  \item Erizos de mar
  \item Estrellas de mar
  \item Pepinos de mar
  \item Caracoles
  \item Almejas
  \item Pulpos
\end{multicols}
\end{enumerate}
\vspace{0.15cm}
\item ¿Qué son los holoturoideos?
\begin{enumerate}[label=\textbf{\Alph*)}]
\begin{multicols}{2}
  \item Pepinos de mar
  \item Estrellas de mar
  \item Erizos de mar
  \item Caracoles
  \item Almejas
  \item Pulpos
\end{multicols}
\end{enumerate}
\vspace{0.15cm}
\item ¿Cuáles son las clases de cnidarios?
\begin{enumerate}[label=\textbf{\Alph*)}]
\begin{multicols}{2}
  \item Hidrozoos, escifozoos y antozoos
  \item Hidrozoos, escifozoos e insectos
  \item Hidrozoos, escifozoos y crustáceos
  \item Hidrozoos, escifozoos y arácnidos
  \item Hidrozoos, escifozoos y moluscos
  \item Hidrozoos, escifozoos y anélidos
\end{multicols}
\end{enumerate}
\vspace{0.15cm}
\item ¿Qué son los hidrozoos?
\begin{enumerate}[label=\textbf{\Alph*)}]
\begin{multicols}{2}
  \item Hidras y medusas pequeñas
  \item Corales
  \item Anémonas
  \item Estrellas de mar
  \item Erizos de mar
  \item Pepinos de mar
\end{multicols}
\end{enumerate}
\vspace{0.15cm}
\item ¿Qué son los escifozoos?
\begin{enumerate}[label=\textbf{\Alph*)}]
\begin{multicols}{2}
  \item Medusas grandes
  \item Hidras
  \item Corales
  \item Anémonas
  \item Estrellas de mar
  \item Erizos de mar
\end{multicols}
\end{enumerate}
\vspace{0.15cm}
\item ¿Qué son los antozoos?
\begin{enumerate}[label=\textbf{\Alph*)}]
\begin{multicols}{2}
  \item Corales y anémonas
  \item Hidras
  \item Medusas
  \item Estrellas de mar
  \item Erizos de mar
  \item Pepinos de mar
\end{multicols}
\end{enumerate}
\vspace{0.15cm}
\item ¿Qué son los poríferos?
\begin{enumerate}[label=\textbf{\Alph*)}]
\begin{multicols}{2}
  \item Esponjas
  \item Medusas
  \item Corales
  \item Anémonas
  \item Estrellas de mar
  \item Erizos de mar
\end{multicols}
\end{enumerate}
\vspace{0.15cm}
\item ¿Qué son los platelmintos?
\begin{enumerate}[label=\textbf{\Alph*)}]
\begin{multicols}{2}
  \item Gusanos planos
  \item Gusanos cilíndricos
  \item Gusanos segmentados
  \item Moluscos
  \item Artrópodos
  \item Equinodermos
\end{multicols}
\end{enumerate}
\vspace{0.15cm}
\item ¿Qué son los nematodos?
\begin{enumerate}[label=\textbf{\Alph*)}]
\begin{multicols}{2}
  \item Gusanos cilíndricos
  \item Gusanos planos
  \item Gusanos segmentados
  \item Moluscos
  \item Artrópodos
  \item Equinodermos
\end{multicols}
\end{enumerate}
\vspace{0.15cm}
\item ¿Qué son las briofitas?
\begin{enumerate}[label=\textbf{\Alph*)}]
\begin{multicols}{2}
  \item Musgos y hepáticas
  \item Helechos
  \item Pinos
  \item Rosas
  \item Algas
  \item Hongos
\end{multicols}
\end{enumerate}
\vspace{0.15cm}
\item ¿Qué son las pteridofitas?
\begin{enumerate}[label=\textbf{\Alph*)}]
\begin{multicols}{2}
  \item Helechos
  \item Musgos
  \item Pinos
  \item Rosas
  \item Algas
  \item Hongos
\end{multicols}
\end{enumerate}
\vspace{0.15cm}
\item ¿Qué son las gimnospermas?
\begin{enumerate}[label=\textbf{\Alph*)}]
\begin{multicols}{2}
  \item Pinos, abetos y cipreses
  \item Musgos
  \item Helechos
  \item Rosas
  \item Algas
  \item Hongos
\end{multicols}
\end{enumerate}
\vspace{0.15cm}
\item ¿Qué son las angiospermas?
\begin{enumerate}[label=\textbf{\Alph*)}]
\begin{multicols}{2}
  \item Plantas con flores y frutos
  \item Pinos
  \item Musgos
  \item Helechos
  \item Algas
  \item Hongos ---
\end{multicols}
\end{enumerate}
\vspace{0.15cm}
\end{enumerate}\n
\subsubsection*{SECCIÓN V: CLASIFICACIÓN Y EJEMPLOS (201–250)}

\begin{enumerate}[label=\textbf{\arabic*.}]
\setcounter{enumi}{200}
\item ¿A qué reino pertenece el ser humano?
\begin{enumerate}[label=\textbf{\Alph*)}]
\begin{multicols}{2}
  \item Animalia
  \item Plantae
  \item Fungi
  \item Protista
  \item Monera
  \item Archaea
\end{multicols}
\end{enumerate}
\vspace{0.15cm}
\item ¿A qué reino pertenece el perro?
\begin{enumerate}[label=\textbf{\Alph*)}]
\begin{multicols}{2}
  \item Animalia
  \item Plantae
  \item Fungi
  \item Protista
  \item Monera
  \item Archaea
\end{multicols}
\end{enumerate}
\vspace{0.15cm}
\item ¿A qué reino pertenece el gato?
\begin{enumerate}[label=\textbf{\Alph*)}]
\begin{multicols}{2}
  \item Animalia
  \item Plantae
  \item Fungi
  \item Protista
  \item Monera
  \item Archaea
\end{multicols}
\end{enumerate}
\vspace{0.15cm}
\item ¿A qué reino pertenece el maíz?
\begin{enumerate}[label=\textbf{\Alph*)}]
\begin{multicols}{2}
  \item Plantae
  \item Animalia
  \item Fungi
  \item Protista
  \item Monera
  \item Archaea
\end{multicols}
\end{enumerate}
\vspace{0.15cm}
\item ¿A qué reino pertenece el trigo?
\begin{enumerate}[label=\textbf{\Alph*)}]
\begin{multicols}{2}
  \item Plantae
  \item Animalia
  \item Fungi
  \item Protista
  \item Monera
  \item Archaea
\end{multicols}
\end{enumerate}
\vspace{0.15cm}
\item ¿A qué reino pertenece el arroz?
\begin{enumerate}[label=\textbf{\Alph*)}]
\begin{multicols}{2}
  \item Plantae
  \item Animalia
  \item Fungi
  \item Protista
  \item Monera
  \item Archaea
\end{multicols}
\end{enumerate}
\vspace{0.15cm}
\item ¿A qué reino pertenece el león?
\begin{enumerate}[label=\textbf{\Alph*)}]
\begin{multicols}{2}
  \item Animalia
  \item Plantae
  \item Fungi
  \item Protista
  \item Monera
  \item Archaea
\end{multicols}
\end{enumerate}
\vspace{0.15cm}
\item ¿A qué reino pertenece el tigre?
\begin{enumerate}[label=\textbf{\Alph*)}]
\begin{multicols}{2}
  \item Animalia
  \item Plantae
  \item Fungi
  \item Protista
  \item Monera
  \item Archaea
\end{multicols}
\end{enumerate}
\vspace{0.15cm}
\item ¿A qué reino pertenece el elefante?
\begin{enumerate}[label=\textbf{\Alph*)}]
\begin{multicols}{2}
  \item Animalia
  \item Plantae
  \item Fungi
  \item Protista
  \item Monera
  \item Archaea
\end{multicols}
\end{enumerate}
\vspace{0.15cm}
\item ¿A qué reino pertenece la ballena?
\begin{enumerate}[label=\textbf{\Alph*)}]
\begin{multicols}{2}
  \item Animalia
  \item Plantae
  \item Fungi
  \item Protista
  \item Monera
  \item Archaea
\end{multicols}
\end{enumerate}
\vspace{0.15cm}
\item ¿A qué reino pertenece el murciélago?
\begin{enumerate}[label=\textbf{\Alph*)}]
\begin{multicols}{2}
  \item Animalia
  \item Plantae
  \item Fungi
  \item Protista
  \item Monera
  \item Archaea
\end{multicols}
\end{enumerate}
\vspace{0.15cm}
\item ¿A qué clase pertenece el águila?
\begin{enumerate}[label=\textbf{\Alph*)}]
\begin{multicols}{2}
  \item Aves
  \item Mamíferos
  \item Reptiles
  \item Anfibios
  \item Peces
  \item Insectos
\end{multicols}
\end{enumerate}
\vspace{0.15cm}
\item ¿A qué clase pertenece la serpiente?
\begin{enumerate}[label=\textbf{\Alph*)}]
\begin{multicols}{2}
  \item Reptiles
  \item Mamíferos
  \item Aves
  \item Anfibios
  \item Peces
  \item Insectos
\end{multicols}
\end{enumerate}
\vspace{0.15cm}
\item ¿A qué clase pertenece la rana?
\begin{enumerate}[label=\textbf{\Alph*)}]
\begin{multicols}{2}
  \item Anfibios
  \item Mamíferos
  \item Aves
  \item Reptiles
  \item Peces
  \item Insectos
\end{multicols}
\end{enumerate}
\vspace{0.15cm}
\item ¿A qué clase pertenece el pez?
\begin{enumerate}[label=\textbf{\Alph*)}]
\begin{multicols}{2}
  \item Peces
  \item Mamíferos
  \item Aves
  \item Reptiles
  \item Anfibios
  \item Insectos
\end{multicols}
\end{enumerate}
\vspace{0.15cm}
\item ¿A qué clase pertenece la abeja?
\begin{enumerate}[label=\textbf{\Alph*)}]
\begin{multicols}{2}
  \item Insecta
  \item Arachnida
  \item Crustacea
  \item Diplopoda
  \item Chilopoda
  \item Merostomata
\end{multicols}
\end{enumerate}
\vspace{0.15cm}
\item ¿A qué clase pertenece la araña?
\begin{enumerate}[label=\textbf{\Alph*)}]
\begin{multicols}{2}
  \item Arachnida
  \item Insecta
  \item Crustacea
  \item Diplopoda
  \item Chilopoda
  \item Merostomata
\end{multicols}
\end{enumerate}
\vspace{0.15cm}
\item ¿A qué clase pertenece el cangrejo?
\begin{enumerate}[label=\textbf{\Alph*)}]
\begin{multicols}{2}
  \item Crustacea
  \item Insecta
  \item Arachnida
  \item Diplopoda
  \item Chilopoda
  \item Merostomata
\end{multicols}
\end{enumerate}
\vspace{0.15cm}
\item ¿A qué clase pertenece el caracol?
\begin{enumerate}[label=\textbf{\Alph*)}]
\begin{multicols}{2}
  \item Gasterópoda
  \item Bivalvia
  \item Cephalopoda
  \item Insecta
  \item Arachnida
  \item Crustacea
\end{multicols}
\end{enumerate}
\vspace{0.15cm}
\item ¿A qué clase pertenece la lombriz de tierra?
\begin{enumerate}[label=\textbf{\Alph*)}]
\begin{multicols}{2}
  \item Oligochaeta
  \item Polychaeta
  \item Hirudinea
  \item Insecta
  \item Arachnida
  \item Crustacea
\end{multicols}
\end{enumerate}
\vspace{0.15cm}
\item ¿A qué clase pertenece la estrella de mar?
\begin{enumerate}[label=\textbf{\Alph*)}]
\begin{multicols}{2}
  \item Asteroidea
  \item Echinoidea
  \item Holothuroidea
  \item Insecta
  \item Arachnida
  \item Crustacea
\end{multicols}
\end{enumerate}
\vspace{0.15cm}
\item ¿A qué clase pertenece la medusa?
\begin{enumerate}[label=\textbf{\Alph*)}]
\begin{multicols}{2}
  \item Scyphozoa
  \item Hydrozoa
  \item Anthozoa
  \item Insecta
  \item Arachnida
  \item Crustacea
\end{multicols}
\end{enumerate}
\vspace{0.15cm}
\item ¿A qué filo pertenece la esponja?
\begin{enumerate}[label=\textbf{\Alph*)}]
\begin{multicols}{2}
  \item Porifera
  \item Cnidaria
  \item Platyhelminthes
  \item Nematoda
  \item Annelida
  \item Mollusca
\end{multicols}
\end{enumerate}
\vspace{0.15cm}
\item ¿A qué filo pertenece la tenia?
\begin{enumerate}[label=\textbf{\Alph*)}]
\begin{multicols}{2}
  \item Platyhelminthes
  \item Nematoda
  \item Annelida
  \item Mollusca
  \item Arthropoda
  \item Echinodermata
\end{multicols}
\end{enumerate}
\vspace{0.15cm}
\item ¿A qué filo pertenece el áscaris?
\begin{enumerate}[label=\textbf{\Alph*)}]
\begin{multicols}{2}
  \item Nematoda
  \item Platyhelminthes
  \item Annelida
  \item Mollusca
  \item Arthropoda
  \item Echinodermata
\end{multicols}
\end{enumerate}
\vspace{0.15cm}
\item ¿A qué reino pertenece el champiñón?
\begin{enumerate}[label=\textbf{\Alph*)}]
\begin{multicols}{2}
  \item Fungi
  \item Plantae
  \item Animalia
  \item Protista
  \item Monera
  \item Archaea
\end{multicols}
\end{enumerate}
\vspace{0.15cm}
\item ¿A qué reino pertenece la levadura?
\begin{enumerate}[label=\textbf{\Alph*)}]
\begin{multicols}{2}
  \item Fungi
  \item Plantae
  \item Animalia
  \item Protista
  \item Monera
  \item Archaea
\end{multicols}
\end{enumerate}
\vspace{0.15cm}
\item ¿A qué dominio pertenece *Escherichia coli*?
\begin{enumerate}[label=\textbf{\Alph*)}]
\begin{multicols}{2}
  \item Bacteria
  \item Archaea
  \item Eukarya
  \item Monera
  \item Protista
  \item Fungi
\end{multicols}
\end{enumerate}
\vspace{0.15cm}
\item ¿A qué dominio pertenece una arquea?
\begin{enumerate}[label=\textbf{\Alph*)}]
\begin{multicols}{2}
  \item Archaea
  \item Bacteria
  \item Eukarya
  \item Monera
  \item Protista
  \item Fungi
\end{multicols}
\end{enumerate}
\vspace{0.15cm}
\item ¿A qué reino pertenece un alga?
\begin{enumerate}[label=\textbf{\Alph*)}]
\begin{multicols}{2}
  \item Protista
  \item Monera
  \item Fungi
  \item Plantae
  \item Animalia
  \item Archaea
\end{multicols}
\end{enumerate}
\vspace{0.15cm}
\item ¿A qué reino pertenece un protozoo?
\begin{enumerate}[label=\textbf{\Alph*)}]
\begin{multicols}{2}
  \item Protista
  \item Monera
  \item Fungi
  \item Plantae
  \item Animalia
  \item Archaea
\end{multicols}
\end{enumerate}
\vspace{0.15cm}
\item ¿A qué grupo pertenece el musgo?
\begin{enumerate}[label=\textbf{\Alph*)}]
\begin{multicols}{2}
  \item Briofitas
  \item Pteridofitas
  \item Gimnospermas
  \item Angiospermas
  \item Algas
  \item Hongos
\end{multicols}
\end{enumerate}
\vspace{0.15cm}
\item ¿A qué grupo pertenece el helecho?
\begin{enumerate}[label=\textbf{\Alph*)}]
\begin{multicols}{2}
  \item Pteridofitas
  \item Briofitas
  \item Gimnospermas
  \item Angiospermas
  \item Algas
  \item Hongos
\end{multicols}
\end{enumerate}
\vspace{0.15cm}
\item ¿A qué grupo pertenece el pino?
\begin{enumerate}[label=\textbf{\Alph*)}]
\begin{multicols}{2}
  \item Gimnospermas
  \item Angiospermas
  \item Briofitas
  \item Pteridofitas
  \item Algas
  \item Hongos
\end{multicols}
\end{enumerate}
\vspace{0.15cm}
\item ¿A qué grupo pertenece la rosa?
\begin{enumerate}[label=\textbf{\Alph*)}]
\begin{multicols}{2}
  \item Angiospermas
  \item Gimnospermas
  \item Briofitas
  \item Pteridofitas
  \item Algas
  \item Hongos
\end{multicols}
\end{enumerate}
\vspace{0.15cm}
\item ¿A qué clase pertenece el ser humano?
\begin{enumerate}[label=\textbf{\Alph*)}]
\begin{multicols}{2}
  \item Mammalia
  \item Aves
  \item Reptilia
  \item Amphibia
  \item Pisces
  \item Insecta
\end{multicols}
\end{enumerate}
\vspace{0.15cm}
\item ¿A qué clase pertenece el perro?
\begin{enumerate}[label=\textbf{\Alph*)}]
\begin{multicols}{2}
  \item Mammalia
  \item Aves
  \item Reptilia
  \item Amphibia
  \item Pisces
  \item Insecta
\end{multicols}
\end{enumerate}
\vspace{0.15cm}
\item ¿A qué clase pertenece el cocodrilo?
\begin{enumerate}[label=\textbf{\Alph*)}]
\begin{multicols}{2}
  \item Reptilia
  \item Mammalia
  \item Aves
  \item Amphibia
  \item Pisces
  \item Insecta
\end{multicols}
\end{enumerate}
\vspace{0.15cm}
\item ¿A qué clase pertenece la salamandra?
\begin{enumerate}[label=\textbf{\Alph*)}]
\begin{multicols}{2}
  \item Amphibia
  \item Mammalia
  \item Aves
  \item Reptilia
  \item Pisces
  \item Insecta
\end{multicols}
\end{enumerate}
\vspace{0.15cm}
\item ¿A qué clase pertenece el tiburón?
\begin{enumerate}[label=\textbf{\Alph*)}]
\begin{multicols}{2}
  \item Chondrichthyes
  \item Osteichthyes
  \item Mammalia
  \item Aves
  \item Reptilia
  \item Amphibia
\end{multicols}
\end{enumerate}
\vspace{0.15cm}
\item ¿Para qué sirve una clave dicotómica?
\begin{enumerate}[label=\textbf{\Alph*)}]
\begin{multicols}{2}
  \item Para identificar organismos
  \item Para clasificar minerales
  \item Para medir organismos
  \item Para contar organismos
  \item Para dibujar organismos
  \item Para fotografiar organismos
\end{multicols}
\end{enumerate}
\vspace{0.15cm}
\item ¿Qué caracteres usa la taxonomía morfológica?
\begin{enumerate}[label=\textbf{\Alph*)}]
\begin{multicols}{2}
  \item Forma y estructura
  \item Color
  \item Tamaño
  \item Peso
  \item Olor
  \item Sabor
\end{multicols}
\end{enumerate}
\vspace{0.15cm}
\item ¿Qué caracteres usa la taxonomía anatómica?
\begin{enumerate}[label=\textbf{\Alph*)}]
\begin{multicols}{2}
  \item Estructura interna
  \item Color
  \item Tamaño
  \item Peso
  \item Olor
  \item Sabor
\end{multicols}
\end{enumerate}
\vspace{0.15cm}
\item ¿Qué caracteres usa la taxonomía fisiológica?
\begin{enumerate}[label=\textbf{\Alph*)}]
\begin{multicols}{2}
  \item Funciones vitales
  \item Color
  \item Tamaño
  \item Peso
  \item Olor
  \item Sabor
\end{multicols}
\end{enumerate}
\vspace{0.15cm}
\item ¿Qué caracteres usa la taxonomía molecular?
\begin{enumerate}[label=\textbf{\Alph*)}]
\begin{multicols}{2}
  \item ADN y proteínas
  \item Color
  \item Tamaño
  \item Peso
  \item Olor
  \item Sabor
\end{multicols}
\end{enumerate}
\vspace{0.15cm}
\item ¿Qué caracteres usa la taxonomía conductual?
\begin{enumerate}[label=\textbf{\Alph*)}]
\begin{multicols}{2}
  \item Comportamiento
  \item Color
  \item Tamaño
  \item Peso
  \item Olor
  \item Sabor
\end{multicols}
\end{enumerate}
\vspace{0.15cm}
\item ¿Qué caracteres usa la taxonomía biogeográfica?
\begin{enumerate}[label=\textbf{\Alph*)}]
\begin{multicols}{2}
  \item Distribución geográfica
  \item Color
  \item Tamaño
  \item Peso
  \item Olor
  \item Sabor
\end{multicols}
\end{enumerate}
\vspace{0.15cm}
\item ¿Qué es la filogenia molecular?
\begin{enumerate}[label=\textbf{\Alph*)}]
\begin{multicols}{2}
  \item El uso de secuencias de ADN para inferir relaciones evolutivas
  \item El uso del color
  \item El uso del tamaño
  \item El uso del peso
  \item El uso del olor
  \item El uso del sabor
\end{multicols}
\end{enumerate}
\vspace{0.15cm}
\item ¿Qué es el reloj molecular?
\begin{enumerate}[label=\textbf{\Alph*)}]
\begin{multicols}{2}
  \item La estimación del tiempo de divergencia usando mutaciones
  \item La medición del tiempo geológico
  \item La medición del tiempo de vida
  \item La medición del tiempo de gestación
  \item La medición del tiempo de crecimiento
  \item La medición del tiempo de reproducción
\end{multicols}
\end{enumerate}
\vspace{0.15cm}
\item ¿Qué es el código de barras de ADN?
\begin{enumerate}[label=\textbf{\Alph*)}]
\begin{multicols}{2}
  \item El uso de una secuencia corta de ADN para identificar especies
  \item El uso del color para identificar especies
  \item El uso del tamaño para identificar especies
  \item El uso del peso para identificar especies
  \item El uso del olor para identificar especies
  \item El uso del sabor para identificar especies ---
\end{multicols}
\end{enumerate}
\vspace{0.15cm}
\end{enumerate}\n
\subsubsection*{SECCIÓN VI: APLICACIONES Y PREGUNTAS MIXTAS (251–300)}

\begin{enumerate}[label=\textbf{\arabic*.}]
\setcounter{enumi}{250}
\item ¿Por qué es importante la taxonomía?
\begin{enumerate}[label=\textbf{\Alph*)}]
\begin{multicols}{2}
  \item Porque organiza y facilita el estudio de la biodiversidad
  \item Porque solo sirve para nombrar
  \item Porque solo sirve para coleccionar
  \item Porque solo sirve para dibujar
  \item Porque solo sirve para contar
  \item Porque solo sirve para medir
\end{multicols}
\end{enumerate}
\vspace{0.15cm}
\item ¿Cómo ayuda la taxonomía a la conservación?
\begin{enumerate}[label=\textbf{\Alph*)}]
\begin{multicols}{2}
  \item Identificando especies y prioridades de conservación
  \item Contando individuos
  \item Midiendo áreas
  \item Dibujando mapas
  \item Fotografiando paisajes
  \item Construyendo zoológicos
\end{multicols}
\end{enumerate}
\vspace{0.15cm}
\item ¿Qué es una especie invasora?
\begin{enumerate}[label=\textbf{\Alph*)}]
\begin{multicols}{2}
  \item Una especie introducida que causa daño al ecosistema
  \item Una especie nativa
  \item Una especie endémica
  \item Una especie migratoria
  \item Una especie domesticada
  \item Una especie extinta
\end{multicols}
\end{enumerate}
\vspace{0.15cm}
\item ¿Qué es un hotspot de biodiversidad?
\begin{enumerate}[label=\textbf{\Alph*)}]
\begin{multicols}{2}
  \item Una región con alta diversidad y amenaza
  \item Una región con baja diversidad
  \item Una región sin especies
  \item Una región desértica
  \item Una región polar
  \item Una región oceánica
\end{multicols}
\end{enumerate}
\vspace{0.15cm}
\item ¿Qué es una especie endémica?
\begin{enumerate}[label=\textbf{\Alph*)}]
\begin{multicols}{2}
  \item La que vive solo en una región determinada
  \item La que vive en todo el mundo
  \item La que vive en varios continentes
  \item La que vive en cautiverio
  \item La que vive en zoológicos
  \item La que vive en jardines
\end{multicols}
\end{enumerate}
\vspace{0.15cm}
\item ¿Qué es una especie en peligro?
\begin{enumerate}[label=\textbf{\Alph*)}]
\begin{multicols}{2}
  \item La que tiene alto riesgo de extinción
  \item La que tiene bajo riesgo
  \item La que está extinta
  \item La que es abundante
  \item La que es doméstica
  \item La que es invasora
\end{multicols}
\end{enumerate}
\vspace{0.15cm}
\item ¿Qué es CITES?
\begin{enumerate}[label=\textbf{\Alph*)}]
\begin{multicols}{2}
  \item La convención que regula el comercio de especies amenazadas
  \item La convención sobre el clima
  \item La convención sobre el mar
  \item La convención sobre el trabajo
  \item La convención sobre los derechos humanos
  \item La convención sobre el patrimonio
\end{multicols}
\end{enumerate}
\vspace{0.15cm}
\item ¿Qué es la UICN?
\begin{enumerate}[label=\textbf{\Alph*)}]
\begin{multicols}{2}
  \item La Unión Internacional para la Conservación de la Naturaleza
  \item La Unión Internacional de Química
  \item La Unión Internacional de Física
  \item La Unión Internacional de Matemáticas
  \item La Unión Internacional de Medicina
  \item La Unión Internacional de Ingeniería
\end{multicols}
\end{enumerate}
\vspace{0.15cm}
\item ¿Qué es la Lista Roja?
\begin{enumerate}[label=\textbf{\Alph*)}]
\begin{multicols}{2}
  \item El inventario de especies amenazadas
  \item El inventario de especies abundantes
  \item El inventario de especies domésticas
  \item El inventario de especies invasoras
  \item El inventario de especies extintas
  \item El inventario de especies endémicas
\end{multicols}
\end{enumerate}
\vspace{0.15cm}
\item ¿Cómo se usa la taxonomía en agricultura?
\begin{enumerate}[label=\textbf{\Alph*)}]
\begin{multicols}{2}
  \item Identificando plagas, cultivos y variedades
  \item Midiendo la producción
  \item Contando semillas
  \item Dibujando parcelas
  \item Fotografiando cultivos
  \item Construyendo invernaderos
\end{multicols}
\end{enumerate}
\vspace{0.15cm}
\item ¿Cómo se usa la taxonomía en medicina?
\begin{enumerate}[label=\textbf{\Alph*)}]
\begin{multicols}{2}
  \item Identificando patógenos y organismos medicinales
  \item Midiendo la presión
  \item Contando células
  \item Dibujando anatomia
  \item Fotografiando pacientes
  \item Construyendo hospitales
\end{multicols}
\end{enumerate}
\vspace{0.15cm}
\item ¿Cómo se usa la taxonomía en biotecnología?
\begin{enumerate}[label=\textbf{\Alph*)}]
\begin{multicols}{2}
  \item Identificando organismos para procesos industriales
  \item Midiendo la producción
  \item Contando colonias
  \item Dibujando biorreactores
  \item Fotografiando laboratorios
  \item Construyendo plantas
\end{multicols}
\end{enumerate}
\vspace{0.15cm}
\item ¿Cómo se usa la taxonomía en ecología?
\begin{enumerate}[label=\textbf{\Alph*)}]
\begin{multicols}{2}
  \item Identificando especies en ecosistemas
  \item Midiendo el clima
  \item Contando individuos
  \item Dibujando mapas
  \item Fotografiando paisajes
  \item Construyendo reservas
\end{multicols}
\end{enumerate}
\vspace{0.15cm}
\item ¿Cómo se usa la taxonomía en evolución?
\begin{enumerate}[label=\textbf{\Alph*)}]
\begin{multicols}{2}
  \item Reconstruyendo relaciones filogenéticas
  \item Midiendo el tiempo
  \item Contando fósiles
  \item Dibujando árboles
  \item Fotografiando rocas
  \item Construyendo museos
\end{multicols}
\end{enumerate}
\vspace{0.15cm}
\item ¿Cómo se usa la taxonomía en biogeografía?
\begin{enumerate}[label=\textbf{\Alph*)}]
\begin{multicols}{2}
  \item Explicando la distribución de las especies
  \item Midiendo distancias
  \item Contando regiones
  \item Dibujando mapas
  \item Fotografiando paisajes
  \item Construyendo rutas
\end{multicols}
\end{enumerate}
\vspace{0.15cm}
\item ¿Cómo se usa la taxonomía en cambio climático?
\begin{enumerate}[label=\textbf{\Alph*)}]
\begin{multicols}{2}
  \item Monitoreando cambios en distribución de especies
  \item Midiendo temperatura
  \item Contando glaciares
  \item Dibujando climas
  \item Fotografiando nubes
  \item Construyendo modelos
\end{multicols}
\end{enumerate}
\vspace{0.15cm}
\item ¿Cómo se usa la taxonomía en salud pública?
\begin{enumerate}[label=\textbf{\Alph*)}]
\begin{multicols}{2}
  \item Identificando vectores y patógenos
  \item Midiendo epidemias
  \item Contando pacientes
  \item Dibujando hospitales
  \item Fotografiando salas
  \item Construyendo clínicas
\end{multicols}
\end{enumerate}
\vspace{0.15cm}
\item ¿Cómo se usa la taxonomía forense?
\begin{enumerate}[label=\textbf{\Alph*)}]
\begin{multicols}{2}
  \item Identificando organismos en investigaciones
  \item Midiendo huellas
  \item Contando evidencias
  \item Dibujando escenas
  \item Fotografiando pistas
  \item Construyendo laboratorios
\end{multicols}
\end{enumerate}
\vspace{0.15cm}
\item ¿Cómo se usa la taxonomía en control biológico?
\begin{enumerate}[label=\textbf{\Alph*)}]
\begin{multicols}{2}
  \item Identificando enemigos naturales de plagas
  \item Midiendo plagas
  \item Contando insectos
  \item Dibujando cultivos
  \item Fotografiando campos
  \item Construyendo trampas
\end{multicols}
\end{enumerate}
\vspace{0.15cm}
\item ¿Cómo se usa la taxonomía en pesca?
\begin{enumerate}[label=\textbf{\Alph*)}]
\begin{multicols}{2}
  \item Identificando especies comerciales y reguladas
  \item Midiendo capturas
  \item Contando peces
  \item Dibujando redes
  \item Fotografiando barcos
  \item Construyendo puertos
\end{multicols}
\end{enumerate}
\vspace{0.15cm}
\item ¿Cómo se usa la taxonomía en silvicultura?
\begin{enumerate}[label=\textbf{\Alph*)}]
\begin{multicols}{2}
  \item Identificando especies forestales
  \item Midiendo árboles
  \item Contando hojas
  \item Dibujando bosques
  \item Fotografiando paisajes
  \item Construyendo aserraderos
\end{multicols}
\end{enumerate}
\vspace{0.15cm}
\item ¿Cómo se usa la taxonomía en ganadería?
\begin{enumerate}[label=\textbf{\Alph*)}]
\begin{multicols}{2}
  \item Identificando razas y especies
  \item Midiendo peso
  \item Contando animales
  \item Dibujando corrales
  \item Fotografiando granjas
  \item Construyendo establos
\end{multicols}
\end{enumerate}
\vspace{0.15cm}
\item ¿Cómo se usa la taxonomía en acuicultura?
\begin{enumerate}[label=\textbf{\Alph*)}]
\begin{multicols}{2}
  \item Identificando especies cultivables
  \item Midiendo estanques
  \item Contando alevines
  \item Dibujando granjas
  \item Fotografiando peces
  \item Construyendo piscinas
\end{multicols}
\end{enumerate}
\vspace{0.15cm}
\item ¿Cómo se usa la taxonomía en jardines botánicos?
\begin{enumerate}[label=\textbf{\Alph*)}]
\begin{multicols}{2}
  \item Etiquetando y conservando especies
  \item Midiendo plantas
  \item Contando flores
  \item Dibujando jardines
  \item Fotografiando paisajes
  \item Construyendo invernaderos
\end{multicols}
\end{enumerate}
\vspace{0.15cm}
\item ¿Cómo se usa la taxonomía en zoológicos?
\begin{enumerate}[label=\textbf{\Alph*)}]
\begin{multicols}{2}
  \item Identificando y conservando especies
  \item Midiendo jaulas
  \item Contando animales
  \item Dibujando recintos
  \item Fotografiando visitantes
  \item Construyendo hábitats
\end{multicols}
\end{enumerate}
\vspace{0.15cm}
\item ¿Qué es una base de datos taxonómica?
\begin{enumerate}[label=\textbf{\Alph*)}]
\begin{multicols}{2}
  \item Un sistema que almacena información de especies
  \item Un sistema que almacena minerales
  \item Un sistema que almacena fósiles
  \item Un sistema que almacena rocas
  \item Un sistema que almacena climas
  \item Un sistema que almacena mapas
\end{multicols}
\end{enumerate}
\vspace{0.15cm}
\item ¿Qué es GBIF?
\begin{enumerate}[label=\textbf{\Alph*)}]
\begin{multicols}{2}
  \item La Infraestructura Mundial de Información en Biodiversidad
  \item La Infraestructura Mundial de Información Geológica
  \item La Infraestructura Mundial de Información Climática
  \item La Infraestructura Mundial de Información Médica
  \item La Infraestructura Mundial de Información Agrícola
  \item La Infraestructura Mundial de Información Minera
\end{multicols}
\end{enumerate}
\vspace{0.15cm}
\item ¿Qué es NCBI?
\begin{enumerate}[label=\textbf{\Alph*)}]
\begin{multicols}{2}
  \item El Centro Nacional para la Información Biotecnológica
  \item El Centro Nacional para la Información Geológica
  \item El Centro Nacional para la Información Climática
  \item El Centro Nacional para la Información Médica
  \item El Centro Nacional para la Información Agrícola
  \item El Centro Nacional para la Información Minera
\end{multicols}
\end{enumerate}
\vspace{0.15cm}
\item ¿Qué es BOLD?
\begin{enumerate}[label=\textbf{\Alph*)}]
\begin{multicols}{2}
  \item El sistema de código de barras de la vida
  \item El sistema de código de barras comercial
  \item El sistema de código de barras geológico
  \item El sistema de código de barras climático
  \item El sistema de código de barras médico
  \item El sistema de código de barras agrícola
\end{multicols}
\end{enumerate}
\vspace{0.15cm}
\item ¿Qué es EOL?
\begin{enumerate}[label=\textbf{\Alph*)}]
\begin{multicols}{2}
  \item La Enciclopedia de la Vida
  \item La Enciclopedia de la Tierra
  \item La Enciclopedia del Mar
  \item La Enciclopedia del Clima
  \item La Enciclopedia de la Medicina
  \item La Enciclopedia de la Agricultura
\end{multicols}
\end{enumerate}
\vspace{0.15cm}
\item ¿Qué es ITIS?
\begin{enumerate}[label=\textbf{\Alph*)}]
\begin{multicols}{2}
  \item El Sistema Integrado de Información Taxonómica
  \item El Sistema Integrado de Información Geológica
  \item El Sistema Integrado de Información Climática
  \item El Sistema Integrado de Información Médica
  \item El Sistema Integrado de Información Agrícola
  \item El Sistema Integrado de Información Minera
\end{multicols}
\end{enumerate}
\vspace{0.15cm}
\item ¿Qué es Catalogue of Life?
\begin{enumerate}[label=\textbf{\Alph*)}]
\begin{multicols}{2}
  \item Un catálogo global de especies
  \item Un catálogo de minerales
  \item Un catálogo de fósiles
  \item Un catálogo de rocas
  \item Un catálogo de climas
  \item Un catálogo de mapas
\end{multicols}
\end{enumerate}
\vspace{0.15cm}
\item ¿Qué es WoRMS?
\begin{enumerate}[label=\textbf{\Alph*)}]
\begin{multicols}{2}
  \item El Registro Mundial de Especies Marinas
  \item El Registro Mundial de Especies Terrestres
  \item El Registro Mundial de Especies Dulceacuícolas
  \item El Registro Mundial de Especies Aéreas
  \item El Registro Mundial de Especies Subterráneas
  \item El Registro Mundial de Especies Extintas
\end{multicols}
\end{enumerate}
\vspace{0.15cm}
\item ¿Qué es AlgaeBase?
\begin{enumerate}[label=\textbf{\Alph*)}]
\begin{multicols}{2}
  \item Una base de datos de algas
  \item Una base de datos de hongos
  \item Una base de datos de plantas
  \item Una base de datos de animales
  \item Una base de datos de bacterias
  \item Una base de datos de virus
\end{multicols}
\end{enumerate}
\vspace{0.15cm}
\item ¿Qué es Index Fungorum?
\begin{enumerate}[label=\textbf{\Alph*)}]
\begin{multicols}{2}
  \item Una base de datos de hongos
  \item Una base de datos de algas
  \item Una base de datos de plantas
  \item Una base de datos de animales
  \item Una base de datos de bacterias
  \item Una base de datos de virus
\end{multicols}
\end{enumerate}
\vspace{0.15cm}
\item ¿Qué es IPNI?
\begin{enumerate}[label=\textbf{\Alph*)}]
\begin{multicols}{2}
  \item El Índice Internacional de Nombres de Plantas
  \item El Índice Internacional de Nombres de Animales
  \item El Índice Internacional de Nombres de Hongos
  \item El Índice Internacional de Nombres de Algas
  \item El Índice Internacional de Nombres de Bacterias
  \item El Índice Internacional de Nombres de Virus
\end{multicols}
\end{enumerate}
\vspace{0.15cm}
\item ¿Qué es ZooBank?
\begin{enumerate}[label=\textbf{\Alph*)}]
\begin{multicols}{2}
  \item El registro oficial de nombres zoológicos
  \item El registro oficial de nombres botánicos
  \item El registro oficial de nombres micológicos
  \item El registro oficial de nombres bacteriológicos
  \item El registro oficial de nombres virales
  \item El registro oficial de nombres algales
\end{multicols}
\end{enumerate}
\vspace{0.15cm}
\item ¿Qué es MycoBank?
\begin{enumerate}[label=\textbf{\Alph*)}]
\begin{multicols}{2}
  \item Una base de datos de nombres micológicos
  \item Una base de datos de nombres zoológicos
  \item Una base de datos de nombres botánicos
  \item Una base de datos de nombres bacteriológicos
  \item Una base de datos de nombres virales
  \item Una base de datos de nombres algales
\end{multicols}
\end{enumerate}
\vspace{0.15cm}
\item ¿Qué es la taxonomía integrativa?
\begin{enumerate}[label=\textbf{\Alph*)}]
\begin{multicols}{2}
  \item La que combina datos morfológicos, moleculares y ecológicos
  \item La que usa solo morfología
  \item La que usa solo moléculas
  \item La que usa solo ecología
  \item La que usa solo comportamiento
  \item La que usa solo biogeografía
\end{multicols}
\end{enumerate}
\vspace{0.15cm}
\item ¿Qué es la taxonomía alfa?
\begin{enumerate}[label=\textbf{\Alph*)}]
\begin{multicols}{2}
  \item La descripción de nuevas especies
  \item La reconstrucción filogenética
  \item La clasificación superior
  \item La nomenclatura
  \item La colección
  \item La conservación
\end{multicols}
\end{enumerate}
\vspace{0.15cm}
\item ¿Qué es la taxonomía beta?
\begin{enumerate}[label=\textbf{\Alph*)}]
\begin{multicols}{2}
  \item La clasificación de niveles superiores
  \item La descripción de nuevas especies
  \item La nomenclatura
  \item La colección
  \item La conservación
  \item La biogeografía
\end{multicols}
\end{enumerate}
\vspace{0.15cm}
\item ¿Qué es la taxonomía gamma?
\begin{enumerate}[label=\textbf{\Alph*)}]
\begin{multicols}{2}
  \item La integración de datos en clasificaciones completas
  \item La descripción de nuevas especies
  \item La nomenclatura
  \item La colección
  \item La conservación
  \item La biogeografía
\end{multicols}
\end{enumerate}
\vspace{0.15cm}
\item ¿Qué es un filograma?
\begin{enumerate}[label=\textbf{\Alph*)}]
\begin{multicols}{2}
  \item Un árbol filogenético con longitudes de rama proporcionales al cambio
  \item Una lista de especies
  \item Una clave dicotómica
  \item Un catálogo de nombres
  \item Un herbario
  \item Un museo
\end{multicols}
\end{enumerate}
\vspace{0.15cm}
\item ¿Qué es un cladograma?
\begin{enumerate}[label=\textbf{\Alph*)}]
\begin{multicols}{2}
  \item Un diagrama que muestra relaciones por caracteres derivados compartidos
  \item Una lista de especies
  \item Una clave dicotómica
  \item Un catálogo de nombres
  \item Un herbario
  \item Un museo
\end{multicols}
\end{enumerate}
\vspace{0.15cm}
\item ¿Qué es la sistemática filogenética?
\begin{enumerate}[label=\textbf{\Alph*)}]
\begin{multicols}{2}
  \item La clasificación basada estrictamente en la filogenia
  \item La clasificación basada en el color
  \item La clasificación basada en el tamaño
  \item La clasificación basada en el hábitat
  \item La clasificación basada en la utilidad
  \item La clasificación basada en la forma
\end{multicols}
\end{enumerate}
\vspace{0.15cm}
\item ¿Qué es la escuela fenética?
\begin{enumerate}[label=\textbf{\Alph*)}]
\begin{multicols}{2}
  \item La clasificación basada en similitud global
  \item La clasificación basada en filogenia
  \item La clasificación basada en cladística
  \item La clasificación basada en evolución
  \item La clasificación basada en genética
  \item La clasificación basada en molecular
\end{multicols}
\end{enumerate}
\vspace{0.15cm}
\item ¿Qué es la escuela cladista?
\begin{enumerate}[label=\textbf{\Alph*)}]
\begin{multicols}{2}
  \item La clasificación basada en caracteres derivados compartidos
  \item La clasificación basada en similitud global
  \item La clasificación basada en color
  \item La clasificación basada en tamaño
  \item La clasificación basada en hábitat
  \item La clasificación basada en utilidad
\end{multicols}
\end{enumerate}
\vspace{0.15cm}
\item ¿Qué es la escuela evolutiva?
\begin{enumerate}[label=\textbf{\Alph*)}]
\begin{multicols}{2}
  \item La clasificación que considera filogenia y divergencia
  \item La clasificación basada solo en similitud
  \item La clasificación basada solo en cladística
  \item La clasificación basada solo en color
  \item La clasificación basada solo en tamaño
  \item La clasificación basada solo en hábitat
\end{multicols}
\end{enumerate}
\vspace{0.15cm}
\item ¿Qué es un carácter taxonómico?
\begin{enumerate}[label=\textbf{\Alph*)}]
\begin{multicols}{2}
  \item Cualquier atributo usado para clasificar
  \item Solo el color
  \item Solo el tamaño
  \item Solo la forma
  \item Solo el hábitat
  \item Solo la utilidad
\end{multicols}
\end{enumerate}
\vspace{0.15cm}
\item ¿Qué es la biodiversidad taxonómica?
\begin{enumerate}[label=\textbf{\Alph*)}]
\begin{multicols}{2}
  \item La variedad de taxones en un área
  \item La variedad de genes
  \item La variedad de ecosistemas
  \item La variedad de individuos
  \item La variedad de poblaciones
  \item La variedad de comunidades ---
\end{multicols}
\end{enumerate}
\vspace{0.15cm}
\end{enumerate}\n
\vspace{0.4cm}
\begin{tcolorbox}[colback=gray!5!white,colframe=gray!60!black,title=\textbf{SOLUCIONARIO -- 5.4 Taxonomía y Ecosistemas}]
\begin{multicols}{4}
\begin{enumerate}[label=\textbf{\arabic*.}]
  \item \textbf{A}
  \item \textbf{A}
  \item \textbf{A}
  \item \textbf{A}
  \item \textbf{A}
  \item \textbf{A}
  \item \textbf{A}
  \item \textbf{A}
  \item \textbf{A}
  \item \textbf{A}
  \item \textbf{A}
  \item \textbf{A}
  \item \textbf{A}
  \item \textbf{A}
  \item \textbf{A}
  \item \textbf{A}
  \item \textbf{A}
  \item \textbf{A}
  \item \textbf{A}
  \item \textbf{A}
  \item \textbf{A}
  \item \textbf{A}
  \item \textbf{A}
  \item \textbf{A}
  \item \textbf{A}
  \item \textbf{A}
  \item \textbf{A}
  \item \textbf{A}
  \item \textbf{A}
  \item \textbf{A}
  \item \textbf{A}
  \item \textbf{A}
  \item \textbf{A}
  \item \textbf{A}
  \item \textbf{A}
  \item \textbf{A}
  \item \textbf{A}
  \item \textbf{A}
  \item \textbf{A}
  \item \textbf{A}
  \item \textbf{A}
  \item \textbf{A}
  \item \textbf{A}
  \item \textbf{A}
  \item \textbf{A}
  \item \textbf{A}
  \item \textbf{A}
  \item \textbf{A}
  \item \textbf{A}
  \item \textbf{A}
  \item \textbf{A}
  \item \textbf{A}
  \item \textbf{A}
  \item \textbf{A}
  \item \textbf{A}
  \item \textbf{A}
  \item \textbf{A}
  \item \textbf{A}
  \item \textbf{A}
  \item \textbf{A}
  \item \textbf{A}
  \item \textbf{A}
  \item \textbf{A}
  \item \textbf{A}
  \item \textbf{A}
  \item \textbf{A}
  \item \textbf{A}
  \item \textbf{A}
  \item \textbf{A}
  \item \textbf{A}
  \item \textbf{A}
  \item \textbf{A}
  \item \textbf{A}
  \item \textbf{A}
  \item \textbf{A}
  \item \textbf{A}
  \item \textbf{A}
  \item \textbf{A}
  \item \textbf{A}
  \item \textbf{A}
  \item \textbf{A}
  \item \textbf{A}
  \item \textbf{A}
  \item \textbf{A}
  \item \textbf{A}
  \item \textbf{A}
  \item \textbf{A}
  \item \textbf{A}
  \item \textbf{A}
  \item \textbf{A}
  \item \textbf{A}
  \item \textbf{A}
  \item \textbf{A}
  \item \textbf{A}
  \item \textbf{A}
  \item \textbf{A}
  \item \textbf{A}
  \item \textbf{A}
  \item \textbf{A}
  \item \textbf{A}
  \item \textbf{A}
  \item \textbf{A}
  \item \textbf{A}
  \item \textbf{A}
  \item \textbf{A}
  \item \textbf{A}
  \item \textbf{A}
  \item \textbf{A}
  \item \textbf{A}
  \item \textbf{A}
  \item \textbf{A}
  \item \textbf{A}
  \item \textbf{A}
  \item \textbf{A}
  \item \textbf{A}
  \item \textbf{A}
  \item \textbf{A}
  \item \textbf{A}
  \item \textbf{A}
  \item \textbf{A}
  \item \textbf{A}
  \item \textbf{A}
  \item \textbf{A}
  \item \textbf{A}
  \item \textbf{A}
  \item \textbf{A}
  \item \textbf{A}
  \item \textbf{A}
  \item \textbf{A}
  \item \textbf{A}
  \item \textbf{A}
  \item \textbf{A}
  \item \textbf{A}
  \item \textbf{A}
  \item \textbf{A}
  \item \textbf{A}
  \item \textbf{A}
  \item \textbf{A}
  \item \textbf{A}
  \item \textbf{A}
  \item \textbf{A}
  \item \textbf{A}
  \item \textbf{A}
  \item \textbf{A}
  \item \textbf{A}
  \item \textbf{A}
  \item \textbf{A}
  \item \textbf{A}
  \item \textbf{A}
  \item \textbf{A}
  \item \textbf{A}
  \item \textbf{A}
  \item \textbf{A}
  \item \textbf{A}
  \item \textbf{A}
  \item \textbf{A}
  \item \textbf{A}
  \item \textbf{A}
  \item \textbf{A}
  \item \textbf{A}
  \item \textbf{A}
  \item \textbf{A}
  \item \textbf{A}
  \item \textbf{A}
  \item \textbf{A}
  \item \textbf{A}
  \item \textbf{A}
  \item \textbf{A}
  \item \textbf{A}
  \item \textbf{A}
  \item \textbf{A}
  \item \textbf{A}
  \item \textbf{A}
  \item \textbf{A}
  \item \textbf{A}
  \item \textbf{A}
  \item \textbf{A}
  \item \textbf{A}
  \item \textbf{A}
  \item \textbf{A}
  \item \textbf{A}
  \item \textbf{A}
  \item \textbf{A}
  \item \textbf{A}
  \item \textbf{A}
  \item \textbf{A}
  \item \textbf{A}
  \item \textbf{A}
  \item \textbf{A}
  \item \textbf{A}
  \item \textbf{A}
  \item \textbf{A}
  \item \textbf{A}
  \item \textbf{A}
  \item \textbf{A}
  \item \textbf{A}
  \item \textbf{A}
  \item \textbf{A}
  \item \textbf{A}
  \item \textbf{A}
  \item \textbf{A}
  \item \textbf{A}
  \item \textbf{A}
  \item \textbf{A}
  \item \textbf{A}
  \item \textbf{A}
  \item \textbf{A}
  \item \textbf{A}
  \item \textbf{A}
  \item \textbf{A}
  \item \textbf{A}
  \item \textbf{A}
  \item \textbf{A}
  \item \textbf{A}
  \item \textbf{A}
  \item \textbf{A}
  \item \textbf{A}
  \item \textbf{A}
  \item \textbf{A}
  \item \textbf{A}
  \item \textbf{A}
  \item \textbf{A}
  \item \textbf{A}
  \item \textbf{A}
  \item \textbf{A}
  \item \textbf{A}
  \item \textbf{A}
  \item \textbf{A}
  \item \textbf{A}
  \item \textbf{A}
  \item \textbf{A}
  \item \textbf{A}
  \item \textbf{A}
  \item \textbf{A}
  \item \textbf{A}
  \item \textbf{A}
  \item \textbf{A}
  \item \textbf{A}
  \item \textbf{A}
  \item \textbf{A}
  \item \textbf{A}
  \item \textbf{A}
  \item \textbf{A}
  \item \textbf{A}
  \item \textbf{A}
  \item \textbf{A}
  \item \textbf{A}
  \item \textbf{A}
  \item \textbf{A}
  \item \textbf{A}
  \item \textbf{A}
  \item \textbf{A}
  \item \textbf{A}
  \item \textbf{A}
  \item \textbf{A}
  \item \textbf{A}
  \item \textbf{A}
  \item \textbf{A}
  \item \textbf{A}
  \item \textbf{A}
  \item \textbf{A}
  \item \textbf{A}
  \item \textbf{A}
  \item \textbf{A}
  \item \textbf{A}
  \item \textbf{A}
  \item \textbf{A}
  \item \textbf{A}
  \item \textbf{A}
  \item \textbf{A}
  \item \textbf{A}
  \item \textbf{A}
  \item \textbf{A}
  \item \textbf{A}
  \item \textbf{A}
  \item \textbf{A}
  \item \textbf{A}
  \item \textbf{A}
  \item \textbf{A}
  \item \textbf{A}
  \item \textbf{A}
  \item \textbf{A}
  \item \textbf{A}
  \item \textbf{A}
  \item \textbf{A}
  \item \textbf{A}
  \item \textbf{A}
  \item \textbf{A}
  \item \textbf{A}
  \item \textbf{A}
  \item \textbf{A}
  \item \textbf{A}
  \item \textbf{A}
  \item \textbf{A}
  \item \textbf{A}
  \item \textbf{A}
  \item \textbf{A}
  \item \textbf{A}
  \item \textbf{A}
  \item \textbf{A}
\end{enumerate}
\end{multicols}
\end{tcolorbox}
\newpage\n
\end{document}